% ============================================================================
%  中国农业周期性分析报告 —— 第二部分（农业上市公司画像）单独成册
%  编译：xelatex main-part2.tex（两遍 + biber），或 make part2
%  内容 = 封面 + 摘要（第二部分口径）+ 目录 + 图表目录
%         + 第七至第十五章 + 参考文献 + 附录 A（数据来源与口径）、附录 B（公司明细表）
%  正文与 main.tex 同源（\input 同一批 tex/ch07--ch15.tex 与 tex/gen/**），
%  不复制、不改写正文，因此章号、节号、图/表编号与主报告完全一致。
%  第二部分正文里有若干指向第一部分的交叉引用（如“第 \ref{sec:hog} 章”“第
%  \ref{sec:soft} 章”）——第一部分不在本册中，故由 scripts/make_part_refs.py 从
%  main.aux 取出这些外部标签的编号，生成 tex/gen/part1-refs.tex 作为桩文件引入，
%  使引用仍印出与主报告一致的编号。（不用 xr-hyper 的 \externaldocument{main}：
%  它会把 main.aux 的全部标签一起引入，其中包含第二部分自身的标签，从而每次编译都
%  报上百条 “Label ... multiply defined”。）
% ============================================================================
\documentclass[11pt, a4paper, fontset=none, zihao=-4]{ctexart}
% ============================================================================
%  preamble.tex —— 中国农业周期性分析报告 导言区
%  编译：xelatex（需 ctex + tikz/pgfplots + Source Han 字体）
% ============================================================================

% ---------- 中文与字体 ----------
% 注意：ctexart 类已加载 ctex，此处不可重复 \usepackage{ctex}
% （重复加载会触发 ctexhook Warning 并破坏 fontspec 参数解析）
\usepackage{fontspec}
\setmainfont[AutoFakeBold=2.5]{Source Han Serif CN}
\setsansfont[AutoFakeBold=2.5]{Source Han Sans CN}
\setmonofont{WenQuanYi Zen Hei Mono}
\setCJKmainfont[AutoFakeBold=2.5]{Source Han Serif CN}
\setCJKsansfont[AutoFakeBold=2.5]{Source Han Sans CN}
\setCJKmonofont{WenQuanYi Zen Hei Mono}
\newCJKfontfamily\heiti{Source Han Sans CN}
\newCJKfontfamily\songti{Source Han Serif CN}
\newCJKfontfamily\heiB[AutoFakeBold=5]{Source Han Sans CN}

% ---------- 版面 ----------
\usepackage[top=2.6cm, bottom=2.6cm, left=2.6cm, right=2.6cm]{geometry}
\usepackage{setspace}
\onehalfspacing
\sloppy
\setlength{\parskip}{2pt plus 1pt}
\setlength{\parindent}{2em}
\usepackage{fancyhdr}
\usepackage{titlesec}

% ---------- 数学与符号 ----------
\usepackage{amsmath, amssymb, amsfonts, mathtools}
\usepackage{bm}

% ---------- 表格 ----------
\usepackage{booktabs}
\usepackage{longtable}
\usepackage{array}
\usepackage{tabularx}
\usepackage{multirow}
\usepackage{makecell}
\usepackage{threeparttable}
\usepackage{siunitx}
\sisetup{group-separator={,}, group-minimum-digits=4, detect-weight=true}

% ---------- 图形 ----------
\usepackage{graphicx}
\usepackage{xcolor}
\usepackage{tikz}
\usepackage{pgfplots}
\usepackage{pgfplotstable}
\pgfplotsset{compat=1.18}
\usepgfplotslibrary{fillbetween, groupplots, statistics, patchplots}
\usetikzlibrary{arrows.meta, positioning, shapes.geometric, calc, fit,
                backgrounds, patterns, decorations.pathreplacing, shadows.blur,
                matrix, chains, fadings, intersections}
% TikZ 外部库（如未安装 external 依赖则自动忽略，不影响编译）
\usetikzlibrary{external}

% ---------- 参考文献（gb7714-2015 国标，biber 后端） ----------
\usepackage[backend=biber, style=gb7714-2015, sorting=none,
            url=true, doi=false, isbn=false, eprint=false]{biblatex}
\addbibresource{refs.bib}
\setlength{\bibitemsep}{2pt}
\renewcommand*{\bibfont}{\footnotesize}

% ---------- 颜色主题（农本绿） ----------
\definecolor{AgriGreen}{HTML}{2E6E4E}
\definecolor{AgriLight}{HTML}{6FA287}
\definecolor{AgriPale}{HTML}{E4EFE8}
\definecolor{AgriGold}{HTML}{C7A252}
\definecolor{HogPink}{HTML}{C97B84}
\definecolor{Grain}{HTML}{D9A441}
\definecolor{OilBlue}{HTML}{3B6E9E}
\definecolor{SoftRed}{HTML}{B5533C}
\definecolor{AquaTeal}{HTML}{2C8C87}
\definecolor{InkGray}{HTML}{3A3F44}
\definecolor{LightGray}{HTML}{F2F3F4}

% 供 pgfplots 循环使用的品种配色
\pgfplotscreateplotcyclelist{agri}{
  {AgriGreen, thick},
  {HogPink, thick},
  {Grain, thick},
  {OilBlue, thick},
  {SoftRed, thick},
  {AquaTeal, thick},
  {AgriGold, thick},
  {InkGray, thick},
  {AgriLight, thick},
  {violet, thick},
}

% 通用折线图样式
\pgfplotsset{
  % 年份坐标轴：x 为十进制年份，刻度标签取整
  yearticks/.style={
    xticklabel style={/pgf/number format/fixed,
                      /pgf/number format/precision=0,
                      /pgf/number format/1000 sep={},
                      font=\footnotesize},
    x tick label as interval=false,
  },
  % 相关系数条形图（y 轴为类别编号）
  catbar/.style={
    ytick=data, y dir=reverse,
    yticklabel style={font=\scriptsize},
    xlabel style={font=\small},
    tick label style={font=\footnotesize},
    axis line style={InkGray, line width=0.5pt},
    grid=major, grid style={LightGray, line width=0.3pt},
    enlarge y limits=0.03,
  },
  agri axis/.style={
    width=13.2cm, height=6.2cm,
    tick label style={font=\footnotesize},
    label style={font=\small},
    title style={font=\small\heiti},
    legend style={font=\scriptsize, fill=white, fill opacity=0.85,
                  text opacity=1, draw=AgriLight!50, rounded corners=1pt},
    grid=major, grid style={LightGray, line width=0.3pt},
    axis line style={InkGray, line width=0.5pt},
    every axis plot/.append style={line width=1pt},
  },
}

% ---------- 列表、盒子 ----------
\usepackage{enumitem}
\setlist[itemize]{itemsep=2pt, topsep=3pt, parsep=0pt, leftmargin=1.6em}
\setlist[enumerate]{itemsep=2pt, topsep=3pt, parsep=0pt, leftmargin=1.8em}
\usepackage[most]{tcolorbox}

\newtcolorbox{keybox}[1][核心结论]{
  enhanced, breakable, colback=AgriPale, colframe=AgriGreen,
  boxrule=0.4pt, left=6pt, right=6pt, top=4pt, bottom=4pt,
  fonttitle=\sffamily\bfseries\small, title={#1},
  attach boxed title to top left={xshift=6pt, yshift=-2.4pt},
  boxed title style={colback=AgriGreen, colframe=AgriGreen, size=small},
  arc=1.5pt, before skip=6pt, after skip=6pt}

\newtcolorbox{noteh}[1][数据说明]{
  enhanced, breakable, colback=LightGray, colframe=InkGray!40,
  boxrule=0.3pt, left=6pt, right=6pt, top=4pt, bottom=4pt,
  fonttitle=\sffamily\bfseries\scriptsize, title={#1},
  arc=1pt, before skip=5pt, after skip=5pt}

% ---------- 标题格式 ----------
\ctexset{
  section/format = {\heiti\zihao{4}\color{AgriGreen}},
  subsection/format = {\heiti\zihao{-4}\color{InkGray}},
  subsubsection/format = {\heiti\zihao{-4}\color{InkGray!80!black}},
  section/aftername = \quad,
  subsection/aftername = \quad,
}
% 部（Part）样式：ctexart 无 chapter，故用 titlesec 定制 \part
\titleclass{\part}{straight}[\section]
\titleformat{\part}[display]
  {\centering\heiti\zihao{2}\color{AgriGreen}}
  {\zihao{5}\sffamily\color{AgriGold} 第\,\zhnumber{\thepart}\,部分}
  {12pt}
  {}
  [\vspace{-8pt}\textcolor{AgriLight}{\rule{0.6\textwidth}{0.8pt}}]
\titlespacing*{\part}{0pt}{0pt}{28pt}
\usepackage{zhnumber}
\titleformat{\paragraph}[runin]{\heiti\zihao{-4}}{}{0pt}{}[：]

% ---------- 代码与超链接 ----------
\usepackage{listings}
\lstset{basicstyle=\ttfamily\scriptsize, breaklines=true,
        frame=single, rulecolor=\color{InkGray!40},
        backgroundcolor=\color{LightGray}, columns=fullflexible}
\usepackage[colorlinks=true, linkcolor=AgriGreen, citecolor=AgriGreen,
            urlcolor=OilBlue, bookmarksnumbered=true,
            CJKbookmarks=true]{hyperref}
\usepackage{footmisc}
\usepackage{caption}
\captionsetup{font={small,sf}, labelfont={bf,color=AgriGreen},
              labelsep=quad, skip=4pt}
\captionsetup[table]{position=top}

% ---------- 自定义宏 ----------
% 图表数据来源脚注（放在图表正下方）
\newcommand{\srcfile}[1]{\par\vspace{-2pt}%
  {\footnotesize\sffamily\color{InkGray!85}资料来源：#1\par}}
% 数据口径说明
\newcommand{\datacal}[1]{{\footnotesize\sffamily\color{InkGray!75}（#1）}}
% 周期阶段标记
\newcommand{\phase}[1]{\textcolor{AgriGreen}{\textbf{【#1】}}}
% 关键数字高亮
\newcommand{\hl}[1]{\textcolor{SoftRed}{\textbf{#1}}}
\newcommand{\up}{\textcolor{SoftRed}{$\blacktriangle$}}
\newcommand{\down}{\textcolor{AgriGreen}{$\blacktriangledown$}}
% 公司代码
\newcommand{\stk}[1]{\texttt{#1}}
% 缺失数据标记
\newcommand{\nadata}{\textcolor{InkGray!60}{---}}

% 由 scripts/make_part_refs.py 生成 —— 勿手工修改。
% 第二部分分册用到的第一部分标签（编号与页码取自 main.aux）。
% 第二部分自身的标签不在此文件中，由本册正常定义，因此不会出现
% “Label ... multiply defined” 警告。
\makeatletter
\@ifundefined{r@sec:corr}{\newlabel{sec:corr}{{6}{24}}}{}
\@ifundefined{r@sec:framework}{\newlabel{sec:framework}{{1}{2}}}{}
\@ifundefined{r@sec:grain}{\newlabel{sec:grain}{{3}{12}}}{}
\@ifundefined{r@sec:hog}{\newlabel{sec:hog}{{2}{6}}}{}
\@ifundefined{r@sec:hog-cost}{\newlabel{sec:hog-cost}{{2.4}{9}}}{}
\@ifundefined{r@sec:hog-shock}{\newlabel{sec:hog-shock}{{2.5}{10}}}{}
\@ifundefined{r@sec:other-aqua}{\newlabel{sec:other-aqua}{{5.5}{22}}}{}
\@ifundefined{r@sec:other-livestock}{\newlabel{sec:other-livestock}{{5}{20}}}{}
\@ifundefined{r@sec:other-poultry}{\newlabel{sec:other-poultry}{{5.4}{22}}}{}
\@ifundefined{r@sec:soft}{\newlabel{sec:soft}{{4}{16}}}{}
\makeatother
        % 第一部分标签的桩（仅这些，避免标签重定义）

\title{中国农业周期性分析报告（第二部分）}
\author{农业行业研究}
\date{2026 年 10 月}

\begin{document}

% ---------------------------------------------------------------- 封面
% ---------------------------------------------------------------- 封面
% 按要求只保留中文大标题、中文副标题与日期三项。
\begin{titlepage}
\centering
\vspace*{4.2cm}

% 标题用 \zihao{1}：\zihao{0} 时 12 个字宽于版心，会折行成“报/告”两行
{\heiti\zihao{1} 中国农业周期性分析报告}

\vspace{1.1cm}
{\heiti\zihao{3} 基于申万行业分类的品种周期机制与上市公司画像}

\vfill
{\heiti\zihao{4} 2026 年 10 月}
\vspace{2.5cm}
\end{titlepage}

\newpage

% ---------------------------------------------------------------- 摘要与目录
\pagenumbering{Roman}
% ---------------------------------------------------------------- 摘要（第二部分分册）
% 与 tex/abstract.tex 同源，仅供“第二部分”分册使用：
%   * 保留全报告定位一句与“第二部分的主要发现”六条；
%   * 删去“第一部分的主要发现”四条（其口径、样本与结论属于第一部分，见另册）；
%   * 本册保留的 \ref 全部指向第二部分自身的章节（第 7 章、第 8--13 章），
%     分册中这些标签均已定义，不会印成 ??。
\section*{摘要}
\addcontentsline{toc}{section}{摘要}

\begin{keybox}[报告要点]
\small
本报告以\textbf{申万行业分类（2021 版）}为统一框架，分为两大部分：
第一部分从供需关系出发分析农业各子行业的周期性，以\hl{猪周期为基准}
测算跨品种、跨行业的相关性与传导时滞；第二部分系统梳理
申万农林牧渔行业\textbf{全部 \num{104} 家 A 股上市公司}近三年的经营与融资画像。

\textbf{本册为第二部分}，对应第七章至第十五章（逐公司画像、融资历史与潜在融资需求）
以及附录 A、附录 B；第一部分（子行业周期性与跨品种相关性）为另一册。

\textbf{第二部分的主要发现}：
（1）\hl{第二部分统一采用 2026 年半年报口径（合并报表、未年化），
市值一律采用总市值口径}（2026 年 10 月 12 日收盘价 $\times$ 总股本，
不使用流通市值）：100 余家公司的总市值合计 \num{12237} 亿元，
前三家占 \num{38.2}\%；
\num{104} 家公司中 \num{46} 家亏损（占 \num{44.2}\%）、\num{44} 家半年 ROE 为负，
ROE 中位数仅 \num{0.75}\%，半年 ROE 超过 \num{5}\%（约合年化 \num{10}\%）的只有 \num{16} 家；
其中生猪养殖 \num{11} 家在 2026 年上半年\hl{全部亏损}、合计 \num{-163.8} 亿元
（上年同期为 \num{+154.1} 亿元）；
（2）子行业分化鲜明：动保（半年 ROE 中位数 \num{2.95}\%、亏损面 \num{11.1}\%、
资产负债率中位数 \num{30.2}\%）与海洋捕捞（\num{9.79}\%）为抗周期方向，
饲料（半年 ROE 中位数 \num{-6.17}\%、亏损面 \num{66.7}\%）为最受损环节；
（3）\num{14} 家公司同时处于 2026 年上半年亏损与 6 月末资产负债率 $>$ \num{65}\% 状态
（其中 \num{6} 家 $>$ \num{80}\%），\num{31} 家连续三年未分红；
据财务约束与公告行为推断，第一档“补血型”融资需求按最低口径为 \num{150}---\num{250} 亿元、
按“负债率压回 \num{65}\%”的重组口径约 \num{305} 亿元
（若改以偿还债务达标则需减少负债约 \num{567} 亿元），
\hl{当前农业上市公司的资本运作主题是修复资产负债表，而非扩大产能}；
（4）与 2025 年年报（\num{34} 家亏损、ROE 中位数 \num{2.64}\%）相比，
\hl{2026 年半年报的亏损面进一步扩大}：归母净利润合计由上年同期的 \num{+270.9} 亿元
转为 \num{-126.1} 亿元，而营业收入合计 \num{6085} 亿元并未同幅萎缩，
\hl{报表层面的底部晚于价格层面的底部}；
（5）逐公司画像给出三处关键事实：\hl{北大荒} 2026 年上半年由盈转亏，
主因是补缴 2021---2025 年企业所得税与滞纳金 \num{14.10} 亿元（一次性计入当期损益、
不追溯调整），剔除后扣非亏损仅 \num{2.09} 亿元，土地发包主业并未恶化；
\hl{种业去产能已经开始}，全国玉米制种面积自 2023 年的 \num{462} 万亩连续三年降至
2026 年的 \num{358.5} 万亩（\num{-24}\%），但 \num{175}\% 的供需比意味着库存出清尚需时间；
\hl{诺普信}的利润已主要来自基质栽培、错季上市的蓝莓（子公司爱莓庄上半年净利润 \num{6.07} 亿元），
农药制剂业务同期仍稳定增长 \num{2.4}\%；
（6）在周期底部真正区分公司的是\hl{\emph{在低点而非高点加杠杆扩张的能力}}——
在景气高点借来的产能会把成本锁在最高处，只有在行业普遍亏损时仍能拿到低成本资金的公司
才有条件承接别人的产能。
第二部分的各行业章（第 \ref{sec:planting} 章至第 \ref{sec:fishery} 章）末尾设有“上市公司画像”小节，
给出全部 \num{104} 家公司的单独画像。画像\hl{逐家检索公开资料}写成，
不使用统一模板：每家公司的沿革与股权变动、融资历史、风险事件与战略动作各不相同，
正文按“基本情况、股权结构与控制权、经营要点、所处行业的基本面（供需）、
与同类公司比较的竞争优势、股价与周期、盈利与财务、战略转型与融资”的顺序展开，
每条事实均可追溯至公司公告、定期报告或公开报道（数据口径见第 \ref{sec:comp-overview} 章）；
每家配股价与行业指数对比图、近年主要财务指标图
（营业收入与归母净利润、ROE 与资产负债率、货币资金与有息负债）、
关键财务指标表与历史融资/资本运作表。
对经历过重整或债务重组的公司（如天邦食品、正邦科技、傲农生物、华英农业等）
单独梳理其融资与重组过程，并按最新财务指标判断未来的潜在融资需求
（以短期债务缺口与年化经营缺口两个口径度量）。
\end{keybox}

\vspace{6pt}
\noindent\textbf{关键词}\quad 农业周期；猪周期；跨品种相关性；申万行业分类；
上市公司画像；融资需求
\clearpage

\newpage
\makeatletter
\renewcommand\tableofcontents{%
  \@mkboth{\contentsname}{\contentsname}%
  \@starttoc{toc}%
}
\makeatother
\tableofcontents
\newpage

% ---------------------------------------------------------------- 图表目录
\phantomsection
\addcontentsline{toc}{section}{图目录}
\listoffigures
\clearpage
\phantomsection
\addcontentsline{toc}{section}{表目录}
\listoftables
\newpage

% ---------------------------------------------------------------- 正文
\pagenumbering{arabic}
\pagestyle{fancy}
\fancyhf{}
\fancyhead[L]{\small\sffamily 中国农业周期性分析报告（第二部分）}
\fancyhead[R]{\small\sffamily\leftmark}
\fancyfoot[C]{\small\thepage}
\renewcommand{\headrulewidth}{0.4pt}
\renewcommand{\headrule}{\hbox to\headwidth{\leaders\hrule height \headrulewidth\hfill}}

% 部分与章的编号与主报告对齐：主报告中第二部分是“第二部分”、首个 \section 为第 7 章
\setcounter{part}{1}
\part{农业上市公司画像（近三年）}
\setcounter{section}{6}
\section{农业上市公司总览：样本、结构与盈利分布}
\label{sec:comp-overview}

第二部分的目标不是评价个股，而是回答一个问题：\hl{在本轮农业下行周期中，
农业上市公司的经营与融资行为呈现出什么共性}？本部分的样本为申万行业分类（2021 版）
一级行业“农林牧渔”下属的全部 A 股上市公司，共 \num{104} 家。

\subsection{样本与数据口径}
\label{sec:comp-sample}

研究样本为\hl{申万行业分类（2021 版）一级行业“农林牧渔”（行业代码 801010）
下属的全部 A 股上市公司}，共 \num{104} 家。数据处理与第 \ref{sec:framework} 部分一致：
先落盘原始数据，再生成可复现的汇总表，正文中的所有数字均由脚本从原始数据计算得出。

\hl{本部分（第 \ref{sec:comp-overview} 章至第 \ref{sec:financing} 章）的财务数据统一采用
2026 年半年报}（合并报表口径、未年化）；凡属时点性质的指标，资产负债率取 2026 年 6 月末数值。
2023---2025 年年报数据在本部分中只作为\emph{趋势对照}出现，不再作为当期口径；
年报与半年报不可直接比较之处（半年报未年化、种植业收入确认的季节性）均在正文与图表题注中标明。

\renewcommand{\arraystretch}{1.15}
\begin{table}[htbp]
\centering
\footnotesize
\caption{农业上市公司画像的数据来源与用途}
\label{tab:comp-source}
\begin{tabular}{@{}p{2.8cm}p{4.0cm}p{7.8cm}@{}}
\toprule
数据类别 & 来源 & 用途 \\
\midrule
公司概况 & 上市公司公告 & 上市日期、所属市场、注册地、主营业务、注册资本 \\
行情与市值 & 公开行情数据与\textbf{总股本} & 最新收盘价、\textbf{总市值}（收盘价 $\times$ 总股本）、近一年涨跌幅、月度股价序列 \\
|财务数据 & 上市公司定期报告（合并报表） & \textbf{2026 年半年报}的营业收入、归母净利润、ROE、毛利率、净利率与资产负债率（2026 年 6 月末）；2023---2025 年年报数据作趋势对照 \\
|债务与偿债 & 上市公司资产负债表、现金流量表 & 2026 年 6 月末的货币资金、短期借款、一年内到期的非流动负债、长期借款与应付债券，并据此计算有息负债与现金短债比；2023 年以来“取得借款收到的现金”“偿还债务支付的现金”用于刻画债务融资的规模 \\
|公司事实底稿 & 公司公告、官网、交易所问询与媒体报道 & 沿革、股权与控制权变动、融资历史、风险事件、战略动作与融资需求线索，\textbf{仅收录有来源链接的事实} \\
股本变动 & 上市公司公告 & 总股本、近三年股本变动次数与原因 \\
分红 & 上市公司公告 & 近三年现金分红次数与每股分红 \\
公告 & 上市公司公告 & 近三年 \num{41266} 条公告标题（2023 年 9 月---2026 年 10 月），按关键词分类为七类事件 \\
行业归属 & 申万行业分类（2021 版） & 一级/二级/三级行业归属 \\
\bottomrule
\end{tabular}
\end{table}

本部分的\emph{有息负债}定义为短期借款、一年内到期的非流动负债、长期借款与应付债券四项之和
（不含租赁负债、应付票据等经营性负债）；\emph{现金短债比}为货币资金除以短期债务
（短期借款加一年内到期的非流动负债）。这两个指标是对“未来融资需求”做量化判断的基础，
其局限在于：账面货币资金未必全部可自由动用（可能包含受限资金或募集资金专户），
短期债务也可能通过续贷滚动而非以现金清偿，因此下文给出的缺口是压力情形下的静态口径，
用于公司之间的横向比较，而不是对违约概率的估计。

需要说明三点口径限制：其一，\hl{本部分的市值一律采用\emph{总市值}口径}
（总市值 $=$ 最新收盘价 $\times$ 总股本，总股本取自巨潮股本结构），
不使用流通市值——流通市值只覆盖已上市流通的股份，
用它衡量一家公司在行业内的相对规模会系统性低估控股股东持股比例高的公司
（例如粮油加工与部分国企，其流通市值仅相当于总市值的十分之一）；
其二，公告分类基于标题关键词，
\hl{同一个事件可能被计入多个类别}（例如“关于向特定对象发行股票募集说明书”同时命中“再融资”），
故各类计数只用于横向比较公司之间的活跃度，不能相加为事件总数。
其三，\hl{半年报数据未年化}——半年度的 ROE 与净利率是 6 个月的经营结果，
与年报口径直接比较会系统性地低估盈利能力，因此本部分的当期判断一律在\emph{同口径}下进行
（2026 年上半年对比 2025 年上半年）；确需引用年报口径时，文中会标明“年报对照”或“年化”。
此外，本部分的“近三年”在公告与股本口径下指 2023 年 9 月至 2026 年 10 月；
“两年营收 CAGR”指 2023$\to$2025 年年报营业收入的年均复合增速，属趋势对照指标。

\subsection{行业分布：三个行业的公司数占了三分之一}
\label{sec:comp-industry}

\begin{figure}[htbp]
\centering
\begin{tikzpicture}
\begin{axis}[
  width=15.2cm, height=8.2cm,
  ybar, bar width=7pt,
  xmin=-0.7, xmax=17.7, ymin=0, ymax=13,
  xtick={0,1,2,3,4,5,6,7,8,9,10,11,12,13,14,15,16,17},
  xticklabels from table={data/clean/comp_ind_count.dat}{label},
  xticklabel style={rotate=35, anchor=east, font=\tiny, align=right, text width=1.9cm},
  ylabel={公司数量（家）}, ylabel style={font=\scriptsize},
  yticklabel style={font=\scriptsize},
  nodes near coords, nodes near coords style={font=\tiny},
  every node near coord/.append style={yshift=1pt},
  agri axis,
]
\addplot[fill=AgriGreen!75, draw=AgriGreen] table[x=i, y=n]{data/clean/comp_ind_count.dat};
\end{axis}
\end{tikzpicture}
\caption{申万农林牧渔行业上市公司的三级行业分布（共 \num{104} 家）}
\label{fig:comp-ind-count}
\end{figure}

三级行业分布呈现明显的\hl{“长尾 + 三个中等集群”}特征：生猪养殖、其他农产品加工、
畜禽饲料各 \num{11} 家，合计占全部公司的 31.7\%；动物保健 \num{9} 家、
肉鸡养殖与种子各 \num{8} 家。分布的另一端是海洋捕捞与粮食种植，各仅 \num{2} 家。
\emph{公司数量的结构本身就是周期结构的反映}——生猪产业链（养殖 + 饲料 + 动保）
合计 \num{35} 家，占 \num{33.7}\%，\hl{农业上市公司的资本回报
在很大程度上由猪周期决定}。

\subsection{市值分布：中间大、两头小}
\label{sec:comp-mktcap}

\begin{figure}[htbp]
\centering
\begin{tikzpicture}
\begin{groupplot}[
  group style={group size=2 by 1, horizontal sep=0.95cm},
  agri axis, width=6.3cm, height=5.0cm,
]
\nextgroupplot[
  % ymax 必须大于数据最大值（总市值分布最大 53 家、ROE 分布最大 22 家），
  % 否则最高的一根柱子被截断、数值标签会压到子图标题上
  ybar, bar width=9pt, ymin=0, ymax=58,
  xtick={0,1,2,3,4,5,6},
  xticklabels from table={data/clean/comp_mktcap.dat}{label},
  xticklabel style={rotate=25, anchor=east, font=\tiny},
  ylabel={公司数量（家）},
  nodes near coords, nodes near coords style={font=\tiny},
  title style={font=\small}, title={总市值分布},
]
\addplot[fill=OilBlue!70, draw=OilBlue] table[x=i, y=n]{data/clean/comp_mktcap.dat};
\nextgroupplot[
  ybar, bar width=9pt, ymin=0, ymax=58,
  xtick={0,1,2,3,4,5,6,7},
  xticklabels from table={data/clean/comp_roe.dat}{label},
  xticklabel style={rotate=30, anchor=east, font=\tiny},
  ylabel={公司数量（家）},
  nodes near coords, nodes near coords style={font=\tiny},
  title style={font=\small}, title={2026 上半年 ROE 分布},
]
\addplot[fill=SoftRed!70, draw=SoftRed] table[x=i, y=n]{data/clean/comp_roe.dat};
\end{groupplot}
\end{tikzpicture}
\caption{农业上市公司的总市值分布（左，亿元；市值时点 2026 年 10 月 12 日，\num{104} 家全部可估）与 2026 年上半年净资产收益率分布（右，\%；半年报披露值、未年化，\num{103} 家可算）}
\label{fig:comp-dist}
\end{figure}

市值分布是典型的\hl{“中间大、两头小”}：\num{104} 家公司中，
总市值在 \num{20}---\num{50} 亿元区间的最多，达 \num{53} 家，
而低于 \num{20} 亿元的仅 \num{5} 家、\num{100} 亿元以上的 \num{20} 家
（其中 \num{200} 亿元以上 \num{10} 家、\num{1000} 亿元以上 \num{3} 家）。
最大的一家是牧原股份，总市值 \num{2283.4} 亿元，
\hl{相当于市值排名后 60 家左右公司的市值之和}；金龙鱼 \num{1385.8} 亿元、
温氏股份 \num{1007.3} 亿元，前三家合计 \num{4676} 亿元，
占全部 \num{104} 家公司合计总市值（\num{12237} 亿元）的 \num{38.2}\%。
\hl{改用总市值口径后，行业规模结构的集中度比此前按流通市值衡量的更高}——
金龙鱼的可流通股本仅占总股本的十分之一，
按流通市值它只有 \num{138.7} 亿元、排在第 \zhnumber{13}位，
按总市值则排第 \zhnumber{2}位。

以 2026 年半年报衡量的盈利分布比年报口径更为严峻：\num{46} 家公司在 2026 年上半年录得亏损
（占 \num{44.2}\%），半年 ROE 的中位数仅 \num{0.75}\%（未年化），\num{44} 家公司半年 ROE 为负，
而半年 ROE 超过 \num{5}\%（约相当于年化 \num{10}\%）的只有 \num{16} 家。
需要提醒的是，半年 ROE 排名靠前的公司中，\textbf{新赛股份}（\num{30.8}\%）与
\textbf{獐子岛}（\num{29.5}\%）的净资产已被前期亏损侵蚀到很薄（2026 年 6 月末资产负债率
分别为 \num{82.7}\% 与 \num{97.0}\%），\emph{比率的分母过小使其失去了可比性}；
\textbf{中基健康}因净资产为负，半年 ROE 无法计算。
\hl{农业上市公司整体的资本回报水平已经低于无风险利率}，
这与第 \ref{sec:hog} 章刻画的全行业深度亏损相互印证——
\emph{上市公司是行业产能的放大器，而非对冲工具}。

\subsection{规模与盈利的组合：谁在下行周期中活得好}
\label{sec:comp-scatter}

\begin{figure}[htbp]
\centering
\begin{tikzpicture}
\begin{axis}[
  width=15.4cm, height=7.6cm,
  xlabel={2026 年上半年营业收入同比（\%）},
  ylabel={2026 年上半年销售净利率（\%）},
  xmin=-45, xmax=90, ymin=-60, ymax=40,
  xlabel style={font=\scriptsize}, ylabel style={font=\scriptsize},
  tick label style={font=\scriptsize},
  grid=both, grid style={InkGray!18},
  agri axis,
]
\addplot[only marks, mark=*, mark size=2.0pt, color=AgriGreen, opacity=0.85]
  table[x=growth, y=margin]{data/clean/comp_growth_scatter.dat};
% 象限分割线
\draw[InkGray, dashed, line width=0.5pt] (axis cs:0,-60) -- (axis cs:0,40);
\draw[InkGray, dashed, line width=0.5pt] (axis cs:-45,0) -- (axis cs:60,0);
\node[font=\tiny, AgriGreen, anchor=north west] at (axis cs:2,38) {成长 + 盈利};
\node[font=\tiny, SoftRed, anchor=north east] at (axis cs:-2,38) {收缩 + 盈利};
\node[font=\tiny, OilBlue, anchor=south west] at (axis cs:2,-58) {成长 + 亏损};
\node[font=\tiny, SoftRed, anchor=south east] at (axis cs:-2,-58) {收缩 + 亏损};
\end{axis}
\end{tikzpicture}
\caption{农业上市公司的成长与盈利组合（2026 年半年报：横轴为营业收入同比，纵轴为销售净利率；每点为一家公司；已剔除同比超出 $[-60,120]\%$、净利率超出 $[-80,40]\%$ 的极端值）}
\label{fig:comp-scatter}
\end{figure}

图上呈现出一条清晰的负相关带：\hl{靠规模扩张换取增长的公司，净利率被压到 0 附近甚至为负}。
营业收入规模最大的五家——金龙鱼（\num{1228.5} 亿元、同比 \num{+6.2}\%）、
海大集团（\num{642.4} 亿元、\num{+9.2}\%）、牧原股份（\num{594.1} 亿元、\num{-22.3}\%）、
新希望（\num{555.9} 亿元、\num{+7.7}\%）、温氏股份（\num{467.6} 亿元、\num{-6.2}\%）
——恰好代表两种结局：\textbf{金龙鱼}与\textbf{海大集团}以个位数的净利率
（\num{1.9}\% 与 \num{2.7}\%）守住薄利；\textbf{牧原}与\textbf{新希望}的净利率分别为
\num{-10.2}\% 与 \num{-3.0}\%，\hl{收入端并未同步收缩，利润端已经先行转负}——
这是“以价换量”在报表上的直接表现。
两者的差别不在于规模，而在于\hl{在周期低点是否还留有加杠杆扩张的能力}：
真正决定一家农业公司能否穿越周期的，不是它在景气高点借了多少，
而是它在行业普遍亏损、资产负债表最难看的时候，
是否还有能力（与意愿）以低成本资金去接别人的产能。
在周期高点加杠杆意味着把成本锁在最高点，而\emph{在周期低点加杠杆才有意义}——
牧原 2025 年年报的 ROE 为 \num{20.57}\%（年报对照），2026 年上半年降至 \num{-7.5}\%；
新希望 2026 年 6 月末的资产负债率为 \num{70.7}\%，
而按半年报口径有 \num{23} 家农业上市公司的资产负债率超过 \num{65}\%
（2025 年末为 \num{21} 家），
\emph{这些公司在本轮周期底部几乎没有再融资以外的选择}。

\subsection{周期读数：亏损面由 27.9\% 扩大至 44.2\%}
\label{sec:comp-h1}

本部分的财务口径统一为 2026 年半年报，因此把 2026 年上半年与 2025 年上半年的\emph{同口径}
数据并列，即可得到本轮周期在公司层面的最新读数：

\begin{itemize}[leftmargin=1.4em, itemsep=2pt]
  \item 亏损公司由 \num{29} 家增至 \hl{\num{46} 家}，占比由 \num{27.9}\% 升至 \num{44.2}\%；
  \item 归母净利润合计由 \num{270.9} 亿元转为 \hl{\num{-126.1} 亿元}，
        亏损公司的亏损额合计由 \num{24.8} 亿元扩大到 \num{226.9} 亿元；
  \item 亏损超过 \num{5} 亿元的公司由 \num{1} 家增至 \num{11} 家；生猪养殖 \num{11} 家公司的
        利润由上年同期的合计 \num{154.1} 亿元转为合计亏损 \num{163.8} 亿元；
  \item 2026 年 6 月末资产负债率超过 \num{65}\% 的公司为 \num{23} 家（2025 年末 \num{21} 家），
        其中同时处于亏损状态的 \num{14} 家、负债率超过 \num{80}\% 的 \num{11} 家（2025 年末 \num{7} 家）；
  \item 营业收入合计 \num{6085} 亿元、同比中位数 \num{+3.6}\%，
        需求端并未出现与利润同等幅度的收缩，
        与第一部分“\hl{本轮下行源于供给而非需求}”的判断一致；
  \item 结构上（两期对比，互斥分解、合计 \num{104} 家）：两期均盈利且同比改善 \num{19} 家、
        两期均盈利但同比下滑 \num{30} 家、扭亏 \num{9} 家、由盈转亏 \num{26} 家、两期均亏 \num{20} 家，
        \emph{分化本身是周期底部的常态}。
\end{itemize}

\begin{figure}[htbp]
\centering
\begin{tikzpicture}
\begin{axis}[
  agri axis, width=15.4cm, height=7.0cm,
  ybar, bar width=8pt,
  xmin=-0.6, xmax=5.6, ymin=0, ymax=42,
  xtick={0,1,2,3,4,5},
  xticklabels from table={data/clean/comp_h1_dist.dat}{label},
  xticklabel style={font=\footnotesize},
  xlabel={归母净利润区间（亿元）}, xlabel style={font=\scriptsize},
  ylabel={公司数量（家）}, ylabel style={font=\scriptsize},
  legend style={font=\scriptsize, at={(0.5,1.02)}, anchor=south, legend columns=2, draw=none},
]
\addplot[fill=AgriLight!70, draw=AgriLight] table[x=i, y=a]{data/clean/comp_h1_dist.dat};
\addplot[fill=SoftRed!70, draw=SoftRed] table[x=i, y=b]{data/clean/comp_h1_dist.dat};
\legend{2025 年上半年, 2026 年上半年}
\end{axis}
\end{tikzpicture}
\caption{农业上市公司 2025 年与 2026 年上半年的归母净利润分布对比（横轴为归母净利润区间，亿元；两期均为半年度数据，口径可比）}
\label{fig:comp-h1-dist}
\end{figure}

分行业看，粮油加工（\num{25.5} 亿元）、水产饲料（\num{15.9} 亿元）、
动物保健（\num{8.8} 亿元）与宠物食品（\num{3.4} 亿元）四个子行业的盈利总额合计
\num{53.5} 亿元；需要说明的是，这四个数字是\emph{子行业口径}的盈利总额，
与前一句 \num{28} 家的\emph{公司口径}盈利改善数不是同一统计对象。
\hl{利润的分布比年报口径更集中、也更极端}，这直接决定了各公司融资需求的紧迫程度——
第 \ref{sec:planting} 章至第 \ref{sec:fishery} 章的逐公司画像小节给出了具体判断。

\begin{keybox}[第二部分的四条主线]
\normalsize
\textbf{① 结构的两个“三分之一”}：生猪产业链公司数与农牧加工+饲料公司数各占约三分之一，
前者决定行业整体的盈利波动，后者决定行业整体的规模。前三家市值合计
\hl{占全部 \num{104} 家公司合计总市值的 \num{38.2}\%}。\par\vspace{3pt}
\textbf{② 盈利的接近一半在亏损}：按 2026 年半年报口径，\num{46} 家亏损（\num{44.2}\%）、
\num{44} 家半年 ROE 为负，ROE 中位数 \num{0.75}\%（未年化）。
这不是个别公司的问题，而是周期底部的\emph{分布特征}。\par\vspace{3pt}
\textbf{③ 资本结构的硬约束}：\num{23} 家公司在 2026 年 6 月末的资产负债率超 \num{65}\%，
其中 \num{11} 家超 \num{80}\%，
\hl{本轮周期中“活下去”取决于资产负债表，而不是成本控制}——
这构成了第 \ref{sec:financing} 节讨论未来融资需求的起点。\par\vspace{3pt}
\textbf{④ 周期底部比年报口径更晚、更深}：与 2025 年年报（\num{34} 家亏损）相比，
2026 年半年报的亏损面扩大到 \num{46} 家，生猪养殖板块由盈利 \num{154.1} 亿元
转为亏损 \num{163.8} 亿元，\hl{报表层面的底部尚未出现}。
\end{keybox}

\section{种植业：种子、粮食种植、食用菌与其他种植}
\label{sec:planting}

种植业板块共 \num{20} 家公司（种子 \num{8}、其他种植业 \num{6}、粮食种植 \num{2}、食用菌 \num{4}），
总市值合计约 \num{1581} 亿元（市值时点 2026 年 10 月 12 日）；按 2026 年半年报口径，\num{6} 家亏损、
半年 ROE 中位数 \num{1.43}\%（未年化），2026 年 6 月末资产负债率超过 \num{65}\% 的仅 \num{2} 家。
\hl{这是农业各子行业中最分散、也最“稳”的一块}——
原因在于种植业的资产以土地与生物资产为主，负债率天然偏低，
但盈利高度依赖单一年度的天气与政策价格。
需要提示的是，\emph{种植业的收入与利润确认季节性较强}：半年报反映的是
上半年（播种与生长期）的经营结果，因此本节的比较一律采用同口径
（2026 年上半年对 2025 年上半年），不与年报数据直接并列。

\begingroup\scriptsize\setlength{\tabcolsep}{2pt}
\begin{longtable}{@{}llrrrrrrr@{}}
\caption{农业上市公司分行业明细（种子、粮食种植、其他种植业、食用菌；市值时点 2026-10-12，总市值口径；财务为 2026 年半年报，ROE 未年化，公告计数窗口 2023-09---2026-10）}\label{tab:comp-planting}\\
\toprule
代码 & 公司简称 & 总市值（亿元） & 营收（亿元） & 同比（\%） & 净利（亿元） & ROE（\%） & 再融资公告 & 重组公告 \\
\midrule
\endfirsthead
\multicolumn{9}{l}{\small（续）}\\
\toprule
代码 & 公司简称 & 总市值（亿元） & 营收（亿元） & 同比（\%） & 净利（亿元） & ROE（\%） & 再融资公告 & 重组公告 \\
\midrule
\endhead
\bottomrule
\endlastfoot
601118 & 海南橡胶 & 319.2 & 151.2 & -25.4 & -1.30 & -1.3 & 2 & 2 \\
600598 & 北大荒 & 226.3 & 21.3 & -29.4 & -5.37 & -7.1 & 0 & 1 \\
601952 & 苏垦农发 & 123.5 & 43.3 & -5.6 & 1.44 & 2.0 & 0 & 3 \\
000998 & 隆平高科 & 113.8 & 17.1 & -21.1 & -2.73 & -4.1 & 47 & 4 \\
002215 & 诺普信 & 94.9 & 39.6 & 7.8 & 7.65 & 16.8 & 7 & 0 \\
002041 & 登海种业 & 87.1 & 4.7 & 27.0 & 0.50 & 1.3 & 0 & 0 \\
600108 & 亚盛集团 & 78.5 & 14.7 & -2.6 & 0.25 & 0.6 & 0 & 2 \\
600313 & 农发种业 & 72.6 & 38.7 & 32.8 & 0.53 & 2.2 & 32 & 3 \\
300189 & 神农种业 & 63.1 & 0.6 & -35.3 & -0.13 & -1.9 & 14 & 1 \\
603336 & 宏辉果蔬 & 59.6 & 4.2 & -11.1 & 0.05 & 0.4 & 3 & 5 \\
300087 & ST荃银 & 52.9 & 11.8 & -17.9 & -0.55 & -3.3 & 1 & 18 \\
600354 & 敦煌种业 & 49.1 & 8.7 & 20.9 & 1.45 & 18.7 & 1 & 7 \\
002772 & 众兴菌业 & 41.4 & 10.9 & 19.6 & 1.91 & 5.3 & 1 & 0 \\
600371 & 万向德农 & 41.0 & 1.0 & -15.1 & 0.14 & 2.6 & 0 & 3 \\
000713 & 国投丰乐 & 38.1 & 11.9 & 3.9 & -0.31 & -1.0 & 45 & 7 \\
600540 & 新赛股份 & 30.4 & 25.7 & 77.2 & 1.25 & 30.8 & 1 & 1 \\
600359 & 新农开发 & 27.3 & 2.0 & -31.1 & 0.19 & 2.8 & 0 & 7 \\
300511 & 雪榕生物 & 23.7 & 11.6 & 46.2 & 0.18 & 1.1 & 29 & 0 \\
300970 & 华绿生物 & 20.0 & 6.9 & 39.2 & 0.46 & 2.8 & 0 & 1 \\
600883 & 博闻科技 & 18.6 & 0.3 & -8.8 & 0.15 & 1.5 & 1 & 0 \\
\end{longtable}
\endgroup


\subsection{盈利结构：土地资产的租金回报}
\label{sec:planting-earnings}

板块内上半年盈利最强的仍是两家以土地为核心资产的公司：\textbf{苏垦农发}
（归母净利润 \num{1.44} 亿元、半年 ROE \num{2.04}\%）与
\textbf{亚盛集团}（\num{0.25} 亿元、\num{0.58}\%）；
而总市值最高的 \textbf{北大荒}（\num{226.3} 亿元）录得 \num{-5.37} 亿元亏损
（上年同期为 \num{+9.84} 亿元），营收 \num{21.28} 亿元、同比 \num{-29.4}\%。
\hl{北大荒的亏损不能读作主业恶化，而应读作一次性的税负冲击}：
2026 年 6 月 18 日公司收到哈尔滨市税务局第一税务分局《税务事项通知书》，
认定下属 \num{16} 家农业分公司从非职工家庭农场取得的土地承包费收入
不应享受企业所得税优惠政策，需补缴 2021---2025 年企业所得税 \num{10.24} 亿元
与滞纳金 \num{3.86} 亿元，合计 \num{14.10} 亿元；
公司说明该项补缴不涉及税务行政处罚、\emph{不属于前期会计差错}、不追溯调整前期财务数据，
\hl{税款与滞纳金一次性计入 2026 年当期损益}，且已于上半年内补缴完毕。
\emph{补缴总额超过了公司 2025 年全年的归母净利润（\num{11.66} 亿元）}，
这一次性支出直接解释了净利从 \num{+9.84} 亿元到 \num{-5.37} 亿元的落差：
2026 年上半年公司扣非净利润为 \num{-2.09} 亿元，
6 月末资产负债率由上年末的 \num{12.9}\% 升至 \num{39.5}\%，也主要来自这笔一次性现金流出
（当期“合同负债”因预收全年土地承包费而大幅增加，账上并不缺经营性现金）。

剥离这一次性因素后，北大荒的主业仍符合\hl{“土地租金 + 规模种植”}的逻辑：
土地承包费收入长期占营业收入七成以上，2026 年上半年占比 \num{91.3}\%；
在年报口径下，2025 年营收 \num{52.3} 亿元对应 \num{11.7} 亿元净利润、
净利率高达 \num{22.3}\%、ROE \num{14.9}\%，远高于行业水平。
这与第 \ref{sec:grain} 章中“玉米、小麦在成本线附近波动、
种植环节利润主要沉淀为地租”的判断一致。
需要提示的是，地租型业务的收入确认方式使单半年度报表容易失真：
2026 年上半年营业收入同比 \num{-29.4}\%，主因是化肥等农业投入品
由购销模式改为代销模式、销售收入因此减少 \num{8.86} 亿元；
而毛利率仍有 \num{54.9}\%、净利率却是 \num{-25.4}\%，
\hl{这两个数字之间的 \num{80} 个百分点的缺口，正是 \num{14.10} 亿元补缴税款的落点}。
因此土地类公司的盈利能力应以年报口径并结合一次性事项判断，
半年报只用于观察边际方向，不能外推全年。

\textbf{诺普信}（\num{7.65} 亿元、半年 ROE \num{16.75}\%）是板块中盈利最强的非土地型公司，
\hl{其盈利来源已经不是传统农药，而是蓝莓}：2026 年上半年特色生鲜消费业务收入
\num{20.70} 亿元（同比 \num{+13.87}\%）、占营业收入的 \num{52.2}\%，
核心子公司爱莓庄农业集团上半年收入 \num{20.18} 亿元、净利润 \num{6.07} 亿元，
是公司利润的绝对主力。
蓝莓能在大宗农产品普遍承压的年份贡献利润，原因在于这条业务线同时改写了三个变量：
\emph{以基质栽培替代露地种植}，把单产与品质从天气与土壤中解放出来；
\emph{以云南等低纬高原基地实现错季上市}，把销售窗口从进口果与国产果的集中上市期错开，
从而拿到时间溢价；\emph{以京东、山姆会员店、胖东来等现代渠道直达消费者}，
压缩流通环节的加价层级。与此同时，农药制剂业务 2026 年上半年收入 \num{15.24} 亿元、
同比 \num{+2.43}\%、占比 \num{38.5}\%，
\hl{在制剂需求整体偏弱的年份保持了稳定}，二者构成
“高弹性的生鲜增量 + 稳定的制剂基本盘”的组合。
这条第二曲线的风险同样明确：实际控制人卢柏强的股权质押规模较大，
公司原计划的 \num{14.5} 亿元定增（蓝莓基地扩建 \num{11.0} 亿元、
小浆果国际研发中心 \num{1.5} 亿元、补流 \num{2.0} 亿元）
\hl{已于 2026 年 7 月终止}，转而以银行长期授信替代股权融资，
实际控制人的个人质押也已清零；这意味着\emph{蓝莓产能的扩张只能靠内生现金流}——
在新增股权资金缺席的条件下，\num{3.6} 万亩存量基地的错季溢价能否维持，
是判断这条第二曲线可持续性的关键。
种子板块中 \textbf{敦煌种业}（\num{1.45} 亿元、\num{18.70}\%）
与 \textbf{新赛股份}（\num{1.25} 亿元、\num{30.85}\%）的半年 ROE 更高，
但前者来自低基数下的利润修复，后者的净资产已被前期亏损侵蚀到很薄
（2026 年 6 月末资产负债率 \num{82.7}\%），\emph{均不构成盈利质量改善的证据}。
种子板块整体承压：\num{8} 家公司中 \num{4} 家亏损，
\textbf{隆平高科}亏损 \num{2.73} 亿元、\textbf{ST荃银}亏损 \num{0.55} 亿元、
\textbf{国投丰乐}亏损 \num{0.31} 亿元、神农种业亏损 \num{0.13} 亿元，
仅有 \textbf{登海种业}（\num{0.50} 亿元）与 \textbf{万向德农}（\num{0.14} 亿元）维持微利。
\hl{种业的困境是结构性的，但供给端的去化已经开始}。
玉米、水稻种子同质化严重，审定品种数量逐年增加而单品种推广面积被稀释，
行业曾长期陷入“品种数量增加、单品种收入下降”的通缩，
这与粮食下游加工环节的价格传导不畅互为因果。
\hl{变化发生在 2025 年}：全国玉米制种面积自 2023 年的 \num{462} 万亩
（十年高点）连续三年调减——2024 年 \num{428} 万亩、2025 年 \num{385} 万亩，
2026 年进一步降至 \num{358.5} 万亩、同比再减 \num{7}\%，较峰值回落约 \num{24}\%
（全国农业技术推广服务中心 2026 年 9 月发布的\emph{2026---2027 年度种子产供需形势}）。
龙头企业同样在收缩制种：\textbf{登海种业} 2025 年度制种产量 \num{7220.32} 万公斤，
同比减少 \num{4382.02} 万公斤，公司在半年报中明确“适当调减制种面积、量入为出、以销定产”；
\textbf{隆平高科} 2026 年上半年亏损扩大的原因被公司归为“主动推进产品结构性优化、
加大去库存消化力度”，即降价去库存。
\hl{但去化是结构性的，而不是齐步走}：2026 年有 \num{3} 家头部企业主推品种的制种面积
逆势扩张，另 \num{3} 家因品种市场萎缩下降超过 \num{50}\%，
中小企业普遍下调 \num{10}\%---\num{35}\%，
\emph{玉米制种继续向西北优势产区集中（占全国 \num{77}\%）}——
供给收缩与集中度提升同时发生，出清之后受益的是有品种优势的头部企业，
而不是所有种子公司。
\hl{但去化不等于出清}：2025 年度全国玉米种子总供给约 \num{22.8} 亿公斤、
市场需求约 \num{13} 亿公斤，供需比仍高达约 \num{175}\%，库存压力巨大；
上市种子公司的库存总额从 2025 年三季度的 \num{116} 亿元高点开始回落，
2026 年一季度库存同比增速才转为负值。
因此正确的判断是：\emph{产能去化已经开始，价格弹性要等到库存出清之后才会显现}。

\subsection{天然橡胶与食用菌：两个极端}
\label{sec:planting-extremes}

\textbf{海南橡胶}（总市值 \num{319.2} 亿元）是“其他种植业”中总市值最大的公司，
2026 年上半年营收 \num{151.2} 亿元（同比 \num{-25.4}\%）、亏损 \num{1.30} 亿元，
2026 年 6 月末资产负债率 \num{70.7}\%。
它对应第 \ref{sec:soft} 章中天然橡胶的全球定价逻辑：
中国天然橡胶种植成本长期高于东南亚产区，
\hl{公司盈利由全球供需缺口与汇率决定，与国内农业周期基本无关}。
天然橡胶主力合约 2021 年 2 月至 2026 年 9 月累计上涨 \num{24.8}\%、
20 号胶上涨 \num{43.3}\%，是第 \ref{sec:soft} 章讨论的软商品中少数趋势性上行的品种，
但截至 2026 年上半年，\hl{价格上涨尚未传导到公司的利润表}。

另一端是食用菌：\num{4} 家公司全部实现盈利，其中 \textbf{众兴菌业}（\num{1.91} 亿元、
半年 ROE \num{5.34}\%）、\textbf{华绿生物}（\num{0.46} 亿元、\num{2.79}\%）、
\textbf{雪榕生物}（\num{0.18} 亿元，上年同期为亏损 \num{1.03} 亿元）
与 \textbf{博闻科技}（\num{0.15} 亿元）；子行业营收同比中位数 \num{+29.4}\%（半年报口径），
是板块中唯一量、利同时改善的方向。
食用菌是板块中\hl{唯一“短周期化”的子行业}——
工厂化生产的产能扩张周期仅 \num{6}---\num{9} 个月，
价格随产能快速调整，龙头企业以成本优势取得稳定的加工利润。

\begin{keybox}[种植业的画像结论]
\normalsize
\textbf{① 土地类公司的利润波动来自确认时点。}苏垦农发上半年归母净利润 \num{1.44} 亿元、
半年 ROE \num{2.04}\%（未年化），北大荒则由上年同期的 \num{+9.84} 亿元转为 \num{-5.37} 亿元；
按年报口径，两家公司 2025 年的 ROE 分别为 \num{7.9}\% 与 \num{14.9}\%（趋势对照）。
\hl{判断土地资产的回报仍应以年报为主，半年报只用于观察边际方向}。\par\vspace{3pt}
\textbf{② 种业的产能去化已经开始，但尚未出清。}\num{8} 家种子公司中 \num{4} 家亏损
（2026 年上半年合计亏损 \num{3.7} 亿元）；全国玉米制种面积已从 2023 年的
\num{462} 万亩连续三年下降至 2026 年的 \num{358.5} 万亩（\num{-24}\%），
龙头企业制种产量同步收缩，\hl{供给端的调整已经启动}；
但 \num{175}\% 的供需比与超过 \num{100} 亿元的公司库存意味着
\hl{去库存仍要一到两个业务年度才能见底}，在这之前种业股缺乏周期弹性。\par\vspace{3pt}
\textbf{③ 天然橡胶是唯一“逆周期”的种植品种}（市值第一、利润仍为负），
2026 年上半年价格的历史涨幅尚未体现在报表中，
其价格弹性要等到 2026---2027 年才可能兑现。\par\vspace{3pt}
\textbf{④ 板块利润继续收缩。}板块 \num{20} 家公司 2026 年上半年合计归母净利润
\num{5.73} 亿元（上年同期 \num{15.87} 亿元），\num{6} 家亏损，
\hl{种业与天然橡胶的疲弱尚未反转}。
\end{keybox}

\subsection{上市公司画像（\num{20} 家）}
\label{sec:co-planting}

本节给出种植业全部 \num{20} 家上市公司的画像，按总市值降序排列。
数据口径与第 \ref{sec:comp-overview} 节一致：\hl{财务数据统一采用 2026 年半年报}
（合并报表口径、未年化，资产负债率为 2026 年 6 月末），
同时列出 2025 年半年报作为同口径对照、2023---2025 年年报作为趋势对照；
股价为月度收盘价，最后一个交易日为 2026 年 10 月 12 日；周期分段使用与第一部分相同的
ZigZag 算法，阈值取 25\%（个股波动大于商品价格，故高于生猪现货的 20\%）；
公告计数窗口为 2023 年 9 月 1 日至 2026 年 10 月 12 日。

与前一版不同，本版画像正文由\hl{逐家检索的公开资料底稿}驱动：底稿只收录有来源链接的事实
（公司沿革、股权与控制权、经营要点、所处行业的基本面、竞争优势、融资历史、
风险事件、战略动作、融资需求线索），
正文按事实生成，各家公司的段落结构与重点并不相同。
每家公司在统一顺序下展开：基本情况与主营结构、股权结构与控制权、经营要点、
所处行业的基本面（供需）、与同类公司比较的竞争优势、历史股价、过去周期分析、
盈利情况、财务分析、重大战略转型、历史融投资情况、未来潜在融投资需求。
\hl{市值一律采用总市值口径}（收盘价 $\times$ 总股本），股权结构取最新披露报告期的
前三大股东与股东户数；行业基本面与竞争优势以逐家检索到的公开资料为依据，
并与三级行业内的同口径指标（营业收入、毛利率、ROE、资产负债率）横向比较。
每家公司配图两张：
股价与行业指数（申万农林牧渔）的对比图，以及近年主要财务指标的组合图
（上排：营业收入与归母净利润、ROE 与资产负债率；下方宽带面板：货币资金与有息负债为左轴柱状、
经营活动现金流净额为其右轴的蓝点折线），并给出关键财务指标表与历史融资/资本运作表。
未来融资需求按两个口径粗略度量：\hl{短期债务缺口 $=$ 期末短期债务 $-$ 货币资金}
与\hl{年化经营缺口 $=$ 半年归母亏损 $\times 2$}；
前者衡量滚动偿付压力，后者衡量维持当期盈利水平的现金消耗，
两者都是压力情形下的静态估算，不构成对融资规模或时点的预测。

% 农业上市公司画像：种子、粮食种植、其他种植业、食用菌（共 20 家，按总市值降序）
% 由 scripts/make_company_profiles.py 自动生成，请勿手工修改

\subsubsection{海南橡胶（601118）}
\label{co:601118}

\pfl{公司基本情况} 2005 年设立至今，海南橡胶是全球最大的天然橡胶种植、加工与贸易企业之一；长期处于扣非亏损状态，盈利对胶价与汇率高度敏感，并以举债和海外并购支撑规模扩张。 公司于 2011-01-07 在上交所首发上市，注册资本 42.79 亿元，总股本 42.79 亿股。 截至 2026 年 10 月 12 日收盘价 7.46 元，对应总市值 319.2 亿元，近一年+37.6\%，总市值在三级行业“其他种植业”的 6 家公司中按由大到小排第\zhnumber{1}位，近一年涨跌幅高于全样本中位数（-10.4\%）48.0 个百分点。 按 2025 年年报披露的行业构成，收入占比前两位为农业 96.8\%；其他(补充) 3.2\%。 股权结构方面，控股股东为海南省农垦投资控股集团有限公司（持股 64.35\%）。

\pfl{股权结构与控制权} 按 2026 年 6 月 30 日的股东名册，前三大股东为海南省农垦投资控股集团有限公司 持股 64.35\%；香港中央结算有限公司（持股比例未披露）；中国银行股份有限公司－大成产业趋势混合型证券投资基金 持股 0.86\%，股东户数 97,198 户。 控制权与股东性质方面，海垦控股集团承诺在 2027 年 12 月 20 日前将体系外经营性橡胶相关资产市场化注入上市公司。

\pfl{经营要点} 从经营要点看，全球拥有 57 家在产橡胶初加工厂（含受托管理的 KM 公司），年加工产能约 217 万吨；2026 年上半年，自产橡胶产量 2.23 万吨、同比增产 1,100 吨。

\pfl{所处行业的基本面} 该子行业按品种各自定价：天然橡胶看全球供给与汽车需求，棉花看新疆产量与配额，农药看原药价格与农业服务需求（第 \ref{sec:framework} 章的供需框架）。 从供需两端的当期数据看，ANRPC 预测 2026 年全球天然橡胶产量 1,531.0 万吨、同比增长 2.3\%；较年初 161.99 万吨下降约 20\%；2026 年 10 月，国内现货均价 1.97 万元/吨、年内涨幅 28.27\%、同比上涨 31.45\%。

\pfl{竞争优势} 把该公司放回同类公司比较，2026 年上半年营业收入 151.2 亿元、在三级行业“其他种植业”的 6 家公司中排第\zhnumber{1}位；毛利率 3.4\%、低于行业内中位数 13.5 个百分点；半年 ROE -1.3\%（未年化）、低于行业内中位数（1.70\%）；资产负债率 70.7\%、高于行业内中位数（59.7\%）。 公司在公开资料中自述与同业比较的相对优势在于：下属经纬公司乳胶丝年产能约 2 万吨、枕头年产能 360 万只。 就当期报表看，公司尚未把经营层面的特征转化为高于行业的资本回报，周期底部的相对优势更多体现在现金流与资产负债表上。

\pfl{历史股价} 以 2021 年 1 月为起点、2026 年 10 月为终点，股价由 5.45 元变为 7.46 元，累计 +36.9\%，区间最高 7.68 元（2026 年 2 月）、最低 4.04 元（2022 年 10 月）；当前价相当于区间最高价的 97.1\%，为区间最低价的 1.85 倍。 月度收益的年化波动率 29.7\%，低于全样本中位数 37.3\%，在三级行业“其他种植业”的 6 家公司中波动率由高到低排第\zhnumber{4}位。 把股价阶段与公司事件对齐看，2022 年 10 月 至 2026 年 2 月股价上涨+90.1\%，同期（2022 年 12 月）公司将所持北京海垦商贸发展有限公司 29.34\% 股权全部转让给海垦控股集团。 从时间对应关系看，公司层面的资本运作与股权、风险事件解释了区间内多数急涨急跌，与行业指数的同步性相对有限。

\begin{center}
\begin{minipage}{\textwidth}\centering
\begin{tikzpicture}
\begin{axis}[
  width=12.6cm, height=3.9cm, scale only axis,
  xmin=2020.922, xmax=2026.972, ymin=3.76, ymax=8.22,
  xtick={2021,2022,2023,2024,2025,2026},
  yearticks,
  yticklabel style={font=\scriptsize},
  ylabel={元/股}, ylabel style={font=\scriptsize},
  grid=major, grid style={LightGray, line width=0.3pt},
  axis line style={InkGray, line width=0.5pt},
  every axis plot/.append style={line width=0.9pt},
]
\addplot[AgriGreen] table[x=x,y=y]{data/clean/profiles/px_601118.dat};
\addplot[SoftRed, dashed, line width=0.7pt] table[x=x,y=y]{data/clean/profiles/pxzz_601118.dat};
\addplot[only marks, mark=*, mark size=1.1pt, SoftRed] table[x=x,y=y]{data/clean/profiles/pxzz_601118.dat};
\end{axis}
\end{tikzpicture}

\captionof{figure}[海南橡胶（601118）股价与周期骨架]{海南橡胶（601118）月度收盘价（元/股）与 ZigZag 周期骨架（阈值 25\%，圆点为周期转折点）}
\end{minipage}
\end{center}

\pfl{过去周期分析} 以 25\% 的 ZigZag 阈值划分，样本区间内股价形成 2 个主升段与 3 个主跌段。 最大主跌段为 2026 年 2 月 至 2026 年 6 月，4 个月-35.9\%；最大主升段出现在 2022 年 10 月 至 2026 年 2 月，40 个月+90.1\%。 最近一次转折出现在 2026 年 10 月（上涨段自 2026 年 6 月起，4 个月+51.6\%），当前处于该段之后的震荡之中。 公司股价较申万农林牧渔指数提前约 43 个月见底，是板块中较早定价周期反转的公司之一。 该子行业的周期来源与猪周期不同（第 \ref{sec:grain} 章、第 \ref{sec:soft} 章），价格更多由自身供给、进口政策与全球定价决定，公司 2026 年 6 月末资产负债率 70.7\%，高于全样本中位数 52.2\%，其股价因此更多反映资金面与信用风险，而不是价格弹性本身。 综合区间形态看，该股单段涨跌幅度偏大、以趋势段为主（最大主升 +90.1\%、最大主跌 -35.9\%），年化波动率 29.7\% 与全样本中位数 37.3\% 相比波动小于样本中位数。

\pfl{盈利情况} 2026 年上半年营业收入 151.23 亿元（同比 -25.4\%），归母净利润 -1.30 亿元，亏损同比收窄 0.46 亿元。 以年报为趋势对照，2023---2025 年营业收入由 334.67 亿元变为 422.37 亿元（两年年均复合增速 +12.3\%），归母净利润由 2.97 亿元变为 -1.03 亿元。 按半年报口径，2026 年上半年的 ROE 为 -1.3\%（未年化）、销售净利率 -1.8\%、毛利率 3.4\%，ROE 在“其他种植业”6 家公司中按数值由高到低排第\zhnumber{6}位。 按同一口径，三级行业“其他种植业”6 家公司的半年 ROE 中位数为 1.70\%，该公司低于其中位数 3.00 个百分点。 从公司披露的原因看，业绩变动主要来自：2024 年，公司亏损主因贷款规模较高、外币借款利息与汇兑损失合计约 5.17 亿元；2026 年 6 月，公司亏损主要受美元汇率下跌导致汇兑损失增加、金融工具平仓损失、合同纠纷判决款等偶发性支出及递延所得税费用增加影响。 公司披露的应对举措集中在：2026 年，公司推进老挝锦森二期项目落地、柬埔寨公司设立，扩大橡胶加工产能。 综合看，公司仍处亏损区间、亏损同比收窄，半年 ROE 在三级行业内按由高到低排第\zhnumber{6}位，单位盈利能力的修复依赖行业景气与自身成本端的改善，报表弹性主要体现为价格与销量的方向性变化。

\begin{center}\small\setlength{\tabcolsep}{3.4pt}
\captionof{table}[海南橡胶（601118）关键财务指标]{海南橡胶（601118）关键财务指标（合并报表口径；两个半年报期均未年化；现金短债比 $=$ 货币资金 $\div$ 短期债务）}
\begin{adjustbox}{max width=\textwidth}
\begin{tabular}{@{}lrrrrr@{}}
\toprule
指标 & 2023 年 & 2024 年 & 2025 年 & 2025 年半年报 & 2026 年半年报 \\
\midrule
营业收入（亿元） & 334.67 & 420.04 & 422.37 & 202.72 & 151.23 \\
归母净利润（亿元） & 2.97 & 1.03 & -1.03 & -1.76 & -1.30 \\
毛利率（\%） & 2.4 & 3.8 & 2.5 & 3.8 & 3.4 \\
ROE（\%） & 3.1 & 1.1 & -1.0 & -1.8 & -1.3 \\
资产负债率（\%） & 65.7 & 68.7 & 70.3 & 69.0 & 70.7 \\
经营活动现金流净额（亿元） & 12.05 & 15.55 & 25.05 & 20.90 & 10.06 \\
货币资金（亿元） & 35.06 & 56.88 & 52.69 & 72.76 & 51.31 \\
有息负债合计（亿元） & 165.56 & 177.92 & 172.30 & 182.33 & 167.80 \\
现金短债比（倍） & 0.33 & 0.55 & 0.50 & 0.66 & 0.55 \\
\bottomrule
\end{tabular}
\end{adjustbox}
\end{center}

\begin{center}
\begin{minipage}{\textwidth}\centering
\begin{tikzpicture}
\begin{groupplot}[
  group style={group size=2 by 1, horizontal sep=0.60cm},
  width=6.85cm, height=4.60cm,
  tick label style={font=\tiny},
  title style={font=\scriptsize\heiti},
  yticklabel style={font=\tiny},
  ylabel style={font=\tiny},
  grid=major, grid style={LightGray, line width=0.3pt},
  axis line style={InkGray, line width=0.5pt},
  enlarge x limits={abs=0.8},
  legend style={font=\scriptsize, at={(0.5,-0.20)}, anchor=north,
                legend columns=2, draw=none, fill=none, column sep=6pt},
]
\nextgroupplot[
  ybar, bar width=4pt,
  ymin=-43.67, ymax=473.21,
  xtick={0,1,2,3,4,5.7},
  xticklabels={2021,2022,2023,2024,2025,2026H1},
  ylabel={亿元},
]
\addplot[fill=AgriGreen!75, draw=AgriGreen, bar shift=-2.6pt] table[x=i, y=rev]{data/clean/profiles/fin_601118.dat};
\addplot[fill=SoftRed!60, draw=SoftRed, bar shift=2.6pt] table[x=i, y=ni]{data/clean/profiles/fin_601118.dat};
\addplot[fill=AgriGreen!30, draw=AgriGreen, pattern=north east lines, bar shift=-2.6pt] table[x=x, y=rev]{data/clean/profiles/finh_601118.dat};
\addplot[fill=SoftRed!20, draw=SoftRed, pattern=north east lines, bar shift=2.6pt] table[x=x, y=ni]{data/clean/profiles/finh_601118.dat};
\legend{营业收入（年/半年）, 归母净利润（年/半年）}
\nextgroupplot[
  ymin=-7.1, ymax=77.9,
  xtick={0,1,2,3,4,5.7},
  xticklabels={2021,2022,2023,2024,2025,2026H1},
  ylabel={\%},
]
\addplot[AgriGold, mark=*, mark size=1.3pt] table[x=i, y=roe]{data/clean/profiles/fin_601118.dat};
\addplot[OilBlue, mark=square*, mark size=1.3pt] table[x=i, y=debt]{data/clean/profiles/fin_601118.dat};
\addplot[only marks, AgriGold, mark=*, mark size=2.0pt] table[x=x, y=roe]{data/clean/profiles/finh_601118.dat};
\addplot[only marks, OilBlue, mark=square*, mark size=2.0pt] table[x=x, y=debt]{data/clean/profiles/finh_601118.dat};
\legend{ROE, 资产负债率}
\end{groupplot}
\end{tikzpicture}
\begin{tikzpicture}
\begin{axis}[
  name=finbot,
  width=13.75cm, height=3.10cm, scale only axis,
  ymin=-17.79, ymax=199.27,
  xtick={0,1,2,3,4,5.7},
  xticklabels={2021,2022,2023,2024,2025,2026H1},
  tick label style={font=\tiny},
  yticklabel style={font=\tiny},
  ylabel={亿元}, ylabel style={font=\tiny},
  ybar, bar width=4pt,
  grid=major, grid style={LightGray, line width=0.3pt},
  axis line style={InkGray, line width=0.5pt},
  enlarge x limits={abs=0.8},
  legend style={font=\scriptsize, at={(0.5,-0.26)}, anchor=north,
                legend columns=2, draw=none, fill=none, column sep=10pt},
]
\addplot[fill=AgriGreen!55, draw=AgriGreen, bar shift=-3.0pt] table[x=i, y=cash]{data/clean/profiles/fin_601118.dat};
\addplot[fill=SoftRed!45, draw=SoftRed, bar shift=3.0pt] table[x=i, y=ibd]{data/clean/profiles/fin_601118.dat};
\addplot[fill=AgriGreen!25, draw=AgriGreen, pattern=north east lines, bar shift=-3.0pt] table[x=x, y=cash]{data/clean/profiles/finh_601118.dat};
\addplot[fill=SoftRed!20, draw=SoftRed, pattern=north east lines, bar shift=3.0pt] table[x=x, y=ibd]{data/clean/profiles/finh_601118.dat};
\legend{货币资金（左轴）, 有息负债（左轴）}
\end{axis}
\begin{axis}[
  at={(finbot.south east)}, anchor=south east,
  width=13.75cm, height=3.10cm, scale only axis,
  axis y line*=right, axis x line*=none, xtick=\empty,
  ymin=-6.51, ymax=32.56,
  tick label style={font=\tiny},
  yticklabel style={font=\tiny}, scaled y ticks=false,
  legend style={font=\scriptsize, at={(0.5,-0.44)}, anchor=north,
                legend columns=1, draw=none, fill=none},
]
\addplot[OilBlue, mark=*, mark size=2.2pt, line width=1.3pt] table[x=i, y=ocf]{data/clean/profiles/fin_601118.dat};
\addplot[only marks, OilBlue, mark=*, mark size=3.2pt] table[x=x, y=ocf]{data/clean/profiles/finh_601118.dat};
\legend{经营活动现金流净额（右轴，亿元，蓝点）}
\end{axis}
\end{tikzpicture}

\captionof{figure}[海南橡胶（601118）近年主要财务指标]{海南橡胶（601118）近年主要财务指标（左上：营业收入与归母净利润；右上：ROE 与资产负债率；下：货币资金、有息负债为左轴柱状，经营活动现金流净额为其右轴的蓝点折线。
2021---2025 年为年报口径，2026H1 为 2026 年半年报/半年末，与其前一年报之间留出间隔、以斜纹或空心点区分；半年报数据未年化）}
\end{minipage}
\end{center}

\pfl{财务分析} 2026 年 6 月末资产负债率 70.7\%，较 2025 年末上升 0.4 个百分点（样本中位数 52.2\%，所处三级行业内按由低到高排第\zhnumber{5}位）。 2026 年 6 月末货币资金 51.31 亿元、短期债务（短期借款与一年内到期非流动负债）93.23 亿元、长期债务 74.57 亿元，有息负债合计 167.80 亿元，现金短债比 0.55 倍——覆盖偏紧。 净有息负债 116.49 亿元，相当于其 2026 年上半年营业收入的 77\%。 流动比率 0.98、速动比率 0.70，短期偿付压力较大。 2026 年上半年经营活动现金流净额 10.06 亿元，净利润为负而经营现金流为正，差额主要来自折旧摊销与营运资本变动，说明亏损尚未完全转化为现金流出。 2026 年 6 月末在建工程 2.62 亿元，较上年末减少 0.11 亿元，在建投入在收缩。 2026 年 6 月末商誉 14.81 亿元，占归母净资产的 15.5\%，是净资产中最脆弱的一块。 综合看，现金短债比 0.55 倍，短期资金安排偏紧；资产负债率高于样本中位数 18.5 个百分点；有息负债中短期债务占比更高，期限结构仍以滚动续作为主。

\pfl{重大战略转型} 从公开披露看，公司的战略推进集中在：并购与重组（2 项）：2022 年 12 月，公司将所持北京海垦商贸发展有限公司 29.34\% 股权全部转让给海垦控股集团；2022 年 11 月，要约收购结束后合计持有 10.86 亿股（约 68.10\%）。 公告口径上，近三年（2023 年 9 月---2026 年 9 月）公司披露重大重组与并购类公告 2 条、对外投资与转型类公告 3 条、回购与股权激励类公告 0 条，合计公告 441 条。 第一大行业仍是“农业”，收入占比 96.8\%，与 2021 年年报相比没有方向性变化。 综合判断，报告期内公司有 5 条与战略相关的公告落地（重大重组与并购 2 条、对外投资与转型 3 条），资本运作与主业调整之间存在直接关联。

\begin{center}\small\setlength{\tabcolsep}{4pt}
\captionof{table}[海南橡胶（601118）历史融资与资本运作]{海南橡胶（601118）历史融资与资本运作（据公司公告与公开报道整理；金额为披露值，含首次公开发行、再融资、债券、银行借款与重整投资等）}
\begin{adjustbox}{max width=\textwidth}
\begin{tabular}{@{}p{1.45cm}p{2.55cm}p{4.05cm}p{4.95cm}@{}}
\toprule
时间 & 方式 & 规模 & 用途与说明 \\
\midrule
2010-12 & IPO（上海证券交易所主板） & 发行 78,\allowbreak{}600 万股，\allowbreak{}发行价 5.99 元/\allowbreak{}股，\allowbreak{}募集资金总额 47.08 亿元 & 招股公告日 2010-12-16；预计募集资金 47.08 亿元 \\
2018-02 & 定向增发（非公开发行） & 发行 34,\allowbreak{}825.62 万股，\allowbreak{}发行价 5.16 元/\allowbreak{}股，\allowbreak{}募集资金总额 17.97 亿元 & 特种胶园更新种植等项目；发行对象共 7 家：华商基金、\allowbreak{}长城（天津）股权投资基金管理、\allowbreak{}太平洋资产管理、\allowbreak{}广发证券资产管理（广东）、\allowbreak{}信泰人寿保险、\allowbreak{}财通基金、\allowbreak{}国信国投基金管理（北京）；2018-02-12 完成新增股份登记，\allowbreak{}总股本增至 427,\allowbreak{}942.78 万股 \\
2025-07 & 中期票据（银行间市场） & “25 海南橡胶 MTN001（科创债）”，\allowbreak{}发行金额 3.50 亿元，\allowbreak{}票面利率 2.07\% & 偿还公司本部及下属企业有息负债；债券代码 102583120 \\
2026-04 & 2026 年度融资额度议案 & 公司及子公司拟通过金融机构借款、\allowbreak{}发行债券、\allowbreak{}银行票据等方式融资，\allowbreak{}融资余额较上年度净增加不超过 47.00 亿元 & 债务性和权益性融资；第六届董事会相关会议审议通过 \\
2026-06 & 银行授信（关联方） & 海南农垦集团财务有限公司授予公司授信总额 60.00 亿元，\allowbreak{}本期实际发生额 4.37 亿元 & 关联方授信业务；截至 2026 年 3 月末，\allowbreak{}公司合并口径共获得金融机构授信总额合计 362.74 亿元 \\
\bottomrule
\end{tabular}
\end{adjustbox}
\end{center}

\pfl{历史融投资情况} 从现金流量表看，2025 年公司取得借款收到的现金 191.94 亿元、偿还债务支付的现金 196.40 亿元、吸收权益性投资 3.51 亿元，筹资活动现金流净额 -24.82 亿元（2023 年为取得借款 150.10 亿元、偿还债务 151.41 亿元）。 2026 年上半年筹资活动现金流净额为 -8.55 亿元（取得借款 68.79 亿元、偿还债务 72.59 亿元，未年化），已转为净流出，处于净还债状态。 2025 年分配股利、利润或偿付利息支付的现金合计 7.07 亿元。 公开披露的融资记录共 5 笔，工具类型以 IPO、定向增发、中期票据、银行授信为主；金额与用途见下表。 其中规模最大的两笔为：2026 年 6 月，银行授信，海南农垦集团财务有限公司授予公司授信总额 60.00 亿元，本期实际发生额 4.37 亿元；2025 年 7 月，中期票据，“25 海南橡胶 MTN001（科创债）”，发行金额 3.50 亿元，票面利率 2.07\%。 股本方面，近三年披露股本变动 3 次；总股本自 2021 年末以来未发生变化（42.79 亿股）。 分红方面，近三年现金分红 2 次、合计每股派现 0.03 元，最近一次实施在 2024 年 12 月。 近三年“再融资”类公告 2 条，涉及定向增发。 综合看，公司融资结构上以债务滚动为主、股权融资为补充；2025 年筹资活动现金流净额 -24.82 亿元，融资节奏仍以维持存量债务周转为主，与前述经营与投资现金流的方向基本一致。

\pfl{未来潜在融资需求分析} 2026 年 6 月末货币资金 51.31 亿元，而短期债务 93.23 亿元（现金短债比 0.55 倍），2026 年上半年经营现金流净额 10.06 亿元。按“短期债务减去账面货币资金”的滚动偿付口径，静态缺口约 41.92 亿元；若再计入上半年亏损的年化额（2.61 亿元），维持一年正常经营所需的外部资金约 41.92---44.53 亿元。 按第 \ref{sec:financing-need} 节的三档划分，公司属于第一档“补血型”：2026 年上半年归母净利润 -1.30 亿元、6 月末资产负债率 70.7\%。 可行的融资路径按现实性排序为：债务重组或展期（与银行和债券持有人协商）、定向增发（需股价配合，且控股股东需放弃优先认购或同步增资）、引入国资或产业投资人（部分公司以控制权让渡为对价）。 公开披露的后续资金安排线索包括：债务与授信安排：2026 年 4 月，公司及子公司融资余额较上年度净增加不超过 47.00 亿元；资本开支安排：2026 年 6 月，特种胶园更新种植项目已累计投入募集资金 12.01 亿元。 资金需求还要叠加在建工程：期末在建工程 2.62 亿元、2026 年上半年购建固定资产等支出 1.50 亿元（未年化），这部分支出在周期底部通常会被主动放缓。 综合看，公司的外部资金需求以债务接续为主，滚动偿付口径下的静态缺口约 41.92 亿元，能否落地取决于授信续作与发行窗口的配合。




\subsubsection{北大荒（600598）}
\label{co:600598}

\pfl{公司基本情况} 北大荒成立于 1998 年，黑龙江垦区核心农业上市公司，以国有耕地承包经营（土地发包）为主业，土地承包费收入长期占营业收入七成左右；2026 年上半年因补缴 2021—2025 年企业所得税及滞纳金 14.10 亿元而由盈转亏。 公司于 2002-03-29 在上交所首发上市，注册资本 17.78 亿元，总股本 17.78 亿股。 截至 2026 年 10 月 12 日收盘价 12.73 元，对应总市值 226.3 亿元，近一年-12.9\%，总市值在三级行业“粮食种植”的 2 家公司中按由大到小排第\zhnumber{1}位，近一年涨跌幅低于全样本中位数（-10.4\%）2.6 个百分点。 按 2026 年半年报披露的产品构成，收入占比前三位为土地承包费 91.3\%；农产品销售 5.2\%；其他商品销售或服务 2.3\%。

\pfl{股权结构与控制权} 按 2026 年 6 月 30 日的股东名册，前三大股东为北大荒农垦集团有限公司 持股 64.14\%；香港中央结算有限公司（持股比例未披露）；大成基金-中国人寿保险股份有限公司-传统险-大成基金 持股 0.83\%，股东户数 104,209 户。 控制权与股东性质方面，公司 2002 年 3 月 29 日在上海证券交易所挂牌交易。

\pfl{经营要点} 经营与产能方面，不涉及税务行政处罚。

\pfl{所处行业的基本面} 粮食种植的供给由播种面积与单产决定、需求由口粮与饲料用粮决定，价格更多由最低收购价与储备政策托底（第 \ref{sec:framework} 章的供需框架）。 与该公司直接相关的行业数据包括：新季玉米集中上市期临近、下游企业以刚需采购为主、库存维持相对低位，国际玉米价格大幅上涨带动新陈玉米价格衔接顺畅。

\pfl{竞争优势} 把该公司放回同类公司比较，2026 年上半年营业收入 21.3 亿元、在三级行业“粮食种植”的 2 家公司中排第\zhnumber{2}位；毛利率 54.9\%、高于行业内中位数 21.8 个百分点；半年 ROE -7.1\%（未年化）、低于行业内中位数（-2.52\%）；资产负债率 39.5\%、低于行业内中位数（43.0\%）。 公司在公开资料中自述与同业比较的相对优势在于：公司是我国目前规模大、现代化水平高的种植业上市公司和重要的商品粮生产基地，在自然资源、基础设施、农机装备、农业科技服务、组织管理等方面具有显著优势。 优势在报表上主要体现为高于行业平均的毛利率，这在周期底部是抵御亏损的主要缓冲。

\pfl{历史股价} 以 2021 年 1 月为起点、2026 年 10 月为终点，股价由 16.37 元变为 12.73 元，累计 -22.2\%，区间最高 17.50 元（2026 年 1 月）、最低 10.84 元（2026 年 6 月）；当前价相当于区间最高价的 72.7\%，为区间最低价的 1.17 倍。 月度收益的年化波动率 24.9\%，低于全样本中位数 37.3\%，在三级行业“粮食种植”的 2 家公司中波动率由高到低排第\zhnumber{2}位。 把股价阶段与公司事件对齐看，2026 年 1 月 至 2026 年 6 月股价下跌-38.1\%，同期（2026 年 6 月）不属于前期会计差错，预计减少 2026 年归母净利润约 14.10 亿元。 把这些阶段与公司事件对齐后可以看到，几轮大涨大跌多由公司自身的资本运作与风险事件触发，行业景气只解释了其中一部分。

\begin{center}
\begin{minipage}{\textwidth}\centering
\begin{tikzpicture}
\begin{axis}[
  width=12.6cm, height=3.9cm, scale only axis,
  xmin=2020.922, xmax=2026.972, ymin=10.08, ymax=18.73,
  xtick={2021,2022,2023,2024,2025,2026},
  yearticks,
  yticklabel style={font=\scriptsize},
  ylabel={元/股}, ylabel style={font=\scriptsize},
  grid=major, grid style={LightGray, line width=0.3pt},
  axis line style={InkGray, line width=0.5pt},
  every axis plot/.append style={line width=0.9pt},
]
\addplot[AgriGreen] table[x=x,y=y]{data/clean/profiles/px_600598.dat};
\addplot[SoftRed, dashed, line width=0.7pt] table[x=x,y=y]{data/clean/profiles/pxzz_600598.dat};
\addplot[only marks, mark=*, mark size=1.1pt, SoftRed] table[x=x,y=y]{data/clean/profiles/pxzz_600598.dat};
\end{axis}
\end{tikzpicture}

\captionof{figure}[北大荒（600598）股价与周期骨架]{北大荒（600598）月度收盘价（元/股）与 ZigZag 周期骨架（阈值 25\%，圆点为周期转折点）}
\end{minipage}
\end{center}

\pfl{过去周期分析} 以 25\% 的 ZigZag 阈值划分，样本区间内股价形成 2 个主升段与 3 个主跌段。 最大主跌段为 2026 年 1 月 至 2026 年 6 月，5 个月-38.1\%；最大主升段出现在 2024 年 1 月 至 2026 年 1 月，24 个月+53.8\%。 最近一次转折出现在 2026 年 6 月（下跌段自 2026 年 1 月起，5 个月-38.1\%），当前处于该段的延续之中。 公司股价与申万农林牧渔指数几乎同步见底，个股并未走出独立行情。 土地经营型企业的收入与地租成本高度相关（第 \ref{sec:planting} 节），粮价与政策补贴是其主要变量，公司 2026 年 6 月末资产负债率 39.5\%、半年 ROE -7.1\%（未年化），在本轮下行中属于随行业同步调整的一类。 综合区间形态看，该股单段涨跌幅度偏大、以趋势段为主（最大主升 +53.8\%、最大主跌 -38.1\%），年化波动率 24.9\% 与全样本中位数 37.3\% 相比波动小于样本中位数。

\pfl{盈利情况} 2026 年上半年营业收入 21.28 亿元（同比 -29.4\%），归母净利润 -5.37 亿元，由上年同期的盈利转为亏损。 以年报为趋势对照，2023---2025 年营业收入由 50.44 亿元变为 52.29 亿元（两年年均复合增速 +1.8\%），归母净利润由 10.64 亿元变为 11.66 亿元。 按半年报口径，2026 年上半年的 ROE 为 -7.1\%（未年化）、销售净利率 -25.4\%、毛利率 54.9\%，ROE 在“粮食种植”2 家公司中按数值由高到低排第\zhnumber{2}位。 按同一口径，三级行业“粮食种植”2 家公司的半年 ROE 中位数为 -2.52\%，该公司低于其中位数 4.56 个百分点。 从公司披露的原因看，业绩变动主要来自：2026 年上半年，导致销售收入同比减少 8.86 亿元。 综合看，公司仍处亏损区间、由盈转亏，半年 ROE 在三级行业内按由高到低排第\zhnumber{2}位，单位盈利能力的修复依赖行业景气与自身成本端的改善，报表弹性主要体现为价格与销量的方向性变化。

\begin{center}\small\setlength{\tabcolsep}{3.4pt}
\captionof{table}[北大荒（600598）关键财务指标]{北大荒（600598）关键财务指标（合并报表口径；两个半年报期均未年化；现金短债比 $=$ 货币资金 $\div$ 短期债务）}
\begin{adjustbox}{max width=\textwidth}
\begin{tabular}{@{}lrrrrr@{}}
\toprule
指标 & 2023 年 & 2024 年 & 2025 年 & 2025 年半年报 & 2026 年半年报 \\
\midrule
营业收入（亿元） & 50.44 & 53.39 & 52.29 & 30.14 & 21.28 \\
归母净利润（亿元） & 10.64 & 10.87 & 11.66 & 9.84 & -5.37 \\
毛利率（\%） & 35.4 & 34.6 & 34.6 & 39.0 & 54.9 \\
ROE（\%） & 14.6 & 14.3 & 14.9 & 12.1 & -7.1 \\
资产负债率（\%） & 16.1 & 15.4 & 12.9 & 34.1 & 39.5 \\
经营活动现金流净额（亿元） & 15.15 & 14.94 & 12.40 & 36.92 & 24.86 \\
货币资金（亿元） & 23.70 & 15.86 & 2.00 & 7.80 & 19.27 \\
有息负债合计（亿元） & 0.57 & 0.73 & 0.10 & 0.16 & 0.10 \\
现金短债比（倍） & 69.01 & 24.33 & 56.25 & 112.61 & 566.76 \\
\bottomrule
\end{tabular}
\end{adjustbox}
\end{center}

\begin{center}
\begin{minipage}{\textwidth}\centering
\begin{tikzpicture}
\begin{groupplot}[
  group style={group size=2 by 1, horizontal sep=0.60cm},
  width=6.85cm, height=4.60cm,
  tick label style={font=\tiny},
  title style={font=\scriptsize\heiti},
  yticklabel style={font=\tiny},
  ylabel style={font=\tiny},
  grid=major, grid style={LightGray, line width=0.3pt},
  axis line style={InkGray, line width=0.5pt},
  enlarge x limits={abs=0.8},
  legend style={font=\scriptsize, at={(0.5,-0.20)}, anchor=north,
                legend columns=2, draw=none, fill=none, column sep=6pt},
]
\nextgroupplot[
  ybar, bar width=4pt,
  ymin=-11.24, ymax=60.44,
  xtick={0,1,2,3,4,5.7},
  xticklabels={2021,2022,2023,2024,2025,2026H1},
  ylabel={亿元},
]
\addplot[fill=AgriGreen!75, draw=AgriGreen, bar shift=-2.6pt] table[x=i, y=rev]{data/clean/profiles/fin_600598.dat};
\addplot[fill=SoftRed!60, draw=SoftRed, bar shift=2.6pt] table[x=i, y=ni]{data/clean/profiles/fin_600598.dat};
\addplot[fill=AgriGreen!30, draw=AgriGreen, pattern=north east lines, bar shift=-2.6pt] table[x=x, y=rev]{data/clean/profiles/finh_600598.dat};
\addplot[fill=SoftRed!20, draw=SoftRed, pattern=north east lines, bar shift=2.6pt] table[x=x, y=ni]{data/clean/profiles/finh_600598.dat};
\legend{营业收入（年/半年）, 归母净利润（年/半年）}
\nextgroupplot[
  ymin=-10.8, ymax=44.2,
  xtick={0,1,2,3,4,5.7},
  xticklabels={2021,2022,2023,2024,2025,2026H1},
  ylabel={\%},
]
\addplot[AgriGold, mark=*, mark size=1.3pt] table[x=i, y=roe]{data/clean/profiles/fin_600598.dat};
\addplot[OilBlue, mark=square*, mark size=1.3pt] table[x=i, y=debt]{data/clean/profiles/fin_600598.dat};
\addplot[only marks, AgriGold, mark=*, mark size=2.0pt] table[x=x, y=roe]{data/clean/profiles/finh_600598.dat};
\addplot[only marks, OilBlue, mark=square*, mark size=2.0pt] table[x=x, y=debt]{data/clean/profiles/finh_600598.dat};
\legend{ROE, 资产负债率}
\end{groupplot}
\end{tikzpicture}
\begin{tikzpicture}
\begin{axis}[
  name=finbot,
  width=13.75cm, height=3.10cm, scale only axis,
  ymin=-3.05, ymax=34.20,
  xtick={0,1,2,3,4,5.7},
  xticklabels={2021,2022,2023,2024,2025,2026H1},
  tick label style={font=\tiny},
  yticklabel style={font=\tiny},
  ylabel={亿元}, ylabel style={font=\tiny},
  ybar, bar width=4pt,
  grid=major, grid style={LightGray, line width=0.3pt},
  axis line style={InkGray, line width=0.5pt},
  enlarge x limits={abs=0.8},
  legend style={font=\scriptsize, at={(0.5,-0.26)}, anchor=north,
                legend columns=2, draw=none, fill=none, column sep=10pt},
]
\addplot[fill=AgriGreen!55, draw=AgriGreen, bar shift=-3.0pt] table[x=i, y=cash]{data/clean/profiles/fin_600598.dat};
\addplot[fill=SoftRed!45, draw=SoftRed, bar shift=3.0pt] table[x=i, y=ibd]{data/clean/profiles/fin_600598.dat};
\addplot[fill=AgriGreen!25, draw=AgriGreen, pattern=north east lines, bar shift=-3.0pt] table[x=x, y=cash]{data/clean/profiles/finh_600598.dat};
\addplot[fill=SoftRed!20, draw=SoftRed, pattern=north east lines, bar shift=3.0pt] table[x=x, y=ibd]{data/clean/profiles/finh_600598.dat};
\legend{货币资金（左轴）, 有息负债（左轴）}
\end{axis}
\begin{axis}[
  at={(finbot.south east)}, anchor=south east,
  width=13.75cm, height=3.10cm, scale only axis,
  axis y line*=right, axis x line*=none, xtick=\empty,
  ymin=-6.46, ymax=32.31,
  tick label style={font=\tiny},
  yticklabel style={font=\tiny}, scaled y ticks=false,
  legend style={font=\scriptsize, at={(0.5,-0.44)}, anchor=north,
                legend columns=1, draw=none, fill=none},
]
\addplot[OilBlue, mark=*, mark size=2.2pt, line width=1.3pt] table[x=i, y=ocf]{data/clean/profiles/fin_600598.dat};
\addplot[only marks, OilBlue, mark=*, mark size=3.2pt] table[x=x, y=ocf]{data/clean/profiles/finh_600598.dat};
\legend{经营活动现金流净额（右轴，亿元，蓝点）}
\end{axis}
\end{tikzpicture}

\captionof{figure}[北大荒（600598）近年主要财务指标]{北大荒（600598）近年主要财务指标（左上：营业收入与归母净利润；右上：ROE 与资产负债率；下：货币资金、有息负债为左轴柱状，经营活动现金流净额为其右轴的蓝点折线。
2021---2025 年为年报口径，2026H1 为 2026 年半年报/半年末，与其前一年报之间留出间隔、以斜纹或空心点区分；半年报数据未年化）}
\end{minipage}
\end{center}

\pfl{财务分析} 2026 年 6 月末资产负债率 39.5\%，较 2025 年末上升 26.6 个百分点（样本中位数 52.2\%，所处三级行业内按由低到高排第\zhnumber{1}位）。 2026 年 6 月末货币资金 19.27 亿元、短期债务（短期借款与一年内到期非流动负债）0.03 亿元、长期债务 0.06 亿元，有息负债合计 0.10 亿元，现金短债比 566.76 倍——覆盖充分。 净有息负债 -19.17 亿元，相当于其 2026 年上半年营业收入的 -90\%。 流动比率 1.46、速动比率 1.41，短期偿债能力处于正常区间。 2026 年上半年经营活动现金流净额 24.86 亿元，净利润为负而经营现金流为正，差额主要来自折旧摊销与营运资本变动，说明亏损尚未完全转化为现金流出。 综合看，现金短债比 566.76 倍，短期偿付压力可控；资产负债率低于样本中位数 12.7 个百分点；有息负债中长期债务占比更高，期限结构对短债滚动的压力较小。 公司披露的监管与合规风险为：2026 年 6 月，不属于前期会计差错，预计减少 2026 年归母净利润约 14.10 亿元。

\pfl{重大战略转型} 从公开披露看，公司的战略推进集中在：并购与重组（1 项）：2014 年，公司剥离北大荒米业集团股权；其他动作（1 项）：2026 年 6 月，报告期末“合同负债”比上期期末增长 25,848.03\%。 公告口径上，近三年（2023 年 9 月---2026 年 9 月）公司披露重大重组与并购类公告 1 条、对外投资与转型类公告 1 条、回购与股权激励类公告 0 条，合计公告 254 条。 第一大产品“土地承包费”的收入占比由 2021 年年报的 77.0\% 变为 91.3\%（14.3 个百分点），收入结构出现再平衡。 综合判断，报告期内公司有 2 条与战略相关的公告落地（重大重组与并购 1 条、对外投资与转型 1 条），资本运作与主业调整之间存在直接关联。

\begin{center}\small\setlength{\tabcolsep}{4pt}
\captionof{table}[北大荒（600598）历史融资与资本运作]{北大荒（600598）历史融资与资本运作（据公司公告与公开报道整理；金额为披露值，含首次公开发行、再融资、债券、银行借款与重整投资等）}
\begin{adjustbox}{max width=\textwidth}
\begin{tabular}{@{}p{1.45cm}p{2.55cm}p{4.05cm}p{4.95cm}@{}}
\toprule
时间 & 方式 & 规模 & 用途与说明 \\
\midrule
2002-03 & IPO（上海证券交易所） & 发行 30,\allowbreak{}000 万股，\allowbreak{}发行价 5.38 元/\allowbreak{}股，\allowbreak{}募集资金总额 16.14 亿元 & 百万吨稻谷精深加工及十万立方米绿色稻草板生产项目、\allowbreak{}农业高新技术大规模综合应用项目、\allowbreak{}浩良河分公司油改煤化肥工程项目；经证监发行字[2002]19 号文核准；募集资金于 2002-03-19 全部划入公司指定账户 \\
2003 & 变更募投并收购资产 & 收购农垦集团所持黑龙江省北大荒米业有限公司 98.55\% 股权 & 由原投资新建改为收购控股股东资产；其中 3.50 亿元来自北大荒上市募集资金 \\
2007-12 & 可转换公司债券 & 发行数量 15.00 亿元（15.00 亿元），\allowbreak{}面值 100.00 元，\allowbreak{}按面值发行 & 经证监发行字[2007]451 号文核准；原股东享有优先配售权；截至 2010-12-31 累计转股 1.43 亿股 \\
2025-05 & 现金分红（2024 年度） & 以总股本 17.78 亿股为基数每股派 0.55 元，\allowbreak{}共派发现金红利 9.78 亿元 & 回报股东；现金红利发放已于 5 月 23 日完成 \\
2026-03 & 现金分红预案（2025 年度） & 每 10 股派发现金股利 5.50 元（含税），\allowbreak{}合计拟派发 977,\allowbreak{}723,\allowbreak{}949.95 元 & 回报股东；不送红股，\allowbreak{}不实施资本公积金转增股本 \\
\bottomrule
\end{tabular}
\end{adjustbox}
\end{center}

\pfl{历史融投资情况} 从现金流量表看，2025 年公司偿还债务支付的现金 0.48 亿元，筹资活动现金流净额 -10.32 亿元。 2026 年上半年筹资活动现金流净额为 -9.78 亿元（取得借款 — 亿元、偿还债务 0.00 亿元，未年化），已转为净流出，处于净还债状态。 2025 年分配股利、利润或偿付利息支付的现金合计 9.76 亿元。 公开披露的融资记录共 5 笔，工具类型以 IPO、可转债为主；金额与用途见下表。 其中规模最大的两笔为：2007 年 12 月，可转债，发行数量 15.00 亿元（15.00 亿元），面值 100.00 元，按面值发行；2003 年，变更募投并收购资产，收购农垦集团所持黑龙江省北大荒米业有限公司 98.55\% 股权。 股本方面，近三年披露股本变动 3 次；总股本自 2021 年末以来未发生变化（17.78 亿股）。 分红方面，近三年现金分红 4 次、合计每股派现 1.96 元，最近一次实施在 2025 年 12 月。 综合看，公司融资结构上股权与债务融资并举；2025 年筹资活动现金流净额 -10.32 亿元，融资节奏仍以维持存量债务周转为主，与前述经营与投资现金流的方向基本一致。

\pfl{未来潜在融资需求分析} 2026 年 6 月末货币资金 19.27 亿元，而短期债务 0.03 亿元（现金短债比 566.76 倍），2026 年上半年经营现金流净额 24.86 亿元。账面货币资金可以覆盖短期债务，缺口不体现在静态偿付层面。 公司 2026 年上半年归母净利润为负（-5.37 亿元），但 6 月末资产负债率 39.5\% 尚未进入第一档（亏损且负债率 $>$65\%）的区间，当前融资需求以补充流动资金、支撑低谷期经营性支出为主。 综合看，公司在静态偿付层面不存在缺口，外部融资更多是配合扩张节奏与并购整合的主动性安排，而非偿债压力驱动。




\subsubsection{苏垦农发（601952）}
\label{co:601952}

\pfl{公司基本情况} 公司前身可追溯至 2011 年，江苏省农垦集团旗下的农业全产业链上市公司，自主经营耕地约 95.5 万亩，形成种植、种子、大米、农资、食用油一体化经营。 公司于 2017-05-15 在上交所首发上市，注册资本 13.78 亿元，总股本 13.78 亿股。 截至 2026 年 10 月 12 日收盘价 8.96 元，对应总市值 123.5 亿元，近一年-5.8\%，总市值在三级行业“粮食种植”的 2 家公司中按由大到小排第\zhnumber{2}位，近一年涨跌幅高于全样本中位数（-10.4\%）4.6 个百分点。 按 2025 年年报披露的行业构成，收入占比前三位为大米加工与销售分部 36.4\%；种植业分部 35.6\%；农资服务与销售分部 21.2\%。 与 2021 年年报相比，第一大行业口径由“米业”（39.7\%）变为“大米加工与销售分部”（36.4\%）；两期披露的分类口径有调整，占比差异中同时包含分类因素。

\pfl{股权结构与控制权} 按 2026 年 6 月 30 日的股东名册，前三大股东为江苏省农垦集团有限公司 持股 67.84\%；全国社保基金四一三组合 持股 1.17\%；中国对外经济贸易信托有限公司一外贸信托-仁桥泽源股票 持股 1.03\%，股东户数 42,420 户。 从控制权角度看，2026 年完成增资扩股、注册资本 20.00 亿元并引入 10.00 亿元外部资金与战略投资者。

\pfl{经营要点} 经营层面，截至 2025 年秋播公司自主经营耕地约 95.5 万亩、土地流转（含土地托管与合作种植）面积约 40.7 万亩。

\pfl{所处行业的基本面} 粮食种植的供给由播种面积与单产决定、需求由口粮与饲料用粮决定，价格更多由最低收购价与储备政策托底（第 \ref{sec:framework} 章的供需框架）。 与该公司直接相关的行业数据包括：2026 年，全国夏粮播种面积 39,798.6 万亩、比 2025 年减少 69.0 万亩、下降 0.2\%；2026 年 8 月，主要粮食价格政策明确；新季玉米上市期下游企业以刚需采购为主、库存维持相对低位。

\pfl{竞争优势} 把该公司放回同类公司比较，2026 年上半年营业收入 43.3 亿元、在三级行业“粮食种植”的 2 家公司中排第\zhnumber{1}位；毛利率 11.3\%、低于行业内中位数 21.8 个百分点；半年 ROE 2.0\%（未年化）、高于行业内中位数（-2.52\%）；资产负债率 46.5\%、高于行业内中位数（43.0\%）。 公司区别于同类公司的经营特征包括：苏垦麦芽建有 15 万吨塔式和 10 万吨箱式生产线各一条、麦芽设计产能 25 万吨；金太阳粮油拥有年 30 万吨精炼产能和年灌装 30 万吨的灌装生产线；报告期内中标国家及省级化肥储备任务。 优势在当期体现为高于行业中位数的资本回报，但优势能否持续仍取决于周期与自身产能投放节奏。

\pfl{历史股价} 自 2021 年 1 月至 2026 年 10 月，股价由 11.62 元变为 8.96 元，累计 -22.9\%，区间最高 16.70 元（2022 年 5 月）、最低 7.49 元（2026 年 6 月）；当前价相当于区间最高价的 53.7\%，为区间最低价的 1.20 倍。 月度收益的年化波动率 26.5\%，低于全样本中位数 37.3\%，在三级行业“粮食种植”的 2 家公司中波动率由高到低排第\zhnumber{1}位。 把股价阶段与公司事件对齐看，2022 年 5 月 至 2026 年 6 月股价下跌-55.1\%，同期（2026 年 6 月）公司完成对江苏大丰华丰种业股份有限公司的投资并购。 整体上看，区间的几处大级别拐点都能在公司的资本运作、股权变动或风险事件上找到对应，股价波动有明显的公司层面驱动，而非单纯的行业 beta。

\begin{center}
\begin{minipage}{\textwidth}\centering
\begin{tikzpicture}
\begin{axis}[
  width=12.6cm, height=3.9cm, scale only axis,
  xmin=2020.922, xmax=2026.972, ymin=6.97, ymax=17.87,
  xtick={2021,2022,2023,2024,2025,2026},
  yearticks,
  yticklabel style={font=\scriptsize},
  ylabel={元/股}, ylabel style={font=\scriptsize},
  grid=major, grid style={LightGray, line width=0.3pt},
  axis line style={InkGray, line width=0.5pt},
  every axis plot/.append style={line width=0.9pt},
]
\addplot[AgriGreen] table[x=x,y=y]{data/clean/profiles/px_601952.dat};
\addplot[SoftRed, dashed, line width=0.7pt] table[x=x,y=y]{data/clean/profiles/pxzz_601952.dat};
\addplot[only marks, mark=*, mark size=1.1pt, SoftRed] table[x=x,y=y]{data/clean/profiles/pxzz_601952.dat};
\end{axis}
\end{tikzpicture}

\captionof{figure}[苏垦农发（601952）股价与周期骨架]{苏垦农发（601952）月度收盘价（元/股）与 ZigZag 周期骨架（阈值 25\%，圆点为周期转折点）}
\end{minipage}
\end{center}

\pfl{过去周期分析} 以 25\% 的 ZigZag 阈值划分，样本区间内股价形成 1 个主升段与 2 个主跌段。 最大主跌段为 2022 年 5 月 至 2026 年 6 月，49 个月-55.1\%；最大主升段出现在 2021 年 9 月 至 2022 年 5 月，8 个月+63.9\%。 最近一次转折出现在 2026 年 6 月（下跌段自 2022 年 5 月起，49 个月-55.1\%），当前处于该段的延续之中。 公司股价与申万农林牧渔指数几乎同步见底，个股并未走出独立行情。 土地经营型企业的收入与地租成本高度相关（第 \ref{sec:planting} 节），粮价与政策补贴是其主要变量，公司 2026 年上半年实现盈利、半年 ROE 2.0\%（未年化），好于全样本中位数（0.75\%），下行周期对其报表的冲击小于同业。 综合区间形态看，该股单段涨跌幅度偏大、以趋势段为主（最大主升 +63.9\%、最大主跌 -55.1\%），年化波动率 26.5\% 与全样本中位数 37.3\% 相比波动小于样本中位数。

\pfl{盈利情况} 2026 年上半年营业收入 43.32 亿元（同比 -5.6\%），归母净利润 1.44 亿元，同比 -32.5\%。 以年报为趋势对照，2023---2025 年营业收入由 121.68 亿元变为 100.71 亿元（两年年均复合增速 -9.0\%），归母净利润由 8.16 亿元变为 5.48 亿元。 按半年报口径，2026 年上半年的 ROE 为 2.0\%（未年化）、销售净利率 3.6\%、毛利率 11.3\%，ROE 在“粮食种植”2 家公司中按数值由高到低排第\zhnumber{1}位。 同期全样本半年 ROE 中位数 0.75\%，公司与之相差 1.29 个百分点（略好于中位数公司）。 综合看，公司在报告期维持盈利（高于全样本半年 ROE 中位数 0.75\%），利润的可持续性取决于粮食种植景气的延续与成本管控的执行。

\begin{center}\small\setlength{\tabcolsep}{3.4pt}
\captionof{table}[苏垦农发（601952）关键财务指标]{苏垦农发（601952）关键财务指标（合并报表口径；两个半年报期均未年化；现金短债比 $=$ 货币资金 $\div$ 短期债务）}
\begin{adjustbox}{max width=\textwidth}
\begin{tabular}{@{}lrrrrr@{}}
\toprule
指标 & 2023 年 & 2024 年 & 2025 年 & 2025 年半年报 & 2026 年半年报 \\
\midrule
营业收入（亿元） & 121.68 & 109.17 & 100.71 & 45.88 & 43.32 \\
归母净利润（亿元） & 8.16 & 7.30 & 5.48 & 2.13 & 1.44 \\
毛利率（\%） & 13.3 & 14.5 & 13.2 & 12.0 & 11.3 \\
ROE（\%） & 12.8 & 10.8 & 7.9 & 3.1 & 2.0 \\
资产负债率（\%） & 50.8 & 49.5 & 48.9 & 50.7 & 46.5 \\
经营活动现金流净额（亿元） & 23.08 & 14.68 & 22.59 & 7.22 & 1.67 \\
货币资金（亿元） & 6.26 & 6.91 & 8.49 & 13.37 & 14.56 \\
有息负债合计（亿元） & 51.18 & 48.89 & 45.57 & 48.97 & 45.25 \\
现金短债比（倍） & 1.10 & 1.16 & 1.67 & 2.49 & 3.40 \\
\bottomrule
\end{tabular}
\end{adjustbox}
\end{center}

\begin{center}
\begin{minipage}{\textwidth}\centering
\begin{tikzpicture}
\begin{groupplot}[
  group style={group size=2 by 1, horizontal sep=0.60cm},
  width=6.85cm, height=4.60cm,
  tick label style={font=\tiny},
  title style={font=\scriptsize\heiti},
  yticklabel style={font=\tiny},
  ylabel style={font=\tiny},
  grid=major, grid style={LightGray, line width=0.3pt},
  axis line style={InkGray, line width=0.5pt},
  enlarge x limits={abs=0.8},
  legend style={font=\scriptsize, at={(0.5,-0.20)}, anchor=north,
                legend columns=2, draw=none, fill=none, column sep=6pt},
]
\nextgroupplot[
  ybar, bar width=4pt,
  ymin=-12.73, ymax=142.55,
  xtick={0,1,2,3,4,5.7},
  xticklabels={2021,2022,2023,2024,2025,2026H1},
  ylabel={亿元},
]
\addplot[fill=AgriGreen!75, draw=AgriGreen, bar shift=-2.6pt] table[x=i, y=rev]{data/clean/profiles/fin_601952.dat};
\addplot[fill=SoftRed!60, draw=SoftRed, bar shift=2.6pt] table[x=i, y=ni]{data/clean/profiles/fin_601952.dat};
\addplot[fill=AgriGreen!30, draw=AgriGreen, pattern=north east lines, bar shift=-2.6pt] table[x=x, y=rev]{data/clean/profiles/finh_601952.dat};
\addplot[fill=SoftRed!20, draw=SoftRed, pattern=north east lines, bar shift=2.6pt] table[x=x, y=ni]{data/clean/profiles/finh_601952.dat};
\legend{营业收入（年/半年）, 归母净利润（年/半年）}
\nextgroupplot[
  ymin=-4.5, ymax=61.5,
  xtick={0,1,2,3,4,5.7},
  xticklabels={2021,2022,2023,2024,2025,2026H1},
  ylabel={\%},
]
\addplot[AgriGold, mark=*, mark size=1.3pt] table[x=i, y=roe]{data/clean/profiles/fin_601952.dat};
\addplot[OilBlue, mark=square*, mark size=1.3pt] table[x=i, y=debt]{data/clean/profiles/fin_601952.dat};
\addplot[only marks, AgriGold, mark=*, mark size=2.0pt] table[x=x, y=roe]{data/clean/profiles/finh_601952.dat};
\addplot[only marks, OilBlue, mark=square*, mark size=2.0pt] table[x=x, y=debt]{data/clean/profiles/finh_601952.dat};
\legend{ROE, 资产负债率}
\end{groupplot}
\end{tikzpicture}
\begin{tikzpicture}
\begin{axis}[
  name=finbot,
  width=13.75cm, height=3.10cm, scale only axis,
  ymin=-5.79, ymax=64.83,
  xtick={0,1,2,3,4,5.7},
  xticklabels={2021,2022,2023,2024,2025,2026H1},
  tick label style={font=\tiny},
  yticklabel style={font=\tiny},
  ylabel={亿元}, ylabel style={font=\tiny},
  ybar, bar width=4pt,
  grid=major, grid style={LightGray, line width=0.3pt},
  axis line style={InkGray, line width=0.5pt},
  enlarge x limits={abs=0.8},
  legend style={font=\scriptsize, at={(0.5,-0.26)}, anchor=north,
                legend columns=2, draw=none, fill=none, column sep=10pt},
]
\addplot[fill=AgriGreen!55, draw=AgriGreen, bar shift=-3.0pt] table[x=i, y=cash]{data/clean/profiles/fin_601952.dat};
\addplot[fill=SoftRed!45, draw=SoftRed, bar shift=3.0pt] table[x=i, y=ibd]{data/clean/profiles/fin_601952.dat};
\addplot[fill=AgriGreen!25, draw=AgriGreen, pattern=north east lines, bar shift=-3.0pt] table[x=x, y=cash]{data/clean/profiles/finh_601952.dat};
\addplot[fill=SoftRed!20, draw=SoftRed, pattern=north east lines, bar shift=3.0pt] table[x=x, y=ibd]{data/clean/profiles/finh_601952.dat};
\legend{货币资金（左轴）, 有息负债（左轴）}
\end{axis}
\begin{axis}[
  at={(finbot.south east)}, anchor=south east,
  width=13.75cm, height=3.10cm, scale only axis,
  axis y line*=right, axis x line*=none, xtick=\empty,
  ymin=-6.00, ymax=30.01,
  tick label style={font=\tiny},
  yticklabel style={font=\tiny}, scaled y ticks=false,
  legend style={font=\scriptsize, at={(0.5,-0.44)}, anchor=north,
                legend columns=1, draw=none, fill=none},
]
\addplot[OilBlue, mark=*, mark size=2.2pt, line width=1.3pt] table[x=i, y=ocf]{data/clean/profiles/fin_601952.dat};
\addplot[only marks, OilBlue, mark=*, mark size=3.2pt] table[x=x, y=ocf]{data/clean/profiles/finh_601952.dat};
\legend{经营活动现金流净额（右轴，亿元，蓝点）}
\end{axis}
\end{tikzpicture}

\captionof{figure}[苏垦农发（601952）近年主要财务指标]{苏垦农发（601952）近年主要财务指标（左上：营业收入与归母净利润；右上：ROE 与资产负债率；下：货币资金、有息负债为左轴柱状，经营活动现金流净额为其右轴的蓝点折线。
2021---2025 年为年报口径，2026H1 为 2026 年半年报/半年末，与其前一年报之间留出间隔、以斜纹或空心点区分；半年报数据未年化）}
\end{minipage}
\end{center}

\pfl{财务分析} 2026 年 6 月末资产负债率 46.5\%，较 2025 年末下降 2.4 个百分点（样本中位数 52.2\%，所处三级行业内按由低到高排第\zhnumber{2}位）。 2026 年 6 月末货币资金 14.56 亿元、短期债务（短期借款与一年内到期非流动负债）4.28 亿元、长期债务 40.98 亿元，有息负债合计 45.25 亿元，现金短债比 3.40 倍——覆盖充分。 净有息负债 30.70 亿元，相当于其 2026 年上半年营业收入的 71\%。 流动比率 3.47、速动比率 2.08，短期指标尚可。 2026 年上半年经营活动现金流净额 1.67 亿元，经营现金流与利润同向为正，内生资金可以支撑运营。 2026 年 6 月末在建工程 3.47 亿元，较上年末增加 1.12 亿元，仍有在建投入。 综合看，现金短债比 3.40 倍，短期偿付压力可控；资产负债率低于样本中位数 5.7 个百分点；有息负债中长期债务占比更高，期限结构对短债滚动的压力较小。

\pfl{重大战略转型} 近三年公司的战略动作集中在：并购与重组（1 项）：2026 年 6 月，公司完成对江苏大丰华丰种业股份有限公司的投资并购；其他动作（1 项）：上述项目，大华种业常规稻麦自主品种审定成果丰硕，完成河南子公司设立。 公告口径上，近三年（2023 年 9 月---2026 年 9 月）公司披露重大重组与并购类公告 3 条、对外投资与转型类公告 4 条、回购与股权激励类公告 0 条，合计公告 289 条。 主营结构的口径变化已经落到报表上：2021 年年报的第一大行业为“米业”（39.7\%），2025 年年报为“大米加工与销售分部”（36.4\%）；公司在这两期的分部划分方式有过调整。 综合判断，报告期内公司有 7 条与战略相关的公告落地（重大重组与并购 3 条、对外投资与转型 4 条），资本运作与主业调整之间存在直接关联。

\begin{center}\small\setlength{\tabcolsep}{4pt}
\captionof{table}[苏垦农发（601952）历史融资与资本运作]{苏垦农发（601952）历史融资与资本运作（据公司公告与公开报道整理；金额为披露值，含首次公开发行、再融资、债券、银行借款与重整投资等）}
\begin{adjustbox}{max width=\textwidth}
\begin{tabular}{@{}p{1.45cm}p{2.55cm}p{4.05cm}p{4.95cm}@{}}
\toprule
时间 & 方式 & 规模 & 用途与说明 \\
\midrule
2017-05 & IPO（上海证券交易所主板） & 公开发行 26,\allowbreak{}000 万股，\allowbreak{}发行价 9.32 元/\allowbreak{}股，\allowbreak{}募集资金总额 24.23 亿元 & 百万亩农田改造建设项目（16.68 亿元）、\allowbreak{}大华种业集团改扩建项目（1.32 亿元）、\allowbreak{}农业科学研究院建设项目、\allowbreak{}农业信息化建设项目、\allowbreak{}补充流动资金（3.00 亿元）；；募集资金 2017-05-09 到账 \\
2022-04 & 变更部分募投项目 & 使用首发募投项目部分现金管理收益及利息 2.46 亿元投资建设“苏垦米业集团改扩建项目”（含江心沙、\allowbreak{}淮海公司、\allowbreak{}宝应湖三个子项目） & 大米加工产能扩建；2024 年 8 月因政策收严、\allowbreak{}淮海项目用地手续耗时较长 \\
2026-04-29 & 子公司增资扩股引进战略投资者（关联交易） & 控股子公司江苏种业完成 10.00 亿元外部资金引入及战略投资者引进；控股股东江苏省农垦集团投资 4.50 亿元认缴新增注册资本 2.90 亿元 & 优化江苏种业股权结构、\allowbreak{}充实种业资本实力；2026-06-17 完成工商变更登记；增资后江苏种业注册资本 20.00 亿元 \\
2026-08 & 中期现金分红 & 以 2026 年半年末总股本 13.78 亿股为基数每 10 股派现金红利 0.40 元 & 回报股东；属 2025 年度股东会授权董事会决策的权限范围，\allowbreak{}无需提交股东会审议 \\
\bottomrule
\end{tabular}
\end{adjustbox}
\end{center}

\pfl{历史融投资情况} 从现金流量表看，2025 年公司取得借款收到的现金 1.05 亿元、偿还债务支付的现金 2.15 亿元，筹资活动现金流净额 -13.17 亿元（2023 年为取得借款 3.68 亿元、偿还债务 7.84 亿元）。 2026 年上半年筹资活动现金流净额为 6.32 亿元（取得借款 0.10 亿元、偿还债务 1.16 亿元，未年化），仍处于净融入状态。 2025 年分配股利、利润或偿付利息支付的现金合计 4.17 亿元。 公开披露的融资记录共 4 笔，工具类型以 IPO、战略投资为主；金额与用途见下表。 其中规模最大的两笔为：2026 年 4 月，战略投资，控股子公司江苏种业完成 10.00 亿元外部资金引入及战略投资者引进，控股股东江苏省农垦集团投资 4.50 亿元认缴新增注册资本 2.90 亿元；2017 年 5 月，IPO，公开发行 26,000 万股，发行价 9.32 元/股，募集资金总额 24.23 亿元。 股本方面，近三年披露股本变动 3 次；总股本自 2021 年末以来未发生变化（13.78 亿股）。 分红方面，近三年现金分红 7 次、合计每股派现 1.12 元，最近一次实施在 2026 年 6 月。 综合看，公司融资结构上股权与债务融资并举；2025 年筹资活动现金流净额 -13.17 亿元，融资节奏仍以维持存量债务周转为主，与前述经营与投资现金流的方向基本一致。

\pfl{未来潜在融资需求分析} 2026 年 6 月末货币资金 14.56 亿元，而短期债务 4.28 亿元（现金短债比 3.40 倍），2026 年上半年经营现金流净额 1.67 亿元。账面货币资金可以覆盖短期债务，缺口不体现在静态偿付层面。 公司 2026 年 6 月末资产负债率 46.5\%、2026 年上半年 ROE 2.0\%（未年化），资产负债表尚有余量，融资需求不是“补血型”的刚性需求，而是配套在建项目投入的主动性安排。 资金需求还要叠加在建工程：期末在建工程 3.47 亿元、2026 年上半年购建固定资产等支出 2.02 亿元（未年化），这部分支出在周期底部通常会被主动放缓。 静态偿付不存在缺口，后续融资更可能服务于产能或并购安排，而不是填补流动性窟窿。




\subsubsection{隆平高科（000998）}
\label{co:000998}

\pfl{公司基本情况} 隆平高科成立于 1999 年，由袁隆平院士发起设立、中信集团为实际控制人的综合性跨国种业集团，全球种业企业前八强，杂交水稻种子国内领先；因国内玉米种业高库存与巴西业务大幅承压，2026 年上半年亏损扩大，商誉规模较大。 公司于 2000-12-11 在深交所主板首发上市，注册资本 14.69 亿元，总股本 13.17 亿股。 截至 2026 年 10 月 12 日收盘价 8.64 元，对应总市值 113.8 亿元，近一年-11.8\%，总市值在三级行业“种子”的 8 家公司中按由大到小排第\zhnumber{1}位，近一年涨跌幅低于全样本中位数（-10.4\%）1.4 个百分点。 按 2026 年半年报披露的行业构成，收入几乎全部来自农业 100.0\%。

\pfl{股权结构与控制权} 按 2026 年 6 月 30 日的股东名册，前三大股东为中信农业科技股份有限公司 持股 25.20\%；湖南杂交水稻研究中心 持股 4.55\%；北大荒中垦(深圳)投资有限公司 持股 4.48\%，股东户数 138,559 户。 股东结构与控制权方面，2025 年 5 月，公司完成向特定对象发行，中信农业认购 1.52 亿股。

\pfl{经营要点} 业务与产能层面，2025 年度，制种产量 34,338.13 万公斤、比 2024 年度增加 2.53\%。

\pfl{所处行业的基本面} 种子的供给由制种面积与库存决定、需求由粮价与种植意愿决定，制种面积在 2025 年起出现萎缩，是本轮供给端最先出清的环节（第 \ref{sec:framework} 章的供需框架）。 从供需两端的当期数据看，2026 年上半年，行业种子售价与销量普遍下行；2026 年，全国玉米制种面积调减至 358.5 万亩、同比下降 7\%。

\pfl{竞争优势} 把该公司放回同类公司比较，2026 年上半年营业收入 17.1 亿元、在三级行业“种子”的 8 家公司中排第\zhnumber{2}位；毛利率 39.1\%、高于行业内中位数 10.6 个百分点；半年 ROE -4.1\%（未年化）、低于行业内中位数（0.17\%）；资产负债率 59.1\%、高于行业内中位数（37.3\%）。 相对优势方面，每年带动农民增收超 130.00 亿元，截至 2025 年末累计有效植物新品种权 1,122 件、有效专利 125 件；成为全国农业领域首家获此荣誉的企业，并连续两年获评“最具社会责任感上市公司”；公司以裕丰 303D、中科玉 505DB 为代表的 19 个次转基因玉米品种已通过审定。 优势在报表上主要体现为高于行业平均的毛利率，这在周期底部是抵御亏损的主要缓冲。

\pfl{历史股价} 以 2021 年 1 月为起点、2026 年 10 月为终点，股价由 17.98 元变为 8.64 元，累计 -51.9\%，区间最高 24.04 元（2021 年 11 月）、最低 7.54 元（2026 年 6 月）；当前价相当于区间最高价的 35.9\%，为区间最低价的 1.15 倍。 月度收益的年化波动率 28.7\%，低于全样本中位数 37.3\%，在三级行业“种子”的 8 家公司中波动率由高到低排第\zhnumber{8}位。 把股价阶段与公司事件对齐看，2021 年 11 月 至 2026 年 6 月股价下跌-68.6\%，同期（2025 年 5 月）发行后中信集团合计持股比例由 17.36\% 提升至不超过 25.93\%。 从时间对应关系看，公司层面的资本运作与股权、风险事件解释了区间内多数急涨急跌，与行业指数的同步性相对有限。

\begin{center}
\begin{minipage}{\textwidth}\centering
\begin{tikzpicture}
\begin{axis}[
  width=12.6cm, height=3.9cm, scale only axis,
  xmin=2020.922, xmax=2026.972, ymin=7.01, ymax=25.72,
  xtick={2021,2022,2023,2024,2025,2026},
  yearticks,
  yticklabel style={font=\scriptsize},
  ylabel={元/股}, ylabel style={font=\scriptsize},
  grid=major, grid style={LightGray, line width=0.3pt},
  axis line style={InkGray, line width=0.5pt},
  every axis plot/.append style={line width=0.9pt},
]
\addplot[AgriGreen] table[x=x,y=y]{data/clean/profiles/px_000998.dat};
\addplot[SoftRed, dashed, line width=0.7pt] table[x=x,y=y]{data/clean/profiles/pxzz_000998.dat};
\addplot[only marks, mark=*, mark size=1.1pt, SoftRed] table[x=x,y=y]{data/clean/profiles/pxzz_000998.dat};
\end{axis}
\end{tikzpicture}

\captionof{figure}[隆平高科（000998）股价与周期骨架]{隆平高科（000998）月度收盘价（元/股）与 ZigZag 周期骨架（阈值 25\%，圆点为周期转折点）}
\end{minipage}
\end{center}

\pfl{过去周期分析} 以 25\% 的 ZigZag 阈值划分，样本区间内股价形成 1 个主升段与 2 个主跌段。 最大主跌段为 2021 年 11 月 至 2026 年 6 月，55 个月-68.6\%；最大主升段出现在 2021 年 6 月 至 2021 年 11 月，5 个月+49.9\%。 最近一次转折出现在 2026 年 6 月（下跌段自 2021 年 11 月起，55 个月-68.6\%），当前处于该段的延续之中。 公司股价与申万农林牧渔指数几乎同步见底，个股并未走出独立行情。 种子行业的价格由制种成本、粮价与政策（转基因商业化、种业振兴）共同决定（第 \ref{sec:planting} 节），与生猪价格的联动有限，公司 2026 年 6 月末资产负债率 59.1\%、半年 ROE -4.1\%（未年化），在本轮下行中属于随行业同步调整的一类。 综合区间形态看，该股单段涨跌幅度偏大、以趋势段为主（最大主升 +49.9\%、最大主跌 -68.6\%），年化波动率 28.7\% 与全样本中位数 37.3\% 相比波动小于样本中位数。

\pfl{盈利情况} 2026 年上半年营业收入 17.10 亿元（同比 -21.1\%），归母净利润 -2.73 亿元，亏损同比扩大 1.09 亿元。 以年报为趋势对照，2023---2025 年营业收入由 92.23 亿元变为 84.77 亿元（两年年均复合增速 -4.1\%），归母净利润由 2.00 亿元变为 1.66 亿元。 按半年报口径，2026 年上半年的 ROE 为 -4.1\%（未年化）、销售净利率 -22.4\%、毛利率 39.1\%，ROE 在“种子”8 家公司中按数值由高到低排第\zhnumber{8}位。 按同一口径，三级行业“种子”8 家公司的半年 ROE 中位数为 0.17\%，该公司低于其中位数 4.26 个百分点。 从公司披露的原因看，业绩变动主要来自：棉花种子受新疆种植面积缩减拖累收入下滑。 针对上述压力，公司披露的应对举措包括：2024 年 9 月，公司商誉账面价值 42.07 亿元。 综合看，公司仍处亏损区间、亏损同比扩大，半年 ROE 在三级行业内按由高到低排第\zhnumber{8}位，单位盈利能力的修复依赖行业景气与自身成本端的改善，报表弹性主要体现为价格与销量的方向性变化。

\begin{center}\small\setlength{\tabcolsep}{3.4pt}
\captionof{table}[隆平高科（000998）关键财务指标]{隆平高科（000998）关键财务指标（合并报表口径；两个半年报期均未年化；现金短债比 $=$ 货币资金 $\div$ 短期债务）}
\begin{adjustbox}{max width=\textwidth}
\begin{tabular}{@{}lrrrrr@{}}
\toprule
指标 & 2023 年 & 2024 年 & 2025 年 & 2025 年半年报 & 2026 年半年报 \\
\midrule
营业收入（亿元） & 92.23 & 85.66 & 84.77 & 21.66 & 17.10 \\
归母净利润（亿元） & 2.00 & 1.14 & 1.66 & -1.64 & -2.73 \\
毛利率（\%） & 39.8 & 36.6 & 39.3 & 36.3 & 39.1 \\
ROE（\%） & 3.2 & 2.2 & 2.9 & -3.1 & -4.1 \\
资产负债率（\%） & 63.5 & 66.5 & 59.5 & 60.0 & 59.1 \\
经营活动现金流净额（亿元） & 7.19 & 4.89 & 7.52 & -9.38 & -8.39 \\
货币资金（亿元） & 39.24 & 27.61 & 29.18 & 16.64 & 24.25 \\
有息负债合计（亿元） & 113.34 & 108.68 & 99.37 & 101.46 & 107.27 \\
现金短债比（倍） & 0.55 & 0.44 & 0.52 & 0.23 & 0.37 \\
\bottomrule
\end{tabular}
\end{adjustbox}
\end{center}

\begin{center}
\begin{minipage}{\textwidth}\centering
\begin{tikzpicture}
\begin{groupplot}[
  group style={group size=2 by 1, horizontal sep=0.60cm},
  width=6.85cm, height=4.60cm,
  tick label style={font=\tiny},
  title style={font=\scriptsize\heiti},
  yticklabel style={font=\tiny},
  ylabel style={font=\tiny},
  grid=major, grid style={LightGray, line width=0.3pt},
  axis line style={InkGray, line width=0.5pt},
  enlarge x limits={abs=0.8},
  legend style={font=\scriptsize, at={(0.5,-0.20)}, anchor=north,
                legend columns=2, draw=none, fill=none, column sep=6pt},
]
\nextgroupplot[
  ybar, bar width=4pt,
  ymin=-18.38, ymax=104.30,
  xtick={0,1,2,3,4,5.7},
  xticklabels={2021,2022,2023,2024,2025,2026H1},
  ylabel={亿元},
]
\addplot[fill=AgriGreen!75, draw=AgriGreen, bar shift=-2.6pt] table[x=i, y=rev]{data/clean/profiles/fin_000998.dat};
\addplot[fill=SoftRed!60, draw=SoftRed, bar shift=2.6pt] table[x=i, y=ni]{data/clean/profiles/fin_000998.dat};
\addplot[fill=AgriGreen!30, draw=AgriGreen, pattern=north east lines, bar shift=-2.6pt] table[x=x, y=rev]{data/clean/profiles/finh_000998.dat};
\addplot[fill=SoftRed!20, draw=SoftRed, pattern=north east lines, bar shift=2.6pt] table[x=x, y=ni]{data/clean/profiles/finh_000998.dat};
\legend{营业收入（年/半年）, 归母净利润（年/半年）}
\nextgroupplot[
  ymin=-20.0, ymax=74.5,
  xtick={0,1,2,3,4,5.7},
  xticklabels={2021,2022,2023,2024,2025,2026H1},
  ylabel={\%},
]
\addplot[AgriGold, mark=*, mark size=1.3pt] table[x=i, y=roe]{data/clean/profiles/fin_000998.dat};
\addplot[OilBlue, mark=square*, mark size=1.3pt] table[x=i, y=debt]{data/clean/profiles/fin_000998.dat};
\addplot[only marks, AgriGold, mark=*, mark size=2.0pt] table[x=x, y=roe]{data/clean/profiles/finh_000998.dat};
\addplot[only marks, OilBlue, mark=square*, mark size=2.0pt] table[x=x, y=debt]{data/clean/profiles/finh_000998.dat};
\legend{ROE, 资产负债率}
\end{groupplot}
\end{tikzpicture}
\begin{tikzpicture}
\begin{axis}[
  name=finbot,
  width=13.75cm, height=3.10cm, scale only axis,
  ymin=-11.33, ymax=126.94,
  xtick={0,1,2,3,4,5.7},
  xticklabels={2021,2022,2023,2024,2025,2026H1},
  tick label style={font=\tiny},
  yticklabel style={font=\tiny},
  ylabel={亿元}, ylabel style={font=\tiny},
  ybar, bar width=4pt,
  grid=major, grid style={LightGray, line width=0.3pt},
  axis line style={InkGray, line width=0.5pt},
  enlarge x limits={abs=0.8},
  legend style={font=\scriptsize, at={(0.5,-0.26)}, anchor=north,
                legend columns=2, draw=none, fill=none, column sep=10pt},
]
\addplot[fill=AgriGreen!55, draw=AgriGreen, bar shift=-3.0pt] table[x=i, y=cash]{data/clean/profiles/fin_000998.dat};
\addplot[fill=SoftRed!45, draw=SoftRed, bar shift=3.0pt] table[x=i, y=ibd]{data/clean/profiles/fin_000998.dat};
\addplot[fill=AgriGreen!25, draw=AgriGreen, pattern=north east lines, bar shift=-3.0pt] table[x=x, y=cash]{data/clean/profiles/finh_000998.dat};
\addplot[fill=SoftRed!20, draw=SoftRed, pattern=north east lines, bar shift=3.0pt] table[x=x, y=ibd]{data/clean/profiles/finh_000998.dat};
\legend{货币资金（左轴）, 有息负债（左轴）}
\end{axis}
\begin{axis}[
  at={(finbot.south east)}, anchor=south east,
  width=13.75cm, height=3.10cm, scale only axis,
  axis y line*=right, axis x line*=none, xtick=\empty,
  ymin=-14.30, ymax=21.15,
  tick label style={font=\tiny},
  yticklabel style={font=\tiny}, scaled y ticks=false,
  legend style={font=\scriptsize, at={(0.5,-0.44)}, anchor=north,
                legend columns=1, draw=none, fill=none},
]
\addplot[OilBlue, mark=*, mark size=2.2pt, line width=1.3pt] table[x=i, y=ocf]{data/clean/profiles/fin_000998.dat};
\addplot[only marks, OilBlue, mark=*, mark size=3.2pt] table[x=x, y=ocf]{data/clean/profiles/finh_000998.dat};
\legend{经营活动现金流净额（右轴，亿元，蓝点）}
\end{axis}
\end{tikzpicture}

\captionof{figure}[隆平高科（000998）近年主要财务指标]{隆平高科（000998）近年主要财务指标（左上：营业收入与归母净利润；右上：ROE 与资产负债率；下：货币资金、有息负债为左轴柱状，经营活动现金流净额为其右轴的蓝点折线。
2021---2025 年为年报口径，2026H1 为 2026 年半年报/半年末，与其前一年报之间留出间隔、以斜纹或空心点区分；半年报数据未年化）}
\end{minipage}
\end{center}

\pfl{财务分析} 2026 年 6 月末资产负债率 59.1\%，较 2025 年末下降 0.4 个百分点（样本中位数 52.2\%，所处三级行业内按由低到高排第\zhnumber{7}位）。 2026 年 6 月末货币资金 24.25 亿元、短期债务（短期借款与一年内到期非流动负债）64.83 亿元、长期债务 42.44 亿元，有息负债合计 107.27 亿元，现金短债比 0.37 倍——覆盖偏紧。 净有息负债 83.02 亿元，相当于其 2026 年上半年营业收入的 486\%。 流动比率 1.06、速动比率 0.60，短期流动性无明显缺口。 2026 年上半年经营活动现金流净额 -8.39 亿元，经营现金流与利润同时为负，现金消耗较快。 2026 年 6 月末在建工程 1.39 亿元，较上年末增加 0.41 亿元，仍有在建投入。 2026 年 6 月末商誉 43.51 亿元，占归母净资产的 42.2\%，是净资产中最脆弱的一块。 综合看，账面货币资金对有息负债与短期债务的覆盖不足，滚动偿付对新增融资的依赖度较高；资产负债率高于样本中位数 6.9 个百分点；有息负债中短期债务占比更高，期限结构仍以滚动续作为主。

\pfl{重大战略转型} 从公开披露看，公司的战略推进集中在：其他动作（3 项）：2024 年 8 月，公司在机构调研中表示正积极应对转基因玉米市场日益增长的机遇；2019 年，成立隆平生物；并购与重组（2 项）：2025 年 3 月，一方面大力处置和剥离非核心资产（含与主业无关的类金融业务）；2017 年，先后收购湖北惠民科技、湖南金稻米业、三瑞农科、河北巡天农业以及巴西陶氏益农特定玉米种子业务。 公告口径上，近三年（2023 年 9 月---2026 年 9 月）公司披露重大重组与并购类公告 4 条、对外投资与转型类公告 0 条、回购与股权激励类公告 4 条，合计公告 399 条。 第一大行业仍是“农业”，收入占比 100.0\%，与 2021 年年报相比没有方向性变化。 综合判断，报告期内公司有 4 条与战略相关的公告落地（重大重组与并购 4 条、对外投资与转型 0 条），资本运作与主业调整之间存在直接关联。

\begin{center}\small\setlength{\tabcolsep}{4pt}
\captionof{table}[隆平高科（000998）历史融资与资本运作]{隆平高科（000998）历史融资与资本运作（据公司公告与公开报道整理；金额为披露值，含首次公开发行、再融资、债券、银行借款与重整投资等）}
\begin{adjustbox}{max width=\textwidth}
\begin{tabular}{@{}p{1.45cm}p{2.55cm}p{4.05cm}p{4.95cm}@{}}
\toprule
时间 & 方式 & 规模 & 用途与说明 \\
\midrule
2016-01 & 定向增发（非公开发行） & 向中信兴业、\allowbreak{}中信建设、\allowbreak{}信农投资、\allowbreak{}现代种业基金和汇添富－优势企业定增计划 5 号资产管理计划合计非公开发行 2.60 亿股 & 发行完成后中信系合计持股 18.72\%，\allowbreak{}中信集团成为实际控制人 \\
2018-03 & 发行股份购买资产 & 向王义波、\allowbreak{}彭泽斌等 45 名自然人发行股份购买联创种业 90\% 股权 & 注入优质玉米种子资产；2018 年 8 月通过证监会审批 \\
2023-08 & 股权受让（现金收购） & 斥资 8.00 亿元受让隆平农业发展股份有限公司 7.14\% 股份 & 获取巴西玉米种子业务，\allowbreak{}扩大营收规模；隆平发展 2023 年实现营收 92.23 亿元 \\
2025-05 & 定向增发（向特定对象发行） & 发行价格 7.87 元/\allowbreak{}股，\allowbreak{}数量不超过 1.52 亿股；实际募集资金人民币 10.35 亿元 & 扣除发行费用后全部用于偿还银行贷款及补充流动资金；发行对象为公司控股股东中信农业科技股份有限公司，\allowbreak{}北京中信信远向中信农业增资 36.00 亿元 \\
\bottomrule
\end{tabular}
\end{adjustbox}
\end{center}

\pfl{历史融投资情况} 从现金流量表看，2025 年公司取得借款收到的现金 102.29 亿元、偿还债务支付的现金 110.19 亿元、吸收权益性投资 22.01 亿元，筹资活动现金流净额 -0.24 亿元（2023 年为取得借款 102.48 亿元、偿还债务 65.61 亿元）。 2026 年上半年筹资活动现金流净额为 3.44 亿元（取得借款 46.34 亿元、偿还债务 38.23 亿元，未年化），仍处于净融入状态。 2025 年分配股利、利润或偿付利息支付的现金合计 6.75 亿元。 公开披露的融资记录共 4 笔，工具类型以定向增发、发行股份购买资产为主；金额与用途见下表。 其中规模最大的两笔为：2025 年 5 月，定向增发，发行价格 7.87 元/股，数量不超过 1.52 亿股，实际募集资金人民币 10.35 亿元；2023 年 8 月，股权受让，斥资 8.00 亿元受让隆平农业发展股份有限公司 7.14\% 股份。 股本方面，近三年披露股本变动 3 次；总股本自 2021 年末以来未发生变化（13.17 亿股）。 分红方面，近三年现金分红 3 次、合计每股派现 0.13 元，最近一次实施在 2025 年 12 月。 近三年“再融资”类公告 47 条，涉及定向增发、募集资金使用。 综合看，公司融资结构上股权与债务融资并举；2025 年筹资活动现金流净额 -0.24 亿元，融资节奏仍以维持存量债务周转为主，与前述经营与投资现金流的方向基本一致。

\pfl{未来潜在融资需求分析} 2026 年 6 月末货币资金 24.25 亿元，而短期债务 64.83 亿元（现金短债比 0.37 倍），2026 年上半年经营现金流净额 -8.39 亿元。按“短期债务减去账面货币资金”的滚动偿付口径，静态缺口约 40.58 亿元；若再计入上半年亏损的年化额（5.45 亿元），维持一年正常经营所需的外部资金约 40.58---46.03 亿元。 公司 2026 年上半年归母净利润为负（-2.73 亿元），但 6 月末资产负债率 59.1\% 尚未进入第一档（亏损且负债率 $>$65\%）的区间，当前融资需求以补充流动资金、支撑低谷期经营性支出为主。 公开披露的后续资金安排线索包括：股权融资线索：2024 年 8 月，公司拟发行不超过 1.52 亿股、募集资金不超过 12.00 亿元，由控股股东中信农业全额认购。 资金需求还要叠加在建工程：期末在建工程 1.39 亿元、2026 年上半年购建固定资产等支出 1.53 亿元（未年化），这部分支出在周期底部通常会被主动放缓。 综合看，公司的外部资金需求以营运资金周转与在建项目投入为主，滚动偿付口径下的静态缺口约 40.58 亿元，能否落地取决于授信续作与发行窗口的配合。




\subsubsection{诺普信（002215）}
\label{co:002215}

\pfl{公司基本情况} 诺普信成立于 1999 年，由农药制剂龙头转型的作物科学公司，以蓝莓为核心的特色生鲜消费业务营收占比已过半，形成“农药制剂＋特色生鲜消费”双轮驱动；14.5 亿元定增已于 2026 年 7 月终止、改以银行长期授信替代，实际控制人个人质押已清零。 公司于 2008-02-18 在深交所主板首发上市，注册资本 10.04 亿元，总股本 10.05 亿股。 截至 2026 年 10 月 12 日收盘价 9.44 元，对应总市值 94.9 亿元，近一年-25.8\%，总市值在三级行业“其他种植业”的 6 家公司中按由大到小排第\zhnumber{2}位，近一年涨跌幅低于全样本中位数（-10.4\%）15.5 个百分点。 按 2026 年半年报披露的行业构成，收入占比前三位为特色生鲜消费业务 52.2\%；农药制剂业务 38.5\%；田田圈业务 9.3\%。 与 2021 年年报相比，第一大行业口径由“主体业务”（59.4\%）变为“特色生鲜消费业务”（52.2\%）；两期披露的分类口径有调整，占比差异中同时包含分类因素。

\pfl{股权结构与控制权} 按 2026 年 6 月 30 日的股东名册，前三大股东为卢柏强 持股 24.42\%；深圳市诺普信投资控股有限公司 持股 4.47\%；华夏基金管理有限公司－社保基金四二二组合（持股比例未披露），股东户数 39,927 户。 股权结构的另一条线索是，实际控制人卢柏强直接持股 24.42\%（2.45 亿股）；2026 年 7 月，原拟募集不超过 14.50 亿元投向云南蓝莓基地扩建 1.35 万亩、小浆果国际研发中心和补充流动资金；2025 年 9 月，公司拟在公告后 15 个交易日后的三个月内减持不超过 3,000 万股。

\pfl{经营要点} 业务与产能层面，蓝莓投产面积由 2021/2022 产季 0.6 万亩、2022/2023 产季 2 万亩、2023/2024 产季 2.8 万亩增至 2024/2025 产季 3.6 万亩；2026 年上半年，研发投入 7,780.36 万元、同比增长 4.89\%。

\pfl{所处行业的基本面} 该子行业按品种各自定价：天然橡胶看全球供给与汽车需求，棉花看新疆产量与配额，农药看原药价格与农业服务需求（第 \ref{sec:framework} 章的供需框架）。 与该公司直接相关的行业数据包括：2020 年，种植面积 158 万亩；2026 年，云南主要产区蓝莓产量同比增 39\%；2024 年，全国蓝莓总产量 78 万吨、比 2020 年增加 197\%，行业市场规模约 500.00 亿元。

\pfl{竞争优势} 把该公司放回同类公司比较，2026 年上半年营业收入 39.6 亿元、在三级行业“其他种植业”的 6 家公司中排第\zhnumber{2}位；毛利率 36.8\%、高于行业内中位数 19.9 个百分点；半年 ROE 16.8\%（未年化）、高于行业内中位数（1.70\%）；资产负债率 60.5\%、高于行业内中位数（59.7\%）。 与同类公司相比，公司的差异化在于：远高于行业平均约 0.8 吨单产有望从 1.5 吨提升至 2 吨；公司实现在价格高峰的 10 月至次年 2 月产果，春节前产量占比提升至约 70\%；KA 渠道出货占比提升至 30\% 以上。 优势在报表上主要体现为高于行业平均的毛利率，这在周期底部是抵御亏损的主要缓冲。

\pfl{历史股价} 本报告样本区间（2021 年 1 月 至 2026 年 10 月）内，股价由 5.30 元变为 9.44 元，累计 +78.1\%，区间最高 12.91 元（2025 年 9 月）、最低 5.15 元（2021 年 7 月）；当前价相当于区间最高价的 73.1\%，为区间最低价的 1.83 倍。 月度收益的年化波动率 33.7\%，低于全样本中位数 37.3\%，在三级行业“其他种植业”的 6 家公司中波动率由高到低排第\zhnumber{3}位。

\begin{center}
\begin{minipage}{\textwidth}\centering
\begin{tikzpicture}
\begin{axis}[
  width=12.6cm, height=3.9cm, scale only axis,
  xmin=2020.922, xmax=2026.972, ymin=4.79, ymax=13.81,
  xtick={2021,2022,2023,2024,2025,2026},
  yearticks,
  yticklabel style={font=\scriptsize},
  ylabel={元/股}, ylabel style={font=\scriptsize},
  grid=major, grid style={LightGray, line width=0.3pt},
  axis line style={InkGray, line width=0.5pt},
  every axis plot/.append style={line width=0.9pt},
]
\addplot[AgriGreen] table[x=x,y=y]{data/clean/profiles/px_002215.dat};
\addplot[SoftRed, dashed, line width=0.7pt] table[x=x,y=y]{data/clean/profiles/pxzz_002215.dat};
\addplot[only marks, mark=*, mark size=1.1pt, SoftRed] table[x=x,y=y]{data/clean/profiles/pxzz_002215.dat};
\end{axis}
\end{tikzpicture}

\captionof{figure}[诺普信（002215）股价与周期骨架]{诺普信（002215）月度收盘价（元/股）与 ZigZag 周期骨架（阈值 25\%，圆点为周期转折点）}
\end{minipage}
\end{center}

\pfl{过去周期分析} 以 25\% 的 ZigZag 阈值划分，样本区间内股价形成 2 个主升段与 3 个主跌段。 最大主跌段为 2025 年 9 月 至 2026 年 6 月，9 个月-36.3\%；最大主升段出现在 2024 年 7 月 至 2025 年 9 月，14 个月+86.3\%。 最近一次转折出现在 2026 年 6 月（下跌段自 2025 年 9 月起，9 个月-36.3\%），当前处于该段的延续之中。 公司股价较申万农林牧渔指数提前约 58 个月见底，是板块中较早定价周期反转的公司之一。 该子行业的周期来源与猪周期不同（第 \ref{sec:grain} 章、第 \ref{sec:soft} 章），价格更多由自身供给、进口政策与全球定价决定，公司 2026 年上半年实现盈利、半年 ROE 16.8\%（未年化），好于全样本中位数（0.75\%），下行周期对其报表的冲击小于同业。 综合区间形态看，该股单段涨跌幅度偏大、以趋势段为主（最大主升 +86.3\%、最大主跌 -36.3\%），年化波动率 33.7\% 与全样本中位数 37.3\% 相比波动与样本中位数接近。

\pfl{盈利情况} 2026 年上半年营业收入 39.64 亿元（同比 +7.8\%），归母净利润 7.65 亿元，同比 +18.0\%。 以年报为趋势对照，2023---2025 年营业收入由 41.20 亿元变为 57.83 亿元（两年年均复合增速 +18.5\%），归母净利润由 2.36 亿元变为 6.50 亿元。 按半年报口径，2026 年上半年的 ROE 为 16.8\%（未年化）、销售净利率 19.2\%、毛利率 36.8\%，ROE 在“其他种植业”6 家公司中按数值由高到低排第\zhnumber{2}位。 同期全样本半年 ROE 中位数 0.75\%，公司与之相差 16.00 个百分点（略好于中位数公司）。 从公司披露的原因看，业绩变动主要来自：2026 年 5 月，公司回应蓝莓行业价格回落：属规模化发展过程中的正常调整。 综合看，公司在报告期维持盈利（高于全样本半年 ROE 中位数 0.75\%），利润的可持续性取决于其他种植业景气的延续与成本管控的执行。

\begin{center}\small\setlength{\tabcolsep}{3.4pt}
\captionof{table}[诺普信（002215）关键财务指标]{诺普信（002215）关键财务指标（合并报表口径；两个半年报期均未年化；现金短债比 $=$ 货币资金 $\div$ 短期债务）}
\begin{adjustbox}{max width=\textwidth}
\begin{tabular}{@{}lrrrrr@{}}
\toprule
指标 & 2023 年 & 2024 年 & 2025 年 & 2025 年半年报 & 2026 年半年报 \\
\midrule
营业收入（亿元） & 41.20 & 52.88 & 57.83 & 36.79 & 39.64 \\
归母净利润（亿元） & 2.36 & 5.85 & 6.50 & 6.48 & 7.65 \\
毛利率（\%） & 29.5 & 37.7 & 37.4 & 37.7 & 36.8 \\
ROE（\%） & 6.4 & 15.4 & 16.4 & 15.8 & 16.8 \\
资产负债率（\%） & 61.1 & 63.7 & 63.3 & 63.7 & 60.5 \\
经营活动现金流净额（亿元） & 6.40 & 13.06 & 14.75 & 2.54 & 0.92 \\
货币资金（亿元） & 13.22 & 14.38 & 11.79 & 15.11 & 12.33 \\
有息负债合计（亿元） & 45.93 & 51.22 & 52.39 & 60.17 & 58.48 \\
现金短债比（倍） & 0.37 & 0.40 & 0.30 & 0.33 & 0.36 \\
\bottomrule
\end{tabular}
\end{adjustbox}
\end{center}

\begin{center}
\begin{minipage}{\textwidth}\centering
\begin{tikzpicture}
\begin{groupplot}[
  group style={group size=2 by 1, horizontal sep=0.60cm},
  width=6.85cm, height=4.60cm,
  tick label style={font=\tiny},
  title style={font=\scriptsize\heiti},
  yticklabel style={font=\tiny},
  ylabel style={font=\tiny},
  grid=major, grid style={LightGray, line width=0.3pt},
  axis line style={InkGray, line width=0.5pt},
  enlarge x limits={abs=0.8},
  legend style={font=\scriptsize, at={(0.5,-0.20)}, anchor=north,
                legend columns=2, draw=none, fill=none, column sep=6pt},
]
\nextgroupplot[
  ybar, bar width=4pt,
  ymin=-5.78, ymax=64.77,
  xtick={0,1,2,3,4,5.7},
  xticklabels={2021,2022,2023,2024,2025,2026H1},
  ylabel={亿元},
]
\addplot[fill=AgriGreen!75, draw=AgriGreen, bar shift=-2.6pt] table[x=i, y=rev]{data/clean/profiles/fin_002215.dat};
\addplot[fill=SoftRed!60, draw=SoftRed, bar shift=2.6pt] table[x=i, y=ni]{data/clean/profiles/fin_002215.dat};
\addplot[fill=AgriGreen!30, draw=AgriGreen, pattern=north east lines, bar shift=-2.6pt] table[x=x, y=rev]{data/clean/profiles/finh_002215.dat};
\addplot[fill=SoftRed!20, draw=SoftRed, pattern=north east lines, bar shift=2.6pt] table[x=x, y=ni]{data/clean/profiles/finh_002215.dat};
\legend{营业收入（年/半年）, 归母净利润（年/半年）}
\nextgroupplot[
  ymin=-5.1, ymax=70.0,
  xtick={0,1,2,3,4,5.7},
  xticklabels={2021,2022,2023,2024,2025,2026H1},
  ylabel={\%},
]
\addplot[AgriGold, mark=*, mark size=1.3pt] table[x=i, y=roe]{data/clean/profiles/fin_002215.dat};
\addplot[OilBlue, mark=square*, mark size=1.3pt] table[x=i, y=debt]{data/clean/profiles/fin_002215.dat};
\addplot[only marks, AgriGold, mark=*, mark size=2.0pt] table[x=x, y=roe]{data/clean/profiles/finh_002215.dat};
\addplot[only marks, OilBlue, mark=square*, mark size=2.0pt] table[x=x, y=debt]{data/clean/profiles/finh_002215.dat};
\legend{ROE, 资产负债率}
\end{groupplot}
\end{tikzpicture}
\begin{tikzpicture}
\begin{axis}[
  name=finbot,
  width=13.75cm, height=3.10cm, scale only axis,
  ymin=-5.85, ymax=65.50,
  xtick={0,1,2,3,4,5.7},
  xticklabels={2021,2022,2023,2024,2025,2026H1},
  tick label style={font=\tiny},
  yticklabel style={font=\tiny},
  ylabel={亿元}, ylabel style={font=\tiny},
  ybar, bar width=4pt,
  grid=major, grid style={LightGray, line width=0.3pt},
  axis line style={InkGray, line width=0.5pt},
  enlarge x limits={abs=0.8},
  legend style={font=\scriptsize, at={(0.5,-0.26)}, anchor=north,
                legend columns=2, draw=none, fill=none, column sep=10pt},
]
\addplot[fill=AgriGreen!55, draw=AgriGreen, bar shift=-3.0pt] table[x=i, y=cash]{data/clean/profiles/fin_002215.dat};
\addplot[fill=SoftRed!45, draw=SoftRed, bar shift=3.0pt] table[x=i, y=ibd]{data/clean/profiles/fin_002215.dat};
\addplot[fill=AgriGreen!25, draw=AgriGreen, pattern=north east lines, bar shift=-3.0pt] table[x=x, y=cash]{data/clean/profiles/finh_002215.dat};
\addplot[fill=SoftRed!20, draw=SoftRed, pattern=north east lines, bar shift=3.0pt] table[x=x, y=ibd]{data/clean/profiles/finh_002215.dat};
\legend{货币资金（左轴）, 有息负债（左轴）}
\end{axis}
\begin{axis}[
  at={(finbot.south east)}, anchor=south east,
  width=13.75cm, height=3.10cm, scale only axis,
  axis y line*=right, axis x line*=none, xtick=\empty,
  ymin=-3.84, ymax=19.18,
  tick label style={font=\tiny},
  yticklabel style={font=\tiny}, scaled y ticks=false,
  legend style={font=\scriptsize, at={(0.5,-0.44)}, anchor=north,
                legend columns=1, draw=none, fill=none},
]
\addplot[OilBlue, mark=*, mark size=2.2pt, line width=1.3pt] table[x=i, y=ocf]{data/clean/profiles/fin_002215.dat};
\addplot[only marks, OilBlue, mark=*, mark size=3.2pt] table[x=x, y=ocf]{data/clean/profiles/finh_002215.dat};
\legend{经营活动现金流净额（右轴，亿元，蓝点）}
\end{axis}
\end{tikzpicture}

\captionof{figure}[诺普信（002215）近年主要财务指标]{诺普信（002215）近年主要财务指标（左上：营业收入与归母净利润；右上：ROE 与资产负债率；下：货币资金、有息负债为左轴柱状，经营活动现金流净额为其右轴的蓝点折线。
2021---2025 年为年报口径，2026H1 为 2026 年半年报/半年末，与其前一年报之间留出间隔、以斜纹或空心点区分；半年报数据未年化）}
\end{minipage}
\end{center}

\pfl{财务分析} 2026 年 6 月末资产负债率 60.5\%，较 2025 年末下降 2.8 个百分点（样本中位数 52.2\%，所处三级行业内按由低到高排第\zhnumber{4}位）。 2026 年 6 月末货币资金 12.33 亿元、短期债务（短期借款与一年内到期非流动负债）33.99 亿元、长期债务 24.49 亿元，有息负债合计 58.48 亿元，现金短债比 0.36 倍——覆盖偏紧。 净有息负债 46.16 亿元，相当于其 2026 年上半年营业收入的 116\%。 流动比率 1.03、速动比率 0.81，短期偿债能力处于正常区间。 2026 年上半年经营活动现金流净额 0.92 亿元，经营现金流与利润同向为正，内生资金可以支撑运营。 2026 年 6 月末在建工程 1.74 亿元，较上年末增加 1.11 亿元，仍有在建投入。 综合看，账面货币资金对有息负债与短期债务的覆盖不足，滚动偿付对新增融资的依赖度较高；资产负债率高于样本中位数 8.3 个百分点；有息负债中短期债务占比更高，期限结构仍以滚动续作为主。 公司披露的监管与合规风险为：2026 年 9 月，公司下属控股公司涉及 2,000.00 万元侵害植物新品种权诉讼。

\pfl{重大战略转型} 从公开披露看，公司的战略推进集中在：融资与资本运作（1 项）：2025 年 11 月，公司启动定增是基于对蓝莓产业长期发展的信心以及把握时间窗口进行战略扩张的内在需求。 公告口径上，近三年（2023 年 9 月---2026 年 9 月）公司披露重大重组与并购类公告 0 条、对外投资与转型类公告 2 条、回购与股权激励类公告 50 条，合计公告 374 条。 主营结构的口径变化已经落到报表上：2021 年年报的第一大行业为“主体业务”（59.4\%），2026 年半年报为“特色生鲜消费业务”（52.2\%）；公司在这两期的分部划分方式有过调整。 综合判断，报告期内公司有 2 条与战略相关的公告落地（重大重组与并购 0 条、对外投资与转型 2 条），资本运作与主业调整之间存在直接关联。

\begin{center}\small\setlength{\tabcolsep}{4pt}
\captionof{table}[诺普信（002215）历史融资与资本运作]{诺普信（002215）历史融资与资本运作（据公司公告与公开报道整理；金额为披露值，含首次公开发行、再融资、债券、银行借款与重整投资等）}
\begin{adjustbox}{max width=\textwidth}
\begin{tabular}{@{}p{1.45cm}p{2.55cm}p{4.05cm}p{4.95cm}@{}}
\toprule
时间 & 方式 & 规模 & 用途与说明 \\
\midrule
2008-02 & IPO（深圳证券交易所中小企业板） & 发行 3,\allowbreak{}000 万股（网下配售 600 万股、\allowbreak{}网上定价发行 2,\allowbreak{}400 万股） & 发行前总股本 9,\allowbreak{}000.00 万股，\allowbreak{}发行后增加 3,\allowbreak{}000.00 万股 \\
2025 & 债务融资规模 & 2025 年筹资活动现金流入小计 33.44 亿元（上年 26.74 亿元） & 日常经营及蓝莓产业投入；公司层面借款规模较大，\allowbreak{}与实控人质押融资为不同口径 \\
2025-11-26 & 定向增发（向特定对象发行）预案 & 募集资金总额不超过 14.50 亿元；拟发行股份数量不超过本次发行前总股本的 30\%（即不超过 3.02 亿股）；发行对象为不超过 35 名特定投资者 & 蓝莓基地新增扩建项目拟用募集资金 11.00 亿元（投资总额 12.21 亿元）、\allowbreak{}小浆果国际研发中心建设项目 1.50 亿元（投资总额 2.74 亿元）、\allowbreak{}补充流动资金 2.00 亿元；已经公司第七届董事会第九次会议（临时）审议通过 \\
2026-04 & 现金分红（2025 年度） & 以 10.05 亿股为基数，\allowbreak{}向全体股东每 10 股派发现金红利 3.00 元（含税） & 回报股东；公司该利润分配预案 \\
\bottomrule
\end{tabular}
\end{adjustbox}
\end{center}

\pfl{历史融投资情况} 从现金流量表看，2025 年公司取得借款收到的现金 65.86 亿元、偿还债务支付的现金 66.78 亿元、吸收权益性投资 0.02 亿元，筹资活动现金流净额 -4.26 亿元（2023 年为取得借款 69.94 亿元、偿还债务 59.81 亿元）。 2026 年上半年筹资活动现金流净额为 -0.20 亿元（取得借款 44.76 亿元、偿还债务 39.85 亿元，未年化），已转为净流出，处于净还债状态。 2025 年分配股利、利润或偿付利息支付的现金合计 4.29 亿元。 公开披露的融资记录共 4 笔，工具类型以 IPO、债务融资规模、定向增发为主；金额与用途见下表。 其中规模最大的两笔为：2025 年 11 月，定向增发，募集资金总额不超过 14.50 亿元，拟发行股份数量不超过本次发行前总股本的 30\%（即不超过 3.02 亿股），发行对象为不超过 35 名特定投资者。 股本方面，近三年披露股本变动 8 次（激励股份解禁 2 次、股份回购 2 次、增发新股上市 1 次、股权激励 1 次）；总股本由 2021 年末的 9.87 亿股变为 10.05 亿股（+1.81\%）。 分红方面，近三年现金分红 4 次、合计每股派现 0.95 元，最近一次实施在 2025 年 12 月。 近三年“再融资”类公告 7 条，涉及定向增发、募集资金使用。 综合看，公司融资结构上以债务滚动为主、股权融资为补充；2025 年筹资活动现金流净额 -4.26 亿元，融资节奏仍以维持存量债务周转为主，与前述经营与投资现金流的方向基本一致。

\pfl{未来潜在融资需求分析} 2026 年 6 月末货币资金 12.33 亿元，而短期债务 33.99 亿元（现金短债比 0.36 倍），2026 年上半年经营现金流净额 0.92 亿元。按“短期债务减去账面货币资金”的滚动偿付口径，静态缺口约 21.66 亿元；上半年盈利或接近盈亏平衡，该缺口不随经营亏损进一步扩大。 按第 \ref{sec:financing-need} 节的三档划分，公司属于第三档“扩张型”：2026 年上半年 ROE 16.8\%（未年化，折合约 33.5\%）、6 月末资产负债率 60.5\%，资产负债表仍具备加杠杆空间；融资用途以产能扩张、品类并购与海外布局为主，单笔规模通常在 5---20 亿元量级，会被市场定价为成长而非补血。 公开披露的后续资金安排线索包括：股权融资线索：2025 年度，公司向特定对象发行 A 股股票预案；2025 年 11 月，公司在投资者关系活动中说明：本次定增为预案。 资金需求还要叠加在建工程：期末在建工程 1.74 亿元、2026 年上半年购建固定资产等支出 3.71 亿元（未年化），这部分支出在周期底部通常会被主动放缓。 综合看，公司的外部资金需求以营运资金周转与在建项目投入为主，滚动偿付口径下的静态缺口约 21.66 亿元，融资的时点与规模更多取决于信用条件与资本市场窗口。




\subsubsection{登海种业（002041）}
\label{co:002041}

\pfl{公司基本情况} 登海种业成立于 2000 年，国内玉米杂交种龙头之一，以“登海”“先玉”系列玉米种为核心，玉米种收入占营业收入约九成，控股股东莱州市农业科学院、李登海为实际控制人；已审定多个转基因玉米品种，2026 年上半年业绩同比明显回升。 公司于 2005-04-18 在深交所主板首发上市，注册资本 8.80 亿元，总股本 8.80 亿股。 截至 2026 年 10 月 12 日收盘价 9.90 元，对应总市值 87.1 亿元，近一年+5.0\%，总市值在三级行业“种子”的 8 家公司中按由大到小排第\zhnumber{2}位，近一年涨跌幅高于全样本中位数（-10.4\%）15.3 个百分点。 按 2026 年半年报披露的行业构成，收入几乎全部来自农业 100.0\%。

\pfl{股权结构与控制权} 按 2026 年 6 月 30 日的股东名册，前三大股东为莱州市农业科学院 持股 53.21\%；李登海 持股 7.35\%；香港中央结算有限公司 持股 0.72\%，股东户数 60,804 户。 股东结构与控制权方面，控股股东莱州市农业科学院持股 4.68 亿股。

\pfl{经营要点} 经营与产能方面，2024---2025 年，业务年度营业总收入 11.81 亿元、同比减少 20.01\%；报告期内公司及各子（孙）公司自主研发并通过国家审定的玉米新品种 12 个，重大在建工程包括登海农科新建厂区项目与育种中心建设项目。

\pfl{所处行业的基本面} 种子的供给由制种面积与库存决定、需求由粮价与种植意愿决定，制种面积在 2025 年起出现萎缩，是本轮供给端最先出清的环节（第 \ref{sec:framework} 章的供需框架）。 从供需两端的当期数据看，2026 年，全国玉米制种面积 358.5 万亩、较上年减少 7\%。

\pfl{竞争优势} 把该公司放回同类公司比较，2026 年上半年营业收入 4.7 亿元、在三级行业“种子”的 8 家公司中排第\zhnumber{6}位；毛利率 30.4\%、高于行业内中位数 1.9 个百分点；半年 ROE 1.3\%（未年化）、高于行业内中位数（0.17\%）；资产负债率 15.8\%、低于行业内中位数（37.3\%）。 与同类公司相比，公司的差异化在于：公司选育的 478、掖 107、DH382、DH351 等 8 个玉米自交系入选“全国杰出贡献玉米自交系”、占 14.55\%。 优势在当期体现为高于行业中位数的资本回报，但优势能否持续仍取决于周期与自身产能投放节奏。

\pfl{历史股价} 本报告样本区间（2021 年 1 月 至 2026 年 10 月）内，股价由 18.52 元变为 9.90 元，累计 -46.5\%，区间最高 26.42 元（2021 年 11 月）、最低 7.87 元（2026 年 6 月）；当前价相当于区间最高价的 37.5\%，为区间最低价的 1.26 倍。 月度收益的年化波动率 33.1\%，低于全样本中位数 37.3\%，在三级行业“种子”的 8 家公司中波动率由高到低排第\zhnumber{7}位。

\begin{center}
\begin{minipage}{\textwidth}\centering
\begin{tikzpicture}
\begin{axis}[
  width=12.6cm, height=3.9cm, scale only axis,
  xmin=2020.922, xmax=2026.972, ymin=7.32, ymax=28.27,
  xtick={2021,2022,2023,2024,2025,2026},
  yearticks,
  yticklabel style={font=\scriptsize},
  ylabel={元/股}, ylabel style={font=\scriptsize},
  grid=major, grid style={LightGray, line width=0.3pt},
  axis line style={InkGray, line width=0.5pt},
  every axis plot/.append style={line width=0.9pt},
]
\addplot[AgriGreen] table[x=x,y=y]{data/clean/profiles/px_002041.dat};
\addplot[SoftRed, dashed, line width=0.7pt] table[x=x,y=y]{data/clean/profiles/pxzz_002041.dat};
\addplot[only marks, mark=*, mark size=1.1pt, SoftRed] table[x=x,y=y]{data/clean/profiles/pxzz_002041.dat};
\end{axis}
\end{tikzpicture}

\captionof{figure}[登海种业（002041）股价与周期骨架]{登海种业（002041）月度收盘价（元/股）与 ZigZag 周期骨架（阈值 25\%，圆点为周期转折点）}
\end{minipage}
\end{center}

\pfl{过去周期分析} 以 25\% 的 ZigZag 阈值划分，样本区间内股价形成 3 个主升段与 3 个主跌段。 最大主升段出现在 2021 年 6 月 至 2021 年 11 月，5 个月+74.9\%；最大主跌段为 2021 年 11 月 至 2024 年 6 月，31 个月-69.4\%。 最近一次转折出现在 2026 年 8 月（上涨段自 2026 年 6 月起，2 个月+37.2\%），当前处于该段之后的震荡之中。 公司股价与申万农林牧渔指数几乎同步见底，个股并未走出独立行情。 种子行业的价格由制种成本、粮价与政策（转基因商业化、种业振兴）共同决定（第 \ref{sec:planting} 节），与生猪价格的联动有限，公司 2026 年上半年实现盈利、半年 ROE 1.3\%（未年化），好于全样本中位数（0.75\%），下行周期对其报表的冲击小于同业。 综合区间形态看，该股单段涨跌幅度偏大、以趋势段为主（最大主升 +74.9\%、最大主跌 -69.4\%），年化波动率 33.1\% 与全样本中位数 37.3\% 相比波动小于样本中位数。

\pfl{盈利情况} 2026 年上半年营业收入 4.69 亿元（同比 +27.0\%），归母净利润 0.50 亿元，同比 +41.0\%。 以年报为趋势对照，2023---2025 年营业收入由 15.52 亿元变为 11.02 亿元（两年年均复合增速 -15.7\%），归母净利润由 2.56 亿元变为 0.92 亿元。 按半年报口径，2026 年上半年的 ROE 为 1.3\%（未年化）、销售净利率 8.9\%、毛利率 30.4\%，ROE 在“种子”8 家公司中按数值由高到低排第\zhnumber{4}位。 按同一口径，三级行业“种子”8 家公司的半年 ROE 中位数为 0.17\%，该公司高于其中位数 1.17 个百分点。 从公司披露的原因看，业绩变动主要来自：2025 年上半年，当期扣非净利润亏损 433.00 万元；2026 年上半年，主要原因是本期借款借贷净额较上年同期减少 2,555.00 万元及上期分红影响。 综合看，公司在报告期维持盈利（高于全样本半年 ROE 中位数 0.75\%），利润的可持续性取决于种子景气的延续与成本管控的执行。

\begin{center}\small\setlength{\tabcolsep}{3.4pt}
\captionof{table}[登海种业（002041）关键财务指标]{登海种业（002041）关键财务指标（合并报表口径；两个半年报期均未年化；现金短债比 $=$ 货币资金 $\div$ 短期债务）}
\begin{adjustbox}{max width=\textwidth}
\begin{tabular}{@{}lrrrrr@{}}
\toprule
指标 & 2023 年 & 2024 年 & 2025 年 & 2025 年半年报 & 2026 年半年报 \\
\midrule
营业收入（亿元） & 15.52 & 12.46 & 11.02 & 3.69 & 4.69 \\
归母净利润（亿元） & 2.56 & 0.57 & 0.92 & 0.35 & 0.50 \\
毛利率（\%） & 28.8 & 28.1 & 29.6 & 27.7 & 30.4 \\
ROE（\%） & 7.4 & 1.6 & 2.5 & 1.0 & 1.3 \\
资产负债率（\%） & 19.1 & 19.1 & 18.7 & 16.9 & 15.8 \\
经营活动现金流净额（亿元） & 1.79 & 1.28 & 2.46 & -0.63 & -0.46 \\
货币资金（亿元） & 13.17 & 15.36 & 23.49 & 17.07 & 21.43 \\
有息负债合计（亿元） & 0.14 & 0.27 & 0.39 & 0.39 & 0.26 \\
现金短债比（倍） & 167.83 & 74.75 & 74.45 & 54.10 & 112.76 \\
\bottomrule
\end{tabular}
\end{adjustbox}
\end{center}

\begin{center}
\begin{minipage}{\textwidth}\centering
\begin{tikzpicture}
\begin{groupplot}[
  group style={group size=2 by 1, horizontal sep=0.60cm},
  width=6.85cm, height=4.60cm,
  tick label style={font=\tiny},
  title style={font=\scriptsize\heiti},
  yticklabel style={font=\tiny},
  ylabel style={font=\tiny},
  grid=major, grid style={LightGray, line width=0.3pt},
  axis line style={InkGray, line width=0.5pt},
  enlarge x limits={abs=0.8},
  legend style={font=\scriptsize, at={(0.5,-0.20)}, anchor=north,
                legend columns=2, draw=none, fill=none, column sep=6pt},
]
\nextgroupplot[
  ybar, bar width=4pt,
  ymin=-1.55, ymax=17.38,
  xtick={0,1,2,3,4,5.7},
  xticklabels={2021,2022,2023,2024,2025,2026H1},
  ylabel={亿元},
]
\addplot[fill=AgriGreen!75, draw=AgriGreen, bar shift=-2.6pt] table[x=i, y=rev]{data/clean/profiles/fin_002041.dat};
\addplot[fill=SoftRed!60, draw=SoftRed, bar shift=2.6pt] table[x=i, y=ni]{data/clean/profiles/fin_002041.dat};
\addplot[fill=AgriGreen!30, draw=AgriGreen, pattern=north east lines, bar shift=-2.6pt] table[x=x, y=rev]{data/clean/profiles/finh_002041.dat};
\addplot[fill=SoftRed!20, draw=SoftRed, pattern=north east lines, bar shift=2.6pt] table[x=x, y=ni]{data/clean/profiles/finh_002041.dat};
\legend{营业收入（年/半年）, 归母净利润（年/半年）}
\nextgroupplot[
  ymin=-1.8, ymax=24.7,
  xtick={0,1,2,3,4,5.7},
  xticklabels={2021,2022,2023,2024,2025,2026H1},
  ylabel={\%},
]
\addplot[AgriGold, mark=*, mark size=1.3pt] table[x=i, y=roe]{data/clean/profiles/fin_002041.dat};
\addplot[OilBlue, mark=square*, mark size=1.3pt] table[x=i, y=debt]{data/clean/profiles/fin_002041.dat};
\addplot[only marks, AgriGold, mark=*, mark size=2.0pt] table[x=x, y=roe]{data/clean/profiles/finh_002041.dat};
\addplot[only marks, OilBlue, mark=square*, mark size=2.0pt] table[x=x, y=debt]{data/clean/profiles/finh_002041.dat};
\legend{ROE, 资产负债率}
\end{groupplot}
\end{tikzpicture}
\begin{tikzpicture}
\begin{axis}[
  name=finbot,
  width=13.75cm, height=3.10cm, scale only axis,
  ymin=-2.35, ymax=26.31,
  xtick={0,1,2,3,4,5.7},
  xticklabels={2021,2022,2023,2024,2025,2026H1},
  tick label style={font=\tiny},
  yticklabel style={font=\tiny},
  ylabel={亿元}, ylabel style={font=\tiny},
  ybar, bar width=4pt,
  grid=major, grid style={LightGray, line width=0.3pt},
  axis line style={InkGray, line width=0.5pt},
  enlarge x limits={abs=0.8},
  legend style={font=\scriptsize, at={(0.5,-0.26)}, anchor=north,
                legend columns=2, draw=none, fill=none, column sep=10pt},
]
\addplot[fill=AgriGreen!55, draw=AgriGreen, bar shift=-3.0pt] table[x=i, y=cash]{data/clean/profiles/fin_002041.dat};
\addplot[fill=SoftRed!45, draw=SoftRed, bar shift=3.0pt] table[x=i, y=ibd]{data/clean/profiles/fin_002041.dat};
\addplot[fill=AgriGreen!25, draw=AgriGreen, pattern=north east lines, bar shift=-3.0pt] table[x=x, y=cash]{data/clean/profiles/finh_002041.dat};
\addplot[fill=SoftRed!20, draw=SoftRed, pattern=north east lines, bar shift=3.0pt] table[x=x, y=ibd]{data/clean/profiles/finh_002041.dat};
\legend{货币资金（左轴）, 有息负债（左轴）}
\end{axis}
\begin{axis}[
  at={(finbot.south east)}, anchor=south east,
  width=13.75cm, height=3.10cm, scale only axis,
  axis y line*=right, axis x line*=none, xtick=\empty,
  ymin=-1.53, ymax=4.89,
  tick label style={font=\tiny},
  yticklabel style={font=\tiny}, scaled y ticks=false,
  legend style={font=\scriptsize, at={(0.5,-0.44)}, anchor=north,
                legend columns=1, draw=none, fill=none},
]
\addplot[OilBlue, mark=*, mark size=2.2pt, line width=1.3pt] table[x=i, y=ocf]{data/clean/profiles/fin_002041.dat};
\addplot[only marks, OilBlue, mark=*, mark size=3.2pt] table[x=x, y=ocf]{data/clean/profiles/finh_002041.dat};
\legend{经营活动现金流净额（右轴，亿元，蓝点）}
\end{axis}
\end{tikzpicture}

\captionof{figure}[登海种业（002041）近年主要财务指标]{登海种业（002041）近年主要财务指标（左上：营业收入与归母净利润；右上：ROE 与资产负债率；下：货币资金、有息负债为左轴柱状，经营活动现金流净额为其右轴的蓝点折线。
2021---2025 年为年报口径，2026H1 为 2026 年半年报/半年末，与其前一年报之间留出间隔、以斜纹或空心点区分；半年报数据未年化）}
\end{minipage}
\end{center}

\pfl{财务分析} 2026 年 6 月末资产负债率 15.8\%，较 2025 年末下降 2.9 个百分点（样本中位数 52.2\%，所处三级行业内按由低到高排第\zhnumber{2}位）。 2026 年 6 月末货币资金 21.43 亿元、短期债务（短期借款与一年内到期非流动负债）0.19 亿元、长期债务 0.07 亿元，有息负债合计 0.26 亿元，现金短债比 112.76 倍——覆盖充分。 净有息负债 -21.16 亿元，相当于其 2026 年上半年营业收入的 -451\%。 流动比率 6.34、速动比率 5.53，短期偿债能力处于正常区间。 2026 年上半年经营活动现金流净额 -0.46 亿元，净利润为正但经营现金流为负，利润的现金含量偏低。 2026 年 6 月末在建工程 0.52 亿元，较上年末增加 0.05 亿元，仍有在建投入。 综合看，现金短债比 112.76 倍，短期偿付压力可控；资产负债率低于样本中位数 36.4 个百分点；有息负债中短期债务占比更高，期限结构仍以滚动续作为主。

\pfl{重大战略转型} 公告口径上，近三年（2023 年 9 月---2026 年 9 月）公司披露重大重组与并购类公告 0 条、对外投资与转型类公告 1 条、回购与股权激励类公告 0 条，合计公告 151 条。 第一大行业仍是“农业”，收入占比 100.0\%，与 2021 年年报相比没有方向性变化。 综合判断，报告期内公司有 1 条与战略相关的公告落地（重大重组与并购 0 条、对外投资与转型 1 条），资本运作与主业调整之间存在直接关联。

\begin{center}\small\setlength{\tabcolsep}{4pt}
\captionof{table}[登海种业（002041）历史融资与资本运作]{登海种业（002041）历史融资与资本运作（据公司公告与公开报道整理；金额为披露值，含首次公开发行、再融资、债券、银行借款与重整投资等）}
\begin{adjustbox}{max width=\textwidth}
\begin{tabular}{@{}p{1.45cm}p{2.55cm}p{4.05cm}p{4.95cm}@{}}
\toprule
时间 & 方式 & 规模 & 用途与说明 \\
\midrule
2005 年以来 & 现金分红（累计） & A 股上市后累计派现 5.84 亿元；近三年累计派现 5,\allowbreak{}720.00 万元 & 回报股东；2026 年上半年不派发现金红利、\allowbreak{}不送红股、\allowbreak{}不以公积金转增股本 \\
2005—2026 & IPO 后的股权／债权再融资 & --- & 本次检索未发现公司 IPO 之后的定增、\allowbreak{}可转债或公司债记录；2015 年年度报告披露“公司报告期无募集资金使用情况” \\
2005-04 & IPO（深圳证券交易所中小企业板） & 发行 2,\allowbreak{}200 万股，\allowbreak{}发行价 16.70 元/\allowbreak{}股 & 发行费用明细：承销费用 1,\allowbreak{}102.20 万元、\allowbreak{}保荐费用 330.00 万元、\allowbreak{}审核费用 3.00 万元、\allowbreak{}会计师费用 180.00 万元、\allowbreak{}律师费用 75.00 万元、\allowbreak{}上网发行费用 97.36 万元；验资确认截至 2005-04-12 募集资金总额 3.67 亿元 \\
\bottomrule
\end{tabular}
\end{adjustbox}
\end{center}

\pfl{历史融投资情况} 从现金流量表看，2025 年公司取得借款收到的现金 0.50 亿元、偿还债务支付的现金 0.37 亿元，筹资活动现金流净额 -0.16 亿元（2023 年为取得借款 0.13 亿元、偿还债务 0.07 亿元）。 2026 年上半年筹资活动现金流净额为 -0.14 亿元（取得借款 0.19 亿元、偿还债务 0.32 亿元，未年化），已转为净流出，处于净还债状态。 2025 年分配股利、利润或偿付利息支付的现金合计 0.27 亿元。 公开披露的融资记录共 3 笔，工具类型以 IPO 为主；金额与用途见下表。 其中规模最大的两笔为：2005 年 4 月，IPO，发行 2,200 万股，发行价 16.70 元/股；2005 年，现金分红，A 股上市后累计派现 5.84 亿元，近三年累计派现 5,720.00 万元。 股本方面，近三年披露股本变动 3 次；总股本自 2021 年末以来未发生变化（8.80 亿股）。 分红方面，近三年现金分红 2 次、合计每股派现 0.06 元，最近一次实施在 2024 年 12 月。 综合看，公司融资结构上股权与债务融资并举；2025 年筹资活动现金流净额 -0.16 亿元，融资节奏仍以维持存量债务周转为主，与前述经营与投资现金流的方向基本一致。

\pfl{未来潜在融资需求分析} 2026 年 6 月末货币资金 21.43 亿元，而短期债务 0.19 亿元（现金短债比 112.76 倍），2026 年上半年经营现金流净额 -0.46 亿元。账面货币资金可以覆盖短期债务，缺口不体现在静态偿付层面。 公司 2026 年 6 月末资产负债率 15.8\%、2026 年上半年 ROE 1.3\%（未年化），资产端与负债端都没有明显压力，融资需求以相机抉择的扩张型用途为主。 公开披露的后续资金安排线索包括：债务与授信安排：2026 年 6 月，本期借款借贷净额较上年同期减少 2,555.00 万元。 资金需求还要叠加在建工程：期末在建工程 0.52 亿元、2026 年上半年购建固定资产等支出 0.39 亿元（未年化），这部分支出在周期底部通常会被主动放缓。 静态偿付不存在缺口，后续融资更可能服务于产能或并购安排，而不是填补流动性窟窿。




\subsubsection{亚盛集团（600108）}
\label{co:600108}

\pfl{公司基本情况} 公司前身可追溯至 1995 年，甘肃省农垦集团控股的现代农业企业集团，种植基地主要分布在甘肃、内蒙，以牧草、马铃薯、制种玉米等规模化种植为基础并延伸加工与农资服务；2026 年上半年盈利微降，经营现金流净流出扩大，资金缺口主要靠借款与中期票据填补。 公司于 1997-08-18 在上交所首发上市，注册资本 19.47 亿元，总股本 19.47 亿股。 截至 2026 年 10 月 12 日收盘价 4.03 元，对应总市值 78.5 亿元，近一年+31.3\%，总市值在三级行业“其他种植业”的 6 家公司中按由大到小排第\zhnumber{3}位，近一年涨跌幅高于全样本中位数（-10.4\%）41.6 个百分点。 按 2026 年半年报披露的行业构成，收入占比前三位为农业 57.1\%；商贸 36.1\%；工业 6.0\%。 股权结构方面，控股股东为甘肃省农垦集团有限责任公司。

\pfl{股权结构与控制权} 按 2026 年 6 月 30 日的股东名册，前三大股东为甘肃省农垦集团有限责任公司 持股 19.08\%；甘肃省农垦资产经营有限公司 持股 9.01\%；香港中央结算有限公司 持股 1.60\%，股东户数 125,017 户。 股权结构的另一条线索是，控股股东所持股份不存在质押、标记或冻结。

\pfl{经营要点} 经营层面，公司主导作物种植面积占总面积的 90\%；子公司中甘肃亚盛田园牧歌草业集团上半年净利润 362.25 万元、甘肃兴农辣椒产业开发 908.79 万元、甘肃亚盛农业综合服务 230.74 万元。

\pfl{所处行业的基本面} 该子行业按品种各自定价：天然橡胶看全球供给与汽车需求，棉花看新疆产量与配额，农药看原药价格与农业服务需求（第 \ref{sec:framework} 章的供需框架）。 从供需两端的当期数据看，截至 2026 年 10 月 7 日甘肃秋粮收获进度 48.38\%，其中玉米收获 802.06 万亩、马铃薯收获 455.96 万亩。

\pfl{竞争优势} 把该公司放回同类公司比较，2026 年上半年营业收入 14.7 亿元、在三级行业“其他种植业”的 6 家公司中排第\zhnumber{4}位；毛利率 17.6\%、高于行业内中位数 0.7 个百分点；半年 ROE 0.6\%（未年化）、低于行业内中位数（1.70\%）；资产负债率 58.9\%、低于行业内中位数（59.7\%）。 与同类公司相比，公司的差异化在于：土地集约化经营程度高，为牧草、制种玉米等作物提供基地支撑，在生产成本和产品品质上具备竞争力；报告期新获授权专利 7 项。 就当期报表看，公司尚未把经营层面的特征转化为高于行业的资本回报，周期底部的相对优势更多体现在现金流与资产负债表上。

\pfl{历史股价} 以 2021 年 1 月为起点、2026 年 10 月为终点，股价由 2.94 元变为 4.03 元，累计 +37.1\%，区间最高 4.55 元（2026 年 3 月）、最低 2.37 元（2024 年 6 月）；当前价相当于区间最高价的 88.6\%，为区间最低价的 1.70 倍。 月度收益的年化波动率 27.8\%，低于全样本中位数 37.3\%，在三级行业“其他种植业”的 6 家公司中波动率由高到低排第\zhnumber{6}位。

\begin{center}
\begin{minipage}{\textwidth}\centering
\begin{tikzpicture}
\begin{axis}[
  width=12.6cm, height=3.9cm, scale only axis,
  xmin=2020.922, xmax=2026.972, ymin=2.20, ymax=4.87,
  xtick={2021,2022,2023,2024,2025,2026},
  yearticks,
  yticklabel style={font=\scriptsize},
  ylabel={元/股}, ylabel style={font=\scriptsize},
  grid=major, grid style={LightGray, line width=0.3pt},
  axis line style={InkGray, line width=0.5pt},
  every axis plot/.append style={line width=0.9pt},
]
\addplot[AgriGreen] table[x=x,y=y]{data/clean/profiles/px_600108.dat};
\addplot[SoftRed, dashed, line width=0.7pt] table[x=x,y=y]{data/clean/profiles/pxzz_600108.dat};
\addplot[only marks, mark=*, mark size=1.1pt, SoftRed] table[x=x,y=y]{data/clean/profiles/pxzz_600108.dat};
\end{axis}
\end{tikzpicture}

\captionof{figure}[亚盛集团（600108）股价与周期骨架]{亚盛集团（600108）月度收盘价（元/股）与 ZigZag 周期骨架（阈值 25\%，圆点为周期转折点）}
\end{minipage}
\end{center}

\pfl{过去周期分析} 以 25\% 的 ZigZag 阈值划分，样本区间内股价形成 3 个主升段与 3 个主跌段。 最大主跌段为 2022 年 5 月 至 2024 年 6 月，25 个月-39.7\%；最大主升段出现在 2024 年 6 月 至 2026 年 3 月，21 个月+92.0\%。 最近一次转折出现在 2026 年 10 月（上涨段自 2026 年 6 月起，4 个月+36.1\%），当前处于该段之后的震荡之中。 公司股价较申万农林牧渔指数提前约 23 个月见底，是板块中较早定价周期反转的公司之一。 该子行业的周期来源与猪周期不同（第 \ref{sec:grain} 章、第 \ref{sec:soft} 章），价格更多由自身供给、进口政策与全球定价决定，公司 2026 年 6 月末资产负债率 58.9\%、半年 ROE 0.6\%（未年化），在本轮下行中属于随行业同步调整的一类。 综合区间形态看，该股单段涨跌幅度偏大、以趋势段为主（最大主升 +92.0\%、最大主跌 -39.7\%），年化波动率 27.8\% 与全样本中位数 37.3\% 相比波动小于样本中位数。

\pfl{盈利情况} 2026 年上半年营业收入 14.72 亿元（同比 -2.6\%），归母净利润 0.25 亿元，同比 -11.0\%。 以年报为趋势对照，2023---2025 年营业收入由 40.05 亿元变为 42.42 亿元（两年年均复合增速 +2.9\%），归母净利润由 1.04 亿元变为 0.57 亿元。 按半年报口径，2026 年上半年的 ROE 为 0.6\%（未年化）、销售净利率 1.7\%、毛利率 17.6\%，ROE 在“其他种植业”6 家公司中按数值由高到低排第\zhnumber{4}位。 同期全样本半年 ROE 中位数 0.75\%，公司与之相差 0.17 个百分点（弱于中位数公司）。 面对这一局面，公司披露的举措包括：公司增资后该公司注册资本将变更为 1.00 亿元，以提升马铃薯加工能力与产能规模效益。 综合看，公司在报告期维持盈利（低于全样本半年 ROE 中位数 0.75\%），利润的可持续性取决于其他种植业景气的延续与成本管控的执行。

\begin{center}\small\setlength{\tabcolsep}{3.4pt}
\captionof{table}[亚盛集团（600108）关键财务指标]{亚盛集团（600108）关键财务指标（合并报表口径；两个半年报期均未年化；现金短债比 $=$ 货币资金 $\div$ 短期债务）}
\begin{adjustbox}{max width=\textwidth}
\begin{tabular}{@{}lrrrrr@{}}
\toprule
指标 & 2023 年 & 2024 年 & 2025 年 & 2025 年半年报 & 2026 年半年报 \\
\midrule
营业收入（亿元） & 40.05 & 41.92 & 42.42 & 15.12 & 14.72 \\
归母净利润（亿元） & 1.04 & 0.96 & 0.57 & 0.28 & 0.25 \\
毛利率（\%） & 17.8 & 17.5 & 15.5 & 17.7 & 17.6 \\
ROE（\%） & 2.5 & 2.3 & 1.3 & 0.7 & 0.6 \\
资产负债率（\%） & 53.5 & 55.3 & 56.3 & 55.8 & 58.9 \\
经营活动现金流净额（亿元） & 1.60 & 1.24 & 1.50 & -1.03 & -1.67 \\
货币资金（亿元） & 6.64 & 4.88 & 3.40 & 3.11 & 5.31 \\
有息负债合计（亿元） & 35.45 & 38.78 & 41.09 & 40.65 & 46.72 \\
现金短债比（倍） & 0.54 & 0.35 & 0.12 & 0.20 & 0.20 \\
\bottomrule
\end{tabular}
\end{adjustbox}
\end{center}

\begin{center}
\begin{minipage}{\textwidth}\centering
\begin{tikzpicture}
\begin{groupplot}[
  group style={group size=2 by 1, horizontal sep=0.60cm},
  width=6.85cm, height=4.60cm,
  tick label style={font=\tiny},
  title style={font=\scriptsize\heiti},
  yticklabel style={font=\tiny},
  ylabel style={font=\tiny},
  grid=major, grid style={LightGray, line width=0.3pt},
  axis line style={InkGray, line width=0.5pt},
  enlarge x limits={abs=0.8},
  legend style={font=\scriptsize, at={(0.5,-0.20)}, anchor=north,
                legend columns=2, draw=none, fill=none, column sep=6pt},
]
\nextgroupplot[
  ybar, bar width=4pt,
  ymin=-4.24, ymax=47.51,
  xtick={0,1,2,3,4,5.7},
  xticklabels={2021,2022,2023,2024,2025,2026H1},
  ylabel={亿元},
]
\addplot[fill=AgriGreen!75, draw=AgriGreen, bar shift=-2.6pt] table[x=i, y=rev]{data/clean/profiles/fin_600108.dat};
\addplot[fill=SoftRed!60, draw=SoftRed, bar shift=2.6pt] table[x=i, y=ni]{data/clean/profiles/fin_600108.dat};
\addplot[fill=AgriGreen!30, draw=AgriGreen, pattern=north east lines, bar shift=-2.6pt] table[x=x, y=rev]{data/clean/profiles/finh_600108.dat};
\addplot[fill=SoftRed!20, draw=SoftRed, pattern=north east lines, bar shift=2.6pt] table[x=x, y=ni]{data/clean/profiles/finh_600108.dat};
\legend{营业收入（年/半年）, 归母净利润（年/半年）}
\nextgroupplot[
  ymin=-4.7, ymax=64.8,
  xtick={0,1,2,3,4,5.7},
  xticklabels={2021,2022,2023,2024,2025,2026H1},
  ylabel={\%},
]
\addplot[AgriGold, mark=*, mark size=1.3pt] table[x=i, y=roe]{data/clean/profiles/fin_600108.dat};
\addplot[OilBlue, mark=square*, mark size=1.3pt] table[x=i, y=debt]{data/clean/profiles/fin_600108.dat};
\addplot[only marks, AgriGold, mark=*, mark size=2.0pt] table[x=x, y=roe]{data/clean/profiles/finh_600108.dat};
\addplot[only marks, OilBlue, mark=square*, mark size=2.0pt] table[x=x, y=debt]{data/clean/profiles/finh_600108.dat};
\legend{ROE, 资产负债率}
\end{groupplot}
\end{tikzpicture}
\begin{tikzpicture}
\begin{axis}[
  name=finbot,
  width=13.75cm, height=3.10cm, scale only axis,
  ymin=-4.67, ymax=52.33,
  xtick={0,1,2,3,4,5.7},
  xticklabels={2021,2022,2023,2024,2025,2026H1},
  tick label style={font=\tiny},
  yticklabel style={font=\tiny},
  ylabel={亿元}, ylabel style={font=\tiny},
  ybar, bar width=4pt,
  grid=major, grid style={LightGray, line width=0.3pt},
  axis line style={InkGray, line width=0.5pt},
  enlarge x limits={abs=0.8},
  legend style={font=\scriptsize, at={(0.5,-0.26)}, anchor=north,
                legend columns=2, draw=none, fill=none, column sep=10pt},
]
\addplot[fill=AgriGreen!55, draw=AgriGreen, bar shift=-3.0pt] table[x=i, y=cash]{data/clean/profiles/fin_600108.dat};
\addplot[fill=SoftRed!45, draw=SoftRed, bar shift=3.0pt] table[x=i, y=ibd]{data/clean/profiles/fin_600108.dat};
\addplot[fill=AgriGreen!25, draw=AgriGreen, pattern=north east lines, bar shift=-3.0pt] table[x=x, y=cash]{data/clean/profiles/finh_600108.dat};
\addplot[fill=SoftRed!20, draw=SoftRed, pattern=north east lines, bar shift=3.0pt] table[x=x, y=ibd]{data/clean/profiles/finh_600108.dat};
\legend{货币资金（左轴）, 有息负债（左轴）}
\end{axis}
\begin{axis}[
  at={(finbot.south east)}, anchor=south east,
  width=13.75cm, height=3.10cm, scale only axis,
  axis y line*=right, axis x line*=none, xtick=\empty,
  ymin=-2.78, ymax=3.86,
  tick label style={font=\tiny},
  yticklabel style={font=\tiny}, scaled y ticks=false,
  legend style={font=\scriptsize, at={(0.5,-0.44)}, anchor=north,
                legend columns=1, draw=none, fill=none},
]
\addplot[OilBlue, mark=*, mark size=2.2pt, line width=1.3pt] table[x=i, y=ocf]{data/clean/profiles/fin_600108.dat};
\addplot[only marks, OilBlue, mark=*, mark size=3.2pt] table[x=x, y=ocf]{data/clean/profiles/finh_600108.dat};
\legend{经营活动现金流净额（右轴，亿元，蓝点）}
\end{axis}
\end{tikzpicture}

\captionof{figure}[亚盛集团（600108）近年主要财务指标]{亚盛集团（600108）近年主要财务指标（左上：营业收入与归母净利润；右上：ROE 与资产负债率；下：货币资金、有息负债为左轴柱状，经营活动现金流净额为其右轴的蓝点折线。
2021---2025 年为年报口径，2026H1 为 2026 年半年报/半年末，与其前一年报之间留出间隔、以斜纹或空心点区分；半年报数据未年化）}
\end{minipage}
\end{center}

\pfl{财务分析} 2026 年 6 月末资产负债率 58.9\%，较 2025 年末上升 2.6 个百分点（样本中位数 52.2\%，所处三级行业内按由低到高排第\zhnumber{3}位）。 2026 年 6 月末货币资金 5.31 亿元、短期债务（短期借款与一年内到期非流动负债）27.04 亿元、长期债务 19.68 亿元，有息负债合计 46.72 亿元，现金短债比 0.20 倍——缺口明显。 净有息负债 41.41 亿元，相当于其 2026 年上半年营业收入的 281\%。 流动比率 1.47、速动比率 0.90，短期偿债能力处于正常区间。 2026 年上半年经营活动现金流净额 -1.67 亿元，净利润为正但经营现金流为负，利润的现金含量偏低。 2026 年 6 月末在建工程 4.70 亿元，较上年末增加 1.00 亿元，仍有在建投入。 综合看，账面货币资金对有息负债与短期债务的覆盖不足，滚动偿付对新增融资的依赖度较高；资产负债率高于样本中位数 6.7 个百分点；有息负债中短期债务占比更高，期限结构仍以滚动续作为主。 公司披露的重整与债务违约风险为：2025 年 6 月，公司对甘肃证券有限公司股权投资成本 2,000.00 万元、甘肃普华甜菊糖有限公司股权投资成本 464.25 万元。

\pfl{重大战略转型} 公司的战略动作可归为以下几类：产能与项目（2 项）：2001 年 11 月，甘肃省财政厅拨付公司国债转贷资金 1,655.00 万元；并购与重组（1 项）：2009 年 3 月，2009-03-09 签署购买资产协议。 公告口径上，近三年（2023 年 9 月---2026 年 9 月）公司披露重大重组与并购类公告 2 条、对外投资与转型类公告 6 条、回购与股权激励类公告 0 条，合计公告 300 条。 第一大行业“农业”的收入占比由 2021 年年报的 81.2\% 变为 57.1\%（24.1 个百分点），收入结构出现再平衡。 综合判断，报告期内公司有 8 条与战略相关的公告落地（重大重组与并购 2 条、对外投资与转型 6 条），资本运作与主业调整之间存在直接关联。

\begin{center}\small\setlength{\tabcolsep}{4pt}
\captionof{table}[亚盛集团（600108）历史融资与资本运作]{亚盛集团（600108）历史融资与资本运作（据公司公告与公开报道整理；金额为披露值，含首次公开发行、再融资、债券、银行借款与重整投资等）}
\begin{adjustbox}{max width=\textwidth}
\begin{tabular}{@{}p{1.45cm}p{2.55cm}p{4.05cm}p{4.95cm}@{}}
\toprule
时间 & 方式 & 规模 & 用途与说明 \\
\midrule
1997-07 & IPO（上海证券交易所） & 公开发行股票 7,\allowbreak{}000 万股 & 以全额预交款、\allowbreak{}比例配售、\allowbreak{}余款即退方式发行；1997-08-18 在上交所挂牌 \\
2009-12 & 发行股份购买资产（定向增发） & 向甘肃省农垦集团有限责任公司非公开发行人民币普通股（A 股）2.96 亿股；同花顺资本运作口径实际募集资金净额 11.57 亿元 & 认购甘肃农垦全资拥有的七家农场全部农业类资产及十宗土地，\allowbreak{}总股本由 14.41 亿股增至 17.37 亿股 \\
2012-05 & 增发 A 股 & 同花顺资本运作口径实际募集资金净额 10.91 亿元（发行起始日 2012-05-08） & 同花顺“资本运作”页披露的历次增发记录 \\
2023-11 & 中期票据（银行间债券市场） & 甘肃亚盛实业（集团）股份有限公司 2023 年度第一期中期票据（乡村振兴）8.00 亿元 & 简称“23 亚盛实业 MTN001（乡村振兴）”，\allowbreak{}实际兑付付息金额 3,\allowbreak{}200.00 万元 \\
2026-04-30 & 现金分红（2025 年度） & 以 2025-12-31 总股本 19.47 亿股为基数 & 回报股东；2025 年度合并报表归属于母公司股东的净利润 5,\allowbreak{}693.36 万元；截至 2025-12-31 母公司报表期末未分配利润 14.48 亿元；该预案尚需 2025 年年度股东会审议通过 \\
2026-06-30 & 计息负债扩张 & 应付债券期末余额 11.10 亿元 & 补充经营资金缺口；公司说明系报告期内借款增加；一年内到期的应付债券 8.19 亿元 \\
\bottomrule
\end{tabular}
\end{adjustbox}
\end{center}

\pfl{历史融投资情况} 从现金流量表看，2025 年公司取得借款收到的现金 19.74 亿元、偿还债务支付的现金 17.39 亿元，筹资活动现金流净额 0.49 亿元（2023 年为取得借款 26.66 亿元、偿还债务 32.02 亿元）。 2026 年上半年筹资活动现金流净额为 4.90 亿元（取得借款 14.51 亿元、偿还债务 15.78 亿元，未年化），仍处于净融入状态。 2025 年分配股利、利润或偿付利息支付的现金合计 1.42 亿元。 公开披露的融资记录共 6 笔，工具类型以 IPO、发行股份购买资产、定向增发、中期票据为主；金额与用途见下表。 其中规模最大的两笔为：2026 年 6 月，计息负债扩张，应付债券期末余额 11.10 亿元；2023 年 11 月，中期票据，甘肃亚盛实业（集团）股份有限公司 2023 年度第一期中期票据（乡村振兴）8.00 亿元。 股本方面，近三年披露股本变动 3 次；总股本自 2021 年末以来未发生变化（19.47 亿股）。 分红方面，近三年现金分红 4 次、合计每股派现 0.02 元，最近一次实施在 2025 年 12 月。 综合看，公司融资结构上股权与债务融资并举；2025 年筹资活动现金流净额 0.49 亿元，年度内仍为净融入，与前述经营与投资现金流的方向基本一致。

\pfl{未来潜在融资需求分析} 2026 年 6 月末货币资金 5.31 亿元，而短期债务 27.04 亿元（现金短债比 0.20 倍），2026 年上半年经营现金流净额 -1.67 亿元。按“短期债务减去账面货币资金”的滚动偿付口径，静态缺口约 21.73 亿元；上半年盈利或接近盈亏平衡，该缺口不随经营亏损进一步扩大。 公司 2026 年 6 月末资产负债率 58.9\%、2026 年上半年 ROE 0.6\%（未年化），现金短债比 0.20 倍、短债覆盖不足，融资需求首先用于置换与滚动存量债务、补充营运资金，属于“补血型”而非扩张型。 公开披露的后续资金安排线索包括：债务与授信安排：2026 年 6 月，筹资活动现金流量净额 4.90 亿元（同比增长 704.58\%）、应付债券余额增加 174.85\% 公司对外部债务融资依赖上升；2026 年 9 月，公司为子公司提供的最高额信用担保额度于 该日到期。 资金需求还要叠加在建工程：期末在建工程 4.70 亿元、2026 年上半年购建固定资产等支出 1.28 亿元（未年化），这部分支出在周期底部通常会被主动放缓。 综合看，公司的外部资金需求以补充营运资金为主、项目投入为辅，滚动偿付口径下的静态缺口约 21.73 亿元，能否落地取决于授信续作与发行窗口的配合。




\subsubsection{农发种业（600313）}
\label{co:600313}

\pfl{公司基本情况} 公司前身可追溯至 1999 年，中国农业发展集团有限公司（中国农发集团）旗下的种业与农资综合企业，主营小麦、玉米、水稻等种子及农药化肥贸易；2026 年 1 月完成向控股股东定增 4.02 亿元，2025 年度实现自 2017 年以来首次现金分红，但母公司累计未分配利润长期为负。 公司于 2001-01-19 在上交所首发上市，注册资本 11.61 亿元，总股本 10.82 亿股。 截至 2026 年 10 月 12 日收盘价 6.71 元，对应总市值 72.6 亿元，近一年+5.0\%，总市值在三级行业“种子”的 8 家公司中按由大到小排第\zhnumber{3}位，近一年涨跌幅高于全样本中位数（-10.4\%）15.4 个百分点。 股权结构方面，实际控制人为中国农业发展集团有限公司。

\pfl{股权结构与控制权} 按 2026 年 6 月 30 日的股东名册，前三大股东为中国农垦集团有限公司 持股 21.55\%；中国农业发展集团有限公司 持股 16.39\%；中国华农资产经营有限公司 持股 5.95\%，股东户数 108,875 户。 控制权与股东性质方面，本次发行前中国农发集团及其一致行动人合计持有公司股份 43,053.23 万股。

\pfl{经营要点} 从经营要点看，2026 年上半年，合同负债 8.73 亿元。

\pfl{所处行业的基本面} 种子的供给由制种面积与库存决定、需求由粮价与种植意愿决定，制种面积在 2025 年起出现萎缩，是本轮供给端最先出清的环节（第 \ref{sec:framework} 章的供需框架）。

\pfl{竞争优势} 把该公司放回同类公司比较，2026 年上半年营业收入 38.7 亿元、在三级行业“种子”的 8 家公司中排第\zhnumber{1}位；毛利率 5.5\%、低于行业内中位数 23.0 个百分点；半年 ROE 2.2\%（未年化）、高于行业内中位数（0.17\%）；资产负债率 37.2\%、低于行业内中位数（37.3\%）。 优势在当期体现为高于行业中位数的资本回报，但优势能否持续仍取决于周期与自身产能投放节奏。

\pfl{历史股价} 以 2021 年 1 月为起点、2026 年 10 月为终点，股价由 5.09 元变为 6.71 元，累计 +31.8\%，区间最高 15.84 元（2022 年 5 月）、最低 4.41 元（2021 年 6 月）；当前价相当于区间最高价的 42.4\%，为区间最低价的 1.52 倍。 月度收益的年化波动率 49.4\%，高于全样本中位数 37.3\%，在三级行业“种子”的 8 家公司中波动率由高到低排第\zhnumber{2}位。

\begin{center}
\begin{minipage}{\textwidth}\centering
\begin{tikzpicture}
\begin{axis}[
  width=12.6cm, height=3.9cm, scale only axis,
  xmin=2020.922, xmax=2026.972, ymin=4.10, ymax=16.95,
  xtick={2021,2022,2023,2024,2025,2026},
  yearticks,
  yticklabel style={font=\scriptsize},
  ylabel={元/股}, ylabel style={font=\scriptsize},
  grid=major, grid style={LightGray, line width=0.3pt},
  axis line style={InkGray, line width=0.5pt},
  every axis plot/.append style={line width=0.9pt},
]
\addplot[AgriGreen] table[x=x,y=y]{data/clean/profiles/px_600313.dat};
\addplot[SoftRed, dashed, line width=0.7pt] table[x=x,y=y]{data/clean/profiles/pxzz_600313.dat};
\addplot[only marks, mark=*, mark size=1.1pt, SoftRed] table[x=x,y=y]{data/clean/profiles/pxzz_600313.dat};
\end{axis}
\end{tikzpicture}

\captionof{figure}[农发种业（600313）股价与周期骨架]{农发种业（600313）月度收盘价（元/股）与 ZigZag 周期骨架（阈值 25\%，圆点为周期转折点）}
\end{minipage}
\end{center}

\pfl{过去周期分析} 以 25\% 的 ZigZag 阈值划分，样本区间内股价形成 3 个主升段与 3 个主跌段。 最大主跌段为 2022 年 5 月 至 2024 年 8 月，27 个月-66.0\%；最大主升段出现在 2021 年 6 月 至 2022 年 5 月，11 个月+259.2\%。 最近一次转折出现在 2026 年 8 月（上涨段自 2026 年 6 月起，2 个月+37.0\%），当前处于该段之后的震荡之中。 公司股价较申万农林牧渔指数提前约 59 个月见底，是板块中较早定价周期反转的公司之一。 种子行业的价格由制种成本、粮价与政策（转基因商业化、种业振兴）共同决定（第 \ref{sec:planting} 节），与生猪价格的联动有限，公司 2026 年上半年实现盈利、半年 ROE 2.2\%（未年化），好于全样本中位数（0.75\%），下行周期对其报表的冲击小于同业。 综合区间形态看，该股单段涨跌幅度偏大、以趋势段为主（最大主升 +259.2\%、最大主跌 -66.0\%），年化波动率 49.4\% 与全样本中位数 37.3\% 相比波动明显大于样本中位数。

\pfl{盈利情况} 2026 年上半年营业收入 38.73 亿元（同比 +32.8\%），归母净利润 0.53 亿元，同比 -16.6\%。 以年报为趋势对照，2023---2025 年营业收入由 67.61 亿元变为 66.54 亿元（两年年均复合增速 -0.8\%），归母净利润由 1.27 亿元变为 0.84 亿元。 按半年报口径，2026 年上半年的 ROE 为 2.2\%（未年化）、销售净利率 1.7\%、毛利率 5.5\%，ROE 在“种子”8 家公司中按数值由高到低排第\zhnumber{3}位。 同期全样本半年 ROE 中位数 0.75\%，公司与之相差 1.48 个百分点（略好于中位数公司）。 从公司披露的原因看，业绩变动主要来自：2026 年上半年，营业总收入 38.73 亿元（同比 +32.84\%）；2024 年，归母净利润同比减少：受国际市场影响农药下游企业出口订单锐减、行业去库存导致国内市场需求走低、农药主要产品价格下跌。 综合看，公司在报告期维持盈利（高于全样本半年 ROE 中位数 0.75\%），利润的可持续性取决于种子景气的延续与成本管控的执行。

\begin{center}\small\setlength{\tabcolsep}{3.4pt}
\captionof{table}[农发种业（600313）关键财务指标]{农发种业（600313）关键财务指标（合并报表口径；两个半年报期均未年化；现金短债比 $=$ 货币资金 $\div$ 短期债务）}
\begin{adjustbox}{max width=\textwidth}
\begin{tabular}{@{}lrrrrr@{}}
\toprule
指标 & 2023 年 & 2024 年 & 2025 年 & 2025 年半年报 & 2026 年半年报 \\
\midrule
营业收入（亿元） & 67.61 & 55.63 & 66.54 & 29.16 & 38.73 \\
归母净利润（亿元） & 1.27 & 0.47 & 0.84 & 0.63 & 0.53 \\
毛利率（\%） & 8.4 & 7.6 & 7.5 & 7.0 & 5.5 \\
ROE（\%） & 7.1 & 2.5 & 4.3 & 3.2 & 2.2 \\
资产负债率（\%） & 40.5 & 44.4 & 43.3 & 39.5 & 37.2 \\
经营活动现金流净额（亿元） & 6.06 & -0.77 & 3.02 & -1.02 & 0.79 \\
货币资金（亿元） & 12.06 & 9.99 & 15.09 & 9.65 & 8.73 \\
有息负债合计（亿元） & 3.75 & 3.87 & 5.05 & 5.97 & 4.77 \\
现金短债比（倍） & 4.57 & 3.06 & 3.83 & 1.98 & 2.44 \\
\bottomrule
\end{tabular}
\end{adjustbox}
\end{center}

\begin{center}
\begin{minipage}{\textwidth}\centering
\begin{tikzpicture}
\begin{groupplot}[
  group style={group size=2 by 1, horizontal sep=0.60cm},
  width=6.85cm, height=4.60cm,
  tick label style={font=\tiny},
  title style={font=\scriptsize\heiti},
  yticklabel style={font=\tiny},
  ylabel style={font=\tiny},
  grid=major, grid style={LightGray, line width=0.3pt},
  axis line style={InkGray, line width=0.5pt},
  enlarge x limits={abs=0.8},
  legend style={font=\scriptsize, at={(0.5,-0.20)}, anchor=north,
                legend columns=2, draw=none, fill=none, column sep=6pt},
]
\nextgroupplot[
  ybar, bar width=4pt,
  ymin=-6.76, ymax=75.72,
  xtick={0,1,2,3,4,5.7},
  xticklabels={2021,2022,2023,2024,2025,2026H1},
  ylabel={亿元},
]
\addplot[fill=AgriGreen!75, draw=AgriGreen, bar shift=-2.6pt] table[x=i, y=rev]{data/clean/profiles/fin_600313.dat};
\addplot[fill=SoftRed!60, draw=SoftRed, bar shift=2.6pt] table[x=i, y=ni]{data/clean/profiles/fin_600313.dat};
\addplot[fill=AgriGreen!30, draw=AgriGreen, pattern=north east lines, bar shift=-2.6pt] table[x=x, y=rev]{data/clean/profiles/finh_600313.dat};
\addplot[fill=SoftRed!20, draw=SoftRed, pattern=north east lines, bar shift=2.6pt] table[x=x, y=ni]{data/clean/profiles/finh_600313.dat};
\legend{营业收入（年/半年）, 归母净利润（年/半年）}
\nextgroupplot[
  ymin=-3.6, ymax=48.9,
  xtick={0,1,2,3,4,5.7},
  xticklabels={2021,2022,2023,2024,2025,2026H1},
  ylabel={\%},
]
\addplot[AgriGold, mark=*, mark size=1.3pt] table[x=i, y=roe]{data/clean/profiles/fin_600313.dat};
\addplot[OilBlue, mark=square*, mark size=1.3pt] table[x=i, y=debt]{data/clean/profiles/fin_600313.dat};
\addplot[only marks, AgriGold, mark=*, mark size=2.0pt] table[x=x, y=roe]{data/clean/profiles/finh_600313.dat};
\addplot[only marks, OilBlue, mark=square*, mark size=2.0pt] table[x=x, y=debt]{data/clean/profiles/finh_600313.dat};
\legend{ROE, 资产负债率}
\end{groupplot}
\end{tikzpicture}
\begin{tikzpicture}
\begin{axis}[
  name=finbot,
  width=13.75cm, height=3.10cm, scale only axis,
  ymin=-1.51, ymax=16.91,
  xtick={0,1,2,3,4,5.7},
  xticklabels={2021,2022,2023,2024,2025,2026H1},
  tick label style={font=\tiny},
  yticklabel style={font=\tiny},
  ylabel={亿元}, ylabel style={font=\tiny},
  ybar, bar width=4pt,
  grid=major, grid style={LightGray, line width=0.3pt},
  axis line style={InkGray, line width=0.5pt},
  enlarge x limits={abs=0.8},
  legend style={font=\scriptsize, at={(0.5,-0.26)}, anchor=north,
                legend columns=2, draw=none, fill=none, column sep=10pt},
]
\addplot[fill=AgriGreen!55, draw=AgriGreen, bar shift=-3.0pt] table[x=i, y=cash]{data/clean/profiles/fin_600313.dat};
\addplot[fill=SoftRed!45, draw=SoftRed, bar shift=3.0pt] table[x=i, y=ibd]{data/clean/profiles/fin_600313.dat};
\addplot[fill=AgriGreen!25, draw=AgriGreen, pattern=north east lines, bar shift=-3.0pt] table[x=x, y=cash]{data/clean/profiles/finh_600313.dat};
\addplot[fill=SoftRed!20, draw=SoftRed, pattern=north east lines, bar shift=3.0pt] table[x=x, y=ibd]{data/clean/profiles/finh_600313.dat};
\legend{货币资金（左轴）, 有息负债（左轴）}
\end{axis}
\begin{axis}[
  at={(finbot.south east)}, anchor=south east,
  width=13.75cm, height=3.10cm, scale only axis,
  axis y line*=right, axis x line*=none, xtick=\empty,
  ymin=-2.54, ymax=8.10,
  tick label style={font=\tiny},
  yticklabel style={font=\tiny}, scaled y ticks=false,
  legend style={font=\scriptsize, at={(0.5,-0.44)}, anchor=north,
                legend columns=1, draw=none, fill=none},
]
\addplot[OilBlue, mark=*, mark size=2.2pt, line width=1.3pt] table[x=i, y=ocf]{data/clean/profiles/fin_600313.dat};
\addplot[only marks, OilBlue, mark=*, mark size=3.2pt] table[x=x, y=ocf]{data/clean/profiles/finh_600313.dat};
\legend{经营活动现金流净额（右轴，亿元，蓝点）}
\end{axis}
\end{tikzpicture}

\captionof{figure}[农发种业（600313）近年主要财务指标]{农发种业（600313）近年主要财务指标（左上：营业收入与归母净利润；右上：ROE 与资产负债率；下：货币资金、有息负债为左轴柱状，经营活动现金流净额为其右轴的蓝点折线。
2021---2025 年为年报口径，2026H1 为 2026 年半年报/半年末，与其前一年报之间留出间隔、以斜纹或空心点区分；半年报数据未年化）}
\end{minipage}
\end{center}

\pfl{财务分析} 2026 年 6 月末资产负债率 37.2\%，较 2025 年末下降 6.1 个百分点（样本中位数 52.2\%，所处三级行业内按由低到高排第\zhnumber{4}位）。 2026 年 6 月末货币资金 8.73 亿元、短期债务（短期借款与一年内到期非流动负债）3.58 亿元、长期债务 1.19 亿元，有息负债合计 4.77 亿元，现金短债比 2.44 倍——覆盖充分。 净有息负债 -3.96 亿元，相当于其 2026 年上半年营业收入的 -10\%。 流动比率 2.24、速动比率 1.69，短期流动性无明显缺口。 2026 年上半年经营活动现金流净额 0.79 亿元，经营现金流与利润同向为正，内生资金可以支撑运营。 2026 年 6 月末在建工程 1.36 亿元，较上年末减少 0.04 亿元，在建投入在收缩。 2026 年 6 月末商誉 1.36 亿元，占归母净资产的 3.9\%，是净资产中最脆弱的一块。 综合看，现金短债比 2.44 倍，短期偿付压力可控；资产负债率低于样本中位数 15.0 个百分点；有息负债中短期债务占比更高，期限结构仍以滚动续作为主。

\pfl{重大战略转型} 近三年公司的战略动作集中在：其他动作（2 项）：2026 年 8 月，公告向全资子公司增资，以现金出资、资金来源为自有资金向全资子公司增资；2022 年 8 月，第七届董事会第十九次会议审议通过投资设立控股子公司。 公告口径上，近三年（2023 年 9 月---2026 年 9 月）公司披露重大重组与并购类公告 3 条、对外投资与转型类公告 5 条、回购与股权激励类公告 0 条，合计公告 291 条。 综合判断，报告期内公司有 8 条与战略相关的公告落地（重大重组与并购 3 条、对外投资与转型 5 条），资本运作与主业调整之间存在直接关联。

\begin{center}\small\setlength{\tabcolsep}{4pt}
\captionof{table}[农发种业（600313）历史融资与资本运作]{农发种业（600313）历史融资与资本运作（据公司公告与公开报道整理；金额为披露值，含首次公开发行、再融资、债券、银行借款与重整投资等）}
\begin{adjustbox}{max width=\textwidth}
\begin{tabular}{@{}p{1.45cm}p{2.55cm}p{4.05cm}p{4.95cm}@{}}
\toprule
时间 & 方式 & 规模 & 用途与说明 \\
\midrule
2000-12 & IPO（上海证券交易所主板） & 向社会公开发行人民币普通股 8,\allowbreak{}000 万股，\allowbreak{}每股发行价格 6.30 元 & 经中国证监会证监发行字〔2000〕178 号文批准发行 \\
2025 & 公积金弥补亏损 & 使用盈余公积弥补亏损 1,\allowbreak{}563.18 万元 & 为提升投资者回报能力 \\
2025-11 至 2026-01 & 定向增发（向特定对象发行 A 股股票并在主板上市） & 发行 7,\allowbreak{}917.53 万股，\allowbreak{}发行价格 5.14 元/\allowbreak{}股，\allowbreak{}募集资金总额 4.07 亿元 & 将中国农发集团以委托贷款方式拨付公司的国拨资金转为对公司的直接投资（股权投资），\allowbreak{}同时补充流动资金；发行对象为公司实际控制人中国农发集团（拟认购 3.07 亿元）及其一致行动人华农资产（拟认购 1.00 亿元）；定价基准日为第七届董事会第五十次会议决议公告日 \\
2026-06 & 现金分红（2025 年度） & 以股权登记日总股本 11.61 亿股为基数 & 回报股东 \\
\bottomrule
\end{tabular}
\end{adjustbox}
\end{center}

\pfl{历史融投资情况} 从现金流量表看，2025 年公司取得借款收到的现金 6.49 亿元、偿还债务支付的现金 6.15 亿元，筹资活动现金流净额 -0.01 亿元（2023 年为取得借款 3.15 亿元、偿还债务 2.49 亿元）。 2026 年上半年筹资活动现金流净额为 0.54 亿元（取得借款 2.69 亿元、偿还债务 2.74 亿元，未年化），仍处于净融入状态。 2025 年分配股利、利润或偿付利息支付的现金合计 0.29 亿元。 公开披露的融资记录共 4 笔，工具类型以 IPO、定向增发为主；金额与用途见下表。 股本方面，近三年披露股本变动 4 次（限售股份上市 1 次）；总股本自 2021 年末以来未发生变化（10.82 亿股）。 分红方面，近三年现金分红 1 次、合计每股派现 0.01 元，最近一次实施在 2026 年 6 月。 近三年“再融资”类公告 32 条，涉及定向增发、发行股份购买资产、募集资金使用。 综合看，公司融资结构上股权与债务融资并举；2025 年筹资活动现金流净额 -0.01 亿元，融资节奏仍以维持存量债务周转为主，与前述经营与投资现金流的方向基本一致。

\pfl{未来潜在融资需求分析} 2026 年 6 月末货币资金 8.73 亿元，而短期债务 3.58 亿元（现金短债比 2.44 倍），2026 年上半年经营现金流净额 0.79 亿元。账面货币资金可以覆盖短期债务，缺口不体现在静态偿付层面。 公司 2026 年 6 月末资产负债率 37.2\%、2026 年上半年 ROE 2.2\%（未年化），资产负债表尚有余量，融资需求不是“补血型”的刚性需求，而是配套在建项目投入的主动性安排。 公开披露的后续资金安排线索包括：股权融资线索：2026 年 1 月，公司完成向中国农发集团及华农资产定增。 资金需求还要叠加在建工程：期末在建工程 1.36 亿元、2026 年上半年购建固定资产等支出 0.71 亿元（未年化），这部分支出在周期底部通常会被主动放缓。 静态偿付不存在缺口，后续融资更可能服务于产能或并购安排，而不是填补流动性窟窿。




\subsubsection{神农种业（300189）}
\label{co:300189}

\pfl{公司基本情况} 神农种业成立于 2000 年，海南起家的水稻、玉米、油菜等农作物种子企业，2011 年创业板上市（原名神农大丰，历经神农基因、神农科技更名）；因信披违法被证监会处罚、主业长期亏损，2025 年靠处置联营企业股权实现账面盈利，2026 年上半年再度亏损。 公司于 2011-03-16 在深交所创业板通过重大资产重组（借壳）上市，注册资本 10.24 亿元，总股本 10.24 亿股。 截至 2026 年 10 月 12 日收盘价 6.16 元，对应总市值 63.1 亿元，近一年+38.1\%，总市值在三级行业“种子”的 8 家公司中按由大到小排第\zhnumber{4}位，近一年涨跌幅高于全样本中位数（-10.4\%）48.5 个百分点。 按 2026 年半年报披露的行业构成，收入占比前两位为农业 95.6\%；其他(补充) 4.4\%。

\pfl{股权结构与控制权} 按 2026 年 6 月 30 日的股东名册，前三大股东为曹欧劼 持股 17.94\%；香港中央结算有限公司 持股 1.05\%；UBSAG 持股 1.02\%，股东户数 80,936 户。 股东结构与控制权方面，持股 17.94\%（1.84 亿股），同时任公司董事长与法定代表人，并为公司实际控制人。

\pfl{经营要点} 经营与产能方面，种子生产主要采取委托制种方式，由公司委托的制种单位组织当地农户集中生产并负责农户管理与集中收购。

\pfl{所处行业的基本面} 种子的供给由制种面积与库存决定、需求由粮价与种植意愿决定，制种面积在 2025 年起出现萎缩，是本轮供给端最先出清的环节（第 \ref{sec:framework} 章的供需框架）。 从供需两端的当期数据看，2026 年，国家继续在稻谷主产区实行最低收购价政策；2026 年，早籼稻、中晚籼稻、粳稻最低收购价维持上年水平，分别为每斤 1 元、1 元、1 元。

\pfl{竞争优势} 把该公司放回同类公司比较，2026 年上半年营业收入 0.6 亿元、在三级行业“种子”的 8 家公司中排第\zhnumber{8}位；毛利率 38.4\%、高于行业内中位数 10.0 个百分点；半年 ROE -1.9\%（未年化）、低于行业内中位数（0.17\%）；资产负债率 39.8\%、高于行业内中位数（37.3\%）。 与同类公司相比，公司的差异化在于：2026 年 9 月，公司及实际控制人曹欧劼为全资孙公司神农种源向中国银行海口龙珠支行申请的 5.48 亿元借款提供连带责任保证。 优势在报表上主要体现为高于行业平均的毛利率，这在周期底部是抵御亏损的主要缓冲。

\pfl{历史股价} 以 2021 年 1 月为起点、2026 年 10 月为终点，股价由 4.15 元变为 6.16 元，累计 +48.4\%，区间最高 8.66 元（2026 年 1 月）、最低 1.79 元（2024 年 6 月）；当前价相当于区间最高价的 71.1\%，为区间最低价的 3.44 倍。 月度收益的年化波动率 44.4\%，高于全样本中位数 37.3\%，在三级行业“种子”的 8 家公司中波动率由高到低排第\zhnumber{3}位。 把股价阶段与公司事件对齐看，2024 年 6 月 至 2026 年 1 月股价上涨+383.8\%，同期（2025 年 1 月）公司转让控股子公司福建神农大丰种业科技有限公司 88.67\% 股权给李坤泰。 从时间对应关系看，公司层面的资本运作与股权、风险事件解释了区间内多数急涨急跌，与行业指数的同步性相对有限。

\begin{center}
\begin{minipage}{\textwidth}\centering
\begin{tikzpicture}
\begin{axis}[
  width=12.6cm, height=3.9cm, scale only axis,
  xmin=2020.922, xmax=2026.972, ymin=1.66, ymax=9.27,
  xtick={2021,2022,2023,2024,2025,2026},
  yearticks,
  yticklabel style={font=\scriptsize},
  ylabel={元/股}, ylabel style={font=\scriptsize},
  grid=major, grid style={LightGray, line width=0.3pt},
  axis line style={InkGray, line width=0.5pt},
  every axis plot/.append style={line width=0.9pt},
]
\addplot[AgriGreen] table[x=x,y=y]{data/clean/profiles/px_300189.dat};
\addplot[SoftRed, dashed, line width=0.7pt] table[x=x,y=y]{data/clean/profiles/pxzz_300189.dat};
\addplot[only marks, mark=*, mark size=1.1pt, SoftRed] table[x=x,y=y]{data/clean/profiles/pxzz_300189.dat};
\end{axis}
\end{tikzpicture}

\captionof{figure}[神农种业（300189）股价与周期骨架]{神农种业（300189）月度收盘价（元/股）与 ZigZag 周期骨架（阈值 25\%，圆点为周期转折点）}
\end{minipage}
\end{center}

\pfl{过去周期分析} 以 25\% 的 ZigZag 阈值划分，样本区间内股价形成 3 个主升段与 2 个主跌段。 最大主跌段为 2022 年 5 月 至 2024 年 6 月，25 个月-68.2\%；最大主升段出现在 2024 年 6 月 至 2026 年 1 月，19 个月+383.8\%。 最近一次转折出现在 2026 年 8 月（上涨段自 2026 年 6 月起，2 个月+45.6\%），当前处于该段之后的震荡之中。 公司股价较申万农林牧渔指数提前约 23 个月见底，是板块中较早定价周期反转的公司之一。 种子行业的价格由制种成本、粮价与政策（转基因商业化、种业振兴）共同决定（第 \ref{sec:planting} 节），与生猪价格的联动有限，公司 2026 年 6 月末资产负债率 39.8\%、半年 ROE -1.9\%（未年化），在本轮下行中属于随行业同步调整的一类。 综合区间形态看，该股单段涨跌幅度偏大、以趋势段为主（最大主升 +383.8\%、最大主跌 -68.2\%），年化波动率 44.4\% 与全样本中位数 37.3\% 相比波动明显大于样本中位数。

\pfl{盈利情况} 2026 年上半年营业收入 0.56 亿元（同比 -35.3\%），归母净利润 -0.13 亿元，亏损同比收窄 0.01 亿元。 以年报为趋势对照，2023---2025 年营业收入由 1.67 亿元变为 2.49 亿元（两年年均复合增速 +22.2\%），归母净利润由 -0.36 亿元变为 1.07 亿元。 按半年报口径，2026 年上半年的 ROE 为 -1.9\%（未年化）、销售净利率 -25.9\%、毛利率 38.4\%，ROE 在“种子”8 家公司中按数值由高到低排第\zhnumber{6}位。 同期全样本半年 ROE 中位数 0.75\%，公司与之相差 2.70 个百分点（弱于中位数公司）。 从公司披露的原因看，业绩变动主要来自：2026 年上半年，毛利率-60.37\% 主要系水产种源受赤潮影响造成损失。 综合看，公司仍处亏损区间、亏损同比收窄，半年 ROE 在三级行业内按由高到低排第\zhnumber{6}位，单位盈利能力的修复依赖行业景气与自身成本端的改善，报表弹性主要体现为价格与销量的方向性变化。

\begin{center}\small\setlength{\tabcolsep}{3.4pt}
\captionof{table}[神农种业（300189）关键财务指标]{神农种业（300189）关键财务指标（合并报表口径；两个半年报期均未年化；现金短债比 $=$ 货币资金 $\div$ 短期债务）}
\begin{adjustbox}{max width=\textwidth}
\begin{tabular}{@{}lrrrrr@{}}
\toprule
指标 & 2023 年 & 2024 年 & 2025 年 & 2025 年半年报 & 2026 年半年报 \\
\midrule
营业收入（亿元） & 1.67 & 1.57 & 2.49 & 0.87 & 0.56 \\
归母净利润（亿元） & -0.36 & -0.51 & 1.07 & -0.14 & -0.13 \\
毛利率（\%） & 35.6 & 37.5 & 35.4 & 30.2 & 38.4 \\
ROE（\%） & -4.6 & -6.8 & 14.0 & -2.0 & -1.9 \\
资产负债率（\%） & 14.5 & 31.5 & 38.9 & 33.4 & 39.8 \\
经营活动现金流净额（亿元） & 0.14 & 0.07 & 0.23 & -0.33 & -0.32 \\
货币资金（亿元） & 1.11 & 0.57 & 0.97 & 0.54 & 0.99 \\
有息负债合计（亿元） & 0.08 & 2.23 & 2.83 & 2.42 & 3.05 \\
现金短债比（倍） & 2,779.09 & 3.57 & 1.63 & 3.84 & 1.11 \\
\bottomrule
\end{tabular}
\end{adjustbox}
\end{center}

\begin{center}
\begin{minipage}{\textwidth}\centering
\begin{tikzpicture}
\begin{groupplot}[
  group style={group size=2 by 1, horizontal sep=0.60cm},
  width=6.85cm, height=4.60cm,
  tick label style={font=\tiny},
  title style={font=\scriptsize\heiti},
  yticklabel style={font=\tiny},
  ylabel style={font=\tiny},
  grid=major, grid style={LightGray, line width=0.3pt},
  axis line style={InkGray, line width=0.5pt},
  enlarge x limits={abs=0.8},
  legend style={font=\scriptsize, at={(0.5,-0.20)}, anchor=north,
                legend columns=2, draw=none, fill=none, column sep=6pt},
]
\nextgroupplot[
  ybar, bar width=4pt,
  ymin=-0.96, ymax=2.87,
  xtick={0,1,2,3,4,5.7},
  xticklabels={2021,2022,2023,2024,2025,2026H1},
  ylabel={亿元},
]
\addplot[fill=AgriGreen!75, draw=AgriGreen, bar shift=-2.6pt] table[x=i, y=rev]{data/clean/profiles/fin_300189.dat};
\addplot[fill=SoftRed!60, draw=SoftRed, bar shift=2.6pt] table[x=i, y=ni]{data/clean/profiles/fin_300189.dat};
\addplot[fill=AgriGreen!30, draw=AgriGreen, pattern=north east lines, bar shift=-2.6pt] table[x=x, y=rev]{data/clean/profiles/finh_300189.dat};
\addplot[fill=SoftRed!20, draw=SoftRed, pattern=north east lines, bar shift=2.6pt] table[x=x, y=ni]{data/clean/profiles/finh_300189.dat};
\legend{营业收入（年/半年）, 归母净利润（年/半年）}
\nextgroupplot[
  ymin=-11.5, ymax=44.6,
  xtick={0,1,2,3,4,5.7},
  xticklabels={2021,2022,2023,2024,2025,2026H1},
  ylabel={\%},
]
\addplot[AgriGold, mark=*, mark size=1.3pt] table[x=i, y=roe]{data/clean/profiles/fin_300189.dat};
\addplot[OilBlue, mark=square*, mark size=1.3pt] table[x=i, y=debt]{data/clean/profiles/fin_300189.dat};
\addplot[only marks, AgriGold, mark=*, mark size=2.0pt] table[x=x, y=roe]{data/clean/profiles/finh_300189.dat};
\addplot[only marks, OilBlue, mark=square*, mark size=2.0pt] table[x=x, y=debt]{data/clean/profiles/finh_300189.dat};
\legend{ROE, 资产负债率}
\end{groupplot}
\end{tikzpicture}
\begin{tikzpicture}
\begin{axis}[
  name=finbot,
  width=13.75cm, height=3.10cm, scale only axis,
  ymin=-0.30, ymax=3.42,
  xtick={0,1,2,3,4,5.7},
  xticklabels={2021,2022,2023,2024,2025,2026H1},
  tick label style={font=\tiny},
  yticklabel style={font=\tiny},
  ylabel={亿元}, ylabel style={font=\tiny},
  ybar, bar width=4pt,
  grid=major, grid style={LightGray, line width=0.3pt},
  axis line style={InkGray, line width=0.5pt},
  enlarge x limits={abs=0.8},
  legend style={font=\scriptsize, at={(0.5,-0.26)}, anchor=north,
                legend columns=2, draw=none, fill=none, column sep=10pt},
]
\addplot[fill=AgriGreen!55, draw=AgriGreen, bar shift=-3.0pt] table[x=i, y=cash]{data/clean/profiles/fin_300189.dat};
\addplot[fill=SoftRed!45, draw=SoftRed, bar shift=3.0pt] table[x=i, y=ibd]{data/clean/profiles/fin_300189.dat};
\addplot[fill=AgriGreen!25, draw=AgriGreen, pattern=north east lines, bar shift=-3.0pt] table[x=x, y=cash]{data/clean/profiles/finh_300189.dat};
\addplot[fill=SoftRed!20, draw=SoftRed, pattern=north east lines, bar shift=3.0pt] table[x=x, y=ibd]{data/clean/profiles/finh_300189.dat};
\legend{货币资金（左轴）, 有息负债（左轴）}
\end{axis}
\begin{axis}[
  at={(finbot.south east)}, anchor=south east,
  width=13.75cm, height=3.10cm, scale only axis,
  axis y line*=right, axis x line*=none, xtick=\empty,
  ymin=-0.51, ymax=0.40,
  tick label style={font=\tiny},
  yticklabel style={font=\tiny}, scaled y ticks=false,
  legend style={font=\scriptsize, at={(0.5,-0.44)}, anchor=north,
                legend columns=1, draw=none, fill=none},
]
\addplot[OilBlue, mark=*, mark size=2.2pt, line width=1.3pt] table[x=i, y=ocf]{data/clean/profiles/fin_300189.dat};
\addplot[only marks, OilBlue, mark=*, mark size=3.2pt] table[x=x, y=ocf]{data/clean/profiles/finh_300189.dat};
\legend{经营活动现金流净额（右轴，亿元，蓝点）}
\end{axis}
\end{tikzpicture}

\captionof{figure}[神农种业（300189）近年主要财务指标]{神农种业（300189）近年主要财务指标（左上：营业收入与归母净利润；右上：ROE 与资产负债率；下：货币资金、有息负债为左轴柱状，经营活动现金流净额为其右轴的蓝点折线。
2021---2025 年为年报口径，2026H1 为 2026 年半年报/半年末，与其前一年报之间留出间隔、以斜纹或空心点区分；半年报数据未年化）}
\end{minipage}
\end{center}

\pfl{财务分析} 2026 年 6 月末资产负债率 39.8\%，较 2025 年末上升 0.9 个百分点（样本中位数 52.2\%，所处三级行业内按由低到高排第\zhnumber{6}位）。 2026 年 6 月末货币资金 0.99 亿元、短期债务（短期借款与一年内到期非流动负债）0.88 亿元、长期债务 2.17 亿元，有息负债合计 3.05 亿元，现金短债比 1.11 倍——覆盖充分。 净有息负债 2.06 亿元，相当于其 2026 年上半年营业收入的 366\%。 流动比率 1.79、速动比率 1.44，短期流动性无明显缺口。 2026 年上半年经营活动现金流净额 -0.32 亿元，经营现金流与利润同时为负，现金消耗较快。 2026 年 6 月末在建工程 2.71 亿元，较上年末增加 0.09 亿元，仍有在建投入。 综合看，现金短债比 1.11 倍，短期偿付压力可控；资产负债率低于样本中位数 12.4 个百分点；有息负债中长期债务占比更高，期限结构对短债滚动的压力较小。

\pfl{重大战略转型} 从公开披露看，公司的战略推进集中在：其他动作（2 项）：2025 年，公司投资收益 1.41 亿元；并购与重组（1 项）：2025 年 1 月，公司转让控股子公司福建神农大丰种业科技有限公司 88.67\% 股权给李坤泰。 公告口径上，近三年（2023 年 9 月---2026 年 9 月）公司披露重大重组与并购类公告 1 条、对外投资与转型类公告 2 条、回购与股权激励类公告 0 条，合计公告 265 条。 第一大行业仍是“农业”，收入占比 95.6\%，与 2021 年年报相比没有方向性变化。 综合判断，报告期内公司有 3 条与战略相关的公告落地（重大重组与并购 1 条、对外投资与转型 2 条），资本运作与主业调整之间存在直接关联。

\begin{center}\small\setlength{\tabcolsep}{4pt}
\captionof{table}[神农种业（300189）历史融资与资本运作]{神农种业（300189）历史融资与资本运作（据公司公告与公开报道整理；金额为披露值，含首次公开发行、再融资、债券、银行借款与重整投资等）}
\begin{adjustbox}{max width=\textwidth}
\begin{tabular}{@{}p{1.45cm}p{2.55cm}p{4.05cm}p{4.95cm}@{}}
\toprule
时间 & 方式 & 规模 & 用途与说明 \\
\midrule
2011-03 & IPO（深交所创业板） & 发行 4,\allowbreak{}000 万股，\allowbreak{}发行价 24.00 元/\allowbreak{}股，\allowbreak{}募集资金总额 9.60 亿元 & 与主营业务相关的项目及主营业务发展所需营运资金；发行后股份总数由 12,\allowbreak{}000 万股增至 16,\allowbreak{}000 万股 \\
2015 & 发行股份购买资产（拟） & 拟购买波莲基因 61.52\% 股权 & 购买标的公司股权；交易文件认定不构成重大资产重组、\allowbreak{}不构成借壳上市 \\
2019 & 行政处罚（罚款） & 公司被罚 60.00 万元，\allowbreak{}黄培劲被罚 70.00 万元 & 海南监管局〔2019〕4 号行政处罚决定书，\allowbreak{}非融资事项，\allowbreak{}列此以说明资本运作约束 \\
2024 & 债务融资 & 2024 年度筹资活动产生的现金流量净额为 2.12 亿元，\allowbreak{}同比增加 1908.04\% & 补充运营资金；2025 年财务费用 316.53 万元 \\
\bottomrule
\end{tabular}
\end{adjustbox}
\end{center}

\pfl{历史融投资情况} 从现金流量表看，2025 年公司取得借款收到的现金 0.83 亿元、偿还债务支付的现金 0.21 亿元，筹资活动现金流净额 0.65 亿元（2023 年为取得借款 0.05 亿元、偿还债务 0.14 亿元）。 2026 年上半年筹资活动现金流净额为 0.03 亿元（取得借款 0.28 亿元、偿还债务 0.17 亿元，未年化），仍处于净融入状态。 2025 年分配股利、利润或偿付利息支付的现金合计 0.08 亿元。 公开披露的融资记录共 4 笔，工具类型以 IPO、发行股份购买资产、债务融资为主；金额与用途见下表。 其中规模最大的两笔为：2024 年，债务融资，2024 年度筹资活动产生的现金流量净额为 2.12 亿元，同比增加 1908.04\%；2019 年，行政处罚，公司被罚 60.00 万元，黄培劲被罚 70.00 万元。 股本方面，近三年披露股本变动 3 次；总股本自 2021 年末以来未发生变化（10.24 亿股）。 分红方面，近三年未实施现金分红，在全部样本中属于分红能力偏弱的一端。 近三年“再融资”类公告 14 条，涉及定向增发、募集资金使用。 综合看，公司融资结构上股权与债务融资并举；2025 年筹资活动现金流净额 0.65 亿元，年度内仍为净融入，与前述经营与投资现金流的方向基本一致。

\pfl{未来潜在融资需求分析} 2026 年 6 月末货币资金 0.99 亿元，而短期债务 0.88 亿元（现金短债比 1.11 倍），2026 年上半年经营现金流净额 -0.32 亿元。账面货币资金可以覆盖短期债务，缺口不体现在静态偿付层面。 公司 2026 年上半年归母净利润为负（-0.13 亿元），但 6 月末资产负债率 39.8\% 尚未进入第一档（亏损且负债率 $>$65\%）的区间，当前融资需求以补充流动资金、支撑低谷期经营性支出为主。 公开披露的后续资金安排线索包括：债务与授信安排：2025 年，主要因本期借款增加、利息支出上升，对外部债务资金依赖上升。 资金需求还要叠加在建工程：期末在建工程 2.71 亿元、2026 年上半年购建固定资产等支出 0.10 亿元（未年化），这部分支出在周期底部通常会被主动放缓。 静态偿付不存在缺口，后续融资更可能服务于产能或并购安排，而不是填补流动性窟窿。




\subsubsection{宏辉果蔬（603336）}
\label{co:603336}

\pfl{公司基本情况} 宏辉果蔬成立于 1992 年，汕头起家的果蔬采后加工与冷链配送服务商，2016 年上交所主板上市；2025 年控制权由创始人黄俊辉转让给苏州国资背景的申泽瑞泰，2026 年以 7 亿元收购施美药业 41.128\% 股权跨界医药，形成果蔬+医药双主业。 公司于 2016-11-24 在上交所首发上市，注册资本 6.09 亿元，总股本 5.70 亿股。 截至 2026 年 10 月 12 日收盘价 10.45 元，对应总市值 59.6 亿元，近一年+14.0\%，总市值在三级行业“其他种植业”的 6 家公司中按由大到小排第\zhnumber{4}位，近一年涨跌幅高于全样本中位数（-10.4\%）24.3 个百分点。 按 2025 年年报披露的行业构成，收入占比前三位为果蔬服务业 72.1\%；供应链管理服务业 22.6\%；农副食品加工业 4.9\%。

\pfl{股权结构与控制权} 按 2026 年 6 月 30 日的股东名册，前三大股东为苏州申泽瑞泰企业管理合伙企业(有限合伙) 持股 24.87\%；黄俊辉 持股 16.54\%；郑幼文 持股 2.19\%，股东户数 21,241 户。 从控制权角度看，控股股东苏州申泽瑞泰企业管理合伙企业（有限合伙）持股 24.87\%；报告期内公司通过全资子公司上海腾晖医药技术有限公司以自有及自筹资金 7.00 亿元收购施美药业 41.128\% 股份并实现控股。

\pfl{经营要点} 经营层面，2025 年 8 月，原大股东黄俊辉为支持公司经营发展，向公司提供不超过 4.60 亿元的无息借款、期限 1 年。

\pfl{所处行业的基本面} 该子行业按品种各自定价：天然橡胶看全球供给与汽车需求，棉花看新疆产量与配额，农药看原药价格与农业服务需求（第 \ref{sec:framework} 章的供需框架）。 从供需两端的当期数据看，果蔬行业长期存在仓储运输损耗高企、产区布局失衡、产销错配、成本承压、同质竞争等结构性问题。

\pfl{竞争优势} 把该公司放回同类公司比较，2026 年上半年营业收入 4.2 亿元、在三级行业“其他种植业”的 6 家公司中排第\zhnumber{5}位；毛利率 16.2\%、低于行业内中位数 0.7 个百分点；半年 ROE 0.4\%（未年化）、低于行业内中位数（1.70\%）；资产负债率 36.3\%、低于行业内中位数（59.7\%）。 公司在公开资料中自述与同业比较的相对优势在于：已在全国设立多个仓储加工配送中心、配送区域覆盖全国各地，并建立果品质量追溯体系。 就当期报表看，公司尚未把经营层面的特征转化为高于行业的资本回报，周期底部的相对优势更多体现在现金流与资产负债表上。

\pfl{历史股价} 以 2021 年 1 月为起点、2026 年 10 月为终点，股价由 10.11 元变为 10.45 元，累计 +3.4\%，区间最高 11.75 元（2026 年 4 月）、最低 3.01 元（2024 年 6 月）；当前价相当于区间最高价的 88.9\%，为区间最低价的 3.47 倍。 月度收益的年化波动率 40.4\%，高于全样本中位数 37.3\%，在三级行业“其他种植业”的 6 家公司中波动率由高到低排第\zhnumber{2}位。 把股价阶段与公司事件对齐看，2024 年 6 月 至 2026 年 4 月股价上涨+290.4\%，同期（2026 年 4 月）第六届董事会第六次会议审议通过收购江西施美药业股份有限公司部分股份，以自有及自筹资金 7.00 亿元受让自然人江鸿持有的施美药业 41.128\% 股份。 把这些阶段与公司事件对齐后可以看到，几轮大涨大跌多由公司自身的资本运作与风险事件触发，行业景气只解释了其中一部分。

\begin{center}
\begin{minipage}{\textwidth}\centering
\begin{tikzpicture}
\begin{axis}[
  width=12.6cm, height=3.9cm, scale only axis,
  xmin=2020.922, xmax=2026.972, ymin=2.80, ymax=12.57,
  xtick={2021,2022,2023,2024,2025,2026},
  yearticks,
  yticklabel style={font=\scriptsize},
  ylabel={元/股}, ylabel style={font=\scriptsize},
  grid=major, grid style={LightGray, line width=0.3pt},
  axis line style={InkGray, line width=0.5pt},
  every axis plot/.append style={line width=0.9pt},
]
\addplot[AgriGreen] table[x=x,y=y]{data/clean/profiles/px_603336.dat};
\addplot[SoftRed, dashed, line width=0.7pt] table[x=x,y=y]{data/clean/profiles/pxzz_603336.dat};
\addplot[only marks, mark=*, mark size=1.1pt, SoftRed] table[x=x,y=y]{data/clean/profiles/pxzz_603336.dat};
\end{axis}
\end{tikzpicture}

\captionof{figure}[宏辉果蔬（603336）股价与周期骨架]{宏辉果蔬（603336）月度收盘价（元/股）与 ZigZag 周期骨架（阈值 25\%，圆点为周期转折点）}
\end{minipage}
\end{center}

\pfl{过去周期分析} 以 25\% 的 ZigZag 阈值划分，样本区间内股价形成 3 个主升段与 3 个主跌段。 最大主跌段为 2021 年 10 月 至 2024 年 6 月，32 个月-62.7\%；最大主升段出现在 2024 年 6 月 至 2026 年 4 月，22 个月+290.4\%。 最近一次转折出现在 2026 年 10 月（上涨段自 2026 年 6 月起，4 个月+31.9\%），当前处于该段之后的震荡之中。 公司股价较申万农林牧渔指数提前约 23 个月见底，是板块中较早定价周期反转的公司之一。 该子行业的周期来源与猪周期不同（第 \ref{sec:grain} 章、第 \ref{sec:soft} 章），价格更多由自身供给、进口政策与全球定价决定，公司 2026 年 6 月末资产负债率 36.3\%、半年 ROE 0.4\%（未年化），在本轮下行中属于随行业同步调整的一类。 综合区间形态看，该股单段涨跌幅度偏大、以趋势段为主（最大主升 +290.4\%、最大主跌 -62.7\%），年化波动率 40.4\% 与全样本中位数 37.3\% 相比波动与样本中位数接近。

\pfl{盈利情况} 2026 年上半年营业收入 4.18 亿元（同比 -11.1\%），归母净利润 0.05 亿元，同比 -25.4\%。 以年报为趋势对照，2023---2025 年营业收入由 10.86 亿元变为 10.52 亿元（两年年均复合增速 -1.6\%），归母净利润由 0.24 亿元变为 0.18 亿元。 按半年报口径，2026 年上半年的 ROE 为 0.4\%（未年化）、销售净利率 5.4\%、毛利率 16.2\%，ROE 在“其他种植业”6 家公司中按数值由高到低排第\zhnumber{5}位。 同期全样本半年 ROE 中位数 0.75\%，公司与之相差 0.39 个百分点（弱于中位数公司）。 从公司披露的原因看，业绩变动主要来自：2025 年，同比减少 44.68\%。 面对这一局面，公司披露的举措包括：2026 年 9 月，主动缩减部分低毛利业务。 综合看，公司在报告期维持盈利（低于全样本半年 ROE 中位数 0.75\%），利润的可持续性取决于其他种植业景气的延续与成本管控的执行。

\begin{center}\small\setlength{\tabcolsep}{3.4pt}
\captionof{table}[宏辉果蔬（603336）关键财务指标]{宏辉果蔬（603336）关键财务指标（合并报表口径；两个半年报期均未年化；现金短债比 $=$ 货币资金 $\div$ 短期债务）}
\begin{adjustbox}{max width=\textwidth}
\begin{tabular}{@{}lrrrrr@{}}
\toprule
指标 & 2023 年 & 2024 年 & 2025 年 & 2025 年半年报 & 2026 年半年报 \\
\midrule
营业收入（亿元） & 10.86 & 10.80 & 10.52 & 4.70 & 4.18 \\
归母净利润（亿元） & 0.24 & 0.18 & 0.18 & 0.07 & 0.05 \\
毛利率（\%） & 8.0 & 9.1 & 9.1 & 9.1 & 16.2 \\
ROE（\%） & 2.1 & 1.6 & 1.4 & 0.6 & 0.4 \\
资产负债率（\%） & 37.6 & 39.5 & 19.8 & 38.9 & 36.3 \\
经营活动现金流净额（亿元） & -0.03 & 0.35 & 2.05 & 0.38 & 0.94 \\
货币资金（亿元） & 0.39 & 0.59 & 4.09 & 0.78 & 4.04 \\
有息负债合计（亿元） & 6.23 & 6.80 & 2.99 & 6.79 & 7.43 \\
现金短债比（倍） & 0.10 & 0.13 & 2.17 & 0.18 & 2.37 \\
\bottomrule
\end{tabular}
\end{adjustbox}
\end{center}

\begin{center}
\begin{minipage}{\textwidth}\centering
\begin{tikzpicture}
\begin{groupplot}[
  group style={group size=2 by 1, horizontal sep=0.60cm},
  width=6.85cm, height=4.60cm,
  tick label style={font=\tiny},
  title style={font=\scriptsize\heiti},
  yticklabel style={font=\tiny},
  ylabel style={font=\tiny},
  grid=major, grid style={LightGray, line width=0.3pt},
  axis line style={InkGray, line width=0.5pt},
  enlarge x limits={abs=0.8},
  legend style={font=\scriptsize, at={(0.5,-0.20)}, anchor=north,
                legend columns=2, draw=none, fill=none, column sep=6pt},
]
\nextgroupplot[
  ybar, bar width=4pt,
  ymin=-1.13, ymax=12.69,
  xtick={0,1,2,3,4,5.7},
  xticklabels={2021,2022,2023,2024,2025,2026H1},
  ylabel={亿元},
]
\addplot[fill=AgriGreen!75, draw=AgriGreen, bar shift=-2.6pt] table[x=i, y=rev]{data/clean/profiles/fin_603336.dat};
\addplot[fill=SoftRed!60, draw=SoftRed, bar shift=2.6pt] table[x=i, y=ni]{data/clean/profiles/fin_603336.dat};
\addplot[fill=AgriGreen!30, draw=AgriGreen, pattern=north east lines, bar shift=-2.6pt] table[x=x, y=rev]{data/clean/profiles/finh_603336.dat};
\addplot[fill=SoftRed!20, draw=SoftRed, pattern=north east lines, bar shift=2.6pt] table[x=x, y=ni]{data/clean/profiles/finh_603336.dat};
\legend{营业收入（年/半年）, 归母净利润（年/半年）}
\nextgroupplot[
  ymin=-3.2, ymax=43.5,
  xtick={0,1,2,3,4,5.7},
  xticklabels={2021,2022,2023,2024,2025,2026H1},
  ylabel={\%},
]
\addplot[AgriGold, mark=*, mark size=1.3pt] table[x=i, y=roe]{data/clean/profiles/fin_603336.dat};
\addplot[OilBlue, mark=square*, mark size=1.3pt] table[x=i, y=debt]{data/clean/profiles/fin_603336.dat};
\addplot[only marks, AgriGold, mark=*, mark size=2.0pt] table[x=x, y=roe]{data/clean/profiles/finh_603336.dat};
\addplot[only marks, OilBlue, mark=square*, mark size=2.0pt] table[x=x, y=debt]{data/clean/profiles/finh_603336.dat};
\legend{ROE, 资产负债率}
\end{groupplot}
\end{tikzpicture}
\begin{tikzpicture}
\begin{axis}[
  name=finbot,
  width=13.75cm, height=3.10cm, scale only axis,
  ymin=-0.74, ymax=8.33,
  xtick={0,1,2,3,4,5.7},
  xticklabels={2021,2022,2023,2024,2025,2026H1},
  tick label style={font=\tiny},
  yticklabel style={font=\tiny},
  ylabel={亿元}, ylabel style={font=\tiny},
  ybar, bar width=4pt,
  grid=major, grid style={LightGray, line width=0.3pt},
  axis line style={InkGray, line width=0.5pt},
  enlarge x limits={abs=0.8},
  legend style={font=\scriptsize, at={(0.5,-0.26)}, anchor=north,
                legend columns=2, draw=none, fill=none, column sep=10pt},
]
\addplot[fill=AgriGreen!55, draw=AgriGreen, bar shift=-3.0pt] table[x=i, y=cash]{data/clean/profiles/fin_603336.dat};
\addplot[fill=SoftRed!45, draw=SoftRed, bar shift=3.0pt] table[x=i, y=ibd]{data/clean/profiles/fin_603336.dat};
\addplot[fill=AgriGreen!25, draw=AgriGreen, pattern=north east lines, bar shift=-3.0pt] table[x=x, y=cash]{data/clean/profiles/finh_603336.dat};
\addplot[fill=SoftRed!20, draw=SoftRed, pattern=north east lines, bar shift=3.0pt] table[x=x, y=ibd]{data/clean/profiles/finh_603336.dat};
\legend{货币资金（左轴）, 有息负债（左轴）}
\end{axis}
\begin{axis}[
  at={(finbot.south east)}, anchor=south east,
  width=13.75cm, height=3.10cm, scale only axis,
  axis y line*=right, axis x line*=none, xtick=\empty,
  ymin=-1.02, ymax=2.78,
  tick label style={font=\tiny},
  yticklabel style={font=\tiny}, scaled y ticks=false,
  legend style={font=\scriptsize, at={(0.5,-0.44)}, anchor=north,
                legend columns=1, draw=none, fill=none},
]
\addplot[OilBlue, mark=*, mark size=2.2pt, line width=1.3pt] table[x=i, y=ocf]{data/clean/profiles/fin_603336.dat};
\addplot[only marks, OilBlue, mark=*, mark size=3.2pt] table[x=x, y=ocf]{data/clean/profiles/finh_603336.dat};
\legend{经营活动现金流净额（右轴，亿元，蓝点）}
\end{axis}
\end{tikzpicture}

\captionof{figure}[宏辉果蔬（603336）近年主要财务指标]{宏辉果蔬（603336）近年主要财务指标（左上：营业收入与归母净利润；右上：ROE 与资产负债率；下：货币资金、有息负债为左轴柱状，经营活动现金流净额为其右轴的蓝点折线。
2021---2025 年为年报口径，2026H1 为 2026 年半年报/半年末，与其前一年报之间留出间隔、以斜纹或空心点区分；半年报数据未年化）}
\end{minipage}
\end{center}

\pfl{财务分析} 2026 年 6 月末资产负债率 36.3\%，较 2025 年末上升 16.6 个百分点（样本中位数 52.2\%，所处三级行业内按由低到高排第\zhnumber{1}位）。 2026 年 6 月末货币资金 4.04 亿元、短期债务（短期借款与一年内到期非流动负债）1.71 亿元、长期债务 5.73 亿元，有息负债合计 7.43 亿元，现金短债比 2.37 倍——覆盖充分。 净有息负债 3.40 亿元，相当于其 2026 年上半年营业收入的 81\%。 流动比率 2.96、速动比率 2.70，短期指标尚可。 2026 年上半年经营活动现金流净额 0.94 亿元，经营现金流与利润同向为正，内生资金可以支撑运营。 2026 年 6 月末商誉 4.07 亿元，占归母净资产的 21.8\%，是净资产中最脆弱的一块。 综合看，现金短债比 2.37 倍，短期偿付压力可控；资产负债率低于样本中位数 15.9 个百分点；有息负债中长期债务占比更高，期限结构对短债滚动的压力较小。 公司披露的股东与质押风险为：2025 年第一季度，报告显示黄俊辉持股 2.52 亿股（44.20\%），其中质押 5,840.00 万股。

\pfl{重大战略转型} 从公开披露看，公司的战略推进集中在：并购与重组（2 项）：2026 年 6 月，支付第三期股权转让款 2.10 亿元（对价的 30\%）。 公告口径上，近三年（2023 年 9 月---2026 年 9 月）公司披露重大重组与并购类公告 5 条、对外投资与转型类公告 1 条、回购与股权激励类公告 0 条，合计公告 385 条。 第一大行业“果蔬服务业”的收入占比由 2021 年年报的 96.8\% 变为 72.1\%（24.7 个百分点），收入结构出现再平衡。 综合判断，报告期内公司有 6 条与战略相关的公告落地（重大重组与并购 5 条、对外投资与转型 1 条），资本运作与主业调整之间存在直接关联。

\begin{center}\small\setlength{\tabcolsep}{4pt}
\captionof{table}[宏辉果蔬（603336）历史融资与资本运作]{宏辉果蔬（603336）历史融资与资本运作（据公司公告与公开报道整理；金额为披露值，含首次公开发行、再融资、债券、银行借款与重整投资等）}
\begin{adjustbox}{max width=\textwidth}
\begin{tabular}{@{}p{1.45cm}p{2.55cm}p{4.05cm}p{4.95cm}@{}}
\toprule
时间 & 方式 & 规模 & 用途与说明 \\
\midrule
2016-11 & IPO（上交所主板） & 发行 3,\allowbreak{}335 万股，\allowbreak{}发行价 9.31 元/\allowbreak{}股，\allowbreak{}募集资金总额 3.10 亿元 & 天津、\allowbreak{}上海、\allowbreak{}广州果蔬加工配送基地建设项目及宏辉果蔬信息化系统建设项目；占发行后总股本 25.01\% \\
2025-06 & 控股股东协议转让（控制权变更，非上市公司融资） & 申泽瑞泰以 5.68 元/\allowbreak{}股受让黄俊辉持有的 1.51 亿股（26.54\%） & 股权转让款由转让方取得，\allowbreak{}不进入上市公司；转让后黄俊辉持股由 2.52 亿股（44.19\%）降至 1.01 亿股（17.66\%） \\
2026-04 & 对外投资（自有及自筹资金） & 7.00 亿元受让江西施美药业股份有限公司 41.128\% 股份 & 布局医药第二增长曲线，\allowbreak{}取得施美药业实质管理控制权与并表；标的 100\% 股权作价 17.02 亿元 \\
\bottomrule
\end{tabular}
\end{adjustbox}
\end{center}

\pfl{历史融投资情况} 从现金流量表看，2025 年公司取得借款收到的现金 3.71 亿元、偿还债务支付的现金 5.18 亿元、吸收权益性投资 0.01 亿元，筹资活动现金流净额 -1.67 亿元（2023 年为取得借款 4.43 亿元、偿还债务 3.89 亿元）。 2026 年上半年筹资活动现金流净额为 4.36 亿元（取得借款 5.47 亿元、偿还债务 1.06 亿元，未年化），仍处于净融入状态。 2025 年分配股利、利润或偿付利息支付的现金合计 0.21 亿元。 公开披露的融资记录共 3 笔，工具类型以 IPO 为主；金额与用途见下表。 其中规模最大的两笔为：2016 年 11 月，IPO，发行 3,335 万股，发行价 9.31 元/股，募集资金总额 3.10 亿元；2026 年 4 月，对外投资，7.00 亿元受让江西施美药业股份有限公司 41.128\% 股份。 股本方面，近三年披露股本变动 10 次（可转债转股 9 次、送股 1 次、转增 1 次）；总股本由 2021 年末的 4.39 亿股变为 5.70 亿股（+30.02\%）。 分红方面，近三年现金分红 3 次、合计每股派现 0.05 元，最近一次实施在 2025 年 12 月。 近三年“再融资”类公告 3 条，涉及可转债、募集资金使用。 综合看，公司融资结构上以债务滚动为主、股权融资为补充；2025 年筹资活动现金流净额 -1.67 亿元，融资节奏仍以维持存量债务周转为主，与前述经营与投资现金流的方向基本一致。

\pfl{未来潜在融资需求分析} 2026 年 6 月末货币资金 4.04 亿元，而短期债务 1.71 亿元（现金短债比 2.37 倍），2026 年上半年经营现金流净额 0.94 亿元。账面货币资金可以覆盖短期债务，缺口不体现在静态偿付层面。 公司 2026 年 6 月末资产负债率 36.3\%、2026 年上半年 ROE 0.4\%（未年化），资产负债表稳健、盈利水平平淡，暂不存在刚性融资需求；潜在的融资需求以技改扩产、并购或被并购为主。 公开披露的后续资金安排线索包括：债务与授信安排：2026 年度，公司及子（孙）公司预计互为提供总额不超过 16.00 亿元的担保，用于满足经营与融资需要。 综合看，公司在静态偿付层面不存在缺口，外部融资更多是配合扩张节奏与并购整合的主动性安排，而非偿债压力驱动。




\subsubsection{ST荃银（300087）}
\label{co:300087}

\pfl{公司基本情况} 合肥的水稻、玉米、小麦种子龙头，2010 年创业板上市；2021 年起由先正达集团系的中化现代农业、中种集团控股，2026 年 1 月中种集团完成部分要约收购后持股 40.51\%；2025 年出现上市以来首次年度亏损，2026 年 6 月因 2024 年年报虚假记载被实施其他风险警示并受行政处罚。 公司于 2010-05-26 在深交所创业板首发上市，注册资本 9.47 亿元，总股本 9.47 亿股。 截至 2026 年 10 月 12 日收盘价 5.58 元，对应总市值 52.9 亿元，近一年-39.5\%，总市值在三级行业“种子”的 8 家公司中按由大到小排第\zhnumber{5}位，近一年涨跌幅低于全样本中位数（-10.4\%）29.1 个百分点。 按 2026 年半年报披露的行业构成，收入占比前两位为农业 97.7\%；其他(补充) 2.3\%。 股权结构方面，控股股东为中种集团控股股东为先正达集团股份有限公司；实际控制人为国务院国资委。

\pfl{股权结构与控制权} 按 2026 年 6 月 30 日的股东名册，前三大股东为中国种子集团有限公司 持股 40.51\%；张琴 持股 4.81\%；贾桂兰 持股 2.52\%，股东户数 34,799 户。 股权结构的另一条线索是，持股 40.51\%；2026 年 1 月，公司与控股股东中种集团签署股权委托管理协议。

\pfl{经营要点} 业务与产能层面，2026 年 1 月，要约收购完成后，中种集团持有荃银高科 40.51\% 股份。

\pfl{所处行业的基本面} 种子的供给由制种面积与库存决定、需求由粮价与种植意愿决定，制种面积在 2025 年起出现萎缩，是本轮供给端最先出清的环节（第 \ref{sec:framework} 章的供需框架）。 行业层面的最新情况是：2026 年，全国杂交稻制种面积继续增加 5.8\%、达到 238.4 万亩。

\pfl{竞争优势} 把该公司放回同类公司比较，2026 年上半年营业收入 11.8 亿元、在三级行业“种子”的 8 家公司中排第\zhnumber{4}位；毛利率 21.1\%、低于行业内中位数 7.3 个百分点；半年 ROE -3.3\%（未年化）、低于行业内中位数（0.17\%）；资产负债率 64.8\%、高于行业内中位数（37.3\%）。 公司区别于同类公司的经营特征包括：公司 2025 年度种子销售收入 27.92 亿元（27.92 亿元）；公司 2025 年研发投入 2.03 亿元、占种子销售收入 7.28\%（上年为 5.65\%）；公司设立荃银农业科学院。 就当期报表看，公司尚未把经营层面的特征转化为高于行业的资本回报，周期底部的相对优势更多体现在现金流与资产负债表上。

\pfl{历史股价} 以 2021 年 1 月为起点、2026 年 10 月为终点，股价由 31.89 元变为 5.58 元，累计 -82.5\%，区间最高 35.18 元（2021 年 8 月）、最低 5.12 元（2026 年 7 月）；当前价相当于区间最高价的 15.9\%，为区间最低价的 1.09 倍。 月度收益的年化波动率 56.8\%，高于全样本中位数 37.3\%，在三级行业“种子”的 8 家公司中波动率由高到低排第\zhnumber{1}位。 把股价阶段与公司事件对齐看，2023 年 1 月 至 2024 年 8 月股价下跌-64.8\%，同期 2026 年上半年，公告显示该笔收购最终交易对价为 1.83 亿元。 整体上看，区间的几处大级别拐点都能在公司的资本运作、股权变动或风险事件上找到对应，股价波动有明显的公司层面驱动，而非单纯的行业 beta。

\begin{center}
\begin{minipage}{\textwidth}\centering
\begin{tikzpicture}
\begin{axis}[
  width=12.6cm, height=3.9cm, scale only axis,
  xmin=2020.922, xmax=2026.972, ymin=4.76, ymax=37.64,
  xtick={2021,2022,2023,2024,2025,2026},
  yearticks,
  yticklabel style={font=\scriptsize},
  ylabel={元/股}, ylabel style={font=\scriptsize},
  grid=major, grid style={LightGray, line width=0.3pt},
  axis line style={InkGray, line width=0.5pt},
  every axis plot/.append style={line width=0.9pt},
]
\addplot[AgriGreen] table[x=x,y=y]{data/clean/profiles/px_300087.dat};
\addplot[SoftRed, dashed, line width=0.7pt] table[x=x,y=y]{data/clean/profiles/pxzz_300087.dat};
\addplot[only marks, mark=*, mark size=1.1pt, SoftRed] table[x=x,y=y]{data/clean/profiles/pxzz_300087.dat};
\end{axis}
\end{tikzpicture}

\captionof{figure}[ST荃银（300087）股价与周期骨架]{ST荃银（300087）月度收盘价（元/股）与 ZigZag 周期骨架（阈值 25\%，圆点为周期转折点）}
\end{minipage}
\end{center}

\pfl{过去周期分析} 以 25\% 的 ZigZag 阈值划分，样本区间内股价形成 4 个主升段与 5 个主跌段。 最大主升段出现在 2024 年 8 月 至 2024 年 10 月，2 个月+121.5\%；最大主跌段为 2023 年 1 月 至 2024 年 8 月，19 个月-64.8\%。 最近一次转折出现在 2026 年 7 月（下跌段自 2025 年 12 月起，7 个月-56.2\%），当前处于该段的延续之中。 与申万农林牧渔指数的最近一次底部相比，该公司见底晚约 2 个月，在本轮下行中属于调整更慢的一类。 种子行业的价格由制种成本、粮价与政策（转基因商业化、种业振兴）共同决定（第 \ref{sec:planting} 节），与生猪价格的联动有限，公司 2026 年 6 月末资产负债率 64.8\%、半年 ROE -3.3\%（未年化），在本轮下行中属于随行业同步调整的一类。 综合区间形态看，该股单段涨跌幅度偏大、以趋势段为主（最大主升 +121.5\%、最大主跌 -64.8\%），年化波动率 56.8\% 与全样本中位数 37.3\% 相比波动明显大于样本中位数。

\pfl{盈利情况} 2026 年上半年营业收入 11.79 亿元（同比 -17.9\%），归母净利润 -0.55 亿元，亏损同比扩大 0.14 亿元。 以年报为趋势对照，2023---2025 年营业收入由 41.03 亿元变为 44.95 亿元（两年年均复合增速 +4.7\%），归母净利润由 2.68 亿元变为 -2.12 亿元。 按半年报口径，2026 年上半年的 ROE 为 -3.3\%（未年化）、销售净利率 -7.3\%、毛利率 21.1\%，ROE 在“种子”8 家公司中按数值由高到低排第\zhnumber{7}位。 同期全样本半年 ROE 中位数 0.75\%，公司与之相差 4.09 个百分点（弱于中位数公司）。 从公司披露的原因看，业绩变动主要来自：2026 年 4 月，公司业绩预告显示 2025 年归母净利润亏损 1.80 亿元—2.70 亿元；2026 年 6 月，导致少计提信用减值损失 1,871.51 万元、虚增利润总额 1,871.51 万元。 综合看，公司仍处亏损区间、亏损同比扩大，半年 ROE 在三级行业内按由高到低排第\zhnumber{7}位，单位盈利能力的修复依赖行业景气与自身成本端的改善，报表弹性主要体现为价格与销量的方向性变化。

\begin{center}\small\setlength{\tabcolsep}{3.4pt}
\captionof{table}[ST荃银（300087）关键财务指标]{ST荃银（300087）关键财务指标（合并报表口径；两个半年报期均未年化；现金短债比 $=$ 货币资金 $\div$ 短期债务）}
\begin{adjustbox}{max width=\textwidth}
\begin{tabular}{@{}lrrrrr@{}}
\toprule
指标 & 2023 年 & 2024 年 & 2025 年 & 2025 年半年报 & 2026 年半年报 \\
\midrule
营业收入（亿元） & 41.03 & 47.09 & 44.95 & 14.36 & 11.79 \\
归母净利润（亿元） & 2.68 & 0.97 & -2.12 & -0.41 & -0.55 \\
毛利率（\%） & 26.9 & 25.0 & 21.5 & 13.6 & 21.1 \\
ROE（\%） & 15.0 & 5.0 & -11.6 & -2.1 & -3.3 \\
资产负债率（\%） & 61.6 & 65.9 & 69.0 & 62.3 & 64.8 \\
经营活动现金流净额（亿元） & 2.84 & 0.65 & 1.99 & -7.17 & -7.03 \\
货币资金（亿元） & 15.31 & 16.00 & 16.14 & 8.37 & 7.67 \\
有息负债合计（亿元） & 14.78 & 21.70 & 22.88 & 21.90 & 22.07 \\
现金短债比（倍） & 1.85 & 0.91 & 1.01 & 0.62 & 0.40 \\
\bottomrule
\end{tabular}
\end{adjustbox}
\end{center}

\begin{center}
\begin{minipage}{\textwidth}\centering
\begin{tikzpicture}
\begin{groupplot}[
  group style={group size=2 by 1, horizontal sep=0.60cm},
  width=6.85cm, height=4.60cm,
  tick label style={font=\tiny},
  title style={font=\scriptsize\heiti},
  yticklabel style={font=\tiny},
  ylabel style={font=\tiny},
  grid=major, grid style={LightGray, line width=0.3pt},
  axis line style={InkGray, line width=0.5pt},
  enlarge x limits={abs=0.8},
  legend style={font=\scriptsize, at={(0.5,-0.20)}, anchor=north,
                legend columns=2, draw=none, fill=none, column sep=6pt},
]
\nextgroupplot[
  ybar, bar width=4pt,
  ymin=-7.04, ymax=52.99,
  xtick={0,1,2,3,4,5.7},
  xticklabels={2021,2022,2023,2024,2025,2026H1},
  ylabel={亿元},
]
\addplot[fill=AgriGreen!75, draw=AgriGreen, bar shift=-2.6pt] table[x=i, y=rev]{data/clean/profiles/fin_300087.dat};
\addplot[fill=SoftRed!60, draw=SoftRed, bar shift=2.6pt] table[x=i, y=ni]{data/clean/profiles/fin_300087.dat};
\addplot[fill=AgriGreen!30, draw=AgriGreen, pattern=north east lines, bar shift=-2.6pt] table[x=x, y=rev]{data/clean/profiles/finh_300087.dat};
\addplot[fill=SoftRed!20, draw=SoftRed, pattern=north east lines, bar shift=2.6pt] table[x=x, y=ni]{data/clean/profiles/finh_300087.dat};
\legend{营业收入（年/半年）, 归母净利润（年/半年）}
\nextgroupplot[
  ymin=-18.1, ymax=77.1,
  xtick={0,1,2,3,4,5.7},
  xticklabels={2021,2022,2023,2024,2025,2026H1},
  ylabel={\%},
]
\addplot[AgriGold, mark=*, mark size=1.3pt] table[x=i, y=roe]{data/clean/profiles/fin_300087.dat};
\addplot[OilBlue, mark=square*, mark size=1.3pt] table[x=i, y=debt]{data/clean/profiles/fin_300087.dat};
\addplot[only marks, AgriGold, mark=*, mark size=2.0pt] table[x=x, y=roe]{data/clean/profiles/finh_300087.dat};
\addplot[only marks, OilBlue, mark=square*, mark size=2.0pt] table[x=x, y=debt]{data/clean/profiles/finh_300087.dat};
\legend{ROE, 资产负债率}
\end{groupplot}
\end{tikzpicture}
\begin{tikzpicture}
\begin{axis}[
  name=finbot,
  width=13.75cm, height=3.10cm, scale only axis,
  ymin=-2.29, ymax=25.63,
  xtick={0,1,2,3,4,5.7},
  xticklabels={2021,2022,2023,2024,2025,2026H1},
  tick label style={font=\tiny},
  yticklabel style={font=\tiny},
  ylabel={亿元}, ylabel style={font=\tiny},
  ybar, bar width=4pt,
  grid=major, grid style={LightGray, line width=0.3pt},
  axis line style={InkGray, line width=0.5pt},
  enlarge x limits={abs=0.8},
  legend style={font=\scriptsize, at={(0.5,-0.26)}, anchor=north,
                legend columns=2, draw=none, fill=none, column sep=10pt},
]
\addplot[fill=AgriGreen!55, draw=AgriGreen, bar shift=-3.0pt] table[x=i, y=cash]{data/clean/profiles/fin_300087.dat};
\addplot[fill=SoftRed!45, draw=SoftRed, bar shift=3.0pt] table[x=i, y=ibd]{data/clean/profiles/fin_300087.dat};
\addplot[fill=AgriGreen!25, draw=AgriGreen, pattern=north east lines, bar shift=-3.0pt] table[x=x, y=cash]{data/clean/profiles/finh_300087.dat};
\addplot[fill=SoftRed!20, draw=SoftRed, pattern=north east lines, bar shift=3.0pt] table[x=x, y=ibd]{data/clean/profiles/finh_300087.dat};
\legend{货币资金（左轴）, 有息负债（左轴）}
\end{axis}
\begin{axis}[
  at={(finbot.south east)}, anchor=south east,
  width=13.75cm, height=3.10cm, scale only axis,
  axis y line*=right, axis x line*=none, xtick=\empty,
  ymin=-9.99, ymax=7.77,
  tick label style={font=\tiny},
  yticklabel style={font=\tiny}, scaled y ticks=false,
  legend style={font=\scriptsize, at={(0.5,-0.44)}, anchor=north,
                legend columns=1, draw=none, fill=none},
]
\addplot[OilBlue, mark=*, mark size=2.2pt, line width=1.3pt] table[x=i, y=ocf]{data/clean/profiles/fin_300087.dat};
\addplot[only marks, OilBlue, mark=*, mark size=3.2pt] table[x=x, y=ocf]{data/clean/profiles/finh_300087.dat};
\legend{经营活动现金流净额（右轴，亿元，蓝点）}
\end{axis}
\end{tikzpicture}

\captionof{figure}[ST荃银（300087）近年主要财务指标]{ST荃银（300087）近年主要财务指标（左上：营业收入与归母净利润；右上：ROE 与资产负债率；下：货币资金、有息负债为左轴柱状，经营活动现金流净额为其右轴的蓝点折线。
2021---2025 年为年报口径，2026H1 为 2026 年半年报/半年末，与其前一年报之间留出间隔、以斜纹或空心点区分；半年报数据未年化）}
\end{minipage}
\end{center}

\pfl{财务分析} 2026 年 6 月末资产负债率 64.8\%，较 2025 年末下降 4.2 个百分点（样本中位数 52.2\%，所处三级行业内按由低到高排第\zhnumber{8}位）。 2026 年 6 月末货币资金 7.67 亿元、短期债务（短期借款与一年内到期非流动负债）19.06 亿元、长期债务 3.01 亿元，有息负债合计 22.07 亿元，现金短债比 0.40 倍——覆盖偏紧。 净有息负债 14.40 亿元，相当于其 2026 年上半年营业收入的 122\%。 流动比率 1.29、速动比率 0.55，短期偿债能力处于正常区间。 2026 年上半年经营活动现金流净额 -7.03 亿元，经营现金流与利润同时为负，日常运营对债务与股东投入的依赖度较高。 2026 年 6 月末在建工程 0.53 亿元，较上年末减少 0.95 亿元，在建投入在收缩。 综合看，账面货币资金对有息负债与短期债务的覆盖不足，滚动偿付对新增融资的依赖度较高；资产负债率高于样本中位数 12.6 个百分点；有息负债中短期债务占比更高，期限结构仍以滚动续作为主。 公司披露的监管与合规风险为：2026 年 1 月，收到中国证监会立案告知书，扣非净利润亏损 2.45 亿元—3.50 亿元。

\pfl{重大战略转型} 近三年公司的战略动作集中在：并购与重组（1 项）：2021 年 1 月，方式包括资产重组、业务调整、委托管理、设立合资公司等。 公告口径上，近三年（2023 年 9 月---2026 年 9 月）公司披露重大重组与并购类公告 18 条、对外投资与转型类公告 0 条、回购与股权激励类公告 8 条，合计公告 318 条。 第一大行业仍是“农业”，收入占比 97.7\%，与 2021 年年报相比没有方向性变化。 综合判断，报告期内公司有 18 条与战略相关的公告落地（重大重组与并购 18 条、对外投资与转型 0 条），资本运作与主业调整之间存在直接关联。

\begin{center}\small\setlength{\tabcolsep}{4pt}
\captionof{table}[ST荃银（300087）历史融资与资本运作]{ST荃银（300087）历史融资与资本运作（据公司公告与公开报道整理；金额为披露值，含首次公开发行、再融资、债券、银行借款与重整投资等）}
\begin{adjustbox}{max width=\textwidth}
\begin{tabular}{@{}p{1.45cm}p{2.55cm}p{4.05cm}p{4.95cm}@{}}
\toprule
时间 & 方式 & 规模 & 用途与说明 \\
\midrule
2010-05 & IPO（深交所创业板） & 发行 1,\allowbreak{}320 万股，\allowbreak{}发行价 35.60 元/\allowbreak{}股，\allowbreak{}实际募集资金 4.70 亿元 & 全部用于与公司主营业务相关的项目及主营业务发展所需的营运资金；占发行后总股本 25\% \\
2015-11 & 发行股份及支付现金购买资产并募集配套资金（未获通过） & 拟购买同路农业 60\% 股权作价 1.44 亿元 & 支付购买资产现金对价及交易费用；交易完成后公司将发行 1,\allowbreak{}049.82 万股 \\
2017-08 & 发行股份及支付现金购买资产并募集配套资金（修订，被否） & 拟购买同路农业 100\% 股权 & 支付购买资产现金对价及本次重组相关费用；2017 年 8 月 25 日临时股东大会 17 项议案全部被否决（中植系、\allowbreak{}大北农系等股东集体反对） \\
2021-07 & 向特定对象发行A股（上市后首次再融资） & 向 6 名特定对象发行 2,\allowbreak{}376.84 万股，\allowbreak{}发行价 23.14 元/\allowbreak{}股 & 研发创新体系建设项目、\allowbreak{}农作物种子海外育繁推一体化建设项目、\allowbreak{}青贮玉米品种产业化及种养结合项目；2021 年 5 月 27 日获证监会同意注册批复，\allowbreak{}2020 年 11 月调减为不超过 5.50 亿元 \\
2023-04-23 & 使用募集资金收购股权 & 第五届董事会第十一次会议通过使用募集资金 2.24 亿元收购河北新纪元种业有限公司股权 & 收购河北新纪元种业股权；2026 年上半年公告显示该笔收购最终交易对价为 1.83 亿元 \\
2025-11-20 & 一致行动关系终止 & 贾桂兰、\allowbreak{}王玉林合计 7,\allowbreak{}435 万股（占总股本 7.85\%）不可撤销承诺预受要约并放弃所持全部股份表决权 & 中种集团表决权比例由 28.36\% 变更为 20.51\%，\allowbreak{}为要约收购前置安排 \\
2026-01 & 控股股东部分要约收购（收购人现金出资，非上市公司融资） & 中种集团以 11.85 元/\allowbreak{}股要约收购 1.89 亿股（占已发行股份 20.00\%） & 提升控股股东持股比例至 40.51\% \\
\bottomrule
\end{tabular}
\end{adjustbox}
\end{center}

\pfl{历史融投资情况} 从现金流量表看，2025 年公司取得借款收到的现金 20.26 亿元、偿还债务支付的现金 19.61 亿元、吸收权益性投资 0.02 亿元，筹资活动现金流净额 -1.16 亿元（2023 年为取得借款 11.99 亿元、偿还债务 8.96 亿元）。 2026 年上半年筹资活动现金流净额为 -0.99 亿元（取得借款 8.82 亿元、偿还债务 9.09 亿元，未年化），已转为净流出，处于净还债状态。 2025 年分配股利、利润或偿付利息支付的现金合计 1.22 亿元。 公开披露的融资记录共 7 笔，工具类型以 IPO、发行股份及支付现金购、定向增发为主；金额与用途见下表。 其中规模最大的两笔为：2023 年 4 月，使用募集资金收购股权，第五届董事会第十一次会议通过使用募集资金 2.24 亿元收购河北新纪元种业有限公司股权；2010 年 5 月，IPO，发行 1,320 万股，发行价 35.60 元/股，实际募集资金 4.70 亿元。 股本方面，近三年披露股本变动 4 次（转增 1 次）；总股本由 2021 年末的 4.54 亿股变为 9.47 亿股（+108.62\%）。 分红方面，近三年现金分红 3 次、合计每股派现 0.25 元，最近一次实施在 2024 年 12 月。 近三年“再融资”类公告 1 条，涉及定向增发。 综合看，公司融资结构上以债务滚动为主、股权融资为补充；2025 年筹资活动现金流净额 -1.16 亿元，融资节奏仍以维持存量债务周转为主，与前述经营与投资现金流的方向基本一致。

\pfl{未来潜在融资需求分析} 2026 年 6 月末货币资金 7.67 亿元，而短期债务 19.06 亿元（现金短债比 0.40 倍），2026 年上半年经营现金流净额 -7.03 亿元。按“短期债务减去账面货币资金”的滚动偿付口径，静态缺口约 11.39 亿元；若再计入上半年亏损的年化额（1.11 亿元），维持一年正常经营所需的外部资金约 11.39---12.50 亿元。 按第 \ref{sec:financing-need} 节的三档划分，公司属于第二档“保壳型”：近三年重大重组与并购类公告 18 条，2026 年上半年归母净利润 -0.55 亿元，融资需求更多体现为“资产置换 + 维持上市地位”，而不是扩产。 资金需求还要叠加在建工程：期末在建工程 0.53 亿元、2026 年上半年购建固定资产等支出 0.73 亿元（未年化），这部分支出在周期底部通常会被主动放缓。 综合看，公司的外部资金需求以补充流动性与项目投入为主，滚动偿付口径下的静态缺口约 11.39 亿元，融资的时点与规模更多取决于信用条件与资本市场窗口。




\subsubsection{敦煌种业（600354）}
\label{co:600354}

\pfl{公司基本情况} 敦煌种业成立于 1998 年，甘肃酒泉的种业与食品加工企业，2004 年上市的老牌种业公司；2018—2019 年连续亏损曾被实施退市风险警示（*ST 敦种），2021 年撤销；2024 年控股股东两度无偿划转、最终变更为酒钢集团，2025—2026 年种子主业量利齐升，但母公司累计未分配利润仍为大额负数，长期不能分红。 公司于 2004-01-15 在上交所首发上市，注册资本 5.28 亿元，总股本 5.28 亿股。 截至 2026 年 10 月 12 日收盘价 9.31 元，对应总市值 49.1 亿元，近一年+46.9\%，总市值在三级行业“种子”的 8 家公司中按由大到小排第\zhnumber{6}位，近一年涨跌幅高于全样本中位数（-10.4\%）57.2 个百分点。 按 2025 年年报披露的行业构成，收入占比前三位为种子 77.1\%；食品与贸易 16.5\%；其他 6.3\%。

\pfl{股权结构与控制权} 按 2026 年 6 月 30 日的股东名册，前三大股东为酒泉钢铁(集团)有限责任公司 持股 11.03\%；江苏融卓投资有限公司 持股 7.94\%；敦煌沙州能源(集团)有限责任公司 持股 5.02\%，股东户数 55,015 户。 股东结构与控制权方面，持股 11.03\%。

\pfl{经营要点} 业务与产能层面，2026 年上半年，分部收入：种子 7.79 亿元、食品与贸易 8,632.72 万元。

\pfl{所处行业的基本面} 种子的供给由制种面积与库存决定、需求由粮价与种植意愿决定，制种面积在 2025 年起出现萎缩，是本轮供给端最先出清的环节（第 \ref{sec:framework} 章的供需框架）。 与该公司直接相关的行业数据包括：2026 年上半年随着制种端主动调减产能；2023 年，2026 年再降至 358.5 万亩、同比下降 7\%。

\pfl{竞争优势} 把该公司放回同类公司比较，2026 年上半年营业收入 8.7 亿元、在三级行业“种子”的 8 家公司中排第\zhnumber{5}位；毛利率 48.6\%、高于行业内中位数 20.2 个百分点；半年 ROE 18.7\%（未年化）、高于行业内中位数（0.17\%）；资产负债率 37.5\%、高于行业内中位数（37.3\%）。 优势在报表上主要体现为高于行业平均的毛利率，这在周期底部是抵御亏损的主要缓冲。

\pfl{历史股价} 本报告样本区间（2021 年 1 月 至 2026 年 10 月）内，股价由 4.44 元变为 9.31 元，累计 +109.7\%，区间最高 9.31 元（2026 年 10 月）、最低 3.95 元（2021 年 3 月）；当前价相当于区间最高价的 100.0\%，为区间最低价的 2.36 倍。 月度收益的年化波动率 38.1\%，高于全样本中位数 37.3\%，在三级行业“种子”的 8 家公司中波动率由高到低排第\zhnumber{5}位。 把股价阶段与公司事件对齐看，2021 年 3 月 至 2023 年 1 月股价上涨+110.1\%，同期（2021 年 5 月）公司撤销退市风险警示。 整体上看，区间的几处大级别拐点都能在公司的资本运作、股权变动或风险事件上找到对应，股价波动有明显的公司层面驱动，而非单纯的行业 beta。

\begin{center}
\begin{minipage}{\textwidth}\centering
\begin{tikzpicture}
\begin{axis}[
  width=12.6cm, height=3.9cm, scale only axis,
  xmin=2020.922, xmax=2026.972, ymin=3.67, ymax=9.96,
  xtick={2021,2022,2023,2024,2025,2026},
  yearticks,
  yticklabel style={font=\scriptsize},
  ylabel={元/股}, ylabel style={font=\scriptsize},
  grid=major, grid style={LightGray, line width=0.3pt},
  axis line style={InkGray, line width=0.5pt},
  every axis plot/.append style={line width=0.9pt},
]
\addplot[AgriGreen] table[x=x,y=y]{data/clean/profiles/px_600354.dat};
\addplot[SoftRed, dashed, line width=0.7pt] table[x=x,y=y]{data/clean/profiles/pxzz_600354.dat};
\addplot[only marks, mark=*, mark size=1.1pt, SoftRed] table[x=x,y=y]{data/clean/profiles/pxzz_600354.dat};
\end{axis}
\end{tikzpicture}

\captionof{figure}[敦煌种业（600354）股价与周期骨架]{敦煌种业（600354）月度收盘价（元/股）与 ZigZag 周期骨架（阈值 25\%，圆点为周期转折点）}
\end{minipage}
\end{center}

\pfl{过去周期分析} 以 25\% 的 ZigZag 阈值划分，样本区间内股价形成 3 个主升段与 3 个主跌段。 最大主跌段为 2023 年 1 月 至 2024 年 8 月，19 个月-43.0\%；最大主升段出现在 2021 年 3 月 至 2023 年 1 月，22 个月+110.1\%。 最近一次转折出现在 2026 年 10 月（上涨段自 2026 年 6 月起，4 个月+95.6\%），当前处于该段之后的震荡之中。 公司股价较申万农林牧渔指数提前约 62 个月见底，是板块中较早定价周期反转的公司之一。 种子行业的价格由制种成本、粮价与政策（转基因商业化、种业振兴）共同决定（第 \ref{sec:planting} 节），与生猪价格的联动有限，公司 2026 年上半年实现盈利、半年 ROE 18.7\%（未年化），好于全样本中位数（0.75\%），下行周期对其报表的冲击小于同业。 综合区间形态看，该股单段涨跌幅度偏大、以趋势段为主（最大主升 +110.1\%、最大主跌 -43.0\%），年化波动率 38.1\% 与全样本中位数 37.3\% 相比波动与样本中位数接近。

\pfl{盈利情况} 2026 年上半年营业收入 8.68 亿元（同比 +20.9\%），归母净利润 1.45 亿元，同比 +165.7\%。 以年报为趋势对照，2023---2025 年营业收入由 11.58 亿元变为 11.61 亿元（两年年均复合增速 +0.1\%），归母净利润由 0.41 亿元变为 0.47 亿元。 按半年报口径，2026 年上半年的 ROE 为 18.7\%（未年化）、销售净利率 33.2\%、毛利率 48.6\%，ROE 在“种子”8 家公司中按数值由高到低排第\zhnumber{1}位。 按同一口径，三级行业“种子”8 家公司的半年 ROE 中位数为 0.17\%，该公司高于其中位数 18.53 个百分点。 公司给出的应对方向是：2026 年 8 月，依托精准市场定位与多渠道推广，自有核心优势品种销量稳步攀升。 综合看，公司在报告期维持盈利（高于全样本半年 ROE 中位数 0.75\%），利润的可持续性取决于种子景气的延续与成本管控的执行。

\begin{center}\small\setlength{\tabcolsep}{3.4pt}
\captionof{table}[敦煌种业（600354）关键财务指标]{敦煌种业（600354）关键财务指标（合并报表口径；两个半年报期均未年化；现金短债比 $=$ 货币资金 $\div$ 短期债务）}
\begin{adjustbox}{max width=\textwidth}
\begin{tabular}{@{}lrrrrr@{}}
\toprule
指标 & 2023 年 & 2024 年 & 2025 年 & 2025 年半年报 & 2026 年半年报 \\
\midrule
营业收入（亿元） & 11.58 & 11.56 & 11.61 & 7.18 & 8.68 \\
归母净利润（亿元） & 0.41 & 0.50 & 0.47 & 0.54 & 1.45 \\
毛利率（\%） & 27.8 & 32.2 & 36.6 & 40.6 & 48.6 \\
ROE（\%） & 7.1 & 8.0 & 7.0 & 8.0 & 18.7 \\
资产负债率（\%） & 59.7 & 58.4 & 57.3 & 44.9 & 37.5 \\
经营活动现金流净额（亿元） & 2.43 & 2.77 & 4.22 & 0.07 & 0.20 \\
货币资金（亿元） & 7.84 & 8.39 & 10.20 & 7.05 & 8.05 \\
有息负债合计（亿元） & 4.65 & 3.14 & 1.94 & 2.66 & 0.65 \\
现金短债比（倍） & 1.87 & 3.08 & 6.34 & 3.16 & 24.20 \\
\bottomrule
\end{tabular}
\end{adjustbox}
\end{center}

\begin{center}
\begin{minipage}{\textwidth}\centering
\begin{tikzpicture}
\begin{groupplot}[
  group style={group size=2 by 1, horizontal sep=0.60cm},
  width=6.85cm, height=4.60cm,
  tick label style={font=\tiny},
  title style={font=\scriptsize\heiti},
  yticklabel style={font=\tiny},
  ylabel style={font=\tiny},
  grid=major, grid style={LightGray, line width=0.3pt},
  axis line style={InkGray, line width=0.5pt},
  enlarge x limits={abs=0.8},
  legend style={font=\scriptsize, at={(0.5,-0.20)}, anchor=north,
                legend columns=2, draw=none, fill=none, column sep=6pt},
]
\nextgroupplot[
  ybar, bar width=4pt,
  ymin=-1.16, ymax=13.00,
  xtick={0,1,2,3,4,5.7},
  xticklabels={2021,2022,2023,2024,2025,2026H1},
  ylabel={亿元},
]
\addplot[fill=AgriGreen!75, draw=AgriGreen, bar shift=-2.6pt] table[x=i, y=rev]{data/clean/profiles/fin_600354.dat};
\addplot[fill=SoftRed!60, draw=SoftRed, bar shift=2.6pt] table[x=i, y=ni]{data/clean/profiles/fin_600354.dat};
\addplot[fill=AgriGreen!30, draw=AgriGreen, pattern=north east lines, bar shift=-2.6pt] table[x=x, y=rev]{data/clean/profiles/finh_600354.dat};
\addplot[fill=SoftRed!20, draw=SoftRed, pattern=north east lines, bar shift=2.6pt] table[x=x, y=ni]{data/clean/profiles/finh_600354.dat};
\legend{营业收入（年/半年）, 归母净利润（年/半年）}
\nextgroupplot[
  ymin=-5.0, ymax=68.8,
  xtick={0,1,2,3,4,5.7},
  xticklabels={2021,2022,2023,2024,2025,2026H1},
  ylabel={\%},
]
\addplot[AgriGold, mark=*, mark size=1.3pt] table[x=i, y=roe]{data/clean/profiles/fin_600354.dat};
\addplot[OilBlue, mark=square*, mark size=1.3pt] table[x=i, y=debt]{data/clean/profiles/fin_600354.dat};
\addplot[only marks, AgriGold, mark=*, mark size=2.0pt] table[x=x, y=roe]{data/clean/profiles/finh_600354.dat};
\addplot[only marks, OilBlue, mark=square*, mark size=2.0pt] table[x=x, y=debt]{data/clean/profiles/finh_600354.dat};
\legend{ROE, 资产负债率}
\end{groupplot}
\end{tikzpicture}
\begin{tikzpicture}
\begin{axis}[
  name=finbot,
  width=13.75cm, height=3.10cm, scale only axis,
  ymin=-1.02, ymax=11.43,
  xtick={0,1,2,3,4,5.7},
  xticklabels={2021,2022,2023,2024,2025,2026H1},
  tick label style={font=\tiny},
  yticklabel style={font=\tiny},
  ylabel={亿元}, ylabel style={font=\tiny},
  ybar, bar width=4pt,
  grid=major, grid style={LightGray, line width=0.3pt},
  axis line style={InkGray, line width=0.5pt},
  enlarge x limits={abs=0.8},
  legend style={font=\scriptsize, at={(0.5,-0.26)}, anchor=north,
                legend columns=2, draw=none, fill=none, column sep=10pt},
]
\addplot[fill=AgriGreen!55, draw=AgriGreen, bar shift=-3.0pt] table[x=i, y=cash]{data/clean/profiles/fin_600354.dat};
\addplot[fill=SoftRed!45, draw=SoftRed, bar shift=3.0pt] table[x=i, y=ibd]{data/clean/profiles/fin_600354.dat};
\addplot[fill=AgriGreen!25, draw=AgriGreen, pattern=north east lines, bar shift=-3.0pt] table[x=x, y=cash]{data/clean/profiles/finh_600354.dat};
\addplot[fill=SoftRed!20, draw=SoftRed, pattern=north east lines, bar shift=3.0pt] table[x=x, y=ibd]{data/clean/profiles/finh_600354.dat};
\legend{货币资金（左轴）, 有息负债（左轴）}
\end{axis}
\begin{axis}[
  at={(finbot.south east)}, anchor=south east,
  width=13.75cm, height=3.10cm, scale only axis,
  axis y line*=right, axis x line*=none, xtick=\empty,
  ymin=-1.10, ymax=5.49,
  tick label style={font=\tiny},
  yticklabel style={font=\tiny}, scaled y ticks=false,
  legend style={font=\scriptsize, at={(0.5,-0.44)}, anchor=north,
                legend columns=1, draw=none, fill=none},
]
\addplot[OilBlue, mark=*, mark size=2.2pt, line width=1.3pt] table[x=i, y=ocf]{data/clean/profiles/fin_600354.dat};
\addplot[only marks, OilBlue, mark=*, mark size=3.2pt] table[x=x, y=ocf]{data/clean/profiles/finh_600354.dat};
\legend{经营活动现金流净额（右轴，亿元，蓝点）}
\end{axis}
\end{tikzpicture}

\captionof{figure}[敦煌种业（600354）近年主要财务指标]{敦煌种业（600354）近年主要财务指标（左上：营业收入与归母净利润；右上：ROE 与资产负债率；下：货币资金、有息负债为左轴柱状，经营活动现金流净额为其右轴的蓝点折线。
2021---2025 年为年报口径，2026H1 为 2026 年半年报/半年末，与其前一年报之间留出间隔、以斜纹或空心点区分；半年报数据未年化）}
\end{minipage}
\end{center}

\pfl{财务分析} 2026 年 6 月末资产负债率 37.5\%，较 2025 年末下降 19.9 个百分点（样本中位数 52.2\%，所处三级行业内按由低到高排第\zhnumber{5}位）。 2026 年 6 月末货币资金 8.05 亿元、短期债务（短期借款与一年内到期非流动负债）0.33 亿元、长期债务 0.32 亿元，有息负债合计 0.65 亿元，现金短债比 24.20 倍——覆盖充分。 净有息负债 -7.39 亿元，相当于其 2026 年上半年营业收入的 -85\%。 流动比率 2.70、速动比率 2.14，短期偿债能力处于正常区间。 2026 年上半年经营活动现金流净额 0.20 亿元，经营现金流与利润同向为正，内生资金可以支撑运营。 综合看，现金短债比 24.20 倍，短期偿付压力可控；资产负债率低于样本中位数 14.7 个百分点；有息负债中短期债务占比更高，期限结构仍以滚动续作为主。

\pfl{重大战略转型} 公司的战略动作可归为以下几类：其他动作（3 项）：2025 年 4 月，与酒钢集团财务有限公司签订金融服务协议，由酒钢财务公司提供存款、贷款及其他综合信贷、结算等服务；2024 年 9 月，预计截至 2024 年末与关联方酒钢集团及其所属公司发生日常关联交易不超过 2,910.00 万元；融资与资本运作（2 项）：2026 年，第一次临时股东会通过：为酒泉敦煌种业百佳食品有限公司贷款担保 6,000.00 万元、为瓜州敦种棉业有限公司贷款担保 5,000.00 万元、为甘肃省敦煌种业集团金昌金从玉农业科技有限公司贷款 1,000.00 万元提供反担保；2026 年 4 月，公司担保总额不高于 1.20 亿元。 公告口径上，近三年（2023 年 9 月---2026 年 9 月）公司披露重大重组与并购类公告 7 条、对外投资与转型类公告 0 条、回购与股权激励类公告 1 条，合计公告 281 条。 第一大行业“种子”的收入占比由 2021 年年报的 65.7\% 变为 77.1\%（11.5 个百分点），收入结构出现再平衡。 综合判断，报告期内公司有 7 条与战略相关的公告落地（重大重组与并购 7 条、对外投资与转型 0 条），资本运作与主业调整之间存在直接关联。

\begin{center}\small\setlength{\tabcolsep}{4pt}
\captionof{table}[敦煌种业（600354）历史融资与资本运作]{敦煌种业（600354）历史融资与资本运作（据公司公告与公开报道整理；金额为披露值，含首次公开发行、再融资、债券、银行借款与重整投资等）}
\begin{adjustbox}{max width=\textwidth}
\begin{tabular}{@{}p{1.45cm}p{2.55cm}p{4.05cm}p{4.95cm}@{}}
\toprule
时间 & 方式 & 规模 & 用途与说明 \\
\midrule
2003-12 & IPO（上交所，二级市场投资者定价市值配售） & 发行 7,\allowbreak{}500 万股，\allowbreak{}发行价 5.68 元/\allowbreak{}股，\allowbreak{}募集资金 42,\allowbreak{}600 万元 & 上市募集资金用于公司主业项目建设；发行数量占发行后总股本 40.33\% \\
2011-02-11 & 非公开发行（向5名特定投资者） & 1,\allowbreak{}758 万股，\allowbreak{}发行价 25.00 元/\allowbreak{}股，\allowbreak{}募集资金净额 4.10 亿元 & 玉米种子烘干生产线建设项目 4,\allowbreak{}900.00 万元、\allowbreak{}购买新疆玛纳斯油脂资产 5,\allowbreak{}100.00 万元、\allowbreak{}新疆玛纳斯 1,\allowbreak{}000 吨/\allowbreak{}日棉籽生产线技术改造 5,\allowbreak{}000.00 万元、\allowbreak{}番茄深加工项目 1.30 亿元、\allowbreak{}脱水蔬菜深加工生产线 1.30 亿元；募集资金已于 2011 年 1 月 26 日存放于专户；后续多次变更募投用途（含购买武威农科院玉米杂交种吉祥 1 号生产经营权 2,\allowbreak{}680.00 万元、\allowbreak{}武威设立全资子公司并建玉米种子生产线） \\
2014-10 & 非公开发行预案 & 拟向特定对象非公开发行 8,\allowbreak{}000 万股，\allowbreak{}发行价格 6.00 元/\allowbreak{}股 & 全部用于补充流动资金；该方案后续实施为 2015 年 10 月的 8,\allowbreak{}000 万股非公开发行 \\
2015-01-13 & 中期票据 & 15 敦煌种业 MTN001，\allowbreak{}发行额度 3.00 亿元 & 偿还银行贷款、\allowbreak{}补充流动资金；报告期内已按募集资金要求使用 \\
2015-10-14 & 非公开发行（向特定投资者） & 8,\allowbreak{}000 万股 & 所募集资金全部用于采购种子等原材料；完成后注册资本变更为 5.28 亿元 \\
2025-04 & 银行综合授信额度（年度议案） & 拟向银行申请不超过人民币 10.60 亿元综合授信额度 & 流动资金贷款等综合授信品种，\allowbreak{}担保总额不高于 1.20 亿元 \\
2026-04-28 & 银行综合授信额度（年度议案） & 拟向银行及各金融机构申请不超过人民币 12.00 亿元的综合授信额度 & 流动资金贷款、\allowbreak{}项目并购贷款、\allowbreak{}中长期贷款、\allowbreak{}银行承兑汇票、\allowbreak{}商业承兑汇票、\allowbreak{}保函、\allowbreak{}信用证、\allowbreak{}抵押贷款、\allowbreak{}担保贷款、\allowbreak{}融资租赁、\allowbreak{}网络供应链融资、\allowbreak{}商业信用等；期限自 2025 年年度股东会审议通过之日起至 2026 年年度股东会召开之日止；授信额度可循环使用 \\
\bottomrule
\end{tabular}
\end{adjustbox}
\end{center}

\pfl{历史融投资情况} 从现金流量表看，2025 年公司取得借款收到的现金 2.64 亿元、偿还债务支付的现金 3.73 亿元，筹资活动现金流净额 -1.88 亿元（2023 年为取得借款 6.28 亿元、偿还债务 6.97 亿元）。 2026 年上半年筹资活动现金流净额为 -2.27 亿元（取得借款 1.31 亿元、偿还债务 2.54 亿元，未年化），已转为净流出，处于净还债状态。 2025 年分配股利、利润或偿付利息支付的现金合计 0.73 亿元。 公开披露的融资记录共 7 笔，工具类型以 IPO、定向增发、中期票据、银行授信为主；金额与用途见下表。 其中规模最大的两笔为：2026 年 4 月，银行授信，拟向银行及各金融机构申请不超过人民币 12.00 亿元的综合授信额度；2025 年 4 月，银行授信，拟向银行申请不超过人民币 10.60 亿元综合授信额度。 股本方面，近三年披露股本变动 3 次；总股本自 2021 年末以来未发生变化（5.28 亿股）。 分红方面，近三年未实施现金分红，在全部样本中属于分红能力偏弱的一端。 近三年“再融资”类公告 1 条，涉及定向增发。 综合看，公司融资结构上股权与债务融资并举；2025 年筹资活动现金流净额 -1.88 亿元，融资节奏仍以维持存量债务周转为主，与前述经营与投资现金流的方向基本一致。

\pfl{未来潜在融资需求分析} 2026 年 6 月末货币资金 8.05 亿元，而短期债务 0.33 亿元（现金短债比 24.20 倍），2026 年上半年经营现金流净额 0.20 亿元。账面货币资金可以覆盖短期债务，缺口不体现在静态偿付层面。 按第 \ref{sec:financing-need} 节的三档划分，公司属于第三档“扩张型”：2026 年上半年 ROE 18.7\%（未年化，折合约 37.4\%）、6 月末资产负债率 37.5\%，资产负债表仍具备加杠杆空间；融资用途以产能扩张、品类并购与海外布局为主，单笔规模通常在 5---20 亿元量级，会被市场定价为成长而非补血。 公开披露的后续资金安排线索包括：债务与授信安排：2026 年度，银行综合授信不超过 12.00 亿元、为子公司担保总额不高于 1.20 亿元继续以银行信贷为主要外部资金来源；资本开支安排：2025 年度，政府补助相关递延收益转入 360.85 万元（与资产相关）、商务局外经贸发展专项资金 309.35 万元、核心种源繁育基地建设项目补助资金 120.00 万元。 以现有货币资金与经营现金流衡量，公司不存在刚性补血的紧迫性，融资动作更多是主动型的。




\subsubsection{众兴菌业（002772）}
\label{co:002772}

\pfl{公司基本情况} 甘肃天水的工厂化食用菌龙头，2015 年深交所上市，以金针菇、双孢菇为核心品种（双孢菇满产后日产能达 420 吨/天）；2025 年下半年起行业企稳回暖，价格修复带来利润大幅弹性，同时以技改方式布局冬虫夏草工厂化生态繁育，推进多品种协同与西南总部基地建设。 公司于 2015-06-26 在深交所主板首发上市，注册资本 3.75 亿元，总股本 3.93 亿股。 截至 2026 年 10 月 12 日收盘价 10.53 元，对应总市值 41.4 亿元，近一年+0.5\%，总市值在三级行业“食用菌”的 4 家公司中按由大到小排第\zhnumber{1}位，近一年涨跌幅高于全样本中位数（-10.4\%）10.8 个百分点。 按 2026 年半年报披露的行业构成，收入占比前两位为农业种植业 99.6\%；其他(补充) 0.4\%。

\pfl{股权结构与控制权} 按 2026 年 6 月 30 日的股东名册，前三大股东为陶军 持股 29.39\%；田德 持股 8.31\%；天水众兴菌业科技股份有限公司－第一期员工持股计划 持股 2.40\%，股东户数 23,472 户。 股东结构与控制权方面，截至 2026 年 6 月末二人累计质押 2,889.00 万股；2026 年，第一次临时股东会审议通过将回购专用证券账户中的 1,849.88 万股用途变更为注销以减少注册资本。

\pfl{经营要点} 经营与产能方面，报告期内投资设立河北众兴、汉中众兴实施冬虫夏草工厂化生态繁育项目。

\pfl{所处行业的基本面} 食用菌的供给是工厂化产能（投放周期 6---9 个月），需求是刚性消费，价格随产能快速调整（第 \ref{sec:framework} 章的供需框架）。 行业层面的最新情况是：行业前三市占率超 40\%；2024 年，全国批发市场鲜金针菇年均价跌至 5.93 元/公斤、同比下跌 9.88\%；2026 年上半年，众兴菌业、华绿生物、雪榕生物、万辰集团四家食用菌上市企业业绩集体回暖。

\pfl{竞争优势} 把该公司放回同类公司比较，2026 年上半年营业收入 10.9 亿元、在三级行业“食用菌”的 4 家公司中排第\zhnumber{2}位；毛利率 30.9\%、高于行业内中位数 14.6 个百分点；半年 ROE 5.3\%（未年化）、高于行业内中位数（2.15\%）；资产负债率 54.0\%、高于行业内中位数（44.7\%）。 相对优势方面，食用菌种植所需农业下脚料资源丰富、原材料成本优势明显；2023 年，公司双孢菇日产能 360 吨、占全国工厂化总产能约 39\% 至 40\%。 优势在报表上主要体现为高于行业平均的毛利率，这在周期底部是抵御亏损的主要缓冲。

\pfl{历史股价} 自 2021 年 1 月至 2026 年 10 月，股价由 8.46 元变为 10.53 元，累计 +24.5\%，区间最高 18.37 元（2026 年 2 月）、最低 5.88 元（2024 年 7 月）；当前价相当于区间最高价的 57.3\%，为区间最低价的 1.79 倍。 月度收益的年化波动率 49.8\%，高于全样本中位数 37.3\%，在三级行业“食用菌”的 4 家公司中波动率由高到低排第\zhnumber{2}位。

\begin{center}
\begin{minipage}{\textwidth}\centering
\begin{tikzpicture}
\begin{axis}[
  width=12.6cm, height=3.9cm, scale only axis,
  xmin=2020.922, xmax=2026.972, ymin=5.47, ymax=19.66,
  xtick={2021,2022,2023,2024,2025,2026},
  yearticks,
  yticklabel style={font=\scriptsize},
  ylabel={元/股}, ylabel style={font=\scriptsize},
  grid=major, grid style={LightGray, line width=0.3pt},
  axis line style={InkGray, line width=0.5pt},
  every axis plot/.append style={line width=0.9pt},
]
\addplot[AgriGreen] table[x=x,y=y]{data/clean/profiles/px_002772.dat};
\addplot[SoftRed, dashed, line width=0.7pt] table[x=x,y=y]{data/clean/profiles/pxzz_002772.dat};
\addplot[only marks, mark=*, mark size=1.1pt, SoftRed] table[x=x,y=y]{data/clean/profiles/pxzz_002772.dat};
\end{axis}
\end{tikzpicture}

\captionof{figure}[众兴菌业（002772）股价与周期骨架]{众兴菌业（002772）月度收盘价（元/股）与 ZigZag 周期骨架（阈值 25\%，圆点为周期转折点）}
\end{minipage}
\end{center}

\pfl{过去周期分析} 以 25\% 的 ZigZag 阈值划分，样本区间内股价形成 3 个主升段与 3 个主跌段。 最大主升段出现在 2024 年 7 月 至 2026 年 2 月，19 个月+212.4\%；最大主跌段为 2021 年 6 月 至 2022 年 4 月，10 个月-59.2\%。 最近一次转折出现在 2026 年 7 月（下跌段自 2026 年 2 月起，5 个月-44.3\%），当前处于该段的延续之中。 公司股价较申万农林牧渔指数提前约 22 个月见底，是板块中较早定价周期反转的公司之一。 食用菌是工厂化生产的制造型农业（第 \ref{sec:soft} 章），产能投放节奏与金针菇等单品价格是盈利的主要变量，与养殖周期基本无关，公司 2026 年上半年实现盈利、半年 ROE 5.3\%（未年化），好于全样本中位数（0.75\%），下行周期对其报表的冲击小于同业。 综合区间形态看，该股单段涨跌幅度偏大、以趋势段为主（最大主升 +212.4\%、最大主跌 -59.2\%），年化波动率 49.8\% 与全样本中位数 37.3\% 相比波动明显大于样本中位数。

\pfl{盈利情况} 2026 年上半年营业收入 10.85 亿元（同比 +19.6\%），归母净利润 1.91 亿元，同比 +176.5\%。 以年报为趋势对照，2023---2025 年营业收入由 19.31 亿元变为 20.98 亿元（两年年均复合增速 +4.2\%），归母净利润由 1.60 亿元变为 3.35 亿元。 按半年报口径，2026 年上半年的 ROE 为 5.3\%（未年化）、销售净利率 17.6\%、毛利率 30.9\%，ROE 在“食用菌”4 家公司中按数值由高到低排第\zhnumber{1}位。 按同一口径，三级行业“食用菌”4 家公司的半年 ROE 中位数为 2.15\%，该公司高于其中位数 3.19 个百分点。 综合看，公司在报告期维持盈利（高于全样本半年 ROE 中位数 0.75\%），利润的可持续性取决于食用菌景气的延续与成本管控的执行。

\begin{center}\small\setlength{\tabcolsep}{3.4pt}
\captionof{table}[众兴菌业（002772）关键财务指标]{众兴菌业（002772）关键财务指标（合并报表口径；两个半年报期均未年化；现金短债比 $=$ 货币资金 $\div$ 短期债务）}
\begin{adjustbox}{max width=\textwidth}
\begin{tabular}{@{}lrrrrr@{}}
\toprule
指标 & 2023 年 & 2024 年 & 2025 年 & 2025 年半年报 & 2026 年半年报 \\
\midrule
营业收入（亿元） & 19.31 & 19.35 & 20.98 & 9.08 & 10.85 \\
归母净利润（亿元） & 1.60 & 1.28 & 3.35 & 0.69 & 1.91 \\
毛利率（\%） & 23.0 & 20.8 & 28.6 & 19.4 & 30.9 \\
ROE（\%） & 4.8 & 3.8 & 9.9 & 2.1 & 5.3 \\
资产负债率（\%） & 43.4 & 46.4 & 50.1 & 49.7 & 54.0 \\
经营活动现金流净额（亿元） & 5.02 & 4.95 & 4.96 & 1.50 & 1.72 \\
货币资金（亿元） & 10.39 & 11.91 & 14.50 & 12.86 & 16.69 \\
有息负债合计（亿元） & 20.84 & 23.70 & 28.90 & 26.90 & 35.65 \\
现金短债比（倍） & 1.03 & 0.91 & 1.57 & 1.17 & 1.35 \\
\bottomrule
\end{tabular}
\end{adjustbox}
\end{center}

\begin{center}
\begin{minipage}{\textwidth}\centering
\begin{tikzpicture}
\begin{groupplot}[
  group style={group size=2 by 1, horizontal sep=0.60cm},
  width=6.85cm, height=4.60cm,
  tick label style={font=\tiny},
  title style={font=\scriptsize\heiti},
  yticklabel style={font=\tiny},
  ylabel style={font=\tiny},
  grid=major, grid style={LightGray, line width=0.3pt},
  axis line style={InkGray, line width=0.5pt},
  enlarge x limits={abs=0.8},
  legend style={font=\scriptsize, at={(0.5,-0.20)}, anchor=north,
                legend columns=2, draw=none, fill=none, column sep=6pt},
]
\nextgroupplot[
  ybar, bar width=4pt,
  ymin=-2.10, ymax=23.50,
  xtick={0,1,2,3,4,5.7},
  xticklabels={2021,2022,2023,2024,2025,2026H1},
  ylabel={亿元},
]
\addplot[fill=AgriGreen!75, draw=AgriGreen, bar shift=-2.6pt] table[x=i, y=rev]{data/clean/profiles/fin_002772.dat};
\addplot[fill=SoftRed!60, draw=SoftRed, bar shift=2.6pt] table[x=i, y=ni]{data/clean/profiles/fin_002772.dat};
\addplot[fill=AgriGreen!30, draw=AgriGreen, pattern=north east lines, bar shift=-2.6pt] table[x=x, y=rev]{data/clean/profiles/finh_002772.dat};
\addplot[fill=SoftRed!20, draw=SoftRed, pattern=north east lines, bar shift=2.6pt] table[x=x, y=ni]{data/clean/profiles/finh_002772.dat};
\legend{营业收入（年/半年）, 归母净利润（年/半年）}
\nextgroupplot[
  ymin=-4.3, ymax=59.4,
  xtick={0,1,2,3,4,5.7},
  xticklabels={2021,2022,2023,2024,2025,2026H1},
  ylabel={\%},
]
\addplot[AgriGold, mark=*, mark size=1.3pt] table[x=i, y=roe]{data/clean/profiles/fin_002772.dat};
\addplot[OilBlue, mark=square*, mark size=1.3pt] table[x=i, y=debt]{data/clean/profiles/fin_002772.dat};
\addplot[only marks, AgriGold, mark=*, mark size=2.0pt] table[x=x, y=roe]{data/clean/profiles/finh_002772.dat};
\addplot[only marks, OilBlue, mark=square*, mark size=2.0pt] table[x=x, y=debt]{data/clean/profiles/finh_002772.dat};
\legend{ROE, 资产负债率}
\end{groupplot}
\end{tikzpicture}
\begin{tikzpicture}
\begin{axis}[
  name=finbot,
  width=13.75cm, height=3.10cm, scale only axis,
  ymin=-3.57, ymax=39.93,
  xtick={0,1,2,3,4,5.7},
  xticklabels={2021,2022,2023,2024,2025,2026H1},
  tick label style={font=\tiny},
  yticklabel style={font=\tiny},
  ylabel={亿元}, ylabel style={font=\tiny},
  ybar, bar width=4pt,
  grid=major, grid style={LightGray, line width=0.3pt},
  axis line style={InkGray, line width=0.5pt},
  enlarge x limits={abs=0.8},
  legend style={font=\scriptsize, at={(0.5,-0.26)}, anchor=north,
                legend columns=2, draw=none, fill=none, column sep=10pt},
]
\addplot[fill=AgriGreen!55, draw=AgriGreen, bar shift=-3.0pt] table[x=i, y=cash]{data/clean/profiles/fin_002772.dat};
\addplot[fill=SoftRed!45, draw=SoftRed, bar shift=3.0pt] table[x=i, y=ibd]{data/clean/profiles/fin_002772.dat};
\addplot[fill=AgriGreen!25, draw=AgriGreen, pattern=north east lines, bar shift=-3.0pt] table[x=x, y=cash]{data/clean/profiles/finh_002772.dat};
\addplot[fill=SoftRed!20, draw=SoftRed, pattern=north east lines, bar shift=3.0pt] table[x=x, y=ibd]{data/clean/profiles/finh_002772.dat};
\legend{货币资金（左轴）, 有息负债（左轴）}
\end{axis}
\begin{axis}[
  at={(finbot.south east)}, anchor=south east,
  width=13.75cm, height=3.10cm, scale only axis,
  axis y line*=right, axis x line*=none, xtick=\empty,
  ymin=-6.52, ymax=7.77,
  tick label style={font=\tiny},
  yticklabel style={font=\tiny}, scaled y ticks=false,
  legend style={font=\scriptsize, at={(0.5,-0.44)}, anchor=north,
                legend columns=1, draw=none, fill=none},
]
\addplot[OilBlue, mark=*, mark size=2.2pt, line width=1.3pt] table[x=i, y=ocf]{data/clean/profiles/fin_002772.dat};
\addplot[only marks, OilBlue, mark=*, mark size=3.2pt] table[x=x, y=ocf]{data/clean/profiles/finh_002772.dat};
\legend{经营活动现金流净额（右轴，亿元，蓝点）}
\end{axis}
\end{tikzpicture}

\captionof{figure}[众兴菌业（002772）近年主要财务指标]{众兴菌业（002772）近年主要财务指标（左上：营业收入与归母净利润；右上：ROE 与资产负债率；下：货币资金、有息负债为左轴柱状，经营活动现金流净额为其右轴的蓝点折线。
2021---2025 年为年报口径，2026H1 为 2026 年半年报/半年末，与其前一年报之间留出间隔、以斜纹或空心点区分；半年报数据未年化）}
\end{minipage}
\end{center}

\pfl{财务分析} 2026 年 6 月末资产负债率 54.0\%，较 2025 年末上升 3.8 个百分点（样本中位数 52.2\%，所处三级行业内按由低到高排第\zhnumber{3}位）。 2026 年 6 月末货币资金 16.69 亿元、短期债务（短期借款与一年内到期非流动负债）12.38 亿元、长期债务 23.27 亿元，有息负债合计 35.65 亿元，现金短债比 1.35 倍——覆盖充分。 净有息负债 18.96 亿元，相当于其 2026 年上半年营业收入的 175\%。 流动比率 2.24、速动比率 1.24，短期指标尚可。 2026 年上半年经营活动现金流净额 1.72 亿元，经营现金流与利润同向为正，内生资金可以支撑运营。 2026 年 6 月末在建工程 5.95 亿元，较上年末增加 1.91 亿元，仍有在建投入。 综合看，现金短债比 1.35 倍，短期偿付压力可控；资产负债率高于样本中位数 1.8 个百分点；有息负债中长期债务占比更高，期限结构对短债滚动的压力较小。

\pfl{重大战略转型} 近三年公司的战略动作集中在：产能与项目（4 项）：2025 年 6 月，第五届董事会第十二次会议通过拟签订项目投资协议暨对外投资；2025 年 2 月，全资子公司陕西众兴高科生物科技对原金针菇生产线停产改造；融资与资本运作（1 项）：2020 年 4 月，非公开发行募投项目年产 7,500 吨蟹味菇生产线、年产 7,500 吨白玉菇生产线募集资金及利息全部变更投向全资子公司安徽众兴菌业科技年产 2 万吨双孢蘑菇及 11 万吨堆肥工厂化生产项目。 公告口径上，近三年（2023 年 9 月---2026 年 9 月）公司披露重大重组与并购类公告 0 条、对外投资与转型类公告 14 条、回购与股权激励类公告 23 条，合计公告 443 条。 第一大行业仍是“农业种植业”，收入占比 99.6\%，与 2021 年年报相比没有方向性变化。 综合判断，报告期内公司有 14 条与战略相关的公告落地（重大重组与并购 0 条、对外投资与转型 14 条），资本运作与主业调整之间存在直接关联。

\begin{center}\small\setlength{\tabcolsep}{4pt}
\captionof{table}[众兴菌业（002772）历史融资与资本运作]{众兴菌业（002772）历史融资与资本运作（据公司公告与公开报道整理；金额为披露值，含首次公开发行、再融资、债券、银行借款与重整投资等）}
\begin{adjustbox}{max width=\textwidth}
\begin{tabular}{@{}p{1.45cm}p{2.55cm}p{4.05cm}p{4.95cm}@{}}
\toprule
时间 & 方式 & 规模 & 用途与说明 \\
\midrule
2015-06 & IPO（深交所） & 发行 3,\allowbreak{}725 万股，\allowbreak{}发行价 13.00 元/\allowbreak{}股，\allowbreak{}募集资金总额 4.84 亿元 & 年产 12,\allowbreak{}600 吨金针菇生产线建设项目、\allowbreak{}年产 7,\allowbreak{}200 吨杏鲍菇生产线建设项目、\allowbreak{}食（药）用菌研发中心建设项目、\allowbreak{}偿还银行贷款；发行价对应发行市盈率 22.98 倍（按询价确定价格除以发行后每股收益计算） \\
2016-03-18 & 限制性股票（股权激励） & 筹资净额 9,\allowbreak{}361.94 万元 & 股权激励认购；公司历次筹资记录中的一项 \\
2016-08-22 & 非公开发行（向特定投资者） & 发行 5,\allowbreak{}368.33 万股，\allowbreak{}每股 21.00 元，\allowbreak{}募集资金总额 11.27 亿元 & 非公开发行募投项目（后部分变更）；截至 2020 年 12 月 31 日，\allowbreak{}非公开发行募集资金本息合计 11.76 亿元 \\
2017-12-13 & 可转换公司债券（众兴转债，128026） & 发行 920 万张、\allowbreak{}每张面值 100.00 元，\allowbreak{}募集资金 9.20 亿元 & 年产 2 万吨双孢蘑菇及 11 万吨堆肥工厂化生产项目、\allowbreak{}年产 32,\allowbreak{}400 吨金针菇生产线建设项目、\allowbreak{}年产 32,\allowbreak{}400 吨金针菇工厂化循环经济产业链建设项目；期限 6 年（2017-12-13 至 2023-12-13）；2018 年第二至四季度合计转股 1.69 万股 \\
2025-02-07 & 股份回购（回购专用证券账户集中竞价） & 累计回购 1,\allowbreak{}849.88 万股，\allowbreak{}占公司总股本 4.7042\% & 股份回购（公司曾披露取得金融机构回购专项贷款承诺函）；回购专户股份不参与 2025 年度利润分配基数 \\
\bottomrule
\end{tabular}
\end{adjustbox}
\end{center}

\pfl{历史融投资情况} 从现金流量表看，2025 年公司取得借款收到的现金 18.26 亿元、偿还债务支付的现金 13.54 亿元、吸收权益性投资 0.07 亿元，筹资活动现金流净额 2.67 亿元（2023 年为取得借款 8.15 亿元、偿还债务 13.59 亿元）。 2026 年上半年筹资活动现金流净额为 3.34 亿元（取得借款 10.96 亿元、偿还债务 5.37 亿元，未年化），仍处于净融入状态。 2025 年分配股利、利润或偿付利息支付的现金合计 1.87 亿元。 公开披露的融资记录共 5 笔，工具类型以 IPO、定向增发、可转债为主；金额与用途见下表。 其中规模最大的两笔为：2017 年 12 月，可转债，发行 920 万张、每张面值 100.00 元，募集资金 9.20 亿元；2016 年 8 月，定向增发，发行 5,368.33 万股，每股 21.00 元，募集资金总额 11.27 亿元。 股本方面，近三年披露股本变动 7 次（可转债转股 5 次、股份回购 1 次）；总股本由 2021 年末的 4.07 亿股变为 3.93 亿股（-3.48\%）。 分红方面，近三年现金分红 6 次、合计每股派现 1.50 元，最近一次实施在 2026 年 6 月。 近三年“再融资”类公告 1 条，涉及可转债。 综合看，公司融资结构上以债务滚动为主、股权融资为补充；2025 年筹资活动现金流净额 2.67 亿元，年度内仍为净融入，与前述经营与投资现金流的方向基本一致。

\pfl{未来潜在融资需求分析} 2026 年 6 月末货币资金 16.69 亿元，而短期债务 12.38 亿元（现金短债比 1.35 倍），2026 年上半年经营现金流净额 1.72 亿元。账面货币资金可以覆盖短期债务，缺口不体现在静态偿付层面。 按第 \ref{sec:financing-need} 节的三档划分，公司属于第三档“扩张型”：2026 年上半年 ROE 5.3\%（未年化，折合约 10.7\%）、6 月末资产负债率 54.0\%，资产负债表仍具备加杠杆空间；融资用途以产能扩张、品类并购与海外布局为主，单笔规模通常在 5---20 亿元量级，会被市场定价为成长而非补血。 公开披露的后续资金安排线索包括：债务与授信安排：2025 年 2 月，股份回购方案资金安排曾取得金融机构回购专项贷款承诺函；2026 年度，公司为合并报表范围内全资子公司及全资孙公司提供担保的预计额度为 11.00 亿元，用于子公司日常经营及项目建设。 资金需求还要叠加在建工程：期末在建工程 5.95 亿元、2026 年上半年购建固定资产等支出 3.02 亿元（未年化），这部分支出在周期底部通常会被主动放缓。 综合看，公司在静态偿付层面不存在缺口，外部融资更多是配合扩张节奏与并购整合的主动性安排，而非偿债压力驱动。




\subsubsection{万向德农（600371）}
\label{co:600371}

\pfl{公司基本情况} 公司前身可追溯至 1995 年，黑龙江起家、2002 年上市的玉米种子企业（原华冠科技）；主营高度集中于玉米杂交种子，收入自 2025 年起大幅下滑，但资产负债率低、无银行借款、分红与增持持续。 公司于 2002-09-16 在上交所首发上市，注册资本 2.93 亿元，总股本 2.93 亿股。 截至 2026 年 10 月 12 日收盘价 14.03 元，对应总市值 41.0 亿元，近一年+53.8\%，总市值在三级行业“种子”的 8 家公司中按由大到小排第\zhnumber{7}位，近一年涨跌幅高于全样本中位数（-10.4\%）64.2 个百分点。 按 2020 年年报披露的行业构成，收入占比前两位为种子行业 97.4\%；其他(补充) 2.6\%。

\pfl{股权结构与控制权} 按 2026 年 6 月 30 日的股东名册，前三大股东为万向三农集团有限公司（持股比例未披露）；陈铭（持股比例未披露）；中国建设银行股份有限公司－诺安多策略混合型证券投资基 持股 0.44\%，股东户数 27,833 户。 控制权与股东性质方面，除万向三农外的前十大股东持股均在 0.73\% 以下。

\pfl{经营要点} 经营与产能方面，2026 年 6 月末短期借款为零。

\pfl{所处行业的基本面} 种子的供给由制种面积与库存决定、需求由粮价与种植意愿决定，制种面积在 2025 年起出现萎缩，是本轮供给端最先出清的环节（第 \ref{sec:framework} 章的供需框架）。 与该公司直接相关的行业数据包括：截至 2026 年 6 月底我国已累计审定转基因玉米品种 160 个，行业竞争格局从“渠道与价格战”向“技术与性状壁垒”重塑；全国农业技术推广服务中心预计 2026 年全国玉米制种面积 358.5 万亩、较上年减少 27 万亩、降幅 7\%。

\pfl{竞争优势} 把该公司放回同类公司比较，2026 年上半年营业收入 1.0 亿元、在三级行业“种子”的 8 家公司中排第\zhnumber{7}位；毛利率 26.5\%、低于行业内中位数 1.9 个百分点；半年 ROE 2.6\%（未年化）、高于行业内中位数（0.17\%）；资产负债率 14.5\%、低于行业内中位数（37.3\%）。 与同类公司相比，公司的差异化在于：公司 2026 年 6 月末总资产 6.66 亿元、归属上市公司股东净资产 5.20 亿元。 优势在当期体现为高于行业中位数的资本回报，但优势能否持续仍取决于周期与自身产能投放节奏。

\pfl{历史股价} 自 2021 年 1 月至 2026 年 10 月，股价由 12.32 元变为 14.03 元，累计 +13.9\%，区间最高 15.72 元（2022 年 5 月）、最低 6.25 元（2026 年 6 月）；当前价相当于区间最高价的 89.2\%，为区间最低价的 2.24 倍。 月度收益的年化波动率 42.8\%，高于全样本中位数 37.3\%，在三级行业“种子”的 8 家公司中波动率由高到低排第\zhnumber{4}位。

\begin{center}
\begin{minipage}{\textwidth}\centering
\begin{tikzpicture}
\begin{axis}[
  width=12.6cm, height=3.9cm, scale only axis,
  xmin=2020.922, xmax=2026.972, ymin=5.81, ymax=16.82,
  xtick={2021,2022,2023,2024,2025,2026},
  yearticks,
  yticklabel style={font=\scriptsize},
  ylabel={元/股}, ylabel style={font=\scriptsize},
  grid=major, grid style={LightGray, line width=0.3pt},
  axis line style={InkGray, line width=0.5pt},
  every axis plot/.append style={line width=0.9pt},
]
\addplot[AgriGreen] table[x=x,y=y]{data/clean/profiles/px_600371.dat};
\addplot[SoftRed, dashed, line width=0.7pt] table[x=x,y=y]{data/clean/profiles/pxzz_600371.dat};
\addplot[only marks, mark=*, mark size=1.1pt, SoftRed] table[x=x,y=y]{data/clean/profiles/pxzz_600371.dat};
\end{axis}
\end{tikzpicture}

\captionof{figure}[万向德农（600371）股价与周期骨架]{万向德农（600371）月度收盘价（元/股）与 ZigZag 周期骨架（阈值 25\%，圆点为周期转折点）}
\end{minipage}
\end{center}

\pfl{过去周期分析} 以 25\% 的 ZigZag 阈值划分，样本区间内股价形成 4 个主升段与 4 个主跌段。 最大主升段出现在 2026 年 6 月 至 2026 年 10 月，4 个月+124.5\%；最大主跌段为 2023 年 1 月 至 2024 年 8 月，19 个月-53.1\%。 最近一次转折出现在 2026 年 10 月（上涨段自 2026 年 6 月起，4 个月+124.5\%），当前处于该段之后的震荡之中。 公司股价与申万农林牧渔指数几乎同步见底，个股并未走出独立行情。 种子行业的价格由制种成本、粮价与政策（转基因商业化、种业振兴）共同决定（第 \ref{sec:planting} 节），与生猪价格的联动有限，公司 2026 年上半年实现盈利、半年 ROE 2.6\%（未年化），好于全样本中位数（0.75\%），下行周期对其报表的冲击小于同业。 综合区间形态看，该股单段涨跌幅度偏大、以趋势段为主（最大主升 +124.5\%、最大主跌 -53.1\%），年化波动率 42.8\% 与全样本中位数 37.3\% 相比波动与样本中位数接近。

\pfl{盈利情况} 2026 年上半年营业收入 0.99 亿元（同比 -15.1\%），归母净利润 0.14 亿元，同比 -44.5\%。 以年报为趋势对照，2023---2025 年营业收入由 3.19 亿元变为 1.98 亿元（两年年均复合增速 -21.3\%），归母净利润由 0.65 亿元变为 0.06 亿元。 按半年报口径，2026 年上半年的 ROE 为 2.6\%（未年化）、销售净利率 15.5\%、毛利率 26.5\%，ROE 在“种子”8 家公司中按数值由高到低排第\zhnumber{2}位。 按同一口径，三级行业“种子”8 家公司的半年 ROE 中位数为 0.17\%，该公司高于其中位数 2.45 个百分点。 从公司披露的原因看，业绩变动主要来自：2025 年，营业收入下降主因玉米种子销量、销售价格均下降，营业成本下降主因制种成本下降。 综合看，公司在报告期维持盈利（高于全样本半年 ROE 中位数 0.75\%），利润的可持续性取决于种子景气的延续与成本管控的执行。

\begin{center}\small\setlength{\tabcolsep}{3.4pt}
\captionof{table}[万向德农（600371）关键财务指标]{万向德农（600371）关键财务指标（合并报表口径；两个半年报期均未年化；现金短债比 $=$ 货币资金 $\div$ 短期债务）}
\begin{adjustbox}{max width=\textwidth}
\begin{tabular}{@{}lrrrrr@{}}
\toprule
指标 & 2023 年 & 2024 年 & 2025 年 & 2025 年半年报 & 2026 年半年报 \\
\midrule
营业收入（亿元） & 3.19 & 3.43 & 1.98 & 1.17 & 0.99 \\
归母净利润（亿元） & 0.65 & 0.53 & 0.06 & 0.25 & 0.14 \\
毛利率（\%） & 34.9 & 27.1 & 27.3 & 33.7 & 26.5 \\
ROE（\%） & 11.3 & 9.1 & 1.0 & 4.4 & 2.6 \\
资产负债率（\%） & 32.4 & 22.5 & 21.0 & 15.0 & 14.5 \\
经营活动现金流净额（亿元） & 0.55 & -0.73 & -0.38 & -0.30 & -0.16 \\
货币资金（亿元） & 4.74 & 3.47 & 2.69 & 2.86 & 2.26 \\
有息负债合计（亿元） & --- & --- & --- & --- & 0.00 \\
现金短债比（倍） & --- & --- & --- & --- & --- \\
\bottomrule
\end{tabular}
\end{adjustbox}
\end{center}

\begin{center}
\begin{minipage}{\textwidth}\centering
\begin{tikzpicture}
\begin{groupplot}[
  group style={group size=2 by 1, horizontal sep=0.60cm},
  width=6.85cm, height=4.60cm,
  tick label style={font=\tiny},
  title style={font=\scriptsize\heiti},
  yticklabel style={font=\tiny},
  ylabel style={font=\tiny},
  grid=major, grid style={LightGray, line width=0.3pt},
  axis line style={InkGray, line width=0.5pt},
  enlarge x limits={abs=0.8},
  legend style={font=\scriptsize, at={(0.5,-0.20)}, anchor=north,
                legend columns=2, draw=none, fill=none, column sep=6pt},
]
\nextgroupplot[
  ybar, bar width=4pt,
  ymin=-0.34, ymax=3.84,
  xtick={0,1,2,3,4,5.7},
  xticklabels={2021,2022,2023,2024,2025,2026H1},
  ylabel={亿元},
]
\addplot[fill=AgriGreen!75, draw=AgriGreen, bar shift=-2.6pt] table[x=i, y=rev]{data/clean/profiles/fin_600371.dat};
\addplot[fill=SoftRed!60, draw=SoftRed, bar shift=2.6pt] table[x=i, y=ni]{data/clean/profiles/fin_600371.dat};
\addplot[fill=AgriGreen!30, draw=AgriGreen, pattern=north east lines, bar shift=-2.6pt] table[x=x, y=rev]{data/clean/profiles/finh_600371.dat};
\addplot[fill=SoftRed!20, draw=SoftRed, pattern=north east lines, bar shift=2.6pt] table[x=x, y=ni]{data/clean/profiles/finh_600371.dat};
\legend{营业收入（年/半年）, 归母净利润（年/半年）}
\nextgroupplot[
  ymin=-2.6, ymax=35.6,
  xtick={0,1,2,3,4,5.7},
  xticklabels={2021,2022,2023,2024,2025,2026H1},
  ylabel={\%},
]
\addplot[AgriGold, mark=*, mark size=1.3pt] table[x=i, y=roe]{data/clean/profiles/fin_600371.dat};
\addplot[OilBlue, mark=square*, mark size=1.3pt] table[x=i, y=debt]{data/clean/profiles/fin_600371.dat};
\addplot[only marks, AgriGold, mark=*, mark size=2.0pt] table[x=x, y=roe]{data/clean/profiles/finh_600371.dat};
\addplot[only marks, OilBlue, mark=square*, mark size=2.0pt] table[x=x, y=debt]{data/clean/profiles/finh_600371.dat};
\legend{ROE, 资产负债率}
\end{groupplot}
\end{tikzpicture}
\begin{tikzpicture}
\begin{axis}[
  name=finbot,
  width=13.75cm, height=3.10cm, scale only axis,
  ymin=-0.47, ymax=5.31,
  xtick={0,1,2,3,4,5.7},
  xticklabels={2021,2022,2023,2024,2025,2026H1},
  tick label style={font=\tiny},
  yticklabel style={font=\tiny},
  ylabel={亿元}, ylabel style={font=\tiny},
  ybar, bar width=4pt,
  grid=major, grid style={LightGray, line width=0.3pt},
  axis line style={InkGray, line width=0.5pt},
  enlarge x limits={abs=0.8},
  legend style={font=\scriptsize, at={(0.5,-0.26)}, anchor=north,
                legend columns=2, draw=none, fill=none, column sep=10pt},
]
\addplot[fill=AgriGreen!55, draw=AgriGreen, bar shift=-3.0pt] table[x=i, y=cash]{data/clean/profiles/fin_600371.dat};
\addplot[fill=SoftRed!45, draw=SoftRed, bar shift=3.0pt] table[x=i, y=ibd]{data/clean/profiles/fin_600371.dat};
\addplot[fill=AgriGreen!25, draw=AgriGreen, pattern=north east lines, bar shift=-3.0pt] table[x=x, y=cash]{data/clean/profiles/finh_600371.dat};
\addplot[fill=SoftRed!20, draw=SoftRed, pattern=north east lines, bar shift=3.0pt] table[x=x, y=ibd]{data/clean/profiles/finh_600371.dat};
\legend{货币资金（左轴）, 有息负债（左轴）}
\end{axis}
\begin{axis}[
  at={(finbot.south east)}, anchor=south east,
  width=13.75cm, height=3.10cm, scale only axis,
  axis y line*=right, axis x line*=none, xtick=\empty,
  ymin=-1.16, ymax=1.41,
  tick label style={font=\tiny},
  yticklabel style={font=\tiny}, scaled y ticks=false,
  legend style={font=\scriptsize, at={(0.5,-0.44)}, anchor=north,
                legend columns=1, draw=none, fill=none},
]
\addplot[OilBlue, mark=*, mark size=2.2pt, line width=1.3pt] table[x=i, y=ocf]{data/clean/profiles/fin_600371.dat};
\addplot[only marks, OilBlue, mark=*, mark size=3.2pt] table[x=x, y=ocf]{data/clean/profiles/finh_600371.dat};
\legend{经营活动现金流净额（右轴，亿元，蓝点）}
\end{axis}
\end{tikzpicture}

\captionof{figure}[万向德农（600371）近年主要财务指标]{万向德农（600371）近年主要财务指标（左上：营业收入与归母净利润；右上：ROE 与资产负债率；下：货币资金、有息负债为左轴柱状，经营活动现金流净额为其右轴的蓝点折线。
2021---2025 年为年报口径，2026H1 为 2026 年半年报/半年末，与其前一年报之间留出间隔、以斜纹或空心点区分；半年报数据未年化）}
\end{minipage}
\end{center}

\pfl{财务分析} 2026 年 6 月末资产负债率 14.5\%，较 2025 年末下降 6.5 个百分点（样本中位数 52.2\%，所处三级行业内按由低到高排第\zhnumber{1}位）。 流动比率 4.96、速动比率 2.78，短期指标尚可。 2026 年上半年经营活动现金流净额 -0.16 亿元，净利润为正但经营现金流为负，利润的现金含量偏低。 综合看，资产负债率低于样本中位数 37.7 个百分点；有息负债中长期债务占比更高，期限结构对短债滚动的压力较小。

\pfl{重大战略转型} 近三年公司的战略动作集中在：并购与重组（2 项）：2026 年 8 月，公司通过出售闲置房产确认资产处置收益 213.89 万元计入非经常性损益，对当期利润形成一定支撑；2013 年 12 月，万向三农集团有限公司通过上交所大宗交易平台累计出售公司 2,046 万股股份；其他动作（1 项）：2024 年 4 月，第九届董事会第二十一次会议审议通过德农种业放弃万向财务公司增资优先认缴出资权暨关联交易议案。 公告口径上，近三年（2023 年 9 月---2026 年 9 月）公司披露重大重组与并购类公告 3 条、对外投资与转型类公告 3 条、回购与股权激励类公告 0 条，合计公告 236 条。 综合判断，报告期内公司有 6 条与战略相关的公告落地（重大重组与并购 3 条、对外投资与转型 3 条），资本运作与主业调整之间存在直接关联。

\begin{center}\small\setlength{\tabcolsep}{4pt}
\captionof{table}[万向德农（600371）历史融资与资本运作]{万向德农（600371）历史融资与资本运作（据公司公告与公开报道整理；金额为披露值，含首次公开发行、再融资、债券、银行借款与重整投资等）}
\begin{adjustbox}{max width=\textwidth}
\begin{tabular}{@{}p{1.45cm}p{2.55cm}p{4.05cm}p{4.95cm}@{}}
\toprule
时间 & 方式 & 规模 & 用途与说明 \\
\midrule
2002-08-29 & IPO（上交所主板） & 发行人民币普通股 4,\allowbreak{}000 万股，\allowbreak{}发行价 3.47 元/\allowbreak{}股 & 公司首发募集资金用于主业项目（招股文件口径）；发行后总股本 10,\allowbreak{}000 万股；公司设立时总股本 6,\allowbreak{}000 万股 \\
2024-04-12 & 现金分红（2023年度） & 以权益分派股权登记日总股本为基数每 10 股派现金股利 2.00 元 & 股东回报；最近三个会计年度累计现金分红金额（含税）1.61 亿元 \\
2025-2026 & 控股股东增持（二级市场集中竞价） & 万向三农集团拟增持总金额不低于人民币 2,\allowbreak{}500.00 万元、\allowbreak{}不超过人民币 5,\allowbreak{}000.00 万元 & 控股股东增持；公司未发行可转换公司债券 \\
\bottomrule
\end{tabular}
\end{adjustbox}
\end{center}

\pfl{历史融投资情况} 从现金流量表看，2025 年公司，筹资活动现金流净额 -0.49 亿元。 2026 年上半年筹资活动现金流净额为 -0.29 亿元（取得借款 — 亿元、偿还债务 — 亿元，未年化），已转为净流出，处于净还债状态。 2025 年分配股利、利润或偿付利息支付的现金合计 0.49 亿元。 公开披露的融资记录共 3 笔，工具类型以 IPO 为主；金额与用途见下表。 其中规模最大的两笔为：2025---2026 年，控股股东增持，万向三农集团拟增持总金额不低于人民币 2,500.00 万元、不超过人民币 5,000.00 万元；2002 年 8 月，IPO，发行人民币普通股 4,000 万股，发行价 3.47 元/股。 股本方面，近三年披露股本变动 3 次；总股本自 2021 年末以来未发生变化（2.93 亿股）。 分红方面，近三年现金分红 4 次、合计每股派现 0.65 元，最近一次实施在 2025 年 12 月。 综合看，公司融资结构上股权与债务融资并举；2025 年筹资活动现金流净额 -0.49 亿元，融资节奏仍以维持存量债务周转为主，与前述经营与投资现金流的方向基本一致。

\pfl{未来潜在融资需求分析} 公司 2026 年 6 月末资产负债率 14.5\%、2026 年上半年 ROE 2.6\%（未年化），资产负债表稳健、盈利水平平淡，暂不存在刚性融资需求；潜在的融资需求以技改扩产、并购或被并购为主。




\subsubsection{国投丰乐（000713）}
\label{co:000713}

\pfl{公司基本情况} 中国种业首家上市公司（1997 年深交所上市，原合肥丰乐种业），主营种子、农化、香料三大板块；2024 年国投种业受让 20\% 股份成为控股股东、实际控制人变更为国务院国资委，2025 年更名为国投丰乐，2026 年 1 月完成向控股股东定增 10.89 亿元，但 2026 年上半年仍亏损。 公司于 1997-04-22 在深交所主板首发上市，注册资本 7.98 亿元，总股本 6.14 亿股。 截至 2026 年 10 月 12 日收盘价 6.21 元，对应总市值 38.1 亿元，近一年-10.1\%，总市值在三级行业“种子”的 8 家公司中按由大到小排第\zhnumber{8}位，近一年涨跌幅高于全样本中位数（-10.4\%）0.2 个百分点。 按 2026 年半年报披露的行业构成，收入占比前三位为农化类 69.6\%；种子类 16.5\%；香料类 13.8\%。

\pfl{股权结构与控制权} 按 2026 年 6 月 30 日的股东名册，前三大股东为国投种业科技有限公司 持股 38.46\%；合肥市建设投资控股(集团)有限公司 持股 7.11\%；夏重阳 持股 1.25\%，股东户数 53,046 户。 股东结构与控制权方面，国投种业以现金认购 2025 年向特定对象发行股票 1.84 亿股；公司实际控制人状态自 2024 年 4 月变更为国务院国资委后未再变化。

\pfl{经营要点} 经营与产能方面，报告期内同玉 105LP 转基因玉米品种通过国家农作物品种审定初审并进入公示阶段。

\pfl{所处行业的基本面} 种子的供给由制种面积与库存决定、需求由粮价与种植意愿决定，制种面积在 2025 年起出现萎缩，是本轮供给端最先出清的环节（第 \ref{sec:framework} 章的供需框架）。 行业层面的最新情况是：3 家头部企业因品种市场萎缩制种面积下降超 50\%。

\pfl{竞争优势} 把该公司放回同类公司比较，2026 年上半年营业收入 11.9 亿元、在三级行业“种子”的 8 家公司中排第\zhnumber{3}位；毛利率 12.0\%、低于行业内中位数 16.4 个百分点；半年 ROE -1.0\%（未年化）、低于行业内中位数（0.17\%）；资产负债率 32.1\%、低于行业内中位数（37.3\%）。 与同类公司相比，公司的差异化在于：公司为中国种业首家上市公司；公司围绕核心品种铁 391K 签署三方联动维权合作协议；2026 年上半年香料产业营业收入 1.65 亿元、同比增长 13.13\%，为公司提供与种业周期弱相关的收入来源。 就当期报表看，公司尚未把经营层面的特征转化为高于行业的资本回报，周期底部的相对优势更多体现在现金流与资产负债表上。

\pfl{历史股价} 以 2021 年 1 月为起点、2026 年 10 月为终点，股价由 15.87 元变为 6.21 元，累计 -60.9\%，区间最高 17.08 元（2021 年 2 月）、最低 5.22 元（2026 年 6 月）；当前价相当于区间最高价的 36.4\%，为区间最低价的 1.19 倍。 月度收益的年化波动率 33.3\%，低于全样本中位数 37.3\%，在三级行业“种子”的 8 家公司中波动率由高到低排第\zhnumber{6}位。 把股价阶段与公司事件对齐看，2022 年 5 月 至 2024 年 6 月股价下跌-50.1\%，同期（2024 年 4 月）完成股份过户登记取得丰乐种业 20\% 股权。 把这些阶段与公司事件对齐后可以看到，几轮大涨大跌多由公司自身的资本运作与风险事件触发，行业景气只解释了其中一部分。

\begin{center}
\begin{minipage}{\textwidth}\centering
\begin{tikzpicture}
\begin{axis}[
  width=12.6cm, height=3.9cm, scale only axis,
  xmin=2020.922, xmax=2026.972, ymin=4.85, ymax=18.28,
  xtick={2021,2022,2023,2024,2025,2026},
  yearticks,
  yticklabel style={font=\scriptsize},
  ylabel={元/股}, ylabel style={font=\scriptsize},
  grid=major, grid style={LightGray, line width=0.3pt},
  axis line style={InkGray, line width=0.5pt},
  every axis plot/.append style={line width=0.9pt},
]
\addplot[AgriGreen] table[x=x,y=y]{data/clean/profiles/px_000713.dat};
\addplot[SoftRed, dashed, line width=0.7pt] table[x=x,y=y]{data/clean/profiles/pxzz_000713.dat};
\addplot[only marks, mark=*, mark size=1.1pt, SoftRed] table[x=x,y=y]{data/clean/profiles/pxzz_000713.dat};
\end{axis}
\end{tikzpicture}

\captionof{figure}[国投丰乐（000713）股价与周期骨架]{国投丰乐（000713）月度收盘价（元/股）与 ZigZag 周期骨架（阈值 25\%，圆点为周期转折点）}
\end{minipage}
\end{center}

\pfl{过去周期分析} 以 25\% 的 ZigZag 阈值划分，样本区间内股价形成 2 个主升段与 3 个主跌段。 最大主升段出现在 2024 年 6 月 至 2024 年 11 月，5 个月+37.6\%；最大主跌段为 2022 年 5 月 至 2024 年 6 月，25 个月-50.1\%。 最近一次转折出现在 2026 年 6 月（下跌段自 2024 年 11 月起，19 个月-30.1\%），当前处于该段的延续之中。 公司股价与申万农林牧渔指数几乎同步见底，个股并未走出独立行情。 种子行业的价格由制种成本、粮价与政策（转基因商业化、种业振兴）共同决定（第 \ref{sec:planting} 节），与生猪价格的联动有限，公司 2026 年 6 月末资产负债率 32.1\%、半年 ROE -1.0\%（未年化），在本轮下行中属于随行业同步调整的一类。 综合区间形态看，该股单段涨跌幅度偏大、以趋势段为主（最大主升 +37.6\%、最大主跌 -50.1\%），年化波动率 33.3\% 与全样本中位数 37.3\% 相比波动小于样本中位数。

\pfl{盈利情况} 2026 年上半年营业收入 11.95 亿元（同比 +3.9\%），归母净利润 -0.31 亿元，亏损同比扩大 0.03 亿元。 以年报为趋势对照，2023---2025 年营业收入由 31.14 亿元变为 29.03 亿元（两年年均复合增速 -3.4\%），归母净利润由 0.40 亿元变为 0.66 亿元。 按半年报口径，2026 年上半年的 ROE 为 -1.0\%（未年化）、销售净利率 -3.0\%、毛利率 12.0\%，ROE 在“种子”8 家公司中按数值由高到低排第\zhnumber{5}位。 同期全样本半年 ROE 中位数 0.75\%，公司与之相差 1.75 个百分点（弱于中位数公司）。 从公司披露的原因看，业绩变动主要来自：2026 年上半年，同比减少 561.36 万元、下降 605.29\%。 综合看，公司仍处亏损区间、亏损同比扩大，半年 ROE 在三级行业内按由高到低排第\zhnumber{5}位，单位盈利能力的修复依赖行业景气与自身成本端的改善，报表弹性主要体现为价格与销量的方向性变化。

\begin{center}\small\setlength{\tabcolsep}{3.4pt}
\captionof{table}[国投丰乐（000713）关键财务指标]{国投丰乐（000713）关键财务指标（合并报表口径；两个半年报期均未年化；现金短债比 $=$ 货币资金 $\div$ 短期债务）}
\begin{adjustbox}{max width=\textwidth}
\begin{tabular}{@{}lrrrrr@{}}
\toprule
指标 & 2023 年 & 2024 年 & 2025 年 & 2025 年半年报 & 2026 年半年报 \\
\midrule
营业收入（亿元） & 31.14 & 29.26 & 29.03 & 11.50 & 11.95 \\
归母净利润（亿元） & 0.40 & 0.70 & 0.66 & -0.28 & -0.31 \\
毛利率（\%） & 12.7 & 15.7 & 16.5 & 10.7 & 12.0 \\
ROE（\%） & 2.1 & 3.5 & 3.2 & -1.4 & -1.0 \\
资产负债率（\%） & 30.6 & 35.8 & 29.7 & 37.1 & 32.1 \\
经营活动现金流净额（亿元） & 1.66 & -0.28 & 2.92 & -1.05 & -0.97 \\
货币资金（亿元） & 1.45 & 2.99 & 16.71 & 2.22 & 15.50 \\
有息负债合计（亿元） & 1.57 & 2.37 & 3.03 & 3.68 & 3.72 \\
现金短债比（倍） & 26.60 & 4.64 & 9.73 & 1.07 & 6.52 \\
\bottomrule
\end{tabular}
\end{adjustbox}
\end{center}

\begin{center}
\begin{minipage}{\textwidth}\centering
\begin{tikzpicture}
\begin{groupplot}[
  group style={group size=2 by 1, horizontal sep=0.60cm},
  width=6.85cm, height=4.60cm,
  tick label style={font=\tiny},
  title style={font=\scriptsize\heiti},
  yticklabel style={font=\tiny},
  ylabel style={font=\tiny},
  grid=major, grid style={LightGray, line width=0.3pt},
  axis line style={InkGray, line width=0.5pt},
  enlarge x limits={abs=0.8},
  legend style={font=\scriptsize, at={(0.5,-0.20)}, anchor=north,
                legend columns=2, draw=none, fill=none, column sep=6pt},
]
\nextgroupplot[
  ybar, bar width=4pt,
  ymin=-3.46, ymax=34.91,
  xtick={0,1,2,3,4,5.7},
  xticklabels={2021,2022,2023,2024,2025,2026H1},
  ylabel={亿元},
]
\addplot[fill=AgriGreen!75, draw=AgriGreen, bar shift=-2.6pt] table[x=i, y=rev]{data/clean/profiles/fin_000713.dat};
\addplot[fill=SoftRed!60, draw=SoftRed, bar shift=2.6pt] table[x=i, y=ni]{data/clean/profiles/fin_000713.dat};
\addplot[fill=AgriGreen!30, draw=AgriGreen, pattern=north east lines, bar shift=-2.6pt] table[x=x, y=rev]{data/clean/profiles/finh_000713.dat};
\addplot[fill=SoftRed!20, draw=SoftRed, pattern=north east lines, bar shift=2.6pt] table[x=x, y=ni]{data/clean/profiles/finh_000713.dat};
\legend{营业收入（年/半年）, 归母净利润（年/半年）}
\nextgroupplot[
  ymin=-3.9, ymax=39.5,
  xtick={0,1,2,3,4,5.7},
  xticklabels={2021,2022,2023,2024,2025,2026H1},
  ylabel={\%},
]
\addplot[AgriGold, mark=*, mark size=1.3pt] table[x=i, y=roe]{data/clean/profiles/fin_000713.dat};
\addplot[OilBlue, mark=square*, mark size=1.3pt] table[x=i, y=debt]{data/clean/profiles/fin_000713.dat};
\addplot[only marks, AgriGold, mark=*, mark size=2.0pt] table[x=x, y=roe]{data/clean/profiles/finh_000713.dat};
\addplot[only marks, OilBlue, mark=square*, mark size=2.0pt] table[x=x, y=debt]{data/clean/profiles/finh_000713.dat};
\legend{ROE, 资产负债率}
\end{groupplot}
\end{tikzpicture}
\begin{tikzpicture}
\begin{axis}[
  name=finbot,
  width=13.75cm, height=3.10cm, scale only axis,
  ymin=-1.67, ymax=18.72,
  xtick={0,1,2,3,4,5.7},
  xticklabels={2021,2022,2023,2024,2025,2026H1},
  tick label style={font=\tiny},
  yticklabel style={font=\tiny},
  ylabel={亿元}, ylabel style={font=\tiny},
  ybar, bar width=4pt,
  grid=major, grid style={LightGray, line width=0.3pt},
  axis line style={InkGray, line width=0.5pt},
  enlarge x limits={abs=0.8},
  legend style={font=\scriptsize, at={(0.5,-0.26)}, anchor=north,
                legend columns=2, draw=none, fill=none, column sep=10pt},
]
\addplot[fill=AgriGreen!55, draw=AgriGreen, bar shift=-3.0pt] table[x=i, y=cash]{data/clean/profiles/fin_000713.dat};
\addplot[fill=SoftRed!45, draw=SoftRed, bar shift=3.0pt] table[x=i, y=ibd]{data/clean/profiles/fin_000713.dat};
\addplot[fill=AgriGreen!25, draw=AgriGreen, pattern=north east lines, bar shift=-3.0pt] table[x=x, y=cash]{data/clean/profiles/finh_000713.dat};
\addplot[fill=SoftRed!20, draw=SoftRed, pattern=north east lines, bar shift=3.0pt] table[x=x, y=ibd]{data/clean/profiles/finh_000713.dat};
\legend{货币资金（左轴）, 有息负债（左轴）}
\end{axis}
\begin{axis}[
  at={(finbot.south east)}, anchor=south east,
  width=13.75cm, height=3.10cm, scale only axis,
  axis y line*=right, axis x line*=none, xtick=\empty,
  ymin=-2.25, ymax=5.44,
  tick label style={font=\tiny},
  yticklabel style={font=\tiny}, scaled y ticks=false,
  legend style={font=\scriptsize, at={(0.5,-0.44)}, anchor=north,
                legend columns=1, draw=none, fill=none},
]
\addplot[OilBlue, mark=*, mark size=2.2pt, line width=1.3pt] table[x=i, y=ocf]{data/clean/profiles/fin_000713.dat};
\addplot[only marks, OilBlue, mark=*, mark size=3.2pt] table[x=x, y=ocf]{data/clean/profiles/finh_000713.dat};
\legend{经营活动现金流净额（右轴，亿元，蓝点）}
\end{axis}
\end{tikzpicture}

\captionof{figure}[国投丰乐（000713）近年主要财务指标]{国投丰乐（000713）近年主要财务指标（左上：营业收入与归母净利润；右上：ROE 与资产负债率；下：货币资金、有息负债为左轴柱状，经营活动现金流净额为其右轴的蓝点折线。
2021---2025 年为年报口径，2026H1 为 2026 年半年报/半年末，与其前一年报之间留出间隔、以斜纹或空心点区分；半年报数据未年化）}
\end{minipage}
\end{center}

\pfl{财务分析} 2026 年 6 月末资产负债率 32.1\%，较 2025 年末上升 2.4 个百分点（样本中位数 52.2\%，所处三级行业内按由低到高排第\zhnumber{3}位）。 2026 年 6 月末货币资金 15.50 亿元、短期债务（短期借款与一年内到期非流动负债）2.38 亿元、长期债务 1.34 亿元，有息负债合计 3.72 亿元，现金短债比 6.52 倍——覆盖充分。 净有息负债 -11.78 亿元，相当于其 2026 年上半年营业收入的 -99\%。 流动比率 2.45、速动比率 1.77，短期流动性无明显缺口。 2026 年上半年经营活动现金流净额 -0.97 亿元，经营现金流与利润同时为负，现金消耗较快。 2026 年 6 月末在建工程 1.95 亿元，较上年末增加 0.51 亿元，仍有在建投入。 2026 年 6 月末商誉 2.21 亿元，占归母净资产的 7.0\%，是净资产中最脆弱的一块。 综合看，现金短债比 6.52 倍，短期偿付压力可控；资产负债率低于样本中位数 20.1 个百分点；有息负债中短期债务占比更高，期限结构仍以滚动续作为主。

\pfl{重大战略转型} 从公开披露看，公司的战略推进集中在：并购与重组（2 项）：2026 年 6 月，与国投种业签署国投（张掖）金种科技有限公司股权转让协议，以 3,376.48 万元现金收购国投金种 60\% 股权；2024 年 4 月，完成股份过户登记取得丰乐种业 20\% 股权；其他动作（2 项）：2026 年上半年，玉米产业销量稳步增长、产品结构不断优化。 公告口径上，近三年（2023 年 9 月---2026 年 9 月）公司披露重大重组与并购类公告 7 条、对外投资与转型类公告 0 条、回购与股权激励类公告 0 条，合计公告 404 条。 第一大行业仍是“农化类”，收入占比 69.6\%，与 2021 年年报相比没有方向性变化。 综合判断，报告期内公司有 7 条与战略相关的公告落地（重大重组与并购 7 条、对外投资与转型 0 条），资本运作与主业调整之间存在直接关联。

\begin{center}\small\setlength{\tabcolsep}{4pt}
\captionof{table}[国投丰乐（000713）历史融资与资本运作]{国投丰乐（000713）历史融资与资本运作（据公司公告与公开报道整理；金额为披露值，含首次公开发行、再融资、债券、银行借款与重整投资等）}
\begin{adjustbox}{max width=\textwidth}
\begin{tabular}{@{}p{1.45cm}p{2.55cm}p{4.05cm}p{4.95cm}@{}}
\toprule
时间 & 方式 & 规模 & 用途与说明 \\
\midrule
1997-04-22 & IPO（深交所） & 发行价 6.75 元/\allowbreak{}股 & 公司为中国种业首家上市公司 \\
2024-04-29 & 控股股东协议受让（控制权变更，非上市公司融资） & 国投种业以现金受让合肥建投持有的 1.23 亿股无限售流通股（占股份总数 20.00\%） & 取得上市公司控股权；2024 年 1 月国投种业收到市场监管总局经营者集中反垄断审查不实施进一步审查决定书（反执二审查决定〔2024〕31 号）；2024 年 3 月 25 日国务院国资委出具国资产权〔2024〕106 号批复原则同意 \\
2025-10-17 & 要约收购义务豁免安排 & 国投种业认购本次发行股票将触发要约收购义务 & 顺利完成向控股股东定增；发行完成后国投种业持有股份比例超过 30\% \\
2025-12-25 & 向特定对象发行A股股票（发行对象为控股股东国投种业，以现金认购） & 发行 1.84 亿股，\allowbreak{}发行价 5.91 元/\allowbreak{}股，\allowbreak{}募集资金总额 10.89 亿元 & 全部用于补充流动资金及偿还银行借款；募集资金于 2025 年 12 月 26 日到账；2026 年 1 月 21 日新增股份上市；发行完成后国投种业持股比例由 20.00\% 升至 38.46\% \\
2026-06-25 & 现金收购关联方股权（关联交易） & 以现金支付方式收购控股股东国投种业持有的国投（张掖）金种科技有限公司 60\% 股权 & 成为国投金种控股股东并纳入合并报表；国投金种另由张掖现代种业集团有限公司持有 40\%；国投种业直接持有公司 38.46\% 股权 \\
\bottomrule
\end{tabular}
\end{adjustbox}
\end{center}

\pfl{历史融投资情况} 从现金流量表看，2025 年公司取得借款收到的现金 2.53 亿元、偿还债务支付的现金 1.70 亿元、吸收权益性投资 10.81 亿元，筹资活动现金流净额 11.42 亿元（2023 年为取得借款 1.65 亿元、偿还债务 1.61 亿元）。 2026 年上半年筹资活动现金流净额为 0.42 亿元（取得借款 1.48 亿元、偿还债务 0.76 亿元，未年化），仍处于净融入状态。 2025 年分配股利、利润或偿付利息支付的现金合计 0.24 亿元。 公开披露的融资记录共 5 笔，工具类型以 IPO、定向增发为主；金额与用途见下表。 其中规模最大的两笔为：2026 年 6 月，现金收购关联方股权，以现金支付方式收购控股股东国投种业持有的国投（张掖）金种科技有限公司 60\% 股权。 股本方面，近三年披露股本变动 3 次；总股本自 2021 年末以来未发生变化（6.14 亿股）。 分红方面，近三年现金分红 4 次、合计每股派现 0.09 元，最近一次实施在 2025 年 12 月。 近三年“再融资”类公告 45 条，涉及定向增发、募集资金使用。 综合看，公司融资结构上以债务滚动为主、股权融资为补充；2025 年筹资活动现金流净额 11.42 亿元，年度内仍为净融入，与前述经营与投资现金流的方向基本一致。

\pfl{未来潜在融资需求分析} 2026 年 6 月末货币资金 15.50 亿元，而短期债务 2.38 亿元（现金短债比 6.52 倍），2026 年上半年经营现金流净额 -0.97 亿元。账面货币资金可以覆盖短期债务，缺口不体现在静态偿付层面。 公司 2026 年上半年归母净利润为负（-0.31 亿元），但 6 月末资产负债率 32.1\% 尚未进入第一档（亏损且负债率 $>$65\%）的区间，当前融资需求以补充流动资金、支撑低谷期经营性支出为主。 公开披露的后续资金安排线索包括：股权融资线索：2026 年 4 月，发行数量不超过 1.84 亿股，募集资金总额不超过 10.89 亿元（含本数）；2025 年度，另支付发行费用合计 307.55 万元（不含税）；债务与授信安排：2026 年 8 月，公司募集资金用途明确为补充流动资金及偿还银行借款。 资金需求还要叠加在建工程：期末在建工程 1.95 亿元、2026 年上半年购建固定资产等支出 0.45 亿元（未年化），这部分支出在周期底部通常会被主动放缓。 静态偿付不存在缺口，后续融资更可能服务于产能或并购安排，而不是填补流动性窟窿。




\subsubsection{新赛股份（600540）}
\label{co:600540}

\pfl{公司基本情况} 2003 年设立至今，新赛股份是新疆兵团系棉花全产业链企业（皮棉、棉籽精深加工、棉油棉蛋白），2004 年上交所上市、历史上曾被实施退市风险警示（*ST 新赛）；2023 年 9 月控股股东变更为中新建物流、实际控制人为兵团国资委；2024—2025 年连续亏损，2026 年上半年扭亏且净利同比增逾 18 倍，同时处于证监会立案调查期间。 公司于 2004-01-07 在上交所首发上市，注册资本 5.81 亿元，总股本 5.81 亿股。 截至 2026 年 10 月 12 日收盘价 5.23 元，对应总市值 30.4 亿元，近一年+8.7\%，总市值在三级行业“其他种植业”的 6 家公司中按由大到小排第\zhnumber{5}位，近一年涨跌幅高于全样本中位数（-10.4\%）19.1 个百分点。 股权结构方面，控股股东为新疆艾比湖农工商联合企业总公司。

\pfl{股权结构与控制权} 按 2026 年 6 月 30 日的股东名册，前三大股东为中新建物流集团有限责任公司 持股 29.16\%；中新建胡杨私募股权基金管理有限公司－新疆中新建胡杨产 持股 5.84\%；新疆艾比湖投资有限公司 持股 5.00\%，股东户数 30,143 户。 控制权与股东性质方面，第二大股东新疆中新建胡杨产业投资基金合伙企业持股 7.83\%。

\pfl{经营要点} 经营层面，上市日期 2004-01-07。

\pfl{所处行业的基本面} 该子行业按品种各自定价：天然橡胶看全球供给与汽车需求，棉花看新疆产量与配额，农药看原药价格与农业服务需求（第 \ref{sec:framework} 章的供需框架）。 行业层面的最新情况是：2026 年年度，全国棉花实播面积同比减少 145.8 万亩、减幅 3.1\%，其中新疆实播面积 4,153.5 万亩；农业农村部 2026 年 6 月将 2026/27 年度新疆棉花种植面积调增至 4,310 万亩、预计总产量 634 万吨。

\pfl{竞争优势} 把该公司放回同类公司比较，2026 年上半年营业收入 25.7 亿元、在三级行业“其他种植业”的 6 家公司中排第\zhnumber{3}位；毛利率 8.7\%、低于行业内中位数 8.2 个百分点；半年 ROE 30.8\%（未年化）、高于行业内中位数（1.70\%）；资产负债率 82.7\%、高于行业内中位数（59.7\%）。 相对优势方面，实施皮棉组货供应并与源头轧花厂、终端纺纱厂合作；上半年该板块收入 2.38 亿元、同比增加 1.31 亿元。 优势在当期体现为高于行业中位数的资本回报，但优势能否持续仍取决于周期与自身产能投放节奏。

\pfl{历史股价} 以 2021 年 1 月为起点、2026 年 10 月为终点，股价由 3.95 元变为 5.23 元，累计 +32.4\%，区间最高 7.61 元（2021 年 3 月）、最低 3.31 元（2026 年 6 月）；当前价相当于区间最高价的 68.7\%，为区间最低价的 1.58 倍。 月度收益的年化波动率 48.0\%，高于全样本中位数 37.3\%，在三级行业“其他种植业”的 6 家公司中波动率由高到低排第\zhnumber{1}位。

\begin{center}
\begin{minipage}{\textwidth}\centering
\begin{tikzpicture}
\begin{axis}[
  width=12.6cm, height=3.9cm, scale only axis,
  xmin=2020.922, xmax=2026.972, ymin=3.08, ymax=8.14,
  xtick={2021,2022,2023,2024,2025,2026},
  yearticks,
  yticklabel style={font=\scriptsize},
  ylabel={元/股}, ylabel style={font=\scriptsize},
  grid=major, grid style={LightGray, line width=0.3pt},
  axis line style={InkGray, line width=0.5pt},
  every axis plot/.append style={line width=0.9pt},
]
\addplot[AgriGreen] table[x=x,y=y]{data/clean/profiles/px_600540.dat};
\addplot[SoftRed, dashed, line width=0.7pt] table[x=x,y=y]{data/clean/profiles/pxzz_600540.dat};
\addplot[only marks, mark=*, mark size=1.1pt, SoftRed] table[x=x,y=y]{data/clean/profiles/pxzz_600540.dat};
\end{axis}
\end{tikzpicture}

\captionof{figure}[新赛股份（600540）股价与周期骨架]{新赛股份（600540）月度收盘价（元/股）与 ZigZag 周期骨架（阈值 25\%，圆点为周期转折点）}
\end{minipage}
\end{center}

\pfl{过去周期分析} 以 25\% 的 ZigZag 阈值划分，样本区间内股价形成 5 个主升段与 4 个主跌段。 最大主跌段为 2023 年 9 月 至 2024 年 6 月，9 个月-47.1\%；最大主升段出现在 2021 年 1 月 至 2021 年 3 月，2 个月+92.7\%。 最近一次转折出现在 2026 年 8 月（上涨段自 2026 年 6 月起，2 个月+68.0\%），当前处于该段之后的震荡之中。 公司股价与申万农林牧渔指数几乎同步见底，个股并未走出独立行情。 该子行业的周期来源与猪周期不同（第 \ref{sec:grain} 章、第 \ref{sec:soft} 章），价格更多由自身供给、进口政策与全球定价决定，公司 2026 年 6 月末资产负债率 82.7\%，高于全样本中位数 52.2\%，其股价因此更多反映资金面与信用风险，而不是价格弹性本身。 综合区间形态看，该股单段涨跌幅度偏大、以趋势段为主（最大主升 +92.7\%、最大主跌 -47.1\%），年化波动率 48.0\% 与全样本中位数 37.3\% 相比波动明显大于样本中位数。

\pfl{盈利情况} 2026 年上半年营业收入 25.71 亿元（同比 +77.2\%），归母净利润 1.25 亿元，同比 +1819.7\%。 以年报为趋势对照，2023---2025 年营业收入由 9.72 亿元变为 30.12 亿元（两年年均复合增速 +76.1\%），归母净利润由 0.15 亿元变为 -1.42 亿元。 按半年报口径，2026 年上半年的 ROE 为 30.8\%（未年化）、销售净利率 4.9\%、毛利率 8.7\%，ROE 在“其他种植业”6 家公司中按数值由高到低排第\zhnumber{1}位。 按同一口径，三级行业“其他种植业”6 家公司的半年 ROE 中位数为 1.70\%，该公司高于其中位数 29.15 个百分点。 综合看，公司在报告期维持盈利（高于全样本半年 ROE 中位数 0.75\%），利润的可持续性取决于其他种植业景气的延续与成本管控的执行。

\begin{center}\small\setlength{\tabcolsep}{3.4pt}
\captionof{table}[新赛股份（600540）关键财务指标]{新赛股份（600540）关键财务指标（合并报表口径；两个半年报期均未年化；现金短债比 $=$ 货币资金 $\div$ 短期债务）}
\begin{adjustbox}{max width=\textwidth}
\begin{tabular}{@{}lrrrrr@{}}
\toprule
指标 & 2023 年 & 2024 年 & 2025 年 & 2025 年半年报 & 2026 年半年报 \\
\midrule
营业收入（亿元） & 9.72 & 15.55 & 30.12 & 14.51 & 25.71 \\
归母净利润（亿元） & 0.15 & -2.44 & -1.42 & 0.07 & 1.25 \\
毛利率（\%） & 3.4 & -4.4 & 0.5 & 2.3 & 8.7 \\
ROE（\%） & 2.0 & -40.2 & -34.4 & 1.3 & 30.8 \\
资产负债率（\%） & 75.1 & 89.3 & 95.1 & 86.2 & 82.7 \\
经营活动现金流净额（亿元） & -4.29 & -8.19 & -13.17 & 18.27 & 24.20 \\
货币资金（亿元） & 8.72 & 10.69 & 17.05 & 16.53 & 8.39 \\
有息负债合计（亿元） & 18.55 & 35.62 & 57.97 & 22.40 & 15.08 \\
现金短债比（倍） & 0.50 & 0.31 & 0.30 & 0.77 & 0.58 \\
\bottomrule
\end{tabular}
\end{adjustbox}
\end{center}

\begin{center}
\begin{minipage}{\textwidth}\centering
\begin{tikzpicture}
\begin{groupplot}[
  group style={group size=2 by 1, horizontal sep=0.60cm},
  width=6.85cm, height=4.60cm,
  tick label style={font=\tiny},
  title style={font=\scriptsize\heiti},
  yticklabel style={font=\tiny},
  ylabel style={font=\tiny},
  grid=major, grid style={LightGray, line width=0.3pt},
  axis line style={InkGray, line width=0.5pt},
  enlarge x limits={abs=0.8},
  legend style={font=\scriptsize, at={(0.5,-0.20)}, anchor=north,
                legend columns=2, draw=none, fill=none, column sep=6pt},
]
\nextgroupplot[
  ybar, bar width=4pt,
  ymin=-6.07, ymax=34.07,
  xtick={0,1,2,3,4,5.7},
  xticklabels={2021,2022,2023,2024,2025,2026H1},
  ylabel={亿元},
]
\addplot[fill=AgriGreen!75, draw=AgriGreen, bar shift=-2.6pt] table[x=i, y=rev]{data/clean/profiles/fin_600540.dat};
\addplot[fill=SoftRed!60, draw=SoftRed, bar shift=2.6pt] table[x=i, y=ni]{data/clean/profiles/fin_600540.dat};
\addplot[fill=AgriGreen!30, draw=AgriGreen, pattern=north east lines, bar shift=-2.6pt] table[x=x, y=rev]{data/clean/profiles/finh_600540.dat};
\addplot[fill=SoftRed!20, draw=SoftRed, pattern=north east lines, bar shift=2.6pt] table[x=x, y=ni]{data/clean/profiles/finh_600540.dat};
\legend{营业收入（年/半年）, 归母净利润（年/半年）}
\nextgroupplot[
  ymin=-51.0, ymax=108.6,
  xtick={0,1,2,3,4,5.7},
  xticklabels={2021,2022,2023,2024,2025,2026H1},
  ylabel={\%},
]
\addplot[AgriGold, mark=*, mark size=1.3pt] table[x=i, y=roe]{data/clean/profiles/fin_600540.dat};
\addplot[OilBlue, mark=square*, mark size=1.3pt] table[x=i, y=debt]{data/clean/profiles/fin_600540.dat};
\addplot[only marks, AgriGold, mark=*, mark size=2.0pt] table[x=x, y=roe]{data/clean/profiles/finh_600540.dat};
\addplot[only marks, OilBlue, mark=square*, mark size=2.0pt] table[x=x, y=debt]{data/clean/profiles/finh_600540.dat};
\legend{ROE, 资产负债率}
\end{groupplot}
\end{tikzpicture}
\begin{tikzpicture}
\begin{axis}[
  name=finbot,
  width=13.75cm, height=3.10cm, scale only axis,
  ymin=-5.80, ymax=64.92,
  xtick={0,1,2,3,4,5.7},
  xticklabels={2021,2022,2023,2024,2025,2026H1},
  tick label style={font=\tiny},
  yticklabel style={font=\tiny},
  ylabel={亿元}, ylabel style={font=\tiny},
  ybar, bar width=4pt,
  grid=major, grid style={LightGray, line width=0.3pt},
  axis line style={InkGray, line width=0.5pt},
  enlarge x limits={abs=0.8},
  legend style={font=\scriptsize, at={(0.5,-0.26)}, anchor=north,
                legend columns=2, draw=none, fill=none, column sep=10pt},
]
\addplot[fill=AgriGreen!55, draw=AgriGreen, bar shift=-3.0pt] table[x=i, y=cash]{data/clean/profiles/fin_600540.dat};
\addplot[fill=SoftRed!45, draw=SoftRed, bar shift=3.0pt] table[x=i, y=ibd]{data/clean/profiles/fin_600540.dat};
\addplot[fill=AgriGreen!25, draw=AgriGreen, pattern=north east lines, bar shift=-3.0pt] table[x=x, y=cash]{data/clean/profiles/finh_600540.dat};
\addplot[fill=SoftRed!20, draw=SoftRed, pattern=north east lines, bar shift=3.0pt] table[x=x, y=ibd]{data/clean/profiles/finh_600540.dat};
\legend{货币资金（左轴）, 有息负债（左轴）}
\end{axis}
\begin{axis}[
  at={(finbot.south east)}, anchor=south east,
  width=13.75cm, height=3.10cm, scale only axis,
  axis y line*=right, axis x line*=none, xtick=\empty,
  ymin=-22.89, ymax=35.42,
  tick label style={font=\tiny},
  yticklabel style={font=\tiny}, scaled y ticks=false,
  legend style={font=\scriptsize, at={(0.5,-0.44)}, anchor=north,
                legend columns=1, draw=none, fill=none},
]
\addplot[OilBlue, mark=*, mark size=2.2pt, line width=1.3pt] table[x=i, y=ocf]{data/clean/profiles/fin_600540.dat};
\addplot[only marks, OilBlue, mark=*, mark size=3.2pt] table[x=x, y=ocf]{data/clean/profiles/finh_600540.dat};
\legend{经营活动现金流净额（右轴，亿元，蓝点）}
\end{axis}
\end{tikzpicture}

\captionof{figure}[新赛股份（600540）近年主要财务指标]{新赛股份（600540）近年主要财务指标（左上：营业收入与归母净利润；右上：ROE 与资产负债率；下：货币资金、有息负债为左轴柱状，经营活动现金流净额为其右轴的蓝点折线。
2021---2025 年为年报口径，2026H1 为 2026 年半年报/半年末，与其前一年报之间留出间隔、以斜纹或空心点区分；半年报数据未年化）}
\end{minipage}
\end{center}

\pfl{财务分析} 2026 年 6 月末资产负债率 82.7\%，较 2025 年末下降 12.4 个百分点（样本中位数 52.2\%，所处三级行业内按由低到高排第\zhnumber{6}位）。 2026 年 6 月末货币资金 8.39 亿元、短期债务（短期借款与一年内到期非流动负债）14.35 亿元、长期债务 0.73 亿元，有息负债合计 15.08 亿元，现金短债比 0.58 倍——覆盖偏紧。 净有息负债 6.69 亿元，相当于其 2026 年上半年营业收入的 26\%。 流动比率 0.79、速动比率 0.70，短期偿债指标偏紧。 2026 年上半年经营活动现金流净额 24.20 亿元，经营现金流与利润同向为正，内生资金可以支撑运营。 综合看，现金短债比 0.58 倍，短期资金安排偏紧；资产负债率高于样本中位数 30.5 个百分点；有息负债中短期债务占比更高，期限结构仍以滚动续作为主。 公司披露的监管与合规风险为：2026 年 6 月，收到中国证监会立案告知书；2026 年 8 月，上海证券交易所对公司 2025 年年度报告信息披露监管问询函回复。

\pfl{重大战略转型} 公司的战略动作可归为以下几类：产能与项目（2 项）：2023 年 7 月，公司披露控股股东和实际控制人拟发生变更的提示性公告；2023 年 9 月，中新建物流增资完成后注册资本至 68.24 亿元；其他动作（2 项）：2026 年上半年，截至 6 月末子企业亏损户数较 2025 年末明显减少、亏损面大幅收窄。 公告口径上，近三年（2023 年 9 月---2026 年 9 月）公司披露重大重组与并购类公告 1 条、对外投资与转型类公告 0 条、回购与股权激励类公告 0 条，合计公告 423 条。 综合判断，报告期内公司有 1 条与战略相关的公告落地（重大重组与并购 1 条、对外投资与转型 0 条），资本运作与主业调整之间存在直接关联。

\begin{center}\small\setlength{\tabcolsep}{4pt}
\captionof{table}[新赛股份（600540）历史融资与资本运作]{新赛股份（600540）历史融资与资本运作（据公司公告与公开报道整理；金额为披露值，含首次公开发行、再融资、债券、银行借款与重整投资等）}
\begin{adjustbox}{max width=\textwidth}
\begin{tabular}{@{}p{1.45cm}p{2.55cm}p{4.05cm}p{4.95cm}@{}}
\toprule
时间 & 方式 & 规模 & 用途与说明 \\
\midrule
2003-12 & IPO（上交所，二级市场投资者定价配售） & 发行 5,\allowbreak{}000 万股，\allowbreak{}发行价 6.67 元/\allowbreak{}股，\allowbreak{}募集资金总额 3.33 亿元 & 首次公开发行募集资金（其中记入股本 5,\allowbreak{}000.00 万元、\allowbreak{}资本公积 2.64 亿元）；网上配号总数 71,\allowbreak{}874,\allowbreak{}879 个 \\
2014-11-20 & 非公开发行（向特定投资者，主承销商广州证券） & 发行人民币普通股（A 股）5,\allowbreak{}954 万股，\allowbreak{}每股发行价格 8.70 元 & 非公开发行募投项目（后续部分用途变更）；经中国证监会核准 \\
2018 & 银行借款与募集资金暂时补流 & 全年通过银行贷款融资 7.92 亿元（其中协助棉业公司 9,\allowbreak{}460.00 万元、\allowbreak{}协助博汇公司 1.62 亿元、\allowbreak{}协助精纺公司 4,\allowbreak{}000.00 万元）；实施募集资金暂时补充流动资金 1.29 亿元 & 满足分子公司在用款高峰期的资金需求、\allowbreak{}缓解企业资金压力；同期全年实际收回各类应收款项 2.44 亿元 \\
2023-08 & 变更部分募集资金用途 & 将原募投项目霍城县可利煤炭物流配送有限公司专用线扩建项目部分募集资金用途变更为偿还总部银行贷款项目，\allowbreak{}金额 3,\allowbreak{}600.00 万元 & 偿还总部银行贷款；其余部分募集资金用途未发生改变 \\
2026-04-28 & 借款及担保额度议案（2026年度） & 预计 2026 年度公司及合并报表范围内子公司所需借款额度 79.00 亿元（可上浮 15\%）；拟提供合计不超过 35.00 亿元的担保（含现有担保、\allowbreak{}展期或续保及新增担保 & 生产经营及子公司资金需求；2026 年 5 月 21 日 2025 年年度股东会审议通过；截至 2026 年 4 月 28 日公司及控股子公司对外担保总额 35.00 亿元 \\
\bottomrule
\end{tabular}
\end{adjustbox}
\end{center}

\pfl{历史融投资情况} 从现金流量表看，2025 年公司取得借款收到的现金 79.41 亿元、偿还债务支付的现金 54.78 亿元，筹资活动现金流净额 23.71 亿元（2023 年为取得借款 25.78 亿元、偿还债务 19.46 亿元）。 2026 年上半年筹资活动现金流净额为 -43.34 亿元（取得借款 8.81 亿元、偿还债务 51.67 亿元，未年化），已转为净流出，处于净还债状态。 2025 年分配股利、利润或偿付利息支付的现金合计 0.91 亿元。 公开披露的融资记录共 5 笔，工具类型以 IPO、定向增发、银行借款为主；金额与用途见下表。 其中规模最大的两笔为：2026 年 4 月，银行借款，预计 2026 年度公司及合并报表范围内子公司所需借款额度 79.00 亿元（可上浮 15\%），拟提供合计不超过 35.00 亿元的担保（含现有担保、展期或续保及新增担保；2023 年 8 月，变更部分募集资金用途，将原募投项目霍城县可利煤炭物流配送有限公司专用线扩建项目部分募集资金用途变更为偿还总部银行贷款项目，金额 3,600.00 万元。 股本方面，近三年披露股本变动 3 次；总股本由 2021 年末的 4.71 亿股变为 5.81 亿股（+23.45\%）。 分红方面，近三年未实施现金分红，在全部样本中属于分红能力偏弱的一端。 近三年“再融资”类公告 1 条，涉及定向增发。 综合看，公司融资结构上股权与债务融资并举；2025 年筹资活动现金流净额 23.71 亿元，年度内仍为净融入，与前述经营与投资现金流的方向基本一致。

\pfl{未来潜在融资需求分析} 2026 年 6 月末货币资金 8.39 亿元，而短期债务 14.35 亿元（现金短债比 0.58 倍），2026 年上半年经营现金流净额 24.20 亿元。按“短期债务减去账面货币资金”的滚动偿付口径，静态缺口约 5.96 亿元；上半年盈利或接近盈亏平衡，该缺口不随经营亏损进一步扩大。 公司 2026 年 6 月末资产负债率 82.7\%、2026 年上半年 ROE 30.8\%（未年化），资产负债率偏高（82.7\%），融资需求首先用于置换与滚动存量债务、补充营运资金，属于“补血型”而非扩张型。 公开披露的后续资金安排线索包括：债务与授信安排：2026 年度，公司及控股子公司借款额度 79.00 亿元、担保额度 35.00 亿元（均可上浮 15\%）公司对外部债务融资依赖度高；2026 年 4 月，公告公司及子公司申请银行综合授信额度部分募集资金投资项目延期对子公司流动资金借款担保计划。 综合看，公司的外部资金需求以补充营运资金为主、项目投入为辅，滚动偿付口径下的静态缺口约 5.96 亿元，融资的时点与规模更多取决于信用条件与资本市场窗口。




\subsubsection{新农开发（600359）}
\label{co:600359}

\pfl{公司基本情况} 1999 年设立至今，新农开发是新疆生产建设兵团第一师国资委控制的农业上市平台，主营作物种子、大米谷物加工、甘草饮品与肉牛养殖；2023 年剥离乳业，2025 年控股股东变更为农垦集团。 公司于 1999-04-29 在上交所首发上市，注册资本 3.82 亿元，总股本 3.82 亿股。 截至 2026 年 10 月 12 日收盘价 7.16 元，对应总市值 27.3 亿元，近一年+3.3\%，总市值在三级行业“其他种植业”的 6 家公司中按由大到小排第\zhnumber{6}位，近一年涨跌幅高于全样本中位数（-10.4\%）13.7 个百分点。 按 2026 年半年报披露的行业构成，收入占比前三位为服务业及其他 72.7\%；农业 20.6\%；工业 6.7\%。 与 2021 年年报相比，第一大行业口径由“农业”（55.1\%）变为“服务业及其他”（72.7\%）；两期披露的分类口径有调整，占比差异中同时包含分类因素。 股权结构方面，控股股东为阿拉尔塔河投资有限责任公司；实际控制人为仍为新疆生产建设兵团第一师国资委。

\pfl{股权结构与控制权} 按 2026 年 6 月 30 日的股东名册，前三大股东为新疆塔里木农垦集团有限公司 持股 40.32\%；高盛公司有限责任公司（持股比例未披露）；中国民生银行股份有限公司－金元顺安元启灵活配置混合型（持股比例未披露），股东户数 36,351 户。 股东结构与控制权方面，本次划转后间接控股股东仍为国投集团。

\pfl{经营要点} 从经营要点看，2023 年，重大资产出售交易价格为 3.18 亿元。

\pfl{所处行业的基本面} 该子行业按品种各自定价：天然橡胶看全球供给与汽车需求，棉花看新疆产量与配额，农药看原药价格与农业服务需求（第 \ref{sec:framework} 章的供需框架）。 从供需两端的当期数据看，2025 年，新疆棉花实际播种面积 3,887.5 万亩、同比增长 5.9\%；目标价格水平为每吨 1.86 万元，以固定产量 510 万吨进行补贴，用于质量补贴部分按当年补贴资金总额的 10\% 左右掌握。

\pfl{竞争优势} 把该公司放回同类公司比较，2026 年上半年营业收入 2.0 亿元、在三级行业“其他种植业”的 6 家公司中排第\zhnumber{6}位；毛利率 30.5\%、高于行业内中位数 13.6 个百分点；半年 ROE 2.8\%（未年化）、高于行业内中位数（1.70\%）；资产负债率 55.6\%、低于行业内中位数（59.7\%）。 与同类公司相比，公司的差异化在于：新农甘草搭建甘草种植、原料收购、精深加工到国内外销售的完整产业链。 优势在报表上主要体现为高于行业平均的毛利率，这在周期底部是抵御亏损的主要缓冲。

\pfl{历史股价} 以 2021 年 1 月为起点、2026 年 10 月为终点，股价由 7.68 元变为 7.16 元，累计 -6.8\%，区间最高 10.11 元（2022 年 5 月）、最低 5.40 元（2026 年 6 月）；当前价相当于区间最高价的 70.8\%，为区间最低价的 1.33 倍。 月度收益的年化波动率 28.9\%，低于全样本中位数 37.3\%，在三级行业“其他种植业”的 6 家公司中波动率由高到低排第\zhnumber{5}位。 把股价阶段与公司事件对齐看，2024 年 6 月 至 2026 年 1 月股价上涨+41.9\%，同期 2025 年半年报，公司披露截至报告期末尚有 3,176.09 万元新农乳业股权转让款未收回。 把这些阶段与公司事件对齐后可以看到，几轮大涨大跌多由公司自身的资本运作与风险事件触发，行业景气只解释了其中一部分。

\begin{center}
\begin{minipage}{\textwidth}\centering
\begin{tikzpicture}
\begin{axis}[
  width=12.6cm, height=3.9cm, scale only axis,
  xmin=2020.922, xmax=2026.972, ymin=5.02, ymax=10.82,
  xtick={2021,2022,2023,2024,2025,2026},
  yearticks,
  yticklabel style={font=\scriptsize},
  ylabel={元/股}, ylabel style={font=\scriptsize},
  grid=major, grid style={LightGray, line width=0.3pt},
  axis line style={InkGray, line width=0.5pt},
  every axis plot/.append style={line width=0.9pt},
]
\addplot[AgriGreen] table[x=x,y=y]{data/clean/profiles/px_600359.dat};
\addplot[SoftRed, dashed, line width=0.7pt] table[x=x,y=y]{data/clean/profiles/pxzz_600359.dat};
\addplot[only marks, mark=*, mark size=1.1pt, SoftRed] table[x=x,y=y]{data/clean/profiles/pxzz_600359.dat};
\end{axis}
\end{tikzpicture}

\captionof{figure}[新农开发（600359）股价与周期骨架]{新农开发（600359）月度收盘价（元/股）与 ZigZag 周期骨架（阈值 25\%，圆点为周期转折点）}
\end{minipage}
\end{center}

\pfl{过去周期分析} 以 25\% 的 ZigZag 阈值划分，样本区间内股价形成 3 个主升段与 2 个主跌段。 最大主跌段为 2022 年 5 月 至 2024 年 6 月，25 个月-44.7\%；最大主升段出现在 2024 年 6 月 至 2026 年 1 月，19 个月+41.9\%。 最近一次转折出现在 2026 年 10 月（上涨段自 2026 年 6 月起，4 个月+32.6\%），当前处于该段之后的震荡之中。 公司股价与申万农林牧渔指数几乎同步见底，个股并未走出独立行情。 该子行业的周期来源与猪周期不同（第 \ref{sec:grain} 章、第 \ref{sec:soft} 章），价格更多由自身供给、进口政策与全球定价决定，公司 2026 年上半年实现盈利、半年 ROE 2.8\%（未年化），好于全样本中位数（0.75\%），下行周期对其报表的冲击小于同业。 综合区间形态看，该股单段涨跌幅度相对温和、以震荡为主（最大主升 +41.9\%、最大主跌 -44.7\%），年化波动率 28.9\% 与全样本中位数 37.3\% 相比波动小于样本中位数。

\pfl{盈利情况} 2026 年上半年营业收入 1.96 亿元（同比 -31.1\%），归母净利润 0.19 亿元，同比 -1.9\%。 以年报为趋势对照，2023---2025 年营业收入由 5.48 亿元变为 4.84 亿元（两年年均复合增速 -6.1\%），归母净利润由 0.69 亿元变为 -0.08 亿元。 按半年报口径，2026 年上半年的 ROE 为 2.8\%（未年化）、销售净利率 10.3\%、毛利率 30.5\%，ROE 在“其他种植业”6 家公司中按数值由高到低排第\zhnumber{3}位。 按同一口径，三级行业“其他种植业”6 家公司的半年 ROE 中位数为 1.70\%，该公司高于其中位数 1.11 个百分点。 综合看，公司在报告期维持盈利（高于全样本半年 ROE 中位数 0.75\%），利润的可持续性取决于其他种植业景气的延续与成本管控的执行。

\begin{center}\small\setlength{\tabcolsep}{3.4pt}
\captionof{table}[新农开发（600359）关键财务指标]{新农开发（600359）关键财务指标（合并报表口径；两个半年报期均未年化；现金短债比 $=$ 货币资金 $\div$ 短期债务）}
\begin{adjustbox}{max width=\textwidth}
\begin{tabular}{@{}lrrrrr@{}}
\toprule
指标 & 2023 年 & 2024 年 & 2025 年 & 2025 年半年报 & 2026 年半年报 \\
\midrule
营业收入（亿元） & 5.48 & 6.22 & 4.84 & 2.84 & 1.96 \\
归母净利润（亿元） & 0.69 & 0.22 & -0.08 & 0.19 & 0.19 \\
毛利率（\%） & 26.7 & 23.0 & 15.6 & 22.8 & 30.5 \\
ROE（\%） & 11.3 & 3.3 & -1.2 & 2.8 & 2.8 \\
资产负债率（\%） & 47.0 & 46.3 & 54.1 & 51.8 & 55.6 \\
经营活动现金流净额（亿元） & -0.47 & 0.95 & 0.35 & 0.76 & -0.26 \\
货币资金（亿元） & 2.53 & 3.42 & 2.74 & 3.84 & 3.42 \\
有息负债合计（亿元） & 2.11 & 4.15 & 5.16 & 5.11 & 6.52 \\
现金短债比（倍） & 1.22 & 0.97 & 0.77 & 0.94 & 0.70 \\
\bottomrule
\end{tabular}
\end{adjustbox}
\end{center}

\begin{center}
\begin{minipage}{\textwidth}\centering
\begin{tikzpicture}
\begin{groupplot}[
  group style={group size=2 by 1, horizontal sep=0.60cm},
  width=6.85cm, height=4.60cm,
  tick label style={font=\tiny},
  title style={font=\scriptsize\heiti},
  yticklabel style={font=\tiny},
  ylabel style={font=\tiny},
  grid=major, grid style={LightGray, line width=0.3pt},
  axis line style={InkGray, line width=0.5pt},
  enlarge x limits={abs=0.8},
  legend style={font=\scriptsize, at={(0.5,-0.20)}, anchor=north,
                legend columns=2, draw=none, fill=none, column sep=6pt},
]
\nextgroupplot[
  ybar, bar width=4pt,
  ymin=-0.77, ymax=7.69,
  xtick={0,1,2,3,4,5.7},
  xticklabels={2021,2022,2023,2024,2025,2026H1},
  ylabel={亿元},
]
\addplot[fill=AgriGreen!75, draw=AgriGreen, bar shift=-2.6pt] table[x=i, y=rev]{data/clean/profiles/fin_600359.dat};
\addplot[fill=SoftRed!60, draw=SoftRed, bar shift=2.6pt] table[x=i, y=ni]{data/clean/profiles/fin_600359.dat};
\addplot[fill=AgriGreen!30, draw=AgriGreen, pattern=north east lines, bar shift=-2.6pt] table[x=x, y=rev]{data/clean/profiles/finh_600359.dat};
\addplot[fill=SoftRed!20, draw=SoftRed, pattern=north east lines, bar shift=2.6pt] table[x=x, y=ni]{data/clean/profiles/finh_600359.dat};
\legend{营业收入（年/半年）, 归母净利润（年/半年）}
\nextgroupplot[
  ymin=-7.0, ymax=78.6,
  xtick={0,1,2,3,4,5.7},
  xticklabels={2021,2022,2023,2024,2025,2026H1},
  ylabel={\%},
]
\addplot[AgriGold, mark=*, mark size=1.3pt] table[x=i, y=roe]{data/clean/profiles/fin_600359.dat};
\addplot[OilBlue, mark=square*, mark size=1.3pt] table[x=i, y=debt]{data/clean/profiles/fin_600359.dat};
\addplot[only marks, AgriGold, mark=*, mark size=2.0pt] table[x=x, y=roe]{data/clean/profiles/finh_600359.dat};
\addplot[only marks, OilBlue, mark=square*, mark size=2.0pt] table[x=x, y=debt]{data/clean/profiles/finh_600359.dat};
\legend{ROE, 资产负债率}
\end{groupplot}
\end{tikzpicture}
\begin{tikzpicture}
\begin{axis}[
  name=finbot,
  width=13.75cm, height=3.10cm, scale only axis,
  ymin=-0.76, ymax=8.53,
  xtick={0,1,2,3,4,5.7},
  xticklabels={2021,2022,2023,2024,2025,2026H1},
  tick label style={font=\tiny},
  yticklabel style={font=\tiny},
  ylabel={亿元}, ylabel style={font=\tiny},
  ybar, bar width=4pt,
  grid=major, grid style={LightGray, line width=0.3pt},
  axis line style={InkGray, line width=0.5pt},
  enlarge x limits={abs=0.8},
  legend style={font=\scriptsize, at={(0.5,-0.26)}, anchor=north,
                legend columns=2, draw=none, fill=none, column sep=10pt},
]
\addplot[fill=AgriGreen!55, draw=AgriGreen, bar shift=-3.0pt] table[x=i, y=cash]{data/clean/profiles/fin_600359.dat};
\addplot[fill=SoftRed!45, draw=SoftRed, bar shift=3.0pt] table[x=i, y=ibd]{data/clean/profiles/fin_600359.dat};
\addplot[fill=AgriGreen!25, draw=AgriGreen, pattern=north east lines, bar shift=-3.0pt] table[x=x, y=cash]{data/clean/profiles/finh_600359.dat};
\addplot[fill=SoftRed!20, draw=SoftRed, pattern=north east lines, bar shift=3.0pt] table[x=x, y=ibd]{data/clean/profiles/finh_600359.dat};
\legend{货币资金（左轴）, 有息负债（左轴）}
\end{axis}
\begin{axis}[
  at={(finbot.south east)}, anchor=south east,
  width=13.75cm, height=3.10cm, scale only axis,
  axis y line*=right, axis x line*=none, xtick=\empty,
  ymin=-0.88, ymax=1.62,
  tick label style={font=\tiny},
  yticklabel style={font=\tiny}, scaled y ticks=false,
  legend style={font=\scriptsize, at={(0.5,-0.44)}, anchor=north,
                legend columns=1, draw=none, fill=none},
]
\addplot[OilBlue, mark=*, mark size=2.2pt, line width=1.3pt] table[x=i, y=ocf]{data/clean/profiles/fin_600359.dat};
\addplot[only marks, OilBlue, mark=*, mark size=3.2pt] table[x=x, y=ocf]{data/clean/profiles/finh_600359.dat};
\legend{经营活动现金流净额（右轴，亿元，蓝点）}
\end{axis}
\end{tikzpicture}

\captionof{figure}[新农开发（600359）近年主要财务指标]{新农开发（600359）近年主要财务指标（左上：营业收入与归母净利润；右上：ROE 与资产负债率；下：货币资金、有息负债为左轴柱状，经营活动现金流净额为其右轴的蓝点折线。
2021---2025 年为年报口径，2026H1 为 2026 年半年报/半年末，与其前一年报之间留出间隔、以斜纹或空心点区分；半年报数据未年化）}
\end{minipage}
\end{center}

\pfl{财务分析} 2026 年 6 月末资产负债率 55.6\%，较 2025 年末上升 1.5 个百分点（样本中位数 52.2\%，所处三级行业内按由低到高排第\zhnumber{2}位）。 2026 年 6 月末货币资金 3.42 亿元、短期债务（短期借款与一年内到期非流动负债）4.87 亿元、长期债务 1.66 亿元，有息负债合计 6.52 亿元，现金短债比 0.70 倍——覆盖偏紧。 净有息负债 3.10 亿元，相当于其 2026 年上半年营业收入的 159\%。 流动比率 1.15、速动比率 0.62，短期流动性无明显缺口。 2026 年上半年经营活动现金流净额 -0.26 亿元，净利润为正但经营现金流为负，利润的现金含量偏低。 2026 年 6 月末在建工程 2.27 亿元，较上年末增加 0.45 亿元，仍有在建投入。 综合看，现金短债比 0.70 倍，短期资金安排偏紧；资产负债率高于样本中位数 3.4 个百分点；有息负债中短期债务占比更高，期限结构仍以滚动续作为主。

\pfl{重大战略转型} 公司的战略动作可归为以下几类：产能与项目（1 项）：2026 年度，公司拟投资建设项目 6 个、总投资 4.14 亿元；并购与重组（1 项）：2026 年 6 月，培育 X 类新兴增长产业、剥离非优势产业。 公告口径上，近三年（2023 年 9 月---2026 年 9 月）公司披露重大重组与并购类公告 7 条、对外投资与转型类公告 0 条、回购与股权激励类公告 0 条，合计公告 312 条。 主营结构的口径变化已经落到报表上：2021 年年报的第一大行业为“农业”（55.1\%），2026 年半年报为“服务业及其他”（72.7\%）；公司在这两期的分部划分方式有过调整。 综合判断，报告期内公司有 7 条与战略相关的公告落地（重大重组与并购 7 条、对外投资与转型 0 条），资本运作与主业调整之间存在直接关联。

\begin{center}\small\setlength{\tabcolsep}{4pt}
\captionof{table}[新农开发（600359）历史融资与资本运作]{新农开发（600359）历史融资与资本运作（据公司公告与公开报道整理；金额为披露值，含首次公开发行、再融资、债券、银行借款与重整投资等）}
\begin{adjustbox}{max width=\textwidth}
\begin{tabular}{@{}p{1.45cm}p{2.55cm}p{4.05cm}p{4.95cm}@{}}
\toprule
时间 & 方式 & 规模 & 用途与说明 \\
\midrule
1999-04 & IPO（上海证券交易所，网上定价发行） & 发行 9,\allowbreak{}000 万股，\allowbreak{}发行价 3.86 元/\allowbreak{}股，\allowbreak{}募集资金 34,\allowbreak{}740 万元 & 司 \\
2001-09 & 配股 & 配售 2,\allowbreak{}700 万股 & 经证监公司字[2001]89 号文核准 \\
2008-10 & 非公开发行 A 股股票（预案） & 拟发行 4,\allowbreak{}000 万—8,\allowbreak{}000 万股，\allowbreak{}发行价格不低于 6.82 元/\allowbreak{}股 & 8 万吨棉浆粕生产线项目 1.59 亿元（与山东海龙按 55\%/\allowbreak{}45\% 共同投入）、\allowbreak{}日处理 900 吨鲜奶奶粉加工生产线建设项目 2.34 亿元；预案经第三届董事会第 24 次会议通过 \\
2014-10 & 非公开发行股票 & 6,\allowbreak{}051.28 万股，\allowbreak{}发行价格 9.75 元/\allowbreak{}股（按此计算约 5.90 亿元） & 2014-12 完成验资并出具 XYZH/\allowbreak{}2014URA1028 号验资报告 \\
2023-04 & 重大资产出售（现金交易） & 3.18 亿元 & 向天润乳业出售所持新农乳业 97.4359\% 股权 \\
2025-01 & 现金收购股权（关联交易） & 3,\allowbreak{}944.76 万元 & 收购控股股东农垦集团所持天山雪 100\% 股权；以自有资金分三期支付（80\%/\allowbreak{}10\%/\allowbreak{}10\%）；2025 年半年报披露已支付 80\% 即 3,\allowbreak{}154.76 万元、\allowbreak{}剩余 790.00 万元未支付；评估值 3,\allowbreak{}944.76 万元 \\
\bottomrule
\end{tabular}
\end{adjustbox}
\end{center}

\pfl{历史融投资情况} 从现金流量表看，2025 年公司取得借款收到的现金 5.60 亿元、偿还债务支付的现金 4.59 亿元，筹资活动现金流净额 0.57 亿元（2023 年为取得借款 2.58 亿元、偿还债务 4.56 亿元）。 2026 年上半年筹资活动现金流净额为 1.23 亿元（取得借款 3.34 亿元、偿还债务 2.03 亿元，未年化），仍处于净融入状态。 2025 年分配股利、利润或偿付利息支付的现金合计 0.12 亿元。 公开披露的融资记录共 6 笔，工具类型以 IPO、配股、定向增发为主；金额与用途见下表。 其中规模最大的两笔为：2023 年 4 月，重大资产出售，3.18 亿元；1999 年 4 月，IPO，发行 9,000 万股，发行价 3.86 元/股，募集资金 34,740 万元。 股本方面，近三年披露股本变动 3 次；总股本自 2021 年末以来未发生变化（3.82 亿股）。 分红方面，近三年未实施现金分红，在全部样本中属于分红能力偏弱的一端。 综合看，公司融资结构上股权与债务融资并举；2025 年筹资活动现金流净额 0.57 亿元，年度内仍为净融入，与前述经营与投资现金流的方向基本一致。

\pfl{未来潜在融资需求分析} 2026 年 6 月末货币资金 3.42 亿元，而短期债务 4.87 亿元（现金短债比 0.70 倍），2026 年上半年经营现金流净额 -0.26 亿元。按“短期债务减去账面货币资金”的滚动偿付口径，静态缺口约 1.44 亿元；上半年盈利或接近盈亏平衡，该缺口不随经营亏损进一步扩大。 公司 2026 年 6 月末资产负债率 55.6\%、2026 年上半年 ROE 2.8\%（未年化），现金短债比 0.70 倍、短债覆盖不足，融资需求首先用于置换与滚动存量债务、补充营运资金，属于“补血型”而非扩张型。 公开披露的后续资金安排线索包括：债务与授信安排：2026 年度，公司投资计划中部分项目建设资金来源明确为“企业自筹、银行贷款”；2026 年，第一次临时股东会审议通过 2026 年度贷款计划 2026 年度预计为子公司借款提供担保 2026 年度为控股子公司提供财务资助预计 2026 年度日常关联交易。 资金需求还要叠加在建工程：期末在建工程 2.27 亿元、2026 年上半年购建固定资产等支出 0.29 亿元（未年化），这部分支出在周期底部通常会被主动放缓。 综合看，公司的外部资金需求以营运资金周转与在建项目投入为主，滚动偿付口径下的静态缺口约 1.44 亿元，融资的时点与规模更多取决于信用条件与资本市场窗口。




\subsubsection{雪榕生物（300511）}
\label{co:300511}

\pfl{公司基本情况} 公司前身可追溯至 2016 年，国内产能最大的食用菌工厂化企业，金针菇日产能全国第一；2021—2024 年连亏、2025 年扭亏，2025 年 11 月控股股东与实控人变更为万紫千鸿、蒋智。 公司于 2016-05-04 在深交所创业板首发上市，注册资本 6.41 亿元，总股本 4.99 亿股。 截至 2026 年 10 月 12 日收盘价 4.75 元，对应总市值 23.7 亿元，近一年-26.8\%，总市值在三级行业“食用菌”的 4 家公司中按由大到小排第\zhnumber{2}位，近一年涨跌幅低于全样本中位数（-10.4\%）16.5 个百分点。 按 2025 年年报披露的行业构成，收入占比前两位为食用菌行业 98.9\%；其他(补充) 1.1\%。

\pfl{股权结构与控制权} 按 2026 年 6 月 30 日的股东名册，前三大股东为杨勇萍 持股 14.79\%；上海万紫千鸿智能科技有限公司 持股 4.93\%；余龙生（持股比例未披露），股东户数 29,257 户。 股权结构的另一条线索是，2025 年 11 月万紫千鸿通过协议受让原实际控制人杨勇萍及其一致行动人共 3,158.54 万股股票后取得控制权；第一大股东杨勇萍持股 14.79\%（9,475.61 万股），其股份 2026 年 6 月 5 日被司法冻结。

\pfl{经营要点} 经营与产能方面，截至 2026 年 6 月末食用菌日产能 1,650 吨，其中金针菇日产能 1,365 吨、位居全国之首。

\pfl{所处行业的基本面} 食用菌的供给是工厂化产能（投放周期 6---9 个月），需求是刚性消费，价格随产能快速调整（第 \ref{sec:framework} 章的供需框架）。 与该公司直接相关的行业数据包括：2025 年下半年，至 2026 年一季度金针菇供给偏紧、市场价格回暖；2026 年上半年，万辰集团食用菌板块营收 3.15 亿元、同比增长 32.61\%。

\pfl{竞争优势} 把该公司放回同类公司比较，2026 年上半年营业收入 11.6 亿元、在三级行业“食用菌”的 4 家公司中排第\zhnumber{1}位；毛利率 17.2\%、高于行业内中位数 0.9 个百分点；半年 ROE 1.1\%（未年化）、低于行业内中位数（2.15\%）；资产负债率 56.5\%、高于行业内中位数（44.7\%）。 相对优势方面，是我国产能最大的食用菌工厂化生产企业；2025 年，2026 年自主选育的白色金针菇新菌种 X139、X129 取得发明专利。 就当期报表看，公司尚未把经营层面的特征转化为高于行业的资本回报，周期底部的相对优势更多体现在现金流与资产负债表上。

\pfl{历史股价} 以 2021 年 1 月为起点、2026 年 10 月为终点，股价由 10.83 元变为 4.75 元，累计 -56.1\%，区间最高 11.29 元（2021 年 2 月）、最低 2.59 元（2024 年 6 月）；当前价相当于区间最高价的 42.1\%，为区间最低价的 1.83 倍。 月度收益的年化波动率 41.5\%，高于全样本中位数 37.3\%，在三级行业“食用菌”的 4 家公司中波动率由高到低排第\zhnumber{3}位。 把股价阶段与公司事件对齐看，2024 年 6 月 至 2026 年 4 月股价上涨+210.0\%，同期（2025 年 11 月）万紫千鸿取得控制权的目的被披露为基于对上市公司内在价值和未来发展前景的信心。 整体上看，区间的几处大级别拐点都能在公司的资本运作、股权变动或风险事件上找到对应，股价波动有明显的公司层面驱动，而非单纯的行业 beta。

\begin{center}
\begin{minipage}{\textwidth}\centering
\begin{tikzpicture}
\begin{axis}[
  width=12.6cm, height=3.9cm, scale only axis,
  xmin=2020.922, xmax=2026.972, ymin=2.41, ymax=12.08,
  xtick={2021,2022,2023,2024,2025,2026},
  yearticks,
  yticklabel style={font=\scriptsize},
  ylabel={元/股}, ylabel style={font=\scriptsize},
  grid=major, grid style={LightGray, line width=0.3pt},
  axis line style={InkGray, line width=0.5pt},
  every axis plot/.append style={line width=0.9pt},
]
\addplot[AgriGreen] table[x=x,y=y]{data/clean/profiles/px_300511.dat};
\addplot[SoftRed, dashed, line width=0.7pt] table[x=x,y=y]{data/clean/profiles/pxzz_300511.dat};
\addplot[only marks, mark=*, mark size=1.1pt, SoftRed] table[x=x,y=y]{data/clean/profiles/pxzz_300511.dat};
\end{axis}
\end{tikzpicture}

\captionof{figure}[雪榕生物（300511）股价与周期骨架]{雪榕生物（300511）月度收盘价（元/股）与 ZigZag 周期骨架（阈值 25\%，圆点为周期转折点）}
\end{minipage}
\end{center}

\pfl{过去周期分析} 以 25\% 的 ZigZag 阈值划分，样本区间内股价形成 2 个主升段与 3 个主跌段。 最大主升段出现在 2024 年 6 月 至 2026 年 4 月，22 个月+210.0\%；最大主跌段为 2023 年 7 月 至 2024 年 6 月，11 个月-66.8\%。 最近一次转折出现在 2026 年 10 月（下跌段自 2026 年 4 月起，6 个月-40.8\%），当前处于该段的延续之中。 公司股价较申万农林牧渔指数提前约 23 个月见底，是板块中较早定价周期反转的公司之一。 食用菌是工厂化生产的制造型农业（第 \ref{sec:soft} 章），产能投放节奏与金针菇等单品价格是盈利的主要变量，与养殖周期基本无关，公司 2026 年上半年实现盈利、半年 ROE 1.1\%（未年化），好于全样本中位数（0.75\%），下行周期对其报表的冲击小于同业。 综合区间形态看，该股单段涨跌幅度偏大、以趋势段为主（最大主升 +210.0\%、最大主跌 -66.8\%），年化波动率 41.5\% 与全样本中位数 37.3\% 相比波动与样本中位数接近。

\pfl{盈利情况} 2026 年上半年营业收入 11.56 亿元（同比 +46.2\%），归母净利润 0.18 亿元，实现扭亏。 以年报为趋势对照，2023---2025 年营业收入由 25.65 亿元变为 19.88 亿元（两年年均复合增速 -12.0\%），归母净利润由 -1.88 亿元变为 0.30 亿元。 按半年报口径，2026 年上半年的 ROE 为 1.1\%（未年化）、销售净利率 1.2\%、毛利率 17.2\%，ROE 在“食用菌”4 家公司中按数值由高到低排第\zhnumber{4}位。 同期全样本半年 ROE 中位数 0.75\%，公司与之相差 0.35 个百分点（略好于中位数公司）。 从公司披露的原因看，业绩变动主要来自：2024 年度，公司主动收缩整体业务规模导致产能利用率下降、单位成本升高，并对相关资产及资产组计提减值（跌价）准备。 综合看，公司在报告期维持盈利（高于全样本半年 ROE 中位数 0.75\%），利润的可持续性取决于食用菌景气的延续与成本管控的执行。

\begin{center}\small\setlength{\tabcolsep}{3.4pt}
\captionof{table}[雪榕生物（300511）关键财务指标]{雪榕生物（300511）关键财务指标（合并报表口径；两个半年报期均未年化；现金短债比 $=$ 货币资金 $\div$ 短期债务）}
\begin{adjustbox}{max width=\textwidth}
\begin{tabular}{@{}lrrrrr@{}}
\toprule
指标 & 2023 年 & 2024 年 & 2025 年 & 2025 年半年报 & 2026 年半年报 \\
\midrule
营业收入（亿元） & 25.65 & 21.62 & 19.88 & 7.91 & 11.56 \\
归母净利润（亿元） & -1.88 & -6.17 & 0.30 & -1.03 & 0.18 \\
毛利率（\%） & 8.4 & 3.6 & 13.6 & 1.2 & 17.2 \\
ROE（\%） & -11.4 & -49.5 & 2.3 & -9.8 & 1.1 \\
资产负债率（\%） & 68.9 & 78.6 & 58.0 & 67.5 & 56.5 \\
经营活动现金流净额（亿元） & 3.58 & 4.37 & 4.16 & 1.38 & 2.21 \\
货币资金（亿元） & 2.53 & 1.68 & 1.61 & 0.76 & 1.76 \\
有息负债合计（亿元） & 22.37 & 21.52 & 12.20 & 15.21 & 12.52 \\
现金短债比（倍） & 0.20 & 0.14 & 0.16 & 0.11 & 0.17 \\
\bottomrule
\end{tabular}
\end{adjustbox}
\end{center}

\begin{center}
\begin{minipage}{\textwidth}\centering
\begin{tikzpicture}
\begin{groupplot}[
  group style={group size=2 by 1, horizontal sep=0.60cm},
  width=6.85cm, height=4.60cm,
  tick label style={font=\tiny},
  title style={font=\scriptsize\heiti},
  yticklabel style={font=\tiny},
  ylabel style={font=\tiny},
  grid=major, grid style={LightGray, line width=0.3pt},
  axis line style={InkGray, line width=0.5pt},
  enlarge x limits={abs=0.8},
  legend style={font=\scriptsize, at={(0.5,-0.20)}, anchor=north,
                legend columns=2, draw=none, fill=none, column sep=6pt},
]
\nextgroupplot[
  ybar, bar width=4pt,
  ymin=-9.35, ymax=29.47,
  xtick={0,1,2,3,4,5.7},
  xticklabels={2021,2022,2023,2024,2025,2026H1},
  ylabel={亿元},
]
\addplot[fill=AgriGreen!75, draw=AgriGreen, bar shift=-2.6pt] table[x=i, y=rev]{data/clean/profiles/fin_300511.dat};
\addplot[fill=SoftRed!60, draw=SoftRed, bar shift=2.6pt] table[x=i, y=ni]{data/clean/profiles/fin_300511.dat};
\addplot[fill=AgriGreen!30, draw=AgriGreen, pattern=north east lines, bar shift=-2.6pt] table[x=x, y=rev]{data/clean/profiles/finh_300511.dat};
\addplot[fill=SoftRed!20, draw=SoftRed, pattern=north east lines, bar shift=2.6pt] table[x=x, y=ni]{data/clean/profiles/finh_300511.dat};
\legend{营业收入（年/半年）, 归母净利润（年/半年）}
\nextgroupplot[
  ymin=-59.8, ymax=91.5,
  xtick={0,1,2,3,4,5.7},
  xticklabels={2021,2022,2023,2024,2025,2026H1},
  ylabel={\%},
]
\addplot[AgriGold, mark=*, mark size=1.3pt] table[x=i, y=roe]{data/clean/profiles/fin_300511.dat};
\addplot[OilBlue, mark=square*, mark size=1.3pt] table[x=i, y=debt]{data/clean/profiles/fin_300511.dat};
\addplot[only marks, AgriGold, mark=*, mark size=2.0pt] table[x=x, y=roe]{data/clean/profiles/finh_300511.dat};
\addplot[only marks, OilBlue, mark=square*, mark size=2.0pt] table[x=x, y=debt]{data/clean/profiles/finh_300511.dat};
\legend{ROE, 资产负债率}
\end{groupplot}
\end{tikzpicture}
\begin{tikzpicture}
\begin{axis}[
  name=finbot,
  width=13.75cm, height=3.10cm, scale only axis,
  ymin=-2.24, ymax=25.06,
  xtick={0,1,2,3,4,5.7},
  xticklabels={2021,2022,2023,2024,2025,2026H1},
  tick label style={font=\tiny},
  yticklabel style={font=\tiny},
  ylabel={亿元}, ylabel style={font=\tiny},
  ybar, bar width=4pt,
  grid=major, grid style={LightGray, line width=0.3pt},
  axis line style={InkGray, line width=0.5pt},
  enlarge x limits={abs=0.8},
  legend style={font=\scriptsize, at={(0.5,-0.26)}, anchor=north,
                legend columns=2, draw=none, fill=none, column sep=10pt},
]
\addplot[fill=AgriGreen!55, draw=AgriGreen, bar shift=-3.0pt] table[x=i, y=cash]{data/clean/profiles/fin_300511.dat};
\addplot[fill=SoftRed!45, draw=SoftRed, bar shift=3.0pt] table[x=i, y=ibd]{data/clean/profiles/fin_300511.dat};
\addplot[fill=AgriGreen!25, draw=AgriGreen, pattern=north east lines, bar shift=-3.0pt] table[x=x, y=cash]{data/clean/profiles/finh_300511.dat};
\addplot[fill=SoftRed!20, draw=SoftRed, pattern=north east lines, bar shift=3.0pt] table[x=x, y=ibd]{data/clean/profiles/finh_300511.dat};
\legend{货币资金（左轴）, 有息负债（左轴）}
\end{axis}
\begin{axis}[
  at={(finbot.south east)}, anchor=south east,
  width=13.75cm, height=3.10cm, scale only axis,
  axis y line*=right, axis x line*=none, xtick=\empty,
  ymin=-1.14, ymax=5.68,
  tick label style={font=\tiny},
  yticklabel style={font=\tiny}, scaled y ticks=false,
  legend style={font=\scriptsize, at={(0.5,-0.44)}, anchor=north,
                legend columns=1, draw=none, fill=none},
]
\addplot[OilBlue, mark=*, mark size=2.2pt, line width=1.3pt] table[x=i, y=ocf]{data/clean/profiles/fin_300511.dat};
\addplot[only marks, OilBlue, mark=*, mark size=3.2pt] table[x=x, y=ocf]{data/clean/profiles/finh_300511.dat};
\legend{经营活动现金流净额（右轴，亿元，蓝点）}
\end{axis}
\end{tikzpicture}

\captionof{figure}[雪榕生物（300511）近年主要财务指标]{雪榕生物（300511）近年主要财务指标（左上：营业收入与归母净利润；右上：ROE 与资产负债率；下：货币资金、有息负债为左轴柱状，经营活动现金流净额为其右轴的蓝点折线。
2021---2025 年为年报口径，2026H1 为 2026 年半年报/半年末，与其前一年报之间留出间隔、以斜纹或空心点区分；半年报数据未年化）}
\end{minipage}
\end{center}

\pfl{财务分析} 2026 年 6 月末资产负债率 56.5\%，较 2025 年末下降 1.4 个百分点（样本中位数 52.2\%，所处三级行业内按由低到高排第\zhnumber{4}位）。 2026 年 6 月末货币资金 1.76 亿元、短期债务（短期借款与一年内到期非流动负债）10.20 亿元、长期债务 2.32 亿元，有息负债合计 12.52 亿元，现金短债比 0.17 倍——缺口明显。 净有息负债 10.76 亿元，相当于其 2026 年上半年营业收入的 93\%。 流动比率 0.37、速动比率 0.15，短期偿付压力较大。 2026 年上半年经营活动现金流净额 2.21 亿元，经营现金流与利润同向为正，内生资金可以支撑运营。 2026 年 6 月末在建工程 0.70 亿元，较上年末增加 0.50 亿元，仍有在建投入。 综合看，账面货币资金对有息负债与短期债务的覆盖不足，滚动偿付对新增融资的依赖度较高；资产负债率高于样本中位数 4.4 个百分点；有息负债中短期债务占比更高，期限结构仍以滚动续作为主。 公司披露的股东与质押风险为：2026 年 9 月，控股股东上海万紫千鸿智能科技有限公司质押公司 1,900 万股。

\pfl{重大战略转型} 公司的战略动作可归为以下几类：产能与项目（2 项）：2025 年 12 月，公司实现扭亏为盈；2024 年，主动收缩整体业务规模以应对行业产能过剩与价格低迷；融资与资本运作（1 项）：2025 年 11 月，万紫千鸿取得控制权的目的被披露为基于对上市公司内在价值和未来发展前景的信心。 公告口径上，近三年（2023 年 9 月---2026 年 9 月）公司披露重大重组与并购类公告 0 条、对外投资与转型类公告 5 条、回购与股权激励类公告 17 条，合计公告 520 条。 第一大行业仍是“食用菌行业”，收入占比 98.9\%，与 2021 年年报相比没有方向性变化。 综合判断，报告期内公司有 5 条与战略相关的公告落地（重大重组与并购 0 条、对外投资与转型 5 条），资本运作与主业调整之间存在直接关联。

\begin{center}\small\setlength{\tabcolsep}{4pt}
\captionof{table}[雪榕生物（300511）历史融资与资本运作]{雪榕生物（300511）历史融资与资本运作（据公司公告与公开报道整理；金额为披露值，含首次公开发行、再融资、债券、银行借款与重整投资等）}
\begin{adjustbox}{max width=\textwidth}
\begin{tabular}{@{}p{1.45cm}p{2.55cm}p{4.05cm}p{4.95cm}@{}}
\toprule
时间 & 方式 & 规模 & 用途与说明 \\
\midrule
2016-05 & IPO（深圳证券交易所创业板） & 发行 3,\allowbreak{}750 万股，\allowbreak{}发行价 16.82 元/\allowbreak{}股，\allowbreak{}募集资金总额 6.31 亿元 & 承销保荐费用 4,\allowbreak{}084.50 万元 \\
2020-06-24 & 公开发行可转换公司债券“雪榕转债”（债券代码 123056） & 85.00 亿元（585 万张，\allowbreak{}每张面值 100.00 元），\allowbreak{}发行信用等级 AA- & 含泰国食用菌工厂化生产车间建设项目等及补充流动资金；2020-10-13 公告变更部分募集资金用途，\allowbreak{}初始转股价 11.89 元/\allowbreak{}股 \\
2025-03 & 向特定对象发行股票（定增，预案） & 发行价格 3.00 元/\allowbreak{}股；若“雪榕转债”未转股 & 扣除发行费用后全部用于补充流动资金；发行对象为控股股东万紫千鸿（现金全额认购） \\
2025-09-03 & 可转债提前赎回 & 赎回价格 100.93 元/\allowbreak{}张（含税）；截至赎回登记日（2025-10-14）收市后 5,\allowbreak{}823,\allowbreak{}234 张已转股、\allowbreak{}剩余 15,\allowbreak{}151 张被赎回 & 终止“雪榕转债”存续 \\
2026-03 & 2026 年限制性股票激励计划（第二类限制性股票） & 拟授予 2,\allowbreak{}704 万股，\allowbreak{}占当时总股本比例 4.22\%，\allowbreak{}每股转让价 6.26 元 & 向激励对象定向发行 A 股普通股股票；2026-03-30 披露向激励对象授予的公告；股权激励计划草案 2026-02-11 披露 \\
\bottomrule
\end{tabular}
\end{adjustbox}
\end{center}

\pfl{历史融投资情况} 从现金流量表看，2025 年公司取得借款收到的现金 4.54 亿元、偿还债务支付的现金 7.79 亿元、吸收权益性投资 0.01 亿元，筹资活动现金流净额 -3.58 亿元（2023 年为取得借款 12.01 亿元、偿还债务 8.96 亿元）。 2026 年上半年筹资活动现金流净额为 -0.90 亿元（取得借款 2.43 亿元、偿还债务 2.81 亿元，未年化），已转为净流出，处于净还债状态。 2025 年分配股利、利润或偿付利息支付的现金合计 0.43 亿元。 公开披露的融资记录共 5 笔，工具类型以 IPO、可转债、定向增发为主；金额与用途见下表。 其中规模最大的两笔为：2025 年 3 月，定向增发，发行价格 3.00 元/股，若“雪榕转债”未转股；2020 年 6 月，可转债，85.00 亿元（585 万张，每张面值 100.00 元），发行信用等级 AA-。 股本方面，近三年披露股本变动 13 次（可转债转股 9 次、其他 1 次、股份回购 1 次）；总股本由 2021 年末的 4.42 亿股变为 4.99 亿股（+12.87\%）。 分红方面，近三年未实施现金分红，在全部样本中属于分红能力偏弱的一端。 近三年“再融资”类公告 29 条，涉及定向增发、可转债、募集资金使用。 综合看，公司融资结构上以债务滚动为主、股权融资为补充；2025 年筹资活动现金流净额 -3.58 亿元，融资节奏仍以维持存量债务周转为主，与前述经营与投资现金流的方向基本一致。

\pfl{未来潜在融资需求分析} 2026 年 6 月末货币资金 1.76 亿元，而短期债务 10.20 亿元（现金短债比 0.17 倍），2026 年上半年经营现金流净额 2.21 亿元。按“短期债务减去账面货币资金”的滚动偿付口径，静态缺口约 8.43 亿元；上半年盈利或接近盈亏平衡，该缺口不随经营亏损进一步扩大。 公司 2026 年 6 月末资产负债率 56.5\%、2026 年上半年 ROE 1.1\%（未年化），现金短债比 0.17 倍、短债覆盖不足，融资需求首先用于置换与滚动存量债务、补充营运资金，属于“补血型”而非扩张型。 公开披露的后续资金安排线索包括：股权融资线索：2025 年度，公司向特定对象发行股票预案：拟向控股股东万紫千鸿发行不超过 1.50 亿股（雪榕转债未转股情形）、发行价 3.00 元/股、募集资金总额不超过 4.49 亿元。 资金需求还要叠加在建工程：期末在建工程 0.70 亿元、2026 年上半年购建固定资产等支出 1.40 亿元（未年化），这部分支出在周期底部通常会被主动放缓。 综合看，公司的外部资金需求以营运资金周转与在建项目投入为主，滚动偿付口径下的静态缺口约 8.43 亿元，能否落地取决于授信续作与发行窗口的配合。




\subsubsection{华绿生物（300970）}
\label{co:300970}

\pfl{公司基本情况} 公司前身可追溯至 2018 年，江苏泗阳的食用菌工厂化企业，以金针菇、真姬菇为主并拓展鹿茸菇等珍稀品种；2021 年创业板上市，2024 年由盈转亏、2025 年随行业回暖大幅回升。 公司于 2021-04-12 在深交所创业板首发上市，注册资本 1.60 亿元，总股本 1.18 亿股。 截至 2026 年 10 月 12 日收盘价 17.00 元，对应总市值 20.0 亿元，近一年+13.8\%，总市值在三级行业“食用菌”的 4 家公司中按由大到小排第\zhnumber{3}位，近一年涨跌幅高于全样本中位数（-10.4\%）24.1 个百分点。 按 2025 年年报披露的行业构成，收入几乎全部来自食用菌 100.0\%。

\pfl{股权结构与控制权} 按 2026 年 6 月 30 日的股东名册，前三大股东为余养朝（持股比例未披露）；阮秀莲（持股比例未披露）；南京极星私募基金管理有限公司－极星北斗 1 号私募证券投（持股比例未披露），股东户数 13,146 户。 从控制权角度看，控股股东、实际控制人余养朝持股 30.04\%（4,766.79 万股），其中 3,575.09 万股为有限售条件股份；前十大股东还包括南京极星私募基金－极星北斗 1 号私募证券投资基金持股 2.46\%、殷宏伟持股 1.55\%；上半年经营活动现金流净额 1.65 亿元、同比增长 249.34\%。

\pfl{经营要点} 经营层面，2024 年下半年，浙江华绿项目金针菇工厂、广西华绿项目一期二厂真姬菇工厂、河南华绿一期金针菇工厂投产，处于产能爬坡与磨合阶段。

\pfl{所处行业的基本面} 食用菌的供给是工厂化产能（投放周期 6---9 个月），需求是刚性消费，价格随产能快速调整（第 \ref{sec:framework} 章的供需框架）。 与该公司直接相关的行业数据包括：2026 年上半年，食用菌行业景气上行源于供给端产能出清与消费端需求回暖双向共振；食用菌行业竞争已从比拼产能、打价格战转向比拼育种技术、产品品类与综合运营能力。

\pfl{竞争优势} 把该公司放回同类公司比较，2026 年上半年营业收入 6.9 亿元、在三级行业“食用菌”的 4 家公司中排第\zhnumber{3}位；毛利率 15.4\%、低于行业内中位数 0.9 个百分点；半年 ROE 2.8\%（未年化）、高于行业内中位数（2.15\%）；资产负债率 35.4\%、低于行业内中位数（44.7\%）。 相对优势方面，2026 年上半年研发费用 897.00 万元、同比增长 50.31\%；2025 年，鹿茸菇收入 0.49 亿元、同比增长 2,536.89\%。 优势在当期体现为高于行业中位数的资本回报，但优势能否持续仍取决于周期与自身产能投放节奏。

\pfl{历史股价} 自 2021 年 4 月至 2026 年 10 月，股价由 57.00 元变为 17.00 元，累计 -70.2\%，区间最高 63.31 元（2021 年 5 月）、最低 11.23 元（2024 年 8 月）；当前价相当于区间最高价的 26.9\%，为区间最低价的 1.51 倍。 月度收益的年化波动率 65.0\%，高于全样本中位数 37.3\%，在三级行业“食用菌”的 4 家公司中波动率由高到低排第\zhnumber{1}位。

\begin{center}
\begin{minipage}{\textwidth}\centering
\begin{tikzpicture}
\begin{axis}[
  width=12.6cm, height=3.9cm, scale only axis,
  xmin=2021.172, xmax=2026.972, ymin=10.44, ymax=67.74,
  xtick={2021,2022,2023,2024,2025,2026},
  yearticks,
  yticklabel style={font=\scriptsize},
  ylabel={元/股}, ylabel style={font=\scriptsize},
  grid=major, grid style={LightGray, line width=0.3pt},
  axis line style={InkGray, line width=0.5pt},
  every axis plot/.append style={line width=0.9pt},
]
\addplot[AgriGreen] table[x=x,y=y]{data/clean/profiles/px_300970.dat};
\addplot[SoftRed, dashed, line width=0.7pt] table[x=x,y=y]{data/clean/profiles/pxzz_300970.dat};
\addplot[only marks, mark=*, mark size=1.1pt, SoftRed] table[x=x,y=y]{data/clean/profiles/pxzz_300970.dat};
\end{axis}
\end{tikzpicture}

\captionof{figure}[华绿生物（300970）股价与周期骨架]{华绿生物（300970）月度收盘价（元/股）与 ZigZag 周期骨架（阈值 25\%，圆点为周期转折点）}
\end{minipage}
\end{center}

\pfl{过去周期分析} 以 25\% 的 ZigZag 阈值划分，样本区间内股价形成 3 个主升段与 4 个主跌段。 最大主跌段为 2021 年 4 月 至 2021 年 7 月，3 个月-63.9\%；最大主升段出现在 2024 年 8 月 至 2026 年 4 月，20 个月+214.3\%。 最近一次转折出现在 2026 年 7 月（下跌段自 2026 年 4 月起，3 个月-54.3\%），当前处于该段的延续之中。 公司股价较申万农林牧渔指数提前约 21 个月见底，是板块中较早定价周期反转的公司之一。 食用菌是工厂化生产的制造型农业（第 \ref{sec:soft} 章），产能投放节奏与金针菇等单品价格是盈利的主要变量，与养殖周期基本无关，公司 2026 年上半年实现盈利、半年 ROE 2.8\%（未年化），好于全样本中位数（0.75\%），下行周期对其报表的冲击小于同业。 综合区间形态看，该股单段涨跌幅度偏大、以趋势段为主（最大主升 +214.3\%、最大主跌 -63.9\%），年化波动率 65.0\% 与全样本中位数 37.3\% 相比波动明显大于样本中位数。

\pfl{盈利情况} 2026 年上半年营业收入 6.94 亿元（同比 +39.2\%），归母净利润 0.46 亿元，实现扭亏。 以年报为趋势对照，2023---2025 年营业收入由 9.96 亿元变为 12.92 亿元（两年年均复合增速 +13.9\%），归母净利润由 0.30 亿元变为 1.10 亿元。 按半年报口径，2026 年上半年的 ROE 为 2.8\%（未年化）、销售净利率 8.0\%、毛利率 15.4\%，ROE 在“食用菌”4 家公司中按数值由高到低排第\zhnumber{2}位。 按同一口径，三级行业“食用菌”4 家公司的半年 ROE 中位数为 2.15\%，该公司高于其中位数 0.64 个百分点。 从公司披露的原因看，业绩变动主要来自：2024 年，由盈转亏：主要产品金针菇整体销售价格较上年同期下降较多、部分项目建成转固导致折旧增加、2022/2023 两期股权激励计划导致报告期累计计提期间费用约 2,000.00 万元；2026 年上半年，营业成本 5.87 亿元。 综合看，公司在报告期维持盈利（高于全样本半年 ROE 中位数 0.75\%），利润的可持续性取决于食用菌景气的延续与成本管控的执行。

\begin{center}\small\setlength{\tabcolsep}{3.4pt}
\captionof{table}[华绿生物（300970）关键财务指标]{华绿生物（300970）关键财务指标（合并报表口径；两个半年报期均未年化；现金短债比 $=$ 货币资金 $\div$ 短期债务）}
\begin{adjustbox}{max width=\textwidth}
\begin{tabular}{@{}lrrrrr@{}}
\toprule
指标 & 2023 年 & 2024 年 & 2025 年 & 2025 年半年报 & 2026 年半年报 \\
\midrule
营业收入（亿元） & 9.96 & 10.32 & 12.92 & 4.99 & 6.94 \\
归母净利润（亿元） & 0.30 & -0.48 & 1.10 & -0.54 & 0.46 \\
毛利率（\%） & 13.5 & 6.0 & 16.1 & -3.4 & 15.4 \\
ROE（\%） & 2.0 & -3.2 & 7.1 & -3.7 & 2.8 \\
资产负债率（\%） & 26.0 & 34.1 & 32.3 & 34.7 & 35.4 \\
经营活动现金流净额（亿元） & 1.67 & 1.38 & 3.45 & 0.47 & 1.65 \\
货币资金（亿元） & 4.10 & 3.61 & 5.90 & 3.52 & 5.89 \\
有息负债合计（亿元） & 3.33 & 5.66 & 5.72 & 5.72 & 7.00 \\
现金短债比（倍） & 6.42 & 7.39 & 8.30 & 13.68 & 19.04 \\
\bottomrule
\end{tabular}
\end{adjustbox}
\end{center}

\begin{center}
\begin{minipage}{\textwidth}\centering
\begin{tikzpicture}
\begin{groupplot}[
  group style={group size=2 by 1, horizontal sep=0.60cm},
  width=6.85cm, height=4.60cm,
  tick label style={font=\tiny},
  title style={font=\scriptsize\heiti},
  yticklabel style={font=\tiny},
  ylabel style={font=\tiny},
  grid=major, grid style={LightGray, line width=0.3pt},
  axis line style={InkGray, line width=0.5pt},
  enlarge x limits={abs=0.8},
  legend style={font=\scriptsize, at={(0.5,-0.20)}, anchor=north,
                legend columns=2, draw=none, fill=none, column sep=6pt},
]
\nextgroupplot[
  ybar, bar width=4pt,
  ymin=-1.82, ymax=14.53,
  xtick={0,1,2,3,4,5.7},
  xticklabels={2021,2022,2023,2024,2025,2026H1},
  ylabel={亿元},
]
\addplot[fill=AgriGreen!75, draw=AgriGreen, bar shift=-2.6pt] table[x=i, y=rev]{data/clean/profiles/fin_300970.dat};
\addplot[fill=SoftRed!60, draw=SoftRed, bar shift=2.6pt] table[x=i, y=ni]{data/clean/profiles/fin_300970.dat};
\addplot[fill=AgriGreen!30, draw=AgriGreen, pattern=north east lines, bar shift=-2.6pt] table[x=x, y=rev]{data/clean/profiles/finh_300970.dat};
\addplot[fill=SoftRed!20, draw=SoftRed, pattern=north east lines, bar shift=2.6pt] table[x=x, y=ni]{data/clean/profiles/finh_300970.dat};
\legend{营业收入（年/半年）, 归母净利润（年/半年）}
\nextgroupplot[
  ymin=-6.3, ymax=39.3,
  xtick={0,1,2,3,4,5.7},
  xticklabels={2021,2022,2023,2024,2025,2026H1},
  ylabel={\%},
]
\addplot[AgriGold, mark=*, mark size=1.3pt] table[x=i, y=roe]{data/clean/profiles/fin_300970.dat};
\addplot[OilBlue, mark=square*, mark size=1.3pt] table[x=i, y=debt]{data/clean/profiles/fin_300970.dat};
\addplot[only marks, AgriGold, mark=*, mark size=2.0pt] table[x=x, y=roe]{data/clean/profiles/finh_300970.dat};
\addplot[only marks, OilBlue, mark=square*, mark size=2.0pt] table[x=x, y=debt]{data/clean/profiles/finh_300970.dat};
\legend{ROE, 资产负债率}
\end{groupplot}
\end{tikzpicture}
\begin{tikzpicture}
\begin{axis}[
  name=finbot,
  width=13.75cm, height=3.10cm, scale only axis,
  ymin=-0.70, ymax=7.84,
  xtick={0,1,2,3,4,5.7},
  xticklabels={2021,2022,2023,2024,2025,2026H1},
  tick label style={font=\tiny},
  yticklabel style={font=\tiny},
  ylabel={亿元}, ylabel style={font=\tiny},
  ybar, bar width=4pt,
  grid=major, grid style={LightGray, line width=0.3pt},
  axis line style={InkGray, line width=0.5pt},
  enlarge x limits={abs=0.8},
  legend style={font=\scriptsize, at={(0.5,-0.26)}, anchor=north,
                legend columns=2, draw=none, fill=none, column sep=10pt},
]
\addplot[fill=AgriGreen!55, draw=AgriGreen, bar shift=-3.0pt] table[x=i, y=cash]{data/clean/profiles/fin_300970.dat};
\addplot[fill=SoftRed!45, draw=SoftRed, bar shift=3.0pt] table[x=i, y=ibd]{data/clean/profiles/fin_300970.dat};
\addplot[fill=AgriGreen!25, draw=AgriGreen, pattern=north east lines, bar shift=-3.0pt] table[x=x, y=cash]{data/clean/profiles/finh_300970.dat};
\addplot[fill=SoftRed!20, draw=SoftRed, pattern=north east lines, bar shift=3.0pt] table[x=x, y=ibd]{data/clean/profiles/finh_300970.dat};
\legend{货币资金（左轴）, 有息负债（左轴）}
\end{axis}
\begin{axis}[
  at={(finbot.south east)}, anchor=south east,
  width=13.75cm, height=3.10cm, scale only axis,
  axis y line*=right, axis x line*=none, xtick=\empty,
  ymin=-0.90, ymax=4.49,
  tick label style={font=\tiny},
  yticklabel style={font=\tiny}, scaled y ticks=false,
  legend style={font=\scriptsize, at={(0.5,-0.44)}, anchor=north,
                legend columns=1, draw=none, fill=none},
]
\addplot[OilBlue, mark=*, mark size=2.2pt, line width=1.3pt] table[x=i, y=ocf]{data/clean/profiles/fin_300970.dat};
\addplot[only marks, OilBlue, mark=*, mark size=3.2pt] table[x=x, y=ocf]{data/clean/profiles/finh_300970.dat};
\legend{经营活动现金流净额（右轴，亿元，蓝点）}
\end{axis}
\end{tikzpicture}

\captionof{figure}[华绿生物（300970）近年主要财务指标]{华绿生物（300970）近年主要财务指标（左上：营业收入与归母净利润；右上：ROE 与资产负债率；下：货币资金、有息负债为左轴柱状，经营活动现金流净额为其右轴的蓝点折线。
2021---2025 年为年报口径，2026H1 为 2026 年半年报/半年末，与其前一年报之间留出间隔、以斜纹或空心点区分；半年报数据未年化）}
\end{minipage}
\end{center}

\pfl{财务分析} 2026 年 6 月末资产负债率 35.4\%，较 2025 年末上升 3.1 个百分点（样本中位数 52.2\%，所处三级行业内按由低到高排第\zhnumber{2}位）。 2026 年 6 月末货币资金 5.89 亿元、短期债务（短期借款与一年内到期非流动负债）0.31 亿元、长期债务 6.69 亿元，有息负债合计 7.00 亿元，现金短债比 19.04 倍——覆盖充分。 净有息负债 1.11 亿元，相当于其 2026 年上半年营业收入的 16\%。 流动比率 4.21、速动比率 2.99，短期偿债能力处于正常区间。 2026 年上半年经营活动现金流净额 1.65 亿元，经营现金流与利润同向为正，内生资金可以支撑运营。 综合看，现金短债比 19.04 倍，短期偿付压力可控；资产负债率低于样本中位数 16.8 个百分点；有息负债中长期债务占比更高，期限结构对短债滚动的压力较小。

\pfl{重大战略转型} 近三年公司的战略动作集中在：产能与项目（3 项）：2026 年 4 月，媒体披露公司年内合计投资 5.00 亿元加码高端产能；2026 年 3 月，公司拟与文山市人民政府签订投资协议，在云南省文山壮族苗族自治州文山市新建日产 50 吨鲜品黄牛肝菌（见手青）工厂化生产项目。 公告口径上，近三年（2023 年 9 月---2026 年 9 月）公司披露重大重组与并购类公告 1 条、对外投资与转型类公告 5 条、回购与股权激励类公告 69 条，合计公告 311 条。 第一大行业仍是“食用菌”，收入占比 100.0\%，与 2021 年年报相比没有方向性变化。 综合判断，报告期内公司有 6 条与战略相关的公告落地（重大重组与并购 1 条、对外投资与转型 5 条），资本运作与主业调整之间存在直接关联。

\begin{center}\small\setlength{\tabcolsep}{4pt}
\captionof{table}[华绿生物（300970）历史融资与资本运作]{华绿生物（300970）历史融资与资本运作（据公司公告与公开报道整理；金额为披露值，含首次公开发行、再融资、债券、银行借款与重整投资等）}
\begin{adjustbox}{max width=\textwidth}
\begin{tabular}{@{}p{1.45cm}p{2.55cm}p{4.05cm}p{4.95cm}@{}}
\toprule
时间 & 方式 & 规模 & 用途与说明 \\
\midrule
2021-04 & IPO（深圳证券交易所创业板） & 发行 1,\allowbreak{}459.00 万股，\allowbreak{}发行价 44.77 元/\allowbreak{}股，\allowbreak{}募集资金总额 6.53 亿元 & 江苏华绿生物科技股份有限公司一厂技术改造项目、\allowbreak{}泗阳华茂农业发展有限公司技术改造项目、\allowbreak{}年产 3 万吨真姬菇工厂化栽培建设项目；；计入资本公积（股本溢价）5.82 亿元 \\
2022-09-30 & 终止募投项目并将剩余募集资金永久补充流动资金 & 剩余募集资金 2.87 亿元及利息 & 终止一厂技术改造项目、\allowbreak{}泗阳华茂农业发展有限公司技术改造项目 \\
2022/\allowbreak{}2023 & 两期股权激励计划（第二类限制性股票，股票来源为向激励对象定向增发） & 2022 年计划以 10.68 元/\allowbreak{}股向 58 名激励对象授予 400.00 万股（授予日 2022-09-30）；2023 年计划以 11.46 元/\allowbreak{}股向 41 名激励对象授予 329.80 万股、\allowbreak{}预留 70.20 万股（首次授予日 2023-08-21） & 员工激励；2023-10-31 第一期 118.86 万股归属上市；2024-10-22 第二期 142.52 万股归属上市 \\
2023-11-30 & 募投项目结项并将节余募集资金永久补充流动资金 & 节余募集资金 856.19 万元（截至 2023-11-30） & 河北华绿之珍年产 5.4 万吨鲜品金针菇工厂化生产项目结项 \\
2024-03-04 & 股份回购 & 计划使用自有资金不低于 3,\allowbreak{}000.00 万元、\allowbreak{}不超过 6,\allowbreak{}000.00 万元；截至 2024-06-30 累计回购 234.66 万股（占总股本 1.95\%） & 在适宜时机用于实施股权激励计划或员工持股计划；2024 年年度报告列示回购股份金额 3,\allowbreak{}001.27 万元；2025 年年度报告以现有总股本 1.23 亿股剔除回购专用账户 234.66 万股为分红基数 \\
2025-04-22 & 2025 年度利润分配预案（含转增） & 以 1.20 亿股为基数，\allowbreak{}每 10 股派发现金红利 5.00 元（含税）、\allowbreak{}送红股 0 股 & 股东回报；2026-04-22 董事会审议通过（公司现有总股本 1.23 亿股 \\
\bottomrule
\end{tabular}
\end{adjustbox}
\end{center}

\pfl{历史融投资情况} 从现金流量表看，2025 年公司取得借款收到的现金 1.49 亿元、偿还债务支付的现金 1.51 亿元、吸收权益性投资 0.26 亿元，筹资活动现金流净额 0.07 亿元（2023 年为取得借款 0.75 亿元、偿还债务 1.30 亿元）。 2026 年上半年筹资活动现金流净额为 -1.16 亿元（取得借款 0.09 亿元、偿还债务 0.44 亿元，未年化），已转为净流出，处于净还债状态。 2025 年分配股利、利润或偿付利息支付的现金合计 0.13 亿元。 公开披露的融资记录共 6 笔，工具类型以 IPO 为主；金额与用途见下表。 其中规模最大的两笔为：2023 年 11 月，募投项目结项并将节余，节余募集资金 856.19 万元（截至 2023-11-30）；2021 年 4 月，IPO，发行 1,459.00 万股，发行价 44.77 元/股，募集资金总额 6.53 亿元。 股本方面，近三年披露股本变动 5 次（增发新股上市 1 次、股权激励 1 次、限售股份上市 1 次）；总股本由 2021 年末的 1.17 亿股变为 1.18 亿股（+1.02\%）。 分红方面，近三年现金分红 3 次、合计每股派现 0.90 元，最近一次实施在 2025 年 12 月。 综合看，公司融资结构上股权与债务融资并举；2025 年筹资活动现金流净额 0.07 亿元，年度内仍为净融入，与前述经营与投资现金流的方向基本一致。

\pfl{未来潜在融资需求分析} 2026 年 6 月末货币资金 5.89 亿元，而短期债务 0.31 亿元（现金短债比 19.04 倍），2026 年上半年经营现金流净额 1.65 亿元。账面货币资金可以覆盖短期债务，缺口不体现在静态偿付层面。 公司 2026 年 6 月末资产负债率 35.4\%、2026 年上半年 ROE 2.8\%（未年化），资产端与负债端都没有明显压力，融资需求以相机抉择的扩张型用途为主。 公开披露的后续资金安排线索包括：债务与授信安排：2026 年 4 月，公司对子公司担保：为南川华绿提供担保额度 1.40 亿元（实际担保余额 9,895.02 万元）、为广西华绿提供担保额度 1.30 亿元（实际担保余额 9,753.75 万元）；2026 年 6 月，在建项目资金来源披露为金融机构贷款：南川鹿茸菇项目预算 5,004.54 万元、工程进度 95.00\%。 静态偿付不存在缺口，后续融资更可能服务于产能或并购安排，而不是填补流动性窟窿。




\subsubsection{博闻科技（600883）}
\label{co:600883}

\pfl{公司基本情况} 博闻科技成立于 1990 年，云南保山起家的老牌上市公司，主业历经水泥—科技—农副食品的多次更迭，现主营云南高原特色农副食品加工（食用菌、火腿、咖啡），并持有新疆众和参股权益与部分证券资产。 公司于 1995-12-08 在上交所首发上市，注册资本 2.36 亿元，总股本 2.36 亿股。 截至 2026 年 10 月 12 日收盘价 7.88 元，对应总市值 18.6 亿元，近一年-3.8\%，总市值在三级行业“食用菌”的 4 家公司中按由大到小排第\zhnumber{4}位，近一年涨跌幅高于全样本中位数（-10.4\%）6.6 个百分点。 按 2025 年年报披露的行业构成，收入占比前三位为食用菌 73.5\%；咖啡 18.0\%；火腿及制品 8.1\%。

\pfl{股权结构与控制权} 按 2026 年 6 月 30 日的股东名册，前三大股东为深圳市得融投资发展有限公司 持股 17.15\%；北京北大资源科技有限公司 持股 11.79\%；云南传奇投资有限公司 持股 10.67\%，股东户数 12,718 户。 控制权与股东性质方面，公司控股云南诺邓金腿食品科技 74.82\%，并持有新疆众和 6.01\% 股份作为联营企业。

\pfl{经营要点} 经营与产能方面，食用菌业务收入 1,742.91 万元、同比减少 10.73\%，火腿业务收入 157.11 万元、同比减少 22.68\%；上半年按权益法确认投资收益 3,135.85 万元、同比增加 19.45\%。

\pfl{所处行业的基本面} 食用菌的供给是工厂化产能（投放周期 6---9 个月），需求是刚性消费，价格随产能快速调整（第 \ref{sec:framework} 章的供需框架）。 与该公司直接相关的行业数据包括：2024 年，全国食用菌总产量 4,419.90 万吨、总产值 4,207.07 亿元；2025 年，松茸产品出口量 729.89 吨、创汇约 2.17 亿元；2024 年，云南省食用菌总产量 119.41 万吨、总产值 472.49 亿元。

\pfl{竞争优势} 把该公司放回同类公司比较，2026 年上半年营业收入 0.3 亿元、在三级行业“食用菌”的 4 家公司中排第\zhnumber{4}位；毛利率 5.4\%、低于行业内中位数 11.0 个百分点；半年 ROE 1.5\%（未年化）、低于行业内中位数（2.15\%）；资产负债率 6.6\%、低于行业内中位数（44.7\%）。 就当期报表看，公司尚未把经营层面的特征转化为高于行业的资本回报，周期底部的相对优势更多体现在现金流与资产负债表上。

\pfl{历史股价} 本报告样本区间（2021 年 1 月 至 2026 年 10 月）内，股价由 7.74 元变为 7.88 元，累计 +1.8\%，区间最高 9.33 元（2022 年 1 月）、最低 5.54 元（2024 年 8 月）；当前价相当于区间最高价的 84.5\%，为区间最低价的 1.42 倍。 月度收益的年化波动率 25.6\%，低于全样本中位数 37.3\%，在三级行业“食用菌”的 4 家公司中波动率由高到低排第\zhnumber{4}位。 把股价阶段与公司事件对齐看，2024 年 8 月 至 2026 年 4 月股价上涨+62.6\%，同期（2025 年 12 月）公司账面同时持有云南白药股份 67.34 万股，并减持众和转债回笼资金。 从时间对应关系看，公司层面的资本运作与股权、风险事件解释了区间内多数急涨急跌，与行业指数的同步性相对有限。

\begin{center}
\begin{minipage}{\textwidth}\centering
\begin{tikzpicture}
\begin{axis}[
  width=12.6cm, height=3.9cm, scale only axis,
  xmin=2020.922, xmax=2026.972, ymin=5.15, ymax=9.98,
  xtick={2021,2022,2023,2024,2025,2026},
  yearticks,
  yticklabel style={font=\scriptsize},
  ylabel={元/股}, ylabel style={font=\scriptsize},
  grid=major, grid style={LightGray, line width=0.3pt},
  axis line style={InkGray, line width=0.5pt},
  every axis plot/.append style={line width=0.9pt},
]
\addplot[AgriGreen] table[x=x,y=y]{data/clean/profiles/px_600883.dat};
\addplot[SoftRed, dashed, line width=0.7pt] table[x=x,y=y]{data/clean/profiles/pxzz_600883.dat};
\addplot[only marks, mark=*, mark size=1.1pt, SoftRed] table[x=x,y=y]{data/clean/profiles/pxzz_600883.dat};
\end{axis}
\end{tikzpicture}

\captionof{figure}[博闻科技（600883）股价与周期骨架]{博闻科技（600883）月度收盘价（元/股）与 ZigZag 周期骨架（阈值 25\%，圆点为周期转折点）}
\end{minipage}
\end{center}

\pfl{过去周期分析} 以 25\% 的 ZigZag 阈值划分，样本区间内股价形成 2 个主升段与 1 个主跌段。 最大主升段出现在 2024 年 8 月 至 2026 年 4 月，20 个月+62.6\%；最大主跌段为 2022 年 1 月 至 2024 年 8 月，31 个月-40.6\%。 最近一次转折出现在 2026 年 4 月（上涨段自 2024 年 8 月起，20 个月+62.6\%），当前处于该段之后的震荡之中。 公司股价较申万农林牧渔指数提前约 21 个月见底，是板块中较早定价周期反转的公司之一。 食用菌是工厂化生产的制造型农业（第 \ref{sec:soft} 章），产能投放节奏与金针菇等单品价格是盈利的主要变量，与养殖周期基本无关，公司 2026 年上半年实现盈利、半年 ROE 1.5\%（未年化），好于全样本中位数（0.75\%），下行周期对其报表的冲击小于同业。 综合区间形态看，该股单段涨跌幅度偏大、以趋势段为主（最大主升 +62.6\%、最大主跌 -40.6\%），年化波动率 25.6\% 与全样本中位数 37.3\% 相比波动小于样本中位数。

\pfl{盈利情况} 2026 年上半年营业收入 0.26 亿元（同比 -8.8\%），归母净利润 0.15 亿元，同比 +5.3\%。 以年报为趋势对照，2023---2025 年营业收入由 0.21 亿元变为 0.56 亿元（两年年均复合增速 +62.2\%），归母净利润由 0.96 亿元变为 0.24 亿元。 按半年报口径，2026 年上半年的 ROE 为 1.5\%（未年化）、销售净利率 56.9\%、毛利率 5.4\%，ROE 在“食用菌”4 家公司中按数值由高到低排第\zhnumber{3}位。 按同一口径，三级行业“食用菌”4 家公司的半年 ROE 中位数为 2.15\%，该公司低于其中位数 0.64 个百分点。 综合看，公司在报告期维持盈利（高于全样本半年 ROE 中位数 0.75\%），利润的可持续性取决于食用菌景气的延续与成本管控的执行。

\begin{center}\small\setlength{\tabcolsep}{3.4pt}
\captionof{table}[博闻科技（600883）关键财务指标]{博闻科技（600883）关键财务指标（合并报表口径；两个半年报期均未年化；现金短债比 $=$ 货币资金 $\div$ 短期债务）}
\begin{adjustbox}{max width=\textwidth}
\begin{tabular}{@{}lrrrrr@{}}
\toprule
指标 & 2023 年 & 2024 年 & 2025 年 & 2025 年半年报 & 2026 年半年报 \\
\midrule
营业收入（亿元） & 0.21 & 0.43 & 0.56 & 0.29 & 0.26 \\
归母净利润（亿元） & 0.96 & 0.72 & 0.24 & 0.14 & 0.15 \\
毛利率（\%） & -1.0 & 0.9 & 8.0 & 5.6 & 5.4 \\
ROE（\%） & 11.2 & 7.7 & 2.4 & 1.5 & 1.5 \\
资产负债率（\%） & 5.6 & 5.7 & 5.9 & 5.8 & 6.6 \\
经营活动现金流净额（亿元） & -0.33 & -0.40 & -0.11 & -0.05 & -0.05 \\
货币资金（亿元） & 1.07 & 0.35 & 1.63 & 1.07 & 1.43 \\
有息负债合计（亿元） & 0.04 & 0.04 & 0.05 & 0.04 & 0.12 \\
现金短债比（倍） & 30.98 & 8.33 & 35.25 & 25.76 & 30.40 \\
\bottomrule
\end{tabular}
\end{adjustbox}
\end{center}

\begin{center}
\begin{minipage}{\textwidth}\centering
\begin{tikzpicture}
\begin{groupplot}[
  group style={group size=2 by 1, horizontal sep=0.60cm},
  width=6.85cm, height=4.60cm,
  tick label style={font=\tiny},
  title style={font=\scriptsize\heiti},
  yticklabel style={font=\tiny},
  ylabel style={font=\tiny},
  grid=major, grid style={LightGray, line width=0.3pt},
  axis line style={InkGray, line width=0.5pt},
  enlarge x limits={abs=0.8},
  legend style={font=\scriptsize, at={(0.5,-0.20)}, anchor=north,
                legend columns=2, draw=none, fill=none, column sep=6pt},
]
\nextgroupplot[
  ybar, bar width=4pt,
  ymin=-0.10, ymax=1.08,
  xtick={0,1,2,3,4,5.7},
  xticklabels={2021,2022,2023,2024,2025,2026H1},
  ylabel={亿元},
]
\addplot[fill=AgriGreen!75, draw=AgriGreen, bar shift=-2.6pt] table[x=i, y=rev]{data/clean/profiles/fin_600883.dat};
\addplot[fill=SoftRed!60, draw=SoftRed, bar shift=2.6pt] table[x=i, y=ni]{data/clean/profiles/fin_600883.dat};
\addplot[fill=AgriGreen!30, draw=AgriGreen, pattern=north east lines, bar shift=-2.6pt] table[x=x, y=rev]{data/clean/profiles/finh_600883.dat};
\addplot[fill=SoftRed!20, draw=SoftRed, pattern=north east lines, bar shift=2.6pt] table[x=x, y=ni]{data/clean/profiles/finh_600883.dat};
\legend{营业收入（年/半年）, 归母净利润（年/半年）}
\nextgroupplot[
  ymin=-0.9, ymax=12.3,
  xtick={0,1,2,3,4,5.7},
  xticklabels={2021,2022,2023,2024,2025,2026H1},
  ylabel={\%},
]
\addplot[AgriGold, mark=*, mark size=1.3pt] table[x=i, y=roe]{data/clean/profiles/fin_600883.dat};
\addplot[OilBlue, mark=square*, mark size=1.3pt] table[x=i, y=debt]{data/clean/profiles/fin_600883.dat};
\addplot[only marks, AgriGold, mark=*, mark size=2.0pt] table[x=x, y=roe]{data/clean/profiles/finh_600883.dat};
\addplot[only marks, OilBlue, mark=square*, mark size=2.0pt] table[x=x, y=debt]{data/clean/profiles/finh_600883.dat};
\legend{ROE, 资产负债率}
\end{groupplot}
\end{tikzpicture}
\begin{tikzpicture}
\begin{axis}[
  name=finbot,
  width=13.75cm, height=3.10cm, scale only axis,
  ymin=-0.16, ymax=1.82,
  xtick={0,1,2,3,4,5.7},
  xticklabels={2021,2022,2023,2024,2025,2026H1},
  tick label style={font=\tiny},
  yticklabel style={font=\tiny},
  ylabel={亿元}, ylabel style={font=\tiny},
  ybar, bar width=4pt,
  grid=major, grid style={LightGray, line width=0.3pt},
  axis line style={InkGray, line width=0.5pt},
  enlarge x limits={abs=0.8},
  legend style={font=\scriptsize, at={(0.5,-0.26)}, anchor=north,
                legend columns=2, draw=none, fill=none, column sep=10pt},
]
\addplot[fill=AgriGreen!55, draw=AgriGreen, bar shift=-3.0pt] table[x=i, y=cash]{data/clean/profiles/fin_600883.dat};
\addplot[fill=SoftRed!45, draw=SoftRed, bar shift=3.0pt] table[x=i, y=ibd]{data/clean/profiles/fin_600883.dat};
\addplot[fill=AgriGreen!25, draw=AgriGreen, pattern=north east lines, bar shift=-3.0pt] table[x=x, y=cash]{data/clean/profiles/finh_600883.dat};
\addplot[fill=SoftRed!20, draw=SoftRed, pattern=north east lines, bar shift=3.0pt] table[x=x, y=ibd]{data/clean/profiles/finh_600883.dat};
\legend{货币资金（左轴）, 有息负债（左轴）}
\end{axis}
\begin{axis}[
  at={(finbot.south east)}, anchor=south east,
  width=13.75cm, height=3.10cm, scale only axis,
  axis y line*=right, axis x line*=none, xtick=\empty,
  ymin=-0.50, ymax=0.12,
  tick label style={font=\tiny},
  yticklabel style={font=\tiny}, scaled y ticks=false,
  legend style={font=\scriptsize, at={(0.5,-0.44)}, anchor=north,
                legend columns=1, draw=none, fill=none},
]
\addplot[OilBlue, mark=*, mark size=2.2pt, line width=1.3pt] table[x=i, y=ocf]{data/clean/profiles/fin_600883.dat};
\addplot[only marks, OilBlue, mark=*, mark size=3.2pt] table[x=x, y=ocf]{data/clean/profiles/finh_600883.dat};
\legend{经营活动现金流净额（右轴，亿元，蓝点）}
\end{axis}
\end{tikzpicture}

\captionof{figure}[博闻科技（600883）近年主要财务指标]{博闻科技（600883）近年主要财务指标（左上：营业收入与归母净利润；右上：ROE 与资产负债率；下：货币资金、有息负债为左轴柱状，经营活动现金流净额为其右轴的蓝点折线。
2021---2025 年为年报口径，2026H1 为 2026 年半年报/半年末，与其前一年报之间留出间隔、以斜纹或空心点区分；半年报数据未年化）}
\end{minipage}
\end{center}

\pfl{财务分析} 2026 年 6 月末资产负债率 6.6\%，较 2025 年末上升 0.7 个百分点（样本中位数 52.2\%，所处三级行业内按由低到高排第\zhnumber{1}位）。 2026 年 6 月末货币资金 1.43 亿元、短期债务（短期借款与一年内到期非流动负债）0.05 亿元、长期债务 0.07 亿元，有息负债合计 0.12 亿元，现金短债比 30.40 倍——覆盖充分。 净有息负债 -1.31 亿元，相当于其 2026 年上半年营业收入的 -502\%。 流动比率 5.89、速动比率 5.13，短期流动性无明显缺口。 2026 年上半年经营活动现金流净额 -0.05 亿元，净利润为正但经营现金流为负，利润的现金含量偏低。 综合看，现金短债比 30.40 倍，短期偿付压力可控；资产负债率低于样本中位数 45.6 个百分点；有息负债中长期债务占比更高，期限结构对短债滚动的压力较小。

\pfl{重大战略转型} 从公开披露看，公司的战略推进集中在：并购与重组（1 项）：2024 年，持股 5\% 以上股东保山智源教育投资发展有限公司通过公开征集转让方式协议转让所持公司全部股份 1,553.17 万股（占公司总股本 6.58\%）。 公告口径上，近三年（2023 年 9 月---2026 年 9 月）公司披露重大重组与并购类公告 0 条、对外投资与转型类公告 4 条、回购与股权激励类公告 0 条，合计公告 239 条。 第一大行业仍是“食用菌”，收入占比 73.5\%，与 2021 年年报相比没有方向性变化。 综合判断，报告期内公司有 4 条与战略相关的公告落地（重大重组与并购 0 条、对外投资与转型 4 条），资本运作与主业调整之间存在直接关联。

\begin{center}\small\setlength{\tabcolsep}{4pt}
\captionof{table}[博闻科技（600883）历史融资与资本运作]{博闻科技（600883）历史融资与资本运作（据公司公告与公开报道整理；金额为披露值，含首次公开发行、再融资、债券、银行借款与重整投资等）}
\begin{adjustbox}{max width=\textwidth}
\begin{tabular}{@{}p{1.45cm}p{2.55cm}p{4.05cm}p{4.95cm}@{}}
\toprule
时间 & 方式 & 规模 & 用途与说明 \\
\midrule
1995-12 & IPO（上海证券交易所挂牌） & 募集资金规模 & 挂牌时公司名为云南省保山水泥股份有限公司 \\
2025-03 & 自有资金出资设立咖啡业务子公司 & 以自有货币资金出资 2,\allowbreak{}000.00 万元人民币 & 在云南省怒江傈僳族自治州泸水市设立子公司，\allowbreak{}作为咖啡种植基地项目投资运营主体；年报同时披露怒江博闻咖啡有限公司实收资本 250.00 万元；截至 2025 年末 \\
2025-07-17 至 2025-08-08 & 减持所持新疆众和可转换公司债券（众和转债） & 减持 500,\allowbreak{}000 张，\allowbreak{}累计成交金额 6,\allowbreak{}491.87 万元 & 资产处置与资金回笼；本次减持完成后公司不再持有众和转债编号临 2025-023 \\
2026-05-22 & 2025 年度利润分配实施 & 按每股人民币 0.04 元（含税）、\allowbreak{}已发行股份 2.36 亿股计算 & 股东回报；2025 年度股东大会批准；母公司 2025 年末可供分配利润 5.48 亿元 \\
\bottomrule
\end{tabular}
\end{adjustbox}
\end{center}

\pfl{历史融投资情况} 从现金流量表看，2025 年公司偿还债务支付的现金 0.00 亿元，筹资活动现金流净额 -0.12 亿元（2023 年为取得借款 0.03 亿元）。 2026 年上半年筹资活动现金流净额为 -0.07 亿元（取得借款 — 亿元、偿还债务 0.00 亿元，未年化），已转为净流出，处于净还债状态。 2025 年分配股利、利润或偿付利息支付的现金合计 0.11 亿元。 公开披露的融资记录共 4 笔，工具类型以 IPO、可转债为主；金额与用途见下表。 其中规模最大的两笔为：2025 年 7 月，可转债，减持 500,000 张，累计成交金额 6,491.87 万元；2025 年 3 月，自有资金出资设立咖啡，以自有货币资金出资 2,000.00 万元人民币。 股本方面，近三年披露股本变动 3 次；总股本自 2021 年末以来未发生变化（2.36 亿股）。 分红方面，近三年现金分红 4 次、合计每股派现 0.25 元，最近一次实施在 2025 年 12 月。 近三年“再融资”类公告 1 条，涉及可转债。 综合看，公司融资结构上股权与债务融资并举；2025 年筹资活动现金流净额 -0.12 亿元，融资节奏仍以维持存量债务周转为主，与前述经营与投资现金流的方向基本一致。

\pfl{未来潜在融资需求分析} 2026 年 6 月末货币资金 1.43 亿元，而短期债务 0.05 亿元（现金短债比 30.40 倍），2026 年上半年经营现金流净额 -0.05 亿元。账面货币资金可以覆盖短期债务，缺口不体现在静态偿付层面。 公司 2026 年 6 月末资产负债率 6.6\%、2026 年上半年 ROE 1.5\%（未年化），资产端与负债端都没有明显压力，融资需求以相机抉择的扩张型用途为主。 综合看，公司在静态偿付层面不存在缺口，外部融资更多是配合扩张节奏与并购整合的主动性安排，而非偿债压力驱动。






\section{养殖业：生猪、肉鸡与其他养殖}
\label{sec:breeding}

养殖业是农业上市公司中市值最大的板块：\num{22} 家公司、总市值合计约 \num{4833} 亿元，
占全部农业上市公司合计总市值的 \hl{39.5}\%。
按 2026 年半年报口径，\num{15} 家亏损、\num{7} 家资产负债率超过 \num{65}\%、
半年 ROE 中位数 \num{-6.26}\%（未年化）。
\hl{这是整个报告中\emph{平均值代表性最低}的板块}：
平均值掩盖了生猪养殖（\num{11} 家全部亏损）与肉鸡养殖（整体仍盈利 \num{8.5} 亿元）
之间的巨大差异。

\begingroup\scriptsize\setlength{\tabcolsep}{2pt}
\begin{longtable}{@{}llrrrrrrr@{}}
\caption{农业上市公司分行业明细（生猪养殖、肉鸡养殖、其他养殖；市值时点 2026-10-12，总市值口径；财务为 2026 年半年报，ROE 未年化，公告计数窗口 2023-09---2026-10）}\label{tab:comp-breeding}\\
\toprule
代码 & 公司简称 & 总市值（亿元） & 营收（亿元） & 同比（\%） & 净利（亿元） & ROE（\%） & 再融资公告 & 重组公告 \\
\midrule
\endfirsthead
\multicolumn{9}{l}{\small（续）}\\
\toprule
代码 & 公司简称 & 总市值（亿元） & 营收（亿元） & 同比（\%） & 净利（亿元） & ROE（\%） & 再融资公告 & 重组公告 \\
\midrule
\endhead
\bottomrule
\endlastfoot
002714 & 牧原股份 & 2,283.4 & 594.1 & -22.3 & -60.78 & -7.5 & 24 & 0 \\
300498 & 温氏股份 & 1,007.3 & 467.6 & -6.2 & -43.66 & -11.1 & 59 & 3 \\
002157 & 正邦科技 & 269.6 & 77.2 & 12.9 & -7.36 & -7.0 & 12 & 0 \\
002299 & 圣农发展 & 197.1 & 107.8 & 21.7 & 5.02 & 4.4 & 0 & 3 \\
300761 & 立华股份 & 168.1 & 91.3 & 9.3 & -0.72 & -0.8 & 3 & 0 \\
605296 & 神农集团 & 159.7 & 24.9 & -11.0 & -8.09 & -17.9 & 27 & 1 \\
000048 & 京基智农 & 91.0 & 21.3 & -10.7 & -4.89 & -12.9 & 9 & 1 \\
002458 & 益生股份 & 87.4 & 17.0 & 28.4 & 3.08 & 6.9 & 10 & 0 \\
603477 & 巨星农牧 & 71.8 & 35.6 & -4.3 & -8.38 & -30.3 & 75 & 0 \\
600975 & 新五丰 & 64.6 & 28.3 & -23.2 & -7.36 & -41.0 & 4 & 3 \\
000735 & 罗牛山 & 64.4 & 12.4 & 14.5 & -3.10 & -7.4 & 1 & 4 \\
002746 & 仙坛股份 & 53.6 & 29.9 & 17.9 & 1.80 & 3.9 & 0 & 2 \\
002124 & 天邦食品 & 43.3 & 36.6 & -22.1 & -15.43 & -96.8 & 29 & 4 \\
002321 & 华英农业 & 43.1 & 23.7 & 14.5 & -0.29 & -3.2 & 1 & 0 \\
301116 & 益客食品 & 41.7 & 99.0 & 13.2 & 0.34 & 2.2 & 56 & 0 \\
300313 & 天山生物 & 37.8 & 2.2 & 427.7 & 0.04 & 12.1 & 35 & 2 \\
001201 & 东瑞股份 & 37.1 & 10.4 & -5.1 & -4.25 & -14.3 & 43 & 0 \\
300967 & 晓鸣股份 & 29.5 & 5.7 & -24.5 & 0.02 & 0.2 & 33 & 0 \\
002234 & 民和股份 & 26.8 & 10.3 & 4.2 & -0.95 & -5.5 & 1 & 0 \\
600965 & *ST福成 & 20.1 & 5.7 & 6.9 & 0.18 & 0.8 & 0 & 0 \\
002982 & 湘佳股份 & 18.3 & 23.2 & 8.6 & -0.12 & -0.7 & 15 & 0 \\
603717 & 天域生物 & 17.2 & 3.5 & -15.0 & -0.54 & -14.6 & 30 & 1 \\
\end{longtable}
\endgroup


\subsection{生猪养殖：成本曲线决定生存}
\label{sec:breeding-hog}

\num{11} 家生猪养殖公司在 2026 年上半年\hl{全部亏损}，
合计归母净利润 $-163.8$ 亿元（上年同期为 $+154.1$ 亿元），
这正是猪周期在报表层面的底部形态：

\begin{itemize}[leftmargin=1.4em, itemsep=2pt]
  \item \textbf{牧原股份}：营收 \num{594.1} 亿元（同比 \num{-22.3}\%）、
        归母净利润 \num{-60.8} 亿元、半年 ROE \num{-7.5}\%，
        其亏损额在全部 \num{104} 家农业上市公司中最大，
        但 2025 年年报仍实现 \num{154.9} 亿元净利润、ROE \num{20.6}\%（趋势对照）；
  \item \textbf{温氏股份}：营收 \num{467.6} 亿元（同比 \num{-6.2}\%）、
        亏损 \num{-43.7} 亿元、半年 ROE \num{-11.1}\%；
  \item \textbf{天邦食品}：营收 \num{36.6} 亿元、亏损 \num{-15.4} 亿元，
        半年 ROE \num{-96.8}\%——\emph{净资产已被亏损消耗殆尽}，
        2026 年 6 月末资产负债率 \num{91.2}\%；
  \item 尾部还有 \textbf{巨星农牧}（\num{-8.4} 亿元）、\textbf{神农集团}（\num{-8.1} 亿元）、
        \textbf{正邦科技}（\num{-7.4} 亿元）、\textbf{新五丰}（\num{-7.4} 亿元），
        亏损额均超过 \num{7} 亿元。
\end{itemize}

\hl{牧原与温氏两家合计亏损 \num{104.4} 亿元，占板块亏损的六成以上}。
这与第 \ref{sec:hog-cost} 节的判断一致：以行业平均完全成本约 \num{13} 元/公斤衡量，
在猪价 \num{10}---\num{12} 元/公斤的深度亏损区间，只有成本处在该水平以下的企业还能维持盈亏平衡。
牧原的 \hl{自繁自养一体化与低豆粕日粮}结构使其成本处于行业最低分位，
2026 年上半年的亏损因而小于其收入规模的收缩幅度；
而“公司 + 农户”模式下外购仔猪比例更高的企业，在仔猪价格从 \num{108.6} 元/公斤
跌至 \num{22.05} 元/公斤的过程中承受了双向亏损（仔猪跌价损失 + 育肥亏损）。

\hl{正邦科技是观察周期破坏力的最典型样本}：总市值仍高达 \num{269.6} 亿元（板块第三），
但 2026 年上半年营收仅 \num{77.2} 亿元、亏损 \num{7.4} 亿元、半年 ROE \num{-7.0}\%，
2025 年年报口径下同样亏损 \num{5.5} 亿元。
市场对其给予高市值的原因是重整预期与产能保留价值，
而不是当前的盈利能力——\emph{这本身是周期底部的典型定价特征}。

\subsection{肉鸡养殖：反转最快的一环}
\label{sec:breeding-poultry}

\num{8} 家肉鸡养殖公司 2026 年上半年合计盈利 \num{8.5} 亿元（上年同期 \num{10.0} 亿元），
是养殖业中唯一整体为正的子行业：\textbf{圣农发展}（\num{5.02} 亿元、半年 ROE \num{4.4}\%）、
\textbf{益生股份}（\num{3.08} 亿元、\num{6.9}\%，上年同期仅 \num{0.06} 亿元）、
\textbf{仙坛股份}（\num{1.80} 亿元、\num{3.9}\%）
与 \textbf{益客食品}（\num{0.34} 亿元，上年同期亏损 \num{1.81} 亿元）实现盈利；
\textbf{立华股份}（\num{-0.72} 亿元）、\textbf{民和股份}（\num{-0.95} 亿元）
与 \textbf{湘佳股份}（\num{-0.12} 亿元）小幅亏损。
这一分化与第 \ref{sec:other-poultry} 节的产业数据相互印证：
2025 年白羽肉鸡祖代更新量创历史新高、出栏均价跌至近十年低位，
\hl{上游种鸡与商品代养殖环节同时承压}；
而圣农发展以全产业链（自有种源 + 自养 + 屠宰 + 深加工）模式
锁定了下游食品端的利润，抵消了商品代的周期波动。

2026 年 1---4 月毛鸡养殖利润同比 +$297.22\%$，
\hl{禽链的盈利反转已经开始}：半年报中益生股份与益客食品的利润已经改善，
但整体仍低于上年同期，\emph{反转尚未完全体现在报表中}。

\begin{keybox}[养殖业的三点判断]
\normalsize
\textbf{① 集中度提升不可逆。}生猪养殖 11 家公司在 2026 年上半年全部亏损，
但牧原、温氏两家合计亏损 \num{104.4} 亿元、占板块亏损六成以上，
行业产能与资本回报向低成本龙头集中的趋势在 2023---2026 年持续强化。\par\vspace{3pt}
\textbf{② 高负债 + 亏损 = 融资依赖。}板块内 \num{7} 家公司在 2026 年 6 月末的
资产负债率超过 \num{65}\%（天邦食品 \num{91.2}\%、新五丰 \num{84.2}\%、天域生物 \num{81.0}\%），
其中多家同时处于亏损状态，\hl{其对股权融资的依赖程度高于其他子行业}
（详见第 \ref{sec:financing} 节）。\par\vspace{3pt}
\textbf{③ 利润修复尚未出现。}板块 \num{22} 家公司 2026 年上半年合计归母净利润
$-155.4$ 亿元（上年同期 \hl{$+164.1$ 亿元}），\num{15} 家亏损：
其中生猪养殖 \num{11} 家全部亏损、合计 $-163.8$ 亿元，
肉鸡养殖 \num{8} 家合计 $+8.5$ 亿元、较上年同期的 $+10.0$ 亿元略有回落。
\hl{猪链仍在磨底，禽链的产能反转也尚未完全体现为报表利润}，
两类差异更可能出现在 2026 年年报中。
\end{keybox}

\subsection{上市公司画像（\num{22} 家）}
\label{sec:co-breeding}

本节给出养殖业全部 \num{22} 家上市公司的画像，按总市值降序排列。
数据口径与第 \ref{sec:comp-overview} 节一致：\hl{财务数据统一采用 2026 年半年报}
（合并报表口径、未年化，资产负债率为 2026 年 6 月末），
同时列出 2025 年半年报作为同口径对照、2023---2025 年年报作为趋势对照；
股价为月度收盘价，最后一个交易日为 2026 年 10 月 12 日；周期分段使用与第一部分相同的
ZigZag 算法，阈值取 25\%（个股波动大于商品价格，故高于生猪现货的 20\%）；
公告计数窗口为 2023 年 9 月 1 日至 2026 年 10 月 12 日。

与前一版不同，本版画像正文由\hl{逐家检索的公开资料底稿}驱动：底稿只收录有来源链接的事实
（公司沿革、股权与控制权、经营要点、所处行业的基本面、竞争优势、融资历史、
风险事件、战略动作、融资需求线索），
正文按事实生成，各家公司的段落结构与重点并不相同。
每家公司在统一顺序下展开：基本情况与主营结构、股权结构与控制权、经营要点、
所处行业的基本面（供需）、与同类公司比较的竞争优势、历史股价、过去周期分析、
盈利情况、财务分析、重大战略转型、历史融投资情况、未来潜在融投资需求；
\hl{市值一律采用总市值口径}，股权结构取最新披露报告期前三大股东与股东户数，
行业基本面与竞争优势以逐家检索的公开资料为依据，并与三级行业内的同口径指标横向比较。
每家公司配图两张：
股价与行业指数（申万农林牧渔）的对比图，以及近年主要财务指标的组合图
（上排：营业收入与归母净利润、ROE 与资产负债率；下方宽带面板：货币资金与有息负债为左轴柱状、
经营活动现金流净额为其右轴的蓝点折线），并给出关键财务指标表与历史融资/资本运作表。
未来融资需求按两个口径粗略度量：\hl{短期债务缺口 $=$ 期末短期债务 $-$ 货币资金}
与\hl{年化经营缺口 $=$ 半年归母亏损 $\times 2$}；
前者衡量滚动偿付压力，后者衡量维持当期盈利水平的现金消耗，
两者都是压力情形下的静态估算，不构成对融资规模或时点的预测。

% 农业上市公司画像：生猪养殖、肉鸡养殖、其他养殖（共 22 家，按总市值降序）
% 由 scripts/make_company_profiles.py 自动生成，请勿手工修改

\subsubsection{牧原股份（002714）}
\label{co:002714}

\pfl{公司基本情况} 公司前身可追溯至 2014 年，自繁自养一体化模式的全球最大生猪养殖企业，2026 年 H 股上市后形成 A+H 双平台；2026 年上半年猪价深跌录得大额亏损，公司主动调减能繁母猪产能并压降养殖成本。 公司于 2014-01-28 在深交所主板首发上市，注册资本 57.73 亿元，总股本 54.65 亿股。 截至 2026 年 10 月 12 日收盘价 41.78 元，对应总市值 2,283.4 亿元，近一年-21.5\%，总市值在三级行业“生猪养殖”的 11 家公司中按由大到小排第\zhnumber{1}位，近一年涨跌幅低于全样本中位数（-10.4\%）11.1 个百分点。 按 2026 年半年报披露的产品构成，收入占比前三位为生猪 88.3\%；屠宰、肉食产品 37.1\%；饲料原料 7.3\%。

\pfl{股权结构与控制权} 按 2026 年 6 月 30 日的股东名册，前三大股东为秦英林 持股 36.14\%；牧原实业集团有限公司 持股 14.70\%；HKSCC NOMINEES LIMITED 持股 5.37\%，股东户数 331,108 户。 控制权与股东性质方面，二人分别控制 48.64\% 和 3.33\% 的股份；2026 年 2 月，公司募资约 120.00 亿港元且含超额配股权，其中 60\% 投向海外业务；2026 年上半年，董事会另审议通过总金额 3 亿至 5.00 亿港元的 H 股回购方案。

\pfl{经营要点} 业务与产能层面，2026 年 6 月，生猪养殖完全成本 11.7 元/公斤，断奶至上市成活率约 92\%；截至 2026 年 6 月末能繁母猪存栏 311.3 万头。

\pfl{所处行业的基本面} 生猪的供给由能繁母猪存栏决定、需求由猪肉消费与替代品价格决定，本轮下行的成因在供给端的产能去化不彻底，而非需求收缩（第 \ref{sec:framework} 章的供需框架）。 从供需两端的当期数据看，2026 年上半年，全国生猪出栏 37246 万头、同比增长 1.7\%，二季度末生猪存栏 42491 万头、同比增长 0.1\%；2026 年二季度，末全国能繁母猪存栏 3780 万头；2026 年 5 月，农业农村部印发生猪产能综合调控实施方案，将能繁母猪正常保有量由 3900 万头下调至 3750 万头。

\pfl{竞争优势} 把该公司放回同类公司比较，2026 年上半年营业收入 594.1 亿元、在三级行业“生猪养殖”的 11 家公司中排第\zhnumber{1}位；毛利率 -1.6\%、与行业内中位数基本持平；半年 ROE -7.5\%（未年化）、高于行业内中位数（-14.32\%）；资产负债率 54.2\%、低于行业内中位数（58.9\%）。 与同类公司相比，公司的差异化在于：2026 年 7 月，完全成本已降至约 11.5 元/公斤，而行业平均完全成本已降至约 12 元/公斤；约 45\% 的商品猪进入自有屠宰体系，当养殖业务毛利率转负时屠宰端仍月月盈利；2026 年 6 月，与阿里云签署战略合作。 优势在当期体现为高于行业中位数的资本回报，但优势能否持续仍取决于周期与自身产能投放节奏。

\pfl{历史股价} 以 2021 年 1 月为起点、2026 年 10 月为终点，股价由 88.90 元变为 41.78 元，累计 -53.0\%，区间最高 113.90 元（2021 年 2 月）、最低 34.99 元（2026 年 6 月）；当前价相当于区间最高价的 36.7\%，为区间最低价的 1.19 倍。 月度收益的年化波动率 38.4\%，高于全样本中位数 37.3\%，在三级行业“生猪养殖”的 11 家公司中波动率由高到低排第\zhnumber{5}位。

\begin{center}
\begin{minipage}{\textwidth}\centering
\begin{tikzpicture}
\begin{axis}[
  width=12.6cm, height=3.9cm, scale only axis,
  xmin=2020.922, xmax=2026.972, ymin=32.54, ymax=121.87,
  xtick={2021,2022,2023,2024,2025,2026},
  yearticks,
  yticklabel style={font=\scriptsize},
  ylabel={元/股}, ylabel style={font=\scriptsize},
  grid=major, grid style={LightGray, line width=0.3pt},
  axis line style={InkGray, line width=0.5pt},
  every axis plot/.append style={line width=0.9pt},
]
\addplot[AgriGreen] table[x=x,y=y]{data/clean/profiles/px_002714.dat};
\addplot[SoftRed, dashed, line width=0.7pt] table[x=x,y=y]{data/clean/profiles/pxzz_002714.dat};
\addplot[only marks, mark=*, mark size=1.1pt, SoftRed] table[x=x,y=y]{data/clean/profiles/pxzz_002714.dat};
\end{axis}
\end{tikzpicture}

\captionof{figure}[牧原股份（002714）股价与周期骨架]{牧原股份（002714）月度收盘价（元/股）与 ZigZag 周期骨架（阈值 25\%，圆点为周期转折点）}
\end{minipage}
\end{center}

\pfl{过去周期分析} 以 25\% 的 ZigZag 阈值划分，样本区间内股价形成 2 个主升段与 3 个主跌段。 最大主升段出现在 2024 年 1 月 至 2025 年 8 月，19 个月+55.9\%；最大主跌段为 2021 年 1 月 至 2021 年 7 月，6 个月-52.5\%。 最近一次转折出现在 2026 年 6 月（下跌段自 2025 年 8 月起，10 个月-36.3\%），当前处于该段的延续之中。 公司股价与申万农林牧渔指数几乎同步见底，个股并未走出独立行情。 生猪养殖的定价主线是产能去化与成本曲线的赛跑（第 \ref{sec:hog} 章）：2026 年三季度行业仍处在“产能去化未完成、价格磨底”的阶段，公司 2026 年 6 月末资产负债率 54.2\%、半年 ROE -7.5\%（未年化），在本轮下行中属于随行业同步调整的一类。 综合区间形态看，该股单段涨跌幅度偏大、以趋势段为主（最大主升 +55.9\%、最大主跌 -52.5\%），年化波动率 38.4\% 与全样本中位数 37.3\% 相比波动与样本中位数接近。

\pfl{盈利情况} 2026 年上半年营业收入 594.10 亿元（同比 -22.3\%），归母净利润 -60.78 亿元，由上年同期的盈利转为亏损。 以年报为趋势对照，2023---2025 年营业收入由 1,108.61 亿元变为 1,441.45 亿元（两年年均复合增速 +14.0\%），归母净利润由 -42.63 亿元变为 154.87 亿元。 按半年报口径，2026 年上半年的 ROE 为 -7.5\%（未年化）、销售净利率 -10.2\%、毛利率 -1.6\%，ROE 在“生猪养殖”11 家公司中按数值由高到低排第\zhnumber{3}位。 按同一口径，三级行业“生猪养殖”11 家公司的半年 ROE 中位数为 -14.32\%，该公司高于其中位数 6.84 个百分点。 综合看，公司仍处亏损区间、由盈转亏，半年 ROE 在三级行业内按由高到低排第\zhnumber{3}位，单位盈利能力的修复依赖行业景气与自身成本端的改善，报表弹性主要体现为价格与销量的方向性变化。

\begin{center}\small\setlength{\tabcolsep}{3.4pt}
\captionof{table}[牧原股份（002714）关键财务指标]{牧原股份（002714）关键财务指标（合并报表口径；两个半年报期均未年化；现金短债比 $=$ 货币资金 $\div$ 短期债务）}
\begin{adjustbox}{max width=\textwidth}
\begin{tabular}{@{}lrrrrr@{}}
\toprule
指标 & 2023 年 & 2024 年 & 2025 年 & 2025 年半年报 & 2026 年半年报 \\
\midrule
营业收入（亿元） & 1,108.61 & 1,379.47 & 1,441.45 & 764.63 & 594.10 \\
归母净利润（亿元） & -42.63 & 178.81 & 154.87 & 105.30 & -60.78 \\
毛利率（\%） & 3.1 & 19.1 & 17.8 & 19.0 & -1.6 \\
ROE（\%） & -6.4 & 25.1 & 20.6 & 13.8 & -7.5 \\
资产负债率（\%） & 62.1 & 58.7 & 54.2 & 56.1 & 54.2 \\
经营活动现金流净额（亿元） & 98.93 & 375.43 & 300.56 & 173.51 & -22.24 \\
货币资金（亿元） & 194.29 & 169.52 & 138.62 & 202.85 & 165.47 \\
有息负债合计（亿元） & 787.62 & 760.80 & 702.49 & 773.01 & 748.02 \\
现金短债比（倍） & 0.35 & 0.32 & 0.28 & 0.37 & 0.29 \\
\bottomrule
\end{tabular}
\end{adjustbox}
\end{center}

\begin{center}
\begin{minipage}{\textwidth}\centering
\begin{tikzpicture}
\begin{groupplot}[
  group style={group size=2 by 1, horizontal sep=0.60cm},
  width=6.85cm, height=4.60cm,
  tick label style={font=\tiny},
  title style={font=\scriptsize\heiti},
  yticklabel style={font=\tiny},
  ylabel style={font=\tiny},
  grid=major, grid style={LightGray, line width=0.3pt},
  axis line style={InkGray, line width=0.5pt},
  enlarge x limits={abs=0.8},
  legend style={font=\scriptsize, at={(0.5,-0.20)}, anchor=north,
                legend columns=2, draw=none, fill=none, column sep=6pt},
]
\nextgroupplot[
  ybar, bar width=4pt,
  ymin=-211.00, ymax=1621.72,
  xtick={0,1,2,3,4,5.7},
  xticklabels={2021,2022,2023,2024,2025,2026H1},
  ylabel={亿元},
]
\addplot[fill=AgriGreen!75, draw=AgriGreen, bar shift=-2.6pt] table[x=i, y=rev]{data/clean/profiles/fin_002714.dat};
\addplot[fill=SoftRed!60, draw=SoftRed, bar shift=2.6pt] table[x=i, y=ni]{data/clean/profiles/fin_002714.dat};
\addplot[fill=AgriGreen!30, draw=AgriGreen, pattern=north east lines, bar shift=-2.6pt] table[x=x, y=rev]{data/clean/profiles/finh_002714.dat};
\addplot[fill=SoftRed!20, draw=SoftRed, pattern=north east lines, bar shift=2.6pt] table[x=x, y=ni]{data/clean/profiles/finh_002714.dat};
\legend{营业收入（年/半年）, 归母净利润（年/半年）}
\nextgroupplot[
  ymin=-13.0, ymax=69.1,
  xtick={0,1,2,3,4,5.7},
  xticklabels={2021,2022,2023,2024,2025,2026H1},
  ylabel={\%},
]
\addplot[AgriGold, mark=*, mark size=1.3pt] table[x=i, y=roe]{data/clean/profiles/fin_002714.dat};
\addplot[OilBlue, mark=square*, mark size=1.3pt] table[x=i, y=debt]{data/clean/profiles/fin_002714.dat};
\addplot[only marks, AgriGold, mark=*, mark size=2.0pt] table[x=x, y=roe]{data/clean/profiles/finh_002714.dat};
\addplot[only marks, OilBlue, mark=square*, mark size=2.0pt] table[x=x, y=debt]{data/clean/profiles/finh_002714.dat};
\legend{ROE, 资产负债率}
\end{groupplot}
\end{tikzpicture}
\begin{tikzpicture}
\begin{axis}[
  name=finbot,
  width=13.75cm, height=3.10cm, scale only axis,
  ymin=-78.76, ymax=882.14,
  xtick={0,1,2,3,4,5.7},
  xticklabels={2021,2022,2023,2024,2025,2026H1},
  tick label style={font=\tiny},
  yticklabel style={font=\tiny},
  ylabel={亿元}, ylabel style={font=\tiny},
  ybar, bar width=4pt,
  grid=major, grid style={LightGray, line width=0.3pt},
  axis line style={InkGray, line width=0.5pt},
  enlarge x limits={abs=0.8},
  legend style={font=\scriptsize, at={(0.5,-0.26)}, anchor=north,
                legend columns=2, draw=none, fill=none, column sep=10pt},
]
\addplot[fill=AgriGreen!55, draw=AgriGreen, bar shift=-3.0pt] table[x=i, y=cash]{data/clean/profiles/fin_002714.dat};
\addplot[fill=SoftRed!45, draw=SoftRed, bar shift=3.0pt] table[x=i, y=ibd]{data/clean/profiles/fin_002714.dat};
\addplot[fill=AgriGreen!25, draw=AgriGreen, pattern=north east lines, bar shift=-3.0pt] table[x=x, y=cash]{data/clean/profiles/finh_002714.dat};
\addplot[fill=SoftRed!20, draw=SoftRed, pattern=north east lines, bar shift=3.0pt] table[x=x, y=ibd]{data/clean/profiles/finh_002714.dat};
\legend{货币资金（左轴）, 有息负债（左轴）}
\end{axis}
\begin{axis}[
  at={(finbot.south east)}, anchor=south east,
  width=13.75cm, height=3.10cm, scale only axis,
  axis y line*=right, axis x line*=none, xtick=\empty,
  ymin=-125.63, ymax=494.73,
  tick label style={font=\tiny},
  yticklabel style={font=\tiny}, scaled y ticks=false,
  legend style={font=\scriptsize, at={(0.5,-0.44)}, anchor=north,
                legend columns=1, draw=none, fill=none},
]
\addplot[OilBlue, mark=*, mark size=2.2pt, line width=1.3pt] table[x=i, y=ocf]{data/clean/profiles/fin_002714.dat};
\addplot[only marks, OilBlue, mark=*, mark size=3.2pt] table[x=x, y=ocf]{data/clean/profiles/finh_002714.dat};
\legend{经营活动现金流净额（右轴，亿元，蓝点）}
\end{axis}
\end{tikzpicture}

\captionof{figure}[牧原股份（002714）近年主要财务指标]{牧原股份（002714）近年主要财务指标（左上：营业收入与归母净利润；右上：ROE 与资产负债率；下：货币资金、有息负债为左轴柱状，经营活动现金流净额为其右轴的蓝点折线。
2021---2025 年为年报口径，2026H1 为 2026 年半年报/半年末，与其前一年报之间留出间隔、以斜纹或空心点区分；半年报数据未年化）}
\end{minipage}
\end{center}

\pfl{财务分析} 2026 年 6 月末资产负债率 54.2\%，与 2025 年末基本持平（样本中位数 52.2\%，所处三级行业内按由低到高排第\zhnumber{5}位）。 2026 年 6 月末货币资金 165.47 亿元、短期债务（短期借款与一年内到期非流动负债）568.32 亿元、长期债务 179.70 亿元，有息负债合计 748.02 亿元，现金短债比 0.29 倍——缺口明显。 净有息负债 582.56 亿元，相当于其 2026 年上半年营业收入的 98\%。 流动比率 0.79、速动比率 0.28，短期偿付压力较大。 2026 年上半年经营活动现金流净额 -22.24 亿元，经营现金流与利润同时为负，现金消耗较快。 2026 年 6 月末在建工程 23.80 亿元，较上年末减少 1.09 亿元，在建投入在收缩。 综合看，账面货币资金对有息负债与短期债务的覆盖不足，滚动偿付对新增融资的依赖度较高；资产负债率高于样本中位数 2.0 个百分点；有息负债中短期债务占比更高，期限结构仍以滚动续作为主。

\pfl{重大战略转型} 从公开披露看，公司的战略推进集中在：其他动作（1 项）：2019 年，公司设立牧原肉食、正阳肉食、商水肉食；股权与股东回报（1 项）：2024 年 9 月，公司回购总金额不低于 30.00 亿元且不超过 40.00 亿元。 公告口径上，近三年（2023 年 9 月---2026 年 9 月）公司披露重大重组与并购类公告 0 条、对外投资与转型类公告 1 条、回购与股权激励类公告 34 条，合计公告 618 条。 第一大产品仍是“生猪”，收入占比 88.3\%，与 2021 年年报相比没有方向性变化。 综合判断，报告期内公司有 1 条与战略相关的公告落地（重大重组与并购 0 条、对外投资与转型 1 条），资本运作与主业调整之间存在直接关联。

\begin{center}\small\setlength{\tabcolsep}{4pt}
\captionof{table}[牧原股份（002714）历史融资与资本运作]{牧原股份（002714）历史融资与资本运作（据公司公告与公开报道整理；金额为披露值，含首次公开发行、再融资、债券、银行借款与重整投资等）}
\begin{adjustbox}{max width=\textwidth}
\begin{tabular}{@{}p{1.45cm}p{2.55cm}p{4.05cm}p{4.95cm}@{}}
\toprule
时间 & 方式 & 规模 & 用途与说明 \\
\midrule
2014-01 & IPO（深交所中小板） & 发行 6,\allowbreak{}050 万股，\allowbreak{}发行价 24.07 元/\allowbreak{}股 & 生猪养殖产能及补充流动资金；其中新股 3,\allowbreak{}000 万股、\allowbreak{}老股转让 3,\allowbreak{}050 万股 \\
2015-12 & 定向增发 & 3,\allowbreak{}287.31 万股，\allowbreak{}发行价 30.42 元/\allowbreak{}股，\allowbreak{}实际募集资金净额 9.93 亿元 & 生猪养殖项目建设；2015 年 12 月 22 日发行 \\
2017-04 & 定向增发 & 1.25 亿股，\allowbreak{}发行价 24.67 元/\allowbreak{}股，\allowbreak{}募集资金总额 30.77 亿元 & 生猪养殖产能扩张；2017 年 4 月 14 日验资 \\
2018 & 优先股 & 6.00 亿元 & 募集资金用于生猪产能扩张等项目；票面利率 6.8\% \\
2019 & 公司债券（扶贫专项） & 8.00 亿元 & 债券简称“19 牧原 01”，\allowbreak{}代码 112849，\allowbreak{}票面利率 7.00\% \\
2019-09 & 定向增发 & 7,\allowbreak{}666.36 万股，\allowbreak{}发行价 65.22 元/\allowbreak{}股 & 新增约 475 万头生猪产能（2018 年 12 月启动预案）；新增股份 2019 年 9 月 3 日在深交所上市 \\
2021-08 & 可转换公司债券（牧原转债，127045） & 5.00 亿元（9,\allowbreak{}550 万张，\allowbreak{}面值 100.00 元） & 51.00 亿元投生猪养殖项目、\allowbreak{}19.00 亿元投生猪屠宰项目、\allowbreak{}25.50 亿元偿还银行贷款及补充流动资金；2021 年 8 月 16 日发行 \\
2022-12 & 定向增发（2021 年度非公开发行） & 1.50 亿股，\allowbreak{}发行价 39.97 元/\allowbreak{}股 & 全部用于补充流动资金；2021 年 10 月 18 日董事会通过预案 \\
2024-03 & 中期票据（银行间市场） & 注册总额度 30.00 亿元，\allowbreak{}2024 年度第一期发行金额上限 5.00 亿元、\allowbreak{}期限 3 年 & 补充流动资金及债务结构调整；在中国银行间市场交易商协会注册 \\
2026-02 & H 股 IPO（香港联交所主板） & 基础发行 2.74 亿股 & 拓展全球商业版图及研发投入，\allowbreak{}不用于国内新增产能建设；基石投资者合计认购约 53.42 亿港元 \\
\bottomrule
\end{tabular}
\end{adjustbox}
\end{center}

\pfl{历史融投资情况} 从现金流量表看，2025 年公司取得借款收到的现金 619.05 亿元、偿还债务支付的现金 679.52 亿元、吸收权益性投资 10.53 亿元，筹资活动现金流净额 -228.78 亿元（2023 年为取得借款 716.49 亿元、偿还债务 533.66 亿元）。 2026 年上半年筹资活动现金流净额为 111.60 亿元（取得借款 471.88 亿元、偿还债务 419.89 亿元，未年化），仍处于净融入状态。 2025 年分配股利、利润或偿付利息支付的现金合计 103.45 亿元。 公开披露的融资记录共 10 笔，工具类型以 IPO、定向增发、优先股、公司债等为主；金额与用途见下表。 其中规模最大的两笔为：2026 年 2 月，IPO，基础发行 2.74 亿股；2022 年 12 月，定向增发，1.50 亿股，发行价 39.97 元/股。 股本方面，近三年披露股本变动 34 次（可转债转股 28 次、激励股份解禁 3 次、股份回购 2 次、限售股份上市 1 次）；总股本由 2021 年末的 52.62 亿股变为 54.65 亿股（+3.86\%）。 分红方面，近三年现金分红 5 次、合计每股派现 3.50 元，最近一次实施在 2025 年 12 月。 近三年“再融资”类公告 24 条，涉及定向增发、可转债。 综合看，公司融资结构上以债务滚动为主、股权融资为补充；2025 年筹资活动现金流净额 -228.78 亿元，融资节奏仍以维持存量债务周转为主，与前述经营与投资现金流的方向基本一致。

\pfl{未来潜在融资需求分析} 2026 年 6 月末货币资金 165.47 亿元，而短期债务 568.32 亿元（现金短债比 0.29 倍），2026 年上半年经营现金流净额 -22.24 亿元。按“短期债务减去账面货币资金”的滚动偿付口径，静态缺口约 402.86 亿元；若再计入上半年亏损的年化额（121.55 亿元），维持一年正常经营所需的外部资金约 402.86---524.41 亿元。 公司 2026 年上半年归母净利润为负（-60.78 亿元），但 6 月末资产负债率 54.2\% 尚未进入第一档（亏损且负债率 $>$65\%）的区间，当前融资需求以补充流动资金、支撑低谷期经营性支出为主。 公开披露的后续资金安排线索包括：股权融资线索：2026 年 2 月，H 股发行已落地，为海外扩张与研发提供资金。 资金需求还要叠加在建工程：期末在建工程 23.80 亿元、2026 年上半年购建固定资产等支出 65.21 亿元（未年化），这部分支出在周期底部通常会被主动放缓。 综合看，公司的外部资金需求以营运资金周转与在建项目投入为主，滚动偿付口径下的静态缺口约 402.86 亿元，融资的时点与规模更多取决于信用条件与资本市场窗口。




\subsubsection{温氏股份（300498）}
\label{co:300498}

\pfl{公司基本情况} 以“公司+农户”模式经营的畜禽养殖龙头，肉鸡规模居全国首位、肉猪规模居行业前列，2015 年换股吸收合并大华农实现整体上市；2026 年上半年猪价深跌导致大额亏损，肉禽板块盈利部分对冲。 公司于 2015-11-02 在深交所创业板首发上市，注册资本 66.54 亿元，总股本 66.53 亿股。 截至 2026 年 10 月 12 日收盘价 15.14 元，对应总市值 1,007.3 亿元，近一年-19.6\%，总市值在三级行业“生猪养殖”的 11 家公司中按由大到小排第\zhnumber{2}位，近一年涨跌幅低于全样本中位数（-10.4\%）9.3 个百分点。 按 2026 年半年报披露的行业分部数据（含内部交易、未抵消至合并口径）为养殖行业 131.1\%；工程及设备制造业 2.5\%；动保行业 2.1\%；该口径与合并营业收入不可直接比较。

\pfl{股权结构与控制权} 按 2026 年 7 月 3 日的股东名册，前三大股东为温鹏程 持股 3.59\%；严居然 持股 2.70\%；梁焕珍 持股 2.58\%。 控制权与股东性质方面，董事长温志芬个人持股 1.526 亿股；2026 年 9 月，惠誉将公司评级展望由稳定调整为负面，同时确认长期外币发行人违约评级为 BB；公司是国内唯一发行美元债且选择公开评级的养殖上市企业。

\pfl{经营要点} 经营层面，2026 年上半年，肉猪上市率 93.1\%、料肉比 2.55；生猪屠宰竣工产能约 500 万头。

\pfl{所处行业的基本面} 生猪的供给由能繁母猪存栏决定、需求由猪肉消费与替代品价格决定，本轮下行的成因在供给端的产能去化不彻底，而非需求收缩（第 \ref{sec:framework} 章的供需框架）。 从供需两端的当期数据看，2026 年上半年，国内生猪市场价格持续下行；2026 年二季度，末全国能繁母猪存栏 3780 万头，上半年全国生猪出栏 37246 万头、同比增长 1.7\%；2026 年 5 月，农业农村部将能繁母猪正常保有量下调至 3750 万头左右，并把能繁母猪存栏总量 10 万头以上的大型集团养殖企业纳入监测名单、实行年度生产备案管理。

\pfl{竞争优势} 把该公司放回同类公司比较，2026 年上半年营业收入 467.6 亿元、在三级行业“生猪养殖”的 11 家公司中排第\zhnumber{2}位；毛利率 -2.9\%、低于行业内中位数 1.3 个百分点；半年 ROE -11.1\%（未年化）、高于行业内中位数（-14.32\%）；资产负债率 58.9\%、高于行业内中位数（58.9\%）。 与同类公司相比，公司的差异化在于：2026 年上半年猪业亏损约 40.00 亿元（含计提肉猪跌价准备约 5.00 亿元）；2026 年，固定资本开支计划为 50 亿至 60.00 亿元。 优势在当期体现为高于行业中位数的资本回报，但优势能否持续仍取决于周期与自身产能投放节奏。

\pfl{历史股价} 自 2021 年 1 月至 2026 年 10 月，股价由 16.96 元变为 15.14 元，累计 -10.7\%，区间最高 24.04 元（2022 年 7 月）、最低 12.32 元（2026 年 6 月）；当前价相当于区间最高价的 63.0\%，为区间最低价的 1.23 倍。 月度收益的年化波动率 29.7\%，低于全样本中位数 37.3\%，在三级行业“生猪养殖”的 11 家公司中波动率由高到低排第\zhnumber{10}位。

\begin{center}
\begin{minipage}{\textwidth}\centering
\begin{tikzpicture}
\begin{axis}[
  width=12.6cm, height=3.9cm, scale only axis,
  xmin=2020.922, xmax=2026.972, ymin=11.46, ymax=25.72,
  xtick={2021,2022,2023,2024,2025,2026},
  yearticks,
  yticklabel style={font=\scriptsize},
  ylabel={元/股}, ylabel style={font=\scriptsize},
  grid=major, grid style={LightGray, line width=0.3pt},
  axis line style={InkGray, line width=0.5pt},
  every axis plot/.append style={line width=0.9pt},
]
\addplot[AgriGreen] table[x=x,y=y]{data/clean/profiles/px_300498.dat};
\addplot[SoftRed, dashed, line width=0.7pt] table[x=x,y=y]{data/clean/profiles/pxzz_300498.dat};
\addplot[only marks, mark=*, mark size=1.1pt, SoftRed] table[x=x,y=y]{data/clean/profiles/pxzz_300498.dat};
\end{axis}
\end{tikzpicture}

\captionof{figure}[温氏股份（300498）股价与周期骨架]{温氏股份（300498）月度收盘价（元/股）与 ZigZag 周期骨架（阈值 25\%，圆点为周期转折点）}
\end{minipage}
\end{center}

\pfl{过去周期分析} 以 25\% 的 ZigZag 阈值划分，样本区间内股价形成 2 个主升段与 3 个主跌段。 最大主跌段为 2024 年 5 月 至 2026 年 6 月，25 个月-42.5\%；最大主升段出现在 2021 年 7 月 至 2022 年 7 月，12 个月+91.1\%。 最近一次转折出现在 2026 年 6 月（下跌段自 2024 年 5 月起，25 个月-42.5\%），当前处于该段的延续之中。 公司股价与申万农林牧渔指数几乎同步见底，个股并未走出独立行情。 生猪养殖的定价主线是产能去化与成本曲线的赛跑（第 \ref{sec:hog} 章）：2026 年三季度行业仍处在“产能去化未完成、价格磨底”的阶段，公司 2026 年 6 月末资产负债率 58.9\%、半年 ROE -11.1\%（未年化），在本轮下行中属于随行业同步调整的一类。 综合区间形态看，该股单段涨跌幅度偏大、以趋势段为主（最大主升 +91.1\%、最大主跌 -42.5\%），年化波动率 29.7\% 与全样本中位数 37.3\% 相比波动小于样本中位数。

\pfl{盈利情况} 2026 年上半年营业收入 467.65 亿元（同比 -6.2\%），归母净利润 -43.66 亿元，由上年同期的盈利转为亏损。 以年报为趋势对照，2023---2025 年营业收入由 899.21 亿元变为 1,038.62 亿元（两年年均复合增速 +7.5\%），归母净利润由 -63.90 亿元变为 52.66 亿元。 按半年报口径，2026 年上半年的 ROE 为 -11.1\%（未年化）、销售净利率 -9.4\%、毛利率 -2.9\%，ROE 在“生猪养殖”11 家公司中按数值由高到低排第\zhnumber{4}位。 同期全样本半年 ROE 中位数 0.75\%，公司与之相差 11.88 个百分点（弱于中位数公司）。 从公司披露的原因看，业绩变动主要来自：2023 年度，主要因生猪销售价格同比大幅下降及计提减值。 综合看，公司仍处亏损区间、由盈转亏，半年 ROE 在三级行业内按由高到低排第\zhnumber{4}位，单位盈利能力的修复依赖行业景气与自身成本端的改善，报表弹性主要体现为价格与销量的方向性变化。

\begin{center}\small\setlength{\tabcolsep}{3.4pt}
\captionof{table}[温氏股份（300498）关键财务指标]{温氏股份（300498）关键财务指标（合并报表口径；两个半年报期均未年化；现金短债比 $=$ 货币资金 $\div$ 短期债务）}
\begin{adjustbox}{max width=\textwidth}
\begin{tabular}{@{}lrrrrr@{}}
\toprule
指标 & 2023 年 & 2024 年 & 2025 年 & 2025 年半年报 & 2026 年半年报 \\
\midrule
营业收入（亿元） & 899.21 & 1,056.45 & 1,038.62 & 498.75 & 467.65 \\
归母净利润（亿元） & -63.90 & 92.80 & 52.66 & 34.75 & -43.66 \\
毛利率（\%） & 0.8 & 17.3 & 11.2 & 13.8 & -2.9 \\
ROE（\%） & -17.6 & 24.7 & 12.6 & 8.3 & -11.1 \\
资产负债率（\%） & 61.4 & 53.1 & 49.8 & 50.6 & 58.9 \\
经营活动现金流净额（亿元） & 75.94 & 196.07 & 118.07 & 64.65 & 4.22 \\
货币资金（亿元） & 42.66 & 47.66 & 10.13 & 24.97 & 36.75 \\
有息负债合计（亿元） & 326.61 & 259.72 & 193.85 & 215.37 & 286.24 \\
现金短债比（倍） & 0.47 & 0.67 & 0.25 & 0.41 & 0.25 \\
\bottomrule
\end{tabular}
\end{adjustbox}
\end{center}

\begin{center}
\begin{minipage}{\textwidth}\centering
\begin{tikzpicture}
\begin{groupplot}[
  group style={group size=2 by 1, horizontal sep=0.60cm},
  width=6.85cm, height=4.60cm,
  tick label style={font=\tiny},
  title style={font=\scriptsize\heiti},
  yticklabel style={font=\tiny},
  ylabel style={font=\tiny},
  grid=major, grid style={LightGray, line width=0.3pt},
  axis line style={InkGray, line width=0.5pt},
  enlarge x limits={abs=0.8},
  legend style={font=\scriptsize, at={(0.5,-0.20)}, anchor=north,
                legend columns=2, draw=none, fill=none, column sep=6pt},
]
\nextgroupplot[
  ybar, bar width=4pt,
  ymin=-253.09, ymax=1199.31,
  xtick={0,1,2,3,4,5.7},
  xticklabels={2021,2022,2023,2024,2025,2026H1},
  ylabel={亿元},
]
\addplot[fill=AgriGreen!75, draw=AgriGreen, bar shift=-2.6pt] table[x=i, y=rev]{data/clean/profiles/fin_300498.dat};
\addplot[fill=SoftRed!60, draw=SoftRed, bar shift=2.6pt] table[x=i, y=ni]{data/clean/profiles/fin_300498.dat};
\addplot[fill=AgriGreen!30, draw=AgriGreen, pattern=north east lines, bar shift=-2.6pt] table[x=x, y=rev]{data/clean/profiles/finh_300498.dat};
\addplot[fill=SoftRed!20, draw=SoftRed, pattern=north east lines, bar shift=2.6pt] table[x=x, y=ni]{data/clean/profiles/finh_300498.dat};
\legend{营业收入（年/半年）, 归母净利润（年/半年）}
\nextgroupplot[
  ymin=-42.0, ymax=73.9,
  xtick={0,1,2,3,4,5.7},
  xticklabels={2021,2022,2023,2024,2025,2026H1},
  ylabel={\%},
]
\addplot[AgriGold, mark=*, mark size=1.3pt] table[x=i, y=roe]{data/clean/profiles/fin_300498.dat};
\addplot[OilBlue, mark=square*, mark size=1.3pt] table[x=i, y=debt]{data/clean/profiles/fin_300498.dat};
\addplot[only marks, AgriGold, mark=*, mark size=2.0pt] table[x=x, y=roe]{data/clean/profiles/finh_300498.dat};
\addplot[only marks, OilBlue, mark=square*, mark size=2.0pt] table[x=x, y=debt]{data/clean/profiles/finh_300498.dat};
\legend{ROE, 资产负债率}
\end{groupplot}
\end{tikzpicture}
\begin{tikzpicture}
\begin{axis}[
  name=finbot,
  width=13.75cm, height=3.10cm, scale only axis,
  ymin=-44.96, ymax=503.57,
  xtick={0,1,2,3,4,5.7},
  xticklabels={2021,2022,2023,2024,2025,2026H1},
  tick label style={font=\tiny},
  yticklabel style={font=\tiny},
  ylabel={亿元}, ylabel style={font=\tiny},
  ybar, bar width=4pt,
  grid=major, grid style={LightGray, line width=0.3pt},
  axis line style={InkGray, line width=0.5pt},
  enlarge x limits={abs=0.8},
  legend style={font=\scriptsize, at={(0.5,-0.26)}, anchor=north,
                legend columns=2, draw=none, fill=none, column sep=10pt},
]
\addplot[fill=AgriGreen!55, draw=AgriGreen, bar shift=-3.0pt] table[x=i, y=cash]{data/clean/profiles/fin_300498.dat};
\addplot[fill=SoftRed!45, draw=SoftRed, bar shift=3.0pt] table[x=i, y=ibd]{data/clean/profiles/fin_300498.dat};
\addplot[fill=AgriGreen!25, draw=AgriGreen, pattern=north east lines, bar shift=-3.0pt] table[x=x, y=cash]{data/clean/profiles/finh_300498.dat};
\addplot[fill=SoftRed!20, draw=SoftRed, pattern=north east lines, bar shift=3.0pt] table[x=x, y=ibd]{data/clean/profiles/finh_300498.dat};
\legend{货币资金（左轴）, 有息负债（左轴）}
\end{axis}
\begin{axis}[
  at={(finbot.south east)}, anchor=south east,
  width=13.75cm, height=3.10cm, scale only axis,
  axis y line*=right, axis x line*=none, xtick=\empty,
  ymin=-50.98, ymax=254.89,
  tick label style={font=\tiny},
  yticklabel style={font=\tiny}, scaled y ticks=false,
  legend style={font=\scriptsize, at={(0.5,-0.44)}, anchor=north,
                legend columns=1, draw=none, fill=none},
]
\addplot[OilBlue, mark=*, mark size=2.2pt, line width=1.3pt] table[x=i, y=ocf]{data/clean/profiles/fin_300498.dat};
\addplot[only marks, OilBlue, mark=*, mark size=3.2pt] table[x=x, y=ocf]{data/clean/profiles/finh_300498.dat};
\legend{经营活动现金流净额（右轴，亿元，蓝点）}
\end{axis}
\end{tikzpicture}

\captionof{figure}[温氏股份（300498）近年主要财务指标]{温氏股份（300498）近年主要财务指标（左上：营业收入与归母净利润；右上：ROE 与资产负债率；下：货币资金、有息负债为左轴柱状，经营活动现金流净额为其右轴的蓝点折线。
2021---2025 年为年报口径，2026H1 为 2026 年半年报/半年末，与其前一年报之间留出间隔、以斜纹或空心点区分；半年报数据未年化）}
\end{minipage}
\end{center}

\pfl{财务分析} 2026 年 6 月末资产负债率 58.9\%，较 2025 年末上升 9.1 个百分点（样本中位数 52.2\%，所处三级行业内按由低到高排第\zhnumber{6}位）。 2026 年 6 月末货币资金 36.75 亿元、短期债务（短期借款与一年内到期非流动负债）146.20 亿元、长期债务 140.04 亿元，有息负债合计 286.24 亿元，现金短债比 0.25 倍——缺口明显。 净有息负债 249.48 亿元，相当于其 2026 年上半年营业收入的 53\%。 流动比率 0.83、速动比率 0.29，短期偿债指标偏紧。 2026 年上半年经营活动现金流净额 4.22 亿元，净利润为负而经营现金流为正，差额主要来自折旧摊销与营运资本变动，说明亏损尚未完全转化为现金流出。 2026 年 6 月末在建工程 28.48 亿元，较上年末增加 8.95 亿元，仍有在建投入。 2026 年 6 月末商誉 1.68 亿元，占归母净资产的 0.4\%，是净资产中最脆弱的一块。 综合看，账面货币资金对有息负债与短期债务的覆盖不足，滚动偿付对新增融资的依赖度较高；资产负债率高于样本中位数 6.7 个百分点；有息负债中短期债务占比更高，期限结构仍以滚动续作为主。

\pfl{重大战略转型} 公司的战略动作可归为以下几类：股权与股东回报（1 项）：2024 年 9 月，公司回购资金总额不低于 9.00 亿元且不超过 18.00 亿元。 公告口径上，近三年（2023 年 9 月---2026 年 9 月）公司披露重大重组与并购类公告 3 条、对外投资与转型类公告 1 条、回购与股权激励类公告 82 条，合计公告 689 条。 第一大行业仍是“养殖行业”，分部口径的收入（含内部交易） 131.1\%，与 2021 年年报相比没有方向性变化。 综合判断，报告期内公司有 4 条与战略相关的公告落地（重大重组与并购 3 条、对外投资与转型 1 条），资本运作与主业调整之间存在直接关联。

\begin{center}\small\setlength{\tabcolsep}{4pt}
\captionof{table}[温氏股份（300498）历史融资与资本运作]{温氏股份（300498）历史融资与资本运作（据公司公告与公开报道整理；金额为披露值，含首次公开发行、再融资、债券、银行借款与重整投资等）}
\begin{adjustbox}{max width=\textwidth}
\begin{tabular}{@{}p{1.45cm}p{2.55cm}p{4.05cm}p{4.95cm}@{}}
\toprule
时间 & 方式 & 规模 & 用途与说明 \\
\midrule
2017-03-20 & 公司债券（17 温氏 01，112506） & 5.00 亿元 & 公司经营资金；票面利率 4.60\% \\
2017-07-04 & 公司债券（17 温氏 02，112539） & 20.00 亿元 & 公司经营资金；票面利率 4.87\% \\
2020-07-01 & 公司债券（20 温氏 01） & 0.00 亿元 & 公司经营资金；起息日 2020 年 7 月 1 日 \\
2020-08-11 & 中期票据（20 温氏食品 MTN001） & 0.00 亿元 & 补充流动资金等；到期兑付日 2023 年 8 月 11 日 \\
2021-03-29 & 可转换公司债券（温氏转债，123107） & 97.00 亿元（92,\allowbreak{}970,\allowbreak{}000 张，\allowbreak{}每张 100 元） & 生猪养殖项目及补充流动资金；截至 2025 年 3 月底募投项目已使用 36.16 亿元、\allowbreak{}永久补充流动资金 27.09 亿元；2021 年 4 月 21 日上市 \\
2025-05 & 公司债券（注册） & 拟发行面值总额不超过 50.00 亿元、\allowbreak{}分期发行 & 包括一般公司债券、\allowbreak{}绿色公司债券、\allowbreak{}乡村振兴公司债券、\allowbreak{}科技创新公司债券等品种；2025 年 4 月 23 日董事会、\allowbreak{}2025 年 5 月 14 日 2024 年度股东大会审议通过；截至募集说明书签署时点尚未取得证监会注册 \\
\bottomrule
\end{tabular}
\end{adjustbox}
\end{center}

\pfl{历史融投资情况} 从现金流量表看，2025 年公司取得借款收到的现金 7.06 亿元、偿还债务支付的现金 77.76 亿元、吸收权益性投资 0.14 亿元，筹资活动现金流净额 -134.76 亿元（2023 年为取得借款 100.67 亿元、偿还债务 127.01 亿元）。 2026 年上半年筹资活动现金流净额为 62.81 亿元（取得借款 128.27 亿元、偿还债务 36.99 亿元，未年化），仍处于净融入状态。 2025 年分配股利、利润或偿付利息支付的现金合计 40.71 亿元。 公开披露的融资记录共 6 笔，工具类型以公司债、中期票据、可转债为主；金额与用途见下表。 其中规模最大的两笔为：2025 年 5 月，公司债，拟发行面值总额不超过 50.00 亿元、分期发行；2021 年 3 月，可转债，97.00 亿元（92,970,000 张，每张 100 元）。 股本方面，近三年披露股本变动 17 次（可转债转股 13 次、股权激励 3 次、增发新股上市 2 次、股份回购 1 次）；总股本由 2021 年末的 63.52 亿股变为 66.53 亿股（+4.74\%）。 分红方面，近三年现金分红 7 次、合计每股派现 1.25 元，最近一次实施在 2025 年 12 月。 近三年“再融资”类公告 59 条，涉及可转债。 综合看，公司融资结构上以债务滚动为主、股权融资为补充；2025 年筹资活动现金流净额 -134.76 亿元，融资节奏仍以维持存量债务周转为主，与前述经营与投资现金流的方向基本一致。

\pfl{未来潜在融资需求分析} 2026 年 6 月末货币资金 36.75 亿元，而短期债务 146.20 亿元（现金短债比 0.25 倍），2026 年上半年经营现金流净额 4.22 亿元。按“短期债务减去账面货币资金”的滚动偿付口径，静态缺口约 109.44 亿元；若再计入上半年亏损的年化额（87.32 亿元），维持一年正常经营所需的外部资金约 109.44---196.76 亿元。 公司 2026 年上半年归母净利润为负（-43.66 亿元），但 6 月末资产负债率 58.9\% 尚未进入第一档（亏损且负债率 $>$65\%）的区间，当前融资需求以补充流动资金、支撑低谷期经营性支出为主。 公开披露的后续资金安排线索包括：股权融资线索：2025---2026 年，公司拟公开发行不超过 50.00 亿元公司债券（分期发行）已进入注册程序。 资金需求还要叠加在建工程：期末在建工程 28.48 亿元、2026 年上半年购建固定资产等支出 56.80 亿元（未年化），这部分支出在周期底部通常会被主动放缓。 综合看，公司的外部资金需求以补充流动性与项目投入为主，滚动偿付口径下的静态缺口约 109.44 亿元，能否落地取决于授信续作与发行窗口的配合。




\subsubsection{正邦科技（002157）}
\label{co:002157}

\pfl{公司基本情况} 2007 年设立至今，正邦科技是曾经的江西养猪龙头企业，2022—2023 年陷入债务危机并被法院裁定重整，由双胞胎集团牵头联合体接盘；2024 年 6 月撤销退市风险警示，现依托双胞胎协同恢复生猪与饲料业务，控股股东承诺 4 年内完成资产整体上市。 公司于 2007-08-17 在深交所主板首发上市，注册资本 92.50 亿元，总股本 92.98 亿股。 截至 2026 年 10 月 12 日收盘价 2.90 元，对应总市值 269.6 亿元，近一年-3.6\%，总市值在三级行业“生猪养殖”的 11 家公司中按由大到小排第\zhnumber{3}位，近一年涨跌幅高于全样本中位数（-10.4\%）6.7 个百分点。 按 2026 年半年报披露的行业构成，收入占比前三位为养殖 53.0\%；饲料 44.3\%；其他(补充) 1.8\%。

\pfl{股权结构与控制权} 按 2026 年 6 月 30 日的股东名册，前三大股东为江西双胞胎农业有限公司 持股 15.13\%；南昌新振邦企业管理中心(有限合伙) 持股 5.95\%；江西正邦科技股份有限公司破产企业财产处置专用账户 持股 3.38\%，股东户数 179,149 户。 从控制权角度看，双胞胎控股持有双胞胎农业 100\% 股权，鲍洪星与华涛合计控制双胞胎控股 57.4094\% 股权；2023 年 7 月，南昌中院裁定公司破产重整，双胞胎农业作为产业投资人认购资本公积金转增的 14 亿股 A 股股票；2023 年 12 月，双胞胎集团由此完整接管上市公司治理体系。

\pfl{经营要点} 业务与产能层面，2026 年，1 至 8 月累计销售生猪 726.76 万头、同比上升 47.07\%，累计销售收入 52.75 亿元、同比下降 3.34\%；2026 年上半年，销售费用同比下降 18.05\%、管理费用同比下降 30.02\%。

\pfl{所处行业的基本面} 生猪的供给由能繁母猪存栏决定、需求由猪肉消费与替代品价格决定，本轮下行的成因在供给端的产能去化不彻底，而非需求收缩（第 \ref{sec:framework} 章的供需框架）。 行业层面的最新情况是：2026 年上半年，全国生猪出栏 37246 万头、同比增长 1.7\%，能繁母猪存栏 3780 万头、同比下降 6.5\%；2026 年 4 月，部分区域价格已先行走低至 8 元/公斤以下，低于行业现金成本；2026 年 5 月，农业农村部将能繁母猪正常保有量下调至 3750 万头左右。

\pfl{竞争优势} 把该公司放回同类公司比较，2026 年上半年营业收入 77.2 亿元、在三级行业“生猪养殖”的 11 家公司中排第\zhnumber{3}位；毛利率 -0.2\%、高于行业内中位数 1.4 个百分点；半年 ROE -7.0\%（未年化）、高于行业内中位数（-14.32\%）；资产负债率 53.5\%、低于行业内中位数（58.9\%）。 公司在公开资料中自述与同业比较的相对优势在于：上半年饲料收入逆势增长 28.47\% 并向外部拓展客户；2023 年，公司持续处于预重整及重整阶段，重整投资人合计提供投资款 43.40 亿元；16 家头部猪企成本分层清晰。 优势在当期体现为高于行业中位数的资本回报，但优势能否持续仍取决于周期与自身产能投放节奏。

\pfl{历史股价} 以 2021 年 1 月为起点、2026 年 10 月为终点，股价由 17.34 元变为 2.90 元，累计 -83.3\%，区间最高 17.34 元（2021 年 1 月）、最低 2.39 元（2023 年 5 月）；当前价相当于区间最高价的 16.7\%，为区间最低价的 1.21 倍。 月度收益的年化波动率 36.0\%，低于全样本中位数 37.3\%，在三级行业“生猪养殖”的 11 家公司中波动率由高到低排第\zhnumber{7}位。 把股价阶段与公司事件对齐看，2021 年 1 月 至 2023 年 5 月股价下跌-86.2\%，同期（2023 年 5 月）公司股票被深交所实施退市风险警示；2023 年 5 月 至 2026 年 4 月股价上涨+45.2\%，同期（2025 年 12 月）公司公告双胞胎已全面启动资产注入前期筹备。 从时间对应关系看，公司层面的资本运作与股权、风险事件解释了区间内多数急涨急跌，与行业指数的同步性相对有限。

\begin{center}
\begin{minipage}{\textwidth}\centering
\begin{tikzpicture}
\begin{axis}[
  width=12.6cm, height=3.9cm, scale only axis,
  xmin=2020.922, xmax=2026.972, ymin=2.22, ymax=18.55,
  xtick={2021,2022,2023,2024,2025,2026},
  yearticks,
  yticklabel style={font=\scriptsize},
  ylabel={元/股}, ylabel style={font=\scriptsize},
  grid=major, grid style={LightGray, line width=0.3pt},
  axis line style={InkGray, line width=0.5pt},
  every axis plot/.append style={line width=0.9pt},
]
\addplot[AgriGreen] table[x=x,y=y]{data/clean/profiles/px_002157.dat};
\addplot[SoftRed, dashed, line width=0.7pt] table[x=x,y=y]{data/clean/profiles/pxzz_002157.dat};
\addplot[only marks, mark=*, mark size=1.1pt, SoftRed] table[x=x,y=y]{data/clean/profiles/pxzz_002157.dat};
\end{axis}
\end{tikzpicture}

\captionof{figure}[正邦科技（002157）股价与周期骨架]{正邦科技（002157）月度收盘价（元/股）与 ZigZag 周期骨架（阈值 25\%，圆点为周期转折点）}
\end{minipage}
\end{center}

\pfl{过去周期分析} 以 25\% 的 ZigZag 阈值划分，样本区间内股价形成 1 个主升段与 1 个主跌段。 最大主升段出现在 2023 年 5 月 至 2026 年 4 月，35 个月+45.2\%；最大主跌段为 2021 年 1 月 至 2023 年 5 月，28 个月-86.2\%。 最近一次转折出现在 2026 年 4 月（上涨段自 2023 年 5 月起，35 个月+45.2\%），当前处于该段之后的震荡之中。 公司股价较申万农林牧渔指数提前约 36 个月见底，是板块中较早定价周期反转的公司之一。 生猪养殖的定价主线是产能去化与成本曲线的赛跑（第 \ref{sec:hog} 章）：2026 年三季度行业仍处在“产能去化未完成、价格磨底”的阶段，公司 2026 年 6 月末资产负债率 53.5\%、半年 ROE -7.0\%（未年化），在本轮下行中属于随行业同步调整的一类。 综合区间形态看，该股单段涨跌幅度偏大、以趋势段为主（最大主升 +45.2\%、最大主跌 -86.2\%），年化波动率 36.0\% 与全样本中位数 37.3\% 相比波动与样本中位数接近。

\pfl{盈利情况} 2026 年上半年营业收入 77.20 亿元（同比 +12.9\%），归母净利润 -7.36 亿元，由上年同期的盈利转为亏损。 以年报为趋势对照，2023---2025 年营业收入由 69.92 亿元变为 147.95 亿元（两年年均复合增速 +45.5\%），归母净利润由 85.29 亿元变为 -5.46 亿元。 按半年报口径，2026 年上半年的 ROE 为 -7.0\%（未年化）、销售净利率 -12.7\%、毛利率 -0.2\%，ROE 在“生猪养殖”11 家公司中按数值由高到低排第\zhnumber{1}位。 同期全样本半年 ROE 中位数 0.75\%，公司与之相差 7.74 个百分点（弱于中位数公司）。 公司披露的应对举措集中在：2025 年 12 月，公司公告双胞胎已全面启动资产注入前期筹备。 综合看，公司仍处亏损区间、由盈转亏，半年 ROE 在三级行业内按由高到低排第\zhnumber{1}位，单位盈利能力的修复依赖行业景气与自身成本端的改善，报表弹性主要体现为价格与销量的方向性变化。

\begin{center}\small\setlength{\tabcolsep}{3.4pt}
\captionof{table}[正邦科技（002157）关键财务指标]{正邦科技（002157）关键财务指标（合并报表口径；两个半年报期均未年化；现金短债比 $=$ 货币资金 $\div$ 短期债务）}
\begin{adjustbox}{max width=\textwidth}
\begin{tabular}{@{}lrrrrr@{}}
\toprule
指标 & 2023 年 & 2024 年 & 2025 年 & 2025 年半年报 & 2026 年半年报 \\
\midrule
营业收入（亿元） & 69.92 & 88.70 & 147.95 & 68.38 & 77.20 \\
归母净利润（亿元） & 85.29 & 2.16 & -5.46 & 2.02 & -7.36 \\
毛利率（\%） & -14.3 & 14.8 & 8.0 & 14.8 & -0.2 \\
ROE（\%） & 202.0 & 2.0 & -4.9 & 1.7 & -7.0 \\
资产负债率（\%） & 54.0 & 46.4 & 50.1 & 45.7 & 53.5 \\
经营活动现金流净额（亿元） & -13.83 & 1.74 & 8.30 & 0.20 & -4.07 \\
货币资金（亿元） & 24.42 & 33.21 & 31.20 & 34.18 & 16.97 \\
有息负债合计（亿元） & 23.19 & 18.28 & 18.17 & 17.17 & 19.10 \\
现金短债比（倍） & 3.42 & 6.05 & 5.29 & 6.80 & 2.43 \\
\bottomrule
\end{tabular}
\end{adjustbox}
\end{center}

\begin{center}
\begin{minipage}{\textwidth}\centering
\begin{tikzpicture}
\begin{groupplot}[
  group style={group size=2 by 1, horizontal sep=0.60cm},
  width=6.85cm, height=4.60cm,
  tick label style={font=\tiny},
  title style={font=\scriptsize\heiti},
  yticklabel style={font=\tiny},
  ylabel style={font=\tiny},
  grid=major, grid style={LightGray, line width=0.3pt},
  axis line style={InkGray, line width=0.5pt},
  enlarge x limits={abs=0.8},
  legend style={font=\scriptsize, at={(0.5,-0.20)}, anchor=north,
                legend columns=2, draw=none, fill=none, column sep=6pt},
]
\nextgroupplot[
  ybar, bar width=4pt,
  ymin=-254.68, ymax=556.49,
  xtick={0,1,2,3,4,5.7},
  xticklabels={2021,2022,2023,2024,2025,2026H1},
  ylabel={亿元},
]
\addplot[fill=AgriGreen!75, draw=AgriGreen, bar shift=-2.6pt] table[x=i, y=rev]{data/clean/profiles/fin_002157.dat};
\addplot[fill=SoftRed!60, draw=SoftRed, bar shift=2.6pt] table[x=i, y=ni]{data/clean/profiles/fin_002157.dat};
\addplot[fill=AgriGreen!30, draw=AgriGreen, pattern=north east lines, bar shift=-2.6pt] table[x=x, y=rev]{data/clean/profiles/finh_002157.dat};
\addplot[fill=SoftRed!20, draw=SoftRed, pattern=north east lines, bar shift=2.6pt] table[x=x, y=ni]{data/clean/profiles/finh_002157.dat};
\legend{营业收入（年/半年）, 归母净利润（年/半年）}
\nextgroupplot[
  ymin=-4409.1, ymax=628.9,
  xtick={0,1,2,3,4,5.7},
  xticklabels={2021,2022,2023,2024,2025,2026H1},
  ylabel={\%},
]
\addplot[AgriGold, mark=*, mark size=1.3pt] table[x=i, y=roe]{data/clean/profiles/fin_002157.dat};
\addplot[OilBlue, mark=square*, mark size=1.3pt] table[x=i, y=debt]{data/clean/profiles/fin_002157.dat};
\addplot[only marks, AgriGold, mark=*, mark size=2.0pt] table[x=x, y=roe]{data/clean/profiles/finh_002157.dat};
\addplot[only marks, OilBlue, mark=square*, mark size=2.0pt] table[x=x, y=debt]{data/clean/profiles/finh_002157.dat};
\legend{ROE, 资产负债率}
\end{groupplot}
\end{tikzpicture}
\begin{tikzpicture}
\begin{axis}[
  name=finbot,
  width=13.75cm, height=3.10cm, scale only axis,
  ymin=-28.73, ymax=321.77,
  xtick={0,1,2,3,4,5.7},
  xticklabels={2021,2022,2023,2024,2025,2026H1},
  tick label style={font=\tiny},
  yticklabel style={font=\tiny},
  ylabel={亿元}, ylabel style={font=\tiny},
  ybar, bar width=4pt,
  grid=major, grid style={LightGray, line width=0.3pt},
  axis line style={InkGray, line width=0.5pt},
  enlarge x limits={abs=0.8},
  legend style={font=\scriptsize, at={(0.5,-0.26)}, anchor=north,
                legend columns=2, draw=none, fill=none, column sep=10pt},
]
\addplot[fill=AgriGreen!55, draw=AgriGreen, bar shift=-3.0pt] table[x=i, y=cash]{data/clean/profiles/fin_002157.dat};
\addplot[fill=SoftRed!45, draw=SoftRed, bar shift=3.0pt] table[x=i, y=ibd]{data/clean/profiles/fin_002157.dat};
\addplot[fill=AgriGreen!25, draw=AgriGreen, pattern=north east lines, bar shift=-3.0pt] table[x=x, y=cash]{data/clean/profiles/finh_002157.dat};
\addplot[fill=SoftRed!20, draw=SoftRed, pattern=north east lines, bar shift=3.0pt] table[x=x, y=ibd]{data/clean/profiles/finh_002157.dat};
\legend{货币资金（左轴）, 有息负债（左轴）}
\end{axis}
\begin{axis}[
  at={(finbot.south east)}, anchor=south east,
  width=13.75cm, height=3.10cm, scale only axis,
  axis y line*=right, axis x line*=none, xtick=\empty,
  ymin=-71.05, ymax=27.19,
  tick label style={font=\tiny},
  yticklabel style={font=\tiny}, scaled y ticks=false,
  legend style={font=\scriptsize, at={(0.5,-0.44)}, anchor=north,
                legend columns=1, draw=none, fill=none},
]
\addplot[OilBlue, mark=*, mark size=2.2pt, line width=1.3pt] table[x=i, y=ocf]{data/clean/profiles/fin_002157.dat};
\addplot[only marks, OilBlue, mark=*, mark size=3.2pt] table[x=x, y=ocf]{data/clean/profiles/finh_002157.dat};
\legend{经营活动现金流净额（右轴，亿元，蓝点）}
\end{axis}
\end{tikzpicture}

\captionof{figure}[正邦科技（002157）近年主要财务指标]{正邦科技（002157）近年主要财务指标（左上：营业收入与归母净利润；右上：ROE 与资产负债率；下：货币资金、有息负债为左轴柱状，经营活动现金流净额为其右轴的蓝点折线。
2021---2025 年为年报口径，2026H1 为 2026 年半年报/半年末，与其前一年报之间留出间隔、以斜纹或空心点区分；半年报数据未年化）}
\end{minipage}
\end{center}

\pfl{财务分析} 2026 年 6 月末资产负债率 53.5\%，较 2025 年末上升 3.4 个百分点（样本中位数 52.2\%，所处三级行业内按由低到高排第\zhnumber{4}位）。 2026 年 6 月末货币资金 16.97 亿元、短期债务（短期借款与一年内到期非流动负债）6.97 亿元、长期债务 12.12 亿元，有息负债合计 19.10 亿元，现金短债比 2.43 倍——覆盖充分。 净有息负债 2.13 亿元，相当于其 2026 年上半年营业收入的 3\%。 流动比率 0.74、速动比率 0.44，短期偿债指标偏紧。 2026 年上半年经营活动现金流净额 -4.07 亿元，经营现金流与利润同时为负，现金消耗较快。 2026 年 6 月末在建工程 26.34 亿元，较上年末减少 2.56 亿元，在建投入在收缩。 综合看，现金短债比 2.43 倍，短期偿付压力可控；资产负债率高于样本中位数 1.3 个百分点；有息负债中长期债务占比更高，期限结构对短债滚动的压力较小。 公司披露的重整与债务违约风险为：2024 年末，已逾期未偿还短期借款 1.68 亿元（重整后遗留部分子公司借款逾期事项）。

\pfl{重大战略转型} 近三年公司的战略动作集中在：并购与重组（4 项）：2024 年 5 月，双胞胎集团承诺：重整完成后 24 个月内逐步启动自身生猪养殖、饲料等业务资产置入上市公司的重组程序；2023 年 11 月，公司通过新设合伙企业以现金及抵债资产出资方式参与重整。 公告口径上，近三年（2023 年 9 月---2026 年 9 月）公司披露重大重组与并购类公告 0 条、对外投资与转型类公告 6 条、回购与股权激励类公告 14 条，合计公告 475 条。 第一大行业仍是“养殖”，收入占比 53.0\%，与 2021 年年报相比没有方向性变化。 综合判断，报告期内公司有 6 条与战略相关的公告落地（重大重组与并购 0 条、对外投资与转型 6 条），资本运作与主业调整之间存在直接关联。

\begin{center}\small\setlength{\tabcolsep}{4pt}
\captionof{table}[正邦科技（002157）历史融资与资本运作]{正邦科技（002157）历史融资与资本运作（据公司公告与公开报道整理；金额为披露值，含首次公开发行、再融资、债券、银行借款与重整投资等）}
\begin{adjustbox}{max width=\textwidth}
\begin{tabular}{@{}p{1.45cm}p{2.55cm}p{4.05cm}p{4.95cm}@{}}
\toprule
时间 & 方式 & 规模 & 用途与说明 \\
\midrule
2007-08 & IPO（深交所） & 1,\allowbreak{}900 万股，\allowbreak{}发行价 11.09 元/\allowbreak{}股（约 2.11 亿元） & 饲料及养殖项目建设；发行后总股本 7,\allowbreak{}529.68 万股 \\
2020 & 可转换公司债券（正邦转债，128114） & 16.00 亿元 & 养殖产能扩张；重整期间债券持有人会议决议：正邦转债最后交易日为 2023 年 8 月 4 日 \\
2021-02-01 & 定向增发（2020 年度非公开发行） & 5.70 亿股，\allowbreak{}发行价 13.16 元/\allowbreak{}股 & 发展生猪养殖和补充流动资金；2020 年 6 月 29 日董事会通过预案 \\
2023-12 & 破产重整投资款（转增股票引入重整投资人） & 重整投资人合计提供资金 43.40 亿元 & 支付重整费用、\allowbreak{}共益债务、\allowbreak{}清偿各类债务及补充上市公司流动资金；重整计划中 31.5 亿股转增股票用于引入重整投资人：产业投资人（双胞胎指定主体）以 1.1 元/\allowbreak{}股受让 14 亿股 \\
2026 & 关联方借款额度 & 任一时点借款余额不超过 15.00 亿元 & 满足短期内、\allowbreak{}临时性补充流动资金需求；向控股股东江西双胞胎农业有限公司同一控制下的关联方双胞胎畜牧集团有限公司申请 \\
\bottomrule
\end{tabular}
\end{adjustbox}
\end{center}

\pfl{历史融投资情况} 从现金流量表看，2025 年公司偿还债务支付的现金 0.01 亿元、吸收权益性投资 0.08 亿元，筹资活动现金流净额 -2.90 亿元（2023 年为取得借款 7.64 亿元、偿还债务 1.91 亿元）。 2026 年上半年筹资活动现金流净额为 1.89 亿元（取得借款 — 亿元、偿还债务 — 亿元，未年化），仍处于净融入状态。 2025 年分配股利、利润或偿付利息支付的现金合计 0.26 亿元。 公开披露的融资记录共 5 笔，工具类型以 IPO、可转债、定向增发、重整投资等为主；金额与用途见下表。 其中规模最大的两笔为：2026 年，银行借款，任一时点借款余额不超过 15.00 亿元；2023 年 12 月，重整投资，重整投资人合计提供资金 43.40 亿元。 股本方面，近三年披露股本变动 10 次（可转债转股 6 次、转增 1 次、限售股份上市 1 次）；总股本由 2021 年末的 31.46 亿股变为 92.98 亿股（+195.51\%）。 分红方面，近三年未实施现金分红，在全部样本中属于分红能力偏弱的一端。 近三年“再融资”类公告 12 条，涉及定向增发、可转债。 综合看，公司融资结构上以债务滚动为主、股权融资为补充；2025 年筹资活动现金流净额 -2.90 亿元，融资节奏仍以维持存量债务周转为主，与前述经营与投资现金流的方向基本一致。

\pfl{未来潜在融资需求分析} 2026 年 6 月末货币资金 16.97 亿元，而短期债务 6.97 亿元（现金短债比 2.43 倍），2026 年上半年经营现金流净额 -4.07 亿元。账面货币资金可以覆盖短期债务，缺口不体现在静态偿付层面。 公司 2026 年上半年归母净利润为负（-7.36 亿元），但 6 月末资产负债率 53.5\% 尚未进入第一档（亏损且负债率 $>$65\%）的区间，当前融资需求以补充流动资金、支撑低谷期经营性支出为主。 公开披露的后续资金安排线索包括：债务与授信安排：2026 年度，资产负债率 70\% 以上的子公司可享担保总额度 20.00 亿元。 资金需求还要叠加在建工程：期末在建工程 26.34 亿元、2026 年上半年购建固定资产等支出 6.94 亿元（未年化），这部分支出在周期底部通常会被主动放缓。 静态偿付不存在缺口，后续融资更可能服务于产能或并购安排，而不是填补流动性窟窿。




\subsubsection{圣农发展（002299）}
\label{co:002299}

\pfl{公司基本情况} 圣农发展成立于 2009 年，国内白羽肉鸡全产业链龙头，覆盖自主育种、饲料、养殖、屠宰与深加工，产能规模居行业前列；2025 年完成对原嘉吉中国白羽鸡项目（太阳谷）的全资控股，2026 年上半年营收突破百亿元、扣非净利同比增长近四成。 公司于 2009-10-21 在深交所主板首发上市，注册资本 12.43 亿元，总股本 12.43 亿股。 截至 2026 年 10 月 12 日收盘价 15.85 元，对应总市值 197.1 亿元，近一年-11.6\%，总市值在三级行业“肉鸡养殖”的 8 家公司中按由大到小排第\zhnumber{1}位，近一年涨跌幅低于全样本中位数（-10.4\%）1.2 个百分点。 按 2026 年半年报披露的行业构成，收入占比前三位为家禽饲养加工行业 49.8\%；食品加工行业 45.5\%；其他(补充) 4.7\%。

\pfl{股权结构与控制权} 按 2026 年 6 月 30 日的股东名册，前三大股东为福建圣农控股集团有限公司（持股比例未披露）；环胜信息技术(上海)有限公司 持股 5.01\%；傅长玉 持股 2.68\%，股东户数 34,847 户。 股东结构与控制权方面，2025 年 7 月，控股股东圣农集团与一致行动人傅露芳通过大宗交易减持 2,003.95 万股、占总股本 1.6117\%；第二大股东环胜信息技术（上海）有限公司持股 5.01\%（6,222.05 万股）。

\pfl{经营要点} 从经营要点看，2026 年上半年，鸡肉生食销售量 73.2 万吨、同比增长 10.3\%；白羽鸡养殖产能近 8 亿羽，已建及在建食品深加工产能合计超 60 万吨、位列全国第一。

\pfl{所处行业的基本面} 白羽与黄羽肉鸡的供给取决于祖代引种量与父母代存栏，需求端受餐饮与替代效应驱动，产能周期只有 6---12 个月（第 \ref{sec:framework} 章的供需框架）。 从供需两端的当期数据看，2026 年，当月父母代鸡苗销量 845.15 万套、同比增长 19.49\%；2021---2025 年，国产占比由 10.6\% 升至 34.5\%；2026 年上半年，全国家禽出栏 87.6 亿只、同比增加 6.3 亿只、增长 7.7\%。

\pfl{竞争优势} 把该公司放回同类公司比较，2026 年上半年营业收入 107.8 亿元、在三级行业“肉鸡养殖”的 8 家公司中排第\zhnumber{1}位；毛利率 13.1\%、高于行业内中位数 2.3 个百分点；半年 ROE 4.4\%（未年化）、高于行业内中位数（1.21\%）；资产负债率 51.9\%、低于行业内中位数（52.1\%）。 公司在公开资料中自述与同业比较的相对优势在于：2022 年 6 月起批量出售父母代种鸡，为国内首个打破国外种源垄断的企业。 优势在当期体现为高于行业中位数的资本回报，但优势能否持续仍取决于周期与自身产能投放节奏。

\pfl{历史股价} 自 2021 年 1 月至 2026 年 10 月，股价由 28.65 元变为 15.85 元，累计 -44.7\%，区间最高 28.80 元（2021 年 2 月）、最低 12.17 元（2024 年 8 月）；当前价相当于区间最高价的 55.0\%，为区间最低价的 1.30 倍。 月度收益的年化波动率 31.2\%，低于全样本中位数 37.3\%，在三级行业“肉鸡养殖”的 8 家公司中波动率由高到低排第\zhnumber{8}位。 把股价阶段与公司事件对齐看，2024 年 8 月 至 2026 年 2 月股价上涨+58.6\%，同期（2025 年 5 月）公司完成太阳谷剩余 54\% 股权收购，公司产能向 10 亿羽/年目标推进。 整体上看，区间的几处大级别拐点都能在公司的资本运作、股权变动或风险事件上找到对应，股价波动有明显的公司层面驱动，而非单纯的行业 beta。

\begin{center}
\begin{minipage}{\textwidth}\centering
\begin{tikzpicture}
\begin{axis}[
  width=12.6cm, height=3.9cm, scale only axis,
  xmin=2020.922, xmax=2026.972, ymin=11.32, ymax=30.82,
  xtick={2021,2022,2023,2024,2025,2026},
  yearticks,
  yticklabel style={font=\scriptsize},
  ylabel={元/股}, ylabel style={font=\scriptsize},
  grid=major, grid style={LightGray, line width=0.3pt},
  axis line style={InkGray, line width=0.5pt},
  every axis plot/.append style={line width=0.9pt},
]
\addplot[AgriGreen] table[x=x,y=y]{data/clean/profiles/px_002299.dat};
\addplot[SoftRed, dashed, line width=0.7pt] table[x=x,y=y]{data/clean/profiles/pxzz_002299.dat};
\addplot[only marks, mark=*, mark size=1.1pt, SoftRed] table[x=x,y=y]{data/clean/profiles/pxzz_002299.dat};
\end{axis}
\end{tikzpicture}

\captionof{figure}[圣农发展（002299）股价与周期骨架]{圣农发展（002299）月度收盘价（元/股）与 ZigZag 周期骨架（阈值 25\%，圆点为周期转折点）}
\end{minipage}
\end{center}

\pfl{过去周期分析} 以 25\% 的 ZigZag 阈值划分，样本区间内股价形成 3 个主升段与 3 个主跌段。 最大主升段出现在 2024 年 8 月 至 2026 年 2 月，18 个月+58.6\%；最大主跌段为 2022 年 11 月 至 2024 年 8 月，21 个月-51.4\%。 最近一次转折出现在 2026 年 2 月（上涨段自 2024 年 8 月起，18 个月+58.6\%），当前处于该段之后的震荡之中。 公司股价较申万农林牧渔指数提前约 21 个月见底，是板块中较早定价周期反转的公司之一。 禽链的产能周期只有 6---12 个月、反转快于生猪（第 \ref{sec:other-livestock} 节），2026 年前四个月毛鸡养殖利润同比大幅回升，链条利润的修复先于生猪，公司 2026 年上半年实现盈利、半年 ROE 4.4\%（未年化），好于全样本中位数（0.75\%），下行周期对其报表的冲击小于同业。 综合区间形态看，该股单段涨跌幅度偏大、以趋势段为主（最大主升 +58.6\%、最大主跌 -51.4\%），年化波动率 31.2\% 与全样本中位数 37.3\% 相比波动小于样本中位数。

\pfl{盈利情况} 2026 年上半年营业收入 107.81 亿元（同比 +21.7\%），归母净利润 5.02 亿元，同比 -44.8\%。 以年报为趋势对照，2023---2025 年营业收入由 184.87 亿元变为 200.94 亿元（两年年均复合增速 +4.3\%），归母净利润由 6.64 亿元变为 13.80 亿元。 按半年报口径，2026 年上半年的 ROE 为 4.4\%（未年化）、销售净利率 4.7\%、毛利率 13.1\%，ROE 在“肉鸡养殖”8 家公司中按数值由高到低排第\zhnumber{2}位。 同期全样本半年 ROE 中位数 0.75\%，公司与之相差 3.65 个百分点（略好于中位数公司）。 综合看，公司在报告期维持盈利（高于全样本半年 ROE 中位数 0.75\%），利润的可持续性取决于肉鸡养殖景气的延续与成本管控的执行。

\begin{center}\small\setlength{\tabcolsep}{3.4pt}
\captionof{table}[圣农发展（002299）关键财务指标]{圣农发展（002299）关键财务指标（合并报表口径；两个半年报期均未年化；现金短债比 $=$ 货币资金 $\div$ 短期债务）}
\begin{adjustbox}{max width=\textwidth}
\begin{tabular}{@{}lrrrrr@{}}
\toprule
指标 & 2023 年 & 2024 年 & 2025 年 & 2025 年半年报 & 2026 年半年报 \\
\midrule
营业收入（亿元） & 184.87 & 185.86 & 200.94 & 88.56 & 107.81 \\
归母净利润（亿元） & 6.64 & 7.24 & 13.80 & 9.10 & 5.02 \\
毛利率（\%） & 10.6 & 11.1 & 12.6 & 12.0 & 13.1 \\
ROE（\%） & 6.6 & 7.0 & 12.6 & 8.4 & 4.4 \\
资产负债率（\%） & 54.1 & 50.0 & 51.7 & 55.2 & 51.9 \\
经营活动现金流净额（亿元） & 22.91 & 30.49 & 40.32 & 14.25 & 9.14 \\
货币资金（亿元） & 7.38 & 7.54 & 14.47 & 27.55 & 11.73 \\
有息负债合计（亿元） & 81.38 & 59.17 & 67.99 & 79.71 & 76.86 \\
现金短债比（倍） & 0.11 & 0.15 & 0.24 & 0.38 & 0.17 \\
\bottomrule
\end{tabular}
\end{adjustbox}
\end{center}

\begin{center}
\begin{minipage}{\textwidth}\centering
\begin{tikzpicture}
\begin{groupplot}[
  group style={group size=2 by 1, horizontal sep=0.60cm},
  width=6.85cm, height=4.60cm,
  tick label style={font=\tiny},
  title style={font=\scriptsize\heiti},
  yticklabel style={font=\tiny},
  ylabel style={font=\tiny},
  grid=major, grid style={LightGray, line width=0.3pt},
  axis line style={InkGray, line width=0.5pt},
  enlarge x limits={abs=0.8},
  legend style={font=\scriptsize, at={(0.5,-0.20)}, anchor=north,
                legend columns=2, draw=none, fill=none, column sep=6pt},
]
\nextgroupplot[
  ybar, bar width=4pt,
  ymin=-20.09, ymax=225.06,
  xtick={0,1,2,3,4,5.7},
  xticklabels={2021,2022,2023,2024,2025,2026H1},
  ylabel={亿元},
]
\addplot[fill=AgriGreen!75, draw=AgriGreen, bar shift=-2.6pt] table[x=i, y=rev]{data/clean/profiles/fin_002299.dat};
\addplot[fill=SoftRed!60, draw=SoftRed, bar shift=2.6pt] table[x=i, y=ni]{data/clean/profiles/fin_002299.dat};
\addplot[fill=AgriGreen!30, draw=AgriGreen, pattern=north east lines, bar shift=-2.6pt] table[x=x, y=rev]{data/clean/profiles/finh_002299.dat};
\addplot[fill=SoftRed!20, draw=SoftRed, pattern=north east lines, bar shift=2.6pt] table[x=x, y=ni]{data/clean/profiles/finh_002299.dat};
\legend{营业收入（年/半年）, 归母净利润（年/半年）}
\nextgroupplot[
  ymin=-4.3, ymax=59.6,
  xtick={0,1,2,3,4,5.7},
  xticklabels={2021,2022,2023,2024,2025,2026H1},
  ylabel={\%},
]
\addplot[AgriGold, mark=*, mark size=1.3pt] table[x=i, y=roe]{data/clean/profiles/fin_002299.dat};
\addplot[OilBlue, mark=square*, mark size=1.3pt] table[x=i, y=debt]{data/clean/profiles/fin_002299.dat};
\addplot[only marks, AgriGold, mark=*, mark size=2.0pt] table[x=x, y=roe]{data/clean/profiles/finh_002299.dat};
\addplot[only marks, OilBlue, mark=square*, mark size=2.0pt] table[x=x, y=debt]{data/clean/profiles/finh_002299.dat};
\legend{ROE, 资产负债率}
\end{groupplot}
\end{tikzpicture}
\begin{tikzpicture}
\begin{axis}[
  name=finbot,
  width=13.75cm, height=3.10cm, scale only axis,
  ymin=-8.14, ymax=91.14,
  xtick={0,1,2,3,4,5.7},
  xticklabels={2021,2022,2023,2024,2025,2026H1},
  tick label style={font=\tiny},
  yticklabel style={font=\tiny},
  ylabel={亿元}, ylabel style={font=\tiny},
  ybar, bar width=4pt,
  grid=major, grid style={LightGray, line width=0.3pt},
  axis line style={InkGray, line width=0.5pt},
  enlarge x limits={abs=0.8},
  legend style={font=\scriptsize, at={(0.5,-0.26)}, anchor=north,
                legend columns=2, draw=none, fill=none, column sep=10pt},
]
\addplot[fill=AgriGreen!55, draw=AgriGreen, bar shift=-3.0pt] table[x=i, y=cash]{data/clean/profiles/fin_002299.dat};
\addplot[fill=SoftRed!45, draw=SoftRed, bar shift=3.0pt] table[x=i, y=ibd]{data/clean/profiles/fin_002299.dat};
\addplot[fill=AgriGreen!25, draw=AgriGreen, pattern=north east lines, bar shift=-3.0pt] table[x=x, y=cash]{data/clean/profiles/finh_002299.dat};
\addplot[fill=SoftRed!20, draw=SoftRed, pattern=north east lines, bar shift=3.0pt] table[x=x, y=ibd]{data/clean/profiles/finh_002299.dat};
\legend{货币资金（左轴）, 有息负债（左轴）}
\end{axis}
\begin{axis}[
  at={(finbot.south east)}, anchor=south east,
  width=13.75cm, height=3.10cm, scale only axis,
  axis y line*=right, axis x line*=none, xtick=\empty,
  ymin=-10.48, ymax=52.42,
  tick label style={font=\tiny},
  yticklabel style={font=\tiny}, scaled y ticks=false,
  legend style={font=\scriptsize, at={(0.5,-0.44)}, anchor=north,
                legend columns=1, draw=none, fill=none},
]
\addplot[OilBlue, mark=*, mark size=2.2pt, line width=1.3pt] table[x=i, y=ocf]{data/clean/profiles/fin_002299.dat};
\addplot[only marks, OilBlue, mark=*, mark size=3.2pt] table[x=x, y=ocf]{data/clean/profiles/finh_002299.dat};
\legend{经营活动现金流净额（右轴，亿元，蓝点）}
\end{axis}
\end{tikzpicture}

\captionof{figure}[圣农发展（002299）近年主要财务指标]{圣农发展（002299）近年主要财务指标（左上：营业收入与归母净利润；右上：ROE 与资产负债率；下：货币资金、有息负债为左轴柱状，经营活动现金流净额为其右轴的蓝点折线。
2021---2025 年为年报口径，2026H1 为 2026 年半年报/半年末，与其前一年报之间留出间隔、以斜纹或空心点区分；半年报数据未年化）}
\end{minipage}
\end{center}

\pfl{财务分析} 2026 年 6 月末资产负债率 51.9\%，较 2025 年末上升 0.3 个百分点（样本中位数 52.2\%，所处三级行业内按由低到高排第\zhnumber{4}位）。 2026 年 6 月末货币资金 11.73 亿元、短期债务（短期借款与一年内到期非流动负债）69.56 亿元、长期债务 7.30 亿元，有息负债合计 76.86 亿元，现金短债比 0.17 倍——缺口明显。 净有息负债 65.13 亿元，相当于其 2026 年上半年营业收入的 60\%。 流动比率 0.67、速动比率 0.36，短期偿债指标偏紧。 2026 年上半年经营活动现金流净额 9.14 亿元，经营现金流与利润同向为正，内生资金可以支撑运营。 2026 年 6 月末在建工程 4.40 亿元，较上年末增加 0.78 亿元，仍有在建投入。 2026 年 6 月末商誉 10.16 亿元，占归母净资产的 8.8\%，是净资产中最脆弱的一块。 综合看，账面货币资金对有息负债与短期债务的覆盖不足，滚动偿付对新增融资的依赖度较高；资产负债率低于样本中位数 0.3 个百分点；有息负债中短期债务占比更高，期限结构仍以滚动续作为主。

\pfl{重大战略转型} 从公开披露看，公司的战略推进集中在：并购与重组（2 项）：2025 年 5 月，公司完成太阳谷剩余 54\% 股权收购，公司产能向 10 亿羽/年目标推进；2023 年 11 月，从德弘资本手中受让 46\% 股权并受让可转股债权，向上游产能扩张。 公告口径上，近三年（2023 年 9 月---2026 年 9 月）公司披露重大重组与并购类公告 3 条、对外投资与转型类公告 3 条、回购与股权激励类公告 32 条，合计公告 382 条。 第一大行业“家禽饲养加工行业”的收入占比由 2021 年年报的 61.7\% 变为 49.8\%（11.9 个百分点），收入结构出现再平衡。 综合判断，报告期内公司有 6 条与战略相关的公告落地（重大重组与并购 3 条、对外投资与转型 3 条），资本运作与主业调整之间存在直接关联。

\begin{center}\small\setlength{\tabcolsep}{4pt}
\captionof{table}[圣农发展（002299）历史融资与资本运作]{圣农发展（002299）历史融资与资本运作（据公司公告与公开报道整理；金额为披露值，含首次公开发行、再融资、债券、银行借款与重整投资等）}
\begin{adjustbox}{max width=\textwidth}
\begin{tabular}{@{}p{1.45cm}p{2.55cm}p{4.05cm}p{4.95cm}@{}}
\toprule
时间 & 方式 & 规模 & 用途与说明 \\
\midrule
2009-10 & IPO（深交所中小板） & 4,\allowbreak{}100 万股，\allowbreak{}发行价 19.75 元/\allowbreak{}股 & 优化各生产环节产能（自养自宰模式下的养殖与屠宰项目）；网下配售 820 万股、\allowbreak{}网上发行 3,\allowbreak{}280 万股 \\
2017-10 & 发行股份购买资产（关联交易） & 交易价格 20.20 亿元，\allowbreak{}发行 1.29 亿股，\allowbreak{}发行价 15.71 元/\allowbreak{}股 & 购买圣农实业等 8 名交易对方合计持有的圣农食品 100\% 股权；圣农食品 100\% 股权评估值 20.20 亿元；2017 年 10 月 23 日完成资产过户 \\
2023-11-13 & 对外投资（收购股权及可转股债权） & 圣农发展以 2.45 亿元受让太阳谷 45.1557\% 股权 & 取得原嘉吉中国白羽鸡项目相关主体的股权与债权；公司第六届董事会第二十三次会议审议通过 \\
2025-05-21 & 现金收购股权 & 转让总价款 11.26 亿元 & 收购德弘钰玺、\allowbreak{}Ancient Steel (HK) Limited 及廖俊杰先生持有的太阳谷合计 54\% 股权 \\
\bottomrule
\end{tabular}
\end{adjustbox}
\end{center}

\pfl{历史融投资情况} 从现金流量表看，2025 年公司取得借款收到的现金 113.31 亿元、偿还债务支付的现金 110.25 亿元、吸收权益性投资 0.56 亿元，筹资活动现金流净额 -11.05 亿元（2023 年为取得借款 77.43 亿元、偿还债务 65.70 亿元）。 2026 年上半年筹资活动现金流净额为 2.88 亿元（取得借款 59.33 亿元、偿还债务 53.29 亿元，未年化），仍处于净融入状态。 2025 年分配股利、利润或偿付利息支付的现金合计 10.70 亿元。 公开披露的融资记录共 4 笔，工具类型以 IPO、发行股份购买资产为主；金额与用途见下表。 其中规模最大的两笔为：2009 年 10 月，IPO，4,100 万股，发行价 19.75 元/股；2025 年 5 月，现金收购股权，转让总价款 11.26 亿元。 股本方面，近三年披露股本变动 9 次（股份回购 3 次、激励股份解禁 2 次、其他 1 次）；总股本由 2021 年末的 12.44 亿股变为 12.43 亿股（-0.05\%）。 分红方面，近三年现金分红 6 次、合计每股派现 1.40 元，最近一次实施在 2025 年 12 月。 综合看，公司融资结构上以债务滚动为主、股权融资为补充；2025 年筹资活动现金流净额 -11.05 亿元，融资节奏仍以维持存量债务周转为主，与前述经营与投资现金流的方向基本一致。

\pfl{未来潜在融资需求分析} 2026 年 6 月末货币资金 11.73 亿元，而短期债务 69.56 亿元（现金短债比 0.17 倍），2026 年上半年经营现金流净额 9.14 亿元。按“短期债务减去账面货币资金”的滚动偿付口径，静态缺口约 57.83 亿元；上半年盈利或接近盈亏平衡，该缺口不随经营亏损进一步扩大。 公司 2026 年 6 月末资产负债率 51.9\%、2026 年上半年 ROE 4.4\%（未年化），现金短债比 0.17 倍、短债覆盖不足，融资需求首先用于置换与滚动存量债务、补充营运资金，属于“补血型”而非扩张型。 资金需求还要叠加在建工程：期末在建工程 4.40 亿元、2026 年上半年购建固定资产等支出 8.64 亿元（未年化），这部分支出在周期底部通常会被主动放缓。 综合看，公司的外部资金需求以营运资金周转与在建项目投入为主，滚动偿付口径下的静态缺口约 57.83 亿元，能否落地取决于授信续作与发行窗口的配合。




\subsubsection{立华股份（300761）}
\label{co:300761}

\pfl{公司基本情况} 公司前身可追溯至 1997 年，全国第二大黄羽鸡养殖企业、江苏首家上市畜禽养殖企业，2019 年创业板上市，近年形成“黄羽鸡＋生猪”双轮驱动；2026 年上半年猪价深跌致整体由盈转亏，黄鸡板块盈利形成对冲。 公司于 2019-02-18 在深交所创业板首发上市，注册资本 8.46 亿元，总股本 8.28 亿股。 截至 2026 年 10 月 12 日收盘价 20.31 元，对应总市值 168.1 亿元，近一年-11.5\%，总市值在三级行业“肉鸡养殖”的 8 家公司中按由大到小排第\zhnumber{2}位，近一年涨跌幅低于全样本中位数（-10.4\%）1.1 个百分点。 按 2025 年年报披露的行业构成，收入几乎全部来自养殖行业 100.0\%。

\pfl{股权结构与控制权} 按 2026 年 6 月 30 日的股东名册，前三大股东为程立力 持股 29.44\%；常州市奔腾牧业技术服务中心(有限合伙) 持股 17.32\%；林州昊鸣技术服务中心(有限合伙) 持股 5.41\%，股东户数 15,343 户。 股东结构与控制权方面，公司质押 8,335.00 万股，该合伙企业持股 17.32\%、质押 6,520.00 万股；报告期内公司控股股东、实际控制人未发生变更。

\pfl{经营要点} 从经营要点看，2026 年上半年，肉鸡完全成本约 11.20 元/公斤、同比微增；黄羽鸡出栏量居全国第二。

\pfl{所处行业的基本面} 白羽与黄羽肉鸡的供给取决于祖代引种量与父母代存栏，需求端受餐饮与替代效应驱动，产能周期只有 6---12 个月（第 \ref{sec:framework} 章的供需框架）。 当期该行业的关键数据是：2026 年上半年，华东地区商品毛鸡出栏均价仅 9.5 元/公斤，而行业平均养殖成本约 12 元/公斤，中小产能加速出清；2026 年，第 8 周黄羽在产父母代存栏降至 1,299.7 万套、处 2022 年以来同期低位；2026 年上半年，全国家禽出栏 87.6 亿只、同比增长 7.7\%。

\pfl{竞争优势} 把该公司放回同类公司比较，2026 年上半年营业收入 91.3 亿元、在三级行业“肉鸡养殖”的 8 家公司中排第\zhnumber{3}位；毛利率 8.0\%、低于行业内中位数 2.8 个百分点；半年 ROE -0.8\%（未年化）、低于行业内中位数（1.21\%）；资产负债率 43.8\%、低于行业内中位数（52.1\%）。 就当期报表看，公司尚未把经营层面的特征转化为高于行业的资本回报，周期底部的相对优势更多体现在现金流与资产负债表上。

\pfl{历史股价} 本报告样本区间（2021 年 1 月 至 2026 年 10 月）内，股价由 33.08 元变为 20.31 元，累计 -38.6\%，区间最高 45.81 元（2022 年 8 月）、最低 16.41 元（2026 年 6 月）；当前价相当于区间最高价的 44.3\%，为区间最低价的 1.24 倍。 月度收益的年化波动率 43.8\%，高于全样本中位数 37.3\%，在三级行业“肉鸡养殖”的 8 家公司中波动率由高到低排第\zhnumber{2}位。 把股价阶段与公司事件对齐看，2022 年 4 月 至 2022 年 8 月股价上涨+63.9\%，同期（2022 年 8 月）实际控制人程立力全额认购 12.70 亿元定增。 从时间对应关系看，公司层面的资本运作与股权、风险事件解释了区间内多数急涨急跌，与行业指数的同步性相对有限。

\begin{center}
\begin{minipage}{\textwidth}\centering
\begin{tikzpicture}
\begin{axis}[
  width=12.6cm, height=3.9cm, scale only axis,
  xmin=2020.922, xmax=2026.972, ymin=15.26, ymax=49.02,
  xtick={2021,2022,2023,2024,2025,2026},
  yearticks,
  yticklabel style={font=\scriptsize},
  ylabel={元/股}, ylabel style={font=\scriptsize},
  grid=major, grid style={LightGray, line width=0.3pt},
  axis line style={InkGray, line width=0.5pt},
  every axis plot/.append style={line width=0.9pt},
]
\addplot[AgriGreen] table[x=x,y=y]{data/clean/profiles/px_300761.dat};
\addplot[SoftRed, dashed, line width=0.7pt] table[x=x,y=y]{data/clean/profiles/pxzz_300761.dat};
\addplot[only marks, mark=*, mark size=1.1pt, SoftRed] table[x=x,y=y]{data/clean/profiles/pxzz_300761.dat};
\end{axis}
\end{tikzpicture}

\captionof{figure}[立华股份（300761）股价与周期骨架]{立华股份（300761）月度收盘价（元/股）与 ZigZag 周期骨架（阈值 25\%，圆点为周期转折点）}
\end{minipage}
\end{center}

\pfl{过去周期分析} 以 25\% 的 ZigZag 阈值划分，样本区间内股价形成 4 个主升段与 5 个主跌段。 最大主升段出现在 2022 年 4 月 至 2022 年 8 月，4 个月+63.9\%；最大主跌段为 2022 年 8 月 至 2023 年 10 月，14 个月-63.0\%。 最近一次转折出现在 2026 年 6 月（下跌段自 2026 年 4 月起，2 个月-27.9\%），当前处于该段的延续之中。 公司股价与申万农林牧渔指数几乎同步见底，个股并未走出独立行情。 禽链的产能周期只有 6---12 个月、反转快于生猪（第 \ref{sec:other-livestock} 节），2026 年前四个月毛鸡养殖利润同比大幅回升，链条利润的修复先于生猪，公司 2026 年 6 月末资产负债率 43.8\%、半年 ROE -0.8\%（未年化），在本轮下行中属于随行业同步调整的一类。 综合区间形态看，该股单段涨跌幅度偏大、以趋势段为主（最大主升 +63.9\%、最大主跌 -63.0\%），年化波动率 43.8\% 与全样本中位数 37.3\% 相比波动明显大于样本中位数。

\pfl{盈利情况} 2026 年上半年营业收入 91.28 亿元（同比 +9.3\%），归母净利润 -0.72 亿元，由上年同期的盈利转为亏损。 以年报为趋势对照，2023---2025 年营业收入由 153.54 亿元变为 187.29 亿元（两年年均复合增速 +10.4\%），归母净利润由 -4.37 亿元变为 5.78 亿元。 按半年报口径，2026 年上半年的 ROE 为 -0.8\%（未年化）、销售净利率 -0.8\%、毛利率 8.0\%，ROE 在“肉鸡养殖”8 家公司中按数值由高到低排第\zhnumber{7}位。 同期全样本半年 ROE 中位数 0.75\%，公司与之相差 1.52 个百分点（弱于中位数公司）。 从公司披露的原因看，业绩变动主要来自：2026 年，业绩同比下滑主要受畜禽价格周期性波动影响。 综合看，公司仍处亏损区间、由盈转亏，半年 ROE 在三级行业内按由高到低排第\zhnumber{7}位，单位盈利能力的修复依赖行业景气与自身成本端的改善，报表弹性主要体现为价格与销量的方向性变化。

\begin{center}\small\setlength{\tabcolsep}{3.4pt}
\captionof{table}[立华股份（300761）关键财务指标]{立华股份（300761）关键财务指标（合并报表口径；两个半年报期均未年化；现金短债比 $=$ 货币资金 $\div$ 短期债务）}
\begin{adjustbox}{max width=\textwidth}
\begin{tabular}{@{}lrrrrr@{}}
\toprule
指标 & 2023 年 & 2024 年 & 2025 年 & 2025 年半年报 & 2026 年半年报 \\
\midrule
营业收入（亿元） & 153.54 & 177.25 & 187.29 & 83.53 & 91.28 \\
归母净利润（亿元） & -4.37 & 15.21 & 5.78 & 1.49 & -0.72 \\
毛利率（\%） & 4.8 & 15.9 & 11.0 & 10.0 & 8.0 \\
ROE（\%） & -5.4 & 18.0 & 6.3 & 1.6 & -0.8 \\
资产负债率（\%） & 43.0 & 37.8 & 42.0 & 41.1 & 43.8 \\
经营活动现金流净额（亿元） & 14.50 & 24.08 & 18.48 & 5.59 & 5.21 \\
货币资金（亿元） & 3.66 & 1.83 & 2.64 & 4.61 & 3.67 \\
有息负债合计（亿元） & 28.50 & 22.31 & 32.04 & 29.53 & 35.68 \\
现金短债比（倍） & 0.17 & 0.11 & 0.10 & 0.19 & 0.12 \\
\bottomrule
\end{tabular}
\end{adjustbox}
\end{center}

\begin{center}
\begin{minipage}{\textwidth}\centering
\begin{tikzpicture}
\begin{groupplot}[
  group style={group size=2 by 1, horizontal sep=0.60cm},
  width=6.85cm, height=4.60cm,
  tick label style={font=\tiny},
  title style={font=\scriptsize\heiti},
  yticklabel style={font=\tiny},
  ylabel style={font=\tiny},
  grid=major, grid style={LightGray, line width=0.3pt},
  axis line style={InkGray, line width=0.5pt},
  enlarge x limits={abs=0.8},
  legend style={font=\scriptsize, at={(0.5,-0.20)}, anchor=north,
                legend columns=2, draw=none, fill=none, column sep=6pt},
]
\nextgroupplot[
  ybar, bar width=4pt,
  ymin=-23.54, ymax=210.29,
  xtick={0,1,2,3,4,5.7},
  xticklabels={2021,2022,2023,2024,2025,2026H1},
  ylabel={亿元},
]
\addplot[fill=AgriGreen!75, draw=AgriGreen, bar shift=-2.6pt] table[x=i, y=rev]{data/clean/profiles/fin_300761.dat};
\addplot[fill=SoftRed!60, draw=SoftRed, bar shift=2.6pt] table[x=i, y=ni]{data/clean/profiles/fin_300761.dat};
\addplot[fill=AgriGreen!30, draw=AgriGreen, pattern=north east lines, bar shift=-2.6pt] table[x=x, y=rev]{data/clean/profiles/finh_300761.dat};
\addplot[fill=SoftRed!20, draw=SoftRed, pattern=north east lines, bar shift=2.6pt] table[x=x, y=ni]{data/clean/profiles/finh_300761.dat};
\legend{营业收入（年/半年）, 归母净利润（年/半年）}
\nextgroupplot[
  ymin=-10.1, ymax=48.8,
  xtick={0,1,2,3,4,5.7},
  xticklabels={2021,2022,2023,2024,2025,2026H1},
  ylabel={\%},
]
\addplot[AgriGold, mark=*, mark size=1.3pt] table[x=i, y=roe]{data/clean/profiles/fin_300761.dat};
\addplot[OilBlue, mark=square*, mark size=1.3pt] table[x=i, y=debt]{data/clean/profiles/fin_300761.dat};
\addplot[only marks, AgriGold, mark=*, mark size=2.0pt] table[x=x, y=roe]{data/clean/profiles/finh_300761.dat};
\addplot[only marks, OilBlue, mark=square*, mark size=2.0pt] table[x=x, y=debt]{data/clean/profiles/finh_300761.dat};
\legend{ROE, 资产负债率}
\end{groupplot}
\end{tikzpicture}
\begin{tikzpicture}
\begin{axis}[
  name=finbot,
  width=13.75cm, height=3.10cm, scale only axis,
  ymin=-3.57, ymax=39.96,
  xtick={0,1,2,3,4,5.7},
  xticklabels={2021,2022,2023,2024,2025,2026H1},
  tick label style={font=\tiny},
  yticklabel style={font=\tiny},
  ylabel={亿元}, ylabel style={font=\tiny},
  ybar, bar width=4pt,
  grid=major, grid style={LightGray, line width=0.3pt},
  axis line style={InkGray, line width=0.5pt},
  enlarge x limits={abs=0.8},
  legend style={font=\scriptsize, at={(0.5,-0.26)}, anchor=north,
                legend columns=2, draw=none, fill=none, column sep=10pt},
]
\addplot[fill=AgriGreen!55, draw=AgriGreen, bar shift=-3.0pt] table[x=i, y=cash]{data/clean/profiles/fin_300761.dat};
\addplot[fill=SoftRed!45, draw=SoftRed, bar shift=3.0pt] table[x=i, y=ibd]{data/clean/profiles/fin_300761.dat};
\addplot[fill=AgriGreen!25, draw=AgriGreen, pattern=north east lines, bar shift=-3.0pt] table[x=x, y=cash]{data/clean/profiles/finh_300761.dat};
\addplot[fill=SoftRed!20, draw=SoftRed, pattern=north east lines, bar shift=3.0pt] table[x=x, y=ibd]{data/clean/profiles/finh_300761.dat};
\legend{货币资金（左轴）, 有息负债（左轴）}
\end{axis}
\begin{axis}[
  at={(finbot.south east)}, anchor=south east,
  width=13.75cm, height=3.10cm, scale only axis,
  axis y line*=right, axis x line*=none, xtick=\empty,
  ymin=-6.26, ymax=31.30,
  tick label style={font=\tiny},
  yticklabel style={font=\tiny}, scaled y ticks=false,
  legend style={font=\scriptsize, at={(0.5,-0.44)}, anchor=north,
                legend columns=1, draw=none, fill=none},
]
\addplot[OilBlue, mark=*, mark size=2.2pt, line width=1.3pt] table[x=i, y=ocf]{data/clean/profiles/fin_300761.dat};
\addplot[only marks, OilBlue, mark=*, mark size=3.2pt] table[x=x, y=ocf]{data/clean/profiles/finh_300761.dat};
\legend{经营活动现金流净额（右轴，亿元，蓝点）}
\end{axis}
\end{tikzpicture}

\captionof{figure}[立华股份（300761）近年主要财务指标]{立华股份（300761）近年主要财务指标（左上：营业收入与归母净利润；右上：ROE 与资产负债率；下：货币资金、有息负债为左轴柱状，经营活动现金流净额为其右轴的蓝点折线。
2021---2025 年为年报口径，2026H1 为 2026 年半年报/半年末，与其前一年报之间留出间隔、以斜纹或空心点区分；半年报数据未年化）}
\end{minipage}
\end{center}

\pfl{财务分析} 2026 年 6 月末资产负债率 43.8\%，较 2025 年末上升 1.8 个百分点（样本中位数 52.2\%，所处三级行业内按由低到高排第\zhnumber{3}位）。 2026 年 6 月末货币资金 3.67 亿元、短期债务（短期借款与一年内到期非流动负债）29.94 亿元、长期债务 5.74 亿元，有息负债合计 35.68 亿元，现金短债比 0.12 倍——缺口明显。 净有息负债 32.00 亿元，相当于其 2026 年上半年营业收入的 35\%。 流动比率 0.76、速动比率 0.24，短期偿付压力较大。 2026 年上半年经营活动现金流净额 5.21 亿元，净利润为负而经营现金流为正，差额主要来自折旧摊销与营运资本变动，说明亏损尚未完全转化为现金流出。 2026 年 6 月末在建工程 2.79 亿元，较上年末增加 0.93 亿元，仍有在建投入。 综合看，账面货币资金对有息负债与短期债务的覆盖不足，滚动偿付对新增融资的依赖度较高；资产负债率低于样本中位数 8.4 个百分点；有息负债中短期债务占比更高，期限结构仍以滚动续作为主。

\pfl{重大战略转型} 公司的战略动作可归为以下几类：融资与资本运作（1 项）：2022 年 8 月，实际控制人程立力全额认购 12.70 亿元定增。 公告口径上，近三年（2023 年 9 月---2026 年 9 月）公司披露重大重组与并购类公告 0 条、对外投资与转型类公告 2 条、回购与股权激励类公告 47 条，合计公告 455 条。 第一大行业仍是“养殖行业”，收入占比 100.0\%，与 2021 年年报相比没有方向性变化。 综合判断，报告期内公司有 2 条与战略相关的公告落地（重大重组与并购 0 条、对外投资与转型 2 条），资本运作与主业调整之间存在直接关联。

\begin{center}\small\setlength{\tabcolsep}{4pt}
\captionof{table}[立华股份（300761）历史融资与资本运作]{立华股份（300761）历史融资与资本运作（据公司公告与公开报道整理；金额为披露值，含首次公开发行、再融资、债券、银行借款与重整投资等）}
\begin{adjustbox}{max width=\textwidth}
\begin{tabular}{@{}p{1.45cm}p{2.55cm}p{4.05cm}p{4.95cm}@{}}
\toprule
时间 & 方式 & 规模 & 用途与说明 \\
\midrule
2019-02-18 & IPO（深交所创业板） & 4,\allowbreak{}128 万股，\allowbreak{}发行价 29.35 元/\allowbreak{}股 & 安庆立华年出栏 1,\allowbreak{}750 万羽一体化养鸡建设项目、\allowbreak{}扬州立华年出栏 1,\allowbreak{}750 万羽一体化养鸡建设项目等；资金 2019 年 2 月 12 日到账；2022 年 4 月董事会通过募集资金项目结项及变更部分募集资金用途并永久性补充流动资金 \\
2021-09 & 可转换公司债券（方案） & 拟向不特定对象发行可转债募集资金不超过 20.45 亿元 & 黄羽鸡一体化养殖、\allowbreak{}屠宰等项目；公司第三届董事会第二次会议通过调整公司向不特定对象发行可转换公司债券方案 \\
2022-08 & 定向增发（向特定对象发行股票） & 5,\allowbreak{}592.25 万股，\allowbreak{}发行价 22.71 元/\allowbreak{}股 & 江苏、\allowbreak{}湖南等地养鸡及屠宰项目、\allowbreak{}总部基地项目及补充流动资金（项目总投资 27.14 亿元）；资金 2022 年 8 月 19 日到账 \\
\bottomrule
\end{tabular}
\end{adjustbox}
\end{center}

\pfl{历史融投资情况} 从现金流量表看，2025 年公司取得借款收到的现金 34.06 亿元、偿还债务支付的现金 25.05 亿元、吸收权益性投资 1.07 亿元，筹资活动现金流净额 3.94 亿元（2023 年为取得借款 33.53 亿元、偿还债务 28.39 亿元）。 2026 年上半年筹资活动现金流净额为 1.33 亿元（取得借款 20.13 亿元、偿还债务 16.88 亿元，未年化），仍处于净融入状态。 2025 年分配股利、利润或偿付利息支付的现金合计 5.56 亿元。 公开披露的融资记录共 3 笔，工具类型以 IPO、可转债、定向增发为主；金额与用途见下表。 其中规模最大的两笔为：2022 年 8 月，定向增发，5,592.25 万股，发行价 22.71 元/股；2019 年 2 月，IPO，4,128 万股，发行价 29.35 元/股。 股本方面，近三年披露股本变动 5 次（转增 1 次、限售股份上市 1 次）；总股本由 2021 年末的 4.04 亿股变为 8.28 亿股（+104.92\%）。 分红方面，近三年现金分红 5 次、合计每股派现 1.95 元，最近一次实施在 2026 年 5 月。 近三年“再融资”类公告 3 条，涉及定向增发。 综合看，公司融资结构上以债务滚动为主、股权融资为补充；2025 年筹资活动现金流净额 3.94 亿元，年度内仍为净融入，与前述经营与投资现金流的方向基本一致。

\pfl{未来潜在融资需求分析} 2026 年 6 月末货币资金 3.67 亿元，而短期债务 29.94 亿元（现金短债比 0.12 倍），2026 年上半年经营现金流净额 5.21 亿元。按“短期债务减去账面货币资金”的滚动偿付口径，静态缺口约 26.26 亿元；若再计入上半年亏损的年化额（1.44 亿元），维持一年正常经营所需的外部资金约 26.26---27.70 亿元。 公司 2026 年上半年归母净利润为负（-0.72 亿元），但 6 月末资产负债率 43.8\% 尚未进入第一档（亏损且负债率 $>$65\%）的区间，当前融资需求以补充流动资金、支撑低谷期经营性支出为主。 公开披露的后续资金安排线索包括：股权融资线索：2021 年 9 月，调整后的 20.45 亿元可转债方案。 资金需求还要叠加在建工程：期末在建工程 2.79 亿元、2026 年上半年购建固定资产等支出 8.32 亿元（未年化），这部分支出在周期底部通常会被主动放缓。 综合看，公司的外部资金需求以补充流动性与项目投入为主，滚动偿付口径下的静态缺口约 26.26 亿元，能否落地取决于授信续作与发行窗口的配合。




\subsubsection{神农集团（605296）}
\label{co:605296}

\pfl{公司基本情况} 公司前身可追溯至 1999 年，云南本土生猪一体化企业，饲料—养殖—屠宰全链条布局，2021 年在沪市主板上市后逆周期扩产、出栏量快速攀升；2026 年上半年猪价深跌致大额亏损，此前两项定增计划先后终止。 公司于 2021-05-28 在上交所首发上市，注册资本 5.26 亿元，总股本 5.25 亿股。 截至 2026 年 10 月 12 日收盘价 30.43 元，对应总市值 159.7 亿元，近一年-1.5\%，总市值在三级行业“生猪养殖”的 11 家公司中按由大到小排第\zhnumber{4}位，近一年涨跌幅高于全样本中位数（-10.4\%）8.9 个百分点。 按 2026 年半年报披露的产品构成，收入占比前三位为生猪养殖分部 90.7\%；饲料加工分部 45.4\%；屠宰及食品加工分部 23.6\%。 与 2021 年年报相比，第一大产品口径由“饲料”（62.9\%）变为“生猪养殖分部”（90.7\%）；两期披露的分类口径有调整，占比差异中同时包含分类因素。

\pfl{股权结构与控制权} 按 2026 年 6 月 30 日的股东名册，前三大股东为何祖训 持股 49.74\%；何乔关 持股 11.48\%；云南正道创业投资合伙企业(有限合伙) 持股 10.03\%，股东户数 18,697 户。 股东结构与控制权方面，控股股东何祖训个人持股 56.75\%；2026 年 9 月，公司控股股东、实际控制人何祖训涉及离婚诉讼；董事、副总经理顿灿于 2026 年 9 月 18 日因个人原因辞去公司全部职务。

\pfl{经营要点} 业务与产能层面，2026 年上半年，出栏商品猪 172.4 万头；2026 年 6 月，单月完全成本进一步优化至 11.9 元/公斤，上半年 PSY 稳步提升至 31 头；2025 年末，合作农户 657 户、年度支付代养费 4.63 亿元，代养出栏比例继续提升。

\pfl{所处行业的基本面} 生猪的供给由能繁母猪存栏决定、需求由猪肉消费与替代品价格决定，本轮下行的成因在供给端的产能去化不彻底，而非需求收缩（第 \ref{sec:framework} 章的供需框架）。 行业层面的最新情况是：2026 年上半年，行业处于全行业深度亏损，头均出栏亏损 300 至 500 元；2026 年二季度，末全国能繁母猪存栏 3780 万头；2026 年 3 月，至 18 日中央储备冻猪肉轮换出库、收储竞价交易挂牌 28400 吨。

\pfl{竞争优势} 把该公司放回同类公司比较，2026 年上半年营业收入 24.9 亿元、在三级行业“生猪养殖”的 11 家公司中排第\zhnumber{7}位；毛利率 0.9\%、高于行业内中位数 2.5 个百分点；半年 ROE -17.9\%（未年化）、低于行业内中位数（-14.32\%）；资产负债率 40.7\%、低于行业内中位数（58.9\%）。 与同类公司相比，公司的差异化在于：全产业链布局有助于降低单一环节价格波动对经营业绩的冲击；PSY 升至 31 头、料肉比优化至 2.45 至 2.48；2026 年代养出栏比例与代养户数量继续提升。 就当期报表看，公司尚未把经营层面的特征转化为高于行业的资本回报，周期底部的相对优势更多体现在现金流与资产负债表上。

\pfl{历史股价} 以 2021 年 5 月为起点、2026 年 10 月为终点，股价由 56.99 元变为 30.43 元，累计 -46.6\%，区间最高 56.99 元（2021 年 5 月）、最低 20.97 元（2023 年 9 月）；当前价相当于区间最高价的 53.4\%，为区间最低价的 1.45 倍。 月度收益的年化波动率 33.0\%，低于全样本中位数 37.3\%，在三级行业“生猪养殖”的 11 家公司中波动率由高到低排第\zhnumber{9}位。 把股价阶段与公司事件对齐看，2021 年 5 月 至 2023 年 9 月股价下跌-63.2\%，同期（2024 年 4 月）公告终止该次定增，同时启动简易程序定增方案。 从时间对应关系看，公司层面的资本运作与股权、风险事件解释了区间内多数急涨急跌，与行业指数的同步性相对有限。

\begin{center}
\begin{minipage}{\textwidth}\centering
\begin{tikzpicture}
\begin{axis}[
  width=12.6cm, height=3.9cm, scale only axis,
  xmin=2021.255, xmax=2026.972, ymin=19.50, ymax=60.98,
  xtick={2021,2022,2023,2024,2025,2026},
  yearticks,
  yticklabel style={font=\scriptsize},
  ylabel={元/股}, ylabel style={font=\scriptsize},
  grid=major, grid style={LightGray, line width=0.3pt},
  axis line style={InkGray, line width=0.5pt},
  every axis plot/.append style={line width=0.9pt},
]
\addplot[AgriGreen] table[x=x,y=y]{data/clean/profiles/px_605296.dat};
\addplot[SoftRed, dashed, line width=0.7pt] table[x=x,y=y]{data/clean/profiles/pxzz_605296.dat};
\addplot[only marks, mark=*, mark size=1.1pt, SoftRed] table[x=x,y=y]{data/clean/profiles/pxzz_605296.dat};
\end{axis}
\end{tikzpicture}

\captionof{figure}[神农集团（605296）股价与周期骨架]{神农集团（605296）月度收盘价（元/股）与 ZigZag 周期骨架（阈值 25\%，圆点为周期转折点）}
\end{minipage}
\end{center}

\pfl{过去周期分析} 以 25\% 的 ZigZag 阈值划分，样本区间内股价形成 2 个主升段与 2 个主跌段。 最大主升段出现在 2023 年 9 月 至 2024 年 5 月，8 个月+93.5\%；最大主跌段为 2021 年 5 月 至 2023 年 9 月，28 个月-63.2\%。 最近一次转折出现在 2026 年 8 月（上涨段自 2025 年 1 月起，19 个月+32.7\%），当前处于该段之后的震荡之中。 公司股价较申万农林牧渔指数提前约 32 个月见底，是板块中较早定价周期反转的公司之一。 生猪养殖的定价主线是产能去化与成本曲线的赛跑（第 \ref{sec:hog} 章）：2026 年三季度行业仍处在“产能去化未完成、价格磨底”的阶段，公司 2026 年 6 月末资产负债率 40.7\%、半年 ROE -17.9\%（未年化），在本轮下行中属于随行业同步调整的一类。 综合区间形态看，该股单段涨跌幅度偏大、以趋势段为主（最大主升 +93.5\%、最大主跌 -63.2\%），年化波动率 33.0\% 与全样本中位数 37.3\% 相比波动小于样本中位数。

\pfl{盈利情况} 2026 年上半年营业收入 24.89 亿元（同比 -11.0\%），归母净利润 -8.09 亿元，由上年同期的盈利转为亏损。 以年报为趋势对照，2023---2025 年营业收入由 38.91 亿元变为 53.52 亿元（两年年均复合增速 +17.3\%），归母净利润由 -4.01 亿元变为 3.39 亿元。 按半年报口径，2026 年上半年的 ROE 为 -17.9\%（未年化）、销售净利率 -32.5\%、毛利率 0.9\%，ROE 在“生猪养殖”11 家公司中按数值由高到低排第\zhnumber{8}位。 按同一口径，三级行业“生猪养殖”11 家公司的半年 ROE 中位数为 -14.32\%，该公司低于其中位数 3.57 个百分点。 针对上述压力，公司披露的应对举措包括：公司为云南区域一体化企业，同时推进饲料及屠宰板块产能利用率提升以摊薄单位成本。 综合看，公司仍处亏损区间、由盈转亏，半年 ROE 在三级行业内按由高到低排第\zhnumber{8}位，单位盈利能力的修复依赖行业景气与自身成本端的改善，报表弹性主要体现为价格与销量的方向性变化。

\begin{center}\small\setlength{\tabcolsep}{3.4pt}
\captionof{table}[神农集团（605296）关键财务指标]{神农集团（605296）关键财务指标（合并报表口径；两个半年报期均未年化；现金短债比 $=$ 货币资金 $\div$ 短期债务）}
\begin{adjustbox}{max width=\textwidth}
\begin{tabular}{@{}lrrrrr@{}}
\toprule
指标 & 2023 年 & 2024 年 & 2025 年 & 2025 年半年报 & 2026 年半年报 \\
\midrule
营业收入（亿元） & 38.91 & 55.84 & 53.52 & 27.98 & 24.89 \\
归母净利润（亿元） & -4.01 & 6.87 & 3.39 & 3.88 & -8.09 \\
毛利率（\%） & 4.4 & 20.8 & 16.3 & 21.9 & 0.9 \\
ROE（\%） & -9.1 & 15.1 & 6.9 & 7.8 & -17.9 \\
资产负债率（\%） & 26.4 & 27.4 & 29.5 & 26.0 & 40.7 \\
经营活动现金流净额（亿元） & 0.19 & 15.34 & 8.10 & 5.57 & -3.75 \\
货币资金（亿元） & 5.44 & 10.42 & 7.15 & 9.39 & 4.50 \\
有息负债合计（亿元） & 7.52 & 7.68 & 7.69 & 7.63 & 13.62 \\
现金短债比（倍） & 1.45 & 2.54 & 3.37 & 3.12 & 0.68 \\
\bottomrule
\end{tabular}
\end{adjustbox}
\end{center}

\begin{center}
\begin{minipage}{\textwidth}\centering
\begin{tikzpicture}
\begin{groupplot}[
  group style={group size=2 by 1, horizontal sep=0.60cm},
  width=6.85cm, height=4.60cm,
  tick label style={font=\tiny},
  title style={font=\scriptsize\heiti},
  yticklabel style={font=\tiny},
  ylabel style={font=\tiny},
  grid=major, grid style={LightGray, line width=0.3pt},
  axis line style={InkGray, line width=0.5pt},
  enlarge x limits={abs=0.8},
  legend style={font=\scriptsize, at={(0.5,-0.20)}, anchor=north,
                legend columns=2, draw=none, fill=none, column sep=6pt},
]
\nextgroupplot[
  ybar, bar width=4pt,
  ymin=-14.48, ymax=63.52,
  xtick={0,1,2,3,4,5.7},
  xticklabels={2021,2022,2023,2024,2025,2026H1},
  ylabel={亿元},
]
\addplot[fill=AgriGreen!75, draw=AgriGreen, bar shift=-2.6pt] table[x=i, y=rev]{data/clean/profiles/fin_605296.dat};
\addplot[fill=SoftRed!60, draw=SoftRed, bar shift=2.6pt] table[x=i, y=ni]{data/clean/profiles/fin_605296.dat};
\addplot[fill=AgriGreen!30, draw=AgriGreen, pattern=north east lines, bar shift=-2.6pt] table[x=x, y=rev]{data/clean/profiles/finh_605296.dat};
\addplot[fill=SoftRed!20, draw=SoftRed, pattern=north east lines, bar shift=2.6pt] table[x=x, y=ni]{data/clean/profiles/finh_605296.dat};
\legend{营业收入（年/半年）, 归母净利润（年/半年）}
\nextgroupplot[
  ymin=-22.6, ymax=46.6,
  xtick={0,1,2,3,4,5.7},
  xticklabels={2021,2022,2023,2024,2025,2026H1},
  ylabel={\%},
]
\addplot[AgriGold, mark=*, mark size=1.3pt] table[x=i, y=roe]{data/clean/profiles/fin_605296.dat};
\addplot[OilBlue, mark=square*, mark size=1.3pt] table[x=i, y=debt]{data/clean/profiles/fin_605296.dat};
\addplot[only marks, AgriGold, mark=*, mark size=2.0pt] table[x=x, y=roe]{data/clean/profiles/finh_605296.dat};
\addplot[only marks, OilBlue, mark=square*, mark size=2.0pt] table[x=x, y=debt]{data/clean/profiles/finh_605296.dat};
\legend{ROE, 资产负债率}
\end{groupplot}
\end{tikzpicture}
\begin{tikzpicture}
\begin{axis}[
  name=finbot,
  width=13.75cm, height=3.10cm, scale only axis,
  ymin=-1.36, ymax=15.26,
  xtick={0,1,2,3,4,5.7},
  xticklabels={2021,2022,2023,2024,2025,2026H1},
  tick label style={font=\tiny},
  yticklabel style={font=\tiny},
  ylabel={亿元}, ylabel style={font=\tiny},
  ybar, bar width=4pt,
  grid=major, grid style={LightGray, line width=0.3pt},
  axis line style={InkGray, line width=0.5pt},
  enlarge x limits={abs=0.8},
  legend style={font=\scriptsize, at={(0.5,-0.26)}, anchor=north,
                legend columns=2, draw=none, fill=none, column sep=10pt},
]
\addplot[fill=AgriGreen!55, draw=AgriGreen, bar shift=-3.0pt] table[x=i, y=cash]{data/clean/profiles/fin_605296.dat};
\addplot[fill=SoftRed!45, draw=SoftRed, bar shift=3.0pt] table[x=i, y=ibd]{data/clean/profiles/fin_605296.dat};
\addplot[fill=AgriGreen!25, draw=AgriGreen, pattern=north east lines, bar shift=-3.0pt] table[x=x, y=cash]{data/clean/profiles/finh_605296.dat};
\addplot[fill=SoftRed!20, draw=SoftRed, pattern=north east lines, bar shift=3.0pt] table[x=x, y=ibd]{data/clean/profiles/finh_605296.dat};
\legend{货币资金（左轴）, 有息负债（左轴）}
\end{axis}
\begin{axis}[
  at={(finbot.south east)}, anchor=south east,
  width=13.75cm, height=3.10cm, scale only axis,
  axis y line*=right, axis x line*=none, xtick=\empty,
  ymin=-8.72, ymax=21.06,
  tick label style={font=\tiny},
  yticklabel style={font=\tiny}, scaled y ticks=false,
  legend style={font=\scriptsize, at={(0.5,-0.44)}, anchor=north,
                legend columns=1, draw=none, fill=none},
]
\addplot[OilBlue, mark=*, mark size=2.2pt, line width=1.3pt] table[x=i, y=ocf]{data/clean/profiles/fin_605296.dat};
\addplot[only marks, OilBlue, mark=*, mark size=3.2pt] table[x=x, y=ocf]{data/clean/profiles/finh_605296.dat};
\legend{经营活动现金流净额（右轴，亿元，蓝点）}
\end{axis}
\end{tikzpicture}

\captionof{figure}[神农集团（605296）近年主要财务指标]{神农集团（605296）近年主要财务指标（左上：营业收入与归母净利润；右上：ROE 与资产负债率；下：货币资金、有息负债为左轴柱状，经营活动现金流净额为其右轴的蓝点折线。
2021---2025 年为年报口径，2026H1 为 2026 年半年报/半年末，与其前一年报之间留出间隔、以斜纹或空心点区分；半年报数据未年化）}
\end{minipage}
\end{center}

\pfl{财务分析} 2026 年 6 月末资产负债率 40.7\%，较 2025 年末上升 11.2 个百分点（样本中位数 52.2\%，所处三级行业内按由低到高排第\zhnumber{1}位）。 2026 年 6 月末货币资金 4.50 亿元、短期债务（短期借款与一年内到期非流动负债）6.59 亿元、长期债务 7.03 亿元，有息负债合计 13.62 亿元，现金短债比 0.68 倍——覆盖偏紧。 净有息负债 9.13 亿元，相当于其 2026 年上半年营业收入的 37\%。 流动比率 1.14、速动比率 0.37，短期指标尚可。 2026 年上半年经营活动现金流净额 -3.75 亿元，经营现金流与利润同时为负，现金消耗较快。 2026 年 6 月末在建工程 3.96 亿元，较上年末增加 2.53 亿元，仍有在建投入。 综合看，现金短债比 0.68 倍，短期资金安排偏紧；资产负债率低于样本中位数 11.4 个百分点；有息负债中长期债务占比更高，期限结构对短债滚动的压力较小。

\pfl{重大战略转型} 近三年公司的战略动作集中在：其他动作（3 项）：2025 年度，公司实现归属于上市公司股东的净利润 3.39 亿元；2023---2025 年，2024 年平均完全成本 13.9 元/公斤。 公告口径上，近三年（2023 年 9 月---2026 年 9 月）公司披露重大重组与并购类公告 1 条、对外投资与转型类公告 5 条、回购与股权激励类公告 85 条，合计公告 563 条。 主营结构的口径变化已经落到报表上：2021 年年报的第一大产品为“饲料”（62.9\%），2026 年半年报为“生猪养殖分部”（90.7\%）；公司在这两期的分部划分方式有过调整。 综合判断，报告期内公司有 6 条与战略相关的公告落地（重大重组与并购 1 条、对外投资与转型 5 条），资本运作与主业调整之间存在直接关联。

\begin{center}\small\setlength{\tabcolsep}{4pt}
\captionof{table}[神农集团（605296）历史融资与资本运作]{神农集团（605296）历史融资与资本运作（据公司公告与公开报道整理；金额为披露值，含首次公开发行、再融资、债券、银行借款与重整投资等）}
\begin{adjustbox}{max width=\textwidth}
\begin{tabular}{@{}p{1.45cm}p{2.55cm}p{4.05cm}p{4.95cm}@{}}
\toprule
时间 & 方式 & 规模 & 用途与说明 \\
\midrule
2021-05-28 & IPO（上交所主板） & 发行 4,\allowbreak{}003 万股，\allowbreak{}发行价 56.08 元/\allowbreak{}股，\allowbreak{}募集资金约 22.00 亿元 & 生猪养殖、\allowbreak{}饲料加工等项目 \\
2023-06 & 定向增发（预案，后终止） & 拟募资不超过 11.50 亿元 & 两个饲料加工建设类项目及一个生猪养殖建设类项目；2024 年 4 月公告终止该次定增 \\
2024-2025 & 以简易程序向特定对象发行股票（后终止） & 拟募资 2.90 亿元（1.37 亿元投红河弥勒五山乡年出栏 24 万头优质仔猪扩繁基地项目、\allowbreak{}2,\allowbreak{}500.00 万元投石林林口铺养殖场生物安全技改、\allowbreak{}4,\allowbreak{}333.06 万元投产业链数字化智能化项目、\allowbreak{}8,\allowbreak{}437.41 万元补充流动资金） & 产能扩张、\allowbreak{}技改与补充流动资金；2025 年 3 月披露预案修订稿 \\
\bottomrule
\end{tabular}
\end{adjustbox}
\end{center}

\pfl{历史融投资情况} 从现金流量表看，2025 年公司取得借款收到的现金 3.80 亿元、偿还债务支付的现金 3.94 亿元、吸收权益性投资 1.18 亿元，筹资活动现金流净额 -3.64 亿元（2023 年为取得借款 6.23 亿元、偿还债务 0.32 亿元）。 2026 年上半年筹资活动现金流净额为 5.08 亿元（取得借款 6.59 亿元、偿还债务 0.76 亿元，未年化），仍处于净融入状态。 2025 年分配股利、利润或偿付利息支付的现金合计 2.14 亿元。 公开披露的融资记录共 3 笔，工具类型以 IPO、定向增发为主；金额与用途见下表。 其中规模最大的两笔为：2021 年 5 月，IPO，发行 4,003 万股，发行价 56.08 元/股，募集资金约 22.00 亿元。 股本方面，近三年披露股本变动 10 次（激励股份解禁 3 次、股份回购 2 次、增发新股上市 1 次、股权激励 1 次）；总股本由 2021 年末的 4.00 亿股变为 5.25 亿股（+31.13\%）。 分红方面，近三年现金分红 2 次、合计每股派现 0.64 元，最近一次实施在 2025 年 6 月。 近三年“再融资”类公告 27 条，涉及定向增发、募集资金使用。 综合看，公司融资结构上以债务滚动为主、股权融资为补充；2025 年筹资活动现金流净额 -3.64 亿元，融资节奏仍以维持存量债务周转为主，与前述经营与投资现金流的方向基本一致。

\pfl{未来潜在融资需求分析} 2026 年 6 月末货币资金 4.50 亿元，而短期债务 6.59 亿元（现金短债比 0.68 倍），2026 年上半年经营现金流净额 -3.75 亿元。按“短期债务减去账面货币资金”的滚动偿付口径，静态缺口约 2.10 亿元；若再计入上半年亏损的年化额（16.18 亿元），维持一年正常经营所需的外部资金约 2.10---18.27 亿元。 公司 2026 年上半年归母净利润为负（-8.09 亿元），但 6 月末资产负债率 40.7\% 尚未进入第一档（亏损且负债率 $>$65\%）的区间，当前融资需求以补充流动资金、支撑低谷期经营性支出为主。 公开披露的后续资金安排线索包括：股权融资线索：2023 年度，定增与 2024 年度简易程序定增均已终止。 资金需求还要叠加在建工程：期末在建工程 3.96 亿元、2026 年上半年购建固定资产等支出 2.81 亿元（未年化），这部分支出在周期底部通常会被主动放缓。 综合看，公司的外部资金需求以补充流动性与项目投入为主，滚动偿付口径下的静态缺口约 2.10 亿元，能否落地取决于授信续作与发行窗口的配合。




\subsubsection{京基智农（000048）}
\label{co:000048}

\pfl{公司基本情况} 1979 年设立至今，京基智农是前身为 1994 年上市的“中国农牧第一股”康达尔，2018—2019 年被京基集团要约收购入主后更名，主业转型为广东、广西区域生猪养殖，房地产转为存量去化业务；2026 年上半年猪价低迷导致大额亏损。 公司于 1994-11-01 在深交所主板首发上市，注册资本 5.30 亿元，总股本 5.31 亿股。 截至 2026 年 10 月 12 日收盘价 17.16 元，对应总市值 91.0 亿元，近一年+2.0\%，总市值在三级行业“生猪养殖”的 11 家公司中按由大到小排第\zhnumber{5}位，近一年涨跌幅高于全样本中位数（-10.4\%）12.3 个百分点。 按 2026 年半年报披露的产品构成，收入占比前三位为猪业产品 71.7\%；饲料产品 14.9\%；商品房 7.1\%。 与 2021 年年报相比，第一大产品口径由“商品房”（68.8\%）变为“猪业产品”（71.7\%）；两期披露的分类口径有调整，占比差异中同时包含分类因素。

\pfl{股权结构与控制权} 按 2026 年 8 月 26 日的股东名册，前三大股东为深圳市京基时代实业有限公司 持股 29.46\%；京基集团有限公司 持股 14.80\%；深圳市中恒通私募证券基金管理有限公司－中恒通启航 1 号 持股 5.00\%。 从控制权角度看，持股 44.26\%，合计控制 39.83\% 的股份；2026 年 10 月初，公司副总裁王鸿鹤因个人原因辞职。

\pfl{经营要点} 从经营要点看，截至 2025 年末生猪养殖项目存栏母猪 11.94 万头，其中能繁母猪 10.29 万头、后备母猪 1.65 万头。

\pfl{所处行业的基本面} 生猪的供给由能繁母猪存栏决定、需求由猪肉消费与替代品价格决定，本轮下行的成因在供给端的产能去化不彻底，而非需求收缩（第 \ref{sec:framework} 章的供需框架）。 行业层面的最新情况是：2026 年上半年，全国生猪出栏 37246 万头、同比增长 1.7\%，二季度末生猪存栏 42491 万头、同比增长 0.1\%；2026 年 5 月，农业农村部将全国能繁母猪正常保有量目标下调至 3750 万头，并新增中期调仔猪的调控思路。

\pfl{竞争优势} 把该公司放回同类公司比较，2026 年上半年营业收入 21.3 亿元、在三级行业“生猪养殖”的 11 家公司中排第\zhnumber{8}位；毛利率 -8.1\%、低于行业内中位数 6.5 个百分点；半年 ROE -12.9\%（未年化）、高于行业内中位数（-14.32\%）；资产负债率 63.6\%、高于行业内中位数（58.9\%）。 相对优势方面，2019 年进军生猪业务后仅 3 年即实现超百万头出栏，旗下四个生猪项目获评国家级生猪产能调控基地及所在地菜篮子生产稳供基地；2026 年上半年，并新增 6,000.00 万元对外投资入股江苏汇博机器人技术股份有限公司、持股 3.61\%，在主业承压期布局风险对冲。 优势在当期体现为高于行业中位数的资本回报，但优势能否持续仍取决于周期与自身产能投放节奏。

\pfl{历史股价} 以 2021 年 1 月为起点、2026 年 10 月为终点，股价由 22.79 元变为 17.16 元，累计 -24.7\%，区间最高 29.29 元（2021 年 4 月）、最低 12.54 元（2026 年 5 月）；当前价相当于区间最高价的 58.6\%，为区间最低价的 1.37 倍。 月度收益的年化波动率 51.1\%，高于全样本中位数 37.3\%，在三级行业“生猪养殖”的 11 家公司中波动率由高到低排第\zhnumber{1}位。 把股价阶段与公司事件对齐看，2026 年 5 月 至 2026 年 8 月股价上涨+82.9\%，同期（2026 年 6 月）公司拟将全资子公司深圳市京基智农酒店管理有限公司 100\% 股权转让给广东承希科汇投资控股有限公司的事项；2025 年 11 月 至 2026 年 3 月股价上涨+75.4\%，同期（2025 年 12 月）公司拟投资 6.50 亿元对江苏汇博机器人技术股份有限公司增资扩股并取得控制权。 把这些阶段与公司事件对齐后可以看到，几轮大涨大跌多由公司自身的资本运作与风险事件触发，行业景气只解释了其中一部分。

\begin{center}
\begin{minipage}{\textwidth}\centering
\begin{tikzpicture}
\begin{axis}[
  width=12.6cm, height=3.9cm, scale only axis,
  xmin=2020.922, xmax=2026.972, ymin=11.66, ymax=31.34,
  xtick={2021,2022,2023,2024,2025,2026},
  yearticks,
  yticklabel style={font=\scriptsize},
  ylabel={元/股}, ylabel style={font=\scriptsize},
  grid=major, grid style={LightGray, line width=0.3pt},
  axis line style={InkGray, line width=0.5pt},
  every axis plot/.append style={line width=0.9pt},
]
\addplot[AgriGreen] table[x=x,y=y]{data/clean/profiles/px_000048.dat};
\addplot[SoftRed, dashed, line width=0.7pt] table[x=x,y=y]{data/clean/profiles/pxzz_000048.dat};
\addplot[only marks, mark=*, mark size=1.1pt, SoftRed] table[x=x,y=y]{data/clean/profiles/pxzz_000048.dat};
\end{axis}
\end{tikzpicture}

\captionof{figure}[京基智农（000048）股价与周期骨架]{京基智农（000048）月度收盘价（元/股）与 ZigZag 周期骨架（阈值 25\%，圆点为周期转折点）}
\end{minipage}
\end{center}

\pfl{过去周期分析} 以 25\% 的 ZigZag 阈值划分，样本区间内股价形成 6 个主升段与 6 个主跌段。 最大主升段出现在 2026 年 5 月 至 2026 年 8 月，3 个月+82.9\%；最大主跌段为 2026 年 3 月 至 2026 年 5 月，2 个月-47.0\%。 最近一次转折出现在 2026 年 10 月（下跌段自 2026 年 8 月起，2 个月-25.2\%），当前处于该段的延续之中。 公司股价与申万农林牧渔指数几乎同步见底，个股并未走出独立行情。 生猪养殖的定价主线是产能去化与成本曲线的赛跑（第 \ref{sec:hog} 章）：2026 年三季度行业仍处在“产能去化未完成、价格磨底”的阶段，公司 2026 年 6 月末资产负债率 63.6\%、半年 ROE -12.9\%（未年化），在本轮下行中属于随行业同步调整的一类。 综合区间形态看，该股单段涨跌幅度偏大、以趋势段为主（最大主升 +82.9\%、最大主跌 -47.0\%），年化波动率 51.1\% 与全样本中位数 37.3\% 相比波动明显大于样本中位数。

\pfl{盈利情况} 2026 年上半年营业收入 21.30 亿元（同比 -10.7\%），归母净利润 -4.89 亿元，由上年同期的盈利转为亏损。 以年报为趋势对照，2023---2025 年营业收入由 124.17 亿元变为 48.72 亿元（两年年均复合增速 -37.4\%），归母净利润由 17.46 亿元变为 1.53 亿元。 按半年报口径，2026 年上半年的 ROE 为 -12.9\%（未年化）、销售净利率 -23.1\%、毛利率 -8.1\%，ROE 在“生猪养殖”11 家公司中按数值由高到低排第\zhnumber{5}位。 按同一口径，三级行业“生猪养殖”11 家公司的半年 ROE 中位数为 -14.32\%，该公司高于其中位数 1.44 个百分点。 从公司披露的原因看，业绩变动主要来自：受市场竞争影响毛利有所下滑。 针对上述压力，公司披露的应对举措包括：2025 年，2026 年完全成本目标 12 元/公斤左右。 综合看，公司仍处亏损区间、由盈转亏，半年 ROE 在三级行业内按由高到低排第\zhnumber{5}位，单位盈利能力的修复依赖行业景气与自身成本端的改善，报表弹性主要体现为价格与销量的方向性变化。

\begin{center}\small\setlength{\tabcolsep}{3.4pt}
\captionof{table}[京基智农（000048）关键财务指标]{京基智农（000048）关键财务指标（合并报表口径；两个半年报期均未年化；现金短债比 $=$ 货币资金 $\div$ 短期债务）}
\begin{adjustbox}{max width=\textwidth}
\begin{tabular}{@{}lrrrrr@{}}
\toprule
指标 & 2023 年 & 2024 年 & 2025 年 & 2025 年半年报 & 2026 年半年报 \\
\midrule
营业收入（亿元） & 124.17 & 59.60 & 48.72 & 23.87 & 21.30 \\
归母净利润（亿元） & 17.46 & 7.14 & 1.53 & 2.26 & -4.89 \\
毛利率（\%） & 38.7 & 28.0 & 13.8 & 19.3 & -8.1 \\
ROE（\%） & 55.5 & 17.4 & 3.6 & 5.2 & -12.9 \\
资产负债率（\%） & 67.8 & 60.3 & 59.0 & 58.3 & 63.6 \\
经营活动现金流净额（亿元） & -4.15 & 9.97 & 5.97 & 3.11 & 0.41 \\
货币资金（亿元） & 11.07 & 8.74 & 4.97 & 6.93 & 6.08 \\
有息负债合计（亿元） & 26.40 & 21.19 & 17.01 & 18.77 & 19.75 \\
现金短债比（倍） & 1.04 & 0.77 & 0.44 & 0.57 & 0.46 \\
\bottomrule
\end{tabular}
\end{adjustbox}
\end{center}

\begin{center}
\begin{minipage}{\textwidth}\centering
\begin{tikzpicture}
\begin{groupplot}[
  group style={group size=2 by 1, horizontal sep=0.60cm},
  width=6.85cm, height=4.60cm,
  tick label style={font=\tiny},
  title style={font=\scriptsize\heiti},
  yticklabel style={font=\tiny},
  ylabel style={font=\tiny},
  grid=major, grid style={LightGray, line width=0.3pt},
  axis line style={InkGray, line width=0.5pt},
  enlarge x limits={abs=0.8},
  legend style={font=\scriptsize, at={(0.5,-0.20)}, anchor=north,
                legend columns=2, draw=none, fill=none, column sep=6pt},
]
\nextgroupplot[
  ybar, bar width=4pt,
  ymin=-17.80, ymax=139.66,
  xtick={0,1,2,3,4,5.7},
  xticklabels={2021,2022,2023,2024,2025,2026H1},
  ylabel={亿元},
]
\addplot[fill=AgriGreen!75, draw=AgriGreen, bar shift=-2.6pt] table[x=i, y=rev]{data/clean/profiles/fin_000048.dat};
\addplot[fill=SoftRed!60, draw=SoftRed, bar shift=2.6pt] table[x=i, y=ni]{data/clean/profiles/fin_000048.dat};
\addplot[fill=AgriGreen!30, draw=AgriGreen, pattern=north east lines, bar shift=-2.6pt] table[x=x, y=rev]{data/clean/profiles/finh_000048.dat};
\addplot[fill=SoftRed!20, draw=SoftRed, pattern=north east lines, bar shift=2.6pt] table[x=x, y=ni]{data/clean/profiles/finh_000048.dat};
\legend{营业收入（年/半年）, 归母净利润（年/半年）}
\nextgroupplot[
  ymin=-20.7, ymax=94.6,
  xtick={0,1,2,3,4,5.7},
  xticklabels={2021,2022,2023,2024,2025,2026H1},
  ylabel={\%},
]
\addplot[AgriGold, mark=*, mark size=1.3pt] table[x=i, y=roe]{data/clean/profiles/fin_000048.dat};
\addplot[OilBlue, mark=square*, mark size=1.3pt] table[x=i, y=debt]{data/clean/profiles/fin_000048.dat};
\addplot[only marks, AgriGold, mark=*, mark size=2.0pt] table[x=x, y=roe]{data/clean/profiles/finh_000048.dat};
\addplot[only marks, OilBlue, mark=square*, mark size=2.0pt] table[x=x, y=debt]{data/clean/profiles/finh_000048.dat};
\legend{ROE, 资产负债率}
\end{groupplot}
\end{tikzpicture}
\begin{tikzpicture}
\begin{axis}[
  name=finbot,
  width=13.75cm, height=3.10cm, scale only axis,
  ymin=-3.53, ymax=39.50,
  xtick={0,1,2,3,4,5.7},
  xticklabels={2021,2022,2023,2024,2025,2026H1},
  tick label style={font=\tiny},
  yticklabel style={font=\tiny},
  ylabel={亿元}, ylabel style={font=\tiny},
  ybar, bar width=4pt,
  grid=major, grid style={LightGray, line width=0.3pt},
  axis line style={InkGray, line width=0.5pt},
  enlarge x limits={abs=0.8},
  legend style={font=\scriptsize, at={(0.5,-0.26)}, anchor=north,
                legend columns=2, draw=none, fill=none, column sep=10pt},
]
\addplot[fill=AgriGreen!55, draw=AgriGreen, bar shift=-3.0pt] table[x=i, y=cash]{data/clean/profiles/fin_000048.dat};
\addplot[fill=SoftRed!45, draw=SoftRed, bar shift=3.0pt] table[x=i, y=ibd]{data/clean/profiles/fin_000048.dat};
\addplot[fill=AgriGreen!25, draw=AgriGreen, pattern=north east lines, bar shift=-3.0pt] table[x=x, y=cash]{data/clean/profiles/finh_000048.dat};
\addplot[fill=SoftRed!20, draw=SoftRed, pattern=north east lines, bar shift=3.0pt] table[x=x, y=ibd]{data/clean/profiles/finh_000048.dat};
\legend{货币资金（左轴）, 有息负债（左轴）}
\end{axis}
\begin{axis}[
  at={(finbot.south east)}, anchor=south east,
  width=13.75cm, height=3.10cm, scale only axis,
  axis y line*=right, axis x line*=none, xtick=\empty,
  ymin=-12.09, ymax=35.57,
  tick label style={font=\tiny},
  yticklabel style={font=\tiny}, scaled y ticks=false,
  legend style={font=\scriptsize, at={(0.5,-0.44)}, anchor=north,
                legend columns=1, draw=none, fill=none},
]
\addplot[OilBlue, mark=*, mark size=2.2pt, line width=1.3pt] table[x=i, y=ocf]{data/clean/profiles/fin_000048.dat};
\addplot[only marks, OilBlue, mark=*, mark size=3.2pt] table[x=x, y=ocf]{data/clean/profiles/finh_000048.dat};
\legend{经营活动现金流净额（右轴，亿元，蓝点）}
\end{axis}
\end{tikzpicture}

\captionof{figure}[京基智农（000048）近年主要财务指标]{京基智农（000048）近年主要财务指标（左上：营业收入与归母净利润；右上：ROE 与资产负债率；下：货币资金、有息负债为左轴柱状，经营活动现金流净额为其右轴的蓝点折线。
2021---2025 年为年报口径，2026H1 为 2026 年半年报/半年末，与其前一年报之间留出间隔、以斜纹或空心点区分；半年报数据未年化）}
\end{minipage}
\end{center}

\pfl{财务分析} 2026 年 6 月末资产负债率 63.6\%，较 2025 年末上升 4.7 个百分点（样本中位数 52.2\%，所处三级行业内按由低到高排第\zhnumber{7}位）。 2026 年 6 月末货币资金 6.08 亿元、短期债务（短期借款与一年内到期非流动负债）13.29 亿元、长期债务 6.46 亿元，有息负债合计 19.75 亿元，现金短债比 0.46 倍——覆盖偏紧。 净有息负债 13.68 亿元，相当于其 2026 年上半年营业收入的 64\%。 流动比率 0.40、速动比率 0.17，短期偿债指标偏紧。 2026 年上半年经营活动现金流净额 0.41 亿元，净利润为负而经营现金流为正，差额主要来自折旧摊销与营运资本变动，说明亏损尚未完全转化为现金流出。 2026 年 6 月末在建工程 1.32 亿元，较上年末增加 0.05 亿元，仍有在建投入。 综合看，账面货币资金对有息负债与短期债务的覆盖不足，滚动偿付对新增融资的依赖度较高；资产负债率高于样本中位数 11.4 个百分点；有息负债中短期债务占比更高，期限结构仍以滚动续作为主。 公司披露的股东与质押风险为：2026 年 8 月，控股股东京基集团、京基时代分别将其持有的公司股份 92.19 万股（占总股本 0.17\%）、9.81 万股（占总股本 0.02\%）质押给自然人郭鉴慷。

\pfl{重大战略转型} 公司的战略动作可归为以下几类：其他动作（2 项）：2026 年 6 月，公司终止原 6.50 亿元控股型增资协议；并购与重组（1 项）：2018---2019 年，京基集团通过要约收购取得控制权并承诺帮助上市公司提升管理效率、优化资源配置。 公告口径上，近三年（2023 年 9 月---2026 年 9 月）公司披露重大重组与并购类公告 1 条、对外投资与转型类公告 0 条、回购与股权激励类公告 61 条，合计公告 464 条。 主营结构的口径变化已经落到报表上：2021 年年报的第一大产品为“商品房”（68.8\%），2026 年半年报为“猪业产品”（71.7\%）；公司在这两期的分部划分方式有过调整。 综合判断，报告期内公司有 1 条与战略相关的公告落地（重大重组与并购 1 条、对外投资与转型 0 条），资本运作与主业调整之间存在直接关联。

\pfl{历史融投资情况} 从现金流量表看，2025 年公司取得借款收到的现金 8.22 亿元、偿还债务支付的现金 11.92 亿元、吸收权益性投资 0.07 亿元，筹资活动现金流净额 -8.56 亿元（2023 年为取得借款 13.80 亿元、偿还债务 5.13 亿元）。 2026 年上半年筹资活动现金流净额为 -0.04 亿元（取得借款 9.33 亿元、偿还债务 7.23 亿元，未年化），已转为净流出，处于净还债状态。 2025 年分配股利、利润或偿付利息支付的现金合计 4.57 亿元。 公开披露的融资记录共 2 笔，工具类型以股权融资、定向增发为主；金额与用途见下表。 其中规模最大的两笔为：2024 年 11 月，定向增发，拟募资不超过 19.30 亿元，拟发行股票数量不超过 1.59 亿股；2019---2024 年，股权融资，上市三十年股权融资总额不过 2.70 亿元。 股本方面，近三年披露股本变动 4 次（增发新股上市 1 次、股权激励 1 次）；总股本由 2021 年末的 5.23 亿股变为 5.31 亿股（+1.40\%）。 分红方面，近三年现金分红 4 次、合计每股派现 2.11 元，最近一次实施在 2025 年 6 月。 近三年“再融资”类公告 9 条，涉及定向增发、募集资金使用。 综合看，公司融资结构上以债务滚动为主、股权融资为补充；2025 年筹资活动现金流净额 -8.56 亿元，融资节奏仍以维持存量债务周转为主，与前述经营与投资现金流的方向基本一致。

\pfl{未来潜在融资需求分析} 2026 年 6 月末货币资金 6.08 亿元，而短期债务 13.29 亿元（现金短债比 0.46 倍），2026 年上半年经营现金流净额 0.41 亿元。按“短期债务减去账面货币资金”的滚动偿付口径，静态缺口约 7.22 亿元；若再计入上半年亏损的年化额（9.78 亿元），维持一年正常经营所需的外部资金约 7.22---17.00 亿元。 公司 2026 年上半年归母净利润为负（-4.89 亿元），但 6 月末资产负债率 63.6\% 尚未进入第一档（亏损且负债率 $>$65\%）的区间，当前融资需求以补充流动资金、支撑低谷期经营性支出为主。 公开披露的后续资金安排线索包括：股权融资线索：2024 年 11 月，3.00 亿元定增预案为上市以来最大规模股权融资计划。 资金需求还要叠加在建工程：期末在建工程 1.32 亿元、2026 年上半年购建固定资产等支出 0.95 亿元（未年化），这部分支出在周期底部通常会被主动放缓。 综合看，公司的外部资金需求以补充流动性与项目投入为主，滚动偿付口径下的静态缺口约 7.22 亿元，融资的时点与规模更多取决于信用条件与资本市场窗口。




\subsubsection{益生股份（002458）}
\label{co:002458}

\pfl{公司基本情况} 1997 年设立至今，益生股份是国内最大的祖代白羽肉种鸡企业，祖代肉种鸡饲养规模连续二十年居全国第一，形成“种鸡＋种猪”双轮驱动；2026 年上半年受益白羽肉鸡行业景气上行与引种受限，净利润同比增长近 49 倍。 公司于 2010-08-10 在深交所主板首发上市，注册资本 14.31 亿元，总股本 11.06 亿股。 截至 2026 年 10 月 12 日收盘价 7.90 元，对应总市值 87.4 亿元，近一年-8.6\%，总市值在三级行业“肉鸡养殖”的 8 家公司中按由大到小排第\zhnumber{3}位，近一年涨跌幅高于全样本中位数（-10.4\%）1.8 个百分点。 按 2026 年半年报披露的行业构成，收入占比前三位为鸡收入 83.8\%；猪收入 13.5\%；乳品收入 1.2\%。

\pfl{股权结构与控制权} 按 2026 年 6 月 30 日的股东名册，前三大股东为曹积生（持股比例未披露）；深圳毕升私募证券基金管理有限公司－必胜年年升 1 号私募（持股比例未披露）；济南泰浩投资合伙企业(有限合伙)（持股比例未披露），股东户数 35,308 户。 股权结构的另一条线索是，控股股东、实际控制人曹积生持股 5.87 亿股、占总股本 41.02\%，累计质押 2.37 亿股；总股本由 11.06 亿股增至 14.31 亿股。

\pfl{经营要点} 经营层面，2025 年度，占全国祖代肉鸡进口总量的比例超 42\%；报告期内父母代肉种鸡销量和销售均价同比均大幅上涨。

\pfl{所处行业的基本面} 白羽与黄羽肉鸡的供给取决于祖代引种量与父母代存栏，需求端受餐饮与替代效应驱动，产能周期只有 6---12 个月（第 \ref{sec:framework} 章的供需框架）。 与该公司直接相关的行业数据包括：2026 年，1—5 月父母代鸡苗市场成交均价 40.63 元/套、同比上涨 21.03\%；报告期内我国祖代白羽肉鸡更新量同比下降 9.34\%。

\pfl{竞争优势} 把该公司放回同类公司比较，2026 年上半年营业收入 17.0 亿元、在三级行业“肉鸡养殖”的 8 家公司中排第\zhnumber{6}位；毛利率 28.0\%、高于行业内中位数 17.3 个百分点；半年 ROE 6.9\%（未年化）、高于行业内中位数（1.21\%）；资产负债率 35.1\%、低于行业内中位数（52.1\%）。 公司在公开资料中自述与同业比较的相对优势在于：2025 年祖代引种占全国进口量超 42\% 居首。 优势在报表上主要体现为高于行业平均的毛利率，这在周期底部是抵御亏损的主要缓冲。

\pfl{历史股价} 自 2021 年 1 月至 2026 年 10 月，股价由 10.20 元变为 7.90 元，累计 -22.5\%，区间最高 15.59 元（2023 年 2 月）、最低 7.77 元（2026 年 6 月）；当前价相当于区间最高价的 50.7\%，已接近区间最低价（1.02 倍）。 月度收益的年化波动率 34.9\%，低于全样本中位数 37.3\%，在三级行业“肉鸡养殖”的 8 家公司中波动率由高到低排第\zhnumber{5}位。 把股价阶段与公司事件对齐看，2023 年 2 月 至 2024 年 8 月股价下跌-48.3\%，同期（2023 年 12 月）6.00 亿元定增完成。 整体上看，区间的几处大级别拐点都能在公司的资本运作、股权变动或风险事件上找到对应，股价波动有明显的公司层面驱动，而非单纯的行业 beta。

\begin{center}
\begin{minipage}{\textwidth}\centering
\begin{tikzpicture}
\begin{axis}[
  width=12.6cm, height=3.9cm, scale only axis,
  xmin=2020.922, xmax=2026.972, ymin=7.23, ymax=16.68,
  xtick={2021,2022,2023,2024,2025,2026},
  yearticks,
  yticklabel style={font=\scriptsize},
  ylabel={元/股}, ylabel style={font=\scriptsize},
  grid=major, grid style={LightGray, line width=0.3pt},
  axis line style={InkGray, line width=0.5pt},
  every axis plot/.append style={line width=0.9pt},
]
\addplot[AgriGreen] table[x=x,y=y]{data/clean/profiles/px_002458.dat};
\addplot[SoftRed, dashed, line width=0.7pt] table[x=x,y=y]{data/clean/profiles/pxzz_002458.dat};
\addplot[only marks, mark=*, mark size=1.1pt, SoftRed] table[x=x,y=y]{data/clean/profiles/pxzz_002458.dat};
\end{axis}
\end{tikzpicture}

\captionof{figure}[益生股份（002458）股价与周期骨架]{益生股份（002458）月度收盘价（元/股）与 ZigZag 周期骨架（阈值 25\%，圆点为周期转折点）}
\end{minipage}
\end{center}

\pfl{过去周期分析} 以 25\% 的 ZigZag 阈值划分，样本区间内股价形成 2 个主升段与 3 个主跌段。 最大主跌段为 2023 年 2 月 至 2024 年 8 月，18 个月-48.3\%；最大主升段出现在 2022 年 5 月 至 2023 年 2 月，9 个月+78.0\%。 最近一次转折出现在 2026 年 6 月（下跌段自 2026 年 4 月起，2 个月-30.7\%），当前处于该段的延续之中。 公司股价与申万农林牧渔指数几乎同步见底，个股并未走出独立行情。 禽链的产能周期只有 6---12 个月、反转快于生猪（第 \ref{sec:other-livestock} 节），2026 年前四个月毛鸡养殖利润同比大幅回升，链条利润的修复先于生猪，公司 2026 年上半年实现盈利、半年 ROE 6.9\%（未年化），好于全样本中位数（0.75\%），下行周期对其报表的冲击小于同业。 综合区间形态看，该股单段涨跌幅度偏大、以趋势段为主（最大主升 +78.0\%、最大主跌 -48.3\%），年化波动率 34.9\% 与全样本中位数 37.3\% 相比波动与样本中位数接近。

\pfl{盈利情况} 2026 年上半年营业收入 16.97 亿元（同比 +28.4\%），归母净利润 3.08 亿元，同比 +4897.3\%。 以年报为趋势对照，2023---2025 年营业收入由 32.25 亿元变为 29.52 亿元（两年年均复合增速 -4.3\%），归母净利润由 5.41 亿元变为 1.65 亿元。 按半年报口径，2026 年上半年的 ROE 为 6.9\%（未年化）、销售净利率 18.2\%、毛利率 28.0\%，ROE 在“肉鸡养殖”8 家公司中按数值由高到低排第\zhnumber{1}位。 同期全样本半年 ROE 中位数 0.75\%，公司与之相差 6.11 个百分点（略好于中位数公司）。 综合看，公司在报告期维持盈利（高于全样本半年 ROE 中位数 0.75\%），利润的可持续性取决于肉鸡养殖景气的延续与成本管控的执行。

\begin{center}\small\setlength{\tabcolsep}{3.4pt}
\captionof{table}[益生股份（002458）关键财务指标]{益生股份（002458）关键财务指标（合并报表口径；两个半年报期均未年化；现金短债比 $=$ 货币资金 $\div$ 短期债务）}
\begin{adjustbox}{max width=\textwidth}
\begin{tabular}{@{}lrrrrr@{}}
\toprule
指标 & 2023 年 & 2024 年 & 2025 年 & 2025 年半年报 & 2026 年半年报 \\
\midrule
营业收入（亿元） & 32.25 & 31.36 & 29.52 & 13.21 & 16.97 \\
归母净利润（亿元） & 5.41 & 5.04 & 1.65 & 0.06 & 3.08 \\
毛利率（\%） & 31.6 & 29.4 & 17.8 & 12.1 & 28.0 \\
ROE（\%） & 17.4 & 11.3 & 3.8 & 0.1 & 6.9 \\
资产负债率（\%） & 31.0 & 39.6 & 35.7 & 42.4 & 35.1 \\
经营活动现金流净额（亿元） & 10.04 & 7.53 & 5.62 & 1.43 & 4.61 \\
货币资金（亿元） & 9.09 & 16.27 & 10.01 & 17.65 & 14.79 \\
有息负债合计（亿元） & 10.67 & 20.47 & 17.52 & 24.50 & 18.88 \\
现金短债比（倍） & 1.08 & 0.88 & 0.65 & 0.79 & 0.90 \\
\bottomrule
\end{tabular}
\end{adjustbox}
\end{center}

\begin{center}
\begin{minipage}{\textwidth}\centering
\begin{tikzpicture}
\begin{groupplot}[
  group style={group size=2 by 1, horizontal sep=0.60cm},
  width=6.85cm, height=4.60cm,
  tick label style={font=\tiny},
  title style={font=\scriptsize\heiti},
  yticklabel style={font=\tiny},
  ylabel style={font=\tiny},
  grid=major, grid style={LightGray, line width=0.3pt},
  axis line style={InkGray, line width=0.5pt},
  enlarge x limits={abs=0.8},
  legend style={font=\scriptsize, at={(0.5,-0.20)}, anchor=north,
                legend columns=2, draw=none, fill=none, column sep=6pt},
]
\nextgroupplot[
  ybar, bar width=4pt,
  ymin=-7.26, ymax=36.56,
  xtick={0,1,2,3,4,5.7},
  xticklabels={2021,2022,2023,2024,2025,2026H1},
  ylabel={亿元},
]
\addplot[fill=AgriGreen!75, draw=AgriGreen, bar shift=-2.6pt] table[x=i, y=rev]{data/clean/profiles/fin_002458.dat};
\addplot[fill=SoftRed!60, draw=SoftRed, bar shift=2.6pt] table[x=i, y=ni]{data/clean/profiles/fin_002458.dat};
\addplot[fill=AgriGreen!30, draw=AgriGreen, pattern=north east lines, bar shift=-2.6pt] table[x=x, y=rev]{data/clean/profiles/finh_002458.dat};
\addplot[fill=SoftRed!20, draw=SoftRed, pattern=north east lines, bar shift=2.6pt] table[x=x, y=ni]{data/clean/profiles/finh_002458.dat};
\legend{营业收入（年/半年）, 归母净利润（年/半年）}
\nextgroupplot[
  ymin=-17.5, ymax=58.4,
  xtick={0,1,2,3,4,5.7},
  xticklabels={2021,2022,2023,2024,2025,2026H1},
  ylabel={\%},
]
\addplot[AgriGold, mark=*, mark size=1.3pt] table[x=i, y=roe]{data/clean/profiles/fin_002458.dat};
\addplot[OilBlue, mark=square*, mark size=1.3pt] table[x=i, y=debt]{data/clean/profiles/fin_002458.dat};
\addplot[only marks, AgriGold, mark=*, mark size=2.0pt] table[x=x, y=roe]{data/clean/profiles/finh_002458.dat};
\addplot[only marks, OilBlue, mark=square*, mark size=2.0pt] table[x=x, y=debt]{data/clean/profiles/finh_002458.dat};
\legend{ROE, 资产负债率}
\end{groupplot}
\end{tikzpicture}
\begin{tikzpicture}
\begin{axis}[
  name=finbot,
  width=13.75cm, height=3.10cm, scale only axis,
  ymin=-2.05, ymax=22.93,
  xtick={0,1,2,3,4,5.7},
  xticklabels={2021,2022,2023,2024,2025,2026H1},
  tick label style={font=\tiny},
  yticklabel style={font=\tiny},
  ylabel={亿元}, ylabel style={font=\tiny},
  ybar, bar width=4pt,
  grid=major, grid style={LightGray, line width=0.3pt},
  axis line style={InkGray, line width=0.5pt},
  enlarge x limits={abs=0.8},
  legend style={font=\scriptsize, at={(0.5,-0.26)}, anchor=north,
                legend columns=2, draw=none, fill=none, column sep=10pt},
]
\addplot[fill=AgriGreen!55, draw=AgriGreen, bar shift=-3.0pt] table[x=i, y=cash]{data/clean/profiles/fin_002458.dat};
\addplot[fill=SoftRed!45, draw=SoftRed, bar shift=3.0pt] table[x=i, y=ibd]{data/clean/profiles/fin_002458.dat};
\addplot[fill=AgriGreen!25, draw=AgriGreen, pattern=north east lines, bar shift=-3.0pt] table[x=x, y=cash]{data/clean/profiles/finh_002458.dat};
\addplot[fill=SoftRed!20, draw=SoftRed, pattern=north east lines, bar shift=3.0pt] table[x=x, y=ibd]{data/clean/profiles/finh_002458.dat};
\legend{货币资金（左轴）, 有息负债（左轴）}
\end{axis}
\begin{axis}[
  at={(finbot.south east)}, anchor=south east,
  width=13.75cm, height=3.10cm, scale only axis,
  axis y line*=right, axis x line*=none, xtick=\empty,
  ymin=-2.61, ymax=13.05,
  tick label style={font=\tiny},
  yticklabel style={font=\tiny}, scaled y ticks=false,
  legend style={font=\scriptsize, at={(0.5,-0.44)}, anchor=north,
                legend columns=1, draw=none, fill=none},
]
\addplot[OilBlue, mark=*, mark size=2.2pt, line width=1.3pt] table[x=i, y=ocf]{data/clean/profiles/fin_002458.dat};
\addplot[only marks, OilBlue, mark=*, mark size=3.2pt] table[x=x, y=ocf]{data/clean/profiles/finh_002458.dat};
\legend{经营活动现金流净额（右轴，亿元，蓝点）}
\end{axis}
\end{tikzpicture}

\captionof{figure}[益生股份（002458）近年主要财务指标]{益生股份（002458）近年主要财务指标（左上：营业收入与归母净利润；右上：ROE 与资产负债率；下：货币资金、有息负债为左轴柱状，经营活动现金流净额为其右轴的蓝点折线。
2021---2025 年为年报口径，2026H1 为 2026 年半年报/半年末，与其前一年报之间留出间隔、以斜纹或空心点区分；半年报数据未年化）}
\end{minipage}
\end{center}

\pfl{财务分析} 2026 年 6 月末资产负债率 35.1\%，较 2025 年末下降 0.5 个百分点（样本中位数 52.2\%，所处三级行业内按由低到高排第\zhnumber{1}位）。 2026 年 6 月末货币资金 14.79 亿元、短期债务（短期借款与一年内到期非流动负债）16.42 亿元、长期债务 2.46 亿元，有息负债合计 18.88 亿元，现金短债比 0.90 倍——覆盖偏紧。 净有息负债 4.09 亿元，相当于其 2026 年上半年营业收入的 24\%。 流动比率 0.91、速动比率 0.73，短期偿付压力较大。 2026 年上半年经营活动现金流净额 4.61 亿元，经营现金流与利润同向为正，内生资金可以支撑运营。 2026 年 6 月末在建工程 0.66 亿元，较上年末减少 0.33 亿元，在建投入在收缩。 综合看，现金短债比 0.90 倍，短期资金安排偏紧；资产负债率低于样本中位数 17.1 个百分点；有息负债中短期债务占比更高，期限结构仍以滚动续作为主。

\pfl{重大战略转型} 从公开披露看，公司的战略推进集中在：并购与重组（1 项）：2020---2021 年，2021 年以 1,360.00 万元收购河北润源养殖 100\% 股权。 公告口径上，近三年（2023 年 9 月---2026 年 9 月）公司披露重大重组与并购类公告 0 条、对外投资与转型类公告 2 条、回购与股权激励类公告 52 条，合计公告 454 条。 第一大行业仍是“鸡收入”，收入占比 83.8\%，与 2021 年年报相比没有方向性变化。 综合判断，报告期内公司有 2 条与战略相关的公告落地（重大重组与并购 0 条、对外投资与转型 2 条），资本运作与主业调整之间存在直接关联。

\begin{center}\small\setlength{\tabcolsep}{4pt}
\captionof{table}[益生股份（002458）历史融资与资本运作]{益生股份（002458）历史融资与资本运作（据公司公告与公开报道整理；金额为披露值，含首次公开发行、再融资、债券、银行借款与重整投资等）}
\begin{adjustbox}{max width=\textwidth}
\begin{tabular}{@{}p{1.45cm}p{2.55cm}p{4.05cm}p{4.95cm}@{}}
\toprule
时间 & 方式 & 规模 & 用途与说明 \\
\midrule
2010-08 & IPO（深交所中小板） & 2,\allowbreak{}700 万股，\allowbreak{}发行价 24.00 元/\allowbreak{}股（约 6.48 亿元） & 种鸡养殖及孵化产能扩张等项目 \\
2016-01 & 定向增发（非公开发行） & 5,\allowbreak{}131.46 万股，\allowbreak{}发行价 11.55 元/\allowbreak{}股 & 种鸡养殖项目等；承销及保荐费用 1,\allowbreak{}000.00 万元；为公司 2016 年至 2021 年之前的最近一次再融资 \\
2021-10 & 可转换公司债券（申请未获核准） & 原拟公开募集不超过 10.40 亿元，\allowbreak{}修订稿缩减为 8.90 亿元 & 4 个种猪养殖项目、\allowbreak{}3 个饲料生产项目和 2 个种鸡孵化场项目（修订后各砍掉一个种猪、\allowbreak{}种鸡项目）；公司 2022 年公告可转债发行申请未获证监会核准 \\
2022-09-29 & 定向增发（2022 年度向特定对象发行股票） & 1.14 亿股，\allowbreak{}发行价 10.22 元/\allowbreak{}股 & 种猪养殖项目（双鸭山祖代种猪场 3.90 亿元、\allowbreak{}山西 3,\allowbreak{}600 头能繁原种母猪场及公猪站 1.60 亿元等）、\allowbreak{}种鸡孵化场项目及补充流动资金、\allowbreak{}偿还银行贷款；2022 年 9 月 29 日披露预案（拟募资不超过 11.60 亿元）；；募集资金 2023 年 11 月 24 日到位 \\
\bottomrule
\end{tabular}
\end{adjustbox}
\end{center}

\pfl{历史融投资情况} 从现金流量表看，2025 年公司取得借款收到的现金 32.68 亿元、偿还债务支付的现金 28.66 亿元，筹资活动现金流净额 -6.30 亿元（2023 年为取得借款 9.28 亿元、偿还债务 17.12 亿元）。 2026 年上半年筹资活动现金流净额为 -0.34 亿元（取得借款 12.30 亿元、偿还债务 9.94 亿元，未年化），已转为净流出，处于净还债状态。 2025 年分配股利、利润或偿付利息支付的现金合计 4.25 亿元。 公开披露的融资记录共 4 笔，工具类型以 IPO、定向增发、可转债为主；金额与用途见下表。 其中规模最大的两笔为：2022 年 9 月，定向增发，1.14 亿股，发行价 10.22 元/股；2010 年 8 月，IPO，2,700 万股，发行价 24.00 元/股（约 6.48 亿元）。 股本方面，近三年披露股本变动 5 次（增发新股上市 1 次、限售股份上市 1 次）；总股本由 2021 年末的 9.93 亿股变为 11.06 亿股（+11.43\%）。 分红方面，近三年现金分红 6 次、合计每股派现 0.70 元，最近一次实施在 2026 年 6 月。 近三年“再融资”类公告 10 条，涉及定向增发、募集资金使用。 综合看，公司融资结构上股权与债务融资并举；2025 年筹资活动现金流净额 -6.30 亿元，融资节奏仍以维持存量债务周转为主，与前述经营与投资现金流的方向基本一致。

\pfl{未来潜在融资需求分析} 2026 年 6 月末货币资金 14.79 亿元，而短期债务 16.42 亿元（现金短债比 0.90 倍），2026 年上半年经营现金流净额 4.61 亿元。按“短期债务减去账面货币资金”的滚动偿付口径，静态缺口约 1.63 亿元；上半年盈利或接近盈亏平衡，该缺口不随经营亏损进一步扩大。 按第 \ref{sec:financing-need} 节的三档划分，公司属于第三档“扩张型”：2026 年上半年 ROE 6.9\%（未年化，折合约 13.7\%）、6 月末资产负债率 35.1\%，资产负债表仍具备加杠杆空间；融资用途以产能扩张、品类并购与海外布局为主，单笔规模通常在 5---20 亿元量级，会被市场定价为成长而非补血。 公开披露的后续资金安排线索包括：股权融资线索：2025 年，累计使用 7.55 亿元（占 66.24\%）。 资金需求还要叠加在建工程：期末在建工程 0.66 亿元、2026 年上半年购建固定资产等支出 1.21 亿元（未年化），这部分支出在周期底部通常会被主动放缓。 综合看，公司的外部资金需求以营运资金周转与在建项目投入为主，滚动偿付口径下的静态缺口约 1.63 亿元，能否落地取决于授信续作与发行窗口的配合。




\subsubsection{巨星农牧（603477）}
\label{co:603477}

\pfl{公司基本情况} 四川乐山的生猪养殖与饲料企业，前身为皮革企业振静股份（2017 年上交所上市），2020 年发行股份购买巨星农牧 100\% 股权后主业转型为生猪养殖，2023 年 9 月控股股东变更为四川巨星集团、实际控制人变更为唐光跃；2026 年上半年出栏 280.77 万头但受猪价下行拖累由盈转亏。 公司于 2017-12-18 在上交所首发上市，注册资本 5.23 亿元，总股本 5.10 亿股。 截至 2026 年 10 月 12 日收盘价 14.08 元，对应总市值 71.8 亿元，近一年-30.3\%，总市值在三级行业“生猪养殖”的 11 家公司中按由大到小排第\zhnumber{6}位，近一年涨跌幅低于全样本中位数（-10.4\%）19.9 个百分点。 按 2026 年半年报披露的行业构成，收入占比前两位为农牧业务 96.9\%；皮革业务 3.1\%。 与 2021 年年报相比，第一大行业口径由“养殖行业”（83.1\%）变为“农牧业务”（96.9\%）；两期披露的分类口径有调整，占比差异中同时包含分类因素。

\pfl{股权结构与控制权} 按 2026 年 6 月 30 日的股东名册，前三大股东为四川巨星企业集团有限公司 持股 29.25\%；四川和邦投资集团有限公司 持股 11.48\%；乐山巨星农牧股份有限公司－第一期员工持股计划 持股 4.27\%，股东户数 41,273 户。 从控制权角度看，2026 年上半年，2.27 亿元定向增发落地实施。

\pfl{经营要点} 从经营要点看，2026 年上半年，商品肥猪完全成本控制在 6.3 元/斤以内；2026 年 7 月，育肥成活率保持 95\%，公司挑战目标为 33 头 PSY、31 头 MSY。

\pfl{所处行业的基本面} 生猪的供给由能繁母猪存栏决定、需求由猪肉消费与替代品价格决定，本轮下行的成因在供给端的产能去化不彻底，而非需求收缩（第 \ref{sec:framework} 章的供需框架）。 与该公司直接相关的行业数据包括：2026 年上半年，9 月 29 日四川生猪市场价格约 10 元/公斤；2026 年二季度，末全国能繁母猪存栏 3780 万头、同比下降 6.5\%；2026 年 5 月，农业农村部将能繁母猪正常保有量下调至 3750 万头左右，要求大型养殖企业带头压减产能。

\pfl{竞争优势} 把该公司放回同类公司比较，2026 年上半年营业收入 35.6 亿元、在三级行业“生猪养殖”的 11 家公司中排第\zhnumber{5}位；毛利率 -9.7\%、低于行业内中位数 8.2 个百分点；半年 ROE -30.3\%（未年化）、低于行业内中位数（-14.32\%）；资产负债率 69.4\%、高于行业内中位数（58.9\%）。 相对优势方面，成本控制能力处于生猪养殖上市公司第一梯队；2025 年带动合作农户增收约 8.65 亿元。 就当期报表看，公司尚未把经营层面的特征转化为高于行业的资本回报，周期底部的相对优势更多体现在现金流与资产负债表上。

\pfl{历史股价} 本报告样本区间（2021 年 1 月 至 2026 年 10 月）内，股价由 10.85 元变为 14.08 元，累计 +29.8\%，区间最高 37.47 元（2023 年 12 月）、最低 9.99 元（2021 年 7 月）；当前价相当于区间最高价的 37.6\%，为区间最低价的 1.41 倍。 月度收益的年化波动率 49.0\%，高于全样本中位数 37.3\%，在三级行业“生猪养殖”的 11 家公司中波动率由高到低排第\zhnumber{2}位。 把股价阶段与公司事件对齐看，2022 年 10 月 至 2023 年 12 月股价上涨+95.7\%，同期（2023 年 8 月）四川和邦投资集团有限公司与四川巨星企业集团有限公司签署股份转让协议，以 27.67 元/股向巨星集团转让 2,600 万股（占总股本 5.14\%）。 整体上看，区间的几处大级别拐点都能在公司的资本运作、股权变动或风险事件上找到对应，股价波动有明显的公司层面驱动，而非单纯的行业 beta。

\begin{center}
\begin{minipage}{\textwidth}\centering
\begin{tikzpicture}
\begin{axis}[
  width=12.6cm, height=3.9cm, scale only axis,
  xmin=2020.922, xmax=2026.972, ymin=9.29, ymax=40.09,
  xtick={2021,2022,2023,2024,2025,2026},
  yearticks,
  yticklabel style={font=\scriptsize},
  ylabel={元/股}, ylabel style={font=\scriptsize},
  grid=major, grid style={LightGray, line width=0.3pt},
  axis line style={InkGray, line width=0.5pt},
  every axis plot/.append style={line width=0.9pt},
]
\addplot[AgriGreen] table[x=x,y=y]{data/clean/profiles/px_603477.dat};
\addplot[SoftRed, dashed, line width=0.7pt] table[x=x,y=y]{data/clean/profiles/pxzz_603477.dat};
\addplot[only marks, mark=*, mark size=1.1pt, SoftRed] table[x=x,y=y]{data/clean/profiles/pxzz_603477.dat};
\end{axis}
\end{tikzpicture}

\captionof{figure}[巨星农牧（603477）股价与周期骨架]{巨星农牧（603477）月度收盘价（元/股）与 ZigZag 周期骨架（阈值 25\%，圆点为周期转折点）}
\end{minipage}
\end{center}

\pfl{过去周期分析} 以 25\% 的 ZigZag 阈值划分，样本区间内股价形成 5 个主升段与 5 个主跌段。 最大主升段出现在 2021 年 7 月 至 2022 年 3 月，8 个月+156.9\%；最大主跌段为 2023 年 12 月 至 2025 年 1 月，13 个月-55.3\%。 最近一次转折出现在 2026 年 10 月（下跌段自 2025 年 5 月起，17 个月-35.6\%），当前处于该段的延续之中。 公司股价较申万农林牧渔指数提前约 58 个月见底，是板块中较早定价周期反转的公司之一。 生猪养殖的定价主线是产能去化与成本曲线的赛跑（第 \ref{sec:hog} 章）：2026 年三季度行业仍处在“产能去化未完成、价格磨底”的阶段，公司 2026 年 6 月末资产负债率 69.4\%，高于全样本中位数 52.2\%，其股价因此更多反映资金面与信用风险，而不是价格弹性本身。 综合区间形态看，该股单段涨跌幅度偏大、以趋势段为主（最大主升 +156.9\%、最大主跌 -55.3\%），年化波动率 49.0\% 与全样本中位数 37.3\% 相比波动明显大于样本中位数。

\pfl{盈利情况} 2026 年上半年营业收入 35.57 亿元（同比 -4.3\%），归母净利润 -8.38 亿元，由上年同期的盈利转为亏损。 以年报为趋势对照，2023---2025 年营业收入由 40.41 亿元变为 80.00 亿元（两年年均复合增速 +40.7\%），归母净利润由 -6.45 亿元变为 0.30 亿元。 按半年报口径，2026 年上半年的 ROE 为 -30.3\%（未年化）、销售净利率 -23.6\%、毛利率 -9.7\%，ROE 在“生猪养殖”11 家公司中按数值由高到低排第\zhnumber{9}位。 同期全样本半年 ROE 中位数 0.75\%，公司与之相差 31.07 个百分点（弱于中位数公司）。 综合看，公司仍处亏损区间、由盈转亏，半年 ROE 在三级行业内按由高到低排第\zhnumber{9}位，单位盈利能力的修复依赖行业景气与自身成本端的改善，报表弹性主要体现为价格与销量的方向性变化。

\begin{center}\small\setlength{\tabcolsep}{3.4pt}
\captionof{table}[巨星农牧（603477）关键财务指标]{巨星农牧（603477）关键财务指标（合并报表口径；两个半年报期均未年化；现金短债比 $=$ 货币资金 $\div$ 短期债务）}
\begin{adjustbox}{max width=\textwidth}
\begin{tabular}{@{}lrrrrr@{}}
\toprule
指标 & 2023 年 & 2024 年 & 2025 年 & 2025 年半年报 & 2026 年半年报 \\
\midrule
营业收入（亿元） & 40.41 & 60.78 & 80.00 & 37.17 & 35.57 \\
归母净利润（亿元） & -6.45 & 5.19 & 0.30 & 1.81 & -8.38 \\
毛利率（\%） & -0.4 & 19.2 & 7.3 & 13.7 & -9.7 \\
ROE（\%） & -19.3 & 15.6 & 0.9 & 5.4 & -30.3 \\
资产负债率（\%） & 59.9 & 61.9 & 67.3 & 64.6 & 69.4 \\
经营活动现金流净额（亿元） & -1.05 & 12.65 & 12.18 & 6.02 & -2.86 \\
货币资金（亿元） & 4.46 & 5.77 & 7.56 & 6.70 & 6.94 \\
有息负债合计（亿元） & 31.42 & 34.47 & 39.84 & 36.21 & 37.45 \\
现金短债比（倍） & 0.29 & 0.30 & 0.36 & 0.33 & 0.34 \\
\bottomrule
\end{tabular}
\end{adjustbox}
\end{center}

\begin{center}
\begin{minipage}{\textwidth}\centering
\begin{tikzpicture}
\begin{groupplot}[
  group style={group size=2 by 1, horizontal sep=0.60cm},
  width=6.85cm, height=4.60cm,
  tick label style={font=\tiny},
  title style={font=\scriptsize\heiti},
  yticklabel style={font=\tiny},
  ylabel style={font=\tiny},
  grid=major, grid style={LightGray, line width=0.3pt},
  axis line style={InkGray, line width=0.5pt},
  enlarge x limits={abs=0.8},
  legend style={font=\scriptsize, at={(0.5,-0.20)}, anchor=north,
                legend columns=2, draw=none, fill=none, column sep=6pt},
]
\nextgroupplot[
  ybar, bar width=4pt,
  ymin=-17.22, ymax=90.61,
  xtick={0,1,2,3,4,5.7},
  xticklabels={2021,2022,2023,2024,2025,2026H1},
  ylabel={亿元},
]
\addplot[fill=AgriGreen!75, draw=AgriGreen, bar shift=-2.6pt] table[x=i, y=rev]{data/clean/profiles/fin_603477.dat};
\addplot[fill=SoftRed!60, draw=SoftRed, bar shift=2.6pt] table[x=i, y=ni]{data/clean/profiles/fin_603477.dat};
\addplot[fill=AgriGreen!30, draw=AgriGreen, pattern=north east lines, bar shift=-2.6pt] table[x=x, y=rev]{data/clean/profiles/finh_603477.dat};
\addplot[fill=SoftRed!20, draw=SoftRed, pattern=north east lines, bar shift=2.6pt] table[x=x, y=ni]{data/clean/profiles/finh_603477.dat};
\legend{营业收入（年/半年）, 归母净利润（年/半年）}
\nextgroupplot[
  ymin=-38.3, ymax=79.4,
  xtick={0,1,2,3,4,5.7},
  xticklabels={2021,2022,2023,2024,2025,2026H1},
  ylabel={\%},
]
\addplot[AgriGold, mark=*, mark size=1.3pt] table[x=i, y=roe]{data/clean/profiles/fin_603477.dat};
\addplot[OilBlue, mark=square*, mark size=1.3pt] table[x=i, y=debt]{data/clean/profiles/fin_603477.dat};
\addplot[only marks, AgriGold, mark=*, mark size=2.0pt] table[x=x, y=roe]{data/clean/profiles/finh_603477.dat};
\addplot[only marks, OilBlue, mark=square*, mark size=2.0pt] table[x=x, y=debt]{data/clean/profiles/finh_603477.dat};
\legend{ROE, 资产负债率}
\end{groupplot}
\end{tikzpicture}
\begin{tikzpicture}
\begin{axis}[
  name=finbot,
  width=13.75cm, height=3.10cm, scale only axis,
  ymin=-3.98, ymax=44.62,
  xtick={0,1,2,3,4,5.7},
  xticklabels={2021,2022,2023,2024,2025,2026H1},
  tick label style={font=\tiny},
  yticklabel style={font=\tiny},
  ylabel={亿元}, ylabel style={font=\tiny},
  ybar, bar width=4pt,
  grid=major, grid style={LightGray, line width=0.3pt},
  axis line style={InkGray, line width=0.5pt},
  enlarge x limits={abs=0.8},
  legend style={font=\scriptsize, at={(0.5,-0.26)}, anchor=north,
                legend columns=2, draw=none, fill=none, column sep=10pt},
]
\addplot[fill=AgriGreen!55, draw=AgriGreen, bar shift=-3.0pt] table[x=i, y=cash]{data/clean/profiles/fin_603477.dat};
\addplot[fill=SoftRed!45, draw=SoftRed, bar shift=3.0pt] table[x=i, y=ibd]{data/clean/profiles/fin_603477.dat};
\addplot[fill=AgriGreen!25, draw=AgriGreen, pattern=north east lines, bar shift=-3.0pt] table[x=x, y=cash]{data/clean/profiles/finh_603477.dat};
\addplot[fill=SoftRed!20, draw=SoftRed, pattern=north east lines, bar shift=3.0pt] table[x=x, y=ibd]{data/clean/profiles/finh_603477.dat};
\legend{货币资金（左轴）, 有息负债（左轴）}
\end{axis}
\begin{axis}[
  at={(finbot.south east)}, anchor=south east,
  width=13.75cm, height=3.10cm, scale only axis,
  axis y line*=right, axis x line*=none, xtick=\empty,
  ymin=-6.89, ymax=17.30,
  tick label style={font=\tiny},
  yticklabel style={font=\tiny}, scaled y ticks=false,
  legend style={font=\scriptsize, at={(0.5,-0.44)}, anchor=north,
                legend columns=1, draw=none, fill=none},
]
\addplot[OilBlue, mark=*, mark size=2.2pt, line width=1.3pt] table[x=i, y=ocf]{data/clean/profiles/fin_603477.dat};
\addplot[only marks, OilBlue, mark=*, mark size=3.2pt] table[x=x, y=ocf]{data/clean/profiles/finh_603477.dat};
\legend{经营活动现金流净额（右轴，亿元，蓝点）}
\end{axis}
\end{tikzpicture}

\captionof{figure}[巨星农牧（603477）近年主要财务指标]{巨星农牧（603477）近年主要财务指标（左上：营业收入与归母净利润；右上：ROE 与资产负债率；下：货币资金、有息负债为左轴柱状，经营活动现金流净额为其右轴的蓝点折线。
2021---2025 年为年报口径，2026H1 为 2026 年半年报/半年末，与其前一年报之间留出间隔、以斜纹或空心点区分；半年报数据未年化）}
\end{minipage}
\end{center}

\pfl{财务分析} 2026 年 6 月末资产负债率 69.4\%，较 2025 年末上升 2.1 个百分点（样本中位数 52.2\%，所处三级行业内按由低到高排第\zhnumber{8}位）。 2026 年 6 月末货币资金 6.94 亿元、短期债务（短期借款与一年内到期非流动负债）20.13 亿元、长期债务 17.32 亿元，有息负债合计 37.45 亿元，现金短债比 0.34 倍——覆盖偏紧。 净有息负债 30.51 亿元，相当于其 2026 年上半年营业收入的 86\%。 流动比率 0.80、速动比率 0.23，短期偿债指标偏紧。 2026 年上半年经营活动现金流净额 -2.86 亿元，经营现金流与利润同时为负，日常运营对债务与股东投入的依赖度较高。 2026 年 6 月末在建工程 4.62 亿元，较上年末增加 0.13 亿元，仍有在建投入。 2026 年 6 月末商誉 5.80 亿元，占归母净资产的 21.1\%，是净资产中最脆弱的一块。 综合看，账面货币资金对有息负债与短期债务的覆盖不足，滚动偿付对新增融资的依赖度较高；资产负债率高于样本中位数 17.2 个百分点；有息负债中短期债务占比更高，期限结构仍以滚动续作为主。 公司披露的股东与质押风险为：2024 年 9 月，控股股东巨星集团持有公司 15,066.50 万股，其中累计质押 8,981 万股。

\pfl{重大战略转型} 近三年公司的战略动作集中在：并购与重组（2 项）：2023 年 8 月，四川和邦投资集团有限公司与四川巨星企业集团有限公司签署股份转让协议，以 27.67 元/股向巨星集团转让 2,600 万股（占总股本 5.14\%）；2020 年 7 月，公司通过发行股份及支付现金收购巨星农牧有限公司 100\% 股权（作价 18.20 亿元）；融资与资本运作（1 项）：2024 年 12 月，预计可新增生猪年出栏量 80 万头、现有产能提升约 25\%。 公告口径上，近三年（2023 年 9 月---2026 年 9 月）公司披露重大重组与并购类公告 0 条、对外投资与转型类公告 0 条、回购与股权激励类公告 33 条，合计公告 589 条。 主营结构的口径变化已经落到报表上：2021 年年报的第一大行业为“养殖行业”（83.1\%），2026 年半年报为“农牧业务”（96.9\%）；公司在这两期的分部划分方式有过调整。 综合判断，公司报告期内的变化主要发生在收入结构层面，属于口径与业务重心的再平衡，而非外延式扩张。

\begin{center}\small\setlength{\tabcolsep}{4pt}
\captionof{table}[巨星农牧（603477）历史融资与资本运作]{巨星农牧（603477）历史融资与资本运作（据公司公告与公开报道整理；金额为披露值，含首次公开发行、再融资、债券、银行借款与重整投资等）}
\begin{adjustbox}{max width=\textwidth}
\begin{tabular}{@{}p{1.45cm}p{2.55cm}p{4.05cm}p{4.95cm}@{}}
\toprule
时间 & 方式 & 规模 & 用途与说明 \\
\midrule
2017-12 & IPO（上海证券交易所主板，发行主体原为四川振静股份有限公司） & 实际募集资金 3.35 亿元（发行 6,\allowbreak{}000 万股，\allowbreak{}发行价 5.58 元/\allowbreak{}股） & 皮革主业相关项目，\allowbreak{}其中补充流动资金 7,\allowbreak{}610.12 万元；招股意向书披露预计募资 3.04 亿元 \\
2020-07 & 发行股份及支付现金购买资产（重大资产重组） & 交易作价 18.20 亿元；发行 2.28 亿股、\allowbreak{}发行价 7.47 元/\allowbreak{}股 & 收购巨星农牧有限公司 100\% 股权（评估值 18.21 亿元）；交易对方为四川巨星企业集团有限公司等 11 名非自然人股东及孙德越等 30 名自然人股东；资产过户于 2020 年 7 月完成 \\
2021-06-28 & 非公开发行股票（向 8 名特定对象） & 募集资金总额 4.20 亿元，\allowbreak{}净额 4.08 亿元（3,\allowbreak{}818.18 万股 & 生猪养殖项目（含古蔺巨星皇华种猪场项目 6,\allowbreak{}000.00 万元）及补流还贷；新增股份于 2021 年 7 月 8 日办理完毕登记托管 \\
2022-04-25 & 公开发行可转换公司债券（债券简称“巨星转债”，代码 113648，期限 6 年） & 面值总额 100,\allowbreak{}000.00 万元（1,\allowbreak{}000 万张，\allowbreak{}每张 100 元） & 德昌巨星生猪繁育一体化项目和补充流动资金；；2022 年 10 月 31 日起可转股 \\
2026-04 & 向特定对象发行股票（2023 年度定增，6 名发行对象） & 募集资金总额 2.27 亿元，\allowbreak{}净额 2.19 亿元（1,\allowbreak{}290.88 万股 & 公猪站建设项目、\allowbreak{}养殖技术研究基地建设项目、\allowbreak{}数字智能化建设项目，\allowbreak{}其中约 75\% 用于数字智能化建设项目；认购资金于 2026 年 4 月 21 日到账 \\
\bottomrule
\end{tabular}
\end{adjustbox}
\end{center}

\pfl{历史融投资情况} 从现金流量表看，2025 年公司取得借款收到的现金 19.43 亿元、偿还债务支付的现金 15.34 亿元，筹资活动现金流净额 -3.04 亿元（2023 年为取得借款 19.59 亿元、偿还债务 13.44 亿元）。 2026 年上半年筹资活动现金流净额为 2.30 亿元（取得借款 10.46 亿元、偿还债务 13.00 亿元，未年化），仍处于净融入状态。 2025 年分配股利、利润或偿付利息支付的现金合计 1.99 亿元。 公开披露的融资记录共 5 笔，工具类型以 IPO、发行股份及支付现金购、定向增发、可转债为主；金额与用途见下表。 其中规模最大的两笔为：2022 年 4 月，可转债，面值总额 100,000.00 万元（1,000 万张，每张 100 元）；2017 年 12 月，IPO，实际募集资金 3.35 亿元（发行 6,000 万股，发行价 5.58 元/股）。 股本方面，近三年披露股本变动 25 次（可转债转股 23 次、限售股份上市 1 次、其他 1 次）；总股本由 2021 年末的 5.06 亿股变为 5.10 亿股（+0.79\%）。 分红方面，近三年现金分红 3 次、合计每股派现 0.22 元，最近一次实施在 2025 年 12 月。 近三年“再融资”类公告 75 条，涉及定向增发、可转债、募集资金使用。 综合看，公司融资结构上股权与债务融资并举；2025 年筹资活动现金流净额 -3.04 亿元，融资节奏仍以维持存量债务周转为主，与前述经营与投资现金流的方向基本一致。

\pfl{未来潜在融资需求分析} 2026 年 6 月末货币资金 6.94 亿元，而短期债务 20.13 亿元（现金短债比 0.34 倍），2026 年上半年经营现金流净额 -2.86 亿元。按“短期债务减去账面货币资金”的滚动偿付口径，静态缺口约 13.19 亿元；若再计入上半年亏损的年化额（16.77 亿元），维持一年正常经营所需的外部资金约 13.19---29.96 亿元。 按第 \ref{sec:financing-need} 节的三档划分，公司属于第一档“补血型”：2026 年上半年归母净利润 -8.38 亿元、6 月末资产负债率 69.4\%。 可行的融资路径按现实性排序为：出售非核心资产与参股股权、债务重组或展期（与银行和债券持有人协商）、引入国资或产业投资人（部分公司以控制权让渡为对价）。 公开披露的后续资金安排线索包括：股权融资线索：2024 年，公司积极推动定增；2026 年 2 月，公司向特定对象发行股票申请获上海证券交易所审核通过。 资金需求还要叠加在建工程：期末在建工程 4.62 亿元、2026 年上半年购建固定资产等支出 1.60 亿元（未年化），这部分支出在周期底部通常会被主动放缓。 综合看，公司的外部资金需求以债务接续为主，滚动偿付口径下的静态缺口约 13.19 亿元，融资节奏将受制于银行续贷意愿与股权融资窗口。




\subsubsection{新五丰（600975）}
\label{co:600975}

\pfl{公司基本情况} 湖南省国资委下辖唯一以生猪产业为主业的国有控股上市公司，2004 年上交所上市，2022 年并购天心种业并连续完成多轮再融资，出栏量由 2021 年 44.15 万头增至 2025 年 542.24 万头；2026 年上半年受猪价下行影响营收 28.33 亿元、归母净亏损 7.36 亿元。 公司于 2004-06-09 在上交所首发上市，注册资本 12.50 亿元，总股本 12.61 亿股。 截至 2026 年 10 月 12 日收盘价 5.12 元，对应总市值 64.6 亿元，近一年-19.9\%，总市值在三级行业“生猪养殖”的 11 家公司中按由大到小排第\zhnumber{7}位，近一年涨跌幅低于全样本中位数（-10.4\%）9.5 个百分点。 按 2025 年年报披露的行业构成，收入占比前三位为畜牧业 84.3\%；批发和零售 14.8\%；饲料加工 0.8\%。 与 2021 年年报相比，第一大行业口径由“批发和零售”（55.6\%）变为“畜牧业”（84.3\%）；两期披露的分类口径有调整，占比差异中同时包含分类因素。 股权结构方面，控股股东为湖南省现代农业产业控股集团有限公司；实际控制人为湖南省人民政府国有资产监督管理委员会。

\pfl{股权结构与控制权} 按 2026 年 6 月 30 日的股东名册，前三大股东为湖南省现代农业产业控股集团有限公司 持股 18.69\%；湖南省粮油食品进出口集团有限公司 持股 16.52\%；湖南兴湘投资控股集团有限公司 持股 3.55\%，股东户数 64,152 户。 控制权与股东性质方面，直接持有 2.34 亿股、持股比例 18.69\%；报告期内公司控股股东与实际控制人均未发生变更；报告期末资产负债率约为 67\%。

\pfl{经营要点} 经营与产能方面，2026 年上半年，生猪出栏 276.02 万头，生猪内销量 254.20 万头、同比增加 17.22 万头；截至 2026 年 6 月末能繁母猪存栏由 2025 年末的 24.93 万头调减至 22.09 万头；已完成 205 万头/年屠宰产能布局，其中广联公司 70 万头、舜新项目 50 万头、双新项目 85 万头已于 2026 年投产。

\pfl{所处行业的基本面} 生猪的供给由能繁母猪存栏决定、需求由猪肉消费与替代品价格决定，本轮下行的成因在供给端的产能去化不彻底，而非需求收缩（第 \ref{sec:framework} 章的供需框架）。 行业层面的最新情况是：2026 年上半年，全国生猪价格同比下降 23.1\%，生猪出栏 37246 万头、同比增长 1.7\%；2026 年二季度，末全国能繁母猪存栏已降至 3780 万头；2026 年 5 月，农业农村部将能繁母猪正常保有量下调至 3750 万头左右，大型养殖企业被要求带头压减产能产量。

\pfl{竞争优势} 把该公司放回同类公司比较，2026 年上半年营业收入 28.3 亿元、在三级行业“生猪养殖”的 11 家公司中排第\zhnumber{6}位；毛利率 -5.4\%、低于行业内中位数 3.8 个百分点；半年 ROE -41.0\%（未年化）、低于行业内中位数（-14.32\%）；资产负债率 84.2\%、高于行业内中位数（58.9\%）。 公司区别于同类公司的经营特征包括：与法国科普利信集团合资设立湖南新五丰利信种业有限公司；采购成本整体降幅超过 5\%。 就当期报表看，公司尚未把经营层面的特征转化为高于行业的资本回报，周期底部的相对优势更多体现在现金流与资产负债表上。

\pfl{历史股价} 本报告样本区间（2021 年 1 月 至 2026 年 10 月）内，股价由 7.78 元变为 5.12 元，累计 -34.2\%，区间最高 11.00 元（2023 年 4 月）、最低 4.31 元（2026 年 6 月）；当前价相当于区间最高价的 46.5\%，为区间最低价的 1.19 倍。 月度收益的年化波动率 35.4\%，低于全样本中位数 37.3\%，在三级行业“生猪养殖”的 11 家公司中波动率由高到低排第\zhnumber{8}位。

\begin{center}
\begin{minipage}{\textwidth}\centering
\begin{tikzpicture}
\begin{axis}[
  width=12.6cm, height=3.9cm, scale only axis,
  xmin=2020.922, xmax=2026.972, ymin=4.01, ymax=11.77,
  xtick={2021,2022,2023,2024,2025,2026},
  yearticks,
  yticklabel style={font=\scriptsize},
  ylabel={元/股}, ylabel style={font=\scriptsize},
  grid=major, grid style={LightGray, line width=0.3pt},
  axis line style={InkGray, line width=0.5pt},
  every axis plot/.append style={line width=0.9pt},
]
\addplot[AgriGreen] table[x=x,y=y]{data/clean/profiles/px_600975.dat};
\addplot[SoftRed, dashed, line width=0.7pt] table[x=x,y=y]{data/clean/profiles/pxzz_600975.dat};
\addplot[only marks, mark=*, mark size=1.1pt, SoftRed] table[x=x,y=y]{data/clean/profiles/pxzz_600975.dat};
\end{axis}
\end{tikzpicture}

\captionof{figure}[新五丰（600975）股价与周期骨架]{新五丰（600975）月度收盘价（元/股）与 ZigZag 周期骨架（阈值 25\%，圆点为周期转折点）}
\end{minipage}
\end{center}

\pfl{过去周期分析} 以 25\% 的 ZigZag 阈值划分，样本区间内股价形成 2 个主升段与 3 个主跌段。 最大主升段出现在 2021 年 7 月 至 2022 年 3 月，8 个月+70.7\%；最大主跌段为 2023 年 4 月 至 2026 年 6 月，38 个月-60.8\%。 最近一次转折出现在 2026 年 6 月（下跌段自 2023 年 4 月起，38 个月-60.8\%），当前处于该段的延续之中。 公司股价与申万农林牧渔指数几乎同步见底，个股并未走出独立行情。 生猪养殖的定价主线是产能去化与成本曲线的赛跑（第 \ref{sec:hog} 章）：2026 年三季度行业仍处在“产能去化未完成、价格磨底”的阶段，公司 2026 年 6 月末资产负债率 84.2\%，高于全样本中位数 52.2\%，其股价因此更多反映资金面与信用风险，而不是价格弹性本身。 综合区间形态看，该股单段涨跌幅度偏大、以趋势段为主（最大主升 +70.7\%、最大主跌 -60.8\%），年化波动率 35.4\% 与全样本中位数 37.3\% 相比波动与样本中位数接近。

\pfl{盈利情况} 2026 年上半年营业收入 28.33 亿元（同比 -23.2\%），归母净利润 -7.36 亿元，由上年同期的盈利转为亏损。 以年报为趋势对照，2023---2025 年营业收入由 56.32 亿元变为 72.58 亿元（两年年均复合增速 +13.5\%），归母净利润由 -12.03 亿元变为 -8.76 亿元。 按半年报口径，2026 年上半年的 ROE 为 -41.0\%（未年化）、销售净利率 -26.5\%、毛利率 -5.4\%，ROE 在“生猪养殖”11 家公司中按数值由高到低排第\zhnumber{10}位。 按同一口径，三级行业“生猪养殖”11 家公司的半年 ROE 中位数为 -14.32\%，该公司低于其中位数 26.68 个百分点。 综合看，公司仍处亏损区间、由盈转亏，半年 ROE 在三级行业内按由高到低排第\zhnumber{10}位，单位盈利能力的修复依赖行业景气与自身成本端的改善，报表弹性主要体现为价格与销量的方向性变化。

\begin{center}\small\setlength{\tabcolsep}{3.4pt}
\captionof{table}[新五丰（600975）关键财务指标]{新五丰（600975）关键财务指标（合并报表口径；两个半年报期均未年化；现金短债比 $=$ 货币资金 $\div$ 短期债务）}
\begin{adjustbox}{max width=\textwidth}
\begin{tabular}{@{}lrrrrr@{}}
\toprule
指标 & 2023 年 & 2024 年 & 2025 年 & 2025 年半年报 & 2026 年半年报 \\
\midrule
营业收入（亿元） & 56.32 & 70.45 & 72.58 & 36.88 & 28.33 \\
归母净利润（亿元） & -12.03 & 0.39 & -8.76 & 0.45 & -7.36 \\
毛利率（\%） & -9.2 & 7.4 & 1.0 & 9.5 & -5.4 \\
ROE（\%） & -42.4 & 1.3 & -33.9 & 1.5 & -41.0 \\
资产负债率（\%） & 75.0 & 74.0 & 78.3 & 72.5 & 84.2 \\
经营活动现金流净额（亿元） & -6.42 & 8.57 & 5.98 & 6.06 & -3.46 \\
货币资金（亿元） & 15.00 & 11.70 & 7.80 & 11.38 & 6.08 \\
有息负债合计（亿元） & 83.49 & 79.19 & 70.34 & 73.86 & 73.49 \\
现金短债比（倍） & 0.64 & 0.66 & 0.43 & 0.76 & 0.50 \\
\bottomrule
\end{tabular}
\end{adjustbox}
\end{center}

\begin{center}
\begin{minipage}{\textwidth}\centering
\begin{tikzpicture}
\begin{groupplot}[
  group style={group size=2 by 1, horizontal sep=0.60cm},
  width=6.85cm, height=4.60cm,
  tick label style={font=\tiny},
  title style={font=\scriptsize\heiti},
  yticklabel style={font=\tiny},
  ylabel style={font=\tiny},
  grid=major, grid style={LightGray, line width=0.3pt},
  axis line style={InkGray, line width=0.5pt},
  enlarge x limits={abs=0.8},
  legend style={font=\scriptsize, at={(0.5,-0.20)}, anchor=north,
                legend columns=2, draw=none, fill=none, column sep=6pt},
]
\nextgroupplot[
  ybar, bar width=4pt,
  ymin=-20.50, ymax=82.73,
  xtick={0,1,2,3,4,5.7},
  xticklabels={2021,2022,2023,2024,2025,2026H1},
  ylabel={亿元},
]
\addplot[fill=AgriGreen!75, draw=AgriGreen, bar shift=-2.6pt] table[x=i, y=rev]{data/clean/profiles/fin_600975.dat};
\addplot[fill=SoftRed!60, draw=SoftRed, bar shift=2.6pt] table[x=i, y=ni]{data/clean/profiles/fin_600975.dat};
\addplot[fill=AgriGreen!30, draw=AgriGreen, pattern=north east lines, bar shift=-2.6pt] table[x=x, y=rev]{data/clean/profiles/finh_600975.dat};
\addplot[fill=SoftRed!20, draw=SoftRed, pattern=north east lines, bar shift=2.6pt] table[x=x, y=ni]{data/clean/profiles/finh_600975.dat};
\legend{营业收入（年/半年）, 归母净利润（年/半年）}
\nextgroupplot[
  ymin=-52.5, ymax=96.8,
  xtick={0,1,2,3,4,5.7},
  xticklabels={2021,2022,2023,2024,2025,2026H1},
  ylabel={\%},
]
\addplot[AgriGold, mark=*, mark size=1.3pt] table[x=i, y=roe]{data/clean/profiles/fin_600975.dat};
\addplot[OilBlue, mark=square*, mark size=1.3pt] table[x=i, y=debt]{data/clean/profiles/fin_600975.dat};
\addplot[only marks, AgriGold, mark=*, mark size=2.0pt] table[x=x, y=roe]{data/clean/profiles/finh_600975.dat};
\addplot[only marks, OilBlue, mark=square*, mark size=2.0pt] table[x=x, y=debt]{data/clean/profiles/finh_600975.dat};
\legend{ROE, 资产负债率}
\end{groupplot}
\end{tikzpicture}
\begin{tikzpicture}
\begin{axis}[
  name=finbot,
  width=13.75cm, height=3.10cm, scale only axis,
  ymin=-8.35, ymax=93.51,
  xtick={0,1,2,3,4,5.7},
  xticklabels={2021,2022,2023,2024,2025,2026H1},
  tick label style={font=\tiny},
  yticklabel style={font=\tiny},
  ylabel={亿元}, ylabel style={font=\tiny},
  ybar, bar width=4pt,
  grid=major, grid style={LightGray, line width=0.3pt},
  axis line style={InkGray, line width=0.5pt},
  enlarge x limits={abs=0.8},
  legend style={font=\scriptsize, at={(0.5,-0.26)}, anchor=north,
                legend columns=2, draw=none, fill=none, column sep=10pt},
]
\addplot[fill=AgriGreen!55, draw=AgriGreen, bar shift=-3.0pt] table[x=i, y=cash]{data/clean/profiles/fin_600975.dat};
\addplot[fill=SoftRed!45, draw=SoftRed, bar shift=3.0pt] table[x=i, y=ibd]{data/clean/profiles/fin_600975.dat};
\addplot[fill=AgriGreen!25, draw=AgriGreen, pattern=north east lines, bar shift=-3.0pt] table[x=x, y=cash]{data/clean/profiles/finh_600975.dat};
\addplot[fill=SoftRed!20, draw=SoftRed, pattern=north east lines, bar shift=3.0pt] table[x=x, y=ibd]{data/clean/profiles/finh_600975.dat};
\legend{货币资金（左轴）, 有息负债（左轴）}
\end{axis}
\begin{axis}[
  at={(finbot.south east)}, anchor=south east,
  width=13.75cm, height=3.10cm, scale only axis,
  axis y line*=right, axis x line*=none, xtick=\empty,
  ymin=-10.32, ymax=13.06,
  tick label style={font=\tiny},
  yticklabel style={font=\tiny}, scaled y ticks=false,
  legend style={font=\scriptsize, at={(0.5,-0.44)}, anchor=north,
                legend columns=1, draw=none, fill=none},
]
\addplot[OilBlue, mark=*, mark size=2.2pt, line width=1.3pt] table[x=i, y=ocf]{data/clean/profiles/fin_600975.dat};
\addplot[only marks, OilBlue, mark=*, mark size=3.2pt] table[x=x, y=ocf]{data/clean/profiles/finh_600975.dat};
\legend{经营活动现金流净额（右轴，亿元，蓝点）}
\end{axis}
\end{tikzpicture}

\captionof{figure}[新五丰（600975）近年主要财务指标]{新五丰（600975）近年主要财务指标（左上：营业收入与归母净利润；右上：ROE 与资产负债率；下：货币资金、有息负债为左轴柱状，经营活动现金流净额为其右轴的蓝点折线。
2021---2025 年为年报口径，2026H1 为 2026 年半年报/半年末，与其前一年报之间留出间隔、以斜纹或空心点区分；半年报数据未年化）}
\end{minipage}
\end{center}

\pfl{财务分析} 2026 年 6 月末资产负债率 84.2\%，较 2025 年末上升 5.9 个百分点（样本中位数 52.2\%，所处三级行业内按由低到高排第\zhnumber{10}位）。 2026 年 6 月末货币资金 6.08 亿元、短期债务（短期借款与一年内到期非流动负债）12.10 亿元、长期债务 61.39 亿元，有息负债合计 73.49 亿元，现金短债比 0.50 倍——覆盖偏紧。 净有息负债 67.40 亿元，相当于其 2026 年上半年营业收入的 238\%。 流动比率 1.13、速动比率 0.38，短期流动性无明显缺口。 2026 年上半年经营活动现金流净额 -3.46 亿元，经营现金流与利润同时为负，日常运营对债务与股东投入的依赖度较高。 综合看，现金短债比 0.50 倍，短期资金安排偏紧；资产负债率高于样本中位数 32.0 个百分点；有息负债中长期债务占比更高，期限结构对短债滚动的压力较小。

\pfl{重大战略转型} 公司的战略动作可归为以下几类：并购与重组（1 项）：2022 年，公司并购天心种业以解决现代农业集团内部同业竞争问题；融资与资本运作（1 项）：2023 年 7 月，公司完成募集配套资金 15.50 亿元，为天心种业 5 个募集资金项目提供 7.11 亿元建设资金。 公告口径上，近三年（2023 年 9 月---2026 年 9 月）公司披露重大重组与并购类公告 3 条、对外投资与转型类公告 2 条、回购与股权激励类公告 3 条，合计公告 426 条。 主营结构的口径变化已经落到报表上：2021 年年报的第一大行业为“批发和零售”（55.6\%），2025 年年报为“畜牧业”（84.3\%）；公司在这两期的分部划分方式有过调整。 综合判断，报告期内公司有 5 条与战略相关的公告落地（重大重组与并购 3 条、对外投资与转型 2 条），资本运作与主业调整之间存在直接关联。

\begin{center}\small\setlength{\tabcolsep}{4pt}
\captionof{table}[新五丰（600975）历史融资与资本运作]{新五丰（600975）历史融资与资本运作（据公司公告与公开报道整理；金额为披露值，含首次公开发行、再融资、债券、银行借款与重整投资等）}
\begin{adjustbox}{max width=\textwidth}
\begin{tabular}{@{}p{1.45cm}p{2.55cm}p{4.05cm}p{4.95cm}@{}}
\toprule
时间 & 方式 & 规模 & 用途与说明 \\
\midrule
2004-06 & IPO（上海证券交易所） & 发行 3,\allowbreak{}500 万股、\allowbreak{}发行价 9.20 元/\allowbreak{}股，\allowbreak{}预计募集资金净额 3.05 亿元 & 发行前净资产 1.46 亿元 \\
2021-10-27 & 非公开发行股票（向特定对象发行，4 名发行对象） & 募集资金总额 10.30 亿元（1.52 亿股，\allowbreak{}发行价 6.76 元/\allowbreak{}股） & 湖南新五丰存栏 4.32 万头母猪场项目（拟投入募集资金 5.11 亿元）及补充流动资金（5.19 亿元）；发行对象为现代农业集团、\allowbreak{}粮油集团、\allowbreak{}兴湘集团、\allowbreak{}湖南建工集团共 4 名 \\
2022 & 发行股份及支付现金购买资产（重大资产重组） & 交易作价合计 22.08 亿元 & 收购天心种业 100\% 股权（含 200.00 万元国有独享资本公积）、\allowbreak{}沅江天心 48.20\% 股权、\allowbreak{}衡东天心 39.00\% 股权、\allowbreak{}荆州湘牧 49.00\% 股权、\allowbreak{}临湘天心 46.70\% 股权、\allowbreak{}湖南天翰牧业 100\% 股权、\allowbreak{}郴州市下思田农业科技 100\% 股权、\allowbreak{}新化县久阳农业开发 100\% 股权、\allowbreak{}衡东鑫邦牧业 100\% 股权等；重组前现代农业集团直接及通过粮油集团间接合计持有 36.59\% \\
2023-07 & 发行股份募集配套资金 & 募集资金总额 15.51 亿元（1.85 亿股，\allowbreak{}发行价 8.39 元/\allowbreak{}股） & 部分用于支付本次重组现金对价和补充流动资金，\allowbreak{}天心种业 5 个募集资金项目获 7.11 亿元建设资金；发行对象 20 名 \\
\bottomrule
\end{tabular}
\end{adjustbox}
\end{center}

\pfl{历史融投资情况} 从现金流量表看，2025 年公司取得借款收到的现金 10.18 亿元、偿还债务支付的现金 15.96 亿元、吸收权益性投资 0.15 亿元，筹资活动现金流净额 -8.67 亿元（2023 年为取得借款 29.57 亿元、偿还债务 20.80 亿元）。 2026 年上半年筹资活动现金流净额为 2.15 亿元（取得借款 15.06 亿元、偿还债务 7.90 亿元，未年化），仍处于净融入状态。 2025 年分配股利、利润或偿付利息支付的现金合计 0.54 亿元。 公开披露的融资记录共 4 笔，工具类型以 IPO、定向增发、发行股份及支付现金购、发行股份募集配套资金为主；金额与用途见下表。 其中规模最大的两笔为：2004 年 6 月，IPO，发行 3,500 万股、发行价 9.20 元/股，预计募集资金净额 3.05 亿元；2022 年，发行股份及支付现金购，交易作价合计 22.08 亿元。 股本方面，近三年披露股本变动 6 次（增发新股上市 2 次、限售股份上市 1 次）；总股本由 2021 年末的 8.05 亿股变为 12.61 亿股（+56.67\%）。 分红方面，近三年未实施现金分红，在全部样本中属于分红能力偏弱的一端。 近三年“再融资”类公告 4 条，涉及定向增发、发行股份购买资产。 综合看，公司融资结构上以债务滚动为主、股权融资为补充；2025 年筹资活动现金流净额 -8.67 亿元，融资节奏仍以维持存量债务周转为主，与前述经营与投资现金流的方向基本一致。

\pfl{未来潜在融资需求分析} 2026 年 6 月末货币资金 6.08 亿元，而短期债务 12.10 亿元（现金短债比 0.50 倍），2026 年上半年经营现金流净额 -3.46 亿元。按“短期债务减去账面货币资金”的滚动偿付口径，静态缺口约 6.02 亿元；若再计入上半年亏损的年化额（14.72 亿元），维持一年正常经营所需的外部资金约 6.02---20.74 亿元。 按第 \ref{sec:financing-need} 节的三档划分，公司属于第一档“补血型”：2026 年上半年归母净利润 -7.36 亿元、6 月末资产负债率 84.2\%。负债率已超过 80\%，净资产被亏损侵蚀的程度较深，属于全部样本中资产负债表最紧张的一组。 可行的融资路径按现实性排序为：出售非核心资产与参股股权、债务重组或展期（与银行和债券持有人协商）、定向增发（需股价配合，且控股股东需放弃优先认购或同步增资）。 公开披露的后续资金安排线索包括：资本开支安排：2025 年 8 月，第六届董事会第十八次会议审议原“湖南新五丰存栏 4.32 万头母猪场项目”相关事项。 综合看，公司的外部资金需求以债务接续为主，滚动偿付口径下的静态缺口约 6.02 亿元，融资的时点与规模更多取决于信用条件与资本市场窗口。




\subsubsection{罗牛山（000735）}
\label{co:000735}

\pfl{公司基本情况} 公司前身可追溯至 1987 年，海南海口的老牌农业上市公司，1997 年深交所上市的中国首家“菜篮子”股份制上市企业，主营生猪养殖与屠宰、冷链仓储物流、房地产开发销售及职业教育服务四大板块；2026 年上半年出栏生猪 45.62 万头、营收逆势增长，但受猪价深跌拖累归母净亏损 3.10 亿元。 公司于 1997-06-11 在深交所主板首发上市，注册资本 11.52 亿元，总股本 11.52 亿股。 截至 2026 年 10 月 12 日收盘价 5.59 元，对应总市值 64.4 亿元，近一年-13.9\%，总市值在三级行业“生猪养殖”的 11 家公司中按由大到小排第\zhnumber{8}位，近一年涨跌幅低于全样本中位数（-10.4\%）3.5 个百分点。 按 2026 年半年报披露的产品构成，收入占比前三位为生猪 39.3\%；其他 16.5\%；饲料 15.2\%。 与 2021 年年报相比，第一大产品口径由“肉制品”（38.6\%）变为“生猪”（39.3\%）；两期披露的分类口径有调整，占比差异中同时包含分类因素。 股权结构方面，控股股东为罗牛山集团有限公司（持股 17.14\%）。

\pfl{股权结构与控制权} 按 2026 年 6 月 30 日的股东名册，前三大股东为罗牛山集团有限公司 持股 17.14\%；海口永盛畜牧机械工程有限公司 持股 2.40\%；柴长茂 持股 1.81\%，股东户数 96,831 户。 股权结构的另一条线索是，个人持股 10.36\%，同时担任公司董事长兼总裁、战略委员会主任委员。

\pfl{经营要点} 经营层面，2026 年上半年，外三元猪出栏同比增长 53.22\% 至 42.16 万头、PSY 超过 28 头。

\pfl{所处行业的基本面} 生猪的供给由能繁母猪存栏决定、需求由猪肉消费与替代品价格决定，本轮下行的成因在供给端的产能去化不彻底，而非需求收缩（第 \ref{sec:framework} 章的供需框架）。 与该公司直接相关的行业数据包括：2026 年上半年，生猪板块毛利率由上年的 8.93\% 骤降至-34.05\%；截至 2026 年 6 月底全国能繁母猪存栏 3780 万头。

\pfl{竞争优势} 把该公司放回同类公司比较，2026 年上半年营业收入 12.4 亿元、在三级行业“生猪养殖”的 11 家公司中排第\zhnumber{9}位；毛利率 0.7\%、高于行业内中位数 2.2 个百分点；半年 ROE -7.4\%（未年化）、高于行业内中位数（-14.32\%）；资产负债率 52.3\%、低于行业内中位数（58.9\%）。 公司在公开资料中自述与同业比较的相对优势在于：2026 年上半年屠宰总量同比增长 65.76\%；设定 2026 年末黑猪料肉比降至 3.5 以下、完全成本降至 18 元/公斤以内的量化目标，并推进黑猪专用屠宰线建设与深加工产品开发。 优势在当期体现为高于行业中位数的资本回报，但优势能否持续仍取决于周期与自身产能投放节奏。

\pfl{历史股价} 本报告样本区间（2021 年 1 月 至 2026 年 10 月）内，股价由 8.18 元变为 5.59 元，累计 -31.7\%，区间最高 9.04 元（2025 年 12 月）、最低 4.33 元（2024 年 6 月）；当前价相当于区间最高价的 61.8\%，为区间最低价的 1.29 倍。 月度收益的年化波动率 27.2\%，低于全样本中位数 37.3\%，在三级行业“生猪养殖”的 11 家公司中波动率由高到低排第\zhnumber{11}位。

\begin{center}
\begin{minipage}{\textwidth}\centering
\begin{tikzpicture}
\begin{axis}[
  width=12.6cm, height=3.9cm, scale only axis,
  xmin=2020.922, xmax=2026.972, ymin=4.03, ymax=9.67,
  xtick={2021,2022,2023,2024,2025,2026},
  yearticks,
  yticklabel style={font=\scriptsize},
  ylabel={元/股}, ylabel style={font=\scriptsize},
  grid=major, grid style={LightGray, line width=0.3pt},
  axis line style={InkGray, line width=0.5pt},
  every axis plot/.append style={line width=0.9pt},
]
\addplot[AgriGreen] table[x=x,y=y]{data/clean/profiles/px_000735.dat};
\addplot[SoftRed, dashed, line width=0.7pt] table[x=x,y=y]{data/clean/profiles/pxzz_000735.dat};
\addplot[only marks, mark=*, mark size=1.1pt, SoftRed] table[x=x,y=y]{data/clean/profiles/pxzz_000735.dat};
\end{axis}
\end{tikzpicture}

\captionof{figure}[罗牛山（000735）股价与周期骨架]{罗牛山（000735）月度收盘价（元/股）与 ZigZag 周期骨架（阈值 25\%，圆点为周期转折点）}
\end{minipage}
\end{center}

\pfl{过去周期分析} 以 25\% 的 ZigZag 阈值划分，样本区间内股价形成 3 个主升段与 3 个主跌段。 最大主升段出现在 2024 年 6 月 至 2025 年 12 月，18 个月+108.8\%；最大主跌段为 2025 年 12 月 至 2026 年 6 月，6 个月-47.5\%。 最近一次转折出现在 2026 年 6 月（下跌段自 2025 年 12 月起，6 个月-47.5\%），当前处于该段的延续之中。 公司股价较申万农林牧渔指数提前约 23 个月见底，是板块中较早定价周期反转的公司之一。 生猪养殖的定价主线是产能去化与成本曲线的赛跑（第 \ref{sec:hog} 章）：2026 年三季度行业仍处在“产能去化未完成、价格磨底”的阶段，公司 2026 年 6 月末资产负债率 52.3\%、半年 ROE -7.4\%（未年化），在本轮下行中属于随行业同步调整的一类。 综合区间形态看，该股单段涨跌幅度偏大、以趋势段为主（最大主升 +108.8\%、最大主跌 -47.5\%），年化波动率 27.2\% 与全样本中位数 37.3\% 相比波动小于样本中位数。

\pfl{盈利情况} 2026 年上半年营业收入 12.41 亿元（同比 +14.5\%），归母净利润 -3.10 亿元，亏损同比扩大 3.03 亿元。 以年报为趋势对照，2023---2025 年营业收入由 42.42 亿元变为 33.77 亿元（两年年均复合增速 -10.8\%），归母净利润由 -4.88 亿元变为 1.54 亿元。 按半年报口径，2026 年上半年的 ROE 为 -7.4\%（未年化）、销售净利率 -26.2\%、毛利率 0.7\%，ROE 在“生猪养殖”11 家公司中按数值由高到低排第\zhnumber{2}位。 同期全样本半年 ROE 中位数 0.75\%，公司与之相差 8.19 个百分点（弱于中位数公司）。 综合看，公司仍处亏损区间、亏损同比扩大，半年 ROE 在三级行业内按由高到低排第\zhnumber{2}位，单位盈利能力的修复依赖行业景气与自身成本端的改善，报表弹性主要体现为价格与销量的方向性变化。

\begin{center}\small\setlength{\tabcolsep}{3.4pt}
\captionof{table}[罗牛山（000735）关键财务指标]{罗牛山（000735）关键财务指标（合并报表口径；两个半年报期均未年化；现金短债比 $=$ 货币资金 $\div$ 短期债务）}
\begin{adjustbox}{max width=\textwidth}
\begin{tabular}{@{}lrrrrr@{}}
\toprule
指标 & 2023 年 & 2024 年 & 2025 年 & 2025 年半年报 & 2026 年半年报 \\
\midrule
营业收入（亿元） & 42.42 & 33.32 & 33.77 & 10.84 & 12.41 \\
归母净利润（亿元） & -4.88 & 2.42 & 1.54 & -0.07 & -3.10 \\
毛利率（\%） & 16.1 & 23.8 & 25.6 & 13.7 & 0.7 \\
ROE（\%） & -11.2 & 5.6 & 3.6 & -0.1 & -7.4 \\
资产负债率（\%） & 63.7 & 54.8 & 50.3 & 54.0 & 52.3 \\
经营活动现金流净额（亿元） & 3.75 & 3.71 & -0.93 & -1.27 & -1.38 \\
货币资金（亿元） & 9.53 & 10.95 & 7.59 & 4.24 & 4.11 \\
有息负债合计（亿元） & 29.80 & 16.62 & 17.41 & 14.57 & 19.18 \\
现金短债比（倍） & 1.30 & 1.95 & 0.93 & 0.88 & 0.41 \\
\bottomrule
\end{tabular}
\end{adjustbox}
\end{center}

\begin{center}
\begin{minipage}{\textwidth}\centering
\begin{tikzpicture}
\begin{groupplot}[
  group style={group size=2 by 1, horizontal sep=0.60cm},
  width=6.85cm, height=4.60cm,
  tick label style={font=\tiny},
  title style={font=\scriptsize\heiti},
  yticklabel style={font=\tiny},
  ylabel style={font=\tiny},
  grid=major, grid style={LightGray, line width=0.3pt},
  axis line style={InkGray, line width=0.5pt},
  enlarge x limits={abs=0.8},
  legend style={font=\scriptsize, at={(0.5,-0.20)}, anchor=north,
                legend columns=2, draw=none, fill=none, column sep=6pt},
]
\nextgroupplot[
  ybar, bar width=4pt,
  ymin=-9.61, ymax=48.09,
  xtick={0,1,2,3,4,5.7},
  xticklabels={2021,2022,2023,2024,2025,2026H1},
  ylabel={亿元},
]
\addplot[fill=AgriGreen!75, draw=AgriGreen, bar shift=-2.6pt] table[x=i, y=rev]{data/clean/profiles/fin_000735.dat};
\addplot[fill=SoftRed!60, draw=SoftRed, bar shift=2.6pt] table[x=i, y=ni]{data/clean/profiles/fin_000735.dat};
\addplot[fill=AgriGreen!30, draw=AgriGreen, pattern=north east lines, bar shift=-2.6pt] table[x=x, y=rev]{data/clean/profiles/finh_000735.dat};
\addplot[fill=SoftRed!20, draw=SoftRed, pattern=north east lines, bar shift=2.6pt] table[x=x, y=ni]{data/clean/profiles/finh_000735.dat};
\legend{营业收入（年/半年）, 归母净利润（年/半年）}
\nextgroupplot[
  ymin=-17.2, ymax=71.2,
  xtick={0,1,2,3,4,5.7},
  xticklabels={2021,2022,2023,2024,2025,2026H1},
  ylabel={\%},
]
\addplot[AgriGold, mark=*, mark size=1.3pt] table[x=i, y=roe]{data/clean/profiles/fin_000735.dat};
\addplot[OilBlue, mark=square*, mark size=1.3pt] table[x=i, y=debt]{data/clean/profiles/fin_000735.dat};
\addplot[only marks, AgriGold, mark=*, mark size=2.0pt] table[x=x, y=roe]{data/clean/profiles/finh_000735.dat};
\addplot[only marks, OilBlue, mark=square*, mark size=2.0pt] table[x=x, y=debt]{data/clean/profiles/finh_000735.dat};
\legend{ROE, 资产负债率}
\end{groupplot}
\end{tikzpicture}
\begin{tikzpicture}
\begin{axis}[
  name=finbot,
  width=13.75cm, height=3.10cm, scale only axis,
  ymin=-2.98, ymax=33.37,
  xtick={0,1,2,3,4,5.7},
  xticklabels={2021,2022,2023,2024,2025,2026H1},
  tick label style={font=\tiny},
  yticklabel style={font=\tiny},
  ylabel={亿元}, ylabel style={font=\tiny},
  ybar, bar width=4pt,
  grid=major, grid style={LightGray, line width=0.3pt},
  axis line style={InkGray, line width=0.5pt},
  enlarge x limits={abs=0.8},
  legend style={font=\scriptsize, at={(0.5,-0.26)}, anchor=north,
                legend columns=2, draw=none, fill=none, column sep=10pt},
]
\addplot[fill=AgriGreen!55, draw=AgriGreen, bar shift=-3.0pt] table[x=i, y=cash]{data/clean/profiles/fin_000735.dat};
\addplot[fill=SoftRed!45, draw=SoftRed, bar shift=3.0pt] table[x=i, y=ibd]{data/clean/profiles/fin_000735.dat};
\addplot[fill=AgriGreen!25, draw=AgriGreen, pattern=north east lines, bar shift=-3.0pt] table[x=x, y=cash]{data/clean/profiles/finh_000735.dat};
\addplot[fill=SoftRed!20, draw=SoftRed, pattern=north east lines, bar shift=3.0pt] table[x=x, y=ibd]{data/clean/profiles/finh_000735.dat};
\legend{货币资金（左轴）, 有息负债（左轴）}
\end{axis}
\begin{axis}[
  at={(finbot.south east)}, anchor=south east,
  width=13.75cm, height=3.10cm, scale only axis,
  axis y line*=right, axis x line*=none, xtick=\empty,
  ymin=-4.82, ymax=15.83,
  tick label style={font=\tiny},
  yticklabel style={font=\tiny}, scaled y ticks=false,
  legend style={font=\scriptsize, at={(0.5,-0.44)}, anchor=north,
                legend columns=1, draw=none, fill=none},
]
\addplot[OilBlue, mark=*, mark size=2.2pt, line width=1.3pt] table[x=i, y=ocf]{data/clean/profiles/fin_000735.dat};
\addplot[only marks, OilBlue, mark=*, mark size=3.2pt] table[x=x, y=ocf]{data/clean/profiles/finh_000735.dat};
\legend{经营活动现金流净额（右轴，亿元，蓝点）}
\end{axis}
\end{tikzpicture}

\captionof{figure}[罗牛山（000735）近年主要财务指标]{罗牛山（000735）近年主要财务指标（左上：营业收入与归母净利润；右上：ROE 与资产负债率；下：货币资金、有息负债为左轴柱状，经营活动现金流净额为其右轴的蓝点折线。
2021---2025 年为年报口径，2026H1 为 2026 年半年报/半年末，与其前一年报之间留出间隔、以斜纹或空心点区分；半年报数据未年化）}
\end{minipage}
\end{center}

\pfl{财务分析} 2026 年 6 月末资产负债率 52.3\%，较 2025 年末上升 2.0 个百分点（样本中位数 52.2\%，所处三级行业内按由低到高排第\zhnumber{2}位）。 2026 年 6 月末货币资金 4.11 亿元、短期债务（短期借款与一年内到期非流动负债）9.97 亿元、长期债务 9.21 亿元，有息负债合计 19.18 亿元，现金短债比 0.41 倍——覆盖偏紧。 净有息负债 15.07 亿元，相当于其 2026 年上半年营业收入的 121\%。 流动比率 1.00、速动比率 0.29，短期流动性无明显缺口。 2026 年上半年经营活动现金流净额 -1.38 亿元，经营现金流与利润同时为负，现金消耗较快。 2026 年 6 月末在建工程 1.69 亿元，较上年末增加 0.06 亿元，仍有在建投入。 综合看，账面货币资金对有息负债与短期债务的覆盖不足，滚动偿付对新增融资的依赖度较高；资产负债率高于样本中位数 0.1 个百分点；有息负债中短期债务占比更高，期限结构仍以滚动续作为主。

\pfl{重大战略转型} 近三年公司的战略动作集中在：融资与资本运作（1 项）：2023 年 4 月，产能从 50 万头发展到近 130 万头。 公告口径上，近三年（2023 年 9 月---2026 年 9 月）公司披露重大重组与并购类公告 4 条、对外投资与转型类公告 2 条、回购与股权激励类公告 0 条，合计公告 299 条。 主营结构的口径变化已经落到报表上：2021 年年报的第一大产品为“肉制品”（38.6\%），2026 年半年报为“生猪”（39.3\%）；公司在这两期的分部划分方式有过调整。 综合判断，报告期内公司有 6 条与战略相关的公告落地（重大重组与并购 4 条、对外投资与转型 2 条），资本运作与主业调整之间存在直接关联。

\begin{center}\small\setlength{\tabcolsep}{4pt}
\captionof{table}[罗牛山（000735）历史融资与资本运作]{罗牛山（000735）历史融资与资本运作（据公司公告与公开报道整理；金额为披露值，含首次公开发行、再融资、债券、银行借款与重整投资等）}
\begin{adjustbox}{max width=\textwidth}
\begin{tabular}{@{}p{1.45cm}p{2.55cm}p{4.05cm}p{4.95cm}@{}}
\toprule
时间 & 方式 & 规模 & 用途与说明 \\
\midrule
1997-05 & IPO（深圳证券交易所主板） & 实际募集资金净额 1.87 亿元 & 发行价 4.35 元/\allowbreak{}股，\allowbreak{}上市日期 1997-06-11 \\
2002-01 & 增发 A 股 & 实际募集资金净额 4.79 亿元 & 募投项目包括海南爱华汽车广场项目（承诺投入 5,\allowbreak{}000.00 万元）等；经中国证监会证监发行字[2002]107 号文核准 \\
2016-04-13 & 非公开发行股票（向 6 名特定对象，均以现金认购） & 募集资金总额 16.50 亿元，\allowbreak{}净额 16.31 亿元（2.71 亿股 & “罗牛山十万头现代化猪场项目”“海南农副产品交易配送中心及产业配套项目”和偿还银行贷款；控股股东罗牛山集团有限公司认购部分限售 36 个月 \\
2020 & 定增（终止） & --- & 原计划扩大生猪养殖规模；2020 年公司为扩大生猪养殖规模启动定增项目 \\
2023-11-24 & 定增（终止） & --- & 公司第十届董事会第十次临时会议及第十届监事会第八次会议审议通过终止公司 2023 年度向特定对象发行 A 股股票事项 \\
\bottomrule
\end{tabular}
\end{adjustbox}
\end{center}

\pfl{历史融投资情况} 从现金流量表看，2025 年公司取得借款收到的现金 7.40 亿元、偿还债务支付的现金 7.05 亿元，筹资活动现金流净额 -1.49 亿元（2023 年为取得借款 8.72 亿元、偿还债务 5.90 亿元）。 2026 年上半年筹资活动现金流净额为 1.22 亿元（取得借款 5.67 亿元、偿还债务 3.79 亿元，未年化），仍处于净融入状态。 2025 年分配股利、利润或偿付利息支付的现金合计 0.74 亿元。 公开披露的融资记录共 5 笔，工具类型以 IPO、定向增发为主；金额与用途见下表。 其中规模最大的两笔为：2002 年 1 月，定向增发，实际募集资金净额 4.79 亿元；1997 年 5 月，IPO，实际募集资金净额 1.87 亿元。 股本方面，近三年披露股本变动 3 次；总股本自 2021 年末以来未发生变化（11.52 亿股）。 分红方面，近三年现金分红 4 次、合计每股派现 0.10 元，最近一次实施在 2025 年 12 月。 近三年“再融资”类公告 1 条，涉及定向增发。 综合看，公司融资结构上股权与债务融资并举；2025 年筹资活动现金流净额 -1.49 亿元，融资节奏仍以维持存量债务周转为主，与前述经营与投资现金流的方向基本一致。

\pfl{未来潜在融资需求分析} 2026 年 6 月末货币资金 4.11 亿元，而短期债务 9.97 亿元（现金短债比 0.41 倍），2026 年上半年经营现金流净额 -1.38 亿元。按“短期债务减去账面货币资金”的滚动偿付口径，静态缺口约 5.86 亿元；若再计入上半年亏损的年化额（6.19 亿元），维持一年正常经营所需的外部资金约 5.86---12.05 亿元。 公司 2026 年上半年归母净利润为负（-3.10 亿元），但 6 月末资产负债率 52.3\% 尚未进入第一档（亏损且负债率 $>$65\%）的区间，当前融资需求以补充流动资金、支撑低谷期经营性支出为主。 公开披露的后续资金安排线索包括：股权融资线索：2023 年 11 月，公司终止 2023 年度向特定对象发行 A 股股票事项，原扩张产能的股权融资渠道关闭。 资金需求还要叠加在建工程：期末在建工程 1.69 亿元、2026 年上半年购建固定资产等支出 0.68 亿元（未年化），这部分支出在周期底部通常会被主动放缓。 综合看，公司的外部资金需求以补充营运资金为主、项目投入为辅，滚动偿付口径下的静态缺口约 5.86 亿元，融资的时点与规模更多取决于信用条件与资本市场窗口。




\subsubsection{仙坛股份（002746）}
\label{co:002746}

\pfl{公司基本情况} 山东烟台白羽肉鸡养殖屠宰一体化龙头企业，2001 年成立、2015 年 2 月深交所上市，商品代肉鸡以“公司+农户”委托养殖为主，2020 年定增逾 10 亿元投建诸城年产 1.2 亿羽肉鸡产业生态项目；2026 年上半年营业收入 29.95 亿元、归母净利润 1.80 亿元，同比增长约三成。 公司于 2015-02-16 在深交所主板首发上市，注册资本 8.61 亿元，总股本 8.61 亿股。 截至 2026 年 10 月 12 日收盘价 6.23 元，对应总市值 53.6 亿元，近一年+0.8\%，总市值在三级行业“肉鸡养殖”的 8 家公司中按由大到小排第\zhnumber{4}位，近一年涨跌幅高于全样本中位数（-10.4\%）11.2 个百分点。 按 2026 年半年报披露的行业构成，收入几乎全部来自肉鸡养殖及屠宰加工 100.0\%。

\pfl{股权结构与控制权} 按 2026 年 6 月 30 日的股东名册，前三大股东为王寿纯 持股 23.37\%；曲立荣 持股 21.34\%；全国社保基金四一三组合 持股 1.86\%，股东户数 26,636 户。 股权结构的另一条线索是，控股股东、实际控制人王寿纯持股 2.01 亿股、占总股本 23.37\%，其中限售 1.51 亿股。

\pfl{经营要点} 经营与产能方面，2026 年上半年诸城项目商品代肉鸡出栏量 4,114.91 万只。

\pfl{所处行业的基本面} 白羽与黄羽肉鸡的供给取决于祖代引种量与父母代存栏，需求端受餐饮与替代效应驱动，产能周期只有 6---12 个月（第 \ref{sec:framework} 章的供需框架）。 与该公司直接相关的行业数据包括：2026 年上半年，全国白羽鸡苗均价 2.70 元/羽、同比涨幅 12.50\%。

\pfl{竞争优势} 把该公司放回同类公司比较，2026 年上半年营业收入 29.9 亿元、在三级行业“肉鸡养殖”的 8 家公司中排第\zhnumber{4}位；毛利率 8.7\%、低于行业内中位数 2.1 个百分点；半年 ROE 3.9\%（未年化）、高于行业内中位数（1.21\%）；资产负债率 43.0\%、低于行业内中位数（52.1\%）。 公司区别于同类公司的经营特征包括：调理品一、二期项目产能持续释放，2026 年上半年调理品销量同比增长 86.02\%。 优势在当期体现为高于行业中位数的资本回报，但优势能否持续仍取决于周期与自身产能投放节奏。

\pfl{历史股价} 以 2021 年 1 月为起点、2026 年 10 月为终点，股价由 10.04 元变为 6.23 元，累计 -37.9\%，区间最高 10.98 元（2021 年 3 月）、最低 5.10 元（2026 年 6 月）；当前价相当于区间最高价的 56.7\%，为区间最低价的 1.22 倍。 月度收益的年化波动率 37.3\%，高于全样本中位数 37.3\%，在三级行业“肉鸡养殖”的 8 家公司中波动率由高到低排第\zhnumber{4}位。 把股价阶段与公司事件对齐看，2021 年 12 月 至 2026 年 6 月股价下跌-51.4\%，同期（2026 年 5 月）公司整体毛利率较 2019 年下降达 23.4 个百分点。 整体上看，区间的几处大级别拐点都能在公司的资本运作、股权变动或风险事件上找到对应，股价波动有明显的公司层面驱动，而非单纯的行业 beta。

\begin{center}
\begin{minipage}{\textwidth}\centering
\begin{tikzpicture}
\begin{axis}[
  width=12.6cm, height=3.9cm, scale only axis,
  xmin=2020.922, xmax=2026.972, ymin=4.74, ymax=11.75,
  xtick={2021,2022,2023,2024,2025,2026},
  yearticks,
  yticklabel style={font=\scriptsize},
  ylabel={元/股}, ylabel style={font=\scriptsize},
  grid=major, grid style={LightGray, line width=0.3pt},
  axis line style={InkGray, line width=0.5pt},
  every axis plot/.append style={line width=0.9pt},
]
\addplot[AgriGreen] table[x=x,y=y]{data/clean/profiles/px_002746.dat};
\addplot[SoftRed, dashed, line width=0.7pt] table[x=x,y=y]{data/clean/profiles/pxzz_002746.dat};
\addplot[only marks, mark=*, mark size=1.1pt, SoftRed] table[x=x,y=y]{data/clean/profiles/pxzz_002746.dat};
\end{axis}
\end{tikzpicture}

\captionof{figure}[仙坛股份（002746）股价与周期骨架]{仙坛股份（002746）月度收盘价（元/股）与 ZigZag 周期骨架（阈值 25\%，圆点为周期转折点）}
\end{minipage}
\end{center}

\pfl{过去周期分析} 以 25\% 的 ZigZag 阈值划分，样本区间内股价形成 3 个主升段与 2 个主跌段。 最大主升段出现在 2021 年 11 月 至 2021 年 12 月，1 个月+57.7\%；最大主跌段为 2021 年 12 月 至 2026 年 6 月，54 个月-51.4\%。 最近一次转折出现在 2026 年 8 月（上涨段自 2026 年 6 月起，2 个月+28.2\%），当前处于该段之后的震荡之中。 公司股价与申万农林牧渔指数几乎同步见底，个股并未走出独立行情。 禽链的产能周期只有 6---12 个月、反转快于生猪（第 \ref{sec:other-livestock} 节），2026 年前四个月毛鸡养殖利润同比大幅回升，链条利润的修复先于生猪，公司 2026 年上半年实现盈利、半年 ROE 3.9\%（未年化），好于全样本中位数（0.75\%），下行周期对其报表的冲击小于同业。 综合区间形态看，该股单段涨跌幅度偏大、以趋势段为主（最大主升 +57.7\%、最大主跌 -51.4\%），年化波动率 37.3\% 与全样本中位数 37.3\% 相比波动与样本中位数接近。

\pfl{盈利情况} 2026 年上半年营业收入 29.95 亿元（同比 +17.9\%），归母净利润 1.80 亿元，同比 +32.0\%。 以年报为趋势对照，2023---2025 年营业收入由 57.81 亿元变为 57.72 亿元（两年年均复合增速 -0.1\%），归母净利润由 2.26 亿元变为 2.43 亿元。 按半年报口径，2026 年上半年的 ROE 为 3.9\%（未年化）、销售净利率 6.2\%、毛利率 8.7\%，ROE 在“肉鸡养殖”8 家公司中按数值由高到低排第\zhnumber{3}位。 同期全样本半年 ROE 中位数 0.75\%，公司与之相差 3.12 个百分点（略好于中位数公司）。 综合看，公司在报告期维持盈利（高于全样本半年 ROE 中位数 0.75\%），利润的可持续性取决于肉鸡养殖景气的延续与成本管控的执行。

\begin{center}\small\setlength{\tabcolsep}{3.4pt}
\captionof{table}[仙坛股份（002746）关键财务指标]{仙坛股份（002746）关键财务指标（合并报表口径；两个半年报期均未年化；现金短债比 $=$ 货币资金 $\div$ 短期债务）}
\begin{adjustbox}{max width=\textwidth}
\begin{tabular}{@{}lrrrrr@{}}
\toprule
指标 & 2023 年 & 2024 年 & 2025 年 & 2025 年半年报 & 2026 年半年报 \\
\midrule
营业收入（亿元） & 57.81 & 52.91 & 57.72 & 25.40 & 29.95 \\
归母净利润（亿元） & 2.26 & 2.06 & 2.43 & 1.37 & 1.80 \\
毛利率（\%） & 3.8 & 5.2 & 6.5 & 7.9 & 8.7 \\
ROE（\%） & 5.0 & 4.5 & 5.3 & 2.9 & 3.9 \\
资产负债率（\%） & 31.4 & 38.6 & 41.8 & 41.8 & 43.0 \\
经营活动现金流净额（亿元） & 3.30 & 1.55 & 6.79 & 3.37 & 0.60 \\
货币资金（亿元） & 25.15 & 39.01 & 36.01 & 39.16 & 36.20 \\
有息负债合计（亿元） & 10.55 & 18.60 & 23.83 & 22.08 & 25.92 \\
现金短债比（倍） & 2.77 & 2.38 & 1.61 & 1.90 & 1.43 \\
\bottomrule
\end{tabular}
\end{adjustbox}
\end{center}

\begin{center}
\begin{minipage}{\textwidth}\centering
\begin{tikzpicture}
\begin{groupplot}[
  group style={group size=2 by 1, horizontal sep=0.60cm},
  width=6.85cm, height=4.60cm,
  tick label style={font=\tiny},
  title style={font=\scriptsize\heiti},
  yticklabel style={font=\tiny},
  ylabel style={font=\tiny},
  grid=major, grid style={LightGray, line width=0.3pt},
  axis line style={InkGray, line width=0.5pt},
  enlarge x limits={abs=0.8},
  legend style={font=\scriptsize, at={(0.5,-0.20)}, anchor=north,
                legend columns=2, draw=none, fill=none, column sep=6pt},
]
\nextgroupplot[
  ybar, bar width=4pt,
  ymin=-5.78, ymax=64.75,
  xtick={0,1,2,3,4,5.7},
  xticklabels={2021,2022,2023,2024,2025,2026H1},
  ylabel={亿元},
]
\addplot[fill=AgriGreen!75, draw=AgriGreen, bar shift=-2.6pt] table[x=i, y=rev]{data/clean/profiles/fin_002746.dat};
\addplot[fill=SoftRed!60, draw=SoftRed, bar shift=2.6pt] table[x=i, y=ni]{data/clean/profiles/fin_002746.dat};
\addplot[fill=AgriGreen!30, draw=AgriGreen, pattern=north east lines, bar shift=-2.6pt] table[x=x, y=rev]{data/clean/profiles/finh_002746.dat};
\addplot[fill=SoftRed!20, draw=SoftRed, pattern=north east lines, bar shift=2.6pt] table[x=x, y=ni]{data/clean/profiles/finh_002746.dat};
\legend{营业收入（年/半年）, 归母净利润（年/半年）}
\nextgroupplot[
  ymin=-3.4, ymax=47.3,
  xtick={0,1,2,3,4,5.7},
  xticklabels={2021,2022,2023,2024,2025,2026H1},
  ylabel={\%},
]
\addplot[AgriGold, mark=*, mark size=1.3pt] table[x=i, y=roe]{data/clean/profiles/fin_002746.dat};
\addplot[OilBlue, mark=square*, mark size=1.3pt] table[x=i, y=debt]{data/clean/profiles/fin_002746.dat};
\addplot[only marks, AgriGold, mark=*, mark size=2.0pt] table[x=x, y=roe]{data/clean/profiles/finh_002746.dat};
\addplot[only marks, OilBlue, mark=square*, mark size=2.0pt] table[x=x, y=debt]{data/clean/profiles/finh_002746.dat};
\legend{ROE, 资产负债率}
\end{groupplot}
\end{tikzpicture}
\begin{tikzpicture}
\begin{axis}[
  name=finbot,
  width=13.75cm, height=3.10cm, scale only axis,
  ymin=-3.90, ymax=43.69,
  xtick={0,1,2,3,4,5.7},
  xticklabels={2021,2022,2023,2024,2025,2026H1},
  tick label style={font=\tiny},
  yticklabel style={font=\tiny},
  ylabel={亿元}, ylabel style={font=\tiny},
  ybar, bar width=4pt,
  grid=major, grid style={LightGray, line width=0.3pt},
  axis line style={InkGray, line width=0.5pt},
  enlarge x limits={abs=0.8},
  legend style={font=\scriptsize, at={(0.5,-0.26)}, anchor=north,
                legend columns=2, draw=none, fill=none, column sep=10pt},
]
\addplot[fill=AgriGreen!55, draw=AgriGreen, bar shift=-3.0pt] table[x=i, y=cash]{data/clean/profiles/fin_002746.dat};
\addplot[fill=SoftRed!45, draw=SoftRed, bar shift=3.0pt] table[x=i, y=ibd]{data/clean/profiles/fin_002746.dat};
\addplot[fill=AgriGreen!25, draw=AgriGreen, pattern=north east lines, bar shift=-3.0pt] table[x=x, y=cash]{data/clean/profiles/finh_002746.dat};
\addplot[fill=SoftRed!20, draw=SoftRed, pattern=north east lines, bar shift=3.0pt] table[x=x, y=ibd]{data/clean/profiles/finh_002746.dat};
\legend{货币资金（左轴）, 有息负债（左轴）}
\end{axis}
\begin{axis}[
  at={(finbot.south east)}, anchor=south east,
  width=13.75cm, height=3.10cm, scale only axis,
  axis y line*=right, axis x line*=none, xtick=\empty,
  ymin=-2.88, ymax=9.09,
  tick label style={font=\tiny},
  yticklabel style={font=\tiny}, scaled y ticks=false,
  legend style={font=\scriptsize, at={(0.5,-0.44)}, anchor=north,
                legend columns=1, draw=none, fill=none},
]
\addplot[OilBlue, mark=*, mark size=2.2pt, line width=1.3pt] table[x=i, y=ocf]{data/clean/profiles/fin_002746.dat};
\addplot[only marks, OilBlue, mark=*, mark size=3.2pt] table[x=x, y=ocf]{data/clean/profiles/finh_002746.dat};
\legend{经营活动现金流净额（右轴，亿元，蓝点）}
\end{axis}
\end{tikzpicture}

\captionof{figure}[仙坛股份（002746）近年主要财务指标]{仙坛股份（002746）近年主要财务指标（左上：营业收入与归母净利润；右上：ROE 与资产负债率；下：货币资金、有息负债为左轴柱状，经营活动现金流净额为其右轴的蓝点折线。
2021---2025 年为年报口径，2026H1 为 2026 年半年报/半年末，与其前一年报之间留出间隔、以斜纹或空心点区分；半年报数据未年化）}
\end{minipage}
\end{center}

\pfl{财务分析} 2026 年 6 月末资产负债率 43.0\%，较 2025 年末上升 1.2 个百分点（样本中位数 52.2\%，所处三级行业内按由低到高排第\zhnumber{2}位）。 2026 年 6 月末货币资金 36.20 亿元、短期债务（短期借款与一年内到期非流动负债）25.39 亿元、长期债务 0.53 亿元，有息负债合计 25.92 亿元，现金短债比 1.43 倍——覆盖充分。 净有息负债 -10.28 亿元，相当于其 2026 年上半年营业收入的 -34\%。 流动比率 1.67、速动比率 1.41，短期偿债能力处于正常区间。 2026 年上半年经营活动现金流净额 0.60 亿元，经营现金流与利润同向为正，内生资金可以支撑运营。 综合看，现金短债比 1.43 倍，短期偿付压力可控；资产负债率低于样本中位数 9.2 个百分点；有息负债中短期债务占比更高，期限结构仍以滚动续作为主。

\pfl{重大战略转型} 公司的战略动作可归为以下几类：产能与项目（2 项）：2026 年 6 月，诸城年产 1.2 亿羽肉鸡产业生态项目一期配套的屠宰加工厂、饲料厂和商品鸡养殖场已全部投产；融资与资本运作（2 项）：2026 年 7 月，全资子公司拟使用 2.40 亿元募集资金对全资孙公司山东仙润鸿食品有限公司进行增资。 公告口径上，近三年（2023 年 9 月---2026 年 9 月）公司披露重大重组与并购类公告 2 条、对外投资与转型类公告 1 条、回购与股权激励类公告 0 条，合计公告 327 条。 第一大行业仍是“肉鸡养殖及屠宰加工”，收入占比 100.0\%，与 2021 年年报相比没有方向性变化。 综合判断，报告期内公司有 3 条与战略相关的公告落地（重大重组与并购 2 条、对外投资与转型 1 条），资本运作与主业调整之间存在直接关联。

\begin{center}\small\setlength{\tabcolsep}{4pt}
\captionof{table}[仙坛股份（002746）历史融资与资本运作]{仙坛股份（002746）历史融资与资本运作（据公司公告与公开报道整理；金额为披露值，含首次公开发行、再融资、债券、银行借款与重整投资等）}
\begin{adjustbox}{max width=\textwidth}
\begin{tabular}{@{}p{1.45cm}p{2.55cm}p{4.05cm}p{4.95cm}@{}}
\toprule
时间 & 方式 & 规模 & 用途与说明 \\
\midrule
2015-02 & IPO（深圳证券交易所中小板） & 3,\allowbreak{}985 万股，\allowbreak{}发行价 5.28 元/\allowbreak{}股，\allowbreak{}募集资金 2.10 亿元 & 发行后总股本 15,\allowbreak{}935 万股 \\
2015-06 & 非公开发行股票（预案，后经董事会、股东大会调整） & 拟发行不超过 2,\allowbreak{}619 万股，\allowbreak{}发行价格不低于 31.98 元/\allowbreak{}股 & 熟食品加工项目、\allowbreak{}商品鸡饲养立体养殖技术改造项目及补充流动资金；2015 年 6 月 23 日披露预案并同日复牌 \\
2020-12 & 非公开发行股票 & 发行 1.10 亿股，\allowbreak{}募集资金总额 10.45 亿元、\allowbreak{}净额 10.31 亿元 & 年产 1.2 亿羽肉鸡产业生态项目，\allowbreak{}预制菜产能 15 万吨 \\
\bottomrule
\end{tabular}
\end{adjustbox}
\end{center}

\pfl{历史融投资情况} 从现金流量表看，2025 年公司取得借款收到的现金 22.71 亿元、偿还债务支付的现金 17.50 亿元，筹资活动现金流净额 -1.32 亿元（2023 年为取得借款 9.38 亿元、偿还债务 4.72 亿元）。 2026 年上半年筹资活动现金流净额为 0.94 亿元（取得借款 20.65 亿元、偿还债务 18.55 亿元，未年化），仍处于净融入状态。 2025 年分配股利、利润或偿付利息支付的现金合计 3.31 亿元。 公开披露的融资记录共 3 笔，工具类型以 IPO、定向增发为主；金额与用途见下表。 其中规模最大的两笔为：2015 年 2 月，IPO，3,985 万股，发行价 5.28 元/股，募集资金 2.10 亿元；2020 年 12 月，定向增发，发行 1.10 亿股，募集资金总额 10.45 亿元、净额 10.31 亿元。 股本方面，近三年披露股本变动 3 次；总股本自 2021 年末以来未发生变化（8.61 亿股）。 分红方面，近三年现金分红 7 次、合计每股派现 0.95 元，最近一次实施在 2026 年 6 月。 综合看，公司融资结构上股权与债务融资并举；2025 年筹资活动现金流净额 -1.32 亿元，融资节奏仍以维持存量债务周转为主，与前述经营与投资现金流的方向基本一致。

\pfl{未来潜在融资需求分析} 2026 年 6 月末货币资金 36.20 亿元，而短期债务 25.39 亿元（现金短债比 1.43 倍），2026 年上半年经营现金流净额 0.60 亿元。账面货币资金可以覆盖短期债务，缺口不体现在静态偿付层面。 公司 2026 年 6 月末资产负债率 43.0\%、2026 年上半年 ROE 3.9\%（未年化），资产负债表稳健、盈利水平平淡，暂不存在刚性融资需求；潜在的融资需求以技改扩产、并购或被并购为主。 公开披露的后续资金安排线索包括：债务与授信安排：2026 年 7 月，公司以募集资金 2.40 亿元增资孙公司并向其提供 1.00 亿元无息借款。 综合看，公司在静态偿付层面不存在缺口，外部融资更多是配合扩张节奏与并购整合的主动性安排，而非偿债压力驱动。




\subsubsection{天邦食品（002124）}
\label{co:002124}

\pfl{公司基本情况} 浙江余姚起家的生猪养殖与食品加工企业（曾称“浙江猪王”），1996 年成立、2007 年深交所上市，2021 年起相继出售疫苗、饲料与史记生物种猪业务以聚焦生猪养殖，2024 年 8 月进入预重整程序，2026 年 9 月与产业投资人解除重整投资协议并重新公开招募投资人，截至 2026 年 6 月末归母净资产 8.23 亿元。 公司于 2007-04-03 在深交所主板首发上市，注册资本 22.22 亿元，总股本 22.22 亿股。 截至 2026 年 10 月 12 日收盘价 1.95 元，对应总市值 43.3 亿元，近一年-32.8\%，总市值在三级行业“生猪养殖”的 11 家公司中按由大到小排第\zhnumber{9}位，近一年涨跌幅低于全样本中位数（-10.4\%）22.4 个百分点。 按 2026 年半年报披露的行业构成，收入占比前三位为养殖 65.4\%；食品 31.4\%；饲料产品 3.2\%。

\pfl{股权结构与控制权} 按 2026 年 6 月 30 日的股东名册，前三大股东为张邦辉 持股 14.94\%；浙商银行股份有限公司－国泰中证畜牧养殖交易型开放式指（持股比例未披露）；潘丽英（持股比例未披露），股东户数 98,434 户。 控制权与股东性质方面，个人直接持有约 3.319 亿股；2024 年 8 月 9 日宁波中院受理预重整申请，截至目前仍处预重整阶段；2026 年 9 月，产业投资人厦门建发物产因参与重整投资未获其国资管理部门核准而行使单方解除权，公司已公开重新招募重整投资人。

\pfl{经营要点} 经营层面，2026 年上半年，出栏各类生猪 388.75 万头，其中仔猪 146.83 万头；2026 年上半年，育肥完全成本降至 12.65 元/公斤、同比下降 0.78 元/公斤，断奶仔猪成本降至 278 元/头、同比下降 7.8\%。

\pfl{所处行业的基本面} 生猪的供给由能繁母猪存栏决定、需求由猪肉消费与替代品价格决定，本轮下行的成因在供给端的产能去化不彻底，而非需求收缩（第 \ref{sec:framework} 章的供需框架）。 从供需两端的当期数据看，2026 年二季度，末全国生猪存栏 42491 万头、同比增长 0.1\%；2026 年 8 月，第 4 周全国活猪价格 11.55 元/公斤；2026 年一季度全国生猪出栏 20026 万头、同比增长 2.8\%。

\pfl{竞争优势} 把该公司放回同类公司比较，2026 年上半年营业收入 36.6 亿元、在三级行业“生猪养殖”的 11 家公司中排第\zhnumber{4}位；毛利率 3.0\%、高于行业内中位数 4.6 个百分点；半年 ROE -96.8\%（未年化）、低于行业内中位数（-14.32\%）；资产负债率 91.2\%、高于行业内中位数（58.9\%）。 公司区别于同类公司的经营特征包括：已建成从养殖场到终端成品的全产业链数字化追溯体系；公司通过产销联动、打造大单品、拓展价值客户提升加工业务销量。 优势在报表上主要体现为高于行业平均的毛利率，这在周期底部是抵御亏损的主要缓冲。

\pfl{历史股价} 本报告样本区间（2021 年 1 月 至 2026 年 10 月）内，股价由 14.27 元变为 1.95 元，累计 -86.3\%，区间最高 17.55 元（2021 年 2 月）、最低 1.95 元（2026 年 10 月）；当前价相当于区间最高价的 11.1\%，已接近区间最低价（1.00 倍）。 月度收益的年化波动率 45.1\%，高于全样本中位数 37.3\%，在三级行业“生猪养殖”的 11 家公司中波动率由高到低排第\zhnumber{3}位。 把股价阶段与公司事件对齐看，2022 年 3 月 至 2024 年 6 月股价下跌-76.1\%，同期（2024 年 4 月）公司申请重整及预重整事项经公司 2024 年第二次临时股东大会审议通过。 从时间对应关系看，公司层面的资本运作与股权、风险事件解释了区间内多数急涨急跌，与行业指数的同步性相对有限。

\begin{center}
\begin{minipage}{\textwidth}\centering
\begin{tikzpicture}
\begin{axis}[
  width=12.6cm, height=3.9cm, scale only axis,
  xmin=2020.922, xmax=2026.972, ymin=1.81, ymax=18.78,
  xtick={2021,2022,2023,2024,2025,2026},
  yearticks,
  yticklabel style={font=\scriptsize},
  ylabel={元/股}, ylabel style={font=\scriptsize},
  grid=major, grid style={LightGray, line width=0.3pt},
  axis line style={InkGray, line width=0.5pt},
  every axis plot/.append style={line width=0.9pt},
]
\addplot[AgriGreen] table[x=x,y=y]{data/clean/profiles/px_002124.dat};
\addplot[SoftRed, dashed, line width=0.7pt] table[x=x,y=y]{data/clean/profiles/pxzz_002124.dat};
\addplot[only marks, mark=*, mark size=1.1pt, SoftRed] table[x=x,y=y]{data/clean/profiles/pxzz_002124.dat};
\end{axis}
\end{tikzpicture}

\captionof{figure}[天邦食品（002124）股价与周期骨架]{天邦食品（002124）月度收盘价（元/股）与 ZigZag 周期骨架（阈值 25\%，圆点为周期转折点）}
\end{minipage}
\end{center}

\pfl{过去周期分析} 以 25\% 的 ZigZag 阈值划分，样本区间内股价形成 2 个主升段与 3 个主跌段。 最大主跌段为 2022 年 3 月 至 2024 年 6 月，27 个月-76.1\%；最大主升段出现在 2021 年 8 月 至 2022 年 3 月，7 个月+59.3\%。 最近一次转折出现在 2026 年 10 月（下跌段自 2024 年 11 月起，23 个月-40.5\%），当前处于该段的延续之中。 与申万农林牧渔指数的最近一次底部相比，该公司见底晚约 5 个月，在本轮下行中属于调整更慢的一类。 生猪养殖的定价主线是产能去化与成本曲线的赛跑（第 \ref{sec:hog} 章）：2026 年三季度行业仍处在“产能去化未完成、价格磨底”的阶段，公司 2026 年 6 月末资产负债率 91.2\%，高于全样本中位数 52.2\%，其股价因此更多反映资金面与信用风险，而不是价格弹性本身。 综合区间形态看，该股单段涨跌幅度偏大、以趋势段为主（最大主升 +59.3\%、最大主跌 -76.1\%），年化波动率 45.1\% 与全样本中位数 37.3\% 相比波动明显大于样本中位数。

\pfl{盈利情况} 2026 年上半年营业收入 36.56 亿元（同比 -22.1\%），归母净利润 -15.43 亿元，由上年同期的盈利转为亏损。 以年报为趋势对照，2023---2025 年营业收入由 102.32 亿元变为 86.89 亿元（两年年均复合增速 -7.8\%），归母净利润由 -28.83 亿元变为 -13.09 亿元。 按半年报口径，2026 年上半年的 ROE 为 -96.8\%（未年化）、销售净利率 -42.2\%、毛利率 3.0\%，ROE 在“生猪养殖”11 家公司中按数值由高到低排第\zhnumber{11}位。 按同一口径，三级行业“生猪养殖”11 家公司的半年 ROE 中位数为 -14.32\%，该公司低于其中位数 82.44 个百分点。 综合看，公司仍处亏损区间、由盈转亏，半年 ROE 在三级行业内按由高到低排第\zhnumber{11}位，单位盈利能力的修复依赖行业景气与自身成本端的改善，报表弹性主要体现为价格与销量的方向性变化。

\begin{center}\small\setlength{\tabcolsep}{3.4pt}
\captionof{table}[天邦食品（002124）关键财务指标]{天邦食品（002124）关键财务指标（合并报表口径；两个半年报期均未年化；现金短债比 $=$ 货币资金 $\div$ 短期债务）}
\begin{adjustbox}{max width=\textwidth}
\begin{tabular}{@{}lrrrrr@{}}
\toprule
指标 & 2023 年 & 2024 年 & 2025 年 & 2025 年半年报 & 2026 年半年报 \\
\midrule
营业收入（亿元） & 102.32 & 96.44 & 86.89 & 46.93 & 36.56 \\
归母净利润（亿元） & -28.83 & 14.59 & -13.09 & 3.52 & -15.43 \\
毛利率（\%） & -8.9 & 16.5 & 8.6 & 17.4 & 3.0 \\
ROE（\%） & -116.3 & 49.5 & -43.4 & 9.1 & -96.8 \\
资产负债率（\%） & 86.7 & 72.6 & 79.8 & 69.8 & 91.2 \\
经营活动现金流净额（亿元） & 2.64 & -9.64 & 9.11 & 4.89 & 2.71 \\
货币资金（亿元） & 12.54 & 1.50 & 2.25 & 2.34 & 1.55 \\
有息负债合计（亿元） & 61.93 & 52.86 & 49.26 & 51.10 & 48.40 \\
现金短债比（倍） & 0.30 & 0.04 & 0.06 & 0.07 & 0.04 \\
\bottomrule
\end{tabular}
\end{adjustbox}
\end{center}

\begin{center}
\begin{minipage}{\textwidth}\centering
\begin{tikzpicture}
\begin{groupplot}[
  group style={group size=2 by 1, horizontal sep=0.60cm},
  width=6.85cm, height=4.60cm,
  tick label style={font=\tiny},
  title style={font=\scriptsize\heiti},
  yticklabel style={font=\tiny},
  ylabel style={font=\tiny},
  grid=major, grid style={LightGray, line width=0.3pt},
  axis line style={InkGray, line width=0.5pt},
  enlarge x limits={abs=0.8},
  legend style={font=\scriptsize, at={(0.5,-0.20)}, anchor=north,
                legend columns=2, draw=none, fill=none, column sep=6pt},
]
\nextgroupplot[
  ybar, bar width=4pt,
  ymin=-59.59, ymax=123.03,
  xtick={0,1,2,3,4,5.7},
  xticklabels={2021,2022,2023,2024,2025,2026H1},
  ylabel={亿元},
]
\addplot[fill=AgriGreen!75, draw=AgriGreen, bar shift=-2.6pt] table[x=i, y=rev]{data/clean/profiles/fin_002124.dat};
\addplot[fill=SoftRed!60, draw=SoftRed, bar shift=2.6pt] table[x=i, y=ni]{data/clean/profiles/fin_002124.dat};
\addplot[fill=AgriGreen!30, draw=AgriGreen, pattern=north east lines, bar shift=-2.6pt] table[x=x, y=rev]{data/clean/profiles/finh_002124.dat};
\addplot[fill=SoftRed!20, draw=SoftRed, pattern=north east lines, bar shift=2.6pt] table[x=x, y=ni]{data/clean/profiles/finh_002124.dat};
\legend{营业收入（年/半年）, 归母净利润（年/半年）}
\nextgroupplot[
  ymin=-132.9, ymax=112.0,
  xtick={0,1,2,3,4,5.7},
  xticklabels={2021,2022,2023,2024,2025,2026H1},
  ylabel={\%},
]
\addplot[AgriGold, mark=*, mark size=1.3pt] table[x=i, y=roe]{data/clean/profiles/fin_002124.dat};
\addplot[OilBlue, mark=square*, mark size=1.3pt] table[x=i, y=debt]{data/clean/profiles/fin_002124.dat};
\addplot[only marks, AgriGold, mark=*, mark size=2.0pt] table[x=x, y=roe]{data/clean/profiles/finh_002124.dat};
\addplot[only marks, OilBlue, mark=square*, mark size=2.0pt] table[x=x, y=debt]{data/clean/profiles/finh_002124.dat};
\legend{ROE, 资产负债率}
\end{groupplot}
\end{tikzpicture}
\begin{tikzpicture}
\begin{axis}[
  name=finbot,
  width=13.75cm, height=3.10cm, scale only axis,
  ymin=-7.95, ymax=89.06,
  xtick={0,1,2,3,4,5.7},
  xticklabels={2021,2022,2023,2024,2025,2026H1},
  tick label style={font=\tiny},
  yticklabel style={font=\tiny},
  ylabel={亿元}, ylabel style={font=\tiny},
  ybar, bar width=4pt,
  grid=major, grid style={LightGray, line width=0.3pt},
  axis line style={InkGray, line width=0.5pt},
  enlarge x limits={abs=0.8},
  legend style={font=\scriptsize, at={(0.5,-0.26)}, anchor=north,
                legend columns=2, draw=none, fill=none, column sep=10pt},
]
\addplot[fill=AgriGreen!55, draw=AgriGreen, bar shift=-3.0pt] table[x=i, y=cash]{data/clean/profiles/fin_002124.dat};
\addplot[fill=SoftRed!45, draw=SoftRed, bar shift=3.0pt] table[x=i, y=ibd]{data/clean/profiles/fin_002124.dat};
\addplot[fill=AgriGreen!25, draw=AgriGreen, pattern=north east lines, bar shift=-3.0pt] table[x=x, y=cash]{data/clean/profiles/finh_002124.dat};
\addplot[fill=SoftRed!20, draw=SoftRed, pattern=north east lines, bar shift=3.0pt] table[x=x, y=ibd]{data/clean/profiles/finh_002124.dat};
\legend{货币资金（左轴）, 有息负债（左轴）}
\end{axis}
\begin{axis}[
  at={(finbot.south east)}, anchor=south east,
  width=13.75cm, height=3.10cm, scale only axis,
  axis y line*=right, axis x line*=none, xtick=\empty,
  ymin=-14.52, ymax=14.73,
  tick label style={font=\tiny},
  yticklabel style={font=\tiny}, scaled y ticks=false,
  legend style={font=\scriptsize, at={(0.5,-0.44)}, anchor=north,
                legend columns=1, draw=none, fill=none},
]
\addplot[OilBlue, mark=*, mark size=2.2pt, line width=1.3pt] table[x=i, y=ocf]{data/clean/profiles/fin_002124.dat};
\addplot[only marks, OilBlue, mark=*, mark size=3.2pt] table[x=x, y=ocf]{data/clean/profiles/finh_002124.dat};
\legend{经营活动现金流净额（右轴，亿元，蓝点）}
\end{axis}
\end{tikzpicture}

\captionof{figure}[天邦食品（002124）近年主要财务指标]{天邦食品（002124）近年主要财务指标（左上：营业收入与归母净利润；右上：ROE 与资产负债率；下：货币资金、有息负债为左轴柱状，经营活动现金流净额为其右轴的蓝点折线。
2021---2025 年为年报口径，2026H1 为 2026 年半年报/半年末，与其前一年报之间留出间隔、以斜纹或空心点区分；半年报数据未年化）}
\end{minipage}
\end{center}

\pfl{财务分析} 2026 年 6 月末资产负债率 91.2\%，较 2025 年末上升 11.4 个百分点（样本中位数 52.2\%，所处三级行业内按由低到高排第\zhnumber{11}位）。 2026 年 6 月末货币资金 1.55 亿元、短期债务（短期借款与一年内到期非流动负债）35.01 亿元、长期债务 13.38 亿元，有息负债合计 48.40 亿元，现金短债比 0.04 倍——缺口明显。 净有息负债 46.85 亿元，相当于其 2026 年上半年营业收入的 128\%。 流动比率 0.18、速动比率 0.05，短期偿债指标偏紧。 2026 年上半年经营活动现金流净额 2.71 亿元，净利润为负而经营现金流为正，差额主要来自折旧摊销与营运资本变动，说明亏损尚未完全转化为现金流出。 2026 年 6 月末在建工程 7.73 亿元，较上年末减少 0.21 亿元，在建投入在收缩。 综合看，账面货币资金对有息负债与短期债务的覆盖不足，滚动偿付对新增融资的依赖度较高；资产负债率高于样本中位数 39.0 个百分点；有息负债中短期债务占比更高，期限结构仍以滚动续作为主。 公司披露的风险集中在退市与持续经营、重整与债务违约两类：一是 2026 年，公司股票将被实施退市风险警示（*ST）；二是 2026 年 4 月，决定延长公司预重整期间至 2026 年 11 月 9 日。

\pfl{重大战略转型} 公司的战略动作可归为以下几类：并购与重组（4 项）：2026 年 9 月，公司继续公开招募重整投资人；2022 年 6 月，公司以 10.20 亿元转让史记生物技术（南京）有限公司 51\% 股权给三亚史记生物科技有限公司；融资与资本运作（1 项）：2023 年度，公司向特定对象发行股票募集资金将剩余的募集资金用途全部变更为新募投项目中的“天邦股份数智化猪场升级项目”。 公告口径上，近三年（2023 年 9 月---2026 年 9 月）公司披露重大重组与并购类公告 4 条、对外投资与转型类公告 5 条、回购与股权激励类公告 0 条，合计公告 493 条。 第一大行业仍是“养殖”，收入占比 65.4\%，与 2021 年年报相比没有方向性变化。 综合判断，公司当前的战略重心不是扩张而是化解债务与完成重整，业务层面的调整让位于司法重整程序。

\begin{center}\small\setlength{\tabcolsep}{4pt}
\captionof{table}[天邦食品（002124）历史融资与资本运作]{天邦食品（002124）历史融资与资本运作（据公司公告与公开报道整理；金额为披露值，含首次公开发行、再融资、债券、银行借款与重整投资等）}
\begin{adjustbox}{max width=\textwidth}
\begin{tabular}{@{}p{1.45cm}p{2.55cm}p{4.05cm}p{4.95cm}@{}}
\toprule
时间 & 方式 & 规模 & 用途与说明 \\
\midrule
2007-04 & IPO（深圳证券交易所中小板） & 实际发行 1,\allowbreak{}850.00 万股 & 发行市盈率（按发行后总股本）29.97 倍 \\
2015-12 & 非公开发行股票 & 发行 8,\allowbreak{}078.99 万股，\allowbreak{}发行价 5.57 元/\allowbreak{}股，\allowbreak{}募集资金总额 4.50 亿元 & 公司于 2015 年 12 月在宁波市市场监督管理局办理工商变更手续 \\
2017-03-16 & 非公开发行股票（向特定对象） & 发行 1.36 亿股，\allowbreak{}发行价 10.85 元/\allowbreak{}股，\allowbreak{}募集资金总额 14.80 亿元 & 募投项目，\allowbreak{}后多次变更用途（累计变更用途的募集资金总额占比曾达 87.55\%）；；新增股份于 2017 年 3 月 16 日上市 \\
2020-08-21 & 非公开发行股票（向 7 家投资者） & 发行 1.54 亿股，\allowbreak{}发行价 17.30 元/\allowbreak{}股，\allowbreak{}募集资金总额 26.66 亿元 & 募集资金于 2020 年 8 月 10 日到账 \\
2021-09 & 出售饲料资产并开展战略合作（向通威股份） & 通威股份以 12.71 亿元收购天邦股份旗下水产饲料全部资产及猪料部分资产 & 回笼资金、\allowbreak{}聚焦生猪养殖与食品加工；2021 年 9 月 20 日公司第七届董事会第四十五次会议审议通过向通威股份出售猪料子公司部分股权并开展长期饲料供应战略合作的议案，\allowbreak{}股权；同时签署猪料长单合作协议 \\
2022-06-30 & 转让子公司股权（资产出售） & 以 10.20 亿元转让史记生物技术（南京）有限公司 51\% 股权（评估值 10.19 亿元） & 改善盈利及资产负债状况、\allowbreak{}聚焦生猪养殖业务；受让方为三亚史记生物科技有限公司（史记生物管理及员工主导） \\
2023-12-29 & 向特定对象发行股票（定增） & 发行 3.83 亿股，\allowbreak{}发行价 3.13 元/\allowbreak{}股，\allowbreak{}募集资金总额 11.98 亿元 & 数智化猪场升级项目及补充流动资金（拟投入数智化猪场升级项目 20.00 亿元 \\
2025-05-08 & 重整投资（预重整投资人出资） & 产业投资人厦门建发物产有限公司、\allowbreak{}南宁漓源粮油饲料有限公司拟以 1.85 元/\allowbreak{}股价格分别取得 2 亿股天邦食品转增股份 & 用于天邦食品重整；协议签订前两家公司未持有天邦食品股份；重整后两家各持 2 亿股 \\
2025-05-08 & 重整投资协议（产业投资人） & 85 元/\allowbreak{}股（不得低于前 20 个交易日股票交易均价 3.22 元/\allowbreak{}股的相应比例） & 受让转增股份、\allowbreak{}参与重整投资；经投资人评审委员会综合评定 \\
\bottomrule
\end{tabular}
\end{adjustbox}
\end{center}

\pfl{历史融投资情况} 从现金流量表看，2025 年公司取得借款收到的现金 24.15 亿元、偿还债务支付的现金 24.60 亿元、吸收权益性投资 0.03 亿元，筹资活动现金流净额 -5.64 亿元（2023 年为取得借款 36.06 亿元、偿还债务 35.15 亿元）。 2026 年上半年筹资活动现金流净额为 -2.88 亿元（取得借款 6.87 亿元、偿还债务 7.19 亿元，未年化），已转为净流出，处于净还债状态。 2025 年分配股利、利润或偿付利息支付的现金合计 1.31 亿元。 公开披露的融资记录共 9 笔，工具类型以 IPO、定向增发、重整投资为主；金额与用途见下表。 其中规模最大的两笔为：2025 年 5 月，重整投资，产业投资人厦门建发物产有限公司、南宁漓源粮油饲料有限公司拟以 1.85 元/股价格分别取得 2 亿股天邦食品转增股份；2007 年 4 月，IPO，实际发行 1,850.00 万股。 股本方面，近三年披露股本变动 5 次（增发新股上市 1 次、限售股份上市 1 次）；总股本由 2021 年末的 18.39 亿股变为 22.22 亿股（+20.81\%）。 分红方面，近三年未实施现金分红，与重整期间的资金约束一致。 近三年“再融资”类公告 29 条，涉及定向增发、募集资金使用。 综合看，公司融资结构上以债务滚动为主、股权融资为补充；2025 年筹资活动现金流净额 -5.64 亿元，融资节奏仍以维持存量债务周转为主，与前述经营与投资现金流的方向基本一致。

\pfl{未来潜在融资需求分析} 2026 年 6 月末货币资金 1.55 亿元，而短期债务 35.01 亿元（现金短债比 0.04 倍），2026 年上半年经营现金流净额 2.71 亿元。按“短期债务减去账面货币资金”的滚动偿付口径，静态缺口约 33.46 亿元；若再计入上半年亏损的年化额（30.86 亿元），维持一年正常经营所需的外部资金约 33.46---64.32 亿元。 按第 \ref{sec:financing-need} 节的三档划分，公司属于第一档“补血型”：2026 年上半年归母净利润 -15.43 亿元、6 月末资产负债率 91.2\%。负债率已超过 80\%，净资产被亏损侵蚀的程度较深，属于全部样本中资产负债表最紧张的一组。 可行的融资路径按现实性排序为：出售非核心资产与参股股权、定向增发（需股价配合，且控股股东需放弃优先认购或同步增资）、引入国资或产业投资人（部分公司以控制权让渡为对价）。 截至目前公司仍处于预重整/重整程序中，重整投资人的招募与投资款到位是未来一年最主要的外部资金来源。 资金需求还要叠加在建工程：期末在建工程 7.73 亿元、2026 年上半年购建固定资产等支出 0.50 亿元（未年化），这部分支出在周期底部通常会被主动放缓。 综合看，公司的外部资金需求以债务接续为主，滚动偿付口径下的静态缺口约 33.46 亿元，融资的时点与规模更多取决于债权人与重整投资人的进度。




\subsubsection{华英农业（002321）}
\label{co:002321}

\pfl{公司基本情况} 公司前身可追溯至 1991 年，河南信阳潢川的鸭肉全产业链企业（曾称“世界鸭王”），2009 年深交所中小板上市，2020 年因大额亏损、控股股东非经营性占用资金及无法表示意见的审计报告被实施退市风险警示，2021 年完成破产重整、2022 年摘帽，现聚焦食品与羽绒双主业，2026 年上半年归母净亏损 2,928.25 万元。 公司于 2009-12-16 在深交所主板首发上市，注册资本 21.33 亿元，总股本 21.33 亿股。 截至 2026 年 10 月 12 日收盘价 2.02 元，对应总市值 43.1 亿元，近一年-30.8\%，总市值在三级行业“其他养殖”的 3 家公司中按由大到小排第\zhnumber{1}位，近一年涨跌幅低于全样本中位数（-10.4\%）20.5 个百分点。 按 2026 年半年报披露的行业构成，收入几乎全部来自畜牧业 100.0\%。 股权结构方面，实际控制人为信阳市财政局。

\pfl{股权结构与控制权} 按 2026 年 9 月 9 日的股东名册，前三大股东为信阳市产业投资集团有限公司 持股 16.90\%；河南光州辰悦实业有限公司 持股 15.30\%；广汉东兴羽绒制品有限公司 持股 6.76\%。 从控制权角度看，控股股东信阳市鼎新兴华产业投资合伙企业（有限合伙）持股 24.00\%（5.12 亿股）；鼎新兴华全体合伙人同意解散、按实缴出资比例将所持股份以非交易过户方式分配。

\pfl{经营要点} 经营层面，2026 年上半年，主要产品价格维持在低位区间波动。

\pfl{所处行业的基本面} 牛、羊的供给由母畜存栏与出栏节奏决定、扩张周期长，水产则由投苗量与消费需求决定（第 \ref{sec:framework} 章的供需框架）。 与该公司直接相关的行业数据包括：肉鸭行业前期产能扩张与后期需求不足叠加；羽绒价格处高位运行，对上游水禽养殖形成一定支撑。

\pfl{竞争优势} 把该公司放回同类公司比较，2026 年上半年营业收入 23.7 亿元、在三级行业“其他养殖”的 3 家公司中排第\zhnumber{1}位；毛利率 5.4\%、低于行业内中位数 3.7 个百分点；半年 ROE -3.2\%（未年化）、低于行业内中位数（0.83\%）；资产负债率 66.2\%、高于行业内中位数（66.2\%）。 就当期报表看，公司尚未把经营层面的特征转化为高于行业的资本回报，周期底部的相对优势更多体现在现金流与资产负债表上。

\pfl{历史股价} 自 2021 年 1 月至 2026 年 10 月，股价由 3.20 元变为 2.02 元，累计 -36.9\%，区间最高 3.96 元（2021 年 11 月）、最低 1.43 元（2024 年 6 月）；当前价相当于区间最高价的 51.0\%，为区间最低价的 1.41 倍。 月度收益的年化波动率 35.0\%，低于全样本中位数 37.3\%，在三级行业“其他养殖”的 3 家公司中波动率由高到低排第\zhnumber{3}位。

\begin{center}
\begin{minipage}{\textwidth}\centering
\begin{tikzpicture}
\begin{axis}[
  width=12.6cm, height=3.9cm, scale only axis,
  xmin=2020.922, xmax=2026.972, ymin=1.33, ymax=4.24,
  xtick={2021,2022,2023,2024,2025,2026},
  yearticks,
  yticklabel style={font=\scriptsize},
  ylabel={元/股}, ylabel style={font=\scriptsize},
  grid=major, grid style={LightGray, line width=0.3pt},
  axis line style={InkGray, line width=0.5pt},
  every axis plot/.append style={line width=0.9pt},
]
\addplot[AgriGreen] table[x=x,y=y]{data/clean/profiles/px_002321.dat};
\addplot[SoftRed, dashed, line width=0.7pt] table[x=x,y=y]{data/clean/profiles/pxzz_002321.dat};
\addplot[only marks, mark=*, mark size=1.1pt, SoftRed] table[x=x,y=y]{data/clean/profiles/pxzz_002321.dat};
\end{axis}
\end{tikzpicture}

\captionof{figure}[华英农业（002321）股价与周期骨架]{华英农业（002321）月度收盘价（元/股）与 ZigZag 周期骨架（阈值 25\%，圆点为周期转折点）}
\end{minipage}
\end{center}

\pfl{过去周期分析} 以 25\% 的 ZigZag 阈值划分，样本区间内股价形成 3 个主升段与 4 个主跌段。 最大主跌段为 2023 年 5 月 至 2024 年 6 月，13 个月-50.0\%；最大主升段出现在 2024 年 6 月 至 2025 年 9 月，15 个月+107.7\%。 最近一次转折出现在 2026 年 6 月（下跌段自 2025 年 9 月起，9 个月-43.8\%），当前处于该段的延续之中。 公司股价较申万农林牧渔指数提前约 23 个月见底，是板块中较早定价周期反转的公司之一。 牛、羊、水产的产能周期长于生猪，第 \ref{sec:other-livestock} 章判断肉牛供给的实质性收缩要到 2027 年之后才能兑现，这一时滞决定了个股的股价弹性，公司 2026 年 6 月末资产负债率 66.2\%，高于全样本中位数 52.2\%，其股价因此更多反映资金面与信用风险，而不是价格弹性本身。 综合区间形态看，该股单段涨跌幅度偏大、以趋势段为主（最大主升 +107.7\%、最大主跌 -50.0\%），年化波动率 35.0\% 与全样本中位数 37.3\% 相比波动与样本中位数接近。

\pfl{盈利情况} 2026 年上半年营业收入 23.72 亿元（同比 +14.5\%），归母净利润 -0.29 亿元，亏损同比扩大 0.09 亿元。 以年报为趋势对照，2023---2025 年营业收入由 37.05 亿元变为 54.47 亿元（两年年均复合增速 +21.3\%），归母净利润由 -0.32 亿元变为 -0.60 亿元。 按半年报口径，2026 年上半年的 ROE 为 -3.2\%（未年化）、销售净利率 -0.2\%、毛利率 5.4\%，ROE 在“其他养殖”3 家公司中按数值由高到低排第\zhnumber{3}位。 同期全样本半年 ROE 中位数 0.75\%，公司与之相差 3.92 个百分点（弱于中位数公司）。 从公司披露的原因看，业绩变动主要来自：2025 年，同比增长 15.15\%，受产品价格承压影响板块毛利率回落。 综合看，公司仍处亏损区间、亏损同比扩大，半年 ROE 在三级行业内按由高到低排第\zhnumber{3}位，单位盈利能力的修复依赖行业景气与自身成本端的改善，报表弹性主要体现为价格与销量的方向性变化。

\begin{center}\small\setlength{\tabcolsep}{3.4pt}
\captionof{table}[华英农业（002321）关键财务指标]{华英农业（002321）关键财务指标（合并报表口径；两个半年报期均未年化；现金短债比 $=$ 货币资金 $\div$ 短期债务）}
\begin{adjustbox}{max width=\textwidth}
\begin{tabular}{@{}lrrrrr@{}}
\toprule
指标 & 2023 年 & 2024 年 & 2025 年 & 2025 年半年报 & 2026 年半年报 \\
\midrule
营业收入（亿元） & 37.05 & 47.31 & 54.47 & 20.71 & 23.72 \\
归母净利润（亿元） & -0.32 & 0.32 & -0.60 & -0.21 & -0.29 \\
毛利率（\%） & 9.1 & 7.5 & 5.8 & 6.8 & 5.4 \\
ROE（\%） & -3.2 & 3.3 & -6.4 & -2.2 & -3.2 \\
资产负债率（\%） & 62.1 & 61.1 & 63.4 & 64.8 & 66.2 \\
经营活动现金流净额（亿元） & -1.33 & 2.00 & 3.76 & 1.86 & -2.96 \\
货币资金（亿元） & 2.42 & 2.75 & 4.14 & 3.29 & 1.92 \\
有息负债合计（亿元） & 3.99 & 4.73 & 5.46 & 4.91 & 6.13 \\
现金短债比（倍） & 1.06 & 0.67 & 0.86 & 0.78 & 0.35 \\
\bottomrule
\end{tabular}
\end{adjustbox}
\end{center}

\begin{center}
\begin{minipage}{\textwidth}\centering
\begin{tikzpicture}
\begin{groupplot}[
  group style={group size=2 by 1, horizontal sep=0.60cm},
  width=6.85cm, height=4.60cm,
  tick label style={font=\tiny},
  title style={font=\scriptsize\heiti},
  yticklabel style={font=\tiny},
  ylabel style={font=\tiny},
  grid=major, grid style={LightGray, line width=0.3pt},
  axis line style={InkGray, line width=0.5pt},
  enlarge x limits={abs=0.8},
  legend style={font=\scriptsize, at={(0.5,-0.20)}, anchor=north,
                legend columns=2, draw=none, fill=none, column sep=6pt},
]
\nextgroupplot[
  ybar, bar width=4pt,
  ymin=-33.37, ymax=64.05,
  xtick={0,1,2,3,4,5.7},
  xticklabels={2021,2022,2023,2024,2025,2026H1},
  ylabel={亿元},
]
\addplot[fill=AgriGreen!75, draw=AgriGreen, bar shift=-2.6pt] table[x=i, y=rev]{data/clean/profiles/fin_002321.dat};
\addplot[fill=SoftRed!60, draw=SoftRed, bar shift=2.6pt] table[x=i, y=ni]{data/clean/profiles/fin_002321.dat};
\addplot[fill=AgriGreen!30, draw=AgriGreen, pattern=north east lines, bar shift=-2.6pt] table[x=x, y=rev]{data/clean/profiles/finh_002321.dat};
\addplot[fill=SoftRed!20, draw=SoftRed, pattern=north east lines, bar shift=2.6pt] table[x=x, y=ni]{data/clean/profiles/finh_002321.dat};
\legend{营业收入（年/半年）, 归母净利润（年/半年）}
\nextgroupplot[
  ymin=-1902.7, ymax=248.5,
  xtick={0,1,2,3,4,5.7},
  xticklabels={2021,2022,2023,2024,2025,2026H1},
  ylabel={\%},
]
\addplot[AgriGold, mark=*, mark size=1.3pt] table[x=i, y=roe]{data/clean/profiles/fin_002321.dat};
\addplot[OilBlue, mark=square*, mark size=1.3pt] table[x=i, y=debt]{data/clean/profiles/fin_002321.dat};
\addplot[only marks, AgriGold, mark=*, mark size=2.0pt] table[x=x, y=roe]{data/clean/profiles/finh_002321.dat};
\addplot[only marks, OilBlue, mark=square*, mark size=2.0pt] table[x=x, y=debt]{data/clean/profiles/finh_002321.dat};
\legend{ROE, 资产负债率}
\end{groupplot}
\end{tikzpicture}
\begin{tikzpicture}
\begin{axis}[
  name=finbot,
  width=13.75cm, height=3.10cm, scale only axis,
  ymin=-0.92, ymax=10.30,
  xtick={0,1,2,3,4,5.7},
  xticklabels={2021,2022,2023,2024,2025,2026H1},
  tick label style={font=\tiny},
  yticklabel style={font=\tiny},
  ylabel={亿元}, ylabel style={font=\tiny},
  ybar, bar width=4pt,
  grid=major, grid style={LightGray, line width=0.3pt},
  axis line style={InkGray, line width=0.5pt},
  enlarge x limits={abs=0.8},
  legend style={font=\scriptsize, at={(0.5,-0.26)}, anchor=north,
                legend columns=2, draw=none, fill=none, column sep=10pt},
]
\addplot[fill=AgriGreen!55, draw=AgriGreen, bar shift=-3.0pt] table[x=i, y=cash]{data/clean/profiles/fin_002321.dat};
\addplot[fill=SoftRed!45, draw=SoftRed, bar shift=3.0pt] table[x=i, y=ibd]{data/clean/profiles/fin_002321.dat};
\addplot[fill=AgriGreen!25, draw=AgriGreen, pattern=north east lines, bar shift=-3.0pt] table[x=x, y=cash]{data/clean/profiles/finh_002321.dat};
\addplot[fill=SoftRed!20, draw=SoftRed, pattern=north east lines, bar shift=3.0pt] table[x=x, y=ibd]{data/clean/profiles/finh_002321.dat};
\legend{货币资金（左轴）, 有息负债（左轴）}
\end{axis}
\begin{axis}[
  at={(finbot.south east)}, anchor=south east,
  width=13.75cm, height=3.10cm, scale only axis,
  axis y line*=right, axis x line*=none, xtick=\empty,
  ymin=-9.78, ymax=6.99,
  tick label style={font=\tiny},
  yticklabel style={font=\tiny}, scaled y ticks=false,
  legend style={font=\scriptsize, at={(0.5,-0.44)}, anchor=north,
                legend columns=1, draw=none, fill=none},
]
\addplot[OilBlue, mark=*, mark size=2.2pt, line width=1.3pt] table[x=i, y=ocf]{data/clean/profiles/fin_002321.dat};
\addplot[only marks, OilBlue, mark=*, mark size=3.2pt] table[x=x, y=ocf]{data/clean/profiles/finh_002321.dat};
\legend{经营活动现金流净额（右轴，亿元，蓝点）}
\end{axis}
\end{tikzpicture}

\captionof{figure}[华英农业（002321）近年主要财务指标]{华英农业（002321）近年主要财务指标（左上：营业收入与归母净利润；右上：ROE 与资产负债率；下：货币资金、有息负债为左轴柱状，经营活动现金流净额为其右轴的蓝点折线。
2021---2025 年为年报口径，2026H1 为 2026 年半年报/半年末，与其前一年报之间留出间隔、以斜纹或空心点区分；半年报数据未年化）}
\end{minipage}
\end{center}

\pfl{财务分析} 2026 年 6 月末资产负债率 66.2\%，较 2025 年末上升 2.8 个百分点（样本中位数 52.2\%，所处三级行业内按由低到高排第\zhnumber{2}位）。 2026 年 6 月末货币资金 1.92 亿元、短期债务（短期借款与一年内到期非流动负债）5.52 亿元、长期债务 0.61 亿元，有息负债合计 6.13 亿元，现金短债比 0.35 倍——覆盖偏紧。 净有息负债 4.21 亿元，相当于其 2026 年上半年营业收入的 18\%。 流动比率 1.49、速动比率 0.77，短期流动性无明显缺口。 2026 年上半年经营活动现金流净额 -2.96 亿元，经营现金流与利润同时为负，日常运营对债务与股东投入的依赖度较高。 2026 年 6 月末商誉 1.12 亿元，占归母净资产的 7.9\%，是净资产中最脆弱的一块。 综合看，账面货币资金对有息负债与短期债务的覆盖不足，滚动偿付对新增融资的依赖度较高；资产负债率高于样本中位数 14.0 个百分点；有息负债中短期债务占比更高，期限结构仍以滚动续作为主。

\pfl{重大战略转型} 近三年公司的战略动作集中在：并购与重组（1 项）：2022 年，重整落定后提出冲刺“世界鸭王+世界羽王”双产业目标；其他动作（1 项）：2022 年，公司拟处置评估价值为 5.70 亿元的长期股权投资以及评估值为 2.77 亿元的其他应收款资产。 公告口径上，近三年（2023 年 9 月---2026 年 9 月）公司披露重大重组与并购类公告 0 条、对外投资与转型类公告 0 条、回购与股权激励类公告 8 条，合计公告 355 条。 第一大行业仍是“畜牧业”，收入占比 100.0\%，与 2021 年年报相比没有方向性变化。 综合判断，公司当前的战略重心不是扩张而是化解债务与完成重整，业务层面的调整让位于司法重整程序。

\begin{center}\small\setlength{\tabcolsep}{4pt}
\captionof{table}[华英农业（002321）历史融资与资本运作]{华英农业（002321）历史融资与资本运作（据公司公告与公开报道整理；金额为披露值，含首次公开发行、再融资、债券、银行借款与重整投资等）}
\begin{adjustbox}{max width=\textwidth}
\begin{tabular}{@{}p{1.45cm}p{2.55cm}p{4.05cm}p{4.95cm}@{}}
\toprule
时间 & 方式 & 规模 & 用途与说明 \\
\midrule
2009-12 & IPO（深圳证券交易所中小板） & 3,\allowbreak{}700 万股，\allowbreak{}发行价 16.98 元/\allowbreak{}股，\allowbreak{}实际募集资金 6.28 亿元 & 依序投向淮阳 4000 万只/\allowbreak{}年商品肉鸡系列加工生产项目、\allowbreak{}丰城 1000 万只/\allowbreak{}年及单县 2000 万只/\allowbreak{}年樱桃谷商品雏鸭生产项目等；本次发行采用网下询价配售（740 万股 \\
2021-12 & 破产重整——重整投资人受让转增股票出资 & 产业投资人鼎新兴华以人民币约 3.62 亿元有条件受让华英农业转增股票约 5.12 亿股 & 受让转增股票、\allowbreak{}解决控股股东资金占用；2021 年 12 月 23 日公司与鼎新兴华、\allowbreak{}光州辰悦、\allowbreak{}信阳华信、\allowbreak{}农投新动能、\allowbreak{}广兴股权签署重整投资协议 \\
2021-12 & 资本公积转增股本（重整出资人权益调整） & 以 2021 年 12 月 31 日总股本 53,\allowbreak{}429.11 万股为基数 & 其中 5.32 亿股以股抵债形式清偿公司债务 \\
\bottomrule
\end{tabular}
\end{adjustbox}
\end{center}

\pfl{历史融投资情况} 从现金流量表看，2025 年公司取得借款收到的现金 5.61 亿元、偿还债务支付的现金 5.21 亿元、吸收权益性投资 0.00 亿元，筹资活动现金流净额 -2.88 亿元（2023 年为取得借款 3.11 亿元、偿还债务 2.35 亿元）。 2026 年上半年筹资活动现金流净额为 1.97 亿元（取得借款 1.83 亿元、偿还债务 0.92 亿元，未年化），仍处于净融入状态。 2025 年分配股利、利润或偿付利息支付的现金合计 0.82 亿元。 公开披露的融资记录共 3 笔，工具类型以 IPO、重整投资为主；金额与用途见下表。 其中规模最大的两笔为：2021 年 12 月，重整投资，产业投资人鼎新兴华以人民币约 3.62 亿元有条件受让华英农业转增股票约 5.12 亿股；2009 年 12 月，IPO，3,700 万股，发行价 16.98 元/股，实际募集资金 6.28 亿元。 股本方面，近三年披露股本变动 3 次；总股本自 2021 年末以来未发生变化（21.33 亿股）。 分红方面，近三年未实施现金分红，与重整期间的资金约束一致。 近三年“再融资”类公告 1 条，涉及定向增发。 综合看，公司融资结构上以债务滚动为主、股权融资为补充；2025 年筹资活动现金流净额 -2.88 亿元，融资节奏仍以维持存量债务周转为主，与前述经营与投资现金流的方向基本一致。

\pfl{未来潜在融资需求分析} 2026 年 6 月末货币资金 1.92 亿元，而短期债务 5.52 亿元（现金短债比 0.35 倍），2026 年上半年经营现金流净额 -2.96 亿元。按“短期债务减去账面货币资金”的滚动偿付口径，静态缺口约 3.60 亿元；若再计入上半年亏损的年化额（0.59 亿元），维持一年正常经营所需的外部资金约 3.60---4.18 亿元。 按第 \ref{sec:financing-need} 节的三档划分，公司属于第一档“补血型”：2026 年上半年归母净利润 -0.29 亿元、6 月末资产负债率 66.2\%。 可行的融资路径按现实性排序为：定向增发（需股价配合，且控股股东需放弃优先认购或同步增资）、引入国资或产业投资人（部分公司以控制权让渡为对价）、债务重组或展期（与银行和债券持有人协商）。 截至目前公司仍处于预重整/重整程序中，重整投资人的招募与投资款到位是未来一年最主要的外部资金来源。 综合看，公司的外部资金需求以债务接续为主，滚动偿付口径下的静态缺口约 3.60 亿元，融资的时点与规模更多取决于债权人与重整投资人的进度。




\subsubsection{益客食品（301116）}
\label{co:301116}

\pfl{公司基本情况} 公司前身可追溯至 2008 年，江苏肉禽（肉鸭、肉鸡）屠宰与加工企业，2022 年 1 月创业板上市，形成肉鸡、肉鸭双产业链协同，产品覆盖鸭产品、鸭苗、鸡产品、鸡苗、饲料、调理品、熟食与羽绒；控股股东益客农牧持股 70.11\%，2026 年上半年营业收入 99.02 亿元、归母净利润 3,370.62 万元实现扭亏。 公司于 2022-01-18 在深交所创业板首发上市，注册资本 4.49 亿元，总股本 4.49 亿股。 截至 2026 年 10 月 12 日收盘价 9.29 元，对应总市值 41.7 亿元，近一年-13.5\%，总市值在三级行业“肉鸡养殖”的 8 家公司中按由大到小排第\zhnumber{5}位，近一年涨跌幅低于全样本中位数（-10.4\%）3.1 个百分点。 按 2025 年年报披露的行业构成，收入占比前三位为屠宰 70.5\%；饲料 15.0\%；羽绒 5.4\%。

\pfl{股权结构与控制权} 按 2026 年 6 月 30 日的股东名册，前三大股东为益客农牧 持股 70.11\%；中央企业乡村产业投资基金股份有限公司 持股 2.99\%；宿迁久德 持股 2.16\%，股东户数 11,703 户。 控制权与股东性质方面，控股股东江苏益客农牧投资有限公司直接持股 70.11\%（3.15 亿股），其中质押 3,900.00 万股；2026 年，公司拟向特定对象发行 A 股不超过 1.35 亿股（不超发行前总股本的 30\%）。

\pfl{经营要点} 经营层面，2025 年，产品价格降幅大于成本降幅，其中鸭产品毛利同比减少约 2.20 亿元、鸭苗毛利同比减少约 2.70 亿元。

\pfl{所处行业的基本面} 白羽与黄羽肉鸡的供给取决于祖代引种量与父母代存栏，需求端受餐饮与替代效应驱动，产能周期只有 6---12 个月（第 \ref{sec:framework} 章的供需框架）。 与该公司直接相关的行业数据包括：8 月初反弹至 0.76 元/羽仍低于 1.0 元/羽孵化成本线；2026 年上半年，白羽肉鸡约占肉鸡出栏量的 68.5\%。

\pfl{竞争优势} 把该公司放回同类公司比较，2026 年上半年营业收入 99.0 亿元、在三级行业“肉鸡养殖”的 8 家公司中排第\zhnumber{2}位；毛利率 4.0\%、低于行业内中位数 6.8 个百分点；半年 ROE 2.2\%（未年化）、高于行业内中位数（1.21\%）；资产负债率 68.9\%、高于行业内中位数（52.1\%）。 相对优势方面，2026 年上半年禽肉销量约占全国禽肉产量的 6.3\%；肉鸡屠宰将新增 10,000 万只/年的屠宰产能，熟食调理品产能同步增加。 优势在当期体现为高于行业中位数的资本回报，但优势能否持续仍取决于周期与自身产能投放节奏。

\pfl{历史股价} 以 2022 年 1 月为起点、2026 年 10 月为终点，股价由 22.05 元变为 9.29 元，累计 -57.9\%，区间最高 22.05 元（2022 年 1 月）、最低 8.56 元（2026 年 6 月）；当前价相当于区间最高价的 42.1\%，为区间最低价的 1.09 倍。 月度收益的年化波动率 34.5\%，低于全样本中位数 37.3\%，在三级行业“肉鸡养殖”的 8 家公司中波动率由高到低排第\zhnumber{6}位。

\begin{center}
\begin{minipage}{\textwidth}\centering
\begin{tikzpicture}
\begin{axis}[
  width=12.6cm, height=3.9cm, scale only axis,
  xmin=2021.922, xmax=2026.972, ymin=7.96, ymax=23.59,
  xtick={2022,2023,2024,2025,2026},
  yearticks,
  yticklabel style={font=\scriptsize},
  ylabel={元/股}, ylabel style={font=\scriptsize},
  grid=major, grid style={LightGray, line width=0.3pt},
  axis line style={InkGray, line width=0.5pt},
  every axis plot/.append style={line width=0.9pt},
]
\addplot[AgriGreen] table[x=x,y=y]{data/clean/profiles/px_301116.dat};
\addplot[SoftRed, dashed, line width=0.7pt] table[x=x,y=y]{data/clean/profiles/pxzz_301116.dat};
\addplot[only marks, mark=*, mark size=1.1pt, SoftRed] table[x=x,y=y]{data/clean/profiles/pxzz_301116.dat};
\end{axis}
\end{tikzpicture}

\captionof{figure}[益客食品（301116）股价与周期骨架]{益客食品（301116）月度收盘价（元/股）与 ZigZag 周期骨架（阈值 25\%，圆点为周期转折点）}
\end{minipage}
\end{center}

\pfl{过去周期分析} 以 25\% 的 ZigZag 阈值划分，样本区间内股价形成 1 个主升段与 2 个主跌段。 最大主跌段为 2022 年 1 月 至 2024 年 6 月，29 个月-59.5\%；最大主升段出现在 2024 年 6 月 至 2024 年 11 月，5 个月+53.3\%。 最近一次转折出现在 2026 年 6 月（下跌段自 2024 年 11 月起，19 个月-37.5\%），当前处于该段的延续之中。 公司股价与申万农林牧渔指数几乎同步见底，个股并未走出独立行情。 禽链的产能周期只有 6---12 个月、反转快于生猪（第 \ref{sec:other-livestock} 节），2026 年前四个月毛鸡养殖利润同比大幅回升，链条利润的修复先于生猪，公司 2026 年 6 月末资产负债率 68.9\%，高于全样本中位数 52.2\%，其股价因此更多反映资金面与信用风险，而不是价格弹性本身。 综合区间形态看，该股单段涨跌幅度偏大、以趋势段为主（最大主升 +53.3\%、最大主跌 -59.5\%），年化波动率 34.5\% 与全样本中位数 37.3\% 相比波动与样本中位数接近。

\pfl{盈利情况} 2026 年上半年营业收入 99.02 亿元（同比 +13.2\%），归母净利润 0.34 亿元，实现扭亏。 以年报为趋势对照，2023---2025 年营业收入由 218.88 亿元变为 189.62 亿元（两年年均复合增速 -6.9\%），归母净利润由 -1.71 亿元变为 -2.80 亿元。 按半年报口径，2026 年上半年的 ROE 为 2.2\%（未年化）、销售净利率 0.4\%、毛利率 4.0\%，ROE 在“肉鸡养殖”8 家公司中按数值由高到低排第\zhnumber{4}位。 同期全样本半年 ROE 中位数 0.75\%，公司与之相差 1.43 个百分点（略好于中位数公司）。 从公司披露的原因看，业绩变动主要来自：2026 年上半年，主要系鸡产品、饲料、鸡苗、鸭苗、调理品、羽绒等产品销量、售价同比升高所致；2025 年，公司主要产品销量同比变化不大。 综合看，公司在报告期维持盈利（高于全样本半年 ROE 中位数 0.75\%），利润的可持续性取决于肉鸡养殖景气的延续与成本管控的执行。

\begin{center}\small\setlength{\tabcolsep}{3.4pt}
\captionof{table}[益客食品（301116）关键财务指标]{益客食品（301116）关键财务指标（合并报表口径；两个半年报期均未年化；现金短债比 $=$ 货币资金 $\div$ 短期债务）}
\begin{adjustbox}{max width=\textwidth}
\begin{tabular}{@{}lrrrrr@{}}
\toprule
指标 & 2023 年 & 2024 年 & 2025 年 & 2025 年半年报 & 2026 年半年报 \\
\midrule
营业收入（亿元） & 218.88 & 208.37 & 189.62 & 87.47 & 99.02 \\
归母净利润（亿元） & -1.71 & 1.00 & -2.80 & -1.81 & 0.34 \\
毛利率（\%） & 2.2 & 4.4 & 2.4 & 1.9 & 4.0 \\
ROE（\%） & -9.1 & 5.7 & -17.2 & -10.8 & 2.2 \\
资产负债率（\%） & 61.3 & 61.3 & 65.2 & 68.2 & 68.9 \\
经营活动现金流净额（亿元） & 6.33 & 9.42 & 3.36 & -2.75 & -0.70 \\
货币资金（亿元） & 4.88 & 2.54 & 2.69 & 3.98 & 4.84 \\
有息负债合计（亿元） & 19.44 & 16.33 & 19.04 & 24.45 & 23.69 \\
现金短债比（倍） & 0.41 & 0.31 & 0.19 & 0.25 & 0.28 \\
\bottomrule
\end{tabular}
\end{adjustbox}
\end{center}

\begin{center}
\begin{minipage}{\textwidth}\centering
\begin{tikzpicture}
\begin{groupplot}[
  group style={group size=2 by 1, horizontal sep=0.60cm},
  width=6.85cm, height=4.60cm,
  tick label style={font=\tiny},
  title style={font=\scriptsize\heiti},
  yticklabel style={font=\tiny},
  ylabel style={font=\tiny},
  grid=major, grid style={LightGray, line width=0.3pt},
  axis line style={InkGray, line width=0.5pt},
  enlarge x limits={abs=0.8},
  legend style={font=\scriptsize, at={(0.5,-0.20)}, anchor=north,
                legend columns=2, draw=none, fill=none, column sep=6pt},
]
\nextgroupplot[
  ybar, bar width=4pt,
  ymin=-24.97, ymax=245.48,
  xtick={0,1,2,3,4,5.7},
  xticklabels={2021,2022,2023,2024,2025,2026H1},
  ylabel={亿元},
]
\addplot[fill=AgriGreen!75, draw=AgriGreen, bar shift=-2.6pt] table[x=i, y=rev]{data/clean/profiles/fin_301116.dat};
\addplot[fill=SoftRed!60, draw=SoftRed, bar shift=2.6pt] table[x=i, y=ni]{data/clean/profiles/fin_301116.dat};
\addplot[fill=AgriGreen!30, draw=AgriGreen, pattern=north east lines, bar shift=-2.6pt] table[x=x, y=rev]{data/clean/profiles/finh_301116.dat};
\addplot[fill=SoftRed!20, draw=SoftRed, pattern=north east lines, bar shift=2.6pt] table[x=x, y=ni]{data/clean/profiles/finh_301116.dat};
\legend{营业收入（年/半年）, 归母净利润（年/半年）}
\nextgroupplot[
  ymin=-24.1, ymax=77.5,
  xtick={0,1,2,3,4,5.7},
  xticklabels={2021,2022,2023,2024,2025,2026H1},
  ylabel={\%},
]
\addplot[AgriGold, mark=*, mark size=1.3pt] table[x=i, y=roe]{data/clean/profiles/fin_301116.dat};
\addplot[OilBlue, mark=square*, mark size=1.3pt] table[x=i, y=debt]{data/clean/profiles/fin_301116.dat};
\addplot[only marks, AgriGold, mark=*, mark size=2.0pt] table[x=x, y=roe]{data/clean/profiles/finh_301116.dat};
\addplot[only marks, OilBlue, mark=square*, mark size=2.0pt] table[x=x, y=debt]{data/clean/profiles/finh_301116.dat};
\legend{ROE, 资产负债率}
\end{groupplot}
\end{tikzpicture}
\begin{tikzpicture}
\begin{axis}[
  name=finbot,
  width=13.75cm, height=3.10cm, scale only axis,
  ymin=-2.37, ymax=26.54,
  xtick={0,1,2,3,4,5.7},
  xticklabels={2021,2022,2023,2024,2025,2026H1},
  tick label style={font=\tiny},
  yticklabel style={font=\tiny},
  ylabel={亿元}, ylabel style={font=\tiny},
  ybar, bar width=4pt,
  grid=major, grid style={LightGray, line width=0.3pt},
  axis line style={InkGray, line width=0.5pt},
  enlarge x limits={abs=0.8},
  legend style={font=\scriptsize, at={(0.5,-0.26)}, anchor=north,
                legend columns=2, draw=none, fill=none, column sep=10pt},
]
\addplot[fill=AgriGreen!55, draw=AgriGreen, bar shift=-3.0pt] table[x=i, y=cash]{data/clean/profiles/fin_301116.dat};
\addplot[fill=SoftRed!45, draw=SoftRed, bar shift=3.0pt] table[x=i, y=ibd]{data/clean/profiles/fin_301116.dat};
\addplot[fill=AgriGreen!25, draw=AgriGreen, pattern=north east lines, bar shift=-3.0pt] table[x=x, y=cash]{data/clean/profiles/finh_301116.dat};
\addplot[fill=SoftRed!20, draw=SoftRed, pattern=north east lines, bar shift=3.0pt] table[x=x, y=ibd]{data/clean/profiles/finh_301116.dat};
\legend{货币资金（左轴）, 有息负债（左轴）}
\end{axis}
\begin{axis}[
  at={(finbot.south east)}, anchor=south east,
  width=13.75cm, height=3.10cm, scale only axis,
  axis y line*=right, axis x line*=none, xtick=\empty,
  ymin=-4.53, ymax=12.74,
  tick label style={font=\tiny},
  yticklabel style={font=\tiny}, scaled y ticks=false,
  legend style={font=\scriptsize, at={(0.5,-0.44)}, anchor=north,
                legend columns=1, draw=none, fill=none},
]
\addplot[OilBlue, mark=*, mark size=2.2pt, line width=1.3pt] table[x=i, y=ocf]{data/clean/profiles/fin_301116.dat};
\addplot[only marks, OilBlue, mark=*, mark size=3.2pt] table[x=x, y=ocf]{data/clean/profiles/finh_301116.dat};
\legend{经营活动现金流净额（右轴，亿元，蓝点）}
\end{axis}
\end{tikzpicture}

\captionof{figure}[益客食品（301116）近年主要财务指标]{益客食品（301116）近年主要财务指标（左上：营业收入与归母净利润；右上：ROE 与资产负债率；下：货币资金、有息负债为左轴柱状，经营活动现金流净额为其右轴的蓝点折线。
2021---2025 年为年报口径，2026H1 为 2026 年半年报/半年末，与其前一年报之间留出间隔、以斜纹或空心点区分；半年报数据未年化）}
\end{minipage}
\end{center}

\pfl{财务分析} 2026 年 6 月末资产负债率 68.9\%，较 2025 年末上升 3.6 个百分点（样本中位数 52.2\%，所处三级行业内按由低到高排第\zhnumber{8}位）。 2026 年 6 月末货币资金 4.84 亿元、短期债务（短期借款与一年内到期非流动负债）17.56 亿元、长期债务 6.14 亿元，有息负债合计 23.69 亿元，现金短债比 0.28 倍——缺口明显。 净有息负债 18.86 亿元，相当于其 2026 年上半年营业收入的 19\%。 流动比率 0.74、速动比率 0.31，短期偿债指标偏紧。 2026 年上半年经营活动现金流净额 -0.70 亿元，净利润为正但经营现金流为负，利润的现金含量偏低。 综合看，账面货币资金对有息负债与短期债务的覆盖不足，滚动偿付对新增融资的依赖度较高；资产负债率高于样本中位数 16.7 个百分点；有息负债中短期债务占比更高，期限结构仍以滚动续作为主。

\pfl{重大战略转型} 从公开披露看，公司的战略推进集中在：融资与资本运作（2 项）：2026 年 6 月，控股股东拟认购 5,000.00 万元至 1.50 亿元；2024 年 2 月，控股股东、实际控制人及家属、董事为公司及子公司银行综合授信提供关联担保；其他动作（1 项）：2026 年上半年，业绩扭亏主要系鸡产品、鸭产品、饲料、鸡苗、调理品、熟食、羽绒等主要产品同比毛利升高。 公告口径上，近三年（2023 年 9 月---2026 年 9 月）公司披露重大重组与并购类公告 0 条、对外投资与转型类公告 14 条、回购与股权激励类公告 23 条，合计公告 488 条。 第一大行业仍是“屠宰”，收入占比 70.5\%，与 2021 年年报相比没有方向性变化。 综合判断，报告期内公司有 14 条与战略相关的公告落地（重大重组与并购 0 条、对外投资与转型 14 条），资本运作与主业调整之间存在直接关联。

\begin{center}\small\setlength{\tabcolsep}{4pt}
\captionof{table}[益客食品（301116）历史融资与资本运作]{益客食品（301116）历史融资与资本运作（据公司公告与公开报道整理；金额为披露值，含首次公开发行、再融资、债券、银行借款与重整投资等）}
\begin{adjustbox}{max width=\textwidth}
\begin{tabular}{@{}p{1.45cm}p{2.55cm}p{4.05cm}p{4.95cm}@{}}
\toprule
时间 & 方式 & 规模 & 用途与说明 \\
\midrule
2022-01 & IPO（深圳证券交易所创业板） & 发行 4,\allowbreak{}489.80 万股、\allowbreak{}发行价 11.40 元/\allowbreak{}股，\allowbreak{}募集资金总额 5.12 亿元 & 肉鸭屠宰线建设、\allowbreak{}禽肉熟食扩建、\allowbreak{}供应链数字化、\allowbreak{}专用饲料研发及产业化项目及补充流动资金；截至 2023 年 6 月 30 日公司已使用募集资金 4.44 亿元 \\
2023-04-07 & 向特定对象发行股票（向控股股东益客农牧，后终止） & 拟募集资金不超过 4.50 亿元（发行不超过 3,\allowbreak{}955.69 万股 & 扣除发行费用后全部用于补充流动资金；2024 年 1 月获中国证监会同意注册批复；2024 年 9 月 11 日公司公告因外部客观环境和市场环境发生变化 \\
2026-06-10 & 向特定对象发行股票（预案，含控股股东认购） & 拟募集资金不超过 9.00 亿元（发行股票数量不超过 1.35 亿股） & 益客食品科技谷项目、\allowbreak{}肉禽产业链综合研发项目、\allowbreak{}数智化升级建设项目、\allowbreak{}补充流动资金项目；发行对象为包括益客农牧在内的不超过 35 名特定对象；公司于 2026 年 6 月 10 日与益客农牧签署附条件生效的股份认购协议；2026 年 8 月 21 日发行申请获深交所受理 \\
\bottomrule
\end{tabular}
\end{adjustbox}
\end{center}

\pfl{历史融投资情况} 从现金流量表看，2025 年公司取得借款收到的现金 18.21 亿元、偿还债务支付的现金 14.96 亿元、吸收权益性投资 0.03 亿元，筹资活动现金流净额 1.52 亿元（2023 年为取得借款 12.26 亿元、偿还债务 10.89 亿元）。 2026 年上半年筹资活动现金流净额为 2.35 亿元（取得借款 12.07 亿元、偿还债务 7.76 亿元，未年化），仍处于净融入状态。 2025 年分配股利、利润或偿付利息支付的现金合计 1.37 亿元。 公开披露的融资记录共 3 笔，工具类型以 IPO、定向增发为主；金额与用途见下表。 其中规模最大的两笔为：2022 年 1 月，IPO，发行 4,489.80 万股、发行价 11.40 元/股，募集资金总额 5.12 亿元；2026 年 6 月，定向增发，拟募集资金不超过 9.00 亿元（发行股票数量不超过 1.35 亿股）。 股本方面，近三年披露股本变动 5 次（限售股份上市 1 次、其他 1 次）；总股本由 2021 年末的 4.04 亿股变为 4.49 亿股（+11.11\%）。 分红方面，近三年现金分红 3 次、合计每股派现 0.20 元，最近一次实施在 2024 年 12 月。 近三年“再融资”类公告 56 条，涉及定向增发、募集资金使用。 综合看，公司融资结构上以债务滚动为主、股权融资为补充；2025 年筹资活动现金流净额 1.52 亿元，年度内仍为净融入，与前述经营与投资现金流的方向基本一致。

\pfl{未来潜在融资需求分析} 2026 年 6 月末货币资金 4.84 亿元，而短期债务 17.56 亿元（现金短债比 0.28 倍），2026 年上半年经营现金流净额 -0.70 亿元。按“短期债务减去账面货币资金”的滚动偿付口径，静态缺口约 12.72 亿元；上半年盈利或接近盈亏平衡，该缺口不随经营亏损进一步扩大。 公司 2026 年 6 月末资产负债率 68.9\%、2026 年上半年 ROE 2.2\%（未年化），资产负债率偏高（68.9\%），融资需求首先用于置换与滚动存量债务、补充营运资金，属于“补血型”而非扩张型。 公开披露的后续资金安排线索包括：股权融资线索：2026 年 6 月，公司拟向包括控股股东益客农牧在内的不超过 35 名特定对象发行股票募资不超过 9.00 亿元；2026 年度，公司向特定对象发行 A 股股票申请获深圳证券交易所受理（募集说明书申报稿同日披露）。 综合看，公司的外部资金需求以补充营运资金为主、项目投入为辅，滚动偿付口径下的静态缺口约 12.72 亿元，融资节奏将受制于银行续贷意愿与股权融资窗口。




\subsubsection{天山生物（300313）}
\label{co:300313}

\pfl{公司基本情况} 天山生物成立于 2003 年，新疆昌吉的肉牛全产业链畜牧生物企业，以牛品种改良、冻精与活畜养殖起家；因 23.73 亿元跨界并购大象广告遭合同诈骗巨亏，2025 年 4 月被实施退市风险警示，2026 年 6 月摘帽并恢复“天山生物”简称。 公司于 2012-04-25 在深交所创业板首发上市，注册资本 2.40 亿元，总股本 3.13 亿股。 截至 2026 年 10 月 12 日收盘价 12.08 元，对应总市值 37.8 亿元，近一年+36.6\%，总市值在三级行业“其他养殖”的 3 家公司中按由大到小排第\zhnumber{2}位，近一年涨跌幅高于全样本中位数（-10.4\%）47.0 个百分点。 按 2026 年半年报披露的产品构成，收入占比前三位为肉类 60.1\%；饲料 13.7\%；活畜 11.6\%。 与 2021 年年报相比，第一大产品口径由“活畜”（63.8\%）变为“肉类”（60.1\%）；两期披露的分类口径有调整，占比差异中同时包含分类因素。 股权结构方面，控股股东为湖州皓辉企业管理咨询有限公司。

\pfl{股权结构与控制权} 按 2026 年 6 月 30 日的股东名册，前三大股东为厦门舍德供应链管理有限公司 持股 22.46\%；陈德宏 持股 14.30\%；新疆畜牧业集团有限公司 持股 13.77\%，股东户数 9,916 户。 控制权与股东性质方面，第二大股东陈德宏持股 14.30\%（3,427.91 万股），其中 3,350.98 万股质押、3,427.91 万股冻结。

\pfl{经营要点} 经营层面，2026 年上半年，主营业务收入 2.17 亿元；2025 年，实际控制人变更后，控股股东厦门舍德向公司提供无息借款 7,000.00 万元。

\pfl{所处行业的基本面} 牛、羊的供给由母畜存栏与出栏节奏决定、扩张周期长，水产则由投苗量与消费需求决定（第 \ref{sec:framework} 章的供需框架）。 从供需两端的当期数据看，全国奶牛存栏降至 577.2 万头；2026 年上半年，全国肉牛出栏 2,119 万头、同比减少 83 万头、下降 3.8\%；2026 年，平均价格 9 美元/剂、同比上涨 43.3\%。

\pfl{竞争优势} 把该公司放回同类公司比较，2026 年上半年营业收入 2.2 亿元、在三级行业“其他养殖”的 3 家公司中排第\zhnumber{3}位；毛利率 9.1\%、与行业内中位数基本持平；半年 ROE 12.1\%（未年化）、高于行业内中位数（0.83\%）；资产负债率 82.0\%、高于行业内中位数（66.2\%）。 优势在当期体现为高于行业中位数的资本回报，但优势能否持续仍取决于周期与自身产能投放节奏。

\pfl{历史股价} 自 2021 年 1 月至 2026 年 10 月，股价由 14.95 元变为 12.08 元，累计 -19.2\%，区间最高 15.12 元（2021 年 2 月）、最低 4.76 元（2022 年 4 月）；当前价相当于区间最高价的 79.9\%，为区间最低价的 2.54 倍。 月度收益的年化波动率 70.1\%，高于全样本中位数 37.3\%，在三级行业“其他养殖”的 3 家公司中波动率由高到低排第\zhnumber{1}位。 把股价阶段与公司事件对齐看，2025 年 4 月 至 2026 年 8 月股价上涨+192.9\%，同期（2026 年 4 月）公司向深圳证券交易所申请撤销对公司股票交易实施的退市风险警示及其他风险警示；2023 年 5 月 至 2023 年 10 月股价上涨+162.6\%，同期（2023 年 10 月）公司拟协议转让 4,071.84 万股。 从时间对应关系看，公司层面的资本运作与股权、风险事件解释了区间内多数急涨急跌，与行业指数的同步性相对有限。

\begin{center}
\begin{minipage}{\textwidth}\centering
\begin{tikzpicture}
\begin{axis}[
  width=12.6cm, height=3.9cm, scale only axis,
  xmin=2020.922, xmax=2026.972, ymin=4.43, ymax=16.18,
  xtick={2021,2022,2023,2024,2025,2026},
  yearticks,
  yticklabel style={font=\scriptsize},
  ylabel={元/股}, ylabel style={font=\scriptsize},
  grid=major, grid style={LightGray, line width=0.3pt},
  axis line style={InkGray, line width=0.5pt},
  every axis plot/.append style={line width=0.9pt},
]
\addplot[AgriGreen] table[x=x,y=y]{data/clean/profiles/px_300313.dat};
\addplot[SoftRed, dashed, line width=0.7pt] table[x=x,y=y]{data/clean/profiles/pxzz_300313.dat};
\addplot[only marks, mark=*, mark size=1.1pt, SoftRed] table[x=x,y=y]{data/clean/profiles/pxzz_300313.dat};
\end{axis}
\end{tikzpicture}

\captionof{figure}[天山生物（300313）股价与周期骨架]{天山生物（300313）月度收盘价（元/股）与 ZigZag 周期骨架（阈值 25\%，圆点为周期转折点）}
\end{minipage}
\end{center}

\pfl{过去周期分析} 以 25\% 的 ZigZag 阈值划分，样本区间内股价形成 4 个主升段与 4 个主跌段。 最大主跌段为 2023 年 10 月 至 2025 年 4 月，18 个月-61.4\%；最大主升段出现在 2025 年 4 月 至 2026 年 8 月，16 个月+192.9\%。 最近一次转折出现在 2026 年 8 月（上涨段自 2025 年 4 月起，16 个月+192.9\%），当前处于该段之后的震荡之中。 公司股价较申万农林牧渔指数提前约 49 个月见底，是板块中较早定价周期反转的公司之一。 牛、羊、水产的产能周期长于生猪，第 \ref{sec:other-livestock} 章判断肉牛供给的实质性收缩要到 2027 年之后才能兑现，这一时滞决定了个股的股价弹性，公司 2026 年 6 月末资产负债率 82.0\%，高于全样本中位数 52.2\%，其股价因此更多反映资金面与信用风险，而不是价格弹性本身。 综合区间形态看，该股单段涨跌幅度偏大、以趋势段为主（最大主升 +192.9\%、最大主跌 -61.4\%），年化波动率 70.1\% 与全样本中位数 37.3\% 相比波动明显大于样本中位数。

\pfl{盈利情况} 2026 年上半年营业收入 2.20 亿元（同比 +427.7\%），归母净利润 0.04 亿元，实现扭亏。 以年报为趋势对照，2023---2025 年营业收入由 1.37 亿元变为 2.00 亿元（两年年均复合增速 +20.7\%），归母净利润由 -0.22 亿元变为 -0.27 亿元。 按半年报口径，2026 年上半年的 ROE 为 12.1\%（未年化）、销售净利率 2.4\%、毛利率 9.1\%，ROE 在“其他养殖”3 家公司中按数值由高到低排第\zhnumber{1}位。 同期全样本半年 ROE 中位数 0.75\%，公司与之相差 11.34 个百分点（略好于中位数公司）。 综合看，公司在报告期维持盈利（高于全样本半年 ROE 中位数 0.75\%），利润的可持续性取决于其他养殖景气的延续与成本管控的执行。

\begin{center}\small\setlength{\tabcolsep}{3.4pt}
\captionof{table}[天山生物（300313）关键财务指标]{天山生物（300313）关键财务指标（合并报表口径；两个半年报期均未年化；现金短债比 $=$ 货币资金 $\div$ 短期债务）}
\begin{adjustbox}{max width=\textwidth}
\begin{tabular}{@{}lrrrrr@{}}
\toprule
指标 & 2023 年 & 2024 年 & 2025 年 & 2025 年半年报 & 2026 年半年报 \\
\midrule
营业收入（亿元） & 1.37 & 1.38 & 2.00 & 0.42 & 2.20 \\
归母净利润（亿元） & -0.22 & -0.66 & -0.27 & -0.08 & 0.04 \\
毛利率（\%） & 17.0 & -5.2 & 7.4 & 18.4 & 9.1 \\
ROE（\%） & -30.1 & 0.0 & --- & 0.0 & 12.1 \\
资产负债率（\%） & 77.2 & 97.7 & 84.1 & 100.8 & 82.0 \\
经营活动现金流净额（亿元） & -0.14 & 0.14 & 0.02 & -0.03 & -0.11 \\
货币资金（亿元） & 0.17 & 0.25 & 0.39 & 0.19 & 0.26 \\
有息负债合计（亿元） & 0.86 & 0.84 & 0.38 & 0.83 & 0.93 \\
现金短债比（倍） & 0.20 & 0.48 & 18.90 & 0.38 & 0.45 \\
\bottomrule
\end{tabular}
\end{adjustbox}
\end{center}

\begin{center}
\begin{minipage}{\textwidth}\centering
\begin{tikzpicture}
\begin{groupplot}[
  group style={group size=2 by 1, horizontal sep=0.60cm},
  width=6.85cm, height=4.60cm,
  tick label style={font=\tiny},
  title style={font=\scriptsize\heiti},
  yticklabel style={font=\tiny},
  ylabel style={font=\tiny},
  grid=major, grid style={LightGray, line width=0.3pt},
  axis line style={InkGray, line width=0.5pt},
  enlarge x limits={abs=0.8},
  legend style={font=\scriptsize, at={(0.5,-0.20)}, anchor=north,
                legend columns=2, draw=none, fill=none, column sep=6pt},
]
\nextgroupplot[
  ybar, bar width=4pt,
  ymin=-0.95, ymax=2.54,
  xtick={0,1,2,3,4,5.7},
  xticklabels={2021,2022,2023,2024,2025,2026H1},
  ylabel={亿元},
]
\addplot[fill=AgriGreen!75, draw=AgriGreen, bar shift=-2.6pt] table[x=i, y=rev]{data/clean/profiles/fin_300313.dat};
\addplot[fill=SoftRed!60, draw=SoftRed, bar shift=2.6pt] table[x=i, y=ni]{data/clean/profiles/fin_300313.dat};
\addplot[fill=AgriGreen!30, draw=AgriGreen, pattern=north east lines, bar shift=-2.6pt] table[x=x, y=rev]{data/clean/profiles/finh_300313.dat};
\addplot[fill=SoftRed!20, draw=SoftRed, pattern=north east lines, bar shift=2.6pt] table[x=x, y=ni]{data/clean/profiles/finh_300313.dat};
\legend{营业收入（年/半年）, 归母净利润（年/半年）}
\nextgroupplot[
  ymin=-44.3, ymax=110.9,
  xtick={0,1,2,3,4,5.7},
  xticklabels={2021,2022,2023,2024,2025,2026H1},
  ylabel={\%},
]
\addplot[AgriGold, mark=*, mark size=1.3pt] table[x=i, y=roe]{data/clean/profiles/fin_300313.dat};
\addplot[OilBlue, mark=square*, mark size=1.3pt] table[x=i, y=debt]{data/clean/profiles/fin_300313.dat};
\addplot[only marks, AgriGold, mark=*, mark size=2.0pt] table[x=x, y=roe]{data/clean/profiles/finh_300313.dat};
\addplot[only marks, OilBlue, mark=square*, mark size=2.0pt] table[x=x, y=debt]{data/clean/profiles/finh_300313.dat};
\legend{ROE, 资产负债率}
\end{groupplot}
\end{tikzpicture}
\begin{tikzpicture}
\begin{axis}[
  name=finbot,
  width=13.75cm, height=3.10cm, scale only axis,
  ymin=-0.15, ymax=1.63,
  xtick={0,1,2,3,4,5.7},
  xticklabels={2021,2022,2023,2024,2025,2026H1},
  tick label style={font=\tiny},
  yticklabel style={font=\tiny},
  ylabel={亿元}, ylabel style={font=\tiny},
  ybar, bar width=4pt,
  grid=major, grid style={LightGray, line width=0.3pt},
  axis line style={InkGray, line width=0.5pt},
  enlarge x limits={abs=0.8},
  legend style={font=\scriptsize, at={(0.5,-0.26)}, anchor=north,
                legend columns=2, draw=none, fill=none, column sep=10pt},
]
\addplot[fill=AgriGreen!55, draw=AgriGreen, bar shift=-3.0pt] table[x=i, y=cash]{data/clean/profiles/fin_300313.dat};
\addplot[fill=SoftRed!45, draw=SoftRed, bar shift=3.0pt] table[x=i, y=ibd]{data/clean/profiles/fin_300313.dat};
\addplot[fill=AgriGreen!25, draw=AgriGreen, pattern=north east lines, bar shift=-3.0pt] table[x=x, y=cash]{data/clean/profiles/finh_300313.dat};
\addplot[fill=SoftRed!20, draw=SoftRed, pattern=north east lines, bar shift=3.0pt] table[x=x, y=ibd]{data/clean/profiles/finh_300313.dat};
\legend{货币资金（左轴）, 有息负债（左轴）}
\end{axis}
\begin{axis}[
  at={(finbot.south east)}, anchor=south east,
  width=13.75cm, height=3.10cm, scale only axis,
  axis y line*=right, axis x line*=none, xtick=\empty,
  ymin=-0.48, ymax=0.29,
  tick label style={font=\tiny},
  yticklabel style={font=\tiny}, scaled y ticks=false,
  legend style={font=\scriptsize, at={(0.5,-0.44)}, anchor=north,
                legend columns=1, draw=none, fill=none},
]
\addplot[OilBlue, mark=*, mark size=2.2pt, line width=1.3pt] table[x=i, y=ocf]{data/clean/profiles/fin_300313.dat};
\addplot[only marks, OilBlue, mark=*, mark size=3.2pt] table[x=x, y=ocf]{data/clean/profiles/finh_300313.dat};
\legend{经营活动现金流净额（右轴，亿元，蓝点）}
\end{axis}
\end{tikzpicture}

\captionof{figure}[天山生物（300313）近年主要财务指标]{天山生物（300313）近年主要财务指标（左上：营业收入与归母净利润；右上：ROE 与资产负债率；下：货币资金、有息负债为左轴柱状，经营活动现金流净额为其右轴的蓝点折线。
2021---2025 年为年报口径，2026H1 为 2026 年半年报/半年末，与其前一年报之间留出间隔、以斜纹或空心点区分；半年报数据未年化）}
\end{minipage}
\end{center}

\pfl{财务分析} 2026 年 6 月末资产负债率 82.0\%，较 2025 年末下降 2.0 个百分点（样本中位数 52.2\%，所处三级行业内按由低到高排第\zhnumber{3}位）。 2026 年 6 月末货币资金 0.26 亿元、短期债务（短期借款与一年内到期非流动负债）0.58 亿元、长期债务 0.35 亿元，有息负债合计 0.93 亿元，现金短债比 0.45 倍——覆盖偏紧。 净有息负债 0.66 亿元，相当于其 2026 年上半年营业收入的 30\%。 流动比率 1.01、速动比率 0.60，短期指标尚可。 2026 年上半年经营活动现金流净额 -0.11 亿元，净利润为正但经营现金流为负，利润的现金含量偏低。 综合看，账面货币资金对有息负债与短期债务的覆盖不足，滚动偿付对新增融资的依赖度较高；资产负债率高于样本中位数 29.8 个百分点；有息负债中短期债务占比更高，期限结构仍以滚动续作为主。 公司披露的退市与持续经营风险为：2026 年 4 月，公司向深圳证券交易所申请撤销对公司股票交易实施的退市风险警示及其他风险警示；2024 年末，归属于上市公司股东的净资产为-377.93 万元。

\pfl{重大战略转型} 从公开披露看，公司的战略推进集中在：并购与重组（2 项）：2025 年 9 月，标的股份与债权转让价款合计 2.87 亿元，并通过表决权委托取得控制权；2023 年 10 月，公司拟协议转让 4,071.84 万股。 公告口径上，近三年（2023 年 9 月---2026 年 9 月）公司披露重大重组与并购类公告 2 条、对外投资与转型类公告 3 条、回购与股权激励类公告 2 条，合计公告 414 条。 主营结构的口径变化已经落到报表上：2021 年年报的第一大产品为“活畜”（63.8\%），2026 年半年报为“肉类”（60.1\%）；公司在这两期的分部划分方式有过调整。 综合判断，报告期内公司有 5 条与战略相关的公告落地（重大重组与并购 2 条、对外投资与转型 3 条），资本运作与主业调整之间存在直接关联。

\begin{center}\small\setlength{\tabcolsep}{4pt}
\captionof{table}[天山生物（300313）历史融资与资本运作]{天山生物（300313）历史融资与资本运作（据公司公告与公开报道整理；金额为披露值，含首次公开发行、再融资、债券、银行借款与重整投资等）}
\begin{adjustbox}{max width=\textwidth}
\begin{tabular}{@{}p{1.45cm}p{2.55cm}p{4.05cm}p{4.95cm}@{}}
\toprule
时间 & 方式 & 规模 & 用途与说明 \\
\midrule
2012-04 & IPO（深交所创业板） & 发行 2,\allowbreak{}273 万股，\allowbreak{}发行价 13.00 元/\allowbreak{}股，\allowbreak{}募集资金总额 2.95 亿元 & 募投项目建设；2012 年年度报告披露口径 \\
2016-06 & 发行股份购买资产并募集配套资金（证监许可〔2015〕2393号） & 募集配套资金 3,\allowbreak{}900.00 万元 & 配套募集资金用途相关项目建设；配套资金发行 240.15 万股 \\
2017-09 & 发行股份及支付现金购买资产（收购大象广告96.21\%股权） & 交易作价 23.73 亿元，\allowbreak{}其中现金对价 5.77 亿元、\allowbreak{}股份对价 17.96 亿元 & 跨界并购户外广告资产；该次交易的会计处理于 2023 年年报中作追溯调整 \\
2026-03 & 向特定对象发行A股股票（预案） & 发行数量不超过 5,\allowbreak{}900 万股，\allowbreak{}发行价格 5.91 元/\allowbreak{}股 & 扣除发行费用后拟全部用于补充流动资金；发行对象为公司实际控制人陈明艺（不超过 900 万股）和控股股东厦门舍德（不超过 5,\allowbreak{}000 万股） \\
\bottomrule
\end{tabular}
\end{adjustbox}
\end{center}

\pfl{历史融投资情况} 从现金流量表看，2025 年公司偿还债务支付的现金 0.50 亿元、吸收权益性投资 0.02 亿元，筹资活动现金流净额 0.17 亿元（2023 年为取得借款 1.07 亿元、偿还债务 1.18 亿元）。 2026 年上半年筹资活动现金流净额为 0.01 亿元（取得借款 — 亿元、偿还债务 0.01 亿元，未年化），仍处于净融入状态。 2025 年分配股利、利润或偿付利息支付的现金合计 0.05 亿元。 公开披露的融资记录共 4 笔，工具类型以 IPO、发行股份及支付现金购、定向增发为主；金额与用途见下表。 其中规模最大的两笔为：2016 年 6 月，募集配套资金 3,900.00 万元；2012 年 4 月，IPO，发行 2,273 万股，发行价 13.00 元/股，募集资金总额 2.95 亿元。 股本方面，近三年披露股本变动 3 次；总股本自 2021 年末以来未发生变化（3.13 亿股）。 分红方面，近三年未实施现金分红，在全部样本中属于分红能力偏弱的一端。 近三年“再融资”类公告 35 条，涉及定向增发、发行股份购买资产、募集资金使用。 综合看，公司融资结构上以债务滚动为主、股权融资为补充；2025 年筹资活动现金流净额 0.17 亿元，年度内仍为净融入，与前述经营与投资现金流的方向基本一致。

\pfl{未来潜在融资需求分析} 2026 年 6 月末货币资金 0.26 亿元，而短期债务 0.58 亿元（现金短债比 0.45 倍），2026 年上半年经营现金流净额 -0.11 亿元。账面货币资金可以覆盖短期债务，缺口不体现在静态偿付层面。 公司 2026 年 6 月末资产负债率 82.0\%、2026 年上半年 ROE 12.1\%（未年化），资产负债率偏高（82.0\%），融资需求首先用于置换与滚动存量债务、补充营运资金，属于“补血型”而非扩张型。 公开披露的后续资金安排线索包括：股权融资线索：2026 年 3 月，公司拟募集资金不超过 3.49 亿元全部用于补充流动资金；2026 年 6 月，董事会审议通过 2026 年度向特定对象发行 A 股股票预案（修订稿）等议案；债务与授信安排：2025 年 4 月，不得不从控股股东借款维持经营。 综合看，公司在静态偿付层面不存在缺口，外部融资更多是配合扩张节奏与并购整合的主动性安排，而非偿债压力驱动。




\subsubsection{东瑞股份（001201）}
\label{co:001201}

\pfl{公司基本情况} 广东河源的供港澳活大猪供应商与粤港澳大湾区“菜篮子”重要生产基地，2002 年成立、2021 年 4 月深交所主板上市（发行价 63.38 元、募资总额 20.07 亿元），形成饲料、育种、扩繁、养殖、屠宰冷链一体化生猪全产业链，2026 年上半年出栏生猪 88.60 万头但归母净亏损 4.25 亿元。 公司于 2021-04-28 在深交所主板首发上市，注册资本 2.58 亿元，总股本 2.58 亿股。 截至 2026 年 10 月 12 日收盘价 14.40 元，对应总市值 37.1 亿元，近一年-9.2\%，总市值在三级行业“生猪养殖”的 11 家公司中按由大到小排第\zhnumber{10}位，近一年涨跌幅高于全样本中位数（-10.4\%）1.1 个百分点。 按 2026 年半年报披露的行业构成，收入占比前三位为养殖行业 95.3\%；饲料加工业 2.7\%；屠宰、肉食业务 1.7\%。

\pfl{股权结构与控制权} 按 2026 年 6 月 30 日的股东名册，前三大股东为袁建康 持股 16.65\%；东莞市东晖实业投资有限公司 持股 10.38\%；东莞市安夏实业投资有限公司 持股 5.68\%，股东户数 20,108 户。 股权结构的另一条线索是，持股 24.96\%；2021 年 4 月，在深交所主板上市，实际募资 20.07 亿元。

\pfl{经营要点} 从经营要点看，2026 年上半年，销售生猪 88.60 万头、同比增长 15.82\%；2026 年 7 月，公司完全成本 13.8 元/公斤。

\pfl{所处行业的基本面} 生猪的供给由能繁母猪存栏决定、需求由猪肉消费与替代品价格决定，本轮下行的成因在供给端的产能去化不彻底，而非需求收缩（第 \ref{sec:framework} 章的供需框架）。 行业层面的最新情况是：2026 年二季度，末全国能繁母猪存栏 3780 万头、同比下降 6.5\%。

\pfl{竞争优势} 把该公司放回同类公司比较，2026 年上半年营业收入 10.4 亿元、在三级行业“生猪养殖”的 11 家公司中排第\zhnumber{10}位；毛利率 -6.6\%、低于行业内中位数 5.0 个百分点；半年 ROE -14.3\%（未年化）、低于行业内中位数（-14.32\%）；资产负债率 52.9\%、低于行业内中位数（58.9\%）。 相对优势方面，是农业产业化国家重点龙头企业、国家生猪核心育种场、国家级生猪产能调控基地。 就当期报表看，公司尚未把经营层面的特征转化为高于行业的资本回报，周期底部的相对优势更多体现在现金流与资产负债表上。

\pfl{历史股价} 自 2021 年 4 月至 2026 年 10 月，股价由 73.93 元变为 14.40 元，累计 -80.5\%，区间最高 73.93 元（2021 年 4 月）、最低 10.86 元（2026 年 6 月）；当前价相当于区间最高价的 19.5\%，为区间最低价的 1.33 倍。 月度收益的年化波动率 37.9\%，高于全样本中位数 37.3\%，在三级行业“生猪养殖”的 11 家公司中波动率由高到低排第\zhnumber{6}位。 把股价阶段与公司事件对齐看，2024 年 5 月 至 2026 年 6 月股价下跌-54.4\%，同期 2026 年上半年，受生猪价格下行影响计提各类减值准备（其中生物资产减值约 1.18 亿元）。 把这些阶段与公司事件对齐后可以看到，几轮大涨大跌多由公司自身的资本运作与风险事件触发，行业景气只解释了其中一部分。

\begin{center}
\begin{minipage}{\textwidth}\centering
\begin{tikzpicture}
\begin{axis}[
  width=12.6cm, height=3.9cm, scale only axis,
  xmin=2021.172, xmax=2026.972, ymin=10.10, ymax=79.11,
  xtick={2021,2022,2023,2024,2025,2026},
  yearticks,
  yticklabel style={font=\scriptsize},
  ylabel={元/股}, ylabel style={font=\scriptsize},
  grid=major, grid style={LightGray, line width=0.3pt},
  axis line style={InkGray, line width=0.5pt},
  every axis plot/.append style={line width=0.9pt},
]
\addplot[AgriGreen] table[x=x,y=y]{data/clean/profiles/px_001201.dat};
\addplot[SoftRed, dashed, line width=0.7pt] table[x=x,y=y]{data/clean/profiles/pxzz_001201.dat};
\addplot[only marks, mark=*, mark size=1.1pt, SoftRed] table[x=x,y=y]{data/clean/profiles/pxzz_001201.dat};
\end{axis}
\end{tikzpicture}

\captionof{figure}[东瑞股份（001201）股价与周期骨架]{东瑞股份（001201）月度收盘价（元/股）与 ZigZag 周期骨架（阈值 25\%，圆点为周期转折点）}
\end{minipage}
\end{center}

\pfl{过去周期分析} 以 25\% 的 ZigZag 阈值划分，样本区间内股价形成 3 个主升段与 3 个主跌段。 最大主升段出现在 2024 年 2 月 至 2024 年 5 月，3 个月+33.6\%；最大主跌段为 2022 年 3 月 至 2024 年 2 月，23 个月-59.2\%。 最近一次转折出现在 2026 年 10 月（上涨段自 2026 年 6 月起，4 个月+32.6\%），当前处于该段之后的震荡之中。 公司股价与申万农林牧渔指数几乎同步见底，个股并未走出独立行情。 生猪养殖的定价主线是产能去化与成本曲线的赛跑（第 \ref{sec:hog} 章）：2026 年三季度行业仍处在“产能去化未完成、价格磨底”的阶段，公司 2026 年 6 月末资产负债率 52.9\%、半年 ROE -14.3\%（未年化），在本轮下行中属于随行业同步调整的一类。 综合区间形态看，该股单段涨跌幅度偏大、以趋势段为主（最大主升 +33.6\%、最大主跌 -59.2\%），年化波动率 37.9\% 与全样本中位数 37.3\% 相比波动与样本中位数接近。

\pfl{盈利情况} 2026 年上半年营业收入 10.39 亿元（同比 -5.1\%），归母净利润 -4.25 亿元，由上年同期的盈利转为亏损。 以年报为趋势对照，2023---2025 年营业收入由 10.37 亿元变为 21.35 亿元（两年年均复合增速 +43.5\%），归母净利润由 -5.16 亿元变为 -1.61 亿元。 按半年报口径，2026 年上半年的 ROE 为 -14.3\%（未年化）、销售净利率 -40.9\%、毛利率 -6.6\%，ROE 在“生猪养殖”11 家公司中按数值由高到低排第\zhnumber{6}位。 同期全样本半年 ROE 中位数 0.75\%，公司与之相差 15.07 个百分点（弱于中位数公司）。 针对上述压力，公司披露的应对举措包括：2024 年 12 月，公司终止部分募集资金投资项目并将募集资金永久补充流动资金——终止“惠州东瑞多层楼房生猪养殖项目”。 综合看，公司仍处亏损区间、由盈转亏，半年 ROE 在三级行业内按由高到低排第\zhnumber{6}位，单位盈利能力的修复依赖行业景气与自身成本端的改善，报表弹性主要体现为价格与销量的方向性变化。

\begin{center}\small\setlength{\tabcolsep}{3.4pt}
\captionof{table}[东瑞股份（001201）关键财务指标]{东瑞股份（001201）关键财务指标（合并报表口径；两个半年报期均未年化；现金短债比 $=$ 货币资金 $\div$ 短期债务）}
\begin{adjustbox}{max width=\textwidth}
\begin{tabular}{@{}lrrrrr@{}}
\toprule
指标 & 2023 年 & 2024 年 & 2025 年 & 2025 年半年报 & 2026 年半年报 \\
\midrule
营业收入（亿元） & 10.37 & 15.20 & 21.35 & 10.95 & 10.39 \\
归母净利润（亿元） & -5.16 & -0.66 & -1.61 & 0.03 & -4.25 \\
毛利率（\%） & -1.4 & 21.3 & 12.3 & 17.3 & -6.6 \\
ROE（\%） & -18.1 & -1.9 & -4.9 & 0.1 & -14.3 \\
资产负债率（\%） & 42.5 & 44.1 & 46.7 & 45.0 & 52.9 \\
经营活动现金流净额（亿元） & -1.78 & 0.58 & 1.74 & 1.23 & -1.38 \\
货币资金（亿元） & 11.75 & 5.07 & 4.58 & 5.45 & 4.55 \\
有息负债合计（亿元） & 19.85 & 20.52 & 21.32 & 21.88 & 24.27 \\
现金短债比（倍） & 1.26 & 0.47 & 0.34 & 0.47 & 0.29 \\
\bottomrule
\end{tabular}
\end{adjustbox}
\end{center}

\begin{center}
\begin{minipage}{\textwidth}\centering
\begin{tikzpicture}
\begin{groupplot}[
  group style={group size=2 by 1, horizontal sep=0.60cm},
  width=6.85cm, height=4.60cm,
  tick label style={font=\tiny},
  title style={font=\scriptsize\heiti},
  yticklabel style={font=\tiny},
  ylabel style={font=\tiny},
  grid=major, grid style={LightGray, line width=0.3pt},
  axis line style={InkGray, line width=0.5pt},
  enlarge x limits={abs=0.8},
  legend style={font=\scriptsize, at={(0.5,-0.20)}, anchor=north,
                legend columns=2, draw=none, fill=none, column sep=6pt},
]
\nextgroupplot[
  ybar, bar width=4pt,
  ymin=-7.81, ymax=24.53,
  xtick={0,1,2,3,4,5.7},
  xticklabels={2021,2022,2023,2024,2025,2026H1},
  ylabel={亿元},
]
\addplot[fill=AgriGreen!75, draw=AgriGreen, bar shift=-2.6pt] table[x=i, y=rev]{data/clean/profiles/fin_001201.dat};
\addplot[fill=SoftRed!60, draw=SoftRed, bar shift=2.6pt] table[x=i, y=ni]{data/clean/profiles/fin_001201.dat};
\addplot[fill=AgriGreen!30, draw=AgriGreen, pattern=north east lines, bar shift=-2.6pt] table[x=x, y=rev]{data/clean/profiles/finh_001201.dat};
\addplot[fill=SoftRed!20, draw=SoftRed, pattern=north east lines, bar shift=2.6pt] table[x=x, y=ni]{data/clean/profiles/finh_001201.dat};
\legend{营业收入（年/半年）, 归母净利润（年/半年）}
\nextgroupplot[
  ymin=-23.8, ymax=60.0,
  xtick={0,1,2,3,4,5.7},
  xticklabels={2021,2022,2023,2024,2025,2026H1},
  ylabel={\%},
]
\addplot[AgriGold, mark=*, mark size=1.3pt] table[x=i, y=roe]{data/clean/profiles/fin_001201.dat};
\addplot[OilBlue, mark=square*, mark size=1.3pt] table[x=i, y=debt]{data/clean/profiles/fin_001201.dat};
\addplot[only marks, AgriGold, mark=*, mark size=2.0pt] table[x=x, y=roe]{data/clean/profiles/finh_001201.dat};
\addplot[only marks, OilBlue, mark=square*, mark size=2.0pt] table[x=x, y=debt]{data/clean/profiles/finh_001201.dat};
\legend{ROE, 资产负债率}
\end{groupplot}
\end{tikzpicture}
\begin{tikzpicture}
\begin{axis}[
  name=finbot,
  width=13.75cm, height=3.10cm, scale only axis,
  ymin=-2.43, ymax=27.18,
  xtick={0,1,2,3,4,5.7},
  xticklabels={2021,2022,2023,2024,2025,2026H1},
  tick label style={font=\tiny},
  yticklabel style={font=\tiny},
  ylabel={亿元}, ylabel style={font=\tiny},
  ybar, bar width=4pt,
  grid=major, grid style={LightGray, line width=0.3pt},
  axis line style={InkGray, line width=0.5pt},
  enlarge x limits={abs=0.8},
  legend style={font=\scriptsize, at={(0.5,-0.26)}, anchor=north,
                legend columns=2, draw=none, fill=none, column sep=10pt},
]
\addplot[fill=AgriGreen!55, draw=AgriGreen, bar shift=-3.0pt] table[x=i, y=cash]{data/clean/profiles/fin_001201.dat};
\addplot[fill=SoftRed!45, draw=SoftRed, bar shift=3.0pt] table[x=i, y=ibd]{data/clean/profiles/fin_001201.dat};
\addplot[fill=AgriGreen!25, draw=AgriGreen, pattern=north east lines, bar shift=-3.0pt] table[x=x, y=cash]{data/clean/profiles/finh_001201.dat};
\addplot[fill=SoftRed!20, draw=SoftRed, pattern=north east lines, bar shift=3.0pt] table[x=x, y=ibd]{data/clean/profiles/finh_001201.dat};
\legend{货币资金（左轴）, 有息负债（左轴）}
\end{axis}
\begin{axis}[
  at={(finbot.south east)}, anchor=south east,
  width=13.75cm, height=3.10cm, scale only axis,
  axis y line*=right, axis x line*=none, xtick=\empty,
  ymin=-2.70, ymax=2.80,
  tick label style={font=\tiny},
  yticklabel style={font=\tiny}, scaled y ticks=false,
  legend style={font=\scriptsize, at={(0.5,-0.44)}, anchor=north,
                legend columns=1, draw=none, fill=none},
]
\addplot[OilBlue, mark=*, mark size=2.2pt, line width=1.3pt] table[x=i, y=ocf]{data/clean/profiles/fin_001201.dat};
\addplot[only marks, OilBlue, mark=*, mark size=3.2pt] table[x=x, y=ocf]{data/clean/profiles/finh_001201.dat};
\legend{经营活动现金流净额（右轴，亿元，蓝点）}
\end{axis}
\end{tikzpicture}

\captionof{figure}[东瑞股份（001201）近年主要财务指标]{东瑞股份（001201）近年主要财务指标（左上：营业收入与归母净利润；右上：ROE 与资产负债率；下：货币资金、有息负债为左轴柱状，经营活动现金流净额为其右轴的蓝点折线。
2021---2025 年为年报口径，2026H1 为 2026 年半年报/半年末，与其前一年报之间留出间隔、以斜纹或空心点区分；半年报数据未年化）}
\end{minipage}
\end{center}

\pfl{财务分析} 2026 年 6 月末资产负债率 52.9\%，较 2025 年末上升 6.2 个百分点（样本中位数 52.2\%，所处三级行业内按由低到高排第\zhnumber{3}位）。 2026 年 6 月末货币资金 4.55 亿元、短期债务（短期借款与一年内到期非流动负债）15.69 亿元、长期债务 8.58 亿元，有息负债合计 24.27 亿元，现金短债比 0.29 倍——缺口明显。 净有息负债 19.72 亿元，相当于其 2026 年上半年营业收入的 190\%。 流动比率 0.67、速动比率 0.26，短期偿付压力较大。 2026 年上半年经营活动现金流净额 -1.38 亿元，经营现金流与利润同时为负，现金消耗较快。 2026 年 6 月末在建工程 0.68 亿元，较上年末增加 0.27 亿元，仍有在建投入。 综合看，账面货币资金对有息负债与短期债务的覆盖不足，滚动偿付对新增融资的依赖度较高；资产负债率高于样本中位数 0.7 个百分点；有息负债中短期债务占比更高，期限结构仍以滚动续作为主。

\pfl{重大战略转型} 从公开披露看，公司的战略推进集中在：其他动作（2 项）：2026 年 6 月，2025 年生猪养殖行业收入占比达 96.33\%。 公告口径上，近三年（2023 年 9 月---2026 年 9 月）公司披露重大重组与并购类公告 0 条、对外投资与转型类公告 2 条、回购与股权激励类公告 18 条，合计公告 413 条。 第一大行业仍是“养殖行业”，收入占比 95.3\%，与 2021 年年报相比没有方向性变化。 综合判断，报告期内公司有 2 条与战略相关的公告落地（重大重组与并购 0 条、对外投资与转型 2 条），资本运作与主业调整之间存在直接关联。

\begin{center}\small\setlength{\tabcolsep}{4pt}
\captionof{table}[东瑞股份（001201）历史融资与资本运作]{东瑞股份（001201）历史融资与资本运作（据公司公告与公开报道整理；金额为披露值，含首次公开发行、再融资、债券、银行借款与重整投资等）}
\begin{adjustbox}{max width=\textwidth}
\begin{tabular}{@{}p{1.45cm}p{2.55cm}p{4.05cm}p{4.95cm}@{}}
\toprule
时间 & 方式 & 规模 & 用途与说明 \\
\midrule
2021-04 & IPO（深圳证券交易所主板） & 公开发行 3,\allowbreak{}167 万股、\allowbreak{}发行价 63.38 元/\allowbreak{}股，\allowbreak{}募集资金总额 20.07 亿元 & 紫金东瑞农业发展有限公司富竹生态养殖项目、\allowbreak{}连平东瑞农牧发展有限公司传统养殖改高床养殖节能减排增效项目、\allowbreak{}和平东瑞农牧发展有限公司高床生态养殖项目等（招股说明书拟募集资金 18.10 亿元） \\
2021-05-21 & 募集资金置换自筹资金 & 以募集资金置换预先投入募投项目的自筹资金 4.40 亿元 & 置换预先投入募投项目的自筹资金；截至 2021 年 12 月 31 日，\allowbreak{}公司募集资金累计投入募投项目 11.23 亿元 \\
2023-12-28 & 向特定对象发行股票（2023 年度定增） & 发行 4,\allowbreak{}497.84 万股 & 东源东瑞农牧发展有限公司东源县船塘现代农业综合体项目（黄沙基地）及补充公司流动资金项目（原拟募集资金总额不超过 10.33 亿元）；；新增股份上市时间 2023 年 12 月 28 日 \\
\bottomrule
\end{tabular}
\end{adjustbox}
\end{center}

\pfl{历史融投资情况} 从现金流量表看，2025 年公司取得借款收到的现金 12.66 亿元、偿还债务支付的现金 12.29 亿元，筹资活动现金流净额 -0.46 亿元（2023 年为取得借款 16.35 亿元、偿还债务 6.26 亿元）。 2026 年上半年筹资活动现金流净额为 2.97 亿元（取得借款 8.78 亿元、偿还债务 5.83 亿元，未年化），仍处于净融入状态。 2025 年分配股利、利润或偿付利息支付的现金合计 0.67 亿元。 公开披露的融资记录共 3 笔，工具类型以 IPO、定向增发为主；金额与用途见下表。 其中规模最大的两笔为：2023 年 12 月，定向增发，发行 4,497.84 万股；2021 年 5 月，募集资金置换自筹资金，以募集资金置换预先投入募投项目的自筹资金 4.40 亿元。 股本方面，近三年披露股本变动 6 次（限售股份上市 2 次、增发新股上市 1 次）；总股本由 2021 年末的 1.77 亿股变为 2.58 亿股（+45.36\%）。 分红方面，近三年未实施现金分红，在全部样本中属于分红能力偏弱的一端。 近三年“再融资”类公告 43 条，涉及定向增发、募集资金使用。 综合看，公司融资结构上股权与债务融资并举；2025 年筹资活动现金流净额 -0.46 亿元，融资节奏仍以维持存量债务周转为主，与前述经营与投资现金流的方向基本一致。

\pfl{未来潜在融资需求分析} 2026 年 6 月末货币资金 4.55 亿元，而短期债务 15.69 亿元（现金短债比 0.29 倍），2026 年上半年经营现金流净额 -1.38 亿元。按“短期债务减去账面货币资金”的滚动偿付口径，静态缺口约 11.14 亿元；若再计入上半年亏损的年化额（8.50 亿元），维持一年正常经营所需的外部资金约 11.14---19.64 亿元。 公司 2026 年上半年归母净利润为负（-4.25 亿元），但 6 月末资产负债率 52.9\% 尚未进入第一档（亏损且负债率 $>$65\%）的区间，当前融资需求以补充流动资金、支撑低谷期经营性支出为主。 公开披露的后续资金安排线索包括：股权融资线索：2024 年 12 月，IPO 与 2023 年度定增的募集资金均已全部使用完毕、募投项目全部结项、专户全部注销。 资金需求还要叠加在建工程：期末在建工程 0.68 亿元、2026 年上半年购建固定资产等支出 0.89 亿元（未年化），这部分支出在周期底部通常会被主动放缓。 综合看，公司的外部资金需求以补充流动性与项目投入为主，滚动偿付口径下的静态缺口约 11.14 亿元，能否落地取决于授信续作与发行窗口的配合。




\subsubsection{晓鸣股份（300967）}
\label{co:300967}

\pfl{公司基本情况} 晓鸣股份成立于 1992 年，宁夏银川的蛋鸡制种龙头，集祖代、父母代蛋种鸡饲养与父母代、商品代雏鸡销售于一体，商品代雏鸡国内市占率约 25\%；2026 年上半年受蛋鸡周期底部拖累，营业收入同比下降约 24\%、净利润同比下降 99\%。 公司于 2021-04-13 在深交所创业板首发上市，注册资本 1.88 亿元，总股本 1.89 亿股。 截至 2026 年 10 月 12 日收盘价 15.55 元，对应总市值 29.5 亿元，近一年-30.7\%，总市值在三级行业“肉鸡养殖”的 8 家公司中按由大到小排第\zhnumber{6}位，近一年涨跌幅低于全样本中位数（-10.4\%）20.3 个百分点。 按 2025 年年报披露的行业构成，收入几乎全部来自养殖行业 100.0\%。

\pfl{股权结构与控制权} 按 2026 年 6 月 30 日的股东名册，前三大股东为魏晓明 持股 42.63\%；正大投资股份有限公司 持股 12.03\%；浙商银行股份有限公司－国泰中证畜牧养殖交易型开放式指 持股 1.09\%，股东户数 14,605 户。

\pfl{经营要点} 经营与产能方面，2026 年 8 月，销售鸡产品 2,212.68 万羽、销售收入 9,979.26 万元。

\pfl{所处行业的基本面} 白羽与黄羽肉鸡的供给取决于祖代引种量与父母代存栏，需求端受餐饮与替代效应驱动，产能周期只有 6---12 个月（第 \ref{sec:framework} 章的供需框架）。 当期该行业的关键数据是：2026 年上半年，二季度起产能去化效果显现、产品价格回升。

\pfl{竞争优势} 把该公司放回同类公司比较，2026 年上半年营业收入 5.7 亿元、在三级行业“肉鸡养殖”的 8 家公司中排第\zhnumber{8}位；毛利率 12.9\%、高于行业内中位数 2.1 个百分点；半年 ROE 0.2\%（未年化）、低于行业内中位数（1.21\%）；资产负债率 52.3\%、高于行业内中位数（52.1\%）。 就当期报表看，公司尚未把经营层面的特征转化为高于行业的资本回报，周期底部的相对优势更多体现在现金流与资产负债表上。

\pfl{历史股价} 以 2021 年 4 月为起点、2026 年 10 月为终点，股价由 18.12 元变为 15.55 元，累计 -14.2\%，区间最高 27.90 元（2021 年 10 月）、最低 9.07 元（2024 年 6 月）；当前价相当于区间最高价的 55.7\%，为区间最低价的 1.71 倍。 月度收益的年化波动率 46.7\%，高于全样本中位数 37.3\%，在三级行业“肉鸡养殖”的 8 家公司中波动率由高到低排第\zhnumber{1}位。

\begin{center}
\begin{minipage}{\textwidth}\centering
\begin{tikzpicture}
\begin{axis}[
  width=12.6cm, height=3.9cm, scale only axis,
  xmin=2021.172, xmax=2026.972, ymin=8.44, ymax=29.85,
  xtick={2021,2022,2023,2024,2025,2026},
  yearticks,
  yticklabel style={font=\scriptsize},
  ylabel={元/股}, ylabel style={font=\scriptsize},
  grid=major, grid style={LightGray, line width=0.3pt},
  axis line style={InkGray, line width=0.5pt},
  every axis plot/.append style={line width=0.9pt},
]
\addplot[AgriGreen] table[x=x,y=y]{data/clean/profiles/px_300967.dat};
\addplot[SoftRed, dashed, line width=0.7pt] table[x=x,y=y]{data/clean/profiles/pxzz_300967.dat};
\addplot[only marks, mark=*, mark size=1.1pt, SoftRed] table[x=x,y=y]{data/clean/profiles/pxzz_300967.dat};
\end{axis}
\end{tikzpicture}

\captionof{figure}[晓鸣股份（300967）股价与周期骨架]{晓鸣股份（300967）月度收盘价（元/股）与 ZigZag 周期骨架（阈值 25\%，圆点为周期转折点）}
\end{minipage}
\end{center}

\pfl{过去周期分析} 以 25\% 的 ZigZag 阈值划分，样本区间内股价形成 3 个主升段与 4 个主跌段。 最大主跌段为 2023 年 3 月 至 2024 年 6 月，15 个月-52.8\%；最大主升段出现在 2024 年 6 月 至 2025 年 8 月，14 个月+169.0\%。 最近一次转折出现在 2026 年 6 月（下跌段自 2025 年 8 月起，10 个月-41.1\%），当前处于该段的延续之中。 公司股价较申万农林牧渔指数提前约 23 个月见底，是板块中较早定价周期反转的公司之一。 禽链的产能周期只有 6---12 个月、反转快于生猪（第 \ref{sec:other-livestock} 节），2026 年前四个月毛鸡养殖利润同比大幅回升，链条利润的修复先于生猪，公司 2026 年 6 月末资产负债率 52.3\%、半年 ROE 0.2\%（未年化），在本轮下行中属于随行业同步调整的一类。 综合区间形态看，该股单段涨跌幅度偏大、以趋势段为主（最大主升 +169.0\%、最大主跌 -52.8\%），年化波动率 46.7\% 与全样本中位数 37.3\% 相比波动明显大于样本中位数。

\pfl{盈利情况} 2026 年上半年营业收入 5.68 亿元（同比 -24.5\%），归母净利润 0.02 亿元，同比 -99.0\%。 以年报为趋势对照，2023---2025 年营业收入由 8.31 亿元变为 12.21 亿元（两年年均复合增速 +21.2\%），归母净利润由 -1.55 亿元变为 0.85 亿元。 按半年报口径，2026 年上半年的 ROE 为 0.2\%（未年化）、销售净利率 0.3\%、毛利率 12.9\%，ROE 在“肉鸡养殖”8 家公司中按数值由高到低排第\zhnumber{5}位。 同期全样本半年 ROE 中位数 0.75\%，公司与之相差 0.51 个百分点（弱于中位数公司）。 综合看，公司在报告期维持盈利（低于全样本半年 ROE 中位数 0.75\%），利润的可持续性取决于肉鸡养殖景气的延续与成本管控的执行。

\begin{center}\small\setlength{\tabcolsep}{3.4pt}
\captionof{table}[晓鸣股份（300967）关键财务指标]{晓鸣股份（300967）关键财务指标（合并报表口径；两个半年报期均未年化；现金短债比 $=$ 货币资金 $\div$ 短期债务）}
\begin{adjustbox}{max width=\textwidth}
\begin{tabular}{@{}lrrrrr@{}}
\toprule
指标 & 2023 年 & 2024 年 & 2025 年 & 2025 年半年报 & 2026 年半年报 \\
\midrule
营业收入（亿元） & 8.31 & 9.71 & 12.21 & 7.52 & 5.68 \\
归母净利润（亿元） & -1.55 & 0.45 & 0.85 & 1.85 & 0.02 \\
毛利率（\%） & 0.9 & 16.5 & 22.0 & 34.0 & 12.9 \\
ROE（\%） & -20.0 & 5.9 & 10.8 & 21.3 & 0.2 \\
资产负债率（\%） & 54.4 & 53.6 & 48.9 & 45.4 & 52.3 \\
经营活动现金流净额（亿元） & 0.61 & 3.06 & 2.40 & 2.24 & 1.40 \\
货币资金（亿元） & 0.71 & 1.80 & 1.18 & 1.96 & 1.43 \\
有息负债合计（亿元） & 7.08 & 7.18 & 6.06 & 6.04 & 6.43 \\
现金短债比（倍） & 0.79 & 0.99 & 1.44 & 2.62 & 1.49 \\
\bottomrule
\end{tabular}
\end{adjustbox}
\end{center}

\begin{center}
\begin{minipage}{\textwidth}\centering
\begin{tikzpicture}
\begin{groupplot}[
  group style={group size=2 by 1, horizontal sep=0.60cm},
  width=6.85cm, height=4.60cm,
  tick label style={font=\tiny},
  title style={font=\scriptsize\heiti},
  yticklabel style={font=\tiny},
  ylabel style={font=\tiny},
  grid=major, grid style={LightGray, line width=0.3pt},
  axis line style={InkGray, line width=0.5pt},
  enlarge x limits={abs=0.8},
  legend style={font=\scriptsize, at={(0.5,-0.20)}, anchor=north,
                legend columns=2, draw=none, fill=none, column sep=6pt},
]
\nextgroupplot[
  ybar, bar width=4pt,
  ymin=-2.92, ymax=13.86,
  xtick={0,1,2,3,4,5.7},
  xticklabels={2021,2022,2023,2024,2025,2026H1},
  ylabel={亿元},
]
\addplot[fill=AgriGreen!75, draw=AgriGreen, bar shift=-2.6pt] table[x=i, y=rev]{data/clean/profiles/fin_300967.dat};
\addplot[fill=SoftRed!60, draw=SoftRed, bar shift=2.6pt] table[x=i, y=ni]{data/clean/profiles/fin_300967.dat};
\addplot[fill=AgriGreen!30, draw=AgriGreen, pattern=north east lines, bar shift=-2.6pt] table[x=x, y=rev]{data/clean/profiles/finh_300967.dat};
\addplot[fill=SoftRed!20, draw=SoftRed, pattern=north east lines, bar shift=2.6pt] table[x=x, y=ni]{data/clean/profiles/finh_300967.dat};
\legend{营业收入（年/半年）, 归母净利润（年/半年）}
\nextgroupplot[
  ymin=-25.9, ymax=61.8,
  xtick={0,1,2,3,4,5.7},
  xticklabels={2021,2022,2023,2024,2025,2026H1},
  ylabel={\%},
]
\addplot[AgriGold, mark=*, mark size=1.3pt] table[x=i, y=roe]{data/clean/profiles/fin_300967.dat};
\addplot[OilBlue, mark=square*, mark size=1.3pt] table[x=i, y=debt]{data/clean/profiles/fin_300967.dat};
\addplot[only marks, AgriGold, mark=*, mark size=2.0pt] table[x=x, y=roe]{data/clean/profiles/finh_300967.dat};
\addplot[only marks, OilBlue, mark=square*, mark size=2.0pt] table[x=x, y=debt]{data/clean/profiles/finh_300967.dat};
\legend{ROE, 资产负债率}
\end{groupplot}
\end{tikzpicture}
\begin{tikzpicture}
\begin{axis}[
  name=finbot,
  width=13.75cm, height=3.10cm, scale only axis,
  ymin=-0.72, ymax=8.04,
  xtick={0,1,2,3,4,5.7},
  xticklabels={2021,2022,2023,2024,2025,2026H1},
  tick label style={font=\tiny},
  yticklabel style={font=\tiny},
  ylabel={亿元}, ylabel style={font=\tiny},
  ybar, bar width=4pt,
  grid=major, grid style={LightGray, line width=0.3pt},
  axis line style={InkGray, line width=0.5pt},
  enlarge x limits={abs=0.8},
  legend style={font=\scriptsize, at={(0.5,-0.26)}, anchor=north,
                legend columns=2, draw=none, fill=none, column sep=10pt},
]
\addplot[fill=AgriGreen!55, draw=AgriGreen, bar shift=-3.0pt] table[x=i, y=cash]{data/clean/profiles/fin_300967.dat};
\addplot[fill=SoftRed!45, draw=SoftRed, bar shift=3.0pt] table[x=i, y=ibd]{data/clean/profiles/fin_300967.dat};
\addplot[fill=AgriGreen!25, draw=AgriGreen, pattern=north east lines, bar shift=-3.0pt] table[x=x, y=cash]{data/clean/profiles/finh_300967.dat};
\addplot[fill=SoftRed!20, draw=SoftRed, pattern=north east lines, bar shift=3.0pt] table[x=x, y=ibd]{data/clean/profiles/finh_300967.dat};
\legend{货币资金（左轴）, 有息负债（左轴）}
\end{axis}
\begin{axis}[
  at={(finbot.south east)}, anchor=south east,
  width=13.75cm, height=3.10cm, scale only axis,
  axis y line*=right, axis x line*=none, xtick=\empty,
  ymin=-0.80, ymax=3.98,
  tick label style={font=\tiny},
  yticklabel style={font=\tiny}, scaled y ticks=false,
  legend style={font=\scriptsize, at={(0.5,-0.44)}, anchor=north,
                legend columns=1, draw=none, fill=none},
]
\addplot[OilBlue, mark=*, mark size=2.2pt, line width=1.3pt] table[x=i, y=ocf]{data/clean/profiles/fin_300967.dat};
\addplot[only marks, OilBlue, mark=*, mark size=3.2pt] table[x=x, y=ocf]{data/clean/profiles/finh_300967.dat};
\legend{经营活动现金流净额（右轴，亿元，蓝点）}
\end{axis}
\end{tikzpicture}

\captionof{figure}[晓鸣股份（300967）近年主要财务指标]{晓鸣股份（300967）近年主要财务指标（左上：营业收入与归母净利润；右上：ROE 与资产负债率；下：货币资金、有息负债为左轴柱状，经营活动现金流净额为其右轴的蓝点折线。
2021---2025 年为年报口径，2026H1 为 2026 年半年报/半年末，与其前一年报之间留出间隔、以斜纹或空心点区分；半年报数据未年化）}
\end{minipage}
\end{center}

\pfl{财务分析} 2026 年 6 月末资产负债率 52.3\%，较 2025 年末上升 3.4 个百分点（样本中位数 52.2\%，所处三级行业内按由低到高排第\zhnumber{5}位）。 2026 年 6 月末货币资金 1.43 亿元、短期债务（短期借款与一年内到期非流动负债）0.96 亿元、长期债务 5.47 亿元，有息负债合计 6.43 亿元，现金短债比 1.49 倍——覆盖充分。 净有息负债 5.00 亿元，相当于其 2026 年上半年营业收入的 88\%。 流动比率 1.08、速动比率 0.69，短期指标尚可。 2026 年上半年经营活动现金流净额 1.40 亿元，经营现金流与利润同向为正，内生资金可以支撑运营。 综合看，现金短债比 1.49 倍，短期偿付压力可控；资产负债率高于样本中位数 0.1 个百分点；有息负债中长期债务占比更高，期限结构对短债滚动的压力较小。

\pfl{重大战略转型} 从公开披露看，公司的战略推进集中在：融资与资本运作（1 项）：2022 年，占前次发行总筹资额的 27.51\%。 公告口径上，近三年（2023 年 9 月---2026 年 9 月）公司披露重大重组与并购类公告 0 条、对外投资与转型类公告 3 条、回购与股权激励类公告 26 条，合计公告 457 条。 第一大行业仍是“养殖行业”，收入占比 100.0\%，与 2021 年年报相比没有方向性变化。 综合判断，报告期内公司有 3 条与战略相关的公告落地（重大重组与并购 0 条、对外投资与转型 3 条），资本运作与主业调整之间存在直接关联。

\begin{center}\small\setlength{\tabcolsep}{4pt}
\captionof{table}[晓鸣股份（300967）历史融资与资本运作]{晓鸣股份（300967）历史融资与资本运作（据公司公告与公开报道整理；金额为披露值，含首次公开发行、再融资、债券、银行借款与重整投资等）}
\begin{adjustbox}{max width=\textwidth}
\begin{tabular}{@{}p{1.45cm}p{2.55cm}p{4.05cm}p{4.95cm}@{}}
\toprule
时间 & 方式 & 规模 & 用途与说明 \\
\midrule
2021-04 & IPO（深交所创业板，证监许可〔2021〕613号） & 发行 4,\allowbreak{}700 万股，\allowbreak{}募集资金总额 2.13 亿元 & 募投项目建设（后部分变更为红寺堡智慧农业产业示范园父母代种业基地项目）；资金于 2021 年 4 月 7 日到位 \\
2023-04 & 向不特定对象发行可转换公司债券（证监许可〔2023〕253号） & 发行 3,\allowbreak{}290,\allowbreak{}000 张，\allowbreak{}每张面值 100.00 元 & 募投项目建设（截至 2025 年底募集资金已使用完毕，\allowbreak{}募投项目均达到预定可使用状态）；截至 2026 年 3 月底转股价格调整为 19.34 元/\allowbreak{}股 \\
2024-12 & 股份回购（用于员工持股计划或股权激励） & 回购总金额不低于 2,\allowbreak{}000.00 万元（含）且不超过 4,\allowbreak{}000.00 万元（含） & 员工持股计划或股权激励；2024 年 12 月 20 日第五届董事会第十七次会议、\allowbreak{}2025 年 1 月 10 日 2025 年第一次临时股东大会审议通过 \\
\bottomrule
\end{tabular}
\end{adjustbox}
\end{center}

\pfl{历史融投资情况} 从现金流量表看，2025 年公司取得借款收到的现金 1.08 亿元、偿还债务支付的现金 2.30 亿元，筹资活动现金流净额 -2.03 亿元（2023 年为取得借款 1.07 亿元、偿还债务 0.94 亿元）。 2026 年上半年筹资活动现金流净额为 -0.29 亿元（取得借款 0.59 亿元、偿还债务 0.32 亿元，未年化），已转为净流出，处于净还债状态。 2025 年分配股利、利润或偿付利息支付的现金合计 0.49 亿元。 公开披露的融资记录共 3 笔，工具类型以 IPO 为主；金额与用途见下表。 其中规模最大的两笔为：2023 年 4 月，发行 3,290,000 张，每张面值 100.00 元；2024 年 12 月，股份回购，回购总金额不低于 2,000.00 万元（含）且不超过 4,000.00 万元（含）。 股本方面，近三年披露股本变动 10 次（可转债转股 4 次、激励股份解禁 1 次、股份回购 1 次、限售股份上市 1 次）；总股本由 2021 年末的 1.88 亿股变为 1.89 亿股（+1.06\%）。 分红方面，近三年现金分红 2 次、合计每股派现 0.45 元，最近一次实施在 2025 年 12 月。 近三年“再融资”类公告 33 条，涉及定向增发、可转债、募集资金使用。 综合看，公司融资结构上股权与债务融资并举；2025 年筹资活动现金流净额 -2.03 亿元，融资节奏仍以维持存量债务周转为主，与前述经营与投资现金流的方向基本一致。

\pfl{未来潜在融资需求分析} 2026 年 6 月末货币资金 1.43 亿元，而短期债务 0.96 亿元（现金短债比 1.49 倍），2026 年上半年经营现金流净额 1.40 亿元。账面货币资金可以覆盖短期债务，缺口不体现在静态偿付层面。 公司 2026 年 6 月末资产负债率 52.3\%、2026 年上半年 ROE 0.2\%（未年化），资产端与负债端都没有明显压力，融资需求以相机抉择的扩张型用途为主。 公开披露的后续资金安排线索包括：债务与授信安排：2029 年 4 月，“晓鸣转债”债券余额 3.29 亿元、到期兑付日 该日。 以现有货币资金与经营现金流衡量，公司不存在刚性补血的紧迫性，融资动作更多是主动型的。




\subsubsection{民和股份（002234）}
\label{co:002234}

\pfl{公司基本情况} 1997 年设立至今，民和股份是山东蓬莱的白羽肉鸡一体化企业，以父母代肉种鸡饲养和商品代鸡苗销售为核心，兼营鸡肉制品与有机废弃物资源化利用；白羽鸡行业已连续亏损多年，公司 2025 年、2026 年上半年持续亏损，母公司存在累计未弥补亏损。 公司于 2008-05-16 在深交所主板首发上市，注册资本 3.49 亿元，总股本 3.49 亿股。 截至 2026 年 10 月 12 日收盘价 7.69 元，对应总市值 26.8 亿元，近一年-10.3\%，总市值在三级行业“肉鸡养殖”的 8 家公司中按由大到小排第\zhnumber{7}位，近一年涨跌幅高于全样本中位数（-10.4\%）0.1 个百分点。 按 2026 年半年报披露的产品构成，收入占比前两位为销售商品 99.7\%；让渡资产使用权 0.3\%。 与 2021 年年报相比，第一大产品口径由“雏鸡”（52.7\%）变为“销售商品”（99.7\%）；两期披露的分类口径有调整，占比差异中同时包含分类因素。

\pfl{股权结构与控制权} 按 2026 年 6 月 30 日的股东名册，前三大股东为孙宪法 持股 18.16\%；孙希民 持股 13.51\%；孙宪秋 持股 5.73\%，股东户数 25,860 户。 控制权与股东性质方面，孙宪法持股 6,335.63 万股、占总股本 18.16\%，其中限售 1,583.91 万股、无质押。

\pfl{经营要点} 经营层面，可适当增加新的产能带动业绩增长。

\pfl{所处行业的基本面} 白羽与黄羽肉鸡的供给取决于祖代引种量与父母代存栏，需求端受餐饮与替代效应驱动，产能周期只有 6---12 个月（第 \ref{sec:framework} 章的供需框架）。 与该公司直接相关的行业数据包括：2026 年，7—8 月白羽鸡苗价格快速上涨；2026 年上半年，全国白羽鸡苗均价 2.70 元/羽、同比涨幅 12.50\%；2026 年，1—5 月父母代鸡苗市场成交均价 40.63 元/套、同比上涨 21.03\%。

\pfl{竞争优势} 把该公司放回同类公司比较，2026 年上半年营业收入 10.3 亿元、在三级行业“肉鸡养殖”的 8 家公司中排第\zhnumber{7}位；毛利率 5.0\%、低于行业内中位数 5.8 个百分点；半年 ROE -5.5\%（未年化）、低于行业内中位数（1.21\%）；资产负债率 53.9\%、高于行业内中位数（52.1\%）。 公司在公开资料中自述与同业比较的相对优势在于：商品代鸡苗为核心产品、毛利最高且价格弹性最大、销售价格波动较大。 就当期报表看，公司尚未把经营层面的特征转化为高于行业的资本回报，周期底部的相对优势更多体现在现金流与资产负债表上。

\pfl{历史股价} 自 2021 年 1 月至 2026 年 10 月，股价由 13.34 元变为 7.69 元，累计 -42.4\%，区间最高 22.65 元（2023 年 1 月）、最低 6.41 元（2026 年 6 月）；当前价相当于区间最高价的 34.0\%，为区间最低价的 1.20 倍。 月度收益的年化波动率 32.7\%，低于全样本中位数 37.3\%，在三级行业“肉鸡养殖”的 8 家公司中波动率由高到低排第\zhnumber{7}位。

\begin{center}
\begin{minipage}{\textwidth}\centering
\begin{tikzpicture}
\begin{axis}[
  width=12.6cm, height=3.9cm, scale only axis,
  xmin=2020.922, xmax=2026.972, ymin=5.96, ymax=24.24,
  xtick={2021,2022,2023,2024,2025,2026},
  yearticks,
  yticklabel style={font=\scriptsize},
  ylabel={元/股}, ylabel style={font=\scriptsize},
  grid=major, grid style={LightGray, line width=0.3pt},
  axis line style={InkGray, line width=0.5pt},
  every axis plot/.append style={line width=0.9pt},
]
\addplot[AgriGreen] table[x=x,y=y]{data/clean/profiles/px_002234.dat};
\addplot[SoftRed, dashed, line width=0.7pt] table[x=x,y=y]{data/clean/profiles/pxzz_002234.dat};
\addplot[only marks, mark=*, mark size=1.1pt, SoftRed] table[x=x,y=y]{data/clean/profiles/pxzz_002234.dat};
\end{axis}
\end{tikzpicture}

\captionof{figure}[民和股份（002234）股价与周期骨架]{民和股份（002234）月度收盘价（元/股）与 ZigZag 周期骨架（阈值 25\%，圆点为周期转折点）}
\end{minipage}
\end{center}

\pfl{过去周期分析} 以 25\% 的 ZigZag 阈值划分，样本区间内股价形成 2 个主升段与 1 个主跌段。 最大主跌段为 2023 年 1 月 至 2026 年 6 月，41 个月-71.7\%；最大主升段出现在 2021 年 1 月 至 2023 年 1 月，24 个月+69.8\%。 最近一次转折出现在 2026 年 8 月（上涨段自 2026 年 6 月起，2 个月+28.5\%），当前处于该段之后的震荡之中。 公司股价与申万农林牧渔指数几乎同步见底，个股并未走出独立行情。 禽链的产能周期只有 6---12 个月、反转快于生猪（第 \ref{sec:other-livestock} 节），2026 年前四个月毛鸡养殖利润同比大幅回升，链条利润的修复先于生猪，公司 2026 年 6 月末资产负债率 53.9\%、半年 ROE -5.5\%（未年化），在本轮下行中属于随行业同步调整的一类。 综合区间形态看，该股单段涨跌幅度偏大、以趋势段为主（最大主升 +69.8\%、最大主跌 -71.7\%），年化波动率 32.7\% 与全样本中位数 37.3\% 相比波动小于样本中位数。

\pfl{盈利情况} 2026 年上半年营业收入 10.27 亿元（同比 +4.2\%），归母净利润 -0.95 亿元，亏损同比收窄 1.19 亿元。 以年报为趋势对照，2023---2025 年营业收入由 20.74 亿元变为 21.62 亿元（两年年均复合增速 +2.1\%），归母净利润由 -3.86 亿元变为 -2.69 亿元。 按半年报口径，2026 年上半年的 ROE 为 -5.5\%（未年化）、销售净利率 -9.7\%、毛利率 5.0\%，ROE 在“肉鸡养殖”8 家公司中按数值由高到低排第\zhnumber{8}位。 按同一口径，三级行业“肉鸡养殖”8 家公司的半年 ROE 中位数为 1.21\%，该公司低于其中位数 6.74 个百分点。 从公司披露的原因看，业绩变动主要来自：2025 年 7 月，原因为白羽肉鸡市场整体疲软、鸡苗和鸡肉产品销售价格较低，并对父母代种鸡、种蛋等相关资产计提减值。 综合看，公司仍处亏损区间、亏损同比收窄，半年 ROE 在三级行业内按由高到低排第\zhnumber{8}位，单位盈利能力的修复依赖行业景气与自身成本端的改善，报表弹性主要体现为价格与销量的方向性变化。

\begin{center}\small\setlength{\tabcolsep}{3.4pt}
\captionof{table}[民和股份（002234）关键财务指标]{民和股份（002234）关键财务指标（合并报表口径；两个半年报期均未年化；现金短债比 $=$ 货币资金 $\div$ 短期债务）}
\begin{adjustbox}{max width=\textwidth}
\begin{tabular}{@{}lrrrrr@{}}
\toprule
指标 & 2023 年 & 2024 年 & 2025 年 & 2025 年半年报 & 2026 年半年报 \\
\midrule
营业收入（亿元） & 20.74 & 21.60 & 21.62 & 9.86 & 10.27 \\
归母净利润（亿元） & -3.86 & -2.49 & -2.69 & -2.14 & -0.95 \\
毛利率（\%） & 3.3 & 8.1 & 4.8 & -0.3 & 5.0 \\
ROE（\%） & -15.6 & -11.5 & -14.1 & -11.1 & -5.5 \\
资产负债率（\%） & 39.4 & 47.3 & 51.8 & 52.0 & 53.9 \\
经营活动现金流净额（亿元） & -0.65 & -0.56 & 1.40 & 0.55 & 0.58 \\
货币资金（亿元） & 5.88 & 6.22 & 5.27 & 6.14 & 5.14 \\
有息负债合计（亿元） & 11.25 & 14.80 & 14.45 & 14.83 & 13.65 \\
现金短债比（倍） & 0.53 & 0.44 & 0.40 & 0.45 & 0.42 \\
\bottomrule
\end{tabular}
\end{adjustbox}
\end{center}

\begin{center}
\begin{minipage}{\textwidth}\centering
\begin{tikzpicture}
\begin{groupplot}[
  group style={group size=2 by 1, horizontal sep=0.60cm},
  width=6.85cm, height=4.60cm,
  tick label style={font=\tiny},
  title style={font=\scriptsize\heiti},
  yticklabel style={font=\tiny},
  ylabel style={font=\tiny},
  grid=major, grid style={LightGray, line width=0.3pt},
  axis line style={InkGray, line width=0.5pt},
  enlarge x limits={abs=0.8},
  legend style={font=\scriptsize, at={(0.5,-0.20)}, anchor=north,
                legend columns=2, draw=none, fill=none, column sep=6pt},
]
\nextgroupplot[
  ybar, bar width=4pt,
  ymin=-7.13, ymax=24.75,
  xtick={0,1,2,3,4,5.7},
  xticklabels={2021,2022,2023,2024,2025,2026H1},
  ylabel={亿元},
]
\addplot[fill=AgriGreen!75, draw=AgriGreen, bar shift=-2.6pt] table[x=i, y=rev]{data/clean/profiles/fin_002234.dat};
\addplot[fill=SoftRed!60, draw=SoftRed, bar shift=2.6pt] table[x=i, y=ni]{data/clean/profiles/fin_002234.dat};
\addplot[fill=AgriGreen!30, draw=AgriGreen, pattern=north east lines, bar shift=-2.6pt] table[x=x, y=rev]{data/clean/profiles/finh_002234.dat};
\addplot[fill=SoftRed!20, draw=SoftRed, pattern=north east lines, bar shift=2.6pt] table[x=x, y=ni]{data/clean/profiles/finh_002234.dat};
\legend{营业收入（年/半年）, 归母净利润（年/半年）}
\nextgroupplot[
  ymin=-21.2, ymax=60.9,
  xtick={0,1,2,3,4,5.7},
  xticklabels={2021,2022,2023,2024,2025,2026H1},
  ylabel={\%},
]
\addplot[AgriGold, mark=*, mark size=1.3pt] table[x=i, y=roe]{data/clean/profiles/fin_002234.dat};
\addplot[OilBlue, mark=square*, mark size=1.3pt] table[x=i, y=debt]{data/clean/profiles/fin_002234.dat};
\addplot[only marks, AgriGold, mark=*, mark size=2.0pt] table[x=x, y=roe]{data/clean/profiles/finh_002234.dat};
\addplot[only marks, OilBlue, mark=square*, mark size=2.0pt] table[x=x, y=debt]{data/clean/profiles/finh_002234.dat};
\legend{ROE, 资产负债率}
\end{groupplot}
\end{tikzpicture}
\begin{tikzpicture}
\begin{axis}[
  name=finbot,
  width=13.75cm, height=3.10cm, scale only axis,
  ymin=-1.48, ymax=16.57,
  xtick={0,1,2,3,4,5.7},
  xticklabels={2021,2022,2023,2024,2025,2026H1},
  tick label style={font=\tiny},
  yticklabel style={font=\tiny},
  ylabel={亿元}, ylabel style={font=\tiny},
  ybar, bar width=4pt,
  grid=major, grid style={LightGray, line width=0.3pt},
  axis line style={InkGray, line width=0.5pt},
  enlarge x limits={abs=0.8},
  legend style={font=\scriptsize, at={(0.5,-0.26)}, anchor=north,
                legend columns=2, draw=none, fill=none, column sep=10pt},
]
\addplot[fill=AgriGreen!55, draw=AgriGreen, bar shift=-3.0pt] table[x=i, y=cash]{data/clean/profiles/fin_002234.dat};
\addplot[fill=SoftRed!45, draw=SoftRed, bar shift=3.0pt] table[x=i, y=ibd]{data/clean/profiles/fin_002234.dat};
\addplot[fill=AgriGreen!25, draw=AgriGreen, pattern=north east lines, bar shift=-3.0pt] table[x=x, y=cash]{data/clean/profiles/finh_002234.dat};
\addplot[fill=SoftRed!20, draw=SoftRed, pattern=north east lines, bar shift=3.0pt] table[x=x, y=ibd]{data/clean/profiles/finh_002234.dat};
\legend{货币资金（左轴）, 有息负债（左轴）}
\end{axis}
\begin{axis}[
  at={(finbot.south east)}, anchor=south east,
  width=13.75cm, height=3.10cm, scale only axis,
  axis y line*=right, axis x line*=none, xtick=\empty,
  ymin=-4.00, ymax=2.68,
  tick label style={font=\tiny},
  yticklabel style={font=\tiny}, scaled y ticks=false,
  legend style={font=\scriptsize, at={(0.5,-0.44)}, anchor=north,
                legend columns=1, draw=none, fill=none},
]
\addplot[OilBlue, mark=*, mark size=2.2pt, line width=1.3pt] table[x=i, y=ocf]{data/clean/profiles/fin_002234.dat};
\addplot[only marks, OilBlue, mark=*, mark size=3.2pt] table[x=x, y=ocf]{data/clean/profiles/finh_002234.dat};
\legend{经营活动现金流净额（右轴，亿元，蓝点）}
\end{axis}
\end{tikzpicture}

\captionof{figure}[民和股份（002234）近年主要财务指标]{民和股份（002234）近年主要财务指标（左上：营业收入与归母净利润；右上：ROE 与资产负债率；下：货币资金、有息负债为左轴柱状，经营活动现金流净额为其右轴的蓝点折线。
2021---2025 年为年报口径，2026H1 为 2026 年半年报/半年末，与其前一年报之间留出间隔、以斜纹或空心点区分；半年报数据未年化）}
\end{minipage}
\end{center}

\pfl{财务分析} 2026 年 6 月末资产负债率 53.9\%，较 2025 年末上升 2.2 个百分点（样本中位数 52.2\%，所处三级行业内按由低到高排第\zhnumber{6}位）。 2026 年 6 月末货币资金 5.14 亿元、短期债务（短期借款与一年内到期非流动负债）12.38 亿元、长期债务 1.28 亿元，有息负债合计 13.65 亿元，现金短债比 0.42 倍——覆盖偏紧。 净有息负债 8.51 亿元，相当于其 2026 年上半年营业收入的 83\%。 流动比率 0.76、速动比率 0.40，短期偿债指标偏紧。 2026 年上半年经营活动现金流净额 0.58 亿元，净利润为负而经营现金流为正，差额主要来自折旧摊销与营运资本变动，说明亏损尚未完全转化为现金流出。 2026 年 6 月末在建工程 1.39 亿元，较上年末减少 0.80 亿元，在建投入在收缩。 综合看，账面货币资金对有息负债与短期债务的覆盖不足，滚动偿付对新增融资的依赖度较高；资产负债率高于样本中位数 1.8 个百分点；有息负债中短期债务占比更高，期限结构仍以滚动续作为主。

\pfl{重大战略转型} 从公开披露看，公司的战略推进集中在：产能与项目（1 项）：2015 年，子公司沼气提纯压缩项目完工转固并开始计提折旧。 公告口径上，近三年（2023 年 9 月---2026 年 9 月）公司披露重大重组与并购类公告 0 条、对外投资与转型类公告 0 条、回购与股权激励类公告 0 条，合计公告 263 条。 主营结构的口径变化已经落到报表上：2021 年年报的第一大产品为“雏鸡”（52.7\%），2026 年半年报为“销售商品”（99.7\%）；公司在这两期的分部划分方式有过调整。 综合判断，公司报告期内的变化主要发生在收入结构层面，属于口径与业务重心的再平衡，而非外延式扩张。

\begin{center}\small\setlength{\tabcolsep}{4pt}
\captionof{table}[民和股份（002234）历史融资与资本运作]{民和股份（002234）历史融资与资本运作（据公司公告与公开报道整理；金额为披露值，含首次公开发行、再融资、债券、银行借款与重整投资等）}
\begin{adjustbox}{max width=\textwidth}
\begin{tabular}{@{}p{1.45cm}p{2.55cm}p{4.05cm}p{4.95cm}@{}}
\toprule
时间 & 方式 & 规模 & 用途与说明 \\
\midrule
2008-05 & IPO（深交所中小板） & 发行 2,\allowbreak{}700 万股（网下配售 540 万股、\allowbreak{}网上定价发行 2,\allowbreak{}160 万股） & 募投项目建设（含粪污处理大型沼气工程及资源化利用项目）；网上定价发行的 2,\allowbreak{}160 万股于 2008 年 5 月 16 日上市交易 \\
2011-11 & 非公开发行（A股） & 发行 4,\allowbreak{}352.33 万股，\allowbreak{}发行价 19.30 元/\allowbreak{}股，\allowbreak{}募集资金总额 8.40 亿元 & 募投项目建设；10 名发行对象认购股份限售期 12 个月 \\
2021-04 & 非公开发行（证监许可〔2020〕3346号） & 发行 4,\allowbreak{}691.36 万股，\allowbreak{}发行价 12.15 元/\allowbreak{}股，\allowbreak{}募集资金总额 5.70 亿元 & 募投项目建设；23 名发行对象以现金认购，\allowbreak{}2024 年度使用募集资金 2,\allowbreak{}757.06 万元投入项目 \\
\bottomrule
\end{tabular}
\end{adjustbox}
\end{center}

\pfl{历史融投资情况} 从现金流量表看，2025 年公司取得借款收到的现金 17.12 亿元、偿还债务支付的现金 14.93 亿元，筹资活动现金流净额 -0.18 亿元（2023 年为取得借款 12.96 亿元、偿还债务 11.98 亿元）。 2026 年上半年筹资活动现金流净额为 0.34 亿元（取得借款 8.75 亿元、偿还债务 8.02 亿元，未年化），仍处于净融入状态。 2025 年分配股利、利润或偿付利息支付的现金合计 0.44 亿元。 公开披露的融资记录共 3 笔，工具类型以 IPO、定向增发为主；金额与用途见下表。 股本方面，近三年披露股本变动 3 次；总股本自 2021 年末以来未发生变化（3.49 亿股）。 分红方面，近三年未实施现金分红，在全部样本中属于分红能力偏弱的一端。 近三年“再融资”类公告 1 条，涉及定向增发。 综合看，公司融资结构上股权与债务融资并举；2025 年筹资活动现金流净额 -0.18 亿元，融资节奏仍以维持存量债务周转为主，与前述经营与投资现金流的方向基本一致。

\pfl{未来潜在融资需求分析} 2026 年 6 月末货币资金 5.14 亿元，而短期债务 12.38 亿元（现金短债比 0.42 倍），2026 年上半年经营现金流净额 0.58 亿元。按“短期债务减去账面货币资金”的滚动偿付口径，静态缺口约 7.24 亿元；若再计入上半年亏损的年化额（1.90 亿元），维持一年正常经营所需的外部资金约 7.24---9.14 亿元。 公司 2026 年上半年归母净利润为负（-0.95 亿元），但 6 月末资产负债率 53.9\% 尚未进入第一档（亏损且负债率 $>$65\%）的区间，当前融资需求以补充流动资金、支撑低谷期经营性支出为主。 资金需求还要叠加在建工程：期末在建工程 1.39 亿元、2026 年上半年购建固定资产等支出 0.43 亿元（未年化），这部分支出在周期底部通常会被主动放缓。 综合看，公司的外部资金需求以补充营运资金为主、项目投入为辅，滚动偿付口径下的静态缺口约 7.24 亿元，能否落地取决于授信续作与发行窗口的配合。




\subsubsection{*ST福成（600965）}
\label{co:600965}

\pfl{公司基本情况} *ST 福成成立于 1998 年，河北三河的牛羊肉屠宰加工与餐饮企业，2015 年发行股份收购大股东旗下灵山宝塔陵园，成为 A 股唯一拥有殡葬业务的公司；2026 年 5 月因 2025 年财务报表被出具无法表示意见，被实施退市风险警示并叠加其他风险警示，简称变更为“*ST 福成”。 公司于 2004-07-13 在上交所首发上市，注册资本 8.19 亿元，总股本 8.19 亿股。 截至 2026 年 10 月 12 日收盘价 2.46 元，对应总市值 20.1 亿元，近一年-49.9\%，总市值在三级行业“其他养殖”的 3 家公司中按由大到小排第\zhnumber{3}位，近一年涨跌幅低于全样本中位数（-10.4\%）39.5 个百分点。

\pfl{股权结构与控制权} 按 2026 年 6 月 30 日的股东名册，前三大股东为福成投资集团有限公司 持股 34.51\%；李福成 持股 15.87\%；三河福生投资有限公司 持股 9.01\%，股东户数 19,969 户。 股权结构的另一条线索是，2015 年，公司向控股股东福成投资集团有限公司发行股份购买三河灵山宝塔陵园有限公司 100\% 股权。

\pfl{经营要点} 经营与产能方面，餐饮门店分布于北京市、河北省和内蒙古自治区。

\pfl{所处行业的基本面} 牛、羊的供给由母畜存栏与出栏节奏决定、扩张周期长，水产则由投苗量与消费需求决定（第 \ref{sec:framework} 章的供需框架）。 与该公司直接相关的行业数据包括：2026 年上半年，全国肉牛出栏 2,119 万头、同比减少 83 万头、下降 3.8\%。

\pfl{竞争优势} 把该公司放回同类公司比较，2026 年上半年营业收入 5.7 亿元、在三级行业“其他养殖”的 3 家公司中排第\zhnumber{2}位；毛利率 19.7\%、高于行业内中位数 10.5 个百分点；半年 ROE 0.8\%（未年化）、低于行业内中位数（0.83\%）；资产负债率 16.7\%、低于行业内中位数（66.2\%）。 相对优势方面，殡葬服务毛利率相对较高。 优势在报表上主要体现为高于行业平均的毛利率，这在周期底部是抵御亏损的主要缓冲。

\pfl{历史股价} 以 2021 年 1 月为起点、2026 年 10 月为终点，股价由 4.82 元变为 2.46 元，累计 -49.0\%，区间最高 8.29 元（2022 年 12 月）、最低 2.46 元（2026 年 10 月）；当前价相当于区间最高价的 29.7\%，已接近区间最低价（1.00 倍）。 月度收益的年化波动率 39.1\%，高于全样本中位数 37.3\%，在三级行业“其他养殖”的 3 家公司中波动率由高到低排第\zhnumber{2}位。 把股价阶段与公司事件对齐看，2021 年 1 月 至 2022 年 12 月股价上涨+72.0\%，同期（2022 年 1 月）犯虚开发票罪判处有期徒刑二年十一个月并处罚金五十万元。 整体上看，区间的几处大级别拐点都能在公司的资本运作、股权变动或风险事件上找到对应，股价波动有明显的公司层面驱动，而非单纯的行业 beta。

\begin{center}
\begin{minipage}{\textwidth}\centering
\begin{tikzpicture}
\begin{axis}[
  width=12.6cm, height=3.9cm, scale only axis,
  xmin=2020.922, xmax=2026.972, ymin=2.29, ymax=8.87,
  xtick={2021,2022,2023,2024,2025,2026},
  yearticks,
  yticklabel style={font=\scriptsize},
  ylabel={元/股}, ylabel style={font=\scriptsize},
  grid=major, grid style={LightGray, line width=0.3pt},
  axis line style={InkGray, line width=0.5pt},
  every axis plot/.append style={line width=0.9pt},
]
\addplot[AgriGreen] table[x=x,y=y]{data/clean/profiles/px_600965.dat};
\addplot[SoftRed, dashed, line width=0.7pt] table[x=x,y=y]{data/clean/profiles/pxzz_600965.dat};
\addplot[only marks, mark=*, mark size=1.1pt, SoftRed] table[x=x,y=y]{data/clean/profiles/pxzz_600965.dat};
\end{axis}
\end{tikzpicture}

\captionof{figure}[*ST福成（600965）股价与周期骨架]{*ST福成（600965）月度收盘价（元/股）与 ZigZag 周期骨架（阈值 25\%，圆点为周期转折点）}
\end{minipage}
\end{center}

\pfl{过去周期分析} 以 25\% 的 ZigZag 阈值划分，样本区间内股价形成 2 个主升段与 2 个主跌段。 最大主升段出现在 2024 年 8 月 至 2026 年 1 月，17 个月+87.9\%；最大主跌段为 2026 年 1 月 至 2026 年 10 月，9 个月-61.5\%。 最近一次转折出现在 2026 年 10 月（下跌段自 2026 年 1 月起，9 个月-61.5\%），当前处于该段的延续之中。 与申万农林牧渔指数的最近一次底部相比，该公司见底晚约 5 个月，在本轮下行中属于调整更慢的一类。 牛、羊、水产的产能周期长于生猪，第 \ref{sec:other-livestock} 章判断肉牛供给的实质性收缩要到 2027 年之后才能兑现，这一时滞决定了个股的股价弹性，公司 2026 年上半年实现盈利、半年 ROE 0.8\%（未年化），好于全样本中位数（0.75\%），下行周期对其报表的冲击小于同业。 综合区间形态看，该股单段涨跌幅度偏大、以趋势段为主（最大主升 +87.9\%、最大主跌 -61.5\%），年化波动率 39.1\% 与全样本中位数 37.3\% 相比波动与样本中位数接近。

\pfl{盈利情况} 2026 年上半年营业收入 5.69 亿元（同比 +6.9\%），归母净利润 0.18 亿元，同比 -48.9\%。 以年报为趋势对照，2023---2025 年营业收入由 10.46 亿元变为 11.63 亿元（两年年均复合增速 +5.4\%），归母净利润由 0.99 亿元变为 0.57 亿元。 按半年报口径，2026 年上半年的 ROE 为 0.8\%（未年化）、销售净利率 3.4\%、毛利率 19.7\%，ROE 在“其他养殖”3 家公司中按数值由高到低排第\zhnumber{2}位。 同期全样本半年 ROE 中位数 0.75\%，公司与之相差 0.08 个百分点（略好于中位数公司）。 综合看，公司在报告期维持盈利（高于全样本半年 ROE 中位数 0.75\%），利润的可持续性取决于其他养殖景气的延续与成本管控的执行。

\begin{center}\small\setlength{\tabcolsep}{3.4pt}
\captionof{table}[*ST福成（600965）关键财务指标]{*ST福成（600965）关键财务指标（合并报表口径；两个半年报期均未年化；现金短债比 $=$ 货币资金 $\div$ 短期债务）}
\begin{adjustbox}{max width=\textwidth}
\begin{tabular}{@{}lrrrrr@{}}
\toprule
指标 & 2023 年 & 2024 年 & 2025 年 & 2025 年半年报 & 2026 年半年报 \\
\midrule
营业收入（亿元） & 10.46 & 10.22 & 11.63 & 5.32 & 5.69 \\
归母净利润（亿元） & 0.99 & 0.52 & 0.57 & 0.36 & 0.18 \\
毛利率（\%） & 37.7 & 25.4 & 22.0 & 25.7 & 19.7 \\
ROE（\%） & 4.7 & 2.4 & 2.6 & 1.6 & 0.8 \\
资产负债率（\%） & 15.4 & 16.1 & 16.5 & 16.3 & 16.7 \\
经营活动现金流净额（亿元） & 0.35 & 1.04 & 2.14 & 0.37 & -0.24 \\
货币资金（亿元） & 1.37 & 0.97 & 2.38 & 0.58 & 2.03 \\
有息负债合计（亿元） & 0.73 & 1.18 & 1.30 & 1.17 & 0.98 \\
现金短债比（倍） & 14.20 & 1.12 & 33.53 & 0.66 & 43.48 \\
\bottomrule
\end{tabular}
\end{adjustbox}
\end{center}

\begin{center}
\begin{minipage}{\textwidth}\centering
\begin{tikzpicture}
\begin{groupplot}[
  group style={group size=2 by 1, horizontal sep=0.60cm},
  width=6.85cm, height=4.60cm,
  tick label style={font=\tiny},
  title style={font=\scriptsize\heiti},
  yticklabel style={font=\tiny},
  ylabel style={font=\tiny},
  grid=major, grid style={LightGray, line width=0.3pt},
  axis line style={InkGray, line width=0.5pt},
  enlarge x limits={abs=0.8},
  legend style={font=\scriptsize, at={(0.5,-0.20)}, anchor=north,
                legend columns=2, draw=none, fill=none, column sep=6pt},
]
\nextgroupplot[
  ybar, bar width=4pt,
  ymin=-1.27, ymax=14.21,
  xtick={0,1,2,3,4,5.7},
  xticklabels={2021,2022,2023,2024,2025,2026H1},
  ylabel={亿元},
]
\addplot[fill=AgriGreen!75, draw=AgriGreen, bar shift=-2.6pt] table[x=i, y=rev]{data/clean/profiles/fin_600965.dat};
\addplot[fill=SoftRed!60, draw=SoftRed, bar shift=2.6pt] table[x=i, y=ni]{data/clean/profiles/fin_600965.dat};
\addplot[fill=AgriGreen!30, draw=AgriGreen, pattern=north east lines, bar shift=-2.6pt] table[x=x, y=rev]{data/clean/profiles/finh_600965.dat};
\addplot[fill=SoftRed!20, draw=SoftRed, pattern=north east lines, bar shift=2.6pt] table[x=x, y=ni]{data/clean/profiles/finh_600965.dat};
\legend{营业收入（年/半年）, 归母净利润（年/半年）}
\nextgroupplot[
  ymin=-1.6, ymax=21.3,
  xtick={0,1,2,3,4,5.7},
  xticklabels={2021,2022,2023,2024,2025,2026H1},
  ylabel={\%},
]
\addplot[AgriGold, mark=*, mark size=1.3pt] table[x=i, y=roe]{data/clean/profiles/fin_600965.dat};
\addplot[OilBlue, mark=square*, mark size=1.3pt] table[x=i, y=debt]{data/clean/profiles/fin_600965.dat};
\addplot[only marks, AgriGold, mark=*, mark size=2.0pt] table[x=x, y=roe]{data/clean/profiles/finh_600965.dat};
\addplot[only marks, OilBlue, mark=square*, mark size=2.0pt] table[x=x, y=debt]{data/clean/profiles/finh_600965.dat};
\legend{ROE, 资产负债率}
\end{groupplot}
\end{tikzpicture}
\begin{tikzpicture}
\begin{axis}[
  name=finbot,
  width=13.75cm, height=3.10cm, scale only axis,
  ymin=-1.00, ymax=11.22,
  xtick={0,1,2,3,4,5.7},
  xticklabels={2021,2022,2023,2024,2025,2026H1},
  tick label style={font=\tiny},
  yticklabel style={font=\tiny},
  ylabel={亿元}, ylabel style={font=\tiny},
  ybar, bar width=4pt,
  grid=major, grid style={LightGray, line width=0.3pt},
  axis line style={InkGray, line width=0.5pt},
  enlarge x limits={abs=0.8},
  legend style={font=\scriptsize, at={(0.5,-0.26)}, anchor=north,
                legend columns=2, draw=none, fill=none, column sep=10pt},
]
\addplot[fill=AgriGreen!55, draw=AgriGreen, bar shift=-3.0pt] table[x=i, y=cash]{data/clean/profiles/fin_600965.dat};
\addplot[fill=SoftRed!45, draw=SoftRed, bar shift=3.0pt] table[x=i, y=ibd]{data/clean/profiles/fin_600965.dat};
\addplot[fill=AgriGreen!25, draw=AgriGreen, pattern=north east lines, bar shift=-3.0pt] table[x=x, y=cash]{data/clean/profiles/finh_600965.dat};
\addplot[fill=SoftRed!20, draw=SoftRed, pattern=north east lines, bar shift=3.0pt] table[x=x, y=ibd]{data/clean/profiles/finh_600965.dat};
\legend{货币资金（左轴）, 有息负债（左轴）}
\end{axis}
\begin{axis}[
  at={(finbot.south east)}, anchor=south east,
  width=13.75cm, height=3.10cm, scale only axis,
  axis y line*=right, axis x line*=none, xtick=\empty,
  ymin=-0.86, ymax=2.85,
  tick label style={font=\tiny},
  yticklabel style={font=\tiny}, scaled y ticks=false,
  legend style={font=\scriptsize, at={(0.5,-0.44)}, anchor=north,
                legend columns=1, draw=none, fill=none},
]
\addplot[OilBlue, mark=*, mark size=2.2pt, line width=1.3pt] table[x=i, y=ocf]{data/clean/profiles/fin_600965.dat};
\addplot[only marks, OilBlue, mark=*, mark size=3.2pt] table[x=x, y=ocf]{data/clean/profiles/finh_600965.dat};
\legend{经营活动现金流净额（右轴，亿元，蓝点）}
\end{axis}
\end{tikzpicture}

\captionof{figure}[*ST福成（600965）近年主要财务指标]{*ST福成（600965）近年主要财务指标（左上：营业收入与归母净利润；右上：ROE 与资产负债率；下：货币资金、有息负债为左轴柱状，经营活动现金流净额为其右轴的蓝点折线。
2021---2025 年为年报口径，2026H1 为 2026 年半年报/半年末，与其前一年报之间留出间隔、以斜纹或空心点区分；半年报数据未年化）}
\end{minipage}
\end{center}

\pfl{财务分析} 2026 年 6 月末资产负债率 16.7\%，较 2025 年末上升 0.2 个百分点（样本中位数 52.2\%，所处三级行业内按由低到高排第\zhnumber{1}位）。 2026 年 6 月末货币资金 2.03 亿元、短期债务（短期借款与一年内到期非流动负债）0.05 亿元、长期债务 0.93 亿元，有息负债合计 0.98 亿元，现金短债比 43.48 倍——覆盖充分。 净有息负债 -1.05 亿元，相当于其 2026 年上半年营业收入的 -18\%。 流动比率 4.82、速动比率 1.15，短期流动性无明显缺口。 2026 年上半年经营活动现金流净额 -0.24 亿元，净利润为正但经营现金流为负，利润的现金含量偏低。 综合看，现金短债比 43.48 倍，短期偿付压力可控；资产负债率低于样本中位数 35.5 个百分点；有息负债中长期债务占比更高，期限结构对短债滚动的压力较小。 公司披露的退市与持续经营风险为：2026 年 5 月，公司股票被实施退市风险警示并叠加其他风险警示。

\pfl{重大战略转型} 近三年公司的战略动作集中在：并购与重组（2 项）：2015---2018 年，先后向控股股东福成集团收购宝塔陵园 100\% 股权、通过并购基金持有杭州钱江陵园 60\% 股权。 公告口径上，近三年（2023 年 9 月---2026 年 9 月）公司披露重大重组与并购类公告 0 条、对外投资与转型类公告 1 条、回购与股权激励类公告 0 条，合计公告 221 条。 综合判断，报告期内公司有 1 条与战略相关的公告落地（重大重组与并购 0 条、对外投资与转型 1 条），资本运作与主业调整之间存在直接关联。

\begin{center}\small\setlength{\tabcolsep}{4pt}
\captionof{table}[*ST福成（600965）历史融资与资本运作]{*ST福成（600965）历史融资与资本运作（据公司公告与公开报道整理；金额为披露值，含首次公开发行、再融资、债券、银行借款与重整投资等）}
\begin{adjustbox}{max width=\textwidth}
\begin{tabular}{@{}p{1.45cm}p{2.55cm}p{4.05cm}p{4.95cm}@{}}
\toprule
时间 & 方式 & 规模 & 用途与说明 \\
\midrule
2004-07 & IPO（上海证券交易所） & 首次公开发行人民币普通股 8,\allowbreak{}000 万股，\allowbreak{}发行价 3.15 元/\allowbreak{}股 & 2004 年 6 月 21 日经中国证监会批准 \\
2013-12 & 发行股份购买资产（证监许可〔2013〕1478号） & 向三河福生投资有限公司等发行股份，\allowbreak{}购买福成肥牛餐饮管理有限公司 100\% 股权、\allowbreak{}三河市福成都市食品有限公司 100\% 股权 & 收购大股东旗下餐饮与食品资产；2013 年 12 月 4 日标的资产过户完成 \\
2015-07 & 发行股份购买资产（证监许可〔2015〕298号） & 向控股股东福成投资集团有限公司发行 2.91 亿股 & 收购三河灵山宝塔陵园有限公司 100\% 股权 \\
\bottomrule
\end{tabular}
\end{adjustbox}
\end{center}

\pfl{历史融投资情况} 从现金流量表看，2025 年公司取得借款收到的现金 1.00 亿元、偿还债务支付的现金 0.80 亿元，筹资活动现金流净额 -0.25 亿元。 2026 年上半年筹资活动现金流净额为 -0.00 亿元（取得借款 0.33 亿元、偿还债务 0.30 亿元，未年化），已转为净流出，处于净还债状态。 2025 年分配股利、利润或偿付利息支付的现金合计 0.34 亿元。 公开披露的融资记录共 3 笔，工具类型以 IPO 为主；金额与用途见下表。 其中规模最大的两笔为：2015 年 7 月，向控股股东福成投资集团有限公司发行 2.91 亿股。 股本方面，近三年披露股本变动 3 次；总股本自 2021 年末以来未发生变化（8.19 亿股）。 分红方面，近三年现金分红 4 次、合计每股派现 0.26 元，最近一次实施在 2025 年 4 月。 综合看，公司融资结构上股权与债务融资并举；2025 年筹资活动现金流净额 -0.25 亿元，融资节奏仍以维持存量债务周转为主，与前述经营与投资现金流的方向基本一致。

\pfl{未来潜在融资需求分析} 2026 年 6 月末货币资金 2.03 亿元，而短期债务 0.05 亿元（现金短债比 43.48 倍），2026 年上半年经营现金流净额 -0.24 亿元。账面货币资金可以覆盖短期债务，缺口不体现在静态偿付层面。 公司 2026 年 6 月末资产负债率 16.7\%、2026 年上半年 ROE 0.8\%（未年化），资产端与负债端都没有明显压力，融资需求以相机抉择的扩张型用途为主。 静态偿付不存在缺口，后续融资更可能服务于产能或并购安排，而不是填补流动性窟窿。




\subsubsection{湘佳股份（002982）}
\label{co:002982}

\pfl{公司基本情况} 2003 年设立至今，湘佳股份是湖南石门的黄羽肉鸡一体化企业，集饲料、种禽、养殖、屠宰、冷链与冰鲜禽肉销售于一体，被称为“中国生鲜家禽第一股”，并延伸至生猪养殖；2025 年归母净利润 3,294.67 万元，2026 年上半年由盈转亏。 公司于 2020-04-24 在深交所主板首发上市，注册资本 2.03 亿元，总股本 1.43 亿股。 截至 2026 年 10 月 12 日收盘价 12.81 元，对应总市值 18.3 亿元，近一年-12.8\%，总市值在三级行业“肉鸡养殖”的 8 家公司中按由大到小排第\zhnumber{8}位，近一年涨跌幅低于全样本中位数（-10.4\%）2.4 个百分点。 按 2026 年半年报披露的行业构成，收入几乎全部来自养殖业 100.0\%。 与 2021 年年报相比，第一大行业口径由“养殖行业”（100.0\%）变为“养殖业”（100.0\%）；两期披露的分类口径有调整，占比差异中同时包含分类因素。

\pfl{股权结构与控制权} 按 2026 年 6 月 30 日的股东名册，前三大股东为邢卫民 持股 22.19\%；喻自文 持股 22.19\%；邢成男 持股 1.93\%，股东户数 18,774 户。 从控制权角度看，各质押 210.00 万股；公司质押比例占总股本的 2.07\%，公司总股本为 20,319.41 万股。

\pfl{经营要点} 从经营要点看，截至 2026 年 3 月 31 日十大流通股东中喻自文、邢卫民各持股 8.66\%。

\pfl{所处行业的基本面} 白羽与黄羽肉鸡的供给取决于祖代引种量与父母代存栏，需求端受餐饮与替代效应驱动，产能周期只有 6---12 个月（第 \ref{sec:framework} 章的供需框架）。 当期该行业的关键数据是：2026 年上半年，黄羽鸡行业持续深度亏损，华东地区商品毛鸡出栏均价仅 9.5 元/公斤，而行业平均养殖成本约 12 元/公斤；2026 年上半年，全国禽肉产量 1,390 万吨、同比增长 9.45\%。

\pfl{竞争优势} 把该公司放回同类公司比较，2026 年上半年营业收入 23.2 亿元、在三级行业“肉鸡养殖”的 8 家公司中排第\zhnumber{5}位；毛利率 12.9\%、高于行业内中位数 2.1 个百分点；半年 ROE -0.7\%（未年化）、低于行业内中位数（1.21\%）；资产负债率 60.9\%、高于行业内中位数（52.1\%）。 就当期报表看，公司尚未把经营层面的特征转化为高于行业的资本回报，周期底部的相对优势更多体现在现金流与资产负债表上。

\pfl{历史股价} 自 2021 年 1 月至 2026 年 10 月，股价由 57.70 元变为 12.81 元，累计 -77.8\%，区间最高 62.98 元（2021 年 2 月）、最低 10.55 元（2026 年 6 月）；当前价相当于区间最高价的 20.3\%，为区间最低价的 1.21 倍。 月度收益的年化波动率 38.4\%，高于全样本中位数 37.3\%，在三级行业“肉鸡养殖”的 8 家公司中波动率由高到低排第\zhnumber{3}位。

\begin{center}
\begin{minipage}{\textwidth}\centering
\begin{tikzpicture}
\begin{axis}[
  width=12.6cm, height=3.9cm, scale only axis,
  xmin=2020.922, xmax=2026.972, ymin=9.81, ymax=67.39,
  xtick={2021,2022,2023,2024,2025,2026},
  yearticks,
  yticklabel style={font=\scriptsize},
  ylabel={元/股}, ylabel style={font=\scriptsize},
  grid=major, grid style={LightGray, line width=0.3pt},
  axis line style={InkGray, line width=0.5pt},
  every axis plot/.append style={line width=0.9pt},
]
\addplot[AgriGreen] table[x=x,y=y]{data/clean/profiles/px_002982.dat};
\addplot[SoftRed, dashed, line width=0.7pt] table[x=x,y=y]{data/clean/profiles/pxzz_002982.dat};
\addplot[only marks, mark=*, mark size=1.1pt, SoftRed] table[x=x,y=y]{data/clean/profiles/pxzz_002982.dat};
\end{axis}
\end{tikzpicture}

\captionof{figure}[湘佳股份（002982）股价与周期骨架]{湘佳股份（002982）月度收盘价（元/股）与 ZigZag 周期骨架（阈值 25\%，圆点为周期转折点）}
\end{minipage}
\end{center}

\pfl{过去周期分析} 以 25\% 的 ZigZag 阈值划分，样本区间内股价形成 3 个主升段与 4 个主跌段。 最大主跌段为 2023 年 2 月 至 2024 年 6 月，16 个月-68.0\%；最大主升段出现在 2022 年 4 月 至 2023 年 2 月，10 个月+40.2\%。 最近一次转折出现在 2026 年 6 月（下跌段自 2025 年 4 月起，14 个月-48.4\%），当前处于该段的延续之中。 公司股价与申万农林牧渔指数几乎同步见底，个股并未走出独立行情。 禽链的产能周期只有 6---12 个月、反转快于生猪（第 \ref{sec:other-livestock} 节），2026 年前四个月毛鸡养殖利润同比大幅回升，链条利润的修复先于生猪，公司 2026 年 6 月末资产负债率 60.9\%、半年 ROE -0.7\%（未年化），在本轮下行中属于随行业同步调整的一类。 综合区间形态看，该股单段涨跌幅度偏大、以趋势段为主（最大主升 +40.2\%、最大主跌 -68.0\%），年化波动率 38.4\% 与全样本中位数 37.3\% 相比波动与样本中位数接近。

\pfl{盈利情况} 2026 年上半年营业收入 23.16 亿元（同比 +8.6\%），归母净利润 -0.12 亿元，由上年同期的盈利转为亏损。 以年报为趋势对照，2023---2025 年营业收入由 38.90 亿元变为 44.71 亿元（两年年均复合增速 +7.2\%），归母净利润由 -1.47 亿元变为 0.33 亿元。 按半年报口径，2026 年上半年的 ROE 为 -0.7\%（未年化）、销售净利率 -0.6\%、毛利率 12.9\%，ROE 在“肉鸡养殖”8 家公司中按数值由高到低排第\zhnumber{6}位。 按同一口径，三级行业“肉鸡养殖”8 家公司的半年 ROE 中位数为 1.21\%，该公司低于其中位数 1.91 个百分点。 综合看，公司仍处亏损区间、由盈转亏，半年 ROE 在三级行业内按由高到低排第\zhnumber{6}位，单位盈利能力的修复依赖行业景气与自身成本端的改善，报表弹性主要体现为价格与销量的方向性变化。

\begin{center}\small\setlength{\tabcolsep}{3.4pt}
\captionof{table}[湘佳股份（002982）关键财务指标]{湘佳股份（002982）关键财务指标（合并报表口径；两个半年报期均未年化；现金短债比 $=$ 货币资金 $\div$ 短期债务）}
\begin{adjustbox}{max width=\textwidth}
\begin{tabular}{@{}lrrrrr@{}}
\toprule
指标 & 2023 年 & 2024 年 & 2025 年 & 2025 年半年报 & 2026 年半年报 \\
\midrule
营业收入（亿元） & 38.90 & 42.04 & 44.71 & 21.34 & 23.16 \\
归母净利润（亿元） & -1.47 & 0.98 & 0.33 & 0.06 & -0.12 \\
毛利率（\%） & 12.6 & 18.0 & 14.3 & 12.6 & 12.9 \\
ROE（\%） & -8.4 & 5.8 & 1.9 & 0.3 & -0.7 \\
资产负债率（\%） & 55.0 & 56.6 & 60.2 & 58.8 & 60.9 \\
经营活动现金流净额（亿元） & 1.62 & 3.80 & 3.77 & 0.56 & 0.84 \\
货币资金（亿元） & 3.42 & 3.85 & 4.39 & 3.28 & 3.94 \\
有息负债合计（亿元） & 13.87 & 15.95 & 18.12 & 18.06 & 19.28 \\
现金短债比（倍） & 1.06 & 1.04 & 0.93 & 0.87 & 0.98 \\
\bottomrule
\end{tabular}
\end{adjustbox}
\end{center}

\begin{center}
\begin{minipage}{\textwidth}\centering
\begin{tikzpicture}
\begin{groupplot}[
  group style={group size=2 by 1, horizontal sep=0.60cm},
  width=6.85cm, height=4.60cm,
  tick label style={font=\tiny},
  title style={font=\scriptsize\heiti},
  yticklabel style={font=\tiny},
  ylabel style={font=\tiny},
  grid=major, grid style={LightGray, line width=0.3pt},
  axis line style={InkGray, line width=0.5pt},
  enlarge x limits={abs=0.8},
  legend style={font=\scriptsize, at={(0.5,-0.20)}, anchor=north,
                legend columns=2, draw=none, fill=none, column sep=6pt},
]
\nextgroupplot[
  ybar, bar width=4pt,
  ymin=-6.09, ymax=50.25,
  xtick={0,1,2,3,4,5.7},
  xticklabels={2021,2022,2023,2024,2025,2026H1},
  ylabel={亿元},
]
\addplot[fill=AgriGreen!75, draw=AgriGreen, bar shift=-2.6pt] table[x=i, y=rev]{data/clean/profiles/fin_002982.dat};
\addplot[fill=SoftRed!60, draw=SoftRed, bar shift=2.6pt] table[x=i, y=ni]{data/clean/profiles/fin_002982.dat};
\addplot[fill=AgriGreen!30, draw=AgriGreen, pattern=north east lines, bar shift=-2.6pt] table[x=x, y=rev]{data/clean/profiles/finh_002982.dat};
\addplot[fill=SoftRed!20, draw=SoftRed, pattern=north east lines, bar shift=2.6pt] table[x=x, y=ni]{data/clean/profiles/finh_002982.dat};
\legend{营业收入（年/半年）, 归母净利润（年/半年）}
\nextgroupplot[
  ymin=-14.0, ymax=67.9,
  xtick={0,1,2,3,4,5.7},
  xticklabels={2021,2022,2023,2024,2025,2026H1},
  ylabel={\%},
]
\addplot[AgriGold, mark=*, mark size=1.3pt] table[x=i, y=roe]{data/clean/profiles/fin_002982.dat};
\addplot[OilBlue, mark=square*, mark size=1.3pt] table[x=i, y=debt]{data/clean/profiles/fin_002982.dat};
\addplot[only marks, AgriGold, mark=*, mark size=2.0pt] table[x=x, y=roe]{data/clean/profiles/finh_002982.dat};
\addplot[only marks, OilBlue, mark=square*, mark size=2.0pt] table[x=x, y=debt]{data/clean/profiles/finh_002982.dat};
\legend{ROE, 资产负债率}
\end{groupplot}
\end{tikzpicture}
\begin{tikzpicture}
\begin{axis}[
  name=finbot,
  width=13.75cm, height=3.10cm, scale only axis,
  ymin=-1.93, ymax=21.59,
  xtick={0,1,2,3,4,5.7},
  xticklabels={2021,2022,2023,2024,2025,2026H1},
  tick label style={font=\tiny},
  yticklabel style={font=\tiny},
  ylabel={亿元}, ylabel style={font=\tiny},
  ybar, bar width=4pt,
  grid=major, grid style={LightGray, line width=0.3pt},
  axis line style={InkGray, line width=0.5pt},
  enlarge x limits={abs=0.8},
  legend style={font=\scriptsize, at={(0.5,-0.26)}, anchor=north,
                legend columns=2, draw=none, fill=none, column sep=10pt},
]
\addplot[fill=AgriGreen!55, draw=AgriGreen, bar shift=-3.0pt] table[x=i, y=cash]{data/clean/profiles/fin_002982.dat};
\addplot[fill=SoftRed!45, draw=SoftRed, bar shift=3.0pt] table[x=i, y=ibd]{data/clean/profiles/fin_002982.dat};
\addplot[fill=AgriGreen!25, draw=AgriGreen, pattern=north east lines, bar shift=-3.0pt] table[x=x, y=cash]{data/clean/profiles/finh_002982.dat};
\addplot[fill=SoftRed!20, draw=SoftRed, pattern=north east lines, bar shift=3.0pt] table[x=x, y=ibd]{data/clean/profiles/finh_002982.dat};
\legend{货币资金（左轴）, 有息负债（左轴）}
\end{axis}
\begin{axis}[
  at={(finbot.south east)}, anchor=south east,
  width=13.75cm, height=3.10cm, scale only axis,
  axis y line*=right, axis x line*=none, xtick=\empty,
  ymin=-0.99, ymax=4.94,
  tick label style={font=\tiny},
  yticklabel style={font=\tiny}, scaled y ticks=false,
  legend style={font=\scriptsize, at={(0.5,-0.44)}, anchor=north,
                legend columns=1, draw=none, fill=none},
]
\addplot[OilBlue, mark=*, mark size=2.2pt, line width=1.3pt] table[x=i, y=ocf]{data/clean/profiles/fin_002982.dat};
\addplot[only marks, OilBlue, mark=*, mark size=3.2pt] table[x=x, y=ocf]{data/clean/profiles/finh_002982.dat};
\legend{经营活动现金流净额（右轴，亿元，蓝点）}
\end{axis}
\end{tikzpicture}

\captionof{figure}[湘佳股份（002982）近年主要财务指标]{湘佳股份（002982）近年主要财务指标（左上：营业收入与归母净利润；右上：ROE 与资产负债率；下：货币资金、有息负债为左轴柱状，经营活动现金流净额为其右轴的蓝点折线。
2021---2025 年为年报口径，2026H1 为 2026 年半年报/半年末，与其前一年报之间留出间隔、以斜纹或空心点区分；半年报数据未年化）}
\end{minipage}
\end{center}

\pfl{财务分析} 2026 年 6 月末资产负债率 60.9\%，较 2025 年末上升 0.8 个百分点（样本中位数 52.2\%，所处三级行业内按由低到高排第\zhnumber{7}位）。 2026 年 6 月末货币资金 3.94 亿元、短期债务（短期借款与一年内到期非流动负债）4.02 亿元、长期债务 15.26 亿元，有息负债合计 19.28 亿元，现金短债比 0.98 倍——覆盖偏紧。 净有息负债 15.34 亿元，相当于其 2026 年上半年营业收入的 66\%。 流动比率 1.35、速动比率 0.80，短期偿债能力处于正常区间。 2026 年上半年经营活动现金流净额 0.84 亿元，净利润为负而经营现金流为正，差额主要来自折旧摊销与营运资本变动，说明亏损尚未完全转化为现金流出。 2026 年 6 月末在建工程 2.35 亿元，较上年末减少 0.04 亿元，在建投入在收缩。 综合看，现金短债比 0.98 倍，短期资金安排偏紧；资产负债率高于样本中位数 8.7 个百分点；有息负债中长期债务占比更高，期限结构对短债滚动的压力较小。 公司披露的股东与质押风险为：2026 年 7 月，公司质押总比例 2.27\%，质押总股数 462.04 万股。

\pfl{重大战略转型} 公司的战略动作可归为以下几类：融资与资本运作（2 项）：2020---2022 年，“湘佳转债”募资投向包含 1 万头种猪养殖基地项目；2020 年，有望新增 1,200 万羽黄羽鸡屠宰产能、合计约 2 万吨产能；产能与项目（1 项）：2026 年，实施主体由湖南湘佳牧业股份有限公司变更为广西湘佳牧业有限公司。 公告口径上，近三年（2023 年 9 月---2026 年 9 月）公司披露重大重组与并购类公告 0 条、对外投资与转型类公告 5 条、回购与股权激励类公告 46 条，合计公告 484 条。 主营结构的口径变化已经落到报表上：2021 年年报的第一大行业为“养殖行业”（100.0\%），2026 年半年报为“养殖业”（100.0\%）；公司在这两期的分部划分方式有过调整。 综合判断，报告期内公司有 5 条与战略相关的公告落地（重大重组与并购 0 条、对外投资与转型 5 条），资本运作与主业调整之间存在直接关联。

\begin{center}\small\setlength{\tabcolsep}{4pt}
\captionof{table}[湘佳股份（002982）历史融资与资本运作]{湘佳股份（002982）历史融资与资本运作（据公司公告与公开报道整理；金额为披露值，含首次公开发行、再融资、债券、银行借款与重整投资等）}
\begin{adjustbox}{max width=\textwidth}
\begin{tabular}{@{}p{1.45cm}p{2.55cm}p{4.05cm}p{4.95cm}@{}}
\toprule
时间 & 方式 & 规模 & 用途与说明 \\
\midrule
2020-04 & IPO（深交所中小板，证监许可〔2020〕41号） & 发行 2,\allowbreak{}563.00 万股，\allowbreak{}发行价 29.63 元/\allowbreak{}股，\allowbreak{}募集资金总额 7.59 亿元 & 1,\allowbreak{}250 万羽优质鸡标准化养殖基地建设项目（计划 2.70 亿元）、\allowbreak{}年屠宰 3,\allowbreak{}000 万羽优质鸡加工厂项目（一期 \\
2022-04 & 公开发行可转换公司债券“湘佳转债”（证监许可〔2021〕1814号） & 发行 640.00 万张，\allowbreak{}每张面值 100.00 元，\allowbreak{}按面值发行 & 1,\allowbreak{}350 万羽优质鸡标准化养殖基地建设项目、\allowbreak{}1 万头种猪养殖基地项目和补充流动资金；期限 6 年 \\
2024 年度 & 利润分配（现金分红＋资本公积转增股本） & 每 10 股派发现金红利 1.00 元（含税）、\allowbreak{}以资本公积金每 10 股转增 4 股 & 股东回报；2025 年度分配预案为每 10 股派 1.00 元（含税） \\
\bottomrule
\end{tabular}
\end{adjustbox}
\end{center}

\pfl{历史融投资情况} 从现金流量表看，2025 年公司取得借款收到的现金 5.91 亿元、偿还债务支付的现金 4.16 亿元，筹资活动现金流净额 0.70 亿元（2023 年为取得借款 3.84 亿元、偿还债务 2.44 亿元）。 2026 年上半年筹资活动现金流净额为 0.63 亿元（取得借款 4.07 亿元、偿还债务 3.27 亿元，未年化），仍处于净融入状态。 2025 年分配股利、利润或偿付利息支付的现金合计 0.84 亿元。 公开披露的融资记录共 3 笔，工具类型以 IPO 为主；金额与用途见下表。 其中规模最大的两笔为：2022 年 4 月，发行 640.00 万张，每张面值 100.00 元，按面值发行。 股本方面，近三年披露股本变动 14 次（可转债转股 8 次、其他 3 次、限售股份上市 1 次、转增 1 次）；总股本由 2021 年末的 1.02 亿股变为 1.43 亿股（+40.00\%）。 分红方面，近三年现金分红 5 次、合计每股派现 0.60 元，最近一次实施在 2025 年 12 月。 近三年“再融资”类公告 15 条，涉及可转债。 综合看，公司融资结构上股权与债务融资并举；2025 年筹资活动现金流净额 0.70 亿元，年度内仍为净融入，与前述经营与投资现金流的方向基本一致。

\pfl{未来潜在融资需求分析} 2026 年 6 月末货币资金 3.94 亿元，而短期债务 4.02 亿元（现金短债比 0.98 倍），2026 年上半年经营现金流净额 0.84 亿元。账面货币资金可以覆盖短期债务，缺口不体现在静态偿付层面。 公司 2026 年上半年归母净利润为负（-0.12 亿元），但 6 月末资产负债率 60.9\% 尚未进入第一档（亏损且负债率 $>$65\%）的区间，当前融资需求以补充流动资金、支撑低谷期经营性支出为主。 公开披露的后续资金安排线索包括：债务与授信安排：2025 年 3 月，“湘佳转债”发行规模 6.40 亿元、债券余额 6.40 亿元；2026 年 8 月，公司公告向银行申请综合授信额度，以满足公司日常生产经营的资金需求。 资金需求还要叠加在建工程：期末在建工程 2.35 亿元、2026 年上半年购建固定资产等支出 1.36 亿元（未年化），这部分支出在周期底部通常会被主动放缓。 静态偿付不存在缺口，后续融资更可能服务于产能或并购安排，而不是填补流动性窟窿。




\subsubsection{天域生物（603717）}
\label{co:603717}

\pfl{公司基本情况} 天域生物成立于 2000 年，重庆起家、总部在上海的园林生态工程企业，2019 年起跨界生猪养殖、2021 年新增光伏新能源并于 2024 年更名“天域生物”；2020—2024 年连续五年亏损，控股股东及主要股东股权高比例质押并被司法冻结，实控人因违规减持被证监会处罚。 公司于 2017-03-27 在上交所首发上市，注册资本 2.90 亿元，总股本 2.90 亿股。 截至 2026 年 10 月 12 日收盘价 5.93 元，对应总市值 17.2 亿元，近一年-34.0\%，总市值在三级行业“生猪养殖”的 11 家公司中按由大到小排第\zhnumber{11}位，近一年涨跌幅低于全样本中位数（-10.4\%）23.7 个百分点。 按 2026 年半年报披露的产品构成，收入占比前三位为生态农牧食品业务 90.6\%；生态环境业务 5.8\%；生态能源业务 3.8\%。 与 2021 年年报相比，第一大产品口径由“园林生态工程”（92.6\%）变为“生态农牧食品业务”（90.6\%）；两期披露的分类口径有调整，占比差异中同时包含分类因素。

\pfl{股权结构与控制权} 按 2026 年 6 月 30 日的股东名册，前三大股东为罗卫国 持股 12.84\%；史东伟 持股 11.15\%；深圳泽源私募证券基金管理有限公司－泽源利旺田 42 号私 持股 7.58\%，股东户数 19,042 户。 股东结构与控制权方面，其所持股份累计质押占其持股的 59.06\%；2026 年 8 月，罗卫国因通过基金份额转让隐匿一致行动关系、绕道违规减持 1,167.77 万股（减持金额 1.64 亿元）。

\pfl{经营要点} 业务与产能层面，2026 年上半年，生猪销售收入 2.95 亿元（同比-5.39\%）。

\pfl{所处行业的基本面} 生猪的供给由能繁母猪存栏决定、需求由猪肉消费与替代品价格决定，本轮下行的成因在供给端的产能去化不彻底，而非需求收缩（第 \ref{sec:framework} 章的供需框架）。 从供需两端的当期数据看，2026 年上半年，全国生猪出栏 37246 万头、同比增长 1.7\%。

\pfl{竞争优势} 把该公司放回同类公司比较，2026 年上半年营业收入 3.5 亿元、在三级行业“生猪养殖”的 11 家公司中排第\zhnumber{11}位；毛利率 -1.0\%、高于行业内中位数 0.6 个百分点；半年 ROE -14.6\%（未年化）、低于行业内中位数（-14.32\%）；资产负债率 81.0\%、高于行业内中位数（58.9\%）。 就当期报表看，公司尚未把经营层面的特征转化为高于行业的资本回报，周期底部的相对优势更多体现在现金流与资产负债表上。

\pfl{历史股价} 自 2021 年 1 月至 2026 年 10 月，股价由 8.91 元变为 5.93 元，累计 -33.4\%，区间最高 15.17 元（2022 年 2 月）、最低 5.43 元（2024 年 8 月）；当前价相当于区间最高价的 39.1\%，为区间最低价的 1.09 倍。 月度收益的年化波动率 42.0\%，高于全样本中位数 37.3\%，在三级行业“生猪养殖”的 11 家公司中波动率由高到低排第\zhnumber{4}位。

\begin{center}
\begin{minipage}{\textwidth}\centering
\begin{tikzpicture}
\begin{axis}[
  width=12.6cm, height=3.9cm, scale only axis,
  xmin=2020.922, xmax=2026.972, ymin=5.05, ymax=16.23,
  xtick={2021,2022,2023,2024,2025,2026},
  yearticks,
  yticklabel style={font=\scriptsize},
  ylabel={元/股}, ylabel style={font=\scriptsize},
  grid=major, grid style={LightGray, line width=0.3pt},
  axis line style={InkGray, line width=0.5pt},
  every axis plot/.append style={line width=0.9pt},
]
\addplot[AgriGreen] table[x=x,y=y]{data/clean/profiles/px_603717.dat};
\addplot[SoftRed, dashed, line width=0.7pt] table[x=x,y=y]{data/clean/profiles/pxzz_603717.dat};
\addplot[only marks, mark=*, mark size=1.1pt, SoftRed] table[x=x,y=y]{data/clean/profiles/pxzz_603717.dat};
\end{axis}
\end{tikzpicture}

\captionof{figure}[天域生物（603717）股价与周期骨架]{天域生物（603717）月度收盘价（元/股）与 ZigZag 周期骨架（阈值 25\%，圆点为周期转折点）}
\end{minipage}
\end{center}

\pfl{过去周期分析} 以 25\% 的 ZigZag 阈值划分，样本区间内股价形成 3 个主升段与 3 个主跌段。 最大主升段出现在 2021 年 1 月 至 2022 年 2 月，13 个月+70.3\%；最大主跌段为 2022 年 2 月 至 2022 年 9 月，7 个月-55.7\%。 最近一次转折出现在 2026 年 7 月（下跌段自 2025 年 10 月起，9 个月-39.2\%），当前处于该段的延续之中。 公司股价较申万农林牧渔指数提前约 21 个月见底，是板块中较早定价周期反转的公司之一。 生猪养殖的定价主线是产能去化与成本曲线的赛跑（第 \ref{sec:hog} 章）：2026 年三季度行业仍处在“产能去化未完成、价格磨底”的阶段，公司 2026 年 6 月末资产负债率 81.0\%，高于全样本中位数 52.2\%，其股价因此更多反映资金面与信用风险，而不是价格弹性本身。 综合区间形态看，该股单段涨跌幅度偏大、以趋势段为主（最大主升 +70.3\%、最大主跌 -55.7\%），年化波动率 42.0\% 与全样本中位数 37.3\% 相比波动与样本中位数接近。

\pfl{盈利情况} 2026 年上半年营业收入 3.51 亿元（同比 -15.0\%），归母净利润 -0.54 亿元，由上年同期的盈利转为亏损。 以年报为趋势对照，2023---2025 年营业收入由 6.70 亿元变为 7.26 亿元（两年年均复合增速 +4.0\%），归母净利润由 -4.62 亿元变为 -1.07 亿元。 按半年报口径，2026 年上半年的 ROE 为 -14.6\%（未年化）、销售净利率 -14.7\%、毛利率 -1.0\%，ROE 在“生猪养殖”11 家公司中按数值由高到低排第\zhnumber{7}位。 同期全样本半年 ROE 中位数 0.75\%，公司与之相差 15.35 个百分点（弱于中位数公司）。 综合看，公司仍处亏损区间、由盈转亏，半年 ROE 在三级行业内按由高到低排第\zhnumber{7}位，单位盈利能力的修复依赖行业景气与自身成本端的改善，报表弹性主要体现为价格与销量的方向性变化。

\begin{center}\small\setlength{\tabcolsep}{3.4pt}
\captionof{table}[天域生物（603717）关键财务指标]{天域生物（603717）关键财务指标（合并报表口径；两个半年报期均未年化；现金短债比 $=$ 货币资金 $\div$ 短期债务）}
\begin{adjustbox}{max width=\textwidth}
\begin{tabular}{@{}lrrrrr@{}}
\toprule
指标 & 2023 年 & 2024 年 & 2025 年 & 2025 年半年报 & 2026 年半年报 \\
\midrule
营业收入（亿元） & 6.70 & 8.02 & 7.26 & 4.12 & 3.51 \\
归母净利润（亿元） & -4.62 & -1.07 & -1.07 & 0.11 & -0.54 \\
毛利率（\%） & -34.4 & 10.8 & -0.1 & 14.6 & -1.0 \\
ROE（\%） & -52.7 & -18.2 & -23.0 & 2.1 & -14.6 \\
资产负债率（\%） & 74.5 & 78.0 & 79.5 & 75.6 & 81.0 \\
经营活动现金流净额（亿元） & -1.76 & 1.35 & 2.58 & 0.90 & -0.05 \\
货币资金（亿元） & 1.52 & 1.78 & 1.53 & 1.58 & 1.35 \\
有息负债合计（亿元） & 11.80 & 11.44 & 10.41 & 11.09 & 10.49 \\
现金短债比（倍） & 0.56 & 0.64 & 0.54 & 0.56 & 0.49 \\
\bottomrule
\end{tabular}
\end{adjustbox}
\end{center}

\begin{center}
\begin{minipage}{\textwidth}\centering
\begin{tikzpicture}
\begin{groupplot}[
  group style={group size=2 by 1, horizontal sep=0.60cm},
  width=6.85cm, height=4.60cm,
  tick label style={font=\tiny},
  title style={font=\scriptsize\heiti},
  yticklabel style={font=\tiny},
  ylabel style={font=\tiny},
  grid=major, grid style={LightGray, line width=0.3pt},
  axis line style={InkGray, line width=0.5pt},
  enlarge x limits={abs=0.8},
  legend style={font=\scriptsize, at={(0.5,-0.20)}, anchor=north,
                legend columns=2, draw=none, fill=none, column sep=6pt},
]
\nextgroupplot[
  ybar, bar width=4pt,
  ymin=-6.02, ymax=11.14,
  xtick={0,1,2,3,4,5.7},
  xticklabels={2021,2022,2023,2024,2025,2026H1},
  ylabel={亿元},
]
\addplot[fill=AgriGreen!75, draw=AgriGreen, bar shift=-2.6pt] table[x=i, y=rev]{data/clean/profiles/fin_603717.dat};
\addplot[fill=SoftRed!60, draw=SoftRed, bar shift=2.6pt] table[x=i, y=ni]{data/clean/profiles/fin_603717.dat};
\addplot[fill=AgriGreen!30, draw=AgriGreen, pattern=north east lines, bar shift=-2.6pt] table[x=x, y=rev]{data/clean/profiles/finh_603717.dat};
\addplot[fill=SoftRed!20, draw=SoftRed, pattern=north east lines, bar shift=2.6pt] table[x=x, y=ni]{data/clean/profiles/finh_603717.dat};
\legend{营业收入（年/半年）, 归母净利润（年/半年）}
\nextgroupplot[
  ymin=-63.4, ymax=94.4,
  xtick={0,1,2,3,4,5.7},
  xticklabels={2021,2022,2023,2024,2025,2026H1},
  ylabel={\%},
]
\addplot[AgriGold, mark=*, mark size=1.3pt] table[x=i, y=roe]{data/clean/profiles/fin_603717.dat};
\addplot[OilBlue, mark=square*, mark size=1.3pt] table[x=i, y=debt]{data/clean/profiles/fin_603717.dat};
\addplot[only marks, AgriGold, mark=*, mark size=2.0pt] table[x=x, y=roe]{data/clean/profiles/finh_603717.dat};
\addplot[only marks, OilBlue, mark=square*, mark size=2.0pt] table[x=x, y=debt]{data/clean/profiles/finh_603717.dat};
\legend{ROE, 资产负债率}
\end{groupplot}
\end{tikzpicture}
\begin{tikzpicture}
\begin{axis}[
  name=finbot,
  width=13.75cm, height=3.10cm, scale only axis,
  ymin=-1.18, ymax=13.21,
  xtick={0,1,2,3,4,5.7},
  xticklabels={2021,2022,2023,2024,2025,2026H1},
  tick label style={font=\tiny},
  yticklabel style={font=\tiny},
  ylabel={亿元}, ylabel style={font=\tiny},
  ybar, bar width=4pt,
  grid=major, grid style={LightGray, line width=0.3pt},
  axis line style={InkGray, line width=0.5pt},
  enlarge x limits={abs=0.8},
  legend style={font=\scriptsize, at={(0.5,-0.26)}, anchor=north,
                legend columns=2, draw=none, fill=none, column sep=10pt},
]
\addplot[fill=AgriGreen!55, draw=AgriGreen, bar shift=-3.0pt] table[x=i, y=cash]{data/clean/profiles/fin_603717.dat};
\addplot[fill=SoftRed!45, draw=SoftRed, bar shift=3.0pt] table[x=i, y=ibd]{data/clean/profiles/fin_603717.dat};
\addplot[fill=AgriGreen!25, draw=AgriGreen, pattern=north east lines, bar shift=-3.0pt] table[x=x, y=cash]{data/clean/profiles/finh_603717.dat};
\addplot[fill=SoftRed!20, draw=SoftRed, pattern=north east lines, bar shift=3.0pt] table[x=x, y=ibd]{data/clean/profiles/finh_603717.dat};
\legend{货币资金（左轴）, 有息负债（左轴）}
\end{axis}
\begin{axis}[
  at={(finbot.south east)}, anchor=south east,
  width=13.75cm, height=3.10cm, scale only axis,
  axis y line*=right, axis x line*=none, xtick=\empty,
  ymin=-3.49, ymax=4.03,
  tick label style={font=\tiny},
  yticklabel style={font=\tiny}, scaled y ticks=false,
  legend style={font=\scriptsize, at={(0.5,-0.44)}, anchor=north,
                legend columns=1, draw=none, fill=none},
]
\addplot[OilBlue, mark=*, mark size=2.2pt, line width=1.3pt] table[x=i, y=ocf]{data/clean/profiles/fin_603717.dat};
\addplot[only marks, OilBlue, mark=*, mark size=3.2pt] table[x=x, y=ocf]{data/clean/profiles/finh_603717.dat};
\legend{经营活动现金流净额（右轴，亿元，蓝点）}
\end{axis}
\end{tikzpicture}

\captionof{figure}[天域生物（603717）近年主要财务指标]{天域生物（603717）近年主要财务指标（左上：营业收入与归母净利润；右上：ROE 与资产负债率；下：货币资金、有息负债为左轴柱状，经营活动现金流净额为其右轴的蓝点折线。
2021---2025 年为年报口径，2026H1 为 2026 年半年报/半年末，与其前一年报之间留出间隔、以斜纹或空心点区分；半年报数据未年化）}
\end{minipage}
\end{center}

\pfl{财务分析} 2026 年 6 月末资产负债率 81.0\%，较 2025 年末上升 1.5 个百分点（样本中位数 52.2\%，所处三级行业内按由低到高排第\zhnumber{9}位）。 2026 年 6 月末货币资金 1.35 亿元、短期债务（短期借款与一年内到期非流动负债）2.75 亿元、长期债务 7.73 亿元，有息负债合计 10.49 亿元，现金短债比 0.49 倍——覆盖偏紧。 净有息负债 9.13 亿元，相当于其 2026 年上半年营业收入的 260\%。 流动比率 0.90、速动比率 0.72，短期偿债指标偏紧。 2026 年上半年经营活动现金流净额 -0.05 亿元，经营现金流与利润同时为负，日常运营对债务与股东投入的依赖度较高。 综合看，账面货币资金对有息负债与短期债务的覆盖不足，滚动偿付对新增融资的依赖度较高；资产负债率高于样本中位数 28.8 个百分点；有息负债中长期债务占比更高，期限结构对短债滚动的压力较小。 公司披露的风险集中在股东与质押、监管与合规两类：一是 2026 年 4 月，控股股东、实际控制人罗卫国所持部分股份被司法冻结 312.93 万股、司法标记 690.00 万股；二是 2026 年 8 月，中国证监会重庆监管局作出行政处罚决定书〔2026〕6 号：罗卫国为实现减持目的通过基金份额转让交易方式隐匿一致行动关系、规避减持比例限制。

\pfl{重大战略转型} 公司的战略动作可归为以下几类：主业转型（1 项）：2019 年，2023 年生猪行情低迷时该业务前三季度毛利亏损超 7,300.00 万元；产能与项目（1 项）：2021 年 1 月，承租猪场改造投产后的年出栏量分别约 12 万头、18 万头。 公告口径上，近三年（2023 年 9 月---2026 年 9 月）公司披露重大重组与并购类公告 1 条、对外投资与转型类公告 2 条、回购与股权激励类公告 0 条，合计公告 536 条。 主营结构的口径变化已经落到报表上：2021 年年报的第一大产品为“园林生态工程”（92.6\%），2026 年半年报为“生态农牧食品业务”（90.6\%）；公司在这两期的分部划分方式有过调整。 综合判断，报告期内公司有 3 条与战略相关的公告落地（重大重组与并购 1 条、对外投资与转型 2 条），资本运作与主业调整之间存在直接关联。

\begin{center}\small\setlength{\tabcolsep}{4pt}
\captionof{table}[天域生物（603717）历史融资与资本运作]{天域生物（603717）历史融资与资本运作（据公司公告与公开报道整理；金额为披露值，含首次公开发行、再融资、债券、银行借款与重整投资等）}
\begin{adjustbox}{max width=\textwidth}
\begin{tabular}{@{}p{1.45cm}p{2.55cm}p{4.05cm}p{4.95cm}@{}}
\toprule
时间 & 方式 & 规模 & 用途与说明 \\
\midrule
2017-03 & IPO（上海证券交易所主板） & 发行不超过 4,\allowbreak{}317.79 万股，\allowbreak{}发行价 14.63 元/\allowbreak{}股 & 募投项目建设；发行前（2016 年 12 月 31 日）净资产 6.02 亿元；上市后 2017 年每 10 股派现 2.00 元（含税） \\
2021 & 非公开发行股票 & 筹资净额 3.94 亿元 & 募投项目建设；历次筹资情况表披露口径（中德证券关于公司 2022 年度向特定对象发行股票的核查文件） \\
2026-01 & 终止募投项目并将剩余募集资金永久补充流动资金 & 项目“天长市龙岗红色古镇文化旅游景区工程总承包项目”募集资金计划投资总额 2.76 亿元 & 剩余募集资金永久补充流动资金；2025 年 12 月 19 日第五届董事会第五次会议、\allowbreak{}2026 年 1 月 5 日 2026 年第一次临时股东会审议通过；2026 年 1 月 7 日剩余募集资金划转至自有资金账户 \\
2026-02 & 撤回向特定对象发行股票申请 & 原拟募集资金总额不超过 4.65 亿元（曾为不超过 5.27 亿元、\allowbreak{}发行不超过 8,\allowbreak{}041.76 万股） & 原拟用于补充流动资金和偿还债务；2026 年 2 月 11 日公司与申请 \\
\bottomrule
\end{tabular}
\end{adjustbox}
\end{center}

\pfl{历史融投资情况} 从现金流量表看，2025 年公司取得借款收到的现金 1.16 亿元、偿还债务支付的现金 1.73 亿元、吸收权益性投资 0.01 亿元，筹资活动现金流净额 -1.45 亿元（2023 年为取得借款 3.28 亿元、偿还债务 2.24 亿元）。 2026 年上半年筹资活动现金流净额为 -0.16 亿元（取得借款 0.40 亿元、偿还债务 0.32 亿元，未年化），已转为净流出，处于净还债状态。 2025 年分配股利、利润或偿付利息支付的现金合计 0.50 亿元。 公开披露的融资记录共 4 笔，工具类型以 IPO、定向增发为主；金额与用途见下表。 其中规模最大的两笔为：2017 年 3 月，IPO，发行不超过 4,317.79 万股，发行价 14.63 元/股；2021 年，定向增发，筹资净额 3.94 亿元。 股本方面，近三年披露股本变动 3 次；总股本自 2021 年末以来未发生变化（2.90 亿股）。 分红方面，近三年未实施现金分红，在全部样本中属于分红能力偏弱的一端。 近三年“再融资”类公告 30 条，涉及定向增发、募集资金使用。 综合看，公司融资结构上以债务滚动为主、股权融资为补充；2025 年筹资活动现金流净额 -1.45 亿元，融资节奏仍以维持存量债务周转为主，与前述经营与投资现金流的方向基本一致。

\pfl{未来潜在融资需求分析} 2026 年 6 月末货币资金 1.35 亿元，而短期债务 2.75 亿元（现金短债比 0.49 倍），2026 年上半年经营现金流净额 -0.05 亿元。按“短期债务减去账面货币资金”的滚动偿付口径，静态缺口约 1.40 亿元；若再计入上半年亏损的年化额（1.09 亿元），维持一年正常经营所需的外部资金约 1.40---2.49 亿元。 按第 \ref{sec:financing-need} 节的三档划分，公司属于第一档“补血型”：2026 年上半年归母净利润 -0.54 亿元、6 月末资产负债率 81.0\%。负债率已超过 80\%，净资产被亏损侵蚀的程度较深，属于全部样本中资产负债表最紧张的一组。 可行的融资路径按现实性排序为：债务重组或展期（与银行和债券持有人协商）、引入国资或产业投资人（部分公司以控制权让渡为对价）、定向增发（需股价配合，且控股股东需放弃优先认购或同步增资）。 公开披露的后续资金安排线索包括：股权融资线索：2025 年 12 月，公司推出向特定对象发行 A 股股票方案，拟募集资金总额不超过 4.65 亿元；债务与授信安排：2025 年年度，股东会审议向控股股东罗卫国申请不超过 1.00 亿元借款额度、向持股 5\% 以上股东史东伟申请不超过 1.00 亿元借款额度；2026 年 5 月，股东罗卫国和史东伟合计拟提供最高 1.50 亿元借款。 综合看，公司的外部资金需求以债务接续为主，滚动偿付口径下的静态缺口约 1.40 亿元，能否落地取决于授信续作与发行窗口的配合。






\section{饲料行业：成本传导的“夹心层”}
\label{sec:feed}

饲料板块 \num{15} 家公司（畜禽饲料 \num{11}、水产饲料 \num{4}），
总市值合计约 \num{1674} 亿元。按 2026 年半年报口径，
\hl{\num{10} 家亏损（亏损比例 \num{66.7}\%，为各板块最高）}，
半年 ROE 中位数 \num{-6.17}\%（\hl{全板块最低、且为唯一远低于零的中位数}），
\num{6} 家 2026 年 6 月末资产负债率超过 \num{65}\%。

\begingroup\scriptsize\setlength{\tabcolsep}{2pt}
\begin{longtable}{@{}llrrrrrrr@{}}
\caption{农业上市公司分行业明细（畜禽饲料、水产饲料；市值时点 2026-10-12，总市值口径；财务为 2026 年半年报，ROE 未年化，公告计数窗口 2023-09---2026-10）}\label{tab:comp-feed}\\
\toprule
代码 & 公司简称 & 总市值（亿元） & 营收（亿元） & 同比（\%） & 净利（亿元） & ROE（\%） & 再融资公告 & 重组公告 \\
\midrule
\endfirsthead
\multicolumn{9}{l}{\small（续）}\\
\toprule
代码 & 公司简称 & 总市值（亿元） & 营收（亿元） & 同比（\%） & 净利（亿元） & ROE（\%） & 再融资公告 & 重组公告 \\
\midrule
\endhead
\bottomrule
\endlastfoot
002311 & 海大集团 & 735.5 & 642.4 & 9.2 & 16.11 & 6.5 & 3 & 0 \\
000876 & 新希望 & 311.4 & 555.9 & 7.7 & -17.02 & -9.1 & 96 & 4 \\
002385 & 大北农 & 134.5 & 138.1 & 1.9 & -6.77 & -9.3 & 16 & 1 \\
002100 & 天康生物 & 105.4 & 83.3 & -5.9 & -4.41 & -6.6 & 2 & 5 \\
603609 & 禾丰股份 & 58.3 & 189.6 & 8.9 & -0.08 & -0.1 & 27 & 5 \\
002567 & 唐人神 & 52.0 & 108.8 & -12.7 & -10.46 & -24.9 & 12 & 4 \\
001313 & 粤海饲料 & 49.6 & 37.8 & 41.7 & 0.48 & 1.9 & 0 & 6 \\
603668 & 天马科技 & 47.3 & 31.3 & 5.3 & 0.12 & 0.6 & 26 & 0 \\
002548 & 金新农 & 30.8 & 23.2 & -2.4 & -3.88 & -37.4 & 9 & 0 \\
603151 & 邦基科技 & 30.7 & 26.8 & 11.8 & -0.00 & -0.0 & 4 & 16 \\
002696 & 百洋股份 & 27.9 & 20.5 & 21.1 & -0.83 & -6.2 & 0 & 0 \\
688156 & 路德科技 & 25.6 & 2.1 & 42.9 & 0.09 & 1.2 & 45 & 0 \\
603363 & 傲农生物 & 24.5 & 45.2 & 14.3 & -1.90 & -6.2 & 23 & 4 \\
001366 & 播恩集团 & 20.8 & 6.7 & 9.0 & 0.10 & 1.3 & 0 & 9 \\
000702 & 正虹科技 & 20.0 & 5.7 & 14.2 & -0.21 & -6.2 & 3 & 5 \\
\end{longtable}
\endgroup


\subsection{为什么饲料是最难的一环}
\label{sec:feed-squeeze}

饲料企业同时面对两条价格曲线：\hl{上游原料价格由全球市场决定（豆粕、玉米），
下游产品价格由养殖户的现金流决定}。当养殖端亏损时，
饲料企业面临“饲料需求下降 + 应收账款回收困难 + 降价促销”的三重挤压，
而原料成本不会同步下跌——这正是 2024---2025 年豆粕与玉米价格
在成本线附近横盘（第 \ref{sec:grain} 章）所导致的后果：
\emph{原料没有给饲料企业让出利润空间}。

板块内部的分化极为清晰：\textbf{海大集团}（水产饲料）2026 年上半年实现净利润 \num{16.11} 亿元、
半年 ROE \num{6.47}\%、营收 \num{642.4} 亿元（同比 \num{+9.2}\%），
总市值 \num{735.5} 亿元、居饲料板块首位，是板块内唯一的高质量成长企业；
\hl{其余 \num{14} 家公司合计亏损 \num{44.8} 亿元}，板块净利润总额因此被压缩至 $-28.7$ 亿元。
海大的成功来自\hl{“配方技术 + 服务营销 + 水产料高毛利”}的组合：
水产饲料的毛利率显著高于畜禽料（海大 2026 年上半年毛利率 \num{9.3}\%），
且海大通过苗种、动保、服务形成黏性，因此在水产养殖投苗量下降的年度仍能维持增长。

反观畜禽饲料：\textbf{新希望} 2026 年上半年营收 \num{555.9} 亿元（同比 \num{+7.7}\%）、
亏损 \num{17.02} 亿元（上年同期为 \num{+7.55} 亿元）；
\textbf{唐人神}亏损 \num{10.46} 亿元；\textbf{大北农}亏损 \num{6.77} 亿元。
三家合计亏损 \num{34.3} 亿元，占板块亏损额的三分之二以上。
\hl{其共同点是“饲料 + 养殖”的一体化扩张}：
在周期高点以债务融资扩大生猪产能，在周期底部同时承担饲料价格竞争
与生猪亏损的双重压力（三家 2026 年 6 月末资产负债率分别为
\num{70.7}\%、\num{71.1}\% 与 \num{69.9}\%）。

\subsection{从公告密度看饲料企业的融资压力}
\label{sec:feed-signals}

从公告行为可以观察到压力的分布：饲料板块近三年公司均“再融资相关公告”
为 \num{17.6} 条（中位数 \num{9} 条），高于种植业（\num{9.2} / \num{1}）
和加工业（\num{10.7} / \num{0}），其中 \textbf{新希望}以 \num{95} 条居全部样本公司之首。
\hl{公告密度是融资需求的先行指标}：定增、可转债、募集说明书、募投项目变更
等文件的密集披露，通常预示着公司正在通过资本市场解决资产负债表问题。

\begin{keybox}[饲料行业的画像结论]
\normalsize
\textbf{① 亏损比例最高的板块（\num{66.7}\%）}，半年 ROE 中位数为 \num{-6.17}\%，
是农业下行周期中\hl{受损最直接的中间环节}：上游成本刚性、下游需求收缩。\par\vspace{3pt}
\textbf{② 唯一的结构性赢家是水产料龙头}：海大集团一家半年利润 \num{16.1} 亿元，
与其余 \num{14} 家合计亏损 \num{44.8} 亿元形成完全的两极分化。\par\vspace{3pt}
\textbf{③ 一体化扩张的代价已经显性化}：新希望、唐人神、大北农三家半年合计亏损 \num{34.3} 亿元，
且均伴随高负债率与密集的再融资公告——\hl{这构成第二部分中最主要的潜在融资需求来源}。\par\vspace{3pt}
\textbf{④ 亏损面已扩大到三分之二}。板块 \num{15} 家公司 2026 年上半年
合计归母净利润 $-28.7$ 亿元（上年同期 $+45.8$ 亿元），\num{10} 家亏损，
\hl{生猪与禽类存栏去化对饲料需求的拖累正在加剧}。
\end{keybox}

\subsection{上市公司画像（\num{15} 家）}
\label{sec:co-feed}

本节给出饲料行业全部 \num{15} 家上市公司的画像，按总市值降序排列。
数据口径与第 \ref{sec:comp-overview} 节一致：\hl{财务数据统一采用 2026 年半年报}
（合并报表口径、未年化，资产负债率为 2026 年 6 月末），
同时列出 2025 年半年报作为同口径对照、2023---2025 年年报作为趋势对照；
股价为月度收盘价，最后一个交易日为 2026 年 10 月 12 日；周期分段使用与第一部分相同的
ZigZag 算法，阈值取 25\%（个股波动大于商品价格，故高于生猪现货的 20\%）；
公告计数窗口为 2023 年 9 月 1 日至 2026 年 10 月 12 日。

与前一版不同，本版画像正文由\hl{逐家检索的公开资料底稿}驱动：底稿只收录有来源链接的事实
（公司沿革、股权与控制权、经营要点、所处行业的基本面、竞争优势、融资历史、
风险事件、战略动作、融资需求线索），
正文按事实生成，各家公司的段落结构与重点并不相同。
每家公司在统一顺序下展开：基本情况与主营结构、股权结构与控制权、经营要点、
所处行业的基本面（供需）、与同类公司比较的竞争优势、历史股价、过去周期分析、
盈利情况、财务分析、重大战略转型、历史融投资情况、未来潜在融投资需求；
\hl{市值一律采用总市值口径}，股权结构取最新披露报告期前三大股东与股东户数，
行业基本面与竞争优势以逐家检索的公开资料为依据，并与三级行业内的同口径指标横向比较。
每家公司配图两张：
股价与行业指数（申万农林牧渔）的对比图，以及近年主要财务指标的组合图
（上排：营业收入与归母净利润、ROE 与资产负债率；下方宽带面板：货币资金与有息负债为左轴柱状、
经营活动现金流净额为其右轴的蓝点折线），并给出关键财务指标表与历史融资/资本运作表。
未来融资需求按两个口径粗略度量：\hl{短期债务缺口 $=$ 期末短期债务 $-$ 货币资金}
与\hl{年化经营缺口 $=$ 半年归母亏损 $\times 2$}；
前者衡量滚动偿付压力，后者衡量维持当期盈利水平的现金消耗，
两者都是压力情形下的静态估算，不构成对融资规模或时点的预测。

% 农业上市公司画像：畜禽饲料、水产饲料（共 15 家，按总市值降序）
% 由 scripts/make_company_profiles.py 自动生成，请勿手工修改

\subsubsection{海大集团（002311）}
\label{co:002311}

\pfl{公司基本情况} 海大集团成立于 2004 年，以水产饲料为核心、覆盖饲料—种苗—动保—养殖的全链饲料龙头，饲料销量与海外扩张领先；2026 年上半年生猪养殖亏损拖累盈利，增收不增利。 公司于 2009-11-27 在深交所主板首发上市，注册资本 16.57 亿元，总股本 16.64 亿股。 截至 2026 年 10 月 12 日收盘价 44.21 元，对应总市值 735.5 亿元，近一年-29.7\%，总市值在三级行业“水产饲料”的 4 家公司中按由大到小排第\zhnumber{1}位，近一年涨跌幅低于全样本中位数（-10.4\%）19.4 个百分点。 按 2026 年半年报披露的行业构成，收入占比前两位为饲料行业 87.6\%；养殖行业 12.4\%。

\pfl{股权结构与控制权} 按 2026 年 6 月 30 日的股东名册，前三大股东为广州市海灏投资有限公司（持股比例未披露）；香港中央结算有限公司（持股比例未披露）；全国社保基金一零六组合（持股比例未披露），股东户数 24,691 户。 股东结构与控制权方面，2025 年，公司分拆所属子公司海大国际控股有限公司至香港联交所主板上市的预案。

\pfl{经营要点} 经营与产能方面，2026 年上半年，饲料外销量 1478 万吨、同比增长约 8\%，其中水产料外销量 340 万吨（同比增约 21\%）、猪料 391 万吨（同比增约 15\%）、禽料 722 万吨。

\pfl{所处行业的基本面} 水产饲料的需求由投苗量与存塘鱼价决定，是饲料中毛利率最高、也最依赖服务的品类（第 \ref{sec:framework} 章的供需框架）。 行业层面的最新情况是：2026 年上半年，鱼粉价格从年初 16000 元/吨飙升至 24000 元/吨的历史最高价。

\pfl{竞争优势} 把该公司放回同类公司比较，2026 年上半年营业收入 642.4 亿元、在三级行业“水产饲料”的 4 家公司中排第\zhnumber{1}位；毛利率 9.3\%、高于行业内中位数 0.5 个百分点；半年 ROE 6.5\%（未年化）、高于行业内中位数（1.22\%）；资产负债率 54.6\%、低于行业内中位数（60.5\%）。 公司在公开资料中自述与同业比较的相对优势在于：2025 年上半年，公司实现饲料总销量 1470 万吨、同比增加 25\%，占全国饲料市场份额由 2024 年的 8.41\% 提升至 9.27\%。 优势在当期体现为高于行业中位数的资本回报，但优势能否持续仍取决于周期与自身产能投放节奏。

\pfl{历史股价} 本报告样本区间（2021 年 1 月 至 2026 年 10 月）内，股价由 65.81 元变为 44.21 元，累计 -32.8\%，区间最高 83.40 元（2021 年 4 月）、最低 37.80 元（2024 年 1 月）；当前价相当于区间最高价的 53.0\%，为区间最低价的 1.17 倍。 月度收益的年化波动率 29.9\%，低于全样本中位数 37.3\%，在三级行业“水产饲料”的 4 家公司中波动率由高到低排第\zhnumber{4}位。

\begin{center}
\begin{minipage}{\textwidth}\centering
\begin{tikzpicture}
\begin{axis}[
  width=12.6cm, height=3.9cm, scale only axis,
  xmin=2020.922, xmax=2026.972, ymin=35.15, ymax=89.24,
  xtick={2021,2022,2023,2024,2025,2026},
  yearticks,
  yticklabel style={font=\scriptsize},
  ylabel={元/股}, ylabel style={font=\scriptsize},
  grid=major, grid style={LightGray, line width=0.3pt},
  axis line style={InkGray, line width=0.5pt},
  every axis plot/.append style={line width=0.9pt},
]
\addplot[AgriGreen] table[x=x,y=y]{data/clean/profiles/px_002311.dat};
\addplot[SoftRed, dashed, line width=0.7pt] table[x=x,y=y]{data/clean/profiles/pxzz_002311.dat};
\addplot[only marks, mark=*, mark size=1.1pt, SoftRed] table[x=x,y=y]{data/clean/profiles/pxzz_002311.dat};
\end{axis}
\end{tikzpicture}

\captionof{figure}[海大集团（002311）股价与周期骨架]{海大集团（002311）月度收盘价（元/股）与 ZigZag 周期骨架（阈值 25\%，圆点为周期转折点）}
\end{minipage}
\end{center}

\pfl{过去周期分析} 以 25\% 的 ZigZag 阈值划分，样本区间内股价形成 2 个主升段与 2 个主跌段。 最大主升段出现在 2024 年 1 月 至 2025 年 9 月，20 个月+68.7\%；最大主跌段为 2021 年 4 月 至 2024 年 1 月，33 个月-54.7\%。 最近一次转折出现在 2026 年 10 月（下跌段自 2025 年 9 月起，13 个月-30.7\%），当前处于该段的延续之中。 公司股价较申万农林牧渔指数提前约 28 个月见底，是板块中较早定价周期反转的公司之一。 水产饲料的需求取决于投苗量与存塘鱼价，第 \ref{sec:other-livestock} 章显示 2026 年淡水鱼价格与投苗量已率先回升，其需求驱动与生猪存栏的联系弱于畜禽饲料，公司 2026 年上半年实现盈利、半年 ROE 6.5\%（未年化），好于全样本中位数（0.75\%），下行周期对其报表的冲击小于同业。 综合区间形态看，该股单段涨跌幅度偏大、以趋势段为主（最大主升 +68.7\%、最大主跌 -54.7\%），年化波动率 29.9\% 与全样本中位数 37.3\% 相比波动小于样本中位数。

\pfl{盈利情况} 2026 年上半年营业收入 642.40 亿元（同比 +9.2\%），归母净利润 16.11 亿元，同比 -39.0\%。 以年报为趋势对照，2023---2025 年营业收入由 1,161.17 亿元变为 1,284.68 亿元（两年年均复合增速 +5.2\%），归母净利润由 27.41 亿元变为 42.80 亿元。 按半年报口径，2026 年上半年的 ROE 为 6.5\%（未年化）、销售净利率 2.7\%、毛利率 9.3\%，ROE 在“水产饲料”4 家公司中按数值由高到低排第\zhnumber{1}位。 同期全样本半年 ROE 中位数 0.75\%，公司与之相差 5.72 个百分点（略好于中位数公司）。 综合看，公司在报告期维持盈利（高于全样本半年 ROE 中位数 0.75\%），利润的可持续性取决于水产饲料景气的延续与成本管控的执行。

\begin{center}\small\setlength{\tabcolsep}{3.4pt}
\captionof{table}[海大集团（002311）关键财务指标]{海大集团（002311）关键财务指标（合并报表口径；两个半年报期均未年化；现金短债比 $=$ 货币资金 $\div$ 短期债务）}
\begin{adjustbox}{max width=\textwidth}
\begin{tabular}{@{}lrrrrr@{}}
\toprule
指标 & 2023 年 & 2024 年 & 2025 年 & 2025 年半年报 & 2026 年半年报 \\
\midrule
营业收入（亿元） & 1,161.17 & 1,146.01 & 1,284.68 & 588.31 & 642.40 \\
归母净利润（亿元） & 27.41 & 45.04 & 42.80 & 26.39 & 16.11 \\
毛利率（\%） & 8.5 & 11.3 & 10.3 & 11.7 & 9.3 \\
ROE（\%） & 14.5 & 20.9 & 17.4 & 10.7 & 6.5 \\
资产负债率（\%） & 53.3 & 47.9 & 46.4 & 49.2 & 54.6 \\
经营活动现金流净额（亿元） & 126.98 & 79.96 & 62.60 & 30.48 & 39.24 \\
货币资金（亿元） & 54.76 & 34.78 & 32.03 & 34.81 & 34.25 \\
有息负债合计（亿元） & 70.99 & 52.67 & 40.82 & 44.63 & 64.39 \\
现金短债比（倍） & 2.23 & 2.35 & 1.91 & 2.19 & 0.84 \\
\bottomrule
\end{tabular}
\end{adjustbox}
\end{center}

\begin{center}
\begin{minipage}{\textwidth}\centering
\begin{tikzpicture}
\begin{groupplot}[
  group style={group size=2 by 1, horizontal sep=0.60cm},
  width=6.85cm, height=4.60cm,
  tick label style={font=\tiny},
  title style={font=\scriptsize\heiti},
  yticklabel style={font=\tiny},
  ylabel style={font=\tiny},
  grid=major, grid style={LightGray, line width=0.3pt},
  axis line style={InkGray, line width=0.5pt},
  enlarge x limits={abs=0.8},
  legend style={font=\scriptsize, at={(0.5,-0.20)}, anchor=north,
                legend columns=2, draw=none, fill=none, column sep=6pt},
]
\nextgroupplot[
  ybar, bar width=4pt,
  ymin=-128.47, ymax=1438.84,
  xtick={0,1,2,3,4,5.7},
  xticklabels={2021,2022,2023,2024,2025,2026H1},
  ylabel={亿元},
]
\addplot[fill=AgriGreen!75, draw=AgriGreen, bar shift=-2.6pt] table[x=i, y=rev]{data/clean/profiles/fin_002311.dat};
\addplot[fill=SoftRed!60, draw=SoftRed, bar shift=2.6pt] table[x=i, y=ni]{data/clean/profiles/fin_002311.dat};
\addplot[fill=AgriGreen!30, draw=AgriGreen, pattern=north east lines, bar shift=-2.6pt] table[x=x, y=rev]{data/clean/profiles/finh_002311.dat};
\addplot[fill=SoftRed!20, draw=SoftRed, pattern=north east lines, bar shift=2.6pt] table[x=x, y=ni]{data/clean/profiles/finh_002311.dat};
\legend{营业收入（年/半年）, 归母净利润（年/半年）}
\nextgroupplot[
  ymin=-4.5, ymax=61.9,
  xtick={0,1,2,3,4,5.7},
  xticklabels={2021,2022,2023,2024,2025,2026H1},
  ylabel={\%},
]
\addplot[AgriGold, mark=*, mark size=1.3pt] table[x=i, y=roe]{data/clean/profiles/fin_002311.dat};
\addplot[OilBlue, mark=square*, mark size=1.3pt] table[x=i, y=debt]{data/clean/profiles/fin_002311.dat};
\addplot[only marks, AgriGold, mark=*, mark size=2.0pt] table[x=x, y=roe]{data/clean/profiles/finh_002311.dat};
\addplot[only marks, OilBlue, mark=square*, mark size=2.0pt] table[x=x, y=debt]{data/clean/profiles/finh_002311.dat};
\legend{ROE, 资产负债率}
\end{groupplot}
\end{tikzpicture}
\begin{tikzpicture}
\begin{axis}[
  name=finbot,
  width=13.75cm, height=3.10cm, scale only axis,
  ymin=-10.85, ymax=121.48,
  xtick={0,1,2,3,4,5.7},
  xticklabels={2021,2022,2023,2024,2025,2026H1},
  tick label style={font=\tiny},
  yticklabel style={font=\tiny},
  ylabel={亿元}, ylabel style={font=\tiny},
  ybar, bar width=4pt,
  grid=major, grid style={LightGray, line width=0.3pt},
  axis line style={InkGray, line width=0.5pt},
  enlarge x limits={abs=0.8},
  legend style={font=\scriptsize, at={(0.5,-0.26)}, anchor=north,
                legend columns=2, draw=none, fill=none, column sep=10pt},
]
\addplot[fill=AgriGreen!55, draw=AgriGreen, bar shift=-3.0pt] table[x=i, y=cash]{data/clean/profiles/fin_002311.dat};
\addplot[fill=SoftRed!45, draw=SoftRed, bar shift=3.0pt] table[x=i, y=ibd]{data/clean/profiles/fin_002311.dat};
\addplot[fill=AgriGreen!25, draw=AgriGreen, pattern=north east lines, bar shift=-3.0pt] table[x=x, y=cash]{data/clean/profiles/finh_002311.dat};
\addplot[fill=SoftRed!20, draw=SoftRed, pattern=north east lines, bar shift=3.0pt] table[x=x, y=ibd]{data/clean/profiles/finh_002311.dat};
\legend{货币资金（左轴）, 有息负债（左轴）}
\end{axis}
\begin{axis}[
  at={(finbot.south east)}, anchor=south east,
  width=13.75cm, height=3.10cm, scale only axis,
  axis y line*=right, axis x line*=none, xtick=\empty,
  ymin=-33.02, ymax=165.08,
  tick label style={font=\tiny},
  yticklabel style={font=\tiny}, scaled y ticks=false,
  legend style={font=\scriptsize, at={(0.5,-0.44)}, anchor=north,
                legend columns=1, draw=none, fill=none},
]
\addplot[OilBlue, mark=*, mark size=2.2pt, line width=1.3pt] table[x=i, y=ocf]{data/clean/profiles/fin_002311.dat};
\addplot[only marks, OilBlue, mark=*, mark size=3.2pt] table[x=x, y=ocf]{data/clean/profiles/finh_002311.dat};
\legend{经营活动现金流净额（右轴，亿元，蓝点）}
\end{axis}
\end{tikzpicture}

\captionof{figure}[海大集团（002311）近年主要财务指标]{海大集团（002311）近年主要财务指标（左上：营业收入与归母净利润；右上：ROE 与资产负债率；下：货币资金、有息负债为左轴柱状，经营活动现金流净额为其右轴的蓝点折线。
2021---2025 年为年报口径，2026H1 为 2026 年半年报/半年末，与其前一年报之间留出间隔、以斜纹或空心点区分；半年报数据未年化）}
\end{minipage}
\end{center}

\pfl{财务分析} 2026 年 6 月末资产负债率 54.6\%，较 2025 年末上升 8.2 个百分点（样本中位数 52.2\%，所处三级行业内按由低到高排第\zhnumber{2}位）。 2026 年 6 月末货币资金 34.25 亿元、短期债务（短期借款与一年内到期非流动负债）40.97 亿元、长期债务 23.42 亿元，有息负债合计 64.39 亿元，现金短债比 0.84 倍——覆盖偏紧。 净有息负债 30.14 亿元，相当于其 2026 年上半年营业收入的 5\%。 流动比率 1.06、速动比率 0.55，短期指标尚可。 2026 年上半年经营活动现金流净额 39.24 亿元，经营现金流与利润同向为正，内生资金可以支撑运营。 2026 年 6 月末在建工程 8.71 亿元，较上年末增加 2.67 亿元，仍有在建投入。 2026 年 6 月末商誉 3.88 亿元，占归母净资产的 1.5\%，是净资产中最脆弱的一块。 综合看，现金短债比 0.84 倍，短期资金安排偏紧；资产负债率高于样本中位数 2.4 个百分点；有息负债中短期债务占比更高，期限结构仍以滚动续作为主。 公司披露的股东与质押风险为：2026 年 7 月，控股股东广州市海灏投资有限公司质押公司 320 万股；2026 年 9 月，控股股东海灏投资再质押 300 万股（占其所持 0.33\%）。

\pfl{重大战略转型} 近三年公司的战略动作集中在：其他动作（2 项）：2024 年，股票期权激励计划推进，公司调整该计划股票期权行权价格；2026 年 8 月，半年报提出“饲料、种苗和动保三体联动”核心业务模式，以全球化布局与技术服务优势对冲原料价格波动；海外与国际化（1 项）：2024---2026 年，2026 年上半年海外饲料销量增速超 25\%。 公告口径上，近三年（2023 年 9 月---2026 年 9 月）公司披露重大重组与并购类公告 0 条、对外投资与转型类公告 0 条、回购与股权激励类公告 54 条，合计公告 448 条。 第一大行业仍是“饲料行业”，收入占比 87.6\%，与 2021 年年报相比没有方向性变化。 综合判断，报告期内公司未披露实质性的重组、并购或转型安排，资本开支以维持现有产能与业务为主。

\begin{center}\small\setlength{\tabcolsep}{4pt}
\captionof{table}[海大集团（002311）历史融资与资本运作]{海大集团（002311）历史融资与资本运作（据公司公告与公开报道整理；金额为披露值，含首次公开发行、再融资、债券、银行借款与重整投资等）}
\begin{adjustbox}{max width=\textwidth}
\begin{tabular}{@{}p{1.45cm}p{2.55cm}p{4.05cm}p{4.95cm}@{}}
\toprule
时间 & 方式 & 规模 & 用途与说明 \\
\midrule
2009-11 & IPO（深交所中小板） & 发行 5,\allowbreak{}600 万股、\allowbreak{}28.00 元/\allowbreak{}股 & 上市日 2009-11-27，\allowbreak{}发行市盈率 48.39 倍 \\
2013-11 & 定向增发（非公开发行） & 6,\allowbreak{}850 万股、\allowbreak{}11.34 元/\allowbreak{}股，\allowbreak{}募集资金净额 7.58 亿元 & 新增股份 2013-12-06 上市，\allowbreak{}限售 12 个月 \\
2020-03 & 公开发行可转换公司债券 & 30.00 亿元（2,\allowbreak{}830 万张，\allowbreak{}面值 100.00 元） & 2020-03-19 发行、\allowbreak{}2020-04-16 上市 \\
2020-03 & 可转债（募集资金专户销户） & 28.30 亿元 & 2023-12-28 公告可转债募集资金专户销户完成，\allowbreak{}募集资金已使用完毕 \\
2024 & 现金分红（2024 年度） & 每 10 股派 11.00 元（含税），\allowbreak{}合计派现 18.30 亿元 & 股东回报；上市以来累计实施 19 次现金分红，\allowbreak{}形成“年度高分红+中期稳回馈”模式；2025 年半年报再推 10 派 2.00 元共 3.33 亿元 \\
2026-09 & 2026 年中期权益分派 & 每 10 股派现金红利 1.50 元（含税），\allowbreak{}回购专户 7.56 万股不参与分派 & 股东回报；2026 年 9 月 10 日发布实施公告 \\
\bottomrule
\end{tabular}
\end{adjustbox}
\end{center}

\pfl{历史融投资情况} 从现金流量表看，2025 年公司取得借款收到的现金 50.54 亿元、偿还债务支付的现金 65.82 亿元、吸收权益性投资 0.73 亿元，筹资活动现金流净额 -51.76 亿元（2023 年为取得借款 98.05 亿元、偿还债务 135.73 亿元）。 2026 年上半年筹资活动现金流净额为 -7.25 亿元（取得借款 65.37 亿元、偿还债务 42.72 亿元，未年化），已转为净流出，处于净还债状态。 2025 年分配股利、利润或偿付利息支付的现金合计 26.24 亿元。 公开披露的融资记录共 6 笔，工具类型以 IPO、定向增发、可转债为主；金额与用途见下表。 其中规模最大的两笔为：2020 年 3 月，可转债，30.00 亿元（2,830 万张，面值 100.00 元）；2020 年 3 月，可转债，28.30 亿元。 股本方面，近三年披露股本变动 6 次（期权行权 2 次、其他 1 次）；总股本由 2021 年末的 16.61 亿股变为 16.64 亿股（+0.16\%）。 分红方面，近三年现金分红 6 次、合计每股派现 3.50 元，最近一次实施在 2026 年 6 月。 近三年“再融资”类公告 3 条，涉及定向增发、可转债、募集资金使用。 综合看，公司融资结构上以债务滚动为主、股权融资为补充；2025 年筹资活动现金流净额 -51.76 亿元，融资节奏仍以维持存量债务周转为主，与前述经营与投资现金流的方向基本一致。

\pfl{未来潜在融资需求分析} 2026 年 6 月末货币资金 34.25 亿元，而短期债务 40.97 亿元（现金短债比 0.84 倍），2026 年上半年经营现金流净额 39.24 亿元。按“短期债务减去账面货币资金”的滚动偿付口径，静态缺口约 6.72 亿元；上半年盈利或接近盈亏平衡，该缺口不随经营亏损进一步扩大。 按第 \ref{sec:financing-need} 节的三档划分，公司属于第三档“扩张型”：2026 年上半年 ROE 6.5\%（未年化，折合约 12.9\%）、6 月末资产负债率 54.6\%，资产负债表仍具备加杠杆空间；融资用途以产能扩张、品类并购与海外布局为主，单笔规模通常在 5---20 亿元量级，会被市场定价为成长而非补血。 公开披露的后续资金安排线索包括：股权融资线索：2022 年 4 月，公司募集资金总额不超过 15.00 亿元；2023 年 3 月，广发证券出具 2022 年度向特定对象发行 A 股股票发行保荐书并取得证监会同意注册决定。 资金需求还要叠加在建工程：期末在建工程 8.71 亿元、2026 年上半年购建固定资产等支出 16.93 亿元（未年化），这部分支出在周期底部通常会被主动放缓。 综合看，公司的外部资金需求以补充流动性与项目投入为主，滚动偿付口径下的静态缺口约 6.72 亿元，能否落地取决于授信续作与发行窗口的配合。




\subsubsection{新希望（000876）}
\label{co:000876}

\pfl{公司基本情况} 1998 年设立至今，新希望是农牧全产业链龙头（饲料+生猪养殖与屠宰），2023 年底白羽肉禽与食品深加工引战出表以聚焦主业；2026 年上半年饲料创新高，猪价探底致亏损 17.02 亿元。 公司于 1998-03-11 在深交所主板通过重大资产重组（借壳）上市，注册资本 49.62 亿元，总股本 45.46 亿股。 截至 2026 年 10 月 12 日收盘价 6.85 元，对应总市值 311.4 亿元，近一年-31.0\%，总市值在三级行业“畜禽饲料”的 11 家公司中按由大到小排第\zhnumber{1}位，近一年涨跌幅低于全样本中位数（-10.4\%）20.7 个百分点。 按 2026 年半年报披露的行业构成，收入占比前三位为饲料 74.5\%；猪产业 23.3\%；其他 2.2\%。

\pfl{股权结构与控制权} 按 2026 年 8 月 21 日的股东名册，前三大股东为南方希望实业有限公司 持股 26.78\%；新希望集团有限公司 持股 22.77\%；西藏思壮投资咨询有限公司 持股 2.14\%。 控制权与股东性质方面，直接持股 25.09\%（11.30 亿股），与其一致行动人南方希望实业（29.51\%、13.29 亿股）合计持有 54.61\%；新希望集团质押 5.90 亿股；同花顺披露其对上市公司合计持股比例为 32.46\%。

\pfl{经营要点} 经营与产能方面，2026 年上半年，饲料外销量 1,316 万吨、同比增长 15\% 以上，其中海外外销 358 万吨、增速超过 25\%。

\pfl{所处行业的基本面} 饲料的供给即产能（门槛低、产能利用率偏低），需求直接来自下游存栏，利润则由玉米、豆粕等原料价格与成品售价的价差决定（第 \ref{sec:framework} 章的供需框架）。 与该公司直接相关的行业数据包括：2026 年上半年，全国工业饲料总产量 16,850 万吨、同比增长 3.9\%。

\pfl{竞争优势} 把该公司放回同类公司比较，2026 年上半年营业收入 555.9 亿元、在三级行业“畜禽饲料”的 11 家公司中排第\zhnumber{1}位；毛利率 4.7\%、低于行业内中位数 0.5 个百分点；半年 ROE -9.1\%（未年化）、低于行业内中位数（-6.25\%）；资产负债率 70.7\%、高于行业内中位数（62.1\%）。 与同类公司相比，公司的差异化在于：15\% 以上的外销增速明显高于全国饲料总产量 3.9\% 的行业平均；饲料业务在多数东道国销量排名前三；持续压降产能利用率低于 80\% 的低效公司，并通过总部集采与区域自采结合的采购体系在原料端锁定成本。 就当期报表看，公司尚未把经营层面的特征转化为高于行业的资本回报，周期底部的相对优势更多体现在现金流与资产负债表上。

\pfl{历史股价} 本报告样本区间（2021 年 1 月 至 2026 年 10 月）内，股价由 21.15 元变为 6.85 元，累计 -67.6\%，区间最高 24.07 元（2021 年 2 月）、最低 6.27 元（2026 年 6 月）；当前价相当于区间最高价的 28.5\%，为区间最低价的 1.09 倍。 月度收益的年化波动率 31.3\%，低于全样本中位数 37.3\%，在三级行业“畜禽饲料”的 11 家公司中波动率由高到低排第\zhnumber{10}位。 把股价阶段与公司事件对齐看，2022 年 3 月 至 2024 年 1 月股价下跌-51.9\%，同期（2023 年 12 月）董事会通过白羽肉禽业务引进战略投资者议案：中国牧工商集团有限公司以现金约 27.00 亿元受让禽产业链运营主体山东中新食品集团 51\% 股权。 整体上看，区间的几处大级别拐点都能在公司的资本运作、股权变动或风险事件上找到对应，股价波动有明显的公司层面驱动，而非单纯的行业 beta。

\begin{center}
\begin{minipage}{\textwidth}\centering
\begin{tikzpicture}
\begin{axis}[
  width=12.6cm, height=3.9cm, scale only axis,
  xmin=2020.922, xmax=2026.972, ymin=5.83, ymax=25.75,
  xtick={2021,2022,2023,2024,2025,2026},
  yearticks,
  yticklabel style={font=\scriptsize},
  ylabel={元/股}, ylabel style={font=\scriptsize},
  grid=major, grid style={LightGray, line width=0.3pt},
  axis line style={InkGray, line width=0.5pt},
  every axis plot/.append style={line width=0.9pt},
]
\addplot[AgriGreen] table[x=x,y=y]{data/clean/profiles/px_000876.dat};
\addplot[SoftRed, dashed, line width=0.7pt] table[x=x,y=y]{data/clean/profiles/pxzz_000876.dat};
\addplot[only marks, mark=*, mark size=1.1pt, SoftRed] table[x=x,y=y]{data/clean/profiles/pxzz_000876.dat};
\end{axis}
\end{tikzpicture}

\captionof{figure}[新希望（000876）股价与周期骨架]{新希望（000876）月度收盘价（元/股）与 ZigZag 周期骨架（阈值 25\%，圆点为周期转折点）}
\end{minipage}
\end{center}

\pfl{过去周期分析} 以 25\% 的 ZigZag 阈值划分，样本区间内股价形成 2 个主升段与 3 个主跌段。 最大主升段出现在 2021 年 8 月 至 2022 年 3 月，7 个月+55.2\%；最大主跌段为 2022 年 3 月 至 2024 年 1 月，22 个月-51.9\%。 最近一次转折出现在 2026 年 6 月（下跌段自 2024 年 9 月起，21 个月-39.8\%），当前处于该段的延续之中。 公司股价与申万农林牧渔指数几乎同步见底，个股并未走出独立行情。 饲料夹在原料成本与养殖需求之间（第 \ref{sec:feed} 节），量的部分取决于存栏、利的部分取决于原料与成品的价差，公司 2026 年 6 月末资产负债率 70.7\%，高于全样本中位数 52.2\%，其股价因此更多反映资金面与信用风险，而不是价格弹性本身。 综合区间形态看，该股单段涨跌幅度偏大、以趋势段为主（最大主升 +55.2\%、最大主跌 -51.9\%），年化波动率 31.3\% 与全样本中位数 37.3\% 相比波动小于样本中位数。

\pfl{盈利情况} 2026 年上半年营业收入 555.89 亿元（同比 +7.7\%），归母净利润 -17.02 亿元，由上年同期的盈利转为亏损。 以年报为趋势对照，2023---2025 年营业收入由 1,417.03 亿元变为 1,068.56 亿元（两年年均复合增速 -13.2\%），归母净利润由 2.49 亿元变为 -17.84 亿元。 按半年报口径，2026 年上半年的 ROE 为 -9.1\%（未年化）、销售净利率 -3.0\%、毛利率 4.7\%，ROE 在“畜禽饲料”11 家公司中按数值由高到低排第\zhnumber{8}位。 按同一口径，三级行业“畜禽饲料”11 家公司的半年 ROE 中位数为 -6.25\%，该公司低于其中位数 2.87 个百分点。 公司披露的应对举措集中在：2026 年 8 月，董事会通过使用向特定对象发行募集资金置换预先投入募投项目的自筹资金 7.95 亿元及已支付发行费用 92.11 万元。 综合看，公司仍处亏损区间、由盈转亏，半年 ROE 在三级行业内按由高到低排第\zhnumber{8}位，单位盈利能力的修复依赖行业景气与自身成本端的改善，报表弹性主要体现为价格与销量的方向性变化。

\begin{center}\small\setlength{\tabcolsep}{3.4pt}
\captionof{table}[新希望（000876）关键财务指标]{新希望（000876）关键财务指标（合并报表口径；两个半年报期均未年化；现金短债比 $=$ 货币资金 $\div$ 短期债务）}
\begin{adjustbox}{max width=\textwidth}
\begin{tabular}{@{}lrrrrr@{}}
\toprule
指标 & 2023 年 & 2024 年 & 2025 年 & 2025 年半年报 & 2026 年半年报 \\
\midrule
营业收入（亿元） & 1,417.03 & 1,030.63 & 1,068.56 & 516.25 & 555.89 \\
归母净利润（亿元） & 2.49 & 4.74 & -17.84 & 7.55 & -17.02 \\
毛利率（\%） & 2.8 & 6.7 & 6.3 & 7.9 & 4.7 \\
ROE（\%） & 0.7 & 1.7 & -8.8 & 2.9 & -9.1 \\
资产负债率（\%） & 72.3 & 69.0 & 69.9 & 68.8 & 70.7 \\
经营活动现金流净额（亿元） & 139.04 & 91.27 & 93.78 & 51.53 & 16.46 \\
货币资金（亿元） & 108.50 & 86.99 & 86.27 & 70.72 & 79.63 \\
有息负债合计（亿元） & 639.65 & 614.09 & 582.06 & 589.73 & 579.09 \\
现金短债比（倍） & 0.41 & 0.30 & 0.37 & 0.25 & 0.39 \\
\bottomrule
\end{tabular}
\end{adjustbox}
\end{center}

\begin{center}
\begin{minipage}{\textwidth}\centering
\begin{tikzpicture}
\begin{groupplot}[
  group style={group size=2 by 1, horizontal sep=0.60cm},
  width=6.85cm, height=4.60cm,
  tick label style={font=\tiny},
  title style={font=\scriptsize\heiti},
  yticklabel style={font=\tiny},
  ylabel style={font=\tiny},
  grid=major, grid style={LightGray, line width=0.3pt},
  axis line style={InkGray, line width=0.5pt},
  enlarge x limits={abs=0.8},
  legend style={font=\scriptsize, at={(0.5,-0.20)}, anchor=north,
                legend columns=2, draw=none, fill=none, column sep=6pt},
]
\nextgroupplot[
  ybar, bar width=4pt,
  ymin=-247.20, ymax=1598.59,
  xtick={0,1,2,3,4,5.7},
  xticklabels={2021,2022,2023,2024,2025,2026H1},
  ylabel={亿元},
]
\addplot[fill=AgriGreen!75, draw=AgriGreen, bar shift=-2.6pt] table[x=i, y=rev]{data/clean/profiles/fin_000876.dat};
\addplot[fill=SoftRed!60, draw=SoftRed, bar shift=2.6pt] table[x=i, y=ni]{data/clean/profiles/fin_000876.dat};
\addplot[fill=AgriGreen!30, draw=AgriGreen, pattern=north east lines, bar shift=-2.6pt] table[x=x, y=rev]{data/clean/profiles/finh_000876.dat};
\addplot[fill=SoftRed!20, draw=SoftRed, pattern=north east lines, bar shift=2.6pt] table[x=x, y=ni]{data/clean/profiles/finh_000876.dat};
\legend{营业收入（年/半年）, 归母净利润（年/半年）}
\nextgroupplot[
  ymin=-36.6, ymax=82.4,
  xtick={0,1,2,3,4,5.7},
  xticklabels={2021,2022,2023,2024,2025,2026H1},
  ylabel={\%},
]
\addplot[AgriGold, mark=*, mark size=1.3pt] table[x=i, y=roe]{data/clean/profiles/fin_000876.dat};
\addplot[OilBlue, mark=square*, mark size=1.3pt] table[x=i, y=debt]{data/clean/profiles/fin_000876.dat};
\addplot[only marks, AgriGold, mark=*, mark size=2.0pt] table[x=x, y=roe]{data/clean/profiles/finh_000876.dat};
\addplot[only marks, OilBlue, mark=square*, mark size=2.0pt] table[x=x, y=debt]{data/clean/profiles/finh_000876.dat};
\legend{ROE, 资产负债率}
\end{groupplot}
\end{tikzpicture}
\begin{tikzpicture}
\begin{axis}[
  name=finbot,
  width=13.75cm, height=3.10cm, scale only axis,
  ymin=-69.36, ymax=776.85,
  xtick={0,1,2,3,4,5.7},
  xticklabels={2021,2022,2023,2024,2025,2026H1},
  tick label style={font=\tiny},
  yticklabel style={font=\tiny},
  ylabel={亿元}, ylabel style={font=\tiny},
  ybar, bar width=4pt,
  grid=major, grid style={LightGray, line width=0.3pt},
  axis line style={InkGray, line width=0.5pt},
  enlarge x limits={abs=0.8},
  legend style={font=\scriptsize, at={(0.5,-0.26)}, anchor=north,
                legend columns=2, draw=none, fill=none, column sep=10pt},
]
\addplot[fill=AgriGreen!55, draw=AgriGreen, bar shift=-3.0pt] table[x=i, y=cash]{data/clean/profiles/fin_000876.dat};
\addplot[fill=SoftRed!45, draw=SoftRed, bar shift=3.0pt] table[x=i, y=ibd]{data/clean/profiles/fin_000876.dat};
\addplot[fill=AgriGreen!25, draw=AgriGreen, pattern=north east lines, bar shift=-3.0pt] table[x=x, y=cash]{data/clean/profiles/finh_000876.dat};
\addplot[fill=SoftRed!20, draw=SoftRed, pattern=north east lines, bar shift=3.0pt] table[x=x, y=ibd]{data/clean/profiles/finh_000876.dat};
\legend{货币资金（左轴）, 有息负债（左轴）}
\end{axis}
\begin{axis}[
  at={(finbot.south east)}, anchor=south east,
  width=13.75cm, height=3.10cm, scale only axis,
  axis y line*=right, axis x line*=none, xtick=\empty,
  ymin=-36.15, ymax=180.75,
  tick label style={font=\tiny},
  yticklabel style={font=\tiny}, scaled y ticks=false,
  legend style={font=\scriptsize, at={(0.5,-0.44)}, anchor=north,
                legend columns=1, draw=none, fill=none},
]
\addplot[OilBlue, mark=*, mark size=2.2pt, line width=1.3pt] table[x=i, y=ocf]{data/clean/profiles/fin_000876.dat};
\addplot[only marks, OilBlue, mark=*, mark size=3.2pt] table[x=x, y=ocf]{data/clean/profiles/finh_000876.dat};
\legend{经营活动现金流净额（右轴，亿元，蓝点）}
\end{axis}
\end{tikzpicture}

\captionof{figure}[新希望（000876）近年主要财务指标]{新希望（000876）近年主要财务指标（左上：营业收入与归母净利润；右上：ROE 与资产负债率；下：货币资金、有息负债为左轴柱状，经营活动现金流净额为其右轴的蓝点折线。
2021---2025 年为年报口径，2026H1 为 2026 年半年报/半年末，与其前一年报之间留出间隔、以斜纹或空心点区分；半年报数据未年化）}
\end{minipage}
\end{center}

\pfl{财务分析} 2026 年 6 月末资产负债率 70.7\%，较 2025 年末上升 0.8 个百分点（样本中位数 52.2\%，所处三级行业内按由低到高排第\zhnumber{9}位）。 2026 年 6 月末货币资金 79.63 亿元、短期债务（短期借款与一年内到期非流动负债）205.48 亿元、长期债务 373.61 亿元，有息负债合计 579.09 亿元，现金短债比 0.39 倍——覆盖偏紧。 净有息负债 499.46 亿元，相当于其 2026 年上半年营业收入的 90\%。 流动比率 0.55、速动比率 0.28，短期偿债指标偏紧。 2026 年上半年经营活动现金流净额 16.46 亿元，净利润为负而经营现金流为正，差额主要来自折旧摊销与营运资本变动，说明亏损尚未完全转化为现金流出。 2026 年 6 月末在建工程 35.71 亿元，较上年末增加 3.71 亿元，仍有在建投入。 2026 年 6 月末商誉 11.27 亿元，占归母净资产的 3.3\%，是净资产中最脆弱的一块。 综合看，账面货币资金对有息负债与短期债务的覆盖不足，滚动偿付对新增融资的依赖度较高；资产负债率高于样本中位数 18.5 个百分点；有息负债中长期债务占比更高，期限结构对短债滚动的压力较小。

\pfl{重大战略转型} 公司的战略动作可归为以下几类：并购与重组（2 项）：2023 年 12 月，董事会通过白羽肉禽业务引进战略投资者议案：中国牧工商集团有限公司以现金约 27.00 亿元受让禽产业链运营主体山东中新食品集团 51\% 股权。 公告口径上，近三年（2023 年 9 月---2026 年 9 月）公司披露重大重组与并购类公告 4 条、对外投资与转型类公告 5 条、回购与股权激励类公告 28 条，合计公告 676 条。 第一大行业“饲料”的收入占比由 2021 年年报的 56.1\% 变为 74.5\%（18.5 个百分点），收入结构出现再平衡。 综合判断，报告期内公司有 9 条与战略相关的公告落地（重大重组与并购 4 条、对外投资与转型 5 条），资本运作与主业调整之间存在直接关联。

\begin{center}\small\setlength{\tabcolsep}{4pt}
\captionof{table}[新希望（000876）历史融资与资本运作]{新希望（000876）历史融资与资本运作（据公司公告与公开报道整理；金额为披露值，含首次公开发行、再融资、债券、银行借款与重整投资等）}
\begin{adjustbox}{max width=\textwidth}
\begin{tabular}{@{}p{1.45cm}p{2.55cm}p{4.05cm}p{4.95cm}@{}}
\toprule
时间 & 方式 & 规模 & 用途与说明 \\
\midrule
1998-03 & IPO（深交所） & 公众股 3,\allowbreak{}600 万股、\allowbreak{}内部职工股 400 万股 & 公司原名四川新希望农业股份有限公司 \\
2011 & 发行股份购买资产 & 新增股份 9.05 亿股、\allowbreak{}8.00 元/\allowbreak{}股 & 收购饲料及农牧相关资产（资产出售、\allowbreak{}资产置换及发行股份购买资产一揽子方案）；2011-10-14 完成股份登记不构成借壳上市 \\
2020-01 & 公开发行可转换公司债券（“希望转债”127015） & 40.00 亿元（4,\allowbreak{}000 万张，\allowbreak{}面值 100.00 元），\allowbreak{}期限 6 年 & 2020-01-03 发行、\allowbreak{}2020-02-04 起深交所挂牌 \\
2021-11 & 公开发行可转换公司债券（“希望转2”127049） & 50.00 亿元（8,\allowbreak{}150 万张，\allowbreak{}面值 100.00 元） & 投资生猪养殖项目 57.05 亿元 \\
2023-12 & 资产出售（白羽肉禽与食品深加工引战） & 两笔交易金额合计约 42.00 亿元（27.00 亿元 + 15.01 亿元） & 聚焦核心主业、\allowbreak{}降低资产负债率；2023 年该两项引战影响归母净利润增加约 51.00 亿元至 52.00 亿元 \\
2026-08 & 向特定对象发行股票（定向增发） & 4.59 亿股、\allowbreak{}5.66 元/\allowbreak{}股，\allowbreak{}募集资金总额 26.00 亿元（26.00 亿元） & 猪场生物安全防控及数智化升级项目、\allowbreak{}偿还银行债务；；26 家机构投资者认购，\allowbreak{}总股本增至 49.62 亿股 \\
\bottomrule
\end{tabular}
\end{adjustbox}
\end{center}

\pfl{历史融投资情况} 从现金流量表看，2025 年公司取得借款收到的现金 464.49 亿元、偿还债务支付的现金 518.94 亿元、吸收权益性投资 37.32 亿元，筹资活动现金流净额 -82.59 亿元（2023 年为取得借款 432.13 亿元、偿还债务 485.38 亿元）。 2026 年上半年筹资活动现金流净额为 -13.37 亿元（取得借款 208.32 亿元、偿还债务 229.84 亿元，未年化），已转为净流出，处于净还债状态。 2025 年分配股利、利润或偿付利息支付的现金合计 22.38 亿元。 公开披露的融资记录共 6 笔，工具类型以 IPO、发行股份购买资产、可转债、定向增发为主；金额与用途见下表。 其中规模最大的两笔为：2026 年 8 月，定向增发，4.59 亿股、5.66 元/股，募集资金总额 26.00 亿元（26.00 亿元）；2021 年 11 月，可转债，50.00 亿元（8,150 万张，面值 100.00 元）。 股本方面，近三年披露股本变动 16 次（可转债转股 11 次、股份回购 2 次、增发新股上市 1 次、股权激励 1 次）；总股本由 2021 年末的 45.05 亿股变为 45.46 亿股（+0.90\%）。 分红方面，近三年现金分红 1 次、合计每股派现 0.02 元，最近一次实施在 2024 年 12 月。 近三年“再融资”类公告 96 条，涉及定向增发、可转债、募集资金使用。 综合看，公司融资结构上以债务滚动为主、股权融资为补充；2025 年筹资活动现金流净额 -82.59 亿元，融资节奏仍以维持存量债务周转为主，与前述经营与投资现金流的方向基本一致。

\pfl{未来潜在融资需求分析} 2026 年 6 月末货币资金 79.63 亿元，而短期债务 205.48 亿元（现金短债比 0.39 倍），2026 年上半年经营现金流净额 16.46 亿元。按“短期债务减去账面货币资金”的滚动偿付口径，静态缺口约 125.85 亿元；若再计入上半年亏损的年化额（34.04 亿元），维持一年正常经营所需的外部资金约 125.85---159.89 亿元。 按第 \ref{sec:financing-need} 节的三档划分，公司属于第一档“补血型”：2026 年上半年归母净利润 -17.02 亿元、6 月末资产负债率 70.7\%。 可行的融资路径按现实性排序为：引入国资或产业投资人（部分公司以控制权让渡为对价）、定向增发（需股价配合，且控股股东需放弃优先认购或同步增资）、出售非核心资产与参股股权。 公开披露的后续资金安排线索包括：股权融资线索：2023 年 12 月，定增方案三次调整：2023 年 12 月预案拟募不超过 73.50 亿元 → 2024 年 8 月修订为不超过 38.00 亿元；2026 年 8 月，规范 26.00 亿元定增募集资金存放与使用；债务与授信安排：2025 年，实际控制人下辖财务公司给予公司综合授信额度 130.00 亿元，为流动性补充；2026 年 6 月，截至评级时公司获得授信总额 798.17 亿元、尚未使用额度 337.96 亿元。 资金需求还要叠加在建工程：期末在建工程 35.71 亿元、2026 年上半年购建固定资产等支出 11.35 亿元（未年化），这部分支出在周期底部通常会被主动放缓。 综合看，公司的外部资金需求以债务接续为主，滚动偿付口径下的静态缺口约 125.85 亿元，能否落地取决于授信续作与发行窗口的配合。




\subsubsection{大北农（002385）}
\label{co:002385}

\pfl{公司基本情况} 2010 年设立至今，大北农是以饲料、种业（含转基因生物育种）、生猪养殖与动保为主业的农业科技公司；2026 年上半年猪价低迷致亏损 6.77 亿元，实控人邵根伙去世后控制权由配偶莫云继承。 公司于 2010-04-09 在深交所主板首发上市，注册资本 42.80 亿元，总股本 43.24 亿股。 截至 2026 年 10 月 12 日收盘价 3.11 元，对应总市值 134.5 亿元，近一年-24.9\%，总市值在三级行业“畜禽饲料”的 11 家公司中按由大到小排第\zhnumber{2}位，近一年涨跌幅低于全样本中位数（-10.4\%）14.5 个百分点。 按 2026 年半年报披露的产品构成，收入占比前三位为饲料产品 67.7\%；养猪产品 21.7\%；其他产品 7.3\%。

\pfl{股权结构与控制权} 按 2026 年 9 月 28 日的股东名册，前三大股东为莫云 持股 16.26\%；邵某某（持股比例未披露）；浙商银行股份有限公司－国泰中证畜牧养殖交易型开放式指 持股 1.75\%。 股东结构与控制权方面，持股 9.28 亿股（占 21.67\%）经法定继承过户；邵根伙名下仍有 7,151.52 万股（占总股本 1.67\%）因质押未解除而未完成过户。

\pfl{经营要点} 经营层面，2026 年上半年，饲料销量 281.21 万吨、同比增长 9.21\%；2026 年上半年，玉米、水稻、大豆合计销售 2,500.92 万公斤、同比增长 12.44\%，其中水稻种子销售 2,097.18 万公斤、收入 3.43 亿元（+6.44\%）；2026 年 6 月，回升至 10.72 元/公斤、当月商品肥猪销量 25.91 万头商品猪销售均价同比下降约 30\%。

\pfl{所处行业的基本面} 饲料的供给即产能（门槛低、产能利用率偏低），需求直接来自下游存栏，利润则由玉米、豆粕等原料价格与成品售价的价差决定（第 \ref{sec:framework} 章的供需框架）。 与该公司直接相关的行业数据包括：2026 年上半年，全国工业饲料总产量 16,850 万吨、同比增长 3.9\%。

\pfl{竞争优势} 把该公司放回同类公司比较，2026 年上半年营业收入 138.1 亿元、在三级行业“畜禽饲料”的 11 家公司中排第\zhnumber{3}位；毛利率 8.7\%、高于行业内中位数 3.4 个百分点；半年 ROE -9.3\%（未年化）、低于行业内中位数（-6.25\%）；资产负债率 69.9\%、高于行业内中位数（62.1\%）。 公司区别于同类公司的经营特征包括：全国已审定的 160 个转基因玉米品种中 95 个使用公司性状、约占 59\%；2026 年 5 月，公司性状审定品种数与除草剂登记数均列行业第一。 优势在报表上主要体现为高于行业平均的毛利率，这在周期底部是抵御亏损的主要缓冲。

\pfl{历史股价} 自 2021 年 1 月至 2026 年 10 月，股价由 9.27 元变为 3.11 元，累计 -66.5\%，区间最高 10.49 元（2021 年 12 月）、最低 2.76 元（2026 年 6 月）；当前价相当于区间最高价的 29.6\%，为区间最低价的 1.13 倍。 月度收益的年化波动率 33.1\%，低于全样本中位数 37.3\%，在三级行业“畜禽饲料”的 11 家公司中波动率由高到低排第\zhnumber{9}位。

\begin{center}
\begin{minipage}{\textwidth}\centering
\begin{tikzpicture}
\begin{axis}[
  width=12.6cm, height=3.9cm, scale only axis,
  xmin=2020.922, xmax=2026.972, ymin=2.57, ymax=11.22,
  xtick={2021,2022,2023,2024,2025,2026},
  yearticks,
  yticklabel style={font=\scriptsize},
  ylabel={元/股}, ylabel style={font=\scriptsize},
  grid=major, grid style={LightGray, line width=0.3pt},
  axis line style={InkGray, line width=0.5pt},
  every axis plot/.append style={line width=0.9pt},
]
\addplot[AgriGreen] table[x=x,y=y]{data/clean/profiles/px_002385.dat};
\addplot[SoftRed, dashed, line width=0.7pt] table[x=x,y=y]{data/clean/profiles/pxzz_002385.dat};
\addplot[only marks, mark=*, mark size=1.1pt, SoftRed] table[x=x,y=y]{data/clean/profiles/pxzz_002385.dat};
\end{axis}
\end{tikzpicture}

\captionof{figure}[大北农（002385）股价与周期骨架]{大北农（002385）月度收盘价（元/股）与 ZigZag 周期骨架（阈值 25\%，圆点为周期转折点）}
\end{minipage}
\end{center}

\pfl{过去周期分析} 以 25\% 的 ZigZag 阈值划分，样本区间内股价形成 2 个主升段与 3 个主跌段。 最大主升段出现在 2021 年 6 月 至 2021 年 12 月，6 个月+54.3\%；最大主跌段为 2022 年 11 月 至 2026 年 6 月，43 个月-70.4\%。 最近一次转折出现在 2026 年 6 月（下跌段自 2022 年 11 月起，43 个月-70.4\%），当前处于该段的延续之中。 公司股价与申万农林牧渔指数几乎同步见底，个股并未走出独立行情。 饲料夹在原料成本与养殖需求之间（第 \ref{sec:feed} 节），量的部分取决于存栏、利的部分取决于原料与成品的价差，公司 2026 年 6 月末资产负债率 69.9\%，高于全样本中位数 52.2\%，其股价因此更多反映资金面与信用风险，而不是价格弹性本身。 综合区间形态看，该股单段涨跌幅度偏大、以趋势段为主（最大主升 +54.3\%、最大主跌 -70.4\%），年化波动率 33.1\% 与全样本中位数 37.3\% 相比波动小于样本中位数。

\pfl{盈利情况} 2026 年上半年营业收入 138.11 亿元（同比 +1.9\%），归母净利润 -6.77 亿元，由上年同期的盈利转为亏损。 以年报为趋势对照，2023---2025 年营业收入由 333.90 亿元变为 291.19 亿元（两年年均复合增速 -6.6\%），归母净利润由 -21.74 亿元变为 -6.42 亿元。 按半年报口径，2026 年上半年的 ROE 为 -9.3\%（未年化）、销售净利率 -6.0\%、毛利率 8.7\%，ROE 在“畜禽饲料”11 家公司中按数值由高到低排第\zhnumber{9}位。 同期全样本半年 ROE 中位数 0.75\%，公司与之相差 10.04 个百分点（弱于中位数公司）。 从公司披露的原因看，业绩变动主要来自：2026 年 8 月，公司投资者关系活动披露：以饲料、种业、动保、食品多板块协同应对周期波动。 面对这一局面，公司披露的举措包括：2026 年 7 月，实施股份回购，并对部分限制性股票回购注销、减少注册资本。 综合看，公司仍处亏损区间、由盈转亏，半年 ROE 在三级行业内按由高到低排第\zhnumber{9}位，单位盈利能力的修复依赖行业景气与自身成本端的改善，报表弹性主要体现为价格与销量的方向性变化。

\begin{center}\small\setlength{\tabcolsep}{3.4pt}
\captionof{table}[大北农（002385）关键财务指标]{大北农（002385）关键财务指标（合并报表口径；两个半年报期均未年化；现金短债比 $=$ 货币资金 $\div$ 短期债务）}
\begin{adjustbox}{max width=\textwidth}
\begin{tabular}{@{}lrrrrr@{}}
\toprule
指标 & 2023 年 & 2024 年 & 2025 年 & 2025 年半年报 & 2026 年半年报 \\
\midrule
营业收入（亿元） & 333.90 & 287.67 & 291.19 & 135.59 & 138.11 \\
归母净利润（亿元） & -21.74 & 3.46 & -6.42 & 2.35 & -6.77 \\
毛利率（\%） & 9.6 & 15.1 & 12.7 & 14.1 & 8.7 \\
ROE（\%） & -23.0 & 4.2 & -7.8 & 2.7 & -9.3 \\
资产负债率（\%） & 66.3 & 63.5 & 67.0 & 64.0 & 69.9 \\
经营活动现金流净额（亿元） & 2.45 & 24.83 & 8.29 & 0.24 & -12.25 \\
货币资金（亿元） & 46.41 & 38.24 & 42.77 & 45.62 & 32.81 \\
有息负债合计（亿元） & 131.20 & 115.30 & 133.70 & 130.82 & 135.45 \\
现金短债比（倍） & 0.46 & 0.42 & 0.41 & 0.45 & 0.31 \\
\bottomrule
\end{tabular}
\end{adjustbox}
\end{center}

\begin{center}
\begin{minipage}{\textwidth}\centering
\begin{tikzpicture}
\begin{groupplot}[
  group style={group size=2 by 1, horizontal sep=0.60cm},
  width=6.85cm, height=4.60cm,
  tick label style={font=\tiny},
  title style={font=\scriptsize\heiti},
  yticklabel style={font=\tiny},
  ylabel style={font=\tiny},
  grid=major, grid style={LightGray, line width=0.3pt},
  axis line style={InkGray, line width=0.5pt},
  enlarge x limits={abs=0.8},
  legend style={font=\scriptsize, at={(0.5,-0.20)}, anchor=north,
                legend columns=2, draw=none, fill=none, column sep=6pt},
]
\nextgroupplot[
  ybar, bar width=4pt,
  ymin=-57.30, ymax=376.58,
  xtick={0,1,2,3,4,5.7},
  xticklabels={2021,2022,2023,2024,2025,2026H1},
  ylabel={亿元},
]
\addplot[fill=AgriGreen!75, draw=AgriGreen, bar shift=-2.6pt] table[x=i, y=rev]{data/clean/profiles/fin_002385.dat};
\addplot[fill=SoftRed!60, draw=SoftRed, bar shift=2.6pt] table[x=i, y=ni]{data/clean/profiles/fin_002385.dat};
\addplot[fill=AgriGreen!30, draw=AgriGreen, pattern=north east lines, bar shift=-2.6pt] table[x=x, y=rev]{data/clean/profiles/finh_002385.dat};
\addplot[fill=SoftRed!20, draw=SoftRed, pattern=north east lines, bar shift=2.6pt] table[x=x, y=ni]{data/clean/profiles/finh_002385.dat};
\legend{营业收入（年/半年）, 归母净利润（年/半年）}
\nextgroupplot[
  ymin=-30.4, ymax=79.2,
  xtick={0,1,2,3,4,5.7},
  xticklabels={2021,2022,2023,2024,2025,2026H1},
  ylabel={\%},
]
\addplot[AgriGold, mark=*, mark size=1.3pt] table[x=i, y=roe]{data/clean/profiles/fin_002385.dat};
\addplot[OilBlue, mark=square*, mark size=1.3pt] table[x=i, y=debt]{data/clean/profiles/fin_002385.dat};
\addplot[only marks, AgriGold, mark=*, mark size=2.0pt] table[x=x, y=roe]{data/clean/profiles/finh_002385.dat};
\addplot[only marks, OilBlue, mark=square*, mark size=2.0pt] table[x=x, y=debt]{data/clean/profiles/finh_002385.dat};
\legend{ROE, 资产负债率}
\end{groupplot}
\end{tikzpicture}
\begin{tikzpicture}
\begin{axis}[
  name=finbot,
  width=13.75cm, height=3.10cm, scale only axis,
  ymin=-13.55, ymax=151.71,
  xtick={0,1,2,3,4,5.7},
  xticklabels={2021,2022,2023,2024,2025,2026H1},
  tick label style={font=\tiny},
  yticklabel style={font=\tiny},
  ylabel={亿元}, ylabel style={font=\tiny},
  ybar, bar width=4pt,
  grid=major, grid style={LightGray, line width=0.3pt},
  axis line style={InkGray, line width=0.5pt},
  enlarge x limits={abs=0.8},
  legend style={font=\scriptsize, at={(0.5,-0.26)}, anchor=north,
                legend columns=2, draw=none, fill=none, column sep=10pt},
]
\addplot[fill=AgriGreen!55, draw=AgriGreen, bar shift=-3.0pt] table[x=i, y=cash]{data/clean/profiles/fin_002385.dat};
\addplot[fill=SoftRed!45, draw=SoftRed, bar shift=3.0pt] table[x=i, y=ibd]{data/clean/profiles/fin_002385.dat};
\addplot[fill=AgriGreen!25, draw=AgriGreen, pattern=north east lines, bar shift=-3.0pt] table[x=x, y=cash]{data/clean/profiles/finh_002385.dat};
\addplot[fill=SoftRed!20, draw=SoftRed, pattern=north east lines, bar shift=3.0pt] table[x=x, y=ibd]{data/clean/profiles/finh_002385.dat};
\legend{货币资金（左轴）, 有息负债（左轴）}
\end{axis}
\begin{axis}[
  at={(finbot.south east)}, anchor=south east,
  width=13.75cm, height=3.10cm, scale only axis,
  axis y line*=right, axis x line*=none, xtick=\empty,
  ymin=-21.89, ymax=35.95,
  tick label style={font=\tiny},
  yticklabel style={font=\tiny}, scaled y ticks=false,
  legend style={font=\scriptsize, at={(0.5,-0.44)}, anchor=north,
                legend columns=1, draw=none, fill=none},
]
\addplot[OilBlue, mark=*, mark size=2.2pt, line width=1.3pt] table[x=i, y=ocf]{data/clean/profiles/fin_002385.dat};
\addplot[only marks, OilBlue, mark=*, mark size=3.2pt] table[x=x, y=ocf]{data/clean/profiles/finh_002385.dat};
\legend{经营活动现金流净额（右轴，亿元，蓝点）}
\end{axis}
\end{tikzpicture}

\captionof{figure}[大北农（002385）近年主要财务指标]{大北农（002385）近年主要财务指标（左上：营业收入与归母净利润；右上：ROE 与资产负债率；下：货币资金、有息负债为左轴柱状，经营活动现金流净额为其右轴的蓝点折线。
2021---2025 年为年报口径，2026H1 为 2026 年半年报/半年末，与其前一年报之间留出间隔、以斜纹或空心点区分；半年报数据未年化）}
\end{minipage}
\end{center}

\pfl{财务分析} 2026 年 6 月末资产负债率 69.9\%，较 2025 年末上升 2.9 个百分点（样本中位数 52.2\%，所处三级行业内按由低到高排第\zhnumber{8}位）。 2026 年 6 月末货币资金 32.81 亿元、短期债务（短期借款与一年内到期非流动负债）104.33 亿元、长期债务 31.12 亿元，有息负债合计 135.45 亿元，现金短债比 0.31 倍——覆盖偏紧。 净有息负债 102.64 亿元，相当于其 2026 年上半年营业收入的 74\%。 流动比率 0.73、速动比率 0.43，短期偿债指标偏紧。 2026 年上半年经营活动现金流净额 -12.25 亿元，经营现金流与利润同时为负，日常运营对债务与股东投入的依赖度较高。 2026 年 6 月末在建工程 0.84 亿元，较上年末减少 0.05 亿元，在建投入在收缩。 2026 年 6 月末商誉 6.40 亿元，占归母净资产的 7.5\%，是净资产中最脆弱的一块。 综合看，账面货币资金对有息负债与短期债务的覆盖不足，滚动偿付对新增融资的依赖度较高；资产负债率高于样本中位数 17.7 个百分点；有息负债中短期债务占比更高，期限结构仍以滚动续作为主。 公司披露的股东与质押风险为：2026 年 2 月，其名下约 9.28 亿股中有近 5 亿股处于质押状态；2026 年 9 月，控股股东莫云将所持股份质押给质权人以置换并解除邵根伙名下同等数量质押股份。

\pfl{重大战略转型} 公司的战略动作可归为以下几类：其他动作（3 项）：2026 年 7 月，公司参股公司增资暨关联交易。 公告口径上，近三年（2023 年 9 月---2026 年 9 月）公司披露重大重组与并购类公告 1 条、对外投资与转型类公告 6 条、回购与股权激励类公告 39 条，合计公告 607 条。 第一大产品仍是“饲料产品”，收入占比 67.7\%，与 2021 年年报相比没有方向性变化。 综合判断，报告期内公司有 7 条与战略相关的公告落地（重大重组与并购 1 条、对外投资与转型 6 条），资本运作与主业调整之间存在直接关联。

\begin{center}\small\setlength{\tabcolsep}{4pt}
\captionof{table}[大北农（002385）历史融资与资本运作]{大北农（002385）历史融资与资本运作（据公司公告与公开报道整理；金额为披露值，含首次公开发行、再融资、债券、银行借款与重整投资等）}
\begin{adjustbox}{max width=\textwidth}
\begin{tabular}{@{}p{1.45cm}p{2.55cm}p{4.05cm}p{4.95cm}@{}}
\toprule
时间 & 方式 & 规模 & 用途与说明 \\
\midrule
2010-03 & IPO（深交所） & 6,\allowbreak{}080 万股、\allowbreak{}35.00 元/\allowbreak{}股，\allowbreak{}募集资金净额 20.18 亿元（20.18 亿元） & 新型高效预混料等项目；2010-04-09 上市；截至 2023 年首发后累计派现金额亦在该文件中列示 \\
2014-06 & 公司债券 & 5.00 亿元，\allowbreak{}“14 北农债”（112203），\allowbreak{}票面利率 7.24\% & 2014-06-13 上市，\allowbreak{}发行完成后累计公司债券余额不超过 5.00 亿元 \\
2023—2024 & 向特定对象发行股票（定向增发） & 2.11 亿股、\allowbreak{}3.31 元/\allowbreak{}股，\allowbreak{}募集资金总额 7.00 亿元（7.00 亿元） & 饲料项目建设、\allowbreak{}补充流动资金等（原预案拟募集不超过 19.43 亿元 \\
2025-01 & 募投项目变更（募集资金内部调整） & 终止“大北农辽宁区核心科技园建设项目”，\allowbreak{}节余募集资金 1.51 亿元（1.51 亿元） & 全部改投饲料项目、\allowbreak{}技改项目和信息化建设项目；2025 年 1 月 24 日董事会、\allowbreak{}2025 年 2 月 10 日股东大会通过 \\
\bottomrule
\end{tabular}
\end{adjustbox}
\end{center}

\pfl{历史融投资情况} 从现金流量表看，2025 年公司取得借款收到的现金 101.26 亿元、偿还债务支付的现金 86.72 亿元、吸收权益性投资 0.12 亿元，筹资活动现金流净额 4.09 亿元（2023 年为取得借款 85.89 亿元、偿还债务 68.23 亿元）。 2026 年上半年筹资活动现金流净额为 -0.25 亿元（取得借款 50.29 亿元、偿还债务 49.38 亿元，未年化），已转为净流出，处于净还债状态。 2025 年分配股利、利润或偿付利息支付的现金合计 6.03 亿元。 公开披露的融资记录共 4 笔，工具类型以 IPO、公司债、定向增发为主；金额与用途见下表。 其中规模最大的两笔为：2025 年 1 月，募投项目变更，终止“大北农辽宁区核心科技园建设项目”，节余募集资金 1.51 亿元（1.51 亿元）（该事项后续已终止或解除，不构成已到位资金）；2023---2024 年，定向增发，2.11 亿股、3.31 元/股，募集资金总额 7.00 亿元（7.00 亿元）。 股本方面，近三年披露股本变动 8 次（股份回购 3 次、增发新股上市 1 次、其他 1 次）；总股本由 2021 年末的 41.41 亿股变为 43.24 亿股（+4.41\%）。 分红方面，近三年现金分红 4 次、合计每股派现 0.19 元，最近一次实施在 2025 年 9 月。 近三年“再融资”类公告 16 条，涉及定向增发、募集资金使用。 综合看，公司融资结构上以债务滚动为主、股权融资为补充；2025 年筹资活动现金流净额 4.09 亿元，年度内仍为净融入，与前述经营与投资现金流的方向基本一致。

\pfl{未来潜在融资需求分析} 2026 年 6 月末货币资金 32.81 亿元，而短期债务 104.33 亿元（现金短债比 0.31 倍），2026 年上半年经营现金流净额 -12.25 亿元。按“短期债务减去账面货币资金”的滚动偿付口径，静态缺口约 71.52 亿元；若再计入上半年亏损的年化额（13.53 亿元），维持一年正常经营所需的外部资金约 71.52---85.05 亿元。 按第 \ref{sec:financing-need} 节的三档划分，公司属于第一档“补血型”：2026 年上半年归母净利润 -6.77 亿元、6 月末资产负债率 69.9\%。 可行的融资路径按现实性排序为：债务重组或展期（与银行和债券持有人协商）、定向增发（需股价配合，且控股股东需放弃优先认购或同步增资）、出售非核心资产与参股股权。 公开披露的后续资金安排线索包括：股权融资线索：2022 年度，公司向特定对象发行股票原拟募集不超过 19.43 亿元；债务与授信安排：2026 年 9 月，发布公司及控股子公司担保进展公告，为子公司经营融资提供担保；资本开支安排：2025 年 10 月，董事会通过部分募集资金投资项目延期。 资金需求还要叠加在建工程：期末在建工程 0.84 亿元、2026 年上半年购建固定资产等支出 2.25 亿元（未年化），这部分支出在周期底部通常会被主动放缓。 综合看，公司的外部资金需求以债务接续为主，滚动偿付口径下的静态缺口约 71.52 亿元，融资的时点与规模更多取决于信用条件与资本市场窗口。




\subsubsection{天康生物（002100）}
\label{co:002100}

\pfl{公司基本情况} 公司前身可追溯至 2006 年，新疆兵团国资背景的农牧企业，主营饲料、生猪养殖与食品、动物疫苗；2026 年 6 月收购羌都畜牧 51\% 股权扩产生猪，2026 年上半年亏损 4.41 亿元。 公司于 2006-12-26 在深交所主板首发上市，注册资本 13.65 亿元，总股本 13.65 亿股。 截至 2026 年 10 月 12 日收盘价 7.72 元，对应总市值 105.4 亿元，近一年+0.0\%，总市值在三级行业“畜禽饲料”的 11 家公司中按由大到小排第\zhnumber{3}位，近一年涨跌幅高于全样本中位数（-10.4\%）10.4 个百分点。 按 2026 年半年报披露的行业构成，收入占比前三位为生猪养殖产业链行业 32.4\%；饲料行业 27.8\%；蛋白油脂加工行业 16.7\%。 与 2021 年年报相比，第一大行业口径由“饲料行业”（32.6\%）变为“生猪养殖产业链行业”（32.4\%）；两期披露的分类口径有调整，占比差异中同时包含分类因素。

\pfl{股权结构与控制权} 按 2026 年 6 月 30 日的股东名册，前三大股东为新疆生产建设兵团国有资产经营有限责任公司 持股 23.36\%；北京融元宝通资产管理合伙企业(有限合伙) 持股 5.04\%；海南天邦鸿康创业投资有限公司 持股 4.03\%，股东户数 53,044 户。 股东结构与控制权方面，持股 23.36\%（3.189 亿股）；控股股东所持股份无质押、标记或冻结；正在推进公开发行股票并在北京证券交易所上市。

\pfl{经营要点} 经营层面，2026 年上半年，生猪出栏 184.84 万头、同比增长 20.95\%（含 6 月并表的新疆羌都畜牧）；2026 年上半年，饲料销量 139.97 万吨、同比增长 4.33\%。

\pfl{所处行业的基本面} 饲料的供给即产能（门槛低、产能利用率偏低），需求直接来自下游存栏，利润则由玉米、豆粕等原料价格与成品售价的价差决定（第 \ref{sec:framework} 章的供需框架）。 与该公司直接相关的行业数据包括：2026 年上半年，全国生猪出栏 37,246 万头、同比增长 1.7\%，二季度末全国能繁母猪存栏 3,780 万头、同比下降 6.5\%。

\pfl{竞争优势} 把该公司放回同类公司比较，2026 年上半年营业收入 83.3 亿元、在三级行业“畜禽饲料”的 11 家公司中排第\zhnumber{5}位；毛利率 1.4\%、低于行业内中位数 3.8 个百分点；半年 ROE -6.6\%（未年化）、低于行业内中位数（-6.25\%）；资产负债率 55.6\%、低于行业内中位数（62.1\%）。 公司在公开资料中自述与同业比较的相对优势在于：公司通过收购新疆羌都畜牧 51\% 股权，在若羌、尉犁等片区形成“饲草种植—饲料加工—生猪繁育育肥—屠宰”一体化闭环。 就当期报表看，公司尚未把经营层面的特征转化为高于行业的资本回报，周期底部的相对优势更多体现在现金流与资产负债表上。

\pfl{历史股价} 以 2021 年 1 月为起点、2026 年 10 月为终点，股价由 10.76 元变为 7.72 元，累计 -28.3\%，区间最高 12.41 元（2021 年 2 月）、最低 6.14 元（2025 年 2 月）；当前价相当于区间最高价的 62.2\%，为区间最低价的 1.26 倍。 月度收益的年化波动率 33.5\%，低于全样本中位数 37.3\%，在三级行业“畜禽饲料”的 11 家公司中波动率由高到低排第\zhnumber{7}位。

\begin{center}
\begin{minipage}{\textwidth}\centering
\begin{tikzpicture}
\begin{axis}[
  width=12.6cm, height=3.9cm, scale only axis,
  xmin=2020.922, xmax=2026.972, ymin=5.71, ymax=13.28,
  xtick={2021,2022,2023,2024,2025,2026},
  yearticks,
  yticklabel style={font=\scriptsize},
  ylabel={元/股}, ylabel style={font=\scriptsize},
  grid=major, grid style={LightGray, line width=0.3pt},
  axis line style={InkGray, line width=0.5pt},
  every axis plot/.append style={line width=0.9pt},
]
\addplot[AgriGreen] table[x=x,y=y]{data/clean/profiles/px_002100.dat};
\addplot[SoftRed, dashed, line width=0.7pt] table[x=x,y=y]{data/clean/profiles/pxzz_002100.dat};
\addplot[only marks, mark=*, mark size=1.1pt, SoftRed] table[x=x,y=y]{data/clean/profiles/pxzz_002100.dat};
\end{axis}
\end{tikzpicture}

\captionof{figure}[天康生物（002100）股价与周期骨架]{天康生物（002100）月度收盘价（元/股）与 ZigZag 周期骨架（阈值 25\%，圆点为周期转折点）}
\end{minipage}
\end{center}

\pfl{过去周期分析} 以 25\% 的 ZigZag 阈值划分，样本区间内股价形成 3 个主升段与 3 个主跌段。 最大主跌段为 2022 年 7 月 至 2025 年 2 月，31 个月-49.4\%；最大主升段出现在 2021 年 7 月 至 2022 年 3 月，8 个月+66.9\%。 最近一次转折出现在 2026 年 4 月（上涨段自 2025 年 2 月起，14 个月+29.5\%），当前处于该段之后的震荡之中。 公司股价较申万农林牧渔指数提前约 15 个月见底，是板块中较早定价周期反转的公司之一。 饲料夹在原料成本与养殖需求之间（第 \ref{sec:feed} 节），量的部分取决于存栏、利的部分取决于原料与成品的价差，公司 2026 年 6 月末资产负债率 55.6\%、半年 ROE -6.6\%（未年化），在本轮下行中属于随行业同步调整的一类。 综合区间形态看，该股单段涨跌幅度偏大、以趋势段为主（最大主升 +66.9\%、最大主跌 -49.4\%），年化波动率 33.5\% 与全样本中位数 37.3\% 相比波动小于样本中位数。

\pfl{盈利情况} 2026 年上半年营业收入 83.25 亿元（同比 -5.9\%），归母净利润 -4.41 亿元，由上年同期的盈利转为亏损。 以年报为趋势对照，2023---2025 年营业收入由 190.26 亿元变为 169.81 亿元（两年年均复合增速 -5.5\%），归母净利润由 -13.63 亿元变为 2.11 亿元。 按半年报口径，2026 年上半年的 ROE 为 -6.6\%（未年化）、销售净利率 -5.9\%、毛利率 1.4\%，ROE 在“畜禽饲料”11 家公司中按数值由高到低排第\zhnumber{7}位。 同期全样本半年 ROE 中位数 0.75\%，公司与之相差 7.32 个百分点（弱于中位数公司）。 从公司披露的原因看，业绩变动主要来自：2026 年上半年，兽药（动物疫苗）收入同比下降 43.95\%，为除猪价外的主要拖累项。 针对上述压力，公司披露的应对举措包括：2026 年 6 月，公司表示本次收购整合以现有存量产能优化配置为核心。 综合看，公司仍处亏损区间、由盈转亏，半年 ROE 在三级行业内按由高到低排第\zhnumber{7}位，单位盈利能力的修复依赖行业景气与自身成本端的改善，报表弹性主要体现为价格与销量的方向性变化。

\begin{center}\small\setlength{\tabcolsep}{3.4pt}
\captionof{table}[天康生物（002100）关键财务指标]{天康生物（002100）关键财务指标（合并报表口径；两个半年报期均未年化；现金短债比 $=$ 货币资金 $\div$ 短期债务）}
\begin{adjustbox}{max width=\textwidth}
\begin{tabular}{@{}lrrrrr@{}}
\toprule
指标 & 2023 年 & 2024 年 & 2025 年 & 2025 年半年报 & 2026 年半年报 \\
\midrule
营业收入（亿元） & 190.26 & 171.76 & 169.81 & 88.47 & 83.25 \\
归母净利润（亿元） & -13.63 & 6.05 & 2.11 & 3.38 & -4.41 \\
毛利率（\%） & 2.9 & 12.1 & 9.9 & 12.1 & 1.4 \\
ROE（\%） & -19.2 & 9.0 & 3.0 & 4.7 & -6.6 \\
资产负债率（\%） & 52.8 & 51.5 & 50.9 & 47.0 & 55.6 \\
经营活动现金流净额（亿元） & 17.52 & 11.62 & 3.81 & 16.19 & 14.96 \\
货币资金（亿元） & 28.22 & 31.02 & 20.81 & 29.73 & 25.35 \\
有息负债合计（亿元） & 63.96 & 64.79 & 63.68 & 52.26 & 65.33 \\
现金短债比（倍） & 0.59 & 0.76 & 0.50 & 0.90 & 0.83 \\
\bottomrule
\end{tabular}
\end{adjustbox}
\end{center}

\begin{center}
\begin{minipage}{\textwidth}\centering
\begin{tikzpicture}
\begin{groupplot}[
  group style={group size=2 by 1, horizontal sep=0.60cm},
  width=6.85cm, height=4.60cm,
  tick label style={font=\tiny},
  title style={font=\scriptsize\heiti},
  yticklabel style={font=\tiny},
  ylabel style={font=\tiny},
  grid=major, grid style={LightGray, line width=0.3pt},
  axis line style={InkGray, line width=0.5pt},
  enlarge x limits={abs=0.8},
  legend style={font=\scriptsize, at={(0.5,-0.20)}, anchor=north,
                legend columns=2, draw=none, fill=none, column sep=6pt},
]
\nextgroupplot[
  ybar, bar width=4pt,
  ymin=-34.02, ymax=214.73,
  xtick={0,1,2,3,4,5.7},
  xticklabels={2021,2022,2023,2024,2025,2026H1},
  ylabel={亿元},
]
\addplot[fill=AgriGreen!75, draw=AgriGreen, bar shift=-2.6pt] table[x=i, y=rev]{data/clean/profiles/fin_002100.dat};
\addplot[fill=SoftRed!60, draw=SoftRed, bar shift=2.6pt] table[x=i, y=ni]{data/clean/profiles/fin_002100.dat};
\addplot[fill=AgriGreen!30, draw=AgriGreen, pattern=north east lines, bar shift=-2.6pt] table[x=x, y=rev]{data/clean/profiles/finh_002100.dat};
\addplot[fill=SoftRed!20, draw=SoftRed, pattern=north east lines, bar shift=2.6pt] table[x=x, y=ni]{data/clean/profiles/finh_002100.dat};
\legend{营业收入（年/半年）, 归母净利润（年/半年）}
\nextgroupplot[
  ymin=-25.2, ymax=63.0,
  xtick={0,1,2,3,4,5.7},
  xticklabels={2021,2022,2023,2024,2025,2026H1},
  ylabel={\%},
]
\addplot[AgriGold, mark=*, mark size=1.3pt] table[x=i, y=roe]{data/clean/profiles/fin_002100.dat};
\addplot[OilBlue, mark=square*, mark size=1.3pt] table[x=i, y=debt]{data/clean/profiles/fin_002100.dat};
\addplot[only marks, AgriGold, mark=*, mark size=2.0pt] table[x=x, y=roe]{data/clean/profiles/finh_002100.dat};
\addplot[only marks, OilBlue, mark=square*, mark size=2.0pt] table[x=x, y=debt]{data/clean/profiles/finh_002100.dat};
\legend{ROE, 资产负债率}
\end{groupplot}
\end{tikzpicture}
\begin{tikzpicture}
\begin{axis}[
  name=finbot,
  width=13.75cm, height=3.10cm, scale only axis,
  ymin=-7.38, ymax=82.62,
  xtick={0,1,2,3,4,5.7},
  xticklabels={2021,2022,2023,2024,2025,2026H1},
  tick label style={font=\tiny},
  yticklabel style={font=\tiny},
  ylabel={亿元}, ylabel style={font=\tiny},
  ybar, bar width=4pt,
  grid=major, grid style={LightGray, line width=0.3pt},
  axis line style={InkGray, line width=0.5pt},
  enlarge x limits={abs=0.8},
  legend style={font=\scriptsize, at={(0.5,-0.26)}, anchor=north,
                legend columns=2, draw=none, fill=none, column sep=10pt},
]
\addplot[fill=AgriGreen!55, draw=AgriGreen, bar shift=-3.0pt] table[x=i, y=cash]{data/clean/profiles/fin_002100.dat};
\addplot[fill=SoftRed!45, draw=SoftRed, bar shift=3.0pt] table[x=i, y=ibd]{data/clean/profiles/fin_002100.dat};
\addplot[fill=AgriGreen!25, draw=AgriGreen, pattern=north east lines, bar shift=-3.0pt] table[x=x, y=cash]{data/clean/profiles/finh_002100.dat};
\addplot[fill=SoftRed!20, draw=SoftRed, pattern=north east lines, bar shift=3.0pt] table[x=x, y=ibd]{data/clean/profiles/finh_002100.dat};
\legend{货币资金（左轴）, 有息负债（左轴）}
\end{axis}
\begin{axis}[
  at={(finbot.south east)}, anchor=south east,
  width=13.75cm, height=3.10cm, scale only axis,
  axis y line*=right, axis x line*=none, xtick=\empty,
  ymin=-4.55, ymax=22.77,
  tick label style={font=\tiny},
  yticklabel style={font=\tiny}, scaled y ticks=false,
  legend style={font=\scriptsize, at={(0.5,-0.44)}, anchor=north,
                legend columns=1, draw=none, fill=none},
]
\addplot[OilBlue, mark=*, mark size=2.2pt, line width=1.3pt] table[x=i, y=ocf]{data/clean/profiles/fin_002100.dat};
\addplot[only marks, OilBlue, mark=*, mark size=3.2pt] table[x=x, y=ocf]{data/clean/profiles/finh_002100.dat};
\legend{经营活动现金流净额（右轴，亿元，蓝点）}
\end{axis}
\end{tikzpicture}

\captionof{figure}[天康生物（002100）近年主要财务指标]{天康生物（002100）近年主要财务指标（左上：营业收入与归母净利润；右上：ROE 与资产负债率；下：货币资金、有息负债为左轴柱状，经营活动现金流净额为其右轴的蓝点折线。
2021---2025 年为年报口径，2026H1 为 2026 年半年报/半年末，与其前一年报之间留出间隔、以斜纹或空心点区分；半年报数据未年化）}
\end{minipage}
\end{center}

\pfl{财务分析} 2026 年 6 月末资产负债率 55.6\%，较 2025 年末上升 4.6 个百分点（样本中位数 52.2\%，所处三级行业内按由低到高排第\zhnumber{3}位）。 2026 年 6 月末货币资金 25.35 亿元、短期债务（短期借款与一年内到期非流动负债）30.51 亿元、长期债务 34.83 亿元，有息负债合计 65.33 亿元，现金短债比 0.83 倍——覆盖偏紧。 净有息负债 39.99 亿元，相当于其 2026 年上半年营业收入的 48\%。 流动比率 1.31、速动比率 0.68，短期指标尚可。 2026 年上半年经营活动现金流净额 14.96 亿元，净利润为负而经营现金流为正，差额主要来自折旧摊销与营运资本变动，说明亏损尚未完全转化为现金流出。 2026 年 6 月末在建工程 5.64 亿元，较上年末增加 3.09 亿元，仍有在建投入。 2026 年 6 月末商誉 7.20 亿元，占归母净资产的 8.4\%，是净资产中最脆弱的一块。 综合看，现金短债比 0.83 倍，短期资金安排偏紧；资产负债率高于样本中位数 3.4 个百分点；有息负债中长期债务占比更高，期限结构对短债滚动的压力较小。

\pfl{重大战略转型} 公司的战略动作可归为以下几类：并购与重组（3 项）：2026 年 1 月，国家市场监督管理总局公示天康生物收购新疆羌都畜牧科技有限公司股权案（经营者集中简易案件）；2025 年 12 月，与新疆七星羌都集团农牧有限公司等股东签署股权收购协议，以现金 12.75 亿元收购羌都畜牧 51\% 股权；融资与资本运作（1 项）：2025 年 6 月，子公司天康制药推进分拆上市：招股说明书显示天康生物直接持股 58.9388\%。 公告口径上，近三年（2023 年 9 月---2026 年 9 月）公司披露重大重组与并购类公告 5 条、对外投资与转型类公告 1 条、回购与股权激励类公告 0 条，合计公告 349 条。 主营结构的口径变化已经落到报表上：2021 年年报的第一大行业为“饲料行业”（32.6\%），2026 年半年报为“生猪养殖产业链行业”（32.4\%）；公司在这两期的分部划分方式有过调整。 综合判断，报告期内公司有 6 条与战略相关的公告落地（重大重组与并购 5 条、对外投资与转型 1 条），资本运作与主业调整之间存在直接关联。

\begin{center}\small\setlength{\tabcolsep}{4pt}
\captionof{table}[天康生物（002100）历史融资与资本运作]{天康生物（002100）历史融资与资本运作（据公司公告与公开报道整理；金额为披露值，含首次公开发行、再融资、债券、银行借款与重整投资等）}
\begin{adjustbox}{max width=\textwidth}
\begin{tabular}{@{}p{1.45cm}p{2.55cm}p{4.05cm}p{4.95cm}@{}}
\toprule
时间 & 方式 & 规模 & 用途与说明 \\
\midrule
2006-12 & IPO（深交所） & 1,\allowbreak{}600 万股、\allowbreak{}10.86 元/\allowbreak{}股 & 牛羊专用饲料生产线、\allowbreak{}专业浮性水产饲料厂、\allowbreak{}酶制剂和益绿素纯植物饲料添加剂等项目；上市日 2006-12-26，\allowbreak{}发行后总股本 6,\allowbreak{}400 万股 \\
2015-07 & 发行股份吸收合并（关联交易） & 发行 3.88 亿股，\allowbreak{}同时注销天康控股所持 2.93 亿股 & 吸收合并新疆天康（控股）集团 \\
2017-12 & 公开发行可转换公司债券 & 10.00 亿元（1,\allowbreak{}000 万张，\allowbreak{}面值 100.00 元，\allowbreak{}期限 6 年） & 联合信用评级给予主体与债项 AA；扣除承销费及保荐费 1,\allowbreak{}450.00 万元后实际收到认购资金 9.86 亿元 \\
2021-12 & 非公开发行股票（2020 年度方案） & 2.77 亿股、\allowbreak{}7.45 元/\allowbreak{}股，\allowbreak{}募集资金总额 20.67 亿元，\allowbreak{}净额 20.43 亿元 & 生猪养殖及饲料等项目；2021-12-24 新增股份上市；2025 年报显示该次募集资金累计已大量使用 \\
2026 & 综合授信额度申请 & 公司及下属公司 2026 年度拟向银行及其他类金融机构申请总额不超过 120.00 亿元综合授信 & 建设项目贷款、\allowbreak{}流动资金贷款及原材料采购等用于降低原材料价格波动风险、\allowbreak{}提高各业务板块盈利能力 \\
\bottomrule
\end{tabular}
\end{adjustbox}
\end{center}

\pfl{历史融投资情况} 从现金流量表看，2025 年公司取得借款收到的现金 53.64 亿元、偿还债务支付的现金 54.08 亿元、吸收权益性投资 0.00 亿元，筹资活动现金流净额 -6.36 亿元（2023 年为取得借款 49.15 亿元、偿还债务 58.82 亿元）。 2026 年上半年筹资活动现金流净额为 -3.13 亿元（取得借款 26.99 亿元、偿还债务 28.40 亿元，未年化），已转为净流出，处于净还债状态。 2025 年分配股利、利润或偿付利息支付的现金合计 4.43 亿元。 公开披露的融资记录共 5 笔，工具类型以 IPO、发行股份吸收合并、可转债、定向增发等为主；金额与用途见下表。 其中规模最大的两笔为：2026 年，银行授信，公司及下属公司 2026 年度拟向银行及其他类金融机构申请总额不超过 120.00 亿元综合授信；2021 年 12 月，定向增发，2.77 亿股、7.45 元/股，募集资金总额 20.67 亿元，净额 20.43 亿元。 股本方面，近三年披露股本变动 7 次（可转债转股 5 次）；总股本由 2021 年末的 13.54 亿股变为 13.65 亿股（+0.84\%）。 分红方面，近三年现金分红 3 次、合计每股派现 0.52 元，最近一次实施在 2025 年 12 月。 近三年“再融资”类公告 2 条，涉及可转债。 综合看，公司融资结构上以债务滚动为主、股权融资为补充；2025 年筹资活动现金流净额 -6.36 亿元，融资节奏仍以维持存量债务周转为主，与前述经营与投资现金流的方向基本一致。

\pfl{未来潜在融资需求分析} 2026 年 6 月末货币资金 25.35 亿元，而短期债务 30.51 亿元（现金短债比 0.83 倍），2026 年上半年经营现金流净额 14.96 亿元。按“短期债务减去账面货币资金”的滚动偿付口径，静态缺口约 5.16 亿元；若再计入上半年亏损的年化额（8.83 亿元），维持一年正常经营所需的外部资金约 5.16---13.99 亿元。 公司 2026 年上半年归母净利润为负（-4.41 亿元），但 6 月末资产负债率 55.6\% 尚未进入第一档（亏损且负债率 $>$65\%）的区间，当前融资需求以补充流动资金、支撑低谷期经营性支出为主。 公开披露的后续资金安排线索包括：股权融资线索：2025---2026 年，天康制药分拆上市持续推进，为子公司股权融资线索；债务与授信安排：2026 年 4 月，董事会/年报披露 2026 年度资金需求计划：拟向银行及其他类金融机构申请总额不超过 120.00 亿元综合授信。 资金需求还要叠加在建工程：期末在建工程 5.64 亿元、2026 年上半年购建固定资产等支出 1.52 亿元（未年化），这部分支出在周期底部通常会被主动放缓。 综合看，公司的外部资金需求以补充流动性与项目投入为主，滚动偿付口径下的静态缺口约 5.16 亿元，融资的时点与规模更多取决于信用条件与资本市场窗口。




\subsubsection{禾丰股份（603609）}
\label{co:603609}

\pfl{公司基本情况} 2014 年设立至今，禾丰股份是以饲料为核心、肉禽（白羽肉鸡屠宰）与生猪业务协同的农牧食品企业，荷兰 De Heus 持股约 9\%；2026 年上半年饲料外销创新高，受猪价拖累归母净利微亏。 公司于 2014-08-08 在上交所首发上市，注册资本 9.04 亿元，总股本 9.19 亿股。 截至 2026 年 10 月 12 日收盘价 6.34 元，对应总市值 58.3 亿元，近一年-21.4\%，总市值在三级行业“畜禽饲料”的 11 家公司中按由大到小排第\zhnumber{4}位，近一年涨跌幅低于全样本中位数（-10.4\%）11.1 个百分点。 按 2025 年年报披露的行业构成，收入占比前三位为肉禽 45.3\%；饲料 37.6\%；生猪 9.9\%。 与 2021 年年报相比，第一大行业口径由“饲料业务”（52.1\%）变为“肉禽”（45.3\%）；两期披露的分类口径有调整，占比差异中同时包含分类因素。

\pfl{股权结构与控制权} 按 2026 年 6 月 30 日的股东名册，前三大股东为金卫东 持股 16.46\%；DE HEUS MAURITIUS 持股 9.06\%；丁云峰 持股 8.84\%，股东户数 22,797 户。 从控制权角度看，直接持股 16.46\%；实际控制人及一致行动人所持股份均未质押，截至 2026 年 3 月底公司总股本为 9.08 亿股。

\pfl{经营要点} 经营与产能方面，2026 年上半年，预混料销量同比增长 18\%、配合料增长 23\%，散装猪饲料外销量增长 16\%、占猪饲料总销量比重接近 50\%；2026 年上半年，6 月末母猪存栏合计 6.2 万头（能繁 5.3 万头、后备 0.9 万头）。

\pfl{所处行业的基本面} 饲料的供给即产能（门槛低、产能利用率偏低），需求直接来自下游存栏，利润则由玉米、豆粕等原料价格与成品售价的价差决定（第 \ref{sec:framework} 章的供需框架）。 与该公司直接相关的行业数据包括：2026 年上半年，全国工业饲料总产量 16,850 万吨、同比增长 3.9\%；2026 年 10 月，上旬东北 14\%—15\% 水分新玉米收购价 2,020—2,150 元/吨、较节前下跌 10—60 元/吨。

\pfl{竞争优势} 把该公司放回同类公司比较，2026 年上半年营业收入 189.6 亿元、在三级行业“畜禽饲料”的 11 家公司中排第\zhnumber{2}位；毛利率 6.3\%、高于行业内中位数 1.1 个百分点；半年 ROE -0.1\%（未年化）、高于行业内中位数（-6.25\%）；资产负债率 57.7\%、低于行业内中位数（62.1\%）。 与同类公司相比，公司的差异化在于：2026 年上半年，饲料外销量增速 18\%；2026 年上半年经营活动现金流净额同比改善 31.40\%，投资活动现金流出同比下降 68.67\%。 优势在当期体现为高于行业中位数的资本回报，但优势能否持续仍取决于周期与自身产能投放节奏。

\pfl{历史股价} 自 2021 年 1 月至 2026 年 10 月，股价由 11.70 元变为 6.34 元，累计 -45.8\%，区间最高 13.85 元（2022 年 11 月）、最低 5.76 元（2026 年 6 月）；当前价相当于区间最高价的 45.8\%，为区间最低价的 1.10 倍。 月度收益的年化波动率 28.5\%，低于全样本中位数 37.3\%，在三级行业“畜禽饲料”的 11 家公司中波动率由高到低排第\zhnumber{11}位。

\begin{center}
\begin{minipage}{\textwidth}\centering
\begin{tikzpicture}
\begin{axis}[
  width=12.6cm, height=3.9cm, scale only axis,
  xmin=2020.922, xmax=2026.972, ymin=5.36, ymax=14.82,
  xtick={2021,2022,2023,2024,2025,2026},
  yearticks,
  yticklabel style={font=\scriptsize},
  ylabel={元/股}, ylabel style={font=\scriptsize},
  grid=major, grid style={LightGray, line width=0.3pt},
  axis line style={InkGray, line width=0.5pt},
  every axis plot/.append style={line width=0.9pt},
]
\addplot[AgriGreen] table[x=x,y=y]{data/clean/profiles/px_603609.dat};
\addplot[SoftRed, dashed, line width=0.7pt] table[x=x,y=y]{data/clean/profiles/pxzz_603609.dat};
\addplot[only marks, mark=*, mark size=1.1pt, SoftRed] table[x=x,y=y]{data/clean/profiles/pxzz_603609.dat};
\end{axis}
\end{tikzpicture}

\captionof{figure}[禾丰股份（603609）股价与周期骨架]{禾丰股份（603609）月度收盘价（元/股）与 ZigZag 周期骨架（阈值 25\%，圆点为周期转折点）}
\end{minipage}
\end{center}

\pfl{过去周期分析} 以 25\% 的 ZigZag 阈值划分，样本区间内股价形成 2 个主升段与 3 个主跌段。 最大主升段出现在 2022 年 4 月 至 2022 年 11 月，7 个月+74.2\%；最大主跌段为 2022 年 11 月 至 2024 年 6 月，19 个月-52.8\%。 最近一次转折出现在 2026 年 6 月（下跌段自 2025 年 5 月起，13 个月-34.9\%），当前处于该段的延续之中。 公司股价与申万农林牧渔指数几乎同步见底，个股并未走出独立行情。 饲料夹在原料成本与养殖需求之间（第 \ref{sec:feed} 节），量的部分取决于存栏、利的部分取决于原料与成品的价差，公司 2026 年 6 月末资产负债率 57.7\%、半年 ROE -0.1\%（未年化），在本轮下行中属于随行业同步调整的一类。 综合区间形态看，该股单段涨跌幅度偏大、以趋势段为主（最大主升 +74.2\%、最大主跌 -52.8\%），年化波动率 28.5\% 与全样本中位数 37.3\% 相比波动小于样本中位数。

\pfl{盈利情况} 2026 年上半年营业收入 189.63 亿元（同比 +8.9\%），归母净利润 -0.08 亿元，由上年同期的盈利转为亏损。 以年报为趋势对照，2023---2025 年营业收入由 359.70 亿元变为 357.62 亿元（两年年均复合增速 -0.3\%），归母净利润由 -4.57 亿元变为 0.53 亿元。 按半年报口径，2026 年上半年的 ROE 为 -0.1\%（未年化）、销售净利率 0.1\%、毛利率 6.3\%，ROE 在“畜禽饲料”11 家公司中按数值由高到低排第\zhnumber{4}位。 同期全样本半年 ROE 中位数 0.75\%，公司与之相差 0.88 个百分点（弱于中位数公司）。 综合看，公司仍处亏损区间、由盈转亏，半年 ROE 在三级行业内按由高到低排第\zhnumber{4}位，单位盈利能力的修复依赖行业景气与自身成本端的改善，报表弹性主要体现为价格与销量的方向性变化。

\begin{center}\small\setlength{\tabcolsep}{3.4pt}
\captionof{table}[禾丰股份（603609）关键财务指标]{禾丰股份（603609）关键财务指标（合并报表口径；两个半年报期均未年化；现金短债比 $=$ 货币资金 $\div$ 短期债务）}
\begin{adjustbox}{max width=\textwidth}
\begin{tabular}{@{}lrrrrr@{}}
\toprule
指标 & 2023 年 & 2024 年 & 2025 年 & 2025 年半年报 & 2026 年半年报 \\
\midrule
营业收入（亿元） & 359.70 & 325.45 & 357.62 & 174.07 & 189.63 \\
归母净利润（亿元） & -4.57 & 3.42 & 0.53 & 2.33 & -0.08 \\
毛利率（\%） & 4.0 & 6.4 & 5.8 & 6.4 & 6.3 \\
ROE（\%） & -6.6 & 5.1 & 0.8 & 3.4 & -0.1 \\
资产负债率（\%） & 49.0 & 49.1 & 56.3 & 56.4 & 57.7 \\
经营活动现金流净额（亿元） & 9.56 & 11.77 & 6.81 & -5.12 & -3.51 \\
货币资金（亿元） & 19.15 & 17.94 & 17.20 & 13.98 & 12.18 \\
有息负债合计（亿元） & 44.98 & 42.79 & 63.44 & 65.26 & 66.24 \\
现金短债比（倍） & 1.01 & 1.69 & 0.74 & 0.56 & 0.53 \\
\bottomrule
\end{tabular}
\end{adjustbox}
\end{center}

\begin{center}
\begin{minipage}{\textwidth}\centering
\begin{tikzpicture}
\begin{groupplot}[
  group style={group size=2 by 1, horizontal sep=0.60cm},
  width=6.85cm, height=4.60cm,
  tick label style={font=\tiny},
  title style={font=\scriptsize\heiti},
  yticklabel style={font=\tiny},
  ylabel style={font=\tiny},
  grid=major, grid style={LightGray, line width=0.3pt},
  axis line style={InkGray, line width=0.5pt},
  enlarge x limits={abs=0.8},
  legend style={font=\scriptsize, at={(0.5,-0.20)}, anchor=north,
                legend columns=2, draw=none, fill=none, column sep=6pt},
]
\nextgroupplot[
  ybar, bar width=4pt,
  ymin=-41.00, ymax=403.42,
  xtick={0,1,2,3,4,5.7},
  xticklabels={2021,2022,2023,2024,2025,2026H1},
  ylabel={亿元},
]
\addplot[fill=AgriGreen!75, draw=AgriGreen, bar shift=-2.6pt] table[x=i, y=rev]{data/clean/profiles/fin_603609.dat};
\addplot[fill=SoftRed!60, draw=SoftRed, bar shift=2.6pt] table[x=i, y=ni]{data/clean/profiles/fin_603609.dat};
\addplot[fill=AgriGreen!30, draw=AgriGreen, pattern=north east lines, bar shift=-2.6pt] table[x=x, y=rev]{data/clean/profiles/finh_603609.dat};
\addplot[fill=SoftRed!20, draw=SoftRed, pattern=north east lines, bar shift=2.6pt] table[x=x, y=ni]{data/clean/profiles/finh_603609.dat};
\legend{营业收入（年/半年）, 归母净利润（年/半年）}
\nextgroupplot[
  ymin=-11.7, ymax=64.1,
  xtick={0,1,2,3,4,5.7},
  xticklabels={2021,2022,2023,2024,2025,2026H1},
  ylabel={\%},
]
\addplot[AgriGold, mark=*, mark size=1.3pt] table[x=i, y=roe]{data/clean/profiles/fin_603609.dat};
\addplot[OilBlue, mark=square*, mark size=1.3pt] table[x=i, y=debt]{data/clean/profiles/fin_603609.dat};
\addplot[only marks, AgriGold, mark=*, mark size=2.0pt] table[x=x, y=roe]{data/clean/profiles/finh_603609.dat};
\addplot[only marks, OilBlue, mark=square*, mark size=2.0pt] table[x=x, y=debt]{data/clean/profiles/finh_603609.dat};
\legend{ROE, 资产负债率}
\end{groupplot}
\end{tikzpicture}
\begin{tikzpicture}
\begin{axis}[
  name=finbot,
  width=13.75cm, height=3.10cm, scale only axis,
  ymin=-6.62, ymax=74.18,
  xtick={0,1,2,3,4,5.7},
  xticklabels={2021,2022,2023,2024,2025,2026H1},
  tick label style={font=\tiny},
  yticklabel style={font=\tiny},
  ylabel={亿元}, ylabel style={font=\tiny},
  ybar, bar width=4pt,
  grid=major, grid style={LightGray, line width=0.3pt},
  axis line style={InkGray, line width=0.5pt},
  enlarge x limits={abs=0.8},
  legend style={font=\scriptsize, at={(0.5,-0.26)}, anchor=north,
                legend columns=2, draw=none, fill=none, column sep=10pt},
]
\addplot[fill=AgriGreen!55, draw=AgriGreen, bar shift=-3.0pt] table[x=i, y=cash]{data/clean/profiles/fin_603609.dat};
\addplot[fill=SoftRed!45, draw=SoftRed, bar shift=3.0pt] table[x=i, y=ibd]{data/clean/profiles/fin_603609.dat};
\addplot[fill=AgriGreen!25, draw=AgriGreen, pattern=north east lines, bar shift=-3.0pt] table[x=x, y=cash]{data/clean/profiles/finh_603609.dat};
\addplot[fill=SoftRed!20, draw=SoftRed, pattern=north east lines, bar shift=3.0pt] table[x=x, y=ibd]{data/clean/profiles/finh_603609.dat};
\legend{货币资金（左轴）, 有息负债（左轴）}
\end{axis}
\begin{axis}[
  at={(finbot.south east)}, anchor=south east,
  width=13.75cm, height=3.10cm, scale only axis,
  axis y line*=right, axis x line*=none, xtick=\empty,
  ymin=-7.49, ymax=16.35,
  tick label style={font=\tiny},
  yticklabel style={font=\tiny}, scaled y ticks=false,
  legend style={font=\scriptsize, at={(0.5,-0.44)}, anchor=north,
                legend columns=1, draw=none, fill=none},
]
\addplot[OilBlue, mark=*, mark size=2.2pt, line width=1.3pt] table[x=i, y=ocf]{data/clean/profiles/fin_603609.dat};
\addplot[only marks, OilBlue, mark=*, mark size=3.2pt] table[x=x, y=ocf]{data/clean/profiles/finh_603609.dat};
\legend{经营活动现金流净额（右轴，亿元，蓝点）}
\end{axis}
\end{tikzpicture}

\captionof{figure}[禾丰股份（603609）近年主要财务指标]{禾丰股份（603609）近年主要财务指标（左上：营业收入与归母净利润；右上：ROE 与资产负债率；下：货币资金、有息负债为左轴柱状，经营活动现金流净额为其右轴的蓝点折线。
2021---2025 年为年报口径，2026H1 为 2026 年半年报/半年末，与其前一年报之间留出间隔、以斜纹或空心点区分；半年报数据未年化）}
\end{minipage}
\end{center}

\pfl{财务分析} 2026 年 6 月末资产负债率 57.7\%，较 2025 年末上升 1.4 个百分点（样本中位数 52.2\%，所处三级行业内按由低到高排第\zhnumber{5}位）。 2026 年 6 月末货币资金 12.18 亿元、短期债务（短期借款与一年内到期非流动负债）22.89 亿元、长期债务 43.35 亿元，有息负债合计 66.24 亿元，现金短债比 0.53 倍——覆盖偏紧。 净有息负债 54.06 亿元，相当于其 2026 年上半年营业收入的 29\%。 流动比率 1.22、速动比率 0.58，短期偿债能力处于正常区间。 2026 年上半年经营活动现金流净额 -3.51 亿元，经营现金流与利润同时为负，现金消耗较快。 2026 年 6 月末在建工程 1.07 亿元，较上年末减少 1.77 亿元，在建投入在收缩。 综合看，现金短债比 0.53 倍，短期资金安排偏紧；资产负债率高于样本中位数 5.5 个百分点；有息负债中长期债务占比更高，期限结构对短债滚动的压力较小。

\pfl{重大战略转型} 从公开披露看，公司的战略推进集中在：其他动作（2 项）：2026 年，经营活动现金流同比改善 31.40\%；2025 年 4 月，公司推动 2026 年上半年肉禽业务收入同比增长 32.5\%；融资与资本运作（1 项）：2026 年 4 月，可转债募集资金结项及终止项目形成结余募集资金本金 7.30 亿元、利息收益 313.79 万元。 公告口径上，近三年（2023 年 9 月---2026 年 9 月）公司披露重大重组与并购类公告 5 条、对外投资与转型类公告 0 条、回购与股权激励类公告 52 条，合计公告 464 条。 主营结构的口径变化已经落到报表上：2021 年年报的第一大行业为“饲料业务”（52.1\%），2025 年年报为“肉禽”（45.3\%）；公司在这两期的分部划分方式有过调整。 综合判断，报告期内公司有 5 条与战略相关的公告落地（重大重组与并购 5 条、对外投资与转型 0 条），资本运作与主业调整之间存在直接关联。

\begin{center}\small\setlength{\tabcolsep}{4pt}
\captionof{table}[禾丰股份（603609）历史融资与资本运作]{禾丰股份（603609）历史融资与资本运作（据公司公告与公开报道整理；金额为披露值，含首次公开发行、再融资、债券、银行借款与重整投资等）}
\begin{adjustbox}{max width=\textwidth}
\begin{tabular}{@{}p{1.45cm}p{2.55cm}p{4.05cm}p{4.95cm}@{}}
\toprule
时间 & 方式 & 规模 & 用途与说明 \\
\midrule
2014-08 & IPO（上交所） & 募集资金总额 4.70 亿元 & 上市日 2014-08-08；2014 年年度报告披露“首次发行”募集资金总额 4.70 亿元 \\
2022-04 & 公开发行可转换公司债券 & 15.00 亿元（1,\allowbreak{}500 万张，\allowbreak{}100.00 元/\allowbreak{}张） & 募投项目；；2022-04-28 资金到位，\allowbreak{}初始转股价 10.22 元/\allowbreak{}股 \\
2024 & 股份回购（视同现金分红） & 两轮回购合计 2.87 亿元、\allowbreak{}回购 3,\allowbreak{}998.77 万股；2024 年度另派发现金红利每股 0.06 元（含税） & 维护股价与优化股本结构、\allowbreak{}股东回报；2024 年度现金分红与回购金额合计约 3.38 亿元；截至 2025-04-20 总股本 9.19 亿股 \\
2025 & 股份回购（视同现金分红） & 以现金为对价集中竞价回购金额 5,\allowbreak{}646.65 万元 & 股东回报；2025 年度不派发现金红利、\allowbreak{}不送红股、\allowbreak{}不以公积金转增股本 \\
2026-01 & 回购股份注销（不调整可转债转股价） & --- & 注销回购股份；2026-01-13 公告回购股份注销不调整“禾丰转债”转股价格 \\
2026-09 & 综合授信额度调整 & 2026 年度拟向金融机构申请综合授信总额由“不超过 88.00 亿元”调整为“不超过 89.70 亿元” & 流动资金贷款、\allowbreak{}项目贷款、\allowbreak{}银行承兑汇票、\allowbreak{}保函、\allowbreak{}信用证、\allowbreak{}融资租赁等；2026 年第一次临时股东会会议材料（2026-09-01 披露） \\
\bottomrule
\end{tabular}
\end{adjustbox}
\end{center}

\pfl{历史融投资情况} 从现金流量表看，2025 年公司取得借款收到的现金 32.51 亿元、偿还债务支付的现金 25.24 亿元、吸收权益性投资 0.97 亿元，筹资活动现金流净额 3.00 亿元（2023 年为取得借款 19.59 亿元、偿还债务 19.13 亿元）。 2026 年上半年筹资活动现金流净额为 0.48 亿元（取得借款 19.36 亿元、偿还债务 16.96 亿元，未年化），仍处于净融入状态。 2025 年分配股利、利润或偿付利息支付的现金合计 2.05 亿元。 公开披露的融资记录共 6 笔，工具类型以 IPO、可转债、银行授信为主；金额与用途见下表。 其中规模最大的两笔为：2026 年 9 月，银行授信，2026 年度拟向金融机构申请综合授信总额由“不超过 88.00 亿元”调整为“不超过 89.70 亿元”；2022 年 4 月，可转债，15.00 亿元（1,500 万张，100.00 元/张）。 股本方面，近三年披露股本变动 13 次（可转债转股 9 次、其他 2 次）；总股本由 2021 年末的 9.22 亿股变为 9.19 亿股（-0.27\%）。 分红方面，近三年现金分红 2 次、合计每股派现 0.18 元，最近一次实施在 2024 年 12 月。 近三年“再融资”类公告 27 条，涉及可转债、募集资金使用。 综合看，公司融资结构上以债务滚动为主、股权融资为补充；2025 年筹资活动现金流净额 3.00 亿元，年度内仍为净融入，与前述经营与投资现金流的方向基本一致。

\pfl{未来潜在融资需求分析} 2026 年 6 月末货币资金 12.18 亿元，而短期债务 22.89 亿元（现金短债比 0.53 倍），2026 年上半年经营现金流净额 -3.51 亿元。按“短期债务减去账面货币资金”的滚动偿付口径，静态缺口约 10.71 亿元；若再计入上半年亏损的年化额（0.17 亿元），维持一年正常经营所需的外部资金约 10.71---10.88 亿元。 公司 2026 年上半年归母净利润为负（-0.08 亿元），但 6 月末资产负债率 57.7\% 尚未进入第一档（亏损且负债率 $>$65\%）的区间，当前融资需求以补充流动资金、支撑低谷期经营性支出为主。 公开披露的后续资金安排线索包括：债务与授信安排：2026 年度，综合授信总额由不超过 88.00 亿元上调至不超过 89.70 亿元；2026 年 8 月，为子公司提供担保情况公告，公司以担保方式支持子公司融资。 资金需求还要叠加在建工程：期末在建工程 1.07 亿元、2026 年上半年购建固定资产等支出 1.80 亿元（未年化），这部分支出在周期底部通常会被主动放缓。 综合看，公司的外部资金需求以补充营运资金为主、项目投入为辅，滚动偿付口径下的静态缺口约 10.71 亿元，能否落地取决于授信续作与发行窗口的配合。




\subsubsection{唐人神（002567）}
\label{co:002567}

\pfl{公司基本情况} 公司前身可追溯至 2011 年，湖南株洲的饲料—生猪养殖—肉品一体化企业；2025 年亏损 11.87 亿元，2026 年上半年亏损扩大至 10.46 亿元，以低成本与降负债为穿越周期主线。 公司于 2011-03-25 在深交所主板首发上市，注册资本 14.33 亿元，总股本 14.33 亿股。 截至 2026 年 10 月 12 日收盘价 3.63 元，对应总市值 52.0 亿元，近一年-24.8\%，总市值在三级行业“畜禽饲料”的 11 家公司中按由大到小排第\zhnumber{5}位，近一年涨跌幅低于全样本中位数（-10.4\%）14.5 个百分点。 按 2026 年半年报披露的行业构成，收入占比前三位为饲料 69.0\%；牲猪养殖等 24.9\%；屠宰及肉类 6.0\%。

\pfl{股权结构与控制权} 按 2026 年 6 月 30 日的股东名册，前三大股东为湖南唐人神控股投资股份有限公司 持股 11.24\%；大生行饲料有限公司 持股 4.85\%；中国农业银行股份有限公司－银华农业产业股票型发起式证（持股比例未披露），股东户数 71,230 户。 从控制权角度看，控股股东所持股份中有 9,600 万股处于质押状态；2026 年 5 月，控股股东唐人神控股解除质押 2,700 万股。

\pfl{经营要点} 经营与产能方面，2026 年上半年，销售生猪 272.19 万头、同比增长 4.95\%，全年计划出栏 500 万头左右；2026 年 8 月，生猪销售收入 4.10 亿元、环比上升 31.69\%、同比下降 39.92\% 收入同比降幅较大主要受生猪市场行情波动和出栏结构调整的影响。

\pfl{所处行业的基本面} 饲料的供给即产能（门槛低、产能利用率偏低），需求直接来自下游存栏，利润则由玉米、豆粕等原料价格与成品售价的价差决定（第 \ref{sec:framework} 章的供需框架）。 当期该行业的关键数据是：2026 年上半年，全国生猪出栏 37,246 万头、同比增长 1.7\%；2026 年上半年，生猪价格创近十年新低：仔猪销售均价同比下跌 33.43\%、生猪出场均价同比下跌 28.95\%。

\pfl{竞争优势} 把该公司放回同类公司比较，2026 年上半年营业收入 108.8 亿元、在三级行业“畜禽饲料”的 11 家公司中排第\zhnumber{4}位；毛利率 2.1\%、低于行业内中位数 3.2 个百分点；半年 ROE -24.9\%（未年化）、低于行业内中位数（-6.25\%）；资产负债率 71.1\%、高于行业内中位数（62.1\%）。 与同类公司相比，公司的差异化在于：持续提升母猪场产能利用率；2026 年上半年在饲料销量同比下降 2.79\% 的情况下实现营业收入 75.09 亿元、与上年基本持平。 就当期报表看，公司尚未把经营层面的特征转化为高于行业的资本回报，周期底部的相对优势更多体现在现金流与资产负债表上。

\pfl{历史股价} 本报告样本区间（2021 年 1 月 至 2026 年 10 月）内，股价由 7.93 元变为 3.63 元，累计 -54.2\%，区间最高 10.50 元（2022 年 3 月）、最低 3.35 元（2026 年 6 月）；当前价相当于区间最高价的 34.6\%，为区间最低价的 1.08 倍。 月度收益的年化波动率 33.1\%，低于全样本中位数 37.3\%，在三级行业“畜禽饲料”的 11 家公司中波动率由高到低排第\zhnumber{8}位。

\begin{center}
\begin{minipage}{\textwidth}\centering
\begin{tikzpicture}
\begin{axis}[
  width=12.6cm, height=3.9cm, scale only axis,
  xmin=2020.922, xmax=2026.972, ymin=3.12, ymax=11.24,
  xtick={2021,2022,2023,2024,2025,2026},
  yearticks,
  yticklabel style={font=\scriptsize},
  ylabel={元/股}, ylabel style={font=\scriptsize},
  grid=major, grid style={LightGray, line width=0.3pt},
  axis line style={InkGray, line width=0.5pt},
  every axis plot/.append style={line width=0.9pt},
]
\addplot[AgriGreen] table[x=x,y=y]{data/clean/profiles/px_002567.dat};
\addplot[SoftRed, dashed, line width=0.7pt] table[x=x,y=y]{data/clean/profiles/pxzz_002567.dat};
\addplot[only marks, mark=*, mark size=1.1pt, SoftRed] table[x=x,y=y]{data/clean/profiles/pxzz_002567.dat};
\end{axis}
\end{tikzpicture}

\captionof{figure}[唐人神（002567）股价与周期骨架]{唐人神（002567）月度收盘价（元/股）与 ZigZag 周期骨架（阈值 25\%，圆点为周期转折点）}
\end{minipage}
\end{center}

\pfl{过去周期分析} 以 25\% 的 ZigZag 阈值划分，样本区间内股价形成 2 个主升段与 3 个主跌段。 最大主跌段为 2022 年 7 月 至 2026 年 6 月，47 个月-66.7\%；最大主升段出现在 2021 年 8 月 至 2022 年 3 月，7 个月+79.5\%。 最近一次转折出现在 2026 年 6 月（下跌段自 2022 年 7 月起，47 个月-66.7\%），当前处于该段的延续之中。 公司股价与申万农林牧渔指数几乎同步见底，个股并未走出独立行情。 饲料夹在原料成本与养殖需求之间（第 \ref{sec:feed} 节），量的部分取决于存栏、利的部分取决于原料与成品的价差，公司 2026 年 6 月末资产负债率 71.1\%，高于全样本中位数 52.2\%，其股价因此更多反映资金面与信用风险，而不是价格弹性本身。 综合区间形态看，该股单段涨跌幅度偏大、以趋势段为主（最大主升 +79.5\%、最大主跌 -66.7\%），年化波动率 33.1\% 与全样本中位数 37.3\% 相比波动小于样本中位数。

\pfl{盈利情况} 2026 年上半年营业收入 108.83 亿元（同比 -12.7\%），归母净利润 -10.46 亿元，亏损同比扩大 9.86 亿元。 以年报为趋势对照，2023---2025 年营业收入由 269.49 亿元变为 242.15 亿元（两年年均复合增速 -5.2\%），归母净利润由 -15.26 亿元变为 -11.87 亿元。 按半年报口径，2026 年上半年的 ROE 为 -24.9\%（未年化）、销售净利率 -9.8\%、毛利率 2.1\%，ROE 在“畜禽饲料”11 家公司中按数值由高到低排第\zhnumber{10}位。 按同一口径，三级行业“畜禽饲料”11 家公司的半年 ROE 中位数为 -6.25\%，该公司低于其中位数 18.60 个百分点。 从公司披露的原因看，业绩变动主要来自：2025 年，公司对存栏生猪计提存货跌价准备 3.76 亿元；2025---2026 年，坚持“养猪以低成本为中心”，通过项目制推进成本改善。 公司给出的应对方向是：具有年出栏生猪 180 万头产能。 综合看，公司仍处亏损区间、亏损同比扩大，半年 ROE 在三级行业内按由高到低排第\zhnumber{10}位，单位盈利能力的修复依赖行业景气与自身成本端的改善，报表弹性主要体现为价格与销量的方向性变化。

\begin{center}\small\setlength{\tabcolsep}{3.4pt}
\captionof{table}[唐人神（002567）关键财务指标]{唐人神（002567）关键财务指标（合并报表口径；两个半年报期均未年化；现金短债比 $=$ 货币资金 $\div$ 短期债务）}
\begin{adjustbox}{max width=\textwidth}
\begin{tabular}{@{}lrrrrr@{}}
\toprule
指标 & 2023 年 & 2024 年 & 2025 年 & 2025 年半年报 & 2026 年半年报 \\
\midrule
营业收入（亿元） & 269.49 & 243.43 & 242.15 & 124.68 & 108.83 \\
归母净利润（亿元） & -15.26 & 3.55 & -11.87 & -0.60 & -10.46 \\
毛利率（\%） & 2.7 & 8.8 & 4.1 & 6.6 & 2.1 \\
ROE（\%） & -25.0 & 6.3 & -22.4 & -1.0 & -24.9 \\
资产负债率（\%） & 64.8 & 63.4 & 66.4 & 64.5 & 71.1 \\
经营活动现金流净额（亿元） & 5.36 & 9.36 & 9.15 & 6.60 & -0.66 \\
货币资金（亿元） & 20.91 & 23.05 & 19.47 & 24.86 & 21.10 \\
有息负债合计（亿元） & 76.23 & 67.63 & 58.64 & 67.03 & 61.18 \\
现金短债比（倍） & 0.83 & 1.01 & 0.76 & 0.97 & 0.85 \\
\bottomrule
\end{tabular}
\end{adjustbox}
\end{center}

\begin{center}
\begin{minipage}{\textwidth}\centering
\begin{tikzpicture}
\begin{groupplot}[
  group style={group size=2 by 1, horizontal sep=0.60cm},
  width=6.85cm, height=4.60cm,
  tick label style={font=\tiny},
  title style={font=\scriptsize\heiti},
  yticklabel style={font=\tiny},
  ylabel style={font=\tiny},
  grid=major, grid style={LightGray, line width=0.3pt},
  axis line style={InkGray, line width=0.5pt},
  enlarge x limits={abs=0.8},
  legend style={font=\scriptsize, at={(0.5,-0.20)}, anchor=north,
                legend columns=2, draw=none, fill=none, column sep=6pt},
]
\nextgroupplot[
  ybar, bar width=4pt,
  ymin=-43.73, ymax=303.66,
  xtick={0,1,2,3,4,5.7},
  xticklabels={2021,2022,2023,2024,2025,2026H1},
  ylabel={亿元},
]
\addplot[fill=AgriGreen!75, draw=AgriGreen, bar shift=-2.6pt] table[x=i, y=rev]{data/clean/profiles/fin_002567.dat};
\addplot[fill=SoftRed!60, draw=SoftRed, bar shift=2.6pt] table[x=i, y=ni]{data/clean/profiles/fin_002567.dat};
\addplot[fill=AgriGreen!30, draw=AgriGreen, pattern=north east lines, bar shift=-2.6pt] table[x=x, y=rev]{data/clean/profiles/finh_002567.dat};
\addplot[fill=SoftRed!20, draw=SoftRed, pattern=north east lines, bar shift=2.6pt] table[x=x, y=ni]{data/clean/profiles/finh_002567.dat};
\legend{营业收入（年/半年）, 归母净利润（年/半年）}
\nextgroupplot[
  ymin=-32.7, ymax=80.7,
  xtick={0,1,2,3,4,5.7},
  xticklabels={2021,2022,2023,2024,2025,2026H1},
  ylabel={\%},
]
\addplot[AgriGold, mark=*, mark size=1.3pt] table[x=i, y=roe]{data/clean/profiles/fin_002567.dat};
\addplot[OilBlue, mark=square*, mark size=1.3pt] table[x=i, y=debt]{data/clean/profiles/fin_002567.dat};
\addplot[only marks, AgriGold, mark=*, mark size=2.0pt] table[x=x, y=roe]{data/clean/profiles/finh_002567.dat};
\addplot[only marks, OilBlue, mark=square*, mark size=2.0pt] table[x=x, y=debt]{data/clean/profiles/finh_002567.dat};
\legend{ROE, 资产负债率}
\end{groupplot}
\end{tikzpicture}
\begin{tikzpicture}
\begin{axis}[
  name=finbot,
  width=13.75cm, height=3.10cm, scale only axis,
  ymin=-7.62, ymax=85.37,
  xtick={0,1,2,3,4,5.7},
  xticklabels={2021,2022,2023,2024,2025,2026H1},
  tick label style={font=\tiny},
  yticklabel style={font=\tiny},
  ylabel={亿元}, ylabel style={font=\tiny},
  ybar, bar width=4pt,
  grid=major, grid style={LightGray, line width=0.3pt},
  axis line style={InkGray, line width=0.5pt},
  enlarge x limits={abs=0.8},
  legend style={font=\scriptsize, at={(0.5,-0.26)}, anchor=north,
                legend columns=2, draw=none, fill=none, column sep=10pt},
]
\addplot[fill=AgriGreen!55, draw=AgriGreen, bar shift=-3.0pt] table[x=i, y=cash]{data/clean/profiles/fin_002567.dat};
\addplot[fill=SoftRed!45, draw=SoftRed, bar shift=3.0pt] table[x=i, y=ibd]{data/clean/profiles/fin_002567.dat};
\addplot[fill=AgriGreen!25, draw=AgriGreen, pattern=north east lines, bar shift=-3.0pt] table[x=x, y=cash]{data/clean/profiles/finh_002567.dat};
\addplot[fill=SoftRed!20, draw=SoftRed, pattern=north east lines, bar shift=3.0pt] table[x=x, y=ibd]{data/clean/profiles/finh_002567.dat};
\legend{货币资金（左轴）, 有息负债（左轴）}
\end{axis}
\begin{axis}[
  at={(finbot.south east)}, anchor=south east,
  width=13.75cm, height=3.10cm, scale only axis,
  axis y line*=right, axis x line*=none, xtick=\empty,
  ymin=-7.21, ymax=13.31,
  tick label style={font=\tiny},
  yticklabel style={font=\tiny}, scaled y ticks=false,
  legend style={font=\scriptsize, at={(0.5,-0.44)}, anchor=north,
                legend columns=1, draw=none, fill=none},
]
\addplot[OilBlue, mark=*, mark size=2.2pt, line width=1.3pt] table[x=i, y=ocf]{data/clean/profiles/fin_002567.dat};
\addplot[only marks, OilBlue, mark=*, mark size=3.2pt] table[x=x, y=ocf]{data/clean/profiles/finh_002567.dat};
\legend{经营活动现金流净额（右轴，亿元，蓝点）}
\end{axis}
\end{tikzpicture}

\captionof{figure}[唐人神（002567）近年主要财务指标]{唐人神（002567）近年主要财务指标（左上：营业收入与归母净利润；右上：ROE 与资产负债率；下：货币资金、有息负债为左轴柱状，经营活动现金流净额为其右轴的蓝点折线。
2021---2025 年为年报口径，2026H1 为 2026 年半年报/半年末，与其前一年报之间留出间隔、以斜纹或空心点区分；半年报数据未年化）}
\end{minipage}
\end{center}

\pfl{财务分析} 2026 年 6 月末资产负债率 71.1\%，较 2025 年末上升 4.7 个百分点（样本中位数 52.2\%，所处三级行业内按由低到高排第\zhnumber{10}位）。 2026 年 6 月末货币资金 21.10 亿元、短期债务（短期借款与一年内到期非流动负债）24.84 亿元、长期债务 36.34 亿元，有息负债合计 61.18 亿元，现金短债比 0.85 倍——覆盖偏紧。 净有息负债 40.08 亿元，相当于其 2026 年上半年营业收入的 37\%。 流动比率 0.82、速动比率 0.51，短期偿付压力较大。 2026 年上半年经营活动现金流净额 -0.66 亿元，经营现金流与利润同时为负，日常运营对债务与股东投入的依赖度较高。 2026 年 6 月末在建工程 1.85 亿元，较上年末增加 0.25 亿元，仍有在建投入。 2026 年 6 月末商誉 1.97 亿元，占归母净资产的 4.2\%，是净资产中最脆弱的一块。 综合看，现金短债比 0.85 倍，短期资金安排偏紧；资产负债率高于样本中位数 18.9 个百分点；有息负债中长期债务占比更高，期限结构对短债滚动的压力较小。

\pfl{重大战略转型} 公司的战略动作可归为以下几类：产能与项目（2 项）：2026 年 4 月，公司推进肉品转型：重点投资建设的茶陵香乡猪生产基地具备年出栏 180 万头产能；融资与资本运作（1 项）：2024 年 12 月，公司变更部分募集资金用途：将“东冲三期生猪养殖基地建设项目”剩余募集资金中的 5,800.00 万元变更至“雅安美神养殖”项目。 公告口径上，近三年（2023 年 9 月---2026 年 9 月）公司披露重大重组与并购类公告 4 条、对外投资与转型类公告 2 条、回购与股权激励类公告 43 条，合计公告 578 条。 第一大行业“饲料”的收入占比由 2021 年年报的 80.5\% 变为 69.0\%（11.5 个百分点），收入结构出现再平衡。 综合判断，报告期内公司有 6 条与战略相关的公告落地（重大重组与并购 4 条、对外投资与转型 2 条），资本运作与主业调整之间存在直接关联。

\begin{center}\small\setlength{\tabcolsep}{4pt}
\captionof{table}[唐人神（002567）历史融资与资本运作]{唐人神（002567）历史融资与资本运作（据公司公告与公开报道整理；金额为披露值，含首次公开发行、再融资、债券、银行借款与重整投资等）}
\begin{adjustbox}{max width=\textwidth}
\begin{tabular}{@{}p{1.45cm}p{2.55cm}p{4.05cm}p{4.95cm}@{}}
\toprule
时间 & 方式 & 规模 & 用途与说明 \\
\midrule
2011-03 & IPO（深交所） & 3,\allowbreak{}500 万股、\allowbreak{}27.00 元/\allowbreak{}股（网下 700 万股、\allowbreak{}网上 2,\allowbreak{}800 万股） & 上市日 2011-03-16 \\
2015-11 & 非公开发行股票 & 7,\allowbreak{}240.95 万股、\allowbreak{}8.01 元/\allowbreak{}股 & 实施 2014 年度权益分配方案后发行数量相应调整；2015-11-12 完成登记 \\
2019-12 & 公开发行可转换公司债券（“唐人转债”128092） & 428.00 亿元（1,\allowbreak{}242.8 万张） & 2020-01-22 上市；截至 2020-10-12（停止转股日）累计转股 1.43 亿股 \\
2021-02 & 非公开发行股票（2020 年度方案） & 2.27 亿股、\allowbreak{}6.83 元/\allowbreak{}股，\allowbreak{}募集资金总额 15.50 亿元、\allowbreak{}净额 15.14 亿元 & 239 股）；2021-02-06 资金到账 \\
2022-12 & 非公开发行股票 & 1.75 亿股、\allowbreak{}6.50 元/\allowbreak{}股，\allowbreak{}募集资金总额 11.40 亿元（净额 11.12 亿元） & 2022-12-22 新增股份上市 \\
2023 & 以简易程序向特定对象发行股票（小额快速定增） & 募集资金 3.00 亿元（净额 2.93 亿元） & 数字化信息系统与智能养殖体系建设；后部分募投项目延期至 2028 年 6 月并调整；2025 年年度报告披露累计使用募集资金 8,\allowbreak{}488.62 万元 \\
2026 & 担保额度 & 2026 年对合并报表内单位担保总额不超过 86.17 亿元（约 86.17 亿元） & 支持子公司保理业务、\allowbreak{}融资租赁、\allowbreak{}采购及履约；额度在 2026-01-01 至 2026-12-31 期间滚动使用 \\
2026-01 & 中期票据（科技创新债券） & 2025 年度第二期科技创新债券：注册金额 5.00 亿元，\allowbreak{}本期发行金额 2.00 亿元 & 募集说明书签署日期 2026 年 1 月 \\
\bottomrule
\end{tabular}
\end{adjustbox}
\end{center}

\pfl{历史融投资情况} 从现金流量表看，2025 年公司取得借款收到的现金 21.96 亿元、偿还债务支付的现金 30.63 亿元、吸收权益性投资 1.92 亿元，筹资活动现金流净额 -12.77 亿元（2023 年为取得借款 44.82 亿元、偿还债务 34.54 亿元）。 2026 年上半年筹资活动现金流净额为 2.91 亿元（取得借款 13.72 亿元、偿还债务 12.04 亿元，未年化），仍处于净融入状态。 2025 年分配股利、利润或偿付利息支付的现金合计 2.35 亿元。 公开披露的融资记录共 8 笔，工具类型以 IPO、定向增发、可转债、中期票据为主；金额与用途见下表。 其中规模最大的两笔为：2026 年 1 月，中期票据，2025 年度第二期科技创新债券：注册金额 5.00 亿元，本期发行金额 2.00 亿元；2023 年，定向增发，募集资金 3.00 亿元（净额 2.93 亿元）。 股本方面，近三年披露股本变动 12 次（限售股份上市 3 次、其他 3 次、期权行权 2 次、增发新股上市 1 次）；总股本由 2021 年末的 12.06 亿股变为 14.33 亿股（+18.83\%）。 分红方面，近三年现金分红 1 次、合计每股派现 0.04 元，最近一次实施在 2022 年 12 月。 近三年“再融资”类公告 12 条，涉及定向增发、募集资金使用。 综合看，公司融资结构上以债务滚动为主、股权融资为补充；2025 年筹资活动现金流净额 -12.77 亿元，融资节奏仍以维持存量债务周转为主，与前述经营与投资现金流的方向基本一致。

\pfl{未来潜在融资需求分析} 2026 年 6 月末货币资金 21.10 亿元，而短期债务 24.84 亿元（现金短债比 0.85 倍），2026 年上半年经营现金流净额 -0.66 亿元。按“短期债务减去账面货币资金”的滚动偿付口径，静态缺口约 3.73 亿元；若再计入上半年亏损的年化额（20.93 亿元），维持一年正常经营所需的外部资金约 3.73---24.66 亿元。 按第 \ref{sec:financing-need} 节的三档划分，公司属于第一档“补血型”：2026 年上半年归母净利润 -10.46 亿元、6 月末资产负债率 71.1\%。 可行的融资路径按现实性排序为：出售非核心资产与参股股权、引入国资或产业投资人（部分公司以控制权让渡为对价）、定向增发（需股价配合，且控股股东需放弃优先认购或同步增资）。 公开披露的后续资金安排线索包括：股权融资线索：2025 年度，第二期科技创新债券（中期票据）2.00 亿元发行；债务与授信安排：2026 年，公司担保总额度 86.17 亿元（报表内）+22.40 亿元（报表外），为子公司融资提供支持。 资金需求还要叠加在建工程：期末在建工程 1.85 亿元、2026 年上半年购建固定资产等支出 1.14 亿元（未年化），这部分支出在周期底部通常会被主动放缓。 综合看，公司的外部资金需求以债务接续为主，滚动偿付口径下的静态缺口约 3.73 亿元，融资的时点与规模更多取决于信用条件与资本市场窗口。




\subsubsection{粤海饲料（001313）}
\label{co:001313}

\pfl{公司基本情况} 1994 年设立至今，粤海饲料是总部位于广东湛江的特种水产饲料龙头，以虾料、海水鱼料为核心，2022 年 2 月登陆深交所主板；2024 年因水产行情低迷陷入亏损，2025 年起业绩修复并推进海外与上游纵向并购。 公司于 2022-02-16 在深交所主板首发上市，注册资本 7.00 亿元，总股本 7.00 亿股。 截至 2026 年 10 月 12 日收盘价 7.08 元，对应总市值 49.6 亿元，近一年-1.1\%，总市值在三级行业“水产饲料”的 4 家公司中按由大到小排第\zhnumber{2}位，近一年涨跌幅高于全样本中位数（-10.4\%）9.2 个百分点。 按 2026 年半年报披露的行业构成，收入几乎全部来自饲料行业 100.0\%。

\pfl{股权结构与控制权} 按 2026 年 6 月 30 日的股东名册，前三大股东为湛江市对虾饲料有限公司 持股 35.41\%；香港煌達實業有限公司 持股 23.44\%；湛江承泽投资中心(有限合伙) 持股 6.98\%，股东户数 21,585 户。

\pfl{经营要点} 业务与产能层面，2025 年实现饲料销量 88 万吨；公司整体设计产能 250 万吨以上。

\pfl{所处行业的基本面} 水产饲料的需求由投苗量与存塘鱼价决定，是饲料中毛利率最高、也最依赖服务的品类（第 \ref{sec:framework} 章的供需框架）。 从供需两端的当期数据看，2026 年上半年，公司销量同比增长近 33\%、增速明显高于行业产量增速。

\pfl{竞争优势} 把该公司放回同类公司比较，2026 年上半年营业收入 37.8 亿元、在三级行业“水产饲料”的 4 家公司中排第\zhnumber{2}位；毛利率 11.0\%、高于行业内中位数 2.2 个百分点；半年 ROE 1.9\%（未年化）、高于行业内中位数（1.22\%）；资产负债率 52.1\%、低于行业内中位数（60.5\%）。 公司区别于同类公司的经营特征包括：特种水产饲料市场份额位居行业前列，在华南、华东、华中等省份以及越南共投资设立 30 余家子公司。 优势在当期体现为高于行业中位数的资本回报，但优势能否持续仍取决于周期与自身产能投放节奏。

\pfl{历史股价} 自 2022 年 2 月至 2026 年 10 月，股价由 13.60 元变为 7.08 元，累计 -47.9\%，区间最高 15.15 元（2022 年 3 月）、最低 6.92 元（2024 年 7 月）；当前价相当于区间最高价的 46.7\%，已接近区间最低价（1.02 倍）。 月度收益的年化波动率 30.1\%，低于全样本中位数 37.3\%，在三级行业“水产饲料”的 4 家公司中波动率由高到低排第\zhnumber{3}位。

\begin{center}
\begin{minipage}{\textwidth}\centering
\begin{tikzpicture}
\begin{axis}[
  width=12.6cm, height=3.9cm, scale only axis,
  xmin=2022.005, xmax=2026.972, ymin=6.44, ymax=16.21,
  xtick={2022,2023,2024,2025,2026},
  yearticks,
  yticklabel style={font=\scriptsize},
  ylabel={元/股}, ylabel style={font=\scriptsize},
  grid=major, grid style={LightGray, line width=0.3pt},
  axis line style={InkGray, line width=0.5pt},
  every axis plot/.append style={line width=0.9pt},
]
\addplot[AgriGreen] table[x=x,y=y]{data/clean/profiles/px_001313.dat};
\addplot[SoftRed, dashed, line width=0.7pt] table[x=x,y=y]{data/clean/profiles/pxzz_001313.dat};
\addplot[only marks, mark=*, mark size=1.1pt, SoftRed] table[x=x,y=y]{data/clean/profiles/pxzz_001313.dat};
\end{axis}
\end{tikzpicture}

\captionof{figure}[粤海饲料（001313）股价与周期骨架]{粤海饲料（001313）月度收盘价（元/股）与 ZigZag 周期骨架（阈值 25\%，圆点为周期转折点）}
\end{minipage}
\end{center}

\pfl{过去周期分析} 以 25\% 的 ZigZag 阈值划分，样本区间内股价形成 1 个主升段与 1 个主跌段。 最大主跌段为 2022 年 2 月 至 2024 年 7 月，29 个月-49.1\%；最大主升段出现在 2024 年 7 月 至 2025 年 2 月，7 个月+31.9\%。 最近一次转折出现在 2025 年 2 月（上涨段自 2024 年 7 月起，7 个月+31.9\%），当前处于该段之后的震荡之中。 公司股价较申万农林牧渔指数提前约 22 个月见底，是板块中较早定价周期反转的公司之一。 水产饲料的需求取决于投苗量与存塘鱼价，第 \ref{sec:other-livestock} 章显示 2026 年淡水鱼价格与投苗量已率先回升，其需求驱动与生猪存栏的联系弱于畜禽饲料，公司 2026 年上半年实现盈利、半年 ROE 1.9\%（未年化），好于全样本中位数（0.75\%），下行周期对其报表的冲击小于同业。 综合区间形态看，该股单段涨跌幅度偏大、以趋势段为主（最大主升 +31.9\%、最大主跌 -49.1\%），年化波动率 30.1\% 与全样本中位数 37.3\% 相比波动小于样本中位数。

\pfl{盈利情况} 2026 年上半年营业收入 37.79 亿元（同比 +41.7\%），归母净利润 0.48 亿元，同比 +1240.5\%。 以年报为趋势对照，2023---2025 年营业收入由 68.72 亿元变为 64.89 亿元（两年年均复合增速 -2.8\%），归母净利润由 0.41 亿元变为 0.18 亿元。 按半年报口径，2026 年上半年的 ROE 为 1.9\%（未年化）、销售净利率 1.4\%、毛利率 11.0\%，ROE 在“水产饲料”4 家公司中按数值由高到低排第\zhnumber{2}位。 同期全样本半年 ROE 中位数 0.75\%，公司与之相差 1.11 个百分点（略好于中位数公司）。 综合看，公司在报告期维持盈利（高于全样本半年 ROE 中位数 0.75\%），利润的可持续性取决于水产饲料景气的延续与成本管控的执行。

\begin{center}\small\setlength{\tabcolsep}{3.4pt}
\captionof{table}[粤海饲料（001313）关键财务指标]{粤海饲料（001313）关键财务指标（合并报表口径；两个半年报期均未年化；现金短债比 $=$ 货币资金 $\div$ 短期债务）}
\begin{adjustbox}{max width=\textwidth}
\begin{tabular}{@{}lrrrrr@{}}
\toprule
指标 & 2023 年 & 2024 年 & 2025 年 & 2025 年半年报 & 2026 年半年报 \\
\midrule
营业收入（亿元） & 68.72 & 59.12 & 64.89 & 26.68 & 37.79 \\
归母净利润（亿元） & 0.41 & -0.85 & 0.18 & 0.04 & 0.48 \\
毛利率（\%） & 10.7 & 10.7 & 9.5 & 9.6 & 11.0 \\
ROE（\%） & 1.5 & -3.2 & 0.7 & 0.1 & 1.9 \\
资产负债率（\%） & 41.2 & 46.3 & 46.6 & 45.3 & 52.1 \\
经营活动现金流净额（亿元） & 4.79 & 5.19 & 2.66 & -2.89 & -4.82 \\
货币资金（亿元） & 10.82 & 12.28 & 11.80 & 8.52 & 6.78 \\
有息负债合计（亿元） & 6.51 & 6.58 & 6.17 & 6.81 & 8.40 \\
现金短债比（倍） & 1.87 & 2.38 & 2.55 & 1.62 & 1.23 \\
\bottomrule
\end{tabular}
\end{adjustbox}
\end{center}

\begin{center}
\begin{minipage}{\textwidth}\centering
\begin{tikzpicture}
\begin{groupplot}[
  group style={group size=2 by 1, horizontal sep=0.60cm},
  width=6.85cm, height=4.60cm,
  tick label style={font=\tiny},
  title style={font=\scriptsize\heiti},
  yticklabel style={font=\tiny},
  ylabel style={font=\tiny},
  grid=major, grid style={LightGray, line width=0.3pt},
  axis line style={InkGray, line width=0.5pt},
  enlarge x limits={abs=0.8},
  legend style={font=\scriptsize, at={(0.5,-0.20)}, anchor=north,
                legend columns=2, draw=none, fill=none, column sep=6pt},
]
\nextgroupplot[
  ybar, bar width=4pt,
  ymin=-8.03, ymax=79.53,
  xtick={0,1,2,3,4,5.7},
  xticklabels={2021,2022,2023,2024,2025,2026H1},
  ylabel={亿元},
]
\addplot[fill=AgriGreen!75, draw=AgriGreen, bar shift=-2.6pt] table[x=i, y=rev]{data/clean/profiles/fin_001313.dat};
\addplot[fill=SoftRed!60, draw=SoftRed, bar shift=2.6pt] table[x=i, y=ni]{data/clean/profiles/fin_001313.dat};
\addplot[fill=AgriGreen!30, draw=AgriGreen, pattern=north east lines, bar shift=-2.6pt] table[x=x, y=rev]{data/clean/profiles/finh_001313.dat};
\addplot[fill=SoftRed!20, draw=SoftRed, pattern=north east lines, bar shift=2.6pt] table[x=x, y=ni]{data/clean/profiles/finh_001313.dat};
\legend{营业收入（年/半年）, 归母净利润（年/半年）}
\nextgroupplot[
  ymin=-7.7, ymax=57.6,
  xtick={0,1,2,3,4,5.7},
  xticklabels={2021,2022,2023,2024,2025,2026H1},
  ylabel={\%},
]
\addplot[AgriGold, mark=*, mark size=1.3pt] table[x=i, y=roe]{data/clean/profiles/fin_001313.dat};
\addplot[OilBlue, mark=square*, mark size=1.3pt] table[x=i, y=debt]{data/clean/profiles/fin_001313.dat};
\addplot[only marks, AgriGold, mark=*, mark size=2.0pt] table[x=x, y=roe]{data/clean/profiles/finh_001313.dat};
\addplot[only marks, OilBlue, mark=square*, mark size=2.0pt] table[x=x, y=debt]{data/clean/profiles/finh_001313.dat};
\legend{ROE, 资产负债率}
\end{groupplot}
\end{tikzpicture}
\begin{tikzpicture}
\begin{axis}[
  name=finbot,
  width=13.75cm, height=3.10cm, scale only axis,
  ymin=-1.23, ymax=13.75,
  xtick={0,1,2,3,4,5.7},
  xticklabels={2021,2022,2023,2024,2025,2026H1},
  tick label style={font=\tiny},
  yticklabel style={font=\tiny},
  ylabel={亿元}, ylabel style={font=\tiny},
  ybar, bar width=4pt,
  grid=major, grid style={LightGray, line width=0.3pt},
  axis line style={InkGray, line width=0.5pt},
  enlarge x limits={abs=0.8},
  legend style={font=\scriptsize, at={(0.5,-0.26)}, anchor=north,
                legend columns=2, draw=none, fill=none, column sep=10pt},
]
\addplot[fill=AgriGreen!55, draw=AgriGreen, bar shift=-3.0pt] table[x=i, y=cash]{data/clean/profiles/fin_001313.dat};
\addplot[fill=SoftRed!45, draw=SoftRed, bar shift=3.0pt] table[x=i, y=ibd]{data/clean/profiles/fin_001313.dat};
\addplot[fill=AgriGreen!25, draw=AgriGreen, pattern=north east lines, bar shift=-3.0pt] table[x=x, y=cash]{data/clean/profiles/finh_001313.dat};
\addplot[fill=SoftRed!20, draw=SoftRed, pattern=north east lines, bar shift=3.0pt] table[x=x, y=ibd]{data/clean/profiles/finh_001313.dat};
\legend{货币资金（左轴）, 有息负债（左轴）}
\end{axis}
\begin{axis}[
  at={(finbot.south east)}, anchor=south east,
  width=13.75cm, height=3.10cm, scale only axis,
  axis y line*=right, axis x line*=none, xtick=\empty,
  ymin=-7.71, ymax=9.64,
  tick label style={font=\tiny},
  yticklabel style={font=\tiny}, scaled y ticks=false,
  legend style={font=\scriptsize, at={(0.5,-0.44)}, anchor=north,
                legend columns=1, draw=none, fill=none},
]
\addplot[OilBlue, mark=*, mark size=2.2pt, line width=1.3pt] table[x=i, y=ocf]{data/clean/profiles/fin_001313.dat};
\addplot[only marks, OilBlue, mark=*, mark size=3.2pt] table[x=x, y=ocf]{data/clean/profiles/finh_001313.dat};
\legend{经营活动现金流净额（右轴，亿元，蓝点）}
\end{axis}
\end{tikzpicture}

\captionof{figure}[粤海饲料（001313）近年主要财务指标]{粤海饲料（001313）近年主要财务指标（左上：营业收入与归母净利润；右上：ROE 与资产负债率；下：货币资金、有息负债为左轴柱状，经营活动现金流净额为其右轴的蓝点折线。
2021---2025 年为年报口径，2026H1 为 2026 年半年报/半年末，与其前一年报之间留出间隔、以斜纹或空心点区分；半年报数据未年化）}
\end{minipage}
\end{center}

\pfl{财务分析} 2026 年 6 月末资产负债率 52.1\%，较 2025 年末上升 5.5 个百分点（样本中位数 52.2\%，所处三级行业内按由低到高排第\zhnumber{1}位）。 2026 年 6 月末货币资金 6.78 亿元、短期债务（短期借款与一年内到期非流动负债）5.52 亿元、长期债务 2.89 亿元，有息负债合计 8.40 亿元，现金短债比 1.23 倍——覆盖充分。 净有息负债 1.62 亿元，相当于其 2026 年上半年营业收入的 4\%。 流动比率 1.46、速动比率 1.12，短期偿债能力处于正常区间。 2026 年上半年经营活动现金流净额 -4.82 亿元，净利润为正但经营现金流为负，利润的现金含量偏低。 综合看，现金短债比 1.23 倍，短期偿付压力可控；资产负债率低于样本中位数 0.1 个百分点；有息负债中短期债务占比更高，期限结构仍以滚动续作为主。

\pfl{重大战略转型} 近三年公司的战略动作集中在：海外与国际化（2 项）：2025 年 10 月，总体规划 20 万吨产能，一期工程产能 10 万吨；其他动作（1 项）：2025 年 7 月，第四届董事会第二次会议审议通过以固定资产对全资子公司进行增资，以固定资产对全资子公司增资。 公告口径上，近三年（2023 年 9 月---2026 年 9 月）公司披露重大重组与并购类公告 6 条、对外投资与转型类公告 3 条、回购与股权激励类公告 48 条，合计公告 465 条。 第一大行业仍是“饲料行业”，收入占比 100.0\%，与 2021 年年报相比没有方向性变化。 综合判断，报告期内公司有 9 条与战略相关的公告落地（重大重组与并购 6 条、对外投资与转型 3 条），资本运作与主业调整之间存在直接关联。

\pfl{历史融投资情况} 从现金流量表看，2025 年公司取得借款收到的现金 4.73 亿元、偿还债务支付的现金 5.36 亿元、吸收权益性投资 0.03 亿元，筹资活动现金流净额 -0.99 亿元（2023 年为取得借款 8.08 亿元、偿还债务 11.11 亿元）。 2026 年上半年筹资活动现金流净额为 0.24 亿元（取得借款 5.42 亿元、偿还债务 4.41 亿元，未年化），仍处于净融入状态。 2025 年分配股利、利润或偿付利息支付的现金合计 1.06 亿元。 公开披露的融资记录共 2 笔，工具类型以 IPO 为主；金额与用途见下表。 其中规模最大的两笔为：2022 年 2 月，IPO，公开发行 10,000 万股，发行价 5.38 元/股，募集资金总额 5.38 亿元；2025 年 11 月，现金收购股权，32.00 亿元收购宜兴市天石饲料有限公司 60\% 股权。 股本方面，近三年披露股本变动 5 次（限售股份上市 1 次、其他 1 次）；总股本由 2021 年末的 6.00 亿股变为 7.00 亿股（+16.67\%）。 分红方面，近三年现金分红 3 次、合计每股派现 0.21 元，最近一次实施在 2025 年 6 月。 综合看，公司融资结构上以债务滚动为主、股权融资为补充；2025 年筹资活动现金流净额 -0.99 亿元，融资节奏仍以维持存量债务周转为主，与前述经营与投资现金流的方向基本一致。

\pfl{未来潜在融资需求分析} 2026 年 6 月末货币资金 6.78 亿元，而短期债务 5.52 亿元（现金短债比 1.23 倍），2026 年上半年经营现金流净额 -4.82 亿元。账面货币资金可以覆盖短期债务，缺口不体现在静态偿付层面。 公司 2026 年 6 月末资产负债率 52.1\%、2026 年上半年 ROE 1.9\%（未年化），账面杠杆不高、盈利能力一般，短期内没有必须依靠外部融资才能维持的经营缺口；后续融资更可能指向技改、扩产或产业并购。 静态偿付不存在缺口，后续融资更可能服务于产能或并购安排，而不是填补流动性窟窿。




\subsubsection{天马科技（603668）}
\label{co:603668}

\pfl{公司基本情况} 2017 年设立至今，天马科技是福建特种水产饲料（鳗鲡料等）起家的农牧食品企业，形成“以鳗鱼为核心、以食品为新蓝海、以饲料为主基石”格局；2026 年上半年鳗价低位致净利降 79.53\%。 公司于 2017-01-17 在上交所首发上市，注册资本 5.06 亿元，总股本 5.02 亿股。 截至 2026 年 10 月 12 日收盘价 9.41 元，对应总市值 47.3 亿元，近一年-41.2\%，总市值在三级行业“水产饲料”的 4 家公司中按由大到小排第\zhnumber{3}位，近一年涨跌幅低于全样本中位数（-10.4\%）30.8 个百分点。 按 2025 年年报披露的行业构成，收入占比前三位为饲料行业 87.4\%；养殖及食品行业 31.1\%；其他(补充) 0.5\%。

\pfl{股权结构与控制权} 按 2026 年 6 月 30 日的股东名册，前三大股东为陈庆堂 持股 19.05\%；福建天马投资发展有限公司 持股 12.13\%；杭州皖翰管理咨询合伙企业(有限合伙)（持股比例未披露），股东户数 20,289 户。 股东结构与控制权方面，2026 年 8 月 4 日陈庆堂与一致行动人秋晟天马 1 号私募证券投资基金解除一致行动关系。

\pfl{经营要点} 经营层面，2026 年上半年，动物营养与饲料业务销量 63.60 万吨、同比增长 5.53\%；2026 年上半年，较 2025 年上半年 9502.51 吨减少约 20\%。

\pfl{所处行业的基本面} 水产饲料的需求由投苗量与存塘鱼价决定，是饲料中毛利率最高、也最依赖服务的品类（第 \ref{sec:framework} 章的供需框架）。 从供需两端的当期数据看，2026 年 7 月，日本鳗鱼加工品进口量冲高至 4784 吨、同比增长 84\%；2026 年上半年，全国水产饲料产量 1110 万吨、同比增长 17.4\%。

\pfl{竞争优势} 把该公司放回同类公司比较，2026 年上半年营业收入 31.3 亿元、在三级行业“水产饲料”的 4 家公司中排第\zhnumber{3}位；毛利率 8.3\%、低于行业内中位数 0.5 个百分点；半年 ROE 0.6\%（未年化）、低于行业内中位数（1.22\%）；资产负债率 72.1\%、高于行业内中位数（60.5\%）。 公司区别于同类公司的经营特征包括：公司是全球规模最大、唯一深耕鳗鱼全产业链的 A 股上市企业；公司畜禽饲料年产销量已连续 5 年突破百万吨、稳居福建省第一梯队；截至 2026 年 6 月末公司总资产 96.36 亿元、有息负债 41.55 亿元，其中短期债务占全部有息负债比重达 83.18\%。 就当期报表看，公司尚未把经营层面的特征转化为高于行业的资本回报，周期底部的相对优势更多体现在现金流与资产负债表上。

\pfl{历史股价} 以 2021 年 1 月为起点、2026 年 10 月为终点，股价由 6.91 元变为 9.41 元，累计 +36.2\%，区间最高 21.35 元（2022 年 7 月）、最低 6.91 元（2021 年 1 月）；当前价相当于区间最高价的 44.1\%，为区间最低价的 1.36 倍。 月度收益的年化波动率 45.6\%，高于全样本中位数 37.3\%，在三级行业“水产饲料”的 4 家公司中波动率由高到低排第\zhnumber{1}位。

\begin{center}
\begin{minipage}{\textwidth}\centering
\begin{tikzpicture}
\begin{axis}[
  width=12.6cm, height=3.9cm, scale only axis,
  xmin=2020.922, xmax=2026.972, ymin=6.43, ymax=22.84,
  xtick={2021,2022,2023,2024,2025,2026},
  yearticks,
  yticklabel style={font=\scriptsize},
  ylabel={元/股}, ylabel style={font=\scriptsize},
  grid=major, grid style={LightGray, line width=0.3pt},
  axis line style={InkGray, line width=0.5pt},
  every axis plot/.append style={line width=0.9pt},
]
\addplot[AgriGreen] table[x=x,y=y]{data/clean/profiles/px_603668.dat};
\addplot[SoftRed, dashed, line width=0.7pt] table[x=x,y=y]{data/clean/profiles/pxzz_603668.dat};
\addplot[only marks, mark=*, mark size=1.1pt, SoftRed] table[x=x,y=y]{data/clean/profiles/pxzz_603668.dat};
\end{axis}
\end{tikzpicture}

\captionof{figure}[天马科技（603668）股价与周期骨架]{天马科技（603668）月度收盘价（元/股）与 ZigZag 周期骨架（阈值 25\%，圆点为周期转折点）}
\end{minipage}
\end{center}

\pfl{过去周期分析} 以 25\% 的 ZigZag 阈值划分，样本区间内股价形成 2 个主升段与 2 个主跌段。 最大主升段出现在 2021 年 1 月 至 2022 年 7 月，18 个月+209.0\%；最大主跌段为 2022 年 7 月 至 2025 年 1 月，30 个月-48.5\%。 最近一次转折出现在 2026 年 10 月（下跌段自 2025 年 12 月起，10 个月-43.5\%），当前处于该段的延续之中。 公司股价较申万农林牧渔指数提前约 64 个月见底，是板块中较早定价周期反转的公司之一。 水产饲料的需求取决于投苗量与存塘鱼价，第 \ref{sec:other-livestock} 章显示 2026 年淡水鱼价格与投苗量已率先回升，其需求驱动与生猪存栏的联系弱于畜禽饲料，公司 2026 年 6 月末资产负债率 72.1\%，高于全样本中位数 52.2\%，其股价因此更多反映资金面与信用风险，而不是价格弹性本身。 综合区间形态看，该股单段涨跌幅度偏大、以趋势段为主（最大主升 +209.0\%、最大主跌 -48.5\%），年化波动率 45.6\% 与全样本中位数 37.3\% 相比波动明显大于样本中位数。

\pfl{盈利情况} 2026 年上半年营业收入 31.32 亿元（同比 +5.3\%），归母净利润 0.12 亿元，同比 -79.5\%。 以年报为趋势对照，2023---2025 年营业收入由 69.98 亿元变为 60.01 亿元（两年年均复合增速 -7.4\%），归母净利润由 -1.88 亿元变为 -1.79 亿元。 按半年报口径，2026 年上半年的 ROE 为 0.6\%（未年化）、销售净利率 0.5\%、毛利率 8.3\%，ROE 在“水产饲料”4 家公司中按数值由高到低排第\zhnumber{3}位。 按同一口径，三级行业“水产饲料”4 家公司的半年 ROE 中位数为 1.22\%，该公司低于其中位数 0.64 个百分点。 从公司披露的原因看，业绩变动主要来自：2026 年，活鳗及烤鳗价格持续低位。 综合看，公司在报告期维持盈利（低于全样本半年 ROE 中位数 0.75\%），利润的可持续性取决于水产饲料景气的延续与成本管控的执行。

\begin{center}\small\setlength{\tabcolsep}{3.4pt}
\captionof{table}[天马科技（603668）关键财务指标]{天马科技（603668）关键财务指标（合并报表口径；两个半年报期均未年化；现金短债比 $=$ 货币资金 $\div$ 短期债务）}
\begin{adjustbox}{max width=\textwidth}
\begin{tabular}{@{}lrrrrr@{}}
\toprule
指标 & 2023 年 & 2024 年 & 2025 年 & 2025 年半年报 & 2026 年半年报 \\
\midrule
营业收入（亿元） & 69.98 & 58.54 & 60.01 & 29.74 & 31.32 \\
归母净利润（亿元） & -1.88 & 0.26 & -1.79 & 0.60 & 0.12 \\
毛利率（\%） & 7.5 & 10.9 & 7.3 & 10.7 & 8.3 \\
ROE（\%） & -9.5 & 1.2 & -8.3 & 2.7 & 0.6 \\
资产负债率（\%） & 70.8 & 68.6 & 72.6 & 69.6 & 72.1 \\
经营活动现金流净额（亿元） & 0.71 & 4.83 & 1.01 & 0.20 & 1.87 \\
货币资金（亿元） & 7.96 & 5.11 & 5.36 & 5.45 & 4.61 \\
有息负债合计（亿元） & 36.26 & 37.19 & 42.42 & 43.64 & 41.55 \\
现金短债比（倍） & 0.30 & 0.16 & 0.15 & 0.14 & 0.13 \\
\bottomrule
\end{tabular}
\end{adjustbox}
\end{center}

\begin{center}
\begin{minipage}{\textwidth}\centering
\begin{tikzpicture}
\begin{groupplot}[
  group style={group size=2 by 1, horizontal sep=0.60cm},
  width=6.85cm, height=4.60cm,
  tick label style={font=\tiny},
  title style={font=\scriptsize\heiti},
  yticklabel style={font=\tiny},
  ylabel style={font=\tiny},
  grid=major, grid style={LightGray, line width=0.3pt},
  axis line style={InkGray, line width=0.5pt},
  enlarge x limits={abs=0.8},
  legend style={font=\scriptsize, at={(0.5,-0.20)}, anchor=north,
                legend columns=2, draw=none, fill=none, column sep=6pt},
]
\nextgroupplot[
  ybar, bar width=4pt,
  ymin=-9.07, ymax=78.71,
  xtick={0,1,2,3,4,5.7},
  xticklabels={2021,2022,2023,2024,2025,2026H1},
  ylabel={亿元},
]
\addplot[fill=AgriGreen!75, draw=AgriGreen, bar shift=-2.6pt] table[x=i, y=rev]{data/clean/profiles/fin_603668.dat};
\addplot[fill=SoftRed!60, draw=SoftRed, bar shift=2.6pt] table[x=i, y=ni]{data/clean/profiles/fin_603668.dat};
\addplot[fill=AgriGreen!30, draw=AgriGreen, pattern=north east lines, bar shift=-2.6pt] table[x=x, y=rev]{data/clean/profiles/finh_603668.dat};
\addplot[fill=SoftRed!20, draw=SoftRed, pattern=north east lines, bar shift=2.6pt] table[x=x, y=ni]{data/clean/profiles/finh_603668.dat};
\legend{营业收入（年/半年）, 归母净利润（年/半年）}
\nextgroupplot[
  ymin=-16.1, ymax=80.8,
  xtick={0,1,2,3,4,5.7},
  xticklabels={2021,2022,2023,2024,2025,2026H1},
  ylabel={\%},
]
\addplot[AgriGold, mark=*, mark size=1.3pt] table[x=i, y=roe]{data/clean/profiles/fin_603668.dat};
\addplot[OilBlue, mark=square*, mark size=1.3pt] table[x=i, y=debt]{data/clean/profiles/fin_603668.dat};
\addplot[only marks, AgriGold, mark=*, mark size=2.0pt] table[x=x, y=roe]{data/clean/profiles/finh_603668.dat};
\addplot[only marks, OilBlue, mark=square*, mark size=2.0pt] table[x=x, y=debt]{data/clean/profiles/finh_603668.dat};
\legend{ROE, 资产负债率}
\end{groupplot}
\end{tikzpicture}
\begin{tikzpicture}
\begin{axis}[
  name=finbot,
  width=13.75cm, height=3.10cm, scale only axis,
  ymin=-4.24, ymax=47.51,
  xtick={0,1,2,3,4,5.7},
  xticklabels={2021,2022,2023,2024,2025,2026H1},
  tick label style={font=\tiny},
  yticklabel style={font=\tiny},
  ylabel={亿元}, ylabel style={font=\tiny},
  ybar, bar width=4pt,
  grid=major, grid style={LightGray, line width=0.3pt},
  axis line style={InkGray, line width=0.5pt},
  enlarge x limits={abs=0.8},
  legend style={font=\scriptsize, at={(0.5,-0.26)}, anchor=north,
                legend columns=2, draw=none, fill=none, column sep=10pt},
]
\addplot[fill=AgriGreen!55, draw=AgriGreen, bar shift=-3.0pt] table[x=i, y=cash]{data/clean/profiles/fin_603668.dat};
\addplot[fill=SoftRed!45, draw=SoftRed, bar shift=3.0pt] table[x=i, y=ibd]{data/clean/profiles/fin_603668.dat};
\addplot[fill=AgriGreen!25, draw=AgriGreen, pattern=north east lines, bar shift=-3.0pt] table[x=x, y=cash]{data/clean/profiles/finh_603668.dat};
\addplot[fill=SoftRed!20, draw=SoftRed, pattern=north east lines, bar shift=3.0pt] table[x=x, y=ibd]{data/clean/profiles/finh_603668.dat};
\legend{货币资金（左轴）, 有息负债（左轴）}
\end{axis}
\begin{axis}[
  at={(finbot.south east)}, anchor=south east,
  width=13.75cm, height=3.10cm, scale only axis,
  axis y line*=right, axis x line*=none, xtick=\empty,
  ymin=-1.26, ymax=6.28,
  tick label style={font=\tiny},
  yticklabel style={font=\tiny}, scaled y ticks=false,
  legend style={font=\scriptsize, at={(0.5,-0.44)}, anchor=north,
                legend columns=1, draw=none, fill=none},
]
\addplot[OilBlue, mark=*, mark size=2.2pt, line width=1.3pt] table[x=i, y=ocf]{data/clean/profiles/fin_603668.dat};
\addplot[only marks, OilBlue, mark=*, mark size=3.2pt] table[x=x, y=ocf]{data/clean/profiles/finh_603668.dat};
\legend{经营活动现金流净额（右轴，亿元，蓝点）}
\end{axis}
\end{tikzpicture}

\captionof{figure}[天马科技（603668）近年主要财务指标]{天马科技（603668）近年主要财务指标（左上：营业收入与归母净利润；右上：ROE 与资产负债率；下：货币资金、有息负债为左轴柱状，经营活动现金流净额为其右轴的蓝点折线。
2021---2025 年为年报口径，2026H1 为 2026 年半年报/半年末，与其前一年报之间留出间隔、以斜纹或空心点区分；半年报数据未年化）}
\end{minipage}
\end{center}

\pfl{财务分析} 2026 年 6 月末资产负债率 72.1\%，较 2025 年末下降 0.6 个百分点（样本中位数 52.2\%，所处三级行业内按由低到高排第\zhnumber{4}位）。 2026 年 6 月末货币资金 4.61 亿元、短期债务（短期借款与一年内到期非流动负债）34.56 亿元、长期债务 7.00 亿元，有息负债合计 41.55 亿元，现金短债比 0.13 倍——缺口明显。 净有息负债 36.95 亿元，相当于其 2026 年上半年营业收入的 118\%。 流动比率 0.88、速动比率 0.24，短期偿债指标偏紧。 2026 年上半年经营活动现金流净额 1.87 亿元，经营现金流与利润同向为正，内生资金可以支撑运营。 2026 年 6 月末在建工程 1.99 亿元，较上年末减少 0.54 亿元，在建投入在收缩。 2026 年 6 月末商誉 1.17 亿元，占归母净资产的 4.3\%，是净资产中最脆弱的一块。 综合看，账面货币资金对有息负债与短期债务的覆盖不足，滚动偿付对新增融资的依赖度较高；资产负债率高于样本中位数 19.9 个百分点；有息负债中短期债务占比更高，期限结构仍以滚动续作为主。 公司披露的股东与质押风险为：2026 年 8 月，实际控制人陈庆堂减持：2026-08-05 减持 207.99 万股（均价 11 元）、2026-08-10 减持 452.66 万股（均价 10 元）；2026 年 9 月，控股股东、实际控制人之一致行动人部分股份解除质押及质押。

\pfl{重大战略转型} 近三年公司的战略动作集中在：其他动作（2 项）：2026 年，食品（烤鳗）销量创新高；2025 年 6 月，股票期权激励计划：第一个行权期行权条件于 2025 年 6 月经董事会审议通过成就；融资与资本运作（1 项）：2023 年，公司以简易程序定增募资建设“鳗鲡生态养殖基地建设项目（二期）”。 公告口径上，近三年（2023 年 9 月---2026 年 9 月）公司披露重大重组与并购类公告 0 条、对外投资与转型类公告 1 条、回购与股权激励类公告 12 条，合计公告 587 条。 第一大行业“饲料行业”的收入占比由 2021 年年报的 99.8\% 变为 87.4\%（12.4 个百分点），收入结构出现再平衡。 综合判断，报告期内公司有 1 条与战略相关的公告落地（重大重组与并购 0 条、对外投资与转型 1 条），资本运作与主业调整之间存在直接关联。

\begin{center}\small\setlength{\tabcolsep}{4pt}
\captionof{table}[天马科技（603668）历史融资与资本运作]{天马科技（603668）历史融资与资本运作（据公司公告与公开报道整理；金额为披露值，含首次公开发行、再融资、债券、银行借款与重整投资等）}
\begin{adjustbox}{max width=\textwidth}
\begin{tabular}{@{}p{1.45cm}p{2.55cm}p{4.05cm}p{4.95cm}@{}}
\toprule
时间 & 方式 & 规模 & 用途与说明 \\
\midrule
2017-01 & IPO（上交所主板） & 募集资金总额 3.29 亿元，\allowbreak{}净额 2.91 亿元 & 特种水产配合饲料生产项目（二期）等；发行价 6.21 元/\allowbreak{}股；发行前陈庆堂直接持股 38.43\%、\allowbreak{}通过天马投资控制 10.54\% \\
2018-04 & 公开发行可转换公司债券 & 5.00 亿元（305 万张，\allowbreak{}面值 100.00 元） & 特种水产配合饲料生产项目（二期）；；采用股份质押+保证担保 \\
2021-07 & 非公开发行股票（2020 年度方案） & 9,\allowbreak{}655.17 万股、\allowbreak{}5.80 元/\allowbreak{}股，\allowbreak{}募集资金总额 5.60 亿元、\allowbreak{}净额 5.52 亿元 & 2021-07-21 资金到账，\allowbreak{}7 月 30 日完成登记，\allowbreak{}总股本由 3.40 亿股增至 4.36 亿股 \\
2023-12 & 以简易程序向特定对象发行股票 & 2,\allowbreak{}053.39 万股、\allowbreak{}14.61 元/\allowbreak{}股 & 鳗鲡生态养殖基地建设项目（二期）及补充流动资金；；2023-12-29 资金到位 \\
2024-10-15 & 接受关联方财务资助 & 控股股东、\allowbreak{}实际控制人陈庆堂之全资控股公司天马投资提供最高不超过 3.50 亿元财务资助 & 公司及合并范围内子公司经营资金；额度有效期至 2025 年 12 月 31 日 \\
2025-12-09 & 接受关联方财务资助 & 天马投资提供最高不超过 4.50 亿元财务资助（无抵押、\allowbreak{}质押等担保措施） & 公司及子公司经营资金；利率不高于 LPR；截至 2026-09-16 公告日，\allowbreak{}天马投资向公司及子公司提供的借款本金余额为 1.88 亿元 \\
\bottomrule
\end{tabular}
\end{adjustbox}
\end{center}

\pfl{历史融投资情况} 从现金流量表看，2025 年公司取得借款收到的现金 35.80 亿元、偿还债务支付的现金 34.74 亿元、吸收权益性投资 0.46 亿元，筹资活动现金流净额 0.38 亿元（2023 年为取得借款 31.76 亿元、偿还债务 29.00 亿元）。 2026 年上半年筹资活动现金流净额为 -1.59 亿元（取得借款 13.55 亿元、偿还债务 12.64 亿元，未年化），已转为净流出，处于净还债状态。 2025 年分配股利、利润或偿付利息支付的现金合计 1.54 亿元。 公开披露的融资记录共 6 笔，工具类型以 IPO、可转债、定向增发为主；金额与用途见下表。 其中规模最大的两笔为：2025 年 12 月，接受关联方财务资助，天马投资提供最高不超过 4.50 亿元财务资助（无抵押、质押等担保措施）；2024 年 10 月，接受关联方财务资助，控股股东、实际控制人陈庆堂之全资控股公司天马投资提供最高不超过 3.50 亿元财务资助。 股本方面，近三年披露股本变动 7 次（限售股份上市 2 次、增发新股上市 1 次、转增 1 次）；总股本由 2021 年末的 4.36 亿股变为 5.02 亿股（+15.17\%）。 分红方面，近三年现金分红 4 次、合计每股派现 0.07 元，最近一次实施在 2024 年 6 月。 近三年“再融资”类公告 26 条，涉及定向增发、募集资金使用。 综合看，公司融资结构上股权与债务融资并举；2025 年筹资活动现金流净额 0.38 亿元，年度内仍为净融入，与前述经营与投资现金流的方向基本一致。

\pfl{未来潜在融资需求分析} 2026 年 6 月末货币资金 4.61 亿元，而短期债务 34.56 亿元（现金短债比 0.13 倍），2026 年上半年经营现金流净额 1.87 亿元。按“短期债务减去账面货币资金”的滚动偿付口径，静态缺口约 29.95 亿元；上半年盈利或接近盈亏平衡，该缺口不随经营亏损进一步扩大。 公司 2026 年 6 月末资产负债率 72.1\%、2026 年上半年 ROE 0.6\%（未年化），资产负债率偏高（72.1\%），融资需求首先用于置换与滚动存量债务、补充营运资金，属于“补血型”而非扩张型。 公开披露的后续资金安排线索包括：债务与授信安排：2025 年 12 月，截至 2026-09-16 借款本金余额 1.88 亿元。 资金需求还要叠加在建工程：期末在建工程 1.99 亿元、2026 年上半年购建固定资产等支出 1.16 亿元（未年化），这部分支出在周期底部通常会被主动放缓。 综合看，公司的外部资金需求以补充流动性与项目投入为主，滚动偿付口径下的静态缺口约 29.95 亿元，能否落地取决于授信续作与发行窗口的配合。




\subsubsection{金新农（002548）}
\label{co:002548}

\pfl{公司基本情况} 公司前身可追溯至 2008 年，深圳的饲料与生猪养殖双主业企业，2011 年上市，历经 2019 年、2020 年两轮控制权与股东结构变动后由广州金农产投控股；2026 年上半年猪价创近十年新低，归母净利润亏损 3.88 亿元，净资产大幅收缩。 公司于 2011-02-18 在深交所主板首发上市，注册资本 8.28 亿元，总股本 8.05 亿股。 截至 2026 年 10 月 12 日收盘价 3.83 元，对应总市值 30.8 亿元，近一年-24.3\%，总市值在三级行业“畜禽饲料”的 11 家公司中按由大到小排第\zhnumber{6}位，近一年涨跌幅低于全样本中位数（-10.4\%）14.0 个百分点。 按 2026 年半年报披露的行业构成，收入占比前三位为饲料加工 69.7\%；畜牧养殖 28.8\%；其他 1.5\%。

\pfl{股权结构与控制权} 按 2026 年 6 月 30 日的股东名册，前三大股东为广州金农产业投资合伙企业(有限合伙) 持股 30.27\%；浙商银行股份有限公司－国泰中证畜牧养殖交易型开放式指（持股比例未披露）；MORGANSTANLEY \& CO. INTERN 持股 0.91\%，股东户数 47,370 户。 控制权与股东性质方面，持股 30.27\%（2.507 亿股），其中被质押 5,700 万股；由于控股股东上层股权分散、任一股东均无法形成控制。

\pfl{经营要点} 业务与产能层面，2026 年上半年，饲料销量 53.28 万吨（不含内销 14.23 万吨）、同比增长 2.27\%，其中猪用饲料 46.40 万吨、占比 87.08\%；2026 年上半年，生猪销量 87.00 万头、同比增长 37.60\%；2026 年 9 月，生猪销量 10.82 万头、环比下降 23.86\%。

\pfl{所处行业的基本面} 饲料的供给即产能（门槛低、产能利用率偏低），需求直接来自下游存栏，利润则由玉米、豆粕等原料价格与成品售价的价差决定（第 \ref{sec:framework} 章的供需框架）。 行业层面的最新情况是：2026 年上半年，全国工业饲料总产量 16,850 万吨、同比增长 3.9\%。

\pfl{竞争优势} 把该公司放回同类公司比较，2026 年上半年营业收入 23.2 亿元、在三级行业“畜禽饲料”的 11 家公司中排第\zhnumber{8}位；毛利率 -2.5\%、低于行业内中位数 7.8 个百分点；半年 ROE -37.4\%（未年化）、低于行业内中位数（-6.25\%）；资产负债率 81.5\%、高于行业内中位数（62.1\%）。 公司在公开资料中自述与同业比较的相对优势在于：猪用饲料占饲料销量的 87.08\%。 就当期报表看，公司尚未把经营层面的特征转化为高于行业的资本回报，周期底部的相对优势更多体现在现金流与资产负债表上。

\pfl{历史股价} 本报告样本区间（2021 年 1 月 至 2026 年 10 月）内，股价由 8.03 元变为 3.83 元，累计 -52.3\%，区间最高 8.16 元（2022 年 3 月）、最低 3.69 元（2026 年 6 月）；当前价相当于区间最高价的 46.9\%，已接近区间最低价（1.04 倍）。 月度收益的年化波动率 39.9\%，高于全样本中位数 37.3\%，在三级行业“畜禽饲料”的 11 家公司中波动率由高到低排第\zhnumber{5}位。 把股价阶段与公司事件对齐看，2024 年 8 月 至 2026 年 2 月股价上涨+93.0\%，同期（2025 年）半年度末资产权利受限合计账面余额 10.90 亿元。 整体上看，区间的几处大级别拐点都能在公司的资本运作、股权变动或风险事件上找到对应，股价波动有明显的公司层面驱动，而非单纯的行业 beta。

\begin{center}
\begin{minipage}{\textwidth}\centering
\begin{tikzpicture}
\begin{axis}[
  width=12.6cm, height=3.9cm, scale only axis,
  xmin=2020.922, xmax=2026.972, ymin=3.43, ymax=8.73,
  xtick={2021,2022,2023,2024,2025,2026},
  yearticks,
  yticklabel style={font=\scriptsize},
  ylabel={元/股}, ylabel style={font=\scriptsize},
  grid=major, grid style={LightGray, line width=0.3pt},
  axis line style={InkGray, line width=0.5pt},
  every axis plot/.append style={line width=0.9pt},
]
\addplot[AgriGreen] table[x=x,y=y]{data/clean/profiles/px_002548.dat};
\addplot[SoftRed, dashed, line width=0.7pt] table[x=x,y=y]{data/clean/profiles/pxzz_002548.dat};
\addplot[only marks, mark=*, mark size=1.1pt, SoftRed] table[x=x,y=y]{data/clean/profiles/pxzz_002548.dat};
\end{axis}
\end{tikzpicture}

\captionof{figure}[金新农（002548）股价与周期骨架]{金新农（002548）月度收盘价（元/股）与 ZigZag 周期骨架（阈值 25\%，圆点为周期转折点）}
\end{minipage}
\end{center}

\pfl{过去周期分析} 以 25\% 的 ZigZag 阈值划分，样本区间内股价形成 3 个主升段与 4 个主跌段。 最大主升段出现在 2024 年 8 月 至 2026 年 2 月，18 个月+93.0\%；最大主跌段为 2023 年 6 月 至 2024 年 8 月，14 个月-51.4\%。 最近一次转折出现在 2026 年 6 月（下跌段自 2026 年 2 月起，4 个月-48.6\%），当前处于该段的延续之中。 公司股价与申万农林牧渔指数几乎同步见底，个股并未走出独立行情。 饲料夹在原料成本与养殖需求之间（第 \ref{sec:feed} 节），量的部分取决于存栏、利的部分取决于原料与成品的价差，公司 2026 年 6 月末资产负债率 81.5\%，高于全样本中位数 52.2\%，其股价因此更多反映资金面与信用风险，而不是价格弹性本身。 综合区间形态看，该股单段涨跌幅度偏大、以趋势段为主（最大主升 +93.0\%、最大主跌 -51.4\%），年化波动率 39.9\% 与全样本中位数 37.3\% 相比波动与样本中位数接近。

\pfl{盈利情况} 2026 年上半年营业收入 23.20 亿元（同比 -2.4\%），归母净利润 -3.88 亿元，亏损同比扩大 3.65 亿元。 以年报为趋势对照，2023---2025 年营业收入由 40.40 亿元变为 47.46 亿元（两年年均复合增速 +8.4\%），归母净利润由 -6.60 亿元变为 -2.93 亿元。 按半年报口径，2026 年上半年的 ROE 为 -37.4\%（未年化）、销售净利率 -16.8\%、毛利率 -2.5\%，ROE 在“畜禽饲料”11 家公司中按数值由高到低排第\zhnumber{11}位。 按同一口径，三级行业“畜禽饲料”11 家公司的半年 ROE 中位数为 -6.25\%，该公司低于其中位数 31.19 个百分点。 综合看，公司仍处亏损区间、亏损同比扩大，半年 ROE 在三级行业内按由高到低排第\zhnumber{11}位，单位盈利能力的修复依赖行业景气与自身成本端的改善，报表弹性主要体现为价格与销量的方向性变化。

\begin{center}\small\setlength{\tabcolsep}{3.4pt}
\captionof{table}[金新农（002548）关键财务指标]{金新农（002548）关键财务指标（合并报表口径；两个半年报期均未年化；现金短债比 $=$ 货币资金 $\div$ 短期债务）}
\begin{adjustbox}{max width=\textwidth}
\begin{tabular}{@{}lrrrrr@{}}
\toprule
指标 & 2023 年 & 2024 年 & 2025 年 & 2025 年半年报 & 2026 年半年报 \\
\midrule
营业收入（亿元） & 40.40 & 45.62 & 47.46 & 23.76 & 23.20 \\
归母净利润（亿元） & -6.60 & 0.30 & -2.93 & -0.24 & -3.88 \\
毛利率（\%） & 3.3 & 10.6 & 4.8 & 7.5 & -2.5 \\
ROE（\%） & -36.5 & 2.0 & -21.4 & -1.6 & -37.4 \\
资产负债率（\%） & 74.1 & 70.4 & 75.7 & 72.1 & 81.5 \\
经营活动现金流净额（亿元） & 1.25 & 4.23 & 3.32 & 1.45 & 0.50 \\
货币资金（亿元） & 7.54 & 3.38 & 3.05 & 5.11 & 2.36 \\
有息负债合计（亿元） & 29.75 & 27.57 & 28.60 & 30.67 & 28.92 \\
现金短债比（倍） & 0.36 & 0.17 & 0.14 & 0.22 & 0.10 \\
\bottomrule
\end{tabular}
\end{adjustbox}
\end{center}

\begin{center}
\begin{minipage}{\textwidth}\centering
\begin{tikzpicture}
\begin{groupplot}[
  group style={group size=2 by 1, horizontal sep=0.60cm},
  width=6.85cm, height=4.60cm,
  tick label style={font=\tiny},
  title style={font=\scriptsize\heiti},
  yticklabel style={font=\tiny},
  ylabel style={font=\tiny},
  grid=major, grid style={LightGray, line width=0.3pt},
  axis line style={InkGray, line width=0.5pt},
  enlarge x limits={abs=0.8},
  legend style={font=\scriptsize, at={(0.5,-0.20)}, anchor=north,
                legend columns=2, draw=none, fill=none, column sep=6pt},
]
\nextgroupplot[
  ybar, bar width=4pt,
  ymin=-15.78, ymax=55.70,
  xtick={0,1,2,3,4,5.7},
  xticklabels={2021,2022,2023,2024,2025,2026H1},
  ylabel={亿元},
]
\addplot[fill=AgriGreen!75, draw=AgriGreen, bar shift=-2.6pt] table[x=i, y=rev]{data/clean/profiles/fin_002548.dat};
\addplot[fill=SoftRed!60, draw=SoftRed, bar shift=2.6pt] table[x=i, y=ni]{data/clean/profiles/fin_002548.dat};
\addplot[fill=AgriGreen!30, draw=AgriGreen, pattern=north east lines, bar shift=-2.6pt] table[x=x, y=rev]{data/clean/profiles/finh_002548.dat};
\addplot[fill=SoftRed!20, draw=SoftRed, pattern=north east lines, bar shift=2.6pt] table[x=x, y=ni]{data/clean/profiles/finh_002548.dat};
\legend{营业收入（年/半年）, 归母净利润（年/半年）}
\nextgroupplot[
  ymin=-52.1, ymax=93.9,
  xtick={0,1,2,3,4,5.7},
  xticklabels={2021,2022,2023,2024,2025,2026H1},
  ylabel={\%},
]
\addplot[AgriGold, mark=*, mark size=1.3pt] table[x=i, y=roe]{data/clean/profiles/fin_002548.dat};
\addplot[OilBlue, mark=square*, mark size=1.3pt] table[x=i, y=debt]{data/clean/profiles/fin_002548.dat};
\addplot[only marks, AgriGold, mark=*, mark size=2.0pt] table[x=x, y=roe]{data/clean/profiles/finh_002548.dat};
\addplot[only marks, OilBlue, mark=square*, mark size=2.0pt] table[x=x, y=debt]{data/clean/profiles/finh_002548.dat};
\legend{ROE, 资产负债率}
\end{groupplot}
\end{tikzpicture}
\begin{tikzpicture}
\begin{axis}[
  name=finbot,
  width=13.75cm, height=3.10cm, scale only axis,
  ymin=-3.12, ymax=34.93,
  xtick={0,1,2,3,4,5.7},
  xticklabels={2021,2022,2023,2024,2025,2026H1},
  tick label style={font=\tiny},
  yticklabel style={font=\tiny},
  ylabel={亿元}, ylabel style={font=\tiny},
  ybar, bar width=4pt,
  grid=major, grid style={LightGray, line width=0.3pt},
  axis line style={InkGray, line width=0.5pt},
  enlarge x limits={abs=0.8},
  legend style={font=\scriptsize, at={(0.5,-0.26)}, anchor=north,
                legend columns=2, draw=none, fill=none, column sep=10pt},
]
\addplot[fill=AgriGreen!55, draw=AgriGreen, bar shift=-3.0pt] table[x=i, y=cash]{data/clean/profiles/fin_002548.dat};
\addplot[fill=SoftRed!45, draw=SoftRed, bar shift=3.0pt] table[x=i, y=ibd]{data/clean/profiles/fin_002548.dat};
\addplot[fill=AgriGreen!25, draw=AgriGreen, pattern=north east lines, bar shift=-3.0pt] table[x=x, y=cash]{data/clean/profiles/finh_002548.dat};
\addplot[fill=SoftRed!20, draw=SoftRed, pattern=north east lines, bar shift=3.0pt] table[x=x, y=ibd]{data/clean/profiles/finh_002548.dat};
\legend{货币资金（左轴）, 有息负债（左轴）}
\end{axis}
\begin{axis}[
  at={(finbot.south east)}, anchor=south east,
  width=13.75cm, height=3.10cm, scale only axis,
  axis y line*=right, axis x line*=none, xtick=\empty,
  ymin=-1.10, ymax=5.50,
  tick label style={font=\tiny},
  yticklabel style={font=\tiny}, scaled y ticks=false,
  legend style={font=\scriptsize, at={(0.5,-0.44)}, anchor=north,
                legend columns=1, draw=none, fill=none},
]
\addplot[OilBlue, mark=*, mark size=2.2pt, line width=1.3pt] table[x=i, y=ocf]{data/clean/profiles/fin_002548.dat};
\addplot[only marks, OilBlue, mark=*, mark size=3.2pt] table[x=x, y=ocf]{data/clean/profiles/finh_002548.dat};
\legend{经营活动现金流净额（右轴，亿元，蓝点）}
\end{axis}
\end{tikzpicture}

\captionof{figure}[金新农（002548）近年主要财务指标]{金新农（002548）近年主要财务指标（左上：营业收入与归母净利润；右上：ROE 与资产负债率；下：货币资金、有息负债为左轴柱状，经营活动现金流净额为其右轴的蓝点折线。
2021---2025 年为年报口径，2026H1 为 2026 年半年报/半年末，与其前一年报之间留出间隔、以斜纹或空心点区分；半年报数据未年化）}
\end{minipage}
\end{center}

\pfl{财务分析} 2026 年 6 月末资产负债率 81.5\%，较 2025 年末上升 5.8 个百分点（样本中位数 52.2\%，所处三级行业内按由低到高排第\zhnumber{11}位）。 2026 年 6 月末货币资金 2.36 亿元、短期债务（短期借款与一年内到期非流动负债）22.72 亿元、长期债务 6.20 亿元，有息负债合计 28.92 亿元，现金短债比 0.10 倍——缺口明显。 净有息负债 26.57 亿元，相当于其 2026 年上半年营业收入的 115\%。 流动比率 0.32、速动比率 0.18，短期偿债指标偏紧。 2026 年上半年经营活动现金流净额 0.50 亿元，净利润为负而经营现金流为正，差额主要来自折旧摊销与营运资本变动，说明亏损尚未完全转化为现金流出。 2026 年 6 月末在建工程 0.76 亿元，较上年末增加 0.10 亿元，仍有在建投入。 综合看，账面货币资金对有息负债与短期债务的覆盖不足，滚动偿付对新增融资的依赖度较高；资产负债率高于样本中位数 29.3 个百分点；有息负债中短期债务占比更高，期限结构仍以滚动续作为主。 公司披露的股东与质押风险为：2025 年，半年度末资产权利受限合计账面余额 10.90 亿元；2024 年度报告，长期股权投资 3.37 亿元质押用于融资借款。

\pfl{重大战略转型} 近三年公司的战略动作集中在：其他动作（2 项）：2025 年，推出 2025 年股票期权与限制性股票激励计划；2020 年，公司拟向 154 名激励对象授予 1,354.9932 万股限制性股票；并购与重组（1 项）：2016 年参股东进农牧 4\% 股权、收购福建一春 60\% 股权、增资赣州东进 20\% 股权。 公告口径上，近三年（2023 年 9 月---2026 年 9 月）公司披露重大重组与并购类公告 0 条、对外投资与转型类公告 4 条、回购与股权激励类公告 25 条，合计公告 433 条。 第一大行业“饲料加工”的收入占比由 2021 年年报的 47.6\% 变为 69.7\%（22.1 个百分点），收入结构出现再平衡。 综合判断，报告期内公司有 4 条与战略相关的公告落地（重大重组与并购 0 条、对外投资与转型 4 条），资本运作与主业调整之间存在直接关联。

\begin{center}\small\setlength{\tabcolsep}{4pt}
\captionof{table}[金新农（002548）历史融资与资本运作]{金新农（002548）历史融资与资本运作（据公司公告与公开报道整理；金额为披露值，含首次公开发行、再融资、债券、银行借款与重整投资等）}
\begin{adjustbox}{max width=\textwidth}
\begin{tabular}{@{}p{1.45cm}p{2.55cm}p{4.05cm}p{4.95cm}@{}}
\toprule
时间 & 方式 & 规模 & 用途与说明 \\
\midrule
2011-02 & IPO（深交所） & 发行 2,\allowbreak{}400 万股，\allowbreak{}发行价 24.00 元/\allowbreak{}股，\allowbreak{}募集资金总额 5.76 亿元 & 饲料产能建设 \\
2015 & 发行股份及支付现金购买资产并募集配套资金 & 以 5.25 亿元价格收购盈华讯方 80\% 股权；募集配套资金总额 3.67 亿元 & 支付现金对价（1.84 亿元）、\allowbreak{}补充流动资金、\allowbreak{}支付本次交易税费等；发行对象为金新农 2015 年第一期员工持股计划和陈俊海等 6 人；；2018 年 10 月经股东大会审议通过以 1.20 亿元有条件分阶段支付现金方式继续收购盈华讯方剩余股权 \\
2018-03 & 公开发行可转换公司债券 & 发行 650 万张，\allowbreak{}每张面值 100.00 元，\allowbreak{}发行总额 6.50 亿元（6.50 亿元） & 补充流动资金及项目建设；债券简称“金农转债”，\allowbreak{}代码 128036 \\
2020-12-30 & 非公开发行（定增） & 向广州湾区金农投资合伙企业（有限合伙）非公开发行 1.28 亿股，\allowbreak{}发行价 5.08 元/\allowbreak{}股 & 生猪养殖项目及补充流动资金；新增股份 2020 年 12 月 30 日在深交所上市；本次认购后收购人持有公司权益超过 30\% \\
2022 & 非公开发行（定增） & 向特定对象非公开发行 1.17 亿股，\allowbreak{}发行价 5.98 元/\allowbreak{}股，\allowbreak{}募集资金总额 7.00 亿元 & 募投项目建设及补充流动资金 \\
2023 & 子公司增资扩股引入投资者 & 广州乡村振兴投资合伙企业（有限合伙）对子公司广州金农增资 8,\allowbreak{}000.00 万元 & 发展生猪养殖业务；投资期为 3+2 年，\allowbreak{}投资期内每年需按投资金额的 6\% 支付投资收益；投资期 3 年结束后投资方有权要求公司回购股权 \\
\bottomrule
\end{tabular}
\end{adjustbox}
\end{center}

\pfl{历史融投资情况} 从现金流量表看，2025 年公司取得借款收到的现金 19.11 亿元、偿还债务支付的现金 18.77 亿元、吸收权益性投资 0.49 亿元，筹资活动现金流净额 -2.41 亿元（2023 年为取得借款 18.74 亿元、偿还债务 20.03 亿元）。 2026 年上半年筹资活动现金流净额为 -1.73 亿元（取得借款 6.61 亿元、偿还债务 7.64 亿元，未年化），已转为净流出，处于净还债状态。 2025 年分配股利、利润或偿付利息支付的现金合计 1.17 亿元。 公开披露的融资记录共 6 笔，工具类型以 IPO、发行股份及支付现金购、可转债、定向增发为主；金额与用途见下表。 其中规模最大的两笔为：2018 年 3 月，可转债，发行 650 万张，每张面值 100.00 元，发行总额 6.50 亿元（6.50 亿元）；2015 年，发行股份及支付现金购，以 5.25 亿元价格收购盈华讯方 80\% 股权，募集配套资金总额 3.67 亿元。 股本方面，近三年披露股本变动 12 次（可转债转股 9 次、限售股份上市 2 次、激励股份解禁 2 次、股份回购 1 次）；总股本由 2021 年末的 6.91 亿股变为 8.05 亿股（+16.56\%）。 分红方面，近三年未实施现金分红，在全部样本中属于分红能力偏弱的一端。 近三年“再融资”类公告 9 条，涉及定向增发、可转债。 综合看，公司融资结构上以债务滚动为主、股权融资为补充；2025 年筹资活动现金流净额 -2.41 亿元，融资节奏仍以维持存量债务周转为主，与前述经营与投资现金流的方向基本一致。

\pfl{未来潜在融资需求分析} 2026 年 6 月末货币资金 2.36 亿元，而短期债务 22.72 亿元（现金短债比 0.10 倍），2026 年上半年经营现金流净额 0.50 亿元。按“短期债务减去账面货币资金”的滚动偿付口径，静态缺口约 20.37 亿元；若再计入上半年亏损的年化额（7.77 亿元），维持一年正常经营所需的外部资金约 20.37---28.14 亿元。 按第 \ref{sec:financing-need} 节的三档划分，公司属于第一档“补血型”：2026 年上半年归母净利润 -3.88 亿元、6 月末资产负债率 81.5\%。负债率已超过 80\%，净资产被亏损侵蚀的程度较深，属于全部样本中资产负债表最紧张的一组。 可行的融资路径按现实性排序为：定向增发（需股价配合，且控股股东需放弃优先认购或同步增资）、债务重组或展期（与银行和债券持有人协商）、引入国资或产业投资人（部分公司以控制权让渡为对价）。 资金需求还要叠加在建工程：期末在建工程 0.76 亿元、2026 年上半年购建固定资产等支出 1.06 亿元（未年化），这部分支出在周期底部通常会被主动放缓。 综合看，公司的外部资金需求以债务接续为主，滚动偿付口径下的静态缺口约 20.37 亿元，融资的时点与规模更多取决于信用条件与资本市场窗口。




\subsubsection{邦基科技（603151）}
\label{co:603151}

\pfl{公司基本情况} 公司前身可追溯至 2007 年，山东淄博的猪饲料专业化企业，以猪预混料、浓缩料、配合料为主业，2022 年 10 月登陆上交所主板；2025—2026 年因向生猪养殖下游延伸并受猪价下行冲击，业绩由盈转亏。 公司于 2022-10-19 在上交所首发上市，注册资本 1.71 亿元，总股本 1.68 亿股。 截至 2026 年 10 月 12 日收盘价 18.30 元，对应总市值 30.7 亿元，近一年-26.8\%，总市值在三级行业“畜禽饲料”的 11 家公司中按由大到小排第\zhnumber{7}位，近一年涨跌幅低于全样本中位数（-10.4\%）16.4 个百分点。 按 2025 年年报披露的行业构成，收入占比前三位为饲料行业 96.1\%；畜牧养殖 2.3\%；其他(补充) 1.6\%。

\pfl{股权结构与控制权} 按 2026 年 6 月 30 日的股东名册，前三大股东为山东邦基集团有限公司 持股 51.21\%；淄博邦盈企业管理合伙企业(有限合伙) 持股 3.20\%；淄博邦智企业管理合伙企业(有限合伙) 持股 3.06\%，股东户数 10,878 户。 从控制权角度看，持股 51.21\%（8,772.12 万股）；公司 2007 年 4 月成立、2022 年 10 月 19 日在上海证券交易所上市。

\pfl{经营要点} 业务与产能层面，公司自 2007 年成立以来一直从事猪饲料的研发、生产和销售。

\pfl{所处行业的基本面} 饲料的供给即产能（门槛低、产能利用率偏低），需求直接来自下游存栏，利润则由玉米、豆粕等原料价格与成品售价的价差决定（第 \ref{sec:framework} 章的供需框架）。 从供需两端的当期数据看，2026 年上半年，合计产量占全国 65.2\%、提高 3.4 个百分点；2026 年上半年，全国工业饲料总产量 16,850 万吨、同比增长 3.9\%。

\pfl{竞争优势} 把该公司放回同类公司比较，2026 年上半年营业收入 26.8 亿元、在三级行业“畜禽饲料”的 11 家公司中排第\zhnumber{7}位；毛利率 5.9\%、高于行业内中位数 0.7 个百分点；半年 ROE -0.0\%（未年化）、高于行业内中位数（-6.25\%）；资产负债率 63.4\%、高于行业内中位数（62.1\%）。 相对优势方面，直接带动饲料产品销量增长、饲料主业收入同比提升。 优势在当期体现为高于行业中位数的资本回报，但优势能否持续仍取决于周期与自身产能投放节奏。

\pfl{历史股价} 本报告样本区间（2022 年 10 月 至 2026 年 10 月）内，股价由 17.82 元变为 18.30 元，累计 +2.7\%，区间最高 25.74 元（2025 年 10 月）、最低 9.08 元（2024 年 8 月）；当前价相当于区间最高价的 71.1\%，为区间最低价的 2.02 倍。 月度收益的年化波动率 45.0\%，高于全样本中位数 37.3\%，在三级行业“畜禽饲料”的 11 家公司中波动率由高到低排第\zhnumber{3}位。 把股价阶段与公司事件对齐看，2024 年 8 月 至 2025 年 10 月股价上涨+183.5\%，同期（2025 年 6 月）发行股份及支付现金购买资产暨关联交易预案；2026 年 6 月 至 2026 年 10 月股价上涨+66.5\%，同期（2026 年）归属于上市公司股东的净利润 -47.97 万元（同比 -100.74\%）。 把这些阶段与公司事件对齐后可以看到，几轮大涨大跌多由公司自身的资本运作与风险事件触发，行业景气只解释了其中一部分。

\begin{center}
\begin{minipage}{\textwidth}\centering
\begin{tikzpicture}
\begin{axis}[
  width=12.6cm, height=3.9cm, scale only axis,
  xmin=2022.672, xmax=2026.972, ymin=8.44, ymax=27.54,
  xtick={2022,2023,2024,2025,2026},
  yearticks,
  yticklabel style={font=\scriptsize},
  ylabel={元/股}, ylabel style={font=\scriptsize},
  grid=major, grid style={LightGray, line width=0.3pt},
  axis line style={InkGray, line width=0.5pt},
  every axis plot/.append style={line width=0.9pt},
]
\addplot[AgriGreen] table[x=x,y=y]{data/clean/profiles/px_603151.dat};
\addplot[SoftRed, dashed, line width=0.7pt] table[x=x,y=y]{data/clean/profiles/pxzz_603151.dat};
\addplot[only marks, mark=*, mark size=1.1pt, SoftRed] table[x=x,y=y]{data/clean/profiles/pxzz_603151.dat};
\end{axis}
\end{tikzpicture}

\captionof{figure}[邦基科技（603151）股价与周期骨架]{邦基科技（603151）月度收盘价（元/股）与 ZigZag 周期骨架（阈值 25\%，圆点为周期转折点）}
\end{minipage}
\end{center}

\pfl{过去周期分析} 以 25\% 的 ZigZag 阈值划分，样本区间内股价形成 3 个主升段与 3 个主跌段。 最大主升段出现在 2024 年 8 月 至 2025 年 10 月，14 个月+183.5\%；最大主跌段为 2025 年 10 月 至 2026 年 6 月，8 个月-57.3\%。 最近一次转折出现在 2026 年 10 月（上涨段自 2026 年 6 月起，4 个月+66.5\%），当前处于该段之后的震荡之中。 公司股价较申万农林牧渔指数提前约 21 个月见底，是板块中较早定价周期反转的公司之一。 饲料夹在原料成本与养殖需求之间（第 \ref{sec:feed} 节），量的部分取决于存栏、利的部分取决于原料与成品的价差，公司 2026 年 6 月末资产负债率 63.4\%、半年 ROE -0.0\%（未年化），在本轮下行中属于随行业同步调整的一类。 综合区间形态看，该股单段涨跌幅度偏大、以趋势段为主（最大主升 +183.5\%、最大主跌 -57.3\%），年化波动率 45.0\% 与全样本中位数 37.3\% 相比波动明显大于样本中位数。

\pfl{盈利情况} 2026 年上半年营业收入 26.75 亿元（同比 +11.8\%），归母净利润 -0.00 亿元，由上年同期的盈利转为亏损。 以年报为趋势对照，2023---2025 年营业收入由 16.47 亿元变为 56.99 亿元（两年年均复合增速 +86.0\%），归母净利润由 0.84 亿元变为 1.11 亿元。 按半年报口径，2026 年上半年的 ROE 为 -0.0\%（未年化）、销售净利率 -1.3\%、毛利率 5.9\%，ROE 在“畜禽饲料”11 家公司中按数值由高到低排第\zhnumber{3}位。 同期全样本半年 ROE 中位数 0.75\%，公司与之相差 0.78 个百分点（弱于中位数公司）。 综合看，公司仍处亏损区间、由盈转亏，半年 ROE 在三级行业内按由高到低排第\zhnumber{3}位，单位盈利能力的修复依赖行业景气与自身成本端的改善，报表弹性主要体现为价格与销量的方向性变化。

\begin{center}\small\setlength{\tabcolsep}{3.4pt}
\captionof{table}[邦基科技（603151）关键财务指标]{邦基科技（603151）关键财务指标（合并报表口径；两个半年报期均未年化；现金短债比 $=$ 货币资金 $\div$ 短期债务）}
\begin{adjustbox}{max width=\textwidth}
\begin{tabular}{@{}lrrrrr@{}}
\toprule
指标 & 2023 年 & 2024 年 & 2025 年 & 2025 年半年报 & 2026 年半年报 \\
\midrule
营业收入（亿元） & 16.47 & 25.42 & 56.99 & 23.93 & 26.75 \\
归母净利润（亿元） & 0.84 & 0.50 & 1.11 & 0.64 & -0.00 \\
毛利率（\%） & 13.0 & 9.2 & 7.6 & 8.7 & 5.9 \\
ROE（\%） & 6.9 & 4.1 & 8.5 & 4.9 & -0.0 \\
资产负债率（\%） & 15.9 & 28.8 & 55.9 & 50.7 & 63.4 \\
经营活动现金流净额（亿元） & -0.54 & -1.10 & -6.01 & -5.48 & 1.44 \\
货币资金（亿元） & 3.91 & 1.68 & 3.40 & 2.53 & 7.22 \\
有息负债合计（亿元） & 1.03 & 2.46 & 11.91 & 8.42 & 15.02 \\
现金短债比（倍） & 4.26 & 0.72 & 0.32 & 0.31 & 0.55 \\
\bottomrule
\end{tabular}
\end{adjustbox}
\end{center}

\begin{center}
\begin{minipage}{\textwidth}\centering
\begin{tikzpicture}
\begin{groupplot}[
  group style={group size=2 by 1, horizontal sep=0.60cm},
  width=6.85cm, height=4.60cm,
  tick label style={font=\tiny},
  title style={font=\scriptsize\heiti},
  yticklabel style={font=\tiny},
  ylabel style={font=\tiny},
  grid=major, grid style={LightGray, line width=0.3pt},
  axis line style={InkGray, line width=0.5pt},
  enlarge x limits={abs=0.8},
  legend style={font=\scriptsize, at={(0.5,-0.20)}, anchor=north,
                legend columns=2, draw=none, fill=none, column sep=6pt},
]
\nextgroupplot[
  ybar, bar width=4pt,
  ymin=-5.70, ymax=63.83,
  xtick={0,1,2,3,4,5.7},
  xticklabels={2021,2022,2023,2024,2025,2026H1},
  ylabel={亿元},
]
\addplot[fill=AgriGreen!75, draw=AgriGreen, bar shift=-2.6pt] table[x=i, y=rev]{data/clean/profiles/fin_603151.dat};
\addplot[fill=SoftRed!60, draw=SoftRed, bar shift=2.6pt] table[x=i, y=ni]{data/clean/profiles/fin_603151.dat};
\addplot[fill=AgriGreen!30, draw=AgriGreen, pattern=north east lines, bar shift=-2.6pt] table[x=x, y=rev]{data/clean/profiles/finh_603151.dat};
\addplot[fill=SoftRed!20, draw=SoftRed, pattern=north east lines, bar shift=2.6pt] table[x=x, y=ni]{data/clean/profiles/finh_603151.dat};
\legend{营业收入（年/半年）, 归母净利润（年/半年）}
\nextgroupplot[
  ymin=-5.1, ymax=69.7,
  xtick={0,1,2,3,4,5.7},
  xticklabels={2021,2022,2023,2024,2025,2026H1},
  ylabel={\%},
]
\addplot[AgriGold, mark=*, mark size=1.3pt] table[x=i, y=roe]{data/clean/profiles/fin_603151.dat};
\addplot[OilBlue, mark=square*, mark size=1.3pt] table[x=i, y=debt]{data/clean/profiles/fin_603151.dat};
\addplot[only marks, AgriGold, mark=*, mark size=2.0pt] table[x=x, y=roe]{data/clean/profiles/finh_603151.dat};
\addplot[only marks, OilBlue, mark=square*, mark size=2.0pt] table[x=x, y=debt]{data/clean/profiles/finh_603151.dat};
\legend{ROE, 资产负债率}
\end{groupplot}
\end{tikzpicture}
\begin{tikzpicture}
\begin{axis}[
  name=finbot,
  width=13.75cm, height=3.10cm, scale only axis,
  ymin=-1.50, ymax=16.82,
  xtick={0,1,2,3,4,5.7},
  xticklabels={2021,2022,2023,2024,2025,2026H1},
  tick label style={font=\tiny},
  yticklabel style={font=\tiny},
  ylabel={亿元}, ylabel style={font=\tiny},
  ybar, bar width=4pt,
  grid=major, grid style={LightGray, line width=0.3pt},
  axis line style={InkGray, line width=0.5pt},
  enlarge x limits={abs=0.8},
  legend style={font=\scriptsize, at={(0.5,-0.26)}, anchor=north,
                legend columns=2, draw=none, fill=none, column sep=10pt},
]
\addplot[fill=AgriGreen!55, draw=AgriGreen, bar shift=-3.0pt] table[x=i, y=cash]{data/clean/profiles/fin_603151.dat};
\addplot[fill=SoftRed!45, draw=SoftRed, bar shift=3.0pt] table[x=i, y=ibd]{data/clean/profiles/fin_603151.dat};
\addplot[fill=AgriGreen!25, draw=AgriGreen, pattern=north east lines, bar shift=-3.0pt] table[x=x, y=cash]{data/clean/profiles/finh_603151.dat};
\addplot[fill=SoftRed!20, draw=SoftRed, pattern=north east lines, bar shift=3.0pt] table[x=x, y=ibd]{data/clean/profiles/finh_603151.dat};
\legend{货币资金（左轴）, 有息负债（左轴）}
\end{axis}
\begin{axis}[
  at={(finbot.south east)}, anchor=south east,
  width=13.75cm, height=3.10cm, scale only axis,
  axis y line*=right, axis x line*=none, xtick=\empty,
  ymin=-7.94, ymax=3.67,
  tick label style={font=\tiny},
  yticklabel style={font=\tiny}, scaled y ticks=false,
  legend style={font=\scriptsize, at={(0.5,-0.44)}, anchor=north,
                legend columns=1, draw=none, fill=none},
]
\addplot[OilBlue, mark=*, mark size=2.2pt, line width=1.3pt] table[x=i, y=ocf]{data/clean/profiles/fin_603151.dat};
\addplot[only marks, OilBlue, mark=*, mark size=3.2pt] table[x=x, y=ocf]{data/clean/profiles/finh_603151.dat};
\legend{经营活动现金流净额（右轴，亿元，蓝点）}
\end{axis}
\end{tikzpicture}

\captionof{figure}[邦基科技（603151）近年主要财务指标]{邦基科技（603151）近年主要财务指标（左上：营业收入与归母净利润；右上：ROE 与资产负债率；下：货币资金、有息负债为左轴柱状，经营活动现金流净额为其右轴的蓝点折线。
2021---2025 年为年报口径，2026H1 为 2026 年半年报/半年末，与其前一年报之间留出间隔、以斜纹或空心点区分；半年报数据未年化）}
\end{minipage}
\end{center}

\pfl{财务分析} 2026 年 6 月末资产负债率 63.4\%，较 2025 年末上升 7.5 个百分点（样本中位数 52.2\%，所处三级行业内按由低到高排第\zhnumber{7}位）。 2026 年 6 月末货币资金 7.22 亿元、短期债务（短期借款与一年内到期非流动负债）13.06 亿元、长期债务 1.96 亿元，有息负债合计 15.02 亿元，现金短债比 0.55 倍——覆盖偏紧。 净有息负债 7.79 亿元，相当于其 2026 年上半年营业收入的 29\%。 流动比率 1.28、速动比率 1.01，短期偿债能力处于正常区间。 2026 年上半年经营活动现金流净额 1.44 亿元，净利润为负而经营现金流为正，差额主要来自折旧摊销与营运资本变动，说明亏损尚未完全转化为现金流出。 综合看，现金短债比 0.55 倍，短期资金安排偏紧；资产负债率高于样本中位数 11.2 个百分点；有息负债中短期债务占比更高，期限结构仍以滚动续作为主。

\pfl{重大战略转型} 公司的战略动作可归为以下几类：其他动作（2 项）：2024 年 6 月，推出 2024 年股票期权激励计划，以 12.85 元/股向 154 名激励对象授予 875.00 万份股票期权；并购与重组（2 项）：2025 年 11 月，公司终止上述重大资产重组；2025 年 6 月，发行股份及支付现金购买资产暨关联交易预案。 公告口径上，近三年（2023 年 9 月---2026 年 9 月）公司披露重大重组与并购类公告 16 条、对外投资与转型类公告 4 条、回购与股权激励类公告 0 条，合计公告 407 条。 第一大行业仍是“饲料行业”，收入占比 96.1\%，与 2021 年年报相比没有方向性变化。 综合判断，报告期内公司有 20 条与战略相关的公告落地（重大重组与并购 16 条、对外投资与转型 4 条），资本运作与主业调整之间存在直接关联。

\begin{center}\small\setlength{\tabcolsep}{4pt}
\captionof{table}[邦基科技（603151）历史融资与资本运作]{邦基科技（603151）历史融资与资本运作（据公司公告与公开报道整理；金额为披露值，含首次公开发行、再融资、债券、银行借款与重整投资等）}
\begin{adjustbox}{max width=\textwidth}
\begin{tabular}{@{}p{1.45cm}p{2.55cm}p{4.05cm}p{4.95cm}@{}}
\toprule
时间 & 方式 & 规模 & 用途与说明 \\
\midrule
2022-10 & IPO（上交所主板） & 公开发行 4,\allowbreak{}200.00 万股，\allowbreak{}发行价 17.95 元/\allowbreak{}股，\allowbreak{}募集资金总额 7.54 亿元 & 投向公司及辽宁、\allowbreak{}长春、\allowbreak{}云南、\allowbreak{}张家口、\allowbreak{}山西等地子公司的高档配合猪饲料、\allowbreak{}浓缩饲料智能生产车间等项目；募集资金 2022 年 10 月 14 日到账；项目总投资 6.90 亿元 \\
2024-06-06 & 股票期权激励（授予） & 以 12.85 元/\allowbreak{}份授予 154 名激励对象 875.00 万份股票期权 & 员工长期激励；2025 年 5 月 13 日首次授予第一个行权期可行权数量 332.00 万股（140 人） \\
2025-06-16 & 发行股份及支付现金购买资产（重大资产重组，后终止） & 拟收购山东北溪农牧、\allowbreak{}山东瑞东伟力农牧、\allowbreak{}山东鑫牧农牧、\allowbreak{}瑞东农牧（利津）、\allowbreak{}瑞东农牧（山东）、\allowbreak{}瑞东威力牧业（滨州）各 100\% 股权及派斯东畜牧技术咨询（上海）80\% 股权 & 由单一猪饲料业务向下游养殖业延伸；交易对方为 Riverstone Farm Pte. Ltd \\
\bottomrule
\end{tabular}
\end{adjustbox}
\end{center}

\pfl{历史融投资情况} 从现金流量表看，2025 年公司取得借款收到的现金 11.95 亿元、偿还债务支付的现金 4.19 亿元、吸收权益性投资 1.42 亿元，筹资活动现金流净额 8.98 亿元（2023 年为取得借款 0.94 亿元、偿还债务 0.50 亿元）。 2026 年上半年筹资活动现金流净额为 2.83 亿元（取得借款 8.68 亿元、偿还债务 6.63 亿元，未年化），仍处于净融入状态。 2025 年分配股利、利润或偿付利息支付的现金合计 0.58 亿元。 公开披露的融资记录共 3 笔，工具类型以 IPO、发行股份及支付现金购为主；金额与用途见下表。 股本方面，近三年披露股本变动 4 次（限售股份上市 1 次）；总股本自 2021 年末以来未发生变化（1.68 亿股）。 分红方面，近三年现金分红 6 次、合计每股派现 1.05 元，最近一次实施在 2026 年 6 月。 近三年“再融资”类公告 4 条，涉及定向增发。 综合看，公司融资结构上以债务滚动为主、股权融资为补充；2025 年筹资活动现金流净额 8.98 亿元，年度内仍为净融入，与前述经营与投资现金流的方向基本一致。

\pfl{未来潜在融资需求分析} 2026 年 6 月末货币资金 7.22 亿元，而短期债务 13.06 亿元（现金短债比 0.55 倍），2026 年上半年经营现金流净额 1.44 亿元。按“短期债务减去账面货币资金”的滚动偿付口径，静态缺口约 5.83 亿元；上半年盈利或接近盈亏平衡，该缺口不随经营亏损进一步扩大。 按第 \ref{sec:financing-need} 节的三档划分，公司属于第二档“保壳型”：近三年重大重组与并购类公告 16 条，2026 年上半年归母净利润 -0.00 亿元，融资需求更多体现为“资产置换 + 维持上市地位”，而不是扩产。 公开披露的后续资金安排线索包括：债务与授信安排：2025 年年度报告，公司披露公司债券（含企业债券）和非金融企业债务融资工具、可转换公司债券均不适用。 综合看，公司的外部资金需求以补充流动性与项目投入为主，滚动偿付口径下的静态缺口约 5.83 亿元，融资节奏将受制于银行续贷意愿与股权融资窗口。




\subsubsection{百洋股份（002696）}
\label{co:002696}

\pfl{公司基本情况} 2000 年设立至今，百洋股份是广西南宁的水产企业，以罗非鱼等水产品加工出口起家，2012 年上市；2017 年以 9.74 亿元收购火星时代跨界数字艺术教育、2019 年原价回购剥离，2020 年 5 月控制权转让给青岛国资，现回归水产与饲料主业，2026 年上半年重新陷入亏损。 公司于 2012-09-05 在深交所主板首发上市，注册资本 3.46 亿元，总股本 3.46 亿股。 截至 2026 年 10 月 12 日收盘价 8.05 元，对应总市值 27.9 亿元，近一年+29.4\%，总市值在三级行业“水产饲料”的 4 家公司中按由大到小排第\zhnumber{4}位，近一年涨跌幅高于全样本中位数（-10.4\%）39.8 个百分点。 按 2026 年半年报披露的行业构成，收入占比前三位为饲料行业 51.8\%；食品加工行业 38.2\%；其他 9.9\%。 股权结构方面，实际控制人为青岛市国资委。

\pfl{股权结构与控制权} 按 2026 年 6 月 30 日的股东名册，前三大股东为青岛海洋发展集团有限公司（持股比例未披露）；青岛国信创新股权投资管理有限公司－青岛市海洋新动能产（持股比例未披露）；浙江宁聚投资管理有限公司－宁聚开阳 13 号私募证券投资 持股 5.31\%，股东户数 13,685 户。

\pfl{经营要点} 业务与产能层面，2026 年上半年，收入增长由饲料及饲料原料业务拉动：饲料业务营收 8.37 亿元（同比 +34.33\%）。

\pfl{所处行业的基本面} 水产饲料的需求由投苗量与存塘鱼价决定，是饲料中毛利率最高、也最依赖服务的品类（第 \ref{sec:framework} 章的供需框架）。

\pfl{竞争优势} 把该公司放回同类公司比较，2026 年上半年营业收入 20.5 亿元、在三级行业“水产饲料”的 4 家公司中排第\zhnumber{4}位；毛利率 6.2\%、低于行业内中位数 2.6 个百分点；半年 ROE -6.2\%（未年化）、低于行业内中位数（1.22\%）；资产负债率 66.4\%、高于行业内中位数（60.5\%）。 相对优势方面，国有控股股东通过委托贷款、统借统贷等方式帮助公司节约资金成本。 就当期报表看，公司尚未把经营层面的特征转化为高于行业的资本回报，周期底部的相对优势更多体现在现金流与资产负债表上。

\pfl{历史股价} 以 2021 年 1 月为起点、2026 年 10 月为终点，股价由 5.44 元变为 8.05 元，累计 +48.0\%，区间最高 8.05 元（2026 年 10 月）、最低 3.88 元（2024 年 8 月）；当前价相当于区间最高价的 100.0\%，为区间最低价的 2.07 倍。 月度收益的年化波动率 33.0\%，低于全样本中位数 37.3\%，在三级行业“水产饲料”的 4 家公司中波动率由高到低排第\zhnumber{2}位。

\begin{center}
\begin{minipage}{\textwidth}\centering
\begin{tikzpicture}
\begin{axis}[
  width=12.6cm, height=3.9cm, scale only axis,
  xmin=2020.922, xmax=2026.972, ymin=3.61, ymax=8.61,
  xtick={2021,2022,2023,2024,2025,2026},
  yearticks,
  yticklabel style={font=\scriptsize},
  ylabel={元/股}, ylabel style={font=\scriptsize},
  grid=major, grid style={LightGray, line width=0.3pt},
  axis line style={InkGray, line width=0.5pt},
  every axis plot/.append style={line width=0.9pt},
]
\addplot[AgriGreen] table[x=x,y=y]{data/clean/profiles/px_002696.dat};
\addplot[SoftRed, dashed, line width=0.7pt] table[x=x,y=y]{data/clean/profiles/pxzz_002696.dat};
\addplot[only marks, mark=*, mark size=1.1pt, SoftRed] table[x=x,y=y]{data/clean/profiles/pxzz_002696.dat};
\end{axis}
\end{tikzpicture}

\captionof{figure}[百洋股份（002696）股价与周期骨架]{百洋股份（002696）月度收盘价（元/股）与 ZigZag 周期骨架（阈值 25\%，圆点为周期转折点）}
\end{minipage}
\end{center}

\pfl{过去周期分析} 以 25\% 的 ZigZag 阈值划分，样本区间内股价形成 3 个主升段与 3 个主跌段。 最大主升段出现在 2024 年 8 月 至 2026 年 2 月，18 个月+102.6\%；最大主跌段为 2022 年 12 月 至 2024 年 8 月，20 个月-48.5\%。 最近一次转折出现在 2026 年 10 月（上涨段自 2026 年 6 月起，4 个月+62.0\%），当前处于该段之后的震荡之中。 公司股价较申万农林牧渔指数提前约 21 个月见底，是板块中较早定价周期反转的公司之一。 水产饲料的需求取决于投苗量与存塘鱼价，第 \ref{sec:other-livestock} 章显示 2026 年淡水鱼价格与投苗量已率先回升，其需求驱动与生猪存栏的联系弱于畜禽饲料，公司 2026 年 6 月末资产负债率 66.4\%，高于全样本中位数 52.2\%，其股价因此更多反映资金面与信用风险，而不是价格弹性本身。 综合区间形态看，该股单段涨跌幅度偏大、以趋势段为主（最大主升 +102.6\%、最大主跌 -48.5\%），年化波动率 33.0\% 与全样本中位数 37.3\% 相比波动小于样本中位数。

\pfl{盈利情况} 2026 年上半年营业收入 20.51 亿元（同比 +21.1\%），归母净利润 -0.83 亿元，由上年同期的盈利转为亏损。 以年报为趋势对照，2023---2025 年营业收入由 27.20 亿元变为 40.55 亿元（两年年均复合增速 +22.1\%），归母净利润由 -0.28 亿元变为 0.41 亿元。 按半年报口径，2026 年上半年的 ROE 为 -6.2\%（未年化）、销售净利率 -4.2\%、毛利率 6.2\%，ROE 在“水产饲料”4 家公司中按数值由高到低排第\zhnumber{4}位。 按同一口径，三级行业“水产饲料”4 家公司的半年 ROE 中位数为 1.22\%，该公司低于其中位数 7.42 个百分点。 综合看，公司仍处亏损区间、由盈转亏，半年 ROE 在三级行业内按由高到低排第\zhnumber{4}位，单位盈利能力的修复依赖行业景气与自身成本端的改善，报表弹性主要体现为价格与销量的方向性变化。

\begin{center}\small\setlength{\tabcolsep}{3.4pt}
\captionof{table}[百洋股份（002696）关键财务指标]{百洋股份（002696）关键财务指标（合并报表口径；两个半年报期均未年化；现金短债比 $=$ 货币资金 $\div$ 短期债务）}
\begin{adjustbox}{max width=\textwidth}
\begin{tabular}{@{}lrrrrr@{}}
\toprule
指标 & 2023 年 & 2024 年 & 2025 年 & 2025 年半年报 & 2026 年半年报 \\
\midrule
营业收入（亿元） & 27.20 & 30.78 & 40.55 & 16.94 & 20.51 \\
归母净利润（亿元） & -0.28 & -0.15 & 0.41 & 0.12 & -0.83 \\
毛利率（\%） & 8.3 & 9.0 & 8.8 & 8.8 & 6.2 \\
ROE（\%） & -1.9 & -1.1 & 2.9 & 0.9 & -6.2 \\
资产负债率（\%） & 49.8 & 53.5 & 60.9 & 60.0 & 66.4 \\
经营活动现金流净额（亿元） & 2.40 & -0.60 & -3.15 & -1.99 & -3.30 \\
货币资金（亿元） & 6.20 & 4.55 & 4.09 & 6.65 & 5.32 \\
有息负债合计（亿元） & 8.10 & 6.35 & 12.19 & 11.48 & 16.70 \\
现金短债比（倍） & 1.08 & 0.84 & 0.46 & 0.78 & 0.35 \\
\bottomrule
\end{tabular}
\end{adjustbox}
\end{center}

\begin{center}
\begin{minipage}{\textwidth}\centering
\begin{tikzpicture}
\begin{groupplot}[
  group style={group size=2 by 1, horizontal sep=0.60cm},
  width=6.85cm, height=4.60cm,
  tick label style={font=\tiny},
  title style={font=\scriptsize\heiti},
  yticklabel style={font=\tiny},
  ylabel style={font=\tiny},
  grid=major, grid style={LightGray, line width=0.3pt},
  axis line style={InkGray, line width=0.5pt},
  enlarge x limits={abs=0.8},
  legend style={font=\scriptsize, at={(0.5,-0.20)}, anchor=north,
                legend columns=2, draw=none, fill=none, column sep=6pt},
]
\nextgroupplot[
  ybar, bar width=4pt,
  ymin=-4.96, ymax=45.52,
  xtick={0,1,2,3,4,5.7},
  xticklabels={2021,2022,2023,2024,2025,2026H1},
  ylabel={亿元},
]
\addplot[fill=AgriGreen!75, draw=AgriGreen, bar shift=-2.6pt] table[x=i, y=rev]{data/clean/profiles/fin_002696.dat};
\addplot[fill=SoftRed!60, draw=SoftRed, bar shift=2.6pt] table[x=i, y=ni]{data/clean/profiles/fin_002696.dat};
\addplot[fill=AgriGreen!30, draw=AgriGreen, pattern=north east lines, bar shift=-2.6pt] table[x=x, y=rev]{data/clean/profiles/finh_002696.dat};
\addplot[fill=SoftRed!20, draw=SoftRed, pattern=north east lines, bar shift=2.6pt] table[x=x, y=ni]{data/clean/profiles/finh_002696.dat};
\legend{营业收入（年/半年）, 归母净利润（年/半年）}
\nextgroupplot[
  ymin=-12.0, ymax=73.7,
  xtick={0,1,2,3,4,5.7},
  xticklabels={2021,2022,2023,2024,2025,2026H1},
  ylabel={\%},
]
\addplot[AgriGold, mark=*, mark size=1.3pt] table[x=i, y=roe]{data/clean/profiles/fin_002696.dat};
\addplot[OilBlue, mark=square*, mark size=1.3pt] table[x=i, y=debt]{data/clean/profiles/fin_002696.dat};
\addplot[only marks, AgriGold, mark=*, mark size=2.0pt] table[x=x, y=roe]{data/clean/profiles/finh_002696.dat};
\addplot[only marks, OilBlue, mark=square*, mark size=2.0pt] table[x=x, y=debt]{data/clean/profiles/finh_002696.dat};
\legend{ROE, 资产负债率}
\end{groupplot}
\end{tikzpicture}
\begin{tikzpicture}
\begin{axis}[
  name=finbot,
  width=13.75cm, height=3.10cm, scale only axis,
  ymin=-1.67, ymax=18.71,
  xtick={0,1,2,3,4,5.7},
  xticklabels={2021,2022,2023,2024,2025,2026H1},
  tick label style={font=\tiny},
  yticklabel style={font=\tiny},
  ylabel={亿元}, ylabel style={font=\tiny},
  ybar, bar width=4pt,
  grid=major, grid style={LightGray, line width=0.3pt},
  axis line style={InkGray, line width=0.5pt},
  enlarge x limits={abs=0.8},
  legend style={font=\scriptsize, at={(0.5,-0.26)}, anchor=north,
                legend columns=2, draw=none, fill=none, column sep=10pt},
]
\addplot[fill=AgriGreen!55, draw=AgriGreen, bar shift=-3.0pt] table[x=i, y=cash]{data/clean/profiles/fin_002696.dat};
\addplot[fill=SoftRed!45, draw=SoftRed, bar shift=3.0pt] table[x=i, y=ibd]{data/clean/profiles/fin_002696.dat};
\addplot[fill=AgriGreen!25, draw=AgriGreen, pattern=north east lines, bar shift=-3.0pt] table[x=x, y=cash]{data/clean/profiles/finh_002696.dat};
\addplot[fill=SoftRed!20, draw=SoftRed, pattern=north east lines, bar shift=3.0pt] table[x=x, y=ibd]{data/clean/profiles/finh_002696.dat};
\legend{货币资金（左轴）, 有息负债（左轴）}
\end{axis}
\begin{axis}[
  at={(finbot.south east)}, anchor=south east,
  width=13.75cm, height=3.10cm, scale only axis,
  axis y line*=right, axis x line*=none, xtick=\empty,
  ymin=-4.78, ymax=4.11,
  tick label style={font=\tiny},
  yticklabel style={font=\tiny}, scaled y ticks=false,
  legend style={font=\scriptsize, at={(0.5,-0.44)}, anchor=north,
                legend columns=1, draw=none, fill=none},
]
\addplot[OilBlue, mark=*, mark size=2.2pt, line width=1.3pt] table[x=i, y=ocf]{data/clean/profiles/fin_002696.dat};
\addplot[only marks, OilBlue, mark=*, mark size=3.2pt] table[x=x, y=ocf]{data/clean/profiles/finh_002696.dat};
\legend{经营活动现金流净额（右轴，亿元，蓝点）}
\end{axis}
\end{tikzpicture}

\captionof{figure}[百洋股份（002696）近年主要财务指标]{百洋股份（002696）近年主要财务指标（左上：营业收入与归母净利润；右上：ROE 与资产负债率；下：货币资金、有息负债为左轴柱状，经营活动现金流净额为其右轴的蓝点折线。
2021---2025 年为年报口径，2026H1 为 2026 年半年报/半年末，与其前一年报之间留出间隔、以斜纹或空心点区分；半年报数据未年化）}
\end{minipage}
\end{center}

\pfl{财务分析} 2026 年 6 月末资产负债率 66.4\%，较 2025 年末上升 5.5 个百分点（样本中位数 52.2\%，所处三级行业内按由低到高排第\zhnumber{3}位）。 2026 年 6 月末货币资金 5.32 亿元、短期债务（短期借款与一年内到期非流动负债）15.11 亿元、长期债务 1.59 亿元，有息负债合计 16.70 亿元，现金短债比 0.35 倍——覆盖偏紧。 净有息负债 11.38 亿元，相当于其 2026 年上半年营业收入的 55\%。 流动比率 1.28、速动比率 0.71，短期流动性无明显缺口。 2026 年上半年经营活动现金流净额 -3.30 亿元，经营现金流与利润同时为负，日常运营对债务与股东投入的依赖度较高。 2026 年 6 月末在建工程 1.18 亿元，较上年末增加 0.64 亿元，仍有在建投入。 综合看，账面货币资金对有息负债与短期债务的覆盖不足，滚动偿付对新增融资的依赖度较高；资产负债率高于样本中位数 14.2 个百分点；有息负债中短期债务占比更高，期限结构仍以滚动续作为主。

\pfl{重大战略转型} 近三年公司的战略动作集中在：并购与重组（2 项）：2019 年，教育文化业务大幅下滑后剥离火星时代；2017 年，当年以 9.74 亿元收购火星时代 100\% 股权。 公告口径上，近三年（2023 年 9 月---2026 年 9 月）公司披露重大重组与并购类公告 0 条、对外投资与转型类公告 1 条、回购与股权激励类公告 39 条，合计公告 337 条。 第一大行业仍是“饲料行业”，收入占比 51.8\%，与 2021 年年报相比没有方向性变化。 综合判断，报告期内公司有 1 条与战略相关的公告落地（重大重组与并购 0 条、对外投资与转型 1 条），资本运作与主业调整之间存在直接关联。

\begin{center}\small\setlength{\tabcolsep}{4pt}
\captionof{table}[百洋股份（002696）历史融资与资本运作]{百洋股份（002696）历史融资与资本运作（据公司公告与公开报道整理；金额为披露值，含首次公开发行、再融资、债券、银行借款与重整投资等）}
\begin{adjustbox}{max width=\textwidth}
\begin{tabular}{@{}p{1.45cm}p{2.55cm}p{4.05cm}p{4.95cm}@{}}
\toprule
时间 & 方式 & 规模 & 用途与说明 \\
\midrule
2012-09 & IPO（深交所） & 公开发行 2,\allowbreak{}200 万股，\allowbreak{}发行价 23.90 元/\allowbreak{}股 & 水产品加工产能建设；发行时公司名称为百洋水产集团股份有限公司 \\
2017-09 & 非公开发行（配套融资） & 向孙忠义等 7 名特定投资者非公开发行 2,\allowbreak{}952.17 万股，\allowbreak{}发行价 18.19 元/\allowbreak{}股 & 支付收购火星时代的现金对价等；新增股份上市日为 2017 年 9 月 27 日 \\
2020-03-24 & 控制权协议转让 & 孙忠义、\allowbreak{}蔡晶将其持有的 1.04 亿股（占公司股份总数 29.90\%）以 9.45 元/\allowbreak{}股协议转让给青岛海洋创新产业投资基金有限公司、\allowbreak{}青岛市海洋新动能产业投资基金（有限合伙） & 公司控制权变更；2020 年 5 月 29 日完成股份交割过户 \\
2024-2025 & 限制性股票激励计划 & 2025 年 7 月 31 日为预留部分授予日，\allowbreak{}以 3.14 元/\allowbreak{}股向 21 名激励对象授予 98.80 万股限制性股票 & 核心团队长期激励；2024 年限制性股票激励计划的预留授予实施 \\
2025 & 股份回购 & 2024 年度支付回购库存股款项 5,\allowbreak{}258.90 万元 & 回购公司股份；2025 年度利润分配以 2025 年 12 月 31 日总股本扣除回购专户库存股后 3.46 亿股为基数 \\
\bottomrule
\end{tabular}
\end{adjustbox}
\end{center}

\pfl{历史融投资情况} 从现金流量表看，2025 年公司取得借款收到的现金 11.82 亿元、偿还债务支付的现金 7.31 亿元、吸收权益性投资 0.22 亿元，筹资活动现金流净额 4.73 亿元（2023 年为取得借款 8.46 亿元、偿还债务 4.82 亿元）。 2026 年上半年筹资活动现金流净额为 5.26 亿元（取得借款 9.31 亿元、偿还债务 4.73 亿元，未年化），仍处于净融入状态。 2025 年分配股利、利润或偿付利息支付的现金合计 0.41 亿元。 公开披露的融资记录共 5 笔，工具类型以 IPO、定向增发为主；金额与用途见下表。 其中规模最大的两笔为：2020 年 3 月，控制权协议转让，孙忠义、蔡晶将其持有的 1.04 亿股（占公司股份总数 29.90\%）以 9.45 元/股协议转让给青岛海洋创新产业投资基金有限公司、青岛市海洋新动能产业投资基金（有限合伙）。 股本方面，近三年披露股本变动 5 次（股份回购 1 次、其他 1 次）；总股本由 2021 年末的 3.49 亿股变为 3.46 亿股（-0.87\%）。 分红方面，近三年现金分红 1 次、合计每股派现 0.05 元，最近一次实施在 2025 年 12 月。 综合看，公司融资结构上以债务滚动为主、股权融资为补充；2025 年筹资活动现金流净额 4.73 亿元，年度内仍为净融入，与前述经营与投资现金流的方向基本一致。

\pfl{未来潜在融资需求分析} 2026 年 6 月末货币资金 5.32 亿元，而短期债务 15.11 亿元（现金短债比 0.35 倍），2026 年上半年经营现金流净额 -3.30 亿元。按“短期债务减去账面货币资金”的滚动偿付口径，静态缺口约 9.79 亿元；若再计入上半年亏损的年化额（1.65 亿元），维持一年正常经营所需的外部资金约 9.79---11.44 亿元。 按第 \ref{sec:financing-need} 节的三档划分，公司属于第一档“补血型”：2026 年上半年归母净利润 -0.83 亿元、6 月末资产负债率 66.4\%。 可行的融资路径按现实性排序为：引入国资或产业投资人（部分公司以控制权让渡为对价）、定向增发（需股价配合，且控股股东需放弃优先认购或同步增资）、债务重组或展期（与银行和债券持有人协商）。 公开披露的后续资金安排线索包括：资本开支安排：2025 年度，购建长期资产支付 1.87 亿元（上年同期 1.10 亿元）。 资金需求还要叠加在建工程：期末在建工程 1.18 亿元、2026 年上半年购建固定资产等支出 0.93 亿元（未年化），这部分支出在周期底部通常会被主动放缓。 综合看，公司的外部资金需求以债务接续为主，滚动偿付口径下的静态缺口约 9.79 亿元，融资节奏将受制于银行续贷意愿与股权融资窗口。




\subsubsection{路德科技（688156）}
\label{co:688156}

\pfl{公司基本情况} 公司前身可追溯至 2020 年，科创板上市的生物发酵饲料企业，前身为以无机固废处理为主的路德环境，2023 年起白酒糟资源化利用成为核心主业并于 2025 年 11 月更名为“路德科技”；连续两年亏损后 2026 年上半年扭亏，正推进简易程序定增扩充产能。 公司于 2020-09-22 在上交所科创板首发上市，注册资本 1.01 亿元，总股本 1.01 亿股。 截至 2026 年 10 月 12 日收盘价 25.40 元，对应总市值 25.6 亿元，近一年+31.9\%，总市值在三级行业“畜禽饲料”的 11 家公司中按由大到小排第\zhnumber{8}位，近一年涨跌幅高于全样本中位数（-10.4\%）42.3 个百分点。 按 2025 年年报披露的行业构成，收入占比前两位为生态保护和环境治理业 95.7\%；其他(补充) 4.3\%。

\pfl{股权结构与控制权} 按 2026 年 6 月 30 日的股东名册，前三大股东为季光明（持股比例未披露）；吕大龙 持股 6.85\%；国信证券股份有限公司（持股比例未披露），股东户数 6,777 户。 从控制权角度看，2026 年 5 月协议转让完成后直接持有公司股份 2,099.77 万股、占总股本 20.85\%；2026 年 3 月，季光明与投资者吕大龙签订股份转让协议，以 19.00 元/股转让 690 万股（占总股本 6.8511\%）、价款 1.31 亿元。

\pfl{经营要点} 经营层面，2026 年上半年，白酒糟生物发酵饲料销量 11 万吨、同比增长 45.92\%。

\pfl{所处行业的基本面} 饲料的供给即产能（门槛低、产能利用率偏低），需求直接来自下游存栏，利润则由玉米、豆粕等原料价格与成品售价的价差决定（第 \ref{sec:framework} 章的供需框架）。 与该公司直接相关的行业数据包括：2026 年上半年，全国工业饲料总产量 16,850 万吨、同比增长 3.9\%。

\pfl{竞争优势} 把该公司放回同类公司比较，2026 年上半年营业收入 2.1 亿元、在三级行业“畜禽饲料”的 11 家公司中排第\zhnumber{11}位；毛利率 30.7\%、高于行业内中位数 25.4 个百分点；半年 ROE 1.2\%（未年化）、高于行业内中位数（-6.25\%）；资产负债率 55.9\%、低于行业内中位数（62.1\%）。 优势在报表上主要体现为高于行业平均的毛利率，这在周期底部是抵御亏损的主要缓冲。

\pfl{历史股价} 本报告样本区间（2021 年 1 月 至 2026 年 10 月）内，股价由 17.19 元变为 25.40 元，累计 +47.8\%，区间最高 37.60 元（2023 年 2 月）、最低 10.78 元（2024 年 8 月）；当前价相当于区间最高价的 67.6\%，为区间最低价的 2.36 倍。 月度收益的年化波动率 47.8\%，高于全样本中位数 37.3\%，在三级行业“畜禽饲料”的 11 家公司中波动率由高到低排第\zhnumber{2}位。

\begin{center}
\begin{minipage}{\textwidth}\centering
\begin{tikzpicture}
\begin{axis}[
  width=12.6cm, height=3.9cm, scale only axis,
  xmin=2020.922, xmax=2026.972, ymin=10.03, ymax=40.23,
  xtick={2021,2022,2023,2024,2025,2026},
  yearticks,
  yticklabel style={font=\scriptsize},
  ylabel={元/股}, ylabel style={font=\scriptsize},
  grid=major, grid style={LightGray, line width=0.3pt},
  axis line style={InkGray, line width=0.5pt},
  every axis plot/.append style={line width=0.9pt},
]
\addplot[AgriGreen] table[x=x,y=y]{data/clean/profiles/px_688156.dat};
\addplot[SoftRed, dashed, line width=0.7pt] table[x=x,y=y]{data/clean/profiles/pxzz_688156.dat};
\addplot[only marks, mark=*, mark size=1.1pt, SoftRed] table[x=x,y=y]{data/clean/profiles/pxzz_688156.dat};
\end{axis}
\end{tikzpicture}

\captionof{figure}[路德科技（688156）股价与周期骨架]{路德科技（688156）月度收盘价（元/股）与 ZigZag 周期骨架（阈值 25\%，圆点为周期转折点）}
\end{minipage}
\end{center}

\pfl{过去周期分析} 以 25\% 的 ZigZag 阈值划分，样本区间内股价形成 5 个主升段与 4 个主跌段。 最大主升段出现在 2022 年 4 月 至 2023 年 2 月，10 个月+145.8\%；最大主跌段为 2023 年 2 月 至 2024 年 8 月，18 个月-71.3\%。 最近一次转折出现在 2026 年 10 月（上涨段自 2026 年 7 月起，3 个月+38.2\%），当前处于该段之后的震荡之中。 公司股价较申万农林牧渔指数提前约 21 个月见底，是板块中较早定价周期反转的公司之一。 饲料夹在原料成本与养殖需求之间（第 \ref{sec:feed} 节），量的部分取决于存栏、利的部分取决于原料与成品的价差，公司 2026 年上半年实现盈利、半年 ROE 1.2\%（未年化），好于全样本中位数（0.75\%），下行周期对其报表的冲击小于同业。 综合区间形态看，该股单段涨跌幅度偏大、以趋势段为主（最大主升 +145.8\%、最大主跌 -71.3\%），年化波动率 47.8\% 与全样本中位数 37.3\% 相比波动明显大于样本中位数。

\pfl{盈利情况} 2026 年上半年营业收入 2.10 亿元（同比 +42.9\%），归母净利润 0.09 亿元，实现扭亏。 以年报为趋势对照，2023---2025 年营业收入由 3.51 亿元变为 3.80 亿元（两年年均复合增速 +4.1\%），归母净利润由 0.27 亿元变为 -0.70 亿元。 按半年报口径，2026 年上半年的 ROE 为 1.2\%（未年化）、销售净利率 4.7\%、毛利率 30.7\%，ROE 在“畜禽饲料”11 家公司中按数值由高到低排第\zhnumber{2}位。 同期全样本半年 ROE 中位数 0.75\%，公司与之相差 0.41 个百分点（略好于中位数公司）。 综合看，公司在报告期维持盈利（高于全样本半年 ROE 中位数 0.75\%），利润的可持续性取决于畜禽饲料景气的延续与成本管控的执行。

\begin{center}\small\setlength{\tabcolsep}{3.4pt}
\captionof{table}[路德科技（688156）关键财务指标]{路德科技（688156）关键财务指标（合并报表口径；两个半年报期均未年化；现金短债比 $=$ 货币资金 $\div$ 短期债务）}
\begin{adjustbox}{max width=\textwidth}
\begin{tabular}{@{}lrrrrr@{}}
\toprule
指标 & 2023 年 & 2024 年 & 2025 年 & 2025 年半年报 & 2026 年半年报 \\
\midrule
营业收入（亿元） & 3.51 & 2.78 & 3.80 & 1.47 & 2.10 \\
归母净利润（亿元） & 0.27 & -0.57 & -0.70 & -0.13 & 0.09 \\
毛利率（\%） & 36.2 & 21.9 & 19.9 & 18.9 & 30.7 \\
ROE（\%） & 3.2 & -6.5 & -8.8 & -1.6 & 1.2 \\
资产负债率（\%） & 40.5 & 49.4 & 58.3 & 47.9 & 55.9 \\
经营活动现金流净额（亿元） & -0.19 & 0.53 & 0.33 & 0.12 & 0.33 \\
货币资金（亿元） & 1.99 & 2.26 & 3.12 & 1.55 & 2.43 \\
有息负债合计（亿元） & 2.90 & 5.02 & 5.85 & 5.22 & 5.50 \\
现金短债比（倍） & 1.27 & 0.86 & 0.83 & 0.51 & 0.72 \\
\bottomrule
\end{tabular}
\end{adjustbox}
\end{center}

\begin{center}
\begin{minipage}{\textwidth}\centering
\begin{tikzpicture}
\begin{groupplot}[
  group style={group size=2 by 1, horizontal sep=0.60cm},
  width=6.85cm, height=4.60cm,
  tick label style={font=\tiny},
  title style={font=\scriptsize\heiti},
  yticklabel style={font=\tiny},
  ylabel style={font=\tiny},
  grid=major, grid style={LightGray, line width=0.3pt},
  axis line style={InkGray, line width=0.5pt},
  enlarge x limits={abs=0.8},
  legend style={font=\scriptsize, at={(0.5,-0.20)}, anchor=north,
                legend columns=2, draw=none, fill=none, column sep=6pt},
]
\nextgroupplot[
  ybar, bar width=4pt,
  ymin=-1.15, ymax=4.36,
  xtick={0,1,2,3,4,5.7},
  xticklabels={2021,2022,2023,2024,2025,2026H1},
  ylabel={亿元},
]
\addplot[fill=AgriGreen!75, draw=AgriGreen, bar shift=-2.6pt] table[x=i, y=rev]{data/clean/profiles/fin_688156.dat};
\addplot[fill=SoftRed!60, draw=SoftRed, bar shift=2.6pt] table[x=i, y=ni]{data/clean/profiles/fin_688156.dat};
\addplot[fill=AgriGreen!30, draw=AgriGreen, pattern=north east lines, bar shift=-2.6pt] table[x=x, y=rev]{data/clean/profiles/finh_688156.dat};
\addplot[fill=SoftRed!20, draw=SoftRed, pattern=north east lines, bar shift=2.6pt] table[x=x, y=ni]{data/clean/profiles/finh_688156.dat};
\legend{营业收入（年/半年）, 归母净利润（年/半年）}
\nextgroupplot[
  ymin=-14.1, ymax=65.0,
  xtick={0,1,2,3,4,5.7},
  xticklabels={2021,2022,2023,2024,2025,2026H1},
  ylabel={\%},
]
\addplot[AgriGold, mark=*, mark size=1.3pt] table[x=i, y=roe]{data/clean/profiles/fin_688156.dat};
\addplot[OilBlue, mark=square*, mark size=1.3pt] table[x=i, y=debt]{data/clean/profiles/fin_688156.dat};
\addplot[only marks, AgriGold, mark=*, mark size=2.0pt] table[x=x, y=roe]{data/clean/profiles/finh_688156.dat};
\addplot[only marks, OilBlue, mark=square*, mark size=2.0pt] table[x=x, y=debt]{data/clean/profiles/finh_688156.dat};
\legend{ROE, 资产负债率}
\end{groupplot}
\end{tikzpicture}
\begin{tikzpicture}
\begin{axis}[
  name=finbot,
  width=13.75cm, height=3.10cm, scale only axis,
  ymin=-0.59, ymax=6.55,
  xtick={0,1,2,3,4,5.7},
  xticklabels={2021,2022,2023,2024,2025,2026H1},
  tick label style={font=\tiny},
  yticklabel style={font=\tiny},
  ylabel={亿元}, ylabel style={font=\tiny},
  ybar, bar width=4pt,
  grid=major, grid style={LightGray, line width=0.3pt},
  axis line style={InkGray, line width=0.5pt},
  enlarge x limits={abs=0.8},
  legend style={font=\scriptsize, at={(0.5,-0.26)}, anchor=north,
                legend columns=2, draw=none, fill=none, column sep=10pt},
]
\addplot[fill=AgriGreen!55, draw=AgriGreen, bar shift=-3.0pt] table[x=i, y=cash]{data/clean/profiles/fin_688156.dat};
\addplot[fill=SoftRed!45, draw=SoftRed, bar shift=3.0pt] table[x=i, y=ibd]{data/clean/profiles/fin_688156.dat};
\addplot[fill=AgriGreen!25, draw=AgriGreen, pattern=north east lines, bar shift=-3.0pt] table[x=x, y=cash]{data/clean/profiles/finh_688156.dat};
\addplot[fill=SoftRed!20, draw=SoftRed, pattern=north east lines, bar shift=3.0pt] table[x=x, y=ibd]{data/clean/profiles/finh_688156.dat};
\legend{货币资金（左轴）, 有息负债（左轴）}
\end{axis}
\begin{axis}[
  at={(finbot.south east)}, anchor=south east,
  width=13.75cm, height=3.10cm, scale only axis,
  axis y line*=right, axis x line*=none, xtick=\empty,
  ymin=-0.38, ymax=0.75,
  tick label style={font=\tiny},
  yticklabel style={font=\tiny}, scaled y ticks=false,
  legend style={font=\scriptsize, at={(0.5,-0.44)}, anchor=north,
                legend columns=1, draw=none, fill=none},
]
\addplot[OilBlue, mark=*, mark size=2.2pt, line width=1.3pt] table[x=i, y=ocf]{data/clean/profiles/fin_688156.dat};
\addplot[only marks, OilBlue, mark=*, mark size=3.2pt] table[x=x, y=ocf]{data/clean/profiles/finh_688156.dat};
\legend{经营活动现金流净额（右轴，亿元，蓝点）}
\end{axis}
\end{tikzpicture}

\captionof{figure}[路德科技（688156）近年主要财务指标]{路德科技（688156）近年主要财务指标（左上：营业收入与归母净利润；右上：ROE 与资产负债率；下：货币资金、有息负债为左轴柱状，经营活动现金流净额为其右轴的蓝点折线。
2021---2025 年为年报口径，2026H1 为 2026 年半年报/半年末，与其前一年报之间留出间隔、以斜纹或空心点区分；半年报数据未年化）}
\end{minipage}
\end{center}

\pfl{财务分析} 2026 年 6 月末资产负债率 55.9\%，较 2025 年末下降 2.4 个百分点（样本中位数 52.2\%，所处三级行业内按由低到高排第\zhnumber{4}位）。 2026 年 6 月末货币资金 2.43 亿元、短期债务（短期借款与一年内到期非流动负债）3.38 亿元、长期债务 2.12 亿元，有息负债合计 5.50 亿元，现金短债比 0.72 倍——覆盖偏紧。 净有息负债 3.06 亿元，相当于其 2026 年上半年营业收入的 146\%。 流动比率 0.82、速动比率 0.69，短期偿债指标偏紧。 2026 年上半年经营活动现金流净额 0.33 亿元，经营现金流与利润同向为正，内生资金可以支撑运营。 2026 年 6 月末在建工程 1.09 亿元，较上年末增加 0.62 亿元，仍有在建投入。 综合看，现金短债比 0.72 倍，短期资金安排偏紧；资产负债率高于样本中位数 3.7 个百分点；有息负债中短期债务占比更高，期限结构仍以滚动续作为主。 公司披露的股东与质押风险为：2026 年 4 月，持股 1.4894\% 的股东武汉德天众享企业管理合伙企业（有限合伙）计划减持不超 1.4894\% 公司股份。

\pfl{重大战略转型} 近三年公司的战略动作集中在：融资与资本运作（2 项）：2026 年 9 月，公司以简易程序定增募集资金主要用于洋河酒糟资源化利用项目建设；2024 年 4 月，提请 2025 年年度股东会授权董事会办理以简易程序向特定对象发行融资总额不超过人民币三亿元且不超过最近一年末净资产 20\% 的股票；主业转型（1 项）：2025 年 11 月，证券简称由“路德环境”变更为“路德科技”。 公告口径上，近三年（2023 年 9 月---2026 年 9 月）公司披露重大重组与并购类公告 0 条、对外投资与转型类公告 3 条、回购与股权激励类公告 20 条，合计公告 411 条。 第一大行业仍是“生态保护和环境治理业”，收入占比 95.7\%，与 2021 年年报相比没有方向性变化。 综合判断，报告期内公司有 3 条与战略相关的公告落地（重大重组与并购 0 条、对外投资与转型 3 条），资本运作与主业调整之间存在直接关联。

\begin{center}\small\setlength{\tabcolsep}{4pt}
\captionof{table}[路德科技（688156）历史融资与资本运作]{路德科技（688156）历史融资与资本运作（据公司公告与公开报道整理；金额为披露值，含首次公开发行、再融资、债券、银行借款与重整投资等）}
\begin{adjustbox}{max width=\textwidth}
\begin{tabular}{@{}p{1.45cm}p{2.55cm}p{4.05cm}p{4.95cm}@{}}
\toprule
时间 & 方式 & 规模 & 用途与说明 \\
\midrule
2020-09 & IPO（上交所科创板） & 首次公开发行 2,\allowbreak{}296 万股，\allowbreak{}发行价 15.91 元/\allowbreak{}股，\allowbreak{}募集资金总额 3.65 亿元 & 研发中心及固废处理相关项目建设；；2020 年 9 月 10 日发行 \\
2023-05 & 向特定对象发行股票（定增） & 发行 834.04 万股，\allowbreak{}发行价 13.57 元/\allowbreak{}股，\allowbreak{}募集资金总额 1.13 亿元 & 补充流动资金及项目建设；发行对象为公司实际控制人季光明 \\
2025-12-22 & 子公司引入战略投资者增资 & 国投聚力旗下聚力杭实基金对永乐路德、\allowbreak{}亳州路德、\allowbreak{}宿迁路德三家子公司战略性增资，\allowbreak{}永乐路德 3,\allowbreak{}700.00 万元） & 加码白酒糟生物发酵饲料核心主业；公司与国投聚力投资管理有限公司签署战略合作协议；聚力杭实基金股东包括国投聚力、\allowbreak{}杭州汇实投资及萧山经开区产业基金 \\
2026-07-27 & 以简易程序向特定对象发行股票（定增，推进中） & 拟募集资金总额不超过 1.50 亿元（含本数） & 洋河酒糟资源化利用项目（拟使用 1.10 亿元）及补充流动资金（4,\allowbreak{}200.00 万元）；发行对象为不超过 35 名特定投资者 \\
\bottomrule
\end{tabular}
\end{adjustbox}
\end{center}

\pfl{历史融投资情况} 从现金流量表看，2025 年公司取得借款收到的现金 3.29 亿元、偿还债务支付的现金 2.77 亿元，筹资活动现金流净额 2.15 亿元（2023 年为取得借款 2.14 亿元、偿还债务 1.00 亿元）。 2026 年上半年筹资活动现金流净额为 -0.36 亿元（取得借款 1.36 亿元、偿还债务 1.39 亿元，未年化），已转为净流出，处于净还债状态。 2025 年分配股利、利润或偿付利息支付的现金合计 0.20 亿元。 公开披露的融资记录共 4 笔，工具类型以 IPO、定向增发、战略投资为主；金额与用途见下表。 其中规模最大的两笔为：2026 年 7 月，定向增发，拟募集资金总额不超过 1.50 亿元（含本数）；2025 年 12 月，战略投资，国投聚力旗下聚力杭实基金对永乐路德、亳州路德、宿迁路德三家子公司战略性增资，永乐路德 3,700.00 万元）。 股本方面，近三年披露股本变动 5 次（增发新股上市 1 次、限售股份上市 1 次）；总股本由 2021 年末的 0.92 亿股变为 1.01 亿股（+9.66\%）。 分红方面，近三年现金分红 2 次、合计每股派现 0.15 元，最近一次实施在 2023 年 12 月。 近三年“再融资”类公告 45 条，涉及定向增发、可转债、募集资金使用。 综合看，公司融资结构上股权与债务融资并举；2025 年筹资活动现金流净额 2.15 亿元，年度内仍为净融入，与前述经营与投资现金流的方向基本一致。

\pfl{未来潜在融资需求分析} 2026 年 6 月末货币资金 2.43 亿元，而短期债务 3.38 亿元（现金短债比 0.72 倍），2026 年上半年经营现金流净额 0.33 亿元。按“短期债务减去账面货币资金”的滚动偿付口径，静态缺口约 0.94 亿元；上半年盈利或接近盈亏平衡，该缺口不随经营亏损进一步扩大。 公司 2026 年 6 月末资产负债率 55.9\%、2026 年上半年 ROE 1.2\%（未年化），现金短债比 0.72 倍、短债覆盖不足，融资需求首先用于置换与滚动存量债务、补充营运资金，属于“补血型”而非扩张型。 公开披露的后续资金安排线索包括：股权融资线索：2026 年 7 月，公司以简易程序向特定对象发行股票募资不超 1.52 亿元，用途为洋河酒糟资源化利用项目与补充流动资金；资本开支安排：2020 年，首次公开发行股票募集资金使用情况。 资金需求还要叠加在建工程：期末在建工程 1.09 亿元、2026 年上半年购建固定资产等支出 0.70 亿元（未年化），这部分支出在周期底部通常会被主动放缓。 综合看，公司的外部资金需求以营运资金周转与在建项目投入为主，滚动偿付口径下的静态缺口约 0.94 亿元，融资节奏将受制于银行续贷意愿与股权融资窗口。




\subsubsection{傲农生物（603363）}
\label{co:603363}

\pfl{公司基本情况} 公司前身可追溯至 2011 年，福建饲料起家、后跨界养猪的农牧企业；2024 年末完成破产重整，控股股东变更为泉州国资牵头联合体且无实际控制人，2026 年上半年猪价低迷亏损 1.90 亿元。 公司于 2017-09-26 在上交所首发上市，注册资本 27.31 亿元，总股本 8.70 亿股。 截至 2026 年 10 月 12 日收盘价 2.81 元，对应总市值 24.5 亿元，近一年-47.7\%，总市值在三级行业“畜禽饲料”的 11 家公司中按由大到小排第\zhnumber{9}位，近一年涨跌幅低于全样本中位数（-10.4\%）37.3 个百分点。 按 2026 年半年报披露的行业构成，收入占比前三位为饲料行业 69.5\%；饲养行业 16.9\%；屠宰食品 13.3\%。 与 2021 年年报相比，第一大行业口径由“饲料及动保行业”（56.4\%）变为“饲料行业”（69.5\%）；两期披露的分类口径有调整，占比差异中同时包含分类因素。

\pfl{股权结构与控制权} 按 2026 年 7 月 27 日的股东名册，前三大股东为福建傲农生物科技集团股份有限公司破产企业财产处置专用 持股 8.87\%；漳州傲农投资有限公司 持股 7.86\%；泉州市傲发同心股权投资管理合伙企业(有限合伙) 持股 7.32\%。 股权结构的另一条线索是，2024 年，底完成重整后，由泉州市国资委下属泉州发展集团牵头、共 5 家产业投资人组成的联合体成为控股股东；破产企业财产处置专用账户持股 8.87\%（2.42 亿股）。

\pfl{经营要点} 业务与产能层面，2026 年，生猪出栏全年目标约 240 万头。

\pfl{所处行业的基本面} 饲料的供给即产能（门槛低、产能利用率偏低），需求直接来自下游存栏，利润则由玉米、豆粕等原料价格与成品售价的价差决定（第 \ref{sec:framework} 章的供需框架）。 当期该行业的关键数据是：2026 年上半年，全国生猪出栏 37,246 万头、同比增长 1.7\%；全国年出栏 500 头以上生猪养殖规模化率由 42\% 提升至 71\%。

\pfl{竞争优势} 把该公司放回同类公司比较，2026 年上半年营业收入 45.2 亿元、在三级行业“畜禽饲料”的 11 家公司中排第\zhnumber{6}位；毛利率 1.7\%、低于行业内中位数 3.6 个百分点；半年 ROE -6.2\%（未年化）、低于行业内中位数（-6.25\%）；资产负债率 62.1\%、高于行业内中位数（62.1\%）。 公司在公开资料中自述与同业比较的相对优势在于：2026 年上半年饲料外销 97.36 万吨、毛利率 5.51\%；重整中清退了部分产能利用率较低的资产。 就当期报表看，公司尚未把经营层面的特征转化为高于行业的资本回报，周期底部的相对优势更多体现在现金流与资产负债表上。

\pfl{历史股价} 本报告样本区间（2021 年 1 月 至 2026 年 10 月）内，股价由 13.40 元变为 2.81 元，累计 -79.0\%，区间最高 23.80 元（2022 年 3 月）、最低 2.68 元（2026 年 6 月）；当前价相当于区间最高价的 11.8\%，已接近区间最低价（1.05 倍）。 月度收益的年化波动率 52.2\%，高于全样本中位数 37.3\%，在三级行业“畜禽饲料”的 11 家公司中波动率由高到低排第\zhnumber{1}位。 把股价阶段与公司事件对齐看，2022 年 7 月 至 2025 年 2 月股价下跌-86.0\%，同期（2025 年 1 月）法院裁定终结重整程序后，公司被撤销因重整实施的退市风险警示。 从时间对应关系看，公司层面的资本运作与股权、风险事件解释了区间内多数急涨急跌，与行业指数的同步性相对有限。

\begin{center}
\begin{minipage}{\textwidth}\centering
\begin{tikzpicture}
\begin{axis}[
  width=12.6cm, height=3.9cm, scale only axis,
  xmin=2020.922, xmax=2026.972, ymin=2.49, ymax=25.47,
  xtick={2021,2022,2023,2024,2025,2026},
  yearticks,
  yticklabel style={font=\scriptsize},
  ylabel={元/股}, ylabel style={font=\scriptsize},
  grid=major, grid style={LightGray, line width=0.3pt},
  axis line style={InkGray, line width=0.5pt},
  every axis plot/.append style={line width=0.9pt},
]
\addplot[AgriGreen] table[x=x,y=y]{data/clean/profiles/px_603363.dat};
\addplot[SoftRed, dashed, line width=0.7pt] table[x=x,y=y]{data/clean/profiles/pxzz_603363.dat};
\addplot[only marks, mark=*, mark size=1.1pt, SoftRed] table[x=x,y=y]{data/clean/profiles/pxzz_603363.dat};
\end{axis}
\end{tikzpicture}

\captionof{figure}[傲农生物（603363）股价与周期骨架]{傲农生物（603363）月度收盘价（元/股）与 ZigZag 周期骨架（阈值 25\%，圆点为周期转折点）}
\end{minipage}
\end{center}

\pfl{过去周期分析} 以 25\% 的 ZigZag 阈值划分，样本区间内股价形成 3 个主升段与 4 个主跌段。 最大主跌段为 2022 年 7 月 至 2025 年 2 月，31 个月-86.0\%；最大主升段出现在 2021 年 7 月 至 2022 年 3 月，8 个月+196.8\%。 最近一次转折出现在 2026 年 6 月（下跌段自 2025 年 10 月起，8 个月-47.6\%），当前处于该段的延续之中。 公司股价与申万农林牧渔指数几乎同步见底，个股并未走出独立行情。 饲料夹在原料成本与养殖需求之间（第 \ref{sec:feed} 节），量的部分取决于存栏、利的部分取决于原料与成品的价差，公司 2026 年 6 月末资产负债率 62.1\%、半年 ROE -6.2\%（未年化），在本轮下行中属于随行业同步调整的一类。 综合区间形态看，该股单段涨跌幅度偏大、以趋势段为主（最大主升 +196.8\%、最大主跌 -86.0\%），年化波动率 52.2\% 与全样本中位数 37.3\% 相比波动明显大于样本中位数。

\pfl{盈利情况} 2026 年上半年营业收入 45.21 亿元（同比 +14.3\%），归母净利润 -1.90 亿元，由上年同期的盈利转为亏损。 以年报为趋势对照，2023---2025 年营业收入由 194.58 亿元变为 84.75 亿元（两年年均复合增速 -34.0\%），归母净利润由 -36.51 亿元变为 0.78 亿元。 按半年报口径，2026 年上半年的 ROE 为 -6.2\%（未年化）、销售净利率 -6.0\%、毛利率 1.7\%，ROE 在“畜禽饲料”11 家公司中按数值由高到低排第\zhnumber{6}位。 按同一口径，三级行业“畜禽饲料”11 家公司的半年 ROE 中位数为 -6.25\%，该公司低于其中位数 0.00 个百分点。 从公司披露的原因看，业绩变动主要来自：2026 年上半年，生猪出栏量同比增加 40.53\%（78.33 万头→110.08 万头）。 综合看，公司仍处亏损区间、由盈转亏，半年 ROE 在三级行业内按由高到低排第\zhnumber{6}位，单位盈利能力的修复依赖行业景气与自身成本端的改善，报表弹性主要体现为价格与销量的方向性变化。

\begin{center}\small\setlength{\tabcolsep}{3.4pt}
\captionof{table}[傲农生物（603363）关键财务指标]{傲农生物（603363）关键财务指标（合并报表口径；两个半年报期均未年化；现金短债比 $=$ 货币资金 $\div$ 短期债务）}
\begin{adjustbox}{max width=\textwidth}
\begin{tabular}{@{}lrrrrr@{}}
\toprule
指标 & 2023 年 & 2024 年 & 2025 年 & 2025 年半年报 & 2026 年半年报 \\
\midrule
营业收入（亿元） & 194.58 & 87.63 & 84.75 & 39.57 & 45.21 \\
归母净利润（亿元） & -36.51 & 5.79 & 0.78 & 3.61 & -1.90 \\
毛利率（\%） & -4.7 & 4.5 & 6.3 & 8.3 & 1.7 \\
ROE（\%） & 0.0 & --- & 2.8 & 12.5 & -6.2 \\
资产负债率（\%） & 103.7 & 67.8 & 59.0 & 60.5 & 62.1 \\
经营活动现金流净额（亿元） & 9.84 & -1.17 & 1.53 & -3.16 & -3.29 \\
货币资金（亿元） & 2.08 & 13.84 & 6.98 & 8.20 & 4.75 \\
有息负债合计（亿元） & 80.52 & 29.08 & 23.29 & 26.64 & 23.80 \\
现金短债比（倍） & 0.04 & 1.09 & 0.80 & 0.93 & 0.49 \\
\bottomrule
\end{tabular}
\end{adjustbox}
\end{center}

\begin{center}
\begin{minipage}{\textwidth}\centering
\begin{tikzpicture}
\begin{groupplot}[
  group style={group size=2 by 1, horizontal sep=0.60cm},
  width=6.85cm, height=4.60cm,
  tick label style={font=\tiny},
  title style={font=\scriptsize\heiti},
  yticklabel style={font=\tiny},
  ylabel style={font=\tiny},
  grid=major, grid style={LightGray, line width=0.3pt},
  axis line style={InkGray, line width=0.5pt},
  enlarge x limits={abs=0.8},
  legend style={font=\scriptsize, at={(0.5,-0.20)}, anchor=north,
                legend columns=2, draw=none, fill=none, column sep=6pt},
]
\nextgroupplot[
  ybar, bar width=4pt,
  ymin=-61.77, ymax=246.45,
  xtick={0,1,2,3,4,5.7},
  xticklabels={2021,2022,2023,2024,2025,2026H1},
  ylabel={亿元},
]
\addplot[fill=AgriGreen!75, draw=AgriGreen, bar shift=-2.6pt] table[x=i, y=rev]{data/clean/profiles/fin_603363.dat};
\addplot[fill=SoftRed!60, draw=SoftRed, bar shift=2.6pt] table[x=i, y=ni]{data/clean/profiles/fin_603363.dat};
\addplot[fill=AgriGreen!30, draw=AgriGreen, pattern=north east lines, bar shift=-2.6pt] table[x=x, y=rev]{data/clean/profiles/finh_603363.dat};
\addplot[fill=SoftRed!20, draw=SoftRed, pattern=north east lines, bar shift=2.6pt] table[x=x, y=ni]{data/clean/profiles/finh_603363.dat};
\legend{营业收入（年/半年）, 归母净利润（年/半年）}
\nextgroupplot[
  ymin=-88.3, ymax=121.5,
  xtick={0,1,2,3,4,5.7},
  xticklabels={2021,2022,2023,2024,2025,2026H1},
  ylabel={\%},
]
\addplot[AgriGold, mark=*, mark size=1.3pt] table[x=i, y=roe]{data/clean/profiles/fin_603363.dat};
\addplot[OilBlue, mark=square*, mark size=1.3pt] table[x=i, y=debt]{data/clean/profiles/fin_603363.dat};
\addplot[only marks, AgriGold, mark=*, mark size=2.0pt] table[x=x, y=roe]{data/clean/profiles/finh_603363.dat};
\addplot[only marks, OilBlue, mark=square*, mark size=2.0pt] table[x=x, y=debt]{data/clean/profiles/finh_603363.dat};
\legend{ROE, 资产负债率}
\end{groupplot}
\end{tikzpicture}
\begin{tikzpicture}
\begin{axis}[
  name=finbot,
  width=13.75cm, height=3.10cm, scale only axis,
  ymin=-9.21, ymax=103.16,
  xtick={0,1,2,3,4,5.7},
  xticklabels={2021,2022,2023,2024,2025,2026H1},
  tick label style={font=\tiny},
  yticklabel style={font=\tiny},
  ylabel={亿元}, ylabel style={font=\tiny},
  ybar, bar width=4pt,
  grid=major, grid style={LightGray, line width=0.3pt},
  axis line style={InkGray, line width=0.5pt},
  enlarge x limits={abs=0.8},
  legend style={font=\scriptsize, at={(0.5,-0.26)}, anchor=north,
                legend columns=2, draw=none, fill=none, column sep=10pt},
]
\addplot[fill=AgriGreen!55, draw=AgriGreen, bar shift=-3.0pt] table[x=i, y=cash]{data/clean/profiles/fin_603363.dat};
\addplot[fill=SoftRed!45, draw=SoftRed, bar shift=3.0pt] table[x=i, y=ibd]{data/clean/profiles/fin_603363.dat};
\addplot[fill=AgriGreen!25, draw=AgriGreen, pattern=north east lines, bar shift=-3.0pt] table[x=x, y=cash]{data/clean/profiles/finh_603363.dat};
\addplot[fill=SoftRed!20, draw=SoftRed, pattern=north east lines, bar shift=3.0pt] table[x=x, y=ibd]{data/clean/profiles/finh_603363.dat};
\legend{货币资金（左轴）, 有息负债（左轴）}
\end{axis}
\begin{axis}[
  at={(finbot.south east)}, anchor=south east,
  width=13.75cm, height=3.10cm, scale only axis,
  axis y line*=right, axis x line*=none, xtick=\empty,
  ymin=-6.70, ymax=13.78,
  tick label style={font=\tiny},
  yticklabel style={font=\tiny}, scaled y ticks=false,
  legend style={font=\scriptsize, at={(0.5,-0.44)}, anchor=north,
                legend columns=1, draw=none, fill=none},
]
\addplot[OilBlue, mark=*, mark size=2.2pt, line width=1.3pt] table[x=i, y=ocf]{data/clean/profiles/fin_603363.dat};
\addplot[only marks, OilBlue, mark=*, mark size=3.2pt] table[x=x, y=ocf]{data/clean/profiles/finh_603363.dat};
\legend{经营活动现金流净额（右轴，亿元，蓝点）}
\end{axis}
\end{tikzpicture}

\captionof{figure}[傲农生物（603363）近年主要财务指标]{傲农生物（603363）近年主要财务指标（左上：营业收入与归母净利润；右上：ROE 与资产负债率；下：货币资金、有息负债为左轴柱状，经营活动现金流净额为其右轴的蓝点折线。
2021---2025 年为年报口径，2026H1 为 2026 年半年报/半年末，与其前一年报之间留出间隔、以斜纹或空心点区分；半年报数据未年化）}
\end{minipage}
\end{center}

\pfl{财务分析} 2026 年 6 月末资产负债率 62.1\%，较 2025 年末上升 3.0 个百分点（样本中位数 52.2\%，所处三级行业内按由低到高排第\zhnumber{6}位）。 2026 年 6 月末货币资金 4.75 亿元、短期债务（短期借款与一年内到期非流动负债）9.72 亿元、长期债务 14.08 亿元，有息负债合计 23.80 亿元，现金短债比 0.49 倍——覆盖偏紧。 净有息负债 19.05 亿元，相当于其 2026 年上半年营业收入的 42\%。 流动比率 0.89、速动比率 0.62，短期偿付压力较大。 2026 年上半年经营活动现金流净额 -3.29 亿元，经营现金流与利润同时为负，日常运营对债务与股东投入的依赖度较高。 2026 年 6 月末在建工程 4.45 亿元，较上年末增加 0.23 亿元，仍有在建投入。 2026 年 6 月末商誉 3.38 亿元，占归母净资产的 10.1\%，是净资产中最脆弱的一块。 综合看，账面货币资金对有息负债与短期债务的覆盖不足，滚动偿付对新增融资的依赖度较高；资产负债率高于样本中位数 9.9 个百分点；有息负债中长期债务占比更高，期限结构对短债滚动的压力较小。 公司披露的重整与债务违约风险为：2026 年上半年，由盈转亏：归母净利润 -1.90 亿元、扣非 -2.52 亿元（上年同期含重整收益基数）；2024 年，共发布 26 份新增部分债务逾期公告。

\pfl{重大战略转型} 从公开披露看，公司的战略推进集中在：并购与重组（3 项）：2025---2026 年，重整后剥离部分闲置资产、保留规模猪场有序复产；2025 年 1 月，法院裁定终结重整程序后，公司被撤销因重整实施的退市风险警示；其他动作（1 项）：2026 年，推出 2026 年限制性股票激励计划。 公告口径上，近三年（2023 年 9 月---2026 年 9 月）公司披露重大重组与并购类公告 4 条、对外投资与转型类公告 5 条、回购与股权激励类公告 31 条，合计公告 752 条。 主营结构的口径变化已经落到报表上：2021 年年报的第一大行业为“饲料及动保行业”（56.4\%），2026 年半年报为“饲料行业”（69.5\%）；公司在这两期的分部划分方式有过调整。 综合判断，公司当前的战略重心不是扩张而是化解债务与完成重整，业务层面的调整让位于司法重整程序。

\begin{center}\small\setlength{\tabcolsep}{4pt}
\captionof{table}[傲农生物（603363）历史融资与资本运作]{傲农生物（603363）历史融资与资本运作（据公司公告与公开报道整理；金额为披露值，含首次公开发行、再融资、债券、银行借款与重整投资等）}
\begin{adjustbox}{max width=\textwidth}
\begin{tabular}{@{}p{1.45cm}p{2.55cm}p{4.05cm}p{4.95cm}@{}}
\toprule
时间 & 方式 & 规模 & 用途与说明 \\
\midrule
2017-09 & IPO（上交所主板） & 6,\allowbreak{}000.00 万股、\allowbreak{}4.79 元/\allowbreak{}股，\allowbreak{}募集资金总额 2.87 亿元 & 饲料及养殖项目等；上市日 2017-09-26 \\
2020—2021 & 公开发行可转换公司债券（公司口径为“2020 年公开发行可转换公司债券”） & 10.00 亿元（1,\allowbreak{}000 万张，\allowbreak{}面值 100.00 元，\allowbreak{}期限 6 年） & 母猪场、\allowbreak{}饲料及生猪项目，\allowbreak{}其中 3.00 亿元偿还银行借款；；方案于 2020 年 12 月处于申报推进阶段 \\
2020-05 & 非公开发行股票 & 募集资金总额 13.90 亿元、\allowbreak{}净额 13.65 亿元 & 9 个生猪养殖项目和 2 个饲料生产项目，\allowbreak{}其中 4.10 亿元用于偿还银行贷款；原拟募集不超过 13.90 亿元 \\
2022-04 & 非公开发行股票 & 1.30 亿股，\allowbreak{}募集资金 14.33 亿元 & 宜丰傲农 15,\allowbreak{}000 头母猪自繁自养猪场项目（一期）、\allowbreak{}永新傲农年产 18 万吨饲料（一期）项目及补充流动资金；由公司控股股东厦门傲农投资、\allowbreak{}实际控制人吴有林及其配偶设立的合伙企业厦门裕泽投资全额认购，\allowbreak{}构成关联交易；原拟募集不超过 15.00 亿元 \\
2024-12 & 重整计划出资人权益调整（转增股票） & 转增股票中的 10.05 亿股由重整投资人有条件受让 & 清偿债务、\allowbreak{}引入战略及财务投资人；重整投资转增股份于 2026-01-29、\allowbreak{}2026-03-27、\allowbreak{}2026-05-27、\allowbreak{}2026-06-25 分批上市流通，\allowbreak{}合计 1.02 亿股 \\
2024-12-25 & 重整投资款 & 管理人账户已收到全体重整投资人或其指定主体支付的全部重整投资款合计 17.09 亿元（约 17.09 亿元） & 执行重整计划、\allowbreak{}清偿债务及补充营运资金；2024-09-13 确定泉发外贸联合体为中选重整投资人 \\
\bottomrule
\end{tabular}
\end{adjustbox}
\end{center}

\pfl{历史融投资情况} 从现金流量表看，2025 年公司取得借款收到的现金 2.97 亿元、偿还债务支付的现金 5.20 亿元、吸收权益性投资 0.32 亿元，筹资活动现金流净额 -5.67 亿元（2023 年为取得借款 55.35 亿元、偿还债务 57.39 亿元）。 2026 年上半年筹资活动现金流净额为 1.68 亿元（取得借款 1.19 亿元、偿还债务 1.18 亿元，未年化），仍处于净融入状态。 2025 年分配股利、利润或偿付利息支付的现金合计 1.07 亿元。 公开披露的融资记录共 6 笔，工具类型以 IPO、可转债、定向增发、重整投资为主；金额与用途见下表。 其中规模最大的两笔为：2024 年 12 月，重整投资，管理人账户已收到全体重整投资人或其指定主体支付的全部重整投资款合计 17.09 亿元（约 17.09 亿元）；2022 年 4 月，定向增发，1.30 亿股，募集资金 14.33 亿元。 股本方面，近三年披露股本变动 8 次（激励股份解禁 4 次、股份回购 1 次）；总股本由 2021 年末的 6.84 亿股变为 8.70 亿股（+27.16\%）。 分红方面，近三年未实施现金分红，与重整期间的资金约束一致。 近三年“再融资”类公告 23 条，涉及定向增发、募集资金使用。 综合看，公司融资结构上以债务滚动为主、股权融资为补充；2025 年筹资活动现金流净额 -5.67 亿元，融资节奏仍以维持存量债务周转为主，与前述经营与投资现金流的方向基本一致。

\pfl{未来潜在融资需求分析} 2026 年 6 月末货币资金 4.75 亿元，而短期债务 9.72 亿元（现金短债比 0.49 倍），2026 年上半年经营现金流净额 -3.29 亿元。按“短期债务减去账面货币资金”的滚动偿付口径，静态缺口约 4.97 亿元；若再计入上半年亏损的年化额（3.80 亿元），维持一年正常经营所需的外部资金约 4.97---8.77 亿元。 公司 2026 年上半年归母净利润为负（-1.90 亿元），但 6 月末资产负债率 62.1\% 尚未进入第一档（亏损且负债率 $>$65\%）的区间，当前融资需求以补充流动资金、支撑低谷期经营性支出为主。 公开披露的后续资金安排线索包括：股权融资线索：2024 年，公司终止 2022 年度向特定对象发行股票，转而通过破产重整化解债务；2026 年，第二次临时股东会审议因 2026 年限制性股票激励计划导致的公司股份总数与注册资本变更。 资金需求还要叠加在建工程：期末在建工程 4.45 亿元、2026 年上半年购建固定资产等支出 0.68 亿元（未年化），这部分支出在周期底部通常会被主动放缓。 综合看，公司的外部资金需求以补充营运资金为主、项目投入为辅，滚动偿付口径下的静态缺口约 4.97 亿元，融资的时点与规模更多取决于债权人与重整投资人的进度。




\subsubsection{播恩集团（001366）}
\label{co:001366}

\pfl{公司基本情况} 2006 年设立至今，播恩集团是总部位于广州、注册在江西赣州的猪饲料企业，聚焦教槽料、乳猪料等幼小动物营养产品，2023 年 3 月登陆深交所主板；上市后连续两年亏损，2026 年起通过终止募投项目、以募集资金改道收购希杰三家饲料公司切入反刍料赛道。 公司于 2023-03-07 在深交所主板首发上市，注册资本 1.61 亿元，总股本 1.61 亿股。 截至 2026 年 10 月 12 日收盘价 12.95 元，对应总市值 20.8 亿元，近一年+3.6\%，总市值在三级行业“畜禽饲料”的 11 家公司中按由大到小排第\zhnumber{10}位，近一年涨跌幅高于全样本中位数（-10.4\%）14.0 个百分点。 按 2026 年半年报披露的行业构成，收入几乎全部来自饲料行业 100.0\%。

\pfl{股权结构与控制权} 按 2026 年 6 月 30 日的股东名册，前三大股东为江西八维生物集团有限公司 持股 64.69\%；赣州爱特威投资合伙企业(有限合伙) 持股 1.72\%；赣州华古投资合伙企业(有限合伙) 持股 0.83\%，股东户数 9,464 户。 股权结构的另一条线索是，持股 64.69\%（1.0395 亿股）。

\pfl{经营要点} 业务与产能层面，2024 年广西钦州、甘肃兰州的生物饲料项目投产。

\pfl{所处行业的基本面} 饲料的供给即产能（门槛低、产能利用率偏低），需求直接来自下游存栏，利润则由玉米、豆粕等原料价格与成品售价的价差决定（第 \ref{sec:framework} 章的供需框架）。 行业层面的最新情况是：2026 年上半年，全国工业饲料总产量 16,850 万吨、同比增长 3.9\%。

\pfl{竞争优势} 把该公司放回同类公司比较，2026 年上半年营业收入 6.7 亿元、在三级行业“畜禽饲料”的 11 家公司中排第\zhnumber{9}位；毛利率 13.7\%、高于行业内中位数 8.5 个百分点；半年 ROE 1.3\%（未年化）、高于行业内中位数（-6.25\%）；资产负债率 38.1\%、低于行业内中位数（62.1\%）。 公司在公开资料中自述与同业比较的相对优势在于：2026 年 9 月，公司以 1.12 亿元收购希杰（沈阳）、希杰（天津）、希杰（聊城）三家饲料公司 100\% 股权；截至 2026 年 6 月 30 日拥有授权发明专利 90 项、实用新型专利 32 项。 优势在报表上主要体现为高于行业平均的毛利率，这在周期底部是抵御亏损的主要缓冲。

\pfl{历史股价} 本报告样本区间（2023 年 3 月 至 2026 年 10 月）内，股价由 18.21 元变为 12.95 元，累计 -28.9\%，区间最高 18.21 元（2023 年 3 月）、最低 9.92 元（2024 年 8 月）；当前价相当于区间最高价的 71.1\%，为区间最低价的 1.31 倍。 月度收益的年化波动率 36.2\%，低于全样本中位数 37.3\%，在三级行业“畜禽饲料”的 11 家公司中波动率由高到低排第\zhnumber{6}位。

\begin{center}
\begin{minipage}{\textwidth}\centering
\begin{tikzpicture}
\begin{axis}[
  width=12.6cm, height=3.9cm, scale only axis,
  xmin=2023.088, xmax=2026.972, ymin=9.23, ymax=19.48,
  xtick={2023,2024,2025,2026},
  yearticks,
  yticklabel style={font=\scriptsize},
  ylabel={元/股}, ylabel style={font=\scriptsize},
  grid=major, grid style={LightGray, line width=0.3pt},
  axis line style={InkGray, line width=0.5pt},
  every axis plot/.append style={line width=0.9pt},
]
\addplot[AgriGreen] table[x=x,y=y]{data/clean/profiles/px_001366.dat};
\addplot[SoftRed, dashed, line width=0.7pt] table[x=x,y=y]{data/clean/profiles/pxzz_001366.dat};
\addplot[only marks, mark=*, mark size=1.1pt, SoftRed] table[x=x,y=y]{data/clean/profiles/pxzz_001366.dat};
\end{axis}
\end{tikzpicture}

\captionof{figure}[播恩集团（001366）股价与周期骨架]{播恩集团（001366）月度收盘价（元/股）与 ZigZag 周期骨架（阈值 25\%，圆点为周期转折点）}
\end{minipage}
\end{center}

\pfl{过去周期分析} 以 25\% 的 ZigZag 阈值划分，样本区间内股价形成 2 个主升段与 3 个主跌段。 最大主跌段为 2023 年 3 月 至 2024 年 2 月，11 个月-38.8\%；最大主升段出现在 2024 年 8 月 至 2026 年 4 月，20 个月+56.2\%。 最近一次转折出现在 2026 年 6 月（下跌段自 2026 年 4 月起，2 个月-30.1\%），当前处于该段的延续之中。 公司股价较申万农林牧渔指数提前约 21 个月见底，是板块中较早定价周期反转的公司之一。 饲料夹在原料成本与养殖需求之间（第 \ref{sec:feed} 节），量的部分取决于存栏、利的部分取决于原料与成品的价差，公司 2026 年上半年实现盈利、半年 ROE 1.3\%（未年化），好于全样本中位数（0.75\%），下行周期对其报表的冲击小于同业。 综合区间形态看，该股单段涨跌幅度偏大、以趋势段为主（最大主升 +56.2\%、最大主跌 -38.8\%），年化波动率 36.2\% 与全样本中位数 37.3\% 相比波动与样本中位数接近。

\pfl{盈利情况} 2026 年上半年营业收入 6.69 亿元（同比 +9.0\%），归母净利润 0.10 亿元，实现扭亏。 以年报为趋势对照，2023---2025 年营业收入由 14.36 亿元变为 13.60 亿元（两年年均复合增速 -2.7\%），归母净利润由 0.49 亿元变为 -0.20 亿元。 按半年报口径，2026 年上半年的 ROE 为 1.3\%（未年化）、销售净利率 1.5\%、毛利率 13.7\%，ROE 在“畜禽饲料”11 家公司中按数值由高到低排第\zhnumber{1}位。 按同一口径，三级行业“畜禽饲料”11 家公司的半年 ROE 中位数为 -6.25\%，该公司高于其中位数 7.52 个百分点。 综合看，公司在报告期维持盈利（高于全样本半年 ROE 中位数 0.75\%），利润的可持续性取决于畜禽饲料景气的延续与成本管控的执行。

\begin{center}\small\setlength{\tabcolsep}{3.4pt}
\captionof{table}[播恩集团（001366）关键财务指标]{播恩集团（001366）关键财务指标（合并报表口径；两个半年报期均未年化；现金短债比 $=$ 货币资金 $\div$ 短期债务）}
\begin{adjustbox}{max width=\textwidth}
\begin{tabular}{@{}lrrrrr@{}}
\toprule
指标 & 2023 年 & 2024 年 & 2025 年 & 2025 年半年报 & 2026 年半年报 \\
\midrule
营业收入（亿元） & 14.36 & 10.17 & 13.60 & 6.14 & 6.69 \\
归母净利润（亿元） & 0.49 & -0.29 & -0.20 & -0.15 & 0.10 \\
毛利率（\%） & 14.2 & 13.8 & 12.5 & 12.1 & 13.7 \\
ROE（\%） & 6.1 & -3.4 & -2.5 & -1.8 & 1.3 \\
资产负债率（\%） & 17.8 & 25.0 & 32.9 & 29.3 & 38.1 \\
经营活动现金流净额（亿元） & 0.46 & -0.31 & -0.28 & -0.95 & -1.39 \\
货币资金（亿元） & 3.52 & 2.91 & 2.93 & 2.21 & 2.60 \\
有息负债合计（亿元） & 0.75 & 1.51 & 2.56 & 2.21 & 3.68 \\
现金短债比（倍） & 35.39 & 3.13 & 1.69 & 1.37 & 1.10 \\
\bottomrule
\end{tabular}
\end{adjustbox}
\end{center}

\begin{center}
\begin{minipage}{\textwidth}\centering
\begin{tikzpicture}
\begin{groupplot}[
  group style={group size=2 by 1, horizontal sep=0.60cm},
  width=6.85cm, height=4.60cm,
  tick label style={font=\tiny},
  title style={font=\scriptsize\heiti},
  yticklabel style={font=\tiny},
  ylabel style={font=\tiny},
  grid=major, grid style={LightGray, line width=0.3pt},
  axis line style={InkGray, line width=0.5pt},
  enlarge x limits={abs=0.8},
  legend style={font=\scriptsize, at={(0.5,-0.20)}, anchor=north,
                legend columns=2, draw=none, fill=none, column sep=6pt},
]
\nextgroupplot[
  ybar, bar width=4pt,
  ymin=-2.09, ymax=19.87,
  xtick={0,1,2,3,4,5.7},
  xticklabels={2021,2022,2023,2024,2025,2026H1},
  ylabel={亿元},
]
\addplot[fill=AgriGreen!75, draw=AgriGreen, bar shift=-2.6pt] table[x=i, y=rev]{data/clean/profiles/fin_001366.dat};
\addplot[fill=SoftRed!60, draw=SoftRed, bar shift=2.6pt] table[x=i, y=ni]{data/clean/profiles/fin_001366.dat};
\addplot[fill=AgriGreen!30, draw=AgriGreen, pattern=north east lines, bar shift=-2.6pt] table[x=x, y=rev]{data/clean/profiles/finh_001366.dat};
\addplot[fill=SoftRed!20, draw=SoftRed, pattern=north east lines, bar shift=2.6pt] table[x=x, y=ni]{data/clean/profiles/finh_001366.dat};
\legend{营业收入（年/半年）, 归母净利润（年/半年）}
\nextgroupplot[
  ymin=-6.8, ymax=42.3,
  xtick={0,1,2,3,4,5.7},
  xticklabels={2021,2022,2023,2024,2025,2026H1},
  ylabel={\%},
]
\addplot[AgriGold, mark=*, mark size=1.3pt] table[x=i, y=roe]{data/clean/profiles/fin_001366.dat};
\addplot[OilBlue, mark=square*, mark size=1.3pt] table[x=i, y=debt]{data/clean/profiles/fin_001366.dat};
\addplot[only marks, AgriGold, mark=*, mark size=2.0pt] table[x=x, y=roe]{data/clean/profiles/finh_001366.dat};
\addplot[only marks, OilBlue, mark=square*, mark size=2.0pt] table[x=x, y=debt]{data/clean/profiles/finh_001366.dat};
\legend{ROE, 资产负债率}
\end{groupplot}
\end{tikzpicture}
\begin{tikzpicture}
\begin{axis}[
  name=finbot,
  width=13.75cm, height=3.10cm, scale only axis,
  ymin=-0.37, ymax=4.12,
  xtick={0,1,2,3,4,5.7},
  xticklabels={2021,2022,2023,2024,2025,2026H1},
  tick label style={font=\tiny},
  yticklabel style={font=\tiny},
  ylabel={亿元}, ylabel style={font=\tiny},
  ybar, bar width=4pt,
  grid=major, grid style={LightGray, line width=0.3pt},
  axis line style={InkGray, line width=0.5pt},
  enlarge x limits={abs=0.8},
  legend style={font=\scriptsize, at={(0.5,-0.26)}, anchor=north,
                legend columns=2, draw=none, fill=none, column sep=10pt},
]
\addplot[fill=AgriGreen!55, draw=AgriGreen, bar shift=-3.0pt] table[x=i, y=cash]{data/clean/profiles/fin_001366.dat};
\addplot[fill=SoftRed!45, draw=SoftRed, bar shift=3.0pt] table[x=i, y=ibd]{data/clean/profiles/fin_001366.dat};
\addplot[fill=AgriGreen!25, draw=AgriGreen, pattern=north east lines, bar shift=-3.0pt] table[x=x, y=cash]{data/clean/profiles/finh_001366.dat};
\addplot[fill=SoftRed!20, draw=SoftRed, pattern=north east lines, bar shift=3.0pt] table[x=x, y=ibd]{data/clean/profiles/finh_001366.dat};
\legend{货币资金（左轴）, 有息负债（左轴）}
\end{axis}
\begin{axis}[
  at={(finbot.south east)}, anchor=south east,
  width=13.75cm, height=3.10cm, scale only axis,
  axis y line*=right, axis x line*=none, xtick=\empty,
  ymin=-1.88, ymax=1.06,
  tick label style={font=\tiny},
  yticklabel style={font=\tiny}, scaled y ticks=false,
  legend style={font=\scriptsize, at={(0.5,-0.44)}, anchor=north,
                legend columns=1, draw=none, fill=none},
]
\addplot[OilBlue, mark=*, mark size=2.2pt, line width=1.3pt] table[x=i, y=ocf]{data/clean/profiles/fin_001366.dat};
\addplot[only marks, OilBlue, mark=*, mark size=3.2pt] table[x=x, y=ocf]{data/clean/profiles/finh_001366.dat};
\legend{经营活动现金流净额（右轴，亿元，蓝点）}
\end{axis}
\end{tikzpicture}

\captionof{figure}[播恩集团（001366）近年主要财务指标]{播恩集团（001366）近年主要财务指标（左上：营业收入与归母净利润；右上：ROE 与资产负债率；下：货币资金、有息负债为左轴柱状，经营活动现金流净额为其右轴的蓝点折线。
2021---2025 年为年报口径，2026H1 为 2026 年半年报/半年末，与其前一年报之间留出间隔、以斜纹或空心点区分；半年报数据未年化）}
\end{minipage}
\end{center}

\pfl{财务分析} 2026 年 6 月末资产负债率 38.1\%，较 2025 年末上升 5.2 个百分点（样本中位数 52.2\%，所处三级行业内按由低到高排第\zhnumber{1}位）。 2026 年 6 月末货币资金 2.60 亿元、短期债务（短期借款与一年内到期非流动负债）2.37 亿元、长期债务 1.31 亿元，有息负债合计 3.68 亿元，现金短债比 1.10 倍——覆盖充分。 净有息负债 1.08 亿元，相当于其 2026 年上半年营业收入的 16\%。 流动比率 2.08、速动比率 1.62，短期指标尚可。 2026 年上半年经营活动现金流净额 -1.39 亿元，净利润为正但经营现金流为负，利润的现金含量偏低。 2026 年 6 月末在建工程 0.63 亿元，较上年末减少 0.06 亿元，在建投入在收缩。 综合看，现金短债比 1.10 倍，短期偿付压力可控；资产负债率低于样本中位数 14.1 个百分点；有息负债中短期债务占比更高，期限结构仍以滚动续作为主。

\pfl{重大战略转型} 公司的战略动作可归为以下几类：并购与重组（3 项）：2026 年 9 月，本次收购可新增牛、羊等反刍动物饲料产品线。 公告口径上，近三年（2023 年 9 月---2026 年 9 月）公司披露重大重组与并购类公告 9 条、对外投资与转型类公告 3 条、回购与股权激励类公告 1 条，合计公告 373 条。 第一大行业仍是“饲料行业”，收入占比 100.0\%，与 2021 年年报相比没有方向性变化。 综合判断，报告期内公司有 12 条与战略相关的公告落地（重大重组与并购 9 条、对外投资与转型 3 条），资本运作与主业调整之间存在直接关联。

\begin{center}\small\setlength{\tabcolsep}{4pt}
\captionof{table}[播恩集团（001366）历史融资与资本运作]{播恩集团（001366）历史融资与资本运作（据公司公告与公开报道整理；金额为披露值，含首次公开发行、再融资、债券、银行借款与重整投资等）}
\begin{adjustbox}{max width=\textwidth}
\begin{tabular}{@{}p{1.45cm}p{2.55cm}p{4.05cm}p{4.95cm}@{}}
\toprule
时间 & 方式 & 规模 & 用途与说明 \\
\midrule
2023-03 & IPO（深交所主板） & 发行 4,\allowbreak{}035 万股新股，\allowbreak{}发行价 9.32 元/\allowbreak{}股，\allowbreak{}募集资金总额 3.76 亿元 & 赣州年产 24 万吨饲料项目、\allowbreak{}浙江播恩年产 12 万吨饲料项目、\allowbreak{}播恩生物健康产业基地-维生素复合预混料项目、\allowbreak{}重庆八维生物年产 12 万吨饲料项目、\allowbreak{}研发中心建设项目；发行市盈率（发行后股本）22.97 倍；截至 2026 年中报累计投入 1.22 亿元 \\
2026-05-26 & 签署股权收购条款清单（Term Sheet） & 支付交易保证金 700.00 万元 & 拟收购希杰国际控股有限公司持有的希杰（沈阳）、\allowbreak{}希杰（天津）、\allowbreak{}希杰（聊城）饲料有限公司各 100\% 股权；暂定交易价格 1.25 亿元；约定若截至 2026 年 8 月 31 日未能签署正式股权转让协议等情形 \\
2026-09-14 & 现金收购股权（正式方案） & 股权转让款合计 1.12 亿元（较 5 月披露的 1.25 亿元暂定价格下调约 1,\allowbreak{}250.00 万元） & 由全资子公司广东播恩投资有限公司收购希杰（沈阳）、\allowbreak{}希杰（天津）、\allowbreak{}希杰（聊城）饲料有限公司各 100\% 股权；三家公司将纳入合并报表范围；2025 年三家标的合计营收约 6.24 亿元 \\
\bottomrule
\end{tabular}
\end{adjustbox}
\end{center}

\pfl{历史融投资情况} 从现金流量表看，2025 年公司取得借款收到的现金 2.40 亿元、偿还债务支付的现金 1.35 亿元，筹资活动现金流净额 0.98 亿元（2023 年为取得借款 0.00 亿元、偿还债务 0.16 亿元）。 2026 年上半年筹资活动现金流净额为 1.35 亿元（取得借款 1.87 亿元、偿还债务 0.75 亿元，未年化），仍处于净融入状态。 2025 年分配股利、利润或偿付利息支付的现金合计 0.07 亿元。 公开披露的融资记录共 3 笔，工具类型以 IPO 为主；金额与用途见下表。 其中规模最大的两笔为：2023 年 3 月，IPO，发行 4,035 万股新股，发行价 9.32 元/股，募集资金总额 3.76 亿元；2026 年 9 月，现金收购股权，股权转让款合计 1.12 亿元（较 5 月披露的 1.25 亿元暂定价格下调约 1,250.00 万元）。 股本方面，近三年披露股本变动 5 次（A 股上市 1 次、限售股份上市 1 次）；总股本自 2021 年末以来未发生变化（1.61 亿股）。 分红方面，近三年现金分红 3 次、合计每股派现 0.35 元，最近一次实施在 2026 年 6 月。 综合看，公司融资结构上股权与债务融资并举；2025 年筹资活动现金流净额 0.98 亿元，年度内仍为净融入，与前述经营与投资现金流的方向基本一致。

\pfl{未来潜在融资需求分析} 2026 年 6 月末货币资金 2.60 亿元，而短期债务 2.37 亿元（现金短债比 1.10 倍），2026 年上半年经营现金流净额 -1.39 亿元。账面货币资金可以覆盖短期债务，缺口不体现在静态偿付层面。 公司 2026 年 6 月末资产负债率 38.1\%、2026 年上半年 ROE 1.3\%（未年化），账面杠杆不高、盈利能力一般，短期内没有必须依靠外部融资才能维持的经营缺口；后续融资更可能指向技改、扩产或产业并购。 公开披露的后续资金安排线索包括：资本开支安排：2026 年 8 月，募集资金结余 2.04 亿元；2023 年 3 月，IPO 募资改道并购受到媒体关注：公司 2023 年 3 月上市募得 3.19 亿元。 资金需求还要叠加在建工程：期末在建工程 0.63 亿元、2026 年上半年购建固定资产等支出 0.37 亿元（未年化），这部分支出在周期底部通常会被主动放缓。 静态偿付不存在缺口，后续融资更可能服务于产能或并购安排，而不是填补流动性窟窿。




\subsubsection{正虹科技（000702）}
\label{co:000702}

\pfl{公司基本情况} 正虹科技成立于 1994 年，湖南岳阳的老牌饲料企业，前身可追溯至 1986 年的国营屈原农场饲料厂，1997 年上市并有“中国饲料第一股”之称，2022 年控制权由屈原农垦转至岳阳国资观盛投资。 公司于 1997-03-18 在深交所主板首发上市，注册资本 3.47 亿元，总股本 3.47 亿股。 截至 2026 年 10 月 12 日收盘价 5.77 元，对应总市值 20.0 亿元，近一年-19.8\%，总市值在三级行业“畜禽饲料”的 11 家公司中按由大到小排第\zhnumber{11}位，近一年涨跌幅低于全样本中位数（-10.4\%）9.4 个百分点。 按 2026 年半年报披露的行业构成，收入占比前三位为饲料行业 89.8\%；饲养行业 8.8\%；其他(补充) 1.2\%。 股权结构方面，控股股东为岳阳观盛投资发展有限公司（持股 23.08\%）。

\pfl{股权结构与控制权} 按 2026 年 6 月 30 日的股东名册，前三大股东为岳阳观盛投资发展有限公司 持股 23.08\%；岳阳观盛农业科技有限责任公司 持股 11.64\%；岳阳市屈原农垦有限责任公司 持股 7.70\%，股东户数 26,283 户。 股权结构的另一条线索是，持股 23.08\%（7,999.04 万股、限售）；2022 年 11 月，观盛农业所持股份中 3,227.34 万股质押，屈原农垦所持股份中 1,330 万股被冻结。

\pfl{经营要点} 经营层面，2026 年上半年，累计销售生猪 5.32 万头、同比下降 15.61\%，实现销售收入 5,013.21 万元、同比下降 39.93\%。

\pfl{所处行业的基本面} 饲料的供给即产能（门槛低、产能利用率偏低），需求直接来自下游存栏，利润则由玉米、豆粕等原料价格与成品售价的价差决定（第 \ref{sec:framework} 章的供需框架）。 行业层面的最新情况是：2026 年上半年，全国工业饲料总产量 16,850 万吨、同比增长 3.9\%。

\pfl{竞争优势} 把该公司放回同类公司比较，2026 年上半年营业收入 5.7 亿元、在三级行业“畜禽饲料”的 11 家公司中排第\zhnumber{10}位；毛利率 5.3\%、与行业内中位数基本持平；半年 ROE -6.2\%（未年化）、高于行业内中位数（-6.25\%）；资产负债率 54.2\%、低于行业内中位数（62.1\%）。 相对优势方面，饲料销售贡献营业收入 87.87\%，在 2026 年上半年生猪销量与销售收入双双下滑。 优势在当期体现为高于行业中位数的资本回报，但优势能否持续仍取决于周期与自身产能投放节奏。

\pfl{历史股价} 以 2021 年 1 月为起点、2026 年 10 月为终点，股价由 5.23 元变为 5.77 元，累计 +10.3\%，区间最高 9.49 元（2024 年 10 月）、最低 4.08 元（2024 年 2 月）；当前价相当于区间最高价的 60.8\%，为区间最低价的 1.41 倍。 月度收益的年化波动率 42.1\%，高于全样本中位数 37.3\%，在三级行业“畜禽饲料”的 11 家公司中波动率由高到低排第\zhnumber{4}位。 把股价阶段与公司事件对齐看，2024 年 10 月 至 2026 年 6 月股价下跌-47.7\%，同期（2024 年 12 月）岳阳市屈原农垦有限责任公司持股 7.70\%（2,667.58 万股），其中冻结 1,330.00 万股。 把这些阶段与公司事件对齐后可以看到，几轮大涨大跌多由公司自身的资本运作与风险事件触发，行业景气只解释了其中一部分。

\begin{center}
\begin{minipage}{\textwidth}\centering
\begin{tikzpicture}
\begin{axis}[
  width=12.6cm, height=3.9cm, scale only axis,
  xmin=2020.922, xmax=2026.972, ymin=3.79, ymax=10.15,
  xtick={2021,2022,2023,2024,2025,2026},
  yearticks,
  yticklabel style={font=\scriptsize},
  ylabel={元/股}, ylabel style={font=\scriptsize},
  grid=major, grid style={LightGray, line width=0.3pt},
  axis line style={InkGray, line width=0.5pt},
  every axis plot/.append style={line width=0.9pt},
]
\addplot[AgriGreen] table[x=x,y=y]{data/clean/profiles/px_000702.dat};
\addplot[SoftRed, dashed, line width=0.7pt] table[x=x,y=y]{data/clean/profiles/pxzz_000702.dat};
\addplot[only marks, mark=*, mark size=1.1pt, SoftRed] table[x=x,y=y]{data/clean/profiles/pxzz_000702.dat};
\end{axis}
\end{tikzpicture}

\captionof{figure}[正虹科技（000702）股价与周期骨架]{正虹科技（000702）月度收盘价（元/股）与 ZigZag 周期骨架（阈值 25\%，圆点为周期转折点）}
\end{minipage}
\end{center}

\pfl{过去周期分析} 以 25\% 的 ZigZag 阈值划分，样本区间内股价形成 2 个主升段与 2 个主跌段。 最大主跌段为 2024 年 10 月 至 2026 年 6 月，20 个月-47.7\%；最大主升段出现在 2024 年 2 月 至 2024 年 10 月，8 个月+132.6\%。 最近一次转折出现在 2026 年 6 月（下跌段自 2024 年 10 月起，20 个月-47.7\%），当前处于该段的延续之中。 公司股价较申万农林牧渔指数提前约 27 个月见底，是板块中较早定价周期反转的公司之一。 饲料夹在原料成本与养殖需求之间（第 \ref{sec:feed} 节），量的部分取决于存栏、利的部分取决于原料与成品的价差，公司 2026 年 6 月末资产负债率 54.2\%、半年 ROE -6.2\%（未年化），在本轮下行中属于随行业同步调整的一类。 综合区间形态看，该股单段涨跌幅度偏大、以趋势段为主（最大主升 +132.6\%、最大主跌 -47.7\%），年化波动率 42.1\% 与全样本中位数 37.3\% 相比波动与样本中位数接近。

\pfl{盈利情况} 2026 年上半年营业收入 5.69 亿元（同比 +14.2\%），归母净利润 -0.21 亿元，亏损同比扩大 0.13 亿元。 以年报为趋势对照，2023---2025 年营业收入由 12.41 亿元变为 10.19 亿元（两年年均复合增速 -9.4\%），归母净利润由 -1.40 亿元变为 -0.14 亿元。 按半年报口径，2026 年上半年的 ROE 为 -6.2\%（未年化）、销售净利率 -4.9\%、毛利率 5.3\%，ROE 在“畜禽饲料”11 家公司中按数值由高到低排第\zhnumber{5}位。 同期全样本半年 ROE 中位数 0.75\%，公司与之相差 6.92 个百分点（弱于中位数公司）。 针对上述压力，公司披露的应对举措包括：2026 年，饲料原料贸易业务适度调整以降低采购成本。 综合看，公司仍处亏损区间、亏损同比扩大，半年 ROE 在三级行业内按由高到低排第\zhnumber{5}位，单位盈利能力的修复依赖行业景气与自身成本端的改善，报表弹性主要体现为价格与销量的方向性变化。

\begin{center}\small\setlength{\tabcolsep}{3.4pt}
\captionof{table}[正虹科技（000702）关键财务指标]{正虹科技（000702）关键财务指标（合并报表口径；两个半年报期均未年化；现金短债比 $=$ 货币资金 $\div$ 短期债务）}
\begin{adjustbox}{max width=\textwidth}
\begin{tabular}{@{}lrrrrr@{}}
\toprule
指标 & 2023 年 & 2024 年 & 2025 年 & 2025 年半年报 & 2026 年半年报 \\
\midrule
营业收入（亿元） & 12.41 & 11.21 & 10.19 & 4.99 & 5.69 \\
归母净利润（亿元） & -1.40 & -0.65 & -0.14 & -0.08 & -0.21 \\
毛利率（\%） & 4.0 & 5.1 & 7.1 & 7.4 & 5.3 \\
ROE（\%） & -46.6 & -16.5 & -4.0 & -2.2 & -6.2 \\
资产负债率（\%） & 53.9 & 57.5 & 47.0 & 47.4 & 54.2 \\
经营活动现金流净额（亿元） & 0.61 & 0.11 & -0.31 & -0.18 & -0.63 \\
货币资金（亿元） & 2.49 & 2.67 & 0.66 & 0.92 & 0.60 \\
有息负债合计（亿元） & 2.84 & 2.82 & 1.41 & 1.33 & 1.97 \\
现金短债比（倍） & 0.91 & 0.96 & 0.51 & 0.71 & 0.33 \\
\bottomrule
\end{tabular}
\end{adjustbox}
\end{center}

\begin{center}
\begin{minipage}{\textwidth}\centering
\begin{tikzpicture}
\begin{groupplot}[
  group style={group size=2 by 1, horizontal sep=0.60cm},
  width=6.85cm, height=4.60cm,
  tick label style={font=\tiny},
  title style={font=\scriptsize\heiti},
  yticklabel style={font=\tiny},
  ylabel style={font=\tiny},
  grid=major, grid style={LightGray, line width=0.3pt},
  axis line style={InkGray, line width=0.5pt},
  enlarge x limits={abs=0.8},
  legend style={font=\scriptsize, at={(0.5,-0.20)}, anchor=north,
                legend columns=2, draw=none, fill=none, column sep=6pt},
]
\nextgroupplot[
  ybar, bar width=4pt,
  ymin=-3.90, ymax=16.30,
  xtick={0,1,2,3,4,5.7},
  xticklabels={2021,2022,2023,2024,2025,2026H1},
  ylabel={亿元},
]
\addplot[fill=AgriGreen!75, draw=AgriGreen, bar shift=-2.6pt] table[x=i, y=rev]{data/clean/profiles/fin_000702.dat};
\addplot[fill=SoftRed!60, draw=SoftRed, bar shift=2.6pt] table[x=i, y=ni]{data/clean/profiles/fin_000702.dat};
\addplot[fill=AgriGreen!30, draw=AgriGreen, pattern=north east lines, bar shift=-2.6pt] table[x=x, y=rev]{data/clean/profiles/finh_000702.dat};
\addplot[fill=SoftRed!20, draw=SoftRed, pattern=north east lines, bar shift=2.6pt] table[x=x, y=ni]{data/clean/profiles/finh_000702.dat};
\legend{营业收入（年/半年）, 归母净利润（年/半年）}
\nextgroupplot[
  ymin=-58.9, ymax=98.6,
  xtick={0,1,2,3,4,5.7},
  xticklabels={2021,2022,2023,2024,2025,2026H1},
  ylabel={\%},
]
\addplot[AgriGold, mark=*, mark size=1.3pt] table[x=i, y=roe]{data/clean/profiles/fin_000702.dat};
\addplot[OilBlue, mark=square*, mark size=1.3pt] table[x=i, y=debt]{data/clean/profiles/fin_000702.dat};
\addplot[only marks, AgriGold, mark=*, mark size=2.0pt] table[x=x, y=roe]{data/clean/profiles/finh_000702.dat};
\addplot[only marks, OilBlue, mark=square*, mark size=2.0pt] table[x=x, y=debt]{data/clean/profiles/finh_000702.dat};
\legend{ROE, 资产负债率}
\end{groupplot}
\end{tikzpicture}
\begin{tikzpicture}
\begin{axis}[
  name=finbot,
  width=13.75cm, height=3.10cm, scale only axis,
  ymin=-0.48, ymax=5.43,
  xtick={0,1,2,3,4,5.7},
  xticklabels={2021,2022,2023,2024,2025,2026H1},
  tick label style={font=\tiny},
  yticklabel style={font=\tiny},
  ylabel={亿元}, ylabel style={font=\tiny},
  ybar, bar width=4pt,
  grid=major, grid style={LightGray, line width=0.3pt},
  axis line style={InkGray, line width=0.5pt},
  enlarge x limits={abs=0.8},
  legend style={font=\scriptsize, at={(0.5,-0.26)}, anchor=north,
                legend columns=2, draw=none, fill=none, column sep=10pt},
]
\addplot[fill=AgriGreen!55, draw=AgriGreen, bar shift=-3.0pt] table[x=i, y=cash]{data/clean/profiles/fin_000702.dat};
\addplot[fill=SoftRed!45, draw=SoftRed, bar shift=3.0pt] table[x=i, y=ibd]{data/clean/profiles/fin_000702.dat};
\addplot[fill=AgriGreen!25, draw=AgriGreen, pattern=north east lines, bar shift=-3.0pt] table[x=x, y=cash]{data/clean/profiles/finh_000702.dat};
\addplot[fill=SoftRed!20, draw=SoftRed, pattern=north east lines, bar shift=3.0pt] table[x=x, y=ibd]{data/clean/profiles/finh_000702.dat};
\legend{货币资金（左轴）, 有息负债（左轴）}
\end{axis}
\begin{axis}[
  at={(finbot.south east)}, anchor=south east,
  width=13.75cm, height=3.10cm, scale only axis,
  axis y line*=right, axis x line*=none, xtick=\empty,
  ymin=-1.78, ymax=1.18,
  tick label style={font=\tiny},
  yticklabel style={font=\tiny}, scaled y ticks=false,
  legend style={font=\scriptsize, at={(0.5,-0.44)}, anchor=north,
                legend columns=1, draw=none, fill=none},
]
\addplot[OilBlue, mark=*, mark size=2.2pt, line width=1.3pt] table[x=i, y=ocf]{data/clean/profiles/fin_000702.dat};
\addplot[only marks, OilBlue, mark=*, mark size=3.2pt] table[x=x, y=ocf]{data/clean/profiles/finh_000702.dat};
\legend{经营活动现金流净额（右轴，亿元，蓝点）}
\end{axis}
\end{tikzpicture}

\captionof{figure}[正虹科技（000702）近年主要财务指标]{正虹科技（000702）近年主要财务指标（左上：营业收入与归母净利润；右上：ROE 与资产负债率；下：货币资金、有息负债为左轴柱状，经营活动现金流净额为其右轴的蓝点折线。
2021---2025 年为年报口径，2026H1 为 2026 年半年报/半年末，与其前一年报之间留出间隔、以斜纹或空心点区分；半年报数据未年化）}
\end{minipage}
\end{center}

\pfl{财务分析} 2026 年 6 月末资产负债率 54.2\%，较 2025 年末上升 7.2 个百分点（样本中位数 52.2\%，所处三级行业内按由低到高排第\zhnumber{2}位）。 2026 年 6 月末货币资金 0.60 亿元、短期债务（短期借款与一年内到期非流动负债）1.81 亿元、长期债务 0.16 亿元，有息负债合计 1.97 亿元，现金短债比 0.33 倍——覆盖偏紧。 净有息负债 1.37 亿元，相当于其 2026 年上半年营业收入的 24\%。 流动比率 1.00、速动比率 0.66，短期偿债能力处于正常区间。 2026 年上半年经营活动现金流净额 -0.63 亿元，经营现金流与利润同时为负，现金消耗较快。 综合看，账面货币资金对有息负债与短期债务的覆盖不足，滚动偿付对新增融资的依赖度较高；资产负债率高于样本中位数 2.0 个百分点；有息负债中短期债务占比更高，期限结构仍以滚动续作为主。 公司披露的股东与质押风险为：2024 年 12 月，岳阳市屈原农垦有限责任公司持股 7.70\%（2,667.58 万股），其中冻结 1,330.00 万股。

\pfl{重大战略转型} 近三年公司的战略动作集中在：融资与资本运作（1 项）：2022 年 6 月，公司引入岳阳国资观盛投资作为控股股东，并筹划向观盛投资非公开发行股票募集 3 亿多元用于补充流动资金；其他动作（1 项）：2023 年生猪生产量同比减少 49.39\%。 公告口径上，近三年（2023 年 9 月---2026 年 9 月）公司披露重大重组与并购类公告 5 条、对外投资与转型类公告 3 条、回购与股权激励类公告 0 条，合计公告 309 条。 第一大行业“饲料行业”的收入占比由 2021 年年报的 61.9\% 变为 89.8\%（28.0 个百分点），收入结构出现再平衡。 综合判断，报告期内公司有 8 条与战略相关的公告落地（重大重组与并购 5 条、对外投资与转型 3 条），资本运作与主业调整之间存在直接关联。

\begin{center}\small\setlength{\tabcolsep}{4pt}
\captionof{table}[正虹科技（000702）历史融资与资本运作]{正虹科技（000702）历史融资与资本运作（据公司公告与公开报道整理；金额为披露值，含首次公开发行、再融资、债券、银行借款与重整投资等）}
\begin{adjustbox}{max width=\textwidth}
\begin{tabular}{@{}p{1.45cm}p{2.55cm}p{4.05cm}p{4.95cm}@{}}
\toprule
时间 & 方式 & 规模 & 用途与说明 \\
\midrule
1997-03 & IPO（深交所） & 发行社会公众股 5,\allowbreak{}000 万股（含职工内部股 500 万股），\allowbreak{}发行价 7.06 元/\allowbreak{}股 & 饲料产能建设；中国证监会证监发(1997)44 号、\allowbreak{}45 号文批准；上市时湖南省岳阳市屈原农垦集团公司代表国家持有 61.59\% 股权 \\
2022-06-17 & 控制权安排（股份协议转让 + 表决权委托 + 定增） & 屈原农垦向观盛农业转让 4,\allowbreak{}034.18 万股（占发行前总股本 15.13\%）；屈原农垦将其剩余 2,\allowbreak{}667.58 万股（10.00\%）表决权委托给观盛投资；公司向观盛投资非公开发行不超过 7,\allowbreak{}999.04 万股 & 巩固控制权并补充流动资金；收购报告书披露变动数量合计 1.47 亿股；认购完成后观盛投资合计将支配公司 42.41\% 的表决权 \\
2022-06-17 & 发行对象与关联方情况 & 君泰控股为公司第一大股东观盛农业的少数股东；君泰科技及相关主体系君泰控股控制的企业 & 发行人实际控制人 2022 年变更后，\allowbreak{}董事、\allowbreak{}监事、\allowbreak{}高级管理人员变动较大 \\
2023-07 & 向特定对象发行股票（定增） & 向控股股东观盛投资发行不超过 7,\allowbreak{}999.04 万股 & 全部用于补充流动资金 \\
\bottomrule
\end{tabular}
\end{adjustbox}
\end{center}

\pfl{历史融投资情况} 从现金流量表看，2025 年公司取得借款收到的现金 2.66 亿元、偿还债务支付的现金 4.10 亿元、吸收权益性投资 0.07 亿元，筹资活动现金流净额 -1.56 亿元（2023 年为取得借款 3.15 亿元、偿还债务 4.80 亿元）。 2026 年上半年筹资活动现金流净额为 0.60 亿元（取得借款 0.54 亿元、偿还债务 0.05 亿元，未年化），仍处于净融入状态。 2025 年分配股利、利润或偿付利息支付的现金合计 0.07 亿元。 公开披露的融资记录共 4 笔，工具类型以 IPO、定向增发为主；金额与用途见下表。 其中规模最大的两笔为：2023 年 7 月，定向增发，向控股股东观盛投资发行不超过 7,999.04 万股。 股本方面，近三年披露股本变动 4 次（增发新股上市 1 次）；总股本由 2021 年末的 2.67 亿股变为 3.47 亿股（+30.00\%）。 分红方面，近三年未实施现金分红，在全部样本中属于分红能力偏弱的一端。 近三年“再融资”类公告 3 条，涉及定向增发。 综合看，公司融资结构上以债务滚动为主、股权融资为补充；2025 年筹资活动现金流净额 -1.56 亿元，融资节奏仍以维持存量债务周转为主，与前述经营与投资现金流的方向基本一致。

\pfl{未来潜在融资需求分析} 2026 年 6 月末货币资金 0.60 亿元，而短期债务 1.81 亿元（现金短债比 0.33 倍），2026 年上半年经营现金流净额 -0.63 亿元。按“短期债务减去账面货币资金”的滚动偿付口径，静态缺口约 1.21 亿元；若再计入上半年亏损的年化额（0.41 亿元），维持一年正常经营所需的外部资金约 1.21---1.63 亿元。 公司 2026 年上半年归母净利润为负（-0.21 亿元），但 6 月末资产负债率 54.2\% 尚未进入第一档（亏损且负债率 $>$65\%）的区间，当前融资需求以补充流动资金、支撑低谷期经营性支出为主。 公开披露的后续资金安排线索包括：股权融资线索：2022 年 6 月，后者正筹划公司控制权变更事项，并同步筹划向观盛投资非公开发行股票；2023 年 7 月，公司向观盛投资的定增（不超过 3.38 亿元）认购资金到位。 综合看，公司的外部资金需求以补充营运资金为主、项目投入为辅，滚动偿付口径下的静态缺口约 1.21 亿元，融资节奏将受制于银行续贷意愿与股权融资窗口。






\section{动物保健：周期中唯一稳定的“卖水人”}
\label{sec:vet}

动物保健板块 \num{9} 家公司（全部归入申万三级“动物保健Ⅲ”），
总市值合计约 \num{510} 亿元。按 2026 年半年报口径，仅 \num{1} 家亏损，
半年 ROE 中位数 \num{2.95}\%（未年化），
\hl{资产负债率中位数 \num{30.2}\%，为全部农业二级行业中最低}，
超过 \num{65}\% 的只有 \num{1} 家。

\begingroup\scriptsize\setlength{\tabcolsep}{2pt}
\begin{longtable}{@{}llrrrrrrr@{}}
\caption{农业上市公司分行业明细（动物保健Ⅲ、动物保健；市值时点 2026-10-12，总市值口径；财务为 2026 年半年报，ROE 未年化，公告计数窗口 2023-09---2026-10）}\label{tab:comp-vet}\\
\toprule
代码 & 公司简称 & 总市值（亿元） & 营收（亿元） & 同比（\%） & 净利（亿元） & ROE（\%） & 再融资公告 & 重组公告 \\
\midrule
\endfirsthead
\multicolumn{9}{l}{\small（续）}\\
\toprule
代码 & 公司简称 & 总市值（亿元） & 营收（亿元） & 同比（\%） & 净利（亿元） & ROE（\%） & 再融资公告 & 重组公告 \\
\midrule
\endhead
\bottomrule
\endlastfoot
600201 & 生物股份 & 133.3 & 7.6 & 22.1 & 1.04 & 1.9 & 43 & 0 \\
300119 & 瑞普生物 & 76.7 & 17.2 & 0.8 & 2.04 & 4.3 & 4 & 13 \\
600195 & 中牧股份 & 67.9 & 31.7 & 13.5 & 0.50 & 0.9 & 3 & 2 \\
688526 & 科前生物 & 56.8 & 3.7 & -23.3 & 1.16 & 2.8 & 0 & 0 \\
002688 & 金河生物 & 39.9 & 14.5 & 4.1 & 1.31 & 5.7 & 26 & 1 \\
603566 & 普莱柯 & 35.4 & 4.9 & -12.2 & 0.77 & 3.0 & 4 & 2 \\
688098 & 申联生物 & 33.7 & 1.1 & -8.1 & -0.19 & -1.4 & 3 & 2 \\
300871 & 回盛生物 & 33.4 & 8.6 & 5.2 & 2.06 & 9.1 & 46 & 2 \\
002868 & 绿康生化 & 33.1 & 2.5 & -15.7 & 0.07 & 10.9 & 8 & 23 \\
\end{longtable}
\endgroup


\subsection{为什么动保能在下行周期中盈利}
\label{sec:vet-mechanism}

动保的需求由\hl{存栏量而非价格}决定：即使猪价跌破成本线，
养殖场也不会减少疫苗接种与保健投入，因为疫病造成的损失远大于药费。
这使动保企业的收入与生猪、家禽存栏量挂钩，
而成本端（原料药、发酵产能）与农产品价格无关。
\emph{在农业下行周期中，动保是唯一同时具备“需求稳定 + 成本独立”特征的环节}——
本报告借用产业内的说法称之为“卖水人”（不承担价格风险、只分享产量增长）。

板块盈利结构印证了这一点：\textbf{回盛生物}（\num{2.06} 亿元、半年 ROE \num{9.06}\%）、
\textbf{瑞普生物}（\num{2.04} 亿元、\num{4.27}\%）、
\textbf{金河生物}（\num{1.31} 亿元、\num{5.74}\%）三家合计利润 \num{5.4} 亿元，
占板块 \num{8.8} 亿元利润总额的 \num{61.8}\%；
\textbf{生物股份}的利润由上年同期的 \num{0.60} 亿元增至 \num{1.04} 亿元，
而 \textbf{科前生物}由 \num{2.20} 亿元降至 \num{1.16} 亿元。
\hl{科前生物与普莱柯以猪用疫苗为主，其客户正是牧原、温氏等逆势扩张的龙头}——
龙头企业在周期底部提升出栏量，反而增加了对疫苗的需求，
这是动保与养殖景气度“脱钩”的微观基础。

板块内唯一的亏损是 \textbf{申联生物}（\num{-0.19} 亿元），
主要受口蹄疫疫苗招标价格下行影响；
\textbf{绿康生化}（\num{0.07} 亿元、半年 ROE \num{10.91}\%）虽然盈利，
但其业务重心已转向微生物发酵与兽药原料药，\hl{不再是纯粹的动保公司}，
且 2026 年 6 月末资产负债率 \num{90.1}\%，是板块内唯一的高负债公司。
即便是处于价格竞争最激烈的口蹄疫疫苗市场，
板块的亏损面（\num{11.1}\%）也远低于饲料（\num{66.7}\%）与养殖业（\num{68.2}\%）。

\subsection{后非洲猪瘟时代的结构性变化}
\label{sec:vet-africa}

非洲猪瘟疫情（2018 年 8 月起，见第 \ref{sec:hog-shock} 节）以来，
养殖行业的生物安全标准被永久性抬高：
\hl{“低猪价 + 高防疫投入”成为新常态}。
这带来两个结构性变化：

\begin{enumerate}[leftmargin=1.5em, itemsep=2pt]
  \item \textbf{需求从“治疗”转向“预防”}：疫苗与诊断试剂占比提升，
        化药占比下降。板块内以疫苗为主的公司（科前生物、普莱柯、生物股份）
        ROE 均高于以化药为主的公司。
  \item \textbf{行业集中度提升}：中小养殖场因生物安全投入不足而退出，
        动保客户结构随之向大型集团场集中。
        \hl{这使动保企业获得了对下游的议价能力}，
        在 2024---2025 年养殖端普遍亏损时仍能维持毛利率。
\end{enumerate}

从融资行为看，动保板块近三年公司均“再融资相关公告” \num{15.2} 条（中位数 \num{4} 条），
明显低于养殖（\num{20.5} / \num{13.5}）与饲料（\num{17.6} / \num{9}）。
\hl{盈利稳定 + 低负债 + 低融资依赖，使动保成为农业板块中资产负债表最健康的子行业}。
其中 \textbf{绿康生化}（\num{23} 条重组并购公告）与 \textbf{瑞普生物}（\num{13} 条）
的公告密度显著高于同业，反映出这两家在主动进行外延式扩张或业务转型。

\begin{keybox}[动物保健的画像结论]
\normalsize
\textbf{① 唯一的“逆周期”子行业}：半年 ROE 中位数 \num{2.95}\%、
亏损面 \num{11.1}\%、资产负债率中位数 \num{30.2}\%，三项指标均为农业最优。\par\vspace{3pt}
\textbf{② 需求的锚是存栏量，不是价格}：龙头企业在下行周期扩张产能，
反而推升了疫苗需求，这是动保与猪周期“脱钩”的机制。\par\vspace{3pt}
\textbf{③ 关注外延扩张的两家}：绿康生化转型原料药、瑞普生物持续并购，
\hl{它们代表动保行业从“单一疫苗”走向“疫苗 + 化药 + 诊断”平台化的路径}。\par\vspace{3pt}
\textbf{④ 稳定性再次得到验证。}板块 \num{9} 家公司 2026 年上半年合计归母净利润
\num{8.8} 亿元（上年同期 \num{9.2} 亿元），仅 \num{1} 家亏损，
\hl{是各子行业中波动最小的一类}。
\end{keybox}

\subsection{上市公司画像（\num{9} 家）}
\label{sec:co-vet}

本节给出动物保健行业全部 \num{9} 家上市公司的画像，按总市值降序排列。
数据口径与第 \ref{sec:comp-overview} 节一致：\hl{财务数据统一采用 2026 年半年报}
（合并报表口径、未年化，资产负债率为 2026 年 6 月末），
同时列出 2025 年半年报作为同口径对照、2023---2025 年年报作为趋势对照；
股价为月度收盘价，最后一个交易日为 2026 年 10 月 12 日；周期分段使用与第一部分相同的
ZigZag 算法，阈值取 25\%（个股波动大于商品价格，故高于生猪现货的 20\%）；
公告计数窗口为 2023 年 9 月 1 日至 2026 年 10 月 12 日。

画像正文由\hl{逐家检索的公开资料}写成：每条事实均可追溯至公开来源
（公司沿革、股权与控制权、经营要点、所处行业的基本面、竞争优势、融资历史、
风险事件、战略动作、融资需求线索），
正文依事实展开，各家公司的段落结构与重点并不相同。
每家公司在统一顺序下展开：基本情况与主营结构、股权结构与控制权、经营要点、
所处行业的基本面（供需）、与同类公司比较的竞争优势、历史股价、过去周期分析、
盈利情况、财务分析、重大战略转型、历史融投资情况、未来潜在融投资需求；
\hl{市值一律采用总市值口径}，股权结构取最新披露报告期前三大股东与股东户数，
行业基本面与竞争优势以逐家检索的公开资料为依据，并与三级行业内的同口径指标横向比较。
每家公司配图两张：
股价与行业指数（申万农林牧渔）的对比图，以及近年主要财务指标的组合图
（上排：营业收入与归母净利润、ROE 与资产负债率；下方宽带面板：货币资金与有息负债为左轴柱状、
经营活动现金流净额为其右轴的蓝点折线），并给出关键财务指标表与历史融资/资本运作表。
未来融资需求按两个口径粗略度量：\hl{短期债务缺口 $=$ 期末短期债务 $-$ 货币资金}
与\hl{年化经营缺口 $=$ 半年归母亏损 $\times 2$}；
前者衡量滚动偿付压力，后者衡量维持当期盈利水平的现金消耗，
两者都是压力情形下的静态估算，不构成对融资规模或时点的预测。

% 农业上市公司画像：动物保健Ⅲ、动物保健（共 9 家，按总市值降序）
% 由 scripts/make_company_profiles.py 自动生成，请勿手工修改

\subsubsection{生物股份（600201）}
\label{co:600201}

\pfl{公司基本情况} 内蒙古呼和浩特的兽用生物制品企业，全称金宇生物技术股份有限公司，1992 年 12 月设立、1999 年 1 月 15 日在上海证券交易所上市；主营兽用生物制品的研发、生产与销售，产品涵盖猪、禽、反刍和宠物类四大系列百余种动物疫苗，拥有口蹄疫、高致病性禽流感和布鲁氏菌病三大强制免疫疫苗农业农村部定点生产资质。2025 年归母净利润 1.00 亿元、同比下降 8.14\%；2026 年上半年归母净利润 1.04 亿元、同比增长 72.28\%。公司两次再融资（2021 年度、2023 年度）均告终止。 公司于 1999-01-15 在上交所首发上市，注册资本 11.00 亿元，总股本 11.20 亿股。 截至 2026 年 10 月 12 日收盘价 11.90 元，对应总市值 133.3 亿元，近一年+17.1\%，总市值在三级行业“动物保健Ⅲ”的 9 家公司中按由大到小排第\zhnumber{1}位，近一年涨跌幅高于全样本中位数（-10.4\%）27.5 个百分点。 按 2025 年年报披露的行业构成，收入占比前三位为生物制药 92.7\%；其他(补充) 3.9\%；其他收入 3.5\%。 股权结构方面，控股股东为内蒙古金宇生物控股有限公司。

\pfl{股权结构与控制权} 按 2026 年 6 月 30 日的股东名册，前三大股东为内蒙古金宇生物控股有限公司 持股 11.29\%；伟星资产管理(上海)有限公司－宁波梅山保税港区星棋道 持股 3.66\%；上海淳韬投资管理有限公司－淳韵价值一号证券投资私募基 持股 3.63\%，股东户数 41,853 户。 股东结构与控制权方面，公司质押数量为 2,035.00 万股；报告期内公司控股股东或第一大股东及其一致行动人累计质押股份数量占其所持公司股份数量比例未达到 80\%；张翀宇、张竞父女通过直接持有公司 2.62\% 股权、间接控制生物控股所持 11.00\% 股权。

\pfl{经营要点} 经营层面，公司第一大股东内蒙古金宇生物控股有限公司注册资本 2,150.00 万元。

\pfl{所处行业的基本面} 动保的需求由下游存栏规模而非养殖盈利决定，供给端则由产品批文与研发能力设卡，疫苗是其中格局最稳定的细分（第 \ref{sec:framework} 章的供需框架）。 行业层面的最新情况是：前十企业市场份额占 45.46\%；竞争从单一价格转向产品加技术服务的综合价值竞争；2026 年 5 月，农业农村部修订印发生猪产能综合调控实施方案，将全国能繁母猪正常保有量目标由 3,900 万头下调至 3,750 万头。

\pfl{竞争优势} 把该公司放回同类公司比较，2026 年上半年营业收入 7.6 亿元、在三级行业“动物保健Ⅲ”的 9 家公司中排第\zhnumber{5}位；毛利率 53.9\%、高于行业内中位数 14.2 个百分点；半年 ROE 1.9\%（未年化）、低于行业内中位数（2.95\%）；资产负债率 16.5\%、低于行业内中位数（30.2\%）。 优势在报表上主要体现为高于行业平均的毛利率，这在周期底部是抵御亏损的主要缓冲。

\pfl{历史股价} 以 2021 年 1 月为起点、2026 年 10 月为终点，股价由 22.77 元变为 11.90 元，累计 -47.7\%，区间最高 22.77 元（2021 年 1 月）、最低 6.26 元（2024 年 8 月）；当前价相当于区间最高价的 52.3\%，为区间最低价的 1.90 倍。 月度收益的年化波动率 37.2\%，低于全样本中位数 37.3\%，在三级行业“动物保健Ⅲ”的 9 家公司中波动率由高到低排第\zhnumber{5}位。 把股价阶段与公司事件对齐看，2024 年 8 月 至 2026 年 1 月股价上涨+188.3\%，同期（2025 年 7 月）原拟向控股股东方非公开发行不超过 8.00 亿元（发行价 8.29 元/股）的计划已于 2025 年 7 月终止；2021 年 1 月 至 2022 年 9 月股价下跌-65.7\%，同期（2022 年 4 月）公司终止 2021 年度非公开发行 A 股股票事项。 整体上看，区间的几处大级别拐点都能在公司的资本运作、股权变动或风险事件上找到对应，股价波动有明显的公司层面驱动，而非单纯的行业 beta。

\begin{center}
\begin{minipage}{\textwidth}\centering
\begin{tikzpicture}
\begin{axis}[
  width=12.6cm, height=3.9cm, scale only axis,
  xmin=2020.922, xmax=2026.972, ymin=5.82, ymax=24.36,
  xtick={2021,2022,2023,2024,2025,2026},
  yearticks,
  yticklabel style={font=\scriptsize},
  ylabel={元/股}, ylabel style={font=\scriptsize},
  grid=major, grid style={LightGray, line width=0.3pt},
  axis line style={InkGray, line width=0.5pt},
  every axis plot/.append style={line width=0.9pt},
]
\addplot[AgriGreen] table[x=x,y=y]{data/clean/profiles/px_600201.dat};
\addplot[SoftRed, dashed, line width=0.7pt] table[x=x,y=y]{data/clean/profiles/pxzz_600201.dat};
\addplot[only marks, mark=*, mark size=1.1pt, SoftRed] table[x=x,y=y]{data/clean/profiles/pxzz_600201.dat};
\end{axis}
\end{tikzpicture}

\captionof{figure}[生物股份（600201）股价与周期骨架]{生物股份（600201）月度收盘价（元/股）与 ZigZag 周期骨架（阈值 25\%，圆点为周期转折点）}
\end{minipage}
\end{center}

\pfl{过去周期分析} 以 25\% 的 ZigZag 阈值划分，样本区间内股价形成 2 个主升段与 3 个主跌段。 最大主升段出现在 2024 年 8 月 至 2026 年 1 月，17 个月+188.3\%；最大主跌段为 2021 年 1 月 至 2022 年 9 月，20 个月-65.7\%。 最近一次转折出现在 2026 年 6 月（下跌段自 2026 年 1 月起，5 个月-35.5\%），当前处于该段的延续之中。 公司股价较申万农林牧渔指数提前约 21 个月见底，是板块中较早定价周期反转的公司之一。 动保的需求由存栏量而非价格决定（第 \ref{sec:vet} 节）：猪价下行阶段龙头养殖企业的逆势扩产反而推升了疫苗与兽药需求，这正是动保股价与猪价“脱钩”的原因，公司 2026 年上半年实现盈利、半年 ROE 1.9\%（未年化），好于全样本中位数（0.75\%），下行周期对其报表的冲击小于同业。 综合区间形态看，该股单段涨跌幅度偏大、以趋势段为主（最大主升 +188.3\%、最大主跌 -65.7\%），年化波动率 37.2\% 与全样本中位数 37.3\% 相比波动与样本中位数接近。

\pfl{盈利情况} 2026 年上半年营业收入 7.56 亿元（同比 +22.1\%），归母净利润 1.04 亿元，同比 +72.3\%。 以年报为趋势对照，2023---2025 年营业收入由 15.98 亿元变为 13.96 亿元（两年年均复合增速 -6.5\%），归母净利润由 2.84 亿元变为 1.00 亿元。 按半年报口径，2026 年上半年的 ROE 为 1.9\%（未年化）、销售净利率 14.7\%、毛利率 53.9\%，ROE 在“动物保健Ⅲ”9 家公司中按数值由高到低排第\zhnumber{7}位。 按同一口径，三级行业“动物保健Ⅲ”9 家公司的半年 ROE 中位数为 2.95\%，该公司低于其中位数 1.03 个百分点。 综合看，公司在报告期维持盈利（高于全样本半年 ROE 中位数 0.75\%），利润的可持续性取决于动物保健Ⅲ景气的延续与成本管控的执行。

\begin{center}\small\setlength{\tabcolsep}{3.4pt}
\captionof{table}[生物股份（600201）关键财务指标]{生物股份（600201）关键财务指标（合并报表口径；两个半年报期均未年化；现金短债比 $=$ 货币资金 $\div$ 短期债务）}
\begin{adjustbox}{max width=\textwidth}
\begin{tabular}{@{}lrrrrr@{}}
\toprule
指标 & 2023 年 & 2024 年 & 2025 年 & 2025 年半年报 & 2026 年半年报 \\
\midrule
营业收入（亿元） & 15.98 & 12.55 & 13.96 & 6.20 & 7.56 \\
归母净利润（亿元） & 2.84 & 1.09 & 1.00 & 0.60 & 1.04 \\
毛利率（\%） & 59.2 & 53.9 & 53.5 & 52.2 & 53.9 \\
ROE（\%） & 5.4 & 1.9 & 1.9 & 1.1 & 1.9 \\
资产负债率（\%） & 17.7 & 19.9 & 16.0 & 16.8 & 16.5 \\
经营活动现金流净额（亿元） & 4.15 & 2.90 & 3.16 & 0.89 & 0.51 \\
货币资金（亿元） & 17.12 & 11.40 & 12.65 & 10.60 & 16.09 \\
有息负债合计（亿元） & 0.04 & 0.18 & 0.02 & 0.17 & 1.01 \\
现金短债比（倍） & 1,117.89 & 69.10 & 884.75 & 64.56 & 1,426.53 \\
\bottomrule
\end{tabular}
\end{adjustbox}
\end{center}

\begin{center}
\begin{minipage}{\textwidth}\centering
\begin{tikzpicture}
\begin{groupplot}[
  group style={group size=2 by 1, horizontal sep=0.60cm},
  width=6.85cm, height=4.60cm,
  tick label style={font=\tiny},
  title style={font=\scriptsize\heiti},
  yticklabel style={font=\tiny},
  ylabel style={font=\tiny},
  grid=major, grid style={LightGray, line width=0.3pt},
  axis line style={InkGray, line width=0.5pt},
  enlarge x limits={abs=0.8},
  legend style={font=\scriptsize, at={(0.5,-0.20)}, anchor=north,
                legend columns=2, draw=none, fill=none, column sep=6pt},
]
\nextgroupplot[
  ybar, bar width=4pt,
  ymin=-1.78, ymax=19.89,
  xtick={0,1,2,3,4,5.7},
  xticklabels={2021,2022,2023,2024,2025,2026H1},
  ylabel={亿元},
]
\addplot[fill=AgriGreen!75, draw=AgriGreen, bar shift=-2.6pt] table[x=i, y=rev]{data/clean/profiles/fin_600201.dat};
\addplot[fill=SoftRed!60, draw=SoftRed, bar shift=2.6pt] table[x=i, y=ni]{data/clean/profiles/fin_600201.dat};
\addplot[fill=AgriGreen!30, draw=AgriGreen, pattern=north east lines, bar shift=-2.6pt] table[x=x, y=rev]{data/clean/profiles/finh_600201.dat};
\addplot[fill=SoftRed!20, draw=SoftRed, pattern=north east lines, bar shift=2.6pt] table[x=x, y=ni]{data/clean/profiles/finh_600201.dat};
\legend{营业收入（年/半年）, 归母净利润（年/半年）}
\nextgroupplot[
  ymin=-1.6, ymax=21.9,
  xtick={0,1,2,3,4,5.7},
  xticklabels={2021,2022,2023,2024,2025,2026H1},
  ylabel={\%},
]
\addplot[AgriGold, mark=*, mark size=1.3pt] table[x=i, y=roe]{data/clean/profiles/fin_600201.dat};
\addplot[OilBlue, mark=square*, mark size=1.3pt] table[x=i, y=debt]{data/clean/profiles/fin_600201.dat};
\addplot[only marks, AgriGold, mark=*, mark size=2.0pt] table[x=x, y=roe]{data/clean/profiles/finh_600201.dat};
\addplot[only marks, OilBlue, mark=square*, mark size=2.0pt] table[x=x, y=debt]{data/clean/profiles/finh_600201.dat};
\legend{ROE, 资产负债率}
\end{groupplot}
\end{tikzpicture}
\begin{tikzpicture}
\begin{axis}[
  name=finbot,
  width=13.75cm, height=3.10cm, scale only axis,
  ymin=-1.71, ymax=19.18,
  xtick={0,1,2,3,4,5.7},
  xticklabels={2021,2022,2023,2024,2025,2026H1},
  tick label style={font=\tiny},
  yticklabel style={font=\tiny},
  ylabel={亿元}, ylabel style={font=\tiny},
  ybar, bar width=4pt,
  grid=major, grid style={LightGray, line width=0.3pt},
  axis line style={InkGray, line width=0.5pt},
  enlarge x limits={abs=0.8},
  legend style={font=\scriptsize, at={(0.5,-0.26)}, anchor=north,
                legend columns=2, draw=none, fill=none, column sep=10pt},
]
\addplot[fill=AgriGreen!55, draw=AgriGreen, bar shift=-3.0pt] table[x=i, y=cash]{data/clean/profiles/fin_600201.dat};
\addplot[fill=SoftRed!45, draw=SoftRed, bar shift=3.0pt] table[x=i, y=ibd]{data/clean/profiles/fin_600201.dat};
\addplot[fill=AgriGreen!25, draw=AgriGreen, pattern=north east lines, bar shift=-3.0pt] table[x=x, y=cash]{data/clean/profiles/finh_600201.dat};
\addplot[fill=SoftRed!20, draw=SoftRed, pattern=north east lines, bar shift=3.0pt] table[x=x, y=ibd]{data/clean/profiles/finh_600201.dat};
\legend{货币资金（左轴）, 有息负债（左轴）}
\end{axis}
\begin{axis}[
  at={(finbot.south east)}, anchor=south east,
  width=13.75cm, height=3.10cm, scale only axis,
  axis y line*=right, axis x line*=none, xtick=\empty,
  ymin=-1.39, ymax=6.93,
  tick label style={font=\tiny},
  yticklabel style={font=\tiny}, scaled y ticks=false,
  legend style={font=\scriptsize, at={(0.5,-0.44)}, anchor=north,
                legend columns=1, draw=none, fill=none},
]
\addplot[OilBlue, mark=*, mark size=2.2pt, line width=1.3pt] table[x=i, y=ocf]{data/clean/profiles/fin_600201.dat};
\addplot[only marks, OilBlue, mark=*, mark size=3.2pt] table[x=x, y=ocf]{data/clean/profiles/finh_600201.dat};
\legend{经营活动现金流净额（右轴，亿元，蓝点）}
\end{axis}
\end{tikzpicture}

\captionof{figure}[生物股份（600201）近年主要财务指标]{生物股份（600201）近年主要财务指标（左上：营业收入与归母净利润；右上：ROE 与资产负债率；下：货币资金、有息负债为左轴柱状，经营活动现金流净额为其右轴的蓝点折线。
2021---2025 年为年报口径，2026H1 为 2026 年半年报/半年末，与其前一年报之间留出间隔、以斜纹或空心点区分；半年报数据未年化）}
\end{minipage}
\end{center}

\pfl{财务分析} 2026 年 6 月末资产负债率 16.5\%，较 2025 年末上升 0.4 个百分点（样本中位数 52.2\%，所处三级行业内按由低到高排第\zhnumber{4}位）。 2026 年 6 月末货币资金 16.09 亿元、短期债务（短期借款与一年内到期非流动负债）0.01 亿元、长期债务 1.00 亿元，有息负债合计 1.01 亿元，现金短债比 1,426.53 倍——覆盖充分。 净有息负债 -15.08 亿元，相当于其 2026 年上半年营业收入的 -199\%。 流动比率 4.21、速动比率 3.82，短期偿债能力处于正常区间。 2026 年上半年经营活动现金流净额 0.51 亿元，经营现金流与利润同向为正，内生资金可以支撑运营。 2026 年 6 月末在建工程 3.63 亿元，较上年末减少 0.02 亿元，在建投入在收缩。 2026 年 6 月末商誉 1.39 亿元，占归母净资产的 2.5\%，是净资产中最脆弱的一块。 综合看，现金短债比 1,426.53 倍，短期偿付压力可控；资产负债率低于样本中位数 35.7 个百分点；有息负债中长期债务占比更高，期限结构对短债滚动的压力较小。

\pfl{重大战略转型} 公司的战略动作可归为以下几类：融资与资本运作（1 项）：2023 年 2 月，公司规划通过向特定对象发行募集资金投入动物 mRNA 疫苗及核酸药物开发项目（2.98 亿元）与动物 mRNA 疫苗及核酸药物生产车间建设项目（2.13 亿元）。 公告口径上，近三年（2023 年 9 月---2026 年 9 月）公司披露重大重组与并购类公告 0 条、对外投资与转型类公告 0 条、回购与股权激励类公告 40 条，合计公告 315 条。 第一大行业仍是“生物制药”，收入占比 92.7\%，与 2021 年年报相比没有方向性变化。 综合判断，报告期内公司未披露实质性的重组、并购或转型安排，资本开支以维持现有产能与业务为主。

\begin{center}\small\setlength{\tabcolsep}{4pt}
\captionof{table}[生物股份（600201）历史融资与资本运作]{生物股份（600201）历史融资与资本运作（据公司公告与公开报道整理；金额为披露值，含首次公开发行、再融资、债券、银行借款与重整投资等）}
\begin{adjustbox}{max width=\textwidth}
\begin{tabular}{@{}p{1.45cm}p{2.55cm}p{4.05cm}p{4.95cm}@{}}
\toprule
时间 & 方式 & 规模 & 用途与说明 \\
\midrule
1999-01 & 首次公开发行A股股票并在上海证券交易所上市（IPO） & 发行价 6.83 元/\allowbreak{}股；首发前总股本 6,\allowbreak{}200.25 万股 & 招股公告日 1998 年 11 月 30 日 \\
2016-09 & 非公开发行A股股票（2015年度非公开发行） & 发行 4,\allowbreak{}032.26 万股，\allowbreak{}发行价格 31.00 元/\allowbreak{}股；募集资金总额 12.50 亿元 & 投资建设金宇生物科技产业园区项目一期工程。；定价基准日为第八届董事会第二十二次会议决议公告日（2016 年 2 月 25 日），\allowbreak{}底价 24.43 元/\allowbreak{}股；发行后实收资本（股本）变更为 6.13 亿元 \\
2021-07-07 & 非公开发行A股股票（2021年度非公开发行，后终止） & 拟非公开发行不超过 6,\allowbreak{}315.79 万股股票（含本数） & 2021 年 7 月 23 日 2021 年第一次临时股东大会审议通过 \\
2022-04-07 & 2021年度非公开发行终止 & 终止拟向内蒙古金宇生物控股有限公司非公开发行不超过 6,\allowbreak{}315.79 万股的计划 & 不适用。；终止原因为鉴于内外部环境变化等因素，\allowbreak{}综合公司整体规划 \\
2023-02-07 & 向特定对象发行A股股票（2023年度向特定对象发行） & 发行价格 8.29 元/\allowbreak{}股；拟募集资金总额不超过人民币 8.00 亿元（含本数）；发行对象为内蒙古金宇生物控股有限公司、\allowbreak{}张翀宇先生和张竞女士 & 拟投向动物 mRNA 疫苗及核酸药物开发项目 2.98 亿元、\allowbreak{}动物 mRNA 疫苗及核酸药物生产车间建设项目 2.13 亿元、\allowbreak{}补充流动资金 2.89 亿元。；2023 年 3 月 16 日 2023 年第二次临时股东大会审议通过；2023 年 5 月 9 日董事会审议通过调整方案；2024 年 1 月 12 日董事会审议通过调减募集资金总额暨调整方案 \\
2024-07-24 & 2023年度向特定对象发行取得中国证监会注册批复 & 中国证监会同意金宇生物批复自同意注册之日（2024 年 7 月 16 日）起 12 个月内有效 & 对应前次发行预案的 mRNA 疫苗及核酸药物开发、\allowbreak{}生产车间建设及补充流动资金项目。；公司于 2024 年 6 月 25 日公告该申请获得上海证券交易所审核通过 \\
2024-10-31 & 股份回购（2024年以集中竞价交易方式回购公司股份） & 拟回购股份数量占总股本比例 0.71\%\textasciitilde{}1.43\% & 用于股权激励及/\allowbreak{}或员工持股计划等（以公司回购方案公告为准）。；2026 年 5 月 19 日公司回购注销部分股票（暨减少注册资本）通知债权人 \\
2025-07-14 & 2023年度向特定对象发行终止 & 终止向特定对象发行 A 股股票 & 不适用 \\
2026-06 & 股份回购并取得银行股票回购专项贷款承诺函 & 回购价格不超过 20.74 元/\allowbreak{}股 & 用于股份回购；2026 年 6 月 13 日公司公告取得该回购专项贷款承诺函 \\
\bottomrule
\end{tabular}
\end{adjustbox}
\end{center}

\pfl{历史融投资情况} 从现金流量表看，2025 年公司偿还债务支付的现金 0.15 亿元、吸收权益性投资 0.07 亿元，筹资活动现金流净额 -3.07 亿元（2023 年为偿还债务 0.10 亿元）。 2026 年上半年筹资活动现金流净额为 -1.13 亿元（取得借款 1.00 亿元、偿还债务 0.20 亿元，未年化），已转为净流出，处于净还债状态。 2025 年分配股利、利润或偿付利息支付的现金合计 0.33 亿元。 公开披露的融资记录共 9 笔，工具类型以 IPO、定向增发、银行借款为主；金额与用途见下表。 其中规模最大的两笔为：2026 年 6 月，银行借款，回购价格不超过 20.74 元/股；2023 年 2 月，定向增发，发行价格 8.29 元/股，拟募集资金总额不超过人民币 8.00 亿元（含本数），发行对象为内蒙古金宇生物控股有限公司、张翀宇先生和张竞女士。 股本方面，近三年披露股本变动 4 次（其他 1 次）；总股本由 2021 年末的 11.26 亿股变为 11.20 亿股（-0.52\%）。 分红方面，近三年现金分红 4 次、合计每股派现 0.21 元，最近一次实施在 2025 年 12 月。 近三年“再融资”类公告 43 条，涉及定向增发、募集资金使用。 综合看，公司融资结构上以债务滚动为主、股权融资为补充；2025 年筹资活动现金流净额 -3.07 亿元，融资节奏仍以维持存量债务周转为主，与前述经营与投资现金流的方向基本一致。

\pfl{未来潜在融资需求分析} 2026 年 6 月末货币资金 16.09 亿元，而短期债务 0.01 亿元（现金短债比 1,426.53 倍），2026 年上半年经营现金流净额 0.51 亿元。账面货币资金可以覆盖短期债务，缺口不体现在静态偿付层面。 公司 2026 年 6 月末资产负债率 16.5\%、2026 年上半年 ROE 1.9\%（未年化），资产负债表尚有余量，融资需求不是“补血型”的刚性需求，而是配套在建项目投入的主动性安排。 公开披露的后续资金安排线索包括：股权融资线索：2026 年 4 月，公司未公告银行综合授信额度、公司债、中票或短融的发行计划，亦无新的定向增发预案；债务与授信安排：2026 年 6 月，贷款金额不超过 3.60 亿元（不超过回购金额的 90\%）。 资金需求还要叠加在建工程：期末在建工程 3.63 亿元、2026 年上半年购建固定资产等支出 0.44 亿元（未年化），这部分支出在周期底部通常会被主动放缓。 以现有货币资金与经营现金流衡量，公司不存在刚性补血的紧迫性，融资动作更多是主动型的。




\subsubsection{瑞普生物（300119）}
\label{co:300119}

\pfl{公司基本情况} 天津自贸试验区的动物保健企业，全称瑞普生物股份有限公司，2008 年 5 月 19 日由天津瑞普生物技术集团有限公司整体变更设立，2010 年 9 月在深圳证券交易所创业板上市；所属行业为医药制造类，产品体系涵盖兽用生物制品、兽用药物制剂、功能性添加剂与兽用原料药，并参股瑞派宠物医院（持股 9.09\%）布局宠物医疗。控股股东、实际控制人为自然人李守军（持股 36.04\%，任董事长）。2025 年归母净利润 4.01 亿元、同比增长 33.18\%；2026 年上半年归母净利润 2.04 亿元、同比下降 20.43\%；2026 年 3 月董事会提请股东会授权以简易程序向特定对象发行融资总额不超过 3 亿元的股票。 公司于 2010-09-17 在深交所创业板首发上市，注册资本 4.65 亿元，总股本 4.66 亿股。 截至 2026 年 10 月 12 日收盘价 16.47 元，对应总市值 76.7 亿元，近一年-24.7\%，总市值在三级行业“动物保健Ⅲ”的 9 家公司中按由大到小排第\zhnumber{2}位，近一年涨跌幅低于全样本中位数（-10.4\%）14.3 个百分点。 按 2026 年半年报披露的行业构成，收入占比前两位为动保业务 75.0\%；宠物供应链 25.0\%。 与 2021 年年报相比，第一大行业口径由“兽用生物制品”（53.3\%）变为“动保业务”（75.0\%）；两期披露的分类口径有调整，占比差异中同时包含分类因素。

\pfl{股权结构与控制权} 按 2026 年 6 月 30 日的股东名册，前三大股东为李守军 持股 36.04\%；梁武 持股 5.04\%；张振（持股比例未披露），股东户数 22,340 户。 控制权与股东性质方面，占总股本 36.04\%，其中质押 4,720.00 万股；截至 2025 年 12 月 31 日公司累计发行股本总数 4.65 亿股；报告期内控股股东与实际控制人均未发生变更。

\pfl{经营要点} 经营与产能方面，2026 年 1 月起，直接减少收入 6,410.82 万元；宠物业务收入 4.79 亿元（同比+15.43\%）；国际业务收入 9,332.01 万元（同比+45.19\%）。

\pfl{所处行业的基本面} 动保的需求由下游存栏规模而非养殖盈利决定，供给端则由产品批文与研发能力设卡，疫苗是其中格局最稳定的细分（第 \ref{sec:framework} 章的供需框架）。 当期该行业的关键数据是：平均毛利率 29.25\%；2025 年，中国兽药产品出口额仅 60.05 亿元（约 8.41 亿美元）；2026 年上半年，全国生猪出栏 37,246 万头（同比+1.7\%），二季度末能繁母猪存栏 3,780 万头（同比-6.5\%）。

\pfl{竞争优势} 把该公司放回同类公司比较，2026 年上半年营业收入 17.2 亿元、在三级行业“动物保健Ⅲ”的 9 家公司中排第\zhnumber{2}位；毛利率 39.7\%、与行业内中位数基本持平；半年 ROE 4.3\%（未年化）、高于行业内中位数（2.95\%）；资产负债率 34.4\%、高于行业内中位数（30.2\%）。 公司区别于同类公司的经营特征包括：拥有 11 个规模化生产基地、92 条生产线；锁定家禽 TOP300、家畜 TOP100、水禽 TOP30、反刍 TOP30 养殖集团，以直销、经销商渠道和政府招标采购三种模式覆盖养殖终端。 优势在当期体现为高于行业中位数的资本回报，但优势能否持续仍取决于周期与自身产能投放节奏。

\pfl{历史股价} 本报告样本区间（2021 年 1 月 至 2026 年 10 月）内，股价由 24.52 元变为 16.47 元，累计 -32.8\%，区间最高 38.05 元（2021 年 5 月）、最低 12.09 元（2024 年 8 月）；当前价相当于区间最高价的 43.3\%，为区间最低价的 1.36 倍。 月度收益的年化波动率 39.8\%，高于全样本中位数 37.3\%，在三级行业“动物保健Ⅲ”的 9 家公司中波动率由高到低排第\zhnumber{3}位。

\begin{center}
\begin{minipage}{\textwidth}\centering
\begin{tikzpicture}
\begin{axis}[
  width=12.6cm, height=3.9cm, scale only axis,
  xmin=2020.922, xmax=2026.972, ymin=11.24, ymax=40.71,
  xtick={2021,2022,2023,2024,2025,2026},
  yearticks,
  yticklabel style={font=\scriptsize},
  ylabel={元/股}, ylabel style={font=\scriptsize},
  grid=major, grid style={LightGray, line width=0.3pt},
  axis line style={InkGray, line width=0.5pt},
  every axis plot/.append style={line width=0.9pt},
]
\addplot[AgriGreen] table[x=x,y=y]{data/clean/profiles/px_300119.dat};
\addplot[SoftRed, dashed, line width=0.7pt] table[x=x,y=y]{data/clean/profiles/pxzz_300119.dat};
\addplot[only marks, mark=*, mark size=1.1pt, SoftRed] table[x=x,y=y]{data/clean/profiles/pxzz_300119.dat};
\end{axis}
\end{tikzpicture}

\captionof{figure}[瑞普生物（300119）股价与周期骨架]{瑞普生物（300119）月度收盘价（元/股）与 ZigZag 周期骨架（阈值 25\%，圆点为周期转折点）}
\end{minipage}
\end{center}

\pfl{过去周期分析} 以 25\% 的 ZigZag 阈值划分，样本区间内股价形成 3 个主升段与 3 个主跌段。 最大主跌段为 2021 年 5 月 至 2022 年 4 月，11 个月-59.9\%；最大主升段出现在 2024 年 8 月 至 2025 年 5 月，9 个月+103.3\%。 最近一次转折出现在 2026 年 6 月（下跌段自 2025 年 5 月起，13 个月-41.3\%），当前处于该段的延续之中。 公司股价较申万农林牧渔指数提前约 21 个月见底，是板块中较早定价周期反转的公司之一。 动保的需求由存栏量而非价格决定（第 \ref{sec:vet} 节）：猪价下行阶段龙头养殖企业的逆势扩产反而推升了疫苗与兽药需求，这正是动保股价与猪价“脱钩”的原因，公司 2026 年上半年实现盈利、半年 ROE 4.3\%（未年化），好于全样本中位数（0.75\%），下行周期对其报表的冲击小于同业。 综合区间形态看，该股单段涨跌幅度偏大、以趋势段为主（最大主升 +103.3\%、最大主跌 -59.9\%），年化波动率 39.8\% 与全样本中位数 37.3\% 相比波动与样本中位数接近。

\pfl{盈利情况} 2026 年上半年营业收入 17.20 亿元（同比 +0.8\%），归母净利润 2.04 亿元，同比 -20.4\%。 以年报为趋势对照，2023---2025 年营业收入由 27.09 亿元变为 33.98 亿元（两年年均复合增速 +12.0\%），归母净利润由 4.54 亿元变为 4.01 亿元。 按半年报口径，2026 年上半年的 ROE 为 4.3\%（未年化）、销售净利率 12.9\%、毛利率 39.7\%，ROE 在“动物保健Ⅲ”9 家公司中按数值由高到低排第\zhnumber{4}位。 同期全样本半年 ROE 中位数 0.75\%，公司与之相差 3.52 个百分点（略好于中位数公司）。 综合看，公司在报告期维持盈利（高于全样本半年 ROE 中位数 0.75\%），利润的可持续性取决于动物保健Ⅲ景气的延续与成本管控的执行。

\begin{center}\small\setlength{\tabcolsep}{3.4pt}
\captionof{table}[瑞普生物（300119）关键财务指标]{瑞普生物（300119）关键财务指标（合并报表口径；两个半年报期均未年化；现金短债比 $=$ 货币资金 $\div$ 短期债务）}
\begin{adjustbox}{max width=\textwidth}
\begin{tabular}{@{}lrrrrr@{}}
\toprule
指标 & 2023 年 & 2024 年 & 2025 年 & 2025 年半年报 & 2026 年半年报 \\
\midrule
营业收入（亿元） & 27.09 & 30.70 & 33.98 & 17.08 & 17.20 \\
归母净利润（亿元） & 4.54 & 3.01 & 4.01 & 2.57 & 2.04 \\
毛利率（\%） & 49.7 & 41.4 & 41.3 & 42.1 & 39.7 \\
ROE（\%） & 10.2 & 6.5 & 8.8 & 5.6 & 4.3 \\
资产负债率（\%） & 27.2 & 37.1 & 32.3 & 36.0 & 34.4 \\
经营活动现金流净额（亿元） & 4.10 & 7.36 & 2.98 & 1.35 & 1.94 \\
货币资金（亿元） & 3.39 & 4.87 & 6.51 & 5.33 & 5.32 \\
有息负债合计（亿元） & 8.97 & 13.81 & 14.60 & 16.22 & 16.16 \\
现金短债比（倍） & 0.59 & 0.57 & 0.60 & 0.56 & 0.42 \\
\bottomrule
\end{tabular}
\end{adjustbox}
\end{center}

\begin{center}
\begin{minipage}{\textwidth}\centering
\begin{tikzpicture}
\begin{groupplot}[
  group style={group size=2 by 1, horizontal sep=0.60cm},
  width=6.85cm, height=4.60cm,
  tick label style={font=\tiny},
  title style={font=\scriptsize\heiti},
  yticklabel style={font=\tiny},
  ylabel style={font=\tiny},
  grid=major, grid style={LightGray, line width=0.3pt},
  axis line style={InkGray, line width=0.5pt},
  enlarge x limits={abs=0.8},
  legend style={font=\scriptsize, at={(0.5,-0.20)}, anchor=north,
                legend columns=2, draw=none, fill=none, column sep=6pt},
]
\nextgroupplot[
  ybar, bar width=4pt,
  ymin=-3.40, ymax=38.06,
  xtick={0,1,2,3,4,5.7},
  xticklabels={2021,2022,2023,2024,2025,2026H1},
  ylabel={亿元},
]
\addplot[fill=AgriGreen!75, draw=AgriGreen, bar shift=-2.6pt] table[x=i, y=rev]{data/clean/profiles/fin_300119.dat};
\addplot[fill=SoftRed!60, draw=SoftRed, bar shift=2.6pt] table[x=i, y=ni]{data/clean/profiles/fin_300119.dat};
\addplot[fill=AgriGreen!30, draw=AgriGreen, pattern=north east lines, bar shift=-2.6pt] table[x=x, y=rev]{data/clean/profiles/finh_300119.dat};
\addplot[fill=SoftRed!20, draw=SoftRed, pattern=north east lines, bar shift=2.6pt] table[x=x, y=ni]{data/clean/profiles/finh_300119.dat};
\legend{营业收入（年/半年）, 归母净利润（年/半年）}
\nextgroupplot[
  ymin=-3.0, ymax=40.9,
  xtick={0,1,2,3,4,5.7},
  xticklabels={2021,2022,2023,2024,2025,2026H1},
  ylabel={\%},
]
\addplot[AgriGold, mark=*, mark size=1.3pt] table[x=i, y=roe]{data/clean/profiles/fin_300119.dat};
\addplot[OilBlue, mark=square*, mark size=1.3pt] table[x=i, y=debt]{data/clean/profiles/fin_300119.dat};
\addplot[only marks, AgriGold, mark=*, mark size=2.0pt] table[x=x, y=roe]{data/clean/profiles/finh_300119.dat};
\addplot[only marks, OilBlue, mark=square*, mark size=2.0pt] table[x=x, y=debt]{data/clean/profiles/finh_300119.dat};
\legend{ROE, 资产负债率}
\end{groupplot}
\end{tikzpicture}
\begin{tikzpicture}
\begin{axis}[
  name=finbot,
  width=13.75cm, height=3.10cm, scale only axis,
  ymin=-1.62, ymax=18.10,
  xtick={0,1,2,3,4,5.7},
  xticklabels={2021,2022,2023,2024,2025,2026H1},
  tick label style={font=\tiny},
  yticklabel style={font=\tiny},
  ylabel={亿元}, ylabel style={font=\tiny},
  ybar, bar width=4pt,
  grid=major, grid style={LightGray, line width=0.3pt},
  axis line style={InkGray, line width=0.5pt},
  enlarge x limits={abs=0.8},
  legend style={font=\scriptsize, at={(0.5,-0.26)}, anchor=north,
                legend columns=2, draw=none, fill=none, column sep=10pt},
]
\addplot[fill=AgriGreen!55, draw=AgriGreen, bar shift=-3.0pt] table[x=i, y=cash]{data/clean/profiles/fin_300119.dat};
\addplot[fill=SoftRed!45, draw=SoftRed, bar shift=3.0pt] table[x=i, y=ibd]{data/clean/profiles/fin_300119.dat};
\addplot[fill=AgriGreen!25, draw=AgriGreen, pattern=north east lines, bar shift=-3.0pt] table[x=x, y=cash]{data/clean/profiles/finh_300119.dat};
\addplot[fill=SoftRed!20, draw=SoftRed, pattern=north east lines, bar shift=3.0pt] table[x=x, y=ibd]{data/clean/profiles/finh_300119.dat};
\legend{货币资金（左轴）, 有息负债（左轴）}
\end{axis}
\begin{axis}[
  at={(finbot.south east)}, anchor=south east,
  width=13.75cm, height=3.10cm, scale only axis,
  axis y line*=right, axis x line*=none, xtick=\empty,
  ymin=-1.91, ymax=9.56,
  tick label style={font=\tiny},
  yticklabel style={font=\tiny}, scaled y ticks=false,
  legend style={font=\scriptsize, at={(0.5,-0.44)}, anchor=north,
                legend columns=1, draw=none, fill=none},
]
\addplot[OilBlue, mark=*, mark size=2.2pt, line width=1.3pt] table[x=i, y=ocf]{data/clean/profiles/fin_300119.dat};
\addplot[only marks, OilBlue, mark=*, mark size=3.2pt] table[x=x, y=ocf]{data/clean/profiles/finh_300119.dat};
\legend{经营活动现金流净额（右轴，亿元，蓝点）}
\end{axis}
\end{tikzpicture}

\captionof{figure}[瑞普生物（300119）近年主要财务指标]{瑞普生物（300119）近年主要财务指标（左上：营业收入与归母净利润；右上：ROE 与资产负债率；下：货币资金、有息负债为左轴柱状，经营活动现金流净额为其右轴的蓝点折线。
2021---2025 年为年报口径，2026H1 为 2026 年半年报/半年末，与其前一年报之间留出间隔、以斜纹或空心点区分；半年报数据未年化）}
\end{minipage}
\end{center}

\pfl{财务分析} 2026 年 6 月末资产负债率 34.4\%，较 2025 年末上升 2.1 个百分点（样本中位数 52.2\%，所处三级行业内按由低到高排第\zhnumber{7}位）。 2026 年 6 月末货币资金 5.32 亿元、短期债务（短期借款与一年内到期非流动负债）12.68 亿元、长期债务 3.49 亿元，有息负债合计 16.16 亿元，现金短债比 0.42 倍——覆盖偏紧。 净有息负债 10.85 亿元，相当于其 2026 年上半年营业收入的 63\%。 流动比率 1.62、速动比率 1.38，短期指标尚可。 2026 年上半年经营活动现金流净额 1.94 亿元，经营现金流与利润同向为正，内生资金可以支撑运营。 2026 年 6 月末在建工程 0.81 亿元，较上年末增加 0.49 亿元，仍有在建投入。 2026 年 6 月末商誉 4.46 亿元，占归母净资产的 8.7\%，是净资产中最脆弱的一块。 综合看，账面货币资金对有息负债与短期债务的覆盖不足，滚动偿付对新增融资的依赖度较高；资产负债率低于样本中位数 17.8 个百分点；有息负债中短期债务占比更高，期限结构仍以滚动续作为主。

\pfl{重大战略转型} 从公开披露看，公司的战略推进集中在：其他动作（3 项）：2026 年 3 月，公司通过参股瑞派宠物医院管理股份有限公司（9.09\%）并推进瑞普生物、瑞派宠物医院、中瑞供应链“三瑞齐发”战略合作；2025 年 9 月，公司与福建圣维生物科技有限公司合资设立内蒙古瑞圣生物科技有限公司；并购与重组（1 项）：2024 年 4 月，公司以 2.30 亿元自有或自筹资金收购保定市收骏科技有限公司 100\% 股权（关联交易）。 公告口径上，近三年（2023 年 9 月---2026 年 9 月）公司披露重大重组与并购类公告 13 条、对外投资与转型类公告 3 条、回购与股权激励类公告 34 条，合计公告 389 条。 主营结构的口径变化已经落到报表上：2021 年年报的第一大行业为“兽用生物制品”（53.3\%），2026 年半年报为“动保业务”（75.0\%）；公司在这两期的分部划分方式有过调整。 综合判断，报告期内公司有 16 条与战略相关的公告落地（重大重组与并购 13 条、对外投资与转型 3 条），资本运作与主业调整之间存在直接关联。

\begin{center}\small\setlength{\tabcolsep}{4pt}
\captionof{table}[瑞普生物（300119）历史融资与资本运作]{瑞普生物（300119）历史融资与资本运作（据公司公告与公开报道整理；金额为披露值，含首次公开发行、再融资、债券、银行借款与重整投资等）}
\begin{adjustbox}{max width=\textwidth}
\begin{tabular}{@{}p{1.45cm}p{2.55cm}p{4.05cm}p{4.95cm}@{}}
\toprule
时间 & 方式 & 规模 & 用途与说明 \\
\midrule
2010-09 & 首次公开发行A股股票并在深圳证券交易所创业板上市（IPO） & 发行价 60.00 元/\allowbreak{}股；首发前总股本 5,\allowbreak{}554.80 万股 & 限责任公司，\allowbreak{}发行方式为网下询价配售与网上定价发行 \\
2016-08 & 创业板非公开发行A股股票（2015年度非公开发行） & 发行 1,\allowbreak{}534.72 万股，\allowbreak{}发行价 14.92 元/\allowbreak{}股；募集资金总额 2.29 亿元 & 市莞城区可园南路一号；发行对象包括李守军、\allowbreak{}梁武、\allowbreak{}李旭东、\allowbreak{}鲍恩东等（以发行情况报告书为准） \\
2021-11 & 向特定对象发行股票（2021年定增） & 发行 6,\allowbreak{}398.47 万股，\allowbreak{}发行价 20.88 元/\allowbreak{}股；募集资金总额 13.36 亿元 & 投向公司向特定对象发行股票募集资金投资项目；其中部分募集资金已于 2025 年变更用途。；新增股份于 2021 年 11 月 19 日在深圳证券交易所上市 \\
2024-04 & 收购保定市收骏科技有限公司100\%股权（关联交易，自有资金或自筹资金） & 交易金额 2.30 亿元 & 2024 年 3 月 28 日第五届董事会第十三次会议、\allowbreak{}第五届监事会第十次会议审议通过；2024 年 4 月公司公告完成工商变更登记 \\
2025-04-01 & 向控股子公司提供财务资助 & 向控股子公司中瑞供应链提供额度不超过人民币 5,\allowbreak{}000.00 万元的借款 & 支持控股子公司经营周转。；2025 年 4 月 1 日第五届董事会第二十二次会议及第五届监事会第十八次会议审议通过编号 2025-024 \\
2025-06-13 & 向控股子公司提供财务资助 & 向中瑞供应链提供额度不超过人民币 1.20 亿元的借款，\allowbreak{}额度范围内可循环使用 & 支持控股子公司经营周转。；2025 年 6 月 13 日第五届董事会第二十四次（临时）会议审议通过编号 2025-047 \\
2025-09-28 & 变更部分募集资金用途 & 变更 2021 年向特定对象发行股票部分募集资金投资项目用途 & 新项目为“生物制造产业化工程建设项目”。；2025 年 9 月 28 日第五届董事会第二十七次（临时）会议审议通过编号 2025-076 \\
2026-03-27 & 提请股东会授权董事会以简易程序向特定对象发行股票 & 融资总额不超过人民币 3.00 亿元且不超过最近一年末净资产 20\% 的中国境内上市人民币普通股（A 股） & 募集资金具体用途由董事会在授权范围内确定 \\
\bottomrule
\end{tabular}
\end{adjustbox}
\end{center}

\pfl{历史融投资情况} 从现金流量表看，2025 年公司取得借款收到的现金 11.86 亿元、偿还债务支付的现金 10.59 亿元、吸收权益性投资 0.01 亿元，筹资活动现金流净额 -2.50 亿元（2023 年为取得借款 7.28 亿元、偿还债务 7.31 亿元）。 2026 年上半年筹资活动现金流净额为 -0.74 亿元（取得借款 8.25 亿元、偿还债务 6.70 亿元，未年化），已转为净流出，处于净还债状态。 2025 年分配股利、利润或偿付利息支付的现金合计 2.25 亿元。 公开披露的融资记录共 8 笔，工具类型以 IPO、定向增发为主；金额与用途见下表。 其中规模最大的两笔为：2010 年 9 月，IPO，发行价 60.00 元/股，首发前总股本 5,554.80 万股；2026 年 3 月，定向增发，融资总额不超过人民币 3.00 亿元且不超过最近一年末净资产 20\% 的中国境内上市人民币普通股（A 股）。 股本方面，近三年披露股本变动 8 次（股份回购 2 次、其他 2 次、激励股份解禁 1 次）；总股本由 2021 年末的 4.68 亿股变为 4.66 亿股（-0.53\%）。 分红方面，近三年现金分红 4 次、合计每股派现 1.40 元，最近一次实施在 2025 年 12 月。 近三年“再融资”类公告 4 条，涉及定向增发。 综合看，公司融资结构上以债务滚动为主、股权融资为补充；2025 年筹资活动现金流净额 -2.50 亿元，融资节奏仍以维持存量债务周转为主，与前述经营与投资现金流的方向基本一致。

\pfl{未来潜在融资需求分析} 2026 年 6 月末货币资金 5.32 亿元，而短期债务 12.68 亿元（现金短债比 0.42 倍），2026 年上半年经营现金流净额 1.94 亿元。按“短期债务减去账面货币资金”的滚动偿付口径，静态缺口约 7.36 亿元；上半年盈利或接近盈亏平衡，该缺口不随经营亏损进一步扩大。 按第 \ref{sec:financing-need} 节的三档划分，公司属于第二档“保壳型”：近三年重大重组与并购类公告 13 条，2026 年上半年归母净利润 2.04 亿元，融资需求更多体现为“资产置换 + 维持上市地位”，而不是扩产。 公开披露的后续资金安排线索包括：股权融资线索：2025 年 9 月，公司变更 2021 年定增部分募集资金用途用于“生物制造产业化工程建设项目”；2026 年 3 月，公司董事会提请股东会授权以简易程序向特定对象发行股票。 资金需求还要叠加在建工程：期末在建工程 0.81 亿元、2026 年上半年购建固定资产等支出 0.97 亿元（未年化），这部分支出在周期底部通常会被主动放缓。 综合看，公司的外部资金需求以补充流动性与项目投入为主，滚动偿付口径下的静态缺口约 7.36 亿元，融资的时点与规模更多取决于信用条件与资本市场窗口。




\subsubsection{中牧股份（600195）}
\label{co:600195}

\pfl{公司基本情况} 北京丰台的央企动保平台，全称中牧实业股份有限公司，1998 年 12 月由中国牧工商（集团）总公司独家发起以募集方式设立，1999 年 1 月 7 日在上海证券交易所上市；专业从事动物保健品和动物营养品研发、生产、销售及技术服务，设立 20 余个生产基地，产品覆盖兽用生物制品、兽用化药、饲料及饲料添加剂与饲料原料贸易，年营收超 60 亿元。2025 年归母净利润 1.71 亿元、同比增长 141.68\%，但扣非净利润仅 0.20 亿元、同比下降 67.42\%；2026 年上半年归母净利润 0.50 亿元、同比下降 39.04\%。 公司于 1999-01-07 在上交所首发上市，注册资本 10.21 亿元，总股本 10.21 亿股。 截至 2026 年 10 月 12 日收盘价 6.65 元，对应总市值 67.9 亿元，近一年-8.5\%，总市值在三级行业“动物保健Ⅲ”的 9 家公司中按由大到小排第\zhnumber{3}位，近一年涨跌幅高于全样本中位数（-10.4\%）1.8 个百分点。 股权结构方面，控股股东为中国牧工商集团有限公司；实际控制人为中国农业发展集团有限公司。

\pfl{股权结构与控制权} 按 2026 年 6 月 30 日的股东名册，前三大股东为中国牧工商集团有限公司 持股 49.37\%；中国华农资产经营有限公司 持股 1.63\%；浙商银行股份有限公司－国泰中证畜牧养殖交易型开放式指（持股比例未披露），股东户数 37,078 户。 股东结构与控制权方面，1.90 亿股、占其总股本 16.39\%；报告期内公司控股股东及实际控制人未发生变更；2026 年 6 月末持股 49.37\%（5.04 亿股）。

\pfl{经营要点} 经营与产能方面，公司共设立 20 余个生产基地，涵盖近千个产品；2026 年上半年，化药收入 11.91 亿元（同比+19.20\%）、贸易 11.02 亿元（同比+41.53\%）、饲料 5.31 亿元（同比-0.95\%）、生物制品 3.41 亿元（同比-24.23\%）。

\pfl{所处行业的基本面} 动保的需求由下游存栏规模而非养殖盈利决定，供给端则由产品批文与研发能力设卡，疫苗是其中格局最稳定的细分（第 \ref{sec:framework} 章的供需框架）。 行业层面的最新情况是：2026 年二季度，末全国生猪存栏 42,491 万头（同比+0.1\%），上半年生猪出栏 37,246 万头（同比+1.7\%）；2026 年上半年，全国家禽出栏 87.6 亿只（同比+7.7\%），牛存栏 9,557 万头（同比-4.3\%）。

\pfl{竞争优势} 把该公司放回同类公司比较，2026 年上半年营业收入 31.7 亿元、在三级行业“动物保健Ⅲ”的 9 家公司中排第\zhnumber{1}位；毛利率 18.4\%、低于行业内中位数 21.3 个百分点；半年 ROE 0.9\%（未年化）、低于行业内中位数（2.95\%）；资产负债率 30.2\%、高于行业内中位数（30.2\%）。 相对优势方面，重大动物疫病疫苗、复合维生素添加剂、蛋鸡预混料、泰妙菌素等产品的品质、产销量与市场占有率位居国内前列；2026 年上半年，综合毛利率由 17.20\% 提升至 18.44\%。 就当期报表看，公司尚未把经营层面的特征转化为高于行业的资本回报，周期底部的相对优势更多体现在现金流与资产负债表上。

\pfl{历史股价} 以 2021 年 1 月为起点、2026 年 10 月为终点，股价由 13.89 元变为 6.65 元，累计 -52.1\%，区间最高 15.92 元（2023 年 2 月）、最低 5.90 元（2026 年 6 月）；当前价相当于区间最高价的 41.8\%，为区间最低价的 1.13 倍。 月度收益的年化波动率 30.5\%，低于全样本中位数 37.3\%，在三级行业“动物保健Ⅲ”的 9 家公司中波动率由高到低排第\zhnumber{8}位。 把股价阶段与公司事件对齐看，2021 年 8 月 至 2023 年 2 月股价上涨+66.7\%，同期（2022 年 12 月）公司终止乾元浩分拆创业板上市后，公司未再公告新的资产分拆或重大资产重组方案。 整体上看，区间的几处大级别拐点都能在公司的资本运作、股权变动或风险事件上找到对应，股价波动有明显的公司层面驱动，而非单纯的行业 beta。

\begin{center}
\begin{minipage}{\textwidth}\centering
\begin{tikzpicture}
\begin{axis}[
  width=12.6cm, height=3.9cm, scale only axis,
  xmin=2020.922, xmax=2026.972, ymin=5.49, ymax=17.03,
  xtick={2021,2022,2023,2024,2025,2026},
  yearticks,
  yticklabel style={font=\scriptsize},
  ylabel={元/股}, ylabel style={font=\scriptsize},
  grid=major, grid style={LightGray, line width=0.3pt},
  axis line style={InkGray, line width=0.5pt},
  every axis plot/.append style={line width=0.9pt},
]
\addplot[AgriGreen] table[x=x,y=y]{data/clean/profiles/px_600195.dat};
\addplot[SoftRed, dashed, line width=0.7pt] table[x=x,y=y]{data/clean/profiles/pxzz_600195.dat};
\addplot[only marks, mark=*, mark size=1.1pt, SoftRed] table[x=x,y=y]{data/clean/profiles/pxzz_600195.dat};
\end{axis}
\end{tikzpicture}

\captionof{figure}[中牧股份（600195）股价与周期骨架]{中牧股份（600195）月度收盘价（元/股）与 ZigZag 周期骨架（阈值 25\%，圆点为周期转折点）}
\end{minipage}
\end{center}

\pfl{过去周期分析} 以 25\% 的 ZigZag 阈值划分，样本区间内股价形成 2 个主升段与 3 个主跌段。 最大主跌段为 2023 年 2 月 至 2024 年 8 月，18 个月-61.5\%；最大主升段出现在 2021 年 8 月 至 2023 年 2 月，18 个月+66.7\%。 最近一次转折出现在 2026 年 6 月（下跌段自 2026 年 2 月起，4 个月-27.9\%），当前处于该段的延续之中。 公司股价与申万农林牧渔指数几乎同步见底，个股并未走出独立行情。 动保的需求由存栏量而非价格决定（第 \ref{sec:vet} 节）：猪价下行阶段龙头养殖企业的逆势扩产反而推升了疫苗与兽药需求，这正是动保股价与猪价“脱钩”的原因，公司 2026 年上半年实现盈利、半年 ROE 0.9\%（未年化），好于全样本中位数（0.75\%），下行周期对其报表的冲击小于同业。 综合区间形态看，该股单段涨跌幅度偏大、以趋势段为主（最大主升 +66.7\%、最大主跌 -61.5\%），年化波动率 30.5\% 与全样本中位数 37.3\% 相比波动小于样本中位数。

\pfl{盈利情况} 2026 年上半年营业收入 31.71 亿元（同比 +13.5\%），归母净利润 0.50 亿元，同比 -39.0\%。 以年报为趋势对照，2023---2025 年营业收入由 54.06 亿元变为 62.05 亿元（两年年均复合增速 +7.1\%），归母净利润由 4.03 亿元变为 1.71 亿元。 按半年报口径，2026 年上半年的 ROE 为 0.9\%（未年化）、销售净利率 2.3\%、毛利率 18.4\%，ROE 在“动物保健Ⅲ”9 家公司中按数值由高到低排第\zhnumber{8}位。 同期全样本半年 ROE 中位数 0.75\%，公司与之相差 0.13 个百分点（略好于中位数公司）。 综合看，公司在报告期维持盈利（高于全样本半年 ROE 中位数 0.75\%），利润的可持续性取决于动物保健Ⅲ景气的延续与成本管控的执行。

\begin{center}\small\setlength{\tabcolsep}{3.4pt}
\captionof{table}[中牧股份（600195）关键财务指标]{中牧股份（600195）关键财务指标（合并报表口径；两个半年报期均未年化；现金短债比 $=$ 货币资金 $\div$ 短期债务）}
\begin{adjustbox}{max width=\textwidth}
\begin{tabular}{@{}lrrrrr@{}}
\toprule
指标 & 2023 年 & 2024 年 & 2025 年 & 2025 年半年报 & 2026 年半年报 \\
\midrule
营业收入（亿元） & 54.06 & 60.17 & 62.05 & 27.93 & 31.71 \\
归母净利润（亿元） & 4.03 & 0.71 & 1.71 & 0.82 & 0.50 \\
毛利率（\%） & 19.8 & 16.4 & 17.5 & 17.2 & 18.4 \\
ROE（\%） & 7.5 & 1.3 & 3.1 & 1.5 & 0.9 \\
资产负债率（\%） & 26.9 & 26.5 & 29.8 & 25.7 & 30.2 \\
经营活动现金流净额（亿元） & 3.72 & 4.20 & 11.58 & -2.32 & -2.42 \\
货币资金（亿元） & 14.93 & 10.91 & 18.73 & 8.96 & 16.65 \\
有息负债合计（亿元） & 13.64 & 10.18 & 10.46 & 11.17 & 11.12 \\
现金短债比（倍） & 3.23 & 3.09 & 2.51 & 1.94 & 1.97 \\
\bottomrule
\end{tabular}
\end{adjustbox}
\end{center}

\begin{center}
\begin{minipage}{\textwidth}\centering
\begin{tikzpicture}
\begin{groupplot}[
  group style={group size=2 by 1, horizontal sep=0.60cm},
  width=6.85cm, height=4.60cm,
  tick label style={font=\tiny},
  title style={font=\scriptsize\heiti},
  yticklabel style={font=\tiny},
  ylabel style={font=\tiny},
  grid=major, grid style={LightGray, line width=0.3pt},
  axis line style={InkGray, line width=0.5pt},
  enlarge x limits={abs=0.8},
  legend style={font=\scriptsize, at={(0.5,-0.20)}, anchor=north,
                legend columns=2, draw=none, fill=none, column sep=6pt},
]
\nextgroupplot[
  ybar, bar width=4pt,
  ymin=-6.20, ymax=69.49,
  xtick={0,1,2,3,4,5.7},
  xticklabels={2021,2022,2023,2024,2025,2026H1},
  ylabel={亿元},
]
\addplot[fill=AgriGreen!75, draw=AgriGreen, bar shift=-2.6pt] table[x=i, y=rev]{data/clean/profiles/fin_600195.dat};
\addplot[fill=SoftRed!60, draw=SoftRed, bar shift=2.6pt] table[x=i, y=ni]{data/clean/profiles/fin_600195.dat};
\addplot[fill=AgriGreen!30, draw=AgriGreen, pattern=north east lines, bar shift=-2.6pt] table[x=x, y=rev]{data/clean/profiles/finh_600195.dat};
\addplot[fill=SoftRed!20, draw=SoftRed, pattern=north east lines, bar shift=2.6pt] table[x=x, y=ni]{data/clean/profiles/finh_600195.dat};
\legend{营业收入（年/半年）, 归母净利润（年/半年）}
\nextgroupplot[
  ymin=-2.4, ymax=33.3,
  xtick={0,1,2,3,4,5.7},
  xticklabels={2021,2022,2023,2024,2025,2026H1},
  ylabel={\%},
]
\addplot[AgriGold, mark=*, mark size=1.3pt] table[x=i, y=roe]{data/clean/profiles/fin_600195.dat};
\addplot[OilBlue, mark=square*, mark size=1.3pt] table[x=i, y=debt]{data/clean/profiles/fin_600195.dat};
\addplot[only marks, AgriGold, mark=*, mark size=2.0pt] table[x=x, y=roe]{data/clean/profiles/finh_600195.dat};
\addplot[only marks, OilBlue, mark=square*, mark size=2.0pt] table[x=x, y=debt]{data/clean/profiles/finh_600195.dat};
\legend{ROE, 资产负债率}
\end{groupplot}
\end{tikzpicture}
\begin{tikzpicture}
\begin{axis}[
  name=finbot,
  width=13.75cm, height=3.10cm, scale only axis,
  ymin=-1.87, ymax=20.98,
  xtick={0,1,2,3,4,5.7},
  xticklabels={2021,2022,2023,2024,2025,2026H1},
  tick label style={font=\tiny},
  yticklabel style={font=\tiny},
  ylabel={亿元}, ylabel style={font=\tiny},
  ybar, bar width=4pt,
  grid=major, grid style={LightGray, line width=0.3pt},
  axis line style={InkGray, line width=0.5pt},
  enlarge x limits={abs=0.8},
  legend style={font=\scriptsize, at={(0.5,-0.26)}, anchor=north,
                legend columns=2, draw=none, fill=none, column sep=10pt},
]
\addplot[fill=AgriGreen!55, draw=AgriGreen, bar shift=-3.0pt] table[x=i, y=cash]{data/clean/profiles/fin_600195.dat};
\addplot[fill=SoftRed!45, draw=SoftRed, bar shift=3.0pt] table[x=i, y=ibd]{data/clean/profiles/fin_600195.dat};
\addplot[fill=AgriGreen!25, draw=AgriGreen, pattern=north east lines, bar shift=-3.0pt] table[x=x, y=cash]{data/clean/profiles/finh_600195.dat};
\addplot[fill=SoftRed!20, draw=SoftRed, pattern=north east lines, bar shift=3.0pt] table[x=x, y=ibd]{data/clean/profiles/finh_600195.dat};
\legend{货币资金（左轴）, 有息负债（左轴）}
\end{axis}
\begin{axis}[
  at={(finbot.south east)}, anchor=south east,
  width=13.75cm, height=3.10cm, scale only axis,
  axis y line*=right, axis x line*=none, xtick=\empty,
  ymin=-6.06, ymax=15.77,
  tick label style={font=\tiny},
  yticklabel style={font=\tiny}, scaled y ticks=false,
  legend style={font=\scriptsize, at={(0.5,-0.44)}, anchor=north,
                legend columns=1, draw=none, fill=none},
]
\addplot[OilBlue, mark=*, mark size=2.2pt, line width=1.3pt] table[x=i, y=ocf]{data/clean/profiles/fin_600195.dat};
\addplot[only marks, OilBlue, mark=*, mark size=3.2pt] table[x=x, y=ocf]{data/clean/profiles/finh_600195.dat};
\legend{经营活动现金流净额（右轴，亿元，蓝点）}
\end{axis}
\end{tikzpicture}

\captionof{figure}[中牧股份（600195）近年主要财务指标]{中牧股份（600195）近年主要财务指标（左上：营业收入与归母净利润；右上：ROE 与资产负债率；下：货币资金、有息负债为左轴柱状，经营活动现金流净额为其右轴的蓝点折线。
2021---2025 年为年报口径，2026H1 为 2026 年半年报/半年末，与其前一年报之间留出间隔、以斜纹或空心点区分；半年报数据未年化）}
\end{minipage}
\end{center}

\pfl{财务分析} 2026 年 6 月末资产负债率 30.2\%，较 2025 年末上升 0.4 个百分点（样本中位数 52.2\%，所处三级行业内按由低到高排第\zhnumber{5}位）。 2026 年 6 月末货币资金 16.65 亿元、短期债务（短期借款与一年内到期非流动负债）8.44 亿元、长期债务 2.69 亿元，有息负债合计 11.12 亿元，现金短债比 1.97 倍——覆盖充分。 净有息负债 -5.52 亿元，相当于其 2026 年上半年营业收入的 -17\%。 流动比率 1.77、速动比率 1.35，短期指标尚可。 2026 年上半年经营活动现金流净额 -2.42 亿元，净利润为正但经营现金流为负，利润的现金含量偏低。 2026 年 6 月末在建工程 0.58 亿元，较上年末增加 0.03 亿元，仍有在建投入。 综合看，现金短债比 1.97 倍，短期偿付压力可控；资产负债率低于样本中位数 22.0 个百分点；有息负债中短期债务占比更高，期限结构仍以滚动续作为主。 公司披露的诉讼与资产冻结风险为：2026 年 6 月，第二大股东中国华农资产经营有限公司所持公司 1,668.35 万股（占总股本 1.63\%）处于冻结状态。

\pfl{重大战略转型} 近三年公司的战略动作集中在：并购与重组（1 项）：2022 年 12 月，公司未再公告新的资产分拆或重大资产重组方案。 公告口径上，近三年（2023 年 9 月---2026 年 9 月）公司披露重大重组与并购类公告 2 条、对外投资与转型类公告 0 条、回购与股权激励类公告 0 条，合计公告 261 条。 综合判断，报告期内公司有 2 条与战略相关的公告落地（重大重组与并购 2 条、对外投资与转型 0 条），资本运作与主业调整之间存在直接关联。

\begin{center}\small\setlength{\tabcolsep}{4pt}
\captionof{table}[中牧股份（600195）历史融资与资本运作]{中牧股份（600195）历史融资与资本运作（据公司公告与公开报道整理；金额为披露值，含首次公开发行、再融资、债券、银行借款与重整投资等）}
\begin{adjustbox}{max width=\textwidth}
\begin{tabular}{@{}p{1.45cm}p{2.55cm}p{4.05cm}p{4.95cm}@{}}
\toprule
时间 & 方式 & 规模 & 用途与说明 \\
\midrule
1999-01 & 首次公开发行A股股票并在上海证券交易所上市（IPO） & 发行价 6.28 元/\allowbreak{}股；首发前总股本 18,\allowbreak{}000.00 万股 & 限责任公司，\allowbreak{}招股公告日 1998 年 11 月 16 日 \\
2013-10 & 非公开发行A股股票（定增） & 发行 3,\allowbreak{}980 万股 & 公司在认购邀请书规定的有效申报时间内共收到 6 家投资者发出的有效申购报价单；新增股份为有限售条件流通股 \\
2015-09-18 & 重大资产重组（停牌筹划并终止） & --- & 不适用。；公司因筹划股权收购事项股票自 2015 年 9 月 18 日起进入重大资产重组停牌程序 \\
2016-02 & 公开发行公司债券（2016年第一期） & 过 12.00 亿元的公司债券；本期基础发行规模 8.00 亿元 & 用于拓宽融资渠道、\allowbreak{}优化债务结构及补充营运资金（以募集说明书为准）。路 66 号 4 号楼；2016 年 2 月 24 日披露发行公告 \\
2019-12-10 & 公开发行公司债券（预案） & 公司董事会审议通过公司债券发行预案，\allowbreak{}拟申请公开发行公司债券 & 为进一步拓宽融资渠道、\allowbreak{}优化债务结构、\allowbreak{}满足资金需求。；该预案公告编号为临 2019-055 \\
2021-07-23 & 分拆控股子公司乾元浩生物股份有限公司至创业板上市 & --- & 不适用。；公司启动分拆乾元浩至创业板上市事项 \\
2022-12-30 & 终止分拆控股子公司乾元浩至创业板上市 & 终止分拆控股子公司乾元浩生物股份有限公司至创业板上市 & 不适用。；公司于 2022 年 12 月 30 日召开第八届董事会 2022 年第二十次临时会议、\allowbreak{}第八届监事会 2022 年第四次临时会议审议通过终止筹划控股子公司分拆上市编号临 2022-045 \\
\bottomrule
\end{tabular}
\end{adjustbox}
\end{center}

\pfl{历史融投资情况} 从现金流量表看，2025 年公司取得借款收到的现金 8.01 亿元、偿还债务支付的现金 7.40 亿元、吸收权益性投资 0.72 亿元，筹资活动现金流净额 0.36 亿元（2023 年为取得借款 13.46 亿元、偿还债务 3.61 亿元）。 2026 年上半年筹资活动现金流净额为 0.55 亿元（取得借款 5.54 亿元、偿还债务 4.87 亿元，未年化），仍处于净融入状态。 2025 年分配股利、利润或偿付利息支付的现金合计 0.93 亿元。 公开披露的融资记录共 7 笔，工具类型以 IPO、定向增发、公司债为主；金额与用途见下表。 其中规模最大的两笔为：2016 年 2 月，公司债，过 12.00 亿元的公司债券，本期基础发行规模 8.00 亿元；2013 年 10 月，定向增发，发行 3,980 万股。 股本方面，近三年披露股本变动 3 次；总股本由 2021 年末的 10.16 亿股变为 10.21 亿股（+0.55\%）。 分红方面，近三年现金分红 5 次、合计每股派现 0.52 元，最近一次实施在 2026 年 7 月。 近三年“再融资”类公告 3 条，涉及可转债。 综合看，公司融资结构上股权与债务融资并举；2025 年筹资活动现金流净额 0.36 亿元，年度内仍为净融入，与前述经营与投资现金流的方向基本一致。

\pfl{未来潜在融资需求分析} 2026 年 6 月末货币资金 16.65 亿元，而短期债务 8.44 亿元（现金短债比 1.97 倍），2026 年上半年经营现金流净额 -2.42 亿元。账面货币资金可以覆盖短期债务，缺口不体现在静态偿付层面。 公司 2026 年 6 月末资产负债率 30.2\%、2026 年上半年 ROE 0.9\%（未年化），账面杠杆不高、盈利能力一般，短期内没有必须依靠外部融资才能维持的经营缺口；后续融资更可能指向技改、扩产或产业并购。 公开披露的后续资金安排线索包括：股权融资线索：2026 年 8 月，公司未公告新的定向增发预案、配股、可转债或公司债发行计划，亦未公告 2026 年度银行综合授信额度；债务与授信安排：2019 年 12 月，公司曾于 2019 年 12 月披露公司债券发行预案。 资金需求还要叠加在建工程：期末在建工程 0.58 亿元、2026 年上半年购建固定资产等支出 0.77 亿元（未年化），这部分支出在周期底部通常会被主动放缓。 以现有货币资金与经营现金流衡量，公司不存在刚性补血的紧迫性，融资动作更多是主动型的。




\subsubsection{科前生物（688526）}
\label{co:688526}

\pfl{公司基本情况} 科前生物成立于 2016 年，武汉的兽用生物制品企业，全称武汉科前生物股份有限公司，2020 年 9 月 22 日在上海证券交易所科创板上市；专注于兽用生物制品研发、生产、销售及动物防疫技术服务，主要产品是非国家强制免疫猪用疫苗和禽用疫苗，可生产 85 个兽用生物制品类产品，是我国兽用生物制品行业内产品种类最丰富的企业之一，并已向禽用疫苗、宠物疫苗、诊断试剂等领域拓展。公司无控股股东，实际控制人为何启盖、吴斌、方六荣、吴美洲四人（一致行动人协议）。2025 年归母净利润 4.19 亿元、同比增长 9.51\%；2026 年上半年归母净利润 1.16 亿元、同比下降 47.37\%；2023 年一次定增撤回、一次定增被上交所终止审核。 公司于 2020-09-22 在上交所科创板首发上市，注册资本 4.66 亿元，总股本 4.66 亿股。 截至 2026 年 10 月 12 日收盘价 12.18 元，对应总市值 56.8 亿元，近一年-31.8\%，总市值在三级行业“动物保健Ⅲ”的 9 家公司中按由大到小排第\zhnumber{4}位，近一年涨跌幅低于全样本中位数（-10.4\%）21.4 个百分点。 按 2025 年年报披露的行业构成，收入占比前三位为兽用生物制品 93.8\%；其他 4.8\%；其他(补充) 1.5\%。 与 2021 年年报相比，第一大行业口径由“兽用生物制造行业”（98.2\%）变为“兽用生物制品”（93.8\%）；两期披露的分类口径有调整，占比差异中同时包含分类因素。

\pfl{股权结构与控制权} 按 2026 年 6 月 30 日的股东名册，前三大股东为武汉华中农大资产经营有限公司 持股 16.73\%；陈焕春 持股 15.20\%；金梅林 持股 7.87\%，股东户数 9,669 户。 股东结构与控制权方面，上述股份均无质押、标记或冻结；吴斌、何启盖、方六荣、吴美洲签订了一致行动人协议；主要业务为授权范围内的经营性资产和对外投资的股权经营管理、转让、投资。

\pfl{经营要点} 业务与产能层面，报告期新增取得 5 项新兽药注册证书或公告，包括猪圆环病毒 2 型重组杆状病毒灭活疫苗；已获得 103 件国家发明专利；2024 年，公司变更募投项目，将原动物生物制品车间技改项目结余募集资金 2.87 亿元投入高级别动物疫苗产业化基地建设项目（一期）。

\pfl{所处行业的基本面} 动保的需求由下游存栏规模而非养殖盈利决定，供给端则由产品批文与研发能力设卡，疫苗是其中格局最稳定的细分（第 \ref{sec:framework} 章的供需框架）。 与该公司直接相关的行业数据包括：2018 年，我国兽用生物制品销售规模从 132.92 亿元攀升至 153.66 亿元，年均复合增长率 2.09\%；2025 年，禽用生物制品销售额 70.25 亿元（占 45.72\%）、猪用 53.51 亿元（占 34.82\%）、牛羊用 19.88 亿元（占 12.94\%）；低效落后产能持续出清。

\pfl{竞争优势} 把该公司放回同类公司比较，2026 年上半年营业收入 3.7 亿元、在三级行业“动物保健Ⅲ”的 9 家公司中排第\zhnumber{7}位；毛利率 59.4\%、高于行业内中位数 19.7 个百分点；半年 ROE 2.8\%（未年化）、低于行业内中位数（2.95\%）；资产负债率 10.7\%、低于行业内中位数（30.2\%）。 公司在公开资料中自述与同业比较的相对优势在于：猪伪狂犬疫苗市场份额 2016 年至 2025 年连续十年国内排名第一。 优势在报表上主要体现为高于行业平均的毛利率，这在周期底部是抵御亏损的主要缓冲。

\pfl{历史股价} 以 2021 年 1 月为起点、2026 年 10 月为终点，股价由 38.08 元变为 12.18 元，累计 -68.0\%，区间最高 41.46 元（2021 年 4 月）、最低 12.10 元（2026 年 6 月）；当前价相当于区间最高价的 29.4\%，已接近区间最低价（1.01 倍）。 月度收益的年化波动率 36.7\%，低于全样本中位数 37.3\%，在三级行业“动物保健Ⅲ”的 9 家公司中波动率由高到低排第\zhnumber{6}位。 把股价阶段与公司事件对齐看，2023 年 2 月 至 2024 年 8 月股价下跌-60.4\%，同期（2024 年 5 月）本次投资不构成关联交易，亦不构成重大资产重组。 把这些阶段与公司事件对齐后可以看到，几轮大涨大跌多由公司自身的资本运作与风险事件触发，行业景气只解释了其中一部分。

\begin{center}
\begin{minipage}{\textwidth}\centering
\begin{tikzpicture}
\begin{axis}[
  width=12.6cm, height=3.9cm, scale only axis,
  xmin=2020.922, xmax=2026.972, ymin=11.25, ymax=44.36,
  xtick={2021,2022,2023,2024,2025,2026},
  yearticks,
  yticklabel style={font=\scriptsize},
  ylabel={元/股}, ylabel style={font=\scriptsize},
  grid=major, grid style={LightGray, line width=0.3pt},
  axis line style={InkGray, line width=0.5pt},
  every axis plot/.append style={line width=0.9pt},
]
\addplot[AgriGreen] table[x=x,y=y]{data/clean/profiles/px_688526.dat};
\addplot[SoftRed, dashed, line width=0.7pt] table[x=x,y=y]{data/clean/profiles/pxzz_688526.dat};
\addplot[only marks, mark=*, mark size=1.1pt, SoftRed] table[x=x,y=y]{data/clean/profiles/pxzz_688526.dat};
\end{axis}
\end{tikzpicture}

\captionof{figure}[科前生物（688526）股价与周期骨架]{科前生物（688526）月度收盘价（元/股）与 ZigZag 周期骨架（阈值 25\%，圆点为周期转折点）}
\end{minipage}
\end{center}

\pfl{过去周期分析} 以 25\% 的 ZigZag 阈值划分，样本区间内股价形成 2 个主升段与 3 个主跌段。 最大主升段出现在 2022 年 4 月 至 2023 年 2 月，10 个月+73.1\%；最大主跌段为 2023 年 2 月 至 2024 年 8 月，18 个月-60.4\%。 最近一次转折出现在 2026 年 6 月（下跌段自 2025 年 8 月起，10 个月-34.2\%），当前处于该段的延续之中。 公司股价与申万农林牧渔指数几乎同步见底，个股并未走出独立行情。 动保的需求由存栏量而非价格决定（第 \ref{sec:vet} 节）：猪价下行阶段龙头养殖企业的逆势扩产反而推升了疫苗与兽药需求，这正是动保股价与猪价“脱钩”的原因，公司 2026 年上半年实现盈利、半年 ROE 2.8\%（未年化），好于全样本中位数（0.75\%），下行周期对其报表的冲击小于同业。 综合区间形态看，该股单段涨跌幅度偏大、以趋势段为主（最大主升 +73.1\%、最大主跌 -60.4\%），年化波动率 36.7\% 与全样本中位数 37.3\% 相比波动与样本中位数接近。

\pfl{盈利情况} 2026 年上半年营业收入 3.74 亿元（同比 -23.3\%），归母净利润 1.16 亿元，同比 -47.4\%。 以年报为趋势对照，2023---2025 年营业收入由 10.64 亿元变为 9.50 亿元（两年年均复合增速 -5.5\%），归母净利润由 3.96 亿元变为 4.19 亿元。 按半年报口径，2026 年上半年的 ROE 为 2.8\%（未年化）、销售净利率 30.7\%、毛利率 59.4\%，ROE 在“动物保健Ⅲ”9 家公司中按数值由高到低排第\zhnumber{6}位。 按同一口径，三级行业“动物保健Ⅲ”9 家公司的半年 ROE 中位数为 2.95\%，该公司低于其中位数 0.18 个百分点。 从公司披露的原因看，业绩变动主要来自：2026 年 4 月，公司商誉账面余额 330.75 万元，系收购武汉诸乐田源生态农业有限公司形成。 综合看，公司在报告期维持盈利（高于全样本半年 ROE 中位数 0.75\%），利润的可持续性取决于动物保健Ⅲ景气的延续与成本管控的执行。

\begin{center}\small\setlength{\tabcolsep}{3.4pt}
\captionof{table}[科前生物（688526）关键财务指标]{科前生物（688526）关键财务指标（合并报表口径；两个半年报期均未年化；现金短债比 $=$ 货币资金 $\div$ 短期债务）}
\begin{adjustbox}{max width=\textwidth}
\begin{tabular}{@{}lrrrrr@{}}
\toprule
指标 & 2023 年 & 2024 年 & 2025 年 & 2025 年半年报 & 2026 年半年报 \\
\midrule
营业收入（亿元） & 10.64 & 9.42 & 9.50 & 4.87 & 3.74 \\
归母净利润（亿元） & 3.96 & 3.82 & 4.19 & 2.20 & 1.16 \\
毛利率（\%） & 72.1 & 63.8 & 67.0 & 67.3 & 59.4 \\
ROE（\%） & 10.8 & 9.9 & 10.3 & 5.4 & 2.8 \\
资产负债率（\%） & 17.8 & 12.3 & 11.5 & 11.4 & 10.7 \\
经营活动现金流净额（亿元） & 5.38 & 3.93 & 4.69 & 1.74 & 1.16 \\
货币资金（亿元） & 5.82 & 5.43 & 1.77 & 5.95 & 1.70 \\
有息负债合计（亿元） & 2.94 & 0.07 & 0.07 & 0.07 & 0.07 \\
现金短债比（倍） & 2.02 & 402.35 & 110.03 & 426.69 & 96.61 \\
\bottomrule
\end{tabular}
\end{adjustbox}
\end{center}

\begin{center}
\begin{minipage}{\textwidth}\centering
\begin{tikzpicture}
\begin{groupplot}[
  group style={group size=2 by 1, horizontal sep=0.60cm},
  width=6.85cm, height=4.60cm,
  tick label style={font=\tiny},
  title style={font=\scriptsize\heiti},
  yticklabel style={font=\tiny},
  ylabel style={font=\tiny},
  grid=major, grid style={LightGray, line width=0.3pt},
  axis line style={InkGray, line width=0.5pt},
  enlarge x limits={abs=0.8},
  legend style={font=\scriptsize, at={(0.5,-0.20)}, anchor=north,
                legend columns=2, draw=none, fill=none, column sep=6pt},
]
\nextgroupplot[
  ybar, bar width=4pt,
  ymin=-1.10, ymax=12.35,
  xtick={0,1,2,3,4,5.7},
  xticklabels={2021,2022,2023,2024,2025,2026H1},
  ylabel={亿元},
]
\addplot[fill=AgriGreen!75, draw=AgriGreen, bar shift=-2.6pt] table[x=i, y=rev]{data/clean/profiles/fin_688526.dat};
\addplot[fill=SoftRed!60, draw=SoftRed, bar shift=2.6pt] table[x=i, y=ni]{data/clean/profiles/fin_688526.dat};
\addplot[fill=AgriGreen!30, draw=AgriGreen, pattern=north east lines, bar shift=-2.6pt] table[x=x, y=rev]{data/clean/profiles/finh_688526.dat};
\addplot[fill=SoftRed!20, draw=SoftRed, pattern=north east lines, bar shift=2.6pt] table[x=x, y=ni]{data/clean/profiles/finh_688526.dat};
\legend{营业收入（年/半年）, 归母净利润（年/半年）}
\nextgroupplot[
  ymin=-1.6, ymax=22.3,
  xtick={0,1,2,3,4,5.7},
  xticklabels={2021,2022,2023,2024,2025,2026H1},
  ylabel={\%},
]
\addplot[AgriGold, mark=*, mark size=1.3pt] table[x=i, y=roe]{data/clean/profiles/fin_688526.dat};
\addplot[OilBlue, mark=square*, mark size=1.3pt] table[x=i, y=debt]{data/clean/profiles/fin_688526.dat};
\addplot[only marks, AgriGold, mark=*, mark size=2.0pt] table[x=x, y=roe]{data/clean/profiles/finh_688526.dat};
\addplot[only marks, OilBlue, mark=square*, mark size=2.0pt] table[x=x, y=debt]{data/clean/profiles/finh_688526.dat};
\legend{ROE, 资产负债率}
\end{groupplot}
\end{tikzpicture}
\begin{tikzpicture}
\begin{axis}[
  name=finbot,
  width=13.75cm, height=3.10cm, scale only axis,
  ymin=-0.58, ymax=6.52,
  xtick={0,1,2,3,4,5.7},
  xticklabels={2021,2022,2023,2024,2025,2026H1},
  tick label style={font=\tiny},
  yticklabel style={font=\tiny},
  ylabel={亿元}, ylabel style={font=\tiny},
  ybar, bar width=4pt,
  grid=major, grid style={LightGray, line width=0.3pt},
  axis line style={InkGray, line width=0.5pt},
  enlarge x limits={abs=0.8},
  legend style={font=\scriptsize, at={(0.5,-0.26)}, anchor=north,
                legend columns=2, draw=none, fill=none, column sep=10pt},
]
\addplot[fill=AgriGreen!55, draw=AgriGreen, bar shift=-3.0pt] table[x=i, y=cash]{data/clean/profiles/fin_688526.dat};
\addplot[fill=SoftRed!45, draw=SoftRed, bar shift=3.0pt] table[x=i, y=ibd]{data/clean/profiles/fin_688526.dat};
\addplot[fill=AgriGreen!25, draw=AgriGreen, pattern=north east lines, bar shift=-3.0pt] table[x=x, y=cash]{data/clean/profiles/finh_688526.dat};
\addplot[fill=SoftRed!20, draw=SoftRed, pattern=north east lines, bar shift=3.0pt] table[x=x, y=ibd]{data/clean/profiles/finh_688526.dat};
\legend{货币资金（左轴）, 有息负债（左轴）}
\end{axis}
\begin{axis}[
  at={(finbot.south east)}, anchor=south east,
  width=13.75cm, height=3.10cm, scale only axis,
  axis y line*=right, axis x line*=none, xtick=\empty,
  ymin=-1.40, ymax=6.99,
  tick label style={font=\tiny},
  yticklabel style={font=\tiny}, scaled y ticks=false,
  legend style={font=\scriptsize, at={(0.5,-0.44)}, anchor=north,
                legend columns=1, draw=none, fill=none},
]
\addplot[OilBlue, mark=*, mark size=2.2pt, line width=1.3pt] table[x=i, y=ocf]{data/clean/profiles/fin_688526.dat};
\addplot[only marks, OilBlue, mark=*, mark size=3.2pt] table[x=x, y=ocf]{data/clean/profiles/finh_688526.dat};
\legend{经营活动现金流净额（右轴，亿元，蓝点）}
\end{axis}
\end{tikzpicture}

\captionof{figure}[科前生物（688526）近年主要财务指标]{科前生物（688526）近年主要财务指标（左上：营业收入与归母净利润；右上：ROE 与资产负债率；下：货币资金、有息负债为左轴柱状，经营活动现金流净额为其右轴的蓝点折线。
2021---2025 年为年报口径，2026H1 为 2026 年半年报/半年末，与其前一年报之间留出间隔、以斜纹或空心点区分；半年报数据未年化）}
\end{minipage}
\end{center}

\pfl{财务分析} 2026 年 6 月末资产负债率 10.7\%，较 2025 年末下降 0.8 个百分点（样本中位数 52.2\%，所处三级行业内按由低到高排第\zhnumber{2}位）。 2026 年 6 月末货币资金 1.70 亿元、短期债务（短期借款与一年内到期非流动负债）0.02 亿元、长期债务 0.05 亿元，有息负债合计 0.07 亿元，现金短债比 96.61 倍——覆盖充分。 净有息负债 -1.63 亿元，相当于其 2026 年上半年营业收入的 -44\%。 流动比率 7.82、速动比率 7.44，短期流动性无明显缺口。 2026 年上半年经营活动现金流净额 1.16 亿元，经营现金流与利润同向为正，内生资金可以支撑运营。 综合看，现金短债比 96.61 倍，短期偿付压力可控；资产负债率低于样本中位数 41.5 个百分点；有息负债中长期债务占比更高，期限结构对短债滚动的压力较小。

\pfl{重大战略转型} 从公开披露看，公司的战略推进集中在：其他动作（2 项）：2026 年上半年，较上年同期增加 6.56 个百分点，在营收下滑背景下继续加大研发投入；主业转型（1 项）：2023 年 9 月，公司实际控制人由陈焕春、金梅林、何启盖、吴斌、方六荣、吴美洲、叶长发七人变更为何启盖、吴斌、方六荣、吴美洲四人本次实际控制人变更不会对公司生产经营产生重大影响。 公告口径上，近三年（2023 年 9 月---2026 年 9 月）公司披露重大重组与并购类公告 0 条、对外投资与转型类公告 1 条、回购与股权激励类公告 35 条，合计公告 343 条。 主营结构的口径变化已经落到报表上：2021 年年报的第一大行业为“兽用生物制造行业”（98.2\%），2025 年年报为“兽用生物制品”（93.8\%）；公司在这两期的分部划分方式有过调整。 综合判断，报告期内公司有 1 条与战略相关的公告落地（重大重组与并购 0 条、对外投资与转型 1 条），资本运作与主业调整之间存在直接关联。

\begin{center}\small\setlength{\tabcolsep}{4pt}
\captionof{table}[科前生物（688526）历史融资与资本运作]{科前生物（688526）历史融资与资本运作（据公司公告与公开报道整理；金额为披露值，含首次公开发行、再融资、债券、银行借款与重整投资等）}
\begin{adjustbox}{max width=\textwidth}
\begin{tabular}{@{}p{1.45cm}p{2.55cm}p{4.05cm}p{4.95cm}@{}}
\toprule
时间 & 方式 & 规模 & 用途与说明 \\
\midrule
2020-09 & 首次公开发行A股股票并在上海证券交易所科创板上市（IPO） & 发行价 11.69 元/\allowbreak{}股；首发前总股本 36,\allowbreak{}000.00 万股 & 投向动物生物制品产业化建设项目（项目投资总额 9.03 亿元、\allowbreak{}募集资金承诺投资额 5.78 亿元）、\allowbreak{}动物生物制品车间技改项目（募集资金承诺投资额 2.87 亿元）、\allowbreak{}科研创新项目及补充流动资金；公司无超募资金。；募集资金于 2020 年 9 月 17 日到位 \\
2022-10-28 & 向特定对象发行A股股票（2022年度向特定对象发行） & 拟向特定对象发行股票数量不超过 876.60 万股（含本数） & 用于“高级别动物疫苗产业化基地建设项目（动物生物安全实验室）”。；2022 年 10 月 28 日董事会、\allowbreak{}2022 年 11 月 14 日股东大会审议通过；2023 年 2 月 27 日第三届董事会第二十四次会议、\allowbreak{}第三届监事会第二十二次会议审议通过 2022 年度向特定对象发行 A 股股票预案（修订稿） \\
2023-06-26 & 终止2022年度向特定对象发行A股股票并撤回申请文件 & 终止拟发行不超过 876.60 万股、\allowbreak{}募集资金总额不超过 1.62 亿元的发行计划 & 不适用。；公司于 2023 年 6 月 26 日召开第三届董事会第二十八次会议和第三届监事会第二十六次会议 \\
2024-05-17 & 变更部分募集资金投资项目 & 拟将首发募投项目“动物生物制品车间技改项目”的募集资金 2.87 亿元（截至 2024 年 4 月 30 日）变更用于新项目“高级别动物疫苗产业化基地建设项目（一期）” & 新项目预计 2028 年 10 月竣工（最终以实际开展情况为准）。；2024 年 5 月 17 日第三届董事会第四十次会议、\allowbreak{}第三届监事会第三十二次会议审议通过 \\
2025-12-31 & 首次公开发行募集资金使用进展 & 截至 2025 年 12 月 31 日，\allowbreak{}首发募集资金净额 11.42 亿元 & 按募投项目计划投入 \\
\bottomrule
\end{tabular}
\end{adjustbox}
\end{center}

\pfl{历史融投资情况} 从现金流量表看，2025 年公司，筹资活动现金流净额 -1.84 亿元（2023 年为取得借款 2.37 亿元、偿还债务 2.22 亿元）。 2026 年上半年筹资活动现金流净额为 -3.00 亿元（取得借款 — 亿元、偿还债务 — 亿元，未年化），已转为净流出，处于净还债状态。 2025 年分配股利、利润或偿付利息支付的现金合计 1.99 亿元。 公开披露的融资记录共 5 笔，工具类型以 IPO、定向增发为主；金额与用途见下表。 其中规模最大的两笔为：2024 年 5 月，变更部分募集资金投资，拟将首发募投项目“动物生物制品车间技改项目”的募集资金 2.87 亿元（截至 2024 年 4 月 30 日）变更用于新项目“高级别动物疫苗产业化基地建设项目（一期）”；2020 年 9 月，IPO，发行价 11.69 元/股，首发前总股本 36,000.00 万股。 股本方面，近三年披露股本变动 7 次（股份回购 2 次、限售股份上市 1 次、其他 1 次）；总股本由 2021 年末的 4.65 亿股变为 4.66 亿股（+0.21\%）。 分红方面，近三年现金分红 4 次、合计每股派现 1.69 元，最近一次实施在 2025 年 12 月。 综合看，公司融资结构上股权与债务融资并举；2025 年筹资活动现金流净额 -1.84 亿元，融资节奏仍以维持存量债务周转为主，与前述经营与投资现金流的方向基本一致。

\pfl{未来潜在融资需求分析} 2026 年 6 月末货币资金 1.70 亿元，而短期债务 0.02 亿元（现金短债比 96.61 倍），2026 年上半年经营现金流净额 1.16 亿元。账面货币资金可以覆盖短期债务，缺口不体现在静态偿付层面。 公司 2026 年 6 月末资产负债率 10.7\%、2026 年上半年 ROE 2.8\%（未年化），财务结构相对宽松而盈利弹性有限，融资属于可选项而非必选项，触发条件多与产能建设或外延并购相关。 公开披露的后续资金安排线索包括：股权融资线索：2026 年 8 月，公司未公告新的定向增发预案、可转债、公司债或银行综合授信额度安排；资本开支安排：2025 年 12 月，首发募集资金累计投入进度为 84.5\%。 静态偿付不存在缺口，后续融资更可能服务于产能或并购安排，而不是填补流动性窟窿。




\subsubsection{金河生物（002688）}
\label{co:002688}

\pfl{公司基本情况} 内蒙古托克托县起家的饲用金霉素龙头，金霉素产销量约占全球一半，2012 年 7 月在深交所中小板上市，现形成兽用化药、兽用疫苗、环保污水处理与玉米深加工业务格局；2016 年、2021 年两次定增合计募资约 12.9 亿元，2025 年因中止非洲猪瘟疫苗研发全额计提吉林佑本商誉减值 1.75 亿元、归母净利润降至 2,755.67 万元，2026 年 8 月再以简易程序定增募资 3 亿元。 公司于 2012-07-13 在深交所主板首发上市，注册资本 8.37 亿元，总股本 7.72 亿股。 截至 2026 年 10 月 12 日收盘价 5.17 元，对应总市值 39.9 亿元，近一年-24.4\%，总市值在三级行业“动物保健Ⅲ”的 9 家公司中按由大到小排第\zhnumber{5}位，近一年涨跌幅低于全样本中位数（-10.4\%）14.1 个百分点。 按 2026 年半年报披露的行业构成，收入占比前三位为兽用化学药品 60.6\%；农产品加工业 22.3\%；兽用生物制品 8.8\%。

\pfl{股权结构与控制权} 按 2026 年 8 月 27 日的股东名册，前三大股东为内蒙古金河控股有限公司 持股 26.18\%；路牡丹 持股 3.10\%；财通基金管理有限公司 持股 2.16\%。 股权结构的另一条线索是，王东晓持有控股股东 82\% 的股权，其配偶路牡丹持有 18\% 的股权；公司回购专用证券账户持有普通股 1,382.14 万股、占总股本 1.80\%。

\pfl{经营要点} 从经营要点看，获批两大美国产能升级项目，明尼苏达州疫苗新扩建项目投资 2,601.56 万美元。

\pfl{所处行业的基本面} 动保的需求由下游存栏规模而非养殖盈利决定，供给端则由产品批文与研发能力设卡，疫苗是其中格局最稳定的细分（第 \ref{sec:framework} 章的供需框架）。 当期该行业的关键数据是：2026 年 8 月，涉及 126 家兽药生产企业和 380 个产品，其中 23 家企业当月批签发超过 25 批、占 126 家企业的 18.3\%；2026 年上半年，兽用化药产品价格走弱。

\pfl{竞争优势} 把该公司放回同类公司比较，2026 年上半年营业收入 14.5 亿元、在三级行业“动物保健Ⅲ”的 9 家公司中排第\zhnumber{3}位；毛利率 37.8\%、低于行业内中位数 1.9 个百分点；半年 ROE 5.7\%（未年化）、高于行业内中位数（2.95\%）；资产负债率 56.1\%、高于行业内中位数（30.2\%）。 公司在公开资料中自述与同业比较的相对优势在于：2026 年上半年，境外收入占比达 39.28\%、境外毛利率 54.86\%；金霉素等主力产品在美国市场份额持续提升。 优势在当期体现为高于行业中位数的资本回报，但优势能否持续仍取决于周期与自身产能投放节奏。

\pfl{历史股价} 本报告样本区间（2021 年 1 月 至 2026 年 10 月）内，股价由 5.05 元变为 5.17 元，累计 +2.4\%，区间最高 7.12 元（2021 年 6 月）、最低 3.63 元（2024 年 6 月）；当前价相当于区间最高价的 72.6\%，为区间最低价的 1.42 倍。 月度收益的年化波动率 30.0\%，低于全样本中位数 37.3\%，在三级行业“动物保健Ⅲ”的 9 家公司中波动率由高到低排第\zhnumber{9}位。 把股价阶段与公司事件对齐看，2024 年 2 月 至 2025 年 8 月股价上涨+94.5\%，同期（2023 年 12 月）公司收购子公司产生的商誉账面价值为 4.99 亿元。 把这些阶段与公司事件对齐后可以看到，几轮大涨大跌多由公司自身的资本运作与风险事件触发，行业景气只解释了其中一部分。

\begin{center}
\begin{minipage}{\textwidth}\centering
\begin{tikzpicture}
\begin{axis}[
  width=12.6cm, height=3.9cm, scale only axis,
  xmin=2020.922, xmax=2026.972, ymin=3.38, ymax=7.62,
  xtick={2021,2022,2023,2024,2025,2026},
  yearticks,
  yticklabel style={font=\scriptsize},
  ylabel={元/股}, ylabel style={font=\scriptsize},
  grid=major, grid style={LightGray, line width=0.3pt},
  axis line style={InkGray, line width=0.5pt},
  every axis plot/.append style={line width=0.9pt},
]
\addplot[AgriGreen] table[x=x,y=y]{data/clean/profiles/px_002688.dat};
\addplot[SoftRed, dashed, line width=0.7pt] table[x=x,y=y]{data/clean/profiles/pxzz_002688.dat};
\addplot[only marks, mark=*, mark size=1.1pt, SoftRed] table[x=x,y=y]{data/clean/profiles/pxzz_002688.dat};
\end{axis}
\end{tikzpicture}

\captionof{figure}[金河生物（002688）股价与周期骨架]{金河生物（002688）月度收盘价（元/股）与 ZigZag 周期骨架（阈值 25\%，圆点为周期转折点）}
\end{minipage}
\end{center}

\pfl{过去周期分析} 以 25\% 的 ZigZag 阈值划分，样本区间内股价形成 3 个主升段与 3 个主跌段。 最大主升段出现在 2024 年 2 月 至 2025 年 8 月，18 个月+94.5\%；最大主跌段为 2021 年 6 月 至 2022 年 4 月，10 个月-40.9\%。 最近一次转折出现在 2026 年 6 月（下跌段自 2025 年 8 月起，10 个月-34.9\%），当前处于该段的延续之中。 公司股价较申万农林牧渔指数提前约 23 个月见底，是板块中较早定价周期反转的公司之一。 动保的需求由存栏量而非价格决定（第 \ref{sec:vet} 节）：猪价下行阶段龙头养殖企业的逆势扩产反而推升了疫苗与兽药需求，这正是动保股价与猪价“脱钩”的原因，公司 2026 年上半年实现盈利、半年 ROE 5.7\%（未年化），好于全样本中位数（0.75\%），下行周期对其报表的冲击小于同业。 综合区间形态看，该股单段涨跌幅度偏大、以趋势段为主（最大主升 +94.5\%、最大主跌 -40.9\%），年化波动率 30.0\% 与全样本中位数 37.3\% 相比波动小于样本中位数。

\pfl{盈利情况} 2026 年上半年营业收入 14.48 亿元（同比 +4.1\%），归母净利润 1.31 亿元，同比 -5.1\%。 以年报为趋势对照，2023---2025 年营业收入由 21.74 亿元变为 28.78 亿元（两年年均复合增速 +15.1\%），归母净利润由 0.86 亿元变为 0.28 亿元。 按半年报口径，2026 年上半年的 ROE 为 5.7\%（未年化）、销售净利率 9.3\%、毛利率 37.8\%，ROE 在“动物保健Ⅲ”9 家公司中按数值由高到低排第\zhnumber{3}位。 同期全样本半年 ROE 中位数 0.75\%，公司与之相差 4.99 个百分点（略好于中位数公司）。 从公司披露的原因看，业绩变动主要来自：2023 年 12 月，公司收购子公司产生的商誉账面价值为 4.99 亿元。 综合看，公司在报告期维持盈利（高于全样本半年 ROE 中位数 0.75\%），利润的可持续性取决于动物保健Ⅲ景气的延续与成本管控的执行。

\begin{center}\small\setlength{\tabcolsep}{3.4pt}
\captionof{table}[金河生物（002688）关键财务指标]{金河生物（002688）关键财务指标（合并报表口径；两个半年报期均未年化；现金短债比 $=$ 货币资金 $\div$ 短期债务）}
\begin{adjustbox}{max width=\textwidth}
\begin{tabular}{@{}lrrrrr@{}}
\toprule
指标 & 2023 年 & 2024 年 & 2025 年 & 2025 年半年报 & 2026 年半年报 \\
\midrule
营业收入（亿元） & 21.74 & 23.71 & 28.78 & 13.90 & 14.48 \\
归母净利润（亿元） & 0.86 & 1.00 & 0.28 & 1.38 & 1.31 \\
毛利率（\%） & 29.6 & 33.5 & 33.9 & 34.6 & 37.8 \\
ROE（\%） & 3.9 & 4.5 & 1.3 & 6.0 & 5.7 \\
资产负债率（\%） & 53.9 & 55.8 & 55.7 & 55.5 & 56.1 \\
经营活动现金流净额（亿元） & 1.99 & 4.98 & 4.69 & 1.34 & 2.54 \\
货币资金（亿元） & 6.08 & 6.36 & 6.18 & 6.25 & 7.00 \\
有息负债合计（亿元） & 23.00 & 23.19 & 23.20 & 24.15 & 23.64 \\
现金短债比（倍） & 0.34 & 0.48 & 0.47 & 0.41 & 0.51 \\
\bottomrule
\end{tabular}
\end{adjustbox}
\end{center}

\begin{center}
\begin{minipage}{\textwidth}\centering
\begin{tikzpicture}
\begin{groupplot}[
  group style={group size=2 by 1, horizontal sep=0.60cm},
  width=6.85cm, height=4.60cm,
  tick label style={font=\tiny},
  title style={font=\scriptsize\heiti},
  yticklabel style={font=\tiny},
  ylabel style={font=\tiny},
  grid=major, grid style={LightGray, line width=0.3pt},
  axis line style={InkGray, line width=0.5pt},
  enlarge x limits={abs=0.8},
  legend style={font=\scriptsize, at={(0.5,-0.20)}, anchor=north,
                legend columns=2, draw=none, fill=none, column sep=6pt},
]
\nextgroupplot[
  ybar, bar width=4pt,
  ymin=-2.88, ymax=32.23,
  xtick={0,1,2,3,4,5.7},
  xticklabels={2021,2022,2023,2024,2025,2026H1},
  ylabel={亿元},
]
\addplot[fill=AgriGreen!75, draw=AgriGreen, bar shift=-2.6pt] table[x=i, y=rev]{data/clean/profiles/fin_002688.dat};
\addplot[fill=SoftRed!60, draw=SoftRed, bar shift=2.6pt] table[x=i, y=ni]{data/clean/profiles/fin_002688.dat};
\addplot[fill=AgriGreen!30, draw=AgriGreen, pattern=north east lines, bar shift=-2.6pt] table[x=x, y=rev]{data/clean/profiles/finh_002688.dat};
\addplot[fill=SoftRed!20, draw=SoftRed, pattern=north east lines, bar shift=2.6pt] table[x=x, y=ni]{data/clean/profiles/finh_002688.dat};
\legend{营业收入（年/半年）, 归母净利润（年/半年）}
\nextgroupplot[
  ymin=-4.5, ymax=61.7,
  xtick={0,1,2,3,4,5.7},
  xticklabels={2021,2022,2023,2024,2025,2026H1},
  ylabel={\%},
]
\addplot[AgriGold, mark=*, mark size=1.3pt] table[x=i, y=roe]{data/clean/profiles/fin_002688.dat};
\addplot[OilBlue, mark=square*, mark size=1.3pt] table[x=i, y=debt]{data/clean/profiles/fin_002688.dat};
\addplot[only marks, AgriGold, mark=*, mark size=2.0pt] table[x=x, y=roe]{data/clean/profiles/finh_002688.dat};
\addplot[only marks, OilBlue, mark=square*, mark size=2.0pt] table[x=x, y=debt]{data/clean/profiles/finh_002688.dat};
\legend{ROE, 资产负债率}
\end{groupplot}
\end{tikzpicture}
\begin{tikzpicture}
\begin{axis}[
  name=finbot,
  width=13.75cm, height=3.10cm, scale only axis,
  ymin=-2.36, ymax=26.47,
  xtick={0,1,2,3,4,5.7},
  xticklabels={2021,2022,2023,2024,2025,2026H1},
  tick label style={font=\tiny},
  yticklabel style={font=\tiny},
  ylabel={亿元}, ylabel style={font=\tiny},
  ybar, bar width=4pt,
  grid=major, grid style={LightGray, line width=0.3pt},
  axis line style={InkGray, line width=0.5pt},
  enlarge x limits={abs=0.8},
  legend style={font=\scriptsize, at={(0.5,-0.26)}, anchor=north,
                legend columns=2, draw=none, fill=none, column sep=10pt},
]
\addplot[fill=AgriGreen!55, draw=AgriGreen, bar shift=-3.0pt] table[x=i, y=cash]{data/clean/profiles/fin_002688.dat};
\addplot[fill=SoftRed!45, draw=SoftRed, bar shift=3.0pt] table[x=i, y=ibd]{data/clean/profiles/fin_002688.dat};
\addplot[fill=AgriGreen!25, draw=AgriGreen, pattern=north east lines, bar shift=-3.0pt] table[x=x, y=cash]{data/clean/profiles/finh_002688.dat};
\addplot[fill=SoftRed!20, draw=SoftRed, pattern=north east lines, bar shift=3.0pt] table[x=x, y=ibd]{data/clean/profiles/finh_002688.dat};
\legend{货币资金（左轴）, 有息负债（左轴）}
\end{axis}
\begin{axis}[
  at={(finbot.south east)}, anchor=south east,
  width=13.75cm, height=3.10cm, scale only axis,
  axis y line*=right, axis x line*=none, xtick=\empty,
  ymin=-1.30, ymax=6.48,
  tick label style={font=\tiny},
  yticklabel style={font=\tiny}, scaled y ticks=false,
  legend style={font=\scriptsize, at={(0.5,-0.44)}, anchor=north,
                legend columns=1, draw=none, fill=none},
]
\addplot[OilBlue, mark=*, mark size=2.2pt, line width=1.3pt] table[x=i, y=ocf]{data/clean/profiles/fin_002688.dat};
\addplot[only marks, OilBlue, mark=*, mark size=3.2pt] table[x=x, y=ocf]{data/clean/profiles/finh_002688.dat};
\legend{经营活动现金流净额（右轴，亿元，蓝点）}
\end{axis}
\end{tikzpicture}

\captionof{figure}[金河生物（002688）近年主要财务指标]{金河生物（002688）近年主要财务指标（左上：营业收入与归母净利润；右上：ROE 与资产负债率；下：货币资金、有息负债为左轴柱状，经营活动现金流净额为其右轴的蓝点折线。
2021---2025 年为年报口径，2026H1 为 2026 年半年报/半年末，与其前一年报之间留出间隔、以斜纹或空心点区分；半年报数据未年化）}
\end{minipage}
\end{center}

\pfl{财务分析} 2026 年 6 月末资产负债率 56.1\%，较 2025 年末上升 0.4 个百分点（样本中位数 52.2\%，所处三级行业内按由低到高排第\zhnumber{8}位）。 2026 年 6 月末货币资金 7.00 亿元、短期债务（短期借款与一年内到期非流动负债）13.82 亿元、长期债务 9.82 亿元，有息负债合计 23.64 亿元，现金短债比 0.51 倍——覆盖偏紧。 净有息负债 16.64 亿元，相当于其 2026 年上半年营业收入的 115\%。 流动比率 1.18、速动比率 0.77，短期流动性无明显缺口。 2026 年上半年经营活动现金流净额 2.54 亿元，经营现金流与利润同向为正，内生资金可以支撑运营。 2026 年 6 月末在建工程 2.83 亿元，较上年末增加 0.87 亿元，仍有在建投入。 2026 年 6 月末商誉 2.78 亿元，占归母净资产的 11.2\%，是净资产中最脆弱的一块。 综合看，现金短债比 0.51 倍，短期资金安排偏紧；资产负债率高于样本中位数 3.9 个百分点；有息负债中短期债务占比更高，期限结构仍以滚动续作为主。 公司披露的风险集中在诉讼与资产冻结、股东与质押两类：一是 2026 年半年报，显示公司货币资金 2,819 元（母公司口径列示于受限资产表）因吉林佑本供应商起诉导致账户被冻结；二是 2026 年，一季报显示控股股东内蒙古金河控股有限公司持有 2.19 亿股（占 28.49\%）。

\pfl{重大战略转型} 近三年公司的战略动作集中在：融资与资本运作（1 项）：2018 年 1 月，公司拟公开发行不超过 6.00 亿元可转债投建动物疫苗生产基地建设项目；并购与重组（1 项）：2023 年 5 月，控股子公司金河佑本以股权转让及增资方式取得吉林百思万可 60\% 股权。 公告口径上，近三年（2023 年 9 月---2026 年 9 月）公司披露重大重组与并购类公告 1 条、对外投资与转型类公告 5 条、回购与股权激励类公告 65 条，合计公告 585 条。 第一大行业仍是“兽用化学药品”，收入占比 60.6\%，与 2021 年年报相比没有方向性变化。 综合判断，报告期内公司有 6 条与战略相关的公告落地（重大重组与并购 1 条、对外投资与转型 5 条），资本运作与主业调整之间存在直接关联。

\begin{center}\small\setlength{\tabcolsep}{4pt}
\captionof{table}[金河生物（002688）历史融资与资本运作]{金河生物（002688）历史融资与资本运作（据公司公告与公开报道整理；金额为披露值，含首次公开发行、再融资、债券、银行借款与重整投资等）}
\begin{adjustbox}{max width=\textwidth}
\begin{tabular}{@{}p{1.45cm}p{2.55cm}p{4.05cm}p{4.95cm}@{}}
\toprule
时间 & 方式 & 规模 & 用途与说明 \\
\midrule
2012-07-05 & IPO（深圳证券交易所中小板） & 发行 2,\allowbreak{}723 万股，\allowbreak{}发行价 18.00 元/\allowbreak{}股，\allowbreak{}募集资金总额 4.90 亿元 & 年产 10,\allowbreak{}000 吨高效饲用金霉素扩建项目、\allowbreak{}金河生物研发中心建设项目；经证监发行字[2012]838 号文核准 \\
2016-02-03 & 非公开发行股票 & 发行 3,\allowbreak{}627.59 万股，\allowbreak{}发行价 13.05 元/\allowbreak{}股，\allowbreak{}认购资金总额 4.73 亿元 & 全部用于补充流动资金；；2016 年 1 月 21 日至 1 月 22 日完成验资 \\
2018-01-04 & 公开发行可转换公司债券（终止） & 拟募集资金总额（含发行费用）不超过人民币 6.00 亿元 & 动物疫苗生产基地建设项目（项目总投资 6.90 亿元 \\
2021-08-10 & 非公开发行股票（向 10 名认购对象） & 发行 1.45 亿股，\allowbreak{}发行价 5.65 元/\allowbreak{}股，\allowbreak{}募集资金总额 8.20 亿元 & 募集资金于 2021 年 7 月 5 日全部到位 \\
2023-05-24 & 现金收购与增资（并购） & 控股子公司金河佑本以 1.63 亿元受让吉林百思万可（后更名金河佑本（吉林）生物科技有限公司）34\% 股权 & 布局非洲猪瘟防控及疫苗产研领域；标的公司原股东为李忠祥、\allowbreak{}张年红、\allowbreak{}张维等 15 名自然人；2026 年 1 月 30 日公司决定中止非洲猪瘟冻干灭活疫苗研发和申报事项 \\
2026-01-30 & 控股子公司售后回租融资租赁（公司提供担保） & 公司按持股比例 85.6071\% 为金河佑本售后回租融资租赁业务提供担保 & 拓宽控股子公司金河佑本融资渠道、\allowbreak{}优化融资结构；经公司第六届董事会第三十六次会议（2026 年 1 月 30 日）审议通过并提交 2026 年第二次临时股东会审议；同日董事会还审议通过变更募集资金用途，\allowbreak{}颗粒剂自动化密闭生产线项目” \\
2026-08-27 & 以简易程序向特定对象发行股票（定增） & 发行 6,\allowbreak{}787.33 万股，\allowbreak{}发行价 4.42 元/\allowbreak{}股，\allowbreak{}募集资金总额 3.00 亿元 & 污水处理扩容及水资源再生利用提标扩建工程项目（拟投入 1.36 亿元）、\allowbreak{}新建产品库房建设项目（粮食储备库 \\
\bottomrule
\end{tabular}
\end{adjustbox}
\end{center}

\pfl{历史融投资情况} 从现金流量表看，2025 年公司取得借款收到的现金 14.17 亿元、偿还债务支付的现金 14.26 亿元、吸收权益性投资 0.02 亿元，筹资活动现金流净额 -1.24 亿元（2023 年为取得借款 20.33 亿元、偿还债务 14.33 亿元）。 2026 年上半年筹资活动现金流净额为 -0.24 亿元（取得借款 5.90 亿元、偿还债务 6.75 亿元，未年化），已转为净流出，处于净还债状态。 2025 年分配股利、利润或偿付利息支付的现金合计 1.48 亿元。 公开披露的融资记录共 7 笔，工具类型以 IPO、定向增发、可转债、融资租赁为主；金额与用途见下表。 其中规模最大的两笔为：2026 年 1 月，融资租赁，公司按持股比例 85.6071\% 为金河佑本售后回租融资租赁业务提供担保；2012 年 7 月，IPO，发行 2,723 万股，发行价 18.00 元/股，募集资金总额 4.90 亿元。 股本方面，近三年披露股本变动 7 次（股份回购 3 次、股权激励 2 次、其他 1 次）；总股本由 2021 年末的 7.80 亿股变为 7.72 亿股（-1.12\%）。 分红方面，近三年现金分红 4 次、合计每股派现 0.40 元，最近一次实施在 2025 年 12 月。 近三年“再融资”类公告 26 条，涉及定向增发、募集资金使用。 综合看，公司融资结构上以债务滚动为主、股权融资为补充；2025 年筹资活动现金流净额 -1.24 亿元，融资节奏仍以维持存量债务周转为主，与前述经营与投资现金流的方向基本一致。

\pfl{未来潜在融资需求分析} 2026 年 6 月末货币资金 7.00 亿元，而短期债务 13.82 亿元（现金短债比 0.51 倍），2026 年上半年经营现金流净额 2.54 亿元。按“短期债务减去账面货币资金”的滚动偿付口径，静态缺口约 6.82 亿元；上半年盈利或接近盈亏平衡，该缺口不随经营亏损进一步扩大。 按第 \ref{sec:financing-need} 节的三档划分，公司属于第三档“扩张型”：2026 年上半年 ROE 5.7\%（未年化，折合约 11.5\%）、6 月末资产负债率 56.1\%，资产负债表仍具备加杠杆空间；融资用途以产能扩张、品类并购与海外布局为主，单笔规模通常在 5---20 亿元量级，会被市场定价为成长而非补血。 公开披露的后续资金安排线索包括：股权融资线索：2026 年 8 月，公司完成 2025 年度以简易程序向特定对象发行股票；债务与授信安排：2026 年 1 月，控股子公司金河佑本开展售后回租融资租赁业务、公司按 85.6071\% 持股比例提供不超过 1.21 亿元担保并由控股股东全额担保。 资金需求还要叠加在建工程：期末在建工程 2.83 亿元、2026 年上半年购建固定资产等支出 1.44 亿元（未年化），这部分支出在周期底部通常会被主动放缓。 综合看，公司的外部资金需求以补充流动性与项目投入为主，滚动偿付口径下的静态缺口约 6.82 亿元，融资的时点与规模更多取决于信用条件与资本市场窗口。




\subsubsection{普莱柯（603566）}
\label{co:603566}

\pfl{公司基本情况} 河南洛阳的综合性兽用疫苗与兽药企业，2002 年成立、2015 年 5 月在上交所上市，产品覆盖猪用、禽用疫苗及兽用化药；2022 年 9 月完成 8.98 亿元非公开发行投向兽用灭活疫苗产能，2025 年归母净利润 1.85 亿元、同比增长 99.33\%，但 2026 年上半年受下游生猪价格低位、动保行业竞争加剧及生物制品增值税率上调影响，营收 4.91 亿元、同比下降 12.18\%，归母净利润同比下降 33.45\%。 公司于 2015-05-18 在上交所首发上市，注册资本 3.46 亿元，总股本 3.46 亿股。 截至 2026 年 10 月 12 日收盘价 10.24 元，对应总市值 35.4 亿元，近一年-27.4\%，总市值在三级行业“动物保健Ⅲ”的 9 家公司中按由大到小排第\zhnumber{6}位，近一年涨跌幅低于全样本中位数（-10.4\%）17.0 个百分点。 按 2025 年年报披露的行业构成，收入占比前三位为兽用药品 94.0\%；技术许可或转让 2.5\%；功能性保健品等(宠物用) 2.3\%。

\pfl{股权结构与控制权} 按 2026 年 6 月 30 日的股东名册，前三大股东为张许科 持股 31.58\%；孙进忠 持股 4.92\%；中信农业科技股份有限公司 持股 2.33\%，股东户数 16,640 户。 股权结构的另一条线索是，2026 年 6 月末持股 31.58\%（1.09 亿股），报告期内持股数量未发生变化且不存在质押、冻结；第二大股东孙进忠持股 4.92\%（1,702.72 万股）；公司控制权保持稳定。

\pfl{经营要点} 经营与产能方面，生物制品板块收入 3.43 亿元（同比-14.44\%）；子公司普莱柯（南京）新建生产基地建成并通过兽药 GMP 验收，缓解高致病性禽流感疫苗产能瓶颈。

\pfl{所处行业的基本面} 动保的需求由下游存栏规模而非养殖盈利决定，供给端则由产品批文与研发能力设卡，疫苗是其中格局最稳定的细分（第 \ref{sec:framework} 章的供需框架）。 与该公司直接相关的行业数据包括：2026 年上半年，全国生猪出栏 37,246 万头（同比+1.7\%）出栏节奏放缓，二季度末生猪存栏 42,491 万头（同比+0.1\%）。

\pfl{竞争优势} 把该公司放回同类公司比较，2026 年上半年营业收入 4.9 亿元、在三级行业“动物保健Ⅲ”的 9 家公司中排第\zhnumber{6}位；毛利率 55.8\%、高于行业内中位数 16.0 个百分点；半年 ROE 3.0\%（未年化）、低于行业内中位数（2.95\%）；资产负债率 15.2\%、低于行业内中位数（30.2\%）。 相对优势方面，利用新技术、新平台积极推进禽腺病毒亚单位系列联苗、鸡传染性支气管炎（悬浮培养）系列联苗等产品的研发和推广上市，为公司发展提供新动能；但报告期口蹄疫疫苗政府采购中标金额同比提升近 80\%，行业监管趋严有利于头部企业提升市场份额。 优势在报表上主要体现为高于行业平均的毛利率，这在周期底部是抵御亏损的主要缓冲。

\pfl{历史股价} 自 2021 年 1 月至 2026 年 10 月，股价由 22.31 元变为 10.24 元，累计 -54.1\%，区间最高 31.89 元（2023 年 2 月）、最低 9.79 元（2026 年 6 月）；当前价相当于区间最高价的 32.1\%，已接近区间最低价（1.05 倍）。 月度收益的年化波动率 34.7\%，低于全样本中位数 37.3\%，在三级行业“动物保健Ⅲ”的 9 家公司中波动率由高到低排第\zhnumber{7}位。 把股价阶段与公司事件对齐看，2023 年 2 月 至 2026 年 6 月股价下跌-69.3\%，同期（2026 年 6 月）非公开发行募集资金净额 8.86 亿元截至 2026 年 6 月末累计投入 6.51 亿元（进度 73.44\%）、变更用途募资总额 33.07 万元。 整体上看，区间的几处大级别拐点都能在公司的资本运作、股权变动或风险事件上找到对应，股价波动有明显的公司层面驱动，而非单纯的行业 beta。

\begin{center}
\begin{minipage}{\textwidth}\centering
\begin{tikzpicture}
\begin{axis}[
  width=12.6cm, height=3.9cm, scale only axis,
  xmin=2020.922, xmax=2026.972, ymin=9.10, ymax=34.12,
  xtick={2021,2022,2023,2024,2025,2026},
  yearticks,
  yticklabel style={font=\scriptsize},
  ylabel={元/股}, ylabel style={font=\scriptsize},
  grid=major, grid style={LightGray, line width=0.3pt},
  axis line style={InkGray, line width=0.5pt},
  every axis plot/.append style={line width=0.9pt},
]
\addplot[AgriGreen] table[x=x,y=y]{data/clean/profiles/px_603566.dat};
\addplot[SoftRed, dashed, line width=0.7pt] table[x=x,y=y]{data/clean/profiles/pxzz_603566.dat};
\addplot[only marks, mark=*, mark size=1.1pt, SoftRed] table[x=x,y=y]{data/clean/profiles/pxzz_603566.dat};
\end{axis}
\end{tikzpicture}

\captionof{figure}[普莱柯（603566）股价与周期骨架]{普莱柯（603566）月度收盘价（元/股）与 ZigZag 周期骨架（阈值 25\%，圆点为周期转折点）}
\end{minipage}
\end{center}

\pfl{过去周期分析} 以 25\% 的 ZigZag 阈值划分，样本区间内股价形成 2 个主升段与 3 个主跌段。 最大主升段出现在 2021 年 10 月 至 2022 年 2 月，4 个月+58.7\%；最大主跌段为 2023 年 2 月 至 2026 年 6 月，40 个月-69.3\%。 最近一次转折出现在 2026 年 6 月（下跌段自 2023 年 2 月起，40 个月-69.3\%），当前处于该段的延续之中。 公司股价与申万农林牧渔指数几乎同步见底，个股并未走出独立行情。 动保的需求由存栏量而非价格决定（第 \ref{sec:vet} 节）：猪价下行阶段龙头养殖企业的逆势扩产反而推升了疫苗与兽药需求，这正是动保股价与猪价“脱钩”的原因，公司 2026 年上半年实现盈利、半年 ROE 3.0\%（未年化），好于全样本中位数（0.75\%），下行周期对其报表的冲击小于同业。 综合区间形态看，该股单段涨跌幅度偏大、以趋势段为主（最大主升 +58.7\%、最大主跌 -69.3\%），年化波动率 34.7\% 与全样本中位数 37.3\% 相比波动与样本中位数接近。

\pfl{盈利情况} 2026 年上半年营业收入 4.91 亿元（同比 -12.2\%），归母净利润 0.77 亿元，同比 -33.4\%。 以年报为趋势对照，2023---2025 年营业收入由 12.53 亿元变为 10.90 亿元（两年年均复合增速 -6.7\%），归母净利润由 1.75 亿元变为 1.85 亿元。 按半年报口径，2026 年上半年的 ROE 为 3.0\%（未年化）、销售净利率 13.0\%、毛利率 55.8\%，ROE 在“动物保健Ⅲ”9 家公司中按数值由高到低排第\zhnumber{5}位。 同期全样本半年 ROE 中位数 0.75\%，公司与之相差 2.20 个百分点（略好于中位数公司）。 从公司披露的原因看，业绩变动主要来自：2026 年上半年，归属于上市公司股东的净利润 7,732.42 万元、同比下降 33.45\%；2026 年半年报，主要系下游生猪出栏价格长期处于低位、动保行业竞争加剧、猪用疫苗销量和价格下降，同时生物制品增值税适用税率由 3\% 上调至 13\%。 面对这一局面，公司披露的举措包括：2024 年 9 月，公司存在因企业合并形成的商誉（收购普莱柯（南京）形成商誉资产组。 综合看，公司在报告期维持盈利（高于全样本半年 ROE 中位数 0.75\%），利润的可持续性取决于动物保健Ⅲ景气的延续与成本管控的执行。

\begin{center}\small\setlength{\tabcolsep}{3.4pt}
\captionof{table}[普莱柯（603566）关键财务指标]{普莱柯（603566）关键财务指标（合并报表口径；两个半年报期均未年化；现金短债比 $=$ 货币资金 $\div$ 短期债务）}
\begin{adjustbox}{max width=\textwidth}
\begin{tabular}{@{}lrrrrr@{}}
\toprule
指标 & 2023 年 & 2024 年 & 2025 年 & 2025 年半年报 & 2026 年半年报 \\
\midrule
营业收入（亿元） & 12.53 & 10.43 & 10.90 & 5.59 & 4.91 \\
归母净利润（亿元） & 1.75 & 0.93 & 1.85 & 1.16 & 0.77 \\
毛利率（\%） & 61.0 & 61.1 & 59.9 & 59.8 & 55.8 \\
ROE（\%） & 6.3 & 3.5 & 7.1 & 4.4 & 3.0 \\
资产负债率（\%） & 16.5 & 15.2 & 15.4 & 14.8 & 15.2 \\
经营活动现金流净额（亿元） & 2.83 & 2.69 & 3.56 & 1.95 & 0.78 \\
货币资金（亿元） & 3.17 & 2.05 & 1.37 & 1.20 & 0.86 \\
有息负债合计（亿元） & 0.02 & 0.01 & 0.27 & 0.01 & 0.26 \\
现金短债比（倍） & 318.06 & 189.62 & 19.77 & 110.46 & 11.86 \\
\bottomrule
\end{tabular}
\end{adjustbox}
\end{center}

\begin{center}
\begin{minipage}{\textwidth}\centering
\begin{tikzpicture}
\begin{groupplot}[
  group style={group size=2 by 1, horizontal sep=0.60cm},
  width=6.85cm, height=4.60cm,
  tick label style={font=\tiny},
  title style={font=\scriptsize\heiti},
  yticklabel style={font=\tiny},
  ylabel style={font=\tiny},
  grid=major, grid style={LightGray, line width=0.3pt},
  axis line style={InkGray, line width=0.5pt},
  enlarge x limits={abs=0.8},
  legend style={font=\scriptsize, at={(0.5,-0.20)}, anchor=north,
                legend columns=2, draw=none, fill=none, column sep=6pt},
]
\nextgroupplot[
  ybar, bar width=4pt,
  ymin=-1.25, ymax=14.03,
  xtick={0,1,2,3,4,5.7},
  xticklabels={2021,2022,2023,2024,2025,2026H1},
  ylabel={亿元},
]
\addplot[fill=AgriGreen!75, draw=AgriGreen, bar shift=-2.6pt] table[x=i, y=rev]{data/clean/profiles/fin_603566.dat};
\addplot[fill=SoftRed!60, draw=SoftRed, bar shift=2.6pt] table[x=i, y=ni]{data/clean/profiles/fin_603566.dat};
\addplot[fill=AgriGreen!30, draw=AgriGreen, pattern=north east lines, bar shift=-2.6pt] table[x=x, y=rev]{data/clean/profiles/finh_603566.dat};
\addplot[fill=SoftRed!20, draw=SoftRed, pattern=north east lines, bar shift=2.6pt] table[x=x, y=ni]{data/clean/profiles/finh_603566.dat};
\legend{营业收入（年/半年）, 归母净利润（年/半年）}
\nextgroupplot[
  ymin=-1.4, ymax=19.1,
  xtick={0,1,2,3,4,5.7},
  xticklabels={2021,2022,2023,2024,2025,2026H1},
  ylabel={\%},
]
\addplot[AgriGold, mark=*, mark size=1.3pt] table[x=i, y=roe]{data/clean/profiles/fin_603566.dat};
\addplot[OilBlue, mark=square*, mark size=1.3pt] table[x=i, y=debt]{data/clean/profiles/fin_603566.dat};
\addplot[only marks, AgriGold, mark=*, mark size=2.0pt] table[x=x, y=roe]{data/clean/profiles/finh_603566.dat};
\addplot[only marks, OilBlue, mark=square*, mark size=2.0pt] table[x=x, y=debt]{data/clean/profiles/finh_603566.dat};
\legend{ROE, 资产负债率}
\end{groupplot}
\end{tikzpicture}
\begin{tikzpicture}
\begin{axis}[
  name=finbot,
  width=13.75cm, height=3.10cm, scale only axis,
  ymin=-0.32, ymax=3.55,
  xtick={0,1,2,3,4,5.7},
  xticklabels={2021,2022,2023,2024,2025,2026H1},
  tick label style={font=\tiny},
  yticklabel style={font=\tiny},
  ylabel={亿元}, ylabel style={font=\tiny},
  ybar, bar width=4pt,
  grid=major, grid style={LightGray, line width=0.3pt},
  axis line style={InkGray, line width=0.5pt},
  enlarge x limits={abs=0.8},
  legend style={font=\scriptsize, at={(0.5,-0.26)}, anchor=north,
                legend columns=2, draw=none, fill=none, column sep=10pt},
]
\addplot[fill=AgriGreen!55, draw=AgriGreen, bar shift=-3.0pt] table[x=i, y=cash]{data/clean/profiles/fin_603566.dat};
\addplot[fill=SoftRed!45, draw=SoftRed, bar shift=3.0pt] table[x=i, y=ibd]{data/clean/profiles/fin_603566.dat};
\addplot[fill=AgriGreen!25, draw=AgriGreen, pattern=north east lines, bar shift=-3.0pt] table[x=x, y=cash]{data/clean/profiles/finh_603566.dat};
\addplot[fill=SoftRed!20, draw=SoftRed, pattern=north east lines, bar shift=3.0pt] table[x=x, y=ibd]{data/clean/profiles/finh_603566.dat};
\legend{货币资金（左轴）, 有息负债（左轴）}
\end{axis}
\begin{axis}[
  at={(finbot.south east)}, anchor=south east,
  width=13.75cm, height=3.10cm, scale only axis,
  axis y line*=right, axis x line*=none, xtick=\empty,
  ymin=-0.93, ymax=4.63,
  tick label style={font=\tiny},
  yticklabel style={font=\tiny}, scaled y ticks=false,
  legend style={font=\scriptsize, at={(0.5,-0.44)}, anchor=north,
                legend columns=1, draw=none, fill=none},
]
\addplot[OilBlue, mark=*, mark size=2.2pt, line width=1.3pt] table[x=i, y=ocf]{data/clean/profiles/fin_603566.dat};
\addplot[only marks, OilBlue, mark=*, mark size=3.2pt] table[x=x, y=ocf]{data/clean/profiles/finh_603566.dat};
\legend{经营活动现金流净额（右轴，亿元，蓝点）}
\end{axis}
\end{tikzpicture}

\captionof{figure}[普莱柯（603566）近年主要财务指标]{普莱柯（603566）近年主要财务指标（左上：营业收入与归母净利润；右上：ROE 与资产负债率；下：货币资金、有息负债为左轴柱状，经营活动现金流净额为其右轴的蓝点折线。
2021---2025 年为年报口径，2026H1 为 2026 年半年报/半年末，与其前一年报之间留出间隔、以斜纹或空心点区分；半年报数据未年化）}
\end{minipage}
\end{center}

\pfl{财务分析} 2026 年 6 月末资产负债率 15.2\%，较 2025 年末下降 0.2 个百分点（样本中位数 52.2\%，所处三级行业内按由低到高排第\zhnumber{3}位）。 2026 年 6 月末货币资金 0.86 亿元、短期债务（短期借款与一年内到期非流动负债）0.07 亿元、长期债务 0.19 亿元，有息负债合计 0.26 亿元，现金短债比 11.86 倍——覆盖充分。 净有息负债 -0.60 亿元，相当于其 2026 年上半年营业收入的 -12\%。 流动比率 3.46、速动比率 2.86，短期偿债能力处于正常区间。 2026 年上半年经营活动现金流净额 0.78 亿元，经营现金流与利润同向为正，内生资金可以支撑运营。 综合看，现金短债比 11.86 倍，短期偿付压力可控；资产负债率低于样本中位数 37.0 个百分点；有息负债中长期债务占比更高，期限结构对短债滚动的压力较小。

\pfl{重大战略转型} 从公开披露看，公司的战略推进集中在：融资与资本运作（1 项）：2020 年 6 月，公司以非公开发行募集资金投向兽用灭活疫苗生产项目。 公告口径上，近三年（2023 年 9 月---2026 年 9 月）公司披露重大重组与并购类公告 2 条、对外投资与转型类公告 0 条、回购与股权激励类公告 20 条，合计公告 306 条。 第一大行业仍是“兽用药品”，收入占比 94.0\%，与 2021 年年报相比没有方向性变化。 综合判断，报告期内公司有 2 条与战略相关的公告落地（重大重组与并购 2 条、对外投资与转型 0 条），资本运作与主业调整之间存在直接关联。

\begin{center}\small\setlength{\tabcolsep}{4pt}
\captionof{table}[普莱柯（603566）历史融资与资本运作]{普莱柯（603566）历史融资与资本运作（据公司公告与公开报道整理；金额为披露值，含首次公开发行、再融资、债券、银行借款与重整投资等）}
\begin{adjustbox}{max width=\textwidth}
\begin{tabular}{@{}p{1.45cm}p{2.55cm}p{4.05cm}p{4.95cm}@{}}
\toprule
时间 & 方式 & 规模 & 用途与说明 \\
\midrule
2015-05-18 & IPO（上海证券交易所主板） & 发行 4,\allowbreak{}000 万股，\allowbreak{}发行价 15.52 元/\allowbreak{}股；2015 年募集资金总额 5.60 亿元 & 动物疫苗产业化建设项目（新增 6 条生产线用于 17 个产品生产）、\allowbreak{}技术研发中心建设等项目；采用网下向投资者询价配售与网上按市值申购定价发行相结合的方式，\allowbreak{}发行后总股本不超 1.6 亿股；2015 年 8 月 20 日披露的募集资金使用情况显示募集资金总额 5.60 亿元 \\
2020-06-20 & 非公开发行股票（获核准批文） & 拟发行股票募集资金总额不超过 8.98 亿元 & 兽用灭活疫苗生产项目等；其中全资子公司普莱柯（南京）拟投资 4.00 亿元建设兽用灭活疫苗生产项目（一期） \\
2022-09-05 & 非公开发行股票（向特定对象） & 发行 3,\allowbreak{}142.06 万股，\allowbreak{}发行价 28.58 元/\allowbreak{}股，\allowbreak{}募集资金总额 8.98 亿元 & 兽用灭活疫苗生产项目等（截至 2026 年 6 月末累计投入 6.51 亿元、\allowbreak{}投入进度 73.44\%）；200 股；投资者于 2022 年 9 月 2 日 17:00 前足额缴款 \\
\bottomrule
\end{tabular}
\end{adjustbox}
\end{center}

\pfl{历史融投资情况} 从现金流量表看，2025 年公司，筹资活动现金流净额 -2.07 亿元（2023 年为偿还债务 0.07 亿元）。 2026 年上半年筹资活动现金流净额为 -1.37 亿元（取得借款 0.25 亿元、偿还债务 0.25 亿元，未年化），已转为净流出，处于净还债状态。 2025 年分配股利、利润或偿付利息支付的现金合计 2.05 亿元。 公开披露的融资记录共 3 笔，工具类型以 IPO、定向增发为主；金额与用途见下表。 其中规模最大的两笔为：2015 年 5 月，IPO，发行 4,000 万股，发行价 15.52 元/股，2015 年募集资金总额 5.60 亿元；2020 年 6 月，定向增发，拟发行股票募集资金总额不超过 8.98 亿元。 股本方面，近三年披露股本变动 8 次（限售股份上市 2 次、其他 2 次、股份回购 1 次）；总股本由 2021 年末的 3.21 亿股变为 3.46 亿股（+7.64\%）。 分红方面，近三年现金分红 6 次、合计每股派现 2.05 元，最近一次实施在 2025 年 12 月。 近三年“再融资”类公告 4 条，涉及定向增发。 综合看，公司融资结构上股权与债务融资并举；2025 年筹资活动现金流净额 -2.07 亿元，融资节奏仍以维持存量债务周转为主，与前述经营与投资现金流的方向基本一致。

\pfl{未来潜在融资需求分析} 2026 年 6 月末货币资金 0.86 亿元，而短期债务 0.07 亿元（现金短债比 11.86 倍），2026 年上半年经营现金流净额 0.78 亿元。账面货币资金可以覆盖短期债务，缺口不体现在静态偿付层面。 公司 2026 年 6 月末资产负债率 15.2\%、2026 年上半年 ROE 3.0\%（未年化），资产端与负债端都没有明显压力，融资需求以相机抉择的扩张型用途为主。 公开披露的后续资金安排线索包括：股权融资线索：2026 年 6 月，非公开发行募集资金净额 8.86 亿元截至 2026 年 6 月末累计投入 6.51 亿元（进度 73.44\%）、变更用途募资总额 33.07 万元。 以现有货币资金与经营现金流衡量，公司不存在刚性补血的紧迫性，融资动作更多是主动型的。




\subsubsection{申联生物（688098）}
\label{co:688098}

\pfl{公司基本情况} 2018 年设立至今，申联生物是专注口蹄疫疫苗的科创板动保企业（上海），2019 年 10 月 28 日上市、为首家科创板动保公司，主要产品为口蹄疫合成肽疫苗与口蹄疫灭活疫苗；上市后业绩持续承压并处于破发状态，2025 年营业收入 2.88 亿元、归母净利润 -2,049.50 万元（亏损收窄约 54\%），2026 年 2 月以 23,744.1789 万元收购扬州世之源控股权、开始构建“人用药品+动物保健”双主业。 公司于 2019-10-28 在上交所科创板首发上市，注册资本 4.11 亿元，总股本 4.11 亿股。 截至 2026 年 10 月 12 日收盘价 8.21 元，对应总市值 33.7 亿元，近一年-20.8\%，总市值在三级行业“动物保健Ⅲ”的 9 家公司中按由大到小排第\zhnumber{7}位，近一年涨跌幅低于全样本中位数（-10.4\%）10.4 个百分点。 按 2025 年年报披露的行业构成，收入占比前三位为兽用生物制品 99.3\%；其他 0.5\%；其他(补充) 0.2\%。

\pfl{股权结构与控制权} 按 2026 年 6 月 30 日的股东名册，前三大股东为杨玉芳 持股 21.93\%；UNITEDBIOMEDICAL,INC. 持股 7.52\%；杨从州 持股 7.31\%，股东户数 9,297 户。 股权结构的另一条线索是，持股 21.93\%（9,006.32 万股）；INC.持股 7.52\%（3,088.07 万股）。

\pfl{经营要点} 从经营要点看，首发募集资金全部投向悬浮培养口蹄疫灭活疫苗项目。

\pfl{所处行业的基本面} 动保的需求由下游存栏规模而非养殖盈利决定，供给端则由产品批文与研发能力设卡，疫苗是其中格局最稳定的细分（第 \ref{sec:framework} 章的供需框架）。 当期该行业的关键数据是：新兽药从研发到获批上市通常需 7 至 8 年；2026 年上半年，生猪价格长期低位运行；2026 年二季度，末全国能繁母猪存栏 3,780 万头（同比-6.5\%），为正常保有量 3,750 万头的 100.8\%。

\pfl{竞争优势} 把该公司放回同类公司比较，2026 年上半年营业收入 1.1 亿元、在三级行业“动物保健Ⅲ”的 9 家公司中排第\zhnumber{9}位；毛利率 57.4\%、高于行业内中位数 17.6 个百分点；半年 ROE -1.4\%（未年化）、低于行业内中位数（2.95\%）；资产负债率 9.8\%、低于行业内中位数（30.2\%）。 优势在报表上主要体现为高于行业平均的毛利率，这在周期底部是抵御亏损的主要缓冲。

\pfl{历史股价} 自 2021 年 1 月至 2026 年 10 月，股价由 14.49 元变为 8.21 元，累计 -43.3\%，区间最高 14.61 元（2021 年 2 月）、最低 3.92 元（2024 年 8 月）；当前价相当于区间最高价的 56.2\%，为区间最低价的 2.09 倍。 月度收益的年化波动率 38.0\%，高于全样本中位数 37.3\%，在三级行业“动物保健Ⅲ”的 9 家公司中波动率由高到低排第\zhnumber{4}位。

\begin{center}
\begin{minipage}{\textwidth}\centering
\begin{tikzpicture}
\begin{axis}[
  width=12.6cm, height=3.9cm, scale only axis,
  xmin=2020.922, xmax=2026.972, ymin=3.65, ymax=15.63,
  xtick={2021,2022,2023,2024,2025,2026},
  yearticks,
  yticklabel style={font=\scriptsize},
  ylabel={元/股}, ylabel style={font=\scriptsize},
  grid=major, grid style={LightGray, line width=0.3pt},
  axis line style={InkGray, line width=0.5pt},
  every axis plot/.append style={line width=0.9pt},
]
\addplot[AgriGreen] table[x=x,y=y]{data/clean/profiles/px_688098.dat};
\addplot[SoftRed, dashed, line width=0.7pt] table[x=x,y=y]{data/clean/profiles/pxzz_688098.dat};
\addplot[only marks, mark=*, mark size=1.1pt, SoftRed] table[x=x,y=y]{data/clean/profiles/pxzz_688098.dat};
\end{axis}
\end{tikzpicture}

\captionof{figure}[申联生物（688098）股价与周期骨架]{申联生物（688098）月度收盘价（元/股）与 ZigZag 周期骨架（阈值 25\%，圆点为周期转折点）}
\end{minipage}
\end{center}

\pfl{过去周期分析} 以 25\% 的 ZigZag 阈值划分，样本区间内股价形成 1 个主升段与 2 个主跌段。 最大主跌段为 2021 年 1 月 至 2024 年 8 月，43 个月-72.9\%；最大主升段出现在 2024 年 8 月 至 2025 年 8 月，12 个月+215.8\%。 最近一次转折出现在 2026 年 5 月（下跌段自 2025 年 8 月起，9 个月-37.2\%），当前处于该段的延续之中。 公司股价较申万农林牧渔指数提前约 21 个月见底，是板块中较早定价周期反转的公司之一。 动保的需求由存栏量而非价格决定（第 \ref{sec:vet} 节）：猪价下行阶段龙头养殖企业的逆势扩产反而推升了疫苗与兽药需求，这正是动保股价与猪价“脱钩”的原因，公司 2026 年 6 月末资产负债率 9.8\%、半年 ROE -1.4\%（未年化），在本轮下行中属于随行业同步调整的一类。 综合区间形态看，该股单段涨跌幅度偏大、以趋势段为主（最大主升 +215.8\%、最大主跌 -72.9\%），年化波动率 38.0\% 与全样本中位数 37.3\% 相比波动与样本中位数接近。

\pfl{盈利情况} 2026 年上半年营业收入 1.12 亿元（同比 -8.1\%），归母净利润 -0.19 亿元，亏损同比扩大 0.06 亿元。 以年报为趋势对照，2023---2025 年营业收入由 3.01 亿元变为 2.88 亿元（两年年均复合增速 -2.2\%），归母净利润由 0.32 亿元变为 -0.20 亿元。 按半年报口径，2026 年上半年的 ROE 为 -1.4\%（未年化）、销售净利率 -18.2\%、毛利率 57.4\%，ROE 在“动物保健Ⅲ”9 家公司中按数值由高到低排第\zhnumber{9}位。 同期全样本半年 ROE 中位数 0.75\%，公司与之相差 2.14 个百分点（弱于中位数公司）。 从公司披露的原因看，业绩变动主要来自：2025 年，归属于上市公司股东的净利润 -2,049.50 万元，原因是母公司存在未弥补亏损。 综合看，公司仍处亏损区间、亏损同比扩大，半年 ROE 在三级行业内按由高到低排第\zhnumber{9}位，单位盈利能力的修复依赖行业景气与自身成本端的改善，报表弹性主要体现为价格与销量的方向性变化。

\begin{center}\small\setlength{\tabcolsep}{3.4pt}
\captionof{table}[申联生物（688098）关键财务指标]{申联生物（688098）关键财务指标（合并报表口径；两个半年报期均未年化；现金短债比 $=$ 货币资金 $\div$ 短期债务）}
\begin{adjustbox}{max width=\textwidth}
\begin{tabular}{@{}lrrrrr@{}}
\toprule
指标 & 2023 年 & 2024 年 & 2025 年 & 2025 年半年报 & 2026 年半年报 \\
\midrule
营业收入（亿元） & 3.01 & 3.03 & 2.88 & 1.22 & 1.12 \\
归母净利润（亿元） & 0.32 & -0.45 & -0.20 & -0.13 & -0.19 \\
毛利率（\%） & 70.8 & 60.2 & 56.8 & 58.8 & 57.4 \\
ROE（\%） & 2.1 & -3.1 & -1.5 & -0.9 & -1.4 \\
资产负债率（\%） & 5.5 & 8.3 & 9.3 & 8.3 & 9.8 \\
经营活动现金流净额（亿元） & 0.04 & 0.89 & 0.33 & -0.39 & -0.65 \\
货币资金（亿元） & 0.45 & 1.46 & 1.22 & 0.53 & 0.65 \\
有息负债合计（亿元） & 0.01 & 0.15 & 0.14 & 0.15 & 0.53 \\
现金短债比（倍） & 134.73 & 10.09 & 9.20 & 3.80 & 1.23 \\
\bottomrule
\end{tabular}
\end{adjustbox}
\end{center}

\begin{center}
\begin{minipage}{\textwidth}\centering
\begin{tikzpicture}
\begin{groupplot}[
  group style={group size=2 by 1, horizontal sep=0.60cm},
  width=6.85cm, height=4.60cm,
  tick label style={font=\tiny},
  title style={font=\scriptsize\heiti},
  yticklabel style={font=\tiny},
  ylabel style={font=\tiny},
  grid=major, grid style={LightGray, line width=0.3pt},
  axis line style={InkGray, line width=0.5pt},
  enlarge x limits={abs=0.8},
  legend style={font=\scriptsize, at={(0.5,-0.20)}, anchor=north,
                legend columns=2, draw=none, fill=none, column sep=6pt},
]
\nextgroupplot[
  ybar, bar width=4pt,
  ymin=-0.85, ymax=4.07,
  xtick={0,1,2,3,4,5.7},
  xticklabels={2021,2022,2023,2024,2025,2026H1},
  ylabel={亿元},
]
\addplot[fill=AgriGreen!75, draw=AgriGreen, bar shift=-2.6pt] table[x=i, y=rev]{data/clean/profiles/fin_688098.dat};
\addplot[fill=SoftRed!60, draw=SoftRed, bar shift=2.6pt] table[x=i, y=ni]{data/clean/profiles/fin_688098.dat};
\addplot[fill=AgriGreen!30, draw=AgriGreen, pattern=north east lines, bar shift=-2.6pt] table[x=x, y=rev]{data/clean/profiles/finh_688098.dat};
\addplot[fill=SoftRed!20, draw=SoftRed, pattern=north east lines, bar shift=2.6pt] table[x=x, y=ni]{data/clean/profiles/finh_688098.dat};
\legend{营业收入（年/半年）, 归母净利润（年/半年）}
\nextgroupplot[
  ymin=-4.1, ymax=11.1,
  xtick={0,1,2,3,4,5.7},
  xticklabels={2021,2022,2023,2024,2025,2026H1},
  ylabel={\%},
]
\addplot[AgriGold, mark=*, mark size=1.3pt] table[x=i, y=roe]{data/clean/profiles/fin_688098.dat};
\addplot[OilBlue, mark=square*, mark size=1.3pt] table[x=i, y=debt]{data/clean/profiles/fin_688098.dat};
\addplot[only marks, AgriGold, mark=*, mark size=2.0pt] table[x=x, y=roe]{data/clean/profiles/finh_688098.dat};
\addplot[only marks, OilBlue, mark=square*, mark size=2.0pt] table[x=x, y=debt]{data/clean/profiles/finh_688098.dat};
\legend{ROE, 资产负债率}
\end{groupplot}
\end{tikzpicture}
\begin{tikzpicture}
\begin{axis}[
  name=finbot,
  width=13.75cm, height=3.10cm, scale only axis,
  ymin=-0.15, ymax=1.64,
  xtick={0,1,2,3,4,5.7},
  xticklabels={2021,2022,2023,2024,2025,2026H1},
  tick label style={font=\tiny},
  yticklabel style={font=\tiny},
  ylabel={亿元}, ylabel style={font=\tiny},
  ybar, bar width=4pt,
  grid=major, grid style={LightGray, line width=0.3pt},
  axis line style={InkGray, line width=0.5pt},
  enlarge x limits={abs=0.8},
  legend style={font=\scriptsize, at={(0.5,-0.26)}, anchor=north,
                legend columns=2, draw=none, fill=none, column sep=10pt},
]
\addplot[fill=AgriGreen!55, draw=AgriGreen, bar shift=-3.0pt] table[x=i, y=cash]{data/clean/profiles/fin_688098.dat};
\addplot[fill=SoftRed!45, draw=SoftRed, bar shift=3.0pt] table[x=i, y=ibd]{data/clean/profiles/fin_688098.dat};
\addplot[fill=AgriGreen!25, draw=AgriGreen, pattern=north east lines, bar shift=-3.0pt] table[x=x, y=cash]{data/clean/profiles/finh_688098.dat};
\addplot[fill=SoftRed!20, draw=SoftRed, pattern=north east lines, bar shift=3.0pt] table[x=x, y=ibd]{data/clean/profiles/finh_688098.dat};
\legend{货币资金（左轴）, 有息负债（左轴）}
\end{axis}
\begin{axis}[
  at={(finbot.south east)}, anchor=south east,
  width=13.75cm, height=3.10cm, scale only axis,
  axis y line*=right, axis x line*=none, xtick=\empty,
  ymin=-1.05, ymax=1.37,
  tick label style={font=\tiny},
  yticklabel style={font=\tiny}, scaled y ticks=false,
  legend style={font=\scriptsize, at={(0.5,-0.44)}, anchor=north,
                legend columns=1, draw=none, fill=none},
]
\addplot[OilBlue, mark=*, mark size=2.2pt, line width=1.3pt] table[x=i, y=ocf]{data/clean/profiles/fin_688098.dat};
\addplot[only marks, OilBlue, mark=*, mark size=3.2pt] table[x=x, y=ocf]{data/clean/profiles/finh_688098.dat};
\legend{经营活动现金流净额（右轴，亿元，蓝点）}
\end{axis}
\end{tikzpicture}

\captionof{figure}[申联生物（688098）近年主要财务指标]{申联生物（688098）近年主要财务指标（左上：营业收入与归母净利润；右上：ROE 与资产负债率；下：货币资金、有息负债为左轴柱状，经营活动现金流净额为其右轴的蓝点折线。
2021---2025 年为年报口径，2026H1 为 2026 年半年报/半年末，与其前一年报之间留出间隔、以斜纹或空心点区分；半年报数据未年化）}
\end{minipage}
\end{center}

\pfl{财务分析} 2026 年 6 月末资产负债率 9.8\%，较 2025 年末上升 0.5 个百分点（样本中位数 52.2\%，所处三级行业内按由低到高排第\zhnumber{1}位）。 2026 年 6 月末货币资金 0.65 亿元、短期债务（短期借款与一年内到期非流动负债）0.53 亿元、长期债务 0.01 亿元，有息负债合计 0.53 亿元，现金短债比 1.23 倍——覆盖充分。 净有息负债 -0.11 亿元，相当于其 2026 年上半年营业收入的 -10\%。 流动比率 3.54、速动比率 2.72，短期流动性无明显缺口。 2026 年上半年经营活动现金流净额 -0.65 亿元，经营现金流与利润同时为负，现金消耗较快。 综合看，现金短债比 1.23 倍，短期偿付压力可控；资产负债率低于样本中位数 42.4 个百分点；有息负债中短期债务占比更高，期限结构仍以滚动续作为主。

\pfl{重大战略转型} 从公开披露看，公司的战略推进集中在：产能与项目（1 项）：2025 年 12 月，公司 2025 年疫苗销售头份数同比增长 2.42\%、亏损幅度收窄约 54\%；并购与重组（1 项）：2026 年 2 月，公司收购扬州世之源控股权，投资方式为股权受让+增资并引入员工持股平台。 公告口径上，近三年（2023 年 9 月---2026 年 9 月）公司披露重大重组与并购类公告 2 条、对外投资与转型类公告 3 条、回购与股权激励类公告 22 条，合计公告 243 条。 第一大行业仍是“兽用生物制品”，收入占比 99.3\%，与 2021 年年报相比没有方向性变化。 综合判断，报告期内公司有 5 条与战略相关的公告落地（重大重组与并购 2 条、对外投资与转型 3 条），资本运作与主业调整之间存在直接关联。

\begin{center}\small\setlength{\tabcolsep}{4pt}
\captionof{table}[申联生物（688098）历史融资与资本运作]{申联生物（688098）历史融资与资本运作（据公司公告与公开报道整理；金额为披露值，含首次公开发行、再融资、债券、银行借款与重整投资等）}
\begin{adjustbox}{max width=\textwidth}
\begin{tabular}{@{}p{1.45cm}p{2.55cm}p{4.05cm}p{4.95cm}@{}}
\toprule
时间 & 方式 & 规模 & 用途与说明 \\
\midrule
2019-10-16 & IPO（上海证券交易所科创板） & 发行 5,\allowbreak{}000 万股（不低于发行后总股本的 10.00\% & 悬浮培养口蹄疫灭活疫苗项目（招股说明书原拟募集资金 4.50 亿元）；发行市盈率（按发行后总股本）43.87 倍 \\
2023-08-30 & 首发募集资金使用（募投项目） & 募集资金 4.40 亿元、\allowbreak{}净额 400,\allowbreak{}175，\allowbreak{}0 元量级 & 增加并扩大公司口蹄疫疫苗种类及生产规模；公司预先以自筹资金开展募投项目 \\
2026-02-25 & 收购联营公司控股权并开展新业务（关联交易，含银行并购贷款安排） & 公司拟使用 2.37 亿元的自有资金及自筹资金（包括银行并购贷款） & 加速推进公司创新药战略布局 \\
\bottomrule
\end{tabular}
\end{adjustbox}
\end{center}

\pfl{历史融投资情况} 从现金流量表看，2025 年公司偿还债务支付的现金 0.01 亿元，筹资活动现金流净额 -0.04 亿元。 2026 年上半年筹资活动现金流净额为 0.55 亿元（取得借款 0.40 亿元、偿还债务 0.01 亿元，未年化），仍处于净融入状态。 2025 年分配股利、利润或偿付利息支付的现金合计 0.00 亿元。 公开披露的融资记录共 3 笔，工具类型以 IPO 为主；金额与用途见下表。 其中规模最大的两笔为：2026 年 2 月，收购联营公司控股权并，公司拟使用 2.37 亿元的自有资金及自筹资金（包括银行并购贷款）；2019 年 10 月，IPO，发行 5,000 万股（不低于发行后总股本的 10.00\%。 股本方面，近三年披露股本变动 3 次；总股本自 2021 年末以来未发生变化（4.11 亿股）。 分红方面，近三年现金分红 2 次、合计每股派现 0.11 元，最近一次实施在 2023 年 12 月。 近三年“再融资”类公告 3 条，涉及定向增发。 综合看，公司融资结构上股权与债务融资并举；2025 年筹资活动现金流净额 -0.04 亿元，融资节奏仍以维持存量债务周转为主，与前述经营与投资现金流的方向基本一致。

\pfl{未来潜在融资需求分析} 2026 年 6 月末货币资金 0.65 亿元，而短期债务 0.53 亿元（现金短债比 1.23 倍），2026 年上半年经营现金流净额 -0.65 亿元。账面货币资金可以覆盖短期债务，缺口不体现在静态偿付层面。 公司 2026 年上半年归母净利润为负（-0.19 亿元），但 6 月末资产负债率 9.8\% 尚未进入第一档（亏损且负债率 $>$65\%）的区间，当前融资需求以补充流动资金、支撑低谷期经营性支出为主。 公开披露的后续资金安排线索包括：债务与授信安排：2026 年 2 月，公司明确收购世之源控股权的 2.37 亿元资金来源于自有资金及自筹资金，其中包含银行并购贷款。 以现有货币资金与经营现金流衡量，公司不存在刚性补血的紧迫性，融资动作更多是主动型的。




\subsubsection{回盛生物（300871）}
\label{co:300871}

\pfl{公司基本情况} 武汉的兽药企业，全称武汉回盛生物科技股份有限公司，2002 年 1 月 25 日成立，2020 年 8 月 24 日在深圳证券交易所创业板上市；主营兽用药品（化药制剂、原料药、中药制剂）及饲料与添加剂，是农业产业化国家重点龙头企业，截至 2025 年 12 月 31 日已取得兽药批准文号 317 个，产品应用于生猪、家禽、水产、反刍及宠物等领域。控股股东为武汉统盛投资有限公司（持股 41.03\%）。2025 年归母净利润 2.53 亿元、同比增长 1356.01\%（2024 年为亏损 0.20 亿元）；2026 年上半年归母净利润 2.06 亿元、同比增长 75.59\%。2026 年 4 月推出逾两年的 2024 年度定增方案到期失效，5 月被深交所终止审核。 公司于 2020-08-24 在深交所创业板首发上市，注册资本 2.02 亿元，总股本 1.66 亿股。 截至 2026 年 10 月 12 日收盘价 20.13 元，对应总市值 33.4 亿元，近一年-5.6\%，总市值在三级行业“动物保健Ⅲ”的 9 家公司中按由大到小排第\zhnumber{8}位，近一年涨跌幅高于全样本中位数（-10.4\%）4.7 个百分点。 按 2025 年年报披露的行业构成，收入占比前两位为兽用原料药及制剂 95.5\%；其他 4.5\%。

\pfl{股权结构与控制权} 按 2026 年 6 月 30 日的股东名册，前三大股东为武汉统盛投资有限公司 持股 41.03\%；梁栋国 持股 5.19\%；浙商银行股份有限公司－国泰中证畜牧养殖交易型开放式指 持股 1.38\%，股东户数 19,274 户。 股东结构与控制权方面，占总股本 41.03\%，其中质押 420.00 万股；报告期内控股股东及实际控制人均未发生变更；报告期内公司控股股东或第一大股东及其一致行动人累计质押股份数量占其所持公司股份数量比例未达到 80\%。

\pfl{经营要点} 从经营要点看，海外销售同比增长 72.79\%。

\pfl{所处行业的基本面} 动保的需求由下游存栏规模而非养殖盈利决定，供给端则由产品批文与研发能力设卡，疫苗是其中格局最稳定的细分（第 \ref{sec:framework} 章的供需框架）。 行业层面的最新情况是：2024 年，兽药原料药销售前十企业销售额 85.16 亿元、占原料药总销售额的 48\%；2024 年，化药制剂销售前十企业销售额 83.40 亿元、占化药制剂总销售额的 28.71\%；2026 年 10 月，兽药原料价格指数（VPi）收于 64.94。

\pfl{竞争优势} 把该公司放回同类公司比较，2026 年上半年营业收入 8.6 亿元、在三级行业“动物保健Ⅲ”的 9 家公司中排第\zhnumber{4}位；毛利率 35.7\%、低于行业内中位数 4.0 个百分点；半年 ROE 9.1\%（未年化）、高于行业内中位数（2.95\%）；资产负债率 31.8\%、高于行业内中位数（30.2\%）。 公司在公开资料中自述与同业比较的相对优势在于：泰万菌素原料药市占率超过 50\%，大环内酯类原料药产能规模处于行业前列。 优势在当期体现为高于行业中位数的资本回报，但优势能否持续仍取决于周期与自身产能投放节奏。

\pfl{历史股价} 以 2021 年 1 月为起点、2026 年 10 月为终点，股价由 44.00 元变为 20.13 元，累计 -54.2\%，区间最高 52.92 元（2021 年 4 月）、最低 8.47 元（2024 年 8 月）；当前价相当于区间最高价的 38.0\%，为区间最低价的 2.38 倍。 月度收益的年化波动率 53.9\%，高于全样本中位数 37.3\%，在三级行业“动物保健Ⅲ”的 9 家公司中波动率由高到低排第\zhnumber{2}位。

\begin{center}
\begin{minipage}{\textwidth}\centering
\begin{tikzpicture}
\begin{axis}[
  width=12.6cm, height=3.9cm, scale only axis,
  xmin=2020.922, xmax=2026.972, ymin=7.88, ymax=56.62,
  xtick={2021,2022,2023,2024,2025,2026},
  yearticks,
  yticklabel style={font=\scriptsize},
  ylabel={元/股}, ylabel style={font=\scriptsize},
  grid=major, grid style={LightGray, line width=0.3pt},
  axis line style={InkGray, line width=0.5pt},
  every axis plot/.append style={line width=0.9pt},
]
\addplot[AgriGreen] table[x=x,y=y]{data/clean/profiles/px_300871.dat};
\addplot[SoftRed, dashed, line width=0.7pt] table[x=x,y=y]{data/clean/profiles/pxzz_300871.dat};
\addplot[only marks, mark=*, mark size=1.1pt, SoftRed] table[x=x,y=y]{data/clean/profiles/pxzz_300871.dat};
\end{axis}
\end{tikzpicture}

\captionof{figure}[回盛生物（300871）股价与周期骨架]{回盛生物（300871）月度收盘价（元/股）与 ZigZag 周期骨架（阈值 25\%，圆点为周期转折点）}
\end{minipage}
\end{center}

\pfl{过去周期分析} 以 25\% 的 ZigZag 阈值划分，样本区间内股价形成 4 个主升段与 4 个主跌段。 最大主跌段为 2021 年 4 月 至 2022 年 4 月，12 个月-69.4\%；最大主升段出现在 2024 年 8 月 至 2026 年 1 月，17 个月+247.8\%。 最近一次转折出现在 2026 年 10 月（下跌段自 2026 年 1 月起，9 个月-31.7\%），当前处于该段的延续之中。 公司股价较申万农林牧渔指数提前约 21 个月见底，是板块中较早定价周期反转的公司之一。 动保的需求由存栏量而非价格决定（第 \ref{sec:vet} 节）：猪价下行阶段龙头养殖企业的逆势扩产反而推升了疫苗与兽药需求，这正是动保股价与猪价“脱钩”的原因，公司 2026 年上半年实现盈利、半年 ROE 9.1\%（未年化），好于全样本中位数（0.75\%），下行周期对其报表的冲击小于同业。 综合区间形态看，该股单段涨跌幅度偏大、以趋势段为主（最大主升 +247.8\%、最大主跌 -69.4\%），年化波动率 53.9\% 与全样本中位数 37.3\% 相比波动明显大于样本中位数。

\pfl{盈利情况} 2026 年上半年营业收入 8.65 亿元（同比 +5.2\%），归母净利润 2.06 亿元，同比 +75.6\%。 以年报为趋势对照，2023---2025 年营业收入由 10.20 亿元变为 16.81 亿元（两年年均复合增速 +28.4\%），归母净利润由 0.17 亿元变为 2.53 亿元。 按半年报口径，2026 年上半年的 ROE 为 9.1\%（未年化）、销售净利率 23.9\%、毛利率 35.7\%，ROE 在“动物保健Ⅲ”9 家公司中按数值由高到低排第\zhnumber{2}位。 同期全样本半年 ROE 中位数 0.75\%，公司与之相差 8.31 个百分点（略好于中位数公司）。 综合看，公司在报告期维持盈利（高于全样本半年 ROE 中位数 0.75\%），利润的可持续性取决于动物保健Ⅲ景气的延续与成本管控的执行。

\begin{center}\small\setlength{\tabcolsep}{3.4pt}
\captionof{table}[回盛生物（300871）关键财务指标]{回盛生物（300871）关键财务指标（合并报表口径；两个半年报期均未年化；现金短债比 $=$ 货币资金 $\div$ 短期债务）}
\begin{adjustbox}{max width=\textwidth}
\begin{tabular}{@{}lrrrrr@{}}
\toprule
指标 & 2023 年 & 2024 年 & 2025 年 & 2025 年半年报 & 2026 年半年报 \\
\midrule
营业收入（亿元） & 10.20 & 12.00 & 16.81 & 8.22 & 8.65 \\
归母净利润（亿元） & 0.17 & -0.20 & 2.53 & 1.17 & 2.06 \\
毛利率（\%） & 20.9 & 16.3 & 31.0 & 27.6 & 35.7 \\
ROE（\%） & 1.1 & -1.4 & 13.4 & 7.0 & 9.1 \\
资产负债率（\%） & 49.3 & 51.4 & 27.3 & 30.4 & 31.8 \\
经营活动现金流净额（亿元） & 0.93 & 0.44 & 2.38 & 1.24 & 1.28 \\
货币资金（亿元） & 2.94 & 3.26 & 3.10 & 2.36 & 3.59 \\
有息负债合计（亿元） & 8.42 & 8.89 & 2.29 & 3.50 & 4.64 \\
现金短债比（倍） & 1.22 & 0.98 & 1.52 & 0.75 & 1.77 \\
\bottomrule
\end{tabular}
\end{adjustbox}
\end{center}

\begin{center}
\begin{minipage}{\textwidth}\centering
\begin{tikzpicture}
\begin{groupplot}[
  group style={group size=2 by 1, horizontal sep=0.60cm},
  width=6.85cm, height=4.60cm,
  tick label style={font=\tiny},
  title style={font=\scriptsize\heiti},
  yticklabel style={font=\tiny},
  ylabel style={font=\tiny},
  grid=major, grid style={LightGray, line width=0.3pt},
  axis line style={InkGray, line width=0.5pt},
  enlarge x limits={abs=0.8},
  legend style={font=\scriptsize, at={(0.5,-0.20)}, anchor=north,
                legend columns=2, draw=none, fill=none, column sep=6pt},
]
\nextgroupplot[
  ybar, bar width=4pt,
  ymin=-1.90, ymax=18.85,
  xtick={0,1,2,3,4,5.7},
  xticklabels={2021,2022,2023,2024,2025,2026H1},
  ylabel={亿元},
]
\addplot[fill=AgriGreen!75, draw=AgriGreen, bar shift=-2.6pt] table[x=i, y=rev]{data/clean/profiles/fin_300871.dat};
\addplot[fill=SoftRed!60, draw=SoftRed, bar shift=2.6pt] table[x=i, y=ni]{data/clean/profiles/fin_300871.dat};
\addplot[fill=AgriGreen!30, draw=AgriGreen, pattern=north east lines, bar shift=-2.6pt] table[x=x, y=rev]{data/clean/profiles/finh_300871.dat};
\addplot[fill=SoftRed!20, draw=SoftRed, pattern=north east lines, bar shift=2.6pt] table[x=x, y=ni]{data/clean/profiles/finh_300871.dat};
\legend{营业收入（年/半年）, 归母净利润（年/半年）}
\nextgroupplot[
  ymin=-5.6, ymax=56.7,
  xtick={0,1,2,3,4,5.7},
  xticklabels={2021,2022,2023,2024,2025,2026H1},
  ylabel={\%},
]
\addplot[AgriGold, mark=*, mark size=1.3pt] table[x=i, y=roe]{data/clean/profiles/fin_300871.dat};
\addplot[OilBlue, mark=square*, mark size=1.3pt] table[x=i, y=debt]{data/clean/profiles/fin_300871.dat};
\addplot[only marks, AgriGold, mark=*, mark size=2.0pt] table[x=x, y=roe]{data/clean/profiles/finh_300871.dat};
\addplot[only marks, OilBlue, mark=square*, mark size=2.0pt] table[x=x, y=debt]{data/clean/profiles/finh_300871.dat};
\legend{ROE, 资产负债率}
\end{groupplot}
\end{tikzpicture}
\begin{tikzpicture}
\begin{axis}[
  name=finbot,
  width=13.75cm, height=3.10cm, scale only axis,
  ymin=-1.07, ymax=11.96,
  xtick={0,1,2,3,4,5.7},
  xticklabels={2021,2022,2023,2024,2025,2026H1},
  tick label style={font=\tiny},
  yticklabel style={font=\tiny},
  ylabel={亿元}, ylabel style={font=\tiny},
  ybar, bar width=4pt,
  grid=major, grid style={LightGray, line width=0.3pt},
  axis line style={InkGray, line width=0.5pt},
  enlarge x limits={abs=0.8},
  legend style={font=\scriptsize, at={(0.5,-0.26)}, anchor=north,
                legend columns=2, draw=none, fill=none, column sep=10pt},
]
\addplot[fill=AgriGreen!55, draw=AgriGreen, bar shift=-3.0pt] table[x=i, y=cash]{data/clean/profiles/fin_300871.dat};
\addplot[fill=SoftRed!45, draw=SoftRed, bar shift=3.0pt] table[x=i, y=ibd]{data/clean/profiles/fin_300871.dat};
\addplot[fill=AgriGreen!25, draw=AgriGreen, pattern=north east lines, bar shift=-3.0pt] table[x=x, y=cash]{data/clean/profiles/finh_300871.dat};
\addplot[fill=SoftRed!20, draw=SoftRed, pattern=north east lines, bar shift=3.0pt] table[x=x, y=ibd]{data/clean/profiles/finh_300871.dat};
\legend{货币资金（左轴）, 有息负债（左轴）}
\end{axis}
\begin{axis}[
  at={(finbot.south east)}, anchor=south east,
  width=13.75cm, height=3.10cm, scale only axis,
  axis y line*=right, axis x line*=none, xtick=\empty,
  ymin=-0.62, ymax=3.10,
  tick label style={font=\tiny},
  yticklabel style={font=\tiny}, scaled y ticks=false,
  legend style={font=\scriptsize, at={(0.5,-0.44)}, anchor=north,
                legend columns=1, draw=none, fill=none},
]
\addplot[OilBlue, mark=*, mark size=2.2pt, line width=1.3pt] table[x=i, y=ocf]{data/clean/profiles/fin_300871.dat};
\addplot[only marks, OilBlue, mark=*, mark size=3.2pt] table[x=x, y=ocf]{data/clean/profiles/finh_300871.dat};
\legend{经营活动现金流净额（右轴，亿元，蓝点）}
\end{axis}
\end{tikzpicture}

\captionof{figure}[回盛生物（300871）近年主要财务指标]{回盛生物（300871）近年主要财务指标（左上：营业收入与归母净利润；右上：ROE 与资产负债率；下：货币资金、有息负债为左轴柱状，经营活动现金流净额为其右轴的蓝点折线。
2021---2025 年为年报口径，2026H1 为 2026 年半年报/半年末，与其前一年报之间留出间隔、以斜纹或空心点区分；半年报数据未年化）}
\end{minipage}
\end{center}

\pfl{财务分析} 2026 年 6 月末资产负债率 31.8\%，较 2025 年末上升 4.4 个百分点（样本中位数 52.2\%，所处三级行业内按由低到高排第\zhnumber{6}位）。 2026 年 6 月末货币资金 3.59 亿元、短期债务（短期借款与一年内到期非流动负债）2.02 亿元、长期债务 2.62 亿元，有息负债合计 4.64 亿元，现金短债比 1.77 倍——覆盖充分。 净有息负债 1.05 亿元，相当于其 2026 年上半年营业收入的 12\%。 流动比率 2.05、速动比率 1.67，短期偿债能力处于正常区间。 2026 年上半年经营活动现金流净额 1.28 亿元，经营现金流与利润同向为正，内生资金可以支撑运营。 2026 年 6 月末在建工程 1.58 亿元，较上年末增加 1.00 亿元，仍有在建投入。 综合看，现金短债比 1.77 倍，短期偿付压力可控；资产负债率低于样本中位数 20.4 个百分点；有息负债中长期债务占比更高，期限结构对短债滚动的压力较小。 公司披露的股东与质押风险为：2024 年 4 月，可能涉及股票质押作为并购贷的担保措施。

\pfl{重大战略转型} 从公开披露看，公司的战略推进集中在：融资与资本运作（2 项）：2026 年 4 月，2021 年可转债募投项目中包含“宠物制剂综合生产线建设项目”（拟投入 9,000.00 万元）；2024 年 4 月，公司拟向楚盛投资及实际控制人发行股票募集不超过 2.50 亿元全部补充流动资金；股权与股东回报（1 项）：2026 年 6 月，公司实际控制人、董事长张卫元提议回购公司股份。 公告口径上，近三年（2023 年 9 月---2026 年 9 月）公司披露重大重组与并购类公告 2 条、对外投资与转型类公告 3 条、回购与股权激励类公告 37 条，合计公告 499 条。 第一大行业仍是“兽用原料药及制剂”，收入占比 95.5\%，与 2021 年年报相比没有方向性变化。 综合判断，报告期内公司有 5 条与战略相关的公告落地（重大重组与并购 2 条、对外投资与转型 3 条），资本运作与主业调整之间存在直接关联。

\begin{center}\small\setlength{\tabcolsep}{4pt}
\captionof{table}[回盛生物（300871）历史融资与资本运作]{回盛生物（300871）历史融资与资本运作（据公司公告与公开报道整理；金额为披露值，含首次公开发行、再融资、债券、银行借款与重整投资等）}
\begin{adjustbox}{max width=\textwidth}
\begin{tabular}{@{}p{1.45cm}p{2.55cm}p{4.05cm}p{4.95cm}@{}}
\toprule
时间 & 方式 & 规模 & 用途与说明 \\
\midrule
2020-08 & 首次公开发行A股股票并在深圳证券交易所创业板上市（IPO） & 发行价 33.61 元/\allowbreak{}股；首发前总股本 8,\allowbreak{}280.70 万股 & 实际募集资金超出发行人预计募集资金。份有限公司 \\
2021-12-17 & 向不特定对象发行可转换公司债券（回盛转债，债券代码123132） & 发行 700.00 万张，\allowbreak{}每张面值 100.00 元 & 年产 1,\allowbreak{}000 吨泰乐菌素和年产 600 吨泰万菌素生产线扩建项目（拟投入 2.85 亿元 \\
2024-04-03 & 向特定对象发行A股股票（2024年度向特定对象发行，后终止） & 拟募集资金总额不超过人民币 2.50 亿元（含本数） & 全部用于补充流动资金。；2024 年 4 月 3 日第三届董事会第十次会议、\allowbreak{}第三届监事会第八次会议审议通过 \\
2024-04-25 & 向金融机构申请综合授信额度（关联担保） & 公司及子公司 2024 年度拟向银行等金融机构申请授信额度不超过人民币 10.00 亿元（含本数） & 融资方式包括银行流动资金贷款、\allowbreak{}固定资产投资贷款、\allowbreak{}保函、\allowbreak{}信用证、\allowbreak{}银行承兑汇票、\allowbreak{}各类商业票据开立及贴现、\allowbreak{}保理等授信业务。；公司公告明确该授信额度不等同于公司实际融资金额 \\
2026-04-17 & 2024年度向特定对象发行A股股票方案到期失效 & 终止拟募集不超过 2.50 亿元（补充流动资金）的发行计划 & 不适用。；公司 2026 年第二次临时股东会审议未通过延长公司 2024 年度向特定对象发行 A 股股票股东会决议有效期和提请股东会延长授权董事会及其授权人士全权办理公司向特定对象发行 A 股股票相关事宜有效期 \\
2026-05 & 深圳证券交易所终止审核2024年度向特定对象发行 & 终止审核公司向特定对象发行股票申请 & 不适用。；公司于 2026 年 4 月 17 日召开的 2026 年第二次临时股东会审议未通过延长决议有效期及授权的议案后 \\
2026-06-05 & 实际控制人、董事长提议回购公司股份 & 公司实际控制人、\allowbreak{}董事长张卫元先生于 2026 年 6 月 5 日出具提议回购公司股份 & 基于对公司价值的判断和未来发展前景的信心提议回购。；该事项处于提议阶段 \\
\bottomrule
\end{tabular}
\end{adjustbox}
\end{center}

\pfl{历史融投资情况} 从现金流量表看，2025 年公司取得借款收到的现金 2.20 亿元、偿还债务支付的现金 3.68 亿元、吸收权益性投资 0.00 亿元，筹资活动现金流净额 -1.60 亿元（2023 年为取得借款 2.11 亿元、偿还债务 2.44 亿元）。 2026 年上半年筹资活动现金流净额为 2.07 亿元（取得借款 4.24 亿元、偿还债务 2.10 亿元，未年化），仍处于净融入状态。 2025 年分配股利、利润或偿付利息支付的现金合计 0.50 亿元。 公开披露的融资记录共 7 笔，工具类型以 IPO、可转债、定向增发、银行授信为主；金额与用途见下表。 其中规模最大的两笔为：2024 年 4 月，银行授信，公司及子公司 2024 年度拟向银行等金融机构申请授信额度不超过人民币 10.00 亿元（含本数）；2021 年 12 月，可转债，发行 700.00 万张，每张面值 100.00 元。 股本方面，近三年披露股本变动 17 次（可转债转股 9 次、其他 3 次、股份回购 2 次、限售股份上市 1 次）；总股本由 2021 年末的 1.66 亿股变为 1.66 亿股（-0.29\%）。 分红方面，近三年现金分红 6 次、合计每股派现 1.28 元，最近一次实施在 2026 年 6 月。 近三年“再融资”类公告 46 条，涉及定向增发、可转债、募集资金使用。 综合看，公司融资结构上以债务滚动为主、股权融资为补充；2025 年筹资活动现金流净额 -1.60 亿元，融资节奏仍以维持存量债务周转为主，与前述经营与投资现金流的方向基本一致。

\pfl{未来潜在融资需求分析} 2026 年 6 月末货币资金 3.59 亿元，而短期债务 2.02 亿元（现金短债比 1.77 倍），2026 年上半年经营现金流净额 1.28 亿元。账面货币资金可以覆盖短期债务，缺口不体现在静态偿付层面。 按第 \ref{sec:financing-need} 节的三档划分，公司属于第三档“扩张型”：2026 年上半年 ROE 9.1\%（未年化，折合约 18.1\%）、6 月末资产负债率 31.8\%，资产负债表仍具备加杠杆空间；融资用途以产能扩张、品类并购与海外布局为主，单笔规模通常在 5---20 亿元量级，会被市场定价为成长而非补血。 公开披露的后续资金安排线索包括：股权融资线索：2026 年 5 月，公司 2024 年度向特定对象发行股票，已因股东会未通过延期而失效；债务与授信安排：2024 年 4 月，公司及子公司向金融机构申请综合授信额度不超过人民币 10.00 亿元。 资金需求还要叠加在建工程：期末在建工程 1.58 亿元、2026 年上半年购建固定资产等支出 2.35 亿元（未年化），这部分支出在周期底部通常会被主动放缓。 静态偿付不存在缺口，后续融资更可能服务于产能或并购安排，而不是填补流动性窟窿。




\subsubsection{绿康生化（002868）}
\label{co:002868}

\pfl{公司基本情况} 福建南平浦城起家的兽药（饲料添加剂、生物发酵）企业，2017 年 5 月深交所上市；2022—2023 年跨界收购光伏胶膜企业江西纬科并投建海宁胶膜项目，2022—2024 年连续三年扣非净利润为负、2024 年末经审计净资产为 -2,453.61 万元，2025 年 4 月 30 日起被实施退市风险警示（*ST 绿康）；2025 年 4 月原控股股东向福建纵腾网络转让 29.99\% 股份、实际控制人变更为王钻，公司以 0 元剥离三家光伏胶膜子公司，2026 年 5 月 13 日撤销退市风险警示及其他风险警示，2026 年上半年归母净利润 731.60 万元、扭亏为盈。 公司于 2017-05-03 在深交所主板首发上市，注册资本 1.55 亿元，总股本 1.55 亿股。 截至 2026 年 10 月 12 日收盘价 21.30 元，对应总市值 33.1 亿元，近一年-21.6\%，总市值在三级行业“动物保健Ⅲ”的 9 家公司中按由大到小排第\zhnumber{9}位，近一年涨跌幅低于全样本中位数（-10.4\%）11.2 个百分点。 按 2026 年半年报披露的行业构成，收入占比前三位为动保系列 91.4\%；其他 5.0\%；植保系列 3.7\%。 与 2021 年年报相比，第一大行业口径由“兽药”（67.9\%）变为“动保系列”（91.4\%）；两期披露的分类口径有调整，占比差异中同时包含分类因素。

\pfl{股权结构与控制权} 按 2026 年 6 月 30 日的股东名册，前三大股东为福建纵腾网络有限公司 持股 29.99\%；上海康怡投资有限公司 持股 12.85\%；董海 持股 5.00\%，股东户数 8,825 户。 股东结构与控制权方面，2026 年 6 月末持股 29.99\%（4,660.84 万股），其中质押 2,796.50 万股；2025 年 11 月，公司完成控制权变更。

\pfl{经营要点} 经营层面，动保系列实现收入 2.26 亿元、占营业总收入 91.37\%（同比+24.66\%）；公司于 2025 年成功剥离当年度处于深度亏损的光伏胶膜子公司。

\pfl{所处行业的基本面} 动保的需求由下游存栏规模而非养殖盈利决定，供给端则由产品批文与研发能力设卡，疫苗是其中格局最稳定的细分（第 \ref{sec:framework} 章的供需框架）。 从供需两端的当期数据看，2026 年上半年，生猪养殖企业普遍亏损；行业份额持续向头部企业集中。

\pfl{竞争优势} 把该公司放回同类公司比较，2026 年上半年营业收入 2.5 亿元、在三级行业“动物保健Ⅲ”的 9 家公司中排第\zhnumber{8}位；毛利率 22.9\%、低于行业内中位数 16.8 个百分点；半年 ROE 10.9\%（未年化）、高于行业内中位数（2.95\%）；资产负债率 90.1\%、高于行业内中位数（30.2\%）。 公司区别于同类公司的经营特征包括：2026 年上半年，动保产品毛利率 22.70\%（同比提升 1.09 个百分点）。 优势在当期体现为高于行业中位数的资本回报，但优势能否持续仍取决于周期与自身产能投放节奏。

\pfl{历史股价} 本报告样本区间（2021 年 1 月 至 2026 年 10 月）内，股价由 11.13 元变为 21.30 元，累计 +91.4\%，区间最高 61.22 元（2023 年 1 月）、最低 9.29 元（2021 年 10 月）；当前价相当于区间最高价的 34.8\%，为区间最低价的 2.29 倍。 月度收益的年化波动率 57.0\%，高于全样本中位数 37.3\%，在三级行业“动物保健Ⅲ”的 9 家公司中波动率由高到低排第\zhnumber{1}位。 把股价阶段与公司事件对齐看，2021 年 1 月 至 2023 年 1 月股价上涨+450.0\%，同期（2023 年 1 月）董事会审议通过以 9,500.00 万元受让江西纬科 100\% 股权；2025 年 3 月 至 2026 年 2 月股价上涨+189.2\%，同期（2025 年）公司完成剥离胶膜业务、回归兽药业务。 整体上看，区间的几处大级别拐点都能在公司的资本运作、股权变动或风险事件上找到对应，股价波动有明显的公司层面驱动，而非单纯的行业 beta。

\begin{center}
\begin{minipage}{\textwidth}\centering
\begin{tikzpicture}
\begin{axis}[
  width=12.6cm, height=3.9cm, scale only axis,
  xmin=2020.922, xmax=2026.972, ymin=8.64, ymax=65.51,
  xtick={2021,2022,2023,2024,2025,2026},
  yearticks,
  yticklabel style={font=\scriptsize},
  ylabel={元/股}, ylabel style={font=\scriptsize},
  grid=major, grid style={LightGray, line width=0.3pt},
  axis line style={InkGray, line width=0.5pt},
  every axis plot/.append style={line width=0.9pt},
]
\addplot[AgriGreen] table[x=x,y=y]{data/clean/profiles/px_002868.dat};
\addplot[SoftRed, dashed, line width=0.7pt] table[x=x,y=y]{data/clean/profiles/pxzz_002868.dat};
\addplot[only marks, mark=*, mark size=1.1pt, SoftRed] table[x=x,y=y]{data/clean/profiles/pxzz_002868.dat};
\end{axis}
\end{tikzpicture}

\captionof{figure}[绿康生化（002868）股价与周期骨架]{绿康生化（002868）月度收盘价（元/股）与 ZigZag 周期骨架（阈值 25\%，圆点为周期转折点）}
\end{minipage}
\end{center}

\pfl{过去周期分析} 以 25\% 的 ZigZag 阈值划分，样本区间内股价形成 3 个主升段与 3 个主跌段。 最大主跌段为 2023 年 1 月 至 2024 年 8 月，19 个月-79.9\%；最大主升段出现在 2021 年 1 月 至 2023 年 1 月，24 个月+450.0\%。 最近一次转折出现在 2026 年 10 月（下跌段自 2026 年 2 月起，8 个月-50.4\%），当前处于该段的延续之中。 公司股价较申万农林牧渔指数提前约 55 个月见底，是板块中较早定价周期反转的公司之一。 动保的需求由存栏量而非价格决定（第 \ref{sec:vet} 节）：猪价下行阶段龙头养殖企业的逆势扩产反而推升了疫苗与兽药需求，这正是动保股价与猪价“脱钩”的原因，公司 2026 年 6 月末资产负债率 90.1\%，高于全样本中位数 52.2\%，其股价因此更多反映资金面与信用风险，而不是价格弹性本身。 综合区间形态看，该股单段涨跌幅度偏大、以趋势段为主（最大主升 +450.0\%、最大主跌 -79.9\%），年化波动率 57.0\% 与全样本中位数 37.3\% 相比波动明显大于样本中位数。

\pfl{盈利情况} 2026 年上半年营业收入 2.48 亿元（同比 -15.7\%），归母净利润 0.07 亿元，实现扭亏。 以年报为趋势对照，2023---2025 年营业收入由 5.07 亿元变为 5.17 亿元（两年年均复合增速 +0.9\%），归母净利润由 -2.22 亿元变为 -1.39 亿元。 按半年报口径，2026 年上半年的 ROE 为 10.9\%（未年化）、销售净利率 2.8\%、毛利率 22.9\%，ROE 在“动物保健Ⅲ”9 家公司中按数值由高到低排第\zhnumber{1}位。 同期全样本半年 ROE 中位数 0.75\%，公司与之相差 10.16 个百分点（略好于中位数公司）。 从公司披露的原因看，业绩变动主要来自：2025 年，公司完成剥离胶膜业务、回归兽药业务；2025 年 12 月，盐霉素、莫能菌素等抗球虫药产能逐步释放，实现新旧产能系统性转换与升级。 综合看，公司在报告期维持盈利（高于全样本半年 ROE 中位数 0.75\%），利润的可持续性取决于动物保健Ⅲ景气的延续与成本管控的执行。

\begin{center}\small\setlength{\tabcolsep}{3.4pt}
\captionof{table}[绿康生化（002868）关键财务指标]{绿康生化（002868）关键财务指标（合并报表口径；两个半年报期均未年化；现金短债比 $=$ 货币资金 $\div$ 短期债务）}
\begin{adjustbox}{max width=\textwidth}
\begin{tabular}{@{}lrrrrr@{}}
\toprule
指标 & 2023 年 & 2024 年 & 2025 年 & 2025 年半年报 & 2026 年半年报 \\
\midrule
营业收入（亿元） & 5.07 & 6.49 & 5.17 & 2.94 & 2.48 \\
归母净利润（亿元） & -2.22 & -4.45 & -1.39 & -0.59 & 0.07 \\
毛利率（\%） & -9.0 & -8.4 & 6.7 & 5.9 & 22.9 \\
ROE（\%） & -45.1 & 0.0 & 170.5 & 0.0 & 10.9 \\
资产负债率（\%） & 78.9 & 101.6 & 91.3 & 105.8 & 90.1 \\
经营活动现金流净额（亿元） & -1.87 & 1.41 & -0.25 & 0.75 & 0.45 \\
货币资金（亿元） & 0.84 & 0.25 & 0.56 & 0.16 & 0.51 \\
有息负债合计（亿元） & 9.63 & 9.16 & 3.60 & 8.66 & 3.19 \\
现金短债比（倍） & 0.19 & 0.03 & 0.17 & 0.02 & 0.17 \\
\bottomrule
\end{tabular}
\end{adjustbox}
\end{center}

\begin{center}
\begin{minipage}{\textwidth}\centering
\begin{tikzpicture}
\begin{groupplot}[
  group style={group size=2 by 1, horizontal sep=0.60cm},
  width=6.85cm, height=4.60cm,
  tick label style={font=\tiny},
  title style={font=\scriptsize\heiti},
  yticklabel style={font=\tiny},
  ylabel style={font=\tiny},
  grid=major, grid style={LightGray, line width=0.3pt},
  axis line style={InkGray, line width=0.5pt},
  enlarge x limits={abs=0.8},
  legend style={font=\scriptsize, at={(0.5,-0.20)}, anchor=north,
                legend columns=2, draw=none, fill=none, column sep=6pt},
]
\nextgroupplot[
  ybar, bar width=4pt,
  ymin=-5.54, ymax=7.81,
  xtick={0,1,2,3,4,5.7},
  xticklabels={2021,2022,2023,2024,2025,2026H1},
  ylabel={亿元},
]
\addplot[fill=AgriGreen!75, draw=AgriGreen, bar shift=-2.6pt] table[x=i, y=rev]{data/clean/profiles/fin_002868.dat};
\addplot[fill=SoftRed!60, draw=SoftRed, bar shift=2.6pt] table[x=i, y=ni]{data/clean/profiles/fin_002868.dat};
\addplot[fill=AgriGreen!30, draw=AgriGreen, pattern=north east lines, bar shift=-2.6pt] table[x=x, y=rev]{data/clean/profiles/finh_002868.dat};
\addplot[fill=SoftRed!20, draw=SoftRed, pattern=north east lines, bar shift=2.6pt] table[x=x, y=ni]{data/clean/profiles/finh_002868.dat};
\legend{营业收入（年/半年）, 归母净利润（年/半年）}
\nextgroupplot[
  ymin=-62.4, ymax=192.0,
  xtick={0,1,2,3,4,5.7},
  xticklabels={2021,2022,2023,2024,2025,2026H1},
  ylabel={\%},
]
\addplot[AgriGold, mark=*, mark size=1.3pt] table[x=i, y=roe]{data/clean/profiles/fin_002868.dat};
\addplot[OilBlue, mark=square*, mark size=1.3pt] table[x=i, y=debt]{data/clean/profiles/fin_002868.dat};
\addplot[only marks, AgriGold, mark=*, mark size=2.0pt] table[x=x, y=roe]{data/clean/profiles/finh_002868.dat};
\addplot[only marks, OilBlue, mark=square*, mark size=2.0pt] table[x=x, y=debt]{data/clean/profiles/finh_002868.dat};
\legend{ROE, 资产负债率}
\end{groupplot}
\end{tikzpicture}
\begin{tikzpicture}
\begin{axis}[
  name=finbot,
  width=13.75cm, height=3.10cm, scale only axis,
  ymin=-0.96, ymax=10.78,
  xtick={0,1,2,3,4,5.7},
  xticklabels={2021,2022,2023,2024,2025,2026H1},
  tick label style={font=\tiny},
  yticklabel style={font=\tiny},
  ylabel={亿元}, ylabel style={font=\tiny},
  ybar, bar width=4pt,
  grid=major, grid style={LightGray, line width=0.3pt},
  axis line style={InkGray, line width=0.5pt},
  enlarge x limits={abs=0.8},
  legend style={font=\scriptsize, at={(0.5,-0.26)}, anchor=north,
                legend columns=2, draw=none, fill=none, column sep=10pt},
]
\addplot[fill=AgriGreen!55, draw=AgriGreen, bar shift=-3.0pt] table[x=i, y=cash]{data/clean/profiles/fin_002868.dat};
\addplot[fill=SoftRed!45, draw=SoftRed, bar shift=3.0pt] table[x=i, y=ibd]{data/clean/profiles/fin_002868.dat};
\addplot[fill=AgriGreen!25, draw=AgriGreen, pattern=north east lines, bar shift=-3.0pt] table[x=x, y=cash]{data/clean/profiles/finh_002868.dat};
\addplot[fill=SoftRed!20, draw=SoftRed, pattern=north east lines, bar shift=3.0pt] table[x=x, y=ibd]{data/clean/profiles/finh_002868.dat};
\legend{货币资金（左轴）, 有息负债（左轴）}
\end{axis}
\begin{axis}[
  at={(finbot.south east)}, anchor=south east,
  width=13.75cm, height=3.10cm, scale only axis,
  axis y line*=right, axis x line*=none, xtick=\empty,
  ymin=-2.72, ymax=2.40,
  tick label style={font=\tiny},
  yticklabel style={font=\tiny}, scaled y ticks=false,
  legend style={font=\scriptsize, at={(0.5,-0.44)}, anchor=north,
                legend columns=1, draw=none, fill=none},
]
\addplot[OilBlue, mark=*, mark size=2.2pt, line width=1.3pt] table[x=i, y=ocf]{data/clean/profiles/fin_002868.dat};
\addplot[only marks, OilBlue, mark=*, mark size=3.2pt] table[x=x, y=ocf]{data/clean/profiles/finh_002868.dat};
\legend{经营活动现金流净额（右轴，亿元，蓝点）}
\end{axis}
\end{tikzpicture}

\captionof{figure}[绿康生化（002868）近年主要财务指标]{绿康生化（002868）近年主要财务指标（左上：营业收入与归母净利润；右上：ROE 与资产负债率；下：货币资金、有息负债为左轴柱状，经营活动现金流净额为其右轴的蓝点折线。
2021---2025 年为年报口径，2026H1 为 2026 年半年报/半年末，与其前一年报之间留出间隔、以斜纹或空心点区分；半年报数据未年化）}
\end{minipage}
\end{center}

\pfl{财务分析} 2026 年 6 月末资产负债率 90.1\%，较 2025 年末下降 1.3 个百分点（样本中位数 52.2\%，所处三级行业内按由低到高排第\zhnumber{9}位）。 2026 年 6 月末货币资金 0.51 亿元、短期债务（短期借款与一年内到期非流动负债）2.94 亿元、长期债务 0.25 亿元，有息负债合计 3.19 亿元，现金短债比 0.17 倍——缺口明显。 净有息负债 2.67 亿元，相当于其 2026 年上半年营业收入的 108\%。 流动比率 0.37、速动比率 0.20，短期偿债指标偏紧。 2026 年上半年经营活动现金流净额 0.45 亿元，经营现金流与利润同向为正，内生资金可以支撑运营。 综合看，账面货币资金对有息负债与短期债务的覆盖不足，滚动偿付对新增融资的依赖度较高；资产负债率高于样本中位数 37.9 个百分点；有息负债中短期债务占比更高，期限结构仍以滚动续作为主。 公司披露的退市与持续经营风险为：2026 年 3 月，公司股票可能被终止上市的第三次风险提示公告；2025 年 4 月，公司股票自 该日起被实施退市风险警示、其他风险警示。

\pfl{重大战略转型} 从公开披露看，公司的战略推进集中在：并购与重组（5 项）：2026 年 5 月，公司股票撤销退市风险警示及其他风险警示；2025 年 9 月，公司深交所问询函回复；产能与项目（2 项）：2023 年 1 月，董事会审议通过以 9,500.00 万元受让江西纬科 100\% 股权。 公告口径上，近三年（2023 年 9 月---2026 年 9 月）公司披露重大重组与并购类公告 23 条、对外投资与转型类公告 5 条、回购与股权激励类公告 11 条，合计公告 645 条。 主营结构的口径变化已经落到报表上：2021 年年报的第一大行业为“兽药”（67.9\%），2026 年半年报为“动保系列”（91.4\%）；公司在这两期的分部划分方式有过调整。 综合判断，报告期内公司有 28 条与战略相关的公告落地（重大重组与并购 23 条、对外投资与转型 5 条），资本运作与主业调整之间存在直接关联。

\begin{center}\small\setlength{\tabcolsep}{4pt}
\captionof{table}[绿康生化（002868）历史融资与资本运作]{绿康生化（002868）历史融资与资本运作（据公司公告与公开报道整理；金额为披露值，含首次公开发行、再融资、债券、银行借款与重整投资等）}
\begin{adjustbox}{max width=\textwidth}
\begin{tabular}{@{}p{1.45cm}p{2.55cm}p{4.05cm}p{4.95cm}@{}}
\toprule
时间 & 方式 & 规模 & 用途与说明 \\
\midrule
2017-05-03 & IPO（深圳证券交易所） & 实际发行 3,\allowbreak{}000.00 万股 & 份有限公司，\allowbreak{}发行市盈率（按发行后总股本）20.50 倍 \\
2023-01-28 & 现金收购股权（跨界收购光伏胶膜企业） & 作价 9,\allowbreak{}500.00 万元受让玉山县旺宏企业管理中心（有限合伙）、\allowbreak{}王梅钧持有的江西纬科科技有限公司 100\% 股权（2022 年 8 月披露预案时预估值为 1.00 亿元、\allowbreak{}溢价率约 632\%） & 由兽药跨界切入光伏胶膜赛道；经第四届董事会第二十一次（临时）会议审议通过；标的后更名为绿康（玉山）胶膜材料有限公司；公司另投资成立绿康（海宁）胶膜材料有限公司，\allowbreak{}计划在海宁投资 60.00 亿元建设年产 8 亿平方米光伏胶膜项目 \\
2023-06-22 & 向特定对象发行 A 股股票（定增预案，未实施） & 发行对象为不超过 35 名特定投资者，\allowbreak{}发行价格为不低于定价基准日前 20 个交易日公司股票交易均价的 80\% & 报告期无募集资金使用情况）；公司于 2023 年 6 月 22 日在巨潮资讯网披露绿康生化 2023 年度向特定对象发行 A 股股票预案；媒体 2025 年报道称公司定增推进不顺 \\
2024-06-27 & 出售全资子公司股权（关联交易） & 以 2,\allowbreak{}730.00 万元的交易价格向关联法人福建浦潭热能有限公司出售全资子公司武汉绿康生化科技有限公司 100\% 股权 & 剥离资产、\allowbreak{}回笼资金；经公司第五届董事会第二十六次会议审议 \\
2024-07-22 & 以简易程序向特定对象发行股票（定增预案，未实施） & 拟募集资金总额不超过 8,\allowbreak{}000.00 万元（符合简易程序募集资金不超过 3.00 亿元且不超过最近一年末净资产 20\% 的规定） & 投向绿康（海宁）胶膜材料有限公司年产 3.2 亿平方米光伏胶膜项目（项目投资总额 7.54 亿元 \\
2025-04-24 & 控制权转让（股份协议转让 29.99\%） & 康怡投资、\allowbreak{}义睿投资、\allowbreak{}长鑫贰号、\allowbreak{}皓赢投资向福建纵腾网络有限公司转让公司 4,\allowbreak{}660.84 万股普通股（占公司总股本的 29.99\%） & 引入产业投资人、\allowbreak{}化解经营与退市风险；2025 年 11 月 21 日该等股份完成过户 \\
2025-11 & 控股股东借款 & 向控股股东福建纵腾网络有限公司借入 1.50 亿元 & 用于公司日常生产经营；借款期限自借款协议签署日起计算；同期公司完成 1.26 亿元银行借款的到期续贷 \\
2025-11-26 & 重大资产出售（0 元剥离光伏胶膜资产） & 以 0 元价格向江西饶信新能材料有限公司出售绿康（玉山）胶膜材料有限公司、\allowbreak{}绿康（海宁）胶膜材料有限公司、\allowbreak{}绿康新能（上海）进出口贸易有限公司共三家标的公司 100\% 股权；以 2024 年 12 月 31 日为评估基准日 & 退出光伏胶膜赛道、\allowbreak{}剥离亏损资产；本次交易构成关联交易；2025 年 9 月 23 日公司深交所问询函回复 \\
2025-12 & 土地收储（政府补偿） & 公司及子公司福建绿安生物农药有限公司与浦城县土地储备中心签订浦城县国有土地使用权收购合同，\allowbreak{}收储价款人民币 1.87 亿元；截至问询函回复日土地收储补偿第一期款项 3,\allowbreak{}000.00 万元已收讫 & 盘活土地资产、\allowbreak{}支持厂区搬迁与产能升级；收储价款将按协议分期结算 \\
\bottomrule
\end{tabular}
\end{adjustbox}
\end{center}

\pfl{历史融投资情况} 从现金流量表看，2025 年公司取得借款收到的现金 1.72 亿元、偿还债务支付的现金 5.07 亿元，筹资活动现金流净额 1.95 亿元（2023 年为取得借款 6.69 亿元、偿还债务 2.14 亿元）。 2026 年上半年筹资活动现金流净额为 -0.54 亿元（取得借款 1.73 亿元、偿还债务 2.04 亿元，未年化），已转为净流出，处于净还债状态。 2025 年分配股利、利润或偿付利息支付的现金合计 0.43 亿元。 公开披露的融资记录共 9 笔，工具类型以 IPO、定向增发、银行借款为主；金额与用途见下表。 其中规模最大的两笔为：2025 年 11 月，银行借款，向控股股东福建纵腾网络有限公司借入 1.50 亿元；2024 年 7 月，定向增发，拟募集资金总额不超过 8,000.00 万元（符合简易程序募集资金不超过 3.00 亿元且不超过最近一年末净资产 20\% 的规定）。 股本方面，近三年披露股本变动 3 次；总股本自 2021 年末以来未发生变化（1.55 亿股）。 分红方面，近三年未实施现金分红，在全部样本中属于分红能力偏弱的一端。 近三年“再融资”类公告 8 条，涉及定向增发、募集资金使用。 综合看，公司融资结构上股权与债务融资并举；2025 年筹资活动现金流净额 1.95 亿元，年度内仍为净融入，与前述经营与投资现金流的方向基本一致。

\pfl{未来潜在融资需求分析} 2026 年 6 月末货币资金 0.51 亿元，而短期债务 2.94 亿元（现金短债比 0.17 倍），2026 年上半年经营现金流净额 0.45 亿元。按“短期债务减去账面货币资金”的滚动偿付口径，静态缺口约 2.42 亿元；上半年盈利或接近盈亏平衡，该缺口不随经营亏损进一步扩大。 按第 \ref{sec:financing-need} 节的三档划分，公司属于第二档“保壳型”：近三年重大重组与并购类公告 23 条，2026 年上半年归母净利润 0.07 亿元，融资需求更多体现为“资产置换 + 维持上市地位”，而不是扩产。 公开披露的后续资金安排线索包括：股权融资线索：2024 年 7 月，公司 2024 年度以简易程序向特定对象发行股票预案；债务与授信安排：2025 年 11 月，同期到期银行借款续贷 1.26 亿元；2026 年 4 月，公司计划 2026 年向浦城农行申请流动资金 5,030.00 万元贷款续期；资本开支安排：2025 年 12 月，公司与浦城县土地储备中心签订浦城县国有土地使用权收购合同。 综合看，公司的外部资金需求以补充营运资金为主、项目投入为辅，滚动偿付口径下的静态缺口约 2.42 亿元，融资的时点与规模更多取决于信用条件与资本市场窗口。






\section{农产品加工：规模最大、利润最薄}
\label{sec:processing}

农产品加工板块 \num{23} 家公司（其他农产品加工 \num{11}、粮油加工 \num{7}、果蔬加工 \num{5}），
\hl{是公司数量最多的板块}，总市值合计约 \num{2887} 亿元（占全部农业上市公司的 \num{23.6}\%）。
按 2026 年半年报口径，\num{8} 家亏损，半年 ROE 中位数 \num{1.40}\%（未年化），
2026 年 6 月末资产负债率超过 \num{65}\% 的仅 \num{2} 家。

\begingroup\scriptsize\setlength{\tabcolsep}{2pt}
\begin{longtable}{@{}llrrrrrrr@{}}
\caption{农业上市公司分行业明细（粮油加工、果蔬加工、其他农产品加工；市值时点 2026-10-12，总市值口径；财务为 2026 年半年报，ROE 未年化，公告计数窗口 2023-09---2026-10）}\label{tab:comp-processing}\\
\toprule
代码 & 公司简称 & 总市值（亿元） & 营收（亿元） & 同比（\%） & 净利（亿元） & ROE（\%） & 再融资公告 & 重组公告 \\
\midrule
\endfirsthead
\multicolumn{9}{l}{\small（续）}\\
\toprule
代码 & 公司简称 & 总市值（亿元） & 营收（亿元） & 同比（\%） & 净利（亿元） & ROE（\%） & 再融资公告 & 重组公告 \\
\midrule
\endhead
\bottomrule
\endlastfoot
300999 & 金龙鱼 & 1,385.8 & 1,228.5 & 6.2 & 22.94 & 2.4 & 0 & 0 \\
600737 & 中粮糖业 & 342.0 & 92.8 & -21.2 & 6.08 & 5.2 & 1 & 0 \\
605198 & 安德利 & 213.3 & 7.2 & -24.4 & 1.23 & 4.3 & 17 & 3 \\
000930 & 中粮科技 & 96.5 & 81.9 & -7.0 & -1.69 & -1.6 & 0 & 0 \\
600127 & 金健米业 & 89.2 & 19.7 & 25.2 & -0.10 & -1.4 & 0 & 4 \\
600251 & 冠农股份 & 81.2 & 13.9 & 0.3 & 5.37 & 13.3 & 0 & 1 \\
000019 & 深粮控股 & 80.7 & 23.9 & 0.1 & 1.12 & 2.2 & 0 & 0 \\
600962 & 国投中鲁 & 62.5 & 8.8 & -12.4 & 0.02 & 0.2 & 85 & 32 \\
300138 & 晨光生物 & 57.1 & 31.6 & -13.6 & 1.79 & 5.2 & 0 & 0 \\
002481 & 双塔食品 & 54.4 & 11.2 & 7.3 & 0.33 & 1.4 & 1 & 2 \\
001338 & 永顺泰 & 46.6 & 19.1 & -13.9 & 0.51 & 1.4 & 1 & 0 \\
000505 & 京粮控股 & 45.7 & 31.1 & -26.0 & 0.08 & 0.3 & 0 & 2 \\
600191 & 华资实业 & 44.0 & 3.5 & 58.0 & 0.36 & 2.1 & 7 & 2 \\
002086 & 东方海洋 & 39.8 & 1.6 & 0.9 & -0.48 & -3.8 & 0 & 0 \\
003030 & 祖名股份 & 32.5 & 10.1 & 10.0 & -0.24 & -2.5 & 0 & 3 \\
603231 & 索宝蛋白 & 32.4 & 9.7 & 24.9 & 0.68 & 3.4 & 0 & 1 \\
000972 & 中基健康 & 31.6 & 2.0 & -19.6 & -0.66 & — & 36 & 16 \\
002286 & 保龄宝 & 31.1 & 15.0 & 6.9 & 0.82 & 3.6 & 9 & 0 \\
002852 & 道道全 & 29.2 & 28.6 & 2.6 & 0.17 & 0.8 & 3 & 0 \\
000911 & *ST广糖 & 24.9 & 15.8 & 17.4 & -0.16 & 0.0 & 78 & 19 \\
300175 & 朗源股份 & 24.6 & 1.7 & 67.7 & -0.13 & -2.7 & 0 & 2 \\
603182 & 嘉华股份 & 21.9 & 8.6 & 25.7 & 0.57 & 5.0 & 9 & 7 \\
300268 & 佳沃食品 & 19.7 & 4.2 & -66.1 & -0.04 & -1.4 & 0 & 20 \\
\end{longtable}
\endgroup


\subsection{规模与利润的严重错配}
\label{sec:processing-scale}

加工板块最直观的对照来自两家公司：\textbf{金龙鱼} 2026 年上半年营收 \num{1228.5} 亿元
（占 \num{104} 家农业上市公司上半年营收总额 \num{6085} 亿元的 \num{20.2}\%、
占本板块营收的 \num{73.5}\%）实现 \num{22.94} 亿元净利润，
\hl{净利率仅 \num{1.94}\%}；而 \textbf{安德利}以 \num{7.16} 亿元营收
实现 \num{1.23} 亿元净利润、\hl{净利率 \num{17.2}\%}，
安德利的总市值 \num{213.3} 亿元\emph{相当于金龙鱼（\num{1385.8} 亿元）的 \num{15}\%}，
而营收只有金龙鱼的 \num{0.6}\%——\hl{市场定价的是利润质量，而不是收入规模}。
\emph{这组对照揭示了加工行业的周期本质}：
以大宗农产品为原料、产品同质化的“量大利薄”模式（粮油加工）
没有定价权，利润被原料价格波动与渠道费用吞噬；
而以浓缩苹果汁、番茄制品等\hl{细分品类 + 出口市场}为核心的加工企业
（安德利半年 ROE \num{4.3}\%、冠农股份 \num{13.3}\%、晨光生物 \num{5.2}\%）
可以获得接近品牌消费品的利润率。

\subsection{食糖：当周期回到加工环节}
\label{sec:processing-sugar}

\textbf{中粮糖业}（总市值 \num{342.0} 亿元，板块第二，仅次于金龙鱼）2026 年上半年实现营收
\num{92.8} 亿元、净利润 \num{6.08} 亿元（半年 ROE \num{5.19}\%），
\hl{营收同比 \num{-21.2}\%，利润却同比 \num{+36.7}\%}（上年同期 \num{4.45} 亿元）。
把这家公司放回第 \ref{sec:soft} 章的框架中看，其业绩完全由白糖周期解释：
2024 年以来国内食糖产量恢复（广西、云南连续两年丰产）、
配额外进口成本下降，2026 年 8 月末广西白糖现货价格下滑至 \num{5420} 元/吨，
\hl{郑商所白糖主力合约 2021 年 2 月至 2026 年 9 月仅累计上涨 \num{0.8}\%}
——基本回到基期水平，明显弱于其他软商品（棉花 \num{3.5}\%、橡胶 \num{24.8}\%）。
对加工企业而言，原料价格下跌在库存周转周期内先体现为\hl{存货减值}，
因此糖价下跌的第一年（2025 年）加工企业的利润往往承压，
而第二年才因成本下降而受益。\emph{这是加工环节与种植环节的典型时滞差异}，
中粮糖业 2026 年上半年“收入降、利润增”的组合正是这一逻辑的兑现。

\subsection{困境公司的集中区}
\label{sec:processing-distress}

加工板块是农业上市公司中\hl{亏损集中度最高}的区域之一（2026 年上半年 \num{8} 家亏损）：
\textbf{中粮科技}（\num{-1.69} 亿元）、\textbf{中基健康}（\num{-0.66} 亿元）、
\textbf{东方海洋}（\num{-0.48} 亿元）、\textbf{祖名股份}（\num{-0.24} 亿元）
四家合计亏损 \num{3.1} 亿元，占板块全部亏损额（\num{3.5} 亿元）的 \num{88.6}\%。
从公告行为看，这批公司与高公告密度群体高度重合：
\textbf{*ST广糖}近三年再融资类公告 \num{78} 条（在全部 \num{104} 家公司中排第三），
2026 年 6 月末资产负债率 \num{100.2}\%；
\textbf{国投中鲁}（果蔬加工，2026 年上半年微利 \num{0.02} 亿元）
以 \num{85} 条再融资类公告居全样本第二，并伴随 \num{32} 条重组并购类公告（全样本第一）；
中基健康的资产负债率更高达 \num{104.1}\%——\hl{净资产已经为负}，
其半年 ROE 因此无法计算。
\hl{高公告密度 + 持续亏损 + ST 标识，构成一组清晰的“重组与保壳”信号}，
这将在第 \ref{sec:financing} 节中作为融资需求分类的依据。

\begin{keybox}[农产品加工的画像结论]
\normalsize
\textbf{① 规模不产生利润}：板块 \num{23} 家公司的上半年营业收入合计 \num{1671} 亿元、
占全样本的 \num{27.5}\%，而净利润合计只有 \num{38.6} 亿元，
其中金龙鱼一家占 \num{59.5}\%。\par\vspace{3pt}
\textbf{② 细分品类 + 出口是唯一的高毛利路径}：冠农股份（半年 ROE \num{13.3}\%）、
晨光生物（\num{5.2}\%）、安德利（\num{4.3}\%）均明显高于板块中位数 \num{1.40}\%，
而粮油加工环节的 ROE 中位数只有 \num{2.22}\%、其他农产品加工仅 \num{1.36}\%。\par\vspace{3pt}
\textbf{③ 食糖加工是“周期放大器”}：\hl{糖价下跌第一年存货减值、
第二年成本改善}，中粮糖业 2026 年上半年营收同比 \num{-21.2}\%、
利润同比 \num{+36.7}\% 的背离正源于此。\par\vspace{3pt}
\textbf{④ 利润逆势回升。}板块 \num{23} 家公司 2026 年上半年合计归母净利润
\num{38.58} 亿元（上年同期 \num{33.46} 亿元），\num{8} 家亏损；
\hl{低毛利、强刚需的加工环节在猪价下行期的相对优势再次显现}。
\end{keybox}

\subsection{上市公司画像（\num{23} 家）}
\label{sec:co-processing}

本节给出农产品加工行业全部 \num{23} 家上市公司的画像，按总市值降序排列。
数据口径与第 \ref{sec:comp-overview} 节一致：\hl{财务数据统一采用 2026 年半年报}
（合并报表口径、未年化，资产负债率为 2026 年 6 月末），
同时列出 2025 年半年报作为同口径对照、2023---2025 年年报作为趋势对照；
股价为月度收盘价，最后一个交易日为 2026 年 10 月 12 日；周期分段使用与第一部分相同的
ZigZag 算法，阈值取 25\%（个股波动大于商品价格，故高于生猪现货的 20\%）；
公告计数窗口为 2023 年 9 月 1 日至 2026 年 10 月 12 日。

与前一版不同，本版画像正文由\hl{逐家检索的公开资料底稿}驱动：底稿只收录有来源链接的事实
（公司沿革、股权与控制权、经营要点、所处行业的基本面、竞争优势、融资历史、
风险事件、战略动作、融资需求线索），
正文按事实生成，各家公司的段落结构与重点并不相同。
每家公司在统一顺序下展开：基本情况与主营结构、股权结构与控制权、经营要点、
所处行业的基本面（供需）、与同类公司比较的竞争优势、历史股价、过去周期分析、
盈利情况、财务分析、重大战略转型、历史融投资情况、未来潜在融投资需求；
\hl{市值一律采用总市值口径}，股权结构取最新披露报告期前三大股东与股东户数，
行业基本面与竞争优势以逐家检索的公开资料为依据，并与三级行业内的同口径指标横向比较。
每家公司配图两张：
股价与行业指数（申万农林牧渔）的对比图，以及近年主要财务指标的组合图
（上排：营业收入与归母净利润、ROE 与资产负债率；下方宽带面板：货币资金与有息负债为左轴柱状、
经营活动现金流净额为其右轴的蓝点折线），并给出关键财务指标表与历史融资/资本运作表。
未来融资需求按两个口径粗略度量：\hl{短期债务缺口 $=$ 期末短期债务 $-$ 货币资金}
与\hl{年化经营缺口 $=$ 半年归母亏损 $\times 2$}；
前者衡量滚动偿付压力，后者衡量维持当期盈利水平的现金消耗，
两者都是压力情形下的静态估算，不构成对融资规模或时点的预测。

% 农业上市公司画像：粮油加工、果蔬加工、其他农产品加工（共 23 家，按总市值降序）
% 由 scripts/make_company_profiles.py 自动生成，请勿手工修改

\subsubsection{金龙鱼（300999）}
\label{co:300999}

\pfl{公司基本情况} 益海嘉里旗下粮油食品巨头，主营厨房食品（食用油、大米、面粉）与饲料原料、油脂科技，旗下拥有金龙鱼、胡姬花、香满园、欧丽薇兰等品牌；2020 年登陆创业板并创下当时创业板 IPO 募资最大规模，2023—2024 年利润连续下滑后于 2025 年触底回暖。 公司于 2020-10-15 在深交所创业板首发上市，注册资本 54.22 亿元，总股本 54.22 亿股。 截至 2026 年 10 月 12 日收盘价 25.56 元，对应总市值 1,385.8 亿元，近一年-19.0\%，总市值在三级行业“粮油加工”的 7 家公司中按由大到小排第\zhnumber{1}位，近一年涨跌幅低于全样本中位数（-10.4\%）8.6 个百分点。 按 2025 年年报披露的行业构成，收入几乎全部来自农副食品加工业 100.0\%。

\pfl{股权结构与控制权} 按 2026 年 6 月 30 日的股东名册，前三大股东为丰益营销(香港)有限公司 持股 89.99\%；香港中央结算有限公司（持股比例未披露）；刘大鹏 持股 0.17\%，股东户数 110,227 户。 控制权与股东性质方面，2026 年 6 月末持股 48.79 亿股、占总股本 89.99\%，为境外法人。

\pfl{经营要点} 从经营要点看，厨房食品、饲料原料及油脂科技产品销量均同比增加。

\pfl{所处行业的基本面} 粮油加工的供给由大豆、菜籽等原料进口与压榨产能决定，需求由食用消费决定，利润来自原料与成品的价差而非价格方向（第 \ref{sec:framework} 章的供需框架）。 行业层面的最新情况是：2026 年，1—8 月中国累计进口大豆 7,410.7 万吨、同比增长 1.1\%，开机率维持 61\%—68\% 高位；2026 年 10 月，USDA 报告上调美豆单产至每英亩 53.1 蒲式耳、产量 45.6 亿蒲式耳、期末库存 3.15 亿蒲式耳。

\pfl{竞争优势} 把该公司放回同类公司比较，2026 年上半年营业收入 1,228.5 亿元、在三级行业“粮油加工”的 7 家公司中排第\zhnumber{1}位；毛利率 6.3\%、与行业内中位数基本持平；半年 ROE 2.4\%（未年化）、高于行业内中位数（2.22\%）；资产负债率 56.2\%、高于行业内中位数（43.5\%）。 相对优势方面，近三年公司在小包装食用植物油、包装米现代渠道市场份额均排名第一，包装米东北大米线上线下全渠道份额也居首位；2026 年上半年，公司来自境外的销售收入占比仅约 2.8\%；公司份额由 2019 年的 39.3\% 微降至 2024 年的 37.7\% 仍居第一。 优势在当期体现为高于行业中位数的资本回报，但优势能否持续仍取决于周期与自身产能投放节奏。

\pfl{历史股价} 本报告样本区间（2021 年 1 月 至 2026 年 10 月）内，股价由 121.92 元变为 25.56 元，累计 -79.0\%，区间最高 121.92 元（2021 年 1 月）、最低 23.26 元（2026 年 6 月）；当前价相当于区间最高价的 21.0\%，为区间最低价的 1.10 倍。 月度收益的年化波动率 32.2\%，低于全样本中位数 37.3\%，在三级行业“粮油加工”的 7 家公司中波动率由高到低排第\zhnumber{4}位。

\begin{center}
\begin{minipage}{\textwidth}\centering
\begin{tikzpicture}
\begin{axis}[
  width=12.6cm, height=3.9cm, scale only axis,
  xmin=2020.922, xmax=2026.972, ymin=21.63, ymax=130.45,
  xtick={2021,2022,2023,2024,2025,2026},
  yearticks,
  yticklabel style={font=\scriptsize},
  ylabel={元/股}, ylabel style={font=\scriptsize},
  grid=major, grid style={LightGray, line width=0.3pt},
  axis line style={InkGray, line width=0.5pt},
  every axis plot/.append style={line width=0.9pt},
]
\addplot[AgriGreen] table[x=x,y=y]{data/clean/profiles/px_300999.dat};
\addplot[SoftRed, dashed, line width=0.7pt] table[x=x,y=y]{data/clean/profiles/pxzz_300999.dat};
\addplot[only marks, mark=*, mark size=1.1pt, SoftRed] table[x=x,y=y]{data/clean/profiles/pxzz_300999.dat};
\end{axis}
\end{tikzpicture}

\captionof{figure}[金龙鱼（300999）股价与周期骨架]{金龙鱼（300999）月度收盘价（元/股）与 ZigZag 周期骨架（阈值 25\%，圆点为周期转折点）}
\end{minipage}
\end{center}

\pfl{过去周期分析} 以 25\% 的 ZigZag 阈值划分，样本区间内股价形成 1 个主升段与 2 个主跌段。 最大主升段出现在 2024 年 8 月 至 2024 年 9 月，1 个月+37.2\%；最大主跌段为 2021 年 1 月 至 2024 年 8 月，43 个月-78.8\%。 最近一次转折出现在 2026 年 6 月（下跌段自 2024 年 9 月起，21 个月-34.5\%），当前处于该段的延续之中。 公司股价与申万农林牧渔指数几乎同步见底，个股并未走出独立行情。 粮油加工是“成本加成”型生意（第 \ref{sec:grain} 章），利润来自原料与成品的价差及规模效率，收入规模大而净利率薄，公司 2026 年上半年实现盈利、半年 ROE 2.4\%（未年化），好于全样本中位数（0.75\%），下行周期对其报表的冲击小于同业。 综合区间形态看，该股单段涨跌幅度偏大、以趋势段为主（最大主升 +37.2\%、最大主跌 -78.8\%），年化波动率 32.2\% 与全样本中位数 37.3\% 相比波动小于样本中位数。

\pfl{盈利情况} 2026 年上半年营业收入 1,228.52 亿元（同比 +6.2\%），归母净利润 22.94 亿元，同比 +30.7\%。 以年报为趋势对照，2023---2025 年营业收入由 2,515.24 亿元变为 2,451.26 亿元（两年年均复合增速 -1.3\%），归母净利润由 28.48 亿元变为 31.53 亿元。 按半年报口径，2026 年上半年的 ROE 为 2.4\%（未年化）、销售净利率 1.9\%、毛利率 6.3\%，ROE 在“粮油加工”7 家公司中按数值由高到低排第\zhnumber{3}位。 按同一口径，三级行业“粮油加工”7 家公司的半年 ROE 中位数为 2.22\%，该公司高于其中位数 0.14 个百分点。 从公司披露的原因看，业绩变动主要来自：2025 年 4 月，公司 2024 年扣非净利润仅 9.72 亿元，同比下降 26.42\%；2025 年 11 月，判处罚金人民币 100.00 万元。 综合看，公司在报告期维持盈利（高于全样本半年 ROE 中位数 0.75\%），利润的可持续性取决于粮油加工景气的延续与成本管控的执行。

\begin{center}\small\setlength{\tabcolsep}{3.4pt}
\captionof{table}[金龙鱼（300999）关键财务指标]{金龙鱼（300999）关键财务指标（合并报表口径；两个半年报期均未年化；现金短债比 $=$ 货币资金 $\div$ 短期债务）}
\begin{adjustbox}{max width=\textwidth}
\begin{tabular}{@{}lrrrrr@{}}
\toprule
指标 & 2023 年 & 2024 年 & 2025 年 & 2025 年半年报 & 2026 年半年报 \\
\midrule
营业收入（亿元） & 2,515.24 & 2,382.79 & 2,451.26 & 1,156.82 & 1,228.52 \\
归母净利润（亿元） & 28.48 & 25.02 & 31.53 & 17.56 & 22.94 \\
毛利率（\%） & 4.8 & 5.4 & 6.4 & 6.5 & 6.3 \\
ROE（\%） & 3.2 & 2.7 & 3.3 & 1.9 & 2.4 \\
资产负债率（\%） & 59.8 & 56.4 & 55.9 & 55.9 & 56.2 \\
经营活动现金流净额（亿元） & 149.21 & 60.38 & 154.45 & 146.25 & 10.01 \\
货币资金（亿元） & 592.96 & 529.06 & 442.04 & 443.72 & 485.84 \\
有息负债合计（亿元） & 1,157.03 & 943.10 & 864.73 & 887.24 & 927.22 \\
现金短债比（倍） & 0.56 & 0.60 & 0.54 & 0.53 & 0.55 \\
\bottomrule
\end{tabular}
\end{adjustbox}
\end{center}

\begin{center}
\begin{minipage}{\textwidth}\centering
\begin{tikzpicture}
\begin{groupplot}[
  group style={group size=2 by 1, horizontal sep=0.60cm},
  width=6.85cm, height=4.60cm,
  tick label style={font=\tiny},
  title style={font=\scriptsize\heiti},
  yticklabel style={font=\tiny},
  ylabel style={font=\tiny},
  grid=major, grid style={LightGray, line width=0.3pt},
  axis line style={InkGray, line width=0.5pt},
  enlarge x limits={abs=0.8},
  legend style={font=\scriptsize, at={(0.5,-0.20)}, anchor=north,
                legend columns=2, draw=none, fill=none, column sep=6pt},
]
\nextgroupplot[
  ybar, bar width=4pt,
  ymin=-257.49, ymax=2883.84,
  xtick={0,1,2,3,4,5.7},
  xticklabels={2021,2022,2023,2024,2025,2026H1},
  ylabel={亿元},
]
\addplot[fill=AgriGreen!75, draw=AgriGreen, bar shift=-2.6pt] table[x=i, y=rev]{data/clean/profiles/fin_300999.dat};
\addplot[fill=SoftRed!60, draw=SoftRed, bar shift=2.6pt] table[x=i, y=ni]{data/clean/profiles/fin_300999.dat};
\addplot[fill=AgriGreen!30, draw=AgriGreen, pattern=north east lines, bar shift=-2.6pt] table[x=x, y=rev]{data/clean/profiles/finh_300999.dat};
\addplot[fill=SoftRed!20, draw=SoftRed, pattern=north east lines, bar shift=2.6pt] table[x=x, y=ni]{data/clean/profiles/finh_300999.dat};
\legend{营业收入（年/半年）, 归母净利润（年/半年）}
\nextgroupplot[
  ymin=-4.8, ymax=65.7,
  xtick={0,1,2,3,4,5.7},
  xticklabels={2021,2022,2023,2024,2025,2026H1},
  ylabel={\%},
]
\addplot[AgriGold, mark=*, mark size=1.3pt] table[x=i, y=roe]{data/clean/profiles/fin_300999.dat};
\addplot[OilBlue, mark=square*, mark size=1.3pt] table[x=i, y=debt]{data/clean/profiles/fin_300999.dat};
\addplot[only marks, AgriGold, mark=*, mark size=2.0pt] table[x=x, y=roe]{data/clean/profiles/finh_300999.dat};
\addplot[only marks, OilBlue, mark=square*, mark size=2.0pt] table[x=x, y=debt]{data/clean/profiles/finh_300999.dat};
\legend{ROE, 资产负债率}
\end{groupplot}
\end{tikzpicture}
\begin{tikzpicture}
\begin{axis}[
  name=finbot,
  width=13.75cm, height=3.10cm, scale only axis,
  ymin=-115.70, ymax=1295.87,
  xtick={0,1,2,3,4,5.7},
  xticklabels={2021,2022,2023,2024,2025,2026H1},
  tick label style={font=\tiny},
  yticklabel style={font=\tiny},
  ylabel={亿元}, ylabel style={font=\tiny},
  ybar, bar width=4pt,
  grid=major, grid style={LightGray, line width=0.3pt},
  axis line style={InkGray, line width=0.5pt},
  enlarge x limits={abs=0.8},
  legend style={font=\scriptsize, at={(0.5,-0.26)}, anchor=north,
                legend columns=2, draw=none, fill=none, column sep=10pt},
]
\addplot[fill=AgriGreen!55, draw=AgriGreen, bar shift=-3.0pt] table[x=i, y=cash]{data/clean/profiles/fin_300999.dat};
\addplot[fill=SoftRed!45, draw=SoftRed, bar shift=3.0pt] table[x=i, y=ibd]{data/clean/profiles/fin_300999.dat};
\addplot[fill=AgriGreen!25, draw=AgriGreen, pattern=north east lines, bar shift=-3.0pt] table[x=x, y=cash]{data/clean/profiles/finh_300999.dat};
\addplot[fill=SoftRed!20, draw=SoftRed, pattern=north east lines, bar shift=3.0pt] table[x=x, y=ibd]{data/clean/profiles/finh_300999.dat};
\legend{货币资金（左轴）, 有息负债（左轴）}
\end{axis}
\begin{axis}[
  at={(finbot.south east)}, anchor=south east,
  width=13.75cm, height=3.10cm, scale only axis,
  axis y line*=right, axis x line*=none, xtick=\empty,
  ymin=-40.16, ymax=200.78,
  tick label style={font=\tiny},
  yticklabel style={font=\tiny}, scaled y ticks=false,
  legend style={font=\scriptsize, at={(0.5,-0.44)}, anchor=north,
                legend columns=1, draw=none, fill=none},
]
\addplot[OilBlue, mark=*, mark size=2.2pt, line width=1.3pt] table[x=i, y=ocf]{data/clean/profiles/fin_300999.dat};
\addplot[only marks, OilBlue, mark=*, mark size=3.2pt] table[x=x, y=ocf]{data/clean/profiles/finh_300999.dat};
\legend{经营活动现金流净额（右轴，亿元，蓝点）}
\end{axis}
\end{tikzpicture}

\captionof{figure}[金龙鱼（300999）近年主要财务指标]{金龙鱼（300999）近年主要财务指标（左上：营业收入与归母净利润；右上：ROE 与资产负债率；下：货币资金、有息负债为左轴柱状，经营活动现金流净额为其右轴的蓝点折线。
2021---2025 年为年报口径，2026H1 为 2026 年半年报/半年末，与其前一年报之间留出间隔、以斜纹或空心点区分；半年报数据未年化）}
\end{minipage}
\end{center}

\pfl{财务分析} 2026 年 6 月末资产负债率 56.2\%，较 2025 年末上升 0.4 个百分点（样本中位数 52.2\%，所处三级行业内按由低到高排第\zhnumber{6}位）。 2026 年 6 月末货币资金 485.84 亿元、短期债务（短期借款与一年内到期非流动负债）887.12 亿元、长期债务 40.10 亿元，有息负债合计 927.22 亿元，现金短债比 0.55 倍——覆盖偏紧。 净有息负债 441.38 亿元，相当于其 2026 年上半年营业收入的 36\%。 流动比率 1.20、速动比率 0.78，短期流动性无明显缺口。 2026 年上半年经营活动现金流净额 10.01 亿元，经营现金流与利润同向为正，内生资金可以支撑运营。 2026 年 6 月末在建工程 42.03 亿元，较上年末增加 0.83 亿元，仍有在建投入。 2026 年 6 月末商誉 61.11 亿元，占归母净资产的 6.0\%，是净资产中最脆弱的一块。 综合看，现金短债比 0.55 倍，短期资金安排偏紧；资产负债率高于样本中位数 4.0 个百分点；有息负债中短期债务占比更高，期限结构仍以滚动续作为主。 公司披露的监管与合规风险为：2026 年 2 月，两起子公司诉讼合计计提预计负债损失约 7.33 亿元，均为一审未生效判决、相关子公司已提起上诉。

\pfl{重大战略转型} 公司的战略动作可归为以下几类：其他动作（1 项）：2024 年 9 月，公司以参股公司股权作价增资鲁花集团。 公告口径上，近三年（2023 年 9 月---2026 年 9 月）公司披露重大重组与并购类公告 0 条、对外投资与转型类公告 3 条、回购与股权激励类公告 14 条，合计公告 306 条。 第一大行业仍是“农副食品加工业”，收入占比 100.0\%，与 2021 年年报相比没有方向性变化。 综合判断，报告期内公司有 3 条与战略相关的公告落地（重大重组与并购 0 条、对外投资与转型 3 条），资本运作与主业调整之间存在直接关联。

\pfl{历史融投资情况} 从现金流量表看，2025 年公司取得借款收到的现金 1,569.21 亿元、偿还债务支付的现金 1,649.00 亿元、吸收权益性投资 0.70 亿元，筹资活动现金流净额 -46.92 亿元（2023 年为取得借款 1,395.70 亿元、偿还债务 1,225.17 亿元）。 2026 年上半年筹资活动现金流净额为 52.24 亿元（取得借款 870.26 亿元、偿还债务 806.36 亿元，未年化），仍处于净融入状态。 2025 年分配股利、利润或偿付利息支付的现金合计 25.46 亿元。 公开披露的融资记录共 2 笔，工具类型以 IPO 为主；金额与用途见下表。 其中规模最大的两笔为：2020 年 10 月，IPO，首次公开发行 5.42 亿股，发行价 25.70 元/股；2024 年 9 月，以参股公司股权作价对，公司出资作价 22.90 亿元、关联方香港嘉银作价 32.79 亿元。 股本方面，近三年披露股本变动 3 次；总股本自 2021 年末以来未发生变化（54.22 亿股）。 分红方面，近三年现金分红 4 次、合计每股派现 0.53 元，最近一次实施在 2025 年 12 月。 综合看，公司融资结构上以债务滚动为主、股权融资为补充；2025 年筹资活动现金流净额 -46.92 亿元，融资节奏仍以维持存量债务周转为主，与前述经营与投资现金流的方向基本一致。

\pfl{未来潜在融资需求分析} 2026 年 6 月末货币资金 485.84 亿元，而短期债务 887.12 亿元（现金短债比 0.55 倍），2026 年上半年经营现金流净额 10.01 亿元。按“短期债务减去账面货币资金”的滚动偿付口径，静态缺口约 401.28 亿元；上半年盈利或接近盈亏平衡，该缺口不随经营亏损进一步扩大。 公司 2026 年 6 月末资产负债率 56.2\%、2026 年上半年 ROE 2.4\%（未年化），现金短债比 0.55 倍、短债覆盖不足，融资需求首先用于置换与滚动存量债务、补充营运资金，属于“补血型”而非扩张型。 公开披露的后续资金安排线索包括：债务与授信安排：2024 年报，公司披露报告期投资额 821.63 亿元，同比增长 21.61\%；资本开支安排：2020---2024 年，IPO 募集资金净额 136.93 亿元。 资金需求还要叠加在建工程：期末在建工程 42.03 亿元、2026 年上半年购建固定资产等支出 20.91 亿元（未年化），这部分支出在周期底部通常会被主动放缓。 综合看，公司的外部资金需求以营运资金周转与在建项目投入为主，滚动偿付口径下的静态缺口约 401.28 亿元，能否落地取决于授信续作与发行窗口的配合。




\subsubsection{中粮糖业（600737）}
\label{co:600737}

\pfl{公司基本情况} 中粮糖业成立于 1993 年，中粮集团旗下食糖与番茄双主业平台，国内食糖行业龙头，拥有从国内外制糖、进口炼糖到贸易、仓储物流的全产业链；番茄制品规模居亚洲第一、世界第二。经历 2004 年德隆系危机后由中粮重组脱困，近年受糖价与番茄酱价格下行影响业绩承压。 公司于 1996-07-31 在上交所首发上市，注册资本 21.39 亿元，总股本 21.39 亿股。 截至 2026 年 10 月 12 日收盘价 15.99 元，对应总市值 342.0 亿元，近一年+0.4\%，总市值在三级行业“其他农产品加工”的 11 家公司中按由大到小排第\zhnumber{1}位，近一年涨跌幅高于全样本中位数（-10.4\%）10.7 个百分点。 按 2025 年年报披露的行业构成，收入占比前三位为工业 66.3\%；贸易 32.9\%；其他(补充) 0.7\%。 与 2021 年年报相比，第一大行业口径由“贸易”（65.9\%）变为“工业”（66.3\%）；两期披露的分类口径有调整，占比差异中同时包含分类因素。 股权结构方面，控股股东为中粮集团有限公司。

\pfl{股权结构与控制权} 按 2026 年 6 月 30 日的股东名册，前三大股东为中粮集团有限公司 持股 50.73\%；中国建设银行股份有限公司－华泰柏瑞中证红利低波动交易（持股比例未披露）；施燕 持股 1.39\%，股东户数 98,457 户。 控制权与股东性质方面，截至 2026 年 6 月 30 日持股 10.85 亿股、占公司总股本 50.73\%；2025 年年报，报告期内公司控制权未发生变更。

\pfl{经营要点} 经营与产能方面，公司炼糖（加工糖）产能已超过 200 万吨/年、稳居国内行业首位；2025 年度，公司自产糖产量 127.56 万吨、同比增长 2.35\%，自产糖销售量约 126.49 万吨、同比增长 6.26\%。

\pfl{所处行业的基本面} 该子行业横跨制糖、番茄制品、糖醇与添加剂，供给受种植面积与进口配额约束，需求端分食用与工业两类（第 \ref{sec:framework} 章的供需框架）。 与该公司直接相关的行业数据包括：2025 年 12 月，国际糖业组织预计 2025/26 年度全球食糖产量增长 3.2\% 至 1.818 亿吨；进口依存度维持在 35\% 左右。

\pfl{竞争优势} 把该公司放回同类公司比较，2026 年上半年营业收入 92.8 亿元、在三级行业“其他农产品加工”的 11 家公司中排第\zhnumber{1}位；毛利率 14.6\%、高于行业内中位数 5.5 个百分点；半年 ROE 5.2\%（未年化）、高于行业内中位数（1.36\%）；资产负债率 39.1\%、高于行业内中位数（39.1\%）。 公司在公开资料中自述与同业比较的相对优势在于：炼糖产能超 200 万吨/年居国内首位；2026 年上半年，在糖价进入下行周期、销量与营收同比下降的背景下，第二季度单季毛利率升至 17.72\%、净利率升至 7.89\%；2026 年上半年新增申请发明专利 7 件、新授权专利 18 件，并正式发布 3 项团体标准。 优势在报表上主要体现为高于行业平均的毛利率，这在周期底部是抵御亏损的主要缓冲。

\pfl{历史股价} 以 2021 年 1 月为起点、2026 年 10 月为终点，股价由 9.19 元变为 15.99 元，累计 +74.0\%，区间最高 17.74 元（2026 年 1 月）、最低 6.62 元（2022 年 10 月）；当前价相当于区间最高价的 90.1\%，为区间最低价的 2.42 倍。 月度收益的年化波动率 33.4\%，低于全样本中位数 37.3\%，在三级行业“其他农产品加工”的 11 家公司中波动率由高到低排第\zhnumber{8}位。

\begin{center}
\begin{minipage}{\textwidth}\centering
\begin{tikzpicture}
\begin{axis}[
  width=12.6cm, height=3.9cm, scale only axis,
  xmin=2020.922, xmax=2026.972, ymin=6.16, ymax=18.98,
  xtick={2021,2022,2023,2024,2025,2026},
  yearticks,
  yticklabel style={font=\scriptsize},
  ylabel={元/股}, ylabel style={font=\scriptsize},
  grid=major, grid style={LightGray, line width=0.3pt},
  axis line style={InkGray, line width=0.5pt},
  every axis plot/.append style={line width=0.9pt},
]
\addplot[AgriGreen] table[x=x,y=y]{data/clean/profiles/px_600737.dat};
\addplot[SoftRed, dashed, line width=0.7pt] table[x=x,y=y]{data/clean/profiles/pxzz_600737.dat};
\addplot[only marks, mark=*, mark size=1.1pt, SoftRed] table[x=x,y=y]{data/clean/profiles/pxzz_600737.dat};
\end{axis}
\end{tikzpicture}

\captionof{figure}[中粮糖业（600737）股价与周期骨架]{中粮糖业（600737）月度收盘价（元/股）与 ZigZag 周期骨架（阈值 25\%，圆点为周期转折点）}
\end{minipage}
\end{center}

\pfl{过去周期分析} 以 25\% 的 ZigZag 阈值划分，样本区间内股价形成 3 个主升段与 2 个主跌段。 最大主升段出现在 2022 年 10 月 至 2026 年 1 月，39 个月+168.0\%；最大主跌段为 2021 年 6 月 至 2022 年 10 月，16 个月-33.9\%。 最近一次转折出现在 2026 年 10 月（上涨段自 2026 年 6 月起，4 个月+29.7\%），当前处于该段之后的震荡之中。 公司股价较申万农林牧渔指数提前约 43 个月见底，是板块中较早定价周期反转的公司之一。 该子行业横跨糖、番茄制品、添加剂等多个品种，各自的供给周期与政策（进口配额、关税）是主要变量，公司 2026 年上半年实现盈利、半年 ROE 5.2\%（未年化），好于全样本中位数（0.75\%），下行周期对其报表的冲击小于同业。 综合区间形态看，该股单段涨跌幅度偏大、以趋势段为主（最大主升 +168.0\%、最大主跌 -33.9\%），年化波动率 33.4\% 与全样本中位数 37.3\% 相比波动小于样本中位数。

\pfl{盈利情况} 2026 年上半年营业收入 92.77 亿元（同比 -21.2\%），归母净利润 6.08 亿元，同比 +36.7\%。 以年报为趋势对照，2023---2025 年营业收入由 331.14 亿元变为 265.11 亿元（两年年均复合增速 -10.5\%），归母净利润由 20.73 亿元变为 9.51 亿元。 按半年报口径，2026 年上半年的 ROE 为 5.2\%（未年化）、销售净利率 6.7\%、毛利率 14.6\%，ROE 在“其他农产品加工”11 家公司中按数值由高到低排第\zhnumber{1}位。 按同一口径，三级行业“其他农产品加工”11 家公司的半年 ROE 中位数为 1.36\%，该公司高于其中位数 3.83 个百分点。 综合看，公司在报告期维持盈利（高于全样本半年 ROE 中位数 0.75\%），利润的可持续性取决于其他农产品加工景气的延续与成本管控的执行。

\begin{center}\small\setlength{\tabcolsep}{3.4pt}
\captionof{table}[中粮糖业（600737）关键财务指标]{中粮糖业（600737）关键财务指标（合并报表口径；两个半年报期均未年化；现金短债比 $=$ 货币资金 $\div$ 短期债务）}
\begin{adjustbox}{max width=\textwidth}
\begin{tabular}{@{}lrrrrr@{}}
\toprule
指标 & 2023 年 & 2024 年 & 2025 年 & 2025 年半年报 & 2026 年半年报 \\
\midrule
营业收入（亿元） & 331.14 & 324.97 & 265.11 & 117.67 & 92.77 \\
归母净利润（亿元） & 20.73 & 17.13 & 9.51 & 4.45 & 6.08 \\
毛利率（\%） & 14.1 & 12.5 & 11.4 & 8.5 & 14.6 \\
ROE（\%） & 19.4 & 14.4 & 8.2 & 3.9 & 5.2 \\
资产负债率（\%） & 44.9 & 41.5 & 44.1 & 48.0 & 39.1 \\
经营活动现金流净额（亿元） & 8.87 & 26.34 & 12.88 & 10.28 & 11.98 \\
货币资金（亿元） & 10.06 & 8.37 & 14.65 & 21.72 & 8.58 \\
有息负债合计（亿元） & 37.58 & 35.66 & 46.37 & 42.49 & 30.71 \\
现金短债比（倍） & 0.27 & 0.24 & 0.32 & 0.52 & 0.28 \\
\bottomrule
\end{tabular}
\end{adjustbox}
\end{center}

\begin{center}
\begin{minipage}{\textwidth}\centering
\begin{tikzpicture}
\begin{groupplot}[
  group style={group size=2 by 1, horizontal sep=0.60cm},
  width=6.85cm, height=4.60cm,
  tick label style={font=\tiny},
  title style={font=\scriptsize\heiti},
  yticklabel style={font=\tiny},
  ylabel style={font=\tiny},
  grid=major, grid style={LightGray, line width=0.3pt},
  axis line style={InkGray, line width=0.5pt},
  enlarge x limits={abs=0.8},
  legend style={font=\scriptsize, at={(0.5,-0.20)}, anchor=north,
                legend columns=2, draw=none, fill=none, column sep=6pt},
]
\nextgroupplot[
  ybar, bar width=4pt,
  ymin=-33.11, ymax=370.88,
  xtick={0,1,2,3,4,5.7},
  xticklabels={2021,2022,2023,2024,2025,2026H1},
  ylabel={亿元},
]
\addplot[fill=AgriGreen!75, draw=AgriGreen, bar shift=-2.6pt] table[x=i, y=rev]{data/clean/profiles/fin_600737.dat};
\addplot[fill=SoftRed!60, draw=SoftRed, bar shift=2.6pt] table[x=i, y=ni]{data/clean/profiles/fin_600737.dat};
\addplot[fill=AgriGreen!30, draw=AgriGreen, pattern=north east lines, bar shift=-2.6pt] table[x=x, y=rev]{data/clean/profiles/finh_600737.dat};
\addplot[fill=SoftRed!20, draw=SoftRed, pattern=north east lines, bar shift=2.6pt] table[x=x, y=ni]{data/clean/profiles/finh_600737.dat};
\legend{营业收入（年/半年）, 归母净利润（年/半年）}
\nextgroupplot[
  ymin=-3.7, ymax=50.4,
  xtick={0,1,2,3,4,5.7},
  xticklabels={2021,2022,2023,2024,2025,2026H1},
  ylabel={\%},
]
\addplot[AgriGold, mark=*, mark size=1.3pt] table[x=i, y=roe]{data/clean/profiles/fin_600737.dat};
\addplot[OilBlue, mark=square*, mark size=1.3pt] table[x=i, y=debt]{data/clean/profiles/fin_600737.dat};
\addplot[only marks, AgriGold, mark=*, mark size=2.0pt] table[x=x, y=roe]{data/clean/profiles/finh_600737.dat};
\addplot[only marks, OilBlue, mark=square*, mark size=2.0pt] table[x=x, y=debt]{data/clean/profiles/finh_600737.dat};
\legend{ROE, 资产负债率}
\end{groupplot}
\end{tikzpicture}
\begin{tikzpicture}
\begin{axis}[
  name=finbot,
  width=13.75cm, height=3.10cm, scale only axis,
  ymin=-4.64, ymax=51.93,
  xtick={0,1,2,3,4,5.7},
  xticklabels={2021,2022,2023,2024,2025,2026H1},
  tick label style={font=\tiny},
  yticklabel style={font=\tiny},
  ylabel={亿元}, ylabel style={font=\tiny},
  ybar, bar width=4pt,
  grid=major, grid style={LightGray, line width=0.3pt},
  axis line style={InkGray, line width=0.5pt},
  enlarge x limits={abs=0.8},
  legend style={font=\scriptsize, at={(0.5,-0.26)}, anchor=north,
                legend columns=2, draw=none, fill=none, column sep=10pt},
]
\addplot[fill=AgriGreen!55, draw=AgriGreen, bar shift=-3.0pt] table[x=i, y=cash]{data/clean/profiles/fin_600737.dat};
\addplot[fill=SoftRed!45, draw=SoftRed, bar shift=3.0pt] table[x=i, y=ibd]{data/clean/profiles/fin_600737.dat};
\addplot[fill=AgriGreen!25, draw=AgriGreen, pattern=north east lines, bar shift=-3.0pt] table[x=x, y=cash]{data/clean/profiles/finh_600737.dat};
\addplot[fill=SoftRed!20, draw=SoftRed, pattern=north east lines, bar shift=3.0pt] table[x=x, y=ibd]{data/clean/profiles/finh_600737.dat};
\legend{货币资金（左轴）, 有息负债（左轴）}
\end{axis}
\begin{axis}[
  at={(finbot.south east)}, anchor=south east,
  width=13.75cm, height=3.10cm, scale only axis,
  axis y line*=right, axis x line*=none, xtick=\empty,
  ymin=-7.69, ymax=34.45,
  tick label style={font=\tiny},
  yticklabel style={font=\tiny}, scaled y ticks=false,
  legend style={font=\scriptsize, at={(0.5,-0.44)}, anchor=north,
                legend columns=1, draw=none, fill=none},
]
\addplot[OilBlue, mark=*, mark size=2.2pt, line width=1.3pt] table[x=i, y=ocf]{data/clean/profiles/fin_600737.dat};
\addplot[only marks, OilBlue, mark=*, mark size=3.2pt] table[x=x, y=ocf]{data/clean/profiles/finh_600737.dat};
\legend{经营活动现金流净额（右轴，亿元，蓝点）}
\end{axis}
\end{tikzpicture}

\captionof{figure}[中粮糖业（600737）近年主要财务指标]{中粮糖业（600737）近年主要财务指标（左上：营业收入与归母净利润；右上：ROE 与资产负债率；下：货币资金、有息负债为左轴柱状，经营活动现金流净额为其右轴的蓝点折线。
2021---2025 年为年报口径，2026H1 为 2026 年半年报/半年末，与其前一年报之间留出间隔、以斜纹或空心点区分；半年报数据未年化）}
\end{minipage}
\end{center}

\pfl{财务分析} 2026 年 6 月末资产负债率 39.1\%，较 2025 年末下降 5.1 个百分点（样本中位数 52.2\%，所处三级行业内按由低到高排第\zhnumber{6}位）。 2026 年 6 月末货币资金 8.58 亿元、短期债务（短期借款与一年内到期非流动负债）30.28 亿元、长期债务 0.43 亿元，有息负债合计 30.71 亿元，现金短债比 0.28 倍——缺口明显。 净有息负债 22.13 亿元，相当于其 2026 年上半年营业收入的 24\%。 流动比率 1.82、速动比率 0.72，短期指标尚可。 2026 年上半年经营活动现金流净额 11.98 亿元，经营现金流与利润同向为正，内生资金可以支撑运营。 2026 年 6 月末在建工程 1.38 亿元，较上年末增加 0.31 亿元，仍有在建投入。 综合看，账面货币资金对有息负债与短期债务的覆盖不足，滚动偿付对新增融资的依赖度较高；资产负债率低于样本中位数 13.1 个百分点；有息负债中短期债务占比更高，期限结构仍以滚动续作为主。

\pfl{重大战略转型} 公司的战略动作可归为以下几类：并购与重组（1 项）：2005---2007 年，中粮集团收购重组后积极布局糖产业。 公告口径上，近三年（2023 年 9 月---2026 年 9 月）公司披露重大重组与并购类公告 0 条、对外投资与转型类公告 0 条、回购与股权激励类公告 0 条，合计公告 285 条。 主营结构的口径变化已经落到报表上：2021 年年报的第一大行业为“贸易”（65.9\%），2025 年年报为“工业”（66.3\%）；公司在这两期的分部划分方式有过调整。 综合判断，公司报告期内的变化主要发生在收入结构层面，属于口径与业务重心的再平衡，而非外延式扩张。

\begin{center}\small\setlength{\tabcolsep}{4pt}
\captionof{table}[中粮糖业（600737）历史融资与资本运作]{中粮糖业（600737）历史融资与资本运作（据公司公告与公开报道整理；金额为披露值，含首次公开发行、再融资、债券、银行借款与重整投资等）}
\begin{adjustbox}{max width=\textwidth}
\begin{tabular}{@{}p{1.45cm}p{2.55cm}p{4.05cm}p{4.95cm}@{}}
\toprule
时间 & 方式 & 规模 & 用途与说明 \\
\midrule
1996-07-31 & IPO（上海证券交易所，前身新疆屯河） & 实际募集资金合计 1.01 亿元 & 招股公告日 1996-06-29，\allowbreak{}发行价 5.78 元/\allowbreak{}股 \\
2007—2008 & 非公开发行（募集资金全部用于偿还对中粮集团的借款） & --- & 偿还对中粮集团的借款；经公司 2007 年第二次临时股东大会审议，\allowbreak{}中粮财务公司将持有的本公司债权 8.80 亿元转让给中粮集团 \\
2019-04 & 非公开发行 A 股（上交所） & 发行 8,\allowbreak{}697.21 万股，\allowbreak{}发行价 7.52 元/\allowbreak{}股 & 甘蔗糖技术升级改造及配套优质高产高糖糖料蔗基地建设项目、\allowbreak{}甜菜糖项目；201 核准非公开发行不超过 4.10 亿股；2019 年 3 月 25 日包括中粮集团在内的 6 名投资者认购全部股份 \\
2022-12 & 子公司收购关联方资产 & 6,\allowbreak{}843.87 万元 & 全资子公司中粮屯河番茄有限公司收购中粮集团持有的内蒙古中粮番茄制品有限公司 100\% 股权；中粮集团为公司控股股东（持股 50.73\%） \\
\bottomrule
\end{tabular}
\end{adjustbox}
\end{center}

\pfl{历史融投资情况} 从现金流量表看，2025 年公司取得借款收到的现金 104.59 亿元、偿还债务支付的现金 95.16 亿元，筹资活动现金流净额 -0.55 亿元（2023 年为取得借款 123.93 亿元、偿还债务 125.08 亿元）。 2026 年上半年筹资活动现金流净额为 -15.84 亿元（取得借款 24.64 亿元、偿还债务 39.79 亿元，未年化），已转为净流出，处于净还债状态。 2025 年分配股利、利润或偿付利息支付的现金合计 9.59 亿元。 公开披露的融资记录共 4 笔，工具类型以 IPO、定向增发为主；金额与用途见下表。 其中规模最大的两笔为：1996 年 7 月，IPO，实际募集资金合计 1.01 亿元；2022 年 12 月，子公司收购关联方资产，6,843.87 万元。 股本方面，近三年披露股本变动 3 次；总股本自 2021 年末以来未发生变化（21.39 亿股）。 分红方面，近三年现金分红 6 次、合计每股派现 2.08 元，最近一次实施在 2025 年 12 月。 近三年“再融资”类公告 1 条，涉及定向增发、募集资金使用。 综合看，公司融资结构上股权与债务融资并举；2025 年筹资活动现金流净额 -0.55 亿元，融资节奏仍以维持存量债务周转为主，与前述经营与投资现金流的方向基本一致。

\pfl{未来潜在融资需求分析} 2026 年 6 月末货币资金 8.58 亿元，而短期债务 30.28 亿元（现金短债比 0.28 倍），2026 年上半年经营现金流净额 11.98 亿元。按“短期债务减去账面货币资金”的滚动偿付口径，静态缺口约 21.70 亿元；上半年盈利或接近盈亏平衡，该缺口不随经营亏损进一步扩大。 按第 \ref{sec:financing-need} 节的三档划分，公司属于第三档“扩张型”：2026 年上半年 ROE 5.2\%（未年化，折合约 10.4\%）、6 月末资产负债率 39.1\%，资产负债表仍具备加杠杆空间；融资用途以产能扩张、品类并购与海外布局为主，单笔规模通常在 5---20 亿元量级，会被市场定价为成长而非补血。 公开披露的后续资金安排线索包括：股权融资线索：2019 年，非公开发行募集资金截至 2023-12-31 剩余 1,078.73 万元。 资金需求还要叠加在建工程：期末在建工程 1.38 亿元、2026 年上半年购建固定资产等支出 2.04 亿元（未年化），这部分支出在周期底部通常会被主动放缓。 综合看，公司的外部资金需求以补充营运资金为主、项目投入为辅，滚动偿付口径下的静态缺口约 21.70 亿元，融资的时点与规模更多取决于信用条件与资本市场窗口。




\subsubsection{安德利（605198）}
\label{co:605198}

\pfl{公司基本情况} 安德利成立于 1996 年，烟台北方安德利果汁股份有限公司，山东安德利集团旗下企业，全球主要浓缩果汁加工企业之一，主营浓缩苹果汁、浓缩梨汁等多品种浓缩果汁；2003 年 H 股在港上市、2020 年回归 A 股，为国内首家饮料类“A+H”双上市企业，2026 年拟跨界收购电子材料资产。 公司于 2020-09-18 在上交所首发上市，注册资本 3.28 亿元，总股本 3.49 亿股。 截至 2026 年 10 月 12 日收盘价 61.12 元，对应总市值 213.3 亿元，近一年+37.1\%，总市值在三级行业“果蔬加工”的 5 家公司中按由大到小排第\zhnumber{1}位，近一年涨跌幅高于全样本中位数（-10.4\%）47.5 个百分点。 按 2025 年年报披露的行业构成，收入占比前三位为饮料制造业 97.2\%；果渣制造业 2.5\%；其他(补充) 0.3\%。

\pfl{股权结构与控制权} 按 2026 年 6 月 30 日的股东名册，前三大股东为香港中央结算(代理人)有限公司（持股比例未披露）；Donghua Fruit Industry Co. 持股 16.57\%；山东安德利集团有限公司 持股 14.55\%，股东户数 9,683 户。 股权结构的另一条线索是，2026 年 8 月，王安将所持安德利集团 45\% 股权转让给女儿王萌，股东杨玉华将 10\% 股权转让给王淑娟。

\pfl{经营要点} 从经营要点看，2026 年上半年末，公司在建工程 3,360.18 万元、较上年末增长 531.52\%。

\pfl{所处行业的基本面} 果蔬加工的价格由产地原料果价格与出口订单决定，需求端受海外需求与关税影响（第 \ref{sec:framework} 章的供需框架）。 与该公司直接相关的行业数据包括：2025 年 12 月，公司在互动平台表示，预计 2026 年浓缩果汁原料果采购价较 2025 年略有下降。

\pfl{竞争优势} 把该公司放回同类公司比较，2026 年上半年营业收入 7.2 亿元、在三级行业“果蔬加工”的 5 家公司中排第\zhnumber{3}位；毛利率 24.3\%、高于行业内中位数 15.7 个百分点；半年 ROE 4.3\%（未年化）、高于行业内中位数（2.26\%）；资产负债率 7.8\%、低于行业内中位数（30.8\%）。 优势在报表上主要体现为高于行业平均的毛利率，这在周期底部是抵御亏损的主要缓冲。

\pfl{历史股价} 自 2021 年 1 月至 2026 年 10 月，股价由 21.10 元变为 61.12 元，累计 +189.7\%，区间最高 67.75 元（2026 年 8 月）、最低 15.56 元（2022 年 4 月）；当前价相当于区间最高价的 90.2\%，为区间最低价的 3.93 倍。 月度收益的年化波动率 53.2\%，高于全样本中位数 37.3\%，在三级行业“果蔬加工”的 5 家公司中波动率由高到低排第\zhnumber{1}位。 把股价阶段与公司事件对齐看，2024 年 8 月 至 2025 年 5 月股价上涨+172.1\%，同期（2025 年 5 月）公司终止不超 3.00 亿元简易程序定增。 把这些阶段与公司事件对齐后可以看到，几轮大涨大跌多由公司自身的资本运作与风险事件触发，行业景气只解释了其中一部分。

\begin{center}
\begin{minipage}{\textwidth}\centering
\begin{tikzpicture}
\begin{axis}[
  width=12.6cm, height=3.9cm, scale only axis,
  xmin=2020.922, xmax=2026.972, ymin=14.47, ymax=72.49,
  xtick={2021,2022,2023,2024,2025,2026},
  yearticks,
  yticklabel style={font=\scriptsize},
  ylabel={元/股}, ylabel style={font=\scriptsize},
  grid=major, grid style={LightGray, line width=0.3pt},
  axis line style={InkGray, line width=0.5pt},
  every axis plot/.append style={line width=0.9pt},
]
\addplot[AgriGreen] table[x=x,y=y]{data/clean/profiles/px_605198.dat};
\addplot[SoftRed, dashed, line width=0.7pt] table[x=x,y=y]{data/clean/profiles/pxzz_605198.dat};
\addplot[only marks, mark=*, mark size=1.1pt, SoftRed] table[x=x,y=y]{data/clean/profiles/pxzz_605198.dat};
\end{axis}
\end{tikzpicture}

\captionof{figure}[安德利（605198）股价与周期骨架]{安德利（605198）月度收盘价（元/股）与 ZigZag 周期骨架（阈值 25\%，圆点为周期转折点）}
\end{minipage}
\end{center}

\pfl{过去周期分析} 以 25\% 的 ZigZag 阈值划分，样本区间内股价形成 5 个主升段与 4 个主跌段。 最大主跌段为 2024 年 4 月 至 2024 年 8 月，4 个月-44.8\%；最大主升段出现在 2024 年 8 月 至 2025 年 5 月，9 个月+172.1\%。 最近一次转折出现在 2026 年 8 月（上涨段自 2025 年 12 月起，8 个月+80.1\%），当前处于该段之后的震荡之中。 公司股价较申万农林牧渔指数提前约 49 个月见底，是板块中较早定价周期反转的公司之一。 果蔬加工的外向型特征明显，出口订单与原料产地价格决定利润率，公司 2026 年上半年实现盈利、半年 ROE 4.3\%（未年化），好于全样本中位数（0.75\%），下行周期对其报表的冲击小于同业。 综合区间形态看，该股单段涨跌幅度偏大、以趋势段为主（最大主升 +172.1\%、最大主跌 -44.8\%），年化波动率 53.2\% 与全样本中位数 37.3\% 相比波动明显大于样本中位数。

\pfl{盈利情况} 2026 年上半年营业收入 7.16 亿元（同比 -24.4\%），归母净利润 1.23 亿元，同比 -38.7\%。 以年报为趋势对照，2023---2025 年营业收入由 8.76 亿元变为 16.77 亿元（两年年均复合增速 +38.4\%），归母净利润由 2.56 亿元变为 3.30 亿元。 按半年报口径，2026 年上半年的 ROE 为 4.3\%（未年化）、销售净利率 17.2\%、毛利率 24.3\%，ROE 在“果蔬加工”5 家公司中按数值由高到低排第\zhnumber{2}位。 按同一口径，三级行业“果蔬加工”5 家公司的半年 ROE 中位数为 2.26\%，该公司高于其中位数 2.07 个百分点。 面对这一局面，公司披露的举措包括：2026 年 9 月，2026 年一季度已扭亏为盈但全年盈利仍存不确定性。 综合看，公司在报告期维持盈利（高于全样本半年 ROE 中位数 0.75\%），利润的可持续性取决于果蔬加工景气的延续与成本管控的执行。

\begin{center}\small\setlength{\tabcolsep}{3.4pt}
\captionof{table}[安德利（605198）关键财务指标]{安德利（605198）关键财务指标（合并报表口径；两个半年报期均未年化；现金短债比 $=$ 货币资金 $\div$ 短期债务）}
\begin{adjustbox}{max width=\textwidth}
\begin{tabular}{@{}lrrrrr@{}}
\toprule
指标 & 2023 年 & 2024 年 & 2025 年 & 2025 年半年报 & 2026 年半年报 \\
\midrule
营业收入（亿元） & 8.76 & 14.18 & 16.77 & 9.48 & 7.16 \\
归母净利润（亿元） & 2.56 & 2.61 & 3.30 & 2.01 & 1.23 \\
毛利率（\%） & 33.0 & 24.8 & 24.3 & 23.9 & 24.3 \\
ROE（\%） & 10.3 & 9.9 & 11.7 & 7.4 & 4.3 \\
资产负债率（\%） & 3.1 & 4.8 & 5.3 & 7.1 & 7.8 \\
经营活动现金流净额（亿元） & -0.30 & -1.09 & 4.92 & 7.92 & 5.38 \\
货币资金（亿元） & 5.44 & 2.36 & 4.09 & 9.59 & 7.39 \\
有息负债合计（亿元） & --- & --- & --- & --- & --- \\
现金短债比（倍） & --- & --- & --- & --- & --- \\
\bottomrule
\end{tabular}
\end{adjustbox}
\end{center}

\begin{center}
\begin{minipage}{\textwidth}\centering
\begin{tikzpicture}
\begin{groupplot}[
  group style={group size=2 by 1, horizontal sep=0.60cm},
  width=6.85cm, height=4.60cm,
  tick label style={font=\tiny},
  title style={font=\scriptsize\heiti},
  yticklabel style={font=\tiny},
  ylabel style={font=\tiny},
  grid=major, grid style={LightGray, line width=0.3pt},
  axis line style={InkGray, line width=0.5pt},
  enlarge x limits={abs=0.8},
  legend style={font=\scriptsize, at={(0.5,-0.20)}, anchor=north,
                legend columns=2, draw=none, fill=none, column sep=6pt},
]
\nextgroupplot[
  ybar, bar width=4pt,
  ymin=-1.68, ymax=18.78,
  xtick={0,1,2,3,4,5.7},
  xticklabels={2021,2022,2023,2024,2025,2026H1},
  ylabel={亿元},
]
\addplot[fill=AgriGreen!75, draw=AgriGreen, bar shift=-2.6pt] table[x=i, y=rev]{data/clean/profiles/fin_605198.dat};
\addplot[fill=SoftRed!60, draw=SoftRed, bar shift=2.6pt] table[x=i, y=ni]{data/clean/profiles/fin_605198.dat};
\addplot[fill=AgriGreen!30, draw=AgriGreen, pattern=north east lines, bar shift=-2.6pt] table[x=x, y=rev]{data/clean/profiles/finh_605198.dat};
\addplot[fill=SoftRed!20, draw=SoftRed, pattern=north east lines, bar shift=2.6pt] table[x=x, y=ni]{data/clean/profiles/finh_605198.dat};
\legend{营业收入（年/半年）, 归母净利润（年/半年）}
\nextgroupplot[
  ymin=-0.9, ymax=12.9,
  xtick={0,1,2,3,4,5.7},
  xticklabels={2021,2022,2023,2024,2025,2026H1},
  ylabel={\%},
]
\addplot[AgriGold, mark=*, mark size=1.3pt] table[x=i, y=roe]{data/clean/profiles/fin_605198.dat};
\addplot[OilBlue, mark=square*, mark size=1.3pt] table[x=i, y=debt]{data/clean/profiles/fin_605198.dat};
\addplot[only marks, AgriGold, mark=*, mark size=2.0pt] table[x=x, y=roe]{data/clean/profiles/finh_605198.dat};
\addplot[only marks, OilBlue, mark=square*, mark size=2.0pt] table[x=x, y=debt]{data/clean/profiles/finh_605198.dat};
\legend{ROE, 资产负债率}
\end{groupplot}
\end{tikzpicture}
\begin{tikzpicture}
\begin{axis}[
  name=finbot,
  width=13.75cm, height=3.10cm, scale only axis,
  ymin=-0.74, ymax=8.28,
  xtick={0,1,2,3,4,5.7},
  xticklabels={2021,2022,2023,2024,2025,2026H1},
  tick label style={font=\tiny},
  yticklabel style={font=\tiny},
  ylabel={亿元}, ylabel style={font=\tiny},
  ybar, bar width=4pt,
  grid=major, grid style={LightGray, line width=0.3pt},
  axis line style={InkGray, line width=0.5pt},
  enlarge x limits={abs=0.8},
  legend style={font=\scriptsize, at={(0.5,-0.26)}, anchor=north,
                legend columns=2, draw=none, fill=none, column sep=10pt},
]
\addplot[fill=AgriGreen!55, draw=AgriGreen, bar shift=-3.0pt] table[x=i, y=cash]{data/clean/profiles/fin_605198.dat};
\addplot[fill=SoftRed!45, draw=SoftRed, bar shift=3.0pt] table[x=i, y=ibd]{data/clean/profiles/fin_605198.dat};
\addplot[fill=AgriGreen!25, draw=AgriGreen, pattern=north east lines, bar shift=-3.0pt] table[x=x, y=cash]{data/clean/profiles/finh_605198.dat};
\addplot[fill=SoftRed!20, draw=SoftRed, pattern=north east lines, bar shift=3.0pt] table[x=x, y=ibd]{data/clean/profiles/finh_605198.dat};
\legend{货币资金（左轴）, 有息负债（左轴）}
\end{axis}
\begin{axis}[
  at={(finbot.south east)}, anchor=south east,
  width=13.75cm, height=3.10cm, scale only axis,
  axis y line*=right, axis x line*=none, xtick=\empty,
  ymin=-2.78, ymax=7.32,
  tick label style={font=\tiny},
  yticklabel style={font=\tiny}, scaled y ticks=false,
  legend style={font=\scriptsize, at={(0.5,-0.44)}, anchor=north,
                legend columns=1, draw=none, fill=none},
]
\addplot[OilBlue, mark=*, mark size=2.2pt, line width=1.3pt] table[x=i, y=ocf]{data/clean/profiles/fin_605198.dat};
\addplot[only marks, OilBlue, mark=*, mark size=3.2pt] table[x=x, y=ocf]{data/clean/profiles/finh_605198.dat};
\legend{经营活动现金流净额（右轴，亿元，蓝点）}
\end{axis}
\end{tikzpicture}

\captionof{figure}[安德利（605198）近年主要财务指标]{安德利（605198）近年主要财务指标（左上：营业收入与归母净利润；右上：ROE 与资产负债率；下：货币资金、有息负债为左轴柱状，经营活动现金流净额为其右轴的蓝点折线。
2021---2025 年为年报口径，2026H1 为 2026 年半年报/半年末，与其前一年报之间留出间隔、以斜纹或空心点区分；半年报数据未年化）}
\end{minipage}
\end{center}

\pfl{财务分析} 2026 年 6 月末资产负债率 7.8\%，较 2025 年末上升 2.5 个百分点（样本中位数 52.2\%，所处三级行业内按由低到高排第\zhnumber{1}位）。 流动比率 8.54、速动比率 5.56，短期流动性无明显缺口。 2026 年上半年经营活动现金流净额 5.38 亿元，经营现金流与利润同向为正，内生资金可以支撑运营。 综合看，资产负债率低于样本中位数 44.4 个百分点。 公司披露的诉讼与资产冻结风险为：2026 年 9 月，标的公司股东丽水毓甬创业投资合伙企业（有限合伙）持有标的公司 1.46\% 股权被冻结。

\pfl{重大战略转型} 从公开披露看，公司的战略推进集中在：并购与重组（2 项）：2026 年 9 月，公司以约 7.93 亿元现金收购宁波甬强科技约 62.06\% 股权；2023 年，扩产与并购并行：新建浓缩桃汁、浓缩山楂汁生产线；融资与资本运作（1 项）：2024 年 5 月，公司拟定增投向脱色脱酸浓缩果汁与 NFC 果汁产线。 公告口径上，近三年（2023 年 9 月---2026 年 9 月）公司披露重大重组与并购类公告 3 条、对外投资与转型类公告 5 条、回购与股权激励类公告 36 条，合计公告 487 条。 第一大行业仍是“饮料制造业”，收入占比 97.2\%，与 2021 年年报相比没有方向性变化。 综合判断，报告期内公司有 8 条与战略相关的公告落地（重大重组与并购 3 条、对外投资与转型 5 条），资本运作与主业调整之间存在直接关联。

\begin{center}\small\setlength{\tabcolsep}{4pt}
\captionof{table}[安德利（605198）历史融资与资本运作]{安德利（605198）历史融资与资本运作（据公司公告与公开报道整理；金额为披露值，含首次公开发行、再融资、债券、银行借款与重整投资等）}
\begin{adjustbox}{max width=\textwidth}
\begin{tabular}{@{}p{1.45cm}p{2.55cm}p{4.05cm}p{4.95cm}@{}}
\toprule
时间 & 方式 & 规模 & 用途与说明 \\
\midrule
2003-04-22 & H 股 IPO（香港联交所创业板，股份代号 8259） & --- & 国内果汁企业首家在香港上市的公司；2011-01-19 转主板 \\
2020-09-18 & A 股 IPO（上海证券交易所） & 发行 2,\allowbreak{}000.00 万股，\allowbreak{}发行价 7.60 元/\allowbreak{}股，\allowbreak{}募集资金总额 1.52 亿元 & 多品种浓缩果汁生产线建设项目（多品种果汁生产项目）；；首发前总股本 35,\allowbreak{}800 万股 \\
2024-05-09 & 拟以简易程序向特定对象发行 A 股股票 & 拟募集资金总额不超过 3.00 亿元（含本数） & 年产 7,\allowbreak{}200 吨脱色脱酸浓缩果汁生产线项目、\allowbreak{}年产 1.2 万吨 NFC 果汁项目、\allowbreak{}补充流动资金项目；2025 年 5 月 20 日第九届董事会第一次会议审议通过终止公司 2024 年度以简易程序向特定对象发行 A 股股票事项 \\
2024-06—2024-12 & H 股回购 & 累计回购 H 股 780 万股，\allowbreak{}支付资金总额 7,\allowbreak{}374.05 万港元（不含佣金等费用） & 注销/\allowbreak{}回购 H 股股份；2024 年 6 月 6 日至 12 月 31 日期间实施 H 股回购 14 次，\allowbreak{}已发行 H 股 9.94\%；2025 年 2 月 13 日公告完成 780 万股 H 股股份注销 \\
2026-09 & 现金收购（跨界取得控制权） & 合计约 7.93 亿元收购宁波甬强科技有限公司约 62.06\% 股权 & 取得甬强科技控制权，\allowbreak{}11 名股东 24.95\% 股权（按 9.00 亿元估值）、\allowbreak{}4.11 亿元收购 14 家机构股东 24.61\% 股权（按年化收益率定价）；整体估值约 12.77 亿元 \\
\bottomrule
\end{tabular}
\end{adjustbox}
\end{center}

\pfl{历史融投资情况} 从现金流量表看，2025 年公司，筹资活动现金流净额 -1.94 亿元（2023 年为取得借款 0.30 亿元、偿还债务 1.60 亿元）。 2026 年上半年筹资活动现金流净额为 -0.97 亿元（取得借款 — 亿元、偿还债务 — 亿元，未年化），已转为净流出，处于净还债状态。 2025 年分配股利、利润或偿付利息支付的现金合计 0.83 亿元。 公开披露的融资记录共 5 笔，工具类型以 IPO、定向增发为主；金额与用途见下表。 其中规模最大的两笔为：2020 年 9 月，IPO，发行 2,000.00 万股，发行价 7.60 元/股，募集资金总额 1.52 亿元；2026 年 9 月，现金收购，合计约 7.93 亿元收购宁波甬强科技有限公司约 62.06\% 股权。 股本方面，近三年披露股本变动 5 次（限售股份上市 1 次、承诺限售 1 次、股份回购 1 次）；总股本由 2021 年末的 3.67 亿股变为 3.49 亿股（-4.98\%）。 分红方面，近三年现金分红 5 次、合计每股派现 0.92 元，最近一次实施在 2026 年 7 月。 近三年“再融资”类公告 17 条，涉及定向增发、募集资金使用。 综合看，公司融资结构上股权与债务融资并举；2025 年筹资活动现金流净额 -1.94 亿元，融资节奏仍以维持存量债务周转为主，与前述经营与投资现金流的方向基本一致。

\pfl{未来潜在融资需求分析} 公司 2026 年 6 月末资产负债率 7.8\%、2026 年上半年 ROE 4.3\%（未年化），资产端与负债端都没有明显压力，融资需求以相机抉择的扩张型用途为主。 公开披露的后续资金安排线索包括：股权融资线索：2025 年 5 月，公司终止不超 3.00 亿元简易程序定增。




\subsubsection{中粮科技（000930）}
\label{co:000930}

\pfl{公司基本情况} 1998 年设立至今，中粮科技是原名安徽丰原生物化学，现为中粮集团旗下玉米深加工与生物能源平台，主营燃料乙醇、淀粉及淀粉糖、味精等产品，国内燃料乙醇市场份额居前；2025 年盈利微薄，2026 年上半年因产品价格下行转为亏损。 公司于 1999-07-12 在深交所主板首发上市，注册资本 18.56 亿元，总股本 18.60 亿股。 截至 2026 年 10 月 12 日收盘价 5.19 元，对应总市值 96.5 亿元，近一年-14.1\%，总市值在三级行业“其他农产品加工”的 11 家公司中按由大到小排第\zhnumber{2}位，近一年涨跌幅低于全样本中位数（-10.4\%）3.7 个百分点。 按 2026 年半年报披露的行业构成，收入占比前两位为农产品加工及销售 98.3\%；其他(补充) 1.7\%。 股权结构方面，实际控制人为中粮集团。

\pfl{股权结构与控制权} 按 2026 年 6 月 30 日的股东名册，前三大股东为 COFCO Bio-chemical Investm 持股 47.58\%；大耀香港有限公司 持股 8.19\%；香港中央结算有限公司 持股 2.08\%，股东户数 102,006 户。 控制权与股东性质方面，2018 年，重大资产重组完成后，Ltd.以 8.83 亿股、47.80\% 成为控股股东；2011 年 9 月中粮集团将所持 2 亿股注入大耀香港。

\pfl{经营要点} 业务与产能层面，食品原料板块已初步建立淀粉糖竞争优势、产销量位居行业前列；公司是国内首家通过酶法工艺实现阿洛酮糖规模化生产并上市销售的企业。

\pfl{所处行业的基本面} 该子行业横跨制糖、番茄制品、糖醇与添加剂，供给受种植面积与进口配额约束，需求端分食用与工业两类（第 \ref{sec:framework} 章的供需框架）。 行业层面的最新情况是：中国玉米淀粉行业产能由 2015 年的 1,766 万吨提升至 2025 年的 2,944 万吨、年均复合增速 5.2\%；政策严控玉米乙醇新增产能。

\pfl{竞争优势} 把该公司放回同类公司比较，2026 年上半年营业收入 81.9 亿元、在三级行业“其他农产品加工”的 11 家公司中排第\zhnumber{2}位；毛利率 5.0\%、低于行业内中位数 4.1 个百分点；半年 ROE -1.6\%（未年化）、低于行业内中位数（1.36\%）；资产负债率 40.7\%、高于行业内中位数（39.1\%）。 公司区别于同类公司的经营特征包括：公司始终保持燃料乙醇行业市场份额领先地位。 就当期报表看，公司尚未把经营层面的特征转化为高于行业的资本回报，周期底部的相对优势更多体现在现金流与资产负债表上。

\pfl{历史股价} 以 2021 年 1 月为起点、2026 年 10 月为终点，股价由 10.89 元变为 5.19 元，累计 -52.3\%，区间最高 11.68 元（2021 年 11 月）、最低 4.25 元（2026 年 6 月）；当前价相当于区间最高价的 44.4\%，为区间最低价的 1.22 倍。 月度收益的年化波动率 28.0\%，低于全样本中位数 37.3\%，在三级行业“其他农产品加工”的 11 家公司中波动率由高到低排第\zhnumber{11}位。

\begin{center}
\begin{minipage}{\textwidth}\centering
\begin{tikzpicture}
\begin{axis}[
  width=12.6cm, height=3.9cm, scale only axis,
  xmin=2020.922, xmax=2026.972, ymin=3.95, ymax=12.50,
  xtick={2021,2022,2023,2024,2025,2026},
  yearticks,
  yticklabel style={font=\scriptsize},
  ylabel={元/股}, ylabel style={font=\scriptsize},
  grid=major, grid style={LightGray, line width=0.3pt},
  axis line style={InkGray, line width=0.5pt},
  every axis plot/.append style={line width=0.9pt},
]
\addplot[AgriGreen] table[x=x,y=y]{data/clean/profiles/px_000930.dat};
\addplot[SoftRed, dashed, line width=0.7pt] table[x=x,y=y]{data/clean/profiles/pxzz_000930.dat};
\addplot[only marks, mark=*, mark size=1.1pt, SoftRed] table[x=x,y=y]{data/clean/profiles/pxzz_000930.dat};
\end{axis}
\end{tikzpicture}

\captionof{figure}[中粮科技（000930）股价与周期骨架]{中粮科技（000930）月度收盘价（元/股）与 ZigZag 周期骨架（阈值 25\%，圆点为周期转折点）}
\end{minipage}
\end{center}

\pfl{过去周期分析} 以 25\% 的 ZigZag 阈值划分，样本区间内股价形成 2 个主升段与 3 个主跌段。 最大主跌段为 2021 年 11 月 至 2024 年 8 月，33 个月-58.3\%；最大主升段出现在 2024 年 8 月 至 2026 年 2 月，18 个月+39.0\%。 最近一次转折出现在 2026 年 6 月（下跌段自 2026 年 2 月起，4 个月-37.2\%），当前处于该段的延续之中。 公司股价与申万农林牧渔指数几乎同步见底，个股并未走出独立行情。 该子行业横跨糖、番茄制品、添加剂等多个品种，各自的供给周期与政策（进口配额、关税）是主要变量，公司 2026 年 6 月末资产负债率 40.7\%、半年 ROE -1.6\%（未年化），在本轮下行中属于随行业同步调整的一类。 综合区间形态看，该股单段涨跌幅度偏大、以趋势段为主（最大主升 +39.0\%、最大主跌 -58.3\%），年化波动率 28.0\% 与全样本中位数 37.3\% 相比波动小于样本中位数。

\pfl{盈利情况} 2026 年上半年营业收入 81.94 亿元（同比 -7.0\%），归母净利润 -1.69 亿元，由上年同期的盈利转为亏损。 以年报为趋势对照，2023---2025 年营业收入由 203.79 亿元变为 175.12 亿元（两年年均复合增速 -7.3\%），归母净利润由 -6.02 亿元变为 0.38 亿元。 按半年报口径，2026 年上半年的 ROE 为 -1.6\%（未年化）、销售净利率 -2.1\%、毛利率 5.0\%，ROE 在“其他农产品加工”11 家公司中按数值由高到低排第\zhnumber{9}位。 按同一口径，三级行业“其他农产品加工”11 家公司的半年 ROE 中位数为 1.36\%，该公司低于其中位数 2.98 个百分点。 面对这一局面，公司披露的举措包括：2025 年，公司通过优化产品结构、扩大淀粉等产品销量，但产品价格下行压力大于销量增长贡献。 综合看，公司仍处亏损区间、由盈转亏，半年 ROE 在三级行业内按由高到低排第\zhnumber{9}位，单位盈利能力的修复依赖行业景气与自身成本端的改善，报表弹性主要体现为价格与销量的方向性变化。

\begin{center}\small\setlength{\tabcolsep}{3.4pt}
\captionof{table}[中粮科技（000930）关键财务指标]{中粮科技（000930）关键财务指标（合并报表口径；两个半年报期均未年化；现金短债比 $=$ 货币资金 $\div$ 短期债务）}
\begin{adjustbox}{max width=\textwidth}
\begin{tabular}{@{}lrrrrr@{}}
\toprule
指标 & 2023 年 & 2024 年 & 2025 年 & 2025 年半年报 & 2026 年半年报 \\
\midrule
营业收入（亿元） & 203.79 & 200.53 & 175.12 & 88.12 & 81.94 \\
归母净利润（亿元） & -6.02 & 0.25 & 0.38 & 1.07 & -1.69 \\
毛利率（\%） & 5.5 & 7.8 & 8.4 & 9.0 & 5.0 \\
ROE（\%） & -5.3 & 0.2 & 0.4 & 1.0 & -1.6 \\
资产负债率（\%） & 34.5 & 37.0 & 34.3 & 37.6 & 40.7 \\
经营活动现金流净额（亿元） & 19.58 & 1.56 & 18.47 & 7.92 & -8.75 \\
货币资金（亿元） & 14.81 & 19.09 & 24.88 & 23.88 & 21.29 \\
有息负债合计（亿元） & 36.43 & 44.21 & 24.33 & 44.90 & 49.14 \\
现金短债比（倍） & 0.45 & 0.47 & 1.30 & 0.60 & 0.62 \\
\bottomrule
\end{tabular}
\end{adjustbox}
\end{center}

\begin{center}
\begin{minipage}{\textwidth}\centering
\begin{tikzpicture}
\begin{groupplot}[
  group style={group size=2 by 1, horizontal sep=0.60cm},
  width=6.85cm, height=4.60cm,
  tick label style={font=\tiny},
  title style={font=\scriptsize\heiti},
  yticklabel style={font=\tiny},
  ylabel style={font=\tiny},
  grid=major, grid style={LightGray, line width=0.3pt},
  axis line style={InkGray, line width=0.5pt},
  enlarge x limits={abs=0.8},
  legend style={font=\scriptsize, at={(0.5,-0.20)}, anchor=north,
                legend columns=2, draw=none, fill=none, column sep=6pt},
]
\nextgroupplot[
  ybar, bar width=4pt,
  ymin=-30.18, ymax=264.59,
  xtick={0,1,2,3,4,5.7},
  xticklabels={2021,2022,2023,2024,2025,2026H1},
  ylabel={亿元},
]
\addplot[fill=AgriGreen!75, draw=AgriGreen, bar shift=-2.6pt] table[x=i, y=rev]{data/clean/profiles/fin_000930.dat};
\addplot[fill=SoftRed!60, draw=SoftRed, bar shift=2.6pt] table[x=i, y=ni]{data/clean/profiles/fin_000930.dat};
\addplot[fill=AgriGreen!30, draw=AgriGreen, pattern=north east lines, bar shift=-2.6pt] table[x=x, y=rev]{data/clean/profiles/finh_000930.dat};
\addplot[fill=SoftRed!20, draw=SoftRed, pattern=north east lines, bar shift=2.6pt] table[x=x, y=ni]{data/clean/profiles/finh_000930.dat};
\legend{营业收入（年/半年）, 归母净利润（年/半年）}
\nextgroupplot[
  ymin=-9.0, ymax=45.3,
  xtick={0,1,2,3,4,5.7},
  xticklabels={2021,2022,2023,2024,2025,2026H1},
  ylabel={\%},
]
\addplot[AgriGold, mark=*, mark size=1.3pt] table[x=i, y=roe]{data/clean/profiles/fin_000930.dat};
\addplot[OilBlue, mark=square*, mark size=1.3pt] table[x=i, y=debt]{data/clean/profiles/fin_000930.dat};
\addplot[only marks, AgriGold, mark=*, mark size=2.0pt] table[x=x, y=roe]{data/clean/profiles/finh_000930.dat};
\addplot[only marks, OilBlue, mark=square*, mark size=2.0pt] table[x=x, y=debt]{data/clean/profiles/finh_000930.dat};
\legend{ROE, 资产负债率}
\end{groupplot}
\end{tikzpicture}
\begin{tikzpicture}
\begin{axis}[
  name=finbot,
  width=13.75cm, height=3.10cm, scale only axis,
  ymin=-4.91, ymax=55.04,
  xtick={0,1,2,3,4,5.7},
  xticklabels={2021,2022,2023,2024,2025,2026H1},
  tick label style={font=\tiny},
  yticklabel style={font=\tiny},
  ylabel={亿元}, ylabel style={font=\tiny},
  ybar, bar width=4pt,
  grid=major, grid style={LightGray, line width=0.3pt},
  axis line style={InkGray, line width=0.5pt},
  enlarge x limits={abs=0.8},
  legend style={font=\scriptsize, at={(0.5,-0.26)}, anchor=north,
                legend columns=2, draw=none, fill=none, column sep=10pt},
]
\addplot[fill=AgriGreen!55, draw=AgriGreen, bar shift=-3.0pt] table[x=i, y=cash]{data/clean/profiles/fin_000930.dat};
\addplot[fill=SoftRed!45, draw=SoftRed, bar shift=3.0pt] table[x=i, y=ibd]{data/clean/profiles/fin_000930.dat};
\addplot[fill=AgriGreen!25, draw=AgriGreen, pattern=north east lines, bar shift=-3.0pt] table[x=x, y=cash]{data/clean/profiles/finh_000930.dat};
\addplot[fill=SoftRed!20, draw=SoftRed, pattern=north east lines, bar shift=3.0pt] table[x=x, y=ibd]{data/clean/profiles/finh_000930.dat};
\legend{货币资金（左轴）, 有息负债（左轴）}
\end{axis}
\begin{axis}[
  at={(finbot.south east)}, anchor=south east,
  width=13.75cm, height=3.10cm, scale only axis,
  axis y line*=right, axis x line*=none, xtick=\empty,
  ymin=-23.62, ymax=65.57,
  tick label style={font=\tiny},
  yticklabel style={font=\tiny}, scaled y ticks=false,
  legend style={font=\scriptsize, at={(0.5,-0.44)}, anchor=north,
                legend columns=1, draw=none, fill=none},
]
\addplot[OilBlue, mark=*, mark size=2.2pt, line width=1.3pt] table[x=i, y=ocf]{data/clean/profiles/fin_000930.dat};
\addplot[only marks, OilBlue, mark=*, mark size=3.2pt] table[x=x, y=ocf]{data/clean/profiles/finh_000930.dat};
\legend{经营活动现金流净额（右轴，亿元，蓝点）}
\end{axis}
\end{tikzpicture}

\captionof{figure}[中粮科技（000930）近年主要财务指标]{中粮科技（000930）近年主要财务指标（左上：营业收入与归母净利润；右上：ROE 与资产负债率；下：货币资金、有息负债为左轴柱状，经营活动现金流净额为其右轴的蓝点折线。
2021---2025 年为年报口径，2026H1 为 2026 年半年报/半年末，与其前一年报之间留出间隔、以斜纹或空心点区分；半年报数据未年化）}
\end{minipage}
\end{center}

\pfl{财务分析} 2026 年 6 月末资产负债率 40.7\%，较 2025 年末上升 6.4 个百分点（样本中位数 52.2\%，所处三级行业内按由低到高排第\zhnumber{7}位）。 2026 年 6 月末货币资金 21.29 亿元、短期债务（短期借款与一年内到期非流动负债）34.61 亿元、长期债务 14.53 亿元，有息负债合计 49.14 亿元，现金短债比 0.62 倍——覆盖偏紧。 净有息负债 27.85 亿元，相当于其 2026 年上半年营业收入的 34\%。 流动比率 1.27、速动比率 0.86，短期偿债能力处于正常区间。 2026 年上半年经营活动现金流净额 -8.75 亿元，经营现金流与利润同时为负，现金消耗较快。 2026 年 6 月末在建工程 7.92 亿元，较上年末增加 0.40 亿元，仍有在建投入。 2026 年 6 月末商誉 3.66 亿元，占归母净资产的 3.4\%，是净资产中最脆弱的一块。 综合看，现金短债比 0.62 倍，短期资金安排偏紧；资产负债率低于样本中位数 11.5 个百分点；有息负债中短期债务占比更高，期限结构仍以滚动续作为主。

\pfl{重大战略转型} 近三年公司的战略动作集中在：并购与重组（1 项）：2018 年，中粮集团将境内玉米深加工与燃料乙醇资产整体注入上市公司；主业转型（1 项）：2025 年，公司实现原料成本最优。 公告口径上，近三年（2023 年 9 月---2026 年 9 月）公司披露重大重组与并购类公告 0 条、对外投资与转型类公告 1 条、回购与股权激励类公告 8 条，合计公告 338 条。 第一大行业仍是“农产品加工及销售”，收入占比 98.3\%，与 2021 年年报相比没有方向性变化。 综合判断，报告期内公司有 1 条与战略相关的公告落地（重大重组与并购 0 条、对外投资与转型 1 条），资本运作与主业调整之间存在直接关联。

\begin{center}\small\setlength{\tabcolsep}{4pt}
\captionof{table}[中粮科技（000930）历史融资与资本运作]{中粮科技（000930）历史融资与资本运作（据公司公告与公开报道整理；金额为披露值，含首次公开发行、再融资、债券、银行借款与重整投资等）}
\begin{adjustbox}{max width=\textwidth}
\begin{tabular}{@{}p{1.45cm}p{2.55cm}p{4.05cm}p{4.95cm}@{}}
\toprule
时间 & 方式 & 规模 & 用途与说明 \\
\midrule
1999-07-12 & IPO（深圳证券交易所，原安徽丰原生物化学股份有限公司） & 实际发行量 6,\allowbreak{}000 万股，\allowbreak{}发行价 6.50 元/\allowbreak{}股，\allowbreak{}实际募集资金合计 3.90 亿元 & 上市地深圳证券交易所，\allowbreak{}发行方式网上定价发行；首发前总股本 8,\allowbreak{}697.88 万股 \\
2003-04-24 & 可转换公司债券（“丰原转债”） & 5.00 亿元（5.00 亿元） & 经中国证监会证监公司字〔2003〕36 号核准；2003 年 5 月 20 日在深交所挂牌，\allowbreak{}代码 125930；截至 2006 年 12 月 31 日累计转股 9,\allowbreak{}860.28 万股 \\
2018 & 发行股份购买资产（重大资产重组暨关联交易） & 标的资产评估值合计 82.85 亿元；新增股份 8.83 亿股 & 收购生化能源、\allowbreak{}生物化学、\allowbreak{}桦力投资 100\% 股权 \\
2022 & 现金分红 & 以 18.65 亿股为基数每 10 股派发现金红利 5.70 元（含税） & 股东回报；该年度公司盈利水平较高，\allowbreak{}分红力度显著大于后续年度 \\
\bottomrule
\end{tabular}
\end{adjustbox}
\end{center}

\pfl{历史融投资情况} 从现金流量表看，2025 年公司取得借款收到的现金 69.69 亿元、偿还债务支付的现金 76.63 亿元，筹资活动现金流净额 -7.88 亿元（2023 年为取得借款 75.65 亿元、偿还债务 72.42 亿元）。 2026 年上半年筹资活动现金流净额为 14.03 亿元（取得借款 39.63 亿元、偿还债务 25.42 亿元，未年化），仍处于净融入状态。 2025 年分配股利、利润或偿付利息支付的现金合计 0.86 亿元。 公开披露的融资记录共 4 笔，工具类型以 IPO、可转债、发行股份购买资产为主；金额与用途见下表。 其中规模最大的两笔为：2003 年 4 月，可转债，5.00 亿元（5.00 亿元）；1999 年 7 月，IPO，实际发行量 6,000 万股，发行价 6.50 元/股，实际募集资金合计 3.90 亿元。 股本方面，近三年披露股本变动 6 次（股份回购 2 次、激励股份解禁 1 次）；总股本由 2021 年末的 18.66 亿股变为 18.60 亿股（-0.32\%）。 分红方面，近三年现金分红 2 次、合计每股派现 0.60 元，最近一次实施在 2025 年 12 月。 综合看，公司融资结构上股权与债务融资并举；2025 年筹资活动现金流净额 -7.88 亿元，融资节奏仍以维持存量债务周转为主，与前述经营与投资现金流的方向基本一致。

\pfl{未来潜在融资需求分析} 2026 年 6 月末货币资金 21.29 亿元，而短期债务 34.61 亿元（现金短债比 0.62 倍），2026 年上半年经营现金流净额 -8.75 亿元。按“短期债务减去账面货币资金”的滚动偿付口径，静态缺口约 13.32 亿元；若再计入上半年亏损的年化额（3.39 亿元），维持一年正常经营所需的外部资金约 13.32---16.70 亿元。 公司 2026 年上半年归母净利润为负（-1.69 亿元），但 6 月末资产负债率 40.7\% 尚未进入第一档（亏损且负债率 $>$65\%）的区间，当前融资需求以补充流动资金、支撑低谷期经营性支出为主。 资金需求还要叠加在建工程：期末在建工程 7.92 亿元、2026 年上半年购建固定资产等支出 5.45 亿元（未年化），这部分支出在周期底部通常会被主动放缓。 综合看，公司的外部资金需求以补充营运资金为主、项目投入为辅，滚动偿付口径下的静态缺口约 13.32 亿元，融资的时点与规模更多取决于信用条件与资本市场窗口。




\subsubsection{金健米业（600127）}
\label{co:600127}

\pfl{公司基本情况} 公司前身可追溯至 1998 年，湖南常德地方国资背景的粮油食品加工企业，主营粮油食品加工、农产品贸易与乳品；收入规模数十亿元但盈利极薄，累计可供分配利润为负，2024 年与控股股东完成资产置换后主业结构调整。 公司于 1998-05-06 在上交所首发上市，注册资本 6.42 亿元，总股本 6.42 亿股。 截至 2026 年 10 月 12 日收盘价 13.90 元，对应总市值 89.2 亿元，近一年+104.7\%，总市值在三级行业“粮油加工”的 7 家公司中按由大到小排第\zhnumber{2}位，近一年涨跌幅高于全样本中位数（-10.4\%）115.1 个百分点。 按 2026 年半年报披露的行业构成，收入占比前三位为粮油业务 94.5\%；食品行业 3.4\%；乳品业务 2.3\%。 与 2021 年年报相比，第一大行业口径由“粮油食品加工业”（48.1\%）变为“粮油业务”（94.5\%）；两期披露的分类口径有调整，占比差异中同时包含分类因素。 股权结构方面，控股股东为湖南粮食集团有限责任公司；实际控制人为湖南省国资委。

\pfl{股权结构与控制权} 按 2026 年 6 月 30 日的股东名册，前三大股东为湖南粮食集团有限责任公司 持股 21.34\%；香港中央结算有限公司 持股 0.38\%；陶学锭 持股 0.27\%，股东户数 94,848 户。 股东结构与控制权方面，2021 年 7 月，划转前湖南粮食集团通过金霞粮食持有公司 1.37 亿股（占总股本 21.34\%）；2022 年 7 月，实际控制人湖南省国资委旗下的湖南省现代农业产业控股集团有限公司与湖南粮食集团有限责任公司合并组建成立湖南农业发展投资集团有限责任公司。

\pfl{经营要点} 经营层面，2024 年，公司置入裕湘食品与中南粮科院 2 家子公司股权、置出 3 家子公司股权至控股股东。

\pfl{所处行业的基本面} 粮油加工的供给由大豆、菜籽等原料进口与压榨产能决定，需求由食用消费决定，利润来自原料与成品的价差而非价格方向（第 \ref{sec:framework} 章的供需框架）。 行业层面的最新情况是：2026 年上半年，我国稻谷和大米进口量 243 万吨、同比增长 67.7\%；限定收购总量 5,000 万吨；截至 2026 年 6 月末全国早籼稻、中晚籼稻和粳稻收购均价分别为 2,624 元/吨、2,594 元/吨和 2,689 元/吨。

\pfl{竞争优势} 把该公司放回同类公司比较，2026 年上半年营业收入 19.7 亿元、在三级行业“粮油加工”的 7 家公司中排第\zhnumber{5}位；毛利率 6.0\%、低于行业内中位数 0.4 个百分点；半年 ROE -1.4\%（未年化）、低于行业内中位数（2.22\%）；资产负债率 62.2\%、高于行业内中位数（43.5\%）。 公司区别于同类公司的经营特征包括：公司为湖南省属国资控股的粮油上市平台。 就当期报表看，公司尚未把经营层面的特征转化为高于行业的资本回报，周期底部的相对优势更多体现在现金流与资产负债表上。

\pfl{历史股价} 以 2021 年 1 月为起点、2026 年 10 月为终点，股价由 9.10 元变为 13.90 元，累计 +52.7\%，区间最高 13.90 元（2026 年 10 月）、最低 5.20 元（2026 年 6 月）；当前价相当于区间最高价的 100.0\%，为区间最低价的 2.67 倍。 月度收益的年化波动率 43.2\%，高于全样本中位数 37.3\%，在三级行业“粮油加工”的 7 家公司中波动率由高到低排第\zhnumber{1}位。

\begin{center}
\begin{minipage}{\textwidth}\centering
\begin{tikzpicture}
\begin{axis}[
  width=12.6cm, height=3.9cm, scale only axis,
  xmin=2020.922, xmax=2026.972, ymin=4.84, ymax=14.87,
  xtick={2021,2022,2023,2024,2025,2026},
  yearticks,
  yticklabel style={font=\scriptsize},
  ylabel={元/股}, ylabel style={font=\scriptsize},
  grid=major, grid style={LightGray, line width=0.3pt},
  axis line style={InkGray, line width=0.5pt},
  every axis plot/.append style={line width=0.9pt},
]
\addplot[AgriGreen] table[x=x,y=y]{data/clean/profiles/px_600127.dat};
\addplot[SoftRed, dashed, line width=0.7pt] table[x=x,y=y]{data/clean/profiles/pxzz_600127.dat};
\addplot[only marks, mark=*, mark size=1.1pt, SoftRed] table[x=x,y=y]{data/clean/profiles/pxzz_600127.dat};
\end{axis}
\end{tikzpicture}

\captionof{figure}[金健米业（600127）股价与周期骨架]{金健米业（600127）月度收盘价（元/股）与 ZigZag 周期骨架（阈值 25\%，圆点为周期转折点）}
\end{minipage}
\end{center}

\pfl{过去周期分析} 以 25\% 的 ZigZag 阈值划分，样本区间内股价形成 3 个主升段与 2 个主跌段。 最大主升段出现在 2026 年 6 月 至 2026 年 10 月，4 个月+167.3\%；最大主跌段为 2022 年 5 月 至 2024 年 7 月，26 个月-42.1\%。 最近一次转折出现在 2026 年 10 月（上涨段自 2026 年 6 月起，4 个月+167.3\%），当前处于该段之后的震荡之中。 公司股价与申万农林牧渔指数几乎同步见底，个股并未走出独立行情。 粮油加工是“成本加成”型生意（第 \ref{sec:grain} 章），利润来自原料与成品的价差及规模效率，收入规模大而净利率薄，公司 2026 年 6 月末资产负债率 62.2\%、半年 ROE -1.4\%（未年化），在本轮下行中属于随行业同步调整的一类。 综合区间形态看，该股单段涨跌幅度偏大、以趋势段为主（最大主升 +167.3\%、最大主跌 -42.1\%），年化波动率 43.2\% 与全样本中位数 37.3\% 相比波动明显大于样本中位数。

\pfl{盈利情况} 2026 年上半年营业收入 19.70 亿元（同比 +25.2\%），归母净利润 -0.10 亿元，由上年同期的盈利转为亏损。 以年报为趋势对照，2023---2025 年营业收入由 52.69 亿元变为 33.58 亿元（两年年均复合增速 -20.2\%），归母净利润由 0.14 亿元变为 0.04 亿元。 按半年报口径，2026 年上半年的 ROE 为 -1.4\%（未年化）、销售净利率 -0.5\%、毛利率 6.0\%，ROE 在“粮油加工”7 家公司中按数值由高到低排第\zhnumber{7}位。 同期全样本半年 ROE 中位数 0.75\%，公司与之相差 2.19 个百分点（弱于中位数公司）。 综合看，公司仍处亏损区间、由盈转亏，半年 ROE 在三级行业内按由高到低排第\zhnumber{7}位，单位盈利能力的修复依赖行业景气与自身成本端的改善，报表弹性主要体现为价格与销量的方向性变化。

\begin{center}\small\setlength{\tabcolsep}{3.4pt}
\captionof{table}[金健米业（600127）关键财务指标]{金健米业（600127）关键财务指标（合并报表口径；两个半年报期均未年化；现金短债比 $=$ 货币资金 $\div$ 短期债务）}
\begin{adjustbox}{max width=\textwidth}
\begin{tabular}{@{}lrrrrr@{}}
\toprule
指标 & 2023 年 & 2024 年 & 2025 年 & 2025 年半年报 & 2026 年半年报 \\
\midrule
营业收入（亿元） & 52.69 & 46.27 & 33.58 & 15.74 & 19.70 \\
归母净利润（亿元） & 0.14 & 0.02 & 0.04 & 0.12 & -0.10 \\
毛利率（\%） & 4.6 & 5.8 & 7.7 & 8.3 & 6.0 \\
ROE（\%） & 1.6 & 0.3 & 0.6 & 1.7 & -1.4 \\
资产负债率（\%） & 54.8 & 59.4 & 61.3 & 57.4 & 62.2 \\
经营活动现金流净额（亿元） & 1.90 & -0.44 & 0.72 & 1.07 & 0.38 \\
货币资金（亿元） & 1.49 & 1.45 & 1.86 & 1.57 & 1.50 \\
有息负债合计（亿元） & 7.12 & 7.37 & 7.78 & 6.75 & 7.31 \\
现金短债比（倍） & 0.23 & 0.20 & 0.24 & 0.23 & 0.21 \\
\bottomrule
\end{tabular}
\end{adjustbox}
\end{center}

\begin{center}
\begin{minipage}{\textwidth}\centering
\begin{tikzpicture}
\begin{groupplot}[
  group style={group size=2 by 1, horizontal sep=0.60cm},
  width=6.85cm, height=4.60cm,
  tick label style={font=\tiny},
  title style={font=\scriptsize\heiti},
  yticklabel style={font=\tiny},
  ylabel style={font=\tiny},
  grid=major, grid style={LightGray, line width=0.3pt},
  axis line style={InkGray, line width=0.5pt},
  enlarge x limits={abs=0.8},
  legend style={font=\scriptsize, at={(0.5,-0.20)}, anchor=north,
                legend columns=2, draw=none, fill=none, column sep=6pt},
]
\nextgroupplot[
  ybar, bar width=4pt,
  ymin=-7.00, ymax=72.20,
  xtick={0,1,2,3,4,5.7},
  xticklabels={2021,2022,2023,2024,2025,2026H1},
  ylabel={亿元},
]
\addplot[fill=AgriGreen!75, draw=AgriGreen, bar shift=-2.6pt] table[x=i, y=rev]{data/clean/profiles/fin_600127.dat};
\addplot[fill=SoftRed!60, draw=SoftRed, bar shift=2.6pt] table[x=i, y=ni]{data/clean/profiles/fin_600127.dat};
\addplot[fill=AgriGreen!30, draw=AgriGreen, pattern=north east lines, bar shift=-2.6pt] table[x=x, y=rev]{data/clean/profiles/finh_600127.dat};
\addplot[fill=SoftRed!20, draw=SoftRed, pattern=north east lines, bar shift=2.6pt] table[x=x, y=ni]{data/clean/profiles/finh_600127.dat};
\legend{营业收入（年/半年）, 归母净利润（年/半年）}
\nextgroupplot[
  ymin=-13.0, ymax=70.7,
  xtick={0,1,2,3,4,5.7},
  xticklabels={2021,2022,2023,2024,2025,2026H1},
  ylabel={\%},
]
\addplot[AgriGold, mark=*, mark size=1.3pt] table[x=i, y=roe]{data/clean/profiles/fin_600127.dat};
\addplot[OilBlue, mark=square*, mark size=1.3pt] table[x=i, y=debt]{data/clean/profiles/fin_600127.dat};
\addplot[only marks, AgriGold, mark=*, mark size=2.0pt] table[x=x, y=roe]{data/clean/profiles/finh_600127.dat};
\addplot[only marks, OilBlue, mark=square*, mark size=2.0pt] table[x=x, y=debt]{data/clean/profiles/finh_600127.dat};
\legend{ROE, 资产负债率}
\end{groupplot}
\end{tikzpicture}
\begin{tikzpicture}
\begin{axis}[
  name=finbot,
  width=13.75cm, height=3.10cm, scale only axis,
  ymin=-0.95, ymax=10.64,
  xtick={0,1,2,3,4,5.7},
  xticklabels={2021,2022,2023,2024,2025,2026H1},
  tick label style={font=\tiny},
  yticklabel style={font=\tiny},
  ylabel={亿元}, ylabel style={font=\tiny},
  ybar, bar width=4pt,
  grid=major, grid style={LightGray, line width=0.3pt},
  axis line style={InkGray, line width=0.5pt},
  enlarge x limits={abs=0.8},
  legend style={font=\scriptsize, at={(0.5,-0.26)}, anchor=north,
                legend columns=2, draw=none, fill=none, column sep=10pt},
]
\addplot[fill=AgriGreen!55, draw=AgriGreen, bar shift=-3.0pt] table[x=i, y=cash]{data/clean/profiles/fin_600127.dat};
\addplot[fill=SoftRed!45, draw=SoftRed, bar shift=3.0pt] table[x=i, y=ibd]{data/clean/profiles/fin_600127.dat};
\addplot[fill=AgriGreen!25, draw=AgriGreen, pattern=north east lines, bar shift=-3.0pt] table[x=x, y=cash]{data/clean/profiles/finh_600127.dat};
\addplot[fill=SoftRed!20, draw=SoftRed, pattern=north east lines, bar shift=3.0pt] table[x=x, y=ibd]{data/clean/profiles/finh_600127.dat};
\legend{货币资金（左轴）, 有息负债（左轴）}
\end{axis}
\begin{axis}[
  at={(finbot.south east)}, anchor=south east,
  width=13.75cm, height=3.10cm, scale only axis,
  axis y line*=right, axis x line*=none, xtick=\empty,
  ymin=-1.18, ymax=3.23,
  tick label style={font=\tiny},
  yticklabel style={font=\tiny}, scaled y ticks=false,
  legend style={font=\scriptsize, at={(0.5,-0.44)}, anchor=north,
                legend columns=1, draw=none, fill=none},
]
\addplot[OilBlue, mark=*, mark size=2.2pt, line width=1.3pt] table[x=i, y=ocf]{data/clean/profiles/fin_600127.dat};
\addplot[only marks, OilBlue, mark=*, mark size=3.2pt] table[x=x, y=ocf]{data/clean/profiles/finh_600127.dat};
\legend{经营活动现金流净额（右轴，亿元，蓝点）}
\end{axis}
\end{tikzpicture}

\captionof{figure}[金健米业（600127）近年主要财务指标]{金健米业（600127）近年主要财务指标（左上：营业收入与归母净利润；右上：ROE 与资产负债率；下：货币资金、有息负债为左轴柱状，经营活动现金流净额为其右轴的蓝点折线。
2021---2025 年为年报口径，2026H1 为 2026 年半年报/半年末，与其前一年报之间留出间隔、以斜纹或空心点区分；半年报数据未年化）}
\end{minipage}
\end{center}

\pfl{财务分析} 2026 年 6 月末资产负债率 62.2\%，较 2025 年末上升 0.9 个百分点（样本中位数 52.2\%，所处三级行业内按由低到高排第\zhnumber{7}位）。 2026 年 6 月末货币资金 1.50 亿元、短期债务（短期借款与一年内到期非流动负债）7.28 亿元、长期债务 0.03 亿元，有息负债合计 7.31 亿元，现金短债比 0.21 倍——缺口明显。 净有息负债 5.81 亿元，相当于其 2026 年上半年营业收入的 30\%。 流动比率 0.97、速动比率 0.47，短期偿付压力较大。 2026 年上半年经营活动现金流净额 0.38 亿元，净利润为负而经营现金流为正，差额主要来自折旧摊销与营运资本变动，说明亏损尚未完全转化为现金流出。 综合看，账面货币资金对有息负债与短期债务的覆盖不足，滚动偿付对新增融资的依赖度较高；资产负债率高于样本中位数 10.0 个百分点；有息负债中短期债务占比更高，期限结构仍以滚动续作为主。 公司披露的监管与合规风险为：2025 年 5 月，收到上海证券交易所金健米业股份有限公司 2024 年年度报告的信息披露监管问询函。

\pfl{重大战略转型} 公司的战略动作可归为以下几类：其他动作（2 项）：2021 年 7 月，从根本上避免和消除与上市公司形成同业竞争的可能性；并购与重组（1 项）：2024 年 11 月，公司以资产置换方式优化产业布局：置出从事贸易业务的进出口与农产品贸易子公司。 公告口径上，近三年（2023 年 9 月---2026 年 9 月）公司披露重大重组与并购类公告 4 条、对外投资与转型类公告 2 条、回购与股权激励类公告 0 条，合计公告 362 条。 主营结构的口径变化已经落到报表上：2021 年年报的第一大行业为“粮油食品加工业”（48.1\%），2026 年半年报为“粮油业务”（94.5\%）；公司在这两期的分部划分方式有过调整。 综合判断，报告期内公司有 6 条与战略相关的公告落地（重大重组与并购 4 条、对外投资与转型 2 条），资本运作与主业调整之间存在直接关联。

\begin{center}\small\setlength{\tabcolsep}{4pt}
\captionof{table}[金健米业（600127）历史融资与资本运作]{金健米业（600127）历史融资与资本运作（据公司公告与公开报道整理；金额为披露值，含首次公开发行、再融资、债券、银行借款与重整投资等）}
\begin{adjustbox}{max width=\textwidth}
\begin{tabular}{@{}p{1.45cm}p{2.55cm}p{4.05cm}p{4.95cm}@{}}
\toprule
时间 & 方式 & 规模 & 用途与说明 \\
\midrule
1998-04 & IPO（上海证券交易所） & 发行公众股 5,\allowbreak{}000 万股，\allowbreak{}发行价 5.30 元/\allowbreak{}股 & 含公司职工股 500 万股；发行方式为全额预缴款、\allowbreak{}比例配售、\allowbreak{}余款即退；首发前总股本 10,\allowbreak{}000 万股、\allowbreak{}首发后 15,\allowbreak{}000 万股；招股公告日 1998-04-06 \\
2000-12—2001-01 & 公募增发 A 股（上海证券交易所） & 增发 5,\allowbreak{}000 万股，\allowbreak{}发行价格 16.10 元/\allowbreak{}股 & 经中国证监会证监公司字[2000]229 号文批准；网上发行的 3,\allowbreak{}000.00 万股于 2001-02-06 上市流通 \\
2024-11—2024-12 & 与控股股东资产置换（关联交易） & 公司向控股股东支付资产置换差价 8,\allowbreak{}361.76 万元 & 置出湖南金健进出口有限责任公司 100\% 股权、\allowbreak{}金健农产品（湖南）有限公司 100\% 股权、\allowbreak{}金健农产品（营口）有限公司 100\% 股权；置入湖南裕湘食品有限公司 100\% 股权、\allowbreak{}中南粮油食品科学研究院有限公司 82\% 股权；公司原已持有中南粮油食品科学研究院 18\% 股权 \\
\bottomrule
\end{tabular}
\end{adjustbox}
\end{center}

\pfl{历史融投资情况} 从现金流量表看，2025 年公司取得借款收到的现金 5.86 亿元、偿还债务支付的现金 5.44 亿元，筹资活动现金流净额 0.28 亿元（2023 年为取得借款 6.21 亿元、偿还债务 8.86 亿元）。 2026 年上半年筹资活动现金流净额为 -0.53 亿元（取得借款 2.62 亿元、偿还债务 3.10 亿元，未年化），已转为净流出，处于净还债状态。 2025 年分配股利、利润或偿付利息支付的现金合计 0.16 亿元。 公开披露的融资记录共 3 笔，工具类型以 IPO、定向增发为主；金额与用途见下表。 其中规模最大的两笔为：1998 年 4 月，IPO，发行公众股 5,000 万股，发行价 5.30 元/股；2024 年 11 月，与控股股东资产置换，公司向控股股东支付资产置换差价 8,361.76 万元。 股本方面，近三年披露股本变动 3 次；总股本自 2021 年末以来未发生变化（6.42 亿股）。 分红方面，近三年未实施现金分红，在全部样本中属于分红能力偏弱的一端。 综合看，公司融资结构上股权与债务融资并举；2025 年筹资活动现金流净额 0.28 亿元，年度内仍为净融入，与前述经营与投资现金流的方向基本一致。

\pfl{未来潜在融资需求分析} 2026 年 6 月末货币资金 1.50 亿元，而短期债务 7.28 亿元（现金短债比 0.21 倍），2026 年上半年经营现金流净额 0.38 亿元。按“短期债务减去账面货币资金”的滚动偿付口径，静态缺口约 5.78 亿元；若再计入上半年亏损的年化额（0.19 亿元），维持一年正常经营所需的外部资金约 5.78---5.97 亿元。 公司 2026 年上半年归母净利润为负（-0.10 亿元），但 6 月末资产负债率 62.2\% 尚未进入第一档（亏损且负债率 $>$65\%）的区间，当前融资需求以补充流动资金、支撑低谷期经营性支出为主。 综合看，公司的外部资金需求以补充流动性与项目投入为主，滚动偿付口径下的静态缺口约 5.78 亿元，融资节奏将受制于银行续贷意愿与股权融资窗口。




\subsubsection{冠农股份（600251）}
\label{co:600251}

\pfl{公司基本情况} 冠农股份成立于 2003 年，新疆生产建设兵团第二师国资背景的农产品精深加工企业，产业涉及棉花加工、番茄制品、甜菜制糖、棉籽蛋白与生物发酵饲料；核心利润来自对国投罗钾（20.3\%）和开都河水电（25.28\%）的权益法投资收益，主业加工板块盈利能力较弱。 公司于 2003-06-09 在上交所首发上市，注册资本 7.77 亿元，总股本 7.77 亿股。 截至 2026 年 10 月 12 日收盘价 10.45 元，对应总市值 81.2 亿元，近一年+13.8\%，总市值在三级行业“果蔬加工”的 5 家公司中按由大到小排第\zhnumber{2}位，近一年涨跌幅高于全样本中位数（-10.4\%）24.2 个百分点。 按 2026 年半年报披露的产品构成，收入占比前三位为油脂产品 38.7\%；棉花产品 34.4\%；果蔬产品 19.7\%。 与 2021 年年报相比，第一大产品口径由“皮棉”（64.8\%）变为“油脂产品”（38.7\%）；两期披露的分类口径有调整，占比差异中同时包含分类因素。

\pfl{股权结构与控制权} 按 2026 年 6 月 30 日的股东名册，前三大股东为新疆冠农集团有限责任公司 持股 41.30\%；新疆绿原鑫融贸易有限公司 持股 3.24\%；张健柏 持股 1.44\%，股东户数 32,881 户。 股权结构的另一条线索是，2017 年 12 月，划转完成后绿原国资间接持有公司 40.89\% 股份。

\pfl{经营要点} 经营层面，2026 年上半年，收入结构为油脂产品 5.37 亿元、棉花产品 4.78 亿元、果蔬产品 2.74 亿元、糖业产品 0.46 亿元。

\pfl{所处行业的基本面} 果蔬加工的价格由产地原料果价格与出口订单决定，需求端受海外需求与关税影响（第 \ref{sec:framework} 章的供需框架）。 当期该行业的关键数据是：2026 年 6 月，全国棉花实播面积 4632 万亩、同比减少 145.8 万亩（减幅 3.1\%），其中新疆实播 4153.5 万亩、同比减少 3.4\%。

\pfl{竞争优势} 把该公司放回同类公司比较，2026 年上半年营业收入 13.9 亿元、在三级行业“果蔬加工”的 5 家公司中排第\zhnumber{1}位；毛利率 8.6\%、与行业内中位数基本持平；半年 ROE 13.3\%（未年化）、高于行业内中位数（2.26\%）；资产负债率 30.8\%、高于行业内中位数（30.8\%）。 相对优势方面，公司持有国投新疆罗布泊钾盐有限责任公司 20.3\% 和开都河水电 25.28\% 股权并采用权益法核算。 优势在当期体现为高于行业中位数的资本回报，但优势能否持续仍取决于周期与自身产能投放节奏。

\pfl{历史股价} 本报告样本区间（2021 年 1 月 至 2026 年 10 月）内，股价由 5.80 元变为 10.45 元，累计 +80.2\%，区间最高 12.95 元（2026 年 4 月）、最低 5.80 元（2021 年 1 月）；当前价相当于区间最高价的 80.7\%，为区间最低价的 1.80 倍。 月度收益的年化波动率 34.8\%，低于全样本中位数 37.3\%，在三级行业“果蔬加工”的 5 家公司中波动率由高到低排第\zhnumber{5}位。

\begin{center}
\begin{minipage}{\textwidth}\centering
\begin{tikzpicture}
\begin{axis}[
  width=12.6cm, height=3.9cm, scale only axis,
  xmin=2020.922, xmax=2026.972, ymin=5.39, ymax=13.86,
  xtick={2021,2022,2023,2024,2025,2026},
  yearticks,
  yticklabel style={font=\scriptsize},
  ylabel={元/股}, ylabel style={font=\scriptsize},
  grid=major, grid style={LightGray, line width=0.3pt},
  axis line style={InkGray, line width=0.5pt},
  every axis plot/.append style={line width=0.9pt},
]
\addplot[AgriGreen] table[x=x,y=y]{data/clean/profiles/px_600251.dat};
\addplot[SoftRed, dashed, line width=0.7pt] table[x=x,y=y]{data/clean/profiles/pxzz_600251.dat};
\addplot[only marks, mark=*, mark size=1.1pt, SoftRed] table[x=x,y=y]{data/clean/profiles/pxzz_600251.dat};
\end{axis}
\end{tikzpicture}

\captionof{figure}[冠农股份（600251）股价与周期骨架]{冠农股份（600251）月度收盘价（元/股）与 ZigZag 周期骨架（阈值 25\%，圆点为周期转折点）}
\end{minipage}
\end{center}

\pfl{过去周期分析} 以 25\% 的 ZigZag 阈值划分，样本区间内股价形成 4 个主升段与 3 个主跌段。 最大主升段出现在 2024 年 8 月 至 2026 年 4 月，20 个月+93.0\%；最大主跌段为 2022 年 7 月 至 2024 年 8 月，25 个月-41.1\%。 最近一次转折出现在 2026 年 8 月（上涨段自 2026 年 6 月起，2 个月+32.3\%），当前处于该段之后的震荡之中。 公司股价较申万农林牧渔指数提前约 64 个月见底，是板块中较早定价周期反转的公司之一。 果蔬加工的外向型特征明显，出口订单与原料产地价格决定利润率，公司 2026 年上半年实现盈利、半年 ROE 13.3\%（未年化），好于全样本中位数（0.75\%），下行周期对其报表的冲击小于同业。 综合区间形态看，该股单段涨跌幅度偏大、以趋势段为主（最大主升 +93.0\%、最大主跌 -41.1\%），年化波动率 34.8\% 与全样本中位数 37.3\% 相比波动与样本中位数接近。

\pfl{盈利情况} 2026 年上半年营业收入 13.89 亿元（同比 +0.3\%），归母净利润 5.37 亿元，同比 +79.1\%。 以年报为趋势对照，2023---2025 年营业收入由 43.48 亿元变为 25.15 亿元（两年年均复合增速 -23.9\%），归母净利润由 7.20 亿元变为 3.48 亿元。 按半年报口径，2026 年上半年的 ROE 为 13.3\%（未年化）、销售净利率 39.5\%、毛利率 8.6\%，ROE 在“果蔬加工”5 家公司中按数值由高到低排第\zhnumber{1}位。 按同一口径，三级行业“果蔬加工”5 家公司的半年 ROE 中位数为 2.26\%，该公司高于其中位数 11.08 个百分点。 综合看，公司在报告期维持盈利（高于全样本半年 ROE 中位数 0.75\%），利润的可持续性取决于果蔬加工景气的延续与成本管控的执行。

\begin{center}\small\setlength{\tabcolsep}{3.4pt}
\captionof{table}[冠农股份（600251）关键财务指标]{冠农股份（600251）关键财务指标（合并报表口径；两个半年报期均未年化；现金短债比 $=$ 货币资金 $\div$ 短期债务）}
\begin{adjustbox}{max width=\textwidth}
\begin{tabular}{@{}lrrrrr@{}}
\toprule
指标 & 2023 年 & 2024 年 & 2025 年 & 2025 年半年报 & 2026 年半年报 \\
\midrule
营业收入（亿元） & 43.48 & 35.13 & 25.15 & 13.86 & 13.89 \\
归母净利润（亿元） & 7.20 & 2.05 & 3.48 & 3.00 & 5.37 \\
毛利率（\%） & 16.8 & 8.1 & 8.8 & 7.0 & 8.6 \\
ROE（\%） & 20.4 & 5.5 & 9.4 & 8.0 & 13.3 \\
资产负债率（\%） & 48.4 & 46.3 & 46.9 & 34.2 & 30.8 \\
经营活动现金流净额（亿元） & 2.05 & 6.13 & 1.26 & 8.36 & 5.12 \\
货币资金（亿元） & 18.77 & 16.92 & 10.80 & 18.78 & 13.50 \\
有息负债合计（亿元） & 26.81 & 23.42 & 22.77 & 14.86 & 11.87 \\
现金短债比（倍） & 0.76 & 0.78 & 0.51 & 1.43 & 1.19 \\
\bottomrule
\end{tabular}
\end{adjustbox}
\end{center}

\begin{center}
\begin{minipage}{\textwidth}\centering
\begin{tikzpicture}
\begin{groupplot}[
  group style={group size=2 by 1, horizontal sep=0.60cm},
  width=6.85cm, height=4.60cm,
  tick label style={font=\tiny},
  title style={font=\scriptsize\heiti},
  yticklabel style={font=\tiny},
  ylabel style={font=\tiny},
  grid=major, grid style={LightGray, line width=0.3pt},
  axis line style={InkGray, line width=0.5pt},
  enlarge x limits={abs=0.8},
  legend style={font=\scriptsize, at={(0.5,-0.20)}, anchor=north,
                legend columns=2, draw=none, fill=none, column sep=6pt},
]
\nextgroupplot[
  ybar, bar width=4pt,
  ymin=-4.35, ymax=48.69,
  xtick={0,1,2,3,4,5.7},
  xticklabels={2021,2022,2023,2024,2025,2026H1},
  ylabel={亿元},
]
\addplot[fill=AgriGreen!75, draw=AgriGreen, bar shift=-2.6pt] table[x=i, y=rev]{data/clean/profiles/fin_600251.dat};
\addplot[fill=SoftRed!60, draw=SoftRed, bar shift=2.6pt] table[x=i, y=ni]{data/clean/profiles/fin_600251.dat};
\addplot[fill=AgriGreen!30, draw=AgriGreen, pattern=north east lines, bar shift=-2.6pt] table[x=x, y=rev]{data/clean/profiles/finh_600251.dat};
\addplot[fill=SoftRed!20, draw=SoftRed, pattern=north east lines, bar shift=2.6pt] table[x=x, y=ni]{data/clean/profiles/finh_600251.dat};
\legend{营业收入（年/半年）, 归母净利润（年/半年）}
\nextgroupplot[
  ymin=-4.5, ymax=61.4,
  xtick={0,1,2,3,4,5.7},
  xticklabels={2021,2022,2023,2024,2025,2026H1},
  ylabel={\%},
]
\addplot[AgriGold, mark=*, mark size=1.3pt] table[x=i, y=roe]{data/clean/profiles/fin_600251.dat};
\addplot[OilBlue, mark=square*, mark size=1.3pt] table[x=i, y=debt]{data/clean/profiles/fin_600251.dat};
\addplot[only marks, AgriGold, mark=*, mark size=2.0pt] table[x=x, y=roe]{data/clean/profiles/finh_600251.dat};
\addplot[only marks, OilBlue, mark=square*, mark size=2.0pt] table[x=x, y=debt]{data/clean/profiles/finh_600251.dat};
\legend{ROE, 资产负债率}
\end{groupplot}
\end{tikzpicture}
\begin{tikzpicture}
\begin{axis}[
  name=finbot,
  width=13.75cm, height=3.10cm, scale only axis,
  ymin=-3.26, ymax=36.47,
  xtick={0,1,2,3,4,5.7},
  xticklabels={2021,2022,2023,2024,2025,2026H1},
  tick label style={font=\tiny},
  yticklabel style={font=\tiny},
  ylabel={亿元}, ylabel style={font=\tiny},
  ybar, bar width=4pt,
  grid=major, grid style={LightGray, line width=0.3pt},
  axis line style={InkGray, line width=0.5pt},
  enlarge x limits={abs=0.8},
  legend style={font=\scriptsize, at={(0.5,-0.26)}, anchor=north,
                legend columns=2, draw=none, fill=none, column sep=10pt},
]
\addplot[fill=AgriGreen!55, draw=AgriGreen, bar shift=-3.0pt] table[x=i, y=cash]{data/clean/profiles/fin_600251.dat};
\addplot[fill=SoftRed!45, draw=SoftRed, bar shift=3.0pt] table[x=i, y=ibd]{data/clean/profiles/fin_600251.dat};
\addplot[fill=AgriGreen!25, draw=AgriGreen, pattern=north east lines, bar shift=-3.0pt] table[x=x, y=cash]{data/clean/profiles/finh_600251.dat};
\addplot[fill=SoftRed!20, draw=SoftRed, pattern=north east lines, bar shift=3.0pt] table[x=x, y=ibd]{data/clean/profiles/finh_600251.dat};
\legend{货币资金（左轴）, 有息负债（左轴）}
\end{axis}
\begin{axis}[
  at={(finbot.south east)}, anchor=south east,
  width=13.75cm, height=3.10cm, scale only axis,
  axis y line*=right, axis x line*=none, xtick=\empty,
  ymin=-2.66, ymax=13.32,
  tick label style={font=\tiny},
  yticklabel style={font=\tiny}, scaled y ticks=false,
  legend style={font=\scriptsize, at={(0.5,-0.44)}, anchor=north,
                legend columns=1, draw=none, fill=none},
]
\addplot[OilBlue, mark=*, mark size=2.2pt, line width=1.3pt] table[x=i, y=ocf]{data/clean/profiles/fin_600251.dat};
\addplot[only marks, OilBlue, mark=*, mark size=3.2pt] table[x=x, y=ocf]{data/clean/profiles/finh_600251.dat};
\legend{经营活动现金流净额（右轴，亿元，蓝点）}
\end{axis}
\end{tikzpicture}

\captionof{figure}[冠农股份（600251）近年主要财务指标]{冠农股份（600251）近年主要财务指标（左上：营业收入与归母净利润；右上：ROE 与资产负债率；下：货币资金、有息负债为左轴柱状，经营活动现金流净额为其右轴的蓝点折线。
2021---2025 年为年报口径，2026H1 为 2026 年半年报/半年末，与其前一年报之间留出间隔、以斜纹或空心点区分；半年报数据未年化）}
\end{minipage}
\end{center}

\pfl{财务分析} 2026 年 6 月末资产负债率 30.8\%，较 2025 年末下降 16.1 个百分点（样本中位数 52.2\%，所处三级行业内按由低到高排第\zhnumber{3}位）。 2026 年 6 月末货币资金 13.50 亿元、短期债务（短期借款与一年内到期非流动负债）11.36 亿元、长期债务 0.50 亿元，有息负债合计 11.87 亿元，现金短债比 1.19 倍——覆盖充分。 净有息负债 -1.63 亿元，相当于其 2026 年上半年营业收入的 -12\%。 流动比率 1.92、速动比率 1.54，短期偿债能力处于正常区间。 2026 年上半年经营活动现金流净额 5.12 亿元，经营现金流与利润同向为正，内生资金可以支撑运营。 综合看，现金短债比 1.19 倍，短期偿付压力可控；资产负债率低于样本中位数 21.4 个百分点；有息负债中短期债务占比更高，期限结构仍以滚动续作为主。

\pfl{重大战略转型} 近三年公司的战略动作集中在：产能与项目（2 项）：2024 年，公司承诺在 36 个月内继续推进绿洲农业将新疆新建番茄制品有限公司 100\% 股权或控股权对外出让；2023 年，公司参股国投罗钾，增资扩产；融资与资本运作（1 项）：2021 年，筹划公开发行可转债以支持主业扩产。 公告口径上，近三年（2023 年 9 月---2026 年 9 月）公司披露重大重组与并购类公告 1 条、对外投资与转型类公告 0 条、回购与股权激励类公告 0 条，合计公告 346 条。 主营结构的口径变化已经落到报表上：2021 年年报的第一大产品为“皮棉”（64.8\%），2026 年半年报为“油脂产品”（38.7\%）；公司在这两期的分部划分方式有过调整。 综合判断，报告期内公司有 1 条与战略相关的公告落地（重大重组与并购 1 条、对外投资与转型 0 条），资本运作与主业调整之间存在直接关联。

\begin{center}\small\setlength{\tabcolsep}{4pt}
\captionof{table}[冠农股份（600251）历史融资与资本运作]{冠农股份（600251）历史融资与资本运作（据公司公告与公开报道整理；金额为披露值，含首次公开发行、再融资、债券、银行借款与重整投资等）}
\begin{adjustbox}{max width=\textwidth}
\begin{tabular}{@{}p{1.45cm}p{2.55cm}p{4.05cm}p{4.95cm}@{}}
\toprule
时间 & 方式 & 规模 & 用途与说明 \\
\midrule
2003-05 & IPO（上海证券交易所） & 发行 4,\allowbreak{}000 万股，\allowbreak{}发行价 5.75 元/\allowbreak{}股 & 经中国证监会证监发行字[2003]46 号核准，\allowbreak{}采用向二级市场投资者定价配售方式 \\
2017-12 & 国有股权无偿划转（间接控股股东变更） & 冠源投资 100\% 国有股权无偿划转 & 兵团国资体系内部整合；第二师国资委将所持冠源投资 100\% 股权无偿划转至绿原国资 \\
2021—2023 & 拟公开发行可转换公司债券（最终终止） & 拟募集资金总额 11.00 亿元 & 公司公开发行可转换公司债券募投项目（公告未在本底稿核实的页面中列明具体项目）；2021 年第七次临时股东大会等审议通过发行方案；2022-01-04 向中国证监会申报；2023-03-21 董事会审议通过终止议案 \\
2023-12 & 参股公司增资（对外投资） & 投资金额 7,\allowbreak{}202.97 万元 & 对参股公司国投新疆罗布泊钾盐有限责任公司增资；不构成重大资产重组，\allowbreak{}不构成关联交易 \\
2025 & 参股公司现金分红（投资收益回收） & 收到国投罗钾 2024 年度现金分红款 4.12 亿元 & 为公司生产经营和项目投资提供资金保障；公司持有国投罗钾 20.3\% 股权 \\
\bottomrule
\end{tabular}
\end{adjustbox}
\end{center}

\pfl{历史融投资情况} 从现金流量表看，2025 年公司取得借款收到的现金 29.56 亿元、偿还债务支付的现金 30.76 亿元，筹资活动现金流净额 -3.87 亿元（2023 年为取得借款 31.99 亿元、偿还债务 29.70 亿元）。 2026 年上半年筹资活动现金流净额为 -11.60 亿元（取得借款 6.47 亿元、偿还债务 17.81 亿元，未年化），已转为净流出，处于净还债状态。 2025 年分配股利、利润或偿付利息支付的现金合计 2.42 亿元。 公开披露的融资记录共 5 笔，工具类型以 IPO、可转债为主；金额与用途见下表。 其中规模最大的两笔为：2025 年，参股公司现金分红，收到国投罗钾 2024 年度现金分红款 4.12 亿元；2023 年 12 月，参股公司增资，投资金额 7,202.97 万元。 股本方面，近三年披露股本变动 5 次（股份回购 2 次）；总股本由 2021 年末的 7.81 亿股变为 7.77 亿股（-0.57\%）。 分红方面，近三年现金分红 5 次、合计每股派现 0.88 元，最近一次实施在 2025 年 12 月。 综合看，公司融资结构上股权与债务融资并举；2025 年筹资活动现金流净额 -3.87 亿元，融资节奏仍以维持存量债务周转为主，与前述经营与投资现金流的方向基本一致。

\pfl{未来潜在融资需求分析} 2026 年 6 月末货币资金 13.50 亿元，而短期债务 11.36 亿元（现金短债比 1.19 倍），2026 年上半年经营现金流净额 5.12 亿元。账面货币资金可以覆盖短期债务，缺口不体现在静态偿付层面。 按第 \ref{sec:financing-need} 节的三档划分，公司属于第三档“扩张型”：2026 年上半年 ROE 13.3\%（未年化，折合约 26.7\%）、6 月末资产负债率 30.8\%，资产负债表仍具备加杠杆空间；融资用途以产能扩张、品类并购与海外布局为主，单笔规模通常在 5---20 亿元量级，会被市场定价为成长而非补血。 公开披露的后续资金安排线索包括：债务与授信安排：2023 年 3 月，公司终止向不特定对象发行可转换公司债券（原拟募集资金 11.00 亿元）；2026 年半年度报告，期后事项包括 2026-06-18 披露 2025 年年度权益分派实施公告、2026-06-27 披露投资者关系活动记录表、为控股子公司提供担保的进展公告；资本开支安排：2025 年 6 月，公司收到国投罗钾 2024 年度现金分红款 4.12 亿元“为公司生产经营和项目投资提供了更加充足的资金保障”。 以现有货币资金与经营现金流衡量，公司不存在刚性补血的紧迫性，融资动作更多是主动型的。




\subsubsection{深粮控股（000019）}
\label{co:000019}

\pfl{公司基本情况} 1991 年设立至今，深粮控股是深圳市属国有大型上市粮食企业（深农投旗下），承担深圳市超百万吨粮油储备任务，是深圳“米袋子”和粤港澳大湾区粮油供应保障重要载体；前身为 1992 年上市的原深圳市深宝实业（深深宝 A），2018 年发行股份购买深粮集团 100\% 股权完成重组转型。 公司于 1992-10-12 在深交所主板首发上市，注册资本 11.53 亿元，总股本 11.53 亿股。 截至 2026 年 10 月 12 日收盘价 7.00 元，对应总市值 80.7 亿元，近一年+3.7\%，总市值在三级行业“粮油加工”的 7 家公司中按由大到小排第\zhnumber{3}位，近一年涨跌幅高于全样本中位数（-10.4\%）14.1 个百分点。 按 2026 年半年报披露的行业构成，收入占比前三位为批发与零售业 63.3\%；租赁及商业服务业 20.6\%；制造业 16.1\%。

\pfl{股权结构与控制权} 按 2026 年 6 月 30 日的股东名册，前三大股东为深圳农业与食品投资控股集团有限公司 持股 63.79\%；深圳市农产品集团股份有限公司 持股 8.23\%；东莞市果菜副食交易市场有限公司 持股 0.75\%，股东户数 39,173 户。 控制权与股东性质方面，2018 年 4 月，深圳市投资控股有限公司将其持有的公司 7,948.43 万股 A 股股份全部过户给深圳市福德国有资本运营有限公司（福德资本）。

\pfl{经营要点} 经营层面，2026 年上半年，经营活动产生的现金流量净额为-5,288.16 万元、上年同期为 7.37 亿元。

\pfl{所处行业的基本面} 粮油加工的供给由大豆、菜籽等原料进口与压榨产能决定，需求由食用消费决定，利润来自原料与成品的价差而非价格方向（第 \ref{sec:framework} 章的供需框架）。 从供需两端的当期数据看，2026 年，国家继续在稻谷主产区实行最低收购价政策；2026 年上半年，我国稻谷和大米进口量 243 万吨、同比增长 67.7\%，国内稻米市场供给总体充裕、价格呈窄幅偏弱走势。

\pfl{竞争优势} 把该公司放回同类公司比较，2026 年上半年营业收入 23.9 亿元、在三级行业“粮油加工”的 7 家公司中排第\zhnumber{4}位；毛利率 16.2\%、高于行业内中位数 9.9 个百分点；半年 ROE 2.2\%（未年化）、低于行业内中位数（2.22\%）；资产负债率 38.1\%、低于行业内中位数（43.5\%）。 与同类公司相比，公司的差异化在于：公司依托东莞物流、华联公司的港贸一体化建设持续提升供应链协同能力。 优势在报表上主要体现为高于行业平均的毛利率，这在周期底部是抵御亏损的主要缓冲。

\pfl{历史股价} 自 2021 年 1 月至 2026 年 10 月，股价由 6.26 元变为 7.00 元，累计 +11.8\%，区间最高 9.05 元（2022 年 5 月）、最低 5.67 元（2026 年 6 月）；当前价相当于区间最高价的 77.3\%，为区间最低价的 1.23 倍。 月度收益的年化波动率 27.1\%，低于全样本中位数 37.3\%，在三级行业“粮油加工”的 7 家公司中波动率由高到低排第\zhnumber{6}位。

\begin{center}
\begin{minipage}{\textwidth}\centering
\begin{tikzpicture}
\begin{axis}[
  width=12.6cm, height=3.9cm, scale only axis,
  xmin=2020.922, xmax=2026.972, ymin=5.27, ymax=9.68,
  xtick={2021,2022,2023,2024,2025,2026},
  yearticks,
  yticklabel style={font=\scriptsize},
  ylabel={元/股}, ylabel style={font=\scriptsize},
  grid=major, grid style={LightGray, line width=0.3pt},
  axis line style={InkGray, line width=0.5pt},
  every axis plot/.append style={line width=0.9pt},
]
\addplot[AgriGreen] table[x=x,y=y]{data/clean/profiles/px_000019.dat};
\addplot[SoftRed, dashed, line width=0.7pt] table[x=x,y=y]{data/clean/profiles/pxzz_000019.dat};
\addplot[only marks, mark=*, mark size=1.1pt, SoftRed] table[x=x,y=y]{data/clean/profiles/pxzz_000019.dat};
\end{axis}
\end{tikzpicture}

\captionof{figure}[深粮控股（000019）股价与周期骨架]{深粮控股（000019）月度收盘价（元/股）与 ZigZag 周期骨架（阈值 25\%，圆点为周期转折点）}
\end{minipage}
\end{center}

\pfl{过去周期分析} 以 25\% 的 ZigZag 阈值划分，样本区间内股价形成 3 个主升段与 2 个主跌段。 最大主跌段为 2022 年 5 月 至 2024 年 6 月，25 个月-36.6\%；最大主升段出现在 2021 年 1 月 至 2022 年 5 月，16 个月+44.6\%。 最近一次转折出现在 2026 年 8 月（上涨段自 2026 年 6 月起，2 个月+27.0\%），当前处于该段之后的震荡之中。 公司股价与申万农林牧渔指数几乎同步见底，个股并未走出独立行情。 粮油加工是“成本加成”型生意（第 \ref{sec:grain} 章），利润来自原料与成品的价差及规模效率，收入规模大而净利率薄，公司 2026 年上半年实现盈利、半年 ROE 2.2\%（未年化），好于全样本中位数（0.75\%），下行周期对其报表的冲击小于同业。 综合区间形态看，该股单段涨跌幅度相对温和、以震荡为主（最大主升 +44.6\%、最大主跌 -36.6\%），年化波动率 27.1\% 与全样本中位数 37.3\% 相比波动小于样本中位数。

\pfl{盈利情况} 2026 年上半年营业收入 23.88 亿元（同比 +0.1\%），归母净利润 1.12 亿元，同比 -36.4\%。 以年报为趋势对照，2023---2025 年营业收入由 61.90 亿元变为 55.05 亿元（两年年均复合增速 -5.7\%），归母净利润由 3.48 亿元变为 2.43 亿元。 按半年报口径，2026 年上半年的 ROE 为 2.2\%（未年化）、销售净利率 4.7\%、毛利率 16.2\%，ROE 在“粮油加工”7 家公司中按数值由高到低排第\zhnumber{4}位。 按同一口径，三级行业“粮油加工”7 家公司的半年 ROE 中位数为 2.22\%，该公司低于其中位数 0.00 个百分点。 综合看，公司在报告期维持盈利（高于全样本半年 ROE 中位数 0.75\%），利润的可持续性取决于粮油加工景气的延续与成本管控的执行。

\begin{center}\small\setlength{\tabcolsep}{3.4pt}
\captionof{table}[深粮控股（000019）关键财务指标]{深粮控股（000019）关键财务指标（合并报表口径；两个半年报期均未年化；现金短债比 $=$ 货币资金 $\div$ 短期债务）}
\begin{adjustbox}{max width=\textwidth}
\begin{tabular}{@{}lrrrrr@{}}
\toprule
指标 & 2023 年 & 2024 年 & 2025 年 & 2025 年半年报 & 2026 年半年报 \\
\midrule
营业收入（亿元） & 61.90 & 53.75 & 55.05 & 23.84 & 23.88 \\
归母净利润（亿元） & 3.48 & 3.25 & 2.43 & 1.76 & 1.12 \\
毛利率（\%） & 15.5 & 16.7 & 16.0 & 18.5 & 16.2 \\
ROE（\%） & 7.2 & 6.7 & 4.9 & 3.5 & 2.2 \\
资产负债率（\%） & 34.1 & 36.2 & 32.3 & 30.2 & 38.1 \\
经营活动现金流净额（亿元） & 5.86 & -0.15 & 11.00 & 7.37 & -0.53 \\
货币资金（亿元） & 2.41 & 1.68 & 0.74 & 1.80 & 0.87 \\
有息负债合计（亿元） & 12.84 & 16.06 & 12.12 & 11.86 & 21.43 \\
现金短债比（倍） & 0.19 & 0.11 & 0.06 & 0.16 & 0.05 \\
\bottomrule
\end{tabular}
\end{adjustbox}
\end{center}

\begin{center}
\begin{minipage}{\textwidth}\centering
\begin{tikzpicture}
\begin{groupplot}[
  group style={group size=2 by 1, horizontal sep=0.60cm},
  width=6.85cm, height=4.60cm,
  tick label style={font=\tiny},
  title style={font=\scriptsize\heiti},
  yticklabel style={font=\tiny},
  ylabel style={font=\tiny},
  grid=major, grid style={LightGray, line width=0.3pt},
  axis line style={InkGray, line width=0.5pt},
  enlarge x limits={abs=0.8},
  legend style={font=\scriptsize, at={(0.5,-0.20)}, anchor=north,
                legend columns=2, draw=none, fill=none, column sep=6pt},
]
\nextgroupplot[
  ybar, bar width=4pt,
  ymin=-10.14, ymax=113.56,
  xtick={0,1,2,3,4,5.7},
  xticklabels={2021,2022,2023,2024,2025,2026H1},
  ylabel={亿元},
]
\addplot[fill=AgriGreen!75, draw=AgriGreen, bar shift=-2.6pt] table[x=i, y=rev]{data/clean/profiles/fin_000019.dat};
\addplot[fill=SoftRed!60, draw=SoftRed, bar shift=2.6pt] table[x=i, y=ni]{data/clean/profiles/fin_000019.dat};
\addplot[fill=AgriGreen!30, draw=AgriGreen, pattern=north east lines, bar shift=-2.6pt] table[x=x, y=rev]{data/clean/profiles/finh_000019.dat};
\addplot[fill=SoftRed!20, draw=SoftRed, pattern=north east lines, bar shift=2.6pt] table[x=x, y=ni]{data/clean/profiles/finh_000019.dat};
\legend{营业收入（年/半年）, 归母净利润（年/半年）}
\nextgroupplot[
  ymin=-3.1, ymax=42.6,
  xtick={0,1,2,3,4,5.7},
  xticklabels={2021,2022,2023,2024,2025,2026H1},
  ylabel={\%},
]
\addplot[AgriGold, mark=*, mark size=1.3pt] table[x=i, y=roe]{data/clean/profiles/fin_000019.dat};
\addplot[OilBlue, mark=square*, mark size=1.3pt] table[x=i, y=debt]{data/clean/profiles/fin_000019.dat};
\addplot[only marks, AgriGold, mark=*, mark size=2.0pt] table[x=x, y=roe]{data/clean/profiles/finh_000019.dat};
\addplot[only marks, OilBlue, mark=square*, mark size=2.0pt] table[x=x, y=debt]{data/clean/profiles/finh_000019.dat};
\legend{ROE, 资产负债率}
\end{groupplot}
\end{tikzpicture}
\begin{tikzpicture}
\begin{axis}[
  name=finbot,
  width=13.75cm, height=3.10cm, scale only axis,
  ymin=-2.14, ymax=24.01,
  xtick={0,1,2,3,4,5.7},
  xticklabels={2021,2022,2023,2024,2025,2026H1},
  tick label style={font=\tiny},
  yticklabel style={font=\tiny},
  ylabel={亿元}, ylabel style={font=\tiny},
  ybar, bar width=4pt,
  grid=major, grid style={LightGray, line width=0.3pt},
  axis line style={InkGray, line width=0.5pt},
  enlarge x limits={abs=0.8},
  legend style={font=\scriptsize, at={(0.5,-0.26)}, anchor=north,
                legend columns=2, draw=none, fill=none, column sep=10pt},
]
\addplot[fill=AgriGreen!55, draw=AgriGreen, bar shift=-3.0pt] table[x=i, y=cash]{data/clean/profiles/fin_000019.dat};
\addplot[fill=SoftRed!45, draw=SoftRed, bar shift=3.0pt] table[x=i, y=ibd]{data/clean/profiles/fin_000019.dat};
\addplot[fill=AgriGreen!25, draw=AgriGreen, pattern=north east lines, bar shift=-3.0pt] table[x=x, y=cash]{data/clean/profiles/finh_000019.dat};
\addplot[fill=SoftRed!20, draw=SoftRed, pattern=north east lines, bar shift=3.0pt] table[x=x, y=ibd]{data/clean/profiles/finh_000019.dat};
\legend{货币资金（左轴）, 有息负债（左轴）}
\end{axis}
\begin{axis}[
  at={(finbot.south east)}, anchor=south east,
  width=13.75cm, height=3.10cm, scale only axis,
  axis y line*=right, axis x line*=none, xtick=\empty,
  ymin=-3.53, ymax=14.46,
  tick label style={font=\tiny},
  yticklabel style={font=\tiny}, scaled y ticks=false,
  legend style={font=\scriptsize, at={(0.5,-0.44)}, anchor=north,
                legend columns=1, draw=none, fill=none},
]
\addplot[OilBlue, mark=*, mark size=2.2pt, line width=1.3pt] table[x=i, y=ocf]{data/clean/profiles/fin_000019.dat};
\addplot[only marks, OilBlue, mark=*, mark size=3.2pt] table[x=x, y=ocf]{data/clean/profiles/finh_000019.dat};
\legend{经营活动现金流净额（右轴，亿元，蓝点）}
\end{axis}
\end{tikzpicture}

\captionof{figure}[深粮控股（000019）近年主要财务指标]{深粮控股（000019）近年主要财务指标（左上：营业收入与归母净利润；右上：ROE 与资产负债率；下：货币资金、有息负债为左轴柱状，经营活动现金流净额为其右轴的蓝点折线。
2021---2025 年为年报口径，2026H1 为 2026 年半年报/半年末，与其前一年报之间留出间隔、以斜纹或空心点区分；半年报数据未年化）}
\end{minipage}
\end{center}

\pfl{财务分析} 2026 年 6 月末资产负债率 38.1\%，较 2025 年末上升 5.8 个百分点（样本中位数 52.2\%，所处三级行业内按由低到高排第\zhnumber{3}位）。 2026 年 6 月末货币资金 0.87 亿元、短期债务（短期借款与一年内到期非流动负债）18.04 亿元、长期债务 3.39 亿元，有息负债合计 21.43 亿元，现金短债比 0.05 倍——缺口明显。 净有息负债 20.57 亿元，相当于其 2026 年上半年营业收入的 86\%。 流动比率 1.81、速动比率 0.18，短期流动性无明显缺口。 2026 年上半年经营活动现金流净额 -0.53 亿元，净利润为正但经营现金流为负，利润的现金含量偏低。 2026 年 6 月末在建工程 0.71 亿元，较上年末增加 0.19 亿元，仍有在建投入。 综合看，账面货币资金对有息负债与短期债务的覆盖不足，滚动偿付对新增融资的依赖度较高；资产负债率低于样本中位数 14.1 个百分点；有息负债中短期债务占比更高，期限结构仍以滚动续作为主。

\pfl{重大战略转型} 从公开披露看，公司的战略推进集中在：其他动作（2 项）：2025 年上半年，归属于上市公司股东的净利润 1.76 亿元、同比增长 36.44\%；并购与重组（1 项）：2018---2019 年，重组后由茶及天然植物精深加工转型为深圳市属国有粮食储备与粮油流通平台。 公告口径上，近三年（2023 年 9 月---2026 年 9 月）公司披露重大重组与并购类公告 0 条、对外投资与转型类公告 0 条、回购与股权激励类公告 0 条，合计公告 235 条。 第一大行业“批发与零售业”的收入占比由 2021 年年报的 82.5\% 变为 63.3\%（19.2 个百分点），收入结构出现再平衡。 综合判断，公司报告期内的变化主要发生在收入结构层面，属于口径与业务重心的再平衡，而非外延式扩张。

\begin{center}\small\setlength{\tabcolsep}{4pt}
\captionof{table}[深粮控股（000019）历史融资与资本运作]{深粮控股（000019）历史融资与资本运作（据公司公告与公开报道整理；金额为披露值，含首次公开发行、再融资、债券、银行借款与重整投资等）}
\begin{adjustbox}{max width=\textwidth}
\begin{tabular}{@{}p{1.45cm}p{2.55cm}p{4.05cm}p{4.95cm}@{}}
\toprule
时间 & 方式 & 规模 & 用途与说明 \\
\midrule
2018-03-23 & 签署重组协议 & 交易双方初步商定标的资产深粮集团 100\% 股权交易作价为 58.59 亿元 & 本次交易完成后深深宝 A 将持有深粮集团 100\% 股权；2018 年 9 月 25 日重大资产重组获证监会审核通过后复牌 \\
2018-11-12 & 发行股份购买资产（重大资产重组） & 向福德资本发行 6.56 亿股 A 股购买其持有的深粮集团 100\% 股权；根据中企华评报字（2018）第 3558 号资产评估报告 & 注入深圳市粮食集团有限公司（深粮集团）100\% 股权 \\
2026-01-22 & 银行综合授信额度计划（董事会决议） & 2026 年度向银行申请综合授信额度不超过人民币 147.50 亿元；其中公司和全资子公司深圳市粮食集团有限公司拟向 21 家银行申请综合授信总额不超过人民币 143.00 亿元 & 日常经营与粮油储备业务资金需求；同日董事会同意公司（含下属企业）2026 年度通过银行流动资金借款、\allowbreak{}法人透支、\allowbreak{}国内信用证、\allowbreak{}银行承兑汇票及商业承兑汇票代理贴现等融资工具对外融资最高余额不超过人民币 20.00 亿元 \\
\bottomrule
\end{tabular}
\end{adjustbox}
\end{center}

\pfl{历史融投资情况} 从现金流量表看，2025 年公司取得借款收到的现金 16.76 亿元、偿还债务支付的现金 26.04 亿元、吸收权益性投资 0.00 亿元，筹资活动现金流净额 -11.55 亿元（2023 年为取得借款 19.31 亿元、偿还债务 18.72 亿元）。 2026 年上半年筹资活动现金流净额为 1.27 亿元（取得借款 12.85 亿元、偿还债务 9.70 亿元，未年化），仍处于净融入状态。 2025 年分配股利、利润或偿付利息支付的现金合计 2.01 亿元。 公开披露的融资记录共 3 笔，工具类型以发行股份购买资产、银行授信为主；金额与用途见下表。 其中规模最大的两笔为：2026 年 1 月，银行授信，2026 年度向银行申请综合授信额度不超过人民币 147.50 亿元，其中公司和全资子公司深圳市粮食集团有限公司拟向 21 家银行申请综合授信总额不超过人民币 143.00 亿元；2018 年 11 月，发行股份购买资产，向福德资本发行 6.56 亿股 A 股购买其持有的深粮集团 100\% 股权，根据中企华评报字（2018）第 3558 号资产评估报告。 股本方面，近三年披露股本变动 3 次；总股本自 2021 年末以来未发生变化（11.53 亿股）。 分红方面，近三年现金分红 4 次、合计每股派现 0.72 元，最近一次实施在 2025 年 12 月。 综合看，公司融资结构上以债务滚动为主、股权融资为补充；2025 年筹资活动现金流净额 -11.55 亿元，融资节奏仍以维持存量债务周转为主，与前述经营与投资现金流的方向基本一致。

\pfl{未来潜在融资需求分析} 2026 年 6 月末货币资金 0.87 亿元，而短期债务 18.04 亿元（现金短债比 0.05 倍），2026 年上半年经营现金流净额 -0.53 亿元。按“短期债务减去账面货币资金”的滚动偿付口径，静态缺口约 17.18 亿元；上半年盈利或接近盈亏平衡，该缺口不随经营亏损进一步扩大。 公司 2026 年 6 月末资产负债率 38.1\%、2026 年上半年 ROE 2.2\%（未年化），现金短债比 0.05 倍、短债覆盖不足，融资需求首先用于置换与滚动存量债务、补充营运资金，属于“补血型”而非扩张型。 公开披露的后续资金安排线索包括：债务与授信安排：2026 年度，银行综合授信额度不超过 147.50 亿元；2026 年第一季度，报告披露短期借款 10.89 亿元；资本开支安排：2026 年 6 月，新增资本开支需求。 资金需求还要叠加在建工程：期末在建工程 0.71 亿元、2026 年上半年购建固定资产等支出 0.50 亿元（未年化），这部分支出在周期底部通常会被主动放缓。 综合看，公司的外部资金需求以营运资金周转与在建项目投入为主，滚动偿付口径下的静态缺口约 17.18 亿元，融资的时点与规模更多取决于信用条件与资本市场窗口。




\subsubsection{国投中鲁（600962）}
\label{co:600962}

\pfl{公司基本情况} 国投中鲁成立于 2001 年，国家开发投资集团有限公司控股的浓缩苹果汁加工企业，是国内主板第一家上市的浓缩果汁企业，主业规模与盈利能力有限、母公司长期存在未弥补亏损；2025 年底起筹划以发行股份方式收购中国电子工程设计院 100\% 股份，跨界转型电子信息产业。 公司于 2004-06-22 在上交所首发上市，注册资本 2.62 亿元，总股本 2.62 亿股。 截至 2026 年 10 月 12 日收盘价 23.84 元，对应总市值 62.5 亿元，近一年+11.4\%，总市值在三级行业“果蔬加工”的 5 家公司中按由大到小排第\zhnumber{3}位，近一年涨跌幅高于全样本中位数（-10.4\%）21.8 个百分点。 按 2025 年年报披露的行业构成，收入占比前两位为饮料制造业 99.5\%；其他(补充) 0.5\%。

\pfl{股权结构与控制权} 按 2026 年 6 月 30 日的股东名册，前三大股东为国家开发投资集团有限公司 持股 44.57\%；郭岩 持股 1.23\%；应志江 持股 1.21\%，股东户数 24,104 户。 控制权与股东性质方面，2025 年 7 月，公司拟通过发行股份购买中国电子工程设计院股份有限公司 100\% 股份并募集配套资金。

\pfl{经营要点} 业务与产能层面，2026 年上半年，公司管理层将收入下降归因于主要产品售价下行、销量下降。

\pfl{所处行业的基本面} 果蔬加工的价格由产地原料果价格与出口订单决定，需求端受海外需求与关税影响（第 \ref{sec:framework} 章的供需框架）。 行业层面的最新情况是：2026 年上半年，同比分别增长 10.51\% 和 2.65\%，出口均价 1475 美元/吨、同比下降 7.06\%。

\pfl{竞争优势} 把该公司放回同类公司比较，2026 年上半年营业收入 8.8 亿元、在三级行业“果蔬加工”的 5 家公司中排第\zhnumber{2}位；毛利率 14.2\%、高于行业内中位数 5.6 个百分点；半年 ROE 0.2\%（未年化）、低于行业内中位数（2.26\%）；资产负债率 52.0\%、高于行业内中位数（30.8\%）。 优势在报表上主要体现为高于行业平均的毛利率，这在周期底部是抵御亏损的主要缓冲。

\pfl{历史股价} 本报告样本区间（2021 年 1 月 至 2026 年 10 月）内，股价由 7.05 元变为 23.84 元，累计 +238.2\%，区间最高 30.17 元（2026 年 6 月）、最低 7.04 元（2021 年 2 月）；当前价相当于区间最高价的 79.0\%，为区间最低价的 3.39 倍。 月度收益的年化波动率 45.1\%，高于全样本中位数 37.3\%，在三级行业“果蔬加工”的 5 家公司中波动率由高到低排第\zhnumber{3}位。 把股价阶段与公司事件对齐看，2025 年 3 月 至 2026 年 6 月股价上涨+137.9\%，同期（2026 年 3 月）公司收到上海证券交易所国投中鲁果汁股份有限公司发行股份购买资产并募集配套资金暨关联交易申请的审核问询函（上证上审（并购重组）。 整体上看，区间的几处大级别拐点都能在公司的资本运作、股权变动或风险事件上找到对应，股价波动有明显的公司层面驱动，而非单纯的行业 beta。

\begin{center}
\begin{minipage}{\textwidth}\centering
\begin{tikzpicture}
\begin{axis}[
  width=12.6cm, height=3.9cm, scale only axis,
  xmin=2020.922, xmax=2026.972, ymin=6.55, ymax=32.28,
  xtick={2021,2022,2023,2024,2025,2026},
  yearticks,
  yticklabel style={font=\scriptsize},
  ylabel={元/股}, ylabel style={font=\scriptsize},
  grid=major, grid style={LightGray, line width=0.3pt},
  axis line style={InkGray, line width=0.5pt},
  every axis plot/.append style={line width=0.9pt},
]
\addplot[AgriGreen] table[x=x,y=y]{data/clean/profiles/px_600962.dat};
\addplot[SoftRed, dashed, line width=0.7pt] table[x=x,y=y]{data/clean/profiles/pxzz_600962.dat};
\addplot[only marks, mark=*, mark size=1.1pt, SoftRed] table[x=x,y=y]{data/clean/profiles/pxzz_600962.dat};
\end{axis}
\end{tikzpicture}

\captionof{figure}[国投中鲁（600962）股价与周期骨架]{国投中鲁（600962）月度收盘价（元/股）与 ZigZag 周期骨架（阈值 25\%，圆点为周期转折点）}
\end{minipage}
\end{center}

\pfl{过去周期分析} 以 25\% 的 ZigZag 阈值划分，样本区间内股价形成 3 个主升段与 2 个主跌段。 最大主跌段为 2023 年 3 月 至 2024 年 2 月，11 个月-37.5\%；最大主升段出现在 2025 年 3 月 至 2026 年 6 月，15 个月+137.9\%。 最近一次转折出现在 2026 年 6 月（上涨段自 2025 年 3 月起，15 个月+137.9\%），当前处于该段之后的震荡之中。 公司股价较申万农林牧渔指数提前约 63 个月见底，是板块中较早定价周期反转的公司之一。 果蔬加工的外向型特征明显，出口订单与原料产地价格决定利润率，公司 2026 年 6 月末资产负债率 52.0\%、半年 ROE 0.2\%（未年化），在本轮下行中属于随行业同步调整的一类。 综合区间形态看，该股单段涨跌幅度偏大、以趋势段为主（最大主升 +137.9\%、最大主跌 -37.5\%），年化波动率 45.1\% 与全样本中位数 37.3\% 相比波动明显大于样本中位数。

\pfl{盈利情况} 2026 年上半年营业收入 8.83 亿元（同比 -12.4\%），归母净利润 0.02 亿元，同比 -92.7\%。 以年报为趋势对照，2023---2025 年营业收入由 14.87 亿元变为 19.84 亿元（两年年均复合增速 +15.5\%），归母净利润由 0.58 亿元变为 0.41 亿元。 按半年报口径，2026 年上半年的 ROE 为 0.2\%（未年化）、销售净利率 0.7\%、毛利率 14.2\%，ROE 在“果蔬加工”5 家公司中按数值由高到低排第\zhnumber{3}位。 同期全样本半年 ROE 中位数 0.75\%，公司与之相差 0.56 个百分点（弱于中位数公司）。 综合看，公司在报告期维持盈利（低于全样本半年 ROE 中位数 0.75\%），利润的可持续性取决于果蔬加工景气的延续与成本管控的执行。

\begin{center}\small\setlength{\tabcolsep}{3.4pt}
\captionof{table}[国投中鲁（600962）关键财务指标]{国投中鲁（600962）关键财务指标（合并报表口径；两个半年报期均未年化；现金短债比 $=$ 货币资金 $\div$ 短期债务）}
\begin{adjustbox}{max width=\textwidth}
\begin{tabular}{@{}lrrrrr@{}}
\toprule
指标 & 2023 年 & 2024 年 & 2025 年 & 2025 年半年报 & 2026 年半年报 \\
\midrule
营业收入（亿元） & 14.87 & 19.87 & 19.84 & 10.07 & 8.83 \\
归母净利润（亿元） & 0.58 & 0.29 & 0.41 & 0.25 & 0.02 \\
毛利率（\%） & 22.7 & 15.5 & 16.7 & 14.6 & 14.2 \\
ROE（\%） & 7.2 & 3.3 & 4.4 & 2.7 & 0.2 \\
资产负债率（\%） & 59.5 & 65.4 & 61.4 & 50.0 & 52.0 \\
经营活动现金流净额（亿元） & -1.66 & -2.83 & 2.42 & 8.75 & 5.35 \\
货币资金（亿元） & 1.73 & 1.62 & 1.68 & 2.56 & 1.57 \\
有息负债合计（亿元） & 13.11 & 16.70 & 15.18 & 9.29 & 10.30 \\
现金短债比（倍） & 0.14 & 0.10 & 0.11 & 0.29 & 0.17 \\
\bottomrule
\end{tabular}
\end{adjustbox}
\end{center}

\begin{center}
\begin{minipage}{\textwidth}\centering
\begin{tikzpicture}
\begin{groupplot}[
  group style={group size=2 by 1, horizontal sep=0.60cm},
  width=6.85cm, height=4.60cm,
  tick label style={font=\tiny},
  title style={font=\scriptsize\heiti},
  yticklabel style={font=\tiny},
  ylabel style={font=\tiny},
  grid=major, grid style={LightGray, line width=0.3pt},
  axis line style={InkGray, line width=0.5pt},
  enlarge x limits={abs=0.8},
  legend style={font=\scriptsize, at={(0.5,-0.20)}, anchor=north,
                legend columns=2, draw=none, fill=none, column sep=6pt},
]
\nextgroupplot[
  ybar, bar width=4pt,
  ymin=-2.20, ymax=22.28,
  xtick={0,1,2,3,4,5.7},
  xticklabels={2021,2022,2023,2024,2025,2026H1},
  ylabel={亿元},
]
\addplot[fill=AgriGreen!75, draw=AgriGreen, bar shift=-2.6pt] table[x=i, y=rev]{data/clean/profiles/fin_600962.dat};
\addplot[fill=SoftRed!60, draw=SoftRed, bar shift=2.6pt] table[x=i, y=ni]{data/clean/profiles/fin_600962.dat};
\addplot[fill=AgriGreen!30, draw=AgriGreen, pattern=north east lines, bar shift=-2.6pt] table[x=x, y=rev]{data/clean/profiles/finh_600962.dat};
\addplot[fill=SoftRed!20, draw=SoftRed, pattern=north east lines, bar shift=2.6pt] table[x=x, y=ni]{data/clean/profiles/finh_600962.dat};
\legend{营业收入（年/半年）, 归母净利润（年/半年）}
\nextgroupplot[
  ymin=-8.3, ymax=72.2,
  xtick={0,1,2,3,4,5.7},
  xticklabels={2021,2022,2023,2024,2025,2026H1},
  ylabel={\%},
]
\addplot[AgriGold, mark=*, mark size=1.3pt] table[x=i, y=roe]{data/clean/profiles/fin_600962.dat};
\addplot[OilBlue, mark=square*, mark size=1.3pt] table[x=i, y=debt]{data/clean/profiles/fin_600962.dat};
\addplot[only marks, AgriGold, mark=*, mark size=2.0pt] table[x=x, y=roe]{data/clean/profiles/finh_600962.dat};
\addplot[only marks, OilBlue, mark=square*, mark size=2.0pt] table[x=x, y=debt]{data/clean/profiles/finh_600962.dat};
\legend{ROE, 资产负债率}
\end{groupplot}
\end{tikzpicture}
\begin{tikzpicture}
\begin{axis}[
  name=finbot,
  width=13.75cm, height=3.10cm, scale only axis,
  ymin=-1.67, ymax=18.71,
  xtick={0,1,2,3,4,5.7},
  xticklabels={2021,2022,2023,2024,2025,2026H1},
  tick label style={font=\tiny},
  yticklabel style={font=\tiny},
  ylabel={亿元}, ylabel style={font=\tiny},
  ybar, bar width=4pt,
  grid=major, grid style={LightGray, line width=0.3pt},
  axis line style={InkGray, line width=0.5pt},
  enlarge x limits={abs=0.8},
  legend style={font=\scriptsize, at={(0.5,-0.26)}, anchor=north,
                legend columns=2, draw=none, fill=none, column sep=10pt},
]
\addplot[fill=AgriGreen!55, draw=AgriGreen, bar shift=-3.0pt] table[x=i, y=cash]{data/clean/profiles/fin_600962.dat};
\addplot[fill=SoftRed!45, draw=SoftRed, bar shift=3.0pt] table[x=i, y=ibd]{data/clean/profiles/fin_600962.dat};
\addplot[fill=AgriGreen!25, draw=AgriGreen, pattern=north east lines, bar shift=-3.0pt] table[x=x, y=cash]{data/clean/profiles/finh_600962.dat};
\addplot[fill=SoftRed!20, draw=SoftRed, pattern=north east lines, bar shift=3.0pt] table[x=x, y=ibd]{data/clean/profiles/finh_600962.dat};
\legend{货币资金（左轴）, 有息负债（左轴）}
\end{axis}
\begin{axis}[
  at={(finbot.south east)}, anchor=south east,
  width=13.75cm, height=3.10cm, scale only axis,
  axis y line*=right, axis x line*=none, xtick=\empty,
  ymin=-4.95, ymax=7.81,
  tick label style={font=\tiny},
  yticklabel style={font=\tiny}, scaled y ticks=false,
  legend style={font=\scriptsize, at={(0.5,-0.44)}, anchor=north,
                legend columns=1, draw=none, fill=none},
]
\addplot[OilBlue, mark=*, mark size=2.2pt, line width=1.3pt] table[x=i, y=ocf]{data/clean/profiles/fin_600962.dat};
\addplot[only marks, OilBlue, mark=*, mark size=3.2pt] table[x=x, y=ocf]{data/clean/profiles/finh_600962.dat};
\legend{经营活动现金流净额（右轴，亿元，蓝点）}
\end{axis}
\end{tikzpicture}

\captionof{figure}[国投中鲁（600962）近年主要财务指标]{国投中鲁（600962）近年主要财务指标（左上：营业收入与归母净利润；右上：ROE 与资产负债率；下：货币资金、有息负债为左轴柱状，经营活动现金流净额为其右轴的蓝点折线。
2021---2025 年为年报口径，2026H1 为 2026 年半年报/半年末，与其前一年报之间留出间隔、以斜纹或空心点区分；半年报数据未年化）}
\end{minipage}
\end{center}

\pfl{财务分析} 2026 年 6 月末资产负债率 52.0\%，较 2025 年末下降 9.4 个百分点（样本中位数 52.2\%，所处三级行业内按由低到高排第\zhnumber{4}位）。 2026 年 6 月末货币资金 1.57 亿元、短期债务（短期借款与一年内到期非流动负债）9.37 亿元、长期债务 0.93 亿元，有息负债合计 10.30 亿元，现金短债比 0.17 倍——缺口明显。 净有息负债 8.73 亿元，相当于其 2026 年上半年营业收入的 99\%。 流动比率 1.37、速动比率 0.59，短期流动性无明显缺口。 2026 年上半年经营活动现金流净额 5.35 亿元，经营现金流与利润同向为正，内生资金可以支撑运营。 综合看，账面货币资金对有息负债与短期债务的覆盖不足，滚动偿付对新增融资的依赖度较高；资产负债率低于样本中位数 0.2 个百分点；有息负债中短期债务占比更高，期限结构仍以滚动续作为主。 公司披露的监管与合规风险为：2026 年 3 月，公司收到上海证券交易所国投中鲁果汁股份有限公司发行股份购买资产并募集配套资金暨关联交易申请的审核问询函（上证上审（并购重组）。

\pfl{重大战略转型} 从公开披露看，公司的战略推进集中在：并购与重组（3 项）：2026 年 8 月，重组获上交所并购重组审核委员会审核通过并取得中国证监会同意注册批复；2025 年 12 月，提升上市公司资产规模和盈利能力；主业转型（1 项）：1992 年 8 月，标的电子院成立于 该日。 公告口径上，近三年（2023 年 9 月---2026 年 9 月）公司披露重大重组与并购类公告 32 条、对外投资与转型类公告 0 条、回购与股权激励类公告 0 条，合计公告 400 条。 第一大行业仍是“饮料制造业”，收入占比 99.5\%，与 2021 年年报相比没有方向性变化。 综合判断，报告期内公司有 32 条与战略相关的公告落地（重大重组与并购 32 条、对外投资与转型 0 条），资本运作与主业调整之间存在直接关联。

\pfl{历史融投资情况} 从现金流量表看，2025 年公司取得借款收到的现金 14.41 亿元、偿还债务支付的现金 16.17 亿元，筹资活动现金流净额 -2.04 亿元（2023 年为取得借款 12.36 亿元、偿还债务 10.40 亿元）。 2026 年上半年筹资活动现金流净额为 -4.99 亿元（取得借款 0.58 亿元、偿还债务 5.47 亿元，未年化），已转为净流出，处于净还债状态。 2025 年分配股利、利润或偿付利息支付的现金合计 0.27 亿元。 公开披露的融资记录共 2 笔，工具类型以 IPO、发行股份购买资产为主；金额与用途见下表。 其中规模最大的两笔为：2025 年 12 月，发行股份购买资产，发行股份购买电子院 100\% 股份的交易价格 60.26 亿元（60.26 亿元），同时拟向不超过 35 名特定投资者发行股份募集配套资金不超过 17.26 亿元；2004 年 6 月，IPO，发行 6,500 万股，发行价 4.80 元/股。 股本方面，近三年披露股本变动 3 次；总股本自 2021 年末以来未发生变化（2.62 亿股）。 分红方面，近三年未实施现金分红，在全部样本中属于分红能力偏弱的一端。 近三年“再融资”类公告 85 条，涉及发行股份购买资产。 综合看，公司融资结构上股权与债务融资并举；2025 年筹资活动现金流净额 -2.04 亿元，融资节奏仍以维持存量债务周转为主，与前述经营与投资现金流的方向基本一致。

\pfl{未来潜在融资需求分析} 2026 年 6 月末货币资金 1.57 亿元，而短期债务 9.37 亿元（现金短债比 0.17 倍），2026 年上半年经营现金流净额 5.35 亿元。按“短期债务减去账面货币资金”的滚动偿付口径，静态缺口约 7.80 亿元；上半年盈利或接近盈亏平衡，该缺口不随经营亏损进一步扩大。 按第 \ref{sec:financing-need} 节的三档划分，公司属于第二档“保壳型”：近三年重大重组与并购类公告 32 条，2026 年上半年归母净利润 0.02 亿元，融资需求更多体现为“资产置换 + 维持上市地位”，而不是扩产。 公开披露的后续资金安排线索包括：股权融资线索：2025 年 12 月，公司拟向不超过 35 名特定投资者发行股份募集配套资金不超过 17.26 亿元，其中补充流动资金 7.00 亿元。 综合看，公司的外部资金需求以补充流动性与项目投入为主，滚动偿付口径下的静态缺口约 7.80 亿元，融资节奏将受制于银行续贷意愿与股权融资窗口。




\subsubsection{晨光生物（300138）}
\label{co:300138}

\pfl{公司基本情况} 晨光生物成立于 2000 年，河北曲周起家的天然植物提取物龙头，辣椒红色素、辣椒精、叶黄素三大主力产品位居世界第一或前列，棉籽类业务规模大但毛利薄；2010 年创业板上市，成本领先与全球份额是核心逻辑。 公司于 2010-11-05 在深交所创业板首发上市，注册资本 4.83 亿元，总股本 4.83 亿股。 截至 2026 年 10 月 12 日收盘价 11.81 元，对应总市值 57.1 亿元，近一年-12.7\%，总市值在三级行业“其他农产品加工”的 11 家公司中按由大到小排第\zhnumber{3}位，近一年涨跌幅低于全样本中位数（-10.4\%）2.4 个百分点。 按 2025 年年报披露的行业构成，收入几乎全部来自天然植物提取行业 100.0\%。

\pfl{股权结构与控制权} 按 2026 年 7 月 15 日的股东名册，前三大股东为卢庆国 持股 20.33\%；李月斋 持股 3.32\%；关庆彬 持股 1.89\%。 股东结构与控制权方面，公司第一大股东卢庆国持股 9,819.61 万股、占 20.33\%；公司曾计划将棉籽板块子公司新疆晟源生物在北交所上市。

\pfl{经营要点} 经营与产能方面，棉籽类业务实现收入 13.79 亿元、同比减少 21.84\%；公司在河北、新疆、云南、印度等地设立原材料基地；公司产品定价采用成本加成模式。

\pfl{所处行业的基本面} 该子行业横跨制糖、番茄制品、糖醇与添加剂，供给受种植面积与进口配额约束，需求端分食用与工业两类（第 \ref{sec:framework} 章的供需框架）。 当期该行业的关键数据是：2026 年，印度辣椒单位价格高于国内。

\pfl{竞争优势} 把该公司放回同类公司比较，2026 年上半年营业收入 31.6 亿元、在三级行业“其他农产品加工”的 11 家公司中排第\zhnumber{3}位；毛利率 13.5\%、高于行业内中位数 4.4 个百分点；半年 ROE 5.2\%（未年化）、高于行业内中位数（1.36\%）；资产负债率 64.1\%、高于行业内中位数（39.1\%）。 与同类公司相比，公司的差异化在于：持续降低原材料成本。 优势在报表上主要体现为高于行业平均的毛利率，这在周期底部是抵御亏损的主要缓冲。

\pfl{历史股价} 以 2021 年 1 月为起点、2026 年 10 月为终点，股价由 16.50 元变为 11.81 元，累计 -28.4\%，区间最高 18.53 元（2023 年 1 月）、最低 7.52 元（2024 年 8 月）；当前价相当于区间最高价的 63.7\%，为区间最低价的 1.57 倍。 月度收益的年化波动率 33.0\%，低于全样本中位数 37.3\%，在三级行业“其他农产品加工”的 11 家公司中波动率由高到低排第\zhnumber{9}位。 把股价阶段与公司事件对齐看，2024 年 8 月 至 2026 年 2 月股价上涨+92.7\%，同期（2025 年 12 月）可转债募集资金已累计使用 3.89 亿元。 从时间对应关系看，公司层面的资本运作与股权、风险事件解释了区间内多数急涨急跌，与行业指数的同步性相对有限。

\begin{center}
\begin{minipage}{\textwidth}\centering
\begin{tikzpicture}
\begin{axis}[
  width=12.6cm, height=3.9cm, scale only axis,
  xmin=2020.922, xmax=2026.972, ymin=6.99, ymax=19.83,
  xtick={2021,2022,2023,2024,2025,2026},
  yearticks,
  yticklabel style={font=\scriptsize},
  ylabel={元/股}, ylabel style={font=\scriptsize},
  grid=major, grid style={LightGray, line width=0.3pt},
  axis line style={InkGray, line width=0.5pt},
  every axis plot/.append style={line width=0.9pt},
]
\addplot[AgriGreen] table[x=x,y=y]{data/clean/profiles/px_300138.dat};
\addplot[SoftRed, dashed, line width=0.7pt] table[x=x,y=y]{data/clean/profiles/pxzz_300138.dat};
\addplot[only marks, mark=*, mark size=1.1pt, SoftRed] table[x=x,y=y]{data/clean/profiles/pxzz_300138.dat};
\end{axis}
\end{tikzpicture}

\captionof{figure}[晨光生物（300138）股价与周期骨架]{晨光生物（300138）月度收盘价（元/股）与 ZigZag 周期骨架（阈值 25\%，圆点为周期转折点）}
\end{minipage}
\end{center}

\pfl{过去周期分析} 以 25\% 的 ZigZag 阈值划分，样本区间内股价形成 3 个主升段与 3 个主跌段。 最大主跌段为 2023 年 1 月 至 2024 年 8 月，19 个月-59.4\%；最大主升段出现在 2024 年 8 月 至 2026 年 2 月，18 个月+92.7\%。 最近一次转折出现在 2026 年 10 月（上涨段自 2026 年 6 月起，4 个月+31.2\%），当前处于该段之后的震荡之中。 公司股价较申万农林牧渔指数提前约 21 个月见底，是板块中较早定价周期反转的公司之一。 该子行业横跨糖、番茄制品、添加剂等多个品种，各自的供给周期与政策（进口配额、关税）是主要变量，公司 2026 年上半年实现盈利、半年 ROE 5.2\%（未年化），好于全样本中位数（0.75\%），下行周期对其报表的冲击小于同业。 综合区间形态看，该股单段涨跌幅度偏大、以趋势段为主（最大主升 +92.7\%、最大主跌 -59.4\%），年化波动率 33.0\% 与全样本中位数 37.3\% 相比波动小于样本中位数。

\pfl{盈利情况} 2026 年上半年营业收入 31.59 亿元（同比 -13.6\%），归母净利润 1.79 亿元，同比 -17.0\%。 以年报为趋势对照，2023---2025 年营业收入由 68.72 亿元变为 65.59 亿元（两年年均复合增速 -2.3\%），归母净利润由 4.80 亿元变为 3.69 亿元。 按半年报口径，2026 年上半年的 ROE 为 5.2\%（未年化）、销售净利率 5.8\%、毛利率 13.5\%，ROE 在“其他农产品加工”11 家公司中按数值由高到低排第\zhnumber{2}位。 按同一口径，三级行业“其他农产品加工”11 家公司的半年 ROE 中位数为 1.36\%，该公司高于其中位数 3.79 个百分点。 从公司披露的原因看，业绩变动主要来自：2025 年 6 月，受原材料成本下降影响，多数植物提取产品价格同比下降。 综合看，公司在报告期维持盈利（高于全样本半年 ROE 中位数 0.75\%），利润的可持续性取决于其他农产品加工景气的延续与成本管控的执行。

\begin{center}\small\setlength{\tabcolsep}{3.4pt}
\captionof{table}[晨光生物（300138）关键财务指标]{晨光生物（300138）关键财务指标（合并报表口径；两个半年报期均未年化；现金短债比 $=$ 货币资金 $\div$ 短期债务）}
\begin{adjustbox}{max width=\textwidth}
\begin{tabular}{@{}lrrrrr@{}}
\toprule
指标 & 2023 年 & 2024 年 & 2025 年 & 2025 年半年报 & 2026 年半年报 \\
\midrule
营业收入（亿元） & 68.72 & 69.94 & 65.59 & 36.58 & 31.59 \\
归母净利润（亿元） & 4.80 & 0.94 & 3.69 & 2.15 & 1.79 \\
毛利率（\%） & 11.6 & 8.0 & 13.1 & 13.9 & 13.5 \\
ROE（\%） & 14.3 & 2.8 & 11.4 & 6.7 & 5.2 \\
资产负债率（\%） & 58.6 & 65.0 & 63.8 & 64.4 & 64.1 \\
经营活动现金流净额（亿元） & -5.88 & -2.23 & 2.13 & 18.30 & 12.18 \\
货币资金（亿元） & 14.22 & 18.02 & 14.59 & 36.20 & 28.02 \\
有息负债合计（亿元） & 42.36 & 49.31 & 51.23 & 53.01 & 53.17 \\
现金短债比（倍） & 0.42 & 0.47 & 0.29 & 0.71 & 0.53 \\
\bottomrule
\end{tabular}
\end{adjustbox}
\end{center}

\begin{center}
\begin{minipage}{\textwidth}\centering
\begin{tikzpicture}
\begin{groupplot}[
  group style={group size=2 by 1, horizontal sep=0.60cm},
  width=6.85cm, height=4.60cm,
  tick label style={font=\tiny},
  title style={font=\scriptsize\heiti},
  yticklabel style={font=\tiny},
  ylabel style={font=\tiny},
  grid=major, grid style={LightGray, line width=0.3pt},
  axis line style={InkGray, line width=0.5pt},
  enlarge x limits={abs=0.8},
  legend style={font=\scriptsize, at={(0.5,-0.20)}, anchor=north,
                legend columns=2, draw=none, fill=none, column sep=6pt},
]
\nextgroupplot[
  ybar, bar width=4pt,
  ymin=-6.99, ymax=78.33,
  xtick={0,1,2,3,4,5.7},
  xticklabels={2021,2022,2023,2024,2025,2026H1},
  ylabel={亿元},
]
\addplot[fill=AgriGreen!75, draw=AgriGreen, bar shift=-2.6pt] table[x=i, y=rev]{data/clean/profiles/fin_300138.dat};
\addplot[fill=SoftRed!60, draw=SoftRed, bar shift=2.6pt] table[x=i, y=ni]{data/clean/profiles/fin_300138.dat};
\addplot[fill=AgriGreen!30, draw=AgriGreen, pattern=north east lines, bar shift=-2.6pt] table[x=x, y=rev]{data/clean/profiles/finh_300138.dat};
\addplot[fill=SoftRed!20, draw=SoftRed, pattern=north east lines, bar shift=2.6pt] table[x=x, y=ni]{data/clean/profiles/finh_300138.dat};
\legend{营业收入（年/半年）, 归母净利润（年/半年）}
\nextgroupplot[
  ymin=-5.2, ymax=71.5,
  xtick={0,1,2,3,4,5.7},
  xticklabels={2021,2022,2023,2024,2025,2026H1},
  ylabel={\%},
]
\addplot[AgriGold, mark=*, mark size=1.3pt] table[x=i, y=roe]{data/clean/profiles/fin_300138.dat};
\addplot[OilBlue, mark=square*, mark size=1.3pt] table[x=i, y=debt]{data/clean/profiles/fin_300138.dat};
\addplot[only marks, AgriGold, mark=*, mark size=2.0pt] table[x=x, y=roe]{data/clean/profiles/finh_300138.dat};
\addplot[only marks, OilBlue, mark=square*, mark size=2.0pt] table[x=x, y=debt]{data/clean/profiles/finh_300138.dat};
\legend{ROE, 资产负债率}
\end{groupplot}
\end{tikzpicture}
\begin{tikzpicture}
\begin{axis}[
  name=finbot,
  width=13.75cm, height=3.10cm, scale only axis,
  ymin=-5.32, ymax=59.55,
  xtick={0,1,2,3,4,5.7},
  xticklabels={2021,2022,2023,2024,2025,2026H1},
  tick label style={font=\tiny},
  yticklabel style={font=\tiny},
  ylabel={亿元}, ylabel style={font=\tiny},
  ybar, bar width=4pt,
  grid=major, grid style={LightGray, line width=0.3pt},
  axis line style={InkGray, line width=0.5pt},
  enlarge x limits={abs=0.8},
  legend style={font=\scriptsize, at={(0.5,-0.26)}, anchor=north,
                legend columns=2, draw=none, fill=none, column sep=10pt},
]
\addplot[fill=AgriGreen!55, draw=AgriGreen, bar shift=-3.0pt] table[x=i, y=cash]{data/clean/profiles/fin_300138.dat};
\addplot[fill=SoftRed!45, draw=SoftRed, bar shift=3.0pt] table[x=i, y=ibd]{data/clean/profiles/fin_300138.dat};
\addplot[fill=AgriGreen!25, draw=AgriGreen, pattern=north east lines, bar shift=-3.0pt] table[x=x, y=cash]{data/clean/profiles/finh_300138.dat};
\addplot[fill=SoftRed!20, draw=SoftRed, pattern=north east lines, bar shift=3.0pt] table[x=x, y=ibd]{data/clean/profiles/finh_300138.dat};
\legend{货币资金（左轴）, 有息负债（左轴）}
\end{axis}
\begin{axis}[
  at={(finbot.south east)}, anchor=south east,
  width=13.75cm, height=3.10cm, scale only axis,
  axis y line*=right, axis x line*=none, xtick=\empty,
  ymin=-10.58, ymax=17.60,
  tick label style={font=\tiny},
  yticklabel style={font=\tiny}, scaled y ticks=false,
  legend style={font=\scriptsize, at={(0.5,-0.44)}, anchor=north,
                legend columns=1, draw=none, fill=none},
]
\addplot[OilBlue, mark=*, mark size=2.2pt, line width=1.3pt] table[x=i, y=ocf]{data/clean/profiles/fin_300138.dat};
\addplot[only marks, OilBlue, mark=*, mark size=3.2pt] table[x=x, y=ocf]{data/clean/profiles/finh_300138.dat};
\legend{经营活动现金流净额（右轴，亿元，蓝点）}
\end{axis}
\end{tikzpicture}

\captionof{figure}[晨光生物（300138）近年主要财务指标]{晨光生物（300138）近年主要财务指标（左上：营业收入与归母净利润；右上：ROE 与资产负债率；下：货币资金、有息负债为左轴柱状，经营活动现金流净额为其右轴的蓝点折线。
2021---2025 年为年报口径，2026H1 为 2026 年半年报/半年末，与其前一年报之间留出间隔、以斜纹或空心点区分；半年报数据未年化）}
\end{minipage}
\end{center}

\pfl{财务分析} 2026 年 6 月末资产负债率 64.1\%，较 2025 年末上升 0.2 个百分点（样本中位数 52.2\%，所处三级行业内按由低到高排第\zhnumber{10}位）。 2026 年 6 月末货币资金 28.02 亿元、短期债务（短期借款与一年内到期非流动负债）52.71 亿元、长期债务 0.46 亿元，有息负债合计 53.17 亿元，现金短债比 0.53 倍——覆盖偏紧。 净有息负债 25.16 亿元，相当于其 2026 年上半年营业收入的 80\%。 流动比率 1.28、速动比率 0.73，短期流动性无明显缺口。 2026 年上半年经营活动现金流净额 12.18 亿元，经营现金流与利润同向为正，内生资金可以支撑运营。 2026 年 6 月末在建工程 0.94 亿元，较上年末增加 0.21 亿元，仍有在建投入。 综合看，现金短债比 0.53 倍，短期资金安排偏紧；资产负债率高于样本中位数 11.9 个百分点；有息负债中短期债务占比更高，期限结构仍以滚动续作为主。

\pfl{重大战略转型} 近三年公司的战略动作集中在：其他动作（3 项）：2025 年全年，归属于上市公司股东的净利润 3.69 亿元；2025 年 6 月，植提类产品收入增长 9.47\%、毛利额同比增长 32\%。 公告口径上，近三年（2023 年 9 月---2026 年 9 月）公司披露重大重组与并购类公告 0 条、对外投资与转型类公告 1 条、回购与股权激励类公告 46 条，合计公告 619 条。 第一大行业仍是“天然植物提取行业”，收入占比 100.0\%，与 2021 年年报相比没有方向性变化。 综合判断，报告期内公司有 1 条与战略相关的公告落地（重大重组与并购 0 条、对外投资与转型 1 条），资本运作与主业调整之间存在直接关联。

\begin{center}\small\setlength{\tabcolsep}{4pt}
\captionof{table}[晨光生物（300138）历史融资与资本运作]{晨光生物（300138）历史融资与资本运作（据公司公告与公开报道整理；金额为披露值，含首次公开发行、再融资、债券、银行借款与重整投资等）}
\begin{adjustbox}{max width=\textwidth}
\begin{tabular}{@{}p{1.45cm}p{2.55cm}p{4.05cm}p{4.95cm}@{}}
\toprule
时间 & 方式 & 规模 & 用途与说明 \\
\midrule
2010-11-05 & IPO（深交所创业板） & 发行 2,\allowbreak{}300 万股，\allowbreak{}发行价 30.00 元/\allowbreak{}股，\allowbreak{}募集资金净额 6.49 亿元 & 招股说明书披露的募投项目；；发行后总股本不超过 8,\allowbreak{}978.5436 万股 \\
2020-06-17 & 公开发行可转换公司债券（创业板） & 面值总额 6.30 亿元（630 万张），\allowbreak{}募集资金总额 6.30 亿元 & 天然植物综合提取一体化项目（一期）4.50 亿元、\allowbreak{}补充流动资金 1.80 亿元；；期限 6 年；原 A 股股东优先配售；2020 年 7 月 13 日上市 \\
2021-02-24 & 可转债提前赎回与转股 & 2020 年 12 月 23 日至停止交易日共 6,\allowbreak{}273,\allowbreak{}758 张转股，\allowbreak{}转股价 12.25 元/\allowbreak{}股；未转股 26,\allowbreak{}242 张按 100.35 元/\allowbreak{}张赎回 & 转股增加股本 \\
2023-03-28 & 对全资子公司增资（自有资金） & “赞比亚土地开发及配套设施建设项目”建设总投资 1,\allowbreak{}500.00 万美元 & 配套增建菊花加工颗粒车间、\allowbreak{}萃取车间、\allowbreak{}菊花存储池 \\
2023-05-29 & 变更部分可转债募集资金用途 & 将“天然植物综合提取一体化项目（一期）”部分募集资金 6,\allowbreak{}500.00 万元变更用于“营养药用综合提取项目（一期）” & 营养药用综合提取项目（一期）、\allowbreak{}赞比亚土地开发及配套设施建设项目；2023 年 7 月 12 日与集资金三方监管协议 \\
2024-06-30 & 子公司首次公开发行终止 & --- & 控股子公司新疆晨光生物科技股份有限公司决定终止向不特定合格投资者公开发行并在北交所上市 \\
\bottomrule
\end{tabular}
\end{adjustbox}
\end{center}

\pfl{历史融投资情况} 从现金流量表看，2025 年公司取得借款收到的现金 59.76 亿元、偿还债务支付的现金 57.66 亿元，筹资活动现金流净额 0.69 亿元（2023 年为取得借款 46.78 亿元、偿还债务 33.80 亿元）。 2026 年上半年筹资活动现金流净额为 -0.42 亿元（取得借款 27.24 亿元、偿还债务 25.28 亿元，未年化），已转为净流出，处于净还债状态。 2025 年分配股利、利润或偿付利息支付的现金合计 1.60 亿元。 公开披露的融资记录共 6 笔，工具类型以 IPO、可转债为主；金额与用途见下表。 其中规模最大的两笔为：2023 年 5 月，可转债，将“天然植物综合提取一体化项目（一期）”部分募集资金 6,500.00 万元变更用于“营养药用综合提取项目（一期）”；2020 年 6 月，可转债，面值总额 6.30 亿元（630 万张），募集资金总额 6.30 亿元。 股本方面，近三年披露股本变动 5 次（其他 1 次、股份回购 1 次）；总股本由 2021 年末的 5.33 亿股变为 4.83 亿股（-9.32\%）。 分红方面，近三年现金分红 4 次、合计每股派现 0.72 元，最近一次实施在 2025 年 12 月。 综合看，公司融资结构上股权与债务融资并举；2025 年筹资活动现金流净额 0.69 亿元，年度内仍为净融入，与前述经营与投资现金流的方向基本一致。

\pfl{未来潜在融资需求分析} 2026 年 6 月末货币资金 28.02 亿元，而短期债务 52.71 亿元（现金短债比 0.53 倍），2026 年上半年经营现金流净额 12.18 亿元。按“短期债务减去账面货币资金”的滚动偿付口径，静态缺口约 24.70 亿元；上半年盈利或接近盈亏平衡，该缺口不随经营亏损进一步扩大。 按第 \ref{sec:financing-need} 节的三档划分，公司属于第三档“扩张型”：2026 年上半年 ROE 5.2\%（未年化，折合约 10.3\%）、6 月末资产负债率 64.1\%，资产负债表仍具备加杠杆空间；融资用途以产能扩张、品类并购与海外布局为主，单笔规模通常在 5---20 亿元量级，会被市场定价为成长而非补血。 公开披露的后续资金安排线索包括：股权融资线索：2025 年 12 月，可转债募集资金已累计使用 3.89 亿元。 资金需求还要叠加在建工程：期末在建工程 0.94 亿元、2026 年上半年购建固定资产等支出 0.60 亿元（未年化），这部分支出在周期底部通常会被主动放缓。 综合看，公司的外部资金需求以补充流动性与项目投入为主，滚动偿付口径下的静态缺口约 24.70 亿元，能否落地取决于授信续作与发行窗口的配合。




\subsubsection{双塔食品（002481）}
\label{co:002481}

\pfl{公司基本情况} 双塔食品成立于 2010 年，山东招远龙口粉丝行业唯一上市公司，主营豌豆蛋白、粉丝、膳食纤维，是全球领先的豌豆蛋白与豌豆淀粉生产基地；近年受豌豆蛋白海外出口及淀粉售价影响盈利大幅回落，控股股东高比例股权质押构成治理与债务风险关注点。 公司于 2010-09-21 在深交所主板首发上市，注册资本 12.34 亿元，总股本 12.34 亿股。 截至 2026 年 10 月 12 日收盘价 4.41 元，对应总市值 54.4 亿元，近一年-23.8\%，总市值在三级行业“其他农产品加工”的 11 家公司中按由大到小排第\zhnumber{4}位，近一年涨跌幅低于全样本中位数（-10.4\%）13.5 个百分点。 按 2026 年半年报披露的行业构成，收入几乎全部来自农副产品加工业 100.0\%。 与 2021 年年报相比，第一大行业口径由“农副食品加工业”（100.0\%）变为“农副产品加工业”（100.0\%）；两期披露的分类口径有调整，占比差异中同时包含分类因素。

\pfl{股权结构与控制权} 按 2026 年 7 月 7 日的股东名册，前三大股东为招远君兴农业发展中心 持股 32.69\%；杨君敏 持股 12.25\%；烟台双塔食品股份有限公司回购专用证券账户 持股 2.94\%。 控制权与股东性质方面，2026 年 9 月，控股股东君兴农业将合计 1669 万股公司股份办理质押；公司 2023 年员工持股计划在报告期内有效，由董事、监事、高级管理人员及核心人员共 36 人持有公司 3,162.64 万股股票。

\pfl{经营要点} 经营层面，营业成本中原材料占比 78.55\%、金额 7.11 亿元。

\pfl{所处行业的基本面} 该子行业横跨制糖、番茄制品、糖醇与添加剂，供给受种植面积与进口配额约束，需求端分食用与工业两类（第 \ref{sec:framework} 章的供需框架）。 与该公司直接相关的行业数据包括：2023 年 7 月，美国商务部对公司豌豆蛋白启动反倾销、反补贴调查；替代蛋白市场份额将由 2020 年的 2\% 增长到 11\%。

\pfl{竞争优势} 把该公司放回同类公司比较，2026 年上半年营业收入 11.2 亿元、在三级行业“其他农产品加工”的 11 家公司中排第\zhnumber{7}位；毛利率 19.4\%、高于行业内中位数 10.3 个百分点；半年 ROE 1.4\%（未年化）、低于行业内中位数（1.36\%）；资产负债率 44.8\%、高于行业内中位数（39.1\%）。 公司在公开资料中自述与同业比较的相对优势在于：豌豆蛋白、粉丝、膳食纤维产能多年稳居行业第一。 优势在报表上主要体现为高于行业平均的毛利率，这在周期底部是抵御亏损的主要缓冲。

\pfl{历史股价} 自 2021 年 1 月至 2026 年 10 月，股价由 12.00 元变为 4.41 元，累计 -63.3\%，区间最高 13.36 元（2021 年 5 月）、最低 3.79 元（2026 年 6 月）；当前价相当于区间最高价的 33.0\%，为区间最低价的 1.16 倍。 月度收益的年化波动率 34.7\%，低于全样本中位数 37.3\%，在三级行业“其他农产品加工”的 11 家公司中波动率由高到低排第\zhnumber{6}位。

\begin{center}
\begin{minipage}{\textwidth}\centering
\begin{tikzpicture}
\begin{axis}[
  width=12.6cm, height=3.9cm, scale only axis,
  xmin=2020.922, xmax=2026.972, ymin=3.52, ymax=14.30,
  xtick={2021,2022,2023,2024,2025,2026},
  yearticks,
  yticklabel style={font=\scriptsize},
  ylabel={元/股}, ylabel style={font=\scriptsize},
  grid=major, grid style={LightGray, line width=0.3pt},
  axis line style={InkGray, line width=0.5pt},
  every axis plot/.append style={line width=0.9pt},
]
\addplot[AgriGreen] table[x=x,y=y]{data/clean/profiles/px_002481.dat};
\addplot[SoftRed, dashed, line width=0.7pt] table[x=x,y=y]{data/clean/profiles/pxzz_002481.dat};
\addplot[only marks, mark=*, mark size=1.1pt, SoftRed] table[x=x,y=y]{data/clean/profiles/pxzz_002481.dat};
\end{axis}
\end{tikzpicture}

\captionof{figure}[双塔食品（002481）股价与周期骨架]{双塔食品（002481）月度收盘价（元/股）与 ZigZag 周期骨架（阈值 25\%，圆点为周期转折点）}
\end{minipage}
\end{center}

\pfl{过去周期分析} 以 25\% 的 ZigZag 阈值划分，样本区间内股价形成 2 个主升段与 3 个主跌段。 最大主升段出现在 2024 年 6 月 至 2025 年 8 月，14 个月+46.5\%；最大主跌段为 2022 年 5 月 至 2024 年 6 月，25 个月-58.5\%。 最近一次转折出现在 2026 年 6 月（下跌段自 2025 年 8 月起，10 个月-33.2\%），当前处于该段的延续之中。 公司股价与申万农林牧渔指数几乎同步见底，个股并未走出独立行情。 该子行业横跨糖、番茄制品、添加剂等多个品种，各自的供给周期与政策（进口配额、关税）是主要变量，公司 2026 年上半年实现盈利、半年 ROE 1.4\%（未年化），好于全样本中位数（0.75\%），下行周期对其报表的冲击小于同业。 综合区间形态看，该股单段涨跌幅度偏大、以趋势段为主（最大主升 +46.5\%、最大主跌 -58.5\%），年化波动率 34.7\% 与全样本中位数 37.3\% 相比波动与样本中位数接近。

\pfl{盈利情况} 2026 年上半年营业收入 11.23 亿元（同比 +7.3\%），归母净利润 0.33 亿元，同比 -39.0\%。 以年报为趋势对照，2023---2025 年营业收入由 24.56 亿元变为 21.21 亿元（两年年均复合增速 -7.1\%），归母净利润由 0.93 亿元变为 0.33 亿元。 按半年报口径，2026 年上半年的 ROE 为 1.4\%（未年化）、销售净利率 2.9\%、毛利率 19.4\%，ROE 在“其他农产品加工”11 家公司中按数值由高到低排第\zhnumber{6}位。 按同一口径，三级行业“其他农产品加工”11 家公司的半年 ROE 中位数为 1.36\%，该公司低于其中位数 0.00 个百分点。 从公司披露的原因看，业绩变动主要来自：2025 年度，业绩预告预计归母净利润 2,800.00 万元至 4,000.00 万元。 综合看，公司在报告期维持盈利（高于全样本半年 ROE 中位数 0.75\%），利润的可持续性取决于其他农产品加工景气的延续与成本管控的执行。

\begin{center}\small\setlength{\tabcolsep}{3.4pt}
\captionof{table}[双塔食品（002481）关键财务指标]{双塔食品（002481）关键财务指标（合并报表口径；两个半年报期均未年化；现金短债比 $=$ 货币资金 $\div$ 短期债务）}
\begin{adjustbox}{max width=\textwidth}
\begin{tabular}{@{}lrrrrr@{}}
\toprule
指标 & 2023 年 & 2024 年 & 2025 年 & 2025 年半年报 & 2026 年半年报 \\
\midrule
营业收入（亿元） & 24.56 & 24.48 & 21.21 & 10.47 & 11.23 \\
归母净利润（亿元） & 0.93 & 0.95 & 0.33 & 0.54 & 0.33 \\
毛利率（\%） & 14.8 & 15.0 & 20.2 & 17.9 & 19.4 \\
ROE（\%） & 3.9 & 3.6 & 1.3 & 2.0 & 1.4 \\
资产负债率（\%） & 32.5 & 43.7 & 47.6 & 41.5 & 44.8 \\
经营活动现金流净额（亿元） & 2.31 & 3.25 & 0.91 & -1.79 & 1.27 \\
货币资金（亿元） & 4.08 & 11.28 & 7.32 & 6.10 & 3.84 \\
有息负债合计（亿元） & 7.20 & 14.30 & 16.02 & 12.94 & 12.29 \\
现金短债比（倍） & 0.69 & 0.86 & 0.56 & 0.54 & 0.43 \\
\bottomrule
\end{tabular}
\end{adjustbox}
\end{center}

\begin{center}
\begin{minipage}{\textwidth}\centering
\begin{tikzpicture}
\begin{groupplot}[
  group style={group size=2 by 1, horizontal sep=0.60cm},
  width=6.85cm, height=4.60cm,
  tick label style={font=\tiny},
  title style={font=\scriptsize\heiti},
  yticklabel style={font=\tiny},
  ylabel style={font=\tiny},
  grid=major, grid style={LightGray, line width=0.3pt},
  axis line style={InkGray, line width=0.5pt},
  enlarge x limits={abs=0.8},
  legend style={font=\scriptsize, at={(0.5,-0.20)}, anchor=north,
                legend columns=2, draw=none, fill=none, column sep=6pt},
]
\nextgroupplot[
  ybar, bar width=4pt,
  ymin=-5.88, ymax=27.88,
  xtick={0,1,2,3,4,5.7},
  xticklabels={2021,2022,2023,2024,2025,2026H1},
  ylabel={亿元},
]
\addplot[fill=AgriGreen!75, draw=AgriGreen, bar shift=-2.6pt] table[x=i, y=rev]{data/clean/profiles/fin_002481.dat};
\addplot[fill=SoftRed!60, draw=SoftRed, bar shift=2.6pt] table[x=i, y=ni]{data/clean/profiles/fin_002481.dat};
\addplot[fill=AgriGreen!30, draw=AgriGreen, pattern=north east lines, bar shift=-2.6pt] table[x=x, y=rev]{data/clean/profiles/finh_002481.dat};
\addplot[fill=SoftRed!20, draw=SoftRed, pattern=north east lines, bar shift=2.6pt] table[x=x, y=ni]{data/clean/profiles/finh_002481.dat};
\legend{营业收入（年/半年）, 归母净利润（年/半年）}
\nextgroupplot[
  ymin=-16.9, ymax=53.5,
  xtick={0,1,2,3,4,5.7},
  xticklabels={2021,2022,2023,2024,2025,2026H1},
  ylabel={\%},
]
\addplot[AgriGold, mark=*, mark size=1.3pt] table[x=i, y=roe]{data/clean/profiles/fin_002481.dat};
\addplot[OilBlue, mark=square*, mark size=1.3pt] table[x=i, y=debt]{data/clean/profiles/fin_002481.dat};
\addplot[only marks, AgriGold, mark=*, mark size=2.0pt] table[x=x, y=roe]{data/clean/profiles/finh_002481.dat};
\addplot[only marks, OilBlue, mark=square*, mark size=2.0pt] table[x=x, y=debt]{data/clean/profiles/finh_002481.dat};
\legend{ROE, 资产负债率}
\end{groupplot}
\end{tikzpicture}
\begin{tikzpicture}
\begin{axis}[
  name=finbot,
  width=13.75cm, height=3.10cm, scale only axis,
  ymin=-1.60, ymax=17.95,
  xtick={0,1,2,3,4,5.7},
  xticklabels={2021,2022,2023,2024,2025,2026H1},
  tick label style={font=\tiny},
  yticklabel style={font=\tiny},
  ylabel={亿元}, ylabel style={font=\tiny},
  ybar, bar width=4pt,
  grid=major, grid style={LightGray, line width=0.3pt},
  axis line style={InkGray, line width=0.5pt},
  enlarge x limits={abs=0.8},
  legend style={font=\scriptsize, at={(0.5,-0.26)}, anchor=north,
                legend columns=2, draw=none, fill=none, column sep=10pt},
]
\addplot[fill=AgriGreen!55, draw=AgriGreen, bar shift=-3.0pt] table[x=i, y=cash]{data/clean/profiles/fin_002481.dat};
\addplot[fill=SoftRed!45, draw=SoftRed, bar shift=3.0pt] table[x=i, y=ibd]{data/clean/profiles/fin_002481.dat};
\addplot[fill=AgriGreen!25, draw=AgriGreen, pattern=north east lines, bar shift=-3.0pt] table[x=x, y=cash]{data/clean/profiles/finh_002481.dat};
\addplot[fill=SoftRed!20, draw=SoftRed, pattern=north east lines, bar shift=3.0pt] table[x=x, y=ibd]{data/clean/profiles/finh_002481.dat};
\legend{货币资金（左轴）, 有息负债（左轴）}
\end{axis}
\begin{axis}[
  at={(finbot.south east)}, anchor=south east,
  width=13.75cm, height=3.10cm, scale only axis,
  axis y line*=right, axis x line*=none, xtick=\empty,
  ymin=-0.84, ymax=4.22,
  tick label style={font=\tiny},
  yticklabel style={font=\tiny}, scaled y ticks=false,
  legend style={font=\scriptsize, at={(0.5,-0.44)}, anchor=north,
                legend columns=1, draw=none, fill=none},
]
\addplot[OilBlue, mark=*, mark size=2.2pt, line width=1.3pt] table[x=i, y=ocf]{data/clean/profiles/fin_002481.dat};
\addplot[only marks, OilBlue, mark=*, mark size=3.2pt] table[x=x, y=ocf]{data/clean/profiles/finh_002481.dat};
\legend{经营活动现金流净额（右轴，亿元，蓝点）}
\end{axis}
\end{tikzpicture}

\captionof{figure}[双塔食品（002481）近年主要财务指标]{双塔食品（002481）近年主要财务指标（左上：营业收入与归母净利润；右上：ROE 与资产负债率；下：货币资金、有息负债为左轴柱状，经营活动现金流净额为其右轴的蓝点折线。
2021---2025 年为年报口径，2026H1 为 2026 年半年报/半年末，与其前一年报之间留出间隔、以斜纹或空心点区分；半年报数据未年化）}
\end{minipage}
\end{center}

\pfl{财务分析} 2026 年 6 月末资产负债率 44.8\%，较 2025 年末下降 2.8 个百分点（样本中位数 52.2\%，所处三级行业内按由低到高排第\zhnumber{8}位）。 2026 年 6 月末货币资金 3.84 亿元、短期债务（短期借款与一年内到期非流动负债）8.95 亿元、长期债务 3.34 亿元，有息负债合计 12.29 亿元，现金短债比 0.43 倍——覆盖偏紧。 净有息负债 8.44 亿元，相当于其 2026 年上半年营业收入的 75\%。 流动比率 1.21、速动比率 0.59，短期指标尚可。 2026 年上半年经营活动现金流净额 1.27 亿元，经营现金流与利润同向为正，内生资金可以支撑运营。 综合看，账面货币资金对有息负债与短期债务的覆盖不足，滚动偿付对新增融资的依赖度较高；资产负债率低于样本中位数 7.4 个百分点；有息负债中短期债务占比更高，期限结构仍以滚动续作为主。 公司披露的股东与质押风险为：2026 年 9 月，累计质押市值 13.98 亿元，控股股东及一致行动人质押 2.46 亿股。

\pfl{重大战略转型} 公司的战略动作可归为以下几类：融资与资本运作（1 项）：2020 年 3 月，发布非公开发行预案，扩张高端豌豆蛋白产能。 公告口径上，近三年（2023 年 9 月---2026 年 9 月）公司披露重大重组与并购类公告 2 条、对外投资与转型类公告 1 条、回购与股权激励类公告 35 条，合计公告 326 条。 主营结构的口径变化已经落到报表上：2021 年年报的第一大行业为“农副食品加工业”（100.0\%），2026 年半年报为“农副产品加工业”（100.0\%）；公司在这两期的分部划分方式有过调整。 综合判断，报告期内公司有 3 条与战略相关的公告落地（重大重组与并购 2 条、对外投资与转型 1 条），资本运作与主业调整之间存在直接关联。

\begin{center}\small\setlength{\tabcolsep}{4pt}
\captionof{table}[双塔食品（002481）历史融资与资本运作]{双塔食品（002481）历史融资与资本运作（据公司公告与公开报道整理；金额为披露值，含首次公开发行、再融资、债券、银行借款与重整投资等）}
\begin{adjustbox}{max width=\textwidth}
\begin{tabular}{@{}p{1.45cm}p{2.55cm}p{4.05cm}p{4.95cm}@{}}
\toprule
时间 & 方式 & 规模 & 用途与说明 \\
\midrule
2010-09-08 & IPO（深圳证券交易所） & 发行 1,\allowbreak{}500 万股，\allowbreak{}发行价 39.80 元/\allowbreak{}股，\allowbreak{}占发行后总股本的 25.00\% & 年产 20,\allowbreak{}000 吨粉丝项目；2010 年 7 月 23 日首发申请获通过；本次发行全部为流通股 \\
2015-01-14 & 非公开发行 A 股（深交所） & 发行 7,\allowbreak{}335.60 万股，\allowbreak{}发行价 17.50 元/\allowbreak{}股 & 新增股份上市首日股价不除权；认购股票限售期十二个月 \\
2020-03-03 & 非公开发行股票预案（未在本底稿核实到实施完成的公开信息） & 拟募集资金总额不超过 21.50 亿元 & 万吨白蛋白提取综合加工项目 5.10 亿元、\allowbreak{}万吨豌豆功能低聚糖提取加工项目 4.45 亿元、\allowbreak{}豌豆精深加工综合利用项目 6.65 亿元、\allowbreak{}补充流动资金 5.30 亿元；2020 年年度报告披露“公司报告期无募集资金使用情况” \\
2020-11—2021 & 股份回购（用于股权激励） & --- & 回购股份用于股权激励；2020 年 11 月 28 日第五届董事会第六次会议审议通过以集中竞价交易方式回购公司股份的方案，\allowbreak{}2020 年 12 月 9 日首次实施回购；2021 年将回购账户中的 969.35 万股授予激励对象 \\
\bottomrule
\end{tabular}
\end{adjustbox}
\end{center}

\pfl{历史融投资情况} 从现金流量表看，2025 年公司取得借款收到的现金 26.46 亿元、偿还债务支付的现金 24.71 亿元、吸收权益性投资 0.01 亿元，筹资活动现金流净额 0.23 亿元（2023 年为取得借款 8.41 亿元、偿还债务 11.03 亿元）。 2026 年上半年筹资活动现金流净额为 -0.97 亿元（取得借款 10.72 亿元、偿还债务 14.41 亿元，未年化），已转为净流出，处于净还债状态。 2025 年分配股利、利润或偿付利息支付的现金合计 0.55 亿元。 公开披露的融资记录共 4 笔，工具类型以 IPO、定向增发为主；金额与用途见下表。 其中规模最大的两笔为：2015 年 1 月，定向增发，发行 7,335.60 万股，发行价 17.50 元/股。 股本方面，近三年披露股本变动 3 次；总股本由 2021 年末的 12.43 亿股变为 12.34 亿股（-0.78\%）。 分红方面，近三年现金分红 4 次、合计每股派现 0.08 元，最近一次实施在 2025 年 6 月。 近三年“再融资”类公告 1 条，涉及定向增发。 综合看，公司融资结构上以债务滚动为主、股权融资为补充；2025 年筹资活动现金流净额 0.23 亿元，年度内仍为净融入，与前述经营与投资现金流的方向基本一致。

\pfl{未来潜在融资需求分析} 2026 年 6 月末货币资金 3.84 亿元，而短期债务 8.95 亿元（现金短债比 0.43 倍），2026 年上半年经营现金流净额 1.27 亿元。按“短期债务减去账面货币资金”的滚动偿付口径，静态缺口约 5.11 亿元；上半年盈利或接近盈亏平衡，该缺口不随经营亏损进一步扩大。 公司 2026 年 6 月末资产负债率 44.8\%、2026 年上半年 ROE 1.4\%（未年化），现金短债比 0.43 倍、短债覆盖不足，融资需求首先用于置换与滚动存量债务、补充营运资金，属于“补血型”而非扩张型。 公开披露的后续资金安排线索包括：股权融资线索：2020 年 3 月，公司拟定增募资不超过 21.50 亿元（含补充流动资金 5.30 亿元），本底稿；债务与授信安排：2026 年 9 月，控股股东及一致行动人质押股份未来半年内到期部分对应融资余额 5,000.00 万元。 综合看，公司的外部资金需求以补充流动性与项目投入为主，滚动偿付口径下的静态缺口约 5.11 亿元，融资节奏将受制于银行续贷意愿与股权融资窗口。




\subsubsection{永顺泰（001338）}
\label{co:001338}

\pfl{公司基本情况} 1987 年设立至今，永顺泰是广东省属国企粤海控股旗下的啤酒麦芽制造商，国内产销规模最大的麦芽制造企业，产能规模居亚洲第一、世界第四；2022 年 11 月登陆深交所主板，业绩受大麦成本与啤酒行业景气度影响显著。 公司于 2022-11-16 在深交所主板首发上市，注册资本 5.02 亿元，总股本 5.02 亿股。 截至 2026 年 10 月 12 日收盘价 9.28 元，对应总市值 46.6 亿元，近一年-19.1\%，总市值在三级行业“其他农产品加工”的 11 家公司中按由大到小排第\zhnumber{5}位，近一年涨跌幅低于全样本中位数（-10.4\%）8.7 个百分点。 按 2026 年半年报披露的行业构成，收入占比前两位为啤酒制造 95.2\%；其他(补充) 4.8\%。

\pfl{股权结构与控制权} 按 2026 年 6 月 30 日的股东名册，前三大股东为广东粤海控股集团有限公司 持股 51.52\%；广州科技金融创新投资控股有限公司 持股 1.85\%；中国建设银行股份有限公司－鹏华中证酒交易型开放式指数 持股 1.67\%，股东户数 28,059 户。 从控制权角度看，广东省人民政府持有粤海集团 90\% 股权、广东省财政厅持有 10\% 股权。

\pfl{经营要点} 经营与产能方面，产能规模位居亚洲第一、世界第四；2026 年上半年，威士忌麦芽和特种麦芽销量创造半年度历史新高，公司总部统筹强化集中采购项目使采购成本进一步下降，并完成永顺泰数字化顶层规划。

\pfl{所处行业的基本面} 该子行业横跨制糖、番茄制品、糖醇与添加剂，供给受种植面积与进口配额约束，需求端分食用与工业两类（第 \ref{sec:framework} 章的供需框架）。 与该公司直接相关的行业数据包括：麦芽行业资金投入、能耗与仓储物流成本较高；2026 年 7 月，我国啤酒产量约 340.3 万千升、同比下降 8.9\%，均低于前三年同期产量。

\pfl{竞争优势} 把该公司放回同类公司比较，2026 年上半年营业收入 19.1 亿元、在三级行业“其他农产品加工”的 11 家公司中排第\zhnumber{4}位；毛利率 8.3\%、低于行业内中位数 0.8 个百分点；半年 ROE 1.4\%（未年化）、高于行业内中位数（1.36\%）；资产负债率 21.1\%、低于行业内中位数（39.1\%）。 与同类公司相比，公司的差异化在于：提高运营效率并降低物流成本。 优势在当期体现为高于行业中位数的资本回报，但优势能否持续仍取决于周期与自身产能投放节奏。

\pfl{历史股价} 以 2022 年 11 月为起点、2026 年 10 月为终点，股价由 21.30 元变为 9.28 元，累计 -56.4\%，区间最高 21.30 元（2022 年 11 月）、最低 8.39 元（2026 年 6 月）；当前价相当于区间最高价的 43.6\%，为区间最低价的 1.11 倍。 月度收益的年化波动率 28.2\%，低于全样本中位数 37.3\%，在三级行业“其他农产品加工”的 11 家公司中波动率由高到低排第\zhnumber{10}位。

\begin{center}
\begin{minipage}{\textwidth}\centering
\begin{tikzpicture}
\begin{axis}[
  width=12.6cm, height=3.9cm, scale only axis,
  xmin=2022.755, xmax=2026.972, ymin=7.80, ymax=22.79,
  xtick={2022,2023,2024,2025,2026},
  yearticks,
  yticklabel style={font=\scriptsize},
  ylabel={元/股}, ylabel style={font=\scriptsize},
  grid=major, grid style={LightGray, line width=0.3pt},
  axis line style={InkGray, line width=0.5pt},
  every axis plot/.append style={line width=0.9pt},
]
\addplot[AgriGreen] table[x=x,y=y]{data/clean/profiles/px_001338.dat};
\addplot[SoftRed, dashed, line width=0.7pt] table[x=x,y=y]{data/clean/profiles/pxzz_001338.dat};
\addplot[only marks, mark=*, mark size=1.1pt, SoftRed] table[x=x,y=y]{data/clean/profiles/pxzz_001338.dat};
\end{axis}
\end{tikzpicture}

\captionof{figure}[永顺泰（001338）股价与周期骨架]{永顺泰（001338）月度收盘价（元/股）与 ZigZag 周期骨架（阈值 25\%，圆点为周期转折点）}
\end{minipage}
\end{center}

\pfl{过去周期分析} 以 25\% 的 ZigZag 阈值划分，样本区间内股价形成 1 个主升段与 2 个主跌段。 最大主跌段为 2022 年 11 月 至 2024 年 8 月，21 个月-55.9\%；最大主升段出现在 2024 年 8 月 至 2026 年 4 月，20 个月+41.2\%。 最近一次转折出现在 2026 年 6 月（下跌段自 2026 年 4 月起，2 个月-36.8\%），当前处于该段的延续之中。 公司股价与申万农林牧渔指数几乎同步见底，个股并未走出独立行情。 该子行业横跨糖、番茄制品、添加剂等多个品种，各自的供给周期与政策（进口配额、关税）是主要变量，公司 2026 年上半年实现盈利、半年 ROE 1.4\%（未年化），好于全样本中位数（0.75\%），下行周期对其报表的冲击小于同业。 综合区间形态看，该股单段涨跌幅度偏大、以趋势段为主（最大主升 +41.2\%、最大主跌 -55.9\%），年化波动率 28.2\% 与全样本中位数 37.3\% 相比波动小于样本中位数。

\pfl{盈利情况} 2026 年上半年营业收入 19.15 亿元（同比 -13.9\%），归母净利润 0.51 亿元，同比 -69.6\%。 以年报为趋势对照，2023---2025 年营业收入由 48.40 亿元变为 39.27 亿元（两年年均复合增速 -9.9\%），归母净利润由 1.74 亿元变为 2.37 亿元。 按半年报口径，2026 年上半年的 ROE 为 1.4\%（未年化）、销售净利率 2.7\%、毛利率 8.3\%，ROE 在“其他农产品加工”11 家公司中按数值由高到低排第\zhnumber{5}位。 按同一口径，三级行业“其他农产品加工”11 家公司的半年 ROE 中位数为 1.36\%，该公司高于其中位数 0.07 个百分点。 综合看，公司在报告期维持盈利（高于全样本半年 ROE 中位数 0.75\%），利润的可持续性取决于其他农产品加工景气的延续与成本管控的执行。

\begin{center}\small\setlength{\tabcolsep}{3.4pt}
\captionof{table}[永顺泰（001338）关键财务指标]{永顺泰（001338）关键财务指标（合并报表口径；两个半年报期均未年化；现金短债比 $=$ 货币资金 $\div$ 短期债务）}
\begin{adjustbox}{max width=\textwidth}
\begin{tabular}{@{}lrrrrr@{}}
\toprule
指标 & 2023 年 & 2024 年 & 2025 年 & 2025 年半年报 & 2026 年半年报 \\
\midrule
营业收入（亿元） & 48.40 & 42.82 & 39.27 & 22.23 & 19.15 \\
归母净利润（亿元） & 1.74 & 2.99 & 2.37 & 1.69 & 0.51 \\
毛利率（\%） & 8.4 & 11.9 & 12.1 & 12.0 & 8.3 \\
ROE（\%） & 5.3 & 8.6 & 6.5 & 4.8 & 1.4 \\
资产负债率（\%） & 25.3 & 14.8 & 10.0 & 15.0 & 21.1 \\
经营活动现金流净额（亿元） & 30.85 & 24.44 & 15.61 & 11.56 & 2.18 \\
货币资金（亿元） & 5.53 & 8.95 & 9.75 & 9.65 & 14.14 \\
有息负债合计（亿元） & 2.20 & 0.64 & 0.04 & 2.28 & 7.41 \\
现金短债比（倍） & 2.52 & 16.51 & 1,606.40 & 4.36 & 1.92 \\
\bottomrule
\end{tabular}
\end{adjustbox}
\end{center}

\begin{center}
\begin{minipage}{\textwidth}\centering
\begin{tikzpicture}
\begin{groupplot}[
  group style={group size=2 by 1, horizontal sep=0.60cm},
  width=6.85cm, height=4.60cm,
  tick label style={font=\tiny},
  title style={font=\scriptsize\heiti},
  yticklabel style={font=\tiny},
  ylabel style={font=\tiny},
  grid=major, grid style={LightGray, line width=0.3pt},
  axis line style={InkGray, line width=0.5pt},
  enlarge x limits={abs=0.8},
  legend style={font=\scriptsize, at={(0.5,-0.20)}, anchor=north,
                legend columns=2, draw=none, fill=none, column sep=6pt},
]
\nextgroupplot[
  ybar, bar width=4pt,
  ymin=-4.84, ymax=54.21,
  xtick={0,1,2,3,4,5.7},
  xticklabels={2021,2022,2023,2024,2025,2026H1},
  ylabel={亿元},
]
\addplot[fill=AgriGreen!75, draw=AgriGreen, bar shift=-2.6pt] table[x=i, y=rev]{data/clean/profiles/fin_001338.dat};
\addplot[fill=SoftRed!60, draw=SoftRed, bar shift=2.6pt] table[x=i, y=ni]{data/clean/profiles/fin_001338.dat};
\addplot[fill=AgriGreen!30, draw=AgriGreen, pattern=north east lines, bar shift=-2.6pt] table[x=x, y=rev]{data/clean/profiles/finh_001338.dat};
\addplot[fill=SoftRed!20, draw=SoftRed, pattern=north east lines, bar shift=2.6pt] table[x=x, y=ni]{data/clean/profiles/finh_001338.dat};
\legend{营业收入（年/半年）, 归母净利润（年/半年）}
\nextgroupplot[
  ymin=-2.2, ymax=29.6,
  xtick={0,1,2,3,4,5.7},
  xticklabels={2021,2022,2023,2024,2025,2026H1},
  ylabel={\%},
]
\addplot[AgriGold, mark=*, mark size=1.3pt] table[x=i, y=roe]{data/clean/profiles/fin_001338.dat};
\addplot[OilBlue, mark=square*, mark size=1.3pt] table[x=i, y=debt]{data/clean/profiles/fin_001338.dat};
\addplot[only marks, AgriGold, mark=*, mark size=2.0pt] table[x=x, y=roe]{data/clean/profiles/finh_001338.dat};
\addplot[only marks, OilBlue, mark=square*, mark size=2.0pt] table[x=x, y=debt]{data/clean/profiles/finh_001338.dat};
\legend{ROE, 资产负债率}
\end{groupplot}
\end{tikzpicture}
\begin{tikzpicture}
\begin{axis}[
  name=finbot,
  width=13.75cm, height=3.10cm, scale only axis,
  ymin=-1.41, ymax=15.84,
  xtick={0,1,2,3,4,5.7},
  xticklabels={2021,2022,2023,2024,2025,2026H1},
  tick label style={font=\tiny},
  yticklabel style={font=\tiny},
  ylabel={亿元}, ylabel style={font=\tiny},
  ybar, bar width=4pt,
  grid=major, grid style={LightGray, line width=0.3pt},
  axis line style={InkGray, line width=0.5pt},
  enlarge x limits={abs=0.8},
  legend style={font=\scriptsize, at={(0.5,-0.26)}, anchor=north,
                legend columns=2, draw=none, fill=none, column sep=10pt},
]
\addplot[fill=AgriGreen!55, draw=AgriGreen, bar shift=-3.0pt] table[x=i, y=cash]{data/clean/profiles/fin_001338.dat};
\addplot[fill=SoftRed!45, draw=SoftRed, bar shift=3.0pt] table[x=i, y=ibd]{data/clean/profiles/fin_001338.dat};
\addplot[fill=AgriGreen!25, draw=AgriGreen, pattern=north east lines, bar shift=-3.0pt] table[x=x, y=cash]{data/clean/profiles/finh_001338.dat};
\addplot[fill=SoftRed!20, draw=SoftRed, pattern=north east lines, bar shift=3.0pt] table[x=x, y=ibd]{data/clean/profiles/finh_001338.dat};
\legend{货币资金（左轴）, 有息负债（左轴）}
\end{axis}
\begin{axis}[
  at={(finbot.south east)}, anchor=south east,
  width=13.75cm, height=3.10cm, scale only axis,
  axis y line*=right, axis x line*=none, xtick=\empty,
  ymin=-8.02, ymax=40.10,
  tick label style={font=\tiny},
  yticklabel style={font=\tiny}, scaled y ticks=false,
  legend style={font=\scriptsize, at={(0.5,-0.44)}, anchor=north,
                legend columns=1, draw=none, fill=none},
]
\addplot[OilBlue, mark=*, mark size=2.2pt, line width=1.3pt] table[x=i, y=ocf]{data/clean/profiles/fin_001338.dat};
\addplot[only marks, OilBlue, mark=*, mark size=3.2pt] table[x=x, y=ocf]{data/clean/profiles/finh_001338.dat};
\legend{经营活动现金流净额（右轴，亿元，蓝点）}
\end{axis}
\end{tikzpicture}

\captionof{figure}[永顺泰（001338）近年主要财务指标]{永顺泰（001338）近年主要财务指标（左上：营业收入与归母净利润；右上：ROE 与资产负债率；下：货币资金、有息负债为左轴柱状，经营活动现金流净额为其右轴的蓝点折线。
2021---2025 年为年报口径，2026H1 为 2026 年半年报/半年末，与其前一年报之间留出间隔、以斜纹或空心点区分；半年报数据未年化）}
\end{minipage}
\end{center}

\pfl{财务分析} 2026 年 6 月末资产负债率 21.1\%，较 2025 年末上升 11.1 个百分点（样本中位数 52.2\%，所处三级行业内按由低到高排第\zhnumber{2}位）。 2026 年 6 月末货币资金 14.14 亿元、短期债务（短期借款与一年内到期非流动负债）7.38 亿元、长期债务 0.03 亿元，有息负债合计 7.41 亿元，现金短债比 1.92 倍——覆盖充分。 净有息负债 -6.73 亿元，相当于其 2026 年上半年营业收入的 -35\%。 流动比率 3.43、速动比率 2.39，短期偿债能力处于正常区间。 2026 年上半年经营活动现金流净额 2.18 亿元，经营现金流与利润同向为正，内生资金可以支撑运营。 2026 年 6 月末在建工程 1.33 亿元，较上年末增加 0.05 亿元，仍有在建投入。 综合看，现金短债比 1.92 倍，短期偿付压力可控；资产负债率低于样本中位数 31.1 个百分点；有息负债中短期债务占比更高，期限结构仍以滚动续作为主。

\pfl{重大战略转型} 近三年公司的战略动作集中在：产能与项目（3 项）：2025---2026 年，公司推进宝麦公司新建 5 万吨/年特制麦芽生产线项目；2025 年 6 月，2025 年上半年麦芽产品产能 56.92 万吨（同比+8.27\%）。 公告口径上，近三年（2023 年 9 月---2026 年 9 月）公司披露重大重组与并购类公告 0 条、对外投资与转型类公告 0 条、回购与股权激励类公告 0 条，合计公告 283 条。 第一大行业仍是“啤酒制造”，收入占比 95.2\%，与 2021 年年报相比没有方向性变化。 综合判断，报告期内公司未披露实质性的重组、并购或转型安排，资本开支以维持现有产能与业务为主。

\begin{center}\small\setlength{\tabcolsep}{4pt}
\captionof{table}[永顺泰（001338）历史融资与资本运作]{永顺泰（001338）历史融资与资本运作（据公司公告与公开报道整理；金额为披露值，含首次公开发行、再融资、债券、银行借款与重整投资等）}
\begin{adjustbox}{max width=\textwidth}
\begin{tabular}{@{}p{1.45cm}p{2.55cm}p{4.05cm}p{4.95cm}@{}}
\toprule
时间 & 方式 & 规模 & 用途与说明 \\
\midrule
2022-11 & IPO募投项目明细 & 广麦 4 期扩建项目投资总额 5.03 亿元、\allowbreak{}拟投入募集资金 4.25 亿元；年产 13 万吨中高档啤酒麦芽项目投资总额 5.22 亿元、\allowbreak{}拟投入募集资金 3.70 亿元；合计投资总额 10.25 亿元、\allowbreak{}拟投入募集资金 7.94 亿元 & 产能扩张；年产 13 万吨中高档啤酒麦芽项目实施地点为江苏省扬州市宝应县安宜工业园 \\
2022-11-16 & IPO（深交所主板） & 发行 1.25 亿股，\allowbreak{}发行价 6.82 元/\allowbreak{}股 & 广麦 4 期扩建项目、\allowbreak{}年产 13 万吨中高档啤酒麦芽项目；；募集资金于 2022 年 11 月 8 日到位；发行后总股本 5.02 亿股 \\
2023-12-08 & 募投项目结项 & 首次公开发行募投项目（广麦 4 期扩建项目、\allowbreak{}年产 13 万吨中高档啤酒麦芽项目）结项 & 补充流动资金；第二届董事会第三次会议、\allowbreak{}第二届监事会第三次会议审议通过（节余金额原文被截断 \\
2026-06-30 & 闲置募集资金现金管理 & 不超过人民币 1,\allowbreak{}450.00 万元 & 购买现金管理产品；使用期限为自董事会审议通过之日起 12 个月内，\allowbreak{}可循环滚动使用；截至 2025 年 12 月 31 日累计使用募集资金 7.94 亿元 \\
\bottomrule
\end{tabular}
\end{adjustbox}
\end{center}

\pfl{历史融投资情况} 从现金流量表看，2025 年公司取得借款收到的现金 0.20 亿元、偿还债务支付的现金 11.19 亿元，筹资活动现金流净额 -12.14 亿元（2023 年为取得借款 16.02 亿元、偿还债务 45.79 亿元）。 2026 年上半年筹资活动现金流净额为 2.37 亿元（取得借款 — 亿元、偿还债务 2.28 亿元，未年化），仍处于净融入状态。 2025 年分配股利、利润或偿付利息支付的现金合计 1.57 亿元。 公开披露的融资记录共 4 笔，工具类型以 IPO 为主；金额与用途见下表。 其中规模最大的两笔为：2026 年 6 月，闲置募集资金现金管理，不超过人民币 1,450.00 万元；2022 年 11 月，IPO，广麦 4 期扩建项目投资总额 5.03 亿元、拟投入募集资金 4.25 亿元，年产 13 万吨中高档啤酒麦芽项目投资总额 5.22 亿元、拟投入募集资金 3.70 亿元，合计投资总额 10.25 亿元、拟投入募集资金 7.94 亿元。 股本方面，近三年披露股本变动 4 次（限售股份上市 1 次）；总股本自 2021 年末以来未发生变化（5.02 亿股）。 分红方面，近三年现金分红 4 次、合计每股派现 0.71 元，最近一次实施在 2025 年 12 月。 近三年“再融资”类公告 1 条，涉及定向增发。 综合看，公司融资结构上股权与债务融资并举；2025 年筹资活动现金流净额 -12.14 亿元，融资节奏仍以维持存量债务周转为主，与前述经营与投资现金流的方向基本一致。

\pfl{未来潜在融资需求分析} 2026 年 6 月末货币资金 14.14 亿元，而短期债务 7.38 亿元（现金短债比 1.92 倍），2026 年上半年经营现金流净额 2.18 亿元。账面货币资金可以覆盖短期债务，缺口不体现在静态偿付层面。 公司 2026 年 6 月末资产负债率 21.1\%、2026 年上半年 ROE 1.4\%（未年化），资产负债表尚有余量，融资需求不是“补血型”的刚性需求，而是配套在建项目投入的主动性安排。 资金需求还要叠加在建工程：期末在建工程 1.33 亿元、2026 年上半年购建固定资产等支出 0.35 亿元（未年化），这部分支出在周期底部通常会被主动放缓。 以现有货币资金与经营现金流衡量，公司不存在刚性补血的紧迫性，融资动作更多是主动型的。




\subsubsection{京粮控股（000505）}
\label{co:000505}

\pfl{公司基本情况} 京粮控股成立于 1992 年，北京市国资委系统（首农食品集团—京粮集团）旗下油脂油料加工与食品加工上市平台，由原海南珠江控股 2017 年重大资产置换并注入京粮食品而来；油脂贸易体量大但毛利薄，2025 年因贸易收缩与原料成本上升转为大额亏损。 公司于 1992-12-21 在深交所主板首发上市，注册资本 7.27 亿元，总股本 7.27 亿股。 截至 2026 年 10 月 12 日收盘价 6.29 元，对应总市值 45.7 亿元，近一年-0.5\%，总市值在三级行业“粮油加工”的 7 家公司中按由大到小排第\zhnumber{4}位，近一年涨跌幅高于全样本中位数（-10.4\%）9.9 个百分点。 按 2026 年半年报披露的行业构成，收入占比前三位为油脂 89.2\%；食品 10.5\%；其他(补充) 0.3\%。 股权结构方面，控股股东为广州珠江实业集团有限公司。

\pfl{股权结构与控制权} 按 2026 年 6 月 30 日的股东名册，前三大股东为北京粮食集团有限责任公司 持股 39.68\%；北京国有资本运营管理有限公司 持股 6.67\%；王岳成 持股 4.74\%，股东户数 51,588 户。 从控制权角度看，第三大股东王岳成持股 4.74\%。

\pfl{经营要点} 经营与产能方面，部分产品销量下滑。

\pfl{所处行业的基本面} 粮油加工的供给由大豆、菜籽等原料进口与压榨产能决定，需求由食用消费决定，利润来自原料与成品的价差而非价格方向（第 \ref{sec:framework} 章的供需框架）。 行业层面的最新情况是：2026 年，1—8 月中国累计进口大豆 7,410.7 万吨、同比增长 1.1\%，进口大豆到港成本持续上行、榨利持续亏损。

\pfl{竞争优势} 把该公司放回同类公司比较，2026 年上半年营业收入 31.1 亿元、在三级行业“粮油加工”的 7 家公司中排第\zhnumber{2}位；毛利率 5.6\%、低于行业内中位数 0.7 个百分点；半年 ROE 0.3\%（未年化）、低于行业内中位数（2.22\%）；资产负债率 43.5\%、高于行业内中位数（43.5\%）。 公司区别于同类公司的经营特征包括：核心加工产能集中于京粮天津；公司能够以较低成本高效完成大豆和大豆油的储备轮换任务，保障工厂开机率并获得稳定压榨利润；产品覆盖北京、天津、河北千余家门店，在京津冀地区具备品牌与市场占有率优势。 就当期报表看，公司尚未把经营层面的特征转化为高于行业的资本回报，周期底部的相对优势更多体现在现金流与资产负债表上。

\pfl{历史股价} 以 2021 年 1 月为起点、2026 年 10 月为终点，股价由 7.12 元变为 6.29 元，累计 -11.7\%，区间最高 12.31 元（2022 年 3 月）、最低 4.97 元（2024 年 8 月）；当前价相当于区间最高价的 51.1\%，为区间最低价的 1.27 倍。 月度收益的年化波动率 39.7\%，高于全样本中位数 37.3\%，在三级行业“粮油加工”的 7 家公司中波动率由高到低排第\zhnumber{2}位。

\begin{center}
\begin{minipage}{\textwidth}\centering
\begin{tikzpicture}
\begin{axis}[
  width=12.6cm, height=3.9cm, scale only axis,
  xmin=2020.922, xmax=2026.972, ymin=4.62, ymax=13.17,
  xtick={2021,2022,2023,2024,2025,2026},
  yearticks,
  yticklabel style={font=\scriptsize},
  ylabel={元/股}, ylabel style={font=\scriptsize},
  grid=major, grid style={LightGray, line width=0.3pt},
  axis line style={InkGray, line width=0.5pt},
  every axis plot/.append style={line width=0.9pt},
]
\addplot[AgriGreen] table[x=x,y=y]{data/clean/profiles/px_000505.dat};
\addplot[SoftRed, dashed, line width=0.7pt] table[x=x,y=y]{data/clean/profiles/pxzz_000505.dat};
\addplot[only marks, mark=*, mark size=1.1pt, SoftRed] table[x=x,y=y]{data/clean/profiles/pxzz_000505.dat};
\end{axis}
\end{tikzpicture}

\captionof{figure}[京粮控股（000505）股价与周期骨架]{京粮控股（000505）月度收盘价（元/股）与 ZigZag 周期骨架（阈值 25\%，圆点为周期转折点）}
\end{minipage}
\end{center}

\pfl{过去周期分析} 以 25\% 的 ZigZag 阈值划分，样本区间内股价形成 3 个主升段与 2 个主跌段。 最大主跌段为 2022 年 3 月 至 2024 年 8 月，29 个月-59.6\%；最大主升段出现在 2021 年 1 月 至 2022 年 3 月，14 个月+72.9\%。 最近一次转折出现在 2026 年 8 月（上涨段自 2026 年 6 月起，2 个月+40.0\%），当前处于该段之后的震荡之中。 公司股价较申万农林牧渔指数提前约 21 个月见底，是板块中较早定价周期反转的公司之一。 粮油加工是“成本加成”型生意（第 \ref{sec:grain} 章），利润来自原料与成品的价差及规模效率，收入规模大而净利率薄，公司 2026 年 6 月末资产负债率 43.5\%、半年 ROE 0.3\%（未年化），在本轮下行中属于随行业同步调整的一类。 综合区间形态看，该股单段涨跌幅度偏大、以趋势段为主（最大主升 +72.9\%、最大主跌 -59.6\%），年化波动率 39.7\% 与全样本中位数 37.3\% 相比波动与样本中位数接近。

\pfl{盈利情况} 2026 年上半年营业收入 31.15 亿元（同比 -26.0\%），归母净利润 0.08 亿元，同比 -55.3\%。 以年报为趋势对照，2023---2025 年营业收入由 116.02 亿元变为 78.59 亿元（两年年均复合增速 -17.7\%），归母净利润由 1.01 亿元变为 -2.66 亿元。 按半年报口径，2026 年上半年的 ROE 为 0.3\%（未年化）、销售净利率 0.3\%、毛利率 5.6\%，ROE 在“粮油加工”7 家公司中按数值由高到低排第\zhnumber{6}位。 同期全样本半年 ROE 中位数 0.75\%，公司与之相差 0.47 个百分点（弱于中位数公司）。 从公司披露的原因看，业绩变动主要来自：2025 年 12 月，公司亏损主因：油脂贸易业务收缩、南美大豆原料品质波动推高加工成本，及休闲食品传统渠道下滑致毛利承压；2026 年 8 月，净利润下降主要受子公司新增食品加工生产线尚处建设期及休闲食品线下渠道受低价竞争冲击等因素影响。 综合看，公司在报告期维持盈利（低于全样本半年 ROE 中位数 0.75\%），利润的可持续性取决于粮油加工景气的延续与成本管控的执行。

\begin{center}\small\setlength{\tabcolsep}{3.4pt}
\captionof{table}[京粮控股（000505）关键财务指标]{京粮控股（000505）关键财务指标（合并报表口径；两个半年报期均未年化；现金短债比 $=$ 货币资金 $\div$ 短期债务）}
\begin{adjustbox}{max width=\textwidth}
\begin{tabular}{@{}lrrrrr@{}}
\toprule
指标 & 2023 年 & 2024 年 & 2025 年 & 2025 年半年报 & 2026 年半年报 \\
\midrule
营业收入（亿元） & 116.02 & 114.35 & 78.59 & 42.08 & 31.15 \\
归母净利润（亿元） & 1.01 & 0.26 & -2.66 & 0.18 & 0.08 \\
毛利率（\%） & 3.0 & 4.5 & 4.5 & 5.6 & 5.6 \\
ROE（\%） & 3.3 & 0.8 & -8.9 & 0.6 & 0.3 \\
资产负债率（\%） & 45.0 & 47.6 & 47.5 & 53.2 & 43.5 \\
经营活动现金流净额（亿元） & 1.09 & -1.09 & 6.47 & 2.22 & 3.90 \\
货币资金（亿元） & 15.43 & 14.17 & 18.22 & 17.61 & 16.90 \\
有息负债合计（亿元） & 21.11 & 22.05 & 22.25 & 23.49 & 17.91 \\
现金短债比（倍） & 1.15 & 0.76 & 1.20 & 1.19 & 1.55 \\
\bottomrule
\end{tabular}
\end{adjustbox}
\end{center}

\begin{center}
\begin{minipage}{\textwidth}\centering
\begin{tikzpicture}
\begin{groupplot}[
  group style={group size=2 by 1, horizontal sep=0.60cm},
  width=6.85cm, height=4.60cm,
  tick label style={font=\tiny},
  title style={font=\scriptsize\heiti},
  yticklabel style={font=\tiny},
  ylabel style={font=\tiny},
  grid=major, grid style={LightGray, line width=0.3pt},
  axis line style={InkGray, line width=0.5pt},
  enlarge x limits={abs=0.8},
  legend style={font=\scriptsize, at={(0.5,-0.20)}, anchor=north,
                legend columns=2, draw=none, fill=none, column sep=6pt},
]
\nextgroupplot[
  ybar, bar width=4pt,
  ymin=-15.78, ymax=144.33,
  xtick={0,1,2,3,4,5.7},
  xticklabels={2021,2022,2023,2024,2025,2026H1},
  ylabel={亿元},
]
\addplot[fill=AgriGreen!75, draw=AgriGreen, bar shift=-2.6pt] table[x=i, y=rev]{data/clean/profiles/fin_000505.dat};
\addplot[fill=SoftRed!60, draw=SoftRed, bar shift=2.6pt] table[x=i, y=ni]{data/clean/profiles/fin_000505.dat};
\addplot[fill=AgriGreen!30, draw=AgriGreen, pattern=north east lines, bar shift=-2.6pt] table[x=x, y=rev]{data/clean/profiles/finh_000505.dat};
\addplot[fill=SoftRed!20, draw=SoftRed, pattern=north east lines, bar shift=2.6pt] table[x=x, y=ni]{data/clean/profiles/finh_000505.dat};
\legend{营业收入（年/半年）, 归母净利润（年/半年）}
\nextgroupplot[
  ymin=-13.4, ymax=53.3,
  xtick={0,1,2,3,4,5.7},
  xticklabels={2021,2022,2023,2024,2025,2026H1},
  ylabel={\%},
]
\addplot[AgriGold, mark=*, mark size=1.3pt] table[x=i, y=roe]{data/clean/profiles/fin_000505.dat};
\addplot[OilBlue, mark=square*, mark size=1.3pt] table[x=i, y=debt]{data/clean/profiles/fin_000505.dat};
\addplot[only marks, AgriGold, mark=*, mark size=2.0pt] table[x=x, y=roe]{data/clean/profiles/finh_000505.dat};
\addplot[only marks, OilBlue, mark=square*, mark size=2.0pt] table[x=x, y=debt]{data/clean/profiles/finh_000505.dat};
\legend{ROE, 资产负债率}
\end{groupplot}
\end{tikzpicture}
\begin{tikzpicture}
\begin{axis}[
  name=finbot,
  width=13.75cm, height=3.10cm, scale only axis,
  ymin=-2.22, ymax=24.91,
  xtick={0,1,2,3,4,5.7},
  xticklabels={2021,2022,2023,2024,2025,2026H1},
  tick label style={font=\tiny},
  yticklabel style={font=\tiny},
  ylabel={亿元}, ylabel style={font=\tiny},
  ybar, bar width=4pt,
  grid=major, grid style={LightGray, line width=0.3pt},
  axis line style={InkGray, line width=0.5pt},
  enlarge x limits={abs=0.8},
  legend style={font=\scriptsize, at={(0.5,-0.26)}, anchor=north,
                legend columns=2, draw=none, fill=none, column sep=10pt},
]
\addplot[fill=AgriGreen!55, draw=AgriGreen, bar shift=-3.0pt] table[x=i, y=cash]{data/clean/profiles/fin_000505.dat};
\addplot[fill=SoftRed!45, draw=SoftRed, bar shift=3.0pt] table[x=i, y=ibd]{data/clean/profiles/fin_000505.dat};
\addplot[fill=AgriGreen!25, draw=AgriGreen, pattern=north east lines, bar shift=-3.0pt] table[x=x, y=cash]{data/clean/profiles/finh_000505.dat};
\addplot[fill=SoftRed!20, draw=SoftRed, pattern=north east lines, bar shift=3.0pt] table[x=x, y=ibd]{data/clean/profiles/finh_000505.dat};
\legend{货币资金（左轴）, 有息负债（左轴）}
\end{axis}
\begin{axis}[
  at={(finbot.south east)}, anchor=south east,
  width=13.75cm, height=3.10cm, scale only axis,
  axis y line*=right, axis x line*=none, xtick=\empty,
  ymin=-8.40, ymax=10.01,
  tick label style={font=\tiny},
  yticklabel style={font=\tiny}, scaled y ticks=false,
  legend style={font=\scriptsize, at={(0.5,-0.44)}, anchor=north,
                legend columns=1, draw=none, fill=none},
]
\addplot[OilBlue, mark=*, mark size=2.2pt, line width=1.3pt] table[x=i, y=ocf]{data/clean/profiles/fin_000505.dat};
\addplot[only marks, OilBlue, mark=*, mark size=3.2pt] table[x=x, y=ocf]{data/clean/profiles/finh_000505.dat};
\legend{经营活动现金流净额（右轴，亿元，蓝点）}
\end{axis}
\end{tikzpicture}

\captionof{figure}[京粮控股（000505）近年主要财务指标]{京粮控股（000505）近年主要财务指标（左上：营业收入与归母净利润；右上：ROE 与资产负债率；下：货币资金、有息负债为左轴柱状，经营活动现金流净额为其右轴的蓝点折线。
2021---2025 年为年报口径，2026H1 为 2026 年半年报/半年末，与其前一年报之间留出间隔、以斜纹或空心点区分；半年报数据未年化）}
\end{minipage}
\end{center}

\pfl{财务分析} 2026 年 6 月末资产负债率 43.5\%，较 2025 年末下降 4.0 个百分点（样本中位数 52.2\%，所处三级行业内按由低到高排第\zhnumber{4}位）。 2026 年 6 月末货币资金 16.90 亿元、短期债务（短期借款与一年内到期非流动负债）10.88 亿元、长期债务 7.03 亿元，有息负债合计 17.91 亿元，现金短债比 1.55 倍——覆盖充分。 净有息负债 1.01 亿元，相当于其 2026 年上半年营业收入的 3\%。 流动比率 2.26、速动比率 1.32，短期流动性无明显缺口。 2026 年上半年经营活动现金流净额 3.90 亿元，经营现金流与利润同向为正，内生资金可以支撑运营。 2026 年 6 月末在建工程 1.31 亿元，较上年末增加 0.42 亿元，仍有在建投入。 2026 年 6 月末商誉 1.26 亿元，占归母净资产的 3.9\%，是净资产中最脆弱的一块。 综合看，现金短债比 1.55 倍，短期偿付压力可控；资产负债率低于样本中位数 8.7 个百分点；有息负债中短期债务占比更高，期限结构仍以滚动续作为主。

\pfl{重大战略转型} 公司的战略动作可归为以下几类：其他动作（3 项）：2026 年上半年，归属于上市公司股东的净利润 803.24 万元；2026 年 6 月，加权平均净资产收益率 0.28\%（上年同期 0.57\%）。 公告口径上，近三年（2023 年 9 月---2026 年 9 月）公司披露重大重组与并购类公告 2 条、对外投资与转型类公告 0 条、回购与股权激励类公告 0 条，合计公告 335 条。 第一大行业仍是“油脂”，收入占比 89.2\%，与 2021 年年报相比没有方向性变化。 综合判断，报告期内公司有 2 条与战略相关的公告落地（重大重组与并购 2 条、对外投资与转型 0 条），资本运作与主业调整之间存在直接关联。

\begin{center}\small\setlength{\tabcolsep}{4pt}
\captionof{table}[京粮控股（000505）历史融资与资本运作]{京粮控股（000505）历史融资与资本运作（据公司公告与公开报道整理；金额为披露值，含首次公开发行、再融资、债券、银行借款与重整投资等）}
\begin{adjustbox}{max width=\textwidth}
\begin{tabular}{@{}p{1.45cm}p{2.55cm}p{4.05cm}p{4.95cm}@{}}
\toprule
时间 & 方式 & 规模 & 用途与说明 \\
\midrule
2017 & 发行股份募集配套资金 & 非公开发行不超过 4,\allowbreak{}896.54 万股，\allowbreak{}发行价 8.82 元/\allowbreak{}股 & 配套资金（招股/\allowbreak{}重组方案约定用途）；股东京粮集团以货币资金认购 \\
2017-09-14 & 重大资产置换＋发行股份购买资产（借壳上市） & 京粮股份（北京京粮食品有限公司）100\% 股权作价 23.09 亿元 & 资产置换、\allowbreak{}购买京粮食品 100\% 股权、\allowbreak{}配套资金；新增股份上市日期为 2017 年 11 月 15 日；发行完成后注册资本 6.86 亿元 \\
2023-08-21 & 公开发行公司债券（第一期） & 发行总额 3.00 亿元（3.00 亿元），\allowbreak{}期限 3 年，\allowbreak{}票面利率 2.88\% & 偿还到期债务和补充公司流动资金；起息日 2023-08-22 \\
2023-11-21 & 变更募集资金用途 & --- & 变更募集资金用途（2023 年 11 月 22 日披露变更募集资金用途编号 2023-046）；募投项目实施方式调整 \\
\bottomrule
\end{tabular}
\end{adjustbox}
\end{center}

\pfl{历史融投资情况} 从现金流量表看，2025 年公司取得借款收到的现金 41.83 亿元、偿还债务支付的现金 42.42 亿元、吸收权益性投资 0.01 亿元，筹资活动现金流净额 -1.79 亿元（2023 年为取得借款 33.59 亿元、偿还债务 27.96 亿元）。 2026 年上半年筹资活动现金流净额为 -4.57 亿元（取得借款 7.24 亿元、偿还债务 11.54 亿元，未年化），已转为净流出，处于净还债状态。 2025 年分配股利、利润或偿付利息支付的现金合计 0.73 亿元。 公开披露的融资记录共 4 笔，工具类型以发行股份募集配套资金、发行股份购买资产、公司债为主；金额与用途见下表。 其中规模最大的两笔为：2023 年 8 月，公司债，发行总额 3.00 亿元（3.00 亿元），期限 3 年，票面利率 2.88\%；2017 年，发行股份募集配套资金，非公开发行不超过 4,896.54 万股，发行价 8.82 元/股。 股本方面，近三年披露股本变动 4 次（限售股份上市 1 次）；总股本自 2021 年末以来未发生变化（7.27 亿股）。 分红方面，近三年现金分红 2 次、合计每股派现 0.09 元，最近一次实施在 2024 年 12 月。 综合看，公司融资结构上以债务滚动为主、股权融资为补充；2025 年筹资活动现金流净额 -1.79 亿元，融资节奏仍以维持存量债务周转为主，与前述经营与投资现金流的方向基本一致。

\pfl{未来潜在融资需求分析} 2026 年 6 月末货币资金 16.90 亿元，而短期债务 10.88 亿元（现金短债比 1.55 倍），2026 年上半年经营现金流净额 3.90 亿元。账面货币资金可以覆盖短期债务，缺口不体现在静态偿付层面。 公司 2026 年 6 月末资产负债率 43.5\%、2026 年上半年 ROE 0.3\%（未年化），资产负债表尚有余量，融资需求不是“补血型”的刚性需求，而是配套在建项目投入的主动性安排。 公开披露的后续资金安排线索包括：债务与授信安排：2026 年 8 月，“23 京粮 01”；资本开支安排：2026 年 8 月，子公司新增食品加工生产线仍在建设期。 资金需求还要叠加在建工程：期末在建工程 1.31 亿元、2026 年上半年购建固定资产等支出 0.76 亿元（未年化），这部分支出在周期底部通常会被主动放缓。 以现有货币资金与经营现金流衡量，公司不存在刚性补血的紧迫性，融资动作更多是主动型的。




\subsubsection{华资实业（600191）}
\label{co:600191}

\pfl{公司基本情况} 1998 年设立至今，华资实业是包头起家的甜菜制糖老牌上市公司，曾是全国甜菜制糖行业唯一上市公司，长期依靠参股华夏银行、恒泰证券（现金融街证券）的投资收益支撑业绩；近年主业转为谷朊粉、食用酒精、黄原胶等粮食深加工与生物发酵，2025 年由山东中裕食品通过增资控股股东间接入主。 公司于 1998-12-10 在上交所首发上市，注册资本 4.85 亿元，总股本 4.85 亿股。 截至 2026 年 10 月 12 日收盘价 9.07 元，对应总市值 44.0 亿元，近一年+15.8\%，总市值在三级行业“其他农产品加工”的 11 家公司中按由大到小排第\zhnumber{6}位，近一年涨跌幅高于全样本中位数（-10.4\%）26.2 个百分点。 按 2025 年年报披露的行业构成，收入占比前三位为食品制造业 91.3\%；农副食品加工业 7.7\%；其他(补充) 1.0\%。 与 2021 年年报相比，第一大行业口径由“制糖业”（84.2\%）变为“食品制造业”（91.3\%）；两期披露的分类口径有调整，占比差异中同时包含分类因素。 股权结构方面，控股股东为海南盛泰创发实业有限公司。

\pfl{股权结构与控制权} 按 2026 年 6 月 30 日的股东名册，前三大股东为海南盛泰创发实业有限公司 持股 29.90\%；世通投资(山东)有限公司 持股 17.61\%；广东南传私募基金管理有限公司－南传烽火 5 号私募证券投 持股 5.82\%，股东户数 20,445 户。 控制权与股东性质方面，2025 年 4 月，中裕科技投资以现金增资取得盛泰创发 55\% 股权并成为其控股股东，公司实际控制人变更为宋民松；中裕科技投资、滨州中裕食品及实际控制人宋民松于 2025 年 4 月 15 日承诺在三十六个月内通过资产转让、委托管理、剥离等方式消除同业竞争。

\pfl{经营要点} 经营层面，加工原糖达到全年满负荷生产需要原糖 20 万至 24 万吨，而公司进口配额仅 8,000 吨左右。

\pfl{所处行业的基本面} 该子行业横跨制糖、番茄制品、糖醇与添加剂，供给受种植面积与进口配额约束，需求端分食用与工业两类（第 \ref{sec:framework} 章的供需框架）。 与该公司直接相关的行业数据包括：2026 年上半年，公司以 95\% 乙醇为代表的主流食用酒精市场呈现工艺分化、结构优化、价格震荡企稳的运行特征。

\pfl{竞争优势} 把该公司放回同类公司比较，2026 年上半年营业收入 3.5 亿元、在三级行业“其他农产品加工”的 11 家公司中排第\zhnumber{10}位；毛利率 4.6\%、低于行业内中位数 4.4 个百分点；半年 ROE 2.1\%（未年化）、高于行业内中位数（1.36\%）；资产负债率 23.8\%、低于行业内中位数（39.1\%）。 与同类公司相比，公司的差异化在于：蛋白含量达 70\% 以上。 优势在当期体现为高于行业中位数的资本回报，但优势能否持续仍取决于周期与自身产能投放节奏。

\pfl{历史股价} 自 2021 年 1 月至 2026 年 10 月，股价由 4.82 元变为 9.07 元，累计 +88.2\%，区间最高 13.86 元（2026 年 1 月）、最低 2.72 元（2021 年 5 月）；当前价相当于区间最高价的 65.4\%，为区间最低价的 3.33 倍。 月度收益的年化波动率 44.2\%，高于全样本中位数 37.3\%，在三级行业“其他农产品加工”的 11 家公司中波动率由高到低排第\zhnumber{2}位。 把股价阶段与公司事件对齐看，2024 年 8 月 至 2026 年 1 月股价上涨+188.1\%，同期（2025 年 4 月）控股股东获得 13.00 亿元现金增资用于取得控制权；2021 年 5 月 至 2022 年 8 月股价上涨+182.4\%，同期（2022 年 8 月）公司撤销退市风险警示恢复“华资实业”，公司走到退市风险边缘。 把这些阶段与公司事件对齐后可以看到，几轮大涨大跌多由公司自身的资本运作与风险事件触发，行业景气只解释了其中一部分。

\begin{center}
\begin{minipage}{\textwidth}\centering
\begin{tikzpicture}
\begin{axis}[
  width=12.6cm, height=3.9cm, scale only axis,
  xmin=2020.922, xmax=2026.972, ymin=2.53, ymax=14.83,
  xtick={2021,2022,2023,2024,2025,2026},
  yearticks,
  yticklabel style={font=\scriptsize},
  ylabel={元/股}, ylabel style={font=\scriptsize},
  grid=major, grid style={LightGray, line width=0.3pt},
  axis line style={InkGray, line width=0.5pt},
  every axis plot/.append style={line width=0.9pt},
]
\addplot[AgriGreen] table[x=x,y=y]{data/clean/profiles/px_600191.dat};
\addplot[SoftRed, dashed, line width=0.7pt] table[x=x,y=y]{data/clean/profiles/pxzz_600191.dat};
\addplot[only marks, mark=*, mark size=1.1pt, SoftRed] table[x=x,y=y]{data/clean/profiles/pxzz_600191.dat};
\end{axis}
\end{tikzpicture}

\captionof{figure}[华资实业（600191）股价与周期骨架]{华资实业（600191）月度收盘价（元/股）与 ZigZag 周期骨架（阈值 25\%，圆点为周期转折点）}
\end{minipage}
\end{center}

\pfl{过去周期分析} 以 25\% 的 ZigZag 阈值划分，样本区间内股价形成 3 个主升段与 4 个主跌段。 最大主升段出现在 2024 年 8 月 至 2026 年 1 月，17 个月+188.1\%；最大主跌段为 2021 年 1 月 至 2021 年 5 月，4 个月-43.6\%。 最近一次转折出现在 2026 年 10 月（下跌段自 2026 年 1 月起，9 个月-34.6\%），当前处于该段的延续之中。 公司股价较申万农林牧渔指数提前约 60 个月见底，是板块中较早定价周期反转的公司之一。 该子行业横跨糖、番茄制品、添加剂等多个品种，各自的供给周期与政策（进口配额、关税）是主要变量，公司 2026 年上半年实现盈利、半年 ROE 2.1\%（未年化），好于全样本中位数（0.75\%），下行周期对其报表的冲击小于同业。 综合区间形态看，该股单段涨跌幅度偏大、以趋势段为主（最大主升 +188.1\%、最大主跌 -43.6\%），年化波动率 44.2\% 与全样本中位数 37.3\% 相比波动明显大于样本中位数。

\pfl{盈利情况} 2026 年上半年营业收入 3.49 亿元（同比 +58.0\%），归母净利润 0.36 亿元，同比 +722.6\%。 以年报为趋势对照，2023---2025 年营业收入由 6.37 亿元变为 6.29 亿元（两年年均复合增速 -0.6\%），归母净利润由 0.19 亿元变为 0.59 亿元。 按半年报口径，2026 年上半年的 ROE 为 2.1\%（未年化）、销售净利率 10.4\%、毛利率 4.6\%，ROE 在“其他农产品加工”11 家公司中按数值由高到低排第\zhnumber{4}位。 同期全样本半年 ROE 中位数 0.75\%，公司与之相差 1.37 个百分点（略好于中位数公司）。 综合看，公司在报告期维持盈利（高于全样本半年 ROE 中位数 0.75\%），利润的可持续性取决于其他农产品加工景气的延续与成本管控的执行。

\begin{center}\small\setlength{\tabcolsep}{3.4pt}
\captionof{table}[华资实业（600191）关键财务指标]{华资实业（600191）关键财务指标（合并报表口径；两个半年报期均未年化；现金短债比 $=$ 货币资金 $\div$ 短期债务）}
\begin{adjustbox}{max width=\textwidth}
\begin{tabular}{@{}lrrrrr@{}}
\toprule
指标 & 2023 年 & 2024 年 & 2025 年 & 2025 年半年报 & 2026 年半年报 \\
\midrule
营业收入（亿元） & 6.37 & 5.16 & 6.29 & 2.21 & 3.49 \\
归母净利润（亿元） & 0.19 & 0.25 & 0.59 & 0.04 & 0.36 \\
毛利率（\%） & 8.6 & 8.3 & 7.5 & 6.0 & 4.6 \\
ROE（\%） & 1.3 & 1.6 & 3.5 & 0.3 & 2.1 \\
资产负债率（\%） & 20.6 & 25.5 & 27.7 & 27.1 & 23.8 \\
经营活动现金流净额（亿元） & 1.10 & 1.39 & -0.06 & 0.19 & 0.31 \\
货币资金（亿元） & 0.47 & 0.42 & 0.90 & 0.46 & 0.46 \\
有息负债合计（亿元） & 2.70 & 1.68 & 2.68 & 2.93 & 2.53 \\
现金短债比（倍） & 0.27 & 0.50 & 0.59 & 0.28 & 0.33 \\
\bottomrule
\end{tabular}
\end{adjustbox}
\end{center}

\begin{center}
\begin{minipage}{\textwidth}\centering
\begin{tikzpicture}
\begin{groupplot}[
  group style={group size=2 by 1, horizontal sep=0.60cm},
  width=6.85cm, height=4.60cm,
  tick label style={font=\tiny},
  title style={font=\scriptsize\heiti},
  yticklabel style={font=\tiny},
  ylabel style={font=\tiny},
  grid=major, grid style={LightGray, line width=0.3pt},
  axis line style={InkGray, line width=0.5pt},
  enlarge x limits={abs=0.8},
  legend style={font=\scriptsize, at={(0.5,-0.20)}, anchor=north,
                legend columns=2, draw=none, fill=none, column sep=6pt},
]
\nextgroupplot[
  ybar, bar width=4pt,
  ymin=-2.45, ymax=7.34,
  xtick={0,1,2,3,4,5.7},
  xticklabels={2021,2022,2023,2024,2025,2026H1},
  ylabel={亿元},
]
\addplot[fill=AgriGreen!75, draw=AgriGreen, bar shift=-2.6pt] table[x=i, y=rev]{data/clean/profiles/fin_600191.dat};
\addplot[fill=SoftRed!60, draw=SoftRed, bar shift=2.6pt] table[x=i, y=ni]{data/clean/profiles/fin_600191.dat};
\addplot[fill=AgriGreen!30, draw=AgriGreen, pattern=north east lines, bar shift=-2.6pt] table[x=x, y=rev]{data/clean/profiles/finh_600191.dat};
\addplot[fill=SoftRed!20, draw=SoftRed, pattern=north east lines, bar shift=2.6pt] table[x=x, y=ni]{data/clean/profiles/finh_600191.dat};
\legend{营业收入（年/半年）, 归母净利润（年/半年）}
\nextgroupplot[
  ymin=-13.5, ymax=31.5,
  xtick={0,1,2,3,4,5.7},
  xticklabels={2021,2022,2023,2024,2025,2026H1},
  ylabel={\%},
]
\addplot[AgriGold, mark=*, mark size=1.3pt] table[x=i, y=roe]{data/clean/profiles/fin_600191.dat};
\addplot[OilBlue, mark=square*, mark size=1.3pt] table[x=i, y=debt]{data/clean/profiles/fin_600191.dat};
\addplot[only marks, AgriGold, mark=*, mark size=2.0pt] table[x=x, y=roe]{data/clean/profiles/finh_600191.dat};
\addplot[only marks, OilBlue, mark=square*, mark size=2.0pt] table[x=x, y=debt]{data/clean/profiles/finh_600191.dat};
\legend{ROE, 资产负债率}
\end{groupplot}
\end{tikzpicture}
\begin{tikzpicture}
\begin{axis}[
  name=finbot,
  width=13.75cm, height=3.10cm, scale only axis,
  ymin=-0.29, ymax=3.24,
  xtick={0,1,2,3,4,5.7},
  xticklabels={2021,2022,2023,2024,2025,2026H1},
  tick label style={font=\tiny},
  yticklabel style={font=\tiny},
  ylabel={亿元}, ylabel style={font=\tiny},
  ybar, bar width=4pt,
  grid=major, grid style={LightGray, line width=0.3pt},
  axis line style={InkGray, line width=0.5pt},
  enlarge x limits={abs=0.8},
  legend style={font=\scriptsize, at={(0.5,-0.26)}, anchor=north,
                legend columns=2, draw=none, fill=none, column sep=10pt},
]
\addplot[fill=AgriGreen!55, draw=AgriGreen, bar shift=-3.0pt] table[x=i, y=cash]{data/clean/profiles/fin_600191.dat};
\addplot[fill=SoftRed!45, draw=SoftRed, bar shift=3.0pt] table[x=i, y=ibd]{data/clean/profiles/fin_600191.dat};
\addplot[fill=AgriGreen!25, draw=AgriGreen, pattern=north east lines, bar shift=-3.0pt] table[x=x, y=cash]{data/clean/profiles/finh_600191.dat};
\addplot[fill=SoftRed!20, draw=SoftRed, pattern=north east lines, bar shift=3.0pt] table[x=x, y=ibd]{data/clean/profiles/finh_600191.dat};
\legend{货币资金（左轴）, 有息负债（左轴）}
\end{axis}
\begin{axis}[
  at={(finbot.south east)}, anchor=south east,
  width=13.75cm, height=3.10cm, scale only axis,
  axis y line*=right, axis x line*=none, xtick=\empty,
  ymin=-3.19, ymax=2.49,
  tick label style={font=\tiny},
  yticklabel style={font=\tiny}, scaled y ticks=false,
  legend style={font=\scriptsize, at={(0.5,-0.44)}, anchor=north,
                legend columns=1, draw=none, fill=none},
]
\addplot[OilBlue, mark=*, mark size=2.2pt, line width=1.3pt] table[x=i, y=ocf]{data/clean/profiles/fin_600191.dat};
\addplot[only marks, OilBlue, mark=*, mark size=3.2pt] table[x=x, y=ocf]{data/clean/profiles/finh_600191.dat};
\legend{经营活动现金流净额（右轴，亿元，蓝点）}
\end{axis}
\end{tikzpicture}

\captionof{figure}[华资实业（600191）近年主要财务指标]{华资实业（600191）近年主要财务指标（左上：营业收入与归母净利润；右上：ROE 与资产负债率；下：货币资金、有息负债为左轴柱状，经营活动现金流净额为其右轴的蓝点折线。
2021---2025 年为年报口径，2026H1 为 2026 年半年报/半年末，与其前一年报之间留出间隔、以斜纹或空心点区分；半年报数据未年化）}
\end{minipage}
\end{center}

\pfl{财务分析} 2026 年 6 月末资产负债率 23.8\%，较 2025 年末下降 3.9 个百分点（样本中位数 52.2\%，所处三级行业内按由低到高排第\zhnumber{4}位）。 2026 年 6 月末货币资金 0.46 亿元、短期债务（短期借款与一年内到期非流动负债）1.38 亿元、长期债务 1.14 亿元，有息负债合计 2.53 亿元，现金短债比 0.33 倍——覆盖偏紧。 净有息负债 2.07 亿元，相当于其 2026 年上半年营业收入的 59\%。 流动比率 0.61、速动比率 0.36，短期偿债指标偏紧。 2026 年上半年经营活动现金流净额 0.31 亿元，经营现金流与利润同向为正，内生资金可以支撑运营。 综合看，账面货币资金对有息负债与短期债务的覆盖不足，滚动偿付对新增融资的依赖度较高；资产负债率低于样本中位数 28.4 个百分点；有息负债中短期债务占比更高，期限结构仍以滚动续作为主。

\pfl{重大战略转型} 近三年公司的战略动作集中在：其他动作（4 项）：2026 年上半年，归属于上市公司股东的净利润 3,616.08 万元、同比增长 722.60\%；2022 年，公司设立全资子公司山东裕维生物科技有限公司；产能与项目（2 项）：2025 年 4 月，年产 2 万吨黄原胶项目正式投产；2023 年，公司启动大规模产能升级计划。 公告口径上，近三年（2023 年 9 月---2026 年 9 月）公司披露重大重组与并购类公告 2 条、对外投资与转型类公告 2 条、回购与股权激励类公告 0 条，合计公告 316 条。 主营结构的口径变化已经落到报表上：2021 年年报的第一大行业为“制糖业”（84.2\%），2025 年年报为“食品制造业”（91.3\%）；公司在这两期的分部划分方式有过调整。 综合判断，报告期内公司有 4 条与战略相关的公告落地（重大重组与并购 2 条、对外投资与转型 2 条），资本运作与主业调整之间存在直接关联。

\begin{center}\small\setlength{\tabcolsep}{4pt}
\captionof{table}[华资实业（600191）历史融资与资本运作]{华资实业（600191）历史融资与资本运作（据公司公告与公开报道整理；金额为披露值，含首次公开发行、再融资、债券、银行借款与重整投资等）}
\begin{adjustbox}{max width=\textwidth}
\begin{tabular}{@{}p{1.45cm}p{2.55cm}p{4.05cm}p{4.95cm}@{}}
\toprule
时间 & 方式 & 规模 & 用途与说明 \\
\midrule
1998-11-02 & 首发A股（上交所） & 发行 7,\allowbreak{}000 万股，\allowbreak{}发行价 4.30 元/\allowbreak{}股 & 新浪财经报道口径为上市募集资金 2.89 亿元；同花顺 F10 登记的“首发 A 股”实际募集资金净额为 2.86 亿元 \\
2015-2016 & 非公开发行股份募集资金（拟用于增资华夏人寿） & 拟募集不超过 316.80 亿元，\allowbreak{}增资完成后持有华夏人寿股权比例不超过 51\% & 对华夏人寿增资；2016 年 7 月证监会发审委审核通过；2017 年 11 月公司公告称鉴于证券市场发生较大变化并综合考虑融资环境和监管政策变化等因素 \\
2025-04-15 & 控股股东层面增资（间接控股，非上市公司融资） & 中裕科技投资以现金 13.00 亿元对盛泰创发增资，\allowbreak{}取得增资后 55\% 股权 & 取得上市公司控制权；控股股东仍为盛泰创发 \\
\bottomrule
\end{tabular}
\end{adjustbox}
\end{center}

\pfl{历史融投资情况} 从现金流量表看，2025 年公司取得借款收到的现金 1.60 亿元、偿还债务支付的现金 0.44 亿元，筹资活动现金流净额 2.36 亿元（2023 年为取得借款 0.70 亿元、偿还债务 0.79 亿元）。 2026 年上半年筹资活动现金流净额为 -0.58 亿元（取得借款 0.88 亿元、偿还债务 0.98 亿元，未年化），已转为净流出，处于净还债状态。 2025 年分配股利、利润或偿付利息支付的现金合计 0.07 亿元。 公开披露的融资记录共 3 笔，工具类型以定向增发为主；金额与用途见下表。 其中规模最大的两笔为：1998 年 11 月，首发 A 股，发行 7,000 万股，发行价 4.30 元/股；2025 年 4 月，控股股东层面增资，中裕科技投资以现金 13.00 亿元对盛泰创发增资，取得增资后 55\% 股权。 股本方面，近三年披露股本变动 3 次；总股本自 2021 年末以来未发生变化（4.85 亿股）。 分红方面，近三年现金分红 1 次、合计每股派现 0.02 元，最近一次实施在 2025 年 12 月。 近三年“再融资”类公告 7 条，涉及定向增发、募集资金使用。 综合看，公司融资结构上股权与债务融资并举；2025 年筹资活动现金流净额 2.36 亿元，年度内仍为净融入，与前述经营与投资现金流的方向基本一致。

\pfl{未来潜在融资需求分析} 2026 年 6 月末货币资金 0.46 亿元，而短期债务 1.38 亿元（现金短债比 0.33 倍），2026 年上半年经营现金流净额 0.31 亿元。按“短期债务减去账面货币资金”的滚动偿付口径，静态缺口约 0.92 亿元；上半年盈利或接近盈亏平衡，该缺口不随经营亏损进一步扩大。 公司 2026 年 6 月末资产负债率 23.8\%、2026 年上半年 ROE 2.1\%（未年化），现金短债比 0.33 倍、短债覆盖不足，融资需求首先用于置换与滚动存量债务、补充营运资金，属于“补血型”而非扩张型。 综合看，公司的外部资金需求以补充营运资金为主、项目投入为辅，滚动偿付口径下的静态缺口约 0.92 亿元，融资的时点与规模更多取决于信用条件与资本市场窗口。




\subsubsection{东方海洋（002086）}
\label{co:002086}

\pfl{公司基本情况} 2001 年设立至今，东方海洋是山东烟台的海水苗种繁育与养殖、水产品加工企业，2006 年深交所上市；2022—2023 年因债务危机、资金占用与违规担保进入破产重整，2023 年底重整计划执行完毕、2024 年 6 月摘帽，现推进“海洋产业＋大健康（体外诊断）”双主业，但仍持续亏损。 公司于 2006-11-28 在深交所主板首发上市，注册资本 19.59 亿元，总股本 19.59 亿股。 截至 2026 年 10 月 12 日收盘价 2.03 元，对应总市值 39.8 亿元，近一年-23.1\%，总市值在三级行业“其他农产品加工”的 11 家公司中按由大到小排第\zhnumber{7}位，近一年涨跌幅低于全样本中位数（-10.4\%）12.8 个百分点。 按 2026 年半年报披露的行业构成，收入占比前三位为水产品加工 68.4\%；体外诊断 12.9\%；海水养殖 12.7\%。

\pfl{股权结构与控制权} 按 2026 年 6 月 30 日的股东名册，前三大股东为五矿金通股权投资基金管理有限公司－金通东洋私募股权投 持股 18.26\%；烟台市正大城市建设发展有限公司 持股 3.43\%；山高(烟台)辰星投资中心(有限合伙) 持股 2.70\%，股东户数 69,288 户。 从控制权角度看，五矿金通－金通东洋私募股权投资基金出具保持上市公司独立性的承诺函与避免同业竞争的承诺函。

\pfl{经营要点} 从经营要点看，具备年 5.8 万吨水产品加工、8 万吨冷藏仓储产能。

\pfl{所处行业的基本面} 该子行业横跨制糖、番茄制品、糖醇与添加剂，供给受种植面积与进口配额约束，需求端分食用与工业两类（第 \ref{sec:framework} 章的供需框架）。 与该公司直接相关的行业数据包括：2024 年，全国海参养殖总产量达 32.6 万吨、全产业链产值突破 1,300.00 亿元。

\pfl{竞争优势} 把该公司放回同类公司比较，2026 年上半年营业收入 1.6 亿元、在三级行业“其他农产品加工”的 11 家公司中排第\zhnumber{11}位；毛利率 0.2\%、低于行业内中位数 8.8 个百分点；半年 ROE -3.8\%（未年化）、低于行业内中位数（1.36\%）；资产负债率 19.6\%、低于行业内中位数（39.1\%）。 公司在公开资料中自述与同业比较的相对优势在于：公司水产品加工拥有从初级加工到精深加工的完整产业链、工厂产能居行业前列。 就当期报表看，公司尚未把经营层面的特征转化为高于行业的资本回报，周期底部的相对优势更多体现在现金流与资产负债表上。

\pfl{历史股价} 以 2021 年 1 月为起点、2026 年 10 月为终点，股价由 1.66 元变为 2.03 元，累计 +22.3\%，区间最高 3.51 元（2024 年 11 月）、最低 1.58 元（2021 年 4 月）；当前价相当于区间最高价的 57.8\%，为区间最低价的 1.28 倍。 月度收益的年化波动率 43.5\%，高于全样本中位数 37.3\%，在三级行业“其他农产品加工”的 11 家公司中波动率由高到低排第\zhnumber{3}位。

\begin{center}
\begin{minipage}{\textwidth}\centering
\begin{tikzpicture}
\begin{axis}[
  width=12.6cm, height=3.9cm, scale only axis,
  xmin=2020.922, xmax=2026.972, ymin=1.47, ymax=3.76,
  xtick={2021,2022,2023,2024,2025,2026},
  yearticks,
  yticklabel style={font=\scriptsize},
  ylabel={元/股}, ylabel style={font=\scriptsize},
  grid=major, grid style={LightGray, line width=0.3pt},
  axis line style={InkGray, line width=0.5pt},
  every axis plot/.append style={line width=0.9pt},
]
\addplot[AgriGreen] table[x=x,y=y]{data/clean/profiles/px_002086.dat};
\addplot[SoftRed, dashed, line width=0.7pt] table[x=x,y=y]{data/clean/profiles/pxzz_002086.dat};
\addplot[only marks, mark=*, mark size=1.1pt, SoftRed] table[x=x,y=y]{data/clean/profiles/pxzz_002086.dat};
\end{axis}
\end{tikzpicture}

\captionof{figure}[东方海洋（002086）股价与周期骨架]{东方海洋（002086）月度收盘价（元/股）与 ZigZag 周期骨架（阈值 25\%，圆点为周期转折点）}
\end{minipage}
\end{center}

\pfl{过去周期分析} 以 25\% 的 ZigZag 阈值划分，样本区间内股价形成 4 个主升段与 3 个主跌段。 最大主升段出现在 2021 年 1 月 至 2021 年 12 月，11 个月+85.5\%；最大主跌段为 2024 年 11 月 至 2026 年 6 月，19 个月-53.3\%。 最近一次转折出现在 2026 年 8 月（上涨段自 2026 年 6 月起，2 个月+29.9\%），当前处于该段之后的震荡之中。 公司股价较申万农林牧渔指数提前约 61 个月见底，是板块中较早定价周期反转的公司之一。 该子行业横跨糖、番茄制品、添加剂等多个品种，各自的供给周期与政策（进口配额、关税）是主要变量，公司 2026 年 6 月末资产负债率 19.6\%、半年 ROE -3.8\%（未年化），在本轮下行中属于随行业同步调整的一类。 综合区间形态看，该股单段涨跌幅度偏大、以趋势段为主（最大主升 +85.5\%、最大主跌 -53.3\%），年化波动率 43.5\% 与全样本中位数 37.3\% 相比波动明显大于样本中位数。

\pfl{盈利情况} 2026 年上半年营业收入 1.57 亿元（同比 +0.9\%），归母净利润 -0.48 亿元，亏损同比收窄 0.23 亿元。 以年报为趋势对照，2023---2025 年营业收入由 4.37 亿元变为 3.56 亿元（两年年均复合增速 -9.8\%），归母净利润由 17.50 亿元变为 -1.53 亿元。 按半年报口径，2026 年上半年的 ROE 为 -3.8\%（未年化）、销售净利率 -30.7\%、毛利率 0.2\%，ROE 在“其他农产品加工”11 家公司中按数值由高到低排第\zhnumber{11}位。 按同一口径，三级行业“其他农产品加工”11 家公司的半年 ROE 中位数为 1.36\%，该公司低于其中位数 5.20 个百分点。 综合看，公司仍处亏损区间、亏损同比收窄，半年 ROE 在三级行业内按由高到低排第\zhnumber{11}位，单位盈利能力的修复依赖行业景气与自身成本端的改善，报表弹性主要体现为价格与销量的方向性变化。

\begin{center}\small\setlength{\tabcolsep}{3.4pt}
\captionof{table}[东方海洋（002086）关键财务指标]{东方海洋（002086）关键财务指标（合并报表口径；两个半年报期均未年化；现金短债比 $=$ 货币资金 $\div$ 短期债务）}
\begin{adjustbox}{max width=\textwidth}
\begin{tabular}{@{}lrrrrr@{}}
\toprule
指标 & 2023 年 & 2024 年 & 2025 年 & 2025 年半年报 & 2026 年半年报 \\
\midrule
营业收入（亿元） & 4.37 & 3.40 & 3.56 & 1.56 & 1.57 \\
归母净利润（亿元） & 17.50 & -1.90 & -1.53 & -0.71 & -0.48 \\
毛利率（\%） & 13.7 & -3.4 & 0.2 & -0.9 & 0.2 \\
ROE（\%） & 0.0 & -12.5 & -11.3 & -5.1 & -3.8 \\
资产负债率（\%） & 39.9 & 20.9 & 20.4 & 18.5 & 19.6 \\
经营活动现金流净额（亿元） & -0.46 & -1.93 & -1.25 & -0.62 & 0.08 \\
货币资金（亿元） & 12.93 & 4.05 & 2.35 & 3.03 & 2.43 \\
有息负债合计（亿元） & 5.26 & 0.63 & 0.53 & 0.63 & 0.49 \\
现金短债比（倍） & 3.17 & 27.25 & 21.44 & 21.92 & 37.76 \\
\bottomrule
\end{tabular}
\end{adjustbox}
\end{center}

\begin{center}
\begin{minipage}{\textwidth}\centering
\begin{tikzpicture}
\begin{groupplot}[
  group style={group size=2 by 1, horizontal sep=0.60cm},
  width=6.85cm, height=4.60cm,
  tick label style={font=\tiny},
  title style={font=\scriptsize\heiti},
  yticklabel style={font=\tiny},
  ylabel style={font=\tiny},
  grid=major, grid style={LightGray, line width=0.3pt},
  axis line style={InkGray, line width=0.5pt},
  enlarge x limits={abs=0.8},
  legend style={font=\scriptsize, at={(0.5,-0.20)}, anchor=north,
                legend columns=2, draw=none, fill=none, column sep=6pt},
]
\nextgroupplot[
  ybar, bar width=4pt,
  ymin=-19.18, ymax=21.50,
  xtick={0,1,2,3,4,5.7},
  xticklabels={2021,2022,2023,2024,2025,2026H1},
  ylabel={亿元},
]
\addplot[fill=AgriGreen!75, draw=AgriGreen, bar shift=-2.6pt] table[x=i, y=rev]{data/clean/profiles/fin_002086.dat};
\addplot[fill=SoftRed!60, draw=SoftRed, bar shift=2.6pt] table[x=i, y=ni]{data/clean/profiles/fin_002086.dat};
\addplot[fill=AgriGreen!30, draw=AgriGreen, pattern=north east lines, bar shift=-2.6pt] table[x=x, y=rev]{data/clean/profiles/finh_002086.dat};
\addplot[fill=SoftRed!20, draw=SoftRed, pattern=north east lines, bar shift=2.6pt] table[x=x, y=ni]{data/clean/profiles/finh_002086.dat};
\legend{营业收入（年/半年）, 归母净利润（年/半年）}
\nextgroupplot[
  ymin=-132.1, ymax=184.5,
  xtick={0,1,2,3,4,5.7},
  xticklabels={2021,2022,2023,2024,2025,2026H1},
  ylabel={\%},
]
\addplot[AgriGold, mark=*, mark size=1.3pt] table[x=i, y=roe]{data/clean/profiles/fin_002086.dat};
\addplot[OilBlue, mark=square*, mark size=1.3pt] table[x=i, y=debt]{data/clean/profiles/fin_002086.dat};
\addplot[only marks, AgriGold, mark=*, mark size=2.0pt] table[x=x, y=roe]{data/clean/profiles/finh_002086.dat};
\addplot[only marks, OilBlue, mark=square*, mark size=2.0pt] table[x=x, y=debt]{data/clean/profiles/finh_002086.dat};
\legend{ROE, 资产负债率}
\end{groupplot}
\end{tikzpicture}
\begin{tikzpicture}
\begin{axis}[
  name=finbot,
  width=13.75cm, height=3.10cm, scale only axis,
  ymin=-1.29, ymax=14.48,
  xtick={0,1,2,3,4,5.7},
  xticklabels={2021,2022,2023,2024,2025,2026H1},
  tick label style={font=\tiny},
  yticklabel style={font=\tiny},
  ylabel={亿元}, ylabel style={font=\tiny},
  ybar, bar width=4pt,
  grid=major, grid style={LightGray, line width=0.3pt},
  axis line style={InkGray, line width=0.5pt},
  enlarge x limits={abs=0.8},
  legend style={font=\scriptsize, at={(0.5,-0.26)}, anchor=north,
                legend columns=2, draw=none, fill=none, column sep=10pt},
]
\addplot[fill=AgriGreen!55, draw=AgriGreen, bar shift=-3.0pt] table[x=i, y=cash]{data/clean/profiles/fin_002086.dat};
\addplot[fill=SoftRed!45, draw=SoftRed, bar shift=3.0pt] table[x=i, y=ibd]{data/clean/profiles/fin_002086.dat};
\addplot[fill=AgriGreen!25, draw=AgriGreen, pattern=north east lines, bar shift=-3.0pt] table[x=x, y=cash]{data/clean/profiles/finh_002086.dat};
\addplot[fill=SoftRed!20, draw=SoftRed, pattern=north east lines, bar shift=3.0pt] table[x=x, y=ibd]{data/clean/profiles/finh_002086.dat};
\legend{货币资金（左轴）, 有息负债（左轴）}
\end{axis}
\begin{axis}[
  at={(finbot.south east)}, anchor=south east,
  width=13.75cm, height=3.10cm, scale only axis,
  axis y line*=right, axis x line*=none, xtick=\empty,
  ymin=-2.51, ymax=0.96,
  tick label style={font=\tiny},
  yticklabel style={font=\tiny}, scaled y ticks=false,
  legend style={font=\scriptsize, at={(0.5,-0.44)}, anchor=north,
                legend columns=1, draw=none, fill=none},
]
\addplot[OilBlue, mark=*, mark size=2.2pt, line width=1.3pt] table[x=i, y=ocf]{data/clean/profiles/fin_002086.dat};
\addplot[only marks, OilBlue, mark=*, mark size=3.2pt] table[x=x, y=ocf]{data/clean/profiles/finh_002086.dat};
\legend{经营活动现金流净额（右轴，亿元，蓝点）}
\end{axis}
\end{tikzpicture}

\captionof{figure}[东方海洋（002086）近年主要财务指标]{东方海洋（002086）近年主要财务指标（左上：营业收入与归母净利润；右上：ROE 与资产负债率；下：货币资金、有息负债为左轴柱状，经营活动现金流净额为其右轴的蓝点折线。
2021---2025 年为年报口径，2026H1 为 2026 年半年报/半年末，与其前一年报之间留出间隔、以斜纹或空心点区分；半年报数据未年化）}
\end{minipage}
\end{center}

\pfl{财务分析} 2026 年 6 月末资产负债率 19.6\%，较 2025 年末下降 0.8 个百分点（样本中位数 52.2\%，所处三级行业内按由低到高排第\zhnumber{1}位）。 2026 年 6 月末货币资金 2.43 亿元、短期债务（短期借款与一年内到期非流动负债）0.06 亿元、长期债务 0.42 亿元，有息负债合计 0.49 亿元，现金短债比 37.76 倍——覆盖充分。 净有息负债 -1.94 亿元，相当于其 2026 年上半年营业收入的 -123\%。 流动比率 3.71、速动比率 2.05，短期指标尚可。 2026 年上半年经营活动现金流净额 0.08 亿元，净利润为负而经营现金流为正，差额主要来自折旧摊销与营运资本变动，说明亏损尚未完全转化为现金流出。 综合看，现金短债比 37.76 倍，短期偿付压力可控；资产负债率低于样本中位数 32.6 个百分点；有息负债中长期债务占比更高，期限结构对短债滚动的压力较小。 公司披露的重整与债务违约风险为：2024 年 4 月，公司股票可能被终止上市的风险提示公告。

\pfl{重大战略转型} 公司的战略动作可归为以下几类：并购与重组（4 项）：2025 年度，公司表示将通过外延并购及内生增长方式发展壮大医疗大健康产业；2024---2025 年，依托重整投资人提供的流动资金支持和业务资源支持。 公告口径上，近三年（2023 年 9 月---2026 年 9 月）公司披露重大重组与并购类公告 0 条、对外投资与转型类公告 4 条、回购与股权激励类公告 0 条，合计公告 464 条。 第一大行业“水产品加工”的收入占比由 2021 年年报的 35.5\% 变为 68.4\%（32.9 个百分点），收入结构出现再平衡。 综合判断，公司当前的战略重心不是扩张而是化解债务与完成重整，业务层面的调整让位于司法重整程序。

\begin{center}\small\setlength{\tabcolsep}{4pt}
\captionof{table}[东方海洋（002086）历史融资与资本运作]{东方海洋（002086）历史融资与资本运作（据公司公告与公开报道整理；金额为披露值，含首次公开发行、再融资、债券、银行借款与重整投资等）}
\begin{adjustbox}{max width=\textwidth}
\begin{tabular}{@{}p{1.45cm}p{2.55cm}p{4.05cm}p{4.95cm}@{}}
\toprule
时间 & 方式 & 规模 & 用途与说明 \\
\midrule
2006-11-28 & IPO（深交所中小板） & 发行 3,\allowbreak{}450 万股，\allowbreak{}发行价 7.49 元/\allowbreak{}股，\allowbreak{}实际募集资金合计 2.58 亿元 & 全部用于水产品加工贸易基地建设项目；发行市盈率（按发行后总股本）23.77 倍；首发前总股本 5,\allowbreak{}180 万股 \\
2023-12 & 破产重整投资款（资本公积金转增股票引入重整投资人） & 截至 2023 年 12 月 29 日重整投资人向管理人指定银行账户共计支付 13.54 亿元投资款 & 偿付债务、\allowbreak{}支付重整费用、\allowbreak{}补充公司流动资金、\allowbreak{}部分解决非经营性资金占用及违规担保问题；以现有总股本 7.56 亿股为基数 \\
2023-12 & 债务豁免与债务重组 & 债权人直接豁免东方海洋清偿义务 2.56 亿元；豁免金额相应抵偿原控股股东对上市公司同等金额的非经营性资金占用 7.77 亿元 & 解决原控股股东非经营性资金占用及违规担保；部分债权人与原控股股东等签订债务重组与化解协议；原控股股东非经营性占用资金已全部清偿完毕（最终以审计机构专项审核报告为准） \\
2023-12-17 & 重整投资协议 & 五矿金通股权投资基金管理有限公司牵头设立专项私募投资基金 & 认购公司资本公积转增股本，\allowbreak{}帮助公司化解债务风险；招商平安有条件认购 4,\allowbreak{}600 万股公司资本公积金转增股票等安排 \\
\bottomrule
\end{tabular}
\end{adjustbox}
\end{center}

\pfl{历史融投资情况} 从现金流量表看，2025 年公司吸收权益性投资 0.71 亿元，筹资活动现金流净额 0.67 亿元（2023 年为取得借款 0.45 亿元、偿还债务 0.66 亿元）。 2026 年上半年筹资活动现金流净额为 -0.00 亿元（取得借款 — 亿元、偿还债务 — 亿元，未年化），已转为净流出，处于净还债状态。 公开披露的融资记录共 4 笔，工具类型以 IPO、重整投资为主；金额与用途见下表。 其中规模最大的两笔为：2023 年 12 月，重整投资，截至 2023 年 12 月 29 日重整投资人向管理人指定银行账户共计支付 13.54 亿元投资款；2006 年 11 月，IPO，发行 3,450 万股，发行价 7.49 元/股，实际募集资金合计 2.58 亿元。 股本方面，近三年披露股本变动 5 次（转增 1 次、限售股份上市 1 次）；总股本由 2021 年末的 7.56 亿股变为 19.59 亿股（+159.00\%）。 分红方面，近三年未实施现金分红，与重整期间的资金约束一致。 综合看，公司融资结构上股权与债务融资并举；2025 年筹资活动现金流净额 0.67 亿元，年度内仍为净融入，与前述经营与投资现金流的方向基本一致。

\pfl{未来潜在融资需求分析} 2026 年 6 月末货币资金 2.43 亿元，而短期债务 0.06 亿元（现金短债比 37.76 倍），2026 年上半年经营现金流净额 0.08 亿元。账面货币资金可以覆盖短期债务，缺口不体现在静态偿付层面。 公司 2026 年上半年归母净利润为负（-0.48 亿元），但 6 月末资产负债率 19.6\% 尚未进入第一档（亏损且负债率 $>$65\%）的区间，当前融资需求以补充流动资金、支撑低谷期经营性支出为主。 公开披露的后续资金安排线索包括：债务与授信安排：2023 年 12 月，重整投资人合计投入 13.54 亿元资金。 静态偿付不存在缺口，后续融资更可能服务于产能或并购安排，而不是填补流动性窟窿。




\subsubsection{祖名股份（003030）}
\label{co:003030}

\pfl{公司基本情况} 祖名股份成立于 2020 年，国内首家登陆资本市场的豆制品企业（豆腐第一股），2021 年 1 月深交所上市，主营生鲜豆制品、植物蛋白饮品、休闲豆制品等 400 余种产品，以浙江（杭州）区域市场及菜篮子保供为主阵地。2025 年营业收入首破 20 亿元、归母净利润 2,529.58 万元实现扭亏，但扣非仍亏损；2026 年上半年受大豆与能源价格上涨影响亏损扩大至 2,435.69 万元，同时以 10.12\% 持股跨界投资芯片设计公司重庆芯联芯，短期刚性债务压力上升。 公司于 2021-01-06 在深交所主板首发上市，注册资本 1.25 亿元，总股本 1.25 亿股。 截至 2026 年 10 月 12 日收盘价 26.03 元，对应总市值 32.5 亿元，近一年+40.8\%，总市值在三级行业“其他农产品加工”的 11 家公司中按由大到小排第\zhnumber{8}位，近一年涨跌幅高于全样本中位数（-10.4\%）51.1 个百分点。 按 2026 年半年报披露的产品构成，收入占比前三位为生鲜豆制品 67.3\%；植物蛋白饮品 16.3\%；其他 12.3\%。

\pfl{股权结构与控制权} 按 2026 年 6 月 30 日的股东名册，前三大股东为蔡祖明 持股 21.08\%；杭州纤品投资有限公司 持股 14.14\%；蔡水埼 持股 9.79\%，股东户数 8,388 户。 从控制权角度看，上市板块为深圳证券交易所；实际控制人王茶英、蔡水埼与蔡祖明均为公司前十大股东，其中蔡水埼直接持股 1,221.12 万股、占 9.79\%。

\pfl{经营要点} 经营层面，截至 2026 年 6 月末公司已拥有 40 项发明专利、49 项实用新型专利；产品保质期短且物流成本占比高。

\pfl{所处行业的基本面} 该子行业横跨制糖、番茄制品、糖醇与添加剂，供给受种植面积与进口配额约束，需求端分食用与工业两类（第 \ref{sec:framework} 章的供需框架）。

\pfl{竞争优势} 把该公司放回同类公司比较，2026 年上半年营业收入 10.1 亿元、在三级行业“其他农产品加工”的 11 家公司中排第\zhnumber{8}位；毛利率 21.6\%、高于行业内中位数 12.6 个百分点；半年 ROE -2.5\%（未年化）、低于行业内中位数（1.36\%）；资产负债率 57.9\%、高于行业内中位数（39.1\%）。 公司在公开资料中自述与同业比较的相对优势在于：2026 年上半年，毛利率 21.65\% 高于行业平均 11.39\% 的水平。 优势在报表上主要体现为高于行业平均的毛利率，这在周期底部是抵御亏损的主要缓冲。

\pfl{历史股价} 自 2021 年 1 月至 2026 年 10 月，股价由 42.60 元变为 26.03 元，累计 -38.9\%，区间最高 42.60 元（2021 年 1 月）、最低 12.99 元（2024 年 8 月）；当前价相当于区间最高价的 61.1\%，为区间最低价的 2.00 倍。 月度收益的年化波动率 33.8\%，低于全样本中位数 37.3\%，在三级行业“其他农产品加工”的 11 家公司中波动率由高到低排第\zhnumber{7}位。 把股价阶段与公司事件对齐看，2024 年 8 月 至 2026 年 10 月股价上涨+100.4\%，同期（2026 年 8 月）公司新增对子公司担保额度预计并提交股东会审议通过公司仍通过担保与银行授信为子公司经营提供资金支持。 把这些阶段与公司事件对齐后可以看到，几轮大涨大跌多由公司自身的资本运作与风险事件触发，行业景气只解释了其中一部分。

\begin{center}
\begin{minipage}{\textwidth}\centering
\begin{tikzpicture}
\begin{axis}[
  width=12.6cm, height=3.9cm, scale only axis,
  xmin=2020.922, xmax=2026.972, ymin=12.08, ymax=45.58,
  xtick={2021,2022,2023,2024,2025,2026},
  yearticks,
  yticklabel style={font=\scriptsize},
  ylabel={元/股}, ylabel style={font=\scriptsize},
  grid=major, grid style={LightGray, line width=0.3pt},
  axis line style={InkGray, line width=0.5pt},
  every axis plot/.append style={line width=0.9pt},
]
\addplot[AgriGreen] table[x=x,y=y]{data/clean/profiles/px_003030.dat};
\addplot[SoftRed, dashed, line width=0.7pt] table[x=x,y=y]{data/clean/profiles/pxzz_003030.dat};
\addplot[only marks, mark=*, mark size=1.1pt, SoftRed] table[x=x,y=y]{data/clean/profiles/pxzz_003030.dat};
\end{axis}
\end{tikzpicture}

\captionof{figure}[祖名股份（003030）股价与周期骨架]{祖名股份（003030）月度收盘价（元/股）与 ZigZag 周期骨架（阈值 25\%，圆点为周期转折点）}
\end{minipage}
\end{center}

\pfl{过去周期分析} 以 25\% 的 ZigZag 阈值划分，样本区间内股价形成 2 个主升段与 2 个主跌段。 最大主跌段为 2022 年 12 月 至 2024 年 8 月，20 个月-55.4\%；最大主升段出现在 2024 年 8 月 至 2026 年 10 月，26 个月+100.4\%。 最近一次转折出现在 2026 年 10 月（上涨段自 2024 年 8 月起，26 个月+100.4\%），当前处于该段之后的震荡之中。 公司股价较申万农林牧渔指数提前约 21 个月见底，是板块中较早定价周期反转的公司之一。 该子行业横跨糖、番茄制品、添加剂等多个品种，各自的供给周期与政策（进口配额、关税）是主要变量，公司 2026 年 6 月末资产负债率 57.9\%、半年 ROE -2.5\%（未年化），在本轮下行中属于随行业同步调整的一类。 综合区间形态看，该股单段涨跌幅度偏大、以趋势段为主（最大主升 +100.4\%、最大主跌 -55.4\%），年化波动率 33.8\% 与全样本中位数 37.3\% 相比波动与样本中位数接近。

\pfl{盈利情况} 2026 年上半年营业收入 10.13 亿元（同比 +10.0\%），归母净利润 -0.24 亿元，亏损同比扩大 0.15 亿元。 以年报为趋势对照，2023---2025 年营业收入由 14.78 亿元变为 20.25 亿元（两年年均复合增速 +17.0\%），归母净利润由 0.40 亿元变为 0.25 亿元。 按半年报口径，2026 年上半年的 ROE 为 -2.5\%（未年化）、销售净利率 -2.8\%、毛利率 21.6\%，ROE 在“其他农产品加工”11 家公司中按数值由高到低排第\zhnumber{10}位。 同期全样本半年 ROE 中位数 0.75\%，公司与之相差 3.20 个百分点（弱于中位数公司）。 面对这一局面，公司披露的举措包括：2024 年 6 月，实施 2024 年股票期权激励计划，向 41 名激励对象授予股票期权 72.00 万份、行权价格调整为 12.21 元/份。 综合看，公司仍处亏损区间、亏损同比扩大，半年 ROE 在三级行业内按由高到低排第\zhnumber{10}位，单位盈利能力的修复依赖行业景气与自身成本端的改善，报表弹性主要体现为价格与销量的方向性变化。

\begin{center}\small\setlength{\tabcolsep}{3.4pt}
\captionof{table}[祖名股份（003030）关键财务指标]{祖名股份（003030）关键财务指标（合并报表口径；两个半年报期均未年化；现金短债比 $=$ 货币资金 $\div$ 短期债务）}
\begin{adjustbox}{max width=\textwidth}
\begin{tabular}{@{}lrrrrr@{}}
\toprule
指标 & 2023 年 & 2024 年 & 2025 年 & 2025 年半年报 & 2026 年半年报 \\
\midrule
营业收入（亿元） & 14.78 & 16.64 & 20.25 & 9.21 & 10.13 \\
归母净利润（亿元） & 0.40 & -0.25 & 0.25 & -0.09 & -0.24 \\
毛利率（\%） & 27.3 & 22.7 & 23.8 & 24.1 & 21.6 \\
ROE（\%） & 3.9 & -2.5 & 2.5 & -0.9 & -2.5 \\
资产负债率（\%） & 47.3 & 55.6 & 54.6 & 58.2 & 57.9 \\
经营活动现金流净额（亿元） & 1.76 & 1.07 & 1.25 & 0.55 & 0.58 \\
货币资金（亿元） & 1.95 & 1.42 & 1.18 & 1.96 & 1.80 \\
有息负债合计（亿元） & 5.90 & 8.81 & 8.49 & 9.51 & 10.02 \\
现金短债比（倍） & 0.58 & 0.22 & 0.19 & 0.33 & 0.21 \\
\bottomrule
\end{tabular}
\end{adjustbox}
\end{center}

\begin{center}
\begin{minipage}{\textwidth}\centering
\begin{tikzpicture}
\begin{groupplot}[
  group style={group size=2 by 1, horizontal sep=0.60cm},
  width=6.85cm, height=4.60cm,
  tick label style={font=\tiny},
  title style={font=\scriptsize\heiti},
  yticklabel style={font=\tiny},
  ylabel style={font=\tiny},
  grid=major, grid style={LightGray, line width=0.3pt},
  axis line style={InkGray, line width=0.5pt},
  enlarge x limits={abs=0.8},
  legend style={font=\scriptsize, at={(0.5,-0.20)}, anchor=north,
                legend columns=2, draw=none, fill=none, column sep=6pt},
]
\nextgroupplot[
  ybar, bar width=4pt,
  ymin=-2.30, ymax=22.71,
  xtick={0,1,2,3,4,5.7},
  xticklabels={2021,2022,2023,2024,2025,2026H1},
  ylabel={亿元},
]
\addplot[fill=AgriGreen!75, draw=AgriGreen, bar shift=-2.6pt] table[x=i, y=rev]{data/clean/profiles/fin_003030.dat};
\addplot[fill=SoftRed!60, draw=SoftRed, bar shift=2.6pt] table[x=i, y=ni]{data/clean/profiles/fin_003030.dat};
\addplot[fill=AgriGreen!30, draw=AgriGreen, pattern=north east lines, bar shift=-2.6pt] table[x=x, y=rev]{data/clean/profiles/finh_003030.dat};
\addplot[fill=SoftRed!20, draw=SoftRed, pattern=north east lines, bar shift=2.6pt] table[x=x, y=ni]{data/clean/profiles/finh_003030.dat};
\legend{营业收入（年/半年）, 归母净利润（年/半年）}
\nextgroupplot[
  ymin=-7.3, ymax=64.0,
  xtick={0,1,2,3,4,5.7},
  xticklabels={2021,2022,2023,2024,2025,2026H1},
  ylabel={\%},
]
\addplot[AgriGold, mark=*, mark size=1.3pt] table[x=i, y=roe]{data/clean/profiles/fin_003030.dat};
\addplot[OilBlue, mark=square*, mark size=1.3pt] table[x=i, y=debt]{data/clean/profiles/fin_003030.dat};
\addplot[only marks, AgriGold, mark=*, mark size=2.0pt] table[x=x, y=roe]{data/clean/profiles/finh_003030.dat};
\addplot[only marks, OilBlue, mark=square*, mark size=2.0pt] table[x=x, y=debt]{data/clean/profiles/finh_003030.dat};
\legend{ROE, 资产负债率}
\end{groupplot}
\end{tikzpicture}
\begin{tikzpicture}
\begin{axis}[
  name=finbot,
  width=13.75cm, height=3.10cm, scale only axis,
  ymin=-1.00, ymax=11.23,
  xtick={0,1,2,3,4,5.7},
  xticklabels={2021,2022,2023,2024,2025,2026H1},
  tick label style={font=\tiny},
  yticklabel style={font=\tiny},
  ylabel={亿元}, ylabel style={font=\tiny},
  ybar, bar width=4pt,
  grid=major, grid style={LightGray, line width=0.3pt},
  axis line style={InkGray, line width=0.5pt},
  enlarge x limits={abs=0.8},
  legend style={font=\scriptsize, at={(0.5,-0.26)}, anchor=north,
                legend columns=2, draw=none, fill=none, column sep=10pt},
]
\addplot[fill=AgriGreen!55, draw=AgriGreen, bar shift=-3.0pt] table[x=i, y=cash]{data/clean/profiles/fin_003030.dat};
\addplot[fill=SoftRed!45, draw=SoftRed, bar shift=3.0pt] table[x=i, y=ibd]{data/clean/profiles/fin_003030.dat};
\addplot[fill=AgriGreen!25, draw=AgriGreen, pattern=north east lines, bar shift=-3.0pt] table[x=x, y=cash]{data/clean/profiles/finh_003030.dat};
\addplot[fill=SoftRed!20, draw=SoftRed, pattern=north east lines, bar shift=3.0pt] table[x=x, y=ibd]{data/clean/profiles/finh_003030.dat};
\legend{货币资金（左轴）, 有息负债（左轴）}
\end{axis}
\begin{axis}[
  at={(finbot.south east)}, anchor=south east,
  width=13.75cm, height=3.10cm, scale only axis,
  axis y line*=right, axis x line*=none, xtick=\empty,
  ymin=-0.50, ymax=2.51,
  tick label style={font=\tiny},
  yticklabel style={font=\tiny}, scaled y ticks=false,
  legend style={font=\scriptsize, at={(0.5,-0.44)}, anchor=north,
                legend columns=1, draw=none, fill=none},
]
\addplot[OilBlue, mark=*, mark size=2.2pt, line width=1.3pt] table[x=i, y=ocf]{data/clean/profiles/fin_003030.dat};
\addplot[only marks, OilBlue, mark=*, mark size=3.2pt] table[x=x, y=ocf]{data/clean/profiles/finh_003030.dat};
\legend{经营活动现金流净额（右轴，亿元，蓝点）}
\end{axis}
\end{tikzpicture}

\captionof{figure}[祖名股份（003030）近年主要财务指标]{祖名股份（003030）近年主要财务指标（左上：营业收入与归母净利润；右上：ROE 与资产负债率；下：货币资金、有息负债为左轴柱状，经营活动现金流净额为其右轴的蓝点折线。
2021---2025 年为年报口径，2026H1 为 2026 年半年报/半年末，与其前一年报之间留出间隔、以斜纹或空心点区分；半年报数据未年化）}
\end{minipage}
\end{center}

\pfl{财务分析} 2026 年 6 月末资产负债率 57.9\%，较 2025 年末上升 3.3 个百分点（样本中位数 52.2\%，所处三级行业内按由低到高排第\zhnumber{9}位）。 2026 年 6 月末货币资金 1.80 亿元、短期债务（短期借款与一年内到期非流动负债）8.74 亿元、长期债务 1.28 亿元，有息负债合计 10.02 亿元，现金短债比 0.21 倍——缺口明显。 净有息负债 8.22 亿元，相当于其 2026 年上半年营业收入的 81\%。 流动比率 0.54、速动比率 0.44，短期偿债指标偏紧。 2026 年上半年经营活动现金流净额 0.58 亿元，净利润为负而经营现金流为正，差额主要来自折旧摊销与营运资本变动，说明亏损尚未完全转化为现金流出。 综合看，账面货币资金对有息负债与短期债务的覆盖不足，滚动偿付对新增融资的依赖度较高；资产负债率高于样本中位数 5.7 个百分点；有息负债中短期债务占比更高，期限结构仍以滚动续作为主。 公司披露的股东与质押风险为：2025 年 12 月，公司控股股东、实际控制人及其一致行动人权益变动触及 5\% 整数倍的提示性公告。

\pfl{重大战略转型} 从公开披露看，公司的战略推进集中在：其他动作（2 项）：2026 年 6 月，紧扣国家十五五规划中生活消费品提质升级、健康消费扩容增效的核心主线，以外延品类（植物蛋白饮品同比增 41.88\%）带动收入增长；主业转型（1 项）：上述项目，公司以 10.12\% 持股比例入股重庆芯联芯。 公告口径上，近三年（2023 年 9 月---2026 年 9 月）公司披露重大重组与并购类公告 3 条、对外投资与转型类公告 2 条、回购与股权激励类公告 27 条，合计公告 349 条。 第一大产品仍是“生鲜豆制品”，收入占比 67.3\%，与 2021 年年报相比没有方向性变化。 综合判断，报告期内公司有 5 条与战略相关的公告落地（重大重组与并购 3 条、对外投资与转型 2 条），资本运作与主业调整之间存在直接关联。

\pfl{历史融投资情况} 从现金流量表看，2025 年公司取得借款收到的现金 7.04 亿元、偿还债务支付的现金 6.67 亿元、吸收权益性投资 0.00 亿元，筹资活动现金流净额 -0.59 亿元（2023 年为取得借款 5.14 亿元、偿还债务 4.00 亿元）。 2026 年上半年筹资活动现金流净额为 1.20 亿元（取得借款 5.02 亿元、偿还债务 3.94 亿元，未年化），仍处于净融入状态。 2025 年分配股利、利润或偿付利息支付的现金合计 0.48 亿元。 公开披露的融资记录共 2 笔，工具类型以 IPO、定向增发为主；金额与用途见下表。 其中规模最大的两笔为：2021 年 1 月，IPO，首次公开发行人民币普通股（A 股）股票 3,120 万股，每股发行价 15.18 元/股；2026 年 4 月，定向增发，拟发行融资总额不超过人民币 3.00 亿元且不超过最近一年末净资产 20\% 的股票。 股本方面，近三年披露股本变动 4 次（限售股份上市 1 次）；总股本自 2021 年末以来未发生变化（1.25 亿股）。 分红方面，近三年现金分红 4 次、合计每股派现 0.60 元，最近一次实施在 2025 年 12 月。 综合看，公司融资结构上以债务滚动为主、股权融资为补充；2025 年筹资活动现金流净额 -0.59 亿元，融资节奏仍以维持存量债务周转为主，与前述经营与投资现金流的方向基本一致。

\pfl{未来潜在融资需求分析} 2026 年 6 月末货币资金 1.80 亿元，而短期债务 8.74 亿元（现金短债比 0.21 倍），2026 年上半年经营现金流净额 0.58 亿元。按“短期债务减去账面货币资金”的滚动偿付口径，静态缺口约 6.94 亿元；若再计入上半年亏损的年化额（0.49 亿元），维持一年正常经营所需的外部资金约 6.94---7.43 亿元。 公司 2026 年上半年归母净利润为负（-0.24 亿元），但 6 月末资产负债率 57.9\% 尚未进入第一档（亏损且负债率 $>$65\%）的区间，当前融资需求以补充流动资金、支撑低谷期经营性支出为主。 公开披露的后续资金安排线索包括：股权融资线索：2026 年 4 月，董事会提请股东会授权董事会办理小额快速融资，向特定对象发行融资总额不超过人民币 3.00 亿元且不超过最近一年末净资产 20\% 的股票；债务与授信安排：2026 年上半年，公司在货币资金 1.80 亿元、短期借款 5.83 亿元的背景下持续依赖银行债务融资；2026 年 8 月，公司新增对子公司担保额度预计并提交股东会审议通过公司仍通过担保与银行授信为子公司经营提供资金支持。 综合看，公司的外部资金需求以补充营运资金为主、项目投入为辅，滚动偿付口径下的静态缺口约 6.94 亿元，能否落地取决于授信续作与发行窗口的配合。




\subsubsection{索宝蛋白（603231）}
\label{co:603231}

\pfl{公司基本情况} 宁波的大豆蛋白深加工企业，2023 年 12 月 15 日在上交所主板上市，产品涵盖大豆分离蛋白、大豆浓缩蛋白、组织化蛋白与非转基因大豆油，主要生产子公司为山东万得福生物科技有限公司，控股股东山东万得福实业集团、实际控制人刘季善。IPO 净额 9.79 亿元、超募 4.25 亿元，2025 年营收 17.48 亿元、归母净利润 1.90 亿元同比增 56.56\%；2026 年上半年增收不增利，且多个募投项目出现延期与变更。 公司于 2023-12-15 在上交所首发上市，注册资本 1.91 亿元，总股本 1.91 亿股。 截至 2026 年 10 月 12 日收盘价 16.92 元，对应总市值 32.4 亿元，近一年-10.2\%，总市值在三级行业“粮油加工”的 7 家公司中按由大到小排第\zhnumber{5}位，近一年涨跌幅高于全样本中位数（-10.4\%）0.1 个百分点。

\pfl{股权结构与控制权} 按 2026 年 6 月 30 日的股东名册，前三大股东为山东万得福实业集团有限公司 持股 43.22\%；上海复星创富投资管理股份有限公司－宁波梅山保税港区复 持股 5.63\%；济南财金复星惟实股权投资基金合伙企业(有限合伙) 持股 5.63\%，股东户数 14,480 户。 股权结构的另一条线索是，公司增持前持有公司股票 8,274.81 万股、占公司总股本 43.22\%；截至 2026 年 6 月 30 日公司股份总数 1.91 亿股；上市时万得福集团由刘季善出资 3,884.61 万元、占其出资比例 53.77\%。

\pfl{经营要点} 经营与产能方面，公司持有山东索宝元启特医食品有限公司 70\% 股权。

\pfl{所处行业的基本面} 粮油加工的供给由大豆、菜籽等原料进口与压榨产能决定，需求由食用消费决定，利润来自原料与成品的价差而非价格方向（第 \ref{sec:framework} 章的供需框架）。 与该公司直接相关的行业数据包括：2026 年，中央一号文件将大豆纳入新一轮千亿斤粮食产能提升行动；大豆蛋白市场份额向优势企业集中。

\pfl{竞争优势} 把该公司放回同类公司比较，2026 年上半年营业收入 9.7 亿元、在三级行业“粮油加工”的 7 家公司中排第\zhnumber{6}位；毛利率 12.9\%、高于行业内中位数 6.6 个百分点；半年 ROE 3.4\%（未年化）、高于行业内中位数（2.22\%）；资产负债率 13.6\%、低于行业内中位数（43.5\%）。 公司在公开资料中自述与同业比较的相对优势在于：具备从大豆采购、低温浸出到分离蛋白、浓缩蛋白、组织化蛋白、膳食纤维的全系列生产线；上市企业中公司大豆蛋白业务处于行业领先地位；采购价格随市场行情波动。 优势在报表上主要体现为高于行业平均的毛利率，这在周期底部是抵御亏损的主要缓冲。

\pfl{历史股价} 以 2023 年 12 月为起点、2026 年 10 月为终点，股价由 23.52 元变为 16.92 元，累计 -28.1\%，区间最高 23.52 元（2023 年 12 月）、最低 14.76 元（2026 年 7 月）；当前价相当于区间最高价的 71.9\%，为区间最低价的 1.15 倍。 月度收益的年化波动率 28.2\%，低于全样本中位数 37.3\%，在三级行业“粮油加工”的 7 家公司中波动率由高到低排第\zhnumber{5}位。

\begin{center}
\begin{minipage}{\textwidth}\centering
\begin{tikzpicture}
\begin{axis}[
  width=12.6cm, height=3.9cm, scale only axis,
  xmin=2023.838, xmax=2026.972, ymin=13.73, ymax=25.17,
  xtick={2023,2024,2025,2026},
  yearticks,
  yticklabel style={font=\scriptsize},
  ylabel={元/股}, ylabel style={font=\scriptsize},
  grid=major, grid style={LightGray, line width=0.3pt},
  axis line style={InkGray, line width=0.5pt},
  every axis plot/.append style={line width=0.9pt},
]
\addplot[AgriGreen] table[x=x,y=y]{data/clean/profiles/px_603231.dat};
\addplot[SoftRed, dashed, line width=0.7pt] table[x=x,y=y]{data/clean/profiles/pxzz_603231.dat};
\addplot[only marks, mark=*, mark size=1.1pt, SoftRed] table[x=x,y=y]{data/clean/profiles/pxzz_603231.dat};
\end{axis}
\end{tikzpicture}

\captionof{figure}[索宝蛋白（603231）股价与周期骨架]{索宝蛋白（603231）月度收盘价（元/股）与 ZigZag 周期骨架（阈值 25\%，圆点为周期转折点）}
\end{minipage}
\end{center}

\pfl{过去周期分析} 以 25\% 的 ZigZag 阈值划分，样本区间内股价形成 1 个主升段与 2 个主跌段。 最大主跌段为 2023 年 12 月 至 2024 年 8 月，8 个月-35.5\%；最大主升段出现在 2024 年 8 月 至 2025 年 10 月，14 个月+33.8\%。 最近一次转折出现在 2026 年 7 月（下跌段自 2025 年 10 月起，9 个月-27.3\%），当前处于该段的延续之中。 与申万农林牧渔指数的最近一次底部相比，该公司见底晚约 2 个月，在本轮下行中属于调整更慢的一类。 粮油加工是“成本加成”型生意（第 \ref{sec:grain} 章），利润来自原料与成品的价差及规模效率，收入规模大而净利率薄，公司 2026 年上半年实现盈利、半年 ROE 3.4\%（未年化），好于全样本中位数（0.75\%），下行周期对其报表的冲击小于同业。 综合区间形态看，该股单段涨跌幅度相对温和、以震荡为主（最大主升 +33.8\%、最大主跌 -35.5\%），年化波动率 28.2\% 与全样本中位数 37.3\% 相比波动小于样本中位数。

\pfl{盈利情况} 2026 年上半年营业收入 9.67 亿元（同比 +24.9\%），归母净利润 0.68 亿元，同比 -30.1\%。 以年报为趋势对照，2023---2025 年营业收入由 17.48 亿元变为 17.48 亿元（两年年均复合增速 -0.0\%），归母净利润由 1.46 亿元变为 1.90 亿元。 按半年报口径，2026 年上半年的 ROE 为 3.4\%（未年化）、销售净利率 6.9\%、毛利率 12.9\%，ROE 在“粮油加工”7 家公司中按数值由高到低排第\zhnumber{2}位。 同期全样本半年 ROE 中位数 0.75\%，公司与之相差 2.62 个百分点（略好于中位数公司）。 从公司披露的原因看，业绩变动主要来自：2025 年 4 月，由于受到宏观经济波动、原材料价格波动等诸多因素的影响。 综合看，公司在报告期维持盈利（高于全样本半年 ROE 中位数 0.75\%），利润的可持续性取决于粮油加工景气的延续与成本管控的执行。

\begin{center}\small\setlength{\tabcolsep}{3.4pt}
\captionof{table}[索宝蛋白（603231）关键财务指标]{索宝蛋白（603231）关键财务指标（合并报表口径；两个半年报期均未年化；现金短债比 $=$ 货币资金 $\div$ 短期债务）}
\begin{adjustbox}{max width=\textwidth}
\begin{tabular}{@{}lrrrrr@{}}
\toprule
指标 & 2023 年 & 2024 年 & 2025 年 & 2025 年半年报 & 2026 年半年报 \\
\midrule
营业收入（亿元） & 17.48 & 15.55 & 17.48 & 7.74 & 9.67 \\
归母净利润（亿元） & 1.46 & 1.21 & 1.90 & 0.97 & 0.68 \\
毛利率（\%） & 13.6 & 12.0 & 16.7 & 18.7 & 12.9 \\
ROE（\%） & 16.3 & 6.4 & 9.8 & 5.1 & 3.4 \\
资产负债率（\%） & 11.7 & 10.1 & 11.6 & 8.7 & 13.6 \\
经营活动现金流净额（亿元） & 1.44 & 1.48 & 3.11 & 0.29 & 0.08 \\
货币资金（亿元） & 10.73 & 3.20 & 2.39 & 2.61 & 2.84 \\
有息负债合计（亿元） & 0.85 & 0.24 & 0.24 & 0.24 & 0.67 \\
现金短债比（倍） & 12.69 & 13.07 & 9.84 & 10.94 & 4.39 \\
\bottomrule
\end{tabular}
\end{adjustbox}
\end{center}

\begin{center}
\begin{minipage}{\textwidth}\centering
\begin{tikzpicture}
\begin{groupplot}[
  group style={group size=2 by 1, horizontal sep=0.60cm},
  width=6.85cm, height=4.60cm,
  tick label style={font=\tiny},
  title style={font=\scriptsize\heiti},
  yticklabel style={font=\tiny},
  ylabel style={font=\tiny},
  grid=major, grid style={LightGray, line width=0.3pt},
  axis line style={InkGray, line width=0.5pt},
  enlarge x limits={abs=0.8},
  legend style={font=\scriptsize, at={(0.5,-0.20)}, anchor=north,
                legend columns=2, draw=none, fill=none, column sep=6pt},
]
\nextgroupplot[
  ybar, bar width=4pt,
  ymin=-1.85, ymax=20.68,
  xtick={0,1,2,3,4,5.7},
  xticklabels={2021,2022,2023,2024,2025,2026H1},
  ylabel={亿元},
]
\addplot[fill=AgriGreen!75, draw=AgriGreen, bar shift=-2.6pt] table[x=i, y=rev]{data/clean/profiles/fin_603231.dat};
\addplot[fill=SoftRed!60, draw=SoftRed, bar shift=2.6pt] table[x=i, y=ni]{data/clean/profiles/fin_603231.dat};
\addplot[fill=AgriGreen!30, draw=AgriGreen, pattern=north east lines, bar shift=-2.6pt] table[x=x, y=rev]{data/clean/profiles/finh_603231.dat};
\addplot[fill=SoftRed!20, draw=SoftRed, pattern=north east lines, bar shift=2.6pt] table[x=x, y=ni]{data/clean/profiles/finh_603231.dat};
\legend{营业收入（年/半年）, 归母净利润（年/半年）}
\nextgroupplot[
  ymin=-2.4, ymax=32.9,
  xtick={0,1,2,3,4,5.7},
  xticklabels={2021,2022,2023,2024,2025,2026H1},
  ylabel={\%},
]
\addplot[AgriGold, mark=*, mark size=1.3pt] table[x=i, y=roe]{data/clean/profiles/fin_603231.dat};
\addplot[OilBlue, mark=square*, mark size=1.3pt] table[x=i, y=debt]{data/clean/profiles/fin_603231.dat};
\addplot[only marks, AgriGold, mark=*, mark size=2.0pt] table[x=x, y=roe]{data/clean/profiles/finh_603231.dat};
\addplot[only marks, OilBlue, mark=square*, mark size=2.0pt] table[x=x, y=debt]{data/clean/profiles/finh_603231.dat};
\legend{ROE, 资产负债率}
\end{groupplot}
\end{tikzpicture}
\begin{tikzpicture}
\begin{axis}[
  name=finbot,
  width=13.75cm, height=3.10cm, scale only axis,
  ymin=-1.07, ymax=12.02,
  xtick={0,1,2,3,4,5.7},
  xticklabels={2021,2022,2023,2024,2025,2026H1},
  tick label style={font=\tiny},
  yticklabel style={font=\tiny},
  ylabel={亿元}, ylabel style={font=\tiny},
  ybar, bar width=4pt,
  grid=major, grid style={LightGray, line width=0.3pt},
  axis line style={InkGray, line width=0.5pt},
  enlarge x limits={abs=0.8},
  legend style={font=\scriptsize, at={(0.5,-0.26)}, anchor=north,
                legend columns=2, draw=none, fill=none, column sep=10pt},
]
\addplot[fill=AgriGreen!55, draw=AgriGreen, bar shift=-3.0pt] table[x=i, y=cash]{data/clean/profiles/fin_603231.dat};
\addplot[fill=SoftRed!45, draw=SoftRed, bar shift=3.0pt] table[x=i, y=ibd]{data/clean/profiles/fin_603231.dat};
\addplot[fill=AgriGreen!25, draw=AgriGreen, pattern=north east lines, bar shift=-3.0pt] table[x=x, y=cash]{data/clean/profiles/finh_603231.dat};
\addplot[fill=SoftRed!20, draw=SoftRed, pattern=north east lines, bar shift=3.0pt] table[x=x, y=ibd]{data/clean/profiles/finh_603231.dat};
\legend{货币资金（左轴）, 有息负债（左轴）}
\end{axis}
\begin{axis}[
  at={(finbot.south east)}, anchor=south east,
  width=13.75cm, height=3.10cm, scale only axis,
  axis y line*=right, axis x line*=none, xtick=\empty,
  ymin=-0.81, ymax=4.04,
  tick label style={font=\tiny},
  yticklabel style={font=\tiny}, scaled y ticks=false,
  legend style={font=\scriptsize, at={(0.5,-0.44)}, anchor=north,
                legend columns=1, draw=none, fill=none},
]
\addplot[OilBlue, mark=*, mark size=2.2pt, line width=1.3pt] table[x=i, y=ocf]{data/clean/profiles/fin_603231.dat};
\addplot[only marks, OilBlue, mark=*, mark size=3.2pt] table[x=x, y=ocf]{data/clean/profiles/finh_603231.dat};
\legend{经营活动现金流净额（右轴，亿元，蓝点）}
\end{axis}
\end{tikzpicture}

\captionof{figure}[索宝蛋白（603231）近年主要财务指标]{索宝蛋白（603231）近年主要财务指标（左上：营业收入与归母净利润；右上：ROE 与资产负债率；下：货币资金、有息负债为左轴柱状，经营活动现金流净额为其右轴的蓝点折线。
2021---2025 年为年报口径，2026H1 为 2026 年半年报/半年末，与其前一年报之间留出间隔、以斜纹或空心点区分；半年报数据未年化）}
\end{minipage}
\end{center}

\pfl{财务分析} 2026 年 6 月末资产负债率 13.6\%，较 2025 年末上升 2.0 个百分点（样本中位数 52.2\%，所处三级行业内按由低到高排第\zhnumber{1}位）。 2026 年 6 月末货币资金 2.84 亿元、短期债务（短期借款与一年内到期非流动负债）0.65 亿元、长期债务 0.02 亿元，有息负债合计 0.67 亿元，现金短债比 4.39 倍——覆盖充分。 净有息负债 -2.17 亿元，相当于其 2026 年上半年营业收入的 -22\%。 流动比率 5.42、速动比率 4.16，短期流动性无明显缺口。 2026 年上半年经营活动现金流净额 0.08 亿元，经营现金流与利润同向为正，内生资金可以支撑运营。 综合看，现金短债比 4.39 倍，短期偿付压力可控；资产负债率低于样本中位数 38.6 个百分点；有息负债中短期债务占比更高，期限结构仍以滚动续作为主。

\pfl{重大战略转型} 公司的战略动作可归为以下几类：融资与资本运作（1 项）：2025 年 1 月，公司就募投项目 3 万吨大豆组织拉丝蛋白生产线建设项目、5,000 吨大豆颗粒蛋白生产线建设项目进行重新论证并继续实施；产能与项目（1 项）：2025 年 11 月，股东大会通过募投项目变更并向全资子公司山东万得福生物科技有限公司增资以实施万得福工业园高端大豆分离蛋白生产线改造项目。 公告口径上，近三年（2023 年 9 月---2026 年 9 月）公司披露重大重组与并购类公告 1 条、对外投资与转型类公告 2 条、回购与股权激励类公告 0 条，合计公告 283 条。 综合判断，报告期内公司有 3 条与战略相关的公告落地（重大重组与并购 1 条、对外投资与转型 2 条），资本运作与主业调整之间存在直接关联。

\begin{center}\small\setlength{\tabcolsep}{4pt}
\captionof{table}[索宝蛋白（603231）历史融资与资本运作]{索宝蛋白（603231）历史融资与资本运作（据公司公告与公开报道整理；金额为披露值，含首次公开发行、再融资、债券、银行借款与重整投资等）}
\begin{adjustbox}{max width=\textwidth}
\begin{tabular}{@{}p{1.45cm}p{2.55cm}p{4.05cm}p{4.95cm}@{}}
\toprule
时间 & 方式 & 规模 & 用途与说明 \\
\midrule
2023-12-15 & IPO（上海证券交易所主板） & 首次公开发行人民币普通股（A 股）4,\allowbreak{}786.48 万股（无老股转让，\allowbreak{}全部为公开发行新股） & 3 万吨大豆组织拉丝蛋白生产线建设项目（募投计划投入 1.20 亿元）、\allowbreak{}5,\allowbreak{}000 吨大豆颗粒蛋白生产线建设项目（1.20 亿元）、\allowbreak{}75T 中温中压高效煤粉锅炉项目（1.45 亿元）、\allowbreak{}补充流动资金（1.70 亿元） \\
2024-12-02 & 闲置募集资金现金管理（非股权融资） & 同意使用最高额度不超过人民币 5.50 亿元（含本数）的暂时闲置募集资金进行现金管理 & 提高募集资金使用效率；第二届董事会第十七次会议和第二届监事会第十三次会议审议通过 \\
2026-04-24 & 银行综合授信额度 & 公司及控股子公司计划向银行申请总额不超过人民币 5.00 亿元的综合授信额度 & 满足公司发展需要、\allowbreak{}拓宽融资渠道；经第三届董事会第八次会议审议通过 2026 年度向金融机构申请综合授信额度；授信期限自公司 2025 年年度股东会召开之日起至 2026 年年度股东会召开之日止；授信额度最终以银行实际审批金额为准 \\
\bottomrule
\end{tabular}
\end{adjustbox}
\end{center}

\pfl{历史融投资情况} 从现金流量表看，2025 年公司取得借款收到的现金 0.48 亿元、偿还债务支付的现金 0.48 亿元，筹资活动现金流净额 -0.59 亿元（2023 年为取得借款 0.94 亿元、偿还债务 0.95 亿元）。 2026 年上半年筹资活动现金流净额为 -0.30 亿元（取得借款 0.84 亿元、偿还债务 0.44 亿元，未年化），已转为净流出，处于净还债状态。 2025 年分配股利、利润或偿付利息支付的现金合计 0.57 亿元。 公开披露的融资记录共 3 笔，工具类型以 IPO、银行授信为主；金额与用途见下表。 其中规模最大的两笔为：2026 年 4 月，银行授信，公司及控股子公司计划向银行申请总额不超过人民币 5.00 亿元的综合授信额度；2024 年 12 月，闲置募集资金现金管理，同意使用最高额度不超过人民币 5.50 亿元（含本数）的暂时闲置募集资金进行现金管理。 股本方面，近三年披露股本变动 4 次（A 股上市 1 次、限售股份上市 1 次）；总股本自 2021 年末以来未发生变化（1.91 亿股）。 分红方面，近三年现金分红 4 次、合计每股派现 1.17 元，最近一次实施在 2026 年 6 月。 综合看，公司融资结构上股权与债务融资并举；2025 年筹资活动现金流净额 -0.59 亿元，融资节奏仍以维持存量债务周转为主，与前述经营与投资现金流的方向基本一致。

\pfl{未来潜在融资需求分析} 2026 年 6 月末货币资金 2.84 亿元，而短期债务 0.65 亿元（现金短债比 4.39 倍），2026 年上半年经营现金流净额 0.08 亿元。账面货币资金可以覆盖短期债务，缺口不体现在静态偿付层面。 公司 2026 年 6 月末资产负债率 13.6\%、2026 年上半年 ROE 3.4\%（未年化），财务结构相对宽松而盈利弹性有限，融资属于可选项而非必选项，触发条件多与产能建设或外延并购相关。 公开披露的后续资金安排线索包括：债务与授信安排：2026 年 4 月，公司及控股子公司计划向银行申请总额不超过人民币 5.00 亿元的综合授信额度；资本开支安排：2026 年 4 月，公司超募资金中尚有 2.98 亿元暂未确定投向，为后续资本开支预留了内部资金。 综合看，公司在静态偿付层面不存在缺口，外部融资更多是配合扩张节奏与并购整合的主动性安排，而非偿债压力驱动。




\subsubsection{中基健康（000972）}
\label{co:000972}

\pfl{公司基本情况} 1994 年设立至今，中基健康是新疆兵团第六师国资控股的番茄制品企业，主营大包装番茄酱、小包装番茄制品与番茄红素；2024 年因净资产为负、扣除后营收不足 3 亿元被实施退市风险警示，2025 年 7 月进入预重整，2025 年底靠政府补助与债务豁免净资产转正并于 2026 年 6 月摘帽，但 2026 年中期净资产再度转负。 公司于 2000-09-26 在深交所主板首发上市，注册资本 7.71 亿元，总股本 7.71 亿股。 截至 2026 年 10 月 12 日收盘价 4.10 元，对应总市值 31.6 亿元，近一年+12.3\%，总市值在三级行业“果蔬加工”的 5 家公司中按由大到小排第\zhnumber{4}位，近一年涨跌幅高于全样本中位数（-10.4\%）22.7 个百分点。 按 2026 年半年报披露的行业构成，收入占比前三位为制造业 95.3\%；其他 4.5\%；租赁服务 0.3\%。

\pfl{股权结构与控制权} 按 2026 年 6 月 30 日的股东名册，前三大股东为新疆生产建设兵团第六师国有资产经营有限责任公司 持股 16.18\%；新疆国恒投资发展集团有限公司 持股 12.97\%；新疆生产建设兵团投资有限责任公司 持股 6.50\%，股东户数 22,630 户。 控制权与股东性质方面，是兵团第六师旗下的番茄制品上市平台；2025 年 7 月，公司及全资子公司新疆中基红色番茄产业有限公司分别收到六师中院预重整通知书。

\pfl{经营要点} 业务与产能层面，2024 年，出口量及价格水平整体下滑。

\pfl{所处行业的基本面} 果蔬加工的价格由产地原料果价格与出口订单决定，需求端受海外需求与关税影响（第 \ref{sec:framework} 章的供需框架）。 从供需两端的当期数据看，2025 年，国内大桶番茄酱价格持续低迷，进一步加剧价格下行压力；2024 年，中国加工番茄产量同比增长 30.63\%、增速显著高于全球。

\pfl{竞争优势} 把该公司放回同类公司比较，2026 年上半年营业收入 2.0 亿元、在三级行业“果蔬加工”的 5 家公司中排第\zhnumber{4}位；毛利率 -19.8\%、低于行业内中位数 28.4 个百分点；资产负债率 104.1\%、高于行业内中位数（30.8\%）。 与同类公司相比，公司的差异化在于：兵团加工番茄种植面积占新疆番茄种植面积的近 60\%。 就当期报表看，公司尚未把经营层面的特征转化为高于行业的资本回报，周期底部的相对优势更多体现在现金流与资产负债表上。

\pfl{历史股价} 自 2021 年 1 月至 2026 年 10 月，股价由 1.80 元变为 4.10 元，累计 +127.8\%，区间最高 4.45 元（2025 年 8 月）、最低 1.68 元（2021 年 2 月）；当前价相当于区间最高价的 92.1\%，为区间最低价的 2.44 倍。 月度收益的年化波动率 43.6\%，高于全样本中位数 37.3\%，在三级行业“果蔬加工”的 5 家公司中波动率由高到低排第\zhnumber{4}位。 把股价阶段与公司事件对齐看，2022 年 4 月 至 2023 年 1 月股价上涨+139.0\%，同期（2022 年 8 月）构成关联交易及重大资产重组，不导致实际控制人变更；2024 年 6 月 至 2025 年 8 月股价上涨+99.6\%，同期（2025 年 8 月）临时管理人公开招募重整投资人，分别引入产业投资人和财务投资人。 把这些阶段与公司事件对齐后可以看到，几轮大涨大跌多由公司自身的资本运作与风险事件触发，行业景气只解释了其中一部分。

\begin{center}
\begin{minipage}{\textwidth}\centering
\begin{tikzpicture}
\begin{axis}[
  width=12.6cm, height=3.9cm, scale only axis,
  xmin=2020.922, xmax=2026.972, ymin=1.56, ymax=4.76,
  xtick={2021,2022,2023,2024,2025,2026},
  yearticks,
  yticklabel style={font=\scriptsize},
  ylabel={元/股}, ylabel style={font=\scriptsize},
  grid=major, grid style={LightGray, line width=0.3pt},
  axis line style={InkGray, line width=0.5pt},
  every axis plot/.append style={line width=0.9pt},
]
\addplot[AgriGreen] table[x=x,y=y]{data/clean/profiles/px_000972.dat};
\addplot[SoftRed, dashed, line width=0.7pt] table[x=x,y=y]{data/clean/profiles/pxzz_000972.dat};
\addplot[only marks, mark=*, mark size=1.1pt, SoftRed] table[x=x,y=y]{data/clean/profiles/pxzz_000972.dat};
\end{axis}
\end{tikzpicture}

\captionof{figure}[中基健康（000972）股价与周期骨架]{中基健康（000972）月度收盘价（元/股）与 ZigZag 周期骨架（阈值 25\%，圆点为周期转折点）}
\end{minipage}
\end{center}

\pfl{过去周期分析} 以 25\% 的 ZigZag 阈值划分，样本区间内股价形成 4 个主升段与 3 个主跌段。 最大主升段出现在 2022 年 4 月 至 2023 年 1 月，9 个月+139.0\%；最大主跌段为 2023 年 1 月 至 2024 年 6 月，17 个月-47.3\%。 最近一次转折出现在 2026 年 10 月（上涨段自 2026 年 6 月起，4 个月+31.0\%），当前处于该段之后的震荡之中。 公司股价较申万农林牧渔指数提前约 63 个月见底，是板块中较早定价周期反转的公司之一。 果蔬加工的外向型特征明显，出口订单与原料产地价格决定利润率，公司 2026 年 6 月末资产负债率 104.1\%，高于全样本中位数 52.2\%，其股价因此更多反映资金面与信用风险，而不是价格弹性本身。 综合区间形态看，该股单段涨跌幅度偏大、以趋势段为主（最大主升 +139.0\%、最大主跌 -47.3\%），年化波动率 43.6\% 与全样本中位数 37.3\% 相比波动明显大于样本中位数。

\pfl{盈利情况} 2026 年上半年营业收入 1.97 亿元（同比 -19.6\%），归母净利润 -0.66 亿元，亏损同比收窄 0.09 亿元。 以年报为趋势对照，2023---2025 年营业收入由 5.76 亿元变为 4.92 亿元（两年年均复合增速 -7.6\%），归母净利润由 1.08 亿元变为 -0.46 亿元。 按半年报口径，2026 年上半年的 ROE 为 —\%（未年化）、销售净利率 -34.6\%、毛利率 -19.8\%。 综合看，公司仍处亏损区间、亏损同比收窄，单位盈利能力的修复依赖行业景气与自身成本端的改善，报表弹性主要体现为价格与销量的方向性变化。

\begin{center}\small\setlength{\tabcolsep}{3.4pt}
\captionof{table}[中基健康（000972）关键财务指标]{中基健康（000972）关键财务指标（合并报表口径；两个半年报期均未年化；现金短债比 $=$ 货币资金 $\div$ 短期债务）}
\begin{adjustbox}{max width=\textwidth}
\begin{tabular}{@{}lrrrrr@{}}
\toprule
指标 & 2023 年 & 2024 年 & 2025 年 & 2025 年半年报 & 2026 年半年报 \\
\midrule
营业收入（亿元） & 5.76 & 3.08 & 4.92 & 2.45 & 1.97 \\
归母净利润（亿元） & 1.08 & -2.31 & -0.46 & -0.75 & -0.66 \\
毛利率（\%） & 25.3 & 5.0 & -15.7 & -13.8 & -19.8 \\
ROE（\%） & 72.4 & 0.0 & --- & 0.0 & --- \\
资产负债率（\%） & 82.9 & 101.0 & 97.8 & 106.8 & 104.1 \\
经营活动现金流净额（亿元） & -1.82 & -0.98 & 2.47 & -1.01 & 1.27 \\
货币资金（亿元） & 1.11 & 1.60 & 3.13 & 0.66 & 3.01 \\
有息负债合计（亿元） & 4.28 & 6.58 & 6.77 & 7.15 & 5.54 \\
现金短债比（倍） & 0.31 & 0.27 & 0.50 & 0.10 & 0.60 \\
\bottomrule
\end{tabular}
\end{adjustbox}
\end{center}

\begin{center}
\begin{minipage}{\textwidth}\centering
\begin{tikzpicture}
\begin{groupplot}[
  group style={group size=2 by 1, horizontal sep=0.60cm},
  width=6.85cm, height=4.60cm,
  tick label style={font=\tiny},
  title style={font=\scriptsize\heiti},
  yticklabel style={font=\tiny},
  ylabel style={font=\tiny},
  grid=major, grid style={LightGray, line width=0.3pt},
  axis line style={InkGray, line width=0.5pt},
  enlarge x limits={abs=0.8},
  legend style={font=\scriptsize, at={(0.5,-0.20)}, anchor=north,
                legend columns=2, draw=none, fill=none, column sep=6pt},
]
\nextgroupplot[
  ybar, bar width=4pt,
  ymin=-3.13, ymax=6.88,
  xtick={0,1,2,3,4,5.7},
  xticklabels={2021,2022,2023,2024,2025,2026H1},
  ylabel={亿元},
]
\addplot[fill=AgriGreen!75, draw=AgriGreen, bar shift=-2.6pt] table[x=i, y=rev]{data/clean/profiles/fin_000972.dat};
\addplot[fill=SoftRed!60, draw=SoftRed, bar shift=2.6pt] table[x=i, y=ni]{data/clean/profiles/fin_000972.dat};
\addplot[fill=AgriGreen!30, draw=AgriGreen, pattern=north east lines, bar shift=-2.6pt] table[x=x, y=rev]{data/clean/profiles/finh_000972.dat};
\addplot[fill=SoftRed!20, draw=SoftRed, pattern=north east lines, bar shift=2.6pt] table[x=x, y=ni]{data/clean/profiles/finh_000972.dat};
\legend{营业收入（年/半年）, 归母净利润（年/半年）}
\nextgroupplot[
  ymin=-102.4, ymax=123.3,
  xtick={0,1,2,3,4,5.7},
  xticklabels={2021,2022,2023,2024,2025,2026H1},
  ylabel={\%},
]
\addplot[AgriGold, mark=*, mark size=1.3pt] table[x=i, y=roe]{data/clean/profiles/fin_000972.dat};
\addplot[OilBlue, mark=square*, mark size=1.3pt] table[x=i, y=debt]{data/clean/profiles/fin_000972.dat};
\addplot[only marks, AgriGold, mark=*, mark size=2.0pt] table[x=x, y=roe]{data/clean/profiles/finh_000972.dat};
\addplot[only marks, OilBlue, mark=square*, mark size=2.0pt] table[x=x, y=debt]{data/clean/profiles/finh_000972.dat};
\legend{ROE, 资产负债率}
\end{groupplot}
\end{tikzpicture}
\begin{tikzpicture}
\begin{axis}[
  name=finbot,
  width=13.75cm, height=3.10cm, scale only axis,
  ymin=-0.68, ymax=7.59,
  xtick={0,1,2,3,4,5.7},
  xticklabels={2021,2022,2023,2024,2025,2026H1},
  tick label style={font=\tiny},
  yticklabel style={font=\tiny},
  ylabel={亿元}, ylabel style={font=\tiny},
  ybar, bar width=4pt,
  grid=major, grid style={LightGray, line width=0.3pt},
  axis line style={InkGray, line width=0.5pt},
  enlarge x limits={abs=0.8},
  legend style={font=\scriptsize, at={(0.5,-0.26)}, anchor=north,
                legend columns=2, draw=none, fill=none, column sep=10pt},
]
\addplot[fill=AgriGreen!55, draw=AgriGreen, bar shift=-3.0pt] table[x=i, y=cash]{data/clean/profiles/fin_000972.dat};
\addplot[fill=SoftRed!45, draw=SoftRed, bar shift=3.0pt] table[x=i, y=ibd]{data/clean/profiles/fin_000972.dat};
\addplot[fill=AgriGreen!25, draw=AgriGreen, pattern=north east lines, bar shift=-3.0pt] table[x=x, y=cash]{data/clean/profiles/finh_000972.dat};
\addplot[fill=SoftRed!20, draw=SoftRed, pattern=north east lines, bar shift=3.0pt] table[x=x, y=ibd]{data/clean/profiles/finh_000972.dat};
\legend{货币资金（左轴）, 有息负债（左轴）}
\end{axis}
\begin{axis}[
  at={(finbot.south east)}, anchor=south east,
  width=13.75cm, height=3.10cm, scale only axis,
  axis y line*=right, axis x line*=none, xtick=\empty,
  ymin=-3.02, ymax=3.77,
  tick label style={font=\tiny},
  yticklabel style={font=\tiny}, scaled y ticks=false,
  legend style={font=\scriptsize, at={(0.5,-0.44)}, anchor=north,
                legend columns=1, draw=none, fill=none},
]
\addplot[OilBlue, mark=*, mark size=2.2pt, line width=1.3pt] table[x=i, y=ocf]{data/clean/profiles/fin_000972.dat};
\addplot[only marks, OilBlue, mark=*, mark size=3.2pt] table[x=x, y=ocf]{data/clean/profiles/finh_000972.dat};
\legend{经营活动现金流净额（右轴，亿元，蓝点）}
\end{axis}
\end{tikzpicture}

\captionof{figure}[中基健康（000972）近年主要财务指标]{中基健康（000972）近年主要财务指标（左上：营业收入与归母净利润；右上：ROE 与资产负债率；下：货币资金、有息负债为左轴柱状，经营活动现金流净额为其右轴的蓝点折线。
2021---2025 年为年报口径，2026H1 为 2026 年半年报/半年末，与其前一年报之间留出间隔、以斜纹或空心点区分；半年报数据未年化）}
\end{minipage}
\end{center}

\pfl{财务分析} 2026 年 6 月末资产负债率 104.1\%，较 2025 年末上升 6.3 个百分点（样本中位数 52.2\%，所处三级行业内按由低到高排第\zhnumber{5}位）。 2026 年 6 月末货币资金 3.01 亿元、短期债务（短期借款与一年内到期非流动负债）5.01 亿元、长期债务 0.53 亿元，有息负债合计 5.54 亿元，现金短债比 0.60 倍——覆盖偏紧。 净有息负债 2.54 亿元，相当于其 2026 年上半年营业收入的 129\%。 流动比率 0.64、速动比率 0.50，短期偿债指标偏紧。 2026 年上半年经营活动现金流净额 1.27 亿元，净利润为负而经营现金流为正，差额主要来自折旧摊销与营运资本变动，说明亏损尚未完全转化为现金流出。 综合看，现金短债比 0.60 倍，短期资金安排偏紧；资产负债率高于样本中位数 51.9 个百分点；有息负债中短期债务占比更高，期限结构仍以滚动续作为主。 公司披露的退市与持续经营风险为：2026 年 4 月，公司向深圳证券交易所申请撤销退市风险警示；2024 年年度报告，2024 年度利润总额、净利润、扣除非经常性损益后的净利润三者均为负值，且扣除后的营业收入低于 3.00 亿元。

\pfl{重大战略转型} 公司的战略动作可归为以下几类：并购与重组（3 项）：2025 年 9 月，公司将借助重整采取一系列系统性改善措施；2025 年 8 月，临时管理人公开招募重整投资人，分别引入产业投资人和财务投资人；其他动作（2 项）：2025 年度，归属于上市公司股东的净资产 2,612.23 万元（转正）；2025 年度，全资子公司红色番茄累计获得 3.17 亿元政府补助。 公告口径上，近三年（2023 年 9 月---2026 年 9 月）公司披露重大重组与并购类公告 16 条、对外投资与转型类公告 5 条、回购与股权激励类公告 8 条，合计公告 497 条。 第一大行业仍是“制造业”，收入占比 95.3\%，与 2021 年年报相比没有方向性变化。 综合判断，公司当前的战略重心不是扩张而是化解债务与完成重整，业务层面的调整让位于司法重整程序。

\begin{center}\small\setlength{\tabcolsep}{4pt}
\captionof{table}[中基健康（000972）历史融资与资本运作]{中基健康（000972）历史融资与资本运作（据公司公告与公开报道整理；金额为披露值，含首次公开发行、再融资、债券、银行借款与重整投资等）}
\begin{adjustbox}{max width=\textwidth}
\begin{tabular}{@{}p{1.45cm}p{2.55cm}p{4.05cm}p{4.95cm}@{}}
\toprule
时间 & 方式 & 规模 & 用途与说明 \\
\midrule
2000-09-26 & IPO（深交所） & 公开发行人民币普通股 4,\allowbreak{}500.00 万股 & 2000 年 4 月 14 日证监会证监发行字[2000]33 号文批准 \\
2022-08 & 发行股份购买资产并募集配套资金（预案，后终止） & 拟以发行股份方式向新疆中泰农业发展有限责任公司、\allowbreak{}新疆粮油集团有限责任公司、\allowbreak{}新疆艳阳天番茄制品有限责任公司购买其合计持有的新疆新粮艳阳天番茄股份有限公司 100.00\% 股份 & 收购新粮艳阳天 100\% 股份，\allowbreak{}配套资金用于相关项目；构成关联交易及重大资产重组 \\
2024-02-28 & 终止发行股份购买资产并募集配套资金 & 原交易对价 4.13 亿元 & 第九届董事会第四十次临时会议、\allowbreak{}第九届监事会第二十二次临时会议审议通过终止议案 \\
2025-09-12 & 重整投资人招募与《重整投资协议》签署 & 产业投资人新疆新业国有资产经营（集团）有限责任公司；财务投资人 7 家（含新疆兵金建投资有限公司；广州昊智科技投资合伙企业（有限合伙）及信风投资管理有限公司联合体；新疆兵投资产管理有限责任公司及青岛鲁创私募基金管理有限公司联合体；北京博雅春芽投资有限公司及长城（天津）股权投资基金管理有限责任公司联合体等）。具体投资金额 & 以重整投资人投资款解决债务、\allowbreak{}优化资产负债结构与股本结构；2025 年 8 月 5 日披露公开招募重整投资人（2025-063 号） \\
2025-12-29 & 债务豁免（子公司红色番茄） & 新疆生产建设兵团第六师国有资产经营有限责任公司豁免债务 4,\allowbreak{}248.00 万元（含借款本金 3,\allowbreak{}942.44 万元、\allowbreak{}利息 305.56 万元）；新疆国恒投资发展集团有限公司豁免债务 5,\allowbreak{}752.00 万元；合计增加子公司红色番茄资本公积 1.00 亿元 & 改善子公司资产负债结构；2025 年 12 月 29 日与六师国资公司签订债务豁免协议、\allowbreak{}2025 年 12 月 30 日与国恒集团公司签订债务豁免协议；该债务豁免被审计机构列为关键审计事项 \\
\bottomrule
\end{tabular}
\end{adjustbox}
\end{center}

\pfl{历史融投资情况} 从现金流量表看，2025 年公司取得借款收到的现金 2.19 亿元、偿还债务支付的现金 2.18 亿元，筹资活动现金流净额 -0.48 亿元（2023 年为取得借款 5.18 亿元、偿还债务 1.04 亿元）。 2026 年上半年筹资活动现金流净额为 -1.29 亿元（取得借款 — 亿元、偿还债务 1.31 亿元，未年化），已转为净流出，处于净还债状态。 2025 年分配股利、利润或偿付利息支付的现金合计 0.23 亿元。 公开披露的融资记录共 5 笔，工具类型以 IPO、发行股份购买资产、重整投资为主；金额与用途见下表。 其中规模最大的两笔为：2022 年 8 月，发行股份购买资产，拟以发行股份方式向新疆中泰农业发展有限责任公司、新疆粮油集团有限责任公司、新疆艳阳天番茄制品有限责任公司购买其合计持有的新疆新粮艳阳天番茄股份有限公司 100.00\% 股份；2000 年 9 月，IPO，公开发行人民币普通股 4,500.00 万股。 股本方面，近三年披露股本变动 3 次；总股本自 2021 年末以来未发生变化（7.71 亿股）。 分红方面，近三年未实施现金分红，与重整期间的资金约束一致。 近三年“再融资”类公告 36 条，涉及定向增发、发行股份购买资产。 综合看，公司融资结构上股权与债务融资并举；2025 年筹资活动现金流净额 -0.48 亿元，融资节奏仍以维持存量债务周转为主，与前述经营与投资现金流的方向基本一致。

\pfl{未来潜在融资需求分析} 2026 年 6 月末货币资金 3.01 亿元，而短期债务 5.01 亿元（现金短债比 0.60 倍），2026 年上半年经营现金流净额 1.27 亿元。按“短期债务减去账面货币资金”的滚动偿付口径，静态缺口约 2.00 亿元；若再计入上半年亏损的年化额（1.32 亿元），维持一年正常经营所需的外部资金约 2.00---3.32 亿元。 按第 \ref{sec:financing-need} 节的三档划分，公司属于第一档“补血型”：2026 年上半年归母净利润 -0.66 亿元、6 月末资产负债率 104.1\%。负债率已超过 80\%，净资产被亏损侵蚀的程度较深，属于全部样本中资产负债表最紧张的一组。 可行的融资路径按现实性排序为：引入国资或产业投资人（部分公司以控制权让渡为对价）、定向增发（需股价配合，且控股股东需放弃优先认购或同步增资）、出售非核心资产与参股股权。 截至目前公司仍处于预重整/重整程序中，重整投资人的招募与投资款到位是未来一年最主要的外部资金来源。 综合看，公司的外部资金需求以债务接续为主，滚动偿付口径下的静态缺口约 2.00 亿元，融资的时点与规模更多取决于债权人与重整投资人的进度。




\subsubsection{保龄宝（002286）}
\label{co:002286}

\pfl{公司基本情况} 山东禹城的功能糖（益生元、膳食纤维、减糖甜味剂）龙头，国内首家完成赤藓糖醇工业化生产的企业，为元气森林等无糖饮品品牌的核心代糖供应商；2009 年上市。 公司于 2009-08-28 在深交所主板首发上市，注册资本 3.80 亿元，总股本 3.70 亿股。 截至 2026 年 10 月 12 日收盘价 8.41 元，对应总市值 31.1 亿元，近一年-14.9\%，总市值在三级行业“其他农产品加工”的 11 家公司中按由大到小排第\zhnumber{9}位，近一年涨跌幅低于全样本中位数（-10.4\%）4.5 个百分点。 按 2026 年半年报披露的行业构成，收入几乎全部来自农副食品加工业 100.0\%。 与 2021 年年报相比，第一大行业口径由“食品制造业”（99.5\%）变为“农副食品加工业”（100.0\%）；两期披露的分类口径有调整，占比差异中同时包含分类因素。

\pfl{股权结构与控制权} 按 2026 年 6 月 30 日的股东名册，前三大股东为刘宗利 持股 6.11\%；北京永裕投资管理有限公司 持股 5.62\%；章木秀（持股比例未披露），股东户数 27,302 户。 股东结构与控制权方面，公司创始人刘宗利持股 2,326.83 万股、占 6.11\%；报告期内公司控股股东未发生其他变更。

\pfl{经营要点} 从经营要点看，2026 年 5 月，公司年产 2 万吨阿洛酮糖智能工厂扩产项目竣工投产，投产后阿洛酮糖总产能达 3 万吨、领跑全国。

\pfl{所处行业的基本面} 该子行业横跨制糖、番茄制品、糖醇与添加剂，供给受种植面积与进口配额约束，需求端分食用与工业两类（第 \ref{sec:framework} 章的供需框架）。 行业层面的最新情况是：2026 年上半年，饮料行业整体销量下滑但无糖饮料热度不减，赤藓糖醇价格自 2022 年下半年从 4 万元/吨跌至 9500 元/吨；2026 年上半年，华康股份收入 24.80 亿元、同比增长 33\% 但归母净利润降幅达 68.6\%。

\pfl{竞争优势} 把该公司放回同类公司比较，2026 年上半年营业收入 15.0 亿元、在三级行业“其他农产品加工”的 11 家公司中排第\zhnumber{6}位；毛利率 14.7\%、高于行业内中位数 5.6 个百分点；半年 ROE 3.6\%（未年化）、高于行业内中位数（1.36\%）；资产负债率 23.1\%、低于行业内中位数（39.1\%）。 公司区别于同类公司的经营特征包括：从事阿洛酮糖研发已十多年，同时具备晶体、液体阿洛酮糖产品生产能力。 优势在报表上主要体现为高于行业平均的毛利率，这在周期底部是抵御亏损的主要缓冲。

\pfl{历史股价} 自 2021 年 1 月至 2026 年 10 月，股价由 8.42 元变为 8.41 元，累计 -0.1\%，区间最高 14.66 元（2021 年 11 月）、最低 5.51 元（2024 年 6 月）；当前价相当于区间最高价的 57.4\%，为区间最低价的 1.53 倍。 月度收益的年化波动率 35.0\%，低于全样本中位数 37.3\%，在三级行业“其他农产品加工”的 11 家公司中波动率由高到低排第\zhnumber{5}位。

\begin{center}
\begin{minipage}{\textwidth}\centering
\begin{tikzpicture}
\begin{axis}[
  width=12.6cm, height=3.9cm, scale only axis,
  xmin=2020.922, xmax=2026.972, ymin=5.12, ymax=15.69,
  xtick={2021,2022,2023,2024,2025,2026},
  yearticks,
  yticklabel style={font=\scriptsize},
  ylabel={元/股}, ylabel style={font=\scriptsize},
  grid=major, grid style={LightGray, line width=0.3pt},
  axis line style={InkGray, line width=0.5pt},
  every axis plot/.append style={line width=0.9pt},
]
\addplot[AgriGreen] table[x=x,y=y]{data/clean/profiles/px_002286.dat};
\addplot[SoftRed, dashed, line width=0.7pt] table[x=x,y=y]{data/clean/profiles/pxzz_002286.dat};
\addplot[only marks, mark=*, mark size=1.1pt, SoftRed] table[x=x,y=y]{data/clean/profiles/pxzz_002286.dat};
\end{axis}
\end{tikzpicture}

\captionof{figure}[保龄宝（002286）股价与周期骨架]{保龄宝（002286）月度收盘价（元/股）与 ZigZag 周期骨架（阈值 25\%，圆点为周期转折点）}
\end{minipage}
\end{center}

\pfl{过去周期分析} 以 25\% 的 ZigZag 阈值划分，样本区间内股价形成 2 个主升段与 2 个主跌段。 最大主跌段为 2021 年 11 月 至 2024 年 6 月，31 个月-62.4\%；最大主升段出现在 2024 年 6 月 至 2025 年 8 月，14 个月+111.6\%。 最近一次转折出现在 2026 年 7 月（下跌段自 2025 年 8 月起，11 个月-32.0\%），当前处于该段的延续之中。 公司股价较申万农林牧渔指数提前约 23 个月见底，是板块中较早定价周期反转的公司之一。 该子行业横跨糖、番茄制品、添加剂等多个品种，各自的供给周期与政策（进口配额、关税）是主要变量，公司 2026 年上半年实现盈利、半年 ROE 3.6\%（未年化），好于全样本中位数（0.75\%），下行周期对其报表的冲击小于同业。 综合区间形态看，该股单段涨跌幅度偏大、以趋势段为主（最大主升 +111.6\%、最大主跌 -62.4\%），年化波动率 35.0\% 与全样本中位数 37.3\% 相比波动与样本中位数接近。

\pfl{盈利情况} 2026 年上半年营业收入 14.96 亿元（同比 +6.9\%），归母净利润 0.82 亿元，同比 -11.2\%。 以年报为趋势对照，2023---2025 年营业收入由 25.24 亿元变为 27.51 亿元（两年年均复合增速 +4.4\%），归母净利润由 0.54 亿元变为 1.51 亿元。 按半年报口径，2026 年上半年的 ROE 为 3.6\%（未年化）、销售净利率 5.5\%、毛利率 14.7\%，ROE 在“其他农产品加工”11 家公司中按数值由高到低排第\zhnumber{3}位。 按同一口径，三级行业“其他农产品加工”11 家公司的半年 ROE 中位数为 1.36\%，该公司高于其中位数 2.26 个百分点。 综合看，公司在报告期维持盈利（高于全样本半年 ROE 中位数 0.75\%），利润的可持续性取决于其他农产品加工景气的延续与成本管控的执行。

\begin{center}\small\setlength{\tabcolsep}{3.4pt}
\captionof{table}[保龄宝（002286）关键财务指标]{保龄宝（002286）关键财务指标（合并报表口径；两个半年报期均未年化；现金短债比 $=$ 货币资金 $\div$ 短期债务）}
\begin{adjustbox}{max width=\textwidth}
\begin{tabular}{@{}lrrrrr@{}}
\toprule
指标 & 2023 年 & 2024 年 & 2025 年 & 2025 年半年报 & 2026 年半年报 \\
\midrule
营业收入（亿元） & 25.24 & 24.02 & 27.51 & 13.99 & 14.96 \\
归母净利润（亿元） & 0.54 & 1.11 & 1.51 & 0.93 & 0.82 \\
毛利率（\%） & 8.1 & 11.9 & 13.3 & 13.2 & 14.7 \\
ROE（\%） & 2.7 & 5.5 & 7.1 & 4.4 & 3.6 \\
资产负债率（\%） & 25.4 & 23.7 & 19.7 & 23.6 & 23.1 \\
经营活动现金流净额（亿元） & 2.82 & 2.49 & 2.59 & 0.23 & 1.19 \\
货币资金（亿元） & 2.43 & 2.55 & 2.74 & 2.71 & 3.13 \\
有息负债合计（亿元） & 2.93 & 2.72 & 1.16 & 3.09 & 1.77 \\
现金短债比（倍） & 0.95 & 0.94 & 2.39 & 0.88 & 2.29 \\
\bottomrule
\end{tabular}
\end{adjustbox}
\end{center}

\begin{center}
\begin{minipage}{\textwidth}\centering
\begin{tikzpicture}
\begin{groupplot}[
  group style={group size=2 by 1, horizontal sep=0.60cm},
  width=6.85cm, height=4.60cm,
  tick label style={font=\tiny},
  title style={font=\scriptsize\heiti},
  yticklabel style={font=\tiny},
  ylabel style={font=\tiny},
  grid=major, grid style={LightGray, line width=0.3pt},
  axis line style={InkGray, line width=0.5pt},
  enlarge x limits={abs=0.8},
  legend style={font=\scriptsize, at={(0.5,-0.20)}, anchor=north,
                legend columns=2, draw=none, fill=none, column sep=6pt},
]
\nextgroupplot[
  ybar, bar width=4pt,
  ymin=-2.76, ymax=30.97,
  xtick={0,1,2,3,4,5.7},
  xticklabels={2021,2022,2023,2024,2025,2026H1},
  ylabel={亿元},
]
\addplot[fill=AgriGreen!75, draw=AgriGreen, bar shift=-2.6pt] table[x=i, y=rev]{data/clean/profiles/fin_002286.dat};
\addplot[fill=SoftRed!60, draw=SoftRed, bar shift=2.6pt] table[x=i, y=ni]{data/clean/profiles/fin_002286.dat};
\addplot[fill=AgriGreen!30, draw=AgriGreen, pattern=north east lines, bar shift=-2.6pt] table[x=x, y=rev]{data/clean/profiles/finh_002286.dat};
\addplot[fill=SoftRed!20, draw=SoftRed, pattern=north east lines, bar shift=2.6pt] table[x=x, y=ni]{data/clean/profiles/finh_002286.dat};
\legend{营业收入（年/半年）, 归母净利润（年/半年）}
\nextgroupplot[
  ymin=-3.2, ymax=44.2,
  xtick={0,1,2,3,4,5.7},
  xticklabels={2021,2022,2023,2024,2025,2026H1},
  ylabel={\%},
]
\addplot[AgriGold, mark=*, mark size=1.3pt] table[x=i, y=roe]{data/clean/profiles/fin_002286.dat};
\addplot[OilBlue, mark=square*, mark size=1.3pt] table[x=i, y=debt]{data/clean/profiles/fin_002286.dat};
\addplot[only marks, AgriGold, mark=*, mark size=2.0pt] table[x=x, y=roe]{data/clean/profiles/finh_002286.dat};
\addplot[only marks, OilBlue, mark=square*, mark size=2.0pt] table[x=x, y=debt]{data/clean/profiles/finh_002286.dat};
\legend{ROE, 资产负债率}
\end{groupplot}
\end{tikzpicture}
\begin{tikzpicture}
\begin{axis}[
  name=finbot,
  width=13.75cm, height=3.10cm, scale only axis,
  ymin=-0.59, ymax=6.56,
  xtick={0,1,2,3,4,5.7},
  xticklabels={2021,2022,2023,2024,2025,2026H1},
  tick label style={font=\tiny},
  yticklabel style={font=\tiny},
  ylabel={亿元}, ylabel style={font=\tiny},
  ybar, bar width=4pt,
  grid=major, grid style={LightGray, line width=0.3pt},
  axis line style={InkGray, line width=0.5pt},
  enlarge x limits={abs=0.8},
  legend style={font=\scriptsize, at={(0.5,-0.26)}, anchor=north,
                legend columns=2, draw=none, fill=none, column sep=10pt},
]
\addplot[fill=AgriGreen!55, draw=AgriGreen, bar shift=-3.0pt] table[x=i, y=cash]{data/clean/profiles/fin_002286.dat};
\addplot[fill=SoftRed!45, draw=SoftRed, bar shift=3.0pt] table[x=i, y=ibd]{data/clean/profiles/fin_002286.dat};
\addplot[fill=AgriGreen!25, draw=AgriGreen, pattern=north east lines, bar shift=-3.0pt] table[x=x, y=cash]{data/clean/profiles/finh_002286.dat};
\addplot[fill=SoftRed!20, draw=SoftRed, pattern=north east lines, bar shift=3.0pt] table[x=x, y=ibd]{data/clean/profiles/finh_002286.dat};
\legend{货币资金（左轴）, 有息负债（左轴）}
\end{axis}
\begin{axis}[
  at={(finbot.south east)}, anchor=south east,
  width=13.75cm, height=3.10cm, scale only axis,
  axis y line*=right, axis x line*=none, xtick=\empty,
  ymin=-0.73, ymax=3.67,
  tick label style={font=\tiny},
  yticklabel style={font=\tiny}, scaled y ticks=false,
  legend style={font=\scriptsize, at={(0.5,-0.44)}, anchor=north,
                legend columns=1, draw=none, fill=none},
]
\addplot[OilBlue, mark=*, mark size=2.2pt, line width=1.3pt] table[x=i, y=ocf]{data/clean/profiles/fin_002286.dat};
\addplot[only marks, OilBlue, mark=*, mark size=3.2pt] table[x=x, y=ocf]{data/clean/profiles/finh_002286.dat};
\legend{经营活动现金流净额（右轴，亿元，蓝点）}
\end{axis}
\end{tikzpicture}

\captionof{figure}[保龄宝（002286）近年主要财务指标]{保龄宝（002286）近年主要财务指标（左上：营业收入与归母净利润；右上：ROE 与资产负债率；下：货币资金、有息负债为左轴柱状，经营活动现金流净额为其右轴的蓝点折线。
2021---2025 年为年报口径，2026H1 为 2026 年半年报/半年末，与其前一年报之间留出间隔、以斜纹或空心点区分；半年报数据未年化）}
\end{minipage}
\end{center}

\pfl{财务分析} 2026 年 6 月末资产负债率 23.1\%，较 2025 年末上升 3.4 个百分点（样本中位数 52.2\%，所处三级行业内按由低到高排第\zhnumber{3}位）。 2026 年 6 月末货币资金 3.13 亿元、短期债务（短期借款与一年内到期非流动负债）1.37 亿元、长期债务 0.41 亿元，有息负债合计 1.77 亿元，现金短债比 2.29 倍——覆盖充分。 净有息负债 -1.35 亿元，相当于其 2026 年上半年营业收入的 -9\%。 流动比率 1.74、速动比率 1.20，短期流动性无明显缺口。 2026 年上半年经营活动现金流净额 1.19 亿元，经营现金流与利润同向为正，内生资金可以支撑运营。 2026 年 6 月末在建工程 2.30 亿元，较上年末增加 0.99 亿元，仍有在建投入。 综合看，现金短债比 2.29 倍，短期偿付压力可控；资产负债率低于样本中位数 29.1 个百分点；有息负债中短期债务占比更高，期限结构仍以滚动续作为主。 公司披露的股东与质押风险为：2026 年 6 月，控股股东北京永裕投资管理有限公司所持公司部分股份被法院司法冻结；2026 年 9 月，公司总质押股份数量 1,879.39 万股。

\pfl{重大战略转型} 从公开披露看，公司的战略推进集中在：其他动作（2 项）：2025 年，限制性股票激励计划首次授予限制性股票第一个解除限售期解除限售条件成就；2025 年度，公司实现归属于上市公司股东的净利润 1.51 亿元、同比增长 35.50\%。 公告口径上，近三年（2023 年 9 月---2026 年 9 月）公司披露重大重组与并购类公告 0 条、对外投资与转型类公告 3 条、回购与股权激励类公告 34 条，合计公告 357 条。 主营结构的口径变化已经落到报表上：2021 年年报的第一大行业为“食品制造业”（99.5\%），2026 年半年报为“农副食品加工业”（100.0\%）；公司在这两期的分部划分方式有过调整。 综合判断，报告期内公司有 3 条与战略相关的公告落地（重大重组与并购 0 条、对外投资与转型 3 条），资本运作与主业调整之间存在直接关联。

\begin{center}\small\setlength{\tabcolsep}{4pt}
\captionof{table}[保龄宝（002286）历史融资与资本运作]{保龄宝（002286）历史融资与资本运作（据公司公告与公开报道整理；金额为披露值，含首次公开发行、再融资、债券、银行借款与重整投资等）}
\begin{adjustbox}{max width=\textwidth}
\begin{tabular}{@{}p{1.45cm}p{2.55cm}p{4.05cm}p{4.95cm}@{}}
\toprule
时间 & 方式 & 规模 & 用途与说明 \\
\midrule
2009-08-28 & IPO（深交所中小板） & 发行 2,\allowbreak{}000 万股，\allowbreak{}发行价 20.56 元/\allowbreak{}股 & 低聚果糖和水溶性膳食纤维产业化项目；补充流动资金 3,\allowbreak{}400.00 万元；归还银行借款 7,\allowbreak{}950.00 万元；发行市盈率 39.54 倍（按 2008 年经审计扣非前后净利润孰低额除以发行后总股本全面摊薄计算）；发行后股本总额 8,\allowbreak{}000 万股 \\
2013-04 & 非公开发行（定向增发） & 发行 4,\allowbreak{}942.80 万股 A 股，\allowbreak{}发行价 12.60 元/\allowbreak{}股 & 不涉及资产过户）；；新增股份 2013 年 4 月 17 日上市 \\
2021-07-13 & 非公开发行A股股票（预案） & 拟发行股票数量不超过 9,\allowbreak{}000 万股，\allowbreak{}发行价格 7.91 元/\allowbreak{}股 & 年产 3 万吨赤藓糖醇晶体项目 2.53 亿元、\allowbreak{}年产 3 万吨阿洛酮糖（干基）项目 1.94 亿元、\allowbreak{}应用解决方案研究中心项目、\allowbreak{}补充流动资金不超过 2.00 亿元；发行对象为控股股东北京永裕投资管理有限公司、\allowbreak{}赣州中科永裕投资管理中心（有限合伙）、\allowbreak{}深圳中科永裕科技合伙企业（有限合伙） \\
2022-11-24 & 取得非公开发行核准批复 & 核 & 年产 3 万吨赤藓糖醇晶体项目等；公司于 2022 年 6 月 14 日收到证监会行政许可申请受理单（受理序号 221232） \\
2023-11 & 非公开发行终止（批复到期未实施） & 原拟募资不超过 7.07 亿元 & 2023 年年度报告披露“非公开发行项目终止，\allowbreak{}本承诺至 2023 年 11 月 17 日结束”；公司总股本未因该次发行增加（2023 年末总股本 3.71 亿股 \\
\bottomrule
\end{tabular}
\end{adjustbox}
\end{center}

\pfl{历史融投资情况} 从现金流量表看，2025 年公司取得借款收到的现金 2.75 亿元、偿还债务支付的现金 4.32 亿元、吸收权益性投资 0.41 亿元，筹资活动现金流净额 -1.50 亿元（2023 年为取得借款 2.90 亿元、偿还债务 4.97 亿元）。 2026 年上半年筹资活动现金流净额为 0.57 亿元（取得借款 1.90 亿元、偿还债务 1.27 亿元，未年化），仍处于净融入状态。 2025 年分配股利、利润或偿付利息支付的现金合计 0.32 亿元。 公开披露的融资记录共 5 笔，工具类型以 IPO、定向增发为主；金额与用途见下表。 其中规模最大的两笔为：2009 年 8 月，IPO，发行 2,000 万股，发行价 20.56 元/股；2021 年 7 月，定向增发，拟发行股票数量不超过 9,000 万股，发行价格 7.91 元/股。 股本方面，近三年披露股本变动 5 次（股份回购 2 次）；总股本由 2021 年末的 3.72 亿股变为 3.70 亿股（-0.58\%）。 分红方面，近三年现金分红 6 次、合计每股派现 0.40 元，最近一次实施在 2026 年 6 月。 近三年“再融资”类公告 9 条，涉及定向增发。 综合看，公司融资结构上以债务滚动为主、股权融资为补充；2025 年筹资活动现金流净额 -1.50 亿元，融资节奏仍以维持存量债务周转为主，与前述经营与投资现金流的方向基本一致。

\pfl{未来潜在融资需求分析} 2026 年 6 月末货币资金 3.13 亿元，而短期债务 1.37 亿元（现金短债比 2.29 倍），2026 年上半年经营现金流净额 1.19 亿元。账面货币资金可以覆盖短期债务，缺口不体现在静态偿付层面。 公司 2026 年 6 月末资产负债率 23.1\%、2026 年上半年 ROE 3.6\%（未年化），资产负债表尚有余量，融资需求不是“补血型”的刚性需求，而是配套在建项目投入的主动性安排。 公开披露的后续资金安排线索包括：股权融资线索：2026 年 4 月，公司拟以简易程序向特定对象发行融资总额不超过人民币 3.00 亿元且不超过最近一年末净资产 20\% 的股票；债务与授信安排：2025 年 4 月，公司及控股子公司拟继续向银行申请不超过人民币 20.00 亿元的银行融资额度。 资金需求还要叠加在建工程：期末在建工程 2.30 亿元、2026 年上半年购建固定资产等支出 1.32 亿元（未年化），这部分支出在周期底部通常会被主动放缓。 综合看，公司在静态偿付层面不存在缺口，外部融资更多是配合扩张节奏与并购整合的主动性安排，而非偿债压力驱动。




\subsubsection{道道全（002852）}
\label{co:002852}

\pfl{公司基本情况} 湖南长沙起家的菜籽油及食用油加工企业（“菜籽油龙头”），2017 年 3 月深交所上市，自 2019 年起以期货套期保值对冲原料价格波动；2020 年因误判原料单边上行行情导致套保平仓亏损约 2.07 亿元、上市以来首亏，此后业绩逐步修复，2025 年营业收入 61.95 亿元、归母净利润 2.34 亿元；2026 年上半年营业收入 28.65 亿元、归母净利润 1,686.06 万元、同比下降 90.68\%，控股股东、实际控制人刘建军持股 21.69\%、质押 1,670 万股。 公司于 2017-03-10 在深交所主板首发上市，注册资本 3.44 亿元，总股本 3.44 亿股。 截至 2026 年 10 月 12 日收盘价 8.49 元，对应总市值 29.2 亿元，近一年-26.9\%，总市值在三级行业“粮油加工”的 7 家公司中按由大到小排第\zhnumber{6}位，近一年涨跌幅低于全样本中位数（-10.4\%）16.5 个百分点。 按 2026 年半年报披露的行业构成，收入占比前两位为食品加工 99.6\%；其他行业 0.4\%。

\pfl{股权结构与控制权} 按 2026 年 6 月 30 日的股东名册，前三大股东为刘建军 持股 21.69\%；湖南兴创投资管理有限公司 持股 13.51\%；中泰证券资管－中泰资管 7721 号 FOF 单一资产管理计 持股 5.03\%，股东户数 20,576 户。 控制权与股东性质方面，2026 年 6 月末，刘建军持股 7,459.56 万股、占总股本 21.69\%。

\pfl{经营要点} 经营层面，截至 2026 年 6 月末公司区域经销商 1,441 家、较 2025 年同期净增 44 家。

\pfl{所处行业的基本面} 粮油加工的供给由大豆、菜籽等原料进口与压榨产能决定，需求由食用消费决定，利润来自原料与成品的价差而非价格方向（第 \ref{sec:framework} 章的供需框架）。 从供需两端的当期数据看，2026 年，1—8 月中国累计进口大豆 7,410.7 万吨、同比增长 1.1\%，三季度油厂开机率维持 61\%—68\%。

\pfl{竞争优势} 把该公司放回同类公司比较，2026 年上半年营业收入 28.6 亿元、在三级行业“粮油加工”的 7 家公司中排第\zhnumber{3}位；毛利率 5.4\%、低于行业内中位数 0.9 个百分点；半年 ROE 0.8\%（未年化）、低于行业内中位数（2.22\%）；资产负债率 49.1\%、高于行业内中位数（43.5\%）。 公司在公开资料中自述与同业比较的相对优势在于：2026 年上半年西南地区收入同比增长三成多至 5.65 亿元，经销商数量与平均单家经销商营收同步提升；扩产主要基于食用油行业集中度仍在提升。 就当期报表看，公司尚未把经营层面的特征转化为高于行业的资本回报，周期底部的相对优势更多体现在现金流与资产负债表上。

\pfl{历史股价} 以 2021 年 1 月为起点、2026 年 10 月为终点，股价由 12.59 元变为 8.49 元，累计 -32.6\%，区间最高 16.75 元（2021 年 12 月）、最低 6.45 元（2024 年 8 月）；当前价相当于区间最高价的 50.7\%，为区间最低价的 1.32 倍。 月度收益的年化波动率 39.3\%，高于全样本中位数 37.3\%，在三级行业“粮油加工”的 7 家公司中波动率由高到低排第\zhnumber{3}位。 把股价阶段与公司事件对齐看，2023 年 4 月 至 2024 年 8 月股价下跌-56.6\%，同期（2024 年 5 月）控股股东一致行动人兴创投资协议转让 1,730.00 万股给中泰证券（上海）资产管理有限公司代表的中泰资管 5600 号单一资产管理计划。 整体上看，区间的几处大级别拐点都能在公司的资本运作、股权变动或风险事件上找到对应，股价波动有明显的公司层面驱动，而非单纯的行业 beta。

\begin{center}
\begin{minipage}{\textwidth}\centering
\begin{tikzpicture}
\begin{axis}[
  width=12.6cm, height=3.9cm, scale only axis,
  xmin=2020.922, xmax=2026.972, ymin=6.00, ymax=17.92,
  xtick={2021,2022,2023,2024,2025,2026},
  yearticks,
  yticklabel style={font=\scriptsize},
  ylabel={元/股}, ylabel style={font=\scriptsize},
  grid=major, grid style={LightGray, line width=0.3pt},
  axis line style={InkGray, line width=0.5pt},
  every axis plot/.append style={line width=0.9pt},
]
\addplot[AgriGreen] table[x=x,y=y]{data/clean/profiles/px_002852.dat};
\addplot[SoftRed, dashed, line width=0.7pt] table[x=x,y=y]{data/clean/profiles/pxzz_002852.dat};
\addplot[only marks, mark=*, mark size=1.1pt, SoftRed] table[x=x,y=y]{data/clean/profiles/pxzz_002852.dat};
\end{axis}
\end{tikzpicture}

\captionof{figure}[道道全（002852）股价与周期骨架]{道道全（002852）月度收盘价（元/股）与 ZigZag 周期骨架（阈值 25\%，圆点为周期转折点）}
\end{minipage}
\end{center}

\pfl{过去周期分析} 以 25\% 的 ZigZag 阈值划分，样本区间内股价形成 4 个主升段与 4 个主跌段。 最大主升段出现在 2024 年 8 月 至 2025 年 7 月，11 个月+88.1\%；最大主跌段为 2023 年 4 月 至 2024 年 8 月，16 个月-56.6\%。 最近一次转折出现在 2026 年 6 月（下跌段自 2025 年 7 月起，11 个月-41.5\%），当前处于该段的延续之中。 公司股价较申万农林牧渔指数提前约 21 个月见底，是板块中较早定价周期反转的公司之一。 粮油加工是“成本加成”型生意（第 \ref{sec:grain} 章），利润来自原料与成品的价差及规模效率，收入规模大而净利率薄，公司 2026 年 6 月末资产负债率 49.1\%、半年 ROE 0.8\%（未年化），在本轮下行中属于随行业同步调整的一类。 综合区间形态看，该股单段涨跌幅度偏大、以趋势段为主（最大主升 +88.1\%、最大主跌 -56.6\%），年化波动率 39.3\% 与全样本中位数 37.3\% 相比波动与样本中位数接近。

\pfl{盈利情况} 2026 年上半年营业收入 28.65 亿元（同比 +2.6\%），归母净利润 0.17 亿元，同比 -90.7\%。 以年报为趋势对照，2023---2025 年营业收入由 70.01 亿元变为 61.95 亿元（两年年均复合增速 -5.9\%），归母净利润由 0.76 亿元变为 2.34 亿元。 按半年报口径，2026 年上半年的 ROE 为 0.8\%（未年化）、销售净利率 0.7\%、毛利率 5.4\%，ROE 在“粮油加工”7 家公司中按数值由高到低排第\zhnumber{5}位。 按同一口径，三级行业“粮油加工”7 家公司的半年 ROE 中位数为 2.22\%，该公司低于其中位数 1.47 个百分点。 从公司披露的原因看，业绩变动主要来自：2025 年度，报告期内套期工具与被套期项目合计实现损益约 3.04 亿元。 综合看，公司在报告期维持盈利（低于全样本半年 ROE 中位数 0.75\%），利润的可持续性取决于粮油加工景气的延续与成本管控的执行。

\begin{center}\small\setlength{\tabcolsep}{3.4pt}
\captionof{table}[道道全（002852）关键财务指标]{道道全（002852）关键财务指标（合并报表口径；两个半年报期均未年化；现金短债比 $=$ 货币资金 $\div$ 短期债务）}
\begin{adjustbox}{max width=\textwidth}
\begin{tabular}{@{}lrrrrr@{}}
\toprule
指标 & 2023 年 & 2024 年 & 2025 年 & 2025 年半年报 & 2026 年半年报 \\
\midrule
营业收入（亿元） & 70.01 & 59.43 & 61.95 & 27.92 & 28.65 \\
归母净利润（亿元） & 0.76 & 1.77 & 2.34 & 1.81 & 0.17 \\
毛利率（\%） & 7.6 & 9.6 & 8.3 & 11.6 & 5.4 \\
ROE（\%） & 3.8 & 8.6 & 10.7 & 8.2 & 0.8 \\
资产负债率（\%） & 57.2 & 46.1 & 44.8 & 48.1 & 49.1 \\
经营活动现金流净额（亿元） & 30.04 & 31.63 & 9.93 & 3.25 & -2.12 \\
货币资金（亿元） & 2.57 & 2.89 & 5.56 & 1.69 & 4.92 \\
有息负债合计（亿元） & 16.11 & 10.08 & 9.96 & 11.43 & 12.69 \\
现金短债比（倍） & 0.16 & 0.29 & 0.56 & 0.15 & 0.39 \\
\bottomrule
\end{tabular}
\end{adjustbox}
\end{center}

\begin{center}
\begin{minipage}{\textwidth}\centering
\begin{tikzpicture}
\begin{groupplot}[
  group style={group size=2 by 1, horizontal sep=0.60cm},
  width=6.85cm, height=4.60cm,
  tick label style={font=\tiny},
  title style={font=\scriptsize\heiti},
  yticklabel style={font=\tiny},
  ylabel style={font=\tiny},
  grid=major, grid style={LightGray, line width=0.3pt},
  axis line style={InkGray, line width=0.5pt},
  enlarge x limits={abs=0.8},
  legend style={font=\scriptsize, at={(0.5,-0.20)}, anchor=north,
                legend columns=2, draw=none, fill=none, column sep=6pt},
]
\nextgroupplot[
  ybar, bar width=4pt,
  ymin=-11.63, ymax=79.22,
  xtick={0,1,2,3,4,5.7},
  xticklabels={2021,2022,2023,2024,2025,2026H1},
  ylabel={亿元},
]
\addplot[fill=AgriGreen!75, draw=AgriGreen, bar shift=-2.6pt] table[x=i, y=rev]{data/clean/profiles/fin_002852.dat};
\addplot[fill=SoftRed!60, draw=SoftRed, bar shift=2.6pt] table[x=i, y=ni]{data/clean/profiles/fin_002852.dat};
\addplot[fill=AgriGreen!30, draw=AgriGreen, pattern=north east lines, bar shift=-2.6pt] table[x=x, y=rev]{data/clean/profiles/finh_002852.dat};
\addplot[fill=SoftRed!20, draw=SoftRed, pattern=north east lines, bar shift=2.6pt] table[x=x, y=ni]{data/clean/profiles/finh_002852.dat};
\legend{营业收入（年/半年）, 归母净利润（年/半年）}
\nextgroupplot[
  ymin=-25.3, ymax=65.9,
  xtick={0,1,2,3,4,5.7},
  xticklabels={2021,2022,2023,2024,2025,2026H1},
  ylabel={\%},
]
\addplot[AgriGold, mark=*, mark size=1.3pt] table[x=i, y=roe]{data/clean/profiles/fin_002852.dat};
\addplot[OilBlue, mark=square*, mark size=1.3pt] table[x=i, y=debt]{data/clean/profiles/fin_002852.dat};
\addplot[only marks, AgriGold, mark=*, mark size=2.0pt] table[x=x, y=roe]{data/clean/profiles/finh_002852.dat};
\addplot[only marks, OilBlue, mark=square*, mark size=2.0pt] table[x=x, y=debt]{data/clean/profiles/finh_002852.dat};
\legend{ROE, 资产负债率}
\end{groupplot}
\end{tikzpicture}
\begin{tikzpicture}
\begin{axis}[
  name=finbot,
  width=13.75cm, height=3.10cm, scale only axis,
  ymin=-1.61, ymax=18.04,
  xtick={0,1,2,3,4,5.7},
  xticklabels={2021,2022,2023,2024,2025,2026H1},
  tick label style={font=\tiny},
  yticklabel style={font=\tiny},
  ylabel={亿元}, ylabel style={font=\tiny},
  ybar, bar width=4pt,
  grid=major, grid style={LightGray, line width=0.3pt},
  axis line style={InkGray, line width=0.5pt},
  enlarge x limits={abs=0.8},
  legend style={font=\scriptsize, at={(0.5,-0.26)}, anchor=north,
                legend columns=2, draw=none, fill=none, column sep=10pt},
]
\addplot[fill=AgriGreen!55, draw=AgriGreen, bar shift=-3.0pt] table[x=i, y=cash]{data/clean/profiles/fin_002852.dat};
\addplot[fill=SoftRed!45, draw=SoftRed, bar shift=3.0pt] table[x=i, y=ibd]{data/clean/profiles/fin_002852.dat};
\addplot[fill=AgriGreen!25, draw=AgriGreen, pattern=north east lines, bar shift=-3.0pt] table[x=x, y=cash]{data/clean/profiles/finh_002852.dat};
\addplot[fill=SoftRed!20, draw=SoftRed, pattern=north east lines, bar shift=3.0pt] table[x=x, y=ibd]{data/clean/profiles/finh_002852.dat};
\legend{货币资金（左轴）, 有息负债（左轴）}
\end{axis}
\begin{axis}[
  at={(finbot.south east)}, anchor=south east,
  width=13.75cm, height=3.10cm, scale only axis,
  axis y line*=right, axis x line*=none, xtick=\empty,
  ymin=-10.90, ymax=41.75,
  tick label style={font=\tiny},
  yticklabel style={font=\tiny}, scaled y ticks=false,
  legend style={font=\scriptsize, at={(0.5,-0.44)}, anchor=north,
                legend columns=1, draw=none, fill=none},
]
\addplot[OilBlue, mark=*, mark size=2.2pt, line width=1.3pt] table[x=i, y=ocf]{data/clean/profiles/fin_002852.dat};
\addplot[only marks, OilBlue, mark=*, mark size=3.2pt] table[x=x, y=ocf]{data/clean/profiles/finh_002852.dat};
\legend{经营活动现金流净额（右轴，亿元，蓝点）}
\end{axis}
\end{tikzpicture}

\captionof{figure}[道道全（002852）近年主要财务指标]{道道全（002852）近年主要财务指标（左上：营业收入与归母净利润；右上：ROE 与资产负债率；下：货币资金、有息负债为左轴柱状，经营活动现金流净额为其右轴的蓝点折线。
2021---2025 年为年报口径，2026H1 为 2026 年半年报/半年末，与其前一年报之间留出间隔、以斜纹或空心点区分；半年报数据未年化）}
\end{minipage}
\end{center}

\pfl{财务分析} 2026 年 6 月末资产负债率 49.1\%，较 2025 年末上升 4.3 个百分点（样本中位数 52.2\%，所处三级行业内按由低到高排第\zhnumber{5}位）。 2026 年 6 月末货币资金 4.92 亿元、短期债务（短期借款与一年内到期非流动负债）12.69 亿元、长期债务 0.00 亿元，有息负债合计 12.69 亿元，现金短债比 0.39 倍——覆盖偏紧。 净有息负债 7.77 亿元，相当于其 2026 年上半年营业收入的 27\%。 流动比率 1.07、速动比率 0.40，短期指标尚可。 2026 年上半年经营活动现金流净额 -2.12 亿元，净利润为正但经营现金流为负，利润的现金含量偏低。 2026 年 6 月末在建工程 2.35 亿元，较上年末增加 0.63 亿元，仍有在建投入。 综合看，账面货币资金对有息负债与短期债务的覆盖不足，滚动偿付对新增融资的依赖度较高；资产负债率低于样本中位数 3.1 个百分点；有息负债中短期债务占比更高，期限结构仍以滚动续作为主。

\pfl{重大战略转型} 公司的战略动作可归为以下几类：其他动作（3 项）：2026 年 8 月，进一步丰富原料价格风险管理工具；2026 年上半年，公司主营原料价格持续上涨、终端产品售价未能充分覆盖采购成本涨幅；并购与重组（1 项）：2024 年 5 月，控股股东一致行动人兴创投资协议转让 1,730.00 万股给中泰证券（上海）资产管理有限公司代表的中泰资管 5600 号单一资产管理计划。 公告口径上，近三年（2023 年 9 月---2026 年 9 月）公司披露重大重组与并购类公告 0 条、对外投资与转型类公告 0 条、回购与股权激励类公告 0 条，合计公告 301 条。 第一大行业仍是“食品加工”，收入占比 99.6\%，与 2021 年年报相比没有方向性变化。 综合判断，报告期内公司未披露实质性的重组、并购或转型安排，资本开支以维持现有产能与业务为主。

\begin{center}\small\setlength{\tabcolsep}{4pt}
\captionof{table}[道道全（002852）历史融资与资本运作]{道道全（002852）历史融资与资本运作（据公司公告与公开报道整理；金额为披露值，含首次公开发行、再融资、债券、银行借款与重整投资等）}
\begin{adjustbox}{max width=\textwidth}
\begin{tabular}{@{}p{1.45cm}p{2.55cm}p{4.05cm}p{4.95cm}@{}}
\toprule
时间 & 方式 & 规模 & 用途与说明 \\
\midrule
2017-03-10 & IPO（深圳证券交易所） & 实际发行 2,\allowbreak{}500.00 万股 & 券股份有限公司，\allowbreak{}发行市盈率（按发行后总股本）22.99 倍 \\
2024-05-09 & 股份协议转让（大股东一致行动人减持） & 兴创投资拟通过协议转让方式将其持有的公司 1,\allowbreak{}730.00 万股无限售流通股转让给中泰证券（上海）资产管理有限公司代表的“中泰资管 5600 号单一资产管理计划” & 大股东一致行动人调整持股结构；转让后中泰证券资管相关产品进入公司前 10 名股东（2025 年末“中泰证券资管－中 7721 号 FOF 单一资产管理计划”持股 5.03\%） \\
2025-08-06 & 大股东增持（含股份增持专项贷款） & 兴创投资拟增持金额不低于 5,\allowbreak{}000.00 万元、\allowbreak{}不超过 1.00 亿元 & 基于对公司未来发展前景的信心和中长期投资价值的认可，\allowbreak{}其中专项贷款资金占比不超过 90\%；2025 年 11 月 5 日实施期限过半时已累计增持 353.62 万股（占总股本 1.03\%） \\
2026 & 商品期货和衍生品交易业务（含保证金额度） & 本次商品期货套期保值业务可循环使用的保证金最高额度不超过公司最近一期经审计净资产的 20\%（不包括交割头寸而支付的全额保证金在内） & 对冲食用植物油原料价格波动风险；公司 2026 年度继续开展商品期货和衍生品交易业务；2026 年 8 月 6 日另公告增加商品期权作为套期保值业务品种 \\
\bottomrule
\end{tabular}
\end{adjustbox}
\end{center}

\pfl{历史融投资情况} 从现金流量表看，2025 年公司取得借款收到的现金 29.40 亿元、偿还债务支付的现金 34.11 亿元，筹资活动现金流净额 -6.18 亿元（2023 年为取得借款 29.54 亿元、偿还债务 57.57 亿元）。 2026 年上半年筹资活动现金流净额为 2.19 亿元（取得借款 11.21 亿元、偿还债务 8.48 亿元，未年化），仍处于净融入状态。 2025 年分配股利、利润或偿付利息支付的现金合计 1.46 亿元。 公开披露的融资记录共 4 笔，工具类型以 IPO 为主；金额与用途见下表。 其中规模最大的两笔为：2025 年 8 月，大股东增持，兴创投资拟增持金额不低于 5,000.00 万元、不超过 1.00 亿元；2017 年 3 月，IPO，实际发行 2,500.00 万股。 股本方面，近三年披露股本变动 4 次（股份回购 1 次）；总股本由 2021 年末的 3.59 亿股变为 3.44 亿股（-4.19\%）。 分红方面，近三年现金分红 7 次、合计每股派现 0.85 元，最近一次实施在 2026 年 6 月。 近三年“再融资”类公告 3 条，涉及定向增发、可转债。 综合看，公司融资结构上股权与债务融资并举；2025 年筹资活动现金流净额 -6.18 亿元，融资节奏仍以维持存量债务周转为主，与前述经营与投资现金流的方向基本一致。

\pfl{未来潜在融资需求分析} 2026 年 6 月末货币资金 4.92 亿元，而短期债务 12.69 亿元（现金短债比 0.39 倍），2026 年上半年经营现金流净额 -2.12 亿元。按“短期债务减去账面货币资金”的滚动偿付口径，静态缺口约 7.77 亿元；上半年盈利或接近盈亏平衡，该缺口不随经营亏损进一步扩大。 公司 2026 年 6 月末资产负债率 49.1\%、2026 年上半年 ROE 0.8\%（未年化），现金短债比 0.39 倍、短债覆盖不足，融资需求首先用于置换与滚动存量债务、补充营运资金，属于“补血型”而非扩张型。 公开披露的后续资金安排线索包括：债务与授信安排：2026 年 2 月，资金来源含股份增持专项贷款（占比不超过 90\%）。 资金需求还要叠加在建工程：期末在建工程 2.35 亿元、2026 年上半年购建固定资产等支出 0.92 亿元（未年化），这部分支出在周期底部通常会被主动放缓。 综合看，公司的外部资金需求以营运资金周转与在建项目投入为主，滚动偿付口径下的静态缺口约 7.77 亿元，融资的时点与规模更多取决于信用条件与资本市场窗口。




\subsubsection{*ST广糖（000911）}
\label{co:000911}

\pfl{公司基本情况} 广西国资控股的制糖企业，所属糖企始建于 1956 年，1999 年 5 月深交所上市，2019 年广西农村投资集团成为控股股东、2023 年 10 月更名为广农糖业。2012—2025 年扣非归母净利润连续 14 年为负，2025 年末归母净资产转为负值，2026 年 4 月 30 日起被实施退市风险警示，并因年报虚假记载被证监会立案处罚而叠加其他风险警示；2026 年 8 月拟向控股股东体系出售物流与种业资产以修复净资产。 公司于 1999-05-27 在深交所主板首发上市，注册资本 4.00 亿元，总股本 4.00 亿股。 截至 2026 年 10 月 12 日收盘价 6.22 元，对应总市值 24.9 亿元，近一年-15.1\%，总市值在三级行业“其他农产品加工”的 11 家公司中按由大到小排第\zhnumber{10}位，近一年涨跌幅低于全样本中位数（-10.4\%）4.8 个百分点。 按 2026 年半年报披露的行业构成，收入占比前两位为工业 88.6\%；物流贸易及其他 11.4\%。

\pfl{股权结构与控制权} 按 2026 年 6 月 30 日的股东名册，前三大股东为广西农村投资集团有限公司 持股 38.23\%；南宁产业投资集团有限责任公司 持股 14.98\%；深圳金湾私募证券基金管理有限公司－金湾长安一号私募证 持股 0.55\%，股东户数 34,230 户。 股东结构与控制权方面，公司总股本 4.00 亿股（全部为无限售条件股份）；公司股票自 2026 年 4 月 30 日起被实施退市风险警示、简称由广农糖业变更为*ST 广糖。

\pfl{经营要点} 业务与产能层面，公司拥有一批制糖核心技术与专利，占全国甘蔗种植面积的 66.32\%；2025 年上半年，同比增加 123.24 万吨（增幅约 19\%）。

\pfl{所处行业的基本面} 该子行业横跨制糖、番茄制品、糖醇与添加剂，供给受种植面积与进口配额约束，需求端分食用与工业两类（第 \ref{sec:framework} 章的供需框架）。 与该公司直接相关的行业数据包括：2025 年，平均产糖率 13.14\%、同比下降 0.25 个百分点；部分主产区现货价格长时间运行于行业平均成本线下方。

\pfl{竞争优势} 把该公司放回同类公司比较，2026 年上半年营业收入 15.8 亿元、在三级行业“其他农产品加工”的 11 家公司中排第\zhnumber{5}位；毛利率 9.1\%、与行业内中位数基本持平；半年 ROE 0.0\%（未年化）、低于行业内中位数（1.36\%）；资产负债率 100.2\%、高于行业内中位数（39.1\%）。 相对优势方面，官网产业布局页列出甘蔗制糖板块日榨能力 500 至 45,000 吨/日、年产糖能力 0.189 至 65 万吨、甘蔗种植面积 3.89 至 80 万亩；2026 年 8 月，公司拟将持有的南宁云鸥物流 99.90\% 股权及广西南糖种业 100\% 股权转让至广西广农供应链集团。 就当期报表看，公司尚未把经营层面的特征转化为高于行业的资本回报，周期底部的相对优势更多体现在现金流与资产负债表上。

\pfl{历史股价} 本报告样本区间（2021 年 1 月 至 2026 年 10 月）内，股价由 7.70 元变为 6.22 元，累计 -19.2\%，区间最高 13.11 元（2021 年 8 月）、最低 3.83 元（2026 年 6 月）；当前价相当于区间最高价的 47.4\%，为区间最低价的 1.62 倍。 月度收益的年化波动率 40.8\%，高于全样本中位数 37.3\%，在三级行业“其他农产品加工”的 11 家公司中波动率由高到低排第\zhnumber{4}位。 把股价阶段与公司事件对齐看，2026 年 6 月 至 2026 年 10 月股价上涨+62.4\%，同期（2026 年 9 月）下一会计年度经审计期末净资产仍为负将直接进入终止上市程序。 从时间对应关系看，公司层面的资本运作与股权、风险事件解释了区间内多数急涨急跌，与行业指数的同步性相对有限。

\begin{center}
\begin{minipage}{\textwidth}\centering
\begin{tikzpicture}
\begin{axis}[
  width=12.6cm, height=3.9cm, scale only axis,
  xmin=2020.922, xmax=2026.972, ymin=3.56, ymax=14.03,
  xtick={2021,2022,2023,2024,2025,2026},
  yearticks,
  yticklabel style={font=\scriptsize},
  ylabel={元/股}, ylabel style={font=\scriptsize},
  grid=major, grid style={LightGray, line width=0.3pt},
  axis line style={InkGray, line width=0.5pt},
  every axis plot/.append style={line width=0.9pt},
]
\addplot[AgriGreen] table[x=x,y=y]{data/clean/profiles/px_000911.dat};
\addplot[SoftRed, dashed, line width=0.7pt] table[x=x,y=y]{data/clean/profiles/pxzz_000911.dat};
\addplot[only marks, mark=*, mark size=1.1pt, SoftRed] table[x=x,y=y]{data/clean/profiles/pxzz_000911.dat};
\end{axis}
\end{tikzpicture}

\captionof{figure}[*ST广糖（000911）股价与周期骨架]{*ST广糖（000911）月度收盘价（元/股）与 ZigZag 周期骨架（阈值 25\%，圆点为周期转折点）}
\end{minipage}
\end{center}

\pfl{过去周期分析} 以 25\% 的 ZigZag 阈值划分，样本区间内股价形成 5 个主升段与 4 个主跌段。 最大主升段出现在 2021 年 1 月 至 2021 年 8 月，7 个月+70.3\%；最大主跌段为 2024 年 11 月 至 2026 年 6 月，19 个月-55.4\%。 最近一次转折出现在 2026 年 10 月（上涨段自 2026 年 6 月起，4 个月+62.4\%），当前处于该段之后的震荡之中。 公司股价与申万农林牧渔指数几乎同步见底，个股并未走出独立行情。 该子行业横跨糖、番茄制品、添加剂等多个品种，各自的供给周期与政策（进口配额、关税）是主要变量，公司 2026 年 6 月末资产负债率 100.2\%，高于全样本中位数 52.2\%，其股价因此更多反映资金面与信用风险，而不是价格弹性本身。 综合区间形态看，该股单段涨跌幅度偏大、以趋势段为主（最大主升 +70.3\%、最大主跌 -55.4\%），年化波动率 40.8\% 与全样本中位数 37.3\% 相比波动与样本中位数接近。

\pfl{盈利情况} 2026 年上半年营业收入 15.82 亿元（同比 +17.4\%），归母净利润 -0.16 亿元，由上年同期的盈利转为亏损。 以年报为趋势对照，2023---2025 年营业收入由 33.42 亿元变为 24.67 亿元（两年年均复合增速 -14.1\%），归母净利润由 0.07 亿元变为 -1.20 亿元。 按半年报口径，2026 年上半年的 ROE 为 0.0\%（未年化）、销售净利率 -1.1\%、毛利率 9.1\%，ROE 在“其他农产品加工”11 家公司中按数值由高到低排第\zhnumber{7}位。 同期全样本半年 ROE 中位数 0.75\%，公司与之相差 0.75 个百分点（弱于中位数公司）。 综合看，公司仍处亏损区间、由盈转亏，半年 ROE 在三级行业内按由高到低排第\zhnumber{7}位，单位盈利能力的修复依赖行业景气与自身成本端的改善，报表弹性主要体现为价格与销量的方向性变化。

\begin{center}\small\setlength{\tabcolsep}{3.4pt}
\captionof{table}[*ST广糖（000911）关键财务指标]{*ST广糖（000911）关键财务指标（合并报表口径；两个半年报期均未年化；现金短债比 $=$ 货币资金 $\div$ 短期债务）}
\begin{adjustbox}{max width=\textwidth}
\begin{tabular}{@{}lrrrrr@{}}
\toprule
指标 & 2023 年 & 2024 年 & 2025 年 & 2025 年半年报 & 2026 年半年报 \\
\midrule
营业收入（亿元） & 33.42 & 32.75 & 24.67 & 13.48 & 15.82 \\
归母净利润（亿元） & 0.07 & 0.08 & -1.20 & 0.09 & -0.16 \\
毛利率（\%） & 12.8 & 11.9 & 9.7 & 12.1 & 9.1 \\
ROE（\%） & 41.6 & 16.4 & 0.0 & 15.6 & 0.0 \\
资产负债率（\%） & 97.0 & 96.4 & 99.8 & 95.7 & 100.2 \\
经营活动现金流净额（亿元） & 5.51 & 4.54 & 3.61 & -2.04 & -11.74 \\
货币资金（亿元） & 14.27 & 7.61 & 7.93 & 2.72 & 2.29 \\
有息负债合计（亿元） & 37.39 & 26.92 & 25.17 & 25.19 & 31.13 \\
现金短债比（倍） & 0.41 & 0.30 & 0.32 & 0.11 & 0.08 \\
\bottomrule
\end{tabular}
\end{adjustbox}
\end{center}

\begin{center}
\begin{minipage}{\textwidth}\centering
\begin{tikzpicture}
\begin{groupplot}[
  group style={group size=2 by 1, horizontal sep=0.60cm},
  width=6.85cm, height=4.60cm,
  tick label style={font=\tiny},
  title style={font=\scriptsize\heiti},
  yticklabel style={font=\tiny},
  ylabel style={font=\tiny},
  grid=major, grid style={LightGray, line width=0.3pt},
  axis line style={InkGray, line width=0.5pt},
  enlarge x limits={abs=0.8},
  legend style={font=\scriptsize, at={(0.5,-0.20)}, anchor=north,
                legend columns=2, draw=none, fill=none, column sep=6pt},
]
\nextgroupplot[
  ybar, bar width=4pt,
  ymin=-8.37, ymax=37.97,
  xtick={0,1,2,3,4,5.7},
  xticklabels={2021,2022,2023,2024,2025,2026H1},
  ylabel={亿元},
]
\addplot[fill=AgriGreen!75, draw=AgriGreen, bar shift=-2.6pt] table[x=i, y=rev]{data/clean/profiles/fin_000911.dat};
\addplot[fill=SoftRed!60, draw=SoftRed, bar shift=2.6pt] table[x=i, y=ni]{data/clean/profiles/fin_000911.dat};
\addplot[fill=AgriGreen!30, draw=AgriGreen, pattern=north east lines, bar shift=-2.6pt] table[x=x, y=rev]{data/clean/profiles/finh_000911.dat};
\addplot[fill=SoftRed!20, draw=SoftRed, pattern=north east lines, bar shift=2.6pt] table[x=x, y=ni]{data/clean/profiles/finh_000911.dat};
\legend{营业收入（年/半年）, 归母净利润（年/半年）}
\nextgroupplot[
  ymin=-191.7, ymax=127.2,
  xtick={0,1,2,3,4,5.7},
  xticklabels={2021,2022,2023,2024,2025,2026H1},
  ylabel={\%},
]
\addplot[AgriGold, mark=*, mark size=1.3pt] table[x=i, y=roe]{data/clean/profiles/fin_000911.dat};
\addplot[OilBlue, mark=square*, mark size=1.3pt] table[x=i, y=debt]{data/clean/profiles/fin_000911.dat};
\addplot[only marks, AgriGold, mark=*, mark size=2.0pt] table[x=x, y=roe]{data/clean/profiles/finh_000911.dat};
\addplot[only marks, OilBlue, mark=square*, mark size=2.0pt] table[x=x, y=debt]{data/clean/profiles/finh_000911.dat};
\legend{ROE, 资产负债率}
\end{groupplot}
\end{tikzpicture}
\begin{tikzpicture}
\begin{axis}[
  name=finbot,
  width=13.75cm, height=3.10cm, scale only axis,
  ymin=-5.31, ymax=59.48,
  xtick={0,1,2,3,4,5.7},
  xticklabels={2021,2022,2023,2024,2025,2026H1},
  tick label style={font=\tiny},
  yticklabel style={font=\tiny},
  ylabel={亿元}, ylabel style={font=\tiny},
  ybar, bar width=4pt,
  grid=major, grid style={LightGray, line width=0.3pt},
  axis line style={InkGray, line width=0.5pt},
  enlarge x limits={abs=0.8},
  legend style={font=\scriptsize, at={(0.5,-0.26)}, anchor=north,
                legend columns=2, draw=none, fill=none, column sep=10pt},
]
\addplot[fill=AgriGreen!55, draw=AgriGreen, bar shift=-3.0pt] table[x=i, y=cash]{data/clean/profiles/fin_000911.dat};
\addplot[fill=SoftRed!45, draw=SoftRed, bar shift=3.0pt] table[x=i, y=ibd]{data/clean/profiles/fin_000911.dat};
\addplot[fill=AgriGreen!25, draw=AgriGreen, pattern=north east lines, bar shift=-3.0pt] table[x=x, y=cash]{data/clean/profiles/finh_000911.dat};
\addplot[fill=SoftRed!20, draw=SoftRed, pattern=north east lines, bar shift=3.0pt] table[x=x, y=ibd]{data/clean/profiles/finh_000911.dat};
\legend{货币资金（左轴）, 有息负债（左轴）}
\end{axis}
\begin{axis}[
  at={(finbot.south east)}, anchor=south east,
  width=13.75cm, height=3.10cm, scale only axis,
  axis y line*=right, axis x line*=none, xtick=\empty,
  ymin=-16.22, ymax=10.69,
  tick label style={font=\tiny},
  yticklabel style={font=\tiny}, scaled y ticks=false,
  legend style={font=\scriptsize, at={(0.5,-0.44)}, anchor=north,
                legend columns=1, draw=none, fill=none},
]
\addplot[OilBlue, mark=*, mark size=2.2pt, line width=1.3pt] table[x=i, y=ocf]{data/clean/profiles/fin_000911.dat};
\addplot[only marks, OilBlue, mark=*, mark size=3.2pt] table[x=x, y=ocf]{data/clean/profiles/finh_000911.dat};
\legend{经营活动现金流净额（右轴，亿元，蓝点）}
\end{axis}
\end{tikzpicture}

\captionof{figure}[*ST广糖（000911）近年主要财务指标]{*ST广糖（000911）近年主要财务指标（左上：营业收入与归母净利润；右上：ROE 与资产负债率；下：货币资金、有息负债为左轴柱状，经营活动现金流净额为其右轴的蓝点折线。
2021---2025 年为年报口径，2026H1 为 2026 年半年报/半年末，与其前一年报之间留出间隔、以斜纹或空心点区分；半年报数据未年化）}
\end{minipage}
\end{center}

\pfl{财务分析} 2026 年 6 月末资产负债率 100.2\%，较 2025 年末上升 0.4 个百分点（样本中位数 52.2\%，所处三级行业内按由低到高排第\zhnumber{11}位）。 2026 年 6 月末货币资金 2.29 亿元、短期债务（短期借款与一年内到期非流动负债）29.24 亿元、长期债务 1.90 亿元，有息负债合计 31.13 亿元，现金短债比 0.08 倍——缺口明显。 净有息负债 28.85 亿元，相当于其 2026 年上半年营业收入的 182\%。 流动比率 0.60、速动比率 0.27，短期偿付压力较大。 2026 年上半年经营活动现金流净额 -11.74 亿元，经营现金流与利润同时为负，日常运营对债务与股东投入的依赖度较高。 综合看，账面货币资金对有息负债与短期债务的覆盖不足，滚动偿付对新增融资的依赖度较高；资产负债率高于样本中位数 48.0 个百分点；有息负债中短期债务占比更高，期限结构仍以滚动续作为主。 公司披露的退市与持续经营风险为：2026 年 4 月，股票交易被实施退市风险警示和其他风险警示，深交所将决定公司股票终止上市；2026 年 9 月，下一会计年度经审计期末净资产仍为负将直接进入终止上市程序。

\pfl{重大战略转型} 近三年公司的战略动作集中在：主业转型（2 项）：2023 年 10 月，公司名称由南宁糖业股份有限公司变更为广西农投糖业集团股份有限公司；其他动作（1 项）：2019 年 4 月，南宁振宁资产经营有限责任公司所持公司 7,681.38 万股（占总股本 23.70\%）国有股份无偿划转给广西农村投资集团。 公告口径上，近三年（2023 年 9 月---2026 年 9 月）公司披露重大重组与并购类公告 19 条、对外投资与转型类公告 1 条、回购与股权激励类公告 1 条，合计公告 554 条。 第一大行业仍是“工业”，收入占比 88.6\%，与 2021 年年报相比没有方向性变化。 综合判断，报告期内公司有 20 条与战略相关的公告落地（重大重组与并购 19 条、对外投资与转型 1 条），资本运作与主业调整之间存在直接关联。

\begin{center}\small\setlength{\tabcolsep}{4pt}
\captionof{table}[*ST广糖（000911）历史融资与资本运作]{*ST广糖（000911）历史融资与资本运作（据公司公告与公开报道整理；金额为披露值，含首次公开发行、再融资、债券、银行借款与重整投资等）}
\begin{adjustbox}{max width=\textwidth}
\begin{tabular}{@{}p{1.45cm}p{2.55cm}p{4.05cm}p{4.95cm}@{}}
\toprule
时间 & 方式 & 规模 & 用途与说明 \\
\midrule
1999 & IPO（深圳证券交易所主板） & 筹资总额 2.36 亿元，\allowbreak{}发行社会公众股 5,\allowbreak{}600 万股，\allowbreak{}发行后总股本 22,\allowbreak{}400 万股 & 经证监发行字[1999]33 号文批准，\allowbreak{}以上网定价方式发行 \\
2001 & 配股 & 筹资总额 1.68 亿元，\allowbreak{}配售 1,\allowbreak{}680 万股 & 经证券发行字[2001]77 号文批准，\allowbreak{}2001 年 11 月 29 日完成配售 \\
2004 & 配股 & 筹资总额 1.65 亿元，\allowbreak{}配售 2,\allowbreak{}184 万股 & 经证监发行字[2004]153 号文批准 \\
2007 & 非公开发行股票 & 筹资总额 2.93 亿元，\allowbreak{}非公开发行 2,\allowbreak{}400 万股 & 扩大业务规模、\allowbreak{}提升盈利能力（拟募集资金约 2.80 亿元）；经证监发行字[2007]311 号文核准；2007 年 11 月披露的非公开发行尽职调查报告显示 \\
2012-09-17 & 公司债券（12 南糖债，代码 112109） & 发行规模 5.40 亿元，\allowbreak{}7 年期固定利率债券，\allowbreak{}票面利率 7.25\% & 偿还借款、\allowbreak{}调整债务结构和补充流动资金，\allowbreak{}2012 年 10 月 23 日起在深交所挂牌；2017 年 9 月 17 日已兑付债券本金 3.52 亿元（含利息） \\
2015 & 非公开发行股票 & 发行 3,\allowbreak{}744.0937 万股，\allowbreak{}每股 13.55 元，\allowbreak{}募集资金总额 5.07 亿元 & 香山糖厂日榨 10,\allowbreak{}000 吨甘蔗技改工程项目 2.05 亿元、\allowbreak{}南宁云鸥物流公司食糖电子商务及智能配送中心项目 1.53 亿元、\allowbreak{}补充流动资金 1.50 亿元；；2015 年 12 月 11 日公司另以股权评估价值扣除过渡期间损益 5,\allowbreak{}241.66 万元为对价完成对广西环江远丰糖业有限责任公司 75\% 股权的收购 \\
2016-06-06 & 非公开发行股票（预案，已终止） & 拟以 16.59 元/\allowbreak{}股非公开发行 6,\allowbreak{}027.73 万股，\allowbreak{}募集资金总额 10.00 亿元 & 偿还银行贷款、\allowbreak{}补充流动资金；公司 2016 年 6 月 6 日晚间发布预案；2018 年 12 月 21 日第六届董事会 2018 年第九次临时会议审议通过终止非公开发行 A 股股票并撤相关申请材料 \\
2021-12-31 & 非公开发行股票（向控股股东农投集团） & 发行 7,\allowbreak{}623.89 万股，\allowbreak{}发行价 7.87 元/\allowbreak{}股，\allowbreak{}募集资金总额 6.00 亿元 & 认购方为控股股东广西农村投资集团有限公司（以现金认购全部股份）；新增股份于 2021 年 12 月 31 日在深交所上市，\allowbreak{}限售期 36 个月责任公司；发行前注册资本 3.24 亿元 \\
2026-04-25 & 2023 年度向特定对象发行 A 股股票（终止） & 拟募集资金总额不超过 2.60 亿元 & 云鸥物流食糖仓储智能配送中心二期扩容项目（项目投资总额 2.08 亿元、\allowbreak{}拟投入募集资金 1.94 亿元）、\allowbreak{}补充流动资金及偿还银行贷款（6,\allowbreak{}600.00 万元）；公司向特定对象发行申请获得深交所上市审核中心审核通过；2025 年 9 月 3 日收到中国证监会同意广西农投糖业；2026 年 4 月 25 日披露终止 2023 年度向特定对象发行股票事项（公告编号 2026-023） \\
\bottomrule
\end{tabular}
\end{adjustbox}
\end{center}

\pfl{历史融投资情况} 从现金流量表看，2025 年公司取得借款收到的现金 29.04 亿元、偿还债务支付的现金 30.13 亿元，筹资活动现金流净额 -2.06 亿元（2023 年为取得借款 45.31 亿元、偿还债务 46.64 亿元）。 2026 年上半年筹资活动现金流净额为 6.23 亿元（取得借款 14.75 亿元、偿还债务 9.75 亿元，未年化），仍处于净融入状态。 2025 年分配股利、利润或偿付利息支付的现金合计 0.86 亿元。 公开披露的融资记录共 9 笔，工具类型以 IPO、配股、定向增发、公司债为主；金额与用途见下表。 其中规模最大的两笔为：2012 年 9 月，公司债，发行规模 5.40 亿元，7 年期固定利率债券，票面利率 7.25\%；2015 年，定向增发，发行 3,744.0937 万股，每股 13.55 元，募集资金总额 5.07 亿元。 股本方面，近三年披露股本变动 3 次；总股本自 2021 年末以来未发生变化（4.00 亿股）。 分红方面，近三年未实施现金分红，在全部样本中属于分红能力偏弱的一端。 近三年“再融资”类公告 78 条，涉及定向增发、募集资金使用。 综合看，公司融资结构上股权与债务融资并举；2025 年筹资活动现金流净额 -2.06 亿元，融资节奏仍以维持存量债务周转为主，与前述经营与投资现金流的方向基本一致。

\pfl{未来潜在融资需求分析} 2026 年 6 月末货币资金 2.29 亿元，而短期债务 29.24 亿元（现金短债比 0.08 倍），2026 年上半年经营现金流净额 -11.74 亿元。按“短期债务减去账面货币资金”的滚动偿付口径，静态缺口约 26.95 亿元；若再计入上半年亏损的年化额（0.32 亿元），维持一年正常经营所需的外部资金约 26.95---27.27 亿元。 按第 \ref{sec:financing-need} 节的三档划分，公司属于第一档“补血型”：2026 年上半年归母净利润 -0.16 亿元、6 月末资产负债率 100.2\%。负债率已超过 80\%，净资产被亏损侵蚀的程度较深，属于全部样本中资产负债表最紧张的一组。 可行的融资路径按现实性排序为：引入国资或产业投资人（部分公司以控制权让渡为对价）、定向增发（需股价配合，且控股股东需放弃优先认购或同步增资）、出售非核心资产与参股股权。 公开披露的后续资金安排线索包括：股权融资线索：2026 年 4 月，主动终止 20.00 万元。 综合看，公司的外部资金需求以债务接续为主，滚动偿付口径下的静态缺口约 26.95 亿元，融资节奏将受制于银行续贷意愿与股权融资窗口。




\subsubsection{朗源股份（300175）}
\label{co:300175}

\pfl{公司基本情况} 山东龙口起家的果品加工出口企业，主营鲜果、干果与烘焙原料的种植管理、加工、仓储及销售，2011 年 2 月创业板上市；2018 年跨界收购数据中心企业广东优世联合，因后者连续四年虚假交易被证监会立案并于 2025 年 4 月被处以 500 万元罚款；2024 年 4 月原大股东通过表决权委托加 2 亿元借款加 100\% 股份质押方式易主杭州东方行知（实际控制人赵征），2025 年 3 月被实施其他风险警示（ST 朗源）、2026 年 6 月 3 日撤销；2025 年营业收入 2.92 亿元、归母净利润 -1,497.10 万元。 公司于 2011-02-15 在深交所创业板首发上市，注册资本 4.71 亿元，总股本 4.71 亿股。 截至 2026 年 10 月 12 日收盘价 5.23 元，对应总市值 24.6 亿元，近一年-16.1\%，总市值在三级行业“果蔬加工”的 5 家公司中按由大到小排第\zhnumber{5}位，近一年涨跌幅低于全样本中位数（-10.4\%）5.7 个百分点。 按 2025 年年报披露的行业构成，收入几乎全部来自制造业 100.0\%。

\pfl{股权结构与控制权} 按 2026 年 6 月 30 日的股东名册，前三大股东为宁波梅山保税港区龙晟行知私募基金管理合伙企业(有限合 持股 14.51\%；王贵美 持股 9.43\%；黄展鹏 持股 1.66\%，股东户数 14,943 户。 股东结构与控制权方面，2024 年 4 月，王贵美、新疆尚龙将其合计持有的 1.13 亿股（占总股本 23.95\%）所享有的表决权等股东权利不可撤销地委托给杭州东方行知股权投资基金合伙企业；2025 年 10 月，东方行知协议受让新疆尚龙及王贵美合计持有的 6,833.43 万股（占总股本 14.51\%）；2026 年二季度，王贵美持股 9.43\% 且所持股份全部冻结，公司股票自 2026 年 6 月 3 日起撤销其他风险警示。

\pfl{经营要点} 经营层面，公司位于太原的数据中心项目已于 2025 年 4 月完成对外出售。

\pfl{所处行业的基本面} 果蔬加工的价格由产地原料果价格与出口订单决定，需求端受海外需求与关税影响（第 \ref{sec:framework} 章的供需框架）。 从供需两端的当期数据看，2025 年，中国苹果总产量 5155.55 万吨、同比增加 0.53\%。

\pfl{竞争优势} 把该公司放回同类公司比较，2026 年上半年营业收入 1.7 亿元、在三级行业“果蔬加工”的 5 家公司中排第\zhnumber{5}位；毛利率 2.4\%、低于行业内中位数 6.2 个百分点；半年 ROE -2.7\%（未年化）、低于行业内中位数（2.26\%）；资产负债率 8.8\%、低于行业内中位数（30.8\%）。 就当期报表看，公司尚未把经营层面的特征转化为高于行业的资本回报，周期底部的相对优势更多体现在现金流与资产负债表上。

\pfl{历史股价} 本报告样本区间（2021 年 1 月 至 2026 年 10 月）内，股价由 4.60 元变为 5.23 元，累计 +13.7\%，区间最高 10.14 元（2024 年 5 月）、最低 3.21 元（2023 年 5 月）；当前价相当于区间最高价的 51.6\%，为区间最低价的 1.63 倍。 月度收益的年化波动率 51.9\%，高于全样本中位数 37.3\%，在三级行业“果蔬加工”的 5 家公司中波动率由高到低排第\zhnumber{2}位。 把股价阶段与公司事件对齐看，2023 年 5 月 至 2024 年 5 月股价上涨+215.9\%，同期（2024 年 5 月）受让方虞长实业为原大股东新疆尚龙指定主体。 把这些阶段与公司事件对齐后可以看到，几轮大涨大跌多由公司自身的资本运作与风险事件触发，行业景气只解释了其中一部分。

\begin{center}
\begin{minipage}{\textwidth}\centering
\begin{tikzpicture}
\begin{axis}[
  width=12.6cm, height=3.9cm, scale only axis,
  xmin=2020.922, xmax=2026.972, ymin=2.99, ymax=10.85,
  xtick={2021,2022,2023,2024,2025,2026},
  yearticks,
  yticklabel style={font=\scriptsize},
  ylabel={元/股}, ylabel style={font=\scriptsize},
  grid=major, grid style={LightGray, line width=0.3pt},
  axis line style={InkGray, line width=0.5pt},
  every axis plot/.append style={line width=0.9pt},
]
\addplot[AgriGreen] table[x=x,y=y]{data/clean/profiles/px_300175.dat};
\addplot[SoftRed, dashed, line width=0.7pt] table[x=x,y=y]{data/clean/profiles/pxzz_300175.dat};
\addplot[only marks, mark=*, mark size=1.1pt, SoftRed] table[x=x,y=y]{data/clean/profiles/pxzz_300175.dat};
\end{axis}
\end{tikzpicture}

\captionof{figure}[朗源股份（300175）股价与周期骨架]{朗源股份（300175）月度收盘价（元/股）与 ZigZag 周期骨架（阈值 25\%，圆点为周期转折点）}
\end{minipage}
\end{center}

\pfl{过去周期分析} 以 25\% 的 ZigZag 阈值划分，样本区间内股价形成 5 个主升段与 4 个主跌段。 最大主跌段为 2024 年 5 月 至 2024 年 8 月，3 个月-47.8\%；最大主升段出现在 2023 年 5 月 至 2024 年 5 月，12 个月+215.9\%。 最近一次转折出现在 2026 年 8 月（上涨段自 2026 年 6 月起，2 个月+30.2\%），当前处于该段之后的震荡之中。 公司股价较申万农林牧渔指数提前约 36 个月见底，是板块中较早定价周期反转的公司之一。 果蔬加工的外向型特征明显，出口订单与原料产地价格决定利润率，公司 2026 年 6 月末资产负债率 8.8\%、半年 ROE -2.7\%（未年化），在本轮下行中属于随行业同步调整的一类。 综合区间形态看，该股单段涨跌幅度偏大、以趋势段为主（最大主升 +215.9\%、最大主跌 -47.8\%），年化波动率 51.9\% 与全样本中位数 37.3\% 相比波动明显大于样本中位数。

\pfl{盈利情况} 2026 年上半年营业收入 1.73 亿元（同比 +67.7\%），归母净利润 -0.13 亿元，亏损同比扩大 0.08 亿元。 以年报为趋势对照，2023---2025 年营业收入由 2.20 亿元变为 2.92 亿元（两年年均复合增速 +15.1\%），归母净利润由 -0.54 亿元变为 -0.15 亿元。 按半年报口径，2026 年上半年的 ROE 为 -2.7\%（未年化）、销售净利率 -7.8\%、毛利率 2.4\%，ROE 在“果蔬加工”5 家公司中按数值由高到低排第\zhnumber{4}位。 按同一口径，三级行业“果蔬加工”5 家公司的半年 ROE 中位数为 2.26\%，该公司低于其中位数 4.97 个百分点。 从公司披露的原因看，业绩变动主要来自：2025 年度，主要原因为报告期内原材料价格上涨等原因导致整体毛利率较上年同期降低，以及因行政处罚缴纳罚款 500.00 万元。 综合看，公司仍处亏损区间、亏损同比扩大，半年 ROE 在三级行业内按由高到低排第\zhnumber{4}位，单位盈利能力的修复依赖行业景气与自身成本端的改善，报表弹性主要体现为价格与销量的方向性变化。

\begin{center}\small\setlength{\tabcolsep}{3.4pt}
\captionof{table}[朗源股份（300175）关键财务指标]{朗源股份（300175）关键财务指标（合并报表口径；两个半年报期均未年化；现金短债比 $=$ 货币资金 $\div$ 短期债务）}
\begin{adjustbox}{max width=\textwidth}
\begin{tabular}{@{}lrrrrr@{}}
\toprule
指标 & 2023 年 & 2024 年 & 2025 年 & 2025 年半年报 & 2026 年半年报 \\
\midrule
营业收入（亿元） & 2.20 & 2.45 & 2.92 & 1.03 & 1.73 \\
归母净利润（亿元） & -0.54 & -0.36 & -0.15 & -0.05 & -0.13 \\
毛利率（\%） & 16.7 & 11.2 & 8.2 & 11.7 & 2.4 \\
ROE（\%） & -10.2 & -7.6 & -2.9 & -1.0 & -2.7 \\
资产负债率（\%） & 29.6 & 28.5 & 10.1 & 6.2 & 8.8 \\
经营活动现金流净额（亿元） & -0.06 & -0.35 & -0.57 & 0.14 & 0.65 \\
货币资金（亿元） & 0.22 & 0.68 & 0.29 & 1.03 & 0.84 \\
有息负债合计（亿元） & 0.62 & 0.15 & 0.33 & 0.20 & 0.28 \\
现金短债比（倍） & 0.70 & 706.92 & 0.94 & 20.35 & 61.19 \\
\bottomrule
\end{tabular}
\end{adjustbox}
\end{center}

\begin{center}
\begin{minipage}{\textwidth}\centering
\begin{tikzpicture}
\begin{groupplot}[
  group style={group size=2 by 1, horizontal sep=0.60cm},
  width=6.85cm, height=4.60cm,
  tick label style={font=\tiny},
  title style={font=\scriptsize\heiti},
  yticklabel style={font=\tiny},
  ylabel style={font=\tiny},
  grid=major, grid style={LightGray, line width=0.3pt},
  axis line style={InkGray, line width=0.5pt},
  enlarge x limits={abs=0.8},
  legend style={font=\scriptsize, at={(0.5,-0.20)}, anchor=north,
                legend columns=2, draw=none, fill=none, column sep=6pt},
]
\nextgroupplot[
  ybar, bar width=4pt,
  ymin=-0.88, ymax=3.33,
  xtick={0,1,2,3,4,5.7},
  xticklabels={2021,2022,2023,2024,2025,2026H1},
  ylabel={亿元},
]
\addplot[fill=AgriGreen!75, draw=AgriGreen, bar shift=-2.6pt] table[x=i, y=rev]{data/clean/profiles/fin_300175.dat};
\addplot[fill=SoftRed!60, draw=SoftRed, bar shift=2.6pt] table[x=i, y=ni]{data/clean/profiles/fin_300175.dat};
\addplot[fill=AgriGreen!30, draw=AgriGreen, pattern=north east lines, bar shift=-2.6pt] table[x=x, y=rev]{data/clean/profiles/finh_300175.dat};
\addplot[fill=SoftRed!20, draw=SoftRed, pattern=north east lines, bar shift=2.6pt] table[x=x, y=ni]{data/clean/profiles/finh_300175.dat};
\legend{营业收入（年/半年）, 归母净利润（年/半年）}
\nextgroupplot[
  ymin=-13.4, ymax=33.5,
  xtick={0,1,2,3,4,5.7},
  xticklabels={2021,2022,2023,2024,2025,2026H1},
  ylabel={\%},
]
\addplot[AgriGold, mark=*, mark size=1.3pt] table[x=i, y=roe]{data/clean/profiles/fin_300175.dat};
\addplot[OilBlue, mark=square*, mark size=1.3pt] table[x=i, y=debt]{data/clean/profiles/fin_300175.dat};
\addplot[only marks, AgriGold, mark=*, mark size=2.0pt] table[x=x, y=roe]{data/clean/profiles/finh_300175.dat};
\addplot[only marks, OilBlue, mark=square*, mark size=2.0pt] table[x=x, y=debt]{data/clean/profiles/finh_300175.dat};
\legend{ROE, 资产负债率}
\end{groupplot}
\end{tikzpicture}
\begin{tikzpicture}
\begin{axis}[
  name=finbot,
  width=13.75cm, height=3.10cm, scale only axis,
  ymin=-0.09, ymax=1.02,
  xtick={0,1,2,3,4,5.7},
  xticklabels={2021,2022,2023,2024,2025,2026H1},
  tick label style={font=\tiny},
  yticklabel style={font=\tiny},
  ylabel={亿元}, ylabel style={font=\tiny},
  ybar, bar width=4pt,
  grid=major, grid style={LightGray, line width=0.3pt},
  axis line style={InkGray, line width=0.5pt},
  enlarge x limits={abs=0.8},
  legend style={font=\scriptsize, at={(0.5,-0.26)}, anchor=north,
                legend columns=2, draw=none, fill=none, column sep=10pt},
]
\addplot[fill=AgriGreen!55, draw=AgriGreen, bar shift=-3.0pt] table[x=i, y=cash]{data/clean/profiles/fin_300175.dat};
\addplot[fill=SoftRed!45, draw=SoftRed, bar shift=3.0pt] table[x=i, y=ibd]{data/clean/profiles/fin_300175.dat};
\addplot[fill=AgriGreen!25, draw=AgriGreen, pattern=north east lines, bar shift=-3.0pt] table[x=x, y=cash]{data/clean/profiles/finh_300175.dat};
\addplot[fill=SoftRed!20, draw=SoftRed, pattern=north east lines, bar shift=3.0pt] table[x=x, y=ibd]{data/clean/profiles/finh_300175.dat};
\legend{货币资金（左轴）, 有息负债（左轴）}
\end{axis}
\begin{axis}[
  at={(finbot.south east)}, anchor=south east,
  width=13.75cm, height=3.10cm, scale only axis,
  axis y line*=right, axis x line*=none, xtick=\empty,
  ymin=-0.88, ymax=1.01,
  tick label style={font=\tiny},
  yticklabel style={font=\tiny}, scaled y ticks=false,
  legend style={font=\scriptsize, at={(0.5,-0.44)}, anchor=north,
                legend columns=1, draw=none, fill=none},
]
\addplot[OilBlue, mark=*, mark size=2.2pt, line width=1.3pt] table[x=i, y=ocf]{data/clean/profiles/fin_300175.dat};
\addplot[only marks, OilBlue, mark=*, mark size=3.2pt] table[x=x, y=ocf]{data/clean/profiles/finh_300175.dat};
\legend{经营活动现金流净额（右轴，亿元，蓝点）}
\end{axis}
\end{tikzpicture}

\captionof{figure}[朗源股份（300175）近年主要财务指标]{朗源股份（300175）近年主要财务指标（左上：营业收入与归母净利润；右上：ROE 与资产负债率；下：货币资金、有息负债为左轴柱状，经营活动现金流净额为其右轴的蓝点折线。
2021---2025 年为年报口径，2026H1 为 2026 年半年报/半年末，与其前一年报之间留出间隔、以斜纹或空心点区分；半年报数据未年化）}
\end{minipage}
\end{center}

\pfl{财务分析} 2026 年 6 月末资产负债率 8.8\%，较 2025 年末下降 1.3 个百分点（样本中位数 52.2\%，所处三级行业内按由低到高排第\zhnumber{2}位）。 2026 年 6 月末货币资金 0.84 亿元、短期债务（短期借款与一年内到期非流动负债）0.01 亿元、长期债务 0.26 亿元，有息负债合计 0.28 亿元，现金短债比 61.19 倍——覆盖充分。 净有息负债 -0.56 亿元，相当于其 2026 年上半年营业收入的 -32\%。 流动比率 14.77、速动比率 8.83，短期偿债能力处于正常区间。 2026 年上半年经营活动现金流净额 0.65 亿元，净利润为负而经营现金流为正，差额主要来自折旧摊销与营运资本变动，说明亏损尚未完全转化为现金流出。 综合看，现金短债比 61.19 倍，短期偿付压力可控；资产负债率低于样本中位数 43.4 个百分点；有息负债中长期债务占比更高，期限结构对短债滚动的压力较小。 公司披露的风险集中在股东与质押、诉讼与资产冻结两类：一是 2026 年 5 月，公司股东部分股份被轮候冻结；二是 2026 年 6 月，公司银行存款中有 192.75 万元被法院冻结。

\pfl{重大战略转型} 近三年公司的战略动作集中在：并购与重组（4 项）：2025 年 4 月，子公司太原市德蓝达科技有限公司完成转让；2024 年 5 月，董事会审议通过出售优世联合资产暨关联交易，业绩补偿款债权整体打包出售给上海虞长翊实业有限公司；其他动作（3 项）：上述项目，公司通过全资子公司杭州朗源君福商贸有限公司成立控股孙公司杭州瞬知源信息科技有限公司。 公告口径上，近三年（2023 年 9 月---2026 年 9 月）公司披露重大重组与并购类公告 2 条、对外投资与转型类公告 3 条、回购与股权激励类公告 0 条，合计公告 294 条。 第一大行业“制造业”的收入占比由 2021 年年报的 85.9\% 变为 100.0\%（14.1 个百分点），收入结构出现再平衡。 综合判断，报告期内公司有 5 条与战略相关的公告落地（重大重组与并购 2 条、对外投资与转型 3 条），资本运作与主业调整之间存在直接关联。

\begin{center}\small\setlength{\tabcolsep}{4pt}
\captionof{table}[朗源股份（300175）历史融资与资本运作]{朗源股份（300175）历史融资与资本运作（据公司公告与公开报道整理；金额为披露值，含首次公开发行、再融资、债券、银行借款与重整投资等）}
\begin{adjustbox}{max width=\textwidth}
\begin{tabular}{@{}p{1.45cm}p{2.55cm}p{4.05cm}p{4.95cm}@{}}
\toprule
时间 & 方式 & 规模 & 用途与说明 \\
\midrule
2011-02-15 & IPO（深圳证券交易所创业板） & 实际发行 2,\allowbreak{}700.00 万股 & 份有限公司，\allowbreak{}发行市盈率（按发行后总股本）51.82 倍 \\
2024-04-30 & 股东借款及股份质押（控制权变更配套安排） & 杭州东方行知股权投资基金合伙企业（有限合伙）向新疆尚龙股权投资管理有限公司、\allowbreak{}王贵美提供 2.00 亿元借款；新疆尚龙、\allowbreak{}王贵美将其合计持有的 1.13 亿股股份（占公司总股本 23.95\%）全部质押给东方行知 & 用于签署表决权委托协议以取得公司控制权；截至 2024 年年度报告问询回复日，\allowbreak{}是否存在杠杆收购及控制权不稳定风险进行问询 \\
2024-05-15 & 出售子公司资产（资产出售暨关联交易） & 将公司持有的广东优世联合 74.63\% 股份、\allowbreak{}对优世联合享有的借款债权本息以及业绩补偿义务人未完成承诺业绩形成的全部业绩补偿款债权整体打包出售给上海虞长翊实业有限公司 & 回笼资金、\allowbreak{}聚焦农副产品加工业务 \\
2025-04-18 & 资产出售回款 & 转让子公司太原市德蓝达科技有限公司，\allowbreak{}截至 2026 年 6 月 30 日累计收回交易价款 7,\allowbreak{}582.69 万元 & 回笼资金；转让后德蓝达不再纳入公司合并报表范围；2025 年 4 月 8 日注销子公司广东横琴君圆商贸有限公司 \\
2025-10-30 & 股份协议转让（控股股东受让股份） & 东方行知受让新疆尚龙、\allowbreak{}王贵美合计持有的 6,\allowbreak{}833.43 万股股份（占公司总股本 14.51\%） & 东方行知由表决权委托方式取得控制权改为直接持股；本次权益变动后东方行知持有公司 14.51\% 股份 \\
2026-06-16 & 以自有资产抵押向银行申请贷款 & --- & 获取银行授信/\allowbreak{}借款；2026 年 6 月 30 日公司受限资产为：货币资金 192.75 万元（被法院冻结）、\allowbreak{}固定资产账面余额 1.83 亿元、\allowbreak{}账面价值 1.30 亿元（抵押获取银行授信/\allowbreak{}借款）、\allowbreak{}无形资产账面余额 711.17 万元、\allowbreak{}账面价值 331.07 万元（抵押获取银行授信/\allowbreak{}借款） \\
\bottomrule
\end{tabular}
\end{adjustbox}
\end{center}

\pfl{历史融投资情况} 从现金流量表看，2025 年公司取得借款收到的现金 0.44 亿元、偿还债务支付的现金 0.29 亿元、吸收权益性投资 0.00 亿元，筹资活动现金流净额 0.13 亿元（2023 年为取得借款 0.25 亿元）。 2026 年上半年筹资活动现金流净额为 -0.05 亿元（取得借款 0.25 亿元、偿还债务 0.30 亿元，未年化），已转为净流出，处于净还债状态。 2025 年分配股利、利润或偿付利息支付的现金合计 0.01 亿元。 公开披露的融资记录共 6 笔，工具类型以 IPO、银行借款为主；金额与用途见下表。 其中规模最大的两笔为：2024 年 4 月，银行借款，杭州东方行知股权投资基金合伙企业（有限合伙）向新疆尚龙股权投资管理有限公司、王贵美提供 2.00 亿元借款，新疆尚龙、王贵美将其合计持有的 1.13 亿股股份（占公司总股本 23.95\%）全部质押给东方行知；2011 年 2 月，IPO，实际发行 2,700.00 万股。 股本方面，近三年披露股本变动 3 次；总股本自 2021 年末以来未发生变化（4.71 亿股）。 分红方面，近三年未实施现金分红，在全部样本中属于分红能力偏弱的一端。 综合看，公司融资结构上以债务滚动为主、股权融资为补充；2025 年筹资活动现金流净额 0.13 亿元，年度内仍为净融入，与前述经营与投资现金流的方向基本一致。

\pfl{未来潜在融资需求分析} 2026 年 6 月末货币资金 0.84 亿元，而短期债务 0.01 亿元（现金短债比 61.19 倍），2026 年上半年经营现金流净额 0.65 亿元。账面货币资金可以覆盖短期债务，缺口不体现在静态偿付层面。 公司 2026 年上半年归母净利润为负（-0.13 亿元），但 6 月末资产负债率 8.8\% 尚未进入第一档（亏损且负债率 $>$65\%）的区间，当前融资需求以补充流动资金、支撑低谷期经营性支出为主。 公开披露的后续资金安排线索包括：债务与授信安排：2026 年 6 月，公司以自有资产抵押向银行申请贷款。 以现有货币资金与经营现金流衡量，公司不存在刚性补血的紧迫性，融资动作更多是主动型的。




\subsubsection{嘉华股份（603182）}
\label{co:603182}

\pfl{公司基本情况} 公司前身可追溯至 2017 年，山东聊城莘县起家的大豆蛋白深加工企业（前身山东嘉华保健品股份有限公司，曾在新三板挂牌），主营大豆蛋白、大豆油、大豆膳食纤维、低温食用豆粕及蒸汽电力，是国内大豆蛋白主要生产企业之一，2022 年 9 月 9 日上交所主板上市；公司原无控股股东、由张冠玲等六名一致行动人共同控制，2026 年 4 月山东省土地发展集团以协议受让加定增认购方式入主，5 月 21 日控股股东变更为土地集团、实际控制人变更为山东省国资委；2025 年营业收入 13.86 亿元、归母净利润 9,470.06 万元，2026 年上半年营业收入 8.56 亿元、归母净利润 5,666.84 万元。 公司于 2022-09-09 在上交所首发上市，注册资本 1.65 亿元，总股本 1.65 亿股。 截至 2026 年 10 月 12 日收盘价 13.30 元，对应总市值 21.9 亿元，近一年-6.5\%，总市值在三级行业“粮油加工”的 7 家公司中按由大到小排第\zhnumber{7}位，近一年涨跌幅高于全样本中位数（-10.4\%）3.8 个百分点。 按 2025 年年报披露的行业构成，收入占比前三位为农副食品加工业 92.9\%；蒸汽、电力 6.8\%；其他(补充) 0.3\%。 股权结构方面，控股股东为山东省土地发展集团有限公司（持股 22.01\%）；实际控制人为山东省人民政府国有资产监督管理委员会。

\pfl{股权结构与控制权} 按 2026 年 6 月 30 日的股东名册，前三大股东为山东省土地发展集团有限公司 持股 22.01\%；高泽林（持股比例未披露）；YUWEI WU 持股 4.41\%，股东户数 8,981 户。 从控制权角度看，2026 年 5 月，公司控股股东、实际控制人变更后，公司成为山东省国资委实际控制的上市公司。

\pfl{经营要点} 经营层面，大豆蛋白业务收入 5.18 亿元、同比增长 31.03\%、占总营收 60.55\%。

\pfl{所处行业的基本面} 粮油加工的供给由大豆、菜籽等原料进口与压榨产能决定，需求由食用消费决定，利润来自原料与成品的价差而非价格方向（第 \ref{sec:framework} 章的供需框架）。 与该公司直接相关的行业数据包括：智能工厂、自动化生产线与节能降耗工艺成为行业高质量发展的重要方向。

\pfl{竞争优势} 把该公司放回同类公司比较，2026 年上半年营业收入 8.6 亿元、在三级行业“粮油加工”的 7 家公司中排第\zhnumber{7}位；毛利率 12.7\%、高于行业内中位数 6.4 个百分点；半年 ROE 5.0\%（未年化）、高于行业内中位数（2.22\%）；资产负债率 18.8\%、低于行业内中位数（43.5\%）。 与同类公司相比，公司的差异化在于：公司大豆蛋白业务处于行业领先地位，与索宝蛋白同为国内上市的大豆蛋白企业；公司与青岛元达生物科技有限公司合资成立聊城嘉华诺康生物科技有限公司，将大豆深加工产业链延伸至后生元生物制造领域，并布局植物基肉制品、特医食品等应用方向。 优势在报表上主要体现为高于行业平均的毛利率，这在周期底部是抵御亏损的主要缓冲。

\pfl{历史股价} 以 2022 年 9 月为起点、2026 年 10 月为终点，股价由 17.23 元变为 13.30 元，累计 -22.8\%，区间最高 19.27 元（2022 年 11 月）、最低 10.28 元（2024 年 8 月）；当前价相当于区间最高价的 69.0\%，为区间最低价的 1.29 倍。 月度收益的年化波动率 25.1\%，低于全样本中位数 37.3\%，在三级行业“粮油加工”的 7 家公司中波动率由高到低排第\zhnumber{7}位。 把股价阶段与公司事件对齐看，2024 年 8 月 至 2026 年 2 月股价上涨+58.6\%，同期 2025 年末，公司定增募资用途为补充流动资金。 把这些阶段与公司事件对齐后可以看到，几轮大涨大跌多由公司自身的资本运作与风险事件触发，行业景气只解释了其中一部分。

\begin{center}
\begin{minipage}{\textwidth}\centering
\begin{tikzpicture}
\begin{axis}[
  width=12.6cm, height=3.9cm, scale only axis,
  xmin=2022.588, xmax=2026.972, ymin=9.56, ymax=20.62,
  xtick={2022,2023,2024,2025,2026},
  yearticks,
  yticklabel style={font=\scriptsize},
  ylabel={元/股}, ylabel style={font=\scriptsize},
  grid=major, grid style={LightGray, line width=0.3pt},
  axis line style={InkGray, line width=0.5pt},
  every axis plot/.append style={line width=0.9pt},
]
\addplot[AgriGreen] table[x=x,y=y]{data/clean/profiles/px_603182.dat};
\addplot[SoftRed, dashed, line width=0.7pt] table[x=x,y=y]{data/clean/profiles/pxzz_603182.dat};
\addplot[only marks, mark=*, mark size=1.1pt, SoftRed] table[x=x,y=y]{data/clean/profiles/pxzz_603182.dat};
\end{axis}
\end{tikzpicture}

\captionof{figure}[嘉华股份（603182）股价与周期骨架]{嘉华股份（603182）月度收盘价（元/股）与 ZigZag 周期骨架（阈值 25\%，圆点为周期转折点）}
\end{minipage}
\end{center}

\pfl{过去周期分析} 以 25\% 的 ZigZag 阈值划分，样本区间内股价形成 2 个主升段与 2 个主跌段。 最大主跌段为 2022 年 11 月 至 2024 年 8 月，21 个月-46.7\%；最大主升段出现在 2024 年 8 月 至 2026 年 2 月，18 个月+58.6\%。 最近一次转折出现在 2026 年 6 月（下跌段自 2026 年 2 月起，4 个月-30.1\%），当前处于该段的延续之中。 公司股价较申万农林牧渔指数提前约 21 个月见底，是板块中较早定价周期反转的公司之一。 粮油加工是“成本加成”型生意（第 \ref{sec:grain} 章），利润来自原料与成品的价差及规模效率，收入规模大而净利率薄，公司 2026 年上半年实现盈利、半年 ROE 5.0\%（未年化），好于全样本中位数（0.75\%），下行周期对其报表的冲击小于同业。 综合区间形态看，该股单段涨跌幅度偏大、以趋势段为主（最大主升 +58.6\%、最大主跌 -46.7\%），年化波动率 25.1\% 与全样本中位数 37.3\% 相比波动小于样本中位数。

\pfl{盈利情况} 2026 年上半年营业收入 8.56 亿元（同比 +25.7\%），归母净利润 0.57 亿元，同比 -9.0\%。 以年报为趋势对照，2023---2025 年营业收入由 16.95 亿元变为 13.86 亿元（两年年均复合增速 -9.6\%），归母净利润由 1.10 亿元变为 0.95 亿元。 按半年报口径，2026 年上半年的 ROE 为 5.0\%（未年化）、销售净利率 6.6\%、毛利率 12.7\%，ROE 在“粮油加工”7 家公司中按数值由高到低排第\zhnumber{1}位。 按同一口径，三级行业“粮油加工”7 家公司的半年 ROE 中位数为 2.22\%，该公司高于其中位数 2.78 个百分点。 综合看，公司在报告期维持盈利（高于全样本半年 ROE 中位数 0.75\%），利润的可持续性取决于粮油加工景气的延续与成本管控的执行。

\begin{center}\small\setlength{\tabcolsep}{3.4pt}
\captionof{table}[嘉华股份（603182）关键财务指标]{嘉华股份（603182）关键财务指标（合并报表口径；两个半年报期均未年化；现金短债比 $=$ 货币资金 $\div$ 短期债务）}
\begin{adjustbox}{max width=\textwidth}
\begin{tabular}{@{}lrrrrr@{}}
\toprule
指标 & 2023 年 & 2024 年 & 2025 年 & 2025 年半年报 & 2026 年半年报 \\
\midrule
营业收入（亿元） & 16.95 & 14.88 & 13.86 & 6.81 & 8.56 \\
归母净利润（亿元） & 1.10 & 1.09 & 0.95 & 0.62 & 0.57 \\
毛利率（\%） & 12.7 & 12.5 & 13.1 & 16.0 & 12.7 \\
ROE（\%） & 11.1 & 10.4 & 8.6 & 5.6 & 5.0 \\
资产负债率（\%） & 19.9 & 24.9 & 19.6 & 17.9 & 18.8 \\
经营活动现金流净额（亿元） & 0.37 & 3.14 & 0.79 & 0.04 & 0.70 \\
货币资金（亿元） & 3.06 & 1.47 & 0.94 & 0.75 & 1.39 \\
有息负债合计（亿元） & 0.35 & 0.25 & 0.09 & --- & 0.14 \\
现金短债比（倍） & 8.75 & 5.89 & 31.33 & --- & --- \\
\bottomrule
\end{tabular}
\end{adjustbox}
\end{center}

\begin{center}
\begin{minipage}{\textwidth}\centering
\begin{tikzpicture}
\begin{groupplot}[
  group style={group size=2 by 1, horizontal sep=0.60cm},
  width=6.85cm, height=4.60cm,
  tick label style={font=\tiny},
  title style={font=\scriptsize\heiti},
  yticklabel style={font=\tiny},
  ylabel style={font=\tiny},
  grid=major, grid style={LightGray, line width=0.3pt},
  axis line style={InkGray, line width=0.5pt},
  enlarge x limits={abs=0.8},
  legend style={font=\scriptsize, at={(0.5,-0.20)}, anchor=north,
                legend columns=2, draw=none, fill=none, column sep=6pt},
]
\nextgroupplot[
  ybar, bar width=4pt,
  ymin=-1.69, ymax=18.98,
  xtick={0,1,2,3,4,5.7},
  xticklabels={2021,2022,2023,2024,2025,2026H1},
  ylabel={亿元},
]
\addplot[fill=AgriGreen!75, draw=AgriGreen, bar shift=-2.6pt] table[x=i, y=rev]{data/clean/profiles/fin_603182.dat};
\addplot[fill=SoftRed!60, draw=SoftRed, bar shift=2.6pt] table[x=i, y=ni]{data/clean/profiles/fin_603182.dat};
\addplot[fill=AgriGreen!30, draw=AgriGreen, pattern=north east lines, bar shift=-2.6pt] table[x=x, y=rev]{data/clean/profiles/finh_603182.dat};
\addplot[fill=SoftRed!20, draw=SoftRed, pattern=north east lines, bar shift=2.6pt] table[x=x, y=ni]{data/clean/profiles/finh_603182.dat};
\legend{营业收入（年/半年）, 归母净利润（年/半年）}
\nextgroupplot[
  ymin=-3.4, ymax=46.1,
  xtick={0,1,2,3,4,5.7},
  xticklabels={2021,2022,2023,2024,2025,2026H1},
  ylabel={\%},
]
\addplot[AgriGold, mark=*, mark size=1.3pt] table[x=i, y=roe]{data/clean/profiles/fin_603182.dat};
\addplot[OilBlue, mark=square*, mark size=1.3pt] table[x=i, y=debt]{data/clean/profiles/fin_603182.dat};
\addplot[only marks, AgriGold, mark=*, mark size=2.0pt] table[x=x, y=roe]{data/clean/profiles/finh_603182.dat};
\addplot[only marks, OilBlue, mark=square*, mark size=2.0pt] table[x=x, y=debt]{data/clean/profiles/finh_603182.dat};
\legend{ROE, 资产负债率}
\end{groupplot}
\end{tikzpicture}
\begin{tikzpicture}
\begin{axis}[
  name=finbot,
  width=13.75cm, height=3.10cm, scale only axis,
  ymin=-0.44, ymax=4.90,
  xtick={0,1,2,3,4,5.7},
  xticklabels={2021,2022,2023,2024,2025,2026H1},
  tick label style={font=\tiny},
  yticklabel style={font=\tiny},
  ylabel={亿元}, ylabel style={font=\tiny},
  ybar, bar width=4pt,
  grid=major, grid style={LightGray, line width=0.3pt},
  axis line style={InkGray, line width=0.5pt},
  enlarge x limits={abs=0.8},
  legend style={font=\scriptsize, at={(0.5,-0.26)}, anchor=north,
                legend columns=2, draw=none, fill=none, column sep=10pt},
]
\addplot[fill=AgriGreen!55, draw=AgriGreen, bar shift=-3.0pt] table[x=i, y=cash]{data/clean/profiles/fin_603182.dat};
\addplot[fill=SoftRed!45, draw=SoftRed, bar shift=3.0pt] table[x=i, y=ibd]{data/clean/profiles/fin_603182.dat};
\addplot[fill=AgriGreen!25, draw=AgriGreen, pattern=north east lines, bar shift=-3.0pt] table[x=x, y=cash]{data/clean/profiles/finh_603182.dat};
\addplot[fill=SoftRed!20, draw=SoftRed, pattern=north east lines, bar shift=3.0pt] table[x=x, y=ibd]{data/clean/profiles/finh_603182.dat};
\legend{货币资金（左轴）, 有息负债（左轴）}
\end{axis}
\begin{axis}[
  at={(finbot.south east)}, anchor=south east,
  width=13.75cm, height=3.10cm, scale only axis,
  axis y line*=right, axis x line*=none, xtick=\empty,
  ymin=-0.82, ymax=4.08,
  tick label style={font=\tiny},
  yticklabel style={font=\tiny}, scaled y ticks=false,
  legend style={font=\scriptsize, at={(0.5,-0.44)}, anchor=north,
                legend columns=1, draw=none, fill=none},
]
\addplot[OilBlue, mark=*, mark size=2.2pt, line width=1.3pt] table[x=i, y=ocf]{data/clean/profiles/fin_603182.dat};
\addplot[only marks, OilBlue, mark=*, mark size=3.2pt] table[x=x, y=ocf]{data/clean/profiles/finh_603182.dat};
\legend{经营活动现金流净额（右轴，亿元，蓝点）}
\end{axis}
\end{tikzpicture}

\captionof{figure}[嘉华股份（603182）近年主要财务指标]{嘉华股份（603182）近年主要财务指标（左上：营业收入与归母净利润；右上：ROE 与资产负债率；下：货币资金、有息负债为左轴柱状，经营活动现金流净额为其右轴的蓝点折线。
2021---2025 年为年报口径，2026H1 为 2026 年半年报/半年末，与其前一年报之间留出间隔、以斜纹或空心点区分；半年报数据未年化）}
\end{minipage}
\end{center}

\pfl{财务分析} 2026 年 6 月末资产负债率 18.8\%，较 2025 年末下降 0.8 个百分点（样本中位数 52.2\%，所处三级行业内按由低到高排第\zhnumber{2}位）。 流动比率 3.27、速动比率 1.98，短期流动性无明显缺口。 2026 年上半年经营活动现金流净额 0.70 亿元，经营现金流与利润同向为正，内生资金可以支撑运营。 综合看，资产负债率低于样本中位数 33.4 个百分点；有息负债中长期债务占比更高，期限结构对短债滚动的压力较小。

\pfl{重大战略转型} 近三年公司的战略动作集中在：融资与资本运作（2 项）：2024 年 4 月，调增“高端大豆蛋白生产基地建设项目”募集资金投资金额 9,833.14 万元；并购与重组（2 项）：2026 年 4 月，公司控股权拟转让给山东省土地发展集团有限公司，并以定增方式引入其现金认购不超过 2.66 亿元。 公告口径上，近三年（2023 年 9 月---2026 年 9 月）公司披露重大重组与并购类公告 7 条、对外投资与转型类公告 0 条、回购与股权激励类公告 0 条，合计公告 294 条。 第一大行业仍是“农副食品加工业”，收入占比 92.9\%，与 2021 年年报相比没有方向性变化。 综合判断，报告期内公司有 7 条与战略相关的公告落地（重大重组与并购 7 条、对外投资与转型 0 条），资本运作与主业调整之间存在直接关联。

\begin{center}\small\setlength{\tabcolsep}{4pt}
\captionof{table}[嘉华股份（603182）历史融资与资本运作]{嘉华股份（603182）历史融资与资本运作（据公司公告与公开报道整理；金额为披露值，含首次公开发行、再融资、债券、银行借款与重整投资等）}
\begin{adjustbox}{max width=\textwidth}
\begin{tabular}{@{}p{1.45cm}p{2.55cm}p{4.05cm}p{4.95cm}@{}}
\toprule
时间 & 方式 & 规模 & 用途与说明 \\
\midrule
2017-05 & 新三板定向股票发行（挂牌期间） & 发行股份 341.00 万股，\allowbreak{}每股价格 6.50 元，\allowbreak{}募集资金 2,\allowbreak{}216.50 万元 & 2017 年 1 月 19 日第三届董事会第四次会议、\allowbreak{}2017 年第一次临时股东大会审议通过，\allowbreak{}融资额不超过 2,\allowbreak{}307.50 万元；经全国中小企业股份转让系统股转系统函【2017】2529 号确认实际发行 341 万股 \\
2022-09-09 & IPO（上海证券交易所主板） & 公开发行新股 4,\allowbreak{}114.00 万股 & 扣除发行费用后的募集资金净额全部用于公司募投项目；主要募投项目为“高端大豆蛋白生产基地建设项目”；发行数量占发行后公司总股本的 25\% \\
2024-04-16 & 变更募投项目募集资金投资金额（调增） & 调增“高端大豆蛋白生产基地建设项目”募集资金投资金额 9,\allowbreak{}833.14 万元（变更前该项目募集资金投资总额 1.13 亿元） & 继续投入高端大豆蛋白生产基地建设项目；该项目为首次公开发行股票承诺投资项目，\allowbreak{}项目未取消；2023 年年度报告及 2024 年半年度报告披露该项目截至 2023 年末累计投入募集资金 1.23 亿元（投入进度未达计划的原因详见公司） \\
2026-04-07 & 股份协议转让（控制权转让） & 山东省土地发展集团有限公司受让张冠玲、\allowbreak{}李广庆、\allowbreak{}贾辉、\allowbreak{}黄瑞华、\allowbreak{}张效伟、\allowbreak{}陈春佳、\allowbreak{}赵冬杰、\allowbreak{}赵珂欣、\allowbreak{}田丰合计持有的 3,\allowbreak{}622.37 万股股份（占公司总股本的 22.01\%） & 取得上市公司控制权；协议转让完成后土地集团持有上市公司 3,\allowbreak{}622.37 万股（22.01\%） \\
2026-04-07 & 向特定对象发行 A 股股票（定增，认购方为拟入主方） & 发行数量不超过 2,\allowbreak{}143.02 万股（含本数，\allowbreak{}不超过本次发行前公司总股本的 30\%） & 扣除发行费用后全部用于补充流动资金；定价基准日为公司第六届董事会第八次会议决议公告日，\allowbreak{}属于上市规则规定的关联方 \\
2026-05-21 & 控制权变更完成（股份过户） & 土地集团协议受让的 3,\allowbreak{}622.37 万股（占 22.01\%）完成过户 & 公司纳入山东省属国资体系；2026 年公司已取得上海证券交易所合规性确认意见；截至 2026 年 6 月 30 日土地集团位列公司第一大股东，\allowbreak{}持股比例 22.01\% \\
\bottomrule
\end{tabular}
\end{adjustbox}
\end{center}

\pfl{历史融投资情况} 从现金流量表看，2025 年公司取得借款收到的现金 0.09 亿元、偿还债务支付的现金 0.25 亿元，筹资活动现金流净额 -0.82 亿元（2023 年为取得借款 1.25 亿元、偿还债务 0.90 亿元）。 2026 年上半年筹资活动现金流净额为 0.05 亿元（取得借款 0.08 亿元、偿还债务 0.03 亿元，未年化），仍处于净融入状态。 2025 年分配股利、利润或偿付利息支付的现金合计 0.66 亿元。 公开披露的融资记录共 6 笔，工具类型以 IPO、定向增发为主；金额与用途见下表。 其中规模最大的两笔为：2026 年 4 月，定向增发，发行数量不超过 2,143.02 万股（含本数，不超过本次发行前公司总股本的 30\%）；2024 年 4 月，变更募投项目募集资金，调增“高端大豆蛋白生产基地建设项目”募集资金投资金额 9,833.14 万元（变更前该项目募集资金投资总额 1.13 亿元）。 股本方面，近三年披露股本变动 4 次（限售股份上市 1 次）；总股本自 2021 年末以来未发生变化（1.65 亿股）。 分红方面，近三年现金分红 6 次、合计每股派现 1.50 元，最近一次实施在 2026 年 6 月。 近三年“再融资”类公告 9 条，涉及定向增发、募集资金使用。 综合看，公司融资结构上股权与债务融资并举；2025 年筹资活动现金流净额 -0.82 亿元，融资节奏仍以维持存量债务周转为主，与前述经营与投资现金流的方向基本一致。

\pfl{未来潜在融资需求分析} 公司 2026 年 6 月末资产负债率 18.8\%、2026 年上半年 ROE 5.0\%（未年化），资产端与负债端都没有明显压力，融资需求以相机抉择的扩张型用途为主。 公开披露的后续资金安排线索包括：股权融资线索：2025 年末，公司定增募资用途为补充流动资金；2026 年 4 月，公司 2026 年度向特定对象发行 A 股股票预案。




\subsubsection{佳沃食品（300268）}
\label{co:300268}

\pfl{公司基本情况} 联想控股旗下佳沃集团控制的动物蛋白上市公司（深交所创业板，2011 年 9 月上市，曾名万福生科、佳沃股份）。2016 年 12 月佳沃集团通过表决权委托取得控制权，2017 年置出大米加工资产、切入水产加工贸易，2019 年以约 60 亿元现金收购智利三文鱼企业 Australis 后陷入高负债与连年巨亏；2022 年末、2024 年末两度因净资产为负被实施退市风险警示，2025 年 6 月以 1 元对价将三文鱼业务子公司北京臻诚 100\% 股权出售给控股股东全资子公司实现出表与净资产转正，2026 年 5 月 13 日撤销退市风险警示及其他风险警示。 公司于 2011-09-27 在深交所创业板首发上市，注册资本 1.74 亿元，总股本 1.74 亿股。 截至 2026 年 10 月 12 日收盘价 11.30 元，对应总市值 19.7 亿元，近一年+10.6\%，总市值在三级行业“其他农产品加工”的 11 家公司中按由大到小排第\zhnumber{11}位，近一年涨跌幅高于全样本中位数（-10.4\%）20.9 个百分点。 按 2025 年年报披露的行业构成，收入占比前两位为动物蛋白 99.9\%；其他 0.1\%。

\pfl{股权结构与控制权} 按 2026 年 6 月 30 日的股东名册，前三大股东为佳沃集团有限公司 持股 46.08\%；陈彦明（持股比例未披露）；呼延合东（持股比例未披露），股东户数 6,197 户。 股东结构与控制权方面，截至 2021 年 8 月公司总股本 1.74 亿股；公司股票交易撤销退市风险警示及其他风险警示的申请获深交所审核同意。

\pfl{经营要点} 经营与产能方面，销量及售价受全球供需市场及中国消费需求影响。

\pfl{所处行业的基本面} 该子行业横跨制糖、番茄制品、糖醇与添加剂，供给受种植面积与进口配额约束，需求端分食用与工业两类（第 \ref{sec:framework} 章的供需框架）。 从供需两端的当期数据看，2026 年前两个月俄罗斯 H\&G 狭鳕产量约 15.3 万吨、同比下滑 13\%；受真鳕、黑线鳕和巴沙鱼等品种价格高企或供应短缺影响，买家转向价格更亲民的狭鳕；2025 年挪威对华北极甜虾出口突破 7.1 万吨、同比增长 8\%，并借电商直播与精品商超渠道加速向零售端渗透。

\pfl{竞争优势} 把该公司放回同类公司比较，2026 年上半年营业收入 4.2 亿元、在三级行业“其他农产品加工”的 11 家公司中排第\zhnumber{9}位；毛利率 8.4\%、低于行业内中位数 0.6 个百分点；半年 ROE -1.4\%（未年化）、低于行业内中位数（1.36\%）；资产负债率 35.5\%、低于行业内中位数（39.1\%）。 公司在公开资料中自述与同业比较的相对优势在于：公司是国内少数以狭鳕鱼、北极甜虾等野生捕捞海产品为核心业务的上市公司；公司通过 2025 年 6 月剥离三文鱼业务子公司北京佳沃臻诚，公司净资产由负转正并于 2026 年 5 月撤销退市风险警示。 就当期报表看，公司尚未把经营层面的特征转化为高于行业的资本回报，周期底部的相对优势更多体现在现金流与资产负债表上。

\pfl{历史股价} 自 2021 年 1 月至 2026 年 10 月，股价由 8.98 元变为 11.30 元，累计 +25.8\%，区间最高 30.33 元（2022 年 5 月）、最低 6.27 元（2024 年 6 月）；当前价相当于区间最高价的 37.3\%，为区间最低价的 1.80 倍。 月度收益的年化波动率 59.0\%，高于全样本中位数 37.3\%，在三级行业“其他农产品加工”的 11 家公司中波动率由高到低排第\zhnumber{1}位。 把股价阶段与公司事件对齐看，2024 年 6 月 至 2026 年 2 月股价上涨+130.3\%，同期（2025 年 12 月）重大资产重组使公司实现净资产转正、净利润大幅减亏。 从时间对应关系看，公司层面的资本运作与股权、风险事件解释了区间内多数急涨急跌，与行业指数的同步性相对有限。

\begin{center}
\begin{minipage}{\textwidth}\centering
\begin{tikzpicture}
\begin{axis}[
  width=12.6cm, height=3.9cm, scale only axis,
  xmin=2020.922, xmax=2026.972, ymin=5.83, ymax=32.45,
  xtick={2021,2022,2023,2024,2025,2026},
  yearticks,
  yticklabel style={font=\scriptsize},
  ylabel={元/股}, ylabel style={font=\scriptsize},
  grid=major, grid style={LightGray, line width=0.3pt},
  axis line style={InkGray, line width=0.5pt},
  every axis plot/.append style={line width=0.9pt},
]
\addplot[AgriGreen] table[x=x,y=y]{data/clean/profiles/px_300268.dat};
\addplot[SoftRed, dashed, line width=0.7pt] table[x=x,y=y]{data/clean/profiles/pxzz_300268.dat};
\addplot[only marks, mark=*, mark size=1.1pt, SoftRed] table[x=x,y=y]{data/clean/profiles/pxzz_300268.dat};
\end{axis}
\end{tikzpicture}

\captionof{figure}[佳沃食品（300268）股价与周期骨架]{佳沃食品（300268）月度收盘价（元/股）与 ZigZag 周期骨架（阈值 25\%，圆点为周期转折点）}
\end{minipage}
\end{center}

\pfl{过去周期分析} 以 25\% 的 ZigZag 阈值划分，样本区间内股价形成 5 个主升段与 4 个主跌段。 最大主升段出现在 2021 年 1 月 至 2021 年 5 月，4 个月+157.5\%；最大主跌段为 2022 年 12 月 至 2024 年 6 月，18 个月-73.3\%。 最近一次转折出现在 2026 年 8 月（上涨段自 2026 年 6 月起，2 个月+32.7\%），当前处于该段之后的震荡之中。 公司股价较申万农林牧渔指数提前约 23 个月见底，是板块中较早定价周期反转的公司之一。 该子行业横跨糖、番茄制品、添加剂等多个品种，各自的供给周期与政策（进口配额、关税）是主要变量，公司 2026 年 6 月末资产负债率 35.5\%、半年 ROE -1.4\%（未年化），在本轮下行中属于随行业同步调整的一类。 综合区间形态看，该股单段涨跌幅度偏大、以趋势段为主（最大主升 +157.5\%、最大主跌 -73.3\%），年化波动率 59.0\% 与全样本中位数 37.3\% 相比波动明显大于样本中位数。

\pfl{盈利情况} 2026 年上半年营业收入 4.21 亿元（同比 -66.1\%），归母净利润 -0.04 亿元，亏损同比收窄 4.15 亿元。 以年报为趋势对照，2023---2025 年营业收入由 45.05 亿元变为 16.57 亿元（两年年均复合增速 -39.4\%），归母净利润由 -11.52 亿元变为 -4.27 亿元。 按半年报口径，2026 年上半年的 ROE 为 -1.4\%（未年化）、销售净利率 -0.8\%、毛利率 8.4\%，ROE 在“其他农产品加工”11 家公司中按数值由高到低排第\zhnumber{8}位。 同期全样本半年 ROE 中位数 0.75\%，公司与之相差 2.14 个百分点（弱于中位数公司）。 从公司披露的原因看，业绩变动主要来自：公司狭鳕鱼业务面临原料价格持续飙升，导致销售收入下滑、毛利率下降。 综合看，公司仍处亏损区间、亏损同比收窄，半年 ROE 在三级行业内按由高到低排第\zhnumber{8}位，单位盈利能力的修复依赖行业景气与自身成本端的改善，报表弹性主要体现为价格与销量的方向性变化。

\begin{center}\small\setlength{\tabcolsep}{3.4pt}
\captionof{table}[佳沃食品（300268）关键财务指标]{佳沃食品（300268）关键财务指标（合并报表口径；两个半年报期均未年化；现金短债比 $=$ 货币资金 $\div$ 短期债务）}
\begin{adjustbox}{max width=\textwidth}
\begin{tabular}{@{}lrrrrr@{}}
\toprule
指标 & 2023 年 & 2024 年 & 2025 年 & 2025 年半年报 & 2026 年半年报 \\
\midrule
营业收入（亿元） & 45.05 & 34.18 & 16.57 & 12.45 & 4.21 \\
归母净利润（亿元） & -11.52 & -9.24 & -4.27 & -4.19 & -0.04 \\
毛利率（\%） & -5.4 & -10.5 & -2.7 & -6.1 & 8.4 \\
ROE（\%） & --- & --- & --- & --- & -1.4 \\
资产负债率（\%） & 95.6 & 104.9 & 10.9 & 23.8 & 35.5 \\
经营活动现金流净额（亿元） & -0.54 & -1.44 & -5.15 & -6.17 & -2.18 \\
货币资金（亿元） & 1.44 & 2.16 & 1.00 & 0.48 & 0.48 \\
有息负债合计（亿元） & 32.03 & 13.68 & 0.02 & 0.42 & 1.72 \\
现金短债比（倍） & 0.05 & 1.39 & 71.83 & 1.16 & 0.28 \\
\bottomrule
\end{tabular}
\end{adjustbox}
\end{center}

\begin{center}
\begin{minipage}{\textwidth}\centering
\begin{tikzpicture}
\begin{groupplot}[
  group style={group size=2 by 1, horizontal sep=0.60cm},
  width=6.85cm, height=4.60cm,
  tick label style={font=\tiny},
  title style={font=\scriptsize\heiti},
  yticklabel style={font=\tiny},
  ylabel style={font=\tiny},
  grid=major, grid style={LightGray, line width=0.3pt},
  axis line style={InkGray, line width=0.5pt},
  enlarge x limits={abs=0.8},
  legend style={font=\scriptsize, at={(0.5,-0.20)}, anchor=north,
                legend columns=2, draw=none, fill=none, column sep=6pt},
]
\nextgroupplot[
  ybar, bar width=4pt,
  ymin=-18.19, ymax=63.29,
  xtick={0,1,2,3,4,5.7},
  xticklabels={2021,2022,2023,2024,2025,2026H1},
  ylabel={亿元},
]
\addplot[fill=AgriGreen!75, draw=AgriGreen, bar shift=-2.6pt] table[x=i, y=rev]{data/clean/profiles/fin_300268.dat};
\addplot[fill=SoftRed!60, draw=SoftRed, bar shift=2.6pt] table[x=i, y=ni]{data/clean/profiles/fin_300268.dat};
\addplot[fill=AgriGreen!30, draw=AgriGreen, pattern=north east lines, bar shift=-2.6pt] table[x=x, y=rev]{data/clean/profiles/finh_300268.dat};
\addplot[fill=SoftRed!20, draw=SoftRed, pattern=north east lines, bar shift=2.6pt] table[x=x, y=ni]{data/clean/profiles/finh_300268.dat};
\legend{营业收入（年/半年）, 归母净利润（年/半年）}
\nextgroupplot[
  ymin=-48.4, ymax=119.1,
  xtick={0,1,2,3,4,5.7},
  xticklabels={2021,2022,2023,2024,2025,2026H1},
  ylabel={\%},
]
\addplot[AgriGold, mark=*, mark size=1.3pt] table[x=i, y=roe]{data/clean/profiles/fin_300268.dat};
\addplot[OilBlue, mark=square*, mark size=1.3pt] table[x=i, y=debt]{data/clean/profiles/fin_300268.dat};
\addplot[only marks, AgriGold, mark=*, mark size=2.0pt] table[x=x, y=roe]{data/clean/profiles/finh_300268.dat};
\addplot[only marks, OilBlue, mark=square*, mark size=2.0pt] table[x=x, y=debt]{data/clean/profiles/finh_300268.dat};
\legend{ROE, 资产负债率}
\end{groupplot}
\end{tikzpicture}
\begin{tikzpicture}
\begin{axis}[
  name=finbot,
  width=13.75cm, height=3.10cm, scale only axis,
  ymin=-4.83, ymax=54.15,
  xtick={0,1,2,3,4,5.7},
  xticklabels={2021,2022,2023,2024,2025,2026H1},
  tick label style={font=\tiny},
  yticklabel style={font=\tiny},
  ylabel={亿元}, ylabel style={font=\tiny},
  ybar, bar width=4pt,
  grid=major, grid style={LightGray, line width=0.3pt},
  axis line style={InkGray, line width=0.5pt},
  enlarge x limits={abs=0.8},
  legend style={font=\scriptsize, at={(0.5,-0.26)}, anchor=north,
                legend columns=2, draw=none, fill=none, column sep=10pt},
]
\addplot[fill=AgriGreen!55, draw=AgriGreen, bar shift=-3.0pt] table[x=i, y=cash]{data/clean/profiles/fin_300268.dat};
\addplot[fill=SoftRed!45, draw=SoftRed, bar shift=3.0pt] table[x=i, y=ibd]{data/clean/profiles/fin_300268.dat};
\addplot[fill=AgriGreen!25, draw=AgriGreen, pattern=north east lines, bar shift=-3.0pt] table[x=x, y=cash]{data/clean/profiles/finh_300268.dat};
\addplot[fill=SoftRed!20, draw=SoftRed, pattern=north east lines, bar shift=3.0pt] table[x=x, y=ibd]{data/clean/profiles/finh_300268.dat};
\legend{货币资金（左轴）, 有息负债（左轴）}
\end{axis}
\begin{axis}[
  at={(finbot.south east)}, anchor=south east,
  width=13.75cm, height=3.10cm, scale only axis,
  axis y line*=right, axis x line*=none, xtick=\empty,
  ymin=-8.25, ymax=10.36,
  tick label style={font=\tiny},
  yticklabel style={font=\tiny}, scaled y ticks=false,
  legend style={font=\scriptsize, at={(0.5,-0.44)}, anchor=north,
                legend columns=1, draw=none, fill=none},
]
\addplot[OilBlue, mark=*, mark size=2.2pt, line width=1.3pt] table[x=i, y=ocf]{data/clean/profiles/fin_300268.dat};
\addplot[only marks, OilBlue, mark=*, mark size=3.2pt] table[x=x, y=ocf]{data/clean/profiles/finh_300268.dat};
\legend{经营活动现金流净额（右轴，亿元，蓝点）}
\end{axis}
\end{tikzpicture}

\captionof{figure}[佳沃食品（300268）近年主要财务指标]{佳沃食品（300268）近年主要财务指标（左上：营业收入与归母净利润；右上：ROE 与资产负债率；下：货币资金、有息负债为左轴柱状，经营活动现金流净额为其右轴的蓝点折线。
2021---2025 年为年报口径，2026H1 为 2026 年半年报/半年末，与其前一年报之间留出间隔、以斜纹或空心点区分；半年报数据未年化）}
\end{minipage}
\end{center}

\pfl{财务分析} 2026 年 6 月末资产负债率 35.5\%，较 2025 年末上升 24.6 个百分点（样本中位数 52.2\%，所处三级行业内按由低到高排第\zhnumber{5}位）。 2026 年 6 月末货币资金 0.48 亿元、短期债务（短期借款与一年内到期非流动负债）1.70 亿元、长期债务 0.01 亿元，有息负债合计 1.72 亿元，现金短债比 0.28 倍——缺口明显。 净有息负债 1.23 亿元，相当于其 2026 年上半年营业收入的 29\%。 流动比率 2.39、速动比率 0.72，短期指标尚可。 2026 年上半年经营活动现金流净额 -2.18 亿元，经营现金流与利润同时为负，现金消耗较快。 综合看，账面货币资金对有息负债与短期债务的覆盖不足，滚动偿付对新增融资的依赖度较高；资产负债率低于样本中位数 16.7 个百分点；有息负债中短期债务占比更高，期限结构仍以滚动续作为主。 公司披露的退市与持续经营风险为：2026 年 5 月，公司股票交易将被实施退市风险警示并继续被实施其他风险警示暨股票停牌。

\pfl{重大战略转型} 近三年公司的战略动作集中在：并购与重组（4 项）：2025 年 12 月，重大资产重组使公司实现净资产转正、净利润大幅减亏；2019 年 7 月，公司完成对智利三文鱼企业 Australis Seafoods S.A. 的现金收购（持股 99.84\%）。 公告口径上，近三年（2023 年 9 月---2026 年 9 月）公司披露重大重组与并购类公告 20 条、对外投资与转型类公告 1 条、回购与股权激励类公告 2 条，合计公告 495 条。 第一大行业仍是“动物蛋白”，收入占比 99.9\%，与 2021 年年报相比没有方向性变化。 综合判断，报告期内公司有 21 条与战略相关的公告落地（重大重组与并购 20 条、对外投资与转型 1 条），资本运作与主业调整之间存在直接关联。

\begin{center}\small\setlength{\tabcolsep}{4pt}
\captionof{table}[佳沃食品（300268）历史融资与资本运作]{佳沃食品（300268）历史融资与资本运作（据公司公告与公开报道整理；金额为披露值，含首次公开发行、再融资、债券、银行借款与重整投资等）}
\begin{adjustbox}{max width=\textwidth}
\begin{tabular}{@{}p{1.45cm}p{2.55cm}p{4.05cm}p{4.95cm}@{}}
\toprule
时间 & 方式 & 规模 & 用途与说明 \\
\midrule
2017-08 & 重大资产重组（资产出售及资产购买） & 以现金 1.16 亿元向佳沃集团出售万福农业 100\% 股权、\allowbreak{}以现金 5,\allowbreak{}706.31 万元向佳沃集团出售万福粮油 100\% 股权；同时以现金 1.74 亿元增资取得国星股份 3,\allowbreak{}626.70 万股 & 置出停产的大米精粗加工资产 \\
2019 & 向特定对象发行股票（预案，已撤回） & 拟以 10.35 元/\allowbreak{}股的价格定增募集资金不超过 10.00 亿元 & 用于收购智利子公司 Australis；2020 年 3 月公司撤回定增申请 \\
2019-07 & 重大资产购买（现金收购智利三文鱼企业） & 收购智利三文鱼公司 Australis Seafoods S.A.，\allowbreak{}持股比例达到 99.84\% & 取得三文鱼上游养殖端资源，\allowbreak{}本次交易不属于重组上市（不构成借壳）；媒体报道另有 58.87 亿元的并购对价口径 \\
2021-07 & 向特定对象发行股票（募集说明书） & 本次发行募集资金总额不超过 9.20 亿元（含本数） & 智利三文鱼智能工厂项目（项目总投资 5.20 亿元、\allowbreak{}拟投入募集资金 1.70 亿元）等；发行对象为包括佳沃集团在内不超过 35 名特定投资者；发行前佳沃集团持有公司 8,\allowbreak{}026.59 万股、\allowbreak{}占发行前股本总额的 46.08\% \\
2021-12 & 向特定对象发行股票（终止） & --- & 2021 年 12 月公司决定终止 2021 年定增再融资并向深交所申请撤回相关申请材料 \\
2022-04-13 & 向特定对象发行股票（向控股股东，预案） & 拟以 14.8 元/\allowbreak{}股向控股股东佳沃集团发行不低于 3,\allowbreak{}484 万股且不超过 5,\allowbreak{}226 万股 & 全部用于补充流动资金；由控股股东全额认购 \\
2023-10-24 & 控股股东借款额度（关联方资金支持） & 在原有已向公司提供不超过 30.00 亿元人民币及不超过 1.45 亿美元的基础上 & 满足公司对日常营运资金周转的需求 \\
2023-11-29 & 控股股东债务豁免（资本性投入） & 佳沃集团同意豁免公司对其负有的债务本金 8.39 亿元及相关应付利息 2.42 亿元；佳沃集团下属子公司同意豁免公司下属子公司 9,\allowbreak{}993.00 万美元的债务本金 & 减轻公司债务压力、\allowbreak{}优化资产负债结构、\allowbreak{}提升持续经营能力；经公司董事会、\allowbreak{}监事会审议通过 \\
2024-12-31 & 控股股东及子公司债务利息豁免（关联方支持） & 报告期内控股股东佳沃集团有限公司及其子公司豁免公司部分借款利息，\allowbreak{}导致公司资本公积增加 1.99 亿元 & 降低财务费用、\allowbreak{}改善资本结构；截至报告期末公司应付关联方 Joyvio (Hong Kong) Investment Limited、\allowbreak{}Rosy Nation Investment Limited、\allowbreak{}Quafood Limited、\allowbreak{}In Food Limited 合计 8.04 亿元债务利率为 7.20\%；报告期末公司部分应付账款处于逾期状态 \\
2025-06-17 & 重大资产出售（向控股股东全资子公司出售三文鱼业务） & 将持有的北京佳沃臻诚科技有限公司 100\% 股权出售给佳沃品鲜（北京）企业管理有限公司；以 2024 年 12 月 31 日为评估基准日 & 剥离从事三文鱼养殖加工销售业务的子公司 \\
\bottomrule
\end{tabular}
\end{adjustbox}
\end{center}

\pfl{历史融投资情况} 从现金流量表看，2025 年公司取得借款收到的现金 0.92 亿元、偿还债务支付的现金 1.31 亿元，筹资活动现金流净额 5.43 亿元（2023 年为取得借款 3.53 亿元、偿还债务 15.33 亿元）。 2026 年上半年筹资活动现金流净额为 1.66 亿元（取得借款 1.94 亿元、偿还债务 0.26 亿元，未年化），仍处于净融入状态。 2025 年分配股利、利润或偿付利息支付的现金合计 0.23 亿元。 公开披露的融资记录共 10 笔，工具类型以定向增发、银行借款、公司债为主；金额与用途见下表。 其中规模最大的两笔为：2023 年 11 月，控股股东债务豁免，佳沃集团同意豁免公司对其负有的债务本金 8.39 亿元及相关应付利息 2.42 亿元，佳沃集团下属子公司同意豁免公司下属子公司 9,993.00 万美元的债务本金；2023 年 10 月，银行借款，在原有已向公司提供不超过 30.00 亿元人民币及不超过 1.45 亿美元的基础上。 股本方面，近三年披露股本变动 3 次；总股本自 2021 年末以来未发生变化（1.74 亿股）。 分红方面，近三年未实施现金分红，在全部样本中属于分红能力偏弱的一端。 综合看，公司融资结构上股权与债务融资并举；2025 年筹资活动现金流净额 5.43 亿元，年度内仍为净融入，与前述经营与投资现金流的方向基本一致。

\pfl{未来潜在融资需求分析} 2026 年 6 月末货币资金 0.48 亿元，而短期债务 1.70 亿元（现金短债比 0.28 倍），2026 年上半年经营现金流净额 -2.18 亿元。按“短期债务减去账面货币资金”的滚动偿付口径，静态缺口约 1.22 亿元；若再计入上半年亏损的年化额（0.08 亿元），维持一年正常经营所需的外部资金约 1.22---1.30 亿元。 按第 \ref{sec:financing-need} 节的三档划分，公司属于第二档“保壳型”：近三年重大重组与并购类公告 20 条，2026 年上半年归母净利润 -0.04 亿元，融资需求更多体现为“资产置换 + 维持上市地位”，而不是扩产。 公开披露的后续资金安排线索包括：股权融资线索：2026 年 6 月，公司 2019 年至 2022 年的四次定增中仅一次成功；债务与授信安排：2025 年 3 月，第五届董事会第十三次会议和第五届监事会第八次会议审议通过 2025 年度申请综合授信、融资额度及提供担保和 2025 年度接受关联方提供担保暨关联交易：同意公司对子公司以及控股子公司之间预计提供不超过人民币 76.00 亿元担保。 综合看，公司的外部资金需求以营运资金周转与在建项目投入为主，滚动偿付口径下的静态缺口约 1.22 亿元，融资节奏将受制于银行续贷意愿与股权融资窗口。






\section{渔业、林业与宠物食品：三个互不相关的子行业}
\label{sec:fishery}

本章合并讨论 \num{13} 家公司：水产养殖 \num{4}、海洋捕捞 \num{2}、林业 \num{4}、宠物食品 \num{3}
（此外三级行业“农业综合Ⅱ”下的 \num{2} 家也并入本章画像，合计 \num{15} 家）。
它们同属申万农林牧渔一级行业，但\hl{周期性来源完全不同}：
按 2026 年半年报口径，这 \num{13} 家公司合计盈利 \num{3.81} 亿元、仍有 \num{5} 家亏损；
2025 年年报口径下则合计净亏损 \num{7.39} 亿元（趋势对照）。

\begingroup\scriptsize\setlength{\tabcolsep}{2pt}
\begin{longtable}{@{}llrrrrrrr@{}}
\caption{农业上市公司分行业明细（水产养殖、海洋捕捞、林业Ⅲ、宠物食品；市值时点 2026-10-12，总市值口径；财务为 2026 年半年报，ROE 未年化，公告计数窗口 2023-09---2026-10）}\label{tab:comp-fishery}\\
\toprule
代码 & 公司简称 & 总市值（亿元） & 营收（亿元） & 同比（\%） & 净利（亿元） & ROE（\%） & 再融资公告 & 重组公告 \\
\midrule
\endfirsthead
\multicolumn{9}{l}{\small（续）}\\
\toprule
代码 & 公司简称 & 总市值（亿元） & 营收（亿元） & 同比（\%） & 净利（亿元） & ROE（\%） & 再融资公告 & 重组公告 \\
\midrule
\endhead
\bottomrule
\endlastfoot
301498 & 乖宝宠物 & 156.8 & 35.4 & 10.0 & 1.91 & 4.1 & 7 & 0 \\
000592 & 平潭发展 & 142.4 & 4.8 & -35.1 & -0.69 & -4.0 & 1 & 3 \\
002891 & 中宠股份 & 81.4 & 32.8 & 34.9 & 1.27 & 4.0 & 7 & 0 \\
000798 & 中水渔业 & 36.7 & 20.8 & 18.9 & 1.02 & 20.9 & 0 & 4 \\
300673 & 佩蒂股份 & 36.2 & 7.3 & 1.0 & 0.24 & 1.4 & 19 & 0 \\
600467 & 好当家 & 34.5 & 8.0 & 13.0 & 0.30 & 0.9 & 0 & 0 \\
300094 & 国联水产 & 33.4 & 14.5 & -12.5 & 0.10 & 10.0 & 1 & 5 \\
600265 & *ST景谷 & 29.3 & 0.6 & -54.6 & -0.00 & -0.0 & 16 & 23 \\
002679 & 福建金森 & 26.3 & 0.5 & -9.5 & -0.14 & -1.8 & 1 & 0 \\
002069 & 獐子岛 & 25.8 & 7.3 & -5.3 & 0.13 & 29.5 & 8 & 5 \\
600257 & 大湖股份 & 25.7 & 4.3 & 2.0 & 0.17 & 2.1 & 2 & 1 \\
600097 & 开创国际 & 21.9 & 11.8 & -5.4 & -0.30 & -1.3 & 0 & 2 \\
000663 & 永安林业 & 19.9 & 1.5 & 14.7 & -0.19 & -2.1 & 2 & 10 \\
\end{longtable}
\endgroup


\subsection{水产养殖：一家公司的亏损抵消了整个子行业}
\label{sec:fishery-aqua}

\num{4} 家水产养殖公司在 2026 年上半年全部实现盈利、合计 \num{0.70} 亿元
（上年同期为 $-5.08$ 亿元）：\textbf{好当家}（\num{0.30} 亿元）、\textbf{大湖股份}（\num{0.17} 亿元）、
\textbf{獐子岛}（\num{0.13} 亿元）与 \textbf{国联水产}（\num{0.10} 亿元）。
其中 \textbf{国联水产}的反转最能说明问题：2025 年它一家亏损 \hl{\num{16.31} 亿元}、
ROE \hl{\num{-160.2}\%}——\emph{亏损额已远超其净资产}，一家公司的亏损抵消了整个子行业；
2026 年上半年虽然转正，但利润仅 \num{0.10} 亿元，
2026 年 6 月末资产负债率仍高达 \num{95.2}\%。
国联水产的问题不在于水产周期，而在于\hl{对虾加工出口业务的
“高库存 + 汇率 + 美国关税”三重冲击}：
其营收结构与第 \ref{sec:other-aqua} 节描述的水产品进口增长、出口承压格局一致
（2025 年水产品贸易逆差扩大至 \num{42.18} 亿美元）。

从周期角度看，水产养殖的正向信号同样明确：2025 年二季度起淡水鱼价格大幅上涨
（草鱼塘口均价同比 +$43\%$）、养殖盈利改善，
且\hl{这一景气已在 2026 年上半年体现为报表利润}——水产养殖 \num{4} 家上市公司
2026 年上半年合计归母净利润 \hl{$+0.70$ 亿元、无一家亏损}（2025 年同期为 $-5.08$ 亿元），
与投苗到上市 \num{12}---\num{18} 个月的时滞一致。
但 \textbf{獐子岛}（半年 ROE \num{29.5}\%，\emph{由极薄的净资产放大}）
与 \textbf{好当家}（\num{0.87}\%）的绝对利润仍很小，
反映出底播养殖类企业的生物资产减值风险尚未完全释放。

\subsection{海洋捕捞：小板块、高 ROE}
\label{sec:fishery-marine}

海洋捕捞是全部农业三级行业中\hl{半年 ROE 中位数最高}的子行业
（两家公司的中位数为 \num{9.79}\%，未年化）：
\textbf{中水渔业}半年 ROE \num{20.9}\%、盈利 \num{1.02} 亿元，
\textbf{开创国际}营收 \num{11.8} 亿元、小幅亏损 \num{0.30} 亿元，
两家合计盈利 \num{0.72} 亿元（上年同期 \num{1.14} 亿元）。
远洋渔业的收益来自配额、船队与\hl{海外渔场资源}，
与国内农产品价格几乎无关，是农业板块中少见的\emph{资源型、非周期}业务。

\subsection{林业：最典型的“壳资源”集群}
\label{sec:fishery-forestry}

\num{4} 家林业公司在 2026 年上半年\hl{全部亏损}（合计 \num{-1.03} 亿元）：
\textbf{平潭发展}（\num{-0.69} 亿元，总市值却达 \num{142.4} 亿元、为本章第二）、
\textbf{永安林业}（\num{-0.19} 亿元，\num{10} 条重组并购公告）、
\textbf{福建金森}（\num{-0.14} 亿元）与
\textbf{*ST景谷}（\num{-0.002} 亿元，\num{23} 条重组并购公告）。
林业的资源属性（林地资产）与其盈利能力长期脱节，
\hl{较高的市值与持续亏损并存，说明市场定价的是林地资产的重估或转型预期，
而不是木材生产的现金流}。\emph{这使林业成为第二部分中融资/重组需求最特殊的板块}：
其潜在融资需求往往与“保壳 + 资产注入”而非产能扩张相关。

\subsection{宠物食品：农业板块中唯一的高成长赛道}
\label{sec:fishery-pet}

\textbf{乖宝宠物}（半年净利润 \num{1.91} 亿元、半年 ROE \num{4.06}\%、营收同比 \num{+10.0}\%）、
\textbf{中宠股份}（\num{1.27} 亿元、\num{4.00}\%、\num{+34.9}\%）、
\textbf{佩蒂股份}（\num{0.24} 亿元、\num{1.36}\%、\num{+1.0}\%）
三家公司在 2026 年上半年全部盈利，但合计利润 \num{3.42} 亿元、
仅为上年同期（\num{6.60} 亿元）的一半左右，
\hl{半年 ROE 中位数 \num{4.00}\%，明显低于 2025 年年报口径的 \num{13.5}\%}。

宠物食品的逻辑与农业周期几乎无关：其需求来自城镇宠物数量与单只消费支出，
上游原料（鸡肉、谷物）价格下跌反而扩大毛利率。
\hl{“农产品价格下行 $\to$ 原料成本下降 $\to$ 毛利率提升”}
构成一条与其他农业子行业完全相反的传导链，
这也是两家公司在 2023---2025 年农业整体下行中保持高增长的原因。
但 2026 年上半年的数据给出了相反的信号：\hl{收入仍在增长（同比中位数 \num{+10.0}\%）、
利润却接近腰斩}，\emph{成长红利正在被竞争与费用侵蚀}，
这是本章最值得关注的边际变化。
其中 \textbf{乖宝宠物}于 2023 年 8 月上市、\textbf{索宝蛋白}于 2023 年 12 月上市，
是全部 \num{104} 家公司中\hl{上市时间最晚的两家}——农业板块的新增融资窗口，
主要留给了这些与周期脱钩的消费型企业。

\subsection{农业综合：农资流通与节水灌溉}
\label{sec:fishery-misc}

申万三级分类中“农业综合Ⅱ”下另有 \num{2} 家公司，
其业务不属于上述任何一类，但恰好代表了农业产业链的两种“服务型”环节：
\textbf{辉隆股份}（半年营收 \num{85.5} 亿元、净利润 \num{1.13} 亿元、半年 ROE \num{3.08}\%）
主营化肥、农药的内外贸分销与自主品牌复合肥生产，
\hl{其收入规模与盈利水平直接反映农资（化肥、农药）价格的波动}——
在化肥价格下行、农民种植收益受挤压的年份，流通环节的毛利被两头压缩；
\textbf{大禹节水}（半年营收 \num{15.9} 亿元、净利润 \num{-0.01} 亿元、半年 ROE \num{-0.04}\%）
从事农业高效节水与农村污水治理，
\hl{其订单来自政府水利项目，具有“逆农业周期”的财政属性}——
2026 年上半年微亏，说明\emph{其业绩节奏更多取决于财政投资而非农业景气本身}。

\begingroup\scriptsize\setlength{\tabcolsep}{2pt}
\begin{longtable}{@{}llrrrrrrr@{}}
\caption{农业上市公司分行业明细（农业综合Ⅱ；市值时点 2026-10-12，总市值口径；财务为 2026 年半年报，ROE 未年化，公告计数窗口 2023-09---2026-10）}\label{tab:comp-misc}\\
\toprule
代码 & 公司简称 & 总市值（亿元） & 营收（亿元） & 同比（\%） & 净利（亿元） & ROE（\%） & 再融资公告 & 重组公告 \\
\midrule
\endfirsthead
\multicolumn{9}{l}{\small（续）}\\
\toprule
代码 & 公司简称 & 总市值（亿元） & 营收（亿元） & 同比（\%） & 净利（亿元） & ROE（\%） & 再融资公告 & 重组公告 \\
\midrule
\endhead
\bottomrule
\endlastfoot
002556 & 辉隆股份 & 48.5 & 85.5 & 3.3 & 1.13 & 3.1 & 0 & 0 \\
300021 & 大禹节水 & 32.7 & 15.9 & 24.3 & -0.01 & -0.0 & 18 & 9 \\
\end{longtable}
\endgroup


\begin{keybox}[渔业、林业与宠物食品的画像结论]
\normalsize
\textbf{① 一个板块，四种周期}：水产养殖看投苗与出口、海洋捕捞看资源与配额、
林业看资产重估、宠物食品看消费升级——\hl{它们唯一的共性是都被归入“农林牧渔”}。\par\vspace{3pt}
\textbf{② 出口加工型的风险与养殖周期不同步}：国联水产 2025 年一家亏损 \num{16.31} 亿元、
ROE \num{-160.2}\%（已超过其净资产），2026 年上半年虽转正（\num{0.10} 亿元），
但 6 月末负债率仍高达 \num{95.2}\%——\hl{加工出口型企业的修复慢于养殖周期}。\par\vspace{3pt}
\textbf{③ 宠物食品的成长红利出现放缓}：三家全部盈利，但合计利润由上年同期的 \num{6.60} 亿元
降至 \num{3.42} 亿元、半年 ROE 中位数 \num{4.00}\%；
\hl{叠加两家新上市公司（乖宝宠物、索宝蛋白），}
说明资本市场的融资功能正从“周期产能”转向“消费升级”，
\emph{但消费型企业的盈利兑现也开始承压}。\par\vspace{3pt}
\textbf{④ 2026 年上半年整体转正。}本组 \num{15} 家公司 2026 年上半年合计归母净利润
\num{4.92} 亿元（上年同期 \num{2.44} 亿元），其中水产养殖 \num{4} 家合计 \num{0.70} 亿元、
\hl{无一家亏损}；改善主要来自水产养殖（上年同期为 \num{-5.08} 亿元），
部分被宠物食品的回落（\num{6.60} $\to$ \num{3.42} 亿元）抵消。
\end{keybox}

\subsection{上市公司画像（\num{15} 家）}
\label{sec:co-fishery}

本节给出本组 \num{15} 家上市公司的画像，按总市值降序排列：渔业、林业与宠物食品 \num{13} 家，
以及三级行业“农业综合Ⅱ”下的 \num{2} 家（农资流通与节水灌溉，
周期来源与农业生产资料价格及财政投资相关，见第 \ref{sec:fishery-misc} 节）。
数据口径与第 \ref{sec:comp-overview} 节一致：\hl{财务数据统一采用 2026 年半年报}
（合并报表口径、未年化，资产负债率为 2026 年 6 月末），
同时列出 2025 年半年报作为同口径对照、2023---2025 年年报作为趋势对照；
股价为月度收盘价，最后一个交易日为 2026 年 10 月 12 日；周期分段使用与第一部分相同的
ZigZag 算法，阈值取 25\%（个股波动大于商品价格，故高于生猪现货的 20\%）；
公告计数窗口为 2023 年 9 月 1 日至 2026 年 10 月 12 日。

与前一版不同，本版画像正文由\hl{逐家检索的公开资料底稿}驱动：底稿只收录有来源链接的事实
（公司沿革、股权与控制权、经营要点、所处行业的基本面、竞争优势、融资历史、
风险事件、战略动作、融资需求线索），
正文按事实生成，各家公司的段落结构与重点并不相同。
每家公司在统一顺序下展开：基本情况与主营结构、股权结构与控制权、经营要点、
所处行业的基本面（供需）、与同类公司比较的竞争优势、历史股价、过去周期分析、
盈利情况、财务分析、重大战略转型、历史融投资情况、未来潜在融投资需求；
\hl{市值一律采用总市值口径}，股权结构取最新披露报告期前三大股东与股东户数，
行业基本面与竞争优势以逐家检索的公开资料为依据，并与三级行业内的同口径指标横向比较。
每家公司配图两张：
股价与行业指数（申万农林牧渔）的对比图，以及近年主要财务指标的组合图
（上排：营业收入与归母净利润、ROE 与资产负债率；下方宽带面板：货币资金与有息负债为左轴柱状、
经营活动现金流净额为其右轴的蓝点折线），并给出关键财务指标表与历史融资/资本运作表。
未来融资需求按两个口径粗略度量：\hl{短期债务缺口 $=$ 期末短期债务 $-$ 货币资金}
与\hl{年化经营缺口 $=$ 半年归母亏损 $\times 2$}；
前者衡量滚动偿付压力，后者衡量维持当期盈利水平的现金消耗，
两者都是压力情形下的静态估算，不构成对融资规模或时点的预测。

% 农业上市公司画像：水产养殖、海洋捕捞、林业Ⅲ、宠物食品（共 13 家，按总市值降序）
% 由 scripts/make_company_profiles.py 自动生成，请勿手工修改

\subsubsection{乖宝宠物（301498）}
\label{co:301498}

\pfl{公司基本情况} 山东聊城的宠物食品企业，全称乖宝宠物食品集团股份有限公司，2006 年 6 月 26 日注册成立，2023 年 8 月 16 日在深圳证券交易所创业板上市；主营犬用和猫用多品类宠物食品的研发、生产和销售，产品包括主粮系列、零食系列与保健品系列，旗下拥有“麦富迪”“弗列加特”等自有品牌，并在泰国、新西兰布局高端产能。控股股东、实际控制人为秦华。2025 年归母净利润 6.73 亿元、同比增长 7.75\%；2026 年上半年营业收入 35.430 亿元、同比增长 10.01\%，但归母净利润 1.91 亿元、同比下降 49.57\%。2026 年 8 月公司推出向不特定对象发行可转债预案，拟募集资金不超过 11 亿元。 公司于 2023-08-16 在深交所创业板首发上市，注册资本 4.00 亿元，总股本 4.00 亿股。 截至 2026 年 10 月 12 日收盘价 39.20 元，对应总市值 156.8 亿元，近一年-58.2\%，总市值在三级行业“宠物食品”的 3 家公司中按由大到小排第\zhnumber{1}位，近一年涨跌幅低于全样本中位数（-10.4\%）47.9 个百分点。 按 2026 年半年报披露的行业构成，收入占比前两位为宠物食品及用品 99.6\%；其他(补充) 0.4\%。

\pfl{股权结构与控制权} 按 2026 年 6 月 30 日的股东名册，前三大股东为秦华 持股 45.72\%；GOLDEN PROSPERITY INVESTME 持股 14.06\%；聊城市海昂企业管理咨询服务中心(有限合伙) 持股 6.93\%，股东户数 26,150 户。 股权结构的另一条线索是，控股股东、实际控制人秦华持有公司股份 1.83 亿股；主要业务为对卢森堡及外国公司进行投资、以收购或认购股权等方式进行并购、从事证券交易及对投资组合公司的持有管理与开发；公司自成立及上市以来在资产、人员、财务、机构和业务等方面独立于控股股东、实际控制人。

\pfl{经营要点} 业务与产能层面，公司拟发行不超过 11.00 亿元可转债。

\pfl{所处行业的基本面} 宠物食品的供给是产能（含海外基地），需求由宠物数量与单宠消费驱动，是农业中与农产品价格周期联动最弱的一块（第 \ref{sec:framework} 章的供需框架）。 当期该行业的关键数据是：2025 年，中国城镇宠物（犬猫）消费市场规模 3,126.00 亿元、同比增长 4.1\%，其中犬 1,606.00 亿元、猫 1,520.00 亿元；2025 年，城镇犬猫数量 1.26 亿只、同比增长 1.8\%；2025 年，中国宠物食品行业零售市场规模 1,084.00 亿元。

\pfl{竞争优势} 把该公司放回同类公司比较，2026 年上半年营业收入 35.4 亿元、在三级行业“宠物食品”的 3 家公司中排第\zhnumber{1}位；毛利率 40.1\%、高于行业内中位数 11.7 个百分点；半年 ROE 4.1\%（未年化）、高于行业内中位数（4.00\%）；资产负债率 31.4\%、低于行业内中位数（35.3\%）。 与同类公司相比，公司的差异化在于：2015 年至 2025 年公司自有品牌“麦富迪”国内市占率从 2.4\% 提升至 6.3\%、实现翻倍；是国内知名的宠物食品上市公司；智能仓储及数智化分拣中心项目完成后预计可达到 2,400 吨/日智能分拣发货能力，以缓解公司现有仓储负荷并提升供应链服务能力。 优势在报表上主要体现为高于行业平均的毛利率，这在周期底部是抵御亏损的主要缓冲。

\pfl{历史股价} 以 2023 年 8 月为起点、2026 年 10 月为终点，股价由 45.24 元变为 39.20 元，累计 -13.4\%，区间最高 109.37 元（2025 年 6 月）、最低 36.30 元（2024 年 1 月）；当前价相当于区间最高价的 35.8\%，为区间最低价的 1.08 倍。 月度收益的年化波动率 43.9\%，高于全样本中位数 37.3\%，在三级行业“宠物食品”的 3 家公司中波动率由高到低排第\zhnumber{1}位。 把股价阶段与公司事件对齐看，2025 年 6 月 至 2026 年 8 月股价下跌-65.0\%，同期（2026 年 8 月）公司同日披露可转债募集资金使用可行性分析报告。 从时间对应关系看，公司层面的资本运作与股权、风险事件解释了区间内多数急涨急跌，与行业指数的同步性相对有限。

\begin{center}
\begin{minipage}{\textwidth}\centering
\begin{tikzpicture}
\begin{axis}[
  width=12.6cm, height=3.9cm, scale only axis,
  xmin=2023.505, xmax=2026.972, ymin=33.76, ymax=117.03,
  xtick={2023,2024,2025,2026},
  yearticks,
  yticklabel style={font=\scriptsize},
  ylabel={元/股}, ylabel style={font=\scriptsize},
  grid=major, grid style={LightGray, line width=0.3pt},
  axis line style={InkGray, line width=0.5pt},
  every axis plot/.append style={line width=0.9pt},
]
\addplot[AgriGreen] table[x=x,y=y]{data/clean/profiles/px_301498.dat};
\addplot[SoftRed, dashed, line width=0.7pt] table[x=x,y=y]{data/clean/profiles/pxzz_301498.dat};
\addplot[only marks, mark=*, mark size=1.1pt, SoftRed] table[x=x,y=y]{data/clean/profiles/pxzz_301498.dat};
\end{axis}
\end{tikzpicture}

\captionof{figure}[乖宝宠物（301498）股价与周期骨架]{乖宝宠物（301498）月度收盘价（元/股）与 ZigZag 周期骨架（阈值 25\%，圆点为周期转折点）}
\end{minipage}
\end{center}

\pfl{过去周期分析} 以 25\% 的 ZigZag 阈值划分，样本区间内股价形成 1 个主升段与 2 个主跌段。 最大主跌段为 2025 年 6 月 至 2026 年 8 月，14 个月-65.0\%；最大主升段出现在 2024 年 1 月 至 2025 年 6 月，17 个月+201.3\%。 最近一次转折出现在 2026 年 8 月（下跌段自 2025 年 6 月起，14 个月-65.0\%），当前处于该段的延续之中。 公司股价较申万农林牧渔指数提前约 28 个月见底，是板块中较早定价周期反转的公司之一。 宠物食品属于消费型农业（第 \ref{sec:financing-strategy} 节归入“向消费端延伸”），收入增长由宠物数量与单宠消费驱动，与农产品价格周期的联动最弱，公司 2026 年上半年实现盈利、半年 ROE 4.1\%（未年化），好于全样本中位数（0.75\%），下行周期对其报表的冲击小于同业。 综合区间形态看，该股单段涨跌幅度偏大、以趋势段为主（最大主升 +201.3\%、最大主跌 -65.0\%），年化波动率 43.9\% 与全样本中位数 37.3\% 相比波动明显大于样本中位数。

\pfl{盈利情况} 2026 年上半年营业收入 35.43 亿元（同比 +10.0\%），归母净利润 1.91 亿元，同比 -49.6\%。 以年报为趋势对照，2023---2025 年营业收入由 43.27 亿元变为 67.69 亿元（两年年均复合增速 +25.1\%），归母净利润由 4.29 亿元变为 6.73 亿元。 按半年报口径，2026 年上半年的 ROE 为 4.1\%（未年化）、销售净利率 5.4\%、毛利率 40.1\%，ROE 在“宠物食品”3 家公司中按数值由高到低排第\zhnumber{1}位。 按同一口径，三级行业“宠物食品”3 家公司的半年 ROE 中位数为 4.00\%，该公司高于其中位数 0.06 个百分点。 从公司披露的原因看，业绩变动主要来自：公司持续扩大国内智能工厂产能；2026 年 4 月，公司加码泰国、新西兰高端产能布局并搭建 WCM 世界级制造体系。 综合看，公司在报告期维持盈利（高于全样本半年 ROE 中位数 0.75\%），利润的可持续性取决于宠物食品景气的延续与成本管控的执行。

\begin{center}\small\setlength{\tabcolsep}{3.4pt}
\captionof{table}[乖宝宠物（301498）关键财务指标]{乖宝宠物（301498）关键财务指标（合并报表口径；两个半年报期均未年化；现金短债比 $=$ 货币资金 $\div$ 短期债务）}
\begin{adjustbox}{max width=\textwidth}
\begin{tabular}{@{}lrrrrr@{}}
\toprule
指标 & 2023 年 & 2024 年 & 2025 年 & 2025 年半年报 & 2026 年半年报 \\
\midrule
营业收入（亿元） & 43.27 & 52.45 & 67.69 & 32.21 & 35.43 \\
归母净利润（亿元） & 4.29 & 6.25 & 6.73 & 3.78 & 1.91 \\
毛利率（\%） & 36.8 & 42.3 & 42.5 & 42.8 & 40.1 \\
ROE（\%） & 17.4 & 15.8 & 15.3 & 8.7 & 4.1 \\
资产负债率（\%） & 11.4 & 17.1 & 20.5 & 14.9 & 31.4 \\
经营活动现金流净额（亿元） & 6.17 & 7.20 & 8.42 & 3.50 & 1.51 \\
货币资金（亿元） & 10.09 & 6.95 & 8.19 & 6.33 & 9.85 \\
有息负债合计（亿元） & 0.25 & 1.86 & 3.70 & 1.48 & 11.94 \\
现金短债比（倍） & 216.75 & 5.71 & 2.61 & 7.68 & 1.06 \\
\bottomrule
\end{tabular}
\end{adjustbox}
\end{center}

\begin{center}
\begin{minipage}{\textwidth}\centering
\begin{tikzpicture}
\begin{groupplot}[
  group style={group size=2 by 1, horizontal sep=0.60cm},
  width=6.85cm, height=4.60cm,
  tick label style={font=\tiny},
  title style={font=\scriptsize\heiti},
  yticklabel style={font=\tiny},
  ylabel style={font=\tiny},
  grid=major, grid style={LightGray, line width=0.3pt},
  axis line style={InkGray, line width=0.5pt},
  enlarge x limits={abs=0.8},
  legend style={font=\scriptsize, at={(0.5,-0.20)}, anchor=north,
                legend columns=2, draw=none, fill=none, column sep=6pt},
]
\nextgroupplot[
  ybar, bar width=4pt,
  ymin=-6.77, ymax=75.81,
  xtick={0,1,2,3,4,5.7},
  xticklabels={2021,2022,2023,2024,2025,2026H1},
  ylabel={亿元},
]
\addplot[fill=AgriGreen!75, draw=AgriGreen, bar shift=-2.6pt] table[x=i, y=rev]{data/clean/profiles/fin_301498.dat};
\addplot[fill=SoftRed!60, draw=SoftRed, bar shift=2.6pt] table[x=i, y=ni]{data/clean/profiles/fin_301498.dat};
\addplot[fill=AgriGreen!30, draw=AgriGreen, pattern=north east lines, bar shift=-2.6pt] table[x=x, y=rev]{data/clean/profiles/finh_301498.dat};
\addplot[fill=SoftRed!20, draw=SoftRed, pattern=north east lines, bar shift=2.6pt] table[x=x, y=ni]{data/clean/profiles/finh_301498.dat};
\legend{营业收入（年/半年）, 归母净利润（年/半年）}
\nextgroupplot[
  ymin=-2.9, ymax=39.5,
  xtick={0,1,2,3,4,5.7},
  xticklabels={2021,2022,2023,2024,2025,2026H1},
  ylabel={\%},
]
\addplot[AgriGold, mark=*, mark size=1.3pt] table[x=i, y=roe]{data/clean/profiles/fin_301498.dat};
\addplot[OilBlue, mark=square*, mark size=1.3pt] table[x=i, y=debt]{data/clean/profiles/fin_301498.dat};
\addplot[only marks, AgriGold, mark=*, mark size=2.0pt] table[x=x, y=roe]{data/clean/profiles/finh_301498.dat};
\addplot[only marks, OilBlue, mark=square*, mark size=2.0pt] table[x=x, y=debt]{data/clean/profiles/finh_301498.dat};
\legend{ROE, 资产负债率}
\end{groupplot}
\end{tikzpicture}
\begin{tikzpicture}
\begin{axis}[
  name=finbot,
  width=13.75cm, height=3.10cm, scale only axis,
  ymin=-1.19, ymax=13.38,
  xtick={0,1,2,3,4,5.7},
  xticklabels={2021,2022,2023,2024,2025,2026H1},
  tick label style={font=\tiny},
  yticklabel style={font=\tiny},
  ylabel={亿元}, ylabel style={font=\tiny},
  ybar, bar width=4pt,
  grid=major, grid style={LightGray, line width=0.3pt},
  axis line style={InkGray, line width=0.5pt},
  enlarge x limits={abs=0.8},
  legend style={font=\scriptsize, at={(0.5,-0.26)}, anchor=north,
                legend columns=2, draw=none, fill=none, column sep=10pt},
]
\addplot[fill=AgriGreen!55, draw=AgriGreen, bar shift=-3.0pt] table[x=i, y=cash]{data/clean/profiles/fin_301498.dat};
\addplot[fill=SoftRed!45, draw=SoftRed, bar shift=3.0pt] table[x=i, y=ibd]{data/clean/profiles/fin_301498.dat};
\addplot[fill=AgriGreen!25, draw=AgriGreen, pattern=north east lines, bar shift=-3.0pt] table[x=x, y=cash]{data/clean/profiles/finh_301498.dat};
\addplot[fill=SoftRed!20, draw=SoftRed, pattern=north east lines, bar shift=3.0pt] table[x=x, y=ibd]{data/clean/profiles/finh_301498.dat};
\legend{货币资金（左轴）, 有息负债（左轴）}
\end{axis}
\begin{axis}[
  at={(finbot.south east)}, anchor=south east,
  width=13.75cm, height=3.10cm, scale only axis,
  axis y line*=right, axis x line*=none, xtick=\empty,
  ymin=-2.19, ymax=10.95,
  tick label style={font=\tiny},
  yticklabel style={font=\tiny}, scaled y ticks=false,
  legend style={font=\scriptsize, at={(0.5,-0.44)}, anchor=north,
                legend columns=1, draw=none, fill=none},
]
\addplot[OilBlue, mark=*, mark size=2.2pt, line width=1.3pt] table[x=i, y=ocf]{data/clean/profiles/fin_301498.dat};
\addplot[only marks, OilBlue, mark=*, mark size=3.2pt] table[x=x, y=ocf]{data/clean/profiles/finh_301498.dat};
\legend{经营活动现金流净额（右轴，亿元，蓝点）}
\end{axis}
\end{tikzpicture}

\captionof{figure}[乖宝宠物（301498）近年主要财务指标]{乖宝宠物（301498）近年主要财务指标（左上：营业收入与归母净利润；右上：ROE 与资产负债率；下：货币资金、有息负债为左轴柱状，经营活动现金流净额为其右轴的蓝点折线。
2021---2025 年为年报口径，2026H1 为 2026 年半年报/半年末，与其前一年报之间留出间隔、以斜纹或空心点区分；半年报数据未年化）}
\end{minipage}
\end{center}

\pfl{财务分析} 2026 年 6 月末资产负债率 31.4\%，较 2025 年末上升 11.0 个百分点（样本中位数 52.2\%，所处三级行业内按由低到高排第\zhnumber{1}位）。 2026 年 6 月末货币资金 9.85 亿元、短期债务（短期借款与一年内到期非流动负债）9.29 亿元、长期债务 2.65 亿元，有息负债合计 11.94 亿元，现金短债比 1.06 倍——覆盖充分。 净有息负债 2.10 亿元，相当于其 2026 年上半年营业收入的 6\%。 流动比率 2.13、速动比率 1.59，短期偿债能力处于正常区间。 2026 年上半年经营活动现金流净额 1.51 亿元，经营现金流与利润同向为正，内生资金可以支撑运营。 2026 年 6 月末在建工程 3.78 亿元，较上年末增加 3.19 亿元，仍有在建投入。 综合看，现金短债比 1.06 倍，短期偿付压力可控；资产负债率低于样本中位数 20.8 个百分点；有息负债中短期债务占比更高，期限结构仍以滚动续作为主。

\pfl{重大战略转型} 从公开披露看，公司的战略推进集中在：其他动作（1 项）：2026 年 4 月，核心战略为自有品牌市占率优先。 公告口径上，近三年（2023 年 9 月---2026 年 9 月）公司披露重大重组与并购类公告 0 条、对外投资与转型类公告 2 条、回购与股权激励类公告 61 条，合计公告 432 条。 第一大行业仍是“宠物食品及用品”，收入占比 99.6\%，与 2021 年年报相比没有方向性变化。 综合判断，报告期内公司有 2 条与战略相关的公告落地（重大重组与并购 0 条、对外投资与转型 2 条），资本运作与主业调整之间存在直接关联。

\begin{center}\small\setlength{\tabcolsep}{4pt}
\captionof{table}[乖宝宠物（301498）历史融资与资本运作]{乖宝宠物（301498）历史融资与资本运作（据公司公告与公开报道整理；金额为披露值，含首次公开发行、再融资、债券、银行借款与重整投资等）}
\begin{adjustbox}{max width=\textwidth}
\begin{tabular}{@{}p{1.45cm}p{2.55cm}p{4.05cm}p{4.95cm}@{}}
\toprule
时间 & 方式 & 规模 & 用途与说明 \\
\midrule
2023-08 & 首次公开发行A股股票并在深圳证券交易所创业板上市（IPO） & 发行价 39.99 元/\allowbreak{}股；首发前总股本 36,\allowbreak{}004.00 万股 & 投向公司首次公开发行募集资金投资项目；超募资金部分用于投资建设新项目及永久补充流动资金。司与中泰证券股份有限公司；发行市盈率 61.28 倍；招股公告日 2023 年 7 月 27 日 \\
2023-09-19 & 超募资金使用（投资建设新项目及永久补充流动资金） & 拟使用部分超募资金投资建设新项目并用部分超募资金永久补充流动资金 & 新增募投项目及补充流动资金。核查意见 \\
2024-04-16 & 超募资金使用（永久补充流动资金） & 使用部分超募资金永久补充流动资金 & 永久补充流动资金。核查意见 \\
2025-08-10 & 股权激励（第二期限制性股票激励计划） & 公司董事会审议通过第二期限制性股票激励计划（草案）及其摘要等议案 & 用于核心人员激励。；2025 年 4 月 19 日、\allowbreak{}7 月 7 日曾两次调整第一期限制性股票激励计划授予价格并向激励对象授予预留部分限制性股票 \\
2026-04-23 & 募集资金使用进展（IPO募投项目） & 截至 2025 年 12 月 31 日公司累计使用募集资金 13.49 亿元 & 用于首发募投项目及超募资金安排项目。；公司同时披露 2025 年度募集资金存放与使用情况的专项报告及中泰证券核查意见；2026 年 4 月 23 日另公告使用部分闲置募集资金进行现金管理 \\
2026-08-27 & 拟向不特定对象发行可转换公司债券 & 募集资金总额不超过人民币 11.00 亿元，\allowbreak{}尚需履行相应审议与注册程序 & 年产 30 万吨高端宠物主粮食品项目（拟投入 4.50 亿元）、\allowbreak{}智能仓储及数智化分拣中心项目（拟投入 4.00 亿元）、\allowbreak{}数智化运营与决策平台建设项目（拟投入 1.00 亿元）、\allowbreak{}补充流动资金（拟投入 1.50 亿元）。；公司同日披露前次募集资金使用情况报告与德勤华永会计师事务所出具的前次募集资金使用情况审核报告 \\
\bottomrule
\end{tabular}
\end{adjustbox}
\end{center}

\pfl{历史融投资情况} 从现金流量表看，2025 年公司取得借款收到的现金 4.20 亿元、偿还债务支付的现金 2.29 亿元、吸收权益性投资 0.11 亿元，筹资活动现金流净额 -0.16 亿元（2023 年为取得借款 0.98 亿元、偿还债务 2.94 亿元）。 2026 年上半年筹资活动现金流净额为 5.06 亿元（取得借款 9.22 亿元、偿还债务 3.01 亿元，未年化），仍处于净融入状态。 2025 年分配股利、利润或偿付利息支付的现金合计 2.02 亿元。 公开披露的融资记录共 6 笔，工具类型以 IPO、可转债为主；金额与用途见下表。 其中规模最大的两笔为：2026 年 8 月，可转债，募集资金总额不超过人民币 11.00 亿元，尚需履行相应审议与注册程序；2026 年 4 月，募集资金使用进展，截至 2025 年 12 月 31 日公司累计使用募集资金 13.49 亿元。 股本方面，近三年披露股本变动 5 次（限售股份上市 2 次、A 股上市 1 次）；总股本自 2021 年末以来未发生变化（4.00 亿股）。 分红方面，近三年现金分红 7 次、合计每股派现 1.42 元，最近一次实施在 2026 年 9 月。 近三年“再融资”类公告 7 条，涉及可转债、募集资金使用。 综合看，公司融资结构上以债务滚动为主、股权融资为补充；2025 年筹资活动现金流净额 -0.16 亿元，融资节奏仍以维持存量债务周转为主，与前述经营与投资现金流的方向基本一致。

\pfl{未来潜在融资需求分析} 2026 年 6 月末货币资金 9.85 亿元，而短期债务 9.29 亿元（现金短债比 1.06 倍），2026 年上半年经营现金流净额 1.51 亿元。账面货币资金可以覆盖短期债务，缺口不体现在静态偿付层面。 公司 2026 年 6 月末资产负债率 31.4\%、2026 年上半年 ROE 4.1\%（未年化），资产负债表尚有余量，融资需求不是“补血型”的刚性需求，而是配套在建项目投入的主动性安排。 公开披露的后续资金安排线索包括：股权融资线索：2026 年 8 月，公司同日披露可转债募集资金使用可行性分析报告；债务与授信安排：2026 年 6 月，公司未公告银行综合授信额度、公司债或中票发行计划；2026 年 8 月，公司已披露向不特定对象发行可转换公司债券预案。 资金需求还要叠加在建工程：期末在建工程 3.78 亿元、2026 年上半年购建固定资产等支出 4.40 亿元（未年化），这部分支出在周期底部通常会被主动放缓。 以现有货币资金与经营现金流衡量，公司不存在刚性补血的紧迫性，融资动作更多是主动型的。




\subsubsection{平潭发展（000592）}
\label{co:000592}

\pfl{公司基本情况} 平潭发展成立于 1993 年，注册地为福建平潭综合实验区的林业与木材加工企业，1996 年 3 月在深圳证券交易所上市；2007 至 2008 年经股权分置改革与重大资产重组，福建山田实业发展有限公司取得控制权并注入林业资产，现由自然人刘平山实际控制。根据《2024 年下半年上市公司行业分类结果》，公司所处行业为木材加工和木、竹、藤、棕、草制品业，主营中纤板生产销售、林木资产的抚育与采伐销售及化肥贸易等业务。2025 年归母净利润 -1.13 亿元，2026 年上半年归母净利润 -0.69 亿元，连续亏损；控股子公司中福海峡（平潭）置业有限公司已进入破产清算程序。 公司于 1996-03-27 在深交所主板首发上市，注册资本 19.32 亿元，总股本 19.32 亿股。 截至 2026 年 10 月 12 日收盘价 7.37 元，对应总市值 142.4 亿元，近一年+104.7\%，总市值在三级行业“林业Ⅲ”的 4 家公司中按由大到小排第\zhnumber{1}位，近一年涨跌幅高于全样本中位数（-10.4\%）115.1 个百分点。 按 2026 年半年报披露的行业构成，收入占比前三位为纤维板 74.9\%；房地产 10.9\%；林业 5.1\%。 股权结构方面，控股股东为福建山田实业发展有限公司。

\pfl{股权结构与控制权} 按 2026 年 6 月 30 日的股东名册，前三大股东为福建山田实业发展有限公司 持股 17.13\%；香港中央结算有限公司 持股 3.38\%；阿布达比投资局－自有资金 持股 0.64\%，股东户数 266,430 户。 股东结构与控制权方面，占总股本 17.13\%，其中质押 9,650.00 万股；报告期内控股股东未发生变更；实际控制人刘平山同时担任公司董事长，并直接持有公司 54.5 万股（占 0.20\%）。

\pfl{经营要点} 从经营要点看，2026 年 6 月，法院裁定受理公司作为债权人申请破产清算的中福置业事项。

\pfl{所处行业的基本面} 林业的供给由林木蓄积与采伐限额决定、需求来自木材加工与基建，变现周期长（第 \ref{sec:framework} 章的供需框架）。 当期该行业的关键数据是：2026 年，1—2 月中国进口原木及锯材 750.3 万立方米、同比下降 11.2\%；2026 年 9 月，下游纸厂订单改善有限、采购以刚需为主。

\pfl{竞争优势} 把该公司放回同类公司比较，2026 年上半年营业收入 4.8 亿元、在三级行业“林业Ⅲ”的 4 家公司中排第\zhnumber{1}位；毛利率 7.9\%、低于行业内中位数 14.5 个百分点；半年 ROE -4.0\%（未年化）、低于行业内中位数（-1.96\%）；资产负债率 39.3\%、高于行业内中位数（38.0\%）。 相对优势方面，年产能超 50 万立方米、年加工饰面板材超 100 万片；公司拥有领先业内水平的 15 万立方米连续压机超薄板生产线以及迪芬巴赫连续压机中纤板生产线。 就当期报表看，公司尚未把经营层面的特征转化为高于行业的资本回报，周期底部的相对优势更多体现在现金流与资产负债表上。

\pfl{历史股价} 以 2021 年 1 月为起点、2026 年 10 月为终点，股价由 2.31 元变为 7.37 元，累计 +219.0\%，区间最高 15.40 元（2025 年 12 月）、最低 1.30 元（2024 年 6 月）；当前价相当于区间最高价的 47.9\%，为区间最低价的 5.67 倍。 月度收益的年化波动率 58.6\%，高于全样本中位数 37.3\%，在三级行业“林业Ⅲ”的 4 家公司中波动率由高到低排第\zhnumber{1}位。

\begin{center}
\begin{minipage}{\textwidth}\centering
\begin{tikzpicture}
\begin{axis}[
  width=12.6cm, height=3.9cm, scale only axis,
  xmin=2020.922, xmax=2026.972, ymin=1.21, ymax=16.48,
  xtick={2021,2022,2023,2024,2025,2026},
  yearticks,
  yticklabel style={font=\scriptsize},
  ylabel={元/股}, ylabel style={font=\scriptsize},
  grid=major, grid style={LightGray, line width=0.3pt},
  axis line style={InkGray, line width=0.5pt},
  every axis plot/.append style={line width=0.9pt},
]
\addplot[AgriGreen] table[x=x,y=y]{data/clean/profiles/px_000592.dat};
\addplot[SoftRed, dashed, line width=0.7pt] table[x=x,y=y]{data/clean/profiles/pxzz_000592.dat};
\addplot[only marks, mark=*, mark size=1.1pt, SoftRed] table[x=x,y=y]{data/clean/profiles/pxzz_000592.dat};
\end{axis}
\end{tikzpicture}

\captionof{figure}[平潭发展（000592）股价与周期骨架]{平潭发展（000592）月度收盘价（元/股）与 ZigZag 周期骨架（阈值 25\%，圆点为周期转折点）}
\end{minipage}
\end{center}

\pfl{过去周期分析} 以 25\% 的 ZigZag 阈值划分，样本区间内股价形成 3 个主升段与 2 个主跌段。 最大主升段出现在 2024 年 6 月 至 2025 年 12 月，18 个月+1084.6\%；最大主跌段为 2021 年 5 月 至 2024 年 6 月，37 个月-67.1\%。 最近一次转折出现在 2026 年 10 月（上涨段自 2026 年 7 月起，3 个月+26.4\%），当前处于该段之后的震荡之中。 公司股价较申万农林牧渔指数提前约 23 个月见底，是板块中较早定价周期反转的公司之一。 林业的现金流依赖林木采伐与政策补贴，周期长、变现慢，公司 2026 年 6 月末资产负债率 39.3\%、半年 ROE -4.0\%（未年化），在本轮下行中属于随行业同步调整的一类。 综合区间形态看，该股单段涨跌幅度偏大、以趋势段为主（最大主升 +1084.6\%、最大主跌 -67.1\%），年化波动率 58.6\% 与全样本中位数 37.3\% 相比波动明显大于样本中位数。

\pfl{盈利情况} 2026 年上半年营业收入 4.82 亿元（同比 -35.1\%），归母净利润 -0.69 亿元，由上年同期的盈利转为亏损。 以年报为趋势对照，2023---2025 年营业收入由 12.33 亿元变为 13.20 亿元（两年年均复合增速 +3.5\%），归母净利润由 -3.08 亿元变为 -1.13 亿元。 按半年报口径，2026 年上半年的 ROE 为 -4.0\%（未年化）、销售净利率 -33.7\%、毛利率 7.9\%，ROE 在“林业Ⅲ”4 家公司中按数值由高到低排第\zhnumber{4}位。 同期全样本半年 ROE 中位数 0.75\%，公司与之相差 4.74 个百分点（弱于中位数公司）。 公司披露的应对举措集中在：2026 年 3 月，公司通过董事会与临时股东会决议，以债权人身份推动控股子公司中福置业破产清算。 综合看，公司仍处亏损区间、由盈转亏，半年 ROE 在三级行业内按由高到低排第\zhnumber{4}位，单位盈利能力的修复依赖行业景气与自身成本端的改善，报表弹性主要体现为价格与销量的方向性变化。

\begin{center}\small\setlength{\tabcolsep}{3.4pt}
\captionof{table}[平潭发展（000592）关键财务指标]{平潭发展（000592）关键财务指标（合并报表口径；两个半年报期均未年化；现金短债比 $=$ 货币资金 $\div$ 短期债务）}
\begin{adjustbox}{max width=\textwidth}
\begin{tabular}{@{}lrrrrr@{}}
\toprule
指标 & 2023 年 & 2024 年 & 2025 年 & 2025 年半年报 & 2026 年半年报 \\
\midrule
营业收入（亿元） & 12.33 & 15.63 & 13.20 & 7.43 & 4.82 \\
归母净利润（亿元） & -3.08 & -1.17 & -1.13 & 0.15 & -0.69 \\
毛利率（\%） & 6.7 & 11.2 & 6.9 & 7.5 & 7.9 \\
ROE（\%） & -13.9 & -5.9 & -6.2 & 0.8 & -4.0 \\
资产负债率（\%） & 48.0 & 46.9 & 45.0 & 41.3 & 39.3 \\
经营活动现金流净额（亿元） & 1.25 & -1.03 & 2.23 & 0.84 & 0.36 \\
货币资金（亿元） & 3.57 & 4.29 & 3.20 & 2.90 & 4.17 \\
有息负债合计（亿元） & 4.17 & 5.22 & 5.95 & 4.96 & 5.30 \\
现金短债比（倍） & 0.92 & 0.83 & 0.54 & 0.59 & 0.79 \\
\bottomrule
\end{tabular}
\end{adjustbox}
\end{center}

\begin{center}
\begin{minipage}{\textwidth}\centering
\begin{tikzpicture}
\begin{groupplot}[
  group style={group size=2 by 1, horizontal sep=0.60cm},
  width=6.85cm, height=4.60cm,
  tick label style={font=\tiny},
  title style={font=\scriptsize\heiti},
  yticklabel style={font=\tiny},
  ylabel style={font=\tiny},
  grid=major, grid style={LightGray, line width=0.3pt},
  axis line style={InkGray, line width=0.5pt},
  enlarge x limits={abs=0.8},
  legend style={font=\scriptsize, at={(0.5,-0.20)}, anchor=north,
                legend columns=2, draw=none, fill=none, column sep=6pt},
]
\nextgroupplot[
  ybar, bar width=4pt,
  ymin=-5.87, ymax=18.46,
  xtick={0,1,2,3,4,5.7},
  xticklabels={2021,2022,2023,2024,2025,2026H1},
  ylabel={亿元},
]
\addplot[fill=AgriGreen!75, draw=AgriGreen, bar shift=-2.6pt] table[x=i, y=rev]{data/clean/profiles/fin_000592.dat};
\addplot[fill=SoftRed!60, draw=SoftRed, bar shift=2.6pt] table[x=i, y=ni]{data/clean/profiles/fin_000592.dat};
\addplot[fill=AgriGreen!30, draw=AgriGreen, pattern=north east lines, bar shift=-2.6pt] table[x=x, y=rev]{data/clean/profiles/finh_000592.dat};
\addplot[fill=SoftRed!20, draw=SoftRed, pattern=north east lines, bar shift=2.6pt] table[x=x, y=ni]{data/clean/profiles/finh_000592.dat};
\legend{营业收入（年/半年）, 归母净利润（年/半年）}
\nextgroupplot[
  ymin=-18.9, ymax=54.2,
  xtick={0,1,2,3,4,5.7},
  xticklabels={2021,2022,2023,2024,2025,2026H1},
  ylabel={\%},
]
\addplot[AgriGold, mark=*, mark size=1.3pt] table[x=i, y=roe]{data/clean/profiles/fin_000592.dat};
\addplot[OilBlue, mark=square*, mark size=1.3pt] table[x=i, y=debt]{data/clean/profiles/fin_000592.dat};
\addplot[only marks, AgriGold, mark=*, mark size=2.0pt] table[x=x, y=roe]{data/clean/profiles/finh_000592.dat};
\addplot[only marks, OilBlue, mark=square*, mark size=2.0pt] table[x=x, y=debt]{data/clean/profiles/finh_000592.dat};
\legend{ROE, 资产负债率}
\end{groupplot}
\end{tikzpicture}
\begin{tikzpicture}
\begin{axis}[
  name=finbot,
  width=13.75cm, height=3.10cm, scale only axis,
  ymin=-0.60, ymax=6.67,
  xtick={0,1,2,3,4,5.7},
  xticklabels={2021,2022,2023,2024,2025,2026H1},
  tick label style={font=\tiny},
  yticklabel style={font=\tiny},
  ylabel={亿元}, ylabel style={font=\tiny},
  ybar, bar width=4pt,
  grid=major, grid style={LightGray, line width=0.3pt},
  axis line style={InkGray, line width=0.5pt},
  enlarge x limits={abs=0.8},
  legend style={font=\scriptsize, at={(0.5,-0.26)}, anchor=north,
                legend columns=2, draw=none, fill=none, column sep=10pt},
]
\addplot[fill=AgriGreen!55, draw=AgriGreen, bar shift=-3.0pt] table[x=i, y=cash]{data/clean/profiles/fin_000592.dat};
\addplot[fill=SoftRed!45, draw=SoftRed, bar shift=3.0pt] table[x=i, y=ibd]{data/clean/profiles/fin_000592.dat};
\addplot[fill=AgriGreen!25, draw=AgriGreen, pattern=north east lines, bar shift=-3.0pt] table[x=x, y=cash]{data/clean/profiles/finh_000592.dat};
\addplot[fill=SoftRed!20, draw=SoftRed, pattern=north east lines, bar shift=3.0pt] table[x=x, y=ibd]{data/clean/profiles/finh_000592.dat};
\legend{货币资金（左轴）, 有息负债（左轴）}
\end{axis}
\begin{axis}[
  at={(finbot.south east)}, anchor=south east,
  width=13.75cm, height=3.10cm, scale only axis,
  axis y line*=right, axis x line*=none, xtick=\empty,
  ymin=-1.87, ymax=3.21,
  tick label style={font=\tiny},
  yticklabel style={font=\tiny}, scaled y ticks=false,
  legend style={font=\scriptsize, at={(0.5,-0.44)}, anchor=north,
                legend columns=1, draw=none, fill=none},
]
\addplot[OilBlue, mark=*, mark size=2.2pt, line width=1.3pt] table[x=i, y=ocf]{data/clean/profiles/fin_000592.dat};
\addplot[only marks, OilBlue, mark=*, mark size=3.2pt] table[x=x, y=ocf]{data/clean/profiles/finh_000592.dat};
\legend{经营活动现金流净额（右轴，亿元，蓝点）}
\end{axis}
\end{tikzpicture}

\captionof{figure}[平潭发展（000592）近年主要财务指标]{平潭发展（000592）近年主要财务指标（左上：营业收入与归母净利润；右上：ROE 与资产负债率；下：货币资金、有息负债为左轴柱状，经营活动现金流净额为其右轴的蓝点折线。
2021---2025 年为年报口径，2026H1 为 2026 年半年报/半年末，与其前一年报之间留出间隔、以斜纹或空心点区分；半年报数据未年化）}
\end{minipage}
\end{center}

\pfl{财务分析} 2026 年 6 月末资产负债率 39.3\%，较 2025 年末下降 5.6 个百分点（样本中位数 52.2\%，所处三级行业内按由低到高排第\zhnumber{3}位）。 2026 年 6 月末货币资金 4.17 亿元、短期债务（短期借款与一年内到期非流动负债）5.27 亿元、长期债务 0.04 亿元，有息负债合计 5.30 亿元，现金短债比 0.79 倍——覆盖偏紧。 净有息负债 1.14 亿元，相当于其 2026 年上半年营业收入的 24\%。 流动比率 2.04、速动比率 1.18，短期指标尚可。 2026 年上半年经营活动现金流净额 0.36 亿元，净利润为负而经营现金流为正，差额主要来自折旧摊销与营运资本变动，说明亏损尚未完全转化为现金流出。 综合看，现金短债比 0.79 倍，短期资金安排偏紧；资产负债率低于样本中位数 12.9 个百分点；有息负债中短期债务占比更高，期限结构仍以滚动续作为主。 公司披露的重整与债务违约风险为：2023 年 12 月，公司已于 2023 年 12 月、2024 年 1 月召开董事会和股东会审议通过作为债权人向法院申请对中福建材城进行破产预重整事项。

\pfl{重大战略转型} 从公开披露看，公司的战略推进集中在：其他动作（1 项）：2025 年 11 月，公司出资 1.50 亿元认购宜昌东数一号投资有限责任公司增资。 公告口径上，近三年（2023 年 9 月---2026 年 9 月）公司披露重大重组与并购类公告 3 条、对外投资与转型类公告 0 条、回购与股权激励类公告 17 条，合计公告 329 条。 第一大行业“纤维板”的收入占比由 2021 年年报的 46.3\% 变为 74.9\%（28.6 个百分点），收入结构出现再平衡。 综合判断，公司当前的战略重心不是扩张而是化解债务与完成重整，业务层面的调整让位于司法重整程序。

\begin{center}\small\setlength{\tabcolsep}{4pt}
\captionof{table}[平潭发展（000592）历史融资与资本运作]{平潭发展（000592）历史融资与资本运作（据公司公告与公开报道整理；金额为披露值，含首次公开发行、再融资、债券、银行借款与重整投资等）}
\begin{adjustbox}{max width=\textwidth}
\begin{tabular}{@{}p{1.45cm}p{2.55cm}p{4.05cm}p{4.95cm}@{}}
\toprule
时间 & 方式 & 规模 & 用途与说明 \\
\midrule
1996-03 & 首次公开发行A股股票并在深圳证券交易所上市（IPO） & 发行价 4.33 元/\allowbreak{}股；首发前总股本 11,\allowbreak{}540.00 万股 & 招股公告日 1996 年 3 月 7 日 \\
2010 & 向特定对象非公开发行A股股票（定增） & 发行 7,\allowbreak{}280.00 万股，\allowbreak{}发行后公司股本变更为 6.52 亿元 & 2011 年度股东大会决议以资本公积转增股本 1.96 亿元 \\
2015-12 & 向特定对象非公开发行A股股票（定增） & 发行 1.18 亿股，\allowbreak{}发行价 16.88 元/\allowbreak{}股；募集资金总额 20.00 亿元 & 新增股份于 2015 年 12 月 18 日在深圳证券交易所上市，\allowbreak{}以资本公积金向全体股东每 10 股转增 10 股 \\
2016-03 & 控股股东以所持公司股票为标的非公开发行可交换公司债券（股东层面融资） & 控股股东福建山田实业发展有限公司质押公司 5,\allowbreak{}800 万股作为担保 & 该事项已获深圳证券交易所深证函【2016】106 号确认 \\
\bottomrule
\end{tabular}
\end{adjustbox}
\end{center}

\pfl{历史融投资情况} 从现金流量表看，2025 年公司取得借款收到的现金 0.26 亿元、偿还债务支付的现金 0.04 亿元、吸收权益性投资 0.00 亿元，筹资活动现金流净额 0.05 亿元。 2026 年上半年筹资活动现金流净额为 -0.01 亿元（取得借款 — 亿元、偿还债务 — 亿元，未年化），已转为净流出，处于净还债状态。 2025 年分配股利、利润或偿付利息支付的现金合计 0.00 亿元。 公开披露的融资记录共 4 笔，工具类型以 IPO、定向增发为主；金额与用途见下表。 其中规模最大的两笔为：1996 年 3 月，IPO，发行价 4.33 元/股，首发前总股本 11,540.00 万股。 股本方面，近三年披露股本变动 4 次（股权分置受限股份上市 1 次）；总股本自 2021 年末以来未发生变化（19.32 亿股）。 分红方面，近三年未实施现金分红，与重整期间的资金约束一致。 近三年“再融资”类公告 1 条，涉及发行股份购买资产。 综合看，公司融资结构上以债务滚动为主、股权融资为补充；2025 年筹资活动现金流净额 0.05 亿元，年度内仍为净融入，与前述经营与投资现金流的方向基本一致。

\pfl{未来潜在融资需求分析} 2026 年 6 月末货币资金 4.17 亿元，而短期债务 5.27 亿元（现金短债比 0.79 倍），2026 年上半年经营现金流净额 0.36 亿元。按“短期债务减去账面货币资金”的滚动偿付口径，静态缺口约 1.10 亿元；若再计入上半年亏损的年化额（1.39 亿元），维持一年正常经营所需的外部资金约 1.10---2.49 亿元。 公司 2026 年上半年归母净利润为负（-0.69 亿元），但 6 月末资产负债率 39.3\% 尚未进入第一档（亏损且负债率 $>$65\%）的区间，当前融资需求以补充流动资金、支撑低谷期经营性支出为主。 公开披露的后续资金安排线索包括：股权融资线索：2026 年 9 月，公司未公告定向增发预案、配股、可转债或公司债等再融资计划。 综合看，公司的外部资金需求以营运资金周转与在建项目投入为主，滚动偿付口径下的静态缺口约 1.10 亿元，融资的时点与规模更多取决于债权人与重整投资人的进度。




\subsubsection{中宠股份（002891）}
\label{co:002891}

\pfl{公司基本情况} 山东烟台的宠物食品企业，全称烟台中宠食品股份有限公司，2014 年 11 月 24 日由烟台中宠食品有限公司整体变更设立，2017 年 8 月 21 日在深圳证券交易所上市；主营宠物食品和用品的研发、生产与销售，采取自主品牌与 OEM/ODM 双轨并行模式，自有品牌包括“Wanpy 顽皮”“TOPTREES 领先”“新西兰 ZEAL 真致”等。2025 年归母净利润 3.65 亿元、同比下降 7.28\%；2026 年上半年营收 32.806 亿元、同比增长 34.88\%，但归母净利润 1.27 亿元、同比下降 37.63\%。公司近年连续通过可转债与非公开发行融资，2026 年 3 月完成可转债提前赎回。 公司于 2017-08-21 在深交所主板首发上市，注册资本 3.22 亿元，总股本 2.94 亿股。 截至 2026 年 10 月 12 日收盘价 27.66 元，对应总市值 81.4 亿元，近一年-49.7\%，总市值在三级行业“宠物食品”的 3 家公司中按由大到小排第\zhnumber{2}位，近一年涨跌幅低于全样本中位数（-10.4\%）39.4 个百分点。 按 2026 年半年报披露的行业构成，收入占比前两位为宠物食品及用品 93.9\%；其他(补充) 6.1\%。

\pfl{股权结构与控制权} 按 2026 年 6 月 30 日的股东名册，前三大股东为烟台中幸生物科技有限公司 持股 22.72\%；烟台和正投资中心(有限合伙) 持股 15.62\%；日本伊藤株式会社 持股 9.32\%，股东户数 45,389 户。 从控制权角度看，报告期内控股股东未发生变更；持股 10\% 以上的法人股东为烟台和正投资中心（有限合伙）；报告期内公司控股股东或第一大股东及其一致行动人累计质押股份数量占其所持公司股份数量比例未达到 80\%。

\pfl{经营要点} 经营层面，公司产能分布在中国、美国、加拿大、墨西哥、新西兰、柬埔寨等多个国家。

\pfl{所处行业的基本面} 宠物食品的供给是产能（含海外基地），需求由宠物数量与单宠消费驱动，是农业中与农产品价格周期联动最弱的一块（第 \ref{sec:framework} 章的供需框架）。 与该公司直接相关的行业数据包括：2025 年，中国城镇宠物（犬猫）消费市场规模 3,126.00 亿元、同比增长 4.1\%，其中宠物食品市场份额 53.7\%；2025 年，中国宠物食品行业零售市场规模 1,084.00 亿元；2024 年，我国宠物零食行业市场规模 474.40 亿元、同比增长 8.50\%。

\pfl{竞争优势} 把该公司放回同类公司比较，2026 年上半年营业收入 32.8 亿元、在三级行业“宠物食品”的 3 家公司中排第\zhnumber{2}位；毛利率 25.9\%、低于行业内中位数 2.5 个百分点；半年 ROE 4.0\%（未年化）、低于行业内中位数（4.00\%）；资产负债率 44.0\%、高于行业内中位数（35.3\%）。 相对优势方面，公司境内品牌业务毛利率 38.06\% 显著高于境外代工业务的 19.37\%。 就当期报表看，公司尚未把经营层面的特征转化为高于行业的资本回报，周期底部的相对优势更多体现在现金流与资产负债表上。

\pfl{历史股价} 自 2021 年 1 月至 2026 年 10 月，股价由 48.72 元变为 27.66 元，累计 -43.2\%，区间最高 62.86 元（2025 年 5 月）、最低 19.16 元（2024 年 7 月）；当前价相当于区间最高价的 44.0\%，为区间最低价的 1.44 倍。 月度收益的年化波动率 40.4\%，高于全样本中位数 37.3\%，在三级行业“宠物食品”的 3 家公司中波动率由高到低排第\zhnumber{2}位。

\begin{center}
\begin{minipage}{\textwidth}\centering
\begin{tikzpicture}
\begin{axis}[
  width=12.6cm, height=3.9cm, scale only axis,
  xmin=2020.922, xmax=2026.972, ymin=17.82, ymax=67.26,
  xtick={2021,2022,2023,2024,2025,2026},
  yearticks,
  yticklabel style={font=\scriptsize},
  ylabel={元/股}, ylabel style={font=\scriptsize},
  grid=major, grid style={LightGray, line width=0.3pt},
  axis line style={InkGray, line width=0.5pt},
  every axis plot/.append style={line width=0.9pt},
]
\addplot[AgriGreen] table[x=x,y=y]{data/clean/profiles/px_002891.dat};
\addplot[SoftRed, dashed, line width=0.7pt] table[x=x,y=y]{data/clean/profiles/pxzz_002891.dat};
\addplot[only marks, mark=*, mark size=1.1pt, SoftRed] table[x=x,y=y]{data/clean/profiles/pxzz_002891.dat};
\end{axis}
\end{tikzpicture}

\captionof{figure}[中宠股份（002891）股价与周期骨架]{中宠股份（002891）月度收盘价（元/股）与 ZigZag 周期骨架（阈值 25\%，圆点为周期转折点）}
\end{minipage}
\end{center}

\pfl{过去周期分析} 以 25\% 的 ZigZag 阈值划分，样本区间内股价形成 4 个主升段与 4 个主跌段。 最大主升段出现在 2024 年 7 月 至 2025 年 5 月，10 个月+228.1\%；最大主跌段为 2021 年 4 月 至 2022 年 4 月，12 个月-59.6\%。 最近一次转折出现在 2026 年 6 月（下跌段自 2025 年 5 月起，13 个月-57.5\%），当前处于该段的延续之中。 公司股价较申万农林牧渔指数提前约 22 个月见底，是板块中较早定价周期反转的公司之一。 宠物食品属于消费型农业（第 \ref{sec:financing-strategy} 节归入“向消费端延伸”），收入增长由宠物数量与单宠消费驱动，与农产品价格周期的联动最弱，公司 2026 年上半年实现盈利、半年 ROE 4.0\%（未年化），好于全样本中位数（0.75\%），下行周期对其报表的冲击小于同业。 综合区间形态看，该股单段涨跌幅度偏大、以趋势段为主（最大主升 +228.1\%、最大主跌 -59.6\%），年化波动率 40.4\% 与全样本中位数 37.3\% 相比波动与样本中位数接近。

\pfl{盈利情况} 2026 年上半年营业收入 32.81 亿元（同比 +34.9\%），归母净利润 1.27 亿元，同比 -37.6\%。 以年报为趋势对照，2023---2025 年营业收入由 37.47 亿元变为 52.21 亿元（两年年均复合增速 +18.0\%），归母净利润由 2.33 亿元变为 3.65 亿元。 按半年报口径，2026 年上半年的 ROE 为 4.0\%（未年化）、销售净利率 4.2\%、毛利率 25.9\%，ROE 在“宠物食品”3 家公司中按数值由高到低排第\zhnumber{2}位。 同期全样本半年 ROE 中位数 0.75\%，公司与之相差 3.25 个百分点（略好于中位数公司）。 针对上述压力，公司披露的应对举措包括：2026 年 3 月，公司实施“中宠转 2”可转债提前赎回。 综合看，公司在报告期维持盈利（高于全样本半年 ROE 中位数 0.75\%），利润的可持续性取决于宠物食品景气的延续与成本管控的执行。

\begin{center}\small\setlength{\tabcolsep}{3.4pt}
\captionof{table}[中宠股份（002891）关键财务指标]{中宠股份（002891）关键财务指标（合并报表口径；两个半年报期均未年化；现金短债比 $=$ 货币资金 $\div$ 短期债务）}
\begin{adjustbox}{max width=\textwidth}
\begin{tabular}{@{}lrrrrr@{}}
\toprule
指标 & 2023 年 & 2024 年 & 2025 年 & 2025 年半年报 & 2026 年半年报 \\
\midrule
营业收入（亿元） & 37.47 & 44.65 & 52.21 & 24.32 & 32.81 \\
归母净利润（亿元） & 2.33 & 3.94 & 3.65 & 2.03 & 1.27 \\
毛利率（\%） & 26.3 & 28.2 & 29.5 & 31.4 & 25.9 \\
ROE（\%） & 10.9 & 16.8 & 13.5 & 7.8 & 4.0 \\
资产负债率（\%） & 44.9 & 41.2 & 42.7 & 40.9 & 44.0 \\
经营活动现金流净额（亿元） & 4.47 & 4.96 & 5.34 & 2.35 & -1.93 \\
货币资金（亿元） & 5.27 & 5.13 & 12.73 & 13.87 & 13.68 \\
有息负债合计（亿元） & 14.34 & 11.68 & 15.19 & 14.94 & 18.56 \\
现金短债比（倍） & 0.77 & 1.22 & 1.27 & 1.38 & 0.85 \\
\bottomrule
\end{tabular}
\end{adjustbox}
\end{center}

\begin{center}
\begin{minipage}{\textwidth}\centering
\begin{tikzpicture}
\begin{groupplot}[
  group style={group size=2 by 1, horizontal sep=0.60cm},
  width=6.85cm, height=4.60cm,
  tick label style={font=\tiny},
  title style={font=\scriptsize\heiti},
  yticklabel style={font=\tiny},
  ylabel style={font=\tiny},
  grid=major, grid style={LightGray, line width=0.3pt},
  axis line style={InkGray, line width=0.5pt},
  enlarge x limits={abs=0.8},
  legend style={font=\scriptsize, at={(0.5,-0.20)}, anchor=north,
                legend columns=2, draw=none, fill=none, column sep=6pt},
]
\nextgroupplot[
  ybar, bar width=4pt,
  ymin=-5.22, ymax=58.48,
  xtick={0,1,2,3,4,5.7},
  xticklabels={2021,2022,2023,2024,2025,2026H1},
  ylabel={亿元},
]
\addplot[fill=AgriGreen!75, draw=AgriGreen, bar shift=-2.6pt] table[x=i, y=rev]{data/clean/profiles/fin_002891.dat};
\addplot[fill=SoftRed!60, draw=SoftRed, bar shift=2.6pt] table[x=i, y=ni]{data/clean/profiles/fin_002891.dat};
\addplot[fill=AgriGreen!30, draw=AgriGreen, pattern=north east lines, bar shift=-2.6pt] table[x=x, y=rev]{data/clean/profiles/finh_002891.dat};
\addplot[fill=SoftRed!20, draw=SoftRed, pattern=north east lines, bar shift=2.6pt] table[x=x, y=ni]{data/clean/profiles/finh_002891.dat};
\legend{营业收入（年/半年）, 归母净利润（年/半年）}
\nextgroupplot[
  ymin=-3.6, ymax=49.4,
  xtick={0,1,2,3,4,5.7},
  xticklabels={2021,2022,2023,2024,2025,2026H1},
  ylabel={\%},
]
\addplot[AgriGold, mark=*, mark size=1.3pt] table[x=i, y=roe]{data/clean/profiles/fin_002891.dat};
\addplot[OilBlue, mark=square*, mark size=1.3pt] table[x=i, y=debt]{data/clean/profiles/fin_002891.dat};
\addplot[only marks, AgriGold, mark=*, mark size=2.0pt] table[x=x, y=roe]{data/clean/profiles/finh_002891.dat};
\addplot[only marks, OilBlue, mark=square*, mark size=2.0pt] table[x=x, y=debt]{data/clean/profiles/finh_002891.dat};
\legend{ROE, 资产负债率}
\end{groupplot}
\end{tikzpicture}
\begin{tikzpicture}
\begin{axis}[
  name=finbot,
  width=13.75cm, height=3.10cm, scale only axis,
  ymin=-1.86, ymax=20.79,
  xtick={0,1,2,3,4,5.7},
  xticklabels={2021,2022,2023,2024,2025,2026H1},
  tick label style={font=\tiny},
  yticklabel style={font=\tiny},
  ylabel={亿元}, ylabel style={font=\tiny},
  ybar, bar width=4pt,
  grid=major, grid style={LightGray, line width=0.3pt},
  axis line style={InkGray, line width=0.5pt},
  enlarge x limits={abs=0.8},
  legend style={font=\scriptsize, at={(0.5,-0.26)}, anchor=north,
                legend columns=2, draw=none, fill=none, column sep=10pt},
]
\addplot[fill=AgriGreen!55, draw=AgriGreen, bar shift=-3.0pt] table[x=i, y=cash]{data/clean/profiles/fin_002891.dat};
\addplot[fill=SoftRed!45, draw=SoftRed, bar shift=3.0pt] table[x=i, y=ibd]{data/clean/profiles/fin_002891.dat};
\addplot[fill=AgriGreen!25, draw=AgriGreen, pattern=north east lines, bar shift=-3.0pt] table[x=x, y=cash]{data/clean/profiles/finh_002891.dat};
\addplot[fill=SoftRed!20, draw=SoftRed, pattern=north east lines, bar shift=3.0pt] table[x=x, y=ibd]{data/clean/profiles/finh_002891.dat};
\legend{货币资金（左轴）, 有息负债（左轴）}
\end{axis}
\begin{axis}[
  at={(finbot.south east)}, anchor=south east,
  width=13.75cm, height=3.10cm, scale only axis,
  axis y line*=right, axis x line*=none, xtick=\empty,
  ymin=-3.82, ymax=7.53,
  tick label style={font=\tiny},
  yticklabel style={font=\tiny}, scaled y ticks=false,
  legend style={font=\scriptsize, at={(0.5,-0.44)}, anchor=north,
                legend columns=1, draw=none, fill=none},
]
\addplot[OilBlue, mark=*, mark size=2.2pt, line width=1.3pt] table[x=i, y=ocf]{data/clean/profiles/fin_002891.dat};
\addplot[only marks, OilBlue, mark=*, mark size=3.2pt] table[x=x, y=ocf]{data/clean/profiles/finh_002891.dat};
\legend{经营活动现金流净额（右轴，亿元，蓝点）}
\end{axis}
\end{tikzpicture}

\captionof{figure}[中宠股份（002891）近年主要财务指标]{中宠股份（002891）近年主要财务指标（左上：营业收入与归母净利润；右上：ROE 与资产负债率；下：货币资金、有息负债为左轴柱状，经营活动现金流净额为其右轴的蓝点折线。
2021---2025 年为年报口径，2026H1 为 2026 年半年报/半年末，与其前一年报之间留出间隔、以斜纹或空心点区分；半年报数据未年化）}
\end{minipage}
\end{center}

\pfl{财务分析} 2026 年 6 月末资产负债率 44.0\%，较 2025 年末上升 1.3 个百分点（样本中位数 52.2\%，所处三级行业内按由低到高排第\zhnumber{3}位）。 2026 年 6 月末货币资金 13.68 亿元、短期债务（短期借款与一年内到期非流动负债）16.07 亿元、长期债务 2.50 亿元，有息负债合计 18.56 亿元，现金短债比 0.85 倍——覆盖偏紧。 净有息负债 4.88 亿元，相当于其 2026 年上半年营业收入的 15\%。 流动比率 1.29、速动比率 0.97，短期偿债能力处于正常区间。 2026 年上半年经营活动现金流净额 -1.93 亿元，净利润为正但经营现金流为负，利润的现金含量偏低。 2026 年 6 月末在建工程 4.39 亿元，较上年末增加 1.75 亿元，仍有在建投入。 2026 年 6 月末商誉 2.26 亿元，占归母净资产的 6.3\%，是净资产中最脆弱的一块。 综合看，现金短债比 0.85 倍，短期资金安排偏紧；资产负债率低于样本中位数 8.2 个百分点；有息负债中短期债务占比更高，期限结构仍以滚动续作为主。 公司披露的股东与质押风险为：2026 年 6 月，控股股东烟台中幸生物科技有限公司所持公司股份中存在质押 1,090.00 万股的情形。

\pfl{重大战略转型} 从公开披露看，公司的战略推进集中在：融资与资本运作（2 项）：2019 年，可转债募集资金投向年产 3 万吨宠物湿粮项目，提升湿粮品类产能；并购与重组（1 项）：2026 年 4 月，公司通过“自主建厂+并购整合”建设横跨五大洲的生产基地网络。 公告口径上，近三年（2023 年 9 月---2026 年 9 月）公司披露重大重组与并购类公告 0 条、对外投资与转型类公告 1 条、回购与股权激励类公告 35 条，合计公告 423 条。 第一大行业仍是“宠物食品及用品”，收入占比 93.9\%，与 2021 年年报相比没有方向性变化。 综合判断，报告期内公司有 1 条与战略相关的公告落地（重大重组与并购 0 条、对外投资与转型 1 条），资本运作与主业调整之间存在直接关联。

\begin{center}\small\setlength{\tabcolsep}{4pt}
\captionof{table}[中宠股份（002891）历史融资与资本运作]{中宠股份（002891）历史融资与资本运作（据公司公告与公开报道整理；金额为披露值，含首次公开发行、再融资、债券、银行借款与重整投资等）}
\begin{adjustbox}{max width=\textwidth}
\begin{tabular}{@{}p{1.45cm}p{2.55cm}p{4.05cm}p{4.95cm}@{}}
\toprule
时间 & 方式 & 规模 & 用途与说明 \\
\midrule
2017-08 & 首次公开发行A股股票并在深圳证券交易所上市（IPO） & 发行价 15.46 元/\allowbreak{}股；首发前总股本 7,\allowbreak{}500.00 万股 & 限责任公司，\allowbreak{}发行方式为网下询价配售与网上定价发行 \\
2019-03 & 公开发行可转换公司债券（中宠转债，债券代码128054） & 发行 1.94 亿元，\allowbreak{}扣除不含税发行费用后实际募集资金净额 1.84 亿元 & 用于年产 3 万吨宠物湿粮项目，\allowbreak{}该项目投资总额 2.70 亿元。；债券于 2019 年 3 月 14 日在深交所上市；公司后续实施提前赎回 \\
2019-07-31 & 重大资产重组（终止） & --- & 不适用。；公司终止重大资产重组并自 2019 年 7 月 31 日开市起复牌 \\
2020-10 & 非公开发行A股股票（2020年非公开发行） & 发行 1,\allowbreak{}737.33 万股，\allowbreak{}发行价格 37.50 元/\allowbreak{}股；募集资金总额 6.51 亿元 & 拟投入年产 6 万吨宠物干粮项目、\allowbreak{}年产 2 万吨宠物湿粮新西兰项目 2.30 亿元、\allowbreak{}营销中心建设及营销渠道智能化升级项目 1.14 亿元、\allowbreak{}补充流动资金 7,\allowbreak{}000.00 万元 \\
2022-11 & 公开发行可转换公司债券（中宠转2，债券代码127076） & 发行 7.69 亿元 & 投向年产 6 万吨高品质宠物干粮项目 2.30 亿元、\allowbreak{}年产 4 万吨新型宠物湿粮项目 2.56 亿元、\allowbreak{}年产 2,\allowbreak{}000 吨冻干宠物食品项目 5,\allowbreak{}254.86 万元、\allowbreak{}平面仓库智能立体化改造项目 3,\allowbreak{}636.68 万元及补充流动资金。；债券于 2022 年 11 月 21 日在深交所上市 \\
2026-01-23 & 股份回购方案 & 拟以集中竞价交易方式回购公司 A 股股票，\allowbreak{}回购价格不超过 78.00 元/\allowbreak{}股（含） & 使用公司自有资金和自筹资金回购公司股份。；公司于 2026 年 5 月 19 日公告取得金融机构股票回购专项贷款承诺函 \\
2026-05-19 & 股票回购专项贷款 & 取得金融机构股票回购专项贷款承诺函 & 专项用于回购公司股份。；回购方案已于 2026 年 1 月 23 日经第四届董事会第二十次会议审议通过 \\
\bottomrule
\end{tabular}
\end{adjustbox}
\end{center}

\pfl{历史融投资情况} 从现金流量表看，2025 年公司取得借款收到的现金 16.85 亿元、偿还债务支付的现金 10.75 亿元、吸收权益性投资 0.04 亿元，筹资活动现金流净额 3.43 亿元（2023 年为取得借款 6.82 亿元、偿还债务 3.66 亿元）。 2026 年上半年筹资活动现金流净额为 6.89 亿元（取得借款 14.30 亿元、偿还债务 7.43 亿元，未年化），仍处于净融入状态。 2025 年分配股利、利润或偿付利息支付的现金合计 1.37 亿元。 公开披露的融资记录共 7 笔，工具类型以 IPO、可转债、定向增发、银行借款为主；金额与用途见下表。 其中规模最大的两笔为：2026 年 1 月，股份回购方案，拟以集中竞价交易方式回购公司 A 股股票，回购价格不超过 78.00 元/股（含）；2022 年 11 月，可转债，发行 7.69 亿元。 股本方面，近三年披露股本变动 16 次（可转债转股 13 次、其他 1 次）；总股本自 2021 年末以来未发生变化（2.94 亿股）。 分红方面，近三年现金分红 7 次、合计每股派现 1.12 元，最近一次实施在 2026 年 6 月。 近三年“再融资”类公告 7 条，涉及可转债。 综合看，公司融资结构上以债务滚动为主、股权融资为补充；2025 年筹资活动现金流净额 3.43 亿元，年度内仍为净融入，与前述经营与投资现金流的方向基本一致。

\pfl{未来潜在融资需求分析} 2026 年 6 月末货币资金 13.68 亿元，而短期债务 16.07 亿元（现金短债比 0.85 倍），2026 年上半年经营现金流净额 -1.93 亿元。按“短期债务减去账面货币资金”的滚动偿付口径，静态缺口约 2.39 亿元；上半年盈利或接近盈亏平衡，该缺口不随经营亏损进一步扩大。 公司 2026 年 6 月末资产负债率 44.0\%、2026 年上半年 ROE 4.0\%（未年化），现金短债比 0.85 倍、短债覆盖不足，融资需求首先用于置换与滚动存量债务、补充营运资金，属于“补血型”而非扩张型。 公开披露的后续资金安排线索包括：股权融资线索：2026 年 4 月，公司 2022 年可转债募投项目中“年产 6 万吨高品质宠物干粮项目”“年产 4 万吨新型宠物湿粮项目”截至 2025 年末投资进度分别为 0.77\% 与 1.06\%；2026 年 8 月，公司未公告新的定向增发预案、配股或债券发行计划；债务与授信安排：2026 年 5 月，公司取得金融机构股票回购专项贷款承诺函，为最新一笔公开披露的银行融资安排。 资金需求还要叠加在建工程：期末在建工程 4.39 亿元、2026 年上半年购建固定资产等支出 2.81 亿元（未年化），这部分支出在周期底部通常会被主动放缓。 综合看，公司的外部资金需求以补充营运资金为主、项目投入为辅，滚动偿付口径下的静态缺口约 2.39 亿元，融资节奏将受制于银行续贷意愿与股权融资窗口。




\subsubsection{中水渔业（000798）}
\label{co:000798}

\pfl{公司基本情况} 公司前身可追溯至 1997 年，中国农发集团旗下的远洋渔业上市平台，1998 年 2 月在深交所上市（原中水集团远洋股份有限公司），以金枪鱼等大洋性捕捞为基本盘；2023 年以 17.15 亿元现金收购中渔环球 51\% 等农发集团渔业板块资产并配套 11 亿元并购贷款，2025 年扭亏（归母净利润 8,375.42 万元），2026 年上半年归母净利润 1.02 亿元、同比增长 18.60\%。 公司于 1998-02-12 在深交所主板首发上市，注册资本 3.66 亿元，总股本 3.66 亿股。 截至 2026 年 10 月 12 日收盘价 10.03 元，对应总市值 36.7 亿元，近一年+27.3\%，总市值在三级行业“海洋捕捞”的 2 家公司中按由大到小排第\zhnumber{1}位，近一年涨跌幅高于全样本中位数（-10.4\%）37.6 个百分点。 按 2026 年半年报披露的行业构成，收入占比前三位为渔业服务 47.2\%；捕捞 41.7\%；零售及加工贸易 11.0\%。 与 2021 年年报相比，第一大行业口径由“捕捞收入”（77.4\%）变为“渔业服务”（47.2\%）；两期披露的分类口径有调整，占比差异中同时包含分类因素。 股权结构方面，实际控制人为中国农业发展集团有限公司。

\pfl{股权结构与控制权} 按 2026 年 6 月 30 日的股东名册，前三大股东为中国水产有限公司 持股 33.90\%；中国水产舟山海洋渔业有限公司 持股 17.78\%；潘立明 持股 3.96\%，股东户数 28,850 户。 从控制权角度看，持股 33.90\%、1.24 亿股。

\pfl{经营要点} 业务与产能层面，2023 年，公司以现金收购中渔环球 51\% 股权后取得航运项目控制权。

\pfl{所处行业的基本面} 海洋捕捞的供给受资源与配额约束，需求端依赖餐饮与出口，成本端跟随燃油与人工（第 \ref{sec:framework} 章的供需框架）。 当期该行业的关键数据是：2025 年，全国远洋渔业产量 247.64 万吨、同比增长 13.12\%，占全国水产品总产量比重提升至 3.24\%。

\pfl{竞争优势} 把该公司放回同类公司比较，2026 年上半年营业收入 20.8 亿元、在三级行业“海洋捕捞”的 2 家公司中排第\zhnumber{1}位；毛利率 5.5\%、低于行业内中位数 8.8 个百分点；半年 ROE 20.9\%（未年化）、高于行业内中位数（9.79\%）；资产负债率 71.6\%、高于行业内中位数（53.0\%）。 与同类公司相比，公司的差异化在于：公司拥有金枪鱼延绳钓、围网、拖网及鱿鱼钓等多类型作业船队并配套海上加油与运输船舶。 优势在当期体现为高于行业中位数的资本回报，但优势能否持续仍取决于周期与自身产能投放节奏。

\pfl{历史股价} 本报告样本区间（2021 年 1 月 至 2026 年 10 月）内，股价由 8.19 元变为 10.03 元，累计 +22.5\%，区间最高 12.14 元（2022 年 8 月）、最低 5.63 元（2024 年 6 月）；当前价相当于区间最高价的 82.6\%，为区间最低价的 1.78 倍。 月度收益的年化波动率 45.1\%，高于全样本中位数 37.3\%，在三级行业“海洋捕捞”的 2 家公司中波动率由高到低排第\zhnumber{1}位。

\begin{center}
\begin{minipage}{\textwidth}\centering
\begin{tikzpicture}
\begin{axis}[
  width=12.6cm, height=3.9cm, scale only axis,
  xmin=2020.922, xmax=2026.972, ymin=5.24, ymax=12.99,
  xtick={2021,2022,2023,2024,2025,2026},
  yearticks,
  yticklabel style={font=\scriptsize},
  ylabel={元/股}, ylabel style={font=\scriptsize},
  grid=major, grid style={LightGray, line width=0.3pt},
  axis line style={InkGray, line width=0.5pt},
  every axis plot/.append style={line width=0.9pt},
]
\addplot[AgriGreen] table[x=x,y=y]{data/clean/profiles/px_000798.dat};
\addplot[SoftRed, dashed, line width=0.7pt] table[x=x,y=y]{data/clean/profiles/pxzz_000798.dat};
\addplot[only marks, mark=*, mark size=1.1pt, SoftRed] table[x=x,y=y]{data/clean/profiles/pxzz_000798.dat};
\end{axis}
\end{tikzpicture}

\captionof{figure}[中水渔业（000798）股价与周期骨架]{中水渔业（000798）月度收盘价（元/股）与 ZigZag 周期骨架（阈值 25\%，圆点为周期转折点）}
\end{minipage}
\end{center}

\pfl{过去周期分析} 以 25\% 的 ZigZag 阈值划分，样本区间内股价形成 3 个主升段与 3 个主跌段。 最大主跌段为 2022 年 8 月 至 2024 年 6 月，22 个月-53.6\%；最大主升段出现在 2024 年 6 月 至 2025 年 11 月，17 个月+115.6\%。 最近一次转折出现在 2026 年 10 月（上涨段自 2026 年 6 月起，4 个月+42.1\%），当前处于该段之后的震荡之中。 公司股价较申万农林牧渔指数提前约 23 个月见底，是板块中较早定价周期反转的公司之一。 远洋与近海捕捞的产量受资源与配额约束，成本端则跟随燃油与人工，公司 2026 年 6 月末资产负债率 71.6\%，高于全样本中位数 52.2\%，其股价因此更多反映资金面与信用风险，而不是价格弹性本身。 综合区间形态看，该股单段涨跌幅度偏大、以趋势段为主（最大主升 +115.6\%、最大主跌 -53.6\%），年化波动率 45.1\% 与全样本中位数 37.3\% 相比波动明显大于样本中位数。

\pfl{盈利情况} 2026 年上半年营业收入 20.78 亿元（同比 +18.9\%），归母净利润 1.02 亿元，同比 +18.6\%。 以年报为趋势对照，2023---2025 年营业收入由 40.42 亿元变为 43.42 亿元（两年年均复合增速 +3.6\%），归母净利润由 -1.18 亿元变为 0.84 亿元。 按半年报口径，2026 年上半年的 ROE 为 20.9\%（未年化）、销售净利率 6.8\%、毛利率 5.5\%，ROE 在“海洋捕捞”2 家公司中按数值由高到低排第\zhnumber{1}位。 同期全样本半年 ROE 中位数 0.75\%，公司与之相差 20.16 个百分点（略好于中位数公司）。 从公司披露的原因看，业绩变动主要来自：2025 年，归属于上市公司股东的净利润 8,375.00 万元、同比增加 1.88 亿元。 综合看，公司在报告期维持盈利（高于全样本半年 ROE 中位数 0.75\%），利润的可持续性取决于海洋捕捞景气的延续与成本管控的执行。

\begin{center}\small\setlength{\tabcolsep}{3.4pt}
\captionof{table}[中水渔业（000798）关键财务指标]{中水渔业（000798）关键财务指标（合并报表口径；两个半年报期均未年化；现金短债比 $=$ 货币资金 $\div$ 短期债务）}
\begin{adjustbox}{max width=\textwidth}
\begin{tabular}{@{}lrrrrr@{}}
\toprule
指标 & 2023 年 & 2024 年 & 2025 年 & 2025 年半年报 & 2026 年半年报 \\
\midrule
营业收入（亿元） & 40.42 & 49.33 & 43.42 & 17.47 & 20.78 \\
归母净利润（亿元） & -1.18 & -1.04 & 0.84 & 0.86 & 1.02 \\
毛利率（\%） & 5.1 & 3.2 & 8.3 & 5.2 & 5.5 \\
ROE（\%） & -8.3 & -25.4 & 21.1 & 21.3 & 20.9 \\
资产负债率（\%） & 72.5 & 73.8 & 72.9 & 72.4 & 71.6 \\
经营活动现金流净额（亿元） & -1.87 & 5.25 & 4.70 & 3.78 & -0.17 \\
货币资金（亿元） & 6.86 & 8.92 & 8.28 & 11.53 & 8.75 \\
有息负债合计（亿元） & 24.60 & 28.07 & 26.28 & 27.57 & 27.86 \\
现金短债比（倍） & 0.81 & 0.56 & 0.55 & 0.70 & 0.52 \\
\bottomrule
\end{tabular}
\end{adjustbox}
\end{center}

\begin{center}
\begin{minipage}{\textwidth}\centering
\begin{tikzpicture}
\begin{groupplot}[
  group style={group size=2 by 1, horizontal sep=0.60cm},
  width=6.85cm, height=4.60cm,
  tick label style={font=\tiny},
  title style={font=\scriptsize\heiti},
  yticklabel style={font=\tiny},
  ylabel style={font=\tiny},
  grid=major, grid style={LightGray, line width=0.3pt},
  axis line style={InkGray, line width=0.5pt},
  enlarge x limits={abs=0.8},
  legend style={font=\scriptsize, at={(0.5,-0.20)}, anchor=north,
                legend columns=2, draw=none, fill=none, column sep=6pt},
]
\nextgroupplot[
  ybar, bar width=4pt,
  ymin=-6.23, ymax=55.39,
  xtick={0,1,2,3,4,5.7},
  xticklabels={2021,2022,2023,2024,2025,2026H1},
  ylabel={亿元},
]
\addplot[fill=AgriGreen!75, draw=AgriGreen, bar shift=-2.6pt] table[x=i, y=rev]{data/clean/profiles/fin_000798.dat};
\addplot[fill=SoftRed!60, draw=SoftRed, bar shift=2.6pt] table[x=i, y=ni]{data/clean/profiles/fin_000798.dat};
\addplot[fill=AgriGreen!30, draw=AgriGreen, pattern=north east lines, bar shift=-2.6pt] table[x=x, y=rev]{data/clean/profiles/finh_000798.dat};
\addplot[fill=SoftRed!20, draw=SoftRed, pattern=north east lines, bar shift=2.6pt] table[x=x, y=ni]{data/clean/profiles/finh_000798.dat};
\legend{营业收入（年/半年）, 归母净利润（年/半年）}
\nextgroupplot[
  ymin=-33.4, ymax=83.7,
  xtick={0,1,2,3,4,5.7},
  xticklabels={2021,2022,2023,2024,2025,2026H1},
  ylabel={\%},
]
\addplot[AgriGold, mark=*, mark size=1.3pt] table[x=i, y=roe]{data/clean/profiles/fin_000798.dat};
\addplot[OilBlue, mark=square*, mark size=1.3pt] table[x=i, y=debt]{data/clean/profiles/fin_000798.dat};
\addplot[only marks, AgriGold, mark=*, mark size=2.0pt] table[x=x, y=roe]{data/clean/profiles/finh_000798.dat};
\addplot[only marks, OilBlue, mark=square*, mark size=2.0pt] table[x=x, y=debt]{data/clean/profiles/finh_000798.dat};
\legend{ROE, 资产负债率}
\end{groupplot}
\end{tikzpicture}
\begin{tikzpicture}
\begin{axis}[
  name=finbot,
  width=13.75cm, height=3.10cm, scale only axis,
  ymin=-2.81, ymax=31.44,
  xtick={0,1,2,3,4,5.7},
  xticklabels={2021,2022,2023,2024,2025,2026H1},
  tick label style={font=\tiny},
  yticklabel style={font=\tiny},
  ylabel={亿元}, ylabel style={font=\tiny},
  ybar, bar width=4pt,
  grid=major, grid style={LightGray, line width=0.3pt},
  axis line style={InkGray, line width=0.5pt},
  enlarge x limits={abs=0.8},
  legend style={font=\scriptsize, at={(0.5,-0.26)}, anchor=north,
                legend columns=2, draw=none, fill=none, column sep=10pt},
]
\addplot[fill=AgriGreen!55, draw=AgriGreen, bar shift=-3.0pt] table[x=i, y=cash]{data/clean/profiles/fin_000798.dat};
\addplot[fill=SoftRed!45, draw=SoftRed, bar shift=3.0pt] table[x=i, y=ibd]{data/clean/profiles/fin_000798.dat};
\addplot[fill=AgriGreen!25, draw=AgriGreen, pattern=north east lines, bar shift=-3.0pt] table[x=x, y=cash]{data/clean/profiles/finh_000798.dat};
\addplot[fill=SoftRed!20, draw=SoftRed, pattern=north east lines, bar shift=3.0pt] table[x=x, y=ibd]{data/clean/profiles/finh_000798.dat};
\legend{货币资金（左轴）, 有息负债（左轴）}
\end{axis}
\begin{axis}[
  at={(finbot.south east)}, anchor=south east,
  width=13.75cm, height=3.10cm, scale only axis,
  axis y line*=right, axis x line*=none, xtick=\empty,
  ymin=-3.72, ymax=7.39,
  tick label style={font=\tiny},
  yticklabel style={font=\tiny}, scaled y ticks=false,
  legend style={font=\scriptsize, at={(0.5,-0.44)}, anchor=north,
                legend columns=1, draw=none, fill=none},
]
\addplot[OilBlue, mark=*, mark size=2.2pt, line width=1.3pt] table[x=i, y=ocf]{data/clean/profiles/fin_000798.dat};
\addplot[only marks, OilBlue, mark=*, mark size=3.2pt] table[x=x, y=ocf]{data/clean/profiles/finh_000798.dat};
\legend{经营活动现金流净额（右轴，亿元，蓝点）}
\end{axis}
\end{tikzpicture}

\captionof{figure}[中水渔业（000798）近年主要财务指标]{中水渔业（000798）近年主要财务指标（左上：营业收入与归母净利润；右上：ROE 与资产负债率；下：货币资金、有息负债为左轴柱状，经营活动现金流净额为其右轴的蓝点折线。
2021---2025 年为年报口径，2026H1 为 2026 年半年报/半年末，与其前一年报之间留出间隔、以斜纹或空心点区分；半年报数据未年化）}
\end{minipage}
\end{center}

\pfl{财务分析} 2026 年 6 月末资产负债率 71.6\%，较 2025 年末下降 1.3 个百分点（样本中位数 52.2\%，所处三级行业内按由低到高排第\zhnumber{2}位）。 2026 年 6 月末货币资金 8.75 亿元、短期债务（短期借款与一年内到期非流动负债）16.75 亿元、长期债务 11.11 亿元，有息负债合计 27.86 亿元，现金短债比 0.52 倍——覆盖偏紧。 净有息负债 19.11 亿元，相当于其 2026 年上半年营业收入的 92\%。 流动比率 1.02、速动比率 0.44，短期指标尚可。 2026 年上半年经营活动现金流净额 -0.17 亿元，净利润为正但经营现金流为负，利润的现金含量偏低。 综合看，现金短债比 0.52 倍，短期资金安排偏紧；资产负债率高于样本中位数 19.4 个百分点；有息负债中短期债务占比更高，期限结构仍以滚动续作为主。 公司披露的重整与债务违约风险为：2026 年 6 月，公司因张福赐未完成 2015 年业绩承诺要求张福赐与张福庆承担违约责任。

\pfl{重大战略转型} 公司的战略动作可归为以下几类：并购与重组（2 项）：2023 年 6 月，上市公司在远洋捕捞、水产品贸易的基础上新增水产品加工、渔业服务等多元化业务；2020 年，公司收购超低温金枪鱼项目、丰汇项目后与原低温金枪鱼项目实现深度融合。 公告口径上，近三年（2023 年 9 月---2026 年 9 月）公司披露重大重组与并购类公告 4 条、对外投资与转型类公告 0 条、回购与股权激励类公告 0 条，合计公告 332 条。 主营结构的口径变化已经落到报表上：2021 年年报的第一大行业为“捕捞收入”（77.4\%），2026 年半年报为“渔业服务”（47.2\%）；公司在这两期的分部划分方式有过调整。 综合判断，报告期内公司有 4 条与战略相关的公告落地（重大重组与并购 4 条、对外投资与转型 0 条），资本运作与主业调整之间存在直接关联。

\begin{center}\small\setlength{\tabcolsep}{4pt}
\captionof{table}[中水渔业（000798）历史融资与资本运作]{中水渔业（000798）历史融资与资本运作（据公司公告与公开报道整理；金额为披露值，含首次公开发行、再融资、债券、银行借款与重整投资等）}
\begin{adjustbox}{max width=\textwidth}
\begin{tabular}{@{}p{1.45cm}p{2.55cm}p{4.05cm}p{4.95cm}@{}}
\toprule
时间 & 方式 & 规模 & 用途与说明 \\
\midrule
1997-11-24 & IPO（深圳证券交易所） & 发行 6,\allowbreak{}300 万股，\allowbreak{}发行价 7.24 元/\allowbreak{}股，\allowbreak{}实际募集资金 4.56 亿元 & 经证监发字[1997]480 号、\allowbreak{}证监发字[1997]481 号文批准 \\
2014-12 & 现金收购股权（并购） & 以现金购买张福赐持有的厦门新阳洲水产品工贸有限公司 55\% 股权 & 延伸水产品加工与贸易业务；交易对方张福赐承诺新阳洲 2014—2017 年净利润分别不低于 3,\allowbreak{}937.00 万元、\allowbreak{}4,\allowbreak{}324.00 万元、\allowbreak{}4,\allowbreak{}555.00 万元、\allowbreak{}4,\allowbreak{}707.00 万元 \\
2018-10 & 对外投资（关联交易） & 项目总投资 4.20 亿元，\allowbreak{}中国农发集团以现金出资 1.00 亿元占 23.81\% 股权 & 与控股股东及其关联方共同投资渔业相关项目；中水渔业与舟渔公司均系中国农发集团直接控制的企业 \\
2023 & 并购贷款 & 拟向合作银行申请总额 11.00 亿元人民币的并购贷款，\allowbreak{}用于支付本次交易的部分转让现金对价 & 支付重大资产购买对价；根据备考审阅报告，\allowbreak{}2022 年末公司模拟账面货币资金为 6.03 亿元；本次交易完成后公司资产负债率将从 28.97\% 大幅上升至 68.57\% \\
2023-06-15 & 重大资产购买（现金）暨关联交易 & 交易价格合计 17.15 亿元 & 整合中国农发集团渔业板块 \\
\bottomrule
\end{tabular}
\end{adjustbox}
\end{center}

\pfl{历史融投资情况} 从现金流量表看，2025 年公司取得借款收到的现金 20.94 亿元、偿还债务支付的现金 21.22 亿元，筹资活动现金流净额 -2.03 亿元（2023 年为取得借款 33.27 亿元、偿还债务 12.32 亿元）。 2026 年上半年筹资活动现金流净额为 1.23 亿元（取得借款 9.68 亿元、偿还债务 7.89 亿元，未年化），仍处于净融入状态。 2025 年分配股利、利润或偿付利息支付的现金合计 1.16 亿元。 公开披露的融资记录共 5 笔，工具类型以 IPO、银行借款为主；金额与用途见下表。 其中规模最大的两笔为：2023 年，银行借款，拟向合作银行申请总额 11.00 亿元人民币的并购贷款，用于支付本次交易的部分转让现金对价；1997 年 11 月，IPO，发行 6,300 万股，发行价 7.24 元/股，实际募集资金 4.56 亿元。 股本方面，近三年披露股本变动 4 次（限售股份上市 1 次）；总股本由 2021 年末的 3.19 亿股变为 3.66 亿股（+14.53\%）。 分红方面，近三年未实施现金分红，在全部样本中属于分红能力偏弱的一端。 综合看，公司融资结构上股权与债务融资并举；2025 年筹资活动现金流净额 -2.03 亿元，融资节奏仍以维持存量债务周转为主，与前述经营与投资现金流的方向基本一致。

\pfl{未来潜在融资需求分析} 2026 年 6 月末货币资金 8.75 亿元，而短期债务 16.75 亿元（现金短债比 0.52 倍），2026 年上半年经营现金流净额 -0.17 亿元。按“短期债务减去账面货币资金”的滚动偿付口径，静态缺口约 8.00 亿元；上半年盈利或接近盈亏平衡，该缺口不随经营亏损进一步扩大。 公司 2026 年 6 月末资产负债率 71.6\%、2026 年上半年 ROE 20.9\%（未年化），资产负债率偏高（71.6\%），融资需求首先用于置换与滚动存量债务、补充营运资金，属于“补血型”而非扩张型。 公开披露的后续资金安排线索包括：债务与授信安排：2023 年，公司以支付现金方式实施 17.15 亿元重大资产购买并申请 11.00 亿元并购贷款。 综合看，公司的外部资金需求以营运资金周转与在建项目投入为主，滚动偿付口径下的静态缺口约 8.00 亿元，融资节奏将受制于银行续贷意愿与股权融资窗口。




\subsubsection{佩蒂股份（300673）}
\label{co:300673}

\pfl{公司基本情况} 公司前身可追溯至 2013 年，浙江温州平阳起家的宠物食品（畜皮咬胶、植物咬胶、营养肉质零食、主粮）出口企业，2017 年 7 月创业板上市，是国内宠物食品行业最早“走出去”经营的企业之一，在越南（好嚼、巴啦啦、德信）、柬埔寨（爵味）、新西兰（北岛小镇、天然纯）建有生产基地，以 ODM 模式向欧美客户提供宠物零食供应链服务；2020 年定增募资 5.31 亿元、2021 年发行 7.2 亿元“佩蒂转债”，2026 年上半年营业收入 7.35 亿元、归母净利润 2,432.07 万元、同比下降 69.25\%。 公司于 2017-07-11 在深交所创业板首发上市，注册资本 2.47 亿元，总股本 2.49 亿股。 截至 2026 年 10 月 12 日收盘价 14.55 元，对应总市值 36.2 亿元，近一年-21.9\%，总市值在三级行业“宠物食品”的 3 家公司中按由大到小排第\zhnumber{3}位，近一年涨跌幅低于全样本中位数（-10.4\%）11.6 个百分点。 按 2026 年半年报披露的行业构成，收入占比前两位为宠物食品及其他 97.4\%；其他(补充) 2.6\%。

\pfl{股权结构与控制权} 按 2026 年 6 月 30 日的股东名册，前三大股东为陈振标 持股 32.55\%；陈振录 持股 10.51\%；郑香兰 持股 3.28\%，股东户数 13,815 户。 股权结构的另一条线索是，2026 年 6 月末，陈振标持股 8,050.28 万股、占总股本 32.55\%；报告期末公司通过回购专用证券账户持有本公司股份 776.68 万股、占公司股本总数的 3.14\%。

\pfl{经营要点} 从经营要点看，2026 年上半年，国外收入 5.15 亿元、约占总营收的七成；越南、柬埔寨工厂合计产能超 3 万吨，上半年东南亚分部收入 4.48 亿元、同比下降 4.38\%；2026 年上半年，财务费用 2,276.00 万元，上年同期为-362.00 万元、同比增加 728.77\%。

\pfl{所处行业的基本面} 宠物食品的供给是产能（含海外基地），需求由宠物数量与单宠消费驱动，是农业中与农产品价格周期联动最弱的一块（第 \ref{sec:framework} 章的供需框架）。 从供需两端的当期数据看，2025 年，中国城镇犬猫数量 1.26 亿只、同比增长 1.8\%；2024 年，我国宠物零食行业市场规模 474.40 亿元、同比增长 8.50\%。

\pfl{竞争优势} 把该公司放回同类公司比较，2026 年上半年营业收入 7.3 亿元、在三级行业“宠物食品”的 3 家公司中排第\zhnumber{3}位；毛利率 28.4\%、与行业内中位数基本持平；半年 ROE 1.4\%（未年化）、低于行业内中位数（4.00\%）；资产负债率 35.3\%、高于行业内中位数（35.3\%）。 公司区别于同类公司的经营特征包括：2013 年，是国内宠物食品企业出海建厂的先行者，以东南亚产能规避对美贸易摩擦的关税风险。 就当期报表看，公司尚未把经营层面的特征转化为高于行业的资本回报，周期底部的相对优势更多体现在现金流与资产负债表上。

\pfl{历史股价} 自 2021 年 1 月至 2026 年 10 月，股价由 31.15 元变为 14.55 元，累计 -53.3\%，区间最高 34.88 元（2021 年 5 月）、最低 11.08 元（2024 年 2 月）；当前价相当于区间最高价的 41.7\%，为区间最低价的 1.31 倍。 月度收益的年化波动率 38.8\%，高于全样本中位数 37.3\%，在三级行业“宠物食品”的 3 家公司中波动率由高到低排第\zhnumber{3}位。 把股价阶段与公司事件对齐看，2024 年 2 月 至 2025 年 11 月股价上涨+69.9\%，同期 2024 年末，陈振录持股 2,599.05 万股（10.51\%），无质押。 从时间对应关系看，公司层面的资本运作与股权、风险事件解释了区间内多数急涨急跌，与行业指数的同步性相对有限。

\begin{center}
\begin{minipage}{\textwidth}\centering
\begin{tikzpicture}
\begin{axis}[
  width=12.6cm, height=3.9cm, scale only axis,
  xmin=2020.922, xmax=2026.972, ymin=10.30, ymax=37.32,
  xtick={2021,2022,2023,2024,2025,2026},
  yearticks,
  yticklabel style={font=\scriptsize},
  ylabel={元/股}, ylabel style={font=\scriptsize},
  grid=major, grid style={LightGray, line width=0.3pt},
  axis line style={InkGray, line width=0.5pt},
  every axis plot/.append style={line width=0.9pt},
]
\addplot[AgriGreen] table[x=x,y=y]{data/clean/profiles/px_300673.dat};
\addplot[SoftRed, dashed, line width=0.7pt] table[x=x,y=y]{data/clean/profiles/pxzz_300673.dat};
\addplot[only marks, mark=*, mark size=1.1pt, SoftRed] table[x=x,y=y]{data/clean/profiles/pxzz_300673.dat};
\end{axis}
\end{tikzpicture}

\captionof{figure}[佩蒂股份（300673）股价与周期骨架]{佩蒂股份（300673）月度收盘价（元/股）与 ZigZag 周期骨架（阈值 25\%，圆点为周期转折点）}
\end{minipage}
\end{center}

\pfl{过去周期分析} 以 25\% 的 ZigZag 阈值划分，样本区间内股价形成 2 个主升段与 3 个主跌段。 最大主升段出现在 2024 年 2 月 至 2025 年 11 月，21 个月+69.9\%；最大主跌段为 2021 年 1 月 至 2022 年 4 月，15 个月-51.0\%。 最近一次转折出现在 2026 年 7 月（下跌段自 2025 年 11 月起，8 个月-31.9\%），当前处于该段的延续之中。 公司股价较申万农林牧渔指数提前约 27 个月见底，是板块中较早定价周期反转的公司之一。 宠物食品属于消费型农业（第 \ref{sec:financing-strategy} 节归入“向消费端延伸”），收入增长由宠物数量与单宠消费驱动，与农产品价格周期的联动最弱，公司 2026 年上半年实现盈利、半年 ROE 1.4\%（未年化），好于全样本中位数（0.75\%），下行周期对其报表的冲击小于同业。 综合区间形态看，该股单段涨跌幅度偏大、以趋势段为主（最大主升 +69.9\%、最大主跌 -51.0\%），年化波动率 38.8\% 与全样本中位数 37.3\% 相比波动与样本中位数接近。

\pfl{盈利情况} 2026 年上半年营业收入 7.35 亿元（同比 +1.0\%），归母净利润 0.24 亿元，同比 -69.3\%。 以年报为趋势对照，2023---2025 年营业收入由 14.11 亿元变为 14.49 亿元（两年年均复合增速 +1.3\%），归母净利润由 -0.11 亿元变为 1.16 亿元。 按半年报口径，2026 年上半年的 ROE 为 1.4\%（未年化）、销售净利率 3.4\%、毛利率 28.4\%，ROE 在“宠物食品”3 家公司中按数值由高到低排第\zhnumber{3}位。 按同一口径，三级行业“宠物食品”3 家公司的半年 ROE 中位数为 4.00\%，该公司低于其中位数 2.64 个百分点。 针对上述压力，公司披露的应对举措包括：产品涉及高端主粮、宠物零食等多个品类，其中越南、柬埔寨工厂产能合计超 3 万吨；2026 年 5 月，公司因注销部分回购股份减少注册资本暨通知债权人，并完成部分回购股份注销暨股份变动。 综合看，公司在报告期维持盈利（高于全样本半年 ROE 中位数 0.75\%），利润的可持续性取决于宠物食品景气的延续与成本管控的执行。

\begin{center}\small\setlength{\tabcolsep}{3.4pt}
\captionof{table}[佩蒂股份（300673）关键财务指标]{佩蒂股份（300673）关键财务指标（合并报表口径；两个半年报期均未年化；现金短债比 $=$ 货币资金 $\div$ 短期债务）}
\begin{adjustbox}{max width=\textwidth}
\begin{tabular}{@{}lrrrrr@{}}
\toprule
指标 & 2023 年 & 2024 年 & 2025 年 & 2025 年半年报 & 2026 年半年报 \\
\midrule
营业收入（亿元） & 14.11 & 16.59 & 14.49 & 7.28 & 7.35 \\
归母净利润（亿元） & -0.11 & 1.82 & 1.16 & 0.79 & 0.24 \\
毛利率（\%） & 19.3 & 29.4 & 31.9 & 32.0 & 28.4 \\
ROE（\%） & -0.6 & 9.6 & 5.8 & 3.9 & 1.4 \\
资产负债率（\%） & 35.9 & 34.3 & 34.0 & 33.0 & 35.3 \\
经营活动现金流净额（亿元） & 1.78 & 3.66 & 1.91 & -0.31 & -0.66 \\
货币资金（亿元） & 8.16 & 6.75 & 7.50 & 5.19 & 4.95 \\
有息负债合计（亿元） & 8.40 & 7.89 & 7.99 & 7.82 & 8.21 \\
现金短债比（倍） & 9.63 & 13.12 & 51.11 & 45.14 & 34.24 \\
\bottomrule
\end{tabular}
\end{adjustbox}
\end{center}

\begin{center}
\begin{minipage}{\textwidth}\centering
\begin{tikzpicture}
\begin{groupplot}[
  group style={group size=2 by 1, horizontal sep=0.60cm},
  width=6.85cm, height=4.60cm,
  tick label style={font=\tiny},
  title style={font=\scriptsize\heiti},
  yticklabel style={font=\tiny},
  ylabel style={font=\tiny},
  grid=major, grid style={LightGray, line width=0.3pt},
  axis line style={InkGray, line width=0.5pt},
  enlarge x limits={abs=0.8},
  legend style={font=\scriptsize, at={(0.5,-0.20)}, anchor=north,
                legend columns=2, draw=none, fill=none, column sep=6pt},
]
\nextgroupplot[
  ybar, bar width=4pt,
  ymin=-1.85, ymax=19.41,
  xtick={0,1,2,3,4,5.7},
  xticklabels={2021,2022,2023,2024,2025,2026H1},
  ylabel={亿元},
]
\addplot[fill=AgriGreen!75, draw=AgriGreen, bar shift=-2.6pt] table[x=i, y=rev]{data/clean/profiles/fin_300673.dat};
\addplot[fill=SoftRed!60, draw=SoftRed, bar shift=2.6pt] table[x=i, y=ni]{data/clean/profiles/fin_300673.dat};
\addplot[fill=AgriGreen!30, draw=AgriGreen, pattern=north east lines, bar shift=-2.6pt] table[x=x, y=rev]{data/clean/profiles/finh_300673.dat};
\addplot[fill=SoftRed!20, draw=SoftRed, pattern=north east lines, bar shift=2.6pt] table[x=x, y=ni]{data/clean/profiles/finh_300673.dat};
\legend{营业收入（年/半年）, 归母净利润（年/半年）}
\nextgroupplot[
  ymin=-3.6, ymax=40.9,
  xtick={0,1,2,3,4,5.7},
  xticklabels={2021,2022,2023,2024,2025,2026H1},
  ylabel={\%},
]
\addplot[AgriGold, mark=*, mark size=1.3pt] table[x=i, y=roe]{data/clean/profiles/fin_300673.dat};
\addplot[OilBlue, mark=square*, mark size=1.3pt] table[x=i, y=debt]{data/clean/profiles/fin_300673.dat};
\addplot[only marks, AgriGold, mark=*, mark size=2.0pt] table[x=x, y=roe]{data/clean/profiles/finh_300673.dat};
\addplot[only marks, OilBlue, mark=square*, mark size=2.0pt] table[x=x, y=debt]{data/clean/profiles/finh_300673.dat};
\legend{ROE, 资产负债率}
\end{groupplot}
\end{tikzpicture}
\begin{tikzpicture}
\begin{axis}[
  name=finbot,
  width=13.75cm, height=3.10cm, scale only axis,
  ymin=-1.03, ymax=11.53,
  xtick={0,1,2,3,4,5.7},
  xticklabels={2021,2022,2023,2024,2025,2026H1},
  tick label style={font=\tiny},
  yticklabel style={font=\tiny},
  ylabel={亿元}, ylabel style={font=\tiny},
  ybar, bar width=4pt,
  grid=major, grid style={LightGray, line width=0.3pt},
  axis line style={InkGray, line width=0.5pt},
  enlarge x limits={abs=0.8},
  legend style={font=\scriptsize, at={(0.5,-0.26)}, anchor=north,
                legend columns=2, draw=none, fill=none, column sep=10pt},
]
\addplot[fill=AgriGreen!55, draw=AgriGreen, bar shift=-3.0pt] table[x=i, y=cash]{data/clean/profiles/fin_300673.dat};
\addplot[fill=SoftRed!45, draw=SoftRed, bar shift=3.0pt] table[x=i, y=ibd]{data/clean/profiles/fin_300673.dat};
\addplot[fill=AgriGreen!25, draw=AgriGreen, pattern=north east lines, bar shift=-3.0pt] table[x=x, y=cash]{data/clean/profiles/finh_300673.dat};
\addplot[fill=SoftRed!20, draw=SoftRed, pattern=north east lines, bar shift=3.0pt] table[x=x, y=ibd]{data/clean/profiles/finh_300673.dat};
\legend{货币资金（左轴）, 有息负债（左轴）}
\end{axis}
\begin{axis}[
  at={(finbot.south east)}, anchor=south east,
  width=13.75cm, height=3.10cm, scale only axis,
  axis y line*=right, axis x line*=none, xtick=\empty,
  ymin=-1.79, ymax=4.96,
  tick label style={font=\tiny},
  yticklabel style={font=\tiny}, scaled y ticks=false,
  legend style={font=\scriptsize, at={(0.5,-0.44)}, anchor=north,
                legend columns=1, draw=none, fill=none},
]
\addplot[OilBlue, mark=*, mark size=2.2pt, line width=1.3pt] table[x=i, y=ocf]{data/clean/profiles/fin_300673.dat};
\addplot[only marks, OilBlue, mark=*, mark size=3.2pt] table[x=x, y=ocf]{data/clean/profiles/finh_300673.dat};
\legend{经营活动现金流净额（右轴，亿元，蓝点）}
\end{axis}
\end{tikzpicture}

\captionof{figure}[佩蒂股份（300673）近年主要财务指标]{佩蒂股份（300673）近年主要财务指标（左上：营业收入与归母净利润；右上：ROE 与资产负债率；下：货币资金、有息负债为左轴柱状，经营活动现金流净额为其右轴的蓝点折线。
2021---2025 年为年报口径，2026H1 为 2026 年半年报/半年末，与其前一年报之间留出间隔、以斜纹或空心点区分；半年报数据未年化）}
\end{minipage}
\end{center}

\pfl{财务分析} 2026 年 6 月末资产负债率 35.3\%，较 2025 年末上升 1.3 个百分点（样本中位数 52.2\%，所处三级行业内按由低到高排第\zhnumber{2}位）。 2026 年 6 月末货币资金 4.95 亿元、短期债务（短期借款与一年内到期非流动负债）0.14 亿元、长期债务 8.06 亿元，有息负债合计 8.21 亿元，现金短债比 34.24 倍——覆盖充分。 净有息负债 3.26 亿元，相当于其 2026 年上半年营业收入的 44\%。 流动比率 8.13、速动比率 6.20，短期偿债能力处于正常区间。 2026 年上半年经营活动现金流净额 -0.66 亿元，净利润为正但经营现金流为负，利润的现金含量偏低。 2026 年 6 月末在建工程 0.79 亿元，较上年末增加 0.12 亿元，仍有在建投入。 综合看，现金短债比 34.24 倍，短期偿付压力可控；资产负债率低于样本中位数 16.9 个百分点；有息负债中长期债务占比更高，期限结构对短债滚动的压力较小。 公司披露的股东与质押风险为：2025 年末，陈振录持股 2,599.05 万股（10.51\%），无质押。

\pfl{重大战略转型} 从公开披露看，公司的战略推进集中在：并购与重组（3 项）：2019 年 6 月，公司变更 6,500.00 万元募集资金用途并增资新西兰北岛小镇宠物食品有限公司 1,800.00 万美元；2019 年 2 月，公司完成对德信皮业（越南）有限公司 100\% 股权收购并改造成越南生产基地；海外与国际化（2 项）：2013 年，公司最早于 2013 年在越南设立第一家海外生产基地。 公告口径上，近三年（2023 年 9 月---2026 年 9 月）公司披露重大重组与并购类公告 0 条、对外投资与转型类公告 5 条、回购与股权激励类公告 60 条，合计公告 427 条。 第一大行业仍是“宠物食品及其他”，收入占比 97.4\%，与 2021 年年报相比没有方向性变化。 综合判断，报告期内公司有 5 条与战略相关的公告落地（重大重组与并购 0 条、对外投资与转型 5 条），资本运作与主业调整之间存在直接关联。

\begin{center}\small\setlength{\tabcolsep}{4pt}
\captionof{table}[佩蒂股份（300673）历史融资与资本运作]{佩蒂股份（300673）历史融资与资本运作（据公司公告与公开报道整理；金额为披露值，含首次公开发行、再融资、债券、银行借款与重整投资等）}
\begin{adjustbox}{max width=\textwidth}
\begin{tabular}{@{}p{1.45cm}p{2.55cm}p{4.05cm}p{4.95cm}@{}}
\toprule
时间 & 方式 & 规模 & 用途与说明 \\
\midrule
2017-07-11 & IPO（深圳证券交易所创业板） & 实际发行 2,\allowbreak{}000.00 万股 & 份有限公司，\allowbreak{}发行市盈率（按发行后总股本）22.99 倍 \\
2019-02-26 & 境外现金收购（越南生产基地） & 收购德信皮业（越南）有限公司 100\% 股权，\allowbreak{}股权取得成本 2,\allowbreak{}384.53 万元 & 扩充海外产能 \\
2019-06 & 境外投资设立子公司 & 通过子公司新西兰北岛小镇宠物食品有限公司出资 500.00 万美元设立新西兰天然纯宠物食品有限公司 & 实施“新西兰年产 4 万吨高品质宠物干粮新建项目”；该公司于 2019 年 6 月 26 日完成设立 \\
2019-06-11 & 变更部分募集资金用途并向全资子公司增资（用于境外收购） & 变更募集资金 6,\allowbreak{}500.00 万元人民币并加自有资金向新西兰北岛小镇宠物食品有限公司增资 1,\allowbreak{}800.00 万美元 & 收购 BOP Industries Limited 已发行的 77.02\% 股份（实缴）；变动前募投项目为江苏康贝宠物食品有限公司实施的“年产 2,\allowbreak{}500 吨植物咬胶、\allowbreak{}500 吨营养肉质零食生产线项目”；经 2019 年第二次临时股东大会审议通过 \\
2021-12-28 & 向不特定对象发行可转换公司债券（佩蒂转债） & 发行 720 万张、\allowbreak{}每张面值 100.00 元，\allowbreak{}期限 6 年，\allowbreak{}募集资金总额 7.20 亿元 & 经深交所审核同意并经中国证监会同意佩蒂动物营养科技股份有限；债券代码 123133 \\
\bottomrule
\end{tabular}
\end{adjustbox}
\end{center}

\pfl{历史融投资情况} 从现金流量表看，2025 年公司取得借款收到的现金 0.13 亿元、偿还债务支付的现金 0.45 亿元，筹资活动现金流净额 -1.83 亿元（2023 年为取得借款 1.60 亿元、偿还债务 1.38 亿元）。 2026 年上半年筹资活动现金流净额为 -0.69 亿元（取得借款 0.10 亿元、偿还债务 0.07 亿元，未年化），已转为净流出，处于净还债状态。 2025 年分配股利、利润或偿付利息支付的现金合计 0.98 亿元。 公开披露的融资记录共 5 笔，工具类型以 IPO、可转债为主；金额与用途见下表。 其中规模最大的两笔为：2019 年 6 月，变更部分募集资金用途，变更募集资金 6,500.00 万元人民币并加自有资金向新西兰北岛小镇宠物食品有限公司增资 1,800.00 万美元；2017 年 7 月，IPO，实际发行 2,000.00 万股。 股本方面，近三年披露股本变动 17 次（可转债转股 10 次、其他 4 次、股份回购 2 次、股权激励 1 次）；总股本由 2021 年末的 2.53 亿股变为 2.49 亿股（-1.81\%）。 分红方面，近三年现金分红 2 次、合计每股派现 0.60 元，最近一次实施在 2025 年 12 月。 近三年“再融资”类公告 19 条，涉及可转债。 综合看，公司融资结构上股权与债务融资并举；2025 年筹资活动现金流净额 -1.83 亿元，融资节奏仍以维持存量债务周转为主，与前述经营与投资现金流的方向基本一致。

\pfl{未来潜在融资需求分析} 2026 年 6 月末货币资金 4.95 亿元，而短期债务 0.14 亿元（现金短债比 34.24 倍），2026 年上半年经营现金流净额 -0.66 亿元。账面货币资金可以覆盖短期债务，缺口不体现在静态偿付层面。 公司 2026 年 6 月末资产负债率 35.3\%、2026 年上半年 ROE 1.4\%（未年化），财务结构相对宽松而盈利弹性有限，融资属于可选项而非必选项，触发条件多与产能建设或外延并购相关。 公开披露的后续资金安排线索包括：股权融资线索：2026 年 6 月，公司 2026 年半年度募集资金存放与使用情况显示：2021 年可转债募集资金累计投入 3.87 亿元。 资金需求还要叠加在建工程：期末在建工程 0.79 亿元、2026 年上半年购建固定资产等支出 0.41 亿元（未年化），这部分支出在周期底部通常会被主动放缓。 综合看，公司在静态偿付层面不存在缺口，外部融资更多是配合扩张节奏与并购整合的主动性安排，而非偿债压力驱动。




\subsubsection{好当家（600467）}
\label{co:600467}

\pfl{公司基本情况} 山东荣成的海参养殖与海洋食品加工企业，2004 年 4 月在上交所上市，拥有全国最大海参育苗与养殖基地之一；2025 年受海参价格与需求影响营业收入 10.51 亿元、同比下降 25.76\%，归母净利润 2,987.23 万元、同比下降 31.21\%，2026 年上半年营收 8.03 亿元、归母净利润 2,997.88 万元，同时推进不超过 1 亿美元的境外债券发行与不超过 35 亿元的年度融资计划。 公司于 2004-04-05 在上交所首发上市，注册资本 14.61 亿元，总股本 14.61 亿股。 截至 2026 年 10 月 12 日收盘价 2.36 元，对应总市值 34.5 亿元，近一年-1.7\%，总市值在三级行业“水产养殖”的 4 家公司中按由大到小排第\zhnumber{1}位，近一年涨跌幅高于全样本中位数（-10.4\%）8.7 个百分点。 按 2025 年年报披露的行业构成，收入占比前三位为食品加工业 57.1\%；养殖业 30.4\%；捕捞业 11.7\%。

\pfl{股权结构与控制权} 按 2026 年 6 月 30 日的股东名册，前三大股东为好当家集团有限公司 持股 36.62\%；李茹 持股 0.65\%；吴绍金 持股 0.45\%，股东户数 66,054 户。 从控制权角度看，2026 年 6 月末持有公司 5.35 亿股、占总股本 36.62\%，其中质押 3.19 亿股；唐传勤并直接持有公司 500.00 万股（占 0.34\%）；2024 年 6 月，唐传勤通过上海证券交易所交易系统以集中竞价方式增持公司股份 500.00 万股、占公司总股本 0.34\%。

\pfl{经营要点} 经营层面，2026 年上半年，公司鲜海参累计捕捞 3,114.78 吨、同比下降 13.11\%，其中对外销售 2,636.62 吨、下降 7.05\%，对内转库存 478.16 吨。

\pfl{所处行业的基本面} 水产养殖的供给由养殖面积与投苗量决定，需求由餐饮与家庭消费决定，价格弹性大于畜禽（第 \ref{sec:framework} 章的供需框架）。 从供需两端的当期数据看，2025 年，全国水产养殖面积 7760.60 千公顷、同比增长 2.55\%，其中海水养殖面积 2408.52 千公顷、增长 7.52\%；2026 年，大连近海养殖区连续三年清退、养殖面积缩减约 22\%，高端底播参与中低端养殖参价格走势背离；提出到 2030 年全省优质海参产量稳定在 10 万吨以上、产值突破 600.00 亿元。

\pfl{竞争优势} 把该公司放回同类公司比较，2026 年上半年营业收入 8.0 亿元、在三级行业“水产养殖”的 4 家公司中排第\zhnumber{2}位；毛利率 17.2\%、低于行业内中位数 0.5 个百分点；半年 ROE 0.9\%（未年化）、低于行业内中位数（6.05\%）；资产负债率 49.1\%、低于行业内中位数（72.2\%）。 与同类公司相比，公司的差异化在于：公司拥有 5 万亩围堰海参养殖面积。 就当期报表看，公司尚未把经营层面的特征转化为高于行业的资本回报，周期底部的相对优势更多体现在现金流与资产负债表上。

\pfl{历史股价} 本报告样本区间（2021 年 1 月 至 2026 年 10 月）内，股价由 2.50 元变为 2.36 元，累计 -5.6\%，区间最高 3.16 元（2022 年 3 月）、最低 1.42 元（2024 年 6 月）；当前价相当于区间最高价的 74.7\%，为区间最低价的 1.66 倍。 月度收益的年化波动率 28.8\%，低于全样本中位数 37.3\%，在三级行业“水产养殖”的 4 家公司中波动率由高到低排第\zhnumber{4}位。

\begin{center}
\begin{minipage}{\textwidth}\centering
\begin{tikzpicture}
\begin{axis}[
  width=12.6cm, height=3.9cm, scale only axis,
  xmin=2020.922, xmax=2026.972, ymin=1.32, ymax=3.38,
  xtick={2021,2022,2023,2024,2025,2026},
  yearticks,
  yticklabel style={font=\scriptsize},
  ylabel={元/股}, ylabel style={font=\scriptsize},
  grid=major, grid style={LightGray, line width=0.3pt},
  axis line style={InkGray, line width=0.5pt},
  every axis plot/.append style={line width=0.9pt},
]
\addplot[AgriGreen] table[x=x,y=y]{data/clean/profiles/px_600467.dat};
\addplot[SoftRed, dashed, line width=0.7pt] table[x=x,y=y]{data/clean/profiles/pxzz_600467.dat};
\addplot[only marks, mark=*, mark size=1.1pt, SoftRed] table[x=x,y=y]{data/clean/profiles/pxzz_600467.dat};
\end{axis}
\end{tikzpicture}

\captionof{figure}[好当家（600467）股价与周期骨架]{好当家（600467）月度收盘价（元/股）与 ZigZag 周期骨架（阈值 25\%，圆点为周期转折点）}
\end{minipage}
\end{center}

\pfl{过去周期分析} 以 25\% 的 ZigZag 阈值划分，样本区间内股价形成 2 个主升段与 2 个主跌段。 最大主跌段为 2022 年 3 月 至 2024 年 6 月，27 个月-55.1\%；最大主升段出现在 2024 年 6 月 至 2026 年 1 月，19 个月+93.7\%。 最近一次转折出现在 2026 年 6 月（下跌段自 2026 年 1 月起，5 个月-28.4\%），当前处于该段的延续之中。 公司股价较申万农林牧渔指数提前约 23 个月见底，是板块中较早定价周期反转的公司之一。 水产养殖的价格周期由投苗量与消费需求共同决定（第 \ref{sec:other-livestock} 节），与猪周期的同步性弱于禽链，公司 2026 年上半年实现盈利、半年 ROE 0.9\%（未年化），好于全样本中位数（0.75\%），下行周期对其报表的冲击小于同业。 综合区间形态看，该股单段涨跌幅度偏大、以趋势段为主（最大主升 +93.7\%、最大主跌 -55.1\%），年化波动率 28.8\% 与全样本中位数 37.3\% 相比波动小于样本中位数。

\pfl{盈利情况} 2026 年上半年营业收入 8.03 亿元（同比 +13.0\%），归母净利润 0.30 亿元，同比 +21.7\%。 以年报为趋势对照，2023---2025 年营业收入由 15.64 亿元变为 10.51 亿元（两年年均复合增速 -18.0\%），归母净利润由 0.49 亿元变为 0.30 亿元。 按半年报口径，2026 年上半年的 ROE 为 0.9\%（未年化）、销售净利率 3.8\%、毛利率 17.2\%，ROE 在“水产养殖”4 家公司中按数值由高到低排第\zhnumber{4}位。 同期全样本半年 ROE 中位数 0.75\%，公司与之相差 0.12 个百分点（略好于中位数公司）。 综合看，公司在报告期维持盈利（高于全样本半年 ROE 中位数 0.75\%），利润的可持续性取决于水产养殖景气的延续与成本管控的执行。

\begin{center}\small\setlength{\tabcolsep}{3.4pt}
\captionof{table}[好当家（600467）关键财务指标]{好当家（600467）关键财务指标（合并报表口径；两个半年报期均未年化；现金短债比 $=$ 货币资金 $\div$ 短期债务）}
\begin{adjustbox}{max width=\textwidth}
\begin{tabular}{@{}lrrrrr@{}}
\toprule
指标 & 2023 年 & 2024 年 & 2025 年 & 2025 年半年报 & 2026 年半年报 \\
\midrule
营业收入（亿元） & 15.64 & 14.16 & 10.51 & 7.11 & 8.03 \\
归母净利润（亿元） & 0.49 & 0.43 & 0.30 & 0.25 & 0.30 \\
毛利率（\%） & 16.5 & 15.6 & 23.9 & 20.1 & 17.2 \\
ROE（\%） & 1.5 & 1.3 & 0.9 & 0.8 & 0.9 \\
资产负债率（\%） & 49.1 & 49.1 & 48.5 & 51.4 & 49.1 \\
经营活动现金流净额（亿元） & 8.68 & 5.88 & 4.51 & 0.59 & 1.08 \\
货币资金（亿元） & 4.80 & 4.14 & 4.40 & 6.25 & 5.64 \\
有息负债合计（亿元） & 23.87 & 25.92 & 27.10 & 29.28 & 27.09 \\
现金短债比（倍） & 0.21 & 0.18 & 0.18 & 0.25 & 0.24 \\
\bottomrule
\end{tabular}
\end{adjustbox}
\end{center}

\begin{center}
\begin{minipage}{\textwidth}\centering
\begin{tikzpicture}
\begin{groupplot}[
  group style={group size=2 by 1, horizontal sep=0.60cm},
  width=6.85cm, height=4.60cm,
  tick label style={font=\tiny},
  title style={font=\scriptsize\heiti},
  yticklabel style={font=\tiny},
  ylabel style={font=\tiny},
  grid=major, grid style={LightGray, line width=0.3pt},
  axis line style={InkGray, line width=0.5pt},
  enlarge x limits={abs=0.8},
  legend style={font=\scriptsize, at={(0.5,-0.20)}, anchor=north,
                legend columns=2, draw=none, fill=none, column sep=6pt},
]
\nextgroupplot[
  ybar, bar width=4pt,
  ymin=-1.56, ymax=17.51,
  xtick={0,1,2,3,4,5.7},
  xticklabels={2021,2022,2023,2024,2025,2026H1},
  ylabel={亿元},
]
\addplot[fill=AgriGreen!75, draw=AgriGreen, bar shift=-2.6pt] table[x=i, y=rev]{data/clean/profiles/fin_600467.dat};
\addplot[fill=SoftRed!60, draw=SoftRed, bar shift=2.6pt] table[x=i, y=ni]{data/clean/profiles/fin_600467.dat};
\addplot[fill=AgriGreen!30, draw=AgriGreen, pattern=north east lines, bar shift=-2.6pt] table[x=x, y=rev]{data/clean/profiles/finh_600467.dat};
\addplot[fill=SoftRed!20, draw=SoftRed, pattern=north east lines, bar shift=2.6pt] table[x=x, y=ni]{data/clean/profiles/finh_600467.dat};
\legend{营业收入（年/半年）, 归母净利润（年/半年）}
\nextgroupplot[
  ymin=-4.0, ymax=55.6,
  xtick={0,1,2,3,4,5.7},
  xticklabels={2021,2022,2023,2024,2025,2026H1},
  ylabel={\%},
]
\addplot[AgriGold, mark=*, mark size=1.3pt] table[x=i, y=roe]{data/clean/profiles/fin_600467.dat};
\addplot[OilBlue, mark=square*, mark size=1.3pt] table[x=i, y=debt]{data/clean/profiles/fin_600467.dat};
\addplot[only marks, AgriGold, mark=*, mark size=2.0pt] table[x=x, y=roe]{data/clean/profiles/finh_600467.dat};
\addplot[only marks, OilBlue, mark=square*, mark size=2.0pt] table[x=x, y=debt]{data/clean/profiles/finh_600467.dat};
\legend{ROE, 资产负债率}
\end{groupplot}
\end{tikzpicture}
\begin{tikzpicture}
\begin{axis}[
  name=finbot,
  width=13.75cm, height=3.10cm, scale only axis,
  ymin=-2.78, ymax=31.17,
  xtick={0,1,2,3,4,5.7},
  xticklabels={2021,2022,2023,2024,2025,2026H1},
  tick label style={font=\tiny},
  yticklabel style={font=\tiny},
  ylabel={亿元}, ylabel style={font=\tiny},
  ybar, bar width=4pt,
  grid=major, grid style={LightGray, line width=0.3pt},
  axis line style={InkGray, line width=0.5pt},
  enlarge x limits={abs=0.8},
  legend style={font=\scriptsize, at={(0.5,-0.26)}, anchor=north,
                legend columns=2, draw=none, fill=none, column sep=10pt},
]
\addplot[fill=AgriGreen!55, draw=AgriGreen, bar shift=-3.0pt] table[x=i, y=cash]{data/clean/profiles/fin_600467.dat};
\addplot[fill=SoftRed!45, draw=SoftRed, bar shift=3.0pt] table[x=i, y=ibd]{data/clean/profiles/fin_600467.dat};
\addplot[fill=AgriGreen!25, draw=AgriGreen, pattern=north east lines, bar shift=-3.0pt] table[x=x, y=cash]{data/clean/profiles/finh_600467.dat};
\addplot[fill=SoftRed!20, draw=SoftRed, pattern=north east lines, bar shift=3.0pt] table[x=x, y=ibd]{data/clean/profiles/finh_600467.dat};
\legend{货币资金（左轴）, 有息负债（左轴）}
\end{axis}
\begin{axis}[
  at={(finbot.south east)}, anchor=south east,
  width=13.75cm, height=3.10cm, scale only axis,
  axis y line*=right, axis x line*=none, xtick=\empty,
  ymin=-2.26, ymax=11.28,
  tick label style={font=\tiny},
  yticklabel style={font=\tiny}, scaled y ticks=false,
  legend style={font=\scriptsize, at={(0.5,-0.44)}, anchor=north,
                legend columns=1, draw=none, fill=none},
]
\addplot[OilBlue, mark=*, mark size=2.2pt, line width=1.3pt] table[x=i, y=ocf]{data/clean/profiles/fin_600467.dat};
\addplot[only marks, OilBlue, mark=*, mark size=3.2pt] table[x=x, y=ocf]{data/clean/profiles/finh_600467.dat};
\legend{经营活动现金流净额（右轴，亿元，蓝点）}
\end{axis}
\end{tikzpicture}

\captionof{figure}[好当家（600467）近年主要财务指标]{好当家（600467）近年主要财务指标（左上：营业收入与归母净利润；右上：ROE 与资产负债率；下：货币资金、有息负债为左轴柱状，经营活动现金流净额为其右轴的蓝点折线。
2021---2025 年为年报口径，2026H1 为 2026 年半年报/半年末，与其前一年报之间留出间隔、以斜纹或空心点区分；半年报数据未年化）}
\end{minipage}
\end{center}

\pfl{财务分析} 2026 年 6 月末资产负债率 49.1\%，较 2025 年末上升 0.6 个百分点（样本中位数 52.2\%，所处三级行业内按由低到高排第\zhnumber{2}位）。 2026 年 6 月末货币资金 5.64 亿元、短期债务（短期借款与一年内到期非流动负债）23.80 亿元、长期债务 3.29 亿元，有息负债合计 27.09 亿元，现金短债比 0.24 倍——缺口明显。 净有息负债 21.45 亿元，相当于其 2026 年上半年营业收入的 267\%。 流动比率 0.51、速动比率 0.26，短期偿债指标偏紧。 2026 年上半年经营活动现金流净额 1.08 亿元，经营现金流与利润同向为正，内生资金可以支撑运营。 综合看，账面货币资金对有息负债与短期债务的覆盖不足，滚动偿付对新增融资的依赖度较高；资产负债率低于样本中位数 3.1 个百分点；有息负债中短期债务占比更高，期限结构仍以滚动续作为主。 公司披露的股东与质押风险为：2026 年 1 月，控股股东好当家集团有限公司持有公司股份累计质押数量（含本次）3.46 亿股。

\pfl{重大战略转型} 近三年公司的战略动作集中在：产能与项目（1 项）：2024 年 12 月，公司建设并投产 130 多万平方水体现代海参育苗车间；融资与资本运作（1 项）：2026 年 4 月，公司拟委任天风国际证券与期货有限公司担任本次境外债券的和牵头安排行。 公告口径上，近三年（2023 年 9 月---2026 年 9 月）公司披露重大重组与并购类公告 0 条、对外投资与转型类公告 0 条、回购与股权激励类公告 0 条，合计公告 196 条。 第一大行业“食品加工业”的收入占比由 2021 年年报的 44.0\% 变为 57.1\%（13.0 个百分点），收入结构出现再平衡。 综合判断，公司报告期内的变化主要发生在收入结构层面，属于口径与业务重心的再平衡，而非外延式扩张。

\begin{center}\small\setlength{\tabcolsep}{4pt}
\captionof{table}[好当家（600467）历史融资与资本运作]{好当家（600467）历史融资与资本运作（据公司公告与公开报道整理；金额为披露值，含首次公开发行、再融资、债券、银行借款与重整投资等）}
\begin{adjustbox}{max width=\textwidth}
\begin{tabular}{@{}p{1.45cm}p{2.55cm}p{4.05cm}p{4.95cm}@{}}
\toprule
时间 & 方式 & 规模 & 用途与说明 \\
\midrule
2004-04-05 & IPO（上海证券交易所主板） & 发行 6,\allowbreak{}000 万股 & 经证监发行字[2004]26 号文批准 \\
2011-11-23 & 非公开发行股票（向特定对象） & 发行 9,\allowbreak{}689.72 万股 & 公司于 2011 年 4 月 29 日至 2011 年 11 月 23 日期间以自有资金先期投入募投项目 \\
2015 & 资本公积转增股本（权益融资之外股本扩张） & 每 10 股转增 10 股，\allowbreak{}共计转增股本 7.30 亿元 & 扩大股本规模；转增后公司注册资本为人民币 14.61 亿元 \\
2026-01-13 & 控股股东股权质押融资 & 好当家集团有限公司持有公司股份 5.73 亿股（占 39.25\%）；2025 年 1 月 10 日其将 730.00 万股无限售流通股质押给交通银行股份有限公司威海分行提供质押担保 & 为控股股东融资提供质押担保；2025 年 12 月控股股东再将部分股份质押给交通银行威海分行（占公司总股本 0.68\%） \\
2026-04-22 & 年度融资计划与境外债券（董事会授权） & 母公司及控股子公司向国内外银行、\allowbreak{}融资租赁公司、\allowbreak{}商业保理公司及国内外资本市场融资静态余额拟不超过 35.00 亿元 & 满足 2026 年度生产经营需要 \\
\bottomrule
\end{tabular}
\end{adjustbox}
\end{center}

\pfl{历史融投资情况} 从现金流量表看，2025 年公司取得借款收到的现金 35.09 亿元、偿还债务支付的现金 33.44 亿元，筹资活动现金流净额 -1.23 亿元（2023 年为取得借款 27.54 亿元、偿还债务 29.74 亿元）。 2026 年上半年筹资活动现金流净额为 -0.11 亿元（取得借款 14.73 亿元、偿还债务 12.95 亿元，未年化），已转为净流出，处于净还债状态。 2025 年分配股利、利润或偿付利息支付的现金合计 1.75 亿元。 公开披露的融资记录共 5 笔，工具类型以 IPO、定向增发为主；金额与用途见下表。 其中规模最大的两笔为：2026 年 4 月，年度融资计划与境外债，母公司及控股子公司向国内外银行、融资租赁公司、商业保理公司及国内外资本市场融资静态余额拟不超过 35.00 亿元；2011 年 11 月，定向增发，发行 9,689.72 万股。 股本方面，近三年披露股本变动 3 次；总股本自 2021 年末以来未发生变化（14.61 亿股）。 分红方面，近三年现金分红 4 次、合计每股派现 0.04 元，最近一次实施在 2025 年 12 月。 综合看，公司融资结构上股权与债务融资并举；2025 年筹资活动现金流净额 -1.23 亿元，融资节奏仍以维持存量债务周转为主，与前述经营与投资现金流的方向基本一致。

\pfl{未来潜在融资需求分析} 2026 年 6 月末货币资金 5.64 亿元，而短期债务 23.80 亿元（现金短债比 0.24 倍），2026 年上半年经营现金流净额 1.08 亿元。按“短期债务减去账面货币资金”的滚动偿付口径，静态缺口约 18.16 亿元；上半年盈利或接近盈亏平衡，该缺口不随经营亏损进一步扩大。 公司 2026 年 6 月末资产负债率 49.1\%、2026 年上半年 ROE 0.9\%（未年化），现金短债比 0.24 倍、短债覆盖不足，融资需求首先用于置换与滚动存量债务、补充营运资金，属于“补血型”而非扩张型。 公开披露的后续资金安排线索包括：股权融资线索：2026 年 4 月，公司 2026-2027 年度融资计划：银行金融机构、融资租赁公司、商业保理公司及国内外资本市场融资静态余额拟不超过 35.00 亿元。 综合看，公司的外部资金需求以补充营运资金为主、项目投入为辅，滚动偿付口径下的静态缺口约 18.16 亿元，融资节奏将受制于银行续贷意愿与股权融资窗口。




\subsubsection{国联水产（300094）}
\label{co:300094}

\pfl{公司基本情况} 湛江水产品加工与预制菜龙头，2010 年 7 月在深交所创业板上市，2010 年 IPO 募资 11.50 亿元、2022 年定增募资 10 亿元；2023—2025 年连续亏损、2025 年归母净利润 -16.31 亿元，2026 年上半年扭亏为盈（归母净利润 1,003.92 万元）但资产负债率达 95.19\%，公司通过出售子公司与资产、引入益阳市国资股东等方式压降负债与补充流动性。 公司于 2010-07-08 在深交所创业板首发上市，注册资本 11.28 亿元，总股本 11.28 亿股。 截至 2026 年 10 月 12 日收盘价 2.96 元，对应总市值 33.4 亿元，近一年-18.9\%，总市值在三级行业“水产养殖”的 4 家公司中按由大到小排第\zhnumber{2}位，近一年涨跌幅低于全样本中位数（-10.4\%）8.5 个百分点。 股权结构方面，控股股东为新余国通投资管理有限公司（持股 14.00\%）。

\pfl{股权结构与控制权} 按 2026 年 6 月 30 日的股东名册，前三大股东为新余国通投资管理有限公司 持股 14.00\%；冠聯國際投資有限公司 持股 6.89\%；李忠 持股 1.59\%，股东户数 52,365 户。 股东结构与控制权方面，新余国通注册资本 5,000.00 万元，股权结构为李忠持股 47.86\%、陈汉与李国通各持股 26.07\%；冠联国际投资有限公司持有公司 6.89\%，为第二大股东；2026 年 8 月，控股股东新余国通拟向益阳高新农业开发有限公司协议转让公司股份 6,402.22 万股、占总股本 5.67\%。

\pfl{经营要点} 经营与产能方面，2025 年，公司水产品出口美国收入同比增长 26.55\%；集中资源推动面包虾产能、产量、利润三重跃升。

\pfl{所处行业的基本面} 水产养殖的供给由养殖面积与投苗量决定，需求由餐饮与家庭消费决定，价格弹性大于畜禽（第 \ref{sec:framework} 章的供需框架）。 当期该行业的关键数据是：农业农村部渔业渔政管理局中国南美白对虾供需报告估计 2025 年我国南美白对虾养殖产量 250 万吨、增长 5.5\%；该报告预测 2026 年南美白对虾养殖面积增至 423 万亩、产量 263 万吨、进口量 90 万吨；2026 年 6 月，美国贸易代表办公室提出对中国商品加征 12.5\% 关税，中国罗非鱼进入美国市场的综合关税可能升至 37.5\%。

\pfl{竞争优势} 把该公司放回同类公司比较，2026 年上半年营业收入 14.5 亿元、在三级行业“水产养殖”的 4 家公司中排第\zhnumber{1}位；毛利率 11.4\%、低于行业内中位数 6.3 个百分点；半年 ROE 10.0\%（未年化）、高于行业内中位数（6.05\%）；资产负债率 95.2\%、高于行业内中位数（72.2\%）。 公司区别于同类公司的经营特征包括：公司构建“龙霸＋小霸龙＋Talassa”多品牌协同矩阵，以对虾为核心品类提升各渠道市场占有率。 优势在当期体现为高于行业中位数的资本回报，但优势能否持续仍取决于周期与自身产能投放节奏。

\pfl{历史股价} 自 2021 年 1 月至 2026 年 10 月，股价由 3.88 元变为 2.96 元，累计 -23.7\%，区间最高 6.62 元（2022 年 3 月）、最低 2.45 元（2026 年 6 月）；当前价相当于区间最高价的 44.7\%，为区间最低价的 1.21 倍。 月度收益的年化波动率 44.3\%，高于全样本中位数 37.3\%，在三级行业“水产养殖”的 4 家公司中波动率由高到低排第\zhnumber{2}位。

\begin{center}
\begin{minipage}{\textwidth}\centering
\begin{tikzpicture}
\begin{axis}[
  width=12.6cm, height=3.9cm, scale only axis,
  xmin=2020.922, xmax=2026.972, ymin=2.28, ymax=7.08,
  xtick={2021,2022,2023,2024,2025,2026},
  yearticks,
  yticklabel style={font=\scriptsize},
  ylabel={元/股}, ylabel style={font=\scriptsize},
  grid=major, grid style={LightGray, line width=0.3pt},
  axis line style={InkGray, line width=0.5pt},
  every axis plot/.append style={line width=0.9pt},
]
\addplot[AgriGreen] table[x=x,y=y]{data/clean/profiles/px_300094.dat};
\addplot[SoftRed, dashed, line width=0.7pt] table[x=x,y=y]{data/clean/profiles/pxzz_300094.dat};
\addplot[only marks, mark=*, mark size=1.1pt, SoftRed] table[x=x,y=y]{data/clean/profiles/pxzz_300094.dat};
\end{axis}
\end{tikzpicture}

\captionof{figure}[国联水产（300094）股价与周期骨架]{国联水产（300094）月度收盘价（元/股）与 ZigZag 周期骨架（阈值 25\%，圆点为周期转折点）}
\end{minipage}
\end{center}

\pfl{过去周期分析} 以 25\% 的 ZigZag 阈值划分，样本区间内股价形成 5 个主升段与 6 个主跌段。 最大主升段出现在 2024 年 6 月 至 2024 年 11 月，5 个月+83.9\%；最大主跌段为 2023 年 11 月 至 2024 年 6 月，7 个月-55.7\%。 最近一次转折出现在 2026 年 6 月（下跌段自 2024 年 11 月起，19 个月-47.5\%），当前处于该段的延续之中。 公司股价与申万农林牧渔指数几乎同步见底，个股并未走出独立行情。 水产养殖的价格周期由投苗量与消费需求共同决定（第 \ref{sec:other-livestock} 节），与猪周期的同步性弱于禽链，公司 2026 年 6 月末资产负债率 95.2\%，高于全样本中位数 52.2\%，其股价因此更多反映资金面与信用风险，而不是价格弹性本身。 综合区间形态看，该股单段涨跌幅度偏大、以趋势段为主（最大主升 +83.9\%、最大主跌 -55.7\%），年化波动率 44.3\% 与全样本中位数 37.3\% 相比波动明显大于样本中位数。

\pfl{盈利情况} 2026 年上半年营业收入 14.45 亿元（同比 -12.5\%），归母净利润 0.10 亿元，实现扭亏。 以年报为趋势对照，2023---2025 年营业收入由 46.17 亿元变为 31.90 亿元（两年年均复合增速 -16.9\%），归母净利润由 -5.32 亿元变为 -16.31 亿元。 按半年报口径，2026 年上半年的 ROE 为 10.0\%（未年化）、销售净利率 0.7\%、毛利率 11.4\%，ROE 在“水产养殖”4 家公司中按数值由高到低排第\zhnumber{2}位。 按同一口径，三级行业“水产养殖”4 家公司的半年 ROE 中位数为 6.05\%，该公司高于其中位数 3.96 个百分点。 针对上述压力，公司披露的应对举措包括：2026 年 6 月，存货规模由前期高位压降至 3.07 亿元（同比下降 80.78\%）。 综合看，公司在报告期维持盈利（高于全样本半年 ROE 中位数 0.75\%），利润的可持续性取决于水产养殖景气的延续与成本管控的执行。

\begin{center}\small\setlength{\tabcolsep}{3.4pt}
\captionof{table}[国联水产（300094）关键财务指标]{国联水产（300094）关键财务指标（合并报表口径；两个半年报期均未年化；现金短债比 $=$ 货币资金 $\div$ 短期债务）}
\begin{adjustbox}{max width=\textwidth}
\begin{tabular}{@{}lrrrrr@{}}
\toprule
指标 & 2023 年 & 2024 年 & 2025 年 & 2025 年半年报 & 2026 年半年报 \\
\midrule
营业收入（亿元） & 46.17 & 34.09 & 31.90 & 16.51 & 14.45 \\
归母净利润（亿元） & -5.32 & -7.42 & -16.31 & -5.40 & 0.10 \\
毛利率（\%） & 7.1 & 9.8 & -6.1 & -4.7 & 11.4 \\
ROE（\%） & -18.7 & -33.6 & -160.2 & -34.6 & 10.0 \\
资产负债率（\%） & 47.6 & 54.1 & 96.0 & 63.2 & 95.2 \\
经营活动现金流净额（亿元） & 1.71 & 0.29 & 1.06 & 0.98 & 1.92 \\
货币资金（亿元） & 4.82 & 3.78 & 2.58 & 3.05 & 1.98 \\
有息负债合计（亿元） & 16.79 & 14.52 & 13.52 & 14.45 & 10.84 \\
现金短债比（倍） & 0.34 & 0.29 & 0.20 & 0.23 & 0.19 \\
\bottomrule
\end{tabular}
\end{adjustbox}
\end{center}

\begin{center}
\begin{minipage}{\textwidth}\centering
\begin{tikzpicture}
\begin{groupplot}[
  group style={group size=2 by 1, horizontal sep=0.60cm},
  width=6.85cm, height=4.60cm,
  tick label style={font=\tiny},
  title style={font=\scriptsize\heiti},
  yticklabel style={font=\tiny},
  ylabel style={font=\tiny},
  grid=major, grid style={LightGray, line width=0.3pt},
  axis line style={InkGray, line width=0.5pt},
  enlarge x limits={abs=0.8},
  legend style={font=\scriptsize, at={(0.5,-0.20)}, anchor=north,
                legend columns=2, draw=none, fill=none, column sep=6pt},
]
\nextgroupplot[
  ybar, bar width=4pt,
  ymin=-22.79, ymax=56.27,
  xtick={0,1,2,3,4,5.7},
  xticklabels={2021,2022,2023,2024,2025,2026H1},
  ylabel={亿元},
]
\addplot[fill=AgriGreen!75, draw=AgriGreen, bar shift=-2.6pt] table[x=i, y=rev]{data/clean/profiles/fin_300094.dat};
\addplot[fill=SoftRed!60, draw=SoftRed, bar shift=2.6pt] table[x=i, y=ni]{data/clean/profiles/fin_300094.dat};
\addplot[fill=AgriGreen!30, draw=AgriGreen, pattern=north east lines, bar shift=-2.6pt] table[x=x, y=rev]{data/clean/profiles/finh_300094.dat};
\addplot[fill=SoftRed!20, draw=SoftRed, pattern=north east lines, bar shift=2.6pt] table[x=x, y=ni]{data/clean/profiles/finh_300094.dat};
\legend{营业收入（年/半年）, 归母净利润（年/半年）}
\nextgroupplot[
  ymin=-180.7, ymax=121.6,
  xtick={0,1,2,3,4,5.7},
  xticklabels={2021,2022,2023,2024,2025,2026H1},
  ylabel={\%},
]
\addplot[AgriGold, mark=*, mark size=1.3pt] table[x=i, y=roe]{data/clean/profiles/fin_300094.dat};
\addplot[OilBlue, mark=square*, mark size=1.3pt] table[x=i, y=debt]{data/clean/profiles/fin_300094.dat};
\addplot[only marks, AgriGold, mark=*, mark size=2.0pt] table[x=x, y=roe]{data/clean/profiles/finh_300094.dat};
\addplot[only marks, OilBlue, mark=square*, mark size=2.0pt] table[x=x, y=debt]{data/clean/profiles/finh_300094.dat};
\legend{ROE, 资产负债率}
\end{groupplot}
\end{tikzpicture}
\begin{tikzpicture}
\begin{axis}[
  name=finbot,
  width=13.75cm, height=3.10cm, scale only axis,
  ymin=-1.78, ymax=19.98,
  xtick={0,1,2,3,4,5.7},
  xticklabels={2021,2022,2023,2024,2025,2026H1},
  tick label style={font=\tiny},
  yticklabel style={font=\tiny},
  ylabel={亿元}, ylabel style={font=\tiny},
  ybar, bar width=4pt,
  grid=major, grid style={LightGray, line width=0.3pt},
  axis line style={InkGray, line width=0.5pt},
  enlarge x limits={abs=0.8},
  legend style={font=\scriptsize, at={(0.5,-0.26)}, anchor=north,
                legend columns=2, draw=none, fill=none, column sep=10pt},
]
\addplot[fill=AgriGreen!55, draw=AgriGreen, bar shift=-3.0pt] table[x=i, y=cash]{data/clean/profiles/fin_300094.dat};
\addplot[fill=SoftRed!45, draw=SoftRed, bar shift=3.0pt] table[x=i, y=ibd]{data/clean/profiles/fin_300094.dat};
\addplot[fill=AgriGreen!25, draw=AgriGreen, pattern=north east lines, bar shift=-3.0pt] table[x=x, y=cash]{data/clean/profiles/finh_300094.dat};
\addplot[fill=SoftRed!20, draw=SoftRed, pattern=north east lines, bar shift=3.0pt] table[x=x, y=ibd]{data/clean/profiles/finh_300094.dat};
\legend{货币资金（左轴）, 有息负债（左轴）}
\end{axis}
\begin{axis}[
  at={(finbot.south east)}, anchor=south east,
  width=13.75cm, height=3.10cm, scale only axis,
  axis y line*=right, axis x line*=none, xtick=\empty,
  ymin=-0.53, ymax=2.66,
  tick label style={font=\tiny},
  yticklabel style={font=\tiny}, scaled y ticks=false,
  legend style={font=\scriptsize, at={(0.5,-0.44)}, anchor=north,
                legend columns=1, draw=none, fill=none},
]
\addplot[OilBlue, mark=*, mark size=2.2pt, line width=1.3pt] table[x=i, y=ocf]{data/clean/profiles/fin_300094.dat};
\addplot[only marks, OilBlue, mark=*, mark size=3.2pt] table[x=x, y=ocf]{data/clean/profiles/finh_300094.dat};
\legend{经营活动现金流净额（右轴，亿元，蓝点）}
\end{axis}
\end{tikzpicture}

\captionof{figure}[国联水产（300094）近年主要财务指标]{国联水产（300094）近年主要财务指标（左上：营业收入与归母净利润；右上：ROE 与资产负债率；下：货币资金、有息负债为左轴柱状，经营活动现金流净额为其右轴的蓝点折线。
2021---2025 年为年报口径，2026H1 为 2026 年半年报/半年末，与其前一年报之间留出间隔、以斜纹或空心点区分；半年报数据未年化）}
\end{minipage}
\end{center}

\pfl{财务分析} 2026 年 6 月末资产负债率 95.2\%，较 2025 年末下降 0.8 个百分点（样本中位数 52.2\%，所处三级行业内按由低到高排第\zhnumber{3}位）。 2026 年 6 月末货币资金 1.98 亿元、短期债务（短期借款与一年内到期非流动负债）10.29 亿元、长期债务 0.56 亿元，有息负债合计 10.84 亿元，现金短债比 0.19 倍——缺口明显。 净有息负债 8.86 亿元，相当于其 2026 年上半年营业收入的 61\%。 流动比率 0.65、速动比率 0.47，短期偿付压力较大。 2026 年上半年经营活动现金流净额 1.92 亿元，经营现金流与利润同向为正，内生资金可以支撑运营。 综合看，账面货币资金对有息负债与短期债务的覆盖不足，滚动偿付对新增融资的依赖度较高；资产负债率高于样本中位数 43.0 个百分点；有息负债中短期债务占比更高，期限结构仍以滚动续作为主。 公司披露的股东与质押风险为：2026 年 5 月，控股股东关联方冠联国际投资有限公司及实际控制人李忠合计质押公司股份 5,896.68 万股；2026 年 9 月，第三方风险评级显示公司整体质押股份 2.08 亿股、占总股本 18.43\%。

\pfl{重大战略转型} 近三年公司的战略动作集中在：融资与资本运作（1 项）：2022 年 12 月，定增募资 10.00 亿元中的 7.00 亿元投向国美水产中央厨房项目等预制菜产能与渠道建设。 公告口径上，近三年（2023 年 9 月---2026 年 9 月）公司披露重大重组与并购类公告 5 条、对外投资与转型类公告 6 条、回购与股权激励类公告 29 条，合计公告 407 条。 综合判断，报告期内公司有 11 条与战略相关的公告落地（重大重组与并购 5 条、对外投资与转型 6 条），资本运作与主业调整之间存在直接关联。

\begin{center}\small\setlength{\tabcolsep}{4pt}
\captionof{table}[国联水产（300094）历史融资与资本运作]{国联水产（300094）历史融资与资本运作（据公司公告与公开报道整理；金额为披露值，含首次公开发行、再融资、债券、银行借款与重整投资等）}
\begin{adjustbox}{max width=\textwidth}
\begin{tabular}{@{}p{1.45cm}p{2.55cm}p{4.05cm}p{4.95cm}@{}}
\toprule
时间 & 方式 & 规模 & 用途与说明 \\
\midrule
2010-06-28 & IPO（深圳证券交易所创业板） & 发行 8,\allowbreak{}000 万股（占发行后总股本 25\%），\allowbreak{}发行价 14.38 元/\allowbreak{}股 & 罗非鱼加工厂建设项目（拟使用募集资金 7,\allowbreak{}923.56 万元）、\allowbreak{}水产品深加工建设项目、\allowbreak{}国内市场营销网络建设项目、\allowbreak{}水产品加工副产品综合利用建设项目、\allowbreak{}科研中心建设项目及其他与主营业务相关的营运资金；股票于 2010 年 7 月 8 日挂牌 \\
2022-12-12 & 向特定对象发行股票（定增） & 向 15 名特定对象发行 2.21 亿股，\allowbreak{}发行价 4.52 元/\allowbreak{}股 & 不超 10.00 亿元募资中 7.00 亿元投向全资子公司国美水产中央厨房项目、\allowbreak{}其余补充流动资金（后两个募投项目均已终止 \\
2025-12-09 & 出售全资子公司股权（关联交易、承债式收购） & 向控股股东新余国通投资管理有限公司转让全资子公司广东新盈食品科技有限公司 100\% 股权 & 盘活存量资产、\allowbreak{}将部分上游闲置资产及存货转为现金资产、\allowbreak{}优化资产结构与存货周转效率；经 2025 年 12 月 8 日第六届董事会第十七次会议审议通过（公告编号：2025-062） \\
2026-05-29 & 出售资产（孙公司资产打包出售） & 以人民币 1,\allowbreak{}600.00 万元将全资孙公司宜昌湘宜水产品有限公司位于湖北省宜昌市枝江市安福寺镇之字溪大道的工业用地使用权、\allowbreak{}地上建筑物、\allowbreak{}构筑物及厂内全部机器设备、\allowbreak{}配套设施一次性打包出售给湖北鸿展仓储服务有限公司 & 优化闲置资产、\allowbreak{}快速回笼资金改善现金流；经 2026 年 5 月 29 日第六届董事会第二十三次会议审议通过 \\
2026-07-15 & 出售境外房产资产 & 全资子公司 Liancheng Investments，\allowbreak{} LLC 以 1,\allowbreak{}700.00 万美元（折合人民币 1.15 亿元）向 LBA-Logistics 出售位于美国加利福尼亚州尤尼恩城法伯街 2910 号的房地产 & 交易所得款项用于公司业务发展；经 2026 年 7 月 15 日第六届董事会第二十四次会议审议通过 \\
\bottomrule
\end{tabular}
\end{adjustbox}
\end{center}

\pfl{历史融投资情况} 从现金流量表看，2025 年公司取得借款收到的现金 22.27 亿元、偿还债务支付的现金 23.79 亿元，筹资活动现金流净额 -2.76 亿元（2023 年为取得借款 22.56 亿元、偿还债务 23.45 亿元）。 2026 年上半年筹资活动现金流净额为 -1.99 亿元（取得借款 4.63 亿元、偿还债务 6.50 亿元，未年化），已转为净流出，处于净还债状态。 2025 年分配股利、利润或偿付利息支付的现金合计 0.53 亿元。 公开披露的融资记录共 5 笔，工具类型以 IPO、定向增发为主；金额与用途见下表。 其中规模最大的两笔为：2010 年 6 月，IPO，发行 8,000 万股（占发行后总股本 25\%），发行价 14.38 元/股；2026 年 7 月，出售境外房产资产，全资子公司 Liancheng Investments， LLC 以 1,700.00 万美元（折合人民币 1.15 亿元）向 LBA-Logistics 出售位于美国加利福尼亚州尤尼恩城法伯街 2910 号的房地产。 股本方面，近三年披露股本变动 5 次（限售股份上市 1 次、股份回购 1 次）；总股本由 2021 年末的 9.12 亿股变为 11.28 亿股（+23.66\%）。 分红方面，近三年未实施现金分红，在全部样本中属于分红能力偏弱的一端。 近三年“再融资”类公告 1 条，涉及定向增发。 综合看，公司融资结构上股权与债务融资并举；2025 年筹资活动现金流净额 -2.76 亿元，融资节奏仍以维持存量债务周转为主，与前述经营与投资现金流的方向基本一致。

\pfl{未来潜在融资需求分析} 2026 年 6 月末货币资金 1.98 亿元，而短期债务 10.29 亿元（现金短债比 0.19 倍），2026 年上半年经营现金流净额 1.92 亿元。按“短期债务减去账面货币资金”的滚动偿付口径，静态缺口约 8.30 亿元；上半年盈利或接近盈亏平衡，该缺口不随经营亏损进一步扩大。 公司 2026 年 6 月末资产负债率 95.2\%、2026 年上半年 ROE 10.0\%（未年化），资产负债率偏高（95.2\%），融资需求首先用于置换与滚动存量债务、补充营运资金，属于“补血型”而非扩张型。 截至目前公司仍处于预重整/重整程序中，重整投资人的招募与投资款到位是未来一年最主要的外部资金来源。 综合看，公司的外部资金需求以营运资金周转与在建项目投入为主，滚动偿付口径下的静态缺口约 8.30 亿元，融资节奏将受制于银行续贷意愿与股权融资窗口。




\subsubsection{*ST景谷（600265）}
\label{co:600265}

\pfl{公司基本情况} 云南景谷的林业与木材加工上市壳资源型公司（云南景谷林业股份有限公司），2000 年 8 月在上交所上市，2008 年后多次易主、三度被实施退市风险警示；2018 年 8 月周大福投资入主（最终控制方郑家纯），2023 年现金 2.703 亿元收购汇银木业 51\% 后业绩暴雷、2025 年以 13,336.60 万元将汇银木业 51\% 股权回售给控股股东，2025 年 10 月起无偿受赠算力服务公司博达数科 100\% 股权；2025 年净利润 -2.46 亿元且扣除后营业收入低于 3 亿元，2026 年 4 月 27 日起简称变更为“*ST 景谷”。 公司于 2000-08-25 在上交所首发上市，注册资本 1.30 亿元，总股本 1.30 亿股。 截至 2026 年 10 月 12 日收盘价 22.56 元，对应总市值 29.3 亿元，近一年+2.2\%，总市值在三级行业“林业Ⅲ”的 4 家公司中按由大到小排第\zhnumber{2}位，近一年涨跌幅高于全样本中位数（-10.4\%）12.6 个百分点。 按 2025 年年报披露的行业构成，收入占比前两位为制造业行业 71.7\%；其他(补充) 28.3\%。 与 2021 年年报相比，第一大行业口径由“林业行业”（94.4\%）变为“制造业行业”（71.7\%）；两期披露的分类口径有调整，占比差异中同时包含分类因素。

\pfl{股权结构与控制权} 按 2026 年 6 月 30 日的股东名册，前三大股东为周大福投资有限公司 持股 55.00\%；杭州磁晅沛曈投资管理合伙企业(有限合伙) 持股 9.31\%；谢志敏 持股 2.06\%，股东户数 4,365 户。 从控制权角度看，2026 年 4 月其将持有的博达数科 49\% 股权无偿赠与公司；触及退市风险警示情形，股票在 2025 年年度报告披露后被实施退市风险警示（简称前加“*ST”）；公司股票被继续实施其他风险警示（“ST”）。

\pfl{经营要点} 经营与产能方面，2023 年半年报显示公司股票简称为“ST 景谷”、变更前简称为“*ST 景谷”。

\pfl{所处行业的基本面} 林业的供给由林木蓄积与采伐限额决定、需求来自木材加工与基建，变现周期长（第 \ref{sec:framework} 章的供需框架）。 从供需两端的当期数据看，2026 年，1—2 月中国进口原木及锯材 750.3 万立方米、同比下降 11.2\%；报告期内人造板产品持续受房地产行业趋势、下游房地产建筑企业及家装行业购买力下降、行业竞争加剧、产品价格下滑等因素影响。

\pfl{竞争优势} 把该公司放回同类公司比较，2026 年上半年营业收入 0.6 亿元、在三级行业“林业Ⅲ”的 4 家公司中排第\zhnumber{3}位；毛利率 32.9\%、高于行业内中位数 10.5 个百分点；半年 ROE -0.0\%（未年化）、高于行业内中位数（-1.96\%）；资产负债率 26.7\%、低于行业内中位数（38.0\%）。 与同类公司相比，公司的差异化在于：公司所在地景谷县所处的普洱市共有林地面积 5,399.61 万亩、森林覆盖率 74.59\%、森林面积 4,955.82 万亩、森林蓄积量 3 亿立方米；销售分为由公司自行采伐的林木销售与转让活立木及林地使用权的活立木销售两种方式。 优势在报表上主要体现为高于行业平均的毛利率，这在周期底部是抵御亏损的主要缓冲。

\pfl{历史股价} 以 2021 年 1 月为起点、2026 年 10 月为终点，股价由 17.27 元变为 22.56 元，累计 +30.6\%，区间最高 24.45 元（2026 年 5 月）、最低 13.30 元（2022 年 6 月）；当前价相当于区间最高价的 92.3\%，为区间最低价的 1.70 倍。 月度收益的年化波动率 29.3\%，低于全样本中位数 37.3\%，在三级行业“林业Ⅲ”的 4 家公司中波动率由高到低排第\zhnumber{4}位。 把股价阶段与公司事件对齐看，2024 年 8 月 至 2026 年 5 月股价上涨+81.1\%，同期（2026 年 5 月）提示除已披露事项外不存在筹划重大资产重组、股份发行、破产重整、引进战略投资者等应披露而未披露的重大事项。 把这些阶段与公司事件对齐后可以看到，几轮大涨大跌多由公司自身的资本运作与风险事件触发，行业景气只解释了其中一部分。

\begin{center}
\begin{minipage}{\textwidth}\centering
\begin{tikzpicture}
\begin{axis}[
  width=12.6cm, height=3.9cm, scale only axis,
  xmin=2020.922, xmax=2026.972, ymin=12.37, ymax=26.16,
  xtick={2021,2022,2023,2024,2025,2026},
  yearticks,
  yticklabel style={font=\scriptsize},
  ylabel={元/股}, ylabel style={font=\scriptsize},
  grid=major, grid style={LightGray, line width=0.3pt},
  axis line style={InkGray, line width=0.5pt},
  every axis plot/.append style={line width=0.9pt},
]
\addplot[AgriGreen] table[x=x,y=y]{data/clean/profiles/px_600265.dat};
\addplot[SoftRed, dashed, line width=0.7pt] table[x=x,y=y]{data/clean/profiles/pxzz_600265.dat};
\addplot[only marks, mark=*, mark size=1.1pt, SoftRed] table[x=x,y=y]{data/clean/profiles/pxzz_600265.dat};
\end{axis}
\end{tikzpicture}

\captionof{figure}[*ST景谷（600265）股价与周期骨架]{*ST景谷（600265）月度收盘价（元/股）与 ZigZag 周期骨架（阈值 25\%，圆点为周期转折点）}
\end{minipage}
\end{center}

\pfl{过去周期分析} 以 25\% 的 ZigZag 阈值划分，样本区间内股价形成 3 个主升段与 3 个主跌段。 最大主升段出现在 2024 年 8 月 至 2026 年 5 月，21 个月+81.1\%；最大主跌段为 2021 年 8 月 至 2022 年 6 月，10 个月-43.9\%。 最近一次转折出现在 2026 年 5 月（上涨段自 2024 年 8 月起，21 个月+81.1\%），当前处于该段之后的震荡之中。 公司股价较申万农林牧渔指数提前约 47 个月见底，是板块中较早定价周期反转的公司之一。 林业的现金流依赖林木采伐与政策补贴，周期长、变现慢，公司 2026 年 6 月末资产负债率 26.7\%、半年 ROE -0.0\%（未年化），在本轮下行中属于随行业同步调整的一类。 综合区间形态看，该股单段涨跌幅度偏大、以趋势段为主（最大主升 +81.1\%、最大主跌 -43.9\%），年化波动率 29.3\% 与全样本中位数 37.3\% 相比波动小于样本中位数。

\pfl{盈利情况} 2026 年上半年营业收入 0.56 亿元（同比 -54.6\%），归母净利润 -0.00 亿元，亏损同比收窄 1.24 亿元。 以年报为趋势对照，2023---2025 年营业收入由 5.90 亿元变为 1.95 亿元（两年年均复合增速 -42.5\%），归母净利润由 0.06 亿元变为 -2.46 亿元。 按半年报口径，2026 年上半年的 ROE 为 -0.0\%（未年化）、销售净利率 5.4\%、毛利率 32.9\%，ROE 在“林业Ⅲ”4 家公司中按数值由高到低排第\zhnumber{1}位。 按同一口径，三级行业“林业Ⅲ”4 家公司的半年 ROE 中位数为 -1.96\%，该公司高于其中位数 1.95 个百分点。 综合看，公司仍处亏损区间、亏损同比收窄，半年 ROE 在三级行业内按由高到低排第\zhnumber{1}位，单位盈利能力的修复依赖行业景气与自身成本端的改善，报表弹性主要体现为价格与销量的方向性变化。

\begin{center}\small\setlength{\tabcolsep}{3.4pt}
\captionof{table}[*ST景谷（600265）关键财务指标]{*ST景谷（600265）关键财务指标（合并报表口径；两个半年报期均未年化；现金短债比 $=$ 货币资金 $\div$ 短期债务）}
\begin{adjustbox}{max width=\textwidth}
\begin{tabular}{@{}lrrrrr@{}}
\toprule
指标 & 2023 年 & 2024 年 & 2025 年 & 2025 年半年报 & 2026 年半年报 \\
\midrule
营业收入（亿元） & 5.90 & 4.47 & 1.95 & 1.23 & 0.56 \\
归母净利润（亿元） & 0.06 & -0.73 & -2.46 & -1.24 & -0.00 \\
毛利率（\%） & 10.3 & 2.3 & 0.1 & -11.9 & 32.9 \\
ROE（\%） & 4.0 & -56.1 & -4.7 & -1.3 & -0.0 \\
资产负债率（\%） & 63.4 & 73.5 & 29.9 & 77.7 & 26.7 \\
经营活动现金流净额（亿元） & 0.92 & 0.03 & -0.26 & -0.08 & 0.29 \\
货币资金（亿元） & 0.59 & 0.40 & 0.71 & 0.25 & 0.92 \\
有息负债合计（亿元） & 2.03 & 2.28 & 0.05 & 2.25 & 0.04 \\
现金短债比（倍） & 2.18 & 0.53 & 34.26 & 0.11 & 55.11 \\
\bottomrule
\end{tabular}
\end{adjustbox}
\end{center}

\begin{center}
\begin{minipage}{\textwidth}\centering
\begin{tikzpicture}
\begin{groupplot}[
  group style={group size=2 by 1, horizontal sep=0.60cm},
  width=6.85cm, height=4.60cm,
  tick label style={font=\tiny},
  title style={font=\scriptsize\heiti},
  yticklabel style={font=\tiny},
  ylabel style={font=\tiny},
  grid=major, grid style={LightGray, line width=0.3pt},
  axis line style={InkGray, line width=0.5pt},
  enlarge x limits={abs=0.8},
  legend style={font=\scriptsize, at={(0.5,-0.20)}, anchor=north,
                legend columns=2, draw=none, fill=none, column sep=6pt},
]
\nextgroupplot[
  ybar, bar width=4pt,
  ymin=-3.29, ymax=6.90,
  xtick={0,1,2,3,4,5.7},
  xticklabels={2021,2022,2023,2024,2025,2026H1},
  ylabel={亿元},
]
\addplot[fill=AgriGreen!75, draw=AgriGreen, bar shift=-2.6pt] table[x=i, y=rev]{data/clean/profiles/fin_600265.dat};
\addplot[fill=SoftRed!60, draw=SoftRed, bar shift=2.6pt] table[x=i, y=ni]{data/clean/profiles/fin_600265.dat};
\addplot[fill=AgriGreen!30, draw=AgriGreen, pattern=north east lines, bar shift=-2.6pt] table[x=x, y=rev]{data/clean/profiles/finh_600265.dat};
\addplot[fill=SoftRed!20, draw=SoftRed, pattern=north east lines, bar shift=2.6pt] table[x=x, y=ni]{data/clean/profiles/finh_600265.dat};
\legend{营业收入（年/半年）, 归母净利润（年/半年）}
\nextgroupplot[
  ymin=-113.3, ymax=90.8,
  xtick={0,1,2,3,4,5.7},
  xticklabels={2021,2022,2023,2024,2025,2026H1},
  ylabel={\%},
]
\addplot[AgriGold, mark=*, mark size=1.3pt] table[x=i, y=roe]{data/clean/profiles/fin_600265.dat};
\addplot[OilBlue, mark=square*, mark size=1.3pt] table[x=i, y=debt]{data/clean/profiles/fin_600265.dat};
\addplot[only marks, AgriGold, mark=*, mark size=2.0pt] table[x=x, y=roe]{data/clean/profiles/finh_600265.dat};
\addplot[only marks, OilBlue, mark=square*, mark size=2.0pt] table[x=x, y=debt]{data/clean/profiles/finh_600265.dat};
\legend{ROE, 资产负债率}
\end{groupplot}
\end{tikzpicture}
\begin{tikzpicture}
\begin{axis}[
  name=finbot,
  width=13.75cm, height=3.10cm, scale only axis,
  ymin=-0.23, ymax=2.55,
  xtick={0,1,2,3,4,5.7},
  xticklabels={2021,2022,2023,2024,2025,2026H1},
  tick label style={font=\tiny},
  yticklabel style={font=\tiny},
  ylabel={亿元}, ylabel style={font=\tiny},
  ybar, bar width=4pt,
  grid=major, grid style={LightGray, line width=0.3pt},
  axis line style={InkGray, line width=0.5pt},
  enlarge x limits={abs=0.8},
  legend style={font=\scriptsize, at={(0.5,-0.26)}, anchor=north,
                legend columns=2, draw=none, fill=none, column sep=10pt},
]
\addplot[fill=AgriGreen!55, draw=AgriGreen, bar shift=-3.0pt] table[x=i, y=cash]{data/clean/profiles/fin_600265.dat};
\addplot[fill=SoftRed!45, draw=SoftRed, bar shift=3.0pt] table[x=i, y=ibd]{data/clean/profiles/fin_600265.dat};
\addplot[fill=AgriGreen!25, draw=AgriGreen, pattern=north east lines, bar shift=-3.0pt] table[x=x, y=cash]{data/clean/profiles/finh_600265.dat};
\addplot[fill=SoftRed!20, draw=SoftRed, pattern=north east lines, bar shift=3.0pt] table[x=x, y=ibd]{data/clean/profiles/finh_600265.dat};
\legend{货币资金（左轴）, 有息负债（左轴）}
\end{axis}
\begin{axis}[
  at={(finbot.south east)}, anchor=south east,
  width=13.75cm, height=3.10cm, scale only axis,
  axis y line*=right, axis x line*=none, xtick=\empty,
  ymin=-0.86, ymax=1.35,
  tick label style={font=\tiny},
  yticklabel style={font=\tiny}, scaled y ticks=false,
  legend style={font=\scriptsize, at={(0.5,-0.44)}, anchor=north,
                legend columns=1, draw=none, fill=none},
]
\addplot[OilBlue, mark=*, mark size=2.2pt, line width=1.3pt] table[x=i, y=ocf]{data/clean/profiles/fin_600265.dat};
\addplot[only marks, OilBlue, mark=*, mark size=3.2pt] table[x=x, y=ocf]{data/clean/profiles/finh_600265.dat};
\legend{经营活动现金流净额（右轴，亿元，蓝点）}
\end{axis}
\end{tikzpicture}

\captionof{figure}[*ST景谷（600265）近年主要财务指标]{*ST景谷（600265）近年主要财务指标（左上：营业收入与归母净利润；右上：ROE 与资产负债率；下：货币资金、有息负债为左轴柱状，经营活动现金流净额为其右轴的蓝点折线。
2021---2025 年为年报口径，2026H1 为 2026 年半年报/半年末，与其前一年报之间留出间隔、以斜纹或空心点区分；半年报数据未年化）}
\end{minipage}
\end{center}

\pfl{财务分析} 2026 年 6 月末资产负债率 26.7\%，较 2025 年末下降 3.2 个百分点（样本中位数 52.2\%，所处三级行业内按由低到高排第\zhnumber{1}位）。 2026 年 6 月末货币资金 0.92 亿元、短期债务（短期借款与一年内到期非流动负债）0.02 亿元、长期债务 0.02 亿元，有息负债合计 0.04 亿元，现金短债比 55.11 倍——覆盖充分。 净有息负债 -0.87 亿元，相当于其 2026 年上半年营业收入的 -156\%。 流动比率 3.58、速动比率 1.72，短期指标尚可。 2026 年上半年经营活动现金流净额 0.29 亿元，净利润为负而经营现金流为正，差额主要来自折旧摊销与营运资本变动，说明亏损尚未完全转化为现金流出。 综合看，现金短债比 55.11 倍，短期偿付压力可控；资产负债率低于样本中位数 25.5 个百分点；有息负债中长期债务占比更高，期限结构对短债滚动的压力较小。 公司披露的退市与持续经营风险为：2026 年 4 月，公司股票被实施退市风险警示；2025 年 8 月，2024 年营业收入占比达 87.02\% 本次出售完成后主营业务规模将急剧下降。

\pfl{重大战略转型} 近三年公司的战略动作集中在：并购与重组（4 项）：2025 年 12 月，公司完成向控股股东出售汇银木业 51\% 股权的重大资产出售；2023 年 2 月，公司以现金 2.70 亿元收购上述项目权；其他动作（1 项）：2026 年 4 月，控股股东无偿赠与上海博达数智科技有限公司 51\% 股权（2026 年 4 月再受赠剩余 49\%）。 公告口径上，近三年（2023 年 9 月---2026 年 9 月）公司披露重大重组与并购类公告 23 条、对外投资与转型类公告 1 条、回购与股权激励类公告 0 条，合计公告 481 条。 主营结构的口径变化已经落到报表上：2021 年年报的第一大行业为“林业行业”（94.4\%），2025 年年报为“制造业行业”（71.7\%）；公司在这两期的分部划分方式有过调整。 综合判断，报告期内公司有 24 条与战略相关的公告落地（重大重组与并购 23 条、对外投资与转型 1 条），资本运作与主业调整之间存在直接关联。

\begin{center}\small\setlength{\tabcolsep}{4pt}
\captionof{table}[*ST景谷（600265）历史融资与资本运作]{*ST景谷（600265）历史融资与资本运作（据公司公告与公开报道整理；金额为披露值，含首次公开发行、再融资、债券、银行借款与重整投资等）}
\begin{adjustbox}{max width=\textwidth}
\begin{tabular}{@{}p{1.45cm}p{2.55cm}p{4.05cm}p{4.95cm}@{}}
\toprule
时间 & 方式 & 规模 & 用途与说明 \\
\midrule
2000-07-21 & IPO（上海证券交易所） & 发行 4,\allowbreak{}000 万股，\allowbreak{}发行价 5.19 元/\allowbreak{}股，\allowbreak{}募集资金总额 1.96 亿元 & 经证监发行字[2000]99 号文核准，\allowbreak{}发行后注册资本 1.05 亿元 \\
2023-02 & 现金收购重大资产（并购） & 以支付现金方式购买崔会军、\allowbreak{}王兰存、\allowbreak{}京保基金、\allowbreak{}技改基金合计持有的唐县汇银木业有限公司 51\% 股权 & 新增人造板制造业务、\allowbreak{}改善上市公司盈利能力；评估机构亚超评估采用收益法评估结论（北京亚超评报字（2023）第 A003 号），\allowbreak{}增值率 83.7\% \\
2025-08-16 & 重大资产出售暨关联交易（向控股股东出售子公司股权） & 以现金交易方式向控股股东周大福投资有限公司转让持有的汇银木业 51\% 股权 & 剥离严重亏损且面临诉讼风险的子公司、\allowbreak{}化解风险、\allowbreak{}为业务转型打基础；2025 年 8 月 16 日披露筹划出售资产涉及重大资产重组（公告编号：2025-066） \\
2025-10 & 无偿受赠资产（控股股东资产赠与） & 控股股东周大福投资有限公司将其持有的上海博达数智科技有限公司 51\% 股权无偿赠与公司 & 提高公司持续经营能力、\allowbreak{}优化业务结构，\allowbreak{}注入算力服务业务；2026 年 4 月控股股东将博达数科剩余 49\% 股权无偿赠与公司 \\
2025-12-05 & 关联方资金占用清收（重大资产出售配套安排） & 截至上交所问询函回复出具日，\allowbreak{}标的公司汇银木业应偿还上市公司拆借资金的本金及利息为 6,\allowbreak{}172.22 万元 & 明确标的公司对上市公司借款的本息偿付安排；公司回复上海证券交易所云南景谷林业股份有限公司重大资产重组草案的信息披露的问询函 \\
\bottomrule
\end{tabular}
\end{adjustbox}
\end{center}

\pfl{历史融投资情况} 从现金流量表看，2025 年公司取得借款收到的现金 0.23 亿元、偿还债务支付的现金 0.70 亿元、吸收权益性投资 1.95 亿元，筹资活动现金流净额 0.39 亿元（2023 年为取得借款 2.00 亿元、偿还债务 0.06 亿元）。 2026 年上半年筹资活动现金流净额为 -0.04 亿元（取得借款 — 亿元、偿还债务 — 亿元，未年化），已转为净流出，处于净还债状态。 2025 年分配股利、利润或偿付利息支付的现金合计 0.17 亿元。 公开披露的融资记录共 5 笔，工具类型以 IPO 为主；金额与用途见下表。 其中规模最大的两笔为：2000 年 7 月，IPO，发行 4,000 万股，发行价 5.19 元/股，募集资金总额 1.96 亿元；2025 年 12 月，关联方资金占用清收，截至上交所问询函回复出具日，标的公司汇银木业应偿还上市公司拆借资金的本金及利息为 6,172.22 万元。 股本方面，近三年披露股本变动 3 次；总股本自 2021 年末以来未发生变化（1.30 亿股）。 分红方面，近三年未实施现金分红，在全部样本中属于分红能力偏弱的一端。 近三年“再融资”类公告 16 条，涉及定向增发、募集资金使用。 综合看，公司融资结构上股权与债务融资并举；2025 年筹资活动现金流净额 0.39 亿元，年度内仍为净融入，与前述经营与投资现金流的方向基本一致。

\pfl{未来潜在融资需求分析} 2026 年 6 月末货币资金 0.92 亿元，而短期债务 0.02 亿元（现金短债比 55.11 倍），2026 年上半年经营现金流净额 0.29 亿元。账面货币资金可以覆盖短期债务，缺口不体现在静态偿付层面。 按第 \ref{sec:financing-need} 节的三档划分，公司属于第二档“保壳型”：近三年重大重组与并购类公告 23 条，2026 年上半年归母净利润 -0.00 亿元，融资需求更多体现为“资产置换 + 维持上市地位”，而不是扩产。 公开披露的后续资金安排线索包括：资本开支安排：2026 年半年报，提出“立足林业主业精益管理+完成博达数科剩余 49\% 股权受赠”的工作主线。 静态偿付不存在缺口，后续融资更可能服务于产能或并购安排，而不是填补流动性窟窿。




\subsubsection{福建金森（002679）}
\label{co:002679}

\pfl{公司基本情况} 福建三明将乐县的国有控股林业上市公司（福建金森林业股份有限公司），2012 年 6 月在深交所中小板上市，是全国首家纯森林资源培育型国有控股上市公司、南方集体林区 FSC 国际森林认证用材林蓄积量最大的企业；主营森林经营与木材采运，2025 年营业收入 1.487 亿元、归母净利润 923.06 万元，2026 年上半年受季节性因素影响归母净利润 -1,358.18 万元（同比减亏 30.03\%），近年无新增股权融资、以存货抵押贷款为主要融资方式。 公司于 2012-06-05 在深交所主板首发上市，注册资本 2.36 亿元，总股本 2.36 亿股。 截至 2026 年 10 月 12 日收盘价 11.15 元，对应总市值 26.3 亿元，近一年+3.7\%，总市值在三级行业“林业Ⅲ”的 4 家公司中按由大到小排第\zhnumber{3}位，近一年涨跌幅高于全样本中位数（-10.4\%）14.1 个百分点。 按 2026 年半年报披露的行业构成，收入占比前三位为林业 74.4\%；生物质燃料 19.2\%；木材加工 4.3\%。 股权结构方面，控股股东为福建金森集团有限公司（持股 63.63\%）；实际控制人为将乐县财政局。

\pfl{股权结构与控制权} 按 2026 年 6 月 30 日的股东名册，前三大股东为福建金森集团有限公司 持股 63.63\%；将乐县交通基础设施建设有限公司 持股 1.74\%；高盛公司有限责任公司 持股 0.74\%，股东户数 17,247 户。 从控制权角度看，持有公司 63.63\% 股份，合计控制公司 63.83\% 股份。

\pfl{经营要点} 经营与产能方面，公司存货账面价值 7.84 亿元全部用于抵押（抵押贷款）。

\pfl{所处行业的基本面} 林业的供给由林木蓄积与采伐限额决定、需求来自木材加工与基建，变现周期长（第 \ref{sec:framework} 章的供需框架）。 从供需两端的当期数据看，2026 年，1—2 月中国进口原木及锯材 750.3 万立方米、同比下降 11.2\%；2026 年 7 月，前 7 个月进口总量同比下降 6.7\%，进口低迷压制国内木材价格；中国人民银行、金融监管总局、国家林草局金融支持林业高质量发展提出加大林权抵押贷款产品创新、提升国家储备林建设金融支持能力并拓展碳汇收益融资。

\pfl{竞争优势} 把该公司放回同类公司比较，2026 年上半年营业收入 0.5 亿元、在三级行业“林业Ⅲ”的 4 家公司中排第\zhnumber{4}位；毛利率 40.5\%、高于行业内中位数 18.1 个百分点；半年 ROE -1.8\%（未年化）、高于行业内中位数（-1.96\%）；资产负债率 61.9\%、高于行业内中位数（38.0\%）。 与同类公司相比，公司的差异化在于：公司为全国首家纯森林资源培育型国有控股上市公司；林木种植经营历史悠久、成本较低，由国有资本控股在收购国有林时处于有利地位；公司通过 FSC-FM/COC 森林可持续经营认证（认证机构为 SGS），并设立金森碳汇科技有限公司发展林草碳汇服务及其他林业技术服务。 优势在报表上主要体现为高于行业平均的毛利率，这在周期底部是抵御亏损的主要缓冲。

\pfl{历史股价} 以 2021 年 1 月为起点、2026 年 10 月为终点，股价由 9.90 元变为 11.15 元，累计 +12.6\%，区间最高 24.22 元（2021 年 5 月）、最低 6.67 元（2024 年 8 月）；当前价相当于区间最高价的 46.0\%，为区间最低价的 1.67 倍。 月度收益的年化波动率 55.4\%，高于全样本中位数 37.3\%，在三级行业“林业Ⅲ”的 4 家公司中波动率由高到低排第\zhnumber{2}位。 把股价阶段与公司事件对齐看，2024 年 8 月 至 2025 年 10 月股价上涨+109.9\%，同期（2025 年 6 月）同一报告期受限资产还包括 4,250 元 ETC 保证金冻结货币资金。 整体上看，区间的几处大级别拐点都能在公司的资本运作、股权变动或风险事件上找到对应，股价波动有明显的公司层面驱动，而非单纯的行业 beta。

\begin{center}
\begin{minipage}{\textwidth}\centering
\begin{tikzpicture}
\begin{axis}[
  width=12.6cm, height=3.9cm, scale only axis,
  xmin=2020.922, xmax=2026.972, ymin=6.20, ymax=25.92,
  xtick={2021,2022,2023,2024,2025,2026},
  yearticks,
  yticklabel style={font=\scriptsize},
  ylabel={元/股}, ylabel style={font=\scriptsize},
  grid=major, grid style={LightGray, line width=0.3pt},
  axis line style={InkGray, line width=0.5pt},
  every axis plot/.append style={line width=0.9pt},
]
\addplot[AgriGreen] table[x=x,y=y]{data/clean/profiles/px_002679.dat};
\addplot[SoftRed, dashed, line width=0.7pt] table[x=x,y=y]{data/clean/profiles/pxzz_002679.dat};
\addplot[only marks, mark=*, mark size=1.1pt, SoftRed] table[x=x,y=y]{data/clean/profiles/pxzz_002679.dat};
\end{axis}
\end{tikzpicture}

\captionof{figure}[福建金森（002679）股价与周期骨架]{福建金森（002679）月度收盘价（元/股）与 ZigZag 周期骨架（阈值 25\%，圆点为周期转折点）}
\end{minipage}
\end{center}

\pfl{过去周期分析} 以 25\% 的 ZigZag 阈值划分，样本区间内股价形成 4 个主升段与 3 个主跌段。 最大主跌段为 2021 年 5 月 至 2022 年 4 月，11 个月-62.3\%；最大主升段出现在 2021 年 1 月 至 2021 年 5 月，4 个月+144.6\%。 最近一次转折出现在 2026 年 8 月（上涨段自 2026 年 6 月起，2 个月+49.0\%），当前处于该段之后的震荡之中。 公司股价较申万农林牧渔指数提前约 21 个月见底，是板块中较早定价周期反转的公司之一。 林业的现金流依赖林木采伐与政策补贴，周期长、变现慢，公司 2026 年 6 月末资产负债率 61.9\%、半年 ROE -1.8\%（未年化），在本轮下行中属于随行业同步调整的一类。 综合区间形态看，该股单段涨跌幅度偏大、以趋势段为主（最大主升 +144.6\%、最大主跌 -62.3\%），年化波动率 55.4\% 与全样本中位数 37.3\% 相比波动明显大于样本中位数。

\pfl{盈利情况} 2026 年上半年营业收入 0.45 亿元（同比 -9.5\%），归母净利润 -0.14 亿元，亏损同比收窄 0.06 亿元。 以年报为趋势对照，2023---2025 年营业收入由 1.48 亿元变为 1.49 亿元（两年年均复合增速 +0.4\%），归母净利润由 0.08 亿元变为 0.09 亿元。 按半年报口径，2026 年上半年的 ROE 为 -1.8\%（未年化）、销售净利率 -29.5\%、毛利率 40.5\%，ROE 在“林业Ⅲ”4 家公司中按数值由高到低排第\zhnumber{2}位。 按同一口径，三级行业“林业Ⅲ”4 家公司的半年 ROE 中位数为 -1.96\%，该公司高于其中位数 0.16 个百分点。 综合看，公司仍处亏损区间、亏损同比收窄，半年 ROE 在三级行业内按由高到低排第\zhnumber{2}位，单位盈利能力的修复依赖行业景气与自身成本端的改善，报表弹性主要体现为价格与销量的方向性变化。

\begin{center}\small\setlength{\tabcolsep}{3.4pt}
\captionof{table}[福建金森（002679）关键财务指标]{福建金森（002679）关键财务指标（合并报表口径；两个半年报期均未年化；现金短债比 $=$ 货币资金 $\div$ 短期债务）}
\begin{adjustbox}{max width=\textwidth}
\begin{tabular}{@{}lrrrrr@{}}
\toprule
指标 & 2023 年 & 2024 年 & 2025 年 & 2025 年半年报 & 2026 年半年报 \\
\midrule
营业收入（亿元） & 1.48 & 1.51 & 1.49 & 0.50 & 0.45 \\
归母净利润（亿元） & 0.08 & 0.10 & 0.09 & -0.19 & -0.14 \\
毛利率（\%） & 52.7 & 57.1 & 38.5 & 38.5 & 40.5 \\
ROE（\%） & 1.1 & 1.4 & 1.2 & -2.6 & -1.8 \\
资产负债率（\%） & 62.8 & 61.2 & 61.7 & 62.6 & 61.9 \\
经营活动现金流净额（亿元） & 0.70 & 0.77 & 0.71 & -0.33 & -0.46 \\
货币资金（亿元） & 2.55 & 1.83 & 1.83 & 1.53 & 1.34 \\
有息负债合计（亿元） & 12.12 & 11.31 & 11.32 & 11.66 & 11.56 \\
现金短债比（倍） & 1.49 & 1.04 & 1.04 & 0.90 & 0.77 \\
\bottomrule
\end{tabular}
\end{adjustbox}
\end{center}

\begin{center}
\begin{minipage}{\textwidth}\centering
\begin{tikzpicture}
\begin{groupplot}[
  group style={group size=2 by 1, horizontal sep=0.60cm},
  width=6.85cm, height=4.60cm,
  tick label style={font=\tiny},
  title style={font=\scriptsize\heiti},
  yticklabel style={font=\tiny},
  ylabel style={font=\tiny},
  grid=major, grid style={LightGray, line width=0.3pt},
  axis line style={InkGray, line width=0.5pt},
  enlarge x limits={abs=0.8},
  legend style={font=\scriptsize, at={(0.5,-0.20)}, anchor=north,
                legend columns=2, draw=none, fill=none, column sep=6pt},
]
\nextgroupplot[
  ybar, bar width=4pt,
  ymin=-0.34, ymax=2.14,
  xtick={0,1,2,3,4,5.7},
  xticklabels={2021,2022,2023,2024,2025,2026H1},
  ylabel={亿元},
]
\addplot[fill=AgriGreen!75, draw=AgriGreen, bar shift=-2.6pt] table[x=i, y=rev]{data/clean/profiles/fin_002679.dat};
\addplot[fill=SoftRed!60, draw=SoftRed, bar shift=2.6pt] table[x=i, y=ni]{data/clean/profiles/fin_002679.dat};
\addplot[fill=AgriGreen!30, draw=AgriGreen, pattern=north east lines, bar shift=-2.6pt] table[x=x, y=rev]{data/clean/profiles/finh_002679.dat};
\addplot[fill=SoftRed!20, draw=SoftRed, pattern=north east lines, bar shift=2.6pt] table[x=x, y=ni]{data/clean/profiles/finh_002679.dat};
\legend{营业收入（年/半年）, 归母净利润（年/半年）}
\nextgroupplot[
  ymin=-7.0, ymax=69.3,
  xtick={0,1,2,3,4,5.7},
  xticklabels={2021,2022,2023,2024,2025,2026H1},
  ylabel={\%},
]
\addplot[AgriGold, mark=*, mark size=1.3pt] table[x=i, y=roe]{data/clean/profiles/fin_002679.dat};
\addplot[OilBlue, mark=square*, mark size=1.3pt] table[x=i, y=debt]{data/clean/profiles/fin_002679.dat};
\addplot[only marks, AgriGold, mark=*, mark size=2.0pt] table[x=x, y=roe]{data/clean/profiles/finh_002679.dat};
\addplot[only marks, OilBlue, mark=square*, mark size=2.0pt] table[x=x, y=debt]{data/clean/profiles/finh_002679.dat};
\legend{ROE, 资产负债率}
\end{groupplot}
\end{tikzpicture}
\begin{tikzpicture}
\begin{axis}[
  name=finbot,
  width=13.75cm, height=3.10cm, scale only axis,
  ymin=-1.21, ymax=13.58,
  xtick={0,1,2,3,4,5.7},
  xticklabels={2021,2022,2023,2024,2025,2026H1},
  tick label style={font=\tiny},
  yticklabel style={font=\tiny},
  ylabel={亿元}, ylabel style={font=\tiny},
  ybar, bar width=4pt,
  grid=major, grid style={LightGray, line width=0.3pt},
  axis line style={InkGray, line width=0.5pt},
  enlarge x limits={abs=0.8},
  legend style={font=\scriptsize, at={(0.5,-0.26)}, anchor=north,
                legend columns=2, draw=none, fill=none, column sep=10pt},
]
\addplot[fill=AgriGreen!55, draw=AgriGreen, bar shift=-3.0pt] table[x=i, y=cash]{data/clean/profiles/fin_002679.dat};
\addplot[fill=SoftRed!45, draw=SoftRed, bar shift=3.0pt] table[x=i, y=ibd]{data/clean/profiles/fin_002679.dat};
\addplot[fill=AgriGreen!25, draw=AgriGreen, pattern=north east lines, bar shift=-3.0pt] table[x=x, y=cash]{data/clean/profiles/finh_002679.dat};
\addplot[fill=SoftRed!20, draw=SoftRed, pattern=north east lines, bar shift=3.0pt] table[x=x, y=ibd]{data/clean/profiles/finh_002679.dat};
\legend{货币资金（左轴）, 有息负债（左轴）}
\end{axis}
\begin{axis}[
  at={(finbot.south east)}, anchor=south east,
  width=13.75cm, height=3.10cm, scale only axis,
  axis y line*=right, axis x line*=none, xtick=\empty,
  ymin=-0.80, ymax=1.22,
  tick label style={font=\tiny},
  yticklabel style={font=\tiny}, scaled y ticks=false,
  legend style={font=\scriptsize, at={(0.5,-0.44)}, anchor=north,
                legend columns=1, draw=none, fill=none},
]
\addplot[OilBlue, mark=*, mark size=2.2pt, line width=1.3pt] table[x=i, y=ocf]{data/clean/profiles/fin_002679.dat};
\addplot[only marks, OilBlue, mark=*, mark size=3.2pt] table[x=x, y=ocf]{data/clean/profiles/finh_002679.dat};
\legend{经营活动现金流净额（右轴，亿元，蓝点）}
\end{axis}
\end{tikzpicture}

\captionof{figure}[福建金森（002679）近年主要财务指标]{福建金森（002679）近年主要财务指标（左上：营业收入与归母净利润；右上：ROE 与资产负债率；下：货币资金、有息负债为左轴柱状，经营活动现金流净额为其右轴的蓝点折线。
2021---2025 年为年报口径，2026H1 为 2026 年半年报/半年末，与其前一年报之间留出间隔、以斜纹或空心点区分；半年报数据未年化）}
\end{minipage}
\end{center}

\pfl{财务分析} 2026 年 6 月末资产负债率 61.9\%，较 2025 年末上升 0.2 个百分点（样本中位数 52.2\%，所处三级行业内按由低到高排第\zhnumber{4}位）。 2026 年 6 月末货币资金 1.34 亿元、短期债务（短期借款与一年内到期非流动负债）1.75 亿元、长期债务 9.81 亿元，有息负债合计 11.56 亿元，现金短债比 0.77 倍——覆盖偏紧。 净有息负债 10.22 亿元，相当于其 2026 年上半年营业收入的 2,247\%。 流动比率 7.63、速动比率 1.15，短期指标尚可。 2026 年上半年经营活动现金流净额 -0.46 亿元，经营现金流与利润同时为负，现金消耗较快。 综合看，现金短债比 0.77 倍，短期资金安排偏紧；资产负债率高于样本中位数 9.7 个百分点；有息负债中长期债务占比更高，期限结构对短债滚动的压力较小。 公司披露的风险集中在监管与合规、股东与质押两类：一是 2025 年，公司未在 2025 年年度报告中披露重大诉讼、仲裁及处罚事项；二是 2026 年，一季报显示控股股东福建金森集团有限公司持股 1.50 亿股、占 63.63\%。

\pfl{重大战略转型} 近三年公司的战略动作集中在：其他动作（3 项）：2023 年 4 月，公司对外投资设立全资子公司及注销分公司；2023 年 2 月，公司对外投资暨设立控股子公司。 公告口径上，近三年（2023 年 9 月---2026 年 9 月）公司披露重大重组与并购类公告 0 条、对外投资与转型类公告 5 条、回购与股权激励类公告 0 条，合计公告 235 条。 第一大行业“林业”的收入占比由 2021 年年报的 98.3\% 变为 74.4\%（23.8 个百分点），收入结构出现再平衡。 综合判断，报告期内公司有 5 条与战略相关的公告落地（重大重组与并购 0 条、对外投资与转型 5 条），资本运作与主业调整之间存在直接关联。

\begin{center}\small\setlength{\tabcolsep}{4pt}
\captionof{table}[福建金森（002679）历史融资与资本运作]{福建金森（002679）历史融资与资本运作（据公司公告与公开报道整理；金额为披露值，含首次公开发行、再融资、债券、银行借款与重整投资等）}
\begin{adjustbox}{max width=\textwidth}
\begin{tabular}{@{}p{1.45cm}p{2.55cm}p{4.05cm}p{4.95cm}@{}}
\toprule
时间 & 方式 & 规模 & 用途与说明 \\
\midrule
2012-06-05 & IPO（深圳证券交易所中小板） & 发行 3,\allowbreak{}468 万股，\allowbreak{}发行价 12.00 元/\allowbreak{}股，\allowbreak{}募集资金总额 4.16 亿元 & 商品材基地建设林木资源资产并购项目（原计划投入 2.57 亿元 \\
2015-12-02 & 资产转让（向控股股东出售资产） & 公司与控股股东福建金森集团有限公司签订资产转让协议（附生效条款），\allowbreak{}交易金额 8,\allowbreak{}423.66 万元 & 优化资产结构、\allowbreak{}收缩非核心经营资产；该公告于 2015 年 12 月 4 日披露；交易对手为公司控股股东 \\
2023-12-31 & 现金分红 & 以 2.36 亿股为基数，\allowbreak{}向全体股东每 10 股派发现金红利 0.26 元（含税） & 回报股东；经公司董事会审议通过的 2023 年度利润分配预案 \\
2025-06-30 & 存货抵押贷款（银行融资） & 账面价值 7.84 亿元的存货全部抵押用于银行借款 & 生产经营周转资金；同一报告期受限资产还包括 4,\allowbreak{}250 元 ETC 保证金冻结货币资金；公司未披露具体授信额度 \\
\bottomrule
\end{tabular}
\end{adjustbox}
\end{center}

\pfl{历史融投资情况} 从现金流量表看，2025 年公司取得借款收到的现金 1.80 亿元、偿还债务支付的现金 1.79 亿元，筹资活动现金流净额 -0.48 亿元（2023 年为取得借款 5.58 亿元、偿还债务 4.19 亿元）。 2026 年上半年筹资活动现金流净额为 0.04 亿元（取得借款 1.00 亿元、偿还债务 0.77 亿元，未年化），仍处于净融入状态。 2025 年分配股利、利润或偿付利息支付的现金合计 0.49 亿元。 公开披露的融资记录共 4 笔，工具类型以 IPO、银行借款为主；金额与用途见下表。 其中规模最大的两笔为：2012 年 6 月，IPO，发行 3,468 万股，发行价 12.00 元/股，募集资金总额 4.16 亿元；2015 年 12 月，资产转让，公司与控股股东福建金森集团有限公司签订资产转让协议（附生效条款），交易金额 8,423.66 万元。 股本方面，近三年披露股本变动 3 次；总股本自 2021 年末以来未发生变化（2.36 亿股）。 分红方面，近三年现金分红 4 次、合计每股派现 0.13 元，最近一次实施在 2025 年 12 月。 近三年“再融资”类公告 1 条，涉及定向增发。 综合看，公司融资结构上股权与债务融资并举；2025 年筹资活动现金流净额 -0.48 亿元，融资节奏仍以维持存量债务周转为主，与前述经营与投资现金流的方向基本一致。

\pfl{未来潜在融资需求分析} 2026 年 6 月末货币资金 1.34 亿元，而短期债务 1.75 亿元（现金短债比 0.77 倍），2026 年上半年经营现金流净额 -0.46 亿元。账面货币资金可以覆盖短期债务，缺口不体现在静态偿付层面。 公司 2026 年上半年归母净利润为负（-0.14 亿元），但 6 月末资产负债率 61.9\% 尚未进入第一档（亏损且负债率 $>$65\%）的区间，当前融资需求以补充流动资金、支撑低谷期经营性支出为主。 公开披露的后续资金安排线索包括：债务与授信安排：2025 年 6 月，公司主要融资方式为以存货（林木资产）抵押取得银行借款。 静态偿付不存在缺口，后续融资更可能服务于产能或并购安排，而不是填补流动性窟窿。




\subsubsection{獐子岛（002069）}
\label{co:002069}

\pfl{公司基本情况} 辽宁大连长海县獐子岛镇起家的海洋牧场企业，2006 年 9 月深交所上市，A 股“水产第一股”，因 2014 年以来多次“扇贝绝收/跑路”及 2016—2017 年年度报告虚假记载于 2020 年 6 月被证监会处罚、原董事长吴厚刚被终身市场禁入并获刑 15 年；2022 年 3 月原控股股东所持股份被拍卖、大连市国资入主，2022 年 5 月被实施其他风险警示（ST 獐子岛）、2023 年 5 月撤销；现由大连市国资委实际控制，公司以“扭亏、瘦身、聚焦、增效”为主线剥离非核心资产，2026 年上半年营业收入 7.31 亿元、归母净利润 1,287.14 万元。 公司于 2006-09-28 在深交所主板首发上市，注册资本 7.11 亿元，总股本 7.11 亿股。 截至 2026 年 10 月 12 日收盘价 3.63 元，对应总市值 25.8 亿元，近一年-9.0\%，总市值在三级行业“水产养殖”的 4 家公司中按由大到小排第\zhnumber{3}位，近一年涨跌幅高于全样本中位数（-10.4\%）1.3 个百分点。 按 2026 年半年报披露的行业构成，收入占比前三位为水产贸易业 41.3\%；水产加工业 40.5\%；水产养殖业 13.9\%。 股权结构方面，实际控制人为大连市国资委。

\pfl{股权结构与控制权} 按 2026 年 6 月 30 日的股东名册，前三大股东为大连盐化集团有限公司 持股 15.46\%；长海县獐子岛褡裢经济发展中心 持股 7.21\%；长海县獐子岛投资发展中心 持股 7.03\%，股东户数 32,489 户。 从控制权角度看，持有公司 1.10 亿股、占总股本 15.46\%；2024 年 11 月，公司解除对大连盐化集团所持 10,996 万股表决权的委托。

\pfl{经营要点} 经营与产能方面，2025 年，公司转让全资子公司大连獐子岛船舶制造有限公司 100\% 股权。

\pfl{所处行业的基本面} 水产养殖的供给由养殖面积与投苗量决定，需求由餐饮与家庭消费决定，价格弹性大于畜禽（第 \ref{sec:framework} 章的供需框架）。 行业层面的最新情况是：2025 年，全国水产养殖面积 7760.60 千公顷、增长 2.55\%，海水养殖面积增长 7.52\%；2026 年，大连连续三年清退近海养殖区、养殖面积缩减约 22\%，高端底播参价格上行而中低端养殖参承压；提出到 2030 年全省优质海参产量稳定在 10 万吨以上、产值突破 600.00 亿元。

\pfl{竞争优势} 把该公司放回同类公司比较，2026 年上半年营业收入 7.3 亿元、在三级行业“水产养殖”的 4 家公司中排第\zhnumber{3}位；毛利率 18.2\%、高于行业内中位数 0.5 个百分点；半年 ROE 29.5\%（未年化）、高于行业内中位数（6.05\%）；资产负债率 97.0\%、高于行业内中位数（72.2\%）。 优势在当期体现为高于行业中位数的资本回报，但优势能否持续仍取决于周期与自身产能投放节奏。

\pfl{历史股价} 自 2021 年 1 月至 2026 年 10 月，股价由 3.23 元变为 3.63 元，累计 +12.4\%，区间最高 4.95 元（2023 年 7 月）、最低 2.42 元（2024 年 6 月）；当前价相当于区间最高价的 73.3\%，为区间最低价的 1.50 倍。 月度收益的年化波动率 34.0\%，低于全样本中位数 37.3\%，在三级行业“水产养殖”的 4 家公司中波动率由高到低排第\zhnumber{3}位。 把股价阶段与公司事件对齐看，2024 年 6 月 至 2025 年 8 月股价上涨+80.2\%，同期（2024 年）公司另收到船厂土地及附着物收购补偿款 5,860.00 万元；2022 年 5 月 至 2023 年 7 月股价上涨+54.2\%，同期（2022 年 10 月）被判处有期徒刑十五年并处罚金 92.00 万元。 整体上看，区间的几处大级别拐点都能在公司的资本运作、股权变动或风险事件上找到对应，股价波动有明显的公司层面驱动，而非单纯的行业 beta。

\begin{center}
\begin{minipage}{\textwidth}\centering
\begin{tikzpicture}
\begin{axis}[
  width=12.6cm, height=3.9cm, scale only axis,
  xmin=2020.922, xmax=2026.972, ymin=2.25, ymax=5.30,
  xtick={2021,2022,2023,2024,2025,2026},
  yearticks,
  yticklabel style={font=\scriptsize},
  ylabel={元/股}, ylabel style={font=\scriptsize},
  grid=major, grid style={LightGray, line width=0.3pt},
  axis line style={InkGray, line width=0.5pt},
  every axis plot/.append style={line width=0.9pt},
]
\addplot[AgriGreen] table[x=x,y=y]{data/clean/profiles/px_002069.dat};
\addplot[SoftRed, dashed, line width=0.7pt] table[x=x,y=y]{data/clean/profiles/pxzz_002069.dat};
\addplot[only marks, mark=*, mark size=1.1pt, SoftRed] table[x=x,y=y]{data/clean/profiles/pxzz_002069.dat};
\end{axis}
\end{tikzpicture}

\captionof{figure}[獐子岛（002069）股价与周期骨架]{獐子岛（002069）月度收盘价（元/股）与 ZigZag 周期骨架（阈值 25\%，圆点为周期转折点）}
\end{minipage}
\end{center}

\pfl{过去周期分析} 以 25\% 的 ZigZag 阈值划分，样本区间内股价形成 4 个主升段与 3 个主跌段。 最大主升段出现在 2024 年 6 月 至 2025 年 8 月，14 个月+80.2\%；最大主跌段为 2023 年 7 月 至 2024 年 6 月，11 个月-51.1\%。 最近一次转折出现在 2026 年 8 月（上涨段自 2026 年 6 月起，2 个月+26.9\%），当前处于该段之后的震荡之中。 公司股价较申万农林牧渔指数提前约 23 个月见底，是板块中较早定价周期反转的公司之一。 水产养殖的价格周期由投苗量与消费需求共同决定（第 \ref{sec:other-livestock} 节），与猪周期的同步性弱于禽链，公司 2026 年 6 月末资产负债率 97.0\%，高于全样本中位数 52.2\%，其股价因此更多反映资金面与信用风险，而不是价格弹性本身。 综合区间形态看，该股单段涨跌幅度偏大、以趋势段为主（最大主升 +80.2\%、最大主跌 -51.1\%），年化波动率 34.0\% 与全样本中位数 37.3\% 相比波动与样本中位数接近。

\pfl{盈利情况} 2026 年上半年营业收入 7.31 亿元（同比 -5.3\%），归母净利润 0.13 亿元，同比 +26.8\%。 以年报为趋势对照，2023---2025 年营业收入由 16.77 亿元变为 12.78 亿元（两年年均复合增速 -12.7\%），归母净利润由 0.09 亿元变为 -0.11 亿元。 按半年报口径，2026 年上半年的 ROE 为 29.5\%（未年化）、销售净利率 2.5\%、毛利率 18.2\%，ROE 在“水产养殖”4 家公司中按数值由高到低排第\zhnumber{1}位。 同期全样本半年 ROE 中位数 0.75\%，公司与之相差 28.77 个百分点（略好于中位数公司）。 综合看，公司在报告期维持盈利（高于全样本半年 ROE 中位数 0.75\%），利润的可持续性取决于水产养殖景气的延续与成本管控的执行。

\begin{center}\small\setlength{\tabcolsep}{3.4pt}
\captionof{table}[獐子岛（002069）关键财务指标]{獐子岛（002069）关键财务指标（合并报表口径；两个半年报期均未年化；现金短债比 $=$ 货币资金 $\div$ 短期债务）}
\begin{adjustbox}{max width=\textwidth}
\begin{tabular}{@{}lrrrrr@{}}
\toprule
指标 & 2023 年 & 2024 年 & 2025 年 & 2025 年半年报 & 2026 年半年报 \\
\midrule
营业收入（亿元） & 16.77 & 15.83 & 12.78 & 7.72 & 7.31 \\
归母净利润（亿元） & 0.09 & -0.22 & -0.11 & 0.10 & 0.13 \\
毛利率（\%） & 17.4 & 11.2 & 15.9 & 16.1 & 18.2 \\
ROE（\%） & 12.3 & -32.7 & -23.2 & 17.9 & 29.5 \\
资产负债率（\%） & 94.7 & 95.8 & 96.8 & 95.2 & 97.0 \\
经营活动现金流净额（亿元） & 1.90 & 0.79 & 0.42 & 0.52 & 0.26 \\
货币资金（亿元） & 5.53 & 6.04 & 5.80 & 6.15 & 5.45 \\
有息负债合计（亿元） & 20.20 & 19.48 & 19.21 & 19.44 & 19.24 \\
现金短债比（倍） & 0.31 & 0.31 & 0.49 & 0.36 & 0.39 \\
\bottomrule
\end{tabular}
\end{adjustbox}
\end{center}

\begin{center}
\begin{minipage}{\textwidth}\centering
\begin{tikzpicture}
\begin{groupplot}[
  group style={group size=2 by 1, horizontal sep=0.60cm},
  width=6.85cm, height=4.60cm,
  tick label style={font=\tiny},
  title style={font=\scriptsize\heiti},
  yticklabel style={font=\tiny},
  ylabel style={font=\tiny},
  grid=major, grid style={LightGray, line width=0.3pt},
  axis line style={InkGray, line width=0.5pt},
  enlarge x limits={abs=0.8},
  legend style={font=\scriptsize, at={(0.5,-0.20)}, anchor=north,
                legend columns=2, draw=none, fill=none, column sep=6pt},
]
\nextgroupplot[
  ybar, bar width=4pt,
  ymin=-2.32, ymax=23.35,
  xtick={0,1,2,3,4,5.7},
  xticklabels={2021,2022,2023,2024,2025,2026H1},
  ylabel={亿元},
]
\addplot[fill=AgriGreen!75, draw=AgriGreen, bar shift=-2.6pt] table[x=i, y=rev]{data/clean/profiles/fin_002069.dat};
\addplot[fill=SoftRed!60, draw=SoftRed, bar shift=2.6pt] table[x=i, y=ni]{data/clean/profiles/fin_002069.dat};
\addplot[fill=AgriGreen!30, draw=AgriGreen, pattern=north east lines, bar shift=-2.6pt] table[x=x, y=rev]{data/clean/profiles/finh_002069.dat};
\addplot[fill=SoftRed!20, draw=SoftRed, pattern=north east lines, bar shift=2.6pt] table[x=x, y=ni]{data/clean/profiles/finh_002069.dat};
\legend{营业收入（年/半年）, 归母净利润（年/半年）}
\nextgroupplot[
  ymin=-43.0, ymax=110.0,
  xtick={0,1,2,3,4,5.7},
  xticklabels={2021,2022,2023,2024,2025,2026H1},
  ylabel={\%},
]
\addplot[AgriGold, mark=*, mark size=1.3pt] table[x=i, y=roe]{data/clean/profiles/fin_002069.dat};
\addplot[OilBlue, mark=square*, mark size=1.3pt] table[x=i, y=debt]{data/clean/profiles/fin_002069.dat};
\addplot[only marks, AgriGold, mark=*, mark size=2.0pt] table[x=x, y=roe]{data/clean/profiles/finh_002069.dat};
\addplot[only marks, OilBlue, mark=square*, mark size=2.0pt] table[x=x, y=debt]{data/clean/profiles/finh_002069.dat};
\legend{ROE, 资产负债率}
\end{groupplot}
\end{tikzpicture}
\begin{tikzpicture}
\begin{axis}[
  name=finbot,
  width=13.75cm, height=3.10cm, scale only axis,
  ymin=-2.06, ymax=23.03,
  xtick={0,1,2,3,4,5.7},
  xticklabels={2021,2022,2023,2024,2025,2026H1},
  tick label style={font=\tiny},
  yticklabel style={font=\tiny},
  ylabel={亿元}, ylabel style={font=\tiny},
  ybar, bar width=4pt,
  grid=major, grid style={LightGray, line width=0.3pt},
  axis line style={InkGray, line width=0.5pt},
  enlarge x limits={abs=0.8},
  legend style={font=\scriptsize, at={(0.5,-0.26)}, anchor=north,
                legend columns=2, draw=none, fill=none, column sep=10pt},
]
\addplot[fill=AgriGreen!55, draw=AgriGreen, bar shift=-3.0pt] table[x=i, y=cash]{data/clean/profiles/fin_002069.dat};
\addplot[fill=SoftRed!45, draw=SoftRed, bar shift=3.0pt] table[x=i, y=ibd]{data/clean/profiles/fin_002069.dat};
\addplot[fill=AgriGreen!25, draw=AgriGreen, pattern=north east lines, bar shift=-3.0pt] table[x=x, y=cash]{data/clean/profiles/finh_002069.dat};
\addplot[fill=SoftRed!20, draw=SoftRed, pattern=north east lines, bar shift=3.0pt] table[x=x, y=ibd]{data/clean/profiles/finh_002069.dat};
\legend{货币资金（左轴）, 有息负债（左轴）}
\end{axis}
\begin{axis}[
  at={(finbot.south east)}, anchor=south east,
  width=13.75cm, height=3.10cm, scale only axis,
  axis y line*=right, axis x line*=none, xtick=\empty,
  ymin=-0.49, ymax=2.47,
  tick label style={font=\tiny},
  yticklabel style={font=\tiny}, scaled y ticks=false,
  legend style={font=\scriptsize, at={(0.5,-0.44)}, anchor=north,
                legend columns=1, draw=none, fill=none},
]
\addplot[OilBlue, mark=*, mark size=2.2pt, line width=1.3pt] table[x=i, y=ocf]{data/clean/profiles/fin_002069.dat};
\addplot[only marks, OilBlue, mark=*, mark size=3.2pt] table[x=x, y=ocf]{data/clean/profiles/finh_002069.dat};
\legend{经营活动现金流净额（右轴，亿元，蓝点）}
\end{axis}
\end{tikzpicture}

\captionof{figure}[獐子岛（002069）近年主要财务指标]{獐子岛（002069）近年主要财务指标（左上：营业收入与归母净利润；右上：ROE 与资产负债率；下：货币资金、有息负债为左轴柱状，经营活动现金流净额为其右轴的蓝点折线。
2021---2025 年为年报口径，2026H1 为 2026 年半年报/半年末，与其前一年报之间留出间隔、以斜纹或空心点区分；半年报数据未年化）}
\end{minipage}
\end{center}

\pfl{财务分析} 2026 年 6 月末资产负债率 97.0\%，较 2025 年末上升 0.2 个百分点（样本中位数 52.2\%，所处三级行业内按由低到高排第\zhnumber{4}位）。 2026 年 6 月末货币资金 5.45 亿元、短期债务（短期借款与一年内到期非流动负债）14.05 亿元、长期债务 5.19 亿元，有息负债合计 19.24 亿元，现金短债比 0.39 倍——覆盖偏紧。 净有息负债 13.79 亿元，相当于其 2026 年上半年营业收入的 189\%。 流动比率 0.99、速动比率 0.47，短期偿付压力较大。 2026 年上半年经营活动现金流净额 0.26 亿元，经营现金流与利润同向为正，内生资金可以支撑运营。 综合看，账面货币资金对有息负债与短期债务的覆盖不足，滚动偿付对新增融资的依赖度较高；资产负债率高于样本中位数 44.8 个百分点；有息负债中短期债务占比更高，期限结构仍以滚动续作为主。 公司披露的诉讼与资产冻结风险为：2025 年 1 月，公司及控股子公司连续 12 个月内累计诉讼、仲裁事项涉案金额共计 1,082.28 万元。

\pfl{重大战略转型} 公司的战略动作可归为以下几类：其他动作（3 项）：2026 年 3 月，公司拟与大连獐子岛海洋发展集团有限公司共同投资设立大连海发生态渔业发展有限公司；2022 年 3 月，原控股股东长海县獐子岛投资发展中心持有的 1.0996 亿股被大连盐化集团以 3.43 亿元拍下；并购与重组（2 项）：2025 年，公司以“扭亏、瘦身、聚焦、增效”为主线推进系统性改革。 公告口径上，近三年（2023 年 9 月---2026 年 9 月）公司披露重大重组与并购类公告 5 条、对外投资与转型类公告 2 条、回购与股权激励类公告 0 条，合计公告 307 条。 第一大行业仍是“水产贸易业”，收入占比 41.3\%，与 2021 年年报相比没有方向性变化。 综合判断，报告期内公司有 7 条与战略相关的公告落地（重大重组与并购 5 条、对外投资与转型 2 条），资本运作与主业调整之间存在直接关联。

\begin{center}\small\setlength{\tabcolsep}{4pt}
\captionof{table}[獐子岛（002069）历史融资与资本运作]{獐子岛（002069）历史融资与资本运作（据公司公告与公开报道整理；金额为披露值，含首次公开发行、再融资、债券、银行借款与重整投资等）}
\begin{adjustbox}{max width=\textwidth}
\begin{tabular}{@{}p{1.45cm}p{2.55cm}p{4.05cm}p{4.95cm}@{}}
\toprule
时间 & 方式 & 规模 & 用途与说明 \\
\midrule
2006-09-28 & IPO（深圳证券交易所） & 实际发行 2,\allowbreak{}830.00 万股 & 限责任公司，\allowbreak{}发行市盈率（按发行后总股本）18.80 倍 \\
2019-08-31 & 重大资产出售（向亚渔食品出售新中海产 100\% 股权及新中日本 90\% 股权） & 交易对价 2.35 亿元（拟出售资产于估值基准日 2019 年 5 月 31 日的估值 2.41 亿元） & 回笼资金用于集中资源加快海洋牧场重构、\allowbreak{}降低资产负债率、\allowbreak{}保障公司安全运营；交易对方亚渔食品为亚洲渔港股份有限公司全资子公司 \\
2020 & 资产处置回笼资金 & 转让位于长海县广鹿岛的 4 宗海域使用的租赁权暨海底存货，\allowbreak{}已按约定收到首付款和进度款合计 9,\allowbreak{}678.00 万元 & 盘活存量资产；经 2020 年 1 月 3 日第七届董事会第九次会议及 2020 年 2 月 3 日 2020 年第一次临时股东大会审议通过公司转让海域使用的租赁权暨海底存货 \\
2022 & 出售分公司资产 & 将大连乌蟒岛分公司的存货、\allowbreak{}固定资产、\allowbreak{}无形资产等相关资产通过大连产权交易所挂牌转让给大连智慧渔业集团有限公司，\allowbreak{}总价款合计 8,\allowbreak{}272.14 万元 & 盘活存量资产、\allowbreak{}改善财务状况；出售存货收益 2,\allowbreak{}078.73 万元 \\
2024-12-12 & 向控股股东出售资产（关联交易） & 向控股股东大连盐化集团有限公司出售公司持有的原鲍鱼厂资产、\allowbreak{}原育苗三厂资产以及部分海域使用权等，\allowbreak{}交易价格 9,\allowbreak{}215.36 万元 & 盘活存量资产；2024 年公司另收到船厂土地及附着物收购补偿款 5,\allowbreak{}860.00 万元 \\
2025-05-22 & 向特定对象发行 A 股股票（定增预案，对象为间接控股股东） & 拟向间接股东大连獐子岛海洋发展集团有限公司定向发行，\allowbreak{}发行价格 3.09 元/\allowbreak{}股 & 扣除相关发行费用后全部用于偿还银行借款及补充流动资金；经公司第八届董事会第十九次会议、\allowbreak{}第八届监事会第十四次会议审议通过；本次发行属于董事会确定全部发行对象的再融资；尚需取得有权国有资产监督管理部门或其授权主体批准、\allowbreak{}股东大会审议通过、\allowbreak{}深交所审核通过并经中国证监会同意注册；若发行落地，\allowbreak{}国资持股比例将升至 31.69\% \\
2026 & 银行融资置换（拟申请综合授信） & 拟向大连银行股份有限公司及大连农村商业银行股份有限公司申请融资贷款 4.95 亿元 & 用于等额置换公司在中国进出口银行辽宁省分行的贷款；公司间接控股股东大连獐子岛海洋发展集团有限公司提供相关支持；2026 年度公司另拟向中国农业银行大连市分行申请综合授信 4.50 亿元、\allowbreak{}向中国建设银行大连市分行申请 3.60 亿元、\allowbreak{}向中国工商银行大连市分行申请 2.90 亿元等 \\
\bottomrule
\end{tabular}
\end{adjustbox}
\end{center}

\pfl{历史融投资情况} 从现金流量表看，2025 年公司取得借款收到的现金 24.07 亿元、偿还债务支付的现金 24.34 亿元，筹资活动现金流净额 -0.97 亿元（2023 年为取得借款 19.84 亿元、偿还债务 19.97 亿元）。 2026 年上半年筹资活动现金流净额为 -0.32 亿元（取得借款 6.51 亿元、偿还债务 6.52 亿元，未年化），已转为净流出，处于净还债状态。 2025 年分配股利、利润或偿付利息支付的现金合计 0.64 亿元。 公开披露的融资记录共 7 笔，工具类型以 IPO、定向增发为主；金额与用途见下表。 其中规模最大的两笔为：2025 年 5 月，定向增发，拟向间接股东大连獐子岛海洋发展集团有限公司定向发行，发行价格 3.09 元/股；2006 年 9 月，IPO，实际发行 2,830.00 万股。 股本方面，近三年披露股本变动 3 次；总股本自 2021 年末以来未发生变化（7.11 亿股）。 分红方面，近三年未实施现金分红，在全部样本中属于分红能力偏弱的一端。 近三年“再融资”类公告 8 条，涉及定向增发、募集资金使用。 综合看，公司融资结构上股权与债务融资并举；2025 年筹资活动现金流净额 -0.97 亿元，融资节奏仍以维持存量债务周转为主，与前述经营与投资现金流的方向基本一致。

\pfl{未来潜在融资需求分析} 2026 年 6 月末货币资金 5.45 亿元，而短期债务 14.05 亿元（现金短债比 0.39 倍），2026 年上半年经营现金流净额 0.26 亿元。按“短期债务减去账面货币资金”的滚动偿付口径，静态缺口约 8.60 亿元；上半年盈利或接近盈亏平衡，该缺口不随经营亏损进一步扩大。 公司 2026 年 6 月末资产负债率 97.0\%、2026 年上半年 ROE 29.5\%（未年化），资产负债率偏高（97.0\%），融资需求首先用于置换与滚动存量债务、补充营运资金，属于“补血型”而非扩张型。 公开披露的后续资金安排线索包括：股权融资线索：2025 年 5 月，公司拟向间接股东大连獐子岛海洋发展集团有限公司定向发行不超过 1.69 亿股；债务与授信安排：2026 年度，公司拟向下列银行申请综合授信额度并根据银行要求提供相应增信措施：中国农业银行股份有限公司大连市分行 4.50 亿元、中国建设银行股份有限公司大连市分行 3.60 亿元、中国工商银行股份有限公司大连市分行 2.90 亿元、中国民生银行股份有限公司大连分行等。 综合看，公司的外部资金需求以补充流动性与项目投入为主，滚动偿付口径下的静态缺口约 8.60 亿元，融资节奏将受制于银行续贷意愿与股权融资窗口。




\subsubsection{大湖股份（600257）}
\label{co:600257}

\pfl{公司基本情况} 湖南常德起家的淡水养殖上市公司（曾名“洞庭水殖”“大湖水殖”，自称“中国淡水渔业第一股”），2000 年 6 月在上交所上市；2016 年定增募资 5.5 亿元后由水产养殖向水产加工、白酒和康复护理医疗延伸，2020—2021 年分步收购并增资东方华康医疗管理有限公司取得 60\% 股权（2026 年 4 月增至 61.1103\%），2025 年 10 月公司全称变更为“大湖健康产业股份有限公司”；2025 年营业收入 9.20 亿元、归母净利润 506.53 万元，2026 年上半年归母净利润 1,688.13 万元、同比扭亏。 公司于 2000-06-12 在上交所首发上市，注册资本 4.81 亿元，总股本 4.81 亿股。 截至 2026 年 10 月 12 日收盘价 5.34 元，对应总市值 25.7 亿元，近一年+0.9\%，总市值在三级行业“水产养殖”的 4 家公司中按由大到小排第\zhnumber{4}位，近一年涨跌幅高于全样本中位数（-10.4\%）11.3 个百分点。 按 2025 年年报披露的行业构成，收入占比前三位为农业 47.7\%；其他 41.9\%；工业 8.8\%。

\pfl{股权结构与控制权} 按 2026 年 6 月 30 日的股东名册，前三大股东为西藏泓杉科技发展有限公司 持股 28.09\%；范远龙 持股 2.90\%；高盛国际－自有资金（持股比例未披露），股东户数 32,879 户。 从控制权角度看，2026 年 6 月末持有公司 1.35 亿股、占总股本 28.09\%，其中质押 2,884.00 万股；实际控制人罗祖亮通过控股股东等主体间接持有公司 26.76\% 股份；公司于 2025 年 10 月由“大湖水殖股份有限公司”更名为“大湖健康产业股份有限公司”。

\pfl{经营要点} 经营层面，淡水养殖面积约 16.8 万亩。

\pfl{所处行业的基本面} 水产养殖的供给由养殖面积与投苗量决定，需求由餐饮与家庭消费决定，价格弹性大于畜禽（第 \ref{sec:framework} 章的供需框架）。 行业层面的最新情况是：2025 年，全国淡水产品产量 3792.20 万吨、同比增长 3.93\%，淡水养殖面积 5352.08 千公顷、增长 0.46\%；2026 年 7 月，湖北省重点监测水产品批发市场 13 个品种加权均价 20.43 元/公斤、同比下跌 2.6\%。

\pfl{竞争优势} 把该公司放回同类公司比较，2026 年上半年营业收入 4.3 亿元、在三级行业“水产养殖”的 4 家公司中排第\zhnumber{4}位；毛利率 29.4\%、高于行业内中位数 11.8 个百分点；半年 ROE 2.1\%（未年化）、低于行业内中位数（6.05\%）；资产负债率 46.4\%、低于行业内中位数（72.2\%）。 公司区别于同类公司的经营特征包括：公司通过并购、新设方式注入并已实际运营医疗机构 7 家。 优势在报表上主要体现为高于行业平均的毛利率，这在周期底部是抵御亏损的主要缓冲。

\pfl{历史股价} 自 2021 年 1 月至 2026 年 10 月，股价由 5.02 元变为 5.34 元，累计 +6.4\%，区间最高 9.23 元（2023 年 8 月）、最低 4.29 元（2024 年 8 月）；当前价相当于区间最高价的 57.9\%，为区间最低价的 1.24 倍。 月度收益的年化波动率 47.8\%，高于全样本中位数 37.3\%，在三级行业“水产养殖”的 4 家公司中波动率由高到低排第\zhnumber{1}位。 把股价阶段与公司事件对齐看，2022 年 10 月 至 2023 年 8 月股价上涨+105.1\%，同期（2023 年 7 月）公司此前已于 2023 年 1 月 19 日、2 月 6 日审议通过延长该次非公开发行股东大会决议有效期及授权有效期至 2024 年 2 月 13 日；2024 年 8 月 至 2026 年 1 月股价上涨+72.3\%，同期（2026 年 1 月）公司以 700.00 万元购买徐培麒持有的嘉兴智康 35\% 股权、以 100.00 万元购买上海麒聚瑆伙机器人科技合伙企业（有限合伙）持有的 5\% 股权（合计 800.00 万元受让 40\% 股权）。 整体上看，区间的几处大级别拐点都能在公司的资本运作、股权变动或风险事件上找到对应，股价波动有明显的公司层面驱动，而非单纯的行业 beta。

\begin{center}
\begin{minipage}{\textwidth}\centering
\begin{tikzpicture}
\begin{axis}[
  width=12.6cm, height=3.9cm, scale only axis,
  xmin=2020.922, xmax=2026.972, ymin=3.99, ymax=9.88,
  xtick={2021,2022,2023,2024,2025,2026},
  yearticks,
  yticklabel style={font=\scriptsize},
  ylabel={元/股}, ylabel style={font=\scriptsize},
  grid=major, grid style={LightGray, line width=0.3pt},
  axis line style={InkGray, line width=0.5pt},
  every axis plot/.append style={line width=0.9pt},
]
\addplot[AgriGreen] table[x=x,y=y]{data/clean/profiles/px_600257.dat};
\addplot[SoftRed, dashed, line width=0.7pt] table[x=x,y=y]{data/clean/profiles/pxzz_600257.dat};
\addplot[only marks, mark=*, mark size=1.1pt, SoftRed] table[x=x,y=y]{data/clean/profiles/pxzz_600257.dat};
\end{axis}
\end{tikzpicture}

\captionof{figure}[大湖股份（600257）股价与周期骨架]{大湖股份（600257）月度收盘价（元/股）与 ZigZag 周期骨架（阈值 25\%，圆点为周期转折点）}
\end{minipage}
\end{center}

\pfl{过去周期分析} 以 25\% 的 ZigZag 阈值划分，样本区间内股价形成 5 个主升段与 5 个主跌段。 最大主跌段为 2023 年 8 月 至 2024 年 2 月，6 个月-45.9\%；最大主升段出现在 2022 年 10 月 至 2023 年 8 月，10 个月+105.1\%。 最近一次转折出现在 2026 年 8 月（上涨段自 2026 年 6 月起，2 个月+33.3\%），当前处于该段之后的震荡之中。 公司股价较申万农林牧渔指数提前约 21 个月见底，是板块中较早定价周期反转的公司之一。 水产养殖的价格周期由投苗量与消费需求共同决定（第 \ref{sec:other-livestock} 节），与猪周期的同步性弱于禽链，公司 2026 年上半年实现盈利、半年 ROE 2.1\%（未年化），好于全样本中位数（0.75\%），下行周期对其报表的冲击小于同业。 综合区间形态看，该股单段涨跌幅度偏大、以趋势段为主（最大主升 +105.1\%、最大主跌 -45.9\%），年化波动率 47.8\% 与全样本中位数 37.3\% 相比波动明显大于样本中位数。

\pfl{盈利情况} 2026 年上半年营业收入 4.34 亿元（同比 +2.0\%），归母净利润 0.17 亿元，实现扭亏。 以年报为趋势对照，2023---2025 年营业收入由 11.93 亿元变为 9.20 亿元（两年年均复合增速 -12.2\%），归母净利润由 -0.08 亿元变为 0.05 亿元。 按半年报口径，2026 年上半年的 ROE 为 2.1\%（未年化）、销售净利率 7.8\%、毛利率 29.4\%，ROE 在“水产养殖”4 家公司中按数值由高到低排第\zhnumber{3}位。 按同一口径，三级行业“水产养殖”4 家公司的半年 ROE 中位数为 6.05\%，该公司低于其中位数 3.96 个百分点。 综合看，公司在报告期维持盈利（高于全样本半年 ROE 中位数 0.75\%），利润的可持续性取决于水产养殖景气的延续与成本管控的执行。

\begin{center}\small\setlength{\tabcolsep}{3.4pt}
\captionof{table}[大湖股份（600257）关键财务指标]{大湖股份（600257）关键财务指标（合并报表口径；两个半年报期均未年化；现金短债比 $=$ 货币资金 $\div$ 短期债务）}
\begin{adjustbox}{max width=\textwidth}
\begin{tabular}{@{}lrrrrr@{}}
\toprule
指标 & 2023 年 & 2024 年 & 2025 年 & 2025 年半年报 & 2026 年半年报 \\
\midrule
营业收入（亿元） & 11.93 & 10.44 & 9.20 & 4.26 & 4.34 \\
归母净利润（亿元） & -0.08 & -0.77 & 0.05 & -0.03 & 0.17 \\
毛利率（\%） & 24.5 & 23.1 & 27.6 & 25.0 & 29.4 \\
ROE（\%） & -0.9 & -9.2 & 0.6 & -0.3 & 2.1 \\
资产负债率（\%） & 51.0 & 48.6 & 46.5 & 48.0 & 46.4 \\
经营活动现金流净额（亿元） & 1.62 & 0.67 & 2.18 & 0.84 & 0.56 \\
货币资金（亿元） & 3.36 & 1.87 & 3.22 & 2.33 & 2.61 \\
有息负债合计（亿元） & 8.68 & 7.30 & 6.90 & 7.32 & 6.87 \\
现金短债比（倍） & 1.04 & 0.67 & 1.08 & 0.74 & 1.04 \\
\bottomrule
\end{tabular}
\end{adjustbox}
\end{center}

\begin{center}
\begin{minipage}{\textwidth}\centering
\begin{tikzpicture}
\begin{groupplot}[
  group style={group size=2 by 1, horizontal sep=0.60cm},
  width=6.85cm, height=4.60cm,
  tick label style={font=\tiny},
  title style={font=\scriptsize\heiti},
  yticklabel style={font=\tiny},
  ylabel style={font=\tiny},
  grid=major, grid style={LightGray, line width=0.3pt},
  axis line style={InkGray, line width=0.5pt},
  enlarge x limits={abs=0.8},
  legend style={font=\scriptsize, at={(0.5,-0.20)}, anchor=north,
                legend columns=2, draw=none, fill=none, column sep=6pt},
]
\nextgroupplot[
  ybar, bar width=4pt,
  ymin=-3.27, ymax=14.69,
  xtick={0,1,2,3,4,5.7},
  xticklabels={2021,2022,2023,2024,2025,2026H1},
  ylabel={亿元},
]
\addplot[fill=AgriGreen!75, draw=AgriGreen, bar shift=-2.6pt] table[x=i, y=rev]{data/clean/profiles/fin_600257.dat};
\addplot[fill=SoftRed!60, draw=SoftRed, bar shift=2.6pt] table[x=i, y=ni]{data/clean/profiles/fin_600257.dat};
\addplot[fill=AgriGreen!30, draw=AgriGreen, pattern=north east lines, bar shift=-2.6pt] table[x=x, y=rev]{data/clean/profiles/finh_600257.dat};
\addplot[fill=SoftRed!20, draw=SoftRed, pattern=north east lines, bar shift=2.6pt] table[x=x, y=ni]{data/clean/profiles/finh_600257.dat};
\legend{营业收入（年/半年）, 归母净利润（年/半年）}
\nextgroupplot[
  ymin=-22.9, ymax=61.2,
  xtick={0,1,2,3,4,5.7},
  xticklabels={2021,2022,2023,2024,2025,2026H1},
  ylabel={\%},
]
\addplot[AgriGold, mark=*, mark size=1.3pt] table[x=i, y=roe]{data/clean/profiles/fin_600257.dat};
\addplot[OilBlue, mark=square*, mark size=1.3pt] table[x=i, y=debt]{data/clean/profiles/fin_600257.dat};
\addplot[only marks, AgriGold, mark=*, mark size=2.0pt] table[x=x, y=roe]{data/clean/profiles/finh_600257.dat};
\addplot[only marks, OilBlue, mark=square*, mark size=2.0pt] table[x=x, y=debt]{data/clean/profiles/finh_600257.dat};
\legend{ROE, 资产负债率}
\end{groupplot}
\end{tikzpicture}
\begin{tikzpicture}
\begin{axis}[
  name=finbot,
  width=13.75cm, height=3.10cm, scale only axis,
  ymin=-0.95, ymax=10.69,
  xtick={0,1,2,3,4,5.7},
  xticklabels={2021,2022,2023,2024,2025,2026H1},
  tick label style={font=\tiny},
  yticklabel style={font=\tiny},
  ylabel={亿元}, ylabel style={font=\tiny},
  ybar, bar width=4pt,
  grid=major, grid style={LightGray, line width=0.3pt},
  axis line style={InkGray, line width=0.5pt},
  enlarge x limits={abs=0.8},
  legend style={font=\scriptsize, at={(0.5,-0.26)}, anchor=north,
                legend columns=2, draw=none, fill=none, column sep=10pt},
]
\addplot[fill=AgriGreen!55, draw=AgriGreen, bar shift=-3.0pt] table[x=i, y=cash]{data/clean/profiles/fin_600257.dat};
\addplot[fill=SoftRed!45, draw=SoftRed, bar shift=3.0pt] table[x=i, y=ibd]{data/clean/profiles/fin_600257.dat};
\addplot[fill=AgriGreen!25, draw=AgriGreen, pattern=north east lines, bar shift=-3.0pt] table[x=x, y=cash]{data/clean/profiles/finh_600257.dat};
\addplot[fill=SoftRed!20, draw=SoftRed, pattern=north east lines, bar shift=3.0pt] table[x=x, y=ibd]{data/clean/profiles/finh_600257.dat};
\legend{货币资金（左轴）, 有息负债（左轴）}
\end{axis}
\begin{axis}[
  at={(finbot.south east)}, anchor=south east,
  width=13.75cm, height=3.10cm, scale only axis,
  axis y line*=right, axis x line*=none, xtick=\empty,
  ymin=-0.57, ymax=2.83,
  tick label style={font=\tiny},
  yticklabel style={font=\tiny}, scaled y ticks=false,
  legend style={font=\scriptsize, at={(0.5,-0.44)}, anchor=north,
                legend columns=1, draw=none, fill=none},
]
\addplot[OilBlue, mark=*, mark size=2.2pt, line width=1.3pt] table[x=i, y=ocf]{data/clean/profiles/fin_600257.dat};
\addplot[only marks, OilBlue, mark=*, mark size=3.2pt] table[x=x, y=ocf]{data/clean/profiles/finh_600257.dat};
\legend{经营活动现金流净额（右轴，亿元，蓝点）}
\end{axis}
\end{tikzpicture}

\captionof{figure}[大湖股份（600257）近年主要财务指标]{大湖股份（600257）近年主要财务指标（左上：营业收入与归母净利润；右上：ROE 与资产负债率；下：货币资金、有息负债为左轴柱状，经营活动现金流净额为其右轴的蓝点折线。
2021---2025 年为年报口径，2026H1 为 2026 年半年报/半年末，与其前一年报之间留出间隔、以斜纹或空心点区分；半年报数据未年化）}
\end{minipage}
\end{center}

\pfl{财务分析} 2026 年 6 月末资产负债率 46.4\%，较 2025 年末下降 0.1 个百分点（样本中位数 52.2\%，所处三级行业内按由低到高排第\zhnumber{1}位）。 2026 年 6 月末货币资金 2.61 亿元、短期债务（短期借款与一年内到期非流动负债）2.50 亿元、长期债务 4.36 亿元，有息负债合计 6.87 亿元，现金短债比 1.04 倍——覆盖充分。 净有息负债 4.25 亿元，相当于其 2026 年上半年营业收入的 98\%。 流动比率 2.24、速动比率 1.20，短期指标尚可。 2026 年上半年经营活动现金流净额 0.56 亿元，经营现金流与利润同向为正，内生资金可以支撑运营。 2026 年 6 月末商誉 1.45 亿元，占归母净资产的 14.8\%，是净资产中最脆弱的一块。 综合看，现金短债比 1.04 倍，短期偿付压力可控；资产负债率低于样本中位数 5.8 个百分点；有息负债中长期债务占比更高，期限结构对短债滚动的压力较小。 公司披露的风险集中在重整与债务违约、股东与质押两类：一是 2024 年 10 月，公司拟追回已支付的违规费用 600.34 万元并收取 30\% 违约金 180.10 万元；二是 2025 年 12 月，子公司净资产份额 2.95 亿元用于银行借款质押。

\pfl{重大战略转型} 近三年公司的战略动作集中在：并购与重组（5 项）：2026 年 1 月，公司以 700.00 万元购买徐培麒持有的嘉兴智康 35\% 股权、以 100.00 万元购买上海麒聚瑆伙机器人科技合伙企业（有限合伙）持有的 5\% 股权（合计 800.00 万元受让 40\% 股权）；2024 年 8 月，下属子公司皂市渔业与汉寿久达公司签订大湖水殖汉寿中华鳖有限公司的股权转让协议，以人民币 47.92 万元转让皂市渔业持有的汉寿中华鳖 72.7778\% 股权；其他动作（2 项）：2025 年 3 月，公司与子公司东方华康共同投资设立上海筷智医疗科技有限公司。 公告口径上，近三年（2023 年 9 月---2026 年 9 月）公司披露重大重组与并购类公告 1 条、对外投资与转型类公告 1 条、回购与股权激励类公告 0 条，合计公告 275 条。 第一大行业“农业”的收入占比由 2021 年年报的 58.3\% 变为 47.7\%（10.6 个百分点），收入结构出现再平衡。 综合判断，报告期内公司有 2 条与战略相关的公告落地（重大重组与并购 1 条、对外投资与转型 1 条），资本运作与主业调整之间存在直接关联。

\begin{center}\small\setlength{\tabcolsep}{4pt}
\captionof{table}[大湖股份（600257）历史融资与资本运作]{大湖股份（600257）历史融资与资本运作（据公司公告与公开报道整理；金额为披露值，含首次公开发行、再融资、债券、银行借款与重整投资等）}
\begin{adjustbox}{max width=\textwidth}
\begin{tabular}{@{}p{1.45cm}p{2.55cm}p{4.05cm}p{4.95cm}@{}}
\toprule
时间 & 方式 & 规模 & 用途与说明 \\
\midrule
2000-06-12 & IPO（上海证券交易所） & 实际发行 4,\allowbreak{}000.00 万股，\allowbreak{}发行价 8.90 元/\allowbreak{}股；实际募集资金合计 3.56 亿元 & 财证券有限责任公司；招股公告日 2000 年 5 月 11 日 \\
2016-10-18 & 非公开发行股票（定向配售、网下定价发行） & 发行价格 10.15 元，\allowbreak{}实际发行数量 5,\allowbreak{}418.7188 万股 & 公司 2025 年年度报告披露本次发行 5,\allowbreak{}419 万股、\allowbreak{}共募集资金 5.50 亿元 \\
2017-12 & 发行股份及支付现金购买资产并募集配套资金（重大资产重组，未完成） & 拟以 10.40 亿元收购西藏深万投实业有限公司 51\% 股权 & 用于支付本次交易中的现金对价、\allowbreak{}相关税费和中介机构费用（原募集配套资金方案）；标的公司核心资产为深圳市万生堂实业有限公司（日本冈本安全套中国内地总经销商） \\
2018-11-05 & 重大资产重组报告书（草案）披露 & 交易总价 10.40 亿元；募集配套资金不超过 3.50 亿元（配套股份数量不超过发行前总股本 4.81 亿股的 20\% & 支付本次交易中的现金对价、\allowbreak{}支付本次交易相关税费和中介机构费用；招商证券作为独立财务顾问披露；交易完成后公司将持有西藏深万投 51\% 股权 \\
2020-04-14 & 现金收购子公司股权及增资 & 以 2.00 亿元收购东方华康原有股东 32\% 股权 & 布局康复护理医疗产业，\allowbreak{}公司持有东方华康 60\% 股权 \\
2023-07-21 & 终止前次非公开发行股票 & --- & 第九届董事会第二次会议、\allowbreak{}第九届监事会第二次会议审议通过公司终止前次非公开发行股票事项 \\
2023-08-31 & 向特定对象发行 A 股股票（预案修订稿，未实施） & 拟募集资金总额（含发行费用）不超过 5.00 亿元 & 大湖水殖股份有限公司冰鲜冻鲜水产品深加工等项目 \\
2026-04-13 & 业绩补偿方以股权抵偿（受让子公司股权） & 业绩补偿方咖辅健康科技（上海）有限公司、\allowbreak{}上海擢英医疗管理合伙企业（有限合伙）、\allowbreak{}李爱川将其合计持有并已依法实缴的东方华康 1.1103\% 股权无偿转让给公司 & 抵偿业绩补偿方应向公司支付的业绩补偿金；交易完成后公司持有东方华康的比例由 60\% 增加至 61.1103\% \\
\bottomrule
\end{tabular}
\end{adjustbox}
\end{center}

\pfl{历史融投资情况} 从现金流量表看，2025 年公司取得借款收到的现金 2.73 亿元、偿还债务支付的现金 2.83 亿元，筹资活动现金流净额 -0.81 亿元（2023 年为取得借款 3.22 亿元、偿还债务 3.21 亿元）。 2026 年上半年筹资活动现金流净额为 -0.56 亿元（取得借款 1.79 亿元、偿还债务 1.62 亿元，未年化），已转为净流出，处于净还债状态。 2025 年分配股利、利润或偿付利息支付的现金合计 0.12 亿元。 公开披露的融资记录共 8 笔，工具类型以 IPO、定向增发、发行股份及支付现金购为主；金额与用途见下表。 其中规模最大的两笔为：2023 年 8 月，定向增发，拟募集资金总额（含发行费用）不超过 5.00 亿元；2016 年 10 月，定向增发，发行价格 10.15 元，实际发行数量 5,418.7188 万股。 股本方面，近三年披露股本变动 3 次；总股本自 2021 年末以来未发生变化（4.81 亿股）。 分红方面，近三年未实施现金分红，在全部样本中属于分红能力偏弱的一端。 近三年“再融资”类公告 2 条，涉及定向增发。 综合看，公司融资结构上股权与债务融资并举；2025 年筹资活动现金流净额 -0.81 亿元，融资节奏仍以维持存量债务周转为主，与前述经营与投资现金流的方向基本一致。

\pfl{未来潜在融资需求分析} 2026 年 6 月末货币资金 2.61 亿元，而短期债务 2.50 亿元（现金短债比 1.04 倍），2026 年上半年经营现金流净额 0.56 亿元。账面货币资金可以覆盖短期债务，缺口不体现在静态偿付层面。 公司 2026 年 6 月末资产负债率 46.4\%、2026 年上半年 ROE 2.1\%（未年化），财务结构相对宽松而盈利弹性有限，融资属于可选项而非必选项，触发条件多与产能建设或外延并购相关。 公开披露的后续资金安排线索包括：股权融资线索：2023 年 8 月，公司 2023 年度向特定对象发行 A 股股票预案（修订稿）；2024 年，第二次临时股东会审议通过 7 项议案。 综合看，公司在静态偿付层面不存在缺口，外部融资更多是配合扩张节奏与并购整合的主动性安排，而非偿债压力驱动。




\subsubsection{开创国际（600097）}
\label{co:600097}

\pfl{公司基本情况} 公司前身可追溯至 1993 年，上海国资（光明食品集团）控制的远洋渔业上市平台。前身为 1997 年上市的海南恒泰芒果，2008 年 12 月经资产置换重组为以远洋渔业为主业的开创国际；旗下拥有 12 艘大型金枪鱼围网船和 2 艘大型拖网加工船，并控股西班牙 ALBO 等境外水产企业。2025 年归母净利润 0.74 亿元，2026 年上半年因渔场资源波动、基里巴斯部分水域关闭与国际油价上涨转为亏损。 公司于 1997-06-19 在上交所首发上市，注册资本 2.41 亿元，总股本 2.41 亿股。 截至 2026 年 10 月 12 日收盘价 9.11 元，对应总市值 21.9 亿元，近一年-18.7\%，总市值在三级行业“海洋捕捞”的 2 家公司中按由大到小排第\zhnumber{2}位，近一年涨跌幅低于全样本中位数（-10.4\%）8.3 个百分点。 按 2026 年半年报披露的行业构成，收入占比前三位为食品加工 52.0\%；远洋捕捞 33.3\%；渔货贸易 14.5\%。 与 2021 年年报相比，第一大行业口径由“海洋捕捞”（56.7\%）变为“食品加工”（52.0\%）；两期披露的分类口径有调整，占比差异中同时包含分类因素。

\pfl{股权结构与控制权} 按 2026 年 6 月 30 日的股东名册，前三大股东为上海远洋渔业有限公司 持股 42.26\%；上海市北高新股份有限公司 持股 1.59\%；徐开东 持股 1.06\%，股东户数 15,618 户。 股权结构的另一条线索是，2017 年上海市国有资产监督管理委员会将所持上海水产集团 100\% 股权划转至光明食品（集团）有限公司。

\pfl{经营要点} 从经营要点看，2026 年上半年，盘中公司股票报 9.17 元/股、总市值 22.09 亿元。

\pfl{所处行业的基本面} 海洋捕捞的供给受资源与配额约束，需求端依赖餐饮与出口，成本端跟随燃油与人工（第 \ref{sec:framework} 章的供需框架）。 当期该行业的关键数据是：资源国调整入渔许可费用与监管规则将影响捕捞成本；2025 年，全国远洋渔业产量 247.64 万吨、同比增长 13.12\%。

\pfl{竞争优势} 把该公司放回同类公司比较，2026 年上半年营业收入 11.8 亿元、在三级行业“海洋捕捞”的 2 家公司中排第\zhnumber{2}位；毛利率 23.1\%、高于行业内中位数 8.8 个百分点；半年 ROE -1.3\%（未年化）、低于行业内中位数（9.79\%）；资产负债率 34.4\%、低于行业内中位数（53.0\%）。 公司在公开资料中自述与同业比较的相对优势在于：核心作业区位于中西太平洋并与多国签订渔业配额协议。 优势在报表上主要体现为高于行业平均的毛利率，这在周期底部是抵御亏损的主要缓冲。

\pfl{历史股价} 以 2021 年 1 月为起点、2026 年 10 月为终点，股价由 9.23 元变为 9.11 元，累计 -1.3\%，区间最高 12.93 元（2022 年 3 月）、最低 7.62 元（2024 年 8 月）；当前价相当于区间最高价的 70.5\%，为区间最低价的 1.20 倍。 月度收益的年化波动率 33.9\%，低于全样本中位数 37.3\%，在三级行业“海洋捕捞”的 2 家公司中波动率由高到低排第\zhnumber{2}位。 把股价阶段与公司事件对齐看，2024 年 8 月 至 2026 年 2 月股价上涨+63.5\%，同期（2025 年 12 月）公司 2017 年完成的非公开发行为上市以来可查证的最后一次股权融资。 从时间对应关系看，公司层面的资本运作与股权、风险事件解释了区间内多数急涨急跌，与行业指数的同步性相对有限。

\begin{center}
\begin{minipage}{\textwidth}\centering
\begin{tikzpicture}
\begin{axis}[
  width=12.6cm, height=3.9cm, scale only axis,
  xmin=2020.922, xmax=2026.972, ymin=7.09, ymax=13.84,
  xtick={2021,2022,2023,2024,2025,2026},
  yearticks,
  yticklabel style={font=\scriptsize},
  ylabel={元/股}, ylabel style={font=\scriptsize},
  grid=major, grid style={LightGray, line width=0.3pt},
  axis line style={InkGray, line width=0.5pt},
  every axis plot/.append style={line width=0.9pt},
]
\addplot[AgriGreen] table[x=x,y=y]{data/clean/profiles/px_600097.dat};
\addplot[SoftRed, dashed, line width=0.7pt] table[x=x,y=y]{data/clean/profiles/pxzz_600097.dat};
\addplot[only marks, mark=*, mark size=1.1pt, SoftRed] table[x=x,y=y]{data/clean/profiles/pxzz_600097.dat};
\end{axis}
\end{tikzpicture}

\captionof{figure}[开创国际（600097）股价与周期骨架]{开创国际（600097）月度收盘价（元/股）与 ZigZag 周期骨架（阈值 25\%，圆点为周期转折点）}
\end{minipage}
\end{center}

\pfl{过去周期分析} 以 25\% 的 ZigZag 阈值划分，样本区间内股价形成 3 个主升段与 4 个主跌段。 最大主升段出现在 2024 年 8 月 至 2026 年 2 月，18 个月+63.5\%；最大主跌段为 2026 年 2 月 至 2026 年 6 月，4 个月-37.6\%。 最近一次转折出现在 2026 年 6 月（下跌段自 2026 年 2 月起，4 个月-37.6\%），当前处于该段的延续之中。 公司股价较申万农林牧渔指数提前约 21 个月见底，是板块中较早定价周期反转的公司之一。 远洋与近海捕捞的产量受资源与配额约束，成本端则跟随燃油与人工，公司 2026 年 6 月末资产负债率 34.4\%、半年 ROE -1.3\%（未年化），在本轮下行中属于随行业同步调整的一类。 综合区间形态看，该股单段涨跌幅度偏大、以趋势段为主（最大主升 +63.5\%、最大主跌 -37.6\%），年化波动率 33.9\% 与全样本中位数 37.3\% 相比波动与样本中位数接近。

\pfl{盈利情况} 2026 年上半年营业收入 11.82 亿元（同比 -5.4\%），归母净利润 -0.30 亿元，由上年同期的盈利转为亏损。 以年报为趋势对照，2023---2025 年营业收入由 20.15 亿元变为 23.37 亿元（两年年均复合增速 +7.7\%），归母净利润由 1.48 亿元变为 0.74 亿元。 按半年报口径，2026 年上半年的 ROE 为 -1.3\%（未年化）、销售净利率 -2.5\%、毛利率 23.1\%，ROE 在“海洋捕捞”2 家公司中按数值由高到低排第\zhnumber{2}位。 按同一口径，三级行业“海洋捕捞”2 家公司的半年 ROE 中位数为 9.79\%，该公司低于其中位数 11.12 个百分点。 从公司披露的原因看，业绩变动主要来自：2025 年，第三季度毛利指标出现异常的核心原因是公司于当期收到上海市农业农村委员会拨付的国际履约能力提升政府补助约 5,353.86 万元；2026 年上半年，国际油价持续上涨、安排部分船只提前回国修理，同时燃油价格上涨导致船只生产成本增加。 针对上述压力，公司披露的应对举措包括：2026 年 6 月，船队通过优化作业天数配置、加大人工网投放、更新网具等措施缓解产量下滑。 综合看，公司仍处亏损区间、由盈转亏，半年 ROE 在三级行业内按由高到低排第\zhnumber{2}位，单位盈利能力的修复依赖行业景气与自身成本端的改善，报表弹性主要体现为价格与销量的方向性变化。

\begin{center}\small\setlength{\tabcolsep}{3.4pt}
\captionof{table}[开创国际（600097）关键财务指标]{开创国际（600097）关键财务指标（合并报表口径；两个半年报期均未年化；现金短债比 $=$ 货币资金 $\div$ 短期债务）}
\begin{adjustbox}{max width=\textwidth}
\begin{tabular}{@{}lrrrrr@{}}
\toprule
指标 & 2023 年 & 2024 年 & 2025 年 & 2025 年半年报 & 2026 年半年报 \\
\midrule
营业收入（亿元） & 20.15 & 23.16 & 23.37 & 12.50 & 11.82 \\
归母净利润（亿元） & 1.48 & 0.61 & 0.74 & 0.28 & -0.30 \\
毛利率（\%） & 32.5 & 28.1 & 29.7 & 31.5 & 23.1 \\
ROE（\%） & 6.7 & 2.7 & 3.2 & 1.2 & -1.3 \\
资产负债率（\%） & 32.6 & 32.3 & 31.5 & 30.0 & 34.4 \\
经营活动现金流净额（亿元） & 1.07 & 3.35 & 1.68 & 1.26 & 0.48 \\
货币资金（亿元） & 5.19 & 5.25 & 4.28 & 5.00 & 3.90 \\
有息负债合计（亿元） & 6.22 & 5.27 & 4.86 & 5.02 & 4.79 \\
现金短债比（倍） & 4.49 & 4.39 & 2.16 & 3.28 & 2.22 \\
\bottomrule
\end{tabular}
\end{adjustbox}
\end{center}

\begin{center}
\begin{minipage}{\textwidth}\centering
\begin{tikzpicture}
\begin{groupplot}[
  group style={group size=2 by 1, horizontal sep=0.60cm},
  width=6.85cm, height=4.60cm,
  tick label style={font=\tiny},
  title style={font=\scriptsize\heiti},
  yticklabel style={font=\tiny},
  ylabel style={font=\tiny},
  grid=major, grid style={LightGray, line width=0.3pt},
  axis line style={InkGray, line width=0.5pt},
  enlarge x limits={abs=0.8},
  legend style={font=\scriptsize, at={(0.5,-0.20)}, anchor=north,
                legend columns=2, draw=none, fill=none, column sep=6pt},
]
\nextgroupplot[
  ybar, bar width=4pt,
  ymin=-2.67, ymax=26.21,
  xtick={0,1,2,3,4,5.7},
  xticklabels={2021,2022,2023,2024,2025,2026H1},
  ylabel={亿元},
]
\addplot[fill=AgriGreen!75, draw=AgriGreen, bar shift=-2.6pt] table[x=i, y=rev]{data/clean/profiles/fin_600097.dat};
\addplot[fill=SoftRed!60, draw=SoftRed, bar shift=2.6pt] table[x=i, y=ni]{data/clean/profiles/fin_600097.dat};
\addplot[fill=AgriGreen!30, draw=AgriGreen, pattern=north east lines, bar shift=-2.6pt] table[x=x, y=rev]{data/clean/profiles/finh_600097.dat};
\addplot[fill=SoftRed!20, draw=SoftRed, pattern=north east lines, bar shift=2.6pt] table[x=x, y=ni]{data/clean/profiles/finh_600097.dat};
\legend{营业收入（年/半年）, 归母净利润（年/半年）}
\nextgroupplot[
  ymin=-4.5, ymax=42.0,
  xtick={0,1,2,3,4,5.7},
  xticklabels={2021,2022,2023,2024,2025,2026H1},
  ylabel={\%},
]
\addplot[AgriGold, mark=*, mark size=1.3pt] table[x=i, y=roe]{data/clean/profiles/fin_600097.dat};
\addplot[OilBlue, mark=square*, mark size=1.3pt] table[x=i, y=debt]{data/clean/profiles/fin_600097.dat};
\addplot[only marks, AgriGold, mark=*, mark size=2.0pt] table[x=x, y=roe]{data/clean/profiles/finh_600097.dat};
\addplot[only marks, OilBlue, mark=square*, mark size=2.0pt] table[x=x, y=debt]{data/clean/profiles/finh_600097.dat};
\legend{ROE, 资产负债率}
\end{groupplot}
\end{tikzpicture}
\begin{tikzpicture}
\begin{axis}[
  name=finbot,
  width=13.75cm, height=3.10cm, scale only axis,
  ymin=-0.76, ymax=8.49,
  xtick={0,1,2,3,4,5.7},
  xticklabels={2021,2022,2023,2024,2025,2026H1},
  tick label style={font=\tiny},
  yticklabel style={font=\tiny},
  ylabel={亿元}, ylabel style={font=\tiny},
  ybar, bar width=4pt,
  grid=major, grid style={LightGray, line width=0.3pt},
  axis line style={InkGray, line width=0.5pt},
  enlarge x limits={abs=0.8},
  legend style={font=\scriptsize, at={(0.5,-0.26)}, anchor=north,
                legend columns=2, draw=none, fill=none, column sep=10pt},
]
\addplot[fill=AgriGreen!55, draw=AgriGreen, bar shift=-3.0pt] table[x=i, y=cash]{data/clean/profiles/fin_600097.dat};
\addplot[fill=SoftRed!45, draw=SoftRed, bar shift=3.0pt] table[x=i, y=ibd]{data/clean/profiles/fin_600097.dat};
\addplot[fill=AgriGreen!25, draw=AgriGreen, pattern=north east lines, bar shift=-3.0pt] table[x=x, y=cash]{data/clean/profiles/finh_600097.dat};
\addplot[fill=SoftRed!20, draw=SoftRed, pattern=north east lines, bar shift=3.0pt] table[x=x, y=ibd]{data/clean/profiles/finh_600097.dat};
\legend{货币资金（左轴）, 有息负债（左轴）}
\end{axis}
\begin{axis}[
  at={(finbot.south east)}, anchor=south east,
  width=13.75cm, height=3.10cm, scale only axis,
  axis y line*=right, axis x line*=none, xtick=\empty,
  ymin=-0.87, ymax=4.36,
  tick label style={font=\tiny},
  yticklabel style={font=\tiny}, scaled y ticks=false,
  legend style={font=\scriptsize, at={(0.5,-0.44)}, anchor=north,
                legend columns=1, draw=none, fill=none},
]
\addplot[OilBlue, mark=*, mark size=2.2pt, line width=1.3pt] table[x=i, y=ocf]{data/clean/profiles/fin_600097.dat};
\addplot[only marks, OilBlue, mark=*, mark size=3.2pt] table[x=x, y=ocf]{data/clean/profiles/finh_600097.dat};
\legend{经营活动现金流净额（右轴，亿元，蓝点）}
\end{axis}
\end{tikzpicture}

\captionof{figure}[开创国际（600097）近年主要财务指标]{开创国际（600097）近年主要财务指标（左上：营业收入与归母净利润；右上：ROE 与资产负债率；下：货币资金、有息负债为左轴柱状，经营活动现金流净额为其右轴的蓝点折线。
2021---2025 年为年报口径，2026H1 为 2026 年半年报/半年末，与其前一年报之间留出间隔、以斜纹或空心点区分；半年报数据未年化）}
\end{minipage}
\end{center}

\pfl{财务分析} 2026 年 6 月末资产负债率 34.4\%，较 2025 年末上升 2.9 个百分点（样本中位数 52.2\%，所处三级行业内按由低到高排第\zhnumber{1}位）。 2026 年 6 月末货币资金 3.90 亿元、短期债务（短期借款与一年内到期非流动负债）1.76 亿元、长期债务 3.03 亿元，有息负债合计 4.79 亿元，现金短债比 2.22 倍——覆盖充分。 净有息负债 0.89 亿元，相当于其 2026 年上半年营业收入的 8\%。 流动比率 2.05、速动比率 1.09，短期指标尚可。 2026 年上半年经营活动现金流净额 0.48 亿元，净利润为负而经营现金流为正，差额主要来自折旧摊销与营运资本变动，说明亏损尚未完全转化为现金流出。 2026 年 6 月末在建工程 1.49 亿元，较上年末增加 0.74 亿元，仍有在建投入。 2026 年 6 月末商誉 1.04 亿元，占归母净资产的 4.7\%，是净资产中最脆弱的一块。 综合看，现金短债比 2.22 倍，短期偿付压力可控；资产负债率低于样本中位数 17.8 个百分点；有息负债中长期债务占比更高，期限结构对短债滚动的压力较小。 公司披露的监管与合规风险为：2026 年 7 月，公司就上海证券交易所对 2025 年年度报告及 2026 年一季度报告的信息披露监管问询函作出回复（致同会计师事务所出具相关回复）。

\pfl{重大战略转型} 从公开披露看，公司的战略推进集中在：并购与重组（3 项）：2017 年，公司以非公开发行募集资金（总额调整至 6.00 亿元）用于收购西班牙 ALBO 公司股权及补充流动资金；2016 年 6 月，S.L. 100\% 股权的现金收购交割。 公告口径上，近三年（2023 年 9 月---2026 年 9 月）公司披露重大重组与并购类公告 2 条、对外投资与转型类公告 1 条、回购与股权激励类公告 0 条，合计公告 262 条。 主营结构的口径变化已经落到报表上：2021 年年报的第一大行业为“海洋捕捞”（56.7\%），2026 年半年报为“食品加工”（52.0\%）；公司在这两期的分部划分方式有过调整。 综合判断，报告期内公司有 3 条与战略相关的公告落地（重大重组与并购 2 条、对外投资与转型 1 条），资本运作与主业调整之间存在直接关联。

\begin{center}\small\setlength{\tabcolsep}{4pt}
\captionof{table}[开创国际（600097）历史融资与资本运作]{开创国际（600097）历史融资与资本运作（据公司公告与公开报道整理；金额为披露值，含首次公开发行、再融资、债券、银行借款与重整投资等）}
\begin{adjustbox}{max width=\textwidth}
\begin{tabular}{@{}p{1.45cm}p{2.55cm}p{4.05cm}p{4.95cm}@{}}
\toprule
时间 & 方式 & 规模 & 用途与说明 \\
\midrule
1997-05-21 & IPO（上海证券交易所） & 首次向社会公众发行人民币普通股 3,\allowbreak{}500 万股 & 经海南省证券管理办公室（琼证办（1997）92 号）文件与中国证券监督管理委员会（证监发审字（1997）134、\allowbreak{}235 号）文件批准 \\
2015-09 & 非公开发行股票（筹划阶段，规模后调整） & 初步方案拟募集资金约 11.00 亿元 & 收购控股股东上海远洋渔业持有的斐济金洋渔业有限公司 20\% 股权并对其增资至持有其 59.375\% 股权；收购远洋渔业持有的基里巴斯渔业公司 20\% 股权；公司股票经四次延期复牌 \\
2016-04 & 非公开发行股票（预案修订稿，拟募集资金总额不超过 84,000 万元） & 拟募集资金总额不超过 8.40 亿元（含发行费用） & 收购西班牙 ALBO 公司 100\% 股权（拟投入募集资金 4.26 亿元）、\allowbreak{}新建舟山金枪鱼食品加工基地项目（拟投入 2.00 亿元）、\allowbreak{}补充流动资金（拟投入 2.14 亿元）；ALBO 公司 100\% 股权购买价款为 60,\allowbreak{}999,\allowbreak{}840 欧元 \\
2016-05 & 重大资产购买（现金收购，募集资金置换） & 以下属全资子公司开创远洋为收购主体，\allowbreak{}S.L. 100\% 股权 & 公司拟以非公开发行募集资金 4.26 亿元投入本次收购 \\
2016-12-15 & 非公开发行股票（反馈与承诺披露） & 控股股东上海远洋渔业有限公司承诺认购不低于本次非公开发行最终确定发行股票总数量的 20\%（含 20\%） & 不适用（承诺事项）；本次非公开发行的董事会决议日为 2015 年 11 月 24 日；上海远洋及其控制的主体上海宇洋人力资源有限公司、\allowbreak{}上海蒂尔远洋渔业有限公司、\allowbreak{}上海金优远洋渔业有限公司、\allowbreak{}上海水锦洋食品有限公司、\allowbreak{}海南双海海洋渔业有限公司等一致行动人承诺自定价基准日前六个月至本次发行完成后六个月不存在主动减持计划 \\
2017-12-07 & 非公开发行股票（完成） & 发行人民币普通股 3,\allowbreak{}833.87 万股，\allowbreak{}发行价 15.65 元/\allowbreak{}股 & 收购西班牙 ALBO 公司 100\% 股权及补充流动资金（募集资金总额后经削减调整至 6.00 亿元）；经公司第七届董事会第十一次（临时）会议、\allowbreak{}第十三次（临时）会议、\allowbreak{}第十八次（临时）会议及 2016 年第一次临时股东会决议 \\
\bottomrule
\end{tabular}
\end{adjustbox}
\end{center}

\pfl{历史融投资情况} 从现金流量表看，2025 年公司取得借款收到的现金 0.24 亿元、偿还债务支付的现金 0.16 亿元，筹资活动现金流净额 -0.86 亿元（2023 年为偿还债务 0.15 亿元）。 2026 年上半年筹资活动现金流净额为 0.40 亿元（取得借款 1.23 亿元、偿还债务 0.10 亿元，未年化），仍处于净融入状态。 2025 年分配股利、利润或偿付利息支付的现金合计 0.22 亿元。 公开披露的融资记录共 6 笔，工具类型以 IPO、定向增发为主；金额与用途见下表。 其中规模最大的两笔为：2016 年 4 月，定向增发，拟募集资金总额不超过 8.40 亿元（含发行费用）；2015 年 9 月，定向增发，初步方案拟募集资金约 11.00 亿元。 股本方面，近三年披露股本变动 3 次；总股本自 2021 年末以来未发生变化（2.41 亿股）。 分红方面，近三年现金分红 4 次、合计每股派现 0.50 元，最近一次实施在 2025 年 12 月。 综合看，公司融资结构上股权与债务融资并举；2025 年筹资活动现金流净额 -0.86 亿元，融资节奏仍以维持存量债务周转为主，与前述经营与投资现金流的方向基本一致。

\pfl{未来潜在融资需求分析} 2026 年 6 月末货币资金 3.90 亿元，而短期债务 1.76 亿元（现金短债比 2.22 倍），2026 年上半年经营现金流净额 0.48 亿元。账面货币资金可以覆盖短期债务，缺口不体现在静态偿付层面。 公司 2026 年上半年归母净利润为负（-0.30 亿元），但 6 月末资产负债率 34.4\% 尚未进入第一档（亏损且负债率 $>$65\%）的区间，当前融资需求以补充流动资金、支撑低谷期经营性支出为主。 公开披露的后续资金安排线索包括：股权融资线索：2025 年 12 月，公司 2017 年完成的非公开发行为上市以来可查证的最后一次股权融资。 资金需求还要叠加在建工程：期末在建工程 1.49 亿元、2026 年上半年购建固定资产等支出 1.20 亿元（未年化），这部分支出在周期底部通常会被主动放缓。 静态偿付不存在缺口，后续融资更可能服务于产能或并购安排，而不是填补流动性窟窿。




\subsubsection{永安林业（000663）}
\label{co:000663}

\pfl{公司基本情况} 1994 年设立至今，永安林业是全国首家林业综合性股份制公司，1996 年 12 月深交所上市，主营森林资源培育与林木采伐经营、人造板与地板家居（林板一体化）；2021 年 2 月中国林业集团完成战略重组、成为间接控股股东。2023—2025 年归母净利润由盈转亏并连续两年亏损，生物质能循环利用项目停建并计提减值、引发诉讼，中林集团避免同业竞争的五年承诺已逾期未履行。 公司于 1996-12-06 在深交所主板首发上市，注册资本 3.36 亿元，总股本 3.37 亿股。 截至 2026 年 10 月 12 日收盘价 5.92 元，对应总市值 19.9 亿元，近一年-10.4\%，总市值在三级行业“林业Ⅲ”的 4 家公司中按由大到小排第\zhnumber{4}位，近一年涨跌幅低于全样本中位数（-10.4\%）0.1 个百分点。 按 2026 年半年报披露的行业构成，收入占比前三位为人造板制造业 93.1\%；其他(补充) 4.2\%；林业 2.7\%。

\pfl{股权结构与控制权} 按 2026 年 6 月 30 日的股东名册，前三大股东为中林(永安)控股有限公司 持股 19.27\%；天风证券股份有限公司 持股 6.40\%；周爽（持股比例未披露），股东户数 18,921 户。 控制权与股东性质方面，截至 2025 年 12 月 31 日公司控股股东中林（永安）控股有限公司持股 19.27\%；2026 年 6 月末持有公司 19.27\% 股份（6,488.46 万股）；2021 年年报，公司作为中林集团旗下唯一上市平台推进福建区域专业化整合。

\pfl{经营要点} 业务与产能层面，2026 年上半年，生产与销售木材均为 0.07 万立方米、同比减少逾 84\%。

\pfl{所处行业的基本面} 林业的供给由林木蓄积与采伐限额决定、需求来自木材加工与基建，变现周期长（第 \ref{sec:framework} 章的供需框架）。 当期该行业的关键数据是：2026 年 7 月，前 7 个月进口总量同比下降 6.7\%，2026 年 1—2 月累计进口 750.3 万立方米、同比下降 11.2\%。

\pfl{竞争优势} 把该公司放回同类公司比较，2026 年上半年营业收入 1.5 亿元、在三级行业“林业Ⅲ”的 4 家公司中排第\zhnumber{2}位；毛利率 11.9\%、低于行业内中位数 10.5 个百分点；半年 ROE -2.1\%（未年化）、低于行业内中位数（-1.96\%）；资产负债率 36.6\%、低于行业内中位数（38.0\%）。 公司区别于同类公司的经营特征包括：作为中林集团唯一上市平台推进国家储备林建设与福建区域专业化整合。 就当期报表看，公司尚未把经营层面的特征转化为高于行业的资本回报，周期底部的相对优势更多体现在现金流与资产负债表上。

\pfl{历史股价} 本报告样本区间（2021 年 1 月 至 2026 年 10 月）内，股价由 4.25 元变为 5.92 元，累计 +39.3\%，区间最高 12.18 元（2023 年 6 月）、最低 4.19 元（2024 年 6 月）；当前价相当于区间最高价的 48.6\%，为区间最低价的 1.41 倍。 月度收益的年化波动率 51.3\%，高于全样本中位数 37.3\%，在三级行业“林业Ⅲ”的 4 家公司中波动率由高到低排第\zhnumber{3}位。 把股价阶段与公司事件对齐看，2024 年 6 月 至 2025 年 11 月股价上涨+108.1\%，同期（2025 年 10 月）公司全资孙公司向永安市人民法院提起对必奥新能源科技有限公司的买卖合同纠纷诉讼。 从时间对应关系看，公司层面的资本运作与股权、风险事件解释了区间内多数急涨急跌，与行业指数的同步性相对有限。

\begin{center}
\begin{minipage}{\textwidth}\centering
\begin{tikzpicture}
\begin{axis}[
  width=12.6cm, height=3.9cm, scale only axis,
  xmin=2020.922, xmax=2026.972, ymin=3.90, ymax=13.03,
  xtick={2021,2022,2023,2024,2025,2026},
  yearticks,
  yticklabel style={font=\scriptsize},
  ylabel={元/股}, ylabel style={font=\scriptsize},
  grid=major, grid style={LightGray, line width=0.3pt},
  axis line style={InkGray, line width=0.5pt},
  every axis plot/.append style={line width=0.9pt},
]
\addplot[AgriGreen] table[x=x,y=y]{data/clean/profiles/px_000663.dat};
\addplot[SoftRed, dashed, line width=0.7pt] table[x=x,y=y]{data/clean/profiles/pxzz_000663.dat};
\addplot[only marks, mark=*, mark size=1.1pt, SoftRed] table[x=x,y=y]{data/clean/profiles/pxzz_000663.dat};
\end{axis}
\end{tikzpicture}

\captionof{figure}[永安林业（000663）股价与周期骨架]{永安林业（000663）月度收盘价（元/股）与 ZigZag 周期骨架（阈值 25\%，圆点为周期转折点）}
\end{minipage}
\end{center}

\pfl{过去周期分析} 以 25\% 的 ZigZag 阈值划分，样本区间内股价形成 4 个主升段与 4 个主跌段。 最大主跌段为 2023 年 6 月 至 2024 年 6 月，12 个月-65.6\%；最大主升段出现在 2021 年 1 月 至 2021 年 6 月，5 个月+183.5\%。 最近一次转折出现在 2026 年 6 月（下跌段自 2025 年 11 月起，7 个月-41.5\%），当前处于该段的延续之中。 公司股价较申万农林牧渔指数提前约 23 个月见底，是板块中较早定价周期反转的公司之一。 林业的现金流依赖林木采伐与政策补贴，周期长、变现慢，公司 2026 年 6 月末资产负债率 36.6\%、半年 ROE -2.1\%（未年化），在本轮下行中属于随行业同步调整的一类。 综合区间形态看，该股单段涨跌幅度偏大、以趋势段为主（最大主升 +183.5\%、最大主跌 -65.6\%），年化波动率 51.3\% 与全样本中位数 37.3\% 相比波动明显大于样本中位数。

\pfl{盈利情况} 2026 年上半年营业收入 1.51 亿元（同比 +14.7\%），归母净利润 -0.19 亿元，亏损同比扩大 0.02 亿元。 以年报为趋势对照，2023---2025 年营业收入由 6.93 亿元变为 3.24 亿元（两年年均复合增速 -31.6\%），归母净利润由 1.90 亿元变为 -0.95 亿元。 按半年报口径，2026 年上半年的 ROE 为 -2.1\%（未年化）、销售净利率 -13.2\%、毛利率 11.9\%，ROE 在“林业Ⅲ”4 家公司中按数值由高到低排第\zhnumber{3}位。 同期全样本半年 ROE 中位数 0.75\%，公司与之相差 2.86 个百分点（弱于中位数公司）。 从公司披露的原因看，业绩变动主要来自：2025 年度，归属于上市公司股东的净利润 -9,510.41 万元（同比下降 11.84\%）；2026 年上半年，主要系生物质项目停建后产生的诉讼费用、资产减值及政府补助同比减少等因素影响。 综合看，公司仍处亏损区间、亏损同比扩大，半年 ROE 在三级行业内按由高到低排第\zhnumber{3}位，单位盈利能力的修复依赖行业景气与自身成本端的改善，报表弹性主要体现为价格与销量的方向性变化。

\begin{center}\small\setlength{\tabcolsep}{3.4pt}
\captionof{table}[永安林业（000663）关键财务指标]{永安林业（000663）关键财务指标（合并报表口径；两个半年报期均未年化；现金短债比 $=$ 货币资金 $\div$ 短期债务）}
\begin{adjustbox}{max width=\textwidth}
\begin{tabular}{@{}lrrrrr@{}}
\toprule
指标 & 2023 年 & 2024 年 & 2025 年 & 2025 年半年报 & 2026 年半年报 \\
\midrule
营业收入（亿元） & 6.93 & 3.23 & 3.24 & 1.32 & 1.51 \\
归母净利润（亿元） & 1.90 & -0.85 & -0.95 & -0.18 & -0.19 \\
毛利率（\%） & 44.7 & 9.9 & 11.4 & 8.0 & 11.9 \\
ROE（\%） & 18.9 & -8.0 & -9.8 & -1.7 & -2.1 \\
资产负债率（\%） & 32.2 & 33.6 & 37.3 & 32.4 & 36.6 \\
经营活动现金流净额（亿元） & 3.49 & 1.17 & -0.17 & 0.09 & -0.31 \\
货币资金（亿元） & 2.76 & 2.78 & 2.13 & 2.60 & 1.68 \\
有息负债合计（亿元） & 1.92 & 0.82 & 0.55 & 0.56 & 0.55 \\
现金短债比（倍） & 12.47 & 9.32 & 6.79 & 10.43 & 5.28 \\
\bottomrule
\end{tabular}
\end{adjustbox}
\end{center}

\begin{center}
\begin{minipage}{\textwidth}\centering
\begin{tikzpicture}
\begin{groupplot}[
  group style={group size=2 by 1, horizontal sep=0.60cm},
  width=6.85cm, height=4.60cm,
  tick label style={font=\tiny},
  title style={font=\scriptsize\heiti},
  yticklabel style={font=\tiny},
  ylabel style={font=\tiny},
  grid=major, grid style={LightGray, line width=0.3pt},
  axis line style={InkGray, line width=0.5pt},
  enlarge x limits={abs=0.8},
  legend style={font=\scriptsize, at={(0.5,-0.20)}, anchor=north,
                legend columns=2, draw=none, fill=none, column sep=6pt},
]
\nextgroupplot[
  ybar, bar width=4pt,
  ymin=-1.80, ymax=8.59,
  xtick={0,1,2,3,4,5.7},
  xticklabels={2021,2022,2023,2024,2025,2026H1},
  ylabel={亿元},
]
\addplot[fill=AgriGreen!75, draw=AgriGreen, bar shift=-2.6pt] table[x=i, y=rev]{data/clean/profiles/fin_000663.dat};
\addplot[fill=SoftRed!60, draw=SoftRed, bar shift=2.6pt] table[x=i, y=ni]{data/clean/profiles/fin_000663.dat};
\addplot[fill=AgriGreen!30, draw=AgriGreen, pattern=north east lines, bar shift=-2.6pt] table[x=x, y=rev]{data/clean/profiles/finh_000663.dat};
\addplot[fill=SoftRed!20, draw=SoftRed, pattern=north east lines, bar shift=2.6pt] table[x=x, y=ni]{data/clean/profiles/finh_000663.dat};
\legend{营业收入（年/半年）, 归母净利润（年/半年）}
\nextgroupplot[
  ymin=-14.6, ymax=56.2,
  xtick={0,1,2,3,4,5.7},
  xticklabels={2021,2022,2023,2024,2025,2026H1},
  ylabel={\%},
]
\addplot[AgriGold, mark=*, mark size=1.3pt] table[x=i, y=roe]{data/clean/profiles/fin_000663.dat};
\addplot[OilBlue, mark=square*, mark size=1.3pt] table[x=i, y=debt]{data/clean/profiles/fin_000663.dat};
\addplot[only marks, AgriGold, mark=*, mark size=2.0pt] table[x=x, y=roe]{data/clean/profiles/finh_000663.dat};
\addplot[only marks, OilBlue, mark=square*, mark size=2.0pt] table[x=x, y=debt]{data/clean/profiles/finh_000663.dat};
\legend{ROE, 资产负债率}
\end{groupplot}
\end{tikzpicture}
\begin{tikzpicture}
\begin{axis}[
  name=finbot,
  width=13.75cm, height=3.10cm, scale only axis,
  ymin=-0.52, ymax=5.86,
  xtick={0,1,2,3,4,5.7},
  xticklabels={2021,2022,2023,2024,2025,2026H1},
  tick label style={font=\tiny},
  yticklabel style={font=\tiny},
  ylabel={亿元}, ylabel style={font=\tiny},
  ybar, bar width=4pt,
  grid=major, grid style={LightGray, line width=0.3pt},
  axis line style={InkGray, line width=0.5pt},
  enlarge x limits={abs=0.8},
  legend style={font=\scriptsize, at={(0.5,-0.26)}, anchor=north,
                legend columns=2, draw=none, fill=none, column sep=10pt},
]
\addplot[fill=AgriGreen!55, draw=AgriGreen, bar shift=-3.0pt] table[x=i, y=cash]{data/clean/profiles/fin_000663.dat};
\addplot[fill=SoftRed!45, draw=SoftRed, bar shift=3.0pt] table[x=i, y=ibd]{data/clean/profiles/fin_000663.dat};
\addplot[fill=AgriGreen!25, draw=AgriGreen, pattern=north east lines, bar shift=-3.0pt] table[x=x, y=cash]{data/clean/profiles/finh_000663.dat};
\addplot[fill=SoftRed!20, draw=SoftRed, pattern=north east lines, bar shift=3.0pt] table[x=x, y=ibd]{data/clean/profiles/finh_000663.dat};
\legend{货币资金（左轴）, 有息负债（左轴）}
\end{axis}
\begin{axis}[
  at={(finbot.south east)}, anchor=south east,
  width=13.75cm, height=3.10cm, scale only axis,
  axis y line*=right, axis x line*=none, xtick=\empty,
  ymin=-1.55, ymax=5.85,
  tick label style={font=\tiny},
  yticklabel style={font=\tiny}, scaled y ticks=false,
  legend style={font=\scriptsize, at={(0.5,-0.44)}, anchor=north,
                legend columns=1, draw=none, fill=none},
]
\addplot[OilBlue, mark=*, mark size=2.2pt, line width=1.3pt] table[x=i, y=ocf]{data/clean/profiles/fin_000663.dat};
\addplot[only marks, OilBlue, mark=*, mark size=3.2pt] table[x=x, y=ocf]{data/clean/profiles/finh_000663.dat};
\legend{经营活动现金流净额（右轴，亿元，蓝点）}
\end{axis}
\end{tikzpicture}

\captionof{figure}[永安林业（000663）近年主要财务指标]{永安林业（000663）近年主要财务指标（左上：营业收入与归母净利润；右上：ROE 与资产负债率；下：货币资金、有息负债为左轴柱状，经营活动现金流净额为其右轴的蓝点折线。
2021---2025 年为年报口径，2026H1 为 2026 年半年报/半年末，与其前一年报之间留出间隔、以斜纹或空心点区分；半年报数据未年化）}
\end{minipage}
\end{center}

\pfl{财务分析} 2026 年 6 月末资产负债率 36.6\%，较 2025 年末下降 0.7 个百分点（样本中位数 52.2\%，所处三级行业内按由低到高排第\zhnumber{2}位）。 2026 年 6 月末货币资金 1.68 亿元、短期债务（短期借款与一年内到期非流动负债）0.32 亿元、长期债务 0.23 亿元，有息负债合计 0.55 亿元，现金短债比 5.28 倍——覆盖充分。 净有息负债 -1.13 亿元，相当于其 2026 年上半年营业收入的 -75\%。 流动比率 1.56、速动比率 0.56，短期流动性无明显缺口。 2026 年上半年经营活动现金流净额 -0.31 亿元，经营现金流与利润同时为负，现金消耗较快。 2026 年 6 月末在建工程 2.79 亿元，较上年末增加 0.06 亿元，仍有在建投入。 综合看，现金短债比 5.28 倍，短期偿付压力可控；资产负债率低于样本中位数 15.6 个百分点；有息负债中短期债务占比更高，期限结构仍以滚动续作为主。 公司披露的诉讼与资产冻结风险为：2026 年 7 月，闽民终 658 号民事判决向最高人民法院申请再审的张强等 17 人诉讼案（涉案金额 3,172.55 万元）被最高人民法院裁定驳回再审申请；2025 年 10 月，公司全资孙公司向永安市人民法院提起对必奥新能源科技有限公司的买卖合同纠纷诉讼。

\pfl{重大战略转型} 近三年公司的战略动作集中在：其他动作（2 项）：2021 年 2 月，林业发展集团完成股份过户登记手续。 公告口径上，近三年（2023 年 9 月---2026 年 9 月）公司披露重大重组与并购类公告 10 条、对外投资与转型类公告 1 条、回购与股权激励类公告 7 条，合计公告 370 条。 第一大行业“人造板制造业”的收入占比由 2021 年年报的 57.4\% 变为 93.1\%（35.7 个百分点），收入结构出现再平衡。 综合判断，报告期内公司有 11 条与战略相关的公告落地（重大重组与并购 10 条、对外投资与转型 1 条），资本运作与主业调整之间存在直接关联。

\begin{center}\small\setlength{\tabcolsep}{4pt}
\captionof{table}[永安林业（000663）历史融资与资本运作]{永安林业（000663）历史融资与资本运作（据公司公告与公开报道整理；金额为披露值，含首次公开发行、再融资、债券、银行借款与重整投资等）}
\begin{adjustbox}{max width=\textwidth}
\begin{tabular}{@{}p{1.45cm}p{2.55cm}p{4.05cm}p{4.95cm}@{}}
\toprule
时间 & 方式 & 规模 & 用途与说明 \\
\midrule
1996-11 & IPO（深圳证券交易所） & 向社会公开募集股票 1,\allowbreak{}950 万股，\allowbreak{}共募得资金 9,\allowbreak{}150.00 万元 & 招股说明书承诺投资项目（后续发生变更，\allowbreak{}变更原因与程序已在 1997 年年报披露）；公司 1996 年 12 月 6 日上市；1997 年 5 月 12 日股东大会决议向全体股东每 10 股派送 1 股红股并以资本公积按 10:7 转增股本 \\
2015-09 & 发行股份及支付现金购买资产并募集配套资金（重大资产重组） & 以发行股份方式支付 12.50 亿元、\allowbreak{}发行价格 11.75 元/\allowbreak{}股 & 本次交易现金对价支付、\allowbreak{}森源股份信息系统升级改造和营销与服务网点建设、\allowbreak{}补充上市公司流动资金；发行股份购买资产所发行新股上市日为 2015 年 10 月 8 日，\allowbreak{}增发后公司总股本 3.09 亿股；配套融资所发行新股上市首日为 2015 年 12 月 23 日 \\
2023-11-20 & 2022 年度向特定对象发行 A 股股票（终止） & 拟募集资金总额不超过人民币 3.00 亿元，\allowbreak{}发行数量不超过 4,\allowbreak{}360.47 万股 & 涿州市生物质能循环利用项目、\allowbreak{}宁晋县生物质能循环利用项目、\allowbreak{}生物质发酵微生物研发中心项目；本次发行方案于 2022 年 9 月 30 日经第九届董事会第二十八次会议审议通过、\allowbreak{}2022 年 10 月 17 日经 2022 年第五次临时股东大会审议通过；2023 年 11 月 20 日第十届董事会第三次会议和第十届监事会第三次会议审议通过终止公司 2022 年度向特定对象发行股票事项并撤回申请文件并向深交所申请撤回相关申请文件 \\
\bottomrule
\end{tabular}
\end{adjustbox}
\end{center}

\pfl{历史融投资情况} 从现金流量表看，2025 年公司取得借款收到的现金 0.23 亿元、偿还债务支付的现金 0.49 亿元，筹资活动现金流净额 -0.29 亿元（2023 年为取得借款 1.88 亿元）。 2026 年上半年筹资活动现金流净额为 -0.01 亿元（取得借款 — 亿元、偿还债务 — 亿元，未年化），已转为净流出，处于净还债状态。 2025 年分配股利、利润或偿付利息支付的现金合计 0.02 亿元。 公开披露的融资记录共 3 笔，工具类型以 IPO、发行股份及支付现金购、定向增发为主；金额与用途见下表。 其中规模最大的两笔为：2015 年 9 月，发行股份及支付现金购，以发行股份方式支付 12.50 亿元、发行价格 11.75 元/股；1996 年 11 月，IPO，向社会公开募集股票 1,950 万股，共募得资金 9,150.00 万元。 股本方面，近三年披露股本变动 3 次；总股本自 2021 年末以来未发生变化（3.37 亿股）。 分红方面，近三年未实施现金分红，在全部样本中属于分红能力偏弱的一端。 近三年“再融资”类公告 2 条，涉及定向增发。 综合看，公司融资结构上股权与债务融资并举；2025 年筹资活动现金流净额 -0.29 亿元，融资节奏仍以维持存量债务周转为主，与前述经营与投资现金流的方向基本一致。

\pfl{未来潜在融资需求分析} 2026 年 6 月末货币资金 1.68 亿元，而短期债务 0.32 亿元（现金短债比 5.28 倍），2026 年上半年经营现金流净额 -0.31 亿元。账面货币资金可以覆盖短期债务，缺口不体现在静态偿付层面。 按第 \ref{sec:financing-need} 节的三档划分，公司属于第二档“保壳型”：近三年重大重组与并购类公告 10 条，2026 年上半年归母净利润 -0.19 亿元，融资需求更多体现为“资产置换 + 维持上市地位”，而不是扩产。 公开披露的后续资金安排线索包括：股权融资线索：2023 年 11 月，原计划由中林集团全额认购、募资不超过 3.00 亿元的股权融资方案落空。 资金需求还要叠加在建工程：期末在建工程 2.79 亿元、2026 年上半年购建固定资产等支出 0.15 亿元（未年化），这部分支出在周期底部通常会被主动放缓。 以现有货币资金与经营现金流衡量，公司不存在刚性补血的紧迫性，融资动作更多是主动型的。






\subsubsection*{农业综合Ⅱ（2 家）}
\addcontentsline{toc}{subsubsection}{农业综合Ⅱ（2 家）}

% 农业上市公司画像：农业综合Ⅱ（共 2 家，按总市值降序）
% 由 scripts/make_company_profiles.py 自动生成，请勿手工修改

\subsubsection{辉隆股份（002556）}
\label{co:002556}

\pfl{公司基本情况} 辉隆股份成立于 2011 年，安徽省供销社系统背景的农资流通龙头，形成化肥、农药等农资流通与精细化工“双主业”，并以 14 座现代农业综合服务中心为载体构建“耕种管收售”全周期农业服务体系。 公司于 2011-03-02 在深交所主板首发上市，注册资本 9.35 亿元，总股本 9.47 亿股。 截至 2026 年 10 月 12 日收盘价 5.12 元，对应总市值 48.5 亿元，近一年-7.8\%，总市值在三级行业“—”的 2 家公司中按由大到小排第\zhnumber{1}位，近一年涨跌幅高于全样本中位数（-10.4\%）2.6 个百分点。 按 2026 年半年报披露的行业构成，收入占比前三位为农资产品 73.2\%；农副产品及其他 12.3\%；化工产品 8.1\%。 与 2021 年年报相比，第一大行业口径由“化肥”（62.1\%）变为“农资产品”（73.2\%）；两期披露的分类口径有调整，占比差异中同时包含分类因素。 股权结构方面，实际控制人为安徽省供销合作社联合社。

\pfl{股权结构与控制权} 按 2026 年 6 月 30 日的股东名册，前三大股东为安徽辉隆投资集团有限公司 持股 37.74\%；香港中央结算有限公司 持股 5.55\%；郑建平 持股 1.68\%，股东户数 31,836 户。 从控制权角度看，2026 年 6 月末持股 37.74\%（3.57 亿股），报告期内控股股东未发生变更；公司是全国供销社系统农资流通行业的第一家上市公司，自 2011 年上市以来累计实施现金分红 18 次、累计分红金额超 17.00 亿元；报告期末公司回购专用证券账户持有 2,738.17 万股。

\pfl{经营要点} 经营层面，拥有六大肥药生产基地和超百万吨自主品牌复合肥产能；全资子公司海华科技薄荷醇产能 5,000 吨/年、BHT 产能 1.5 万吨/年。

\pfl{所处行业的基本面} 农业综合类公司的收入分布在多个品种上，供需由各自品种决定，整体上更接近农资与服务的需求驱动（第 \ref{sec:framework} 章的供需框架）。 行业层面的最新情况是：2026 年 9 月，中旬全国尿素（中小颗粒）价格 1,792.5 元/吨（较上期+1.2\%）。

\pfl{竞争优势} 把该公司放回同类公司比较，2026 年上半年营业收入 85.5 亿元、在三级行业“农业综合Ⅱ”的 2 家公司中排第\zhnumber{1}位；毛利率 6.2\%、低于行业内中位数 7.2 个百分点；半年 ROE 3.1\%（未年化）、高于行业内中位数（1.52\%）；资产负债率 67.0\%、低于行业内中位数（69.3\%）。 优势在当期体现为高于行业中位数的资本回报，但优势能否持续仍取决于周期与自身产能投放节奏。

\pfl{历史股价} 自 2021 年 1 月至 2026 年 10 月，股价由 8.11 元变为 5.12 元，累计 -36.9\%，区间最高 12.87 元（2021 年 12 月）、最低 4.17 元（2024 年 6 月）；当前价相当于区间最高价的 39.8\%，为区间最低价的 1.23 倍。 月度收益的年化波动率 30.3\%，低于全样本中位数 37.3\%，在三级行业“—”的 2 家公司中波动率由高到低排第\zhnumber{1}位。 把股价阶段与公司事件对齐看，2024 年 6 月 至 2026 年 2 月股价上涨+61.2\%，同期（2025 年 12 月）受限货币资金 3.84 亿元。 从时间对应关系看，公司层面的资本运作与股权、风险事件解释了区间内多数急涨急跌，与行业指数的同步性相对有限。

\begin{center}
\begin{minipage}{\textwidth}\centering
\begin{tikzpicture}
\begin{axis}[
  width=12.6cm, height=3.9cm, scale only axis,
  xmin=2020.922, xmax=2026.972, ymin=3.88, ymax=13.77,
  xtick={2021,2022,2023,2024,2025,2026},
  yearticks,
  yticklabel style={font=\scriptsize},
  ylabel={元/股}, ylabel style={font=\scriptsize},
  grid=major, grid style={LightGray, line width=0.3pt},
  axis line style={InkGray, line width=0.5pt},
  every axis plot/.append style={line width=0.9pt},
]
\addplot[AgriGreen] table[x=x,y=y]{data/clean/profiles/px_002556.dat};
\addplot[SoftRed, dashed, line width=0.7pt] table[x=x,y=y]{data/clean/profiles/pxzz_002556.dat};
\addplot[only marks, mark=*, mark size=1.1pt, SoftRed] table[x=x,y=y]{data/clean/profiles/pxzz_002556.dat};
\end{axis}
\end{tikzpicture}

\captionof{figure}[辉隆股份（002556）股价与周期骨架]{辉隆股份（002556）月度收盘价（元/股）与 ZigZag 周期骨架（阈值 25\%，圆点为周期转折点）}
\end{minipage}
\end{center}

\pfl{过去周期分析} 以 25\% 的 ZigZag 阈值划分，样本区间内股价形成 2 个主升段与 2 个主跌段。 最大主升段出现在 2024 年 6 月 至 2026 年 2 月，20 个月+61.2\%；最大主跌段为 2021 年 12 月 至 2024 年 6 月，30 个月-67.6\%。 最近一次转折出现在 2026 年 6 月（下跌段自 2026 年 2 月起，4 个月-33.3\%），当前处于该段的延续之中。 公司股价较申万农林牧渔指数提前约 23 个月见底，是板块中较早定价周期反转的公司之一。 农业综合类公司业务分散，估值更多取决于资产与转型预期而非单一品种价格，公司 2026 年 6 月末资产负债率 67.0\%，高于全样本中位数 52.2\%，其股价因此更多反映资金面与信用风险，而不是价格弹性本身。 综合区间形态看，该股单段涨跌幅度偏大、以趋势段为主（最大主升 +61.2\%、最大主跌 -67.6\%），年化波动率 30.3\% 与全样本中位数 37.3\% 相比波动小于样本中位数。

\pfl{盈利情况} 2026 年上半年营业收入 85.47 亿元（同比 +3.3\%），归母净利润 1.13 亿元，同比 +1.5\%。 以年报为趋势对照，2023---2025 年营业收入由 178.33 亿元变为 151.32 亿元（两年年均复合增速 -7.9\%），归母净利润由 0.77 亿元变为 1.94 亿元。 按半年报口径，2026 年上半年的 ROE 为 3.1\%（未年化）、销售净利率 1.4\%、毛利率 6.2\%，ROE 在“—”2 家公司中按数值由高到低排第\zhnumber{1}位。 同期全样本半年 ROE 中位数 0.75\%，公司与之相差 2.33 个百分点（略好于中位数公司）。 综合看，公司在报告期维持盈利（高于全样本半年 ROE 中位数 0.75\%），利润的可持续性取决于本行业景气的延续与成本管控的执行。

\begin{center}\small\setlength{\tabcolsep}{3.4pt}
\captionof{table}[辉隆股份（002556）关键财务指标]{辉隆股份（002556）关键财务指标（合并报表口径；两个半年报期均未年化；现金短债比 $=$ 货币资金 $\div$ 短期债务）}
\begin{adjustbox}{max width=\textwidth}
\begin{tabular}{@{}lrrrrr@{}}
\toprule
指标 & 2023 年 & 2024 年 & 2025 年 & 2025 年半年报 & 2026 年半年报 \\
\midrule
营业收入（亿元） & 178.33 & 156.52 & 151.32 & 82.77 & 85.47 \\
归母净利润（亿元） & 0.77 & 1.69 & 1.94 & 1.11 & 1.13 \\
毛利率（\%） & 4.8 & 5.8 & 6.3 & 5.8 & 6.2 \\
ROE（\%） & 2.0 & 4.6 & 5.2 & 3.0 & 3.1 \\
资产负债率（\%） & 66.2 & 64.6 & 64.9 & 67.9 & 67.0 \\
经营活动现金流净额（亿元） & -3.73 & -1.54 & 10.71 & -1.09 & -4.17 \\
货币资金（亿元） & 11.16 & 9.62 & 11.59 & 11.67 & 11.20 \\
有息负债合计（亿元） & 31.71 & 29.92 & 25.60 & 34.29 & 31.42 \\
现金短债比（倍） & 0.70 & 0.79 & 0.82 & 0.65 & 0.40 \\
\bottomrule
\end{tabular}
\end{adjustbox}
\end{center}

\begin{center}
\begin{minipage}{\textwidth}\centering
\begin{tikzpicture}
\begin{groupplot}[
  group style={group size=2 by 1, horizontal sep=0.60cm},
  width=6.85cm, height=4.60cm,
  tick label style={font=\tiny},
  title style={font=\scriptsize\heiti},
  yticklabel style={font=\tiny},
  ylabel style={font=\tiny},
  grid=major, grid style={LightGray, line width=0.3pt},
  axis line style={InkGray, line width=0.5pt},
  enlarge x limits={abs=0.8},
  legend style={font=\scriptsize, at={(0.5,-0.20)}, anchor=north,
                legend columns=2, draw=none, fill=none, column sep=6pt},
]
\nextgroupplot[
  ybar, bar width=4pt,
  ymin=-19.08, ymax=213.71,
  xtick={0,1,2,3,4,5.7},
  xticklabels={2021,2022,2023,2024,2025,2026H1},
  ylabel={亿元},
]
\addplot[fill=AgriGreen!75, draw=AgriGreen, bar shift=-2.6pt] table[x=i, y=rev]{data/clean/profiles/fin_002556.dat};
\addplot[fill=SoftRed!60, draw=SoftRed, bar shift=2.6pt] table[x=i, y=ni]{data/clean/profiles/fin_002556.dat};
\addplot[fill=AgriGreen!30, draw=AgriGreen, pattern=north east lines, bar shift=-2.6pt] table[x=x, y=rev]{data/clean/profiles/finh_002556.dat};
\addplot[fill=SoftRed!20, draw=SoftRed, pattern=north east lines, bar shift=2.6pt] table[x=x, y=ni]{data/clean/profiles/finh_002556.dat};
\legend{营业收入（年/半年）, 归母净利润（年/半年）}
\nextgroupplot[
  ymin=-5.4, ymax=73.7,
  xtick={0,1,2,3,4,5.7},
  xticklabels={2021,2022,2023,2024,2025,2026H1},
  ylabel={\%},
]
\addplot[AgriGold, mark=*, mark size=1.3pt] table[x=i, y=roe]{data/clean/profiles/fin_002556.dat};
\addplot[OilBlue, mark=square*, mark size=1.3pt] table[x=i, y=debt]{data/clean/profiles/fin_002556.dat};
\addplot[only marks, AgriGold, mark=*, mark size=2.0pt] table[x=x, y=roe]{data/clean/profiles/finh_002556.dat};
\addplot[only marks, OilBlue, mark=square*, mark size=2.0pt] table[x=x, y=debt]{data/clean/profiles/finh_002556.dat};
\legend{ROE, 资产负债率}
\end{groupplot}
\end{tikzpicture}
\begin{tikzpicture}
\begin{axis}[
  name=finbot,
  width=13.75cm, height=3.10cm, scale only axis,
  ymin=-3.17, ymax=35.52,
  xtick={0,1,2,3,4,5.7},
  xticklabels={2021,2022,2023,2024,2025,2026H1},
  tick label style={font=\tiny},
  yticklabel style={font=\tiny},
  ylabel={亿元}, ylabel style={font=\tiny},
  ybar, bar width=4pt,
  grid=major, grid style={LightGray, line width=0.3pt},
  axis line style={InkGray, line width=0.5pt},
  enlarge x limits={abs=0.8},
  legend style={font=\scriptsize, at={(0.5,-0.26)}, anchor=north,
                legend columns=2, draw=none, fill=none, column sep=10pt},
]
\addplot[fill=AgriGreen!55, draw=AgriGreen, bar shift=-3.0pt] table[x=i, y=cash]{data/clean/profiles/fin_002556.dat};
\addplot[fill=SoftRed!45, draw=SoftRed, bar shift=3.0pt] table[x=i, y=ibd]{data/clean/profiles/fin_002556.dat};
\addplot[fill=AgriGreen!25, draw=AgriGreen, pattern=north east lines, bar shift=-3.0pt] table[x=x, y=cash]{data/clean/profiles/finh_002556.dat};
\addplot[fill=SoftRed!20, draw=SoftRed, pattern=north east lines, bar shift=3.0pt] table[x=x, y=ibd]{data/clean/profiles/finh_002556.dat};
\legend{货币资金（左轴）, 有息负债（左轴）}
\end{axis}
\begin{axis}[
  at={(finbot.south east)}, anchor=south east,
  width=13.75cm, height=3.10cm, scale only axis,
  axis y line*=right, axis x line*=none, xtick=\empty,
  ymin=-8.49, ymax=17.41,
  tick label style={font=\tiny},
  yticklabel style={font=\tiny}, scaled y ticks=false,
  legend style={font=\scriptsize, at={(0.5,-0.44)}, anchor=north,
                legend columns=1, draw=none, fill=none},
]
\addplot[OilBlue, mark=*, mark size=2.2pt, line width=1.3pt] table[x=i, y=ocf]{data/clean/profiles/fin_002556.dat};
\addplot[only marks, OilBlue, mark=*, mark size=3.2pt] table[x=x, y=ocf]{data/clean/profiles/finh_002556.dat};
\legend{经营活动现金流净额（右轴，亿元，蓝点）}
\end{axis}
\end{tikzpicture}

\captionof{figure}[辉隆股份（002556）近年主要财务指标]{辉隆股份（002556）近年主要财务指标（左上：营业收入与归母净利润；右上：ROE 与资产负债率；下：货币资金、有息负债为左轴柱状，经营活动现金流净额为其右轴的蓝点折线。
2021---2025 年为年报口径，2026H1 为 2026 年半年报/半年末，与其前一年报之间留出间隔、以斜纹或空心点区分；半年报数据未年化）}
\end{minipage}
\end{center}

\pfl{财务分析} 2026 年 6 月末资产负债率 67.0\%，较 2025 年末上升 2.1 个百分点（样本中位数 52.2\%，所处三级行业内按由低到高排第\zhnumber{1}位）。 2026 年 6 月末货币资金 11.20 亿元、短期债务（短期借款与一年内到期非流动负债）27.68 亿元、长期债务 3.74 亿元，有息负债合计 31.42 亿元，现金短债比 0.40 倍——覆盖偏紧。 净有息负债 20.22 亿元，相当于其 2026 年上半年营业收入的 24\%。 流动比率 1.01、速动比率 0.55，短期指标尚可。 2026 年上半年经营活动现金流净额 -4.17 亿元，净利润为正但经营现金流为负，利润的现金含量偏低。 综合看，账面货币资金对有息负债与短期债务的覆盖不足，滚动偿付对新增融资的依赖度较高；资产负债率高于样本中位数 14.8 个百分点；有息负债中短期债务占比更高，期限结构仍以滚动续作为主。 公司披露的风险集中在诉讼与资产冻结、股东与质押两类：一是 2025 年 12 月，受限货币资金 3.84 亿元；二是 2026 年 1 月，媒体报道控股股东解除 410 万股质押。

\pfl{重大战略转型} 近三年公司的战略动作集中在：股权与股东回报（1 项）：2026 年 4 月，2026-08-13 披露进展：注销回购股份 1,157.72 万股、对应回购金额 1.00 亿元。 公告口径上，近三年（2023 年 9 月---2026 年 9 月）公司披露重大重组与并购类公告 0 条、对外投资与转型类公告 1 条、回购与股权激励类公告 16 条，合计公告 346 条。 主营结构的口径变化已经落到报表上：2021 年年报的第一大行业为“化肥”（62.1\%），2026 年半年报为“农资产品”（73.2\%）；公司在这两期的分部划分方式有过调整。 综合判断，报告期内公司有 1 条与战略相关的公告落地（重大重组与并购 0 条、对外投资与转型 1 条），资本运作与主业调整之间存在直接关联。

\begin{center}\small\setlength{\tabcolsep}{4pt}
\captionof{table}[辉隆股份（002556）历史融资与资本运作]{辉隆股份（002556）历史融资与资本运作（据公司公告与公开报道整理；金额为披露值，含首次公开发行、再融资、债券、银行借款与重整投资等）}
\begin{adjustbox}{max width=\textwidth}
\begin{tabular}{@{}p{1.45cm}p{2.55cm}p{4.05cm}p{4.95cm}@{}}
\toprule
时间 & 方式 & 规模 & 用途与说明 \\
\midrule
2011-02 & IPO（深圳证券交易所中小板） & 发行 3,\allowbreak{}750 万股，\allowbreak{}发行价 37.50 元/\allowbreak{}股 & 配送中心建设项目和信息化系统建设项目；深证上[2011]67 号文同意上市限责任公司 \\
2018-2019 /\allowbreak{} 2020-02 & 发行股份、可转换公司债券及支付现金购买资产（关联交易） & 标的为海华科技 100\% 股权 & 收购海华科技 100\% 股权（交易对方为辉隆投资、\allowbreak{}蚌埠隆海及解凤贤、\allowbreak{}解凤苗、\allowbreak{}解凤祥等 40 名自然人）；可转债初始转股价 5.13 元/\allowbreak{}股 \\
2020-04 & 非公开发行股份及可转换公司债券募集配套资金 & 募集配套资金 6.44 亿元 & 向交易对方支付现金对价、\allowbreak{}海华科技项目建设及支付本次交易相关费用；配套可转债于 2020-04-21 登记完成 \\
2023-07-28 & 中期票据（乡村振兴） & 4.00 亿元，\allowbreak{}票面利率 3.50\% & 2025 年度报告披露期末余额 4.06 亿元 \\
2024-03-22 & 中期票据（乡村振兴） & 4.00 亿元，\allowbreak{}票面利率 2.80\% & 2025 年度报告披露期末余额 4.09 亿元 \\
2026-04-17 & 终止部分超募资金投资项目 & 剩余超募资金 4,\allowbreak{}343.59 万元 & 终止首发超募资金投资项目“40 万吨/\allowbreak{}年新型肥料项目”，\allowbreak{}剩余超募资金继续存放原募集资金专户；截至 2026-06-30 公司对募集资金项目累计投入 13.06 亿元 \\
\bottomrule
\end{tabular}
\end{adjustbox}
\end{center}

\pfl{历史融投资情况} 从现金流量表看，2025 年公司取得借款收到的现金 18.70 亿元、偿还债务支付的现金 25.51 亿元、吸收权益性投资 0.01 亿元，筹资活动现金流净额 -10.40 亿元（2023 年为取得借款 36.24 亿元、偿还债务 28.20 亿元）。 2026 年上半年筹资活动现金流净额为 0.00 亿元（取得借款 14.36 亿元、偿还债务 13.04 亿元，未年化），仍处于净融入状态。 2025 年分配股利、利润或偿付利息支付的现金合计 3.18 亿元。 公开披露的融资记录共 6 笔，工具类型以 IPO、可转债、中期票据为主；金额与用途见下表。 其中规模最大的两笔为：2026 年 4 月，终止部分超募资金投资，剩余超募资金 4,343.59 万元（该事项后续已终止或解除，不构成已到位资金）；2020 年 4 月，可转债，募集配套资金 6.44 亿元。 股本方面，近三年披露股本变动 9 次（其他 3 次、限售股份上市 2 次、股份回购 1 次）；总股本由 2021 年末的 9.54 亿股变为 9.47 亿股（-0.78\%）。 分红方面，近三年现金分红 5 次、合计每股派现 0.85 元，最近一次实施在 2025 年 12 月。 综合看，公司融资结构上以债务滚动为主、股权融资为补充；2025 年筹资活动现金流净额 -10.40 亿元，融资节奏仍以维持存量债务周转为主，与前述经营与投资现金流的方向基本一致。

\pfl{未来潜在融资需求分析} 2026 年 6 月末货币资金 11.20 亿元，而短期债务 27.68 亿元（现金短债比 0.40 倍），2026 年上半年经营现金流净额 -4.17 亿元。按“短期债务减去账面货币资金”的滚动偿付口径，静态缺口约 16.47 亿元；上半年盈利或接近盈亏平衡，该缺口不随经营亏损进一步扩大。 公司 2026 年 6 月末资产负债率 67.0\%、2026 年上半年 ROE 3.1\%（未年化），资产负债率偏高（67.0\%），融资需求首先用于置换与滚动存量债务、补充营运资金，属于“补血型”而非扩张型。 公开披露的后续资金安排线索包括：股权融资线索：2026 年 5 月，公司披露变更回购股份用途并注销、减少注册资本及修订公司章程；债务与授信安排：2026 年 6 月，公司对控股公司提供担保及募集资金存放与实际使用情况专项报告。 综合看，公司的外部资金需求以补充营运资金为主、项目投入为辅，滚动偿付口径下的静态缺口约 16.47 亿元，融资节奏将受制于银行续贷意愿与股权融资窗口。




\subsubsection{大禹节水（300021）}
\label{co:300021}

\pfl{公司基本情况} 大禹节水成立于 2009 年，国内节水灌溉与农田水利领域主要上市公司，业务覆盖节水材料装备制造、灌区 EPC/EPCO 总承包与农水设计服务，正向“设计引领+全产业链”模式转型。 公司于 2009-10-30 在深交所创业板首发上市，注册资本 10.22 亿元，总股本 8.54 亿股。 截至 2026 年 10 月 12 日收盘价 3.83 元，对应总市值 32.7 亿元，近一年-22.6\%，总市值在三级行业“—”的 2 家公司中按由大到小排第\zhnumber{2}位，近一年涨跌幅低于全样本中位数（-10.4\%）12.3 个百分点。 按 2025 年年报披露的行业构成，收入几乎全部来自智慧水利行业 100.0\%。 与 2021 年年报相比，第一大行业口径由“节水灌溉行业”（81.0\%）变为“智慧水利行业”（100.0\%）；两期披露的分类口径有调整，占比差异中同时包含分类因素。

\pfl{股权结构与控制权} 按 2026 年 6 月 30 日的股东名册，前三大股东为王浩宇 持股 18.17\%；仇玲 持股 17.74\%；高雅萍（持股比例未披露），股东户数 38,882 户。 从控制权角度看，2026 年 6 月末持股 18.17\%（1.86 亿股），其中有限售条件股份 1.39 亿股、质押 4,500.00 万股；第二大股东仇玲持股 17.74\%（1.81 亿股），与王浩宇为母子关系、二人为一致行动人；高雅萍、蒋仕波构成一致行动关系，合计持有公司 2.71\% 股份。

\pfl{经营要点} 业务与产能层面，上半年累计实现新签合同额 21.66 亿元、完成回款 18.70 亿元；获新疆阿克苏塔里木河干流盖孜库木片区生态综合治理等标志性订单。

\pfl{所处行业的基本面} 农业综合类公司的收入分布在多个品种上，供需由各自品种决定，整体上更接近农资与服务的需求驱动（第 \ref{sec:framework} 章的供需框架）。 行业层面的最新情况是：农业综合生产能力明显提高。

\pfl{竞争优势} 把该公司放回同类公司比较，2026 年上半年营业收入 15.9 亿元、在三级行业“农业综合Ⅱ”的 2 家公司中排第\zhnumber{2}位；毛利率 20.6\%、高于行业内中位数 7.2 个百分点；半年 ROE -0.0\%（未年化）、低于行业内中位数（1.52\%）；资产负债率 71.7\%、高于行业内中位数（69.3\%）。 优势在报表上主要体现为高于行业平均的毛利率，这在周期底部是抵御亏损的主要缓冲。

\pfl{历史股价} 本报告样本区间（2021 年 1 月 至 2026 年 10 月）内，股价由 4.76 元变为 3.83 元，累计 -19.5\%，区间最高 6.32 元（2025 年 7 月）、最低 3.25 元（2024 年 8 月）；当前价相当于区间最高价的 60.6\%，为区间最低价的 1.18 倍。 月度收益的年化波动率 27.0\%，低于全样本中位数 37.3\%，在三级行业“—”的 2 家公司中波动率由高到低排第\zhnumber{2}位。 把股价阶段与公司事件对齐看，2024 年 8 月 至 2025 年 7 月股价上涨+94.5\%，同期（2025 年 6 月）可转债募集资金：本报告期投入 8,418.05 万元；2021 年 12 月 至 2024 年 8 月股价下跌-44.9\%，同期（2024 年 2 月）股东大会通过变更部分募集资金用途，将可转债募投项目尚未使用募集资金及专户存款利息用于新项目。 整体上看，区间的几处大级别拐点都能在公司的资本运作、股权变动或风险事件上找到对应，股价波动有明显的公司层面驱动，而非单纯的行业 beta。

\begin{center}
\begin{minipage}{\textwidth}\centering
\begin{tikzpicture}
\begin{axis}[
  width=12.6cm, height=3.9cm, scale only axis,
  xmin=2020.922, xmax=2026.972, ymin=3.02, ymax=6.76,
  xtick={2021,2022,2023,2024,2025,2026},
  yearticks,
  yticklabel style={font=\scriptsize},
  ylabel={元/股}, ylabel style={font=\scriptsize},
  grid=major, grid style={LightGray, line width=0.3pt},
  axis line style={InkGray, line width=0.5pt},
  every axis plot/.append style={line width=0.9pt},
]
\addplot[AgriGreen] table[x=x,y=y]{data/clean/profiles/px_300021.dat};
\addplot[SoftRed, dashed, line width=0.7pt] table[x=x,y=y]{data/clean/profiles/pxzz_300021.dat};
\addplot[only marks, mark=*, mark size=1.1pt, SoftRed] table[x=x,y=y]{data/clean/profiles/pxzz_300021.dat};
\end{axis}
\end{tikzpicture}

\captionof{figure}[大禹节水（300021）股价与周期骨架]{大禹节水（300021）月度收盘价（元/股）与 ZigZag 周期骨架（阈值 25\%，圆点为周期转折点）}
\end{minipage}
\end{center}

\pfl{过去周期分析} 以 25\% 的 ZigZag 阈值划分，样本区间内股价形成 2 个主升段与 2 个主跌段。 最大主升段出现在 2024 年 8 月 至 2025 年 7 月，11 个月+94.5\%；最大主跌段为 2021 年 12 月 至 2024 年 8 月，32 个月-44.9\%。 最近一次转折出现在 2026 年 10 月（下跌段自 2025 年 7 月起，15 个月-39.4\%），当前处于该段的延续之中。 公司股价较申万农林牧渔指数提前约 21 个月见底，是板块中较早定价周期反转的公司之一。 农业综合类公司业务分散，估值更多取决于资产与转型预期而非单一品种价格，公司 2026 年 6 月末资产负债率 71.7\%，高于全样本中位数 52.2\%，其股价因此更多反映资金面与信用风险，而不是价格弹性本身。 综合区间形态看，该股单段涨跌幅度偏大、以趋势段为主（最大主升 +94.5\%、最大主跌 -44.9\%），年化波动率 27.0\% 与全样本中位数 37.3\% 相比波动小于样本中位数。

\pfl{盈利情况} 2026 年上半年营业收入 15.86 亿元（同比 +24.3\%），归母净利润 -0.01 亿元，由上年同期的盈利转为亏损。 以年报为趋势对照，2023---2025 年营业收入由 34.53 亿元变为 37.61 亿元（两年年均复合增速 +4.4\%），归母净利润由 0.50 亿元变为 0.49 亿元。 按半年报口径，2026 年上半年的 ROE 为 -0.0\%（未年化）、销售净利率 -0.4\%、毛利率 20.6\%，ROE 在“—”2 家公司中按数值由高到低排第\zhnumber{2}位。 同期全样本半年 ROE 中位数 0.75\%，公司与之相差 0.79 个百分点（弱于中位数公司）。 从公司披露的原因看，业绩变动主要来自：2026 年上半年，归属于上市公司股东的净利润转亏 -110.79 万元。 综合看，公司仍处亏损区间、由盈转亏，半年 ROE 在三级行业内按由高到低排第\zhnumber{2}位，单位盈利能力的修复依赖行业景气与自身成本端的改善，报表弹性主要体现为价格与销量的方向性变化。

\begin{center}\small\setlength{\tabcolsep}{3.4pt}
\captionof{table}[大禹节水（300021）关键财务指标]{大禹节水（300021）关键财务指标（合并报表口径；两个半年报期均未年化；现金短债比 $=$ 货币资金 $\div$ 短期债务）}
\begin{adjustbox}{max width=\textwidth}
\begin{tabular}{@{}lrrrrr@{}}
\toprule
指标 & 2023 年 & 2024 年 & 2025 年 & 2025 年半年报 & 2026 年半年报 \\
\midrule
营业收入（亿元） & 34.53 & 43.79 & 37.61 & 12.76 & 15.86 \\
归母净利润（亿元） & 0.50 & 0.81 & 0.49 & 0.13 & -0.01 \\
毛利率（\%） & 24.3 & 17.8 & 21.4 & 23.3 & 20.6 \\
ROE（\%） & 2.5 & 3.9 & 2.2 & 0.6 & -0.0 \\
资产负债率（\%） & 71.2 & 75.1 & 70.6 & 73.6 & 71.7 \\
经营活动现金流净额（亿元） & -0.94 & 5.97 & -2.76 & -5.49 & -5.69 \\
货币资金（亿元） & 11.51 & 16.88 & 14.69 & 12.02 & 12.85 \\
有息负债合计（亿元） & 25.14 & 27.10 & 25.96 & 29.00 & 29.12 \\
现金短债比（倍） & 1.14 & 1.22 & 0.98 & 0.72 & 0.70 \\
\bottomrule
\end{tabular}
\end{adjustbox}
\end{center}

\begin{center}
\begin{minipage}{\textwidth}\centering
\begin{tikzpicture}
\begin{groupplot}[
  group style={group size=2 by 1, horizontal sep=0.60cm},
  width=6.85cm, height=4.60cm,
  tick label style={font=\tiny},
  title style={font=\scriptsize\heiti},
  yticklabel style={font=\tiny},
  ylabel style={font=\tiny},
  grid=major, grid style={LightGray, line width=0.3pt},
  axis line style={InkGray, line width=0.5pt},
  enlarge x limits={abs=0.8},
  legend style={font=\scriptsize, at={(0.5,-0.20)}, anchor=north,
                legend columns=2, draw=none, fill=none, column sep=6pt},
]
\nextgroupplot[
  ybar, bar width=4pt,
  ymin=-4.39, ymax=49.04,
  xtick={0,1,2,3,4,5.7},
  xticklabels={2021,2022,2023,2024,2025,2026H1},
  ylabel={亿元},
]
\addplot[fill=AgriGreen!75, draw=AgriGreen, bar shift=-2.6pt] table[x=i, y=rev]{data/clean/profiles/fin_300021.dat};
\addplot[fill=SoftRed!60, draw=SoftRed, bar shift=2.6pt] table[x=i, y=ni]{data/clean/profiles/fin_300021.dat};
\addplot[fill=AgriGreen!30, draw=AgriGreen, pattern=north east lines, bar shift=-2.6pt] table[x=x, y=rev]{data/clean/profiles/finh_300021.dat};
\addplot[fill=SoftRed!20, draw=SoftRed, pattern=north east lines, bar shift=2.6pt] table[x=x, y=ni]{data/clean/profiles/finh_300021.dat};
\legend{营业收入（年/半年）, 归母净利润（年/半年）}
\nextgroupplot[
  ymin=-6.0, ymax=82.6,
  xtick={0,1,2,3,4,5.7},
  xticklabels={2021,2022,2023,2024,2025,2026H1},
  ylabel={\%},
]
\addplot[AgriGold, mark=*, mark size=1.3pt] table[x=i, y=roe]{data/clean/profiles/fin_300021.dat};
\addplot[OilBlue, mark=square*, mark size=1.3pt] table[x=i, y=debt]{data/clean/profiles/fin_300021.dat};
\addplot[only marks, AgriGold, mark=*, mark size=2.0pt] table[x=x, y=roe]{data/clean/profiles/finh_300021.dat};
\addplot[only marks, OilBlue, mark=square*, mark size=2.0pt] table[x=x, y=debt]{data/clean/profiles/finh_300021.dat};
\legend{ROE, 资产负债率}
\end{groupplot}
\end{tikzpicture}
\begin{tikzpicture}
\begin{axis}[
  name=finbot,
  width=13.75cm, height=3.10cm, scale only axis,
  ymin=-2.91, ymax=32.62,
  xtick={0,1,2,3,4,5.7},
  xticklabels={2021,2022,2023,2024,2025,2026H1},
  tick label style={font=\tiny},
  yticklabel style={font=\tiny},
  ylabel={亿元}, ylabel style={font=\tiny},
  ybar, bar width=4pt,
  grid=major, grid style={LightGray, line width=0.3pt},
  axis line style={InkGray, line width=0.5pt},
  enlarge x limits={abs=0.8},
  legend style={font=\scriptsize, at={(0.5,-0.26)}, anchor=north,
                legend columns=2, draw=none, fill=none, column sep=10pt},
]
\addplot[fill=AgriGreen!55, draw=AgriGreen, bar shift=-3.0pt] table[x=i, y=cash]{data/clean/profiles/fin_300021.dat};
\addplot[fill=SoftRed!45, draw=SoftRed, bar shift=3.0pt] table[x=i, y=ibd]{data/clean/profiles/fin_300021.dat};
\addplot[fill=AgriGreen!25, draw=AgriGreen, pattern=north east lines, bar shift=-3.0pt] table[x=x, y=cash]{data/clean/profiles/finh_300021.dat};
\addplot[fill=SoftRed!20, draw=SoftRed, pattern=north east lines, bar shift=3.0pt] table[x=x, y=ibd]{data/clean/profiles/finh_300021.dat};
\legend{货币资金（左轴）, 有息负债（左轴）}
\end{axis}
\begin{axis}[
  at={(finbot.south east)}, anchor=south east,
  width=13.75cm, height=3.10cm, scale only axis,
  axis y line*=right, axis x line*=none, xtick=\empty,
  ymin=-8.72, ymax=9.46,
  tick label style={font=\tiny},
  yticklabel style={font=\tiny}, scaled y ticks=false,
  legend style={font=\scriptsize, at={(0.5,-0.44)}, anchor=north,
                legend columns=1, draw=none, fill=none},
]
\addplot[OilBlue, mark=*, mark size=2.2pt, line width=1.3pt] table[x=i, y=ocf]{data/clean/profiles/fin_300021.dat};
\addplot[only marks, OilBlue, mark=*, mark size=3.2pt] table[x=x, y=ocf]{data/clean/profiles/finh_300021.dat};
\legend{经营活动现金流净额（右轴，亿元，蓝点）}
\end{axis}
\end{tikzpicture}

\captionof{figure}[大禹节水（300021）近年主要财务指标]{大禹节水（300021）近年主要财务指标（左上：营业收入与归母净利润；右上：ROE 与资产负债率；下：货币资金、有息负债为左轴柱状，经营活动现金流净额为其右轴的蓝点折线。
2021---2025 年为年报口径，2026H1 为 2026 年半年报/半年末，与其前一年报之间留出间隔、以斜纹或空心点区分；半年报数据未年化）}
\end{minipage}
\end{center}

\pfl{财务分析} 2026 年 6 月末资产负债率 71.7\%，较 2025 年末上升 1.1 个百分点（样本中位数 52.2\%，所处三级行业内按由低到高排第\zhnumber{2}位）。 2026 年 6 月末货币资金 12.85 亿元、短期债务（短期借款与一年内到期非流动负债）18.48 亿元、长期债务 10.64 亿元，有息负债合计 29.12 亿元，现金短债比 0.70 倍——覆盖偏紧。 净有息负债 16.27 亿元，相当于其 2026 年上半年营业收入的 103\%。 流动比率 1.25、速动比率 1.20，短期流动性无明显缺口。 2026 年上半年经营活动现金流净额 -5.69 亿元，经营现金流与利润同时为负，日常运营对债务与股东投入的依赖度较高。 2026 年 6 月末在建工程 0.61 亿元，较上年末增加 0.54 亿元，仍有在建投入。 2026 年 6 月末商誉 1.34 亿元，占归母净资产的 4.6\%，是净资产中最脆弱的一块。 综合看，现金短债比 0.70 倍，短期资金安排偏紧；资产负债率高于样本中位数 19.5 个百分点；有息负债中短期债务占比更高，期限结构仍以滚动续作为主。

\pfl{重大战略转型} 近三年公司的战略动作集中在：并购与重组（1 项）：2026 年 3 月，公司转让所持慧图科技 28.5\% 股权给嘉兴乾瞻衡远创业投资合伙企业（有限合伙）。 公告口径上，近三年（2023 年 9 月---2026 年 9 月）公司披露重大重组与并购类公告 9 条、对外投资与转型类公告 35 条、回购与股权激励类公告 42 条，合计公告 635 条。 主营结构的口径变化已经落到报表上：2021 年年报的第一大行业为“节水灌溉行业”（81.0\%），2025 年年报为“智慧水利行业”（100.0\%）；公司在这两期的分部划分方式有过调整。 综合判断，报告期内公司有 44 条与战略相关的公告落地（重大重组与并购 9 条、对外投资与转型 35 条），资本运作与主业调整之间存在直接关联。

\begin{center}\small\setlength{\tabcolsep}{4pt}
\captionof{table}[大禹节水（300021）历史融资与资本运作]{大禹节水（300021）历史融资与资本运作（据公司公告与公开报道整理；金额为披露值，含首次公开发行、再融资、债券、银行借款与重整投资等）}
\begin{adjustbox}{max width=\textwidth}
\begin{tabular}{@{}p{1.45cm}p{2.55cm}p{4.05cm}p{4.95cm}@{}}
\toprule
时间 & 方式 & 规模 & 用途与说明 \\
\midrule
2009-10 & IPO（深圳证券交易所创业板） & 发行 1,\allowbreak{}800 万股，\allowbreak{}发行价 14.00 元/\allowbreak{}股 & 在酒泉总部及全资子公司新疆公司、\allowbreak{}武威公司建设滴灌管（带）及配套节水器材生产线技改扩建项目；超过 1.53 亿元的部分补充流动资金 \\
2020-07-28 & 公开发行可转换公司债券“大禹转债” & 638 万张、\allowbreak{}每张面值 100.00 元，\allowbreak{}发行总额 6.38 亿元 & 高端节水灌溉产品智能工厂建设项目（拟配备以色列先进压力补偿式滴灌生产线及核心智能制造设备）、\allowbreak{}现代农业运营服务和科技研发中心建设项目等；；转股期自 2021-02-03 起至 2026-07-27 止；2024 年 2 月经股东大会审议变更部分募集资金用途 \\
2022-02-24 & 以简易程序向特定对象发行股票 & 5,\allowbreak{}859.38 万股，\allowbreak{}发行价 5.12 元/\allowbreak{}股，\allowbreak{}募集资金总额 3.00 亿元 & 立信会计师事务所出具验资报告；新增股份自 2022-03-08 上市 \\
2025-08-21 & 可转债提前赎回 & 赎回价格 100.20 元/\allowbreak{}张（含税）；公司本年度实际赎回 10,\allowbreak{}724.00 张大禹转债 & 终止“大禹转债”存续并在深交所摘牌；本年度因可转债转股增加股本 17,\allowbreak{}234.74 万股 \\
2026-07 & 2025 年度利润分配 & 以 10.22 亿股为基数，\allowbreak{}每 10 股派发现金红利 0.20 元（含税） & 股东回报；2026 年 7 月完成派发 \\
\bottomrule
\end{tabular}
\end{adjustbox}
\end{center}

\pfl{历史融投资情况} 从现金流量表看，2025 年公司取得借款收到的现金 21.76 亿元、偿还债务支付的现金 16.85 亿元、吸收权益性投资 0.08 亿元，筹资活动现金流净额 2.39 亿元（2023 年为取得借款 14.48 亿元、偿还债务 9.61 亿元）。 2026 年上半年筹资活动现金流净额为 2.82 亿元（取得借款 10.47 亿元、偿还债务 7.10 亿元，未年化），仍处于净融入状态。 2025 年分配股利、利润或偿付利息支付的现金合计 1.46 亿元。 公开披露的融资记录共 5 笔，工具类型以 IPO、可转债、定向增发为主；金额与用途见下表。 其中规模最大的两笔为：2020 年 7 月，可转债，638 万张、每张面值 100.00 元，发行总额 6.38 亿元；2009 年 10 月，IPO，发行 1,800 万股，发行价 14.00 元/股。 股本方面，近三年披露股本变动 17 次（可转债转股 11 次、股份回购 4 次、其他 1 次）；总股本由 2021 年末的 8.01 亿股变为 8.54 亿股（+6.59\%）。 分红方面，近三年现金分红 4 次、合计每股派现 0.21 元，最近一次实施在 2025 年 12 月。 近三年“再融资”类公告 18 条，涉及定向增发、可转债。 综合看，公司融资结构上以债务滚动为主、股权融资为补充；2025 年筹资活动现金流净额 2.39 亿元，年度内仍为净融入，与前述经营与投资现金流的方向基本一致。

\pfl{未来潜在融资需求分析} 2026 年 6 月末货币资金 12.85 亿元，而短期债务 18.48 亿元（现金短债比 0.70 倍），2026 年上半年经营现金流净额 -5.69 亿元。按“短期债务减去账面货币资金”的滚动偿付口径，静态缺口约 5.62 亿元；上半年盈利或接近盈亏平衡，该缺口不随经营亏损进一步扩大。 按第 \ref{sec:financing-need} 节的三档划分，公司属于第一档“补血型”：2026 年上半年归母净利润 -0.01 亿元、6 月末资产负债率 71.7\%。 可行的融资路径按现实性排序为：定向增发（需股价配合，且控股股东需放弃优先认购或同步增资）、引入国资或产业投资人（部分公司以控制权让渡为对价）、出售非核心资产与参股股权。 公开披露的后续资金安排线索包括：股权融资线索：2025 年 6 月，可转债募集资金：本报告期投入 8,418.05 万元；资本开支安排：2025 年 8 月，公司提出高标准农田建设由整改前单体项目拆分招标转为整县整装推进，将显著提升单个订单体量；2025 年 12 月，简易程序发行募集资金：本报告期投入 4,375.47 万元。 资金需求还要叠加在建工程：期末在建工程 0.61 亿元、2026 年上半年购建固定资产等支出 0.45 亿元（未年化），这部分支出在周期底部通常会被主动放缓。 综合看，公司的外部资金需求以债务接续为主，滚动偿付口径下的静态缺口约 5.62 亿元，能否落地取决于授信续作与发行窗口的配合。






\section{战略转型、融资历史与潜在融资需求}
\label{sec:financing}

本章回答第二部分最后、也是最实际的一个问题：\hl{在周期底部，
这 \num{104} 家公司各自需要什么，又能从哪里得到它}？
分析方法是从两类可观测的行为数据出发——\emph{公告}（公司自己披露的意图）
与\emph{财务结构}（公司实际承受的约束），
再把两者交叉，得到融资需求的三个层次。

\subsection{公告中的行为信号}
\label{sec:financing-signals}

近三年（2023 年 9 月---2026 年 10 月）\num{104} 家公司累计披露公告 \num{41266} 条。
按标题关键词分类后得到的事件分布如下：

\begin{table}[htbp]
\centering
\footnotesize
\caption{农业上市公司近三年公告的事件分类统计（按标题关键词归类，共 \num{41266} 条；同一公告可命中多个类别，各行数值不可相加）}
\label{tab:fin-events}
\begin{tabular}{@{}p{2.9cm}rp{7.3cm}@{}}
\toprule
事件类别 & 公告数 & 典型标题特征与解读 \\
\midrule
担保与关联交易 & \num{3211} & “为子公司提供担保”“关联交易”——集团内输血，隐含资金压力 \\
回购与激励 & \num{1872} & “回购股份”“股权激励”“员工持股”——管理层对估值与留人的判断 \\
再融资 & \num{1383} & “向特定对象发行”“募集说明书”“募集资金存放与使用”——融资动作 \\
股东行为 & \num{1364} & “减持”“增持”“质押”——股东层面的流动性需求 \\
业绩与风险 & \num{1264} & “业绩预告”“风险提示”“立案/问询”——经营与合规风险披露 \\
重大重组与并购 & \num{369} & “重大资产重组”“收购”“股权转让”——外延扩张或保壳 \\
对外投资与转型 & \num{274} & “设立子公司”“战略合作”“框架协议”“扩产”——产能与赛道调整 \\
\bottomrule
\end{tabular}
\end{table}

\hl{“担保与关联交易”居首（\num{3211} 条）具有明确的经济含义}：
在农业集团企业内部，母公司为亏损子公司提供担保是维持其
银行授信与生产经营的常规手段。担保公告越密集，通常意味着\emph{集团内部
对单一板块（多为生猪养殖）的资金依赖度越高}。
“再融资”类公告 \num{1383} 条则直接对应资本市场融资动作，
其中 \textbf{新希望}（\num{96} 条）、\textbf{国投中鲁}（\num{85} 条）、
\textbf{*ST广糖}（\num{78} 条）、\textbf{巨星农牧}（\num{75} 条）居前四位，
\hl{这四家公司在 2026 年上半年均处于亏损或微利状态}
（归母净利润分别为 \num{-17.02}、\num{0.02}、\num{-0.16} 与 \num{-8.38} 亿元）。

按三级行业看，公告密度最高的是\textbf{生猪养殖}（公司均再融资公告 \num{28.3} 条、
中位数 \num{27} 条，居全部三级行业首位），其次是\textbf{果蔬加工}（\num{27.6} 条）；
若看重组与并购类公告，密度最高的是\textbf{果蔬加工}（\num{10.6} 条）与\textbf{林业}（\num{9.0} 条）。
前者的高密度对应全行业性的补血需求，后两者则集中在\hl{重组与壳资源运作}上——
国投中鲁、\textbf{*ST景谷}、中基健康、永安林业均属此类，
从净资产与公告结构看，其共同特征是主业萎缩、壳价值可能高于经营价值。

\subsection{融资历史：股本与分红的两个事实}
\label{sec:financing-history}

\begin{itemize}[leftmargin=1.4em, itemsep=2pt]
  \item \textbf{股本变动是普遍事实}：\num{104} 家公司\hl{全部在近三年存在股本变更记录}
        （含送转、股权激励行权、定增、可转债转股）——\emph{没有任何一家的股本在三年内保持不变}。
  \item \textbf{分红能力分化}：\num{73} 家公司在近三年实施过现金分红，
        \hl{\num{31} 家连续三年未分红}（占 \num{29.8}\%），
        其中多数同时处于亏损状态。
\end{itemize}

这两个事实合起来说明一个规律：\hl{农业上市公司的股本变动频率
远高于其分红频率}。在周期底部，公司倾向于\emph{用股权而不是现金流来解决问题}——
定增、可转债、股权激励行权都会增加股本，而现金分红是唯一
真正返还股东的渠道，因此在下行周期中最先被削减。
\emph{这构成了“未来潜在融资需求”分析的基本背景}：
融资的动机在周期底部是刚性（补血），而不是可选（扩张）。

\subsection{重大战略转型的三条路径}
\label{sec:financing-strategy}

本轮周期中，\hl{破产重整与重大债务重组不再是边缘事件}：逐家核对公开资料后，全样本 \num{104} 家
中有 \num{8} 家披露过自身的破产重整、预重整或重大债务重组事项——天邦食品、正邦科技、
傲农生物、华英农业、东方海洋、中基健康、*ST景谷（三度被实施退市风险警示，靠“出售资产＋
债务重组”保壳），以及 2005 年为化解德隆危机由中粮集团受让股权并完成债务重组的中粮糖业。
其中正邦科技（2023 年 12 月）、华英农业（2022 年 4 月）、东方海洋（2023 年 12 月）、
傲农生物（2024 年 12 月）的重整计划已由法院裁定确认执行完毕，正邦科技、华英农业、
东方海洋并已撤销退市风险警示；天邦食品仍在预重整程序中，与投资人的谈判自 2025 年以来
已四次延期。这与第一部分的价格周期判断相互印证：\emph{价格见底并不等于报表见底，
报表见底也不等于资产负债表修复完毕}。

从公告与财务数据的组合中，可以归纳出本轮周期中农业上市公司转型的三条路径。

\paragraph{路径一：一体化收缩（“把亏损的业务砍掉”）。}
以\textbf{新希望}为代表：公司以饲料起家、在 2020---2021 年猪价高点
举债扩张生猪产能，2023---2025 年连续三年将生猪板块作为亏损主体，
2026 年上半年仍亏损 \num{17.02} 亿元、2026 年 6 月末资产负债率 \num{70.7}\%、
再融资相关公告 \num{95} 条（全样本第一）。
其战略实质是\hl{“以融资换时间”}：通过定增、出售资产、引入战投
维持资产负债表，等待猪价回升。
同类公司包括\textbf{唐人神}（\num{-10.46} 亿元）、
\textbf{大北农}（\num{-6.77} 亿元）、\textbf{天邦食品}（\num{-15.43} 亿元）。

\paragraph{路径二：跨界与保壳（“换一个行业重新开始”）。}
以\textbf{国投中鲁}（\num{85} 条再融资公告 + \num{32} 条重组公告，全样本最高重组密度）、
\textbf{绿康生化}（\num{23} 条重组公告，从动保转向微生物发酵与原料药）、
\textbf{*ST景谷}（\num{23} 条重组公告）为代表。
这类公司的特征是：\hl{主业规模小、资产轻、市值仍有一定支撑}，
转型方向多为与农业无关的新兴产业。
\emph{它们的融资需求是“保壳型”的}——不是为了扩产，而是为了维持上市地位。

\paragraph{路径三：向消费端与高附加值环节延伸（“把周期变成品牌”）。}
以\textbf{海大集团}（服务型饲料，半年 ROE \num{6.5}\%、未年化）、
\textbf{乖宝宠物}（半年 ROE \num{4.1}\%、半年营收同比 \num{+10.0}\%）、
\textbf{中宠股份}（半年 ROE \num{4.0}\%、半年营收同比 \num{+34.9}\%）、
\textbf{瑞普生物}（动保平台化，半年 ROE \num{4.3}\%）
为代表。它们的共同点是\hl{把收入结构从“价格 $\times$ 数量”
转向“品牌 $\times$ 复购”}，从而在农产品价格下行时保持利润率。
这类公司是本轮周期中唯一\emph{以扩张为目的}融资的群体——
乖宝宠物（2023 年 8 月）、索宝蛋白（2023 年 12 月）
是全部 \num{104} 家公司中上市最晚的两家。

\subsection{未来潜在融资需求的分层}
\label{sec:financing-need}

把财务约束与转型路径交叉，可以把潜在融资需求分为三档。

\begin{figure}[htbp]
\centering
\begin{tikzpicture}
\begin{groupplot}[
  group style={group size=2 by 1, horizontal sep=1.6cm},
  agri axis, width=5.5cm, height=5.4cm,
]
\nextgroupplot[
  xlabel={2026 年 6 月末资产负债率（\%）}, ylabel={2026 年上半年 ROE（\%）},
  xlabel style={font=\scriptsize}, ylabel style={font=\scriptsize},
  tick label style={font=\scriptsize},
  % 允许超出坐标范围：天邦食品的半年 ROE（-96.8%）已在图注中说明，
  % 若纳入坐标轴会把其余公司的主体分布压缩到不足图高的十分之一
  xmin=0, xmax=105, ymin=-70, ymax=35,
  grid=both, grid style={InkGray!18},
  title style={font=\small}, title={负债率与回报的负相关},
]
\addplot[only marks, mark=*, mark size=1.9pt, color=OilBlue, opacity=0.85]
  table[x=debt, y=roe]{data/clean/comp_debt_roe.dat};
\draw[SoftRed, dashed, line width=0.6pt] (axis cs:65,-70) -- (axis cs:65,35);
\node[font=\tiny, SoftRed, rotate=90, anchor=south] at (axis cs:63,-20) {65\% 警戒线};
\nextgroupplot[
  xbar, xmin=0, xmax=105, ymin=-0.7, ymax=13.7,
  ytick={0,1,2,3,4,5,6,7,8,9,10,11,12,13},
  yticklabels from table={data/clean/comp_risk.dat}{label},
  yticklabel style={font=\tiny, align=right, text width=1.25cm},
  xlabel={2026 年 6 月末资产负债率（\%）}, xlabel style={font=\scriptsize},
  tick label style={font=\scriptsize},
  bar width=5.5pt,
  title style={font=\small}, title={亏损且负债率$>$65\%的公司},
]
\addplot[fill=SoftRed!70, draw=SoftRed] table[y=i, x=debt]{data/clean/comp_risk.dat};
\end{groupplot}
\end{tikzpicture}
\caption{农业上市公司的财务约束：资产负债率与 ROE 的关系（左，2026 年半年报口径）与高负债亏损公司（右；左图 \num{103} 家可算 ROE，其中 \num{1} 家（天邦食品，\num{-96.8}\%）超出纵轴范围未显示）}
\label{fig:fin-risk}
\end{figure}

左图给出一条清晰的经验规律：\hl{资产负债率超过 \num{65}\% 后，
ROE 几乎全部落在负区间}；\num{103} 家可算 ROE 的公司中，
半年 ROE 超过 \num{5}\%（约相当于年化 \num{10}\%）的有 \num{16} 家，
其中 \num{10} 家的资产负债率低于 \num{65}\%——\emph{高负债与高回报在样本中几乎互斥}。
右图则列出同时满足“2026 年上半年亏损 + 6 月末负债率 $>$ \num{65}\%”的 \num{14} 家公司，
其中 \hl{\num{6} 家负债率超过 \num{80}\%}：
中基健康（\num{104.1}\%）、*ST广糖（\num{100.2}\%）、天邦食品（\num{91.2}\%）、
新五丰（\num{84.2}\%）、金新农（\num{81.5}\%）、天域生物（\num{81.0}\%）。
\emph{中基健康与*ST广糖的净资产已经为负，其余四家的负债率也已接近
“净资产被亏损侵蚀殆尽”的临界点}。

据此，未来潜在融资需求可分为三档：

\begin{table}[htbp]
\centering
\footnotesize
\caption{农业上市公司未来潜在融资需求的三层结构（基于 2026 年半年报与近三年公告数据）}
\label{tab:fin-layers}
\begin{tabular}{@{}p{1.7cm}p{3.3cm}p{4.6cm}p{4.6cm}@{}}
\toprule
层次 & 特征与公司数 & 融资动机 & 可行的融资方式 \\
\midrule
第一档\newline 补血型 & 2026 年上半年亏损且 6 月末负债率 $>$65\%：\newline \num{14} 家（其中 $>$80\% \num{6} 家）
  & 维持银行授信、偿还到期债务、避免退市
  & 定增（需股价配合）、债务重组、出售资产、引入国资或战投；\hl{股权融资是最现实选择} \\
第二档\newline 保壳型 & 重组并购公告 $\geq$ \num{10} 条：\newline \num{8} 家（另有 \num{1} 家\emph{中基健康}已因亏损与高负债计入第一档）
  & 资产置换、跨界注入新业务、维持上市地位
  & 重大资产重组、发行股份购买资产、控制权转让；\hl{融资与“易主”同时发生} \\
第三档\newline 扩张型 & 半年 ROE $>$ \num{5}\%（约合年化 \num{10}\%）且 6 月末负债率 $<$ \num{65}\%：\newline \num{10} 家（动保、细分加工与农资为主）
  & 产能扩张、海外布局、品类并购
  & 定增 + 可转债、并购贷款；\hl{这类融资会被市场定价为成长而非补血} \\
\bottomrule
\end{tabular}
\end{table}

三档\hl{互不重复}：分类时按“补血—保壳—扩张”的顺序依次判定，
一家公司只归入最紧迫的一档（例如中基健康既有 \num{10} 条以上重组公告、
又在 2026 年上半年亏损且负债率超过 \num{65}\%（\num{104.1}\%），故计入第一档）。

需要说明的是三档之间的\emph{迁移}：与 2025 年年报口径相比，第一档的公司数由 \num{18} 家
减少到 \num{14} 家（其中 \num{6} 家负债率超过 \num{80}\%），
原因是个别高负债公司在 2026 年上半年实现扭亏或完成了资产处置，
但全样本亏损面却由 \num{34} 家扩大到 \num{46} 家（见第 \ref{sec:comp-h1} 节）——
\hl{“头部补血、尾部出清”是本轮周期底部最典型的资产负债表形态}。

按金额量级估算，第一档的融资需求最为迫切：这 \num{14} 家公司的
2026 年上半年亏损合计约 \hl{\num{73.1} 亿元}
（此处为该档成员的亏损合计数，并非板块或全样本的亏损合计）。
若只按“补足一年经营亏损 + 覆盖一年利息支出”的最低口径，
年度资金缺口约在 \hl{\num{150}---\num{250} 亿元}量级；
若按“把资产负债率压回 \num{65}\% 警戒线”的重组口径，
这 \num{14} 家公司 2026 年 6 月末的负债合计为 \num{1855} 亿元、
净资产合计仅 \num{694} 亿元，\hl{需要补充的净资产将超过 \num{2000} 亿元}
（按 $\text{目标净资产} = \text{负债} \div 0.65$ 粗略推算）。
而第三档公司的融资规模取决于其扩张节奏，单笔通常在
\num{5}---\num{20} 亿元之间。\emph{两者的根本差别在于：
第一档融资的用途是“活下去”，第三档融资的用途是“长大”}。

上面是“把负债率压回警戒线”的重组口径。若换成更贴近日常经营的\emph{滚动偿付}口径，
可以用债务结构直接度量压力：截至 2026 年 6 月末，\num{104} 家农业上市公司的有息负债合计
\hl{\num{4680} 亿元}（短期借款、一年内到期的非流动负债、长期借款与应付债券之和），
而账面货币资金合计仅 \hl{\num{1584} 亿元}；\hl{有 \num{65} 家公司的短期债务超过货币资金
（现金短债比小于 \num{1}），其中 \num{42} 家低于 \num{0.5}}。
以“短期债务 $-$ 货币资金”衡量的静态缺口，全样本合计约 \hl{\num{1781} 亿元}，
其中第一档 \num{14} 家合计约 \num{365} 亿元——这 \num{14} 家的货币资金 \num{228.5} 亿元
对应短期债务 \num{594} 亿元、有息负债 \num{1231} 亿元，
\hl{平均每 1 元货币资金要对应 2.6 元短期债务}。
从债务融资的流量看，2025 年全年全样本“取得借款收到的现金”合计 \num{4531} 亿元、
“偿还债务支付的现金”合计 \num{4773} 亿元，\hl{净偿还约 \num{242} 亿元}——
在本轮下行中，行业整体是在收缩债务而不是扩张债务，
这也解释了为什么“再融资”类公告的数量创出近年新高：\emph{债务端在收缩，}
\emph{资金缺口只能转向股权与重组}。

需要强调的是，这一缺口是\emph{压力情形}：现实中大部分短期债务通过银行续贷滚动，
货币资金中也可能包含募集资金专户与受限资金，
因此它可以用于公司之间的横向比较，却不能直接相加为一个“待融资总额”。
逐家公司的货币资金、短期债务、有息负债、现金短债比与静态缺口见附录表
\ref{tab:company-financing}（按现金短债比升序排列），
每家公司的画像里也给出了对应的债务结构与融资线索。

\begin{keybox}[融资需求分析的三个结论]
\normalsize
\textbf{① 融资动机由周期位置决定，而非行业属性}。同一家公司在 2021 年
募资用于扩产，在 2026 年募资用于补血——\hl{公告主题的变化就是周期的可见痕迹}。\par\vspace{3pt}
\textbf{② 高负债与亏损高度重合}：\num{46} 家亏损公司中有 \num{14} 家同时负债率超过 \num{65}\%，
其中 \num{6} 家超过 \num{80}\%。
\hl{负债率是比净利润更早的预警指标}。\par\vspace{3pt}
\textbf{③ 潜在融资需求的规模在数百亿元量级}，
且集中在生猪产业链（新希望、天邦食品、新五丰、唐人神、巨星农牧）
与困境加工企业（*ST广糖、中基健康、大禹节水、百洋股份）。
\hl{资本市场对农业板块的融资功能，在未来两三年内将主要被用于修复资产负债表，
而不是扩大产能}。
\end{keybox}

\section{结论与启示}
\label{sec:conclusion}

本报告用两个互补的视角考察中国农业的周期性：第一部分的\emph{品种视角}
回答“周期从哪里来、由什么决定、如何相互传导”；
第二部分的\emph{公司视角}回答“在周期中，谁赚到了钱、谁需要什么”。

\subsection{五个结论}

\paragraph{结论一：农业周期的驱动力量已经从需求转向供给与政策。}
在 2024---2025 年的这一轮下行中，中国的粮食与肉类\hl{需求并未出现显著萎缩}：
CPI 食品分项保持温和、水产品消费仍在增长。
价格下跌的主要原因是供给的恢复与产能的过冲：
玉米连续增产、谷物进口大幅收缩，生猪出栏量在产能去化未完成时仍在增长，
白羽肉鸡祖代更新量创历史新高。\emph{这意味着分析农业周期时，
供给端的产能数据（能繁母猪、祖代种鸡、投苗量、棉花种植面积）
比需求端的宏观数据更重要}。

\paragraph{结论二：不同品种的周期长度由“产能调整的生物学时间”决定。}
本报告给出了一个可比较的时间刻度：蛋鸡与肉鸡约 \num{6}---\num{12} 个月、
生猪约 \num{10}---\num{12} 个月、肉羊约 \num{18}---\num{24} 个月、
肉牛 \num{24}---\num{44} 个月以上、多年生作物（苹果、橡胶、糖料）\num{3}---\num{7} 年。
\hl{时间越长的品种，价格调整的幅度越大、拐点的持续性越强}——
肉牛的供给收缩要等到 2027 年之后，而禽类的产能已经完成了一轮出清与反转。

\paragraph{结论三：跨品种的相关性远弱于直觉。}
在第 \ref{sec:corr} 章的统一口径（月度收益率、同一时间窗口）下，
5\% 显著性阈值 $|\rho|>0.24$ 之上只有三个品种：花生 \num{+0.48}、粳米 \num{+0.29}
与豆一 \num{+0.26}（后两者仅略高于阈值）；
玉米 \num{+0.09}、豆粕 \num{+0.13} 等主力饲料原料品种的相关系数均不显著。
\hl{“猪价影响一切农产品”是一种叙事上的便利，而非统计事实}——
真正的共性来自共同的成本（饲料）与共同的宏观环境，
而这些因素在月度频率上已经被品种自身的供需波动淹没。

\paragraph{结论四：上市公司的盈利分布比行业均值更极端。}
按 2026 年半年报口径，\num{104} 家农业上市公司中 \num{46} 家亏损（\num{44.2}\%）、
\num{44} 家半年 ROE 为负，ROE 中位数仅 \num{0.75}\%（未年化）；
\hl{生猪养殖 \num{11} 家公司在 2026 年上半年全部亏损}，
合计亏损 \num{163.8} 亿元（上年同期为合计盈利 \num{154.1} 亿元）。
\hl{周期底部的行业整合不是渐进的，而是在同一时间、以“头部失血、尾部出清”
的方式一次性完成}。
但把逐公司画像放在一起看，\hl{同一板块内部的差别并不在规模，而在杠杆的时点}：
在景气高点借来的钱会把产能与成本锁在最高处，
而在行业普遍亏损的低点仍能拿到低成本资金的公司，
才有条件以低于重置成本的价格承接别人的产能——
\emph{决定谁能穿越周期的，是在\textbf{周期低点}（而不是高点）加杠杆扩张的能力}。
与 2025 年年报（\num{34} 家亏损、ROE 中位数 \num{2.64}\%）相比，
亏损面进一步扩大至 \num{46} 家，归母净利润合计由上年同期的 \num{+270.9} 亿元
转为 \num{-126.1} 亿元，\emph{报表层面的底部比价格层面的底部更晚出现}。

\paragraph{结论五：周期底部的融资需求是刚性的、而且以补血为主。}
按 2026 年半年报口径，\num{14} 家公司同时处于亏损与高负债（$>$ \num{65}\%）状态，
其中 \num{6} 家负债率超过 \num{80}\%；
\num{31} 家公司连续三年未分红。
这 \num{14} 家公司的负债合计 \num{1855} 亿元、净资产合计仅 \num{694} 亿元，
若要把资产负债率压回 \num{65}\% 的警戒线，需要补充的净资产约 \num{305} 亿元
（若改以偿还债务达标，则需减少负债约 \num{567} 亿元），
\hl{相当于这一档公司现有净资产的四成以上}。
近三年公告中“担保与关联交易”类公告数（\num{3184}）远高于
“对外投资与转型”类（\num{274}），\hl{说明农业上市公司当前的资本运作主题
是维持资产负债表，而不是扩张}。

\paragraph{结论六：周期底部最稀缺的能力，是在低点而不是高点加杠杆扩张。}
把 \num{104} 家公司的资产负债表与扩产节奏交叉起来，可以看到三种类型：
\emph{在 2026 年上半年仍能新增资本开支、并以自有资金或低成本融资完成扩张的公司}
（典型如海大集团、牧原股份的养殖成本优势、众兴菌业与华绿生物的工厂化产能）；
\emph{资产负债表已经不允许任何扩张、只能靠再融资维持的公司}
（\num{14} 家同时亏损且负债率超过 \num{65}\%）；
以及\emph{借助“补缴税款、主动去库存”等一次性事项出清历史包袱、
把报表重新做干净的公司}（北大荒补缴 \num{14.10} 亿元税款、
隆平高科主动去库存、种业龙头企业连续三年调减制种面积）。
\hl{第三类公司往往被市场当作“业绩暴雷”卖在第一轮，
但它们恰恰是下一轮周期里成本最低的产能承接者}。

\subsection{对投资者的启示}

\begin{enumerate}[leftmargin=1.5em, itemsep=2pt]
  \item \textbf{在周期底部，优先选择“供给刚性 + 政策底”的品种。}
        肉牛（进口保障措施、\num{44} 个月产能周期）与蛋鸡（存栏同比下降 \num{6.1}\%）
        是当前供给端最确定的两个方向；而玉米、小麦等政策收储品种
        在成本线附近的波动更适合区间交易而非趋势投资。
  \item \textbf{在同一产业链中，选择需求与价格脱钩的环节。}
        动保（半年 ROE 中位数 \num{2.95}\%、亏损面 \num{11.1}\%、
        资产负债率中位数 \num{30.2}\%）在农业整体下行中保持了正回报，
        海洋捕捞（\num{9.79}\%）则是资源属性带来的另一类“脱钩”。
        \hl{“卖水人”与“资源型”是本轮周期中最有效的两个抗周期逻辑}。
        \hl{种业是本轮供给端出清最早的环节}：全国玉米制种面积已从 2023 年的
        \num{462} 万亩降至 2026 年的 \num{358.5} 万亩（\num{-24}\%），
        但 \num{175}\% 的供需比说明库存尚未出清，
        因此种业的周期弹性更可能在 2027 年之后出现，而非当下。
        需要提示的是，宠物食品的半年 ROE 中位数已回落至 \num{4.00}\%，
        高成长赛道同样开始承受竞争压力。
  \item \textbf{对高负债亏损公司保持警觉，而不是抄底}。
        \num{6} 家负债率超过 \num{80}\% 且亏损的公司中，多数需要
        通过股权融资或资产重组才能维持经营，
        \hl{本报告推断这类公司的市值中包含相当比例的“壳价值”与“重整预期”，
        而非盈利能力的现值}——该判断基于市值与盈利的背离，
        仍需公司层面的债务结构与重组进展进一步验证。
  \item \textbf{用公告与产能数据做交叉验证，而不是依赖单一指标。}
        本报告的方法论支持一点：\emph{产能数据（能繁母猪、祖代存栏、投苗量）
        决定方向，公告数据（定增、重组、担保）决定风险，
        财务数据（负债率、ROE）决定谁能等到周期反转}。
\end{enumerate}

\subsection{尚待解决的问题}

本报告在数据上仍有三处限制，也构成后续研究的方向：
\textbf{其一}，由于公开数据源缺少牛、羊、禽、水产的连续现货价格序列，
这些品种的周期判断更多依赖产业统计（存栏、产量、均价）而非高频价格，
无法与生猪期货做严格的月度频率对比；
\textbf{其二}，水产品与部分经济作物的数据粒度只到年度，
限制了周期相位的精确划分；
\textbf{其三}，上市公司画像中的“未来融资需求”是基于财务约束
与公告行为的推断，尚不包含实际的债券到期结构与银行贷款条款等内部信息，
若需精确测算需引入更细的债务数据。

\begin{keybox}[一句话总结]
\normalsize
农业的周期不是一条曲线，而是\hl{一组由生物学时间、政策边界与资产负债表
共同决定的时钟}——\emph{看懂每个品种的时钟现在指在几点，
比预测下个月的价格更有意义}。
\end{keybox}


% ---------------------------------------------------------------- 参考文献
% 第二部分的正文与附录不使用 \cite（逐公司来源以 URL 形式保存在事实底稿
% data/facts/*.yaml 中，见附录 A 的口径说明），因此本册的参考文献列表用 \nocite{*}
% 输出与主报告完全相同的那一份（51 条），保证分册自成一体、不出现空书目。
\nocite{*}
\newpage
\printbibliography[heading=bibintoc, title={参考文献}]

% ---------------------------------------------------------------- 附录
% 附录 C（期货合约参数）属第一部分，本册不含；附录 A、B 均属第二部分的数据口径与明细
\appendix
\section{数据来源与口径说明}
\label{app:sources}

\subsection{数据来源}

本报告使用的数据分为四类，全部来自公开渠道：

\begin{enumerate}[leftmargin=1.5em, itemsep=2pt]
  \item \textbf{价格数据}：期货交易所公开的农产品期货主力连续合约行情；
        生猪现货价格与产业指标（猪易数据、玄田数据）；
        国家统计局的生猪存栏、出栏与猪肉产量统计。
  \item \textbf{行业数据}：申万宏源《申万行业分类标准（2021 版）》的一级、二级、三级行业
        及其成分股名单。
  \item \textbf{公司数据}：上市公司定期报告（合并报表口径）、公司公告（巨潮资讯网）、公司官网与公开报道，
        含公司概况、股本变动、分红送配与公告全文检索。
  \item \textbf{宏观数据}：国家统计局的 CPI、产量与进出口数据；中价指数的价格指数。
\end{enumerate}

\subsection{已知的数据缺口}

为使结论可追溯，报告在使用数据时明确标注了以下缺口：

\begin{enumerate}[leftmargin=1.5em, itemsep=2pt]
  \item \textbf{养殖品种缺少连续现货价格序列}：公开数据源没有肉牛、肉羊、白羽肉鸡、
        蛋鸡与水产品的连续现货价，这些品种的周期判断依赖农业农村部与产业技术体系的
        定期统计（存栏、出栏、产量、均价），无法参与月度频率的相关性检验。
  \item \textbf{期货库存数据不完整}：部分农产品期货的库存与仓单数据缺失，
        故正文未将库存作为周期的核心判据。
  \item \textbf{公司公告分类为关键词口径}：\num{41266} 条公告按标题关键词归入七类事件，
        存在同一公告命中多类的情况，各类计数用于公司之间的比较而非加总。
  \item \textbf{市值一律采用总市值口径}：以“报告数据时点收盘价 $\times$ 总股本”计算
        （总股本取自巨潮股本结构），\emph{不使用流通市值}；
        流通市值只覆盖已流通股份，会低估控股股东持股比例高的公司的相对规模。
  \item \textbf{榨季与市场年度口径}：白糖按榨季（当年 10 月至次年 9 月）、
        棉花与油脂按市场年度，与日历年不一致，正文在引用时已注明。
  \item \textbf{公司画像中的财务分析以公开披露信息为限}：2026 年数据为半年报口径（未年化），
        融资需求的分析基于财务约束与公告行为的推断，不含债券到期结构与银行贷款条款等
        非公开信息。
\end{enumerate}

\subsection{外部资料的核实原则}

除价格与财务数据外，报告引用的产业事实（产量、存栏、政策文件、进口数据）
均来自政府统计公报、部委公告与产业协会报告，
其中无法在原始来源核实的内容未写入正文。例如：\hl{“能繁母猪调控目标下调至 3650 万头”一说
未能在农业农村部原始文件中核实，报告统一采用已核实的 3750 万头}。

\section{农业上市公司明细表}
\label{app:companies}

本附录列出申万农林牧渔行业（801010）全部 \num{104} 家 A 股上市公司的画像数据。
第一张表为\hl{上市公司画像总表}：财务指标统一采用 \textbf{2026 年半年报}（合并口径、未年化），
资产负债率为 2026 年 6 月末数值，\hl{总市值时点为 2026 年 10 月 12 日}（收盘价 $\times$ 总股本）；
“再融资/重组并购”两列为公司近三年（2023 年 9 月---2026 年 9 月）相应类别公告的数量，
用于衡量公司在融资与并购上的活跃程度。
第二张表按申万三级行业汇总各子行业的公司数量与整体经营水平，口径与第一张表一致；
第三张表\hl{债务结构与滚动偿付压力}（\ref{tab:company-financing}）把 2026 年 6 月末的
货币资金、短期债务（短期借款 $+$ 一年内到期的非流动负债）、有息负债
（短期债务 $+$ 长期借款 $+$ 应付债券）与\emph{现金短债比}列在一起，按现金短债比升序排列，
并在末列标出第 \ref{sec:financing-need} 节的分层归属；“静态缺口”为
$\max(0,\ \text{短期债务} - \text{货币资金})$，即假设不续贷、不动用其他资产时的滚动偿付缺口。
需要强调的是，短期债务在实务中大多通过银行续贷滚动，账面货币资金也可能包含受限资金与
募集资金专户，因此该表用于公司之间的横向比较，不能相加为全行业的“待融资总额”。

\begingroup\scriptsize\setlength{\tabcolsep}{1.5pt}
\begin{longtable}{@{}llp{1.9cm}rrrrrr@{}}
\caption{申万农林牧渔行业 A 股上市公司画像总表（按总市值降序；总市值时点 2026-10-12，总市值 $=$ 收盘价 $\times$ 总股本；财务为 2026 年半年报（未年化），资产负债率为 2026 年 6 月末）}\label{tab:company-master}\\
\toprule
代码 & 公司简称 & 申万三级行业 & 总市值（亿元） & 营收（亿元） & 同比（\%） & 净利（亿元） & ROE（\%） & 负债率（\%） \\
\midrule
\endfirsthead
\multicolumn{9}{l}{\small（续）}\\
\toprule
代码 & 公司简称 & 申万三级行业 & 总市值（亿元） & 营收（亿元） & 同比（\%） & 净利（亿元） & ROE（\%） & 负债率（\%） \\
\midrule
\endhead
\bottomrule
\endlastfoot
002714 & 牧原股份 & 生猪养殖 & 2,283.4 & 594.1 & -22.3 & -60.78 & -7.5 & 54.2 \\
300999 & 金龙鱼 & 粮油加工 & 1,385.8 & 1,228.5 & 6.2 & 22.94 & 2.4 & 56.2 \\
300498 & 温氏股份 & 生猪养殖 & 1,007.3 & 467.6 & -6.2 & -43.66 & -11.1 & 58.9 \\
002311 & 海大集团 & 水产饲料 & 735.5 & 642.4 & 9.2 & 16.11 & 6.5 & 54.6 \\
600737 & 中粮糖业 & 其他农产品加工 & 342.0 & 92.8 & -21.2 & 6.08 & 5.2 & 39.1 \\
601118 & 海南橡胶 & 其他种植业 & 319.2 & 151.2 & -25.4 & -1.30 & -1.3 & 70.7 \\
000876 & 新希望 & 畜禽饲料 & 311.4 & 555.9 & 7.7 & -17.02 & -9.1 & 70.7 \\
002157 & 正邦科技 & 生猪养殖 & 269.6 & 77.2 & 12.9 & -7.36 & -7.0 & 53.5 \\
600598 & 北大荒 & 粮食种植 & 226.3 & 21.3 & -29.4 & -5.37 & -7.1 & 39.5 \\
605198 & 安德利 & 果蔬加工 & 213.3 & 7.2 & -24.4 & 1.23 & 4.3 & 7.8 \\
002299 & 圣农发展 & 肉鸡养殖 & 197.1 & 107.8 & 21.7 & 5.02 & 4.4 & 51.9 \\
300761 & 立华股份 & 肉鸡养殖 & 168.1 & 91.3 & 9.3 & -0.72 & -0.8 & 43.8 \\
605296 & 神农集团 & 生猪养殖 & 159.7 & 24.9 & -11.0 & -8.09 & -17.9 & 40.7 \\
301498 & 乖宝宠物 & 宠物食品 & 156.8 & 35.4 & 10.0 & 1.91 & 4.1 & 31.4 \\
000592 & 平潭发展 & 林业Ⅲ & 142.4 & 4.8 & -35.1 & -0.69 & -4.0 & 39.3 \\
002385 & 大北农 & 畜禽饲料 & 134.5 & 138.1 & 1.9 & -6.77 & -9.3 & 69.9 \\
600201 & 生物股份 & 动物保健Ⅲ & 133.3 & 7.6 & 22.1 & 1.04 & 1.9 & 16.5 \\
601952 & 苏垦农发 & 粮食种植 & 123.5 & 43.3 & -5.6 & 1.44 & 2.0 & 46.5 \\
000998 & 隆平高科 & 种子 & 113.8 & 17.1 & -21.1 & -2.73 & -4.1 & 59.1 \\
002100 & 天康生物 & 畜禽饲料 & 105.4 & 83.3 & -5.9 & -4.41 & -6.6 & 55.6 \\
000930 & 中粮科技 & 其他农产品加工 & 96.5 & 81.9 & -7.0 & -1.69 & -1.6 & 40.7 \\
002215 & 诺普信 & 其他种植业 & 94.9 & 39.6 & 7.8 & 7.65 & 16.8 & 60.5 \\
000048 & 京基智农 & 生猪养殖 & 91.0 & 21.3 & -10.7 & -4.89 & -12.9 & 63.6 \\
600127 & 金健米业 & 粮油加工 & 89.2 & 19.7 & 25.2 & -0.10 & -1.4 & 62.2 \\
002458 & 益生股份 & 肉鸡养殖 & 87.4 & 17.0 & 28.4 & 3.08 & 6.9 & 35.1 \\
002041 & 登海种业 & 种子 & 87.1 & 4.7 & 27.0 & 0.50 & 1.3 & 15.8 \\
002891 & 中宠股份 & 宠物食品 & 81.4 & 32.8 & 34.9 & 1.27 & 4.0 & 44.0 \\
600251 & 冠农股份 & 果蔬加工 & 81.2 & 13.9 & 0.3 & 5.37 & 13.3 & 30.8 \\
000019 & 深粮控股 & 粮油加工 & 80.7 & 23.9 & 0.1 & 1.12 & 2.2 & 38.1 \\
600108 & 亚盛集团 & 其他种植业 & 78.5 & 14.7 & -2.6 & 0.25 & 0.6 & 58.9 \\
300119 & 瑞普生物 & 动物保健Ⅲ & 76.7 & 17.2 & 0.8 & 2.04 & 4.3 & 34.4 \\
600313 & 农发种业 & 种子 & 72.6 & 38.7 & 32.8 & 0.53 & 2.2 & 37.2 \\
603477 & 巨星农牧 & 生猪养殖 & 71.8 & 35.6 & -4.3 & -8.38 & -30.3 & 69.4 \\
600195 & 中牧股份 & 动物保健Ⅲ & 67.9 & 31.7 & 13.5 & 0.50 & 0.9 & 30.2 \\
600975 & 新五丰 & 生猪养殖 & 64.6 & 28.3 & -23.2 & -7.36 & -41.0 & 84.2 \\
000735 & 罗牛山 & 生猪养殖 & 64.4 & 12.4 & 14.5 & -3.10 & -7.4 & 52.3 \\
300189 & 神农种业 & 种子 & 63.1 & 0.6 & -35.3 & -0.13 & -1.9 & 39.8 \\
600962 & 国投中鲁 & 果蔬加工 & 62.5 & 8.8 & -12.4 & 0.02 & 0.2 & 52.0 \\
603336 & 宏辉果蔬 & 其他种植业 & 59.6 & 4.2 & -11.1 & 0.05 & 0.4 & 36.3 \\
603609 & 禾丰股份 & 畜禽饲料 & 58.3 & 189.6 & 8.9 & -0.08 & -0.1 & 57.7 \\
300138 & 晨光生物 & 其他农产品加工 & 57.1 & 31.6 & -13.6 & 1.79 & 5.2 & 64.1 \\
688526 & 科前生物 & 动物保健Ⅲ & 56.8 & 3.7 & -23.3 & 1.16 & 2.8 & 10.7 \\
002481 & 双塔食品 & 其他农产品加工 & 54.4 & 11.2 & 7.3 & 0.33 & 1.4 & 44.8 \\
002746 & 仙坛股份 & 肉鸡养殖 & 53.6 & 29.9 & 17.9 & 1.80 & 3.9 & 43.0 \\
300087 & ST荃银 & 种子 & 52.9 & 11.8 & -17.9 & -0.55 & -3.3 & 64.8 \\
002567 & 唐人神 & 畜禽饲料 & 52.0 & 108.8 & -12.7 & -10.46 & -24.9 & 71.1 \\
001313 & 粤海饲料 & 水产饲料 & 49.6 & 37.8 & 41.7 & 0.48 & 1.9 & 52.1 \\
600354 & 敦煌种业 & 种子 & 49.1 & 8.7 & 20.9 & 1.45 & 18.7 & 37.5 \\
002556 & 辉隆股份 & 农业综合Ⅱ & 48.5 & 85.5 & 3.3 & 1.13 & 3.1 & 67.0 \\
603668 & 天马科技 & 水产饲料 & 47.3 & 31.3 & 5.3 & 0.12 & 0.6 & 72.1 \\
001338 & 永顺泰 & 其他农产品加工 & 46.6 & 19.1 & -13.9 & 0.51 & 1.4 & 21.1 \\
000505 & 京粮控股 & 粮油加工 & 45.7 & 31.1 & -26.0 & 0.08 & 0.3 & 43.5 \\
600191 & 华资实业 & 其他农产品加工 & 44.0 & 3.5 & 58.0 & 0.36 & 2.1 & 23.8 \\
002124 & 天邦食品 & 生猪养殖 & 43.3 & 36.6 & -22.1 & -15.43 & -96.8 & 91.2 \\
002321 & 华英农业 & 其他养殖 & 43.1 & 23.7 & 14.5 & -0.29 & -3.2 & 66.2 \\
301116 & 益客食品 & 肉鸡养殖 & 41.7 & 99.0 & 13.2 & 0.34 & 2.2 & 68.9 \\
002772 & 众兴菌业 & 食用菌 & 41.4 & 10.9 & 19.6 & 1.91 & 5.3 & 54.0 \\
600371 & 万向德农 & 种子 & 41.0 & 1.0 & -15.1 & 0.14 & 2.6 & 14.5 \\
002688 & 金河生物 & 动物保健Ⅲ & 39.9 & 14.5 & 4.1 & 1.31 & 5.7 & 56.1 \\
002086 & 东方海洋 & 其他农产品加工 & 39.8 & 1.6 & 0.9 & -0.48 & -3.8 & 19.6 \\
000713 & 国投丰乐 & 种子 & 38.1 & 11.9 & 3.9 & -0.31 & -1.0 & 32.1 \\
300313 & 天山生物 & 其他养殖 & 37.8 & 2.2 & 427.7 & 0.04 & 12.1 & 82.0 \\
001201 & 东瑞股份 & 生猪养殖 & 37.1 & 10.4 & -5.1 & -4.25 & -14.3 & 52.9 \\
000798 & 中水渔业 & 海洋捕捞 & 36.7 & 20.8 & 18.9 & 1.02 & 20.9 & 71.6 \\
300673 & 佩蒂股份 & 宠物食品 & 36.2 & 7.3 & 1.0 & 0.24 & 1.4 & 35.3 \\
603566 & 普莱柯 & 动物保健Ⅲ & 35.4 & 4.9 & -12.2 & 0.77 & 3.0 & 15.2 \\
600467 & 好当家 & 水产养殖 & 34.5 & 8.0 & 13.0 & 0.30 & 0.9 & 49.1 \\
688098 & 申联生物 & 动物保健Ⅲ & 33.7 & 1.1 & -8.1 & -0.19 & -1.4 & 9.8 \\
300094 & 国联水产 & 水产养殖 & 33.4 & 14.5 & -12.5 & 0.10 & 10.0 & 95.2 \\
300871 & 回盛生物 & 动物保健Ⅲ & 33.4 & 8.6 & 5.2 & 2.06 & 9.1 & 31.8 \\
002868 & 绿康生化 & 动物保健Ⅲ & 33.1 & 2.5 & -15.7 & 0.07 & 10.9 & 90.1 \\
300021 & 大禹节水 & 农业综合Ⅱ & 32.7 & 15.9 & 24.3 & -0.01 & -0.0 & 71.7 \\
003030 & 祖名股份 & 其他农产品加工 & 32.5 & 10.1 & 10.0 & -0.24 & -2.5 & 57.9 \\
603231 & 索宝蛋白 & 粮油加工 & 32.4 & 9.7 & 24.9 & 0.68 & 3.4 & 13.6 \\
000972 & 中基健康 & 果蔬加工 & 31.6 & 2.0 & -19.6 & -0.66 & — & 104.1 \\
002286 & 保龄宝 & 其他农产品加工 & 31.1 & 15.0 & 6.9 & 0.82 & 3.6 & 23.1 \\
002548 & 金新农 & 畜禽饲料 & 30.8 & 23.2 & -2.4 & -3.88 & -37.4 & 81.5 \\
603151 & 邦基科技 & 畜禽饲料 & 30.7 & 26.8 & 11.8 & -0.00 & -0.0 & 63.4 \\
600540 & 新赛股份 & 其他种植业 & 30.4 & 25.7 & 77.2 & 1.25 & 30.8 & 82.7 \\
300967 & 晓鸣股份 & 肉鸡养殖 & 29.5 & 5.7 & -24.5 & 0.02 & 0.2 & 52.3 \\
600265 & *ST景谷 & 林业Ⅲ & 29.3 & 0.6 & -54.6 & -0.00 & -0.0 & 26.7 \\
002852 & 道道全 & 粮油加工 & 29.2 & 28.6 & 2.6 & 0.17 & 0.8 & 49.1 \\
002696 & 百洋股份 & 水产饲料 & 27.9 & 20.5 & 21.1 & -0.83 & -6.2 & 66.4 \\
600359 & 新农开发 & 其他种植业 & 27.3 & 2.0 & -31.1 & 0.19 & 2.8 & 55.6 \\
002234 & 民和股份 & 肉鸡养殖 & 26.8 & 10.3 & 4.2 & -0.95 & -5.5 & 53.9 \\
002679 & 福建金森 & 林业Ⅲ & 26.3 & 0.5 & -9.5 & -0.14 & -1.8 & 61.9 \\
002069 & 獐子岛 & 水产养殖 & 25.8 & 7.3 & -5.3 & 0.13 & 29.5 & 97.0 \\
600257 & 大湖股份 & 水产养殖 & 25.7 & 4.3 & 2.0 & 0.17 & 2.1 & 46.4 \\
688156 & 路德科技 & 畜禽饲料 & 25.6 & 2.1 & 42.9 & 0.09 & 1.2 & 55.9 \\
000911 & *ST广糖 & 其他农产品加工 & 24.9 & 15.8 & 17.4 & -0.16 & 0.0 & 100.2 \\
300175 & 朗源股份 & 果蔬加工 & 24.6 & 1.7 & 67.7 & -0.13 & -2.7 & 8.8 \\
603363 & 傲农生物 & 畜禽饲料 & 24.5 & 45.2 & 14.3 & -1.90 & -6.2 & 62.1 \\
300511 & 雪榕生物 & 食用菌 & 23.7 & 11.6 & 46.2 & 0.18 & 1.1 & 56.5 \\
600097 & 开创国际 & 海洋捕捞 & 21.9 & 11.8 & -5.4 & -0.30 & -1.3 & 34.4 \\
603182 & 嘉华股份 & 粮油加工 & 21.9 & 8.6 & 25.7 & 0.57 & 5.0 & 18.8 \\
001366 & 播恩集团 & 畜禽饲料 & 20.8 & 6.7 & 9.0 & 0.10 & 1.3 & 38.1 \\
600965 & *ST福成 & 其他养殖 & 20.1 & 5.7 & 6.9 & 0.18 & 0.8 & 16.7 \\
300970 & 华绿生物 & 食用菌 & 20.0 & 6.9 & 39.2 & 0.46 & 2.8 & 35.4 \\
000702 & 正虹科技 & 畜禽饲料 & 20.0 & 5.7 & 14.2 & -0.21 & -6.2 & 54.2 \\
000663 & 永安林业 & 林业Ⅲ & 19.9 & 1.5 & 14.7 & -0.19 & -2.1 & 36.6 \\
300268 & 佳沃食品 & 其他农产品加工 & 19.7 & 4.2 & -66.1 & -0.04 & -1.4 & 35.5 \\
600883 & 博闻科技 & 食用菌 & 18.6 & 0.3 & -8.8 & 0.15 & 1.5 & 6.6 \\
002982 & 湘佳股份 & 肉鸡养殖 & 18.3 & 23.2 & 8.6 & -0.12 & -0.7 & 60.9 \\
603717 & 天域生物 & 生猪养殖 & 17.2 & 3.5 & -15.0 & -0.54 & -14.6 & 81.0 \\
\end{longtable}
\endgroup


\vspace{10pt}
\begin{table}[htbp]\centering\footnotesize
\caption{按申万三级行业汇总的农业上市公司数量与经营指标（2026 年半年报口径；ROE 为半年报披露的算术平均值、未年化，总市值时点 2026-10-12）}
\label{tab:company-by-industry}
\begin{tabular}{@{}lrrrrrr@{}}
\toprule
申万三级行业 & 公司数 & 总市值 & 平均营收 & 营收同比 & 亏损 & 平均 ROE \\
 & & （亿元） & 2026H1（亿元） & 中位数（\%） & 公司数 & 2026H1（\%） \\
\midrule
其他农产品加工 & 11 & 788 & 26.1 & 0.9 & 5 & 0.9 \\
畜禽饲料 & 11 & 814 & 107.8 & 8.9 & 9 & -8.9 \\
生猪养殖 & 11 & 4,109 & 119.3 & -10.7 & 11 & -23.7 \\
动物保健Ⅲ & 9 & 510 & 10.2 & 0.8 & 1 & 4.1 \\
肉鸡养殖 & 8 & 622 & 48.0 & 11.2 & 3 & 1.3 \\
种子 & 8 & 518 & 11.8 & -5.6 & 4 & 1.8 \\
粮油加工 & 7 & 1,685 & 192.9 & 6.2 & 1 & 1.8 \\
其他种植业 & 6 & 610 & 39.6 & -6.9 & 1 & 8.3 \\
果蔬加工 & 5 & 413 & 6.7 & -12.4 & 2 & 3.8 \\
林业Ⅲ & 4 & 218 & 1.8 & -22.3 & 4 & -2.0 \\
水产饲料 & 4 & 860 & 183.0 & 15.1 & 1 & 0.7 \\
食用菌 & 4 & 104 & 7.4 & 29.4 & 0 & 2.7 \\
水产养殖 & 4 & 119 & 8.5 & -1.7 & 0 & 10.6 \\
其他养殖 & 3 & 101 & 10.5 & 14.5 & 1 & 3.3 \\
宠物食品 & 3 & 274 & 25.2 & 10.0 & 0 & 3.1 \\
农业综合Ⅱ & 2 & 81 & 50.7 & 13.8 & 1 & 1.5 \\
海洋捕捞 & 2 & 59 & 16.3 & 6.8 & 1 & 9.8 \\
粮食种植 & 2 & 350 & 32.3 & -17.5 & 1 & -2.5 \\
\bottomrule
\end{tabular}
\end{table}


\vspace{10pt}
\begingroup\scriptsize\setlength{\tabcolsep}{1.5pt}
\begin{longtable}{@{}llp{1.9cm}rrrrrl@{}}
\caption{农业上市公司债务结构与滚动偿付压力（2026 年 6 月末；按现金短债比升序；单位：亿元，现金短债比为倍）}\label{tab:company-financing}\\
\toprule
代码 & 公司简称 & 申万三级行业 & 货币资金 & 短期债务 & 有息负债 & 现金短债比 & 静态缺口 & 融资需求档 \\
\midrule
\endfirsthead
\multicolumn{9}{l}{\small（续）}\\
\toprule
代码 & 公司简称 & 申万三级行业 & 货币资金 & 短期债务 & 有息负债 & 现金短债比 & 静态缺口 & 融资需求档 \\
\midrule
\endhead
\bottomrule
\endlastfoot
002124 & 天邦食品 & 生猪养殖 & 1.55 & 35.01 & 48.40 & 0.04 & 33.46 & 第一档 \\
000019 & 深粮控股 & 粮油加工 & 0.87 & 18.04 & 21.43 & 0.05 & 17.18 & --- \\
000911 & *ST广糖 & 其他农产品加工 & 2.29 & 29.24 & 31.13 & 0.08 & 26.95 & 第一档 \\
002548 & 金新农 & 畜禽饲料 & 2.36 & 22.72 & 28.92 & 0.10 & 20.37 & 第一档 \\
300761 & 立华股份 & 肉鸡养殖 & 3.67 & 29.94 & 35.68 & 0.12 & 26.26 & --- \\
603668 & 天马科技 & 水产饲料 & 4.61 & 34.56 & 41.55 & 0.13 & 29.95 & --- \\
600962 & 国投中鲁 & 果蔬加工 & 1.57 & 9.37 & 10.30 & 0.17 & 7.80 & 第二档 \\
002299 & 圣农发展 & 肉鸡养殖 & 11.73 & 69.56 & 76.86 & 0.17 & 57.83 & --- \\
300511 & 雪榕生物 & 食用菌 & 1.76 & 10.20 & 12.52 & 0.17 & 8.43 & --- \\
002868 & 绿康生化 & 动物保健Ⅲ & 0.51 & 2.94 & 3.19 & 0.17 & 2.42 & 第二档 \\
300094 & 国联水产 & 水产养殖 & 1.98 & 10.29 & 10.84 & 0.19 & 8.30 & --- \\
600108 & 亚盛集团 & 其他种植业 & 5.31 & 27.04 & 46.72 & 0.20 & 21.73 & --- \\
600127 & 金健米业 & 粮油加工 & 1.50 & 7.28 & 7.31 & 0.21 & 5.78 & --- \\
003030 & 祖名股份 & 其他农产品加工 & 1.80 & 8.74 & 10.02 & 0.21 & 6.94 & --- \\
600467 & 好当家 & 水产养殖 & 5.64 & 23.80 & 27.09 & 0.24 & 18.16 & --- \\
300498 & 温氏股份 & 生猪养殖 & 36.75 & 146.20 & 286.24 & 0.25 & 109.44 & --- \\
301116 & 益客食品 & 肉鸡养殖 & 4.84 & 17.56 & 23.69 & 0.28 & 12.72 & --- \\
300268 & 佳沃食品 & 其他农产品加工 & 0.48 & 1.70 & 1.72 & 0.28 & 1.22 & 第二档 \\
600737 & 中粮糖业 & 其他农产品加工 & 8.58 & 30.28 & 30.71 & 0.28 & 21.70 & 第三档 \\
001201 & 东瑞股份 & 生猪养殖 & 4.55 & 15.69 & 24.27 & 0.29 & 11.14 & --- \\
002714 & 牧原股份 & 生猪养殖 & 165.47 & 568.32 & 748.02 & 0.29 & 402.86 & --- \\
002385 & 大北农 & 畜禽饲料 & 32.81 & 104.33 & 135.45 & 0.31 & 71.52 & 第一档 \\
000702 & 正虹科技 & 畜禽饲料 & 0.60 & 1.81 & 1.97 & 0.33 & 1.21 & --- \\
600191 & 华资实业 & 其他农产品加工 & 0.46 & 1.38 & 2.53 & 0.33 & 0.92 & --- \\
603477 & 巨星农牧 & 生猪养殖 & 6.94 & 20.13 & 37.45 & 0.34 & 13.19 & 第一档 \\
002321 & 华英农业 & 其他养殖 & 1.92 & 5.52 & 6.13 & 0.35 & 3.60 & 第一档 \\
002696 & 百洋股份 & 水产饲料 & 5.32 & 15.11 & 16.70 & 0.35 & 9.79 & 第一档 \\
002215 & 诺普信 & 其他种植业 & 12.33 & 33.99 & 58.48 & 0.36 & 21.66 & 第三档 \\
000998 & 隆平高科 & 种子 & 24.25 & 64.83 & 107.27 & 0.37 & 40.58 & --- \\
000876 & 新希望 & 畜禽饲料 & 79.63 & 205.48 & 579.09 & 0.39 & 125.85 & 第一档 \\
002852 & 道道全 & 粮油加工 & 4.92 & 12.69 & 12.69 & 0.39 & 7.77 & --- \\
002069 & 獐子岛 & 水产养殖 & 5.45 & 14.05 & 19.24 & 0.39 & 8.60 & --- \\
300087 & ST荃银 & 种子 & 7.67 & 19.06 & 22.07 & 0.40 & 11.39 & 第二档 \\
002556 & 辉隆股份 & --- & 11.20 & 27.68 & 31.42 & 0.40 & 16.47 & --- \\
000735 & 罗牛山 & 生猪养殖 & 4.11 & 9.97 & 19.18 & 0.41 & 5.86 & --- \\
002234 & 民和股份 & 肉鸡养殖 & 5.14 & 12.38 & 13.65 & 0.42 & 7.24 & --- \\
300119 & 瑞普生物 & 动物保健Ⅲ & 5.32 & 12.68 & 16.16 & 0.42 & 7.36 & 第二档 \\
002481 & 双塔食品 & 其他农产品加工 & 3.84 & 8.95 & 12.29 & 0.43 & 5.11 & --- \\
300313 & 天山生物 & 其他养殖 & 0.26 & 0.58 & 0.93 & 0.45 & 0.31 & --- \\
000048 & 京基智农 & 生猪养殖 & 6.08 & 13.29 & 19.75 & 0.46 & 7.22 & --- \\
603363 & 傲农生物 & 畜禽饲料 & 4.75 & 9.72 & 23.80 & 0.49 & 4.97 & --- \\
603717 & 天域生物 & 生猪养殖 & 1.35 & 2.75 & 10.49 & 0.49 & 1.40 & 第一档 \\
600975 & 新五丰 & 生猪养殖 & 6.08 & 12.10 & 73.49 & 0.50 & 6.02 & 第一档 \\
002688 & 金河生物 & 动物保健Ⅲ & 7.00 & 13.82 & 23.64 & 0.51 & 6.82 & 第三档 \\
000798 & 中水渔业 & 海洋捕捞 & 8.75 & 16.75 & 27.86 & 0.52 & 8.00 & --- \\
300138 & 晨光生物 & 其他农产品加工 & 28.02 & 52.71 & 53.17 & 0.53 & 24.70 & 第三档 \\
603609 & 禾丰股份 & 畜禽饲料 & 12.18 & 22.89 & 66.24 & 0.53 & 10.71 & --- \\
300999 & 金龙鱼 & 粮油加工 & 485.84 & 887.12 & 927.22 & 0.55 & 401.28 & --- \\
601118 & 海南橡胶 & 其他种植业 & 51.31 & 93.23 & 167.80 & 0.55 & 41.92 & 第一档 \\
603151 & 邦基科技 & 畜禽饲料 & 7.22 & 13.06 & 15.02 & 0.55 & 5.83 & 第二档 \\
600540 & 新赛股份 & 其他种植业 & 8.39 & 14.35 & 15.08 & 0.58 & 5.96 & --- \\
000972 & 中基健康 & 果蔬加工 & 3.01 & 5.01 & 5.54 & 0.60 & 2.00 & 第一档 \\
000930 & 中粮科技 & 其他农产品加工 & 21.29 & 34.61 & 49.14 & 0.62 & 13.32 & --- \\
605296 & 神农集团 & 生猪养殖 & 4.50 & 6.59 & 13.62 & 0.68 & 2.10 & --- \\
300021 & 大禹节水 & --- & 12.85 & 18.48 & 29.12 & 0.70 & 5.62 & 第一档 \\
600359 & 新农开发 & 其他种植业 & 3.42 & 4.87 & 6.52 & 0.70 & 1.44 & --- \\
688156 & 路德科技 & 畜禽饲料 & 2.43 & 3.38 & 5.50 & 0.72 & 0.94 & --- \\
002679 & 福建金森 & 林业Ⅲ & 1.34 & 1.75 & 11.56 & 0.77 & 0.41 & --- \\
000592 & 平潭发展 & 林业Ⅲ & 4.17 & 5.27 & 5.30 & 0.79 & 1.10 & --- \\
002100 & 天康生物 & 畜禽饲料 & 25.35 & 30.51 & 65.33 & 0.83 & 5.16 & --- \\
002311 & 海大集团 & 水产饲料 & 34.25 & 40.97 & 64.39 & 0.84 & 6.72 & 第三档 \\
002567 & 唐人神 & 畜禽饲料 & 21.10 & 24.84 & 61.18 & 0.85 & 3.73 & 第一档 \\
002891 & 中宠股份 & 宠物食品 & 13.68 & 16.07 & 18.56 & 0.85 & 2.39 & --- \\
002458 & 益生股份 & 肉鸡养殖 & 14.79 & 16.42 & 18.88 & 0.90 & 1.63 & 第三档 \\
002982 & 湘佳股份 & 肉鸡养殖 & 3.94 & 4.02 & 19.28 & 0.98 & 0.08 & --- \\
600257 & 大湖股份 & 水产养殖 & 2.61 & 2.50 & 6.87 & 1.04 & -0.11 & --- \\
301498 & 乖宝宠物 & 宠物食品 & 9.85 & 9.29 & 11.94 & 1.06 & -0.55 & --- \\
001366 & 播恩集团 & 畜禽饲料 & 2.60 & 2.37 & 3.68 & 1.10 & -0.23 & --- \\
300189 & 神农种业 & 种子 & 0.99 & 0.88 & 3.05 & 1.11 & -0.10 & --- \\
600251 & 冠农股份 & 果蔬加工 & 13.50 & 11.36 & 11.87 & 1.19 & -2.14 & 第三档 \\
001313 & 粤海饲料 & 水产饲料 & 6.78 & 5.52 & 8.40 & 1.23 & -1.26 & --- \\
688098 & 申联生物 & 动物保健Ⅲ & 0.65 & 0.53 & 0.53 & 1.23 & -0.12 & --- \\
002772 & 众兴菌业 & 食用菌 & 16.69 & 12.38 & 35.65 & 1.35 & -4.31 & 第三档 \\
002746 & 仙坛股份 & 肉鸡养殖 & 36.20 & 25.39 & 25.92 & 1.43 & -10.80 & --- \\
300967 & 晓鸣股份 & 肉鸡养殖 & 1.43 & 0.96 & 6.43 & 1.49 & -0.47 & --- \\
000505 & 京粮控股 & 粮油加工 & 16.90 & 10.88 & 17.91 & 1.55 & -6.02 & --- \\
300871 & 回盛生物 & 动物保健Ⅲ & 3.59 & 2.02 & 4.64 & 1.77 & -1.57 & 第三档 \\
001338 & 永顺泰 & 其他农产品加工 & 14.14 & 7.38 & 7.41 & 1.92 & -6.76 & --- \\
600195 & 中牧股份 & 动物保健Ⅲ & 16.65 & 8.44 & 11.12 & 1.97 & -8.21 & --- \\
600097 & 开创国际 & 海洋捕捞 & 3.90 & 1.76 & 4.79 & 2.22 & -2.14 & --- \\
002286 & 保龄宝 & 其他农产品加工 & 3.13 & 1.37 & 1.77 & 2.29 & -1.76 & --- \\
603336 & 宏辉果蔬 & 其他种植业 & 4.04 & 1.71 & 7.43 & 2.37 & -2.33 & --- \\
002157 & 正邦科技 & 生猪养殖 & 16.97 & 6.97 & 19.10 & 2.43 & -9.99 & --- \\
600313 & 农发种业 & 种子 & 8.73 & 3.58 & 4.77 & 2.44 & -5.15 & --- \\
601952 & 苏垦农发 & 粮食种植 & 14.56 & 4.28 & 45.25 & 3.40 & -10.28 & --- \\
603231 & 索宝蛋白 & 粮油加工 & 2.84 & 0.65 & 0.67 & 4.39 & -2.19 & --- \\
000663 & 永安林业 & 林业Ⅲ & 1.68 & 0.32 & 0.55 & 5.28 & -1.36 & 第二档 \\
000713 & 国投丰乐 & 种子 & 15.50 & 2.38 & 3.72 & 6.52 & -13.12 & --- \\
603566 & 普莱柯 & 动物保健Ⅲ & 0.86 & 0.07 & 0.26 & 11.86 & -0.78 & --- \\
300970 & 华绿生物 & 食用菌 & 5.89 & 0.31 & 7.00 & 19.04 & -5.58 & --- \\
600354 & 敦煌种业 & 种子 & 8.05 & 0.33 & 0.65 & 24.20 & -7.72 & 第三档 \\
600883 & 博闻科技 & 食用菌 & 1.43 & 0.05 & 0.12 & 30.40 & -1.39 & --- \\
300673 & 佩蒂股份 & 宠物食品 & 4.95 & 0.14 & 8.21 & 34.24 & -4.80 & --- \\
002086 & 东方海洋 & 其他农产品加工 & 2.43 & 0.06 & 0.49 & 37.76 & -2.36 & --- \\
600965 & *ST福成 & 其他养殖 & 2.03 & 0.05 & 0.98 & 43.48 & -1.98 & --- \\
600265 & *ST景谷 & 林业Ⅲ & 0.92 & 0.02 & 0.04 & 55.11 & -0.90 & 第二档 \\
300175 & 朗源股份 & 果蔬加工 & 0.84 & 0.01 & 0.28 & 61.19 & -0.82 & --- \\
688526 & 科前生物 & 动物保健Ⅲ & 1.70 & 0.02 & 0.07 & 96.61 & -1.68 & --- \\
002041 & 登海种业 & 种子 & 21.43 & 0.19 & 0.26 & 112.76 & -21.24 & --- \\
600598 & 北大荒 & 粮食种植 & 19.27 & 0.03 & 0.10 & 566.76 & -19.24 & --- \\
600201 & 生物股份 & 动物保健Ⅲ & 16.09 & 0.01 & 1.01 & 1,426.53 & -16.08 & --- \\
605198 & 安德利 & 果蔬加工 & 7.39 & --- & --- & --- & --- & --- \\
600371 & 万向德农 & 种子 & 2.26 & --- & 0.00 & --- & --- & --- \\
603182 & 嘉华股份 & 粮油加工 & 1.39 & --- & 0.14 & --- & --- & --- \\
\end{longtable}
\endgroup


\vspace{10pt}
\noindent\textbf{使用说明}：表中“同比”为 2026 年上半年营业收入相对 2025 年上半年的增速；
“ROE”为半年报披露的净资产收益率，\hl{未年化}（如需与年报口径比较，可近似乘以 2），
净资产为负的公司该指标记为“—”；
“—”表示该数据在所获取的公开接口中缺失。
全部数据由 \texttt{scripts/make\_company.py} 生成（正文统计数据由
\texttt{scripts/company\_stats.py} 汇总），
可通过 \texttt{data/clean/company\_master.csv} 复核任意单元格。


\end{document}
